%    \iffalse
%
% texshade.dtx
% Docstrip archive,  to extract documentation  run twice through 
% LaTeX.
% To  successfully  extract the documentation it is necessary to
% first run the file `texshade.ins' through LaTeX. This produces 
% the needed style  file  `texshade.sty'  and the parameter file
% `texshade.def' as well as several example files.  See the file
% `README.txt' for further information!
% 
%
% Copyright (C) 1999-2021 Eric Beitz
% See the file texshade.txt
% 
%    \fi
%
%    \changes{1.0}{1999-5-12}{First release}
%    \changes{1.1}{1999-5-26}{%
%      Corrections:  `emphregion'  is not extending to the next
%                                  alignment any more;
%                    `namecolor' & `numbercolor' are now correctly
%                                                reordered;
%                    sequence gaps at the beginning or the end are
%                    now treated correctly, i.e. no symbols are shown.
%      Introduction: `seqlength'
%                    `gapcolors'
%                    `rulersteps'
%                    `hideresidues'
%                    `showresidues'
%                    `fingerprint'.}
%    \changes{1.2}{1999-6-12}{%
%      Corrections:  functional shading error in funcgroup no. 8.
%      Introduction: `includeDSSP'
%                    `includeSTRIDE'
%                    `includePHDsec'
%                    `includePHDtopo'
%                    `appearance'
%                    `numcount'
%                    `alphacount'
%                    `Alphacount'
%                    `showonDSSP'
%                    `hideonDSSP'
%                    `showonSTRIDE'
%                    `hideonSTRIDE'
%                    `showonPHDsec'
%                    `hideonPHDsec'
%                    `showonPHDtopo'
%                    `hideonPHDtopo'.}
%    \changes{1.2a}{1999-6-24}{%
%      Corrections:  `namecolor' & `numbercolor' are now really
%                                                correctly reordered;
%                    in sequence names ( and ) are now allowed;
%                    option `case' in `funcshadingstyle' works now.}
%    \changes{1.3}{2000-3-3}{%
%      Corrections:  features in the ttop row do not produce line
%                                                scrambling any more;
%                    `language' replaced by `germanlanguage'
%                                       and `englishlanguage' due to
%                                       incompatibilities with babel;
%                    incompatibility with amsmath's text command fixed.
%      Introduction: new feature option `translate'
%                    `codon'
%                    `geneticcode'
%                    `backtranslabel'
%                    `backtranstext'
%                    `romancount'
%                    `Romancount'
%                    TeXtopo compatibility.}
%    \changes{1.3a}{2000-7-28}{%
%      Introduction: `showleadinggaps'
%                    `hideleadinggaps' to hide or show gap symbols 
%                                      before the actual seq start.}
%    \changes{1.3b}{2000-7-30}{%
%      Corrections:  `showleadinggaps'
%                    `hideleadinggaps' were extended to `ending' gaps;
%                    sequence names input routine now accepts special 
%                    characters.}
%                    
%    \changes{1.4}{2000-9-12}{%
%      Introduction: `movelegend'  allows one to move the legend;
%                    series of sequence numbers, e.g. in `orderseqs',
%                    can now be written with a dash, e.g. {1-3,6-4,7} instead
%                    of {1,2,3,6,5,4,7}.}
%
%    \changes{1.4a}{2000-10-3}{Documentation and FAQ additions}
%    
%    \changes{1.5}{2001-02-22}{%
%      Corrections:  `X's in the alignment lead to a run-time error;
%      Introduction: `ttopspace'
%                    `topspace'
%                    `bottomspace'
%                    `bbottomspace' for controlling vertical space 
%                                   between feature lines.
%                    `showcaption'  for adding a caption to the alignment.
%                    the sequence lengths are now stored in the .aux 
%                    in order to have correct gap breaks after the seqs.}
%                    
%    \changes{1.5a}{2001-03-08}{%
%      Corrections:  Eckhart Guth\"orlein noticed a sorting problem 
%                    when in addition some sequences where `killed' and
%                    the consensus was set to a particular sequence.
%                    This update fixes this problem.}
%                    
%    \changes{1.6}{2002-03-26}{%
%      Corrections:  There is no restriction to `dvips' anymore. One
%                    can add an option to the \usepackage{texshade}
%                    call which is passed to color.sty, e.g. `dvipdf';
%                    `noblockskip' led to over-printing of lines;
%                    `namecolor' and `numbercolor' did not support 
%                    sequence lists - fixed (thanks to Denys Bashtovyy).
%      Introduction: The FASTA file format is now supported;
%                    references to sequences can be made by name in
%                    addition to number;
%                    sequences can be refered to by their name in 
%                    addition to their number in the input file
%                    (suggested by Christoph Gille);
%                    `flexblockspace'   optimizes the space between 
%                                       sequence blocks to be minimal (is
%                                       default as before);
%                    `fixblockspace'    leads to an equal separation of
%                                       sequence blocks independent of
%                                       feature lines;
%                    `firstcolumnDSSP'  lets you choose the first numbering
%                                       column in DSSP input files and
%                    `secondcolumnDSSP' the second column.}
%                    
%    \changes{1.7}{2004-01-05}{%
%      Corrections:  Several bugs were fixed.
%                    In gaps the wrong character was plotted in `donotshade'
%                    mode. Gaps were colored incorrectly when a single 
%                    sequence was set as consensus. (thanks to Jeferson J. 
%                    Arenzon). Another `donotshade' problem was solved 
%                    which led to a halt of the LaTeX run (thanks go to
%                    Naomi Siew). The gap and match labels in diverse
%                    mode were switched (`-' in gaps; `.' at matching
%                    positions) in order to follow convention.
%      Introduction: Spanish labels (contributed by Mikel Egana Aranguren);
%                    New feature label `helix'.}
%   
%    \changes{1.8}{2004-08-26}{%
%      Corrections:  Minor bugs were fixed.
%      Introduction: Definition of "light" versions of all colors;
%                    Definition of three color ramps: 
%                        Red-Blue, Green-Red and Cold-Hot;
%                    New feature labels `bar' and `color'.}
%   
%    \changes{1.9}{2005-02-08}{%
%      Corrections:  Fixed TeXtopo incompatibility introduced with v1.8.
%                    (Thanks to Meike Schmedt)
%      Introduction: Implementation of HMMTOP topology prediction.                   
%                    `includeHMMTOP'
%                    `showonHMMTOP'
%                    `hideonHMMTOP'
%                    new `appearance' option {HMMTOP} with {internal}
%                                                          {external}
%                                                          {TM};
%                    new arrow look with scalable line thickness;
%                    new arrow option `ball';
%                    `frameblock'   colored frame around sequence block;
%                    `shortcaption' allows one to define short caption
%                                   versions for the List of Figures.}
%
%    \changes{1.10}{2005-03-29}{%
%      Corrections:  Sped up drawing of color scales and bar graphs by
%                    by more than 10fold!
%                    (Thanks, Christoph Gille, for asking for it)
%      Introduction: Definition of even lighter versions of all colors;
%                    implementation of a new labeling mode 'tint':                   
%                    `tintregion'
%                    `tintblock'
%                    `tintdefault';
%                    new `feature' option {restriction} for putting a
%                    triangle label pointing between two residues;
%                    data files for color scales and bar graphs can
%                    now contain 'NaN' (not a number) values
%                    (Also requested by Christoph Gille.)}
%
%    \changes{1.11}{2005-04-13}{%
%      Corrections:  Frames were drawn with the wrong height when
%                    separation lines were used. Fixed.
%                    Spacing between bar graph feature line and 
%                    sequence block was wrong after `bargraphstretch'.
%      Introduction: Additional optional parameter for feature rule
%                    thickness;                   
%                    additional optional parameters for feature box
%                    frame color and frame thickness;
%                    definition of three more color scales:
%                    {RedBlue}, {RedGreen}, and {HotCold};
%                    plotting of amino acid features as bar graphs
%                    or color scales:
%                    `hydrophobicity'
%                    `molweight'
%                    `charge';
%                    plotting of protein sequence conservation as
%                    bar graph or color scale:
%                    `conservation';
%                    separate command for stretching color scales:
%                    `colorscalestretch';
%                    color scales on consensus sequence according
%                    to sequence conservation.
%                    }
%
%    \changes{1.12}{2005-09-20}{%
%      Corrections:  Combination of 'setends' with regional labeling
%                    using 'shaderegion', 'frameblock', 'emphregion'
%                    or 'tintregion' produced incorrect output
%                    (thanks to Chris Page). Fixed.
%      Introduction: Optional colors for `showconsensus' foreground
%                    and background.
%                    }
%
%    \changes{1.13}{2006-02-23}{%
%      Corrections:  Helix symbols in feature lines were not drawn
%                    correctly if the standard Computer Modern Font
%                    was changed to another one, e.g. Palatino (thanks
%                    to Markus Heller). Fixed.
%                    Unintended gaps occurred due to numbers at the
%                    end of lines in Clustal W alignment files. Fixed.
%                    Frames were too tall when sequences were hidden
%                    or killed. Fixed.
%                    The limitations in the number of sequences per
%                    alignment have finally been overcome by a more
%                    restrictive use of counter variables.
%      Introduction: The numbering can now be displayed on both sides
%                    of the alignment with the optional parameter
%                    {leftright};
%                    TeXshade tries to guess the sequence type (protein 
%                    or nucleotide) if not defined by the user;
%                    Implementation of sequence logos:
%                    `showsequencelogo', `hidesequencelogo', 
%                    `namesequencelogo', `logostretch',
%                    `logocolor', `clearlogocolors',
%                    `showlogoscale', `hidelogoscale'
%                    `dofrequencycorrection', `undofrequencycorrection';
%                    The ruler numbering can now be rotated with
%                    `rotateruler' and back with `unrotateruler',
%                    this way every position can be numbered which is
%                    often wanted when e.g. sequence logos are plotted;
%                    the font family (sf, rm, tt) can be set for the
%                    ruler, e.g. `rulertt' or `setfamily{ruler}{tt}';
%                    `hideseqs' and `showseqs' in order to hide/show
%                    all sequences, esp. useful with sequence logos;
%                    `allowzero' and `disallowzero' - use (or do not 
%                    use) the number `0' in the sequence numbering as
%                    sometimes wanted in sequence logos;
%                    Implementation of a new way to visualize residues  
%                    which are characteristic for protein subfamilies, 
%                    i.e. subfamily logos:
%                    `showsubfamilylogo', `hidesubfamilylogo', 
%                    `namesubfamilylogo', 
%                    `setsubfamily', 
%                    `shownegatives', `hidenegatives'.
%                    }
%
%    \changes{1.14}{2006-05-11}{%
%      Introduction: `showrelevance', `hiderelevance',
%                    `relevance': commands to set a bit-value above
%                    which subfamily deviations are considered relevant
%                    and to label such positions in the subfamily logo.
%                    }
%
%    \changes{1.15}{2006-06-27}{%
%      Correction:   Logos can now be plotted with pdflatex; pstricks is
%                    not needed anymore.
%                    }
%
%    \changes{1.16}{2007-02-18}{%
%      Corrections:  TeXshade crashed when doing conservation
%                    calculations with sequences containing untypical
%                    residues symbols, such as X. Fixed.
%                    Shading of the reference sequence in diverse mode
%                    is now achieved with `conservedresidues' and
%                    `allmatchresidues' instead of `nomatchresidues'.
%      Introduction: `exportconsensus' produces a pymol script for
%                    coloring according to the TeXshade conservation
%                    calculation;
%                    `namerulerpos' allows one to change labels of
%                    the ruler individually;
%                    `hideblock' allows one to hide parts of the
%                    alignment (still in experimental stage!).
%      New home:     TeXshade, TeXtopo, and BioTeX have a new home:
%                    `www.pharmazie.uni-kiel.de/chem/Prof_Beitz/biotex.html'
%                    }
%
%    \changes{1.17}{2007-06-19}{%
%      Corrections:  .
%      Introduction: `allmatchspecial' now accepts an optional threshold
%                    percentage allowing one to set two levels off 
%                    conservation;
%                    the same effect can be achieved by using a
%                    number as an optional parameter in
%                    `shadingmode',
%                    or by setting an additional parameter in 
%                    `threshold';
%                    names can be displayed left or right of feature
%                    lines using:
%                    `showfeaturename', `showfeaturestylename';
%                    `hidefeaturename', `hidefeaturestylename';
%                    `hidefeaturenames', `hidefeaturestylenames';
%                    the color of such names can be changed with
%                    `featurenamescolor'
%                    `featurestylenamescolor'
%                    font styles can be set as usual, e.g.
%                    `setsize{featurenames}{large}' or
%                    `featurestylenamesrm' etc.
%                    }
%
%
%    \changes{1.18}{2008-04-15}{%
%      Corrections:  bug fixes (featurenames, ordering, numbering).
%      Introduction: T-Coffee shading in alignment, consensus and
%                    feature color scales and bar graphs
%                    `shadingmode[ASCII-file]{T-Coffee}'
%                    `includeTCoffee';
%                    two more feature lines on top and at bottom:
%                    `ttttop', `tttop', `bbbottom', and `bbbbottom';
%                    fusion of the `startnumber' and `setends' commands
%                    via optional parameters.
%                    }
%
%    \changes{1.19}{2009-03-09}{%
%      Corrections:  horizontal scaling of logo characters with `charstretch';
%                    further bug fixes.
%      Introduction: selection of residues depending on PDB file coordinates;
%                    works with `feature', `shaderegion', `shadeblock',
%                    `tintregion', `tintblock', `emphregion',
%                    `emphblock', `frameblock';
%                    3D selection possible around a point, along a line, or 
%                    within a plane;
%                    list of selected residues printable with `printPDBlist',
%                    or viewable during the TeX run with `messagePDBlist';
%                    show only selected residues in alignment with `setdomain',
%                    `domaingaprule' sets thickness of a domain separator rule,
%                    `domaingapcolors' sets fg/bg colors for the separator rule,
%                    these commands replace the unstable `hideblock' command.
%                    }
%
%    \changes{1.19a}{2009-06-11}{%
%      Corrections:  gaps are now treated correctly when setting domains.
%      Introduction: a special character can be displayed for stop-positions
%                    in protein sequences indicated by `*' in the input file;
%                    the stop character can be set using `stopchar'.
%                    }
%
%    \changes{1.20}{2009-10-05}{%
%      Corrections:  landscape mode works now correctly;
%                    annoying `Overfull \hbox ...' signals appearing after
%                    `residuesperline*' commands are ignored;
%                    domains work now with frequency corrected sequence logos.
%      Introduction: new commands for setting residue weight tables: 
%                    `weighttable' (with parameters: identity, structural,
%                                   PAM250, PAM100, and BLOSUM62),
%                    `gappenalty', and `setweight';
%                    similarity/identity percentage tables can be printed with
%                    `similaritytable' and values for specific sequence pairs 
%                    can be utilized with `percentsimilarity{seq1}{seq2}' and 
%                    `percentidentity{seq1}{seq2}'.
%                    }
%
%    \changes{1.20a}{2009-11-30}{%
%      Corrections:  identity/similarity tables work now without the need to
%                    set a conservation plot in the feature line.
%                    }
%
%    \changes{1.21}{2010-03-01}{%
%      Corrections:  `setdomain' ignored a changed `startnumber', corrected;
%                    further bug fixes. 
%      Introduction: in several commands sequence motifs can be given instead
%                    of numeral sequence stretch definitions;
%                    new type of sequence emphasis by lowercase characters:
%                    `lowerregion' and `lowerblock'.
%                    }
%
%    \changes{1.22}{2010-10-11}{%
%      Corrections:  minor bug fixes. 
%      Introduction: `defshadingcolors' defines and names sets of shading  
%                    colors for later re-use within the alignment;
%                    `changeshadingcolors' allows one to use different sets
%                    of shading colors in different sequence blocks.
%                    }
%
%    \changes{1.23}{2011-05-13}{%
%      Introduction: `hideallmatchpositions' removes all positions in the
%                    alignment where all residues match, i.e. a handy feature for
%                    diverse mode to condense the output to the relevant sites.
%                    }
%
%    \changes{1.24}{2011-12-01}{%
%      Introduction: `rulerspace' allows one to adjust the space between the
%                    ruler and the top or bottom sequence row;
%                    the textsize of the ruler numbering can now be set;
%                    Postscript color samples are shown in the manual.
%                    }
%
%    \changes{1.25}{2018-12-01}{%
%      Corrections:  aligments with many seqs produced wrong calculation
%                    of threshold shading, corrected. 
%      Introduction: new feature style |S-S| for disulfide bridges;
%                    hooks in top feature lines can be drawn down to alignment
%                    by adding an underscore in the definition, e.g. |{_,-,}|.
%                    }
%
%    \changes{1.26}{2021-04-01}{%
%      Corrections:  minor fixes;
%      Introduction: Chimera command file export for 'structure memes',
%                    i.e. mapping of sequence logo and subfamily logo data
%                    onto 3D protein structures;
%                    new commands `structurememe', `chimeraballScale',
%                    `chimeraaxisdistance', `memeRed', `memeBlue',
%                    `memeYellow', `memeBlack', `memeWhite',
%                    `memeStandardcolors', `chimerachain',
%                    `memelabelcutoff', `echostructurefile'.
%                    }
%
%    \changes{1.26a}{2022-08-29}{%
%      Corrections:  frames and line ends remain open to indicate
%                    continuity of the marked region; the reference to the
%                    structure meme paper was added;
%                     `par' line endings now treated correctly.}
%
%
%
% \CharacterTable
%  {Upper-case    \A\B\C\D\E\F\G\H\I\J\K\L\M\N\O\P\Q\R\S\T\U\V\W\X\Y\Z
%   Lower-case    \a\b\c\d\e\f\g\h\i\j\k\l\m\n\o\p\q\r\s\t\u\v\w\x\y\z
%   Digits        \0\1\2\3\4\5\6\7\8\9
%   Exclamation   \!     Double quote  \"     Hash (number) \#
%   Dollar        \$     Percent       \%     Ampersand     \&
%   Acute accent  \'     Left paren    \(     Right paren   \)
%   Asterisk      \*     Plus          \+     Comma         \,
%   Minus         \-     Point         \.     Solidus       \/
%   Colon         \:     Semicolon     \;     Less than     \<
%   Equals        \=     Greater than  \>     Question mark \?
%   Commercial at \@     Left bracket  \[     Backslash     \\
%   Right bracket \]     Circumflex    \^     Underscore    \_
%   Grave accent  \`     Left brace    \{     Vertical bar  \|
%   Right brace   \}     Tilde         \~}
%   
%
%    \newsavebox{\mybox}
%    \newenvironment{fmpage}[1][0.975\textwidth]{%
%       \begin{lrbox}{\mybox}\begin{minipage}{#1}}
%      {\end{minipage}\end{lrbox}\fbox{\usebox{\mybox}}}
%        
%    \parindent0mm
%            
%     
%    \title{The \TeXshade{} package\footnote{Please cite: Eric Beitz (2000), 
%           \TeX{}shade: 
%           shading and labeling multiple sequence alignments using \LaTeXe. 
%           \textit{Bioinformatics}: \textbf{16}, 135--139.}\\[2mm] \large
%           Typesetting \\ nucleotide and peptide alignments}
%    \author{Eric Beitz\footnote{University of Kiel, 
%            Pharmaceutical Chemistry, Gutenbergstrasse 8, 
%            D-24118 Kiel, Germany; 
%            send electronic mail to \texttt{ebeitz@pharmazie.uni-kiel.de};
%            for further information, updates and on-line documentation 
%            see my homepage at
%            \texttt{www.pharmazie.uni-kiel.de/chem/Prof\_Beitz/biotex.html} }}
%    \date{\small v1.26a; 2022/08/29\\[2pt]}
%    \maketitle
%    \begin{abstract}
%      Setting alignments of nucleotides and peptides for publication 
%      or presentation purposes is usually a time consuming two-step process. 
%      First, a scientific software is used for the calculation of the
%      alignment. This
%      is done in a few minutes. Then, in order to highlight special sequence
%      relationships and to label positions and regions of interest a 
%      second software with high quality output capability is needed.
%      Manipulating sequence alignments with standard word processing
%      or graphics programs takes its time---often several hours---and 
%      simple layout changes such as
%      re-breaking lines, say from 50 to 40 residues per line, 
%      elongate the working time considerably.
%      
%      \TeXshade{} is an alignment shading software  
%      written in \TeX/\LaTeX{} which can process
%      multiple sequence alignments in the MSF, ALN 
%      and FASTA file format.
%      It provides in addition to common shading algorithms special
%      shading modes featuring functional aspects, e.\,g.\ charge or
%      hydropathy, and a plenitude of commands for handling
%      shading colors, text styles, labels, legends and even allows 
%      the user to define completely new shading modes. \TeXshade{} 
%      combines highest flexibility and the habitual \TeX{} output 
%      quality---with reasonable time expenditure.
%      
%    \end{abstract}
%        
%    \thispagestyle{empty}
%    
%    \tableofcontents
%    \newpage
%    
%    \section{Package Overview}
%    
%    \label{over}
%   
%    After |texshade.ins| is run through \TeX{} the following files
%    should appear in the directory:
%    
%    \begin{tabbing}
%    \quad|texshade.sty|\quad\= the style file with all \TeXshade{} 
%                                      commands\\
%    \quad|texshade.def|\>      an example parameter file with the 
%                                      standard \\
%                              \>      parameter settings\\
%    \quad|AQPDNA.MSF|  \>      an example nucleotide alignment
%                                      (MSF-format)\\
%    \quad|AQPpro.MSF|  \>      an example protein alignment
%                                      (MSF-format)\\
%    \quad|AQP_TC.asc|  \>      an example T-Coffee shading file
%                                      (|score_ascii|-format)\\
%    \quad|AQP2spec.ALN|\>      a further protein alignment
%                                      (minimal ALN-file)\\
%    \quad|AQP1.phd|\>          secondary structure information
%                                      (PHD-format)\\
%    \quad|AQP1.top|\>          topology data extracted 
%                                      from |AQP1.phd|\\
%    \quad|AQP1_HMM.sgl|\>      topology information (single line,
%                                      HMMTOP-format)\\
%    \quad|AQP1_HMM.ext|\>      topology information (extended,
%                                      HMMTOP-format)\\
%    \quad|AQP_sequence_logo_1FX8.cmd|\\ \>      Chimera structure meme file
%                                     \\
%    \quad|AQP_subfamily_logo_1FX8.cmd|\\ \>      Chimera structure meme file
%                                     \\
%    \quad|meme.eps|\>       structure meme example image\\
%    \quad|standard.cod|\>      standard genetic code definitions\\
%    \quad|ciliate.cod|\>       ciliate macronuclear genetic code\\
%    \end{tabbing}
%    The alignment file examples as well as the topology data file are 
%    needed for \TeX{}ing this documentation
%    and can serve as illustrations for the MSF and ALN 
%    file format.
%    
%    The following subsections give an overview on the capabilities of
%    the \TeXshade{} package. All commands are described in detail 
%    later on.
%    
%    
%    \subsection{Version History}
%    
%    \textbf{v1.26(a)  2021/04/01 (2022/08/29)}
%    \medskip
%
%    \emph{Corrections}:  Frames at line ends remain open to indicate
%    continuity of the marked region, and |\par| line endings are now treated
%    correctly.
%
%    \emph{Introductions:}
%     Chimera command file export option for 'structure memes'
%    was introduced,
%     i.e. mapping of sequence logo and subfamily logo data
%     onto 3D protein structures.
%     The new commands |\structurememe|,
%     |\memeRed|, |\memeBlue|,
%     |\memeYellow|, |\memeBlack|, |\memeWhite|,
%     |\memeStandardcolors|, |\chimeraaxisdistance|, |\chimerachain|,  |\chimeraballScale|,
%     |\memelabelcutoff|, and |\echostructurefile|
%     handle the new feature.
%    \medskip
%
%    \textbf{v1.25  2018/01/17}
%    \medskip
%
%    \emph{Corrections}:  Aligments with many seqs produced wrong calculation
%    of threshold shading,\footnote{Noted by Kathryn Parker.} corrected; minor bug fixes. 
%
%    \emph{Introduction}: A new feature style |S-S| for disulfide bridges is
%    implemented; hooks in top feature lines can be drawn down to alignment
%    by adding an underscore in the definition, e.g. |{_,-,}|.
%    \medskip
%
%    \textbf{v1.24  2011/12/01}
%    \medskip
%
%    \emph{Introductions:} 
%    The space between the ruler and the top or bottom sequence row
%    can now be adjusted using |\rulerspace| and the textsize of the 
%    ruler numbering can be set with |\setsize| or |\rulerLarge| etc. 
%    Postscript color samples are shown in the manual next to the 
%    color names list.
%    \medskip
%
%    \textbf{v1.23  2011/05/13}
%    \medskip
%
%    \emph{Introductions:} 
%    In diverse mode sequence positions where all residues match do not
%    contain much information. A new command, 
%    |\hideallmatchpositions|,\footnote{Requested by Matt Russell.}
%    will remove all such positions from the alignment and hence condense
%    the output considerably.
%    \medskip
%
%    \textbf{v1.22  2010/10/11}
%    \medskip
%
%    \emph{Introductions:} 
%    Sets of shading colors can be defined and named using
%    |\defshadingcolors| (p.\ \pageref{Ldefshadingcolors}) and different
%    sets of colors can be selected for different sequence blocks by 
%    |\changeshadingcolors| (\ref{Lchangeshadingcolors}).
%    \medskip
%
%
%    \textbf{v1.21 2010/03/01}
%    \medskip
%
%    \emph{Corrections:} |\setdomain| now recognizes a changed |\startnumber|. 
%
%    \emph{Introductions:} 
%    (a) Sequence motifs (\ref{Lshaderegion}) can be given as a selection pattern
%    in sequence stretch labeling commands, such as |\shaderegion|,
%    |\emphregion| or |\frameblock|, instead of numeral stretch definitions.
%    
%    (b) A new type of sequence emphasis by lowercase letters was implemented: 
%    |\lowerregion| and |\lowerblock| (\ref{Llowerregion}).
%    \medskip
%
%    \textbf{v1.20a 2009/11/30}
%    \medskip
%
%    \emph{Corrections:}  Identity/similarity tables
%    work now without having to set a conservation plot in the feature
%    line. 
%    \medskip
%
%
%    \textbf{v1.20 2009/10/05}
%    \medskip
%
%    \emph{Corrections:}  Landscape mode works now correctly.
%    Annoying |Overfull \hbox ...| signals appearing after
%    |\residuesperline*| commands are ignored;
%    domains work now with frequency corrected sequence logos.
%
%    \emph{Introductions:} 
%    (a) New commands for setting residue weight tables were introduced 
%    |\weighttable| (with parameters: |identity|, |structural|,
%    |PAM250|, |PAM100|, and |BLOSUM62|),
%    |\gappenalty|, and |\setweight|.
%    (b) Similarity/identity percentage tables can be printed with
%    |\similaritytable|\footnote{Asked for by Giovanni M.\ Lesa.} and  
%    values for specific sequence pairs can be utilized with 
%    |\percentsimilarity{|\meta{seqref1}|}{|\meta{seqref2}|}| and 
%    |\percentidentity{|\meta{seqref1}|}{|\meta{seqref2}|}|.
%    \medskip  
%
%    \textbf{v1.19a 2009/06/11}
%    \medskip
%
%    \emph{Correction:} Gaps are now treated correctly when setting 
%    domains.
%    
%    \emph{Introduction:}
%    Stop-positions in protein sequences (|*| in the input file) can be
%    indicated by a special character (|\stopchar|) in the alignment 
%    output.\footnote{Suggestion by Yun He.}
%    \medskip  
%
%
%    \textbf{v1.19 2009/03/09}
%    \medskip
%
%    \emph{Correction:} Logo characters are now horizontally 
%    scalable using |\charstretch|. Minor bugs were fixed.
%    
%    \emph{Introduction:}
%    (a) Selection of residue positions was enhanced (\ref{Lshaderegion}) 
%    \TeXshade{} can
%    now select residues based on their 3D coordinates provided in a 
%    PDB structure file. 3D selection can be due to a certain distance 
%    around a point, along a line, or above and below a plane. It
%    works with |\feature|, |\shaderegion|, |\shadeblock|, |\tintregion|,
%    |\tintblock|, |\emphregion|, |\emphblock|, |\frameblock|.
%    (b) The list of selected residues is printable with |\printPDBlist| or 
%    viewable during the \TeX{} run with |\messagePDBlist|.
%    (c) |\hideblock| and related commands were replaced by |\setdomain|
%    (\ref{Lsetdomain}). 
%    This command will display only selected residues in the alignment.
%    Thickness and colors of a domain separator rule can be set using
%    |\domaingaprule| and |\domaingapcolors|.
%    \bigskip
%
%    \textbf{v1.18 2008/04/15}
%    \medskip
%
%    \emph{Correction:} several bug fixed concerning featurename
%    display, sequence ordering and numbering.
%    
%    \emph{Introduction:}
%    (a) T-Coffee shading information can be loaded and put on the
%    alignment.\footnote{Suggestion by Florian Mertes.} 
%    The conservation data can also be displayed in the 
%    consensus as well as feature color scales and bar plots.
%    (b) Two more feature lines were added on the top and at the
%    bottom (|ttttop|, |tttop|, |bbbottom|, |bbbbottom|).
%    (c) The |startnumber| and |setends| commands have been fused;
%    either command can set both, a new start number as well as
%    end definitions of the sequence section to be displayed.
%    \bigskip
%
%    \textbf{v1.17 2007/06/19}
%    \medskip
%    
%    \emph{Introduction:}\footnote{Asked for by Marat Kazanov.}
%    (a) A second threshold percentage was introduced in order to label
%    two levels of conservation in `identical' and `similar' mode. This is 
%    achieved by setting an optional parameter in |\threshold| or in 
%    |\allmatchspecial|, or by using a number as an optional parameter in 
%    |\shadingmode|. (b) The feature lines can be additionally labeled with
%    a name left or right of the feature. This is handled using
%    |\showfeaturename|, |\showfeaturestylename|,
%    |\hidefeaturename|, |\hidefeaturestylename|,
%    |\hidefeaturenames|, |\hidefeaturestylenames|.
%    The color of such names can be changed with
%    |\featurenamecolor|, |\featurestylenamecolor|,
%    |\featurenamescolor|, |\featurestylenamescolor|,
%    Font styles in feature names can be set as usual, e.g.
%    |\setsize{featurenames}{large}| or
%    |\featurestylenamesrm|.
%    \bigskip
%
%
%    \textbf{v1.16 2007/02/18}
%    \medskip
%    
%    \emph{Correction:} \TeXshade{} crashed when calculating conservation
%    using sequences with untypical residue characters, such as "X".
%    Fixed. The reference sequence in diverse mode can now be shaded with
%    |\conservedresidues| and, if active, |\allmatchresidues|.\footnote{For this
%    and suggesting |namerulerpos| credit to Marco Pasi.}
%
%    \emph{Introduction:}\footnote{Both extensions were suggested by Phillip Hahn.}
%    (a) A command was introduced, i.e. |\exportconsensus| 
%    which produces a pymol script file for coloring a 3D model according to 
%    \TeXshade's conservation calculation. (b) With |namerulerpos| labels of the
%    ruler can be exchanged by a string. (c) Various parts of the alignment 
%    can now be hidden by |\hideblock|. 
%
%    \emph{New home:} \TeXshade, \TeXtopo, and \BioTeX{} have a new home:
%    |www.pharmazie.uni-kiel.de/chem/Prof_Beitz/biotex.html|.
%    \bigskip
%
%
%    \textbf{v1.15 2006/06/27}
%    \medskip
%    
%    \emph{Correction:} Sequence and subfamily logos can now be plotted
%    with pdflatex; pstricks is not needed anymore.
%    \bigskip
%
%
%    \textbf{v1.14 2006/05/11}
%    \medskip
%    
%    \emph{Introduction:} In order to better recognize relevant positions
%    in a subfamily logo [14], a bit-value can now be set by |\relevance|
%    above which a deviation is considered relevant. Such positions
%    can be labeled with a symbol by |\showrelevance| and hidden
%    by |\hiderelevance|.
%    \bigskip
%
%
%    \textbf{v1.13 2006/02/23}
%    \medskip
%    
%    \emph{Corrections:} Helix symbols in feature lines were not drawn
%    correctly if the standard Computer Modern Font was changed to 
%    another one, e.g. Palatino.\footnote{Thanks to Markus Heller} 
%    Fixed. Unintended gaps occurred due to numbers at the
%    end of lines in Clustal W alignment files. Fixed. The limitations
%    in the number of sequences per alignment have finally been overcome
%    by a more restrictive use of counter variables.
%
%    \emph{Introductions:} (a) The numbering can now be displayed---in 
%    addition to left or right---on both sides of the alignment with 
%    the optional parameter |{leftright}| in the |\shownumbering| 
%    command (p.\pageref{Lshownumbering}). (b) TeXshade tries to guess 
%    the sequence type, i.\,e.\ protein or nucleotide, if not defined 
%    by the user. (c) Plotting of sequence logos has been implemented
%    (p.\pageref{Lshowsequencelogo}).
%    Logos can be shown in addition to or together with the consensus,
%    or alone without any alignment sequences. (d) The ruler numbering
%    can be rotated in order to make labeling of every position possible.
%    (e) A new way to visualize subfamily characteristics has been
%    implemented, i.e. subfamily logos (p.\pageref{Lshowsubfamilylogo}) [14].
%    \bigskip
%
%
%    \textbf{v1.12 2005/09/20}
%    \medskip
%    
%    \emph{Corrections:} When regional labeling with |\shaderegion|,
%    |\emphregion|, |\tintregion|, or |\frameblock| was combined with
%    |\setends| incorrect output was produced lacking the
%    labeling.\footnote{Discovered by Chris Page.} Other minor fixes.
%
%    \emph{Introductions:} An additional optional parameter for setting 
%    consensus colors was implemented in the |\showconsensus| command
%    (p.\pageref{Lshowconsensus}). This even allows one to use color 
%    scales illustrating sequence conservation in the consensus line.
%    \bigskip
%
%    \textbf{v1.11 2005/04/13}
%    \medskip
%    
%    \emph{Corrections:} Bounding boxes with |\frameblock| had a wrong
%    height when |\separationline|s were used. Other minor fixes.
%
%    \emph{Introductions:} (a) An additional parameter for setting 
%    individual bar and arrow thicknesses in feature lines has been 
%    introduced. (b) Additional parameters for setting the frame color
%    and thickness of boxes in feature lines have been implemented. (c)
%    Three more color scales have been defined: |RedBlue|, |RedGreen|,
%    and |HotCold|. (d) Plotting of amino acid features (|hydrophobicity|,
%    |molweight|, |charge|) as bar graphs or color scales. (e) Plotting
%    of protein sequence |conservation| as bar graph or color
%    scale\footnote{Ahmad Mirza asked for (e) and (f), great suggestion!}.
%    (f) Color scales can be used for shading the consensus sequence
%    according to protein sequence conservation.
%    (g) Separate command for stretching color scales |\colorscalestretch|.
%    \bigskip
%    
%    \textbf{v1.10 2005/03/29}
%    \medskip
%    
%    \emph{Corrections:} Plotting of color scales and bar graphs has
%    been sped up by more than a factor of 10.\footnote{This and (d)
%    I owe again to Christoph Gille.}
%
%    \emph{Introductions:} (a) More colors have been introduced, i.e. 
%    even lighter versions of the existing PostScript colors 
%    `LightLight' plus color name and `LightLightLight' plus color 
%    name. (b) Sequence stretches and blocks can be tinted for
%    labeling purposes |\tintreqion|, |\tintblock| and |\tintdefault|. 
%    (c) A new feature label style |{restriction}| has been introduced. 
%    (d) Java-typical `NaN' values are now allowed in data files for
%    bar graphs and color scales.
%    \bigskip
%    
%    
%
%    \textbf{v1.9 2005/02/08}
%    \medskip
%    
%    \emph{Corrections:} \TeXshade{} version 1.8 introduced an
%    incompatibility with \TeXtopo{}. This problem was identified
%    by Meike Schmedt and has been fixed.
%
%    \emph{Introductions:} (a) A short version of the figure caption
%    can now be defined for display in the list of figures\footnote{%
%    Meike, here you go \dots} |\shortcaption{|\meta{text}|}|. (b) A
%    colored frame can be drawn around a sequence block for labeling
%    purposes with the command |\frameblock|.\footnote{Alan Robinson, 
%    this is for you.} (c) A new look for feature arrows has been 
%    implemented with scalable line thickness and a new end style
%    `ball'. (d) HMMTOP topology predictions can
%    now be included for plotting feature lines with information on
%    the location of the transmembrane domains.\footnote{Implemented 
%    after a request by Steffen Moeller.} 
%    \bigskip
%
%    \textbf{v1.8 2004/08/26}
%    \medskip
%    
%    \emph{Corrections:} Only minor bugs were fixed.
%
%    \emph{Introductions:} (a) More colors have been designed, i.e. 
%    `light' versions of the existing PostScript colors. (b) 
%    Three color ramps in 5\% steps have been introduced: 
%    i) Blue-Red, ii) Green-Red and iii) Cold-Hot.
%    (c) Two new feature label styles |bar| and |color| have been 
%    introduced which allow one to display number
%    values as bar graphs or color scales along the
%    alignment\footnote{Inspired by Christoph Gille's {\tt STRAP}}. 
%    \bigskip
%
%
%    \textbf{v1.7 2004/01/05}
%    \medskip
%    
%    \emph{Corrections:} Several bugs were fixed.
%    In gaps the wrong character was plotted in `donotshade'
%    mode. Gaps were colored incorrectly when a single 
%    sequence was set as consensus. Another `donotshade' problem was 
%    solved which led to a halt of the LaTeX 
%    run\footnote{Thanks to Jeferson J.\ Arenzon and Naomi Siew}. 
%    Due to several requests, the gap and match labels in |diverse|
%    mode were switched (`|-|' in gaps; `|.|' at matching
%    positions) in order to follow convention.
%
%    \emph{Introduction:} \TeXshade{} speaks spanish (|\spanishlanguage|).
%    Necessary translations were contributed by Mikel Ega\~na Aranguren.
%    A new feature label style |helix| has been introduced.
%    \bigskip
%
%
%    \textbf{v1.6 2002/03/26}
%    \medskip
%    
%    \emph{Corrections:} The unnecessary restriction to the DVIPS
%    driver for |color.sty| has been removed\footnote{As suggested by
%    Eckhart Guth\"ohrlein.}. Any color.sty compatible
%    driver option can be given with the |\usepackage{texshade}| call
%    and is then passed to the |color| package. The `|\namecolor|' and 
%    `|\numbercolor|' commands do now support sequence 
%    lists.\footnote{Thanks to Denys Bashtovyy.}
%        
%    \emph{Introductions:} (a) The FASTA file format is supported by
%    \TeXshade{} as alignment inputs. (b) Two commands set the space
%    between sequence blocks either to be flexible (as so far) 
%    `|\flexblockspace|' or the be fixed `|\fixblockspace|'. (c) One 
%    can now refer to sequences by their name in addition to the number
%    in the input file. (d) Using
%    `|\firstcolumnDSSP|' and `|\secondcolumnDSSP|' one can choose 
%    which of the first to columns should refer to the sequence numbering 
%    (the second column remains default setting)\footnote{c and d were
%    suggested by Christoph Gille.}.
%    \bigskip
%    
%    \textbf{v1.5a 2001/03/08}
%    \medskip
%    
%    \emph{Corrections:} `X's in the alignment file caused a run-time
%    error. Fixed.
%        
%    \emph{Introductions:} (a) The vertical space between feature
%    lines can be controlled by four new commands: |\ttopspace|,
%    |\topspace|, |\bottomspace| and 
%    |\bbottomspace|\footnote{Suggested by Ulrike Folkers.}. (b) It is 
%    now easily possible to add a caption to the alignment with 
%    the |\showcaption| command. (c) \TeXshade{} stores the 
%    sequence lengths in the |.aux| file in order to have correct
%    breaks of the gaps after the sequences.
%    \bigskip
%
%
%    \textbf{v1.4\&4a 2000/9/12 \& 2000/10/3}
%    \medskip
%    
%    \emph{Introductions:} (a) The alignment legend can now be moved 
%    by the command `|\movelegend|'. (b) In commands with parameters 
%    that contain series of sequence numbers, e.\,g. |\orderseqs|, a
%    dash can be used, e.\,g. |{1-3,6-4,7}| instead of 
%    |{1,2,3,6,5,4,7}|.
%    \bigskip
%    
%
%    \textbf{v1.3a\&b 2000/7/28 \& 2000/7/30}
%    \medskip
%    
%    \emph{Introductions:} (a) It is now possible to force \TeXshade{} to 
%    display gap symbols before and after the actual sequence 
%    by the commands `|\showleadinggaps|' and `|\hideleadinggaps|'
%    (\ref{Lshowleadinggaps}).
%    (b) The sequence names input routine is now more tolerant concerning
%    special characters.
%    \bigskip
%    
%    \textbf{v1.3 2000/3/3}
%    \medskip
%    
%    \emph{Corrections:} Line scrambling occured when features where
%    set in the |ttop| row without a feature in the |top| row. Fixed.
%    The incompatible command `|\language|' with the |babel| package has been
%    replaced by `|\germanlanguage|' and `|\englishlanguage|'\footnote%
%    {Thanks to Eckhart Guth\"ohrlein.}.
%    
%    \emph{Introductions:} (a) Now, translations of sequence stretches
%    are possible. Either nucleotide or amino acid sources can be 
%    translated. This is done by the new |{translate}| option for the 
%    feature command. (b) The codons are defined by the new command
%    `|\codon|'. Complete codon sets can be loaded by `|\geneticcode|'.
%    (c) Further, the size and style of the nucleotide triplets of
%    backtranslations can be set by `|\backtranslabel|' and 
%    `|\backtranstext|'. (d) Two more feature counter styles were introduced:
%    `|\Romancount|' and `|\romancount|'. (e) \TeXshade{} is now 
%    compatible with \TeXtopo, a new \TeX{} software
%    for drawing and shading topology plots of membrane proteins.
%    \bigskip
%    
%    \textbf{v1.2a 1999/6/24 (not released)}
%    \medskip
%    
%    \emph{Minor corrections:} `|\namecolor|' and `|\numbercolor|' are 
%    now really correctly reordered. Brackets ( and ) are now allowed
%    in sequence names. The option |{case}| in `|\funcshadingstyle|' 
%    works now.
%    \bigskip
%    
%
%    \textbf{v1.2 1999/6/12}
%    \medskip
%    
%    \emph{Corrections:} (a) Functional group definitions of more than
%    seven groups produced an error when displaying group number
%    eight. These residues where skipped in the alignment. Fixed.
%    
%    \emph{Introductions:} (a) Protein secondary structure files in the DSSP,
%    STRIDE and PHD format can be included and displayed auto\-matically
%    within the alignment by `|\includeDSSP|' (and similar commands for 
%    STRIDE, PHDsec and PHDtopo, \ref{structure}).
%    (b) Which types of secondary structures are to be included or
%    skipped in the alignment is chosen by `|\showonDSSP|' and 
%    `|\hideonDSSP|' (and respective commands for STRIDE, PHDsec and PHDtopo).
%    (c) The appearance of the labels is defined by `|\appearance|'.
%    (d) Internal counters for repeatedly occuring structure types
%    can be activated by `|\numcount|', `|\alphacount|' and
%    `|\Alphacount|'. All commands are described in \ref{structure}.
%    \bigskip
%        
%
%    \textbf{v1.1 1999/5/26}
%    \medskip
%    
%    \emph{Corrections:} (a) The activation of `|emphregion|' lead to 
%    an em\-pha\-sized following alignment. This has been 
%    corrected. (b) `|\namecolor|' and `|\numbercolor|' were not 
%    reordered with the command `|orderseqs|'. Fixed. (c) Sequence 
%    gaps at the beginning or the end of a sequence, i.\,e. before
%    the first and after the last residue where labeled with the
%    gap symbol. Now these positions are left blank.
%    
%    \emph{Introductions:} (a) In order to treat the preceeding and 
%    sequence following gaps correctly, \TeXshade{} needs to know the
%    length of the sequences. Therefore, the command `|\seqlength|' was
%    introduced (\ref{seqlines}). (b) With `|\gapcolors|' (also 
%    \ref{seqlines}) the
%    color selection for gap symbols is independent from non-conserved
%    residues. (c) The divisions of the ruler where so far fixed to
%    10. Now, this value is changeable by `|\rulersteps|' (again 
%    \ref{seqlines}). (d) `|\hideresidues|' and `|\showresidues|' turn
%    off or on the residue names, i.\,e. one can choose between a
%    display of shaded boxes only or with letters in the boxes 
%    (\ref{kill}). (e) The changes (c) through (d) were necessary 
%    for the introduction of `|\fingerprint|'. This command allows one to
%    display the complete sequence in one line for an easy survey of
%    the alignment (\ref{fingerprint}).
%    \bigskip
%    
%    \textbf{v1.0 1999/5/12}
%    \medskip
%    
%    First release.
%    \bigskip
%    
%    
%    \subsection{\LaTeX{} basics}
%    
%    \subsubsection{Typesetting documents using \LaTeX}
%    
%    In order to use any of the macros provided by the 
%    \BioTeX-project 
%    (\TeXshade/\TeXtopo) efficiently a basic understanding of the \TeX{} 
%    typesetting system and its usage is required. Several books are
%    available on this topic, but a rather quick and easy introduction
%    is the \emph{Not so short introduction to \LaTeX}. This document
%    is available from all Comprehensive \TeX{} Archive Network 
%    (CTAN) servers,
%    e.\,g. from |ftp://ftp.dante.de/pub/tex/documentation/lshort/|,
%    in many different languages and formats besides \LaTeX{}, such 
%    as \textsc{PostScript} and on-line viewable PDF.
%    I also put a link from the \BioTeX{} (\TeXshade/\TeXtopo) homepage 
%    to the document collection
%    (|http://pharmazie.uni-kiel.de/chem/Prof_Beitz/BioTeX|).
%    
%    
%    \subsubsection{Memory shortness when using \TeX{}shade}
%    
%    If you are using \TeXshade{} to align several large sequences (about 1000
%    residues/sequence), LaTeX will probably stop compiling and quit with one
%    of the following messages: 
%    
%    |!\ TeX capacity exceeded, sorry [main memory size=384000]| 
%    
%    or 
%    
%    |!\ TeX capacity exceeded, sorry [stack size=300]|.
%    
%    \TeX{} allocates space for different kinds of internal variables. 
%    Setting alignments needs lots of memory,
%    usually more than for typesetting plain text.
%    Thus, the parameter settings of a standard \TeX{} installation might not
%    be sufficient for certain projects. This manifests
%    in \TeX{} error messages about insufficient memory  
%    and the setting process is interrupted. There is no reason to be 
%    concerned. The parameters can be set by hand. Unfortunately, 
%    each \TeX{} system hides its default parameter file in a different
%    place in the system. 
%    
%    
%    
%    
%
%    \subsection{System requirements} \label{require}
%    
%    \TeXshade{} requires \LaTeXe{} with |color.sty| and |graphics.sty|
%    for shading. For arrows in the feature line (p.\pageref{Lfeature})  
%    the AMS Math style is needed.
%    David Carlisle's |color.sty| is part of the Standard \LaTeX{}
%    `Graphics Bundle' [1]. This and the other packages can be downloaded 
%    from any \TeX{} archive, e.g.\ |ftp.dante.de|; usually they are 
%    included in a comprehensive \TeX{} installation. 
%
%    The |color| style allows one to use several |[|\meta{options}|]|, e.\,g.
%    |dvips|, |pdftex| or |dviwin|. These provide the commands which
%    different devices/programs need to display colored output. It is 
%    advisable to make yourself familiar with the |color.sty| manual. 
%    You should define a default driver in the file |color.cfg|.
%    Since there is no direct call of |color.sty| by the user, the 
%    option can be stated when \TeXshade{} is loaded, see next 
%    subsection. If no option is stated the |DVIPS| driver will be 
%    loaded. 
%     
%    With the |[dvips]| option the output DVI-file 
%    can be converted to \textsc{PostScript} using the |DVIPS| program
%    and can later be viewed or printed with the public domain 
%    {\sc GhostView} program which is
%    available for almost all computer platforms. Further, more and more
%    standard \TeX{} viewers are to a certain extent \textsc{PostScript} 
%    compatible. 
%    \bigskip
%    
%    \subsection{The \texttt{texshade} environment}
%    
%    \label{tsenvironment}
%    
%    The commands provided by the \TeXshade{} package are enabled by
%    the following command in the document header section:
%    \medskip
%    
%    \quad |\usepackage[|\meta{option}|]{texshade}|
%    
%    \medskip
%    Make sure that the file `|texshade.sty|' is present in a directory
%    searched by \TeX{} (see the installation notes in the file
%    `|texshade.txt|').
%    
%    The \meta{option} given here is passed to |color.sty| which 
%    handles the color commands for a particular output device, see 
%    previous subsection and the |color.sty| manual.
%    
%    The \TeXshade{} package provides only one single new environment:
%    |texshade|. This environment has one mandatory and
%    one optional argument, both of them designating file names which
%    must be present in a directory searched by \TeX. The
%    required file \meta{alignmentfile} contains the aligned nucleotide 
%    or peptide sequences
%    (see section~\ref{alignfilestruc}). This file is needed, because 
%    \TeXshade{} does no alignment by 
%    itself, it has to take a preprocessed alignment as input.
%    The optional file is a parameter file (section~\ref{paramfilestruc})
%    with definitions for the
%    customized calculation of the consensus, special sequence features
%    or labels etc. In this parameter file all \TeXshade{} commands 
%    which are allowed in the |texshade| environment can be used and are 
%    fully functional.
%    Within the environment further \TeXshade{} commands can be given
%    to replace or complete settings from the parameter file.
%    
%    Thus, setting an alignment with \TeXshade{} is as simple as
%    this:
%    
%    \begin{quote}
%    |\begin{texshade}[|\meta{parameterfile}|]|
%         |{|\meta{alignmentfile}|}|
%         
%    \quad\emph{further \emph{\TeXshade} commands, if needed}
%    
%    |\end{texshade}|
%    \end{quote}
%    
%    \subsection{Shading modes predefined in this package}
%    
%    \subsubsection{Identity mode}
%    
%    \label{ident}
%     
%    This basic type of shading is provided by almost any alignment
%    program. All identical residues at a position are shaded if the
%    number of matching residues is higher than a given threshold
%    (default is 50\%).\medskip
%    
%    \begin{texshade}{AQPpro.MSF}
%       \setends{1}{80..112}
%       \hideconsensus
%    \end{texshade}
%
%    Code:\medskip
%
%    \vbox{%
%    \begin{verbatim}    
%    \begin{texshade}{AQPpro.MSF}
%       \setends{1}{80..112}
%       \hideconsensus
%    \end{texshade}
%    \end{verbatim}  
%    }  
%    
%    Quite uncommon for an alignment shading program is the possibility
%    to display only selected sequence domains, e.\,g.\ to eliminate uninteresting
%    positions from the output:
%    \medskip
%    
%    \begin{texshade}{AQPpro.MSF}
%       \setdomain{1}{80..90,100..110,120..130}
%       \showruler{1}{top}
%       \hidenumbering
%       \hideconsensus
%    \end{texshade}
%
%    Code:\medskip
%
%    \vbox{%
%    \begin{verbatim}    
%    \begin{texshade}{AQPpro.MSF}
%       \setdomain{1}{80..90,100..110,120..130}
%       \showruler{1}{top}
%       \hidenumbering
%       \hideconsensus
%    \end{texshade}
%    \end{verbatim}  
%    }  
%    
%    This goes even further. You can have \TeXshade{} select positions
%    based on the 3D coordinates provided by a PDB file, e.\,g.\ show
%    all residues that are within an 8 \AA{} radius around the
%    $\alpha$-carbon of the residue at position 81:
%    \bigskip
%    
%    \begin{texshade}{AQPpro.MSF}
%       \setdomain{1}{75..86,103..103,148..148,151..153,155..156,193..195,218..219,222..222}
%       \showruler{top}{1} \rulersteps{1}
%       \hidenumbering
%       \hideconsensus
%    \end{texshade}
%
%    Code:\medskip
%
%    \vbox{%
%    \begin{verbatim}    
%    \begin{texshade}{AQPpro.MSF}
%       \setdomain{1}{point[8]:1J4N.pdb,81[CA]}
%       \showruler{top}{1} \rulersteps{1}
%       \hidenumbering
%       \hideconsensus
%    \end{texshade}
%    \end{verbatim}  
%    }  
%    
%    If you like, positions where conservation is very high (here $\ge$ 80\%) can
%    be shaded in a special color and the consensus can be shown with
%    or without shading according to the degree of conservation:
%    \medskip\label{shadecons}
%    
%    \begin{texshade}{AQPpro.MSF}
%       \threshold[80]{50}
%       \setends{1}{80..112}
%       \showconsensus[ColdHot]{bottom}
%       \defconsensus{.}{lower}{upper}
%       \showlegend
%    \end{texshade}
%    
%    Code:\medskip
%
%    \vbox{%
%    \begin{verbatim}    
%    \begin{texshade}{AQPpro.MSF}
%       \threshold[80]{50}
%       \setends{1}{80..112}
%       \showconsensus[ColdHot]{bottom}
%       \defconsensus{.}{lower}{upper}
%       \showlegend
%    \end{texshade}
%    \end{verbatim}}
%    
%    \subsubsection{Similarity mode}
%    
%    \label{similar}
%    
%    In many cases it is expedient---mostly when comparing protein
%    sequences---to shade also residues
%    which are not identical but similar to the consensus sequence.
%    Consider a position where three out of five residues are basic
%    arginines and two more residues are also basic but lysines.
%    In similarity mode \TeXshade{} shades similar residues in a different 
%    color to distinguish them from the consensus residue. Even when 
%    none of the residues alone reaches the 
%    threshold but a group of similar residues does these are shaded
%    in the `similarity' color. This case is given for instance
%    when at a position in a five sequence alignment two aliphatic 
%    valines and two also aliphatic isoleucins are present and the 
%    threshold is set to 50\%. Neither residue exceeds this percentage 
%    but as a group of similars they do. 
%    
%    In grayscale printouts some colors of the following alignment may appear 
%    undistinguishable. Don't worry if you usually use grayscale---all 
%    colors/grays can be selected freely (see \ref{colors}).
%    \medskip 
%    
%    \begin{texshade}{AQPpro.MSF}
%       \shadingmode{similar}
%       \threshold[80]{50}
%       \setends{1}{80..112}
%       \hideconsensus
%       \feature{top}{1}{93..93}{fill:$\downarrow$}{first case (see text)}
%       \feature{bottom}{1}{98..98}{fill:$\uparrow$}{second case (see text)}
%    \end{texshade}    
%    
%    Code:\medskip
%
%    \vbox{%
%    \begin{verbatim}    
%    \begin{texshade}{AQPpro.MSF}
%       \shadingmode{similar}
%       \threshold[80]{50}
%       \setends{1}{80..112}
%       \hideconsensus
%       \feature{top}{1}{93..93}{fill:$\downarrow$}{first case (see text)}
%       \feature{bottom}{1}{98..98}{fill:$\uparrow$}{second case (see text)}
%    \end{texshade}    
%    \end{verbatim}}
%    
%    Probably you know 
%    this kind of shading from the public domain program 
%    |BoxShade| 
%    by \textsc{Kay Hofmann} or from the Macintosh version 
%    |MacBoxShade| by \textsc{Michael D. Barron}. \TeXshade{}
%    provides the same functionality---and goes truly beyond---for the 
%    \TeX{} community.
%
%    
%    \subsubsection{T-Coffee shading}
%    
%    \label{TCoffee}
%    
%    \TeXshade{}'s capabilities of calculating alignment shadings are 
%    limited. |T-Coffee| (|www.tcoffee.org|) is a sophisticated alignment/shading 
%    software. You can apply shading from |T-Coffee| in \TeXshade{}
%    by loading the shading information file (|score_ascii|) generated by
%    |T-Coffee|.
%    \medskip 
%    
%
%    \begin{texshade}{AQPpro.MSF}
%       \shadingmode[AQP_TC.asc]{T-Coffee}
%       \setends{1}{30..63}
%       \feature{top}{1}{30..63}{color:conservation[T-Coffee]}{}
%       \showfeaturestylename{top}{feat-cons}
%       \showconsensus{bottom}
%    \end{texshade}    
%    
%    Code:\medskip
%
%    \vbox{%
%    \begin{verbatim}    
%    \begin{texshade}{AQPpro.MSF}
%       \shadingmode[AQP_TC.asc]{T-Coffee}
%       \setends{1}{30..63}
%       \feature{top}{1}{30..63}{color:conservation[T-Coffee]}{}
%       \showfeaturestylename{top}{feat-cons}
%       \showconsensus{bottom}
%    \end{texshade}    
%    \end{verbatim}}
%    
%    
%    
%    
%    
%    
%    
%    \subsubsection{Diversity mode}
%    
%    \label{diverse}
%    
%    Contrary to the above described modes this shading style displays
%    sequence differences. Thus, it is most suitable for comparing very
%    similar sequences, e.\,g.\ species variants of a protein.
%    
%    One sequence is used as consensus. 
%    Matching residues in other sequences are blanked out, 
%    mismatches are shown in lowercase.
%    \medskip
%        
%    \begin{texshade}{AQP2spec.ALN}
%       \shadingmode{diverse}
%       \setends{1}{77..109} \residuesperline*{33}
%       \featureslarge
%       \feature{top}{1}{77..109}{}{AQP2 species variants}
%       \namesrm\namessl
%       \hidenumbering
%       \showruler{top}{1}
%       \shownames{left}
%       \nameseq{1}{Bos taurus}
%       \nameseq{2}{Canis familiaris}
%       \nameseq{3}{Dugong dugong}
%       \nameseq{4}{Equus caballus}
%       \nameseq{5}{Elephas maximus}
%       \frameblock{1}{82..82,106..106}{Red[1pt]}
%    \end{texshade}\label{frame}
%    
%    Code:\medskip
%
%    \vbox{%
%    \begin{verbatim}    
%    \begin{texshade}{AQP2spec.ALN}
%       \shadingmode{diverse}
%       \setends{1}{77..109}
%       \featureslarge
%       \feature{top}{1}{77..109}{}{AQP2 species variants}
%       \namesrm\namessl
%       \hidenumbering\showruler{top}{1}
%       \shownames{left}
%       \nameseq{1}{Bos taurus}
%       \nameseq{2}{Canis familiaris}
%       \nameseq{3}{Dugong dugong}
%       \nameseq{4}{Equus caballus}
%       \nameseq{5}{Elephas maximus}
%       \frameblock{1}{82..82,106..106}{Red[1pt]}
%    \end{texshade}\label{frame}
%    \end{verbatim}}   
%    
%    
%    \subsubsection{Functionality modes}
%    
%    \label{func}
%    
%    Displaying functional peptide similarities is one of \TeXshade's
%    strong capabilities. Six functional shading modes are predefined;
%    further user specific modes can easily be created. The examples 
%    may not look very impressive when printed in grayscale. Enjoy 
%    them on your screen or use color printouts. As mentioned before,
%    all colors can be changed to others or to grays without restrictions
%    (see chapter \ref{colors}).
%    
%    \begin{itemize}
%    \item [\textbf{charge}:] residues which are charged at physiological pH
%                    (7.4) are shaded if their number at a position
%                    is higher than the threshold \label{charge}
%                    
%    \begin{texshade}{AQPpro.MSF}
%       \residuesperline*{34}
%       \shadingmode[charge]{functional}
%       \setends{1}{138..170}
%       \feature{top}{3}{153..165}{bar[-50,50]:-50,-45,%
%          -40,-30,-20,-10,0,10,20,30,40,45,50}{}
%       \feature{top}{3}{167..186}{color:5,10,15,20,25,30,35,%
%          40,45,50,55,60,65,70,75,80,85,90,95,100[ColdHot]}{}
%       \showlegend
%    \end{texshade}
%    
%    Code:\medskip
%
%    \vbox{%
%    \begin{verbatim}    
%    \begin{texshade}{AQPpro.MSF}
%       \shadingmode[charge]{functional}
%       \setends{1}{138..170}
%       \feature{top}{3}{153..165}{bar[-50,50]:-50,-45,%
%          -40,-30,-20,-10,0,10,20,30,40,45,50}{}
%       \feature{top}{3}{167..186}{color:5,10,15,20,25,30,35,%
%          40,45,50,55,60,65,70,75,80,85,90,95,100[ColdHot]}{}
%       \showlegend
%    \end{texshade}
%    \end{verbatim}}
%    
%    \item [\textbf{hydropathy}:] discrimination between acidic and
%                        basic, polar uncharged and hydrophobic nonpolar residues
%                        \label{hydro}
%    \bigskip    
%    
%    \begin{texshade}{AQPpro.MSF}
%       \residuesperline*{34}
%       \shadingmode[hydropathy]{functional}
%       \feature{top}{1}{158..163}{brace}{tinted}
%       \tintblock{1}{158..163}
%       \feature{top}{1}{QLVLC}{brace}{lowercased}
%       \lowerblock{1}{QLVLC}
%       \setends{1}{138..170}
%       \showlegend
%    \end{texshade}
%    
%    Code:\medskip
%
%    \vbox{%
%    \begin{verbatim}    
%    \begin{texshade}{AQPpro.MSF}
%       \shadingmode[hydropathy]{functional}
%       \feature{top}{1}{158..163}{brace}{tinted}
%       \tintblock{1}{158..163}
%       \feature{top}{1}{QLVLC}{brace}{lowercased}
%       \lowerblock{1}{QLVLC}
%       \setends{1}{138..170}
%       \showlegend
%    \end{texshade}
%    \end{verbatim}}    
%    
%    
%    \item [\textbf{structure}:] displays the potential
%                        localization within the tertiary structure of
%                        the protein \label{struc}
%    \bigskip    
%    
%    \begin{texshade}{AQPpro.MSF}
%       \residuesperline*{34}
%       \shadingmode[structure]{functional}
%       \setends{1}{138..170}
%       \feature{top}{1}{138..157}{box[Blue,Red][0.5pt]: %
%           $\alpha$-helix[Yellow]}{transmembrane domain 4}
%       \feature{top}{1}{158..163}{translate[Blue]}{}
%       \backtranslabel{oblique}
%       \feature{bottom}{1}{[DE]RXXR[DE]}{brace[Blue]}{loop D [Blue]}
%       \feature{top}{1}{164..170}{o->[Red]}{trans. dom. 5}
%       \showlegend
%    \end{texshade}
%          
%    Code:\medskip
%
%    \vbox{%
%    \begin{verbatim}    
%    \begin{texshade}{AQPpro.MSF}
%       \shadingmode[structure]{functional}
%       \setends{1}{138..170}
%       \feature{top}{1}{138..157}{box[Blue,Red][0.5pt]: %
%           $\alpha$-helix[Yellow]}{transmembrane domain 4}
%       \feature{top}{1}{158..163}{translate[Blue]}{}
%       \backtranslabel{oblique}
%       \feature{bottom}{1}{[DE]RXXR[DE]}{brace[Blue]}{loop D [Blue]}
%       \feature{top}{1}{164..170}{o->[Red]}{trans. dom. 5}
%       \showlegend
%    \end{texshade}
%    \end{verbatim}}   
%          
%    
%    \item [\textbf{chemical}:] residues are shaded due to chemical 
%                      properties of
%                      their functional groups \label{chem}
%    \bigskip    
%                      
%    \begin{texshade}{AQPpro.MSF}
%       \residuesperline*{34}
%       \shadingmode[chemical]{functional}
%       \setends{1}{138..170}
%       \showlegend
%    \end{texshade}
%          
%    Code:\medskip
%
%    \vbox{%
%    \begin{verbatim}    
%    \begin{texshade}{AQPpro.MSF}
%       \shadingmode[chemical]{functional}
%       \setends{1}{138..170}
%       \showlegend
%    \end{texshade}
%    \end{verbatim}}    
%          
%          With |\shadeallresidues| the threshold is ignored and
%          all residues are shaded due to their group assignment.
%          This is \emph{not} identical to a threshold of 0\%
%          where only the majority group would be shaded. See the
%          difference:
%    \bigskip    
%          
%    \begin{texshade}{AQPpro.MSF}
%       \residuesperline*{34}
%       \shadingmode[chemical]{functional}
%       \setends{1}{138..170}
%       \shadeallresidues
%    \end{texshade}
%          
%    Code:\medskip
%
%    \vbox{%
%    \begin{verbatim}    
%    \begin{texshade}{AQPpro.MSF}
%       \shadingmode[chemical]{functional}
%       \setends{1}{138..170}
%       \shadeallresidues
%    \end{texshade}
%    \end{verbatim}}    
%          
%    
%    \item [\textbf{rasmol}:] similar to |[chemical]| but with
%                      shading following the rasmol
%                      color scheme \label{ras}
%    \bigskip
%    \bigskip
%    \bigskip
%                      
%    \begin{texshade}{AQPpro.MSF}
%       \residuesperline*{34}
%       \shadingmode[rasmol]{functional}
%       \setends{1}{138..170}
%       \showruler{top}{1}
%       \rulersteps{1}
%       \namerulerpos{150}{site A[Red]}
%       \namerulerpos{155}{site B[Green]}
%       \shadeallresidues
%       \showlegend
%    \end{texshade}
%          
%    Code:\medskip
%
%    \vbox{%
%    \begin{verbatim}    
%    \begin{texshade}{AQPpro.MSF}
%       \shadingmode[rasmol]{functional}
%       \setends{1}{138..170}
%       \showruler{bottom}{1}
%       \rulersteps{1}
%       \namerulerpos{150}{site A[Red]}
%       \namerulerpos{155}{site B[Green]}
%       \shadeallresidues
%       \showlegend
%    \end{texshade}
%    \end{verbatim}}    
%          
%          
%          
%    \item [\textbf{standard area}:] this shading displays the 
%                      differences in the surface 
%                      area \label{starea}
%                      of the different amino acid's sidechains
%    \bigskip    
%                      
%    \begin{texshade}{AQPpro.MSF}
%       \residuesperline*{34}
%       \shadingmode[standard area]{functional}
%       \setends{1}{138..170}
%       \showlegend
%       \shadeallresidues
%    \end{texshade}
%          
%    Code:\medskip
%
%    \vbox{%
%    \begin{verbatim}    
%    \begin{texshade}{AQPpro.MSF}
%       \shadingmode[standard area]{functional}
%       \setends{1}{138..170}
%       \showlegend
%       \shadeallresidues
%    \end{texshade}
%    \end{verbatim}}    
%          
%    \item [\textbf{accessible area}:] \label{accarea} 
%                     here, the surface area which can
%                     be accessed by solvent molecules is used as a 
%                     basis for shading; low accessibility means 
%                     hydrophobic (i.\,e.\ strongly buried 
%                     residues), whereas highly accessible 
%                     sidechains are hydrophilic (compare to
%                     \textbf{hydropathy} and \textbf{structure})
%    \bigskip    
%                   
%    \begin{texshade}{AQPpro.MSF}
%       \residuesperline*{34}
%       \shadingmode[accessible area]{functional}
%       \setends{1}{138..170}
%       \showlegend
%       \feature{top}{1}{138..157,164..170}{helix}{membr.}
%       \feature{top}{1}{158..163}{---}{loop}
%       \featurerule{1mm}
%       \shadeallresidues
%    \end{texshade}
%
%    Code:\medskip
%
%    \vbox{%
%    \begin{verbatim}    
%    \begin{texshade}{AQPpro.MSF}
%       \shadingmode[accessible area]{functional}
%       \setends{1}{138..170}
%       \showlegend
%       \feature{top}{1}{138..157,164..170}{helix}{membr.}
%       \feature{top}{1}{158..163}{---}{loop}
%       \featurerule{1mm}
%       \shadeallresidues
%    \end{texshade}
%    \end{verbatim}}    
%
%    \end{itemize}
%    
%    
%    
%    \subsection{Bar graphs and color scales}
%    
%   \label{graphs}
%    
%    Amino acid properties, such as hydrophobicity, molecular weight,
%    or charge can be shown as bar graphs or color scales along the
%    alignment. Further, the degree of protein sequence conservation 
%    can be indicated. As an example, in the following  
%    aquaporin alignment plots of residue conservation (bars, top), 
%    are shown as well as properties of the AQP1 sequence: charge (scale, top), 
%    molecular weight are shown (scale, bottom), and hydrophobicity (bars, bottom).
%    
%    
%    \begin{texshade}{AQPpro.MSF}
%       \residuesperline*{34}
%       \setends{1}{138..170}
%       \feature{ttop}{1}{138..170}{bar:conservation}{}
%       \showfeaturestylename{ttop}{conserv.}
%       \ttopspace{-\baselineskip}
%       \feature{top}{1}{138..170}{color:charge}{}
%       \showfeaturestylename{top}{charge}
%       \feature{bottom}{1}{138..170}{color:molweight[ColdHot]}{}
%       \showfeaturestylename{bottom}{weight}
%       \bbottomspace{-\baselineskip}
%       \feature{bbottom}{1}{138..170}{bar:hydrophobicity[Red,Gray10]}{}
%       \showfeaturestylename{bbottom}{hydrophob.}
%       \bargraphstretch{3}
%       \featurestylenamescolor{Red}
%       \featurestylenamesrm  \featurestylenamesit
%       \hideconsensus
%    \end{texshade}
%          
%    Code:\medskip
%
%    \vbox{%
%    \begin{verbatim}    
%    \begin{texshade}{AQPpro.MSF}
%       \setends{1}{138..170}
%       \feature{ttop}{1}{138..170}{bar:conservation}{}
%       \showfeaturestylename{ttop}{conserv.}
%       \ttopspace{-\baselineskip}
%       \feature{top}{1}{138..170}{color:charge}{}
%       \showfeaturestylename{top}{charge}
%       \feature{bottom}{1}{138..170}{color:molweight[ColdHot]}{}
%       \showfeaturestylename{bottom}{weight}
%       \bbottomspace{-\baselineskip}
%       \feature{bbottom}{1}{138..170}{bar:hydrophobicity[Red,Gray10]}{}
%       \showfeaturestylename{bbottom}{hydrophob.}
%       \bargraphstretch{3}
%       \featurestylenamescolor{Red}
%       \featurestylenamesrm  \featurestylenamesit
%       \hideconsensus
%    \end{texshade}
%    \end{verbatim}}    
%    \medskip      
%    
%    The degree of similarity and identity between all sequences in the alignment 
%    can be shown in a table using |\similaritytable| (\ref{simtable}) outside the |texshade|
%    environment: \label{simtableEx}
%    
%    \DeleteShortVerb{\|}
%    \begin{center}\similaritytable\end{center}
%    \MakeShortVerb{\|}
%    
%    
%    \subsection{Secondary structures}
%    
%   \label{sec}
%    
%    Predicted protein secondary structures in the DSSP, STRIDE
%    PHD or HMMTOP file format can be included and displayed in the
%    alignment. As an example, the following few commands show an 
%    aquaporin alignment with the PHD topology data for aquaporin 
%    type 1 (top sequence).
%    
%    \vbox{%
%          \begin{verbatim}
%          \begin{texshade}{AQPpro.MSF}
%             \shadingmode[allmatchspecial]{similar}
%             \includePHDtopo{1}{AQP1.phd}
%          \end{texshade}
%          \end{verbatim}
%    }         
%    
%    Abbr.: \emph{int.} -- internal; \emph{ext.} -- external; \emph{TM} --
%    transmembrane domain
%    
%          \begin{texshade}{AQPpro.MSF}
%             \shadingmode[allmatchspecial]{similar}
%             \includePHDtopo{1}{AQP1.phd}
%          \end{texshade}
%    
%    \subsection{Sequence fingerprints}
%    
%    \label{finger}
%    
%    To gain a quick overview of sequence similarities or properties 
%    the |\fingerprint| command has been implemented. It can depict the 
%    complete sequence in one single line. The residues are presented 
%    as colored vertical lines. The implementation of this kind of output
%    was inspired by a publication by \textsc{Kai-Uwe Fr\"ohlich} [6].
%    \medskip
%    
%    \begin{texshade}{AQPpro.MSF}
%       \shadingmode[allmatchspecial]{similar}
%       \shadingcolors{grays}
%       \fingerprint{360}
%       \showlegend
%       \feature{top}{1}{13..36,51..68,94..112,138..156,%
%                        165..185,211..232}{,-,}{TM}
%    \end{texshade}
%    
%    Code:\medskip
%
%    \vbox{%
%    \begin{verbatim}    
%    \begin{texshade}{AQPpro.MSF}
%       \shadingmode[allmatchspecial]{similar}
%       \shadingcolors{grays}
%       \fingerprint{360}
%       \showlegend
%       \feature{top}{1}{13..36,51..68,94..112,138..156,%
%                        165..185,211..232}{,-,}{TM}
%    \end{texshade}
%    \end{verbatim}}    
%    
%    The higher the similarity the darker the vertical lines. In this 
%    overview it becomes obvious that the transmembrane regions of the 
%    aquaporin isoforms are most conserved.
%    \medskip
%    
%    A fingerprint of charge distribution on different aquaporins is shown. 
%    below. Sequence gaps can be left blank (example above) or drawn as lines
%    between the sequence blocks.
%    
%    \begin{texshade}{AQPpro.MSF}
%       \shadingmode[charge]{functional}
%       \shadeallresidues
%       \fingerprint{360}
%       \gapchar{rule}
%       \showlegend
%    \end{texshade}
%
%
%    Code:\medskip
%
%    \vbox{%
%    \begin{verbatim}    
%    \begin{texshade}{AQPpro.MSF}
%       \shadingmode[charge]{functional}
%       \shadeallresidues
%       \fingerprint{360}
%       \gapchar{rule}
%       \showlegend
%    \end{texshade}
%    \end{verbatim}}    
%
%    
%    \subsection{Sequence logos}
%    
%    \label{logo}
%    
%    Sequence logos represent the information content of the aligned 
%    sequences at a position in bit (max.\ 2 bit for DNA, i.\,e.
%    log$_2$4, and 4.322 bit for proteins, i.\,e. log$_2$20) and the 
%    relative frequency of a base or amino acid at this 
%    position [7]. Thus, more information is contained in logos than in
%    a standard consensus sequence.
%    The example below shows a DNA sequence alignment with the logo on the  
%    top.
%
%    It must be remarked that a logo from only five sequences does not
%    produce meaningful results - it rather illustrates the technique.
%
%    \medskip
%    
%    \begin{texshade}{AQPDNA.MSF}
%       \setends{1}{414..443}
%       \showsequencelogo{top}
%    \end{texshade}
%    
%    Code:\medskip
%
%    \vbox{%
%    \begin{verbatim}    
%    \begin{texshade}{AQPDNA.MSF}
%       \setends{1}{414..443}
%       \showsequencelogo{top}
%    \end{texshade}
%    \end{verbatim}}    
%    
%    
%    Next, only the logo of a protein alignment is displayed plus the 
%    degree of sequence conservation as a color scale in the consensus 
%    line. Note, that the full functionality of the feature lines remains.
%    \medskip
%    
%    \begin{texshade}{AQPpro.MSF}
%       \setends{AQP3.PRO}{203..235}
%       \showsequencelogo{top} \showlogoscale{leftright}
%       \hideseqs
%       \residuesperline*{33}
%       \defconsensus{{$\bullet$}}{{$\bullet$}}{{$\bullet$}}
%       \showconsensus[ColdHot]{bottom}
%       \nameconsensus{conservation} \namessf\namessl
%       \showruler{bottom}{AQP3.PRO} \rulersteps{1}
%       \feature{top}{AQP3.PRO}{208..210}{---}{NPA}
%       \feature{top}{AQP3.PRO}{211..219}{helix}{}
%       \feature{top}{AQP3.PRO}{220..232}{brace}{loop E}
%       \feature{top}{AQP3.PRO}{233..235}{helix}{TM6}
%       \feature{bottom}{AQP3.PRO}{203..235}{brace}{1-step numbering}
%    \end{texshade}
%    
%    Code:\medskip
%
%    \vbox{%
%    \begin{verbatim}    
%    \begin{texshade}{AQPpro.MSF}
%       \setends{AQP3.PRO}{203..235}
%       \showsequencelogo{top} \showlogoscale{leftright}
%       \hideseqs
%       \residuesperline*{33}
%       \defconsensus{{$\bullet$}}{{$\bullet$}}{{$\bullet$}}
%       \showconsensus[ColdHot]{bottom}
%       \nameconsensus{conservation} \namessf\namessl
%       \showruler{bottom}{AQP3.PRO} \rulersteps{1}
%       \feature{top}{AQP3.PRO}{208..210}{---}{NPA}
%       \feature{top}{AQP3.PRO}{211..219}{helix}{}
%       \feature{top}{AQP3.PRO}{220..232}{brace}{loop E}
%       \feature{top}{AQP3.PRO}{233..235}{helix}{TM6}
%       \feature{bottom}{AQP3.PRO}{203..235}{brace}{1-step numbering}
%    \end{texshade}
%    \end{verbatim}} 
%    
%    The same logo is shown below but with frequency correction turned
%    on (|\dofrequencycorrection|), see p.\pageref{Lshowsequencelogo}.
%    This takes into account the difference between the amino acid 
%    distribution in the alignment and the equal distribution of
%    5\% for each residue.
%    \medskip
%
%    \begin{texshade}{AQPpro.MSF}
%       \setends{AQP3.PRO}{203..235}
%       \showsequencelogo{top} \showlogoscale{leftright}
%       \hideseqs
%       \residuesperline*{33}
%       \defconsensus{{$\bullet$}}{{$\bullet$}}{{$\bullet$}}
%       \showconsensus[ColdHot]{bottom}
%       \nameconsensus{conservation} \namessf\namessl
%       \showruler{bottom}{AQP3.PRO} \rulersteps{1}
%       \feature{top}{AQP3.PRO}{208..210}{---}{NPA}
%       \feature{top}{AQP3.PRO}{211..219}{helix}{}
%       \feature{top}{AQP3.PRO}{220..232}{brace}{loop E}
%       \feature{top}{AQP3.PRO}{233..235}{helix}{TM6}
%       \feature{bottom}{AQP3.PRO}{203..235}{brace}{1-step numbering}
%       \dofrequencycorrection
%    \end{texshade}
%    
%
%    \subsection{Subfamily logos}
%    
%    \label{sublogo}
%    
%    The following output is derived from the calculation of a
%    subfamily logo [14]. Such logos display relevant deviations of a
%    subfamily compared to the remaining set of sequences. Here,
%    typical residues of AQP3 are shown (upright) which deviate 
%    from the remaining  four aquaporins of this alignment (upside-down).
%    The output can be directly compared to the sequence logo above,
%    which displays the same section of the alignment. 
%    Note, that five sequences are far too few to obtain meaningful
%    results with this method. This is just to illustrate the 
%    approach.
%    \medskip
%
%    \begin{texshade}{AQPpro.MSF}
%       \setends{AQP3.PRO}{203..235}
%       \residuesperline*{33}
%       \setsubfamily{3} 
%       \showsubfamilylogo{top} \showlogoscale{leftright}
%       \namesubfamilylogo[others]{AQP3}
%       \namessf \namessl
%       \showruler{bottom}{AQP3.PRO} \rulersteps{1}
%       \hideseqs
%       \hideconsensus
%       \dofrequencycorrection
%    \end{texshade}
%
%    Code:\medskip
%
%    \vbox{%
%    \begin{verbatim}    
%    \begin{texshade}{AQPpro.MSF}
%       \setends{AQP3.PRO}{203..235}
%       \residuesperline*{33}
%       \setsubfamily{3} 
%       \showsubfamilylogo{top} \showlogoscale{leftright}
%       \namesubfamilylogo[others]{AQP3}
%       \namessf \namessl
%       \showruler{bottom}{AQP3.PRO} \rulersteps{1}
%       \hideseqs
%       \hideconsensus
%       \dofrequencycorrection
%    \end{texshade}
%    \end{verbatim}} 
%    
%    
%
%    \subsection{Structure memes}
%    
%    The calculated sequence logo and subfamily logo data can
%    be used to generate \emph{Structure memes}. This a type of 3D protein
%    display that visualizes the residue composition, their frequency, and the
%    information content of each sequence position. Therefore, in the structure file,
%    a red-yellow-blue plus black and white color code is projected onto a sphere at
%    the $\alpha$-carbon position, 
%    which is scaled according to the information content. \TeXshade\ 
%    generates an output file for the molecular structure display software
%    Chimera. Use the files |AQP_sequence_logo_1FX8.cmd| and
%    |AQP_subfamily_logo_1FX8.cmd| generated from |texshade.dtx| to test it out in Chimera.
%
%    The structure meme example below shows an aquaporin protomer in top down 
%    view with projected sequence logo information. The image was made with Chimera from 
%    |AQP_sequence_logo_1FX8.cmd|
%    using \textit{E.\ coli} GlpF aquaglyceroporin structure data (PDB no.\ 1FX8). The standard residue colors are 
%    |CGNQST| (yellow), |HKR| (blue), |DE| (red), |APV| (white), and |FILMW| (black).
%
%   \vspace*{-2cm}
%
%    \begin{center}\includegraphics[width=0.5\textwidth]{meme.eps}\end{center}
%
%    Code to generate the respective Chimera command file:\medskip
%
%    \vbox{%
%    \begin{verbatim}    
%     \begin{texshade}{AQPall.aln}
%        \seqtype{P}
%        \showsequencelogo{top}
%        \hideseqs
%        \structurememe[AQP_sequence_logo_1FX8.cmd]{14}
%     \end{texshade}
%    \end{verbatim}} 
%
%    \subsection{Customization of the alignment output}
%    
%    Extensive possibilities are given to the user to customize 
%    the final output of an alignment. Thus, all parameters defining the
%    appearance of letters can be changed individually for sequence 
%    residues, names and numbering or the describing feature texts. 
%    Additional manual shading can be applied to any region or
%    block of residues. Sequences are easily re-ordered, separated, hidden 
%    or blanked out without recalculation of the entire alignment; 
%    sections of the alignment can also be shown. 
%    Numbering and rulers can be displayed and set to any value. 
%    A powerful tool is the |\feature|
%    command which allows one to label stretches of residues with bars, 
%    arrows, braces or any fill character and describing text. 
%    Legends are set automatically if desired, but user commands
%    are also provided to build individual legends.
%    
%    
%    \newpage
%    \section{Format of alignment input files}
%    
%    \label{alignfilestruc}
%    
%    \TeXshade{} can handle two common alignment input formats, i.\,e.\
%    the MSF format (\underline{m}ultiple \underline{s}equence 
%    \underline{f}ormat) and the ALN format 
%    (\underline{al}ig\underline{n}ment format). The MSF 
%    format is used by |PILEUP| of the Unix GCG sequence
%    analysis package\footnote{For a description see
%    |http://gene.md.huji.ac.il/Computer/GCG9doc|}. Files in the
%    ALN format are produced by |CLUSTAL| which is 
%    available for free for Unix, DOS and Macintosh. Further, upon 
%    request, the FASTA format is supported since version 1.6.
%    In addition to the mentioned software many alignment programs have 
%    export filters for the MSF, ALN or FASTA
%    format, e.\,g.\ |MACAW| produces ALN files. If
%    you are not sure whether your favorite sequence aligner 
%    produces one of the required formats compare its output to
%    the following examples. \TeXshade{} determines the format from
%    the internal file structure, thus extensions like MSF, ALN
%    or FASTA
%    are not required. If you can choose the alignment format 
%    MSF is recommended, because this format gives information
%    about the sequence type, i.\,e.\ peptide or nucleotide sequences, 
%    and length (for the correct setting of gaps at the sequence end).
%    
%    \subsection{The MSF file format}
%    Files of this type are divided into a header section and the 
%    multiple sequence alignment. The header may contain the
%    following components:
%    
%    
%    \begin{itemize}
%    \item[\textbf{File Type}:] (optional) The first header line 
%              reads for nucleic acids alignments
%              |!!NA_MULTIPLE_ALIGNMENT 1.0| and for amino acid sequences
%              |!!AA_MULTIPLE_ALIGNMENT 1.0| (all uppercase).
%    \item[\textbf{Description}:] (optional) Informative text
%              describing what is in the file.
%    \item[\textbf{Dividing line}:] (required!) Must include the 
%              following attributes:
%       \begin{itemize}
%          \item[|MSF|:] Displays the number of bases or residues in 
%                        the multiple sequence alignment.
%          \item[|Type|:] Displays the sequence type, `P' for a peptide
%                         and `N' for a nucleotide alignment.
%          \item[|Checksum|:] Displays an integer value that 
%                        characterizes the contents of the file.
%          \item[|..|] The two periods act as a divider between the 
%                        descriptive information and the following
%                        sequence information.
%       \end{itemize}
%    \item[\textbf{Name/Weight}:] (required!) Must include the name of
%             each sequence included in the alignment, as well as its
%             length, checksum and weight.
%    \item[\textbf{Two slashes} (|//|):] (required!) This separating 
%             line divides the name/weight information from the 
%             sequence alignment
%    \end{itemize}
%    
%    The alignment section consists of sequence blocks divided by an
%    empty line. Each sequence line starts out with the sequence name.
%    An example file is shown here:
%    \medskip
%    
%    \parindent-1mm
%    \begin{fmpage}
%    \begin{verbatim}
%    
%    AQP.MSF  MSF: 87  Type: P  May 1st, 1998 Check: 2586 ..
%    Name: AQP1.PRO  Len: 66  Check:  1367    Weight: 1.00
%    Name: AQP2.PRO  Len: 58  Check:  2176    Weight: 1.00
%    Name: AQP3.PRO  Len: 83  Check:  1893    Weight: 1.00
%    Name: AQP4.PRO  Len: 63  Check:  3737    Weight: 1.00
%    Name: AQP5.PRO  Len: 59  Check:  3413    Weight: 1.00
%    //
%              1                                           45
%    AQP1.PRO  MAS........................EIKKKLFWRAVVAEFLAM
%    AQP2.PRO  MW.........................ELRSIAFSRAVLAEFLAT
%    AQP3.PRO  M.........NRCG.....EMLHIRYR......LLRQALAECLGT
%    AQP4.PRO  MSDGAAARRWGKCGPPCSRESIMVAFKGVWTQAFWKAVTAEFLAM
%    AQP5.PRO  MK........................KEVCSLAFFKAVFAEFLAT
%    
%              45                                       87
%    AQP1.PRO  TLFVFISIGSALGFNYPLERNQTLVQDNVKVSLAFGLSIATL
%    AQP2.PRO  LLFVFFGLGSALQWA...SS....PPSVLQIAVAFGLGIGIL
%    AQP3.PRO  LILVMFGCGSVAQVVLSRGTHGGF....LTINLAFGFAVTLA
%    AQP4.PRO  LIFVLLSVGSTINWG...GSENPLPVDMVLISLCFGLSIATM
%    AQP5.PRO  LIFVFFGLGSALKWP...SA....LPTILQISIAFGLAIGTL
%    \end{verbatim}
%    \end{fmpage}
%    \bigskip
%    
%    \parindent0mm
%    \TeXshade{} extracts only the information from the file it
%    really needs. So, do not mind all the checksums listed
%    in the file---\TeXshade{} does not either. The same is true 
%    for |Weight|. Required are the string |MSF:| 
%    for the identification of the file format and |Type:| for the 
%    determination of the sequence type (both in the dividing line), 
%    further all |Name:| definitions and finally |//|. The MSF format 
%    allows one to comment out sequences. This is done 
%    by putting an exclamation point directly infront of the respective 
%    |Name|. These sequences are neither displayed nor used for the 
%    calculation of the consensus. This works for \TeXshade, too.
%    To comment out sequences without changing
%    the input file use the \TeXshade{} command 
%    |\killseq{|\meta{seqref}|}| (\ref{kill}). 
%    \medskip
%    
%    \parindent-1mm
%    \begin{fmpage}\label{commout}
%    \begin{verbatim}
%    
%    AQP.MSF  MSF: 87  Type: P  May 1st, 1998 Check: 2586 ..
%    Name: AQP1.PRO  Len: 66  Check:  1367    Weight: 1.00
%    !Name: AQP2.PRO  Len: 58  Check:  2176    Weight: 1.00
%    !Name: AQP3.PRO  Len: 83  Check:  1893    Weight: 1.00
%    Name: AQP4.PRO  Len: 63  Check:  3737    Weight: 1.00
%    Name: AQP5.PRO  Len: 59  Check:  3413    Weight: 1.00
%    //    
%              1                                           45
%    AQP1.PRO  MAS........................EIKKKLFWRAVVAEFLAM
%    AQP2.PRO  MW.........................ELRSIAFSRAVLAEFLAT
%    AQP3.PRO  M.........NRCG.....EMLHIRYR......LLRQALAECLGT
%    AQP4.PRO  MSDGAAARRWGKCGPPCSRESIMVAFKGVWTQAFWKAVTAEFLAM
%    AQP5.PRO  MK........................KEVCSLAFFKAVFAEFLAT
%    
%              45                                       87
%    AQP1.PRO  TLFVFISIGSALGFNYPLERNQTLVQDNVKVSLAFGLSIATL
%    AQP2.PRO  LLFVFFGLGSALQWA...SS....PPSVLQIAVAFGLGIGIL
%    AQP3.PRO  LILVMFGCGSVAQVVLSRGTHGGF....LTINLAFGFAVTLA
%    AQP4.PRO  LIFVLLSVGSTINWG...GSENPLPVDMVLISLCFGLSIATM
%    AQP5.PRO  LIFVFFGLGSALKWP...SA....LPTILQISIAFGLAIGTL
%    \end{verbatim}
%    \end{fmpage}
%    \parindent0mm
%    \bigskip
%    
%    The sequence lengths given after |Len:| are not used by 
%    \TeXshade. Due to the fact that most alignment programms calculate the 
%    sequence length by summing up residues and additionally gaps which 
%    is not really correct. In order to have the sequence break right 
%    after the last residue without printing further gap symbols 
%    \TeXshade{} counts the number of residues by itself. You can  
%    also use the command |\seqlength| in the \TeXshade{}
%    environment to set the values manually if you do not trust a machine.
%    
%    \subsection{The ALN file format}
%    ALN files are quite similar to the above described MSF files.
%    They simply lack a defined header section. Nevertheless, 
%    describing text is allowed before the alignment part. \TeXshade{}
%    determines the number of sequences and their names from the last
%    sequence block---so, no further text lines are allowed after this block!
%    Due to a lacking declaration in the file the sequence type has 
%    to be set in the |texshade| environment by |\seqtype{|\meta{type}|}|
%    \label{Lseqtype} with `P' for peptide and `N' for nucleotide sequences;
%    for the example below: |\seqtype{P}|. If no |\seqtype| command
%    is used \TeXshade{} assumes a nucleotide sequence.
%    \bigskip
%    
%    \parindent-1mm
%    \begin{fmpage}
%    \begin{verbatim}
%    
%                  profalign  May 1st, 1998, 16:58
%                  
%    of AQPpro.MSF{}
%    
%    Muliple alignment parameter:
%    
%    Gap Penalty (fixed):           10.00
%    Gap Penalty (varying):         .05
%    Gap separation penalty range:  8
%    Percent. identity for delay:   0%
%    List of hydrophilic residue:   GPSNDQEKRH
%    Protein Weight Matrix:         blosom
%    
%                      10        20        30        40
%                       .         .         .         .
%    AQP1.PRO  MAS........................EIKKKLFWRAVVAEFLAM
%    AQP2.PRO  MW.........................ELRSIAFSRAVLAEFLAT
%    AQP3.PRO  M.........NRCG.....EMLHIRYR......LLRQALAECLGT
%    AQP4.PRO  MSDGAAARRWGKCGPPCSRESIMVAFKGVWTQAFWKAVTAEFLAM
%    AQP5.PRO  MK........................KEVCSLAFFKAVFAEFLAT
%              *                                  .   ** *.
%    
%    AQP1.PRO  TLFVFISIGSALGFNYPLERNQTLVQDNVKVSLAFGLSIATL
%    AQP2.PRO  LLFVFFGLGSALQWA...SS....PPSVLQIAVAFGLGIGIL
%    AQP3.PRO  LILVMFGCGSVAQVVLSRGTHGGF....LTINLAFGFAVTLA
%    AQP4.PRO  LIFVLLSVGSTINWG...GSENPLPVDMVLISLCFGLSIATM
%    AQP5.PRO  LIFVFFGLGSALKWP...SA....LPTILQISIAFGLAIGTL
%              .. *   .**                    .   **  .
%    \end{verbatim}
%    \end{fmpage}
%    \bigskip
%    
%    The minimal contents of an ALN file are shown below; this 
%    is fully sufficient. Many sequence alignment programs can
%    produce such an output. Have a look at |seqpup| by 
%    \textsc{Don Gilbert} if you need a comprehensive conversion
%    program\footnote{Sorry, |seqpup| is much more!}.
%    \bigskip 
%   
%    \parindent-1mm
%    \begin{fmpage}
%    \begin{verbatim}
%    
%    AQP1.PRO  MAS........................EIKKKLFWRAVVAEFLAM
%    AQP2.PRO  MW.........................ELRSIAFSRAVLAEFLAT
%    AQP3.PRO  M.........NRCG.....EMLHIRYR......LLRQALAECLGT
%    AQP4.PRO  MSDGAAARRWGKCGPPCSRESIMVAFKGVWTQAFWKAVTAEFLAM
%    AQP5.PRO  MK........................KEVCSLAFFKAVFAEFLAT
%    
%    AQP1.PRO  TLFVFISIGSALGFNYPLERNQTLVQDNVKVSLAFGLSIATL
%    AQP2.PRO  LLFVFFGLGSALQWA...SS....PPSVLQIAVAFGLGIGIL
%    AQP3.PRO  LILVMFGCGSVAQVVLSRGTHGGF....LTINLAFGFAVTLA
%    AQP4.PRO  LIFVLLSVGSTINWG...GSENPLPVDMVLISLCFGLSIATM
%    AQP5.PRO  LIFVFFGLGSALKWP...SA....LPTILQISIAFGLAIGTL
%    \end{verbatim}
%    \end{fmpage}
%    \bigskip
%    
%    \subsection{The FASTA file format}
%    In FASTA files each sequence is led
%    by a single description line starting with a `|>|'. \TeXshade{} uses
%    the first word delimited by the leading `|>|' and a space as 
%    the sequence name. If no descriptive text is present \TeXshade{}
%    generates a sequence name consisting of `|seq|' plus a consecutive
%    number. The lines following the description line 
%    contain the sequence.
%    \bigskip
%    
%    \begin{fmpage}
%    \begin{verbatim}
%    
%    >AQP1.PRO  
%    MAS........................EIKKKLFWRAVVAEFLAM
%    TLFVFISIGSALGFNYPLERNQTLVQDNVKVSLAFGLSIATL
%    
%    >AQP2.PRO  
%    MW.........................ELRSIAFSRAVLAEFLAT
%    LLFVFFGLGSALQWA...SS....PPSVLQIAVAFGLGIGIL
%    
%    >AQP3.PRO  
%    M.........NRCG.....EMLHIRYR......LLRQALAECLGT
%    LILVMFGCGSVAQVVLSRGTHGGF....LTINLAFGFAVTLA
%    
%    >AQP4.PRO  
%    MSDGAAARRWGKCGPPCSRESIMVAFKGVWTQAFWKAVTAEFLAM
%    LIFVLLSVGSTINWG...GSENPLPVDMVLISLCFGLSIATM
%    
%    >AQP5.PRO  
%    MK........................KEVCSLAFFKAVFAEFLAT
%    LIFVFFGLGSALKWP...SA....LPTILQISIAFGLAIGTL
%    \end{verbatim}
%    \end{fmpage}
%    \bigskip
%    
%    
%    \parindent0mm
%    \newpage
%    \section{Use of a \TeX{}shade parameter file}
%    
%    \label{paramfilestruc}
%    
%    Using predefined parameter files for repeatedly occuring situations
%    can save a lot of typing and makes the output throughout the 
%    publication or presentation more consistent. Further, such
%    files are an easy way to exchange self-defined shading 
%    modes or new color schemes (i.\,e.\ for a satisfying grayscale output) 
%    with other users. If you have created a
%    parameter file, which you think is of interest for others, please 
%    submit it to me\footnote{|ebeitz@pharmazie.uni-kiel.de|} as an e-mail 
%    attachment together with a short 
%    description. I will take care of those files and post them---with 
%    a reference to the author---together with the next \TeXshade{} 
%    distribution to make them available for all interested users.
%    
%    No special file format is required for parameter 
%    files. \TeXshade{} simply calls the file using the |\input|
%    command right after resetting all parameters to default. An
%    example parameter file is present containing the standard
%    parameters of \TeXshade{} called |texshade.def|. This file can be
%    changed freely and can be used as a template for the creation of
%    personal parameter files. 
%    
%    Five steps are executed by \TeXshade{} when 
%    processing the |texshade| environment:
%    
%    \bigskip
%    \begin{minipage}{12cm}
%    |\begin{texshade}[|\meta{parameterfile}|]{|\meta{alignmentfile}|}|
%    
%    \begin{enumerate}
%    \item Analysis of the \meta{alignmentfile}; determination of 
%    the number of sequences and sequence names
%    
%    \item Setting parameters to default
%    
%    \item Setting parameters to the definitions of the 
%    \meta{parameterfile}, if existent
%    
%    \item Execution of further \TeXshade{} commands within the
%    evironment, if existent
%    
%    \parindent-1cm
%    \medskip
%    |\end{texshade}|
%    
%    \parindent0cm
%    \item Loading and setting the alignment on a line by line basis
%    \end{enumerate}
%    \end{minipage}
%    
%    \newpage
%    \section{\texttt{texshade} user commands}
%    
%    The \TeXshade{} package must be loaded by the |\usepackage|
%    command in the document header section.
%    \medskip
%    
%    \quad|\usepackage[|\meta{option}|]{texshade}|
%    \medskip
%    
%    Then, the |texshade| environment is ready to use as described
%    in \ref{tsenvironment}. See also section \ref{paramfilestruc} for
%    a description of the optional parameter file. All other 
%    commands provided by \TeXshade{} (except |\shadebox| [\ref{Lshadebox}], |\molweight| and
%    |\charge| [\ref{molcharge}], and |\percentsimilarity|, |\percentidentity| and
%    |\similaritytable| [\ref{Lpercentidentity}]) must 
%    be used within the |texshade| environment.
%    
%    
%    
%    \subsection{Using predefined shading modes}
%    
%    \label{predef}
%    
%    \label{Lshadingmode}
%    If no |\shadingmode| command is given in the |texshade| 
%    environment the default shading mode (\emph{identical}, see 
%    \ref{ident}) is active. For the selection of one of the other
%    predefined shading modes the following command is provided.
%    \bigskip
%    
%    \quad |\shadingmode[|\meta{option}|]{|\meta{mode}|}|
%    \bigskip
%    
%    You can choose from four shading modes and declare one option 
%    which depends on the selected mode. 
%    
%    \begin{enumerate}
%    
%    \item |\shadingmode[|\meta{allmatchspecial/number}|]{identical}|
%    
%       See \ref{ident} for examples. Use the
%       option |allmatchspecial| to shade positions with a special color
%       where all residues are identical. Or use a percentage number
%       (0--100) as an option to set an additional threshold for highly
%       conserved residues, e.\,g.\ |\shadingmode[90]{identical}|. 
%       \label{Lallmatchspecial}|\allmatchspecial| can also be 
%       used as a command with or without an optional parameter for
%       setting the high conservation threshold. As both, option or command, 
%       |allmatchspecial| is only active in the \emph{identical} and
%       \emph{similar} shading modes.
%       
%       \label{Lshadingcolors}
%       One can choose from five predefined shading color schemes with
%       the command 
%       |\shadingcolors{|\meta{scheme}|}|. The sets are named `blues' 
%       (used in the example, \ref{ident}), `reds', `greens', 
%       `grays' and `black'. Default is |\shadingcolors{blues}|. Further, the colors 
%       for the non matching, the 
%       conserved and all matching (or highly conserved) residues can be set individually 
%       plus the letter case (lower or upper) or any character
%       can be chosen: \label{Lnomatchresidues}
%       \label{Lconservedresidues}
%       \label{Lallmatchresidues} 
%       \bigskip
%       
%       |\nomatchresidues{|\meta{res.col.}|}{|\meta{shad.col.}|}{|\meta{case}|}{|\meta{style}|}|
%       
%       |\conservedresidues{|\meta{res.col.}|}{|\meta{shad.col.}|}{|\meta{case}|}{|\meta{style}|}|
%       
%       |\allmatchresidues{|\meta{res.col.}|}{|\meta{shad.col.}|}{|\meta{case}|}{|\meta{style}|}|
%       \bigskip
%       
%       For how to handle colors for the foreground \meta{res.col.} and 
%       the background \meta{shad.col.} see section \ref{colors}. 
%       The third parameter \meta{case} tells \TeXshade{} to print the
%       corresponding residue as a lowercase or an uppercase letter or
%       even to print any other character. Finally, the \meta{style} 
%       parameter tells \TeXshade{} which shape to use for the letters.
%       Use one of the following styles 
%       for \meta{style}.
%    
%       \begin{center}
%       \begin{tabular}{cl}
%       \meta{style} & \emph{effect} \\ \hline
%       |bf| & bold face series\\
%       |md| & normal series \\
%       |up| & upright shape (normal shape)\\
%       |it| & italics shape \\
%       |sl| & slanted shape \\
%       |rm| & modern roman family \\
%       |sf| & sans serif family \\
%       |tt| & typewriter family \\
%       \end{tabular}
%       \end{center}
%       \medskip    
%    
%       In order to change only some
%       of the parameters it is sufficient to declare these
%       and use empty braces for the others. Examples:
%       \bigskip
%       
%       \quad|\conservedresidues{White}{Blue}{upper}{bf}|: the conserved
%       residues are printed as bold face white uppercase letters on blue.
%       \bigskip
%       
%       \quad|\nomatchresidues{}{}{{$\bullet$}}{}|: instead of the non 
%       matching residues a `$\bullet$' is printed. The colors and style 
%       are not changed.
%       Note the double curly braces which make \TeXshade{} 
%       interpret this complex symbol description as one single 
%       character.
%
%       Once a set of shading colors is redefined, it can be saved using
%       |\defshadingcolors{|\meta{name}|}|\label{Ldefshadingcolors} for
%       later use in the document (see \ref{shaderegion}).
%
%       \bigskip
%       
%       
%    \item |\shadingmode[|\meta{allmatchspecial/number}|]{similar}|
%    
%       \label{Lsimilarresidues}
%       See \ref{similar} for an example output and an explanation
%       of the shading. In addition to the described commands
%       for changing shading colors this shading mode provides 
%       the command |\similarresidues|.
%       Use it in analogy to the commands above.
%       
%       \label{Lpepsims}\label{Lpepgroups}
%       \label{LDNAsims}\label{LDNAgroups}
%       How does \TeXshade{} know which residues are
%       considered to be similar? These definitions are set by two command
%       couples, i.\,e.\ 
%       |\pepsims|,|\pepgroups| for peptides and
%       |\DNAsims|,|\DNAgroups| for nucleotides. With |\pepsims| and
%       |\DNAsims| residues are defined which are similar to the 
%       consensus residue. Examples:
%       
%       \quad |\pepsims{S}{TA}|\quad If a serine is the consensus
%             residue then all threonins and alanines at this
%             position are shaded in the color for similars. This
%             definition does \emph{not} imply that threonine and
%             alanine are similar to each other! This becomes 
%             obvious when you inspect the next definition:
%             
%       \quad |\pepsims{T}{S}|\quad  Serine but not alanine is declared
%             to be similar to threonine.
%             
%       What happens if there is no consensus residue? How does
%       \TeXshade{} decide if a group of similars is greater than
%       the threshold? For this groups are pre-defined:
%       
%       \quad |\pepgroups{FYW,ILVM,RK,DE,GA,ST,NQ}| This command allows
%             one to set up to nine groups of similars, separated by commas. 
%             Each residue can belong to only one group. If one residue
%             is assigned to several groups only the last assignment is 
%             carried out.
%             
%       \quad |\DNAgroups{GAR,CTY}| This command is used in analogy to
%             the amino acid groups. Here, two ambiguity codes (`R' for
%             pu\underline{r}ine base and `Y' for p\underline{y}rimidine
%             base) are assigned in addition.
%             
%       Residues which do not appear in any of the four commands are 
%       considered not to belong to a group. The default 
%       settings for similars are listed below:
%       \bigskip
%       
%       \begin{verbatim}
% \pepgroups{FYW,ILVM,RK,DE,GA,ST,NQ}
%       
% \pepsims{F}{YW}  % Y and W are similar to F
% \pepsims{Y}{WF}  % W and F are similar to Y
% \pepsims{W}{YF}  % Y and F are similar to W
%    
% \pepsims{I}{LVM} % L, V and M are similar to I
% \pepsims{L}{VMI} % V, M and I are similar to L
% \pepsims{V}{MIL} % M, I and L are similar to V
%    
% \pepsims{R}{KH}  % K and H are similar to R
% \pepsims{K}{HR}  % H and R are similar to K
% \pepsims{H}{RK}  % R and K are similar to H
%       
% \pepsims{A}{GS}  % G and S are similar to A
% \pepsims{G}{A}   % A (but not S) is similar to G
%       
% \pepsims{S}{TA}  % T and A are similar to S
% \pepsims{T}{S}   % S (but not A) is similar to T
%    
% \pepsims{D}{EN}  % E and N (but not Q) are similar to D
% \pepsims{E}{DQ}  % D and Q (but not N) are similar to E
% \pepsims{N}{QD}  % Q and D (but not E) are similar to N
% \pepsims{Q}{NE}  % N and E (but not D) are similar to Q
%    
% \DNAgroups{GAR,CTY}
%       
% \DNAsims{A}{GR}  % G and R are similar to A
% \DNAsims{G}{AR}  % A and R are similar to G
% \DNAsims{R}{AG}  % A and G are similar to R
%       
% \DNAsims{C}{TY}  % T and Y are similar to C
% \DNAsims{T}{CY}  % C and Y are similar to T
% \DNAsims{Y}{CT}  % C and T are similar to Y
%       \end{verbatim}
%
%       
%    \item |\shadingmode[|\meta{filename}|]{T-Coffee}|
%       
%       Enter a \meta{filename} to load the shading
%       information from a |T-Coffee| |score_ascii| file (|www.tcoffee.org|); 
%       see example in \ref{TCoffee}. Make sure
%       that the alignment file specified in the |\texshade| command
%       and this shading file correspond to each other.
%    
%       If you do not enter a \meta{filename} here, a separate
%       command |\includeTCoffee{|\meta{filename}|}| must be used.
%    
%       |T-Coffee| shading can also be used in the consensus
%       p.\,\pageref{Lshowconsensus} and in the feature lines,  
%       in particular color scales and bar plots p.\,\pageref{Lgraphs}, 
%       for the display of shading information.
%    
%       
%    \item |\shadingmode[|\meta{seqref}|]{diverse}|
%    
%       \ref{diverse} depicts an example alignment. Choose the 
%       number or the name of the sequence \meta{seqref} which will be treated
%       as the consensus and to which the other sequences are compared. 
%       If no \meta{seqref} is declared the first sequence is set as 
%       consensus (\meta{seqref} = 1).
%       
%       Standard definitions for |diverse| 
%       mode are:
%
%       \begin{verbatim}
%          \nomatchresidues{Black}{White}{lower}{up}
%          \similarresidues{Black}{White}{lower}{up}
%          \conservedresidues{Black}{White}{{.}}{up}
%          \allmatchresidues{Black}{White}{{.}}{up}
%          \gapchar{-}
%       \end{verbatim}
%   
%       After calling |\shadingmode{diverse}| these commands can be 
%       used to redefine the |diverse| mode settings (mind the double 
%       curly braces around the dot-symbol!). 
%    
%       \label{Lhideallmatchpositions}
%       Since alignment positions where all residues match do not contain
%       much information, those sites can be blanked out using
%      Ê|\hideallmatchpositions|. The resulting break in the alignment is
%       indicated by a gap and a vertical line. See the |\setdomain| 
%       command (\ref{Lsetdomain}) 
%       for further information on how to change the gap and ruler colors.  
%       A single-stepped ruler is also recommended (\ref{Lshowruler}).
%      Ê|\hideallmatchpositions| can be combined with |\setends| 
%       (\ref{Lsetends}).
%   
%    
%    \item |\shadingmode[|\meta{type}|]{functional}|\label{funcdef}
%       There are seven different functional shading modes available for
%       peptide sequences; nucleotide sequences can not be shaded due
%       to functional aspects. Five of \TeXshade's functional modes
%       correspond to the four `alphabets' employed by \textsc{Karlin} 
%       and \textsc{Ghandour} for peptide alignments [2] or by the
%       rasmol software. Additional 
%       `alphabets' to the standard 20-letter array of amino acids
%       can highlight peptide similarities which were otherwise not visible.
%       For the `alphabet' definitions see below:
%       
%       \begin{itemize}
%       \item \meta{type} = |charge|\quad Acidic (D, E) and basic (H, 
%                           K, R).
%    
%       \item \meta{type} = |hydropathy|\quad Acidic and basic (as 
%                           above), polar uncharged (C, G, N, Q, S, 
%                           T, Y) and hydrophobic nonpolar (A, F, I, L, M, 
%                           P, V, W), see also \textsc{Kyte} \& 
%                           \textsc{Doolittle} [3].
%    
%       \item \meta{type} = |structure|\quad External (D, E, H, K, N, Q, R),
%                           internal (F, I, L, M, V) and ambivalent (A, C,
%                           G, P, S, T, W, Y).
%       
%       \item \meta{type} = |chemical|\quad Acidic (D, E), aliphatic
%                           (I, L, V), aliphatic (small) (A, G), 
%                           amide (N, Q), aromatic
%                           (F, W, Y), basic (H, K, R), hydroxyl
%                           (S, T), imino (P) and sulfur (C, M).
%                           
%       \item \meta{type} = |rasmol|\quad (D, E), (K, R, H), (F, Y, W),
%                           (A, G), (C, M), (S, T), (N, Q), (I, L, V), 
%                           (P).
%                           
%       \end{itemize}
%                           
%       The two modes described below highlight sidechain sizes and 
%       hydrophobicity, respectively, according to \textsc{Rose} 
%       \emph{et al.}\ [4,5]. Standard area stands for the surface area
%       of the residue in \AA$^2$, i.\,e. it is a measure for the size
%       of a residue's sidechain. The accessible area value (also in 
%       \AA$^2$) gives information about the size of the surface area
%       which is accessible by solvent molecules within the folded 
%       protein. A very small area means that the residue is
%       strongly buried and is thus very hydrophobic. Hydrophilic
%       residues in turn possess large accessible areas due
%       to their prefered location at the protein surface. Therefore,
%       this kind of shading provides another method, in addition
%       to |hydropathy| and |structure|, for the
%       visualization of structural protein properties.
%       
%       \begin{itemize}
%       
%       \item \meta{type} = |standard area|\quad for the area values
%                           see legend of the alignment in \ref{starea}
%                           
%       \item \meta{type} = |accessible area|\quad for values see
%                           \ref{accarea}
%       
%       \end{itemize}
%       
%       \label{Lclearfuncgroups}
%       If no \meta{type} or an unknown \meta{type} is designated as option
%       all functional groups and shading colors are cleared. This is
%       also achieved by the command 
%       |\clearfuncgroups|. With all groups cleared one can start to
%       build new shading modes from scratch. How to do this is explained
%       in the next section.
%       
%    \label{Lfuncshadingstyle} 
%    In order to exchange the colors but to keep the group definitions
%    and descriptions the command 
%    |\funcshadingstyle| can be
%    employed. Usage:
%    \medskip
%    
%    \quad|\funcshadingstyle{|\meta{residue}|}{|\meta{res.col.}|}{|\meta{shad.col.}|}|
%    
%    \hfill|{|\meta{case}|}{|\meta{style}|}|
%    \medskip
%    
%    \meta{residue} is one representative of the whole amino acid group. The 
%    colors which are declared by the next four parameters are used 
%    for all residues in this group. \meta{case} and \meta{style} are 
%    as described for example in |\nomatchresidues|.
%    \end{enumerate}
%    
%       With |\shadeallresidues| \label{Lshadeallresidues} the 
%       threshold is ignored and
%       all residues are shaded due to their group assignment.
%
%    \subsection{Creating new functional shading modes}
%    
%    The grouping of amino acids due to other properties can make sense as
%    suggested by \textsc{Karlin} and \textsc{Ghandour} [2], e.\,g.\ 
%    physical properties (molecular weight, shape), kinetic properties
%    (reaction velocity, Michaelis-Menton constant), or structure
%    ($\alpha$-helices, $\beta$-sheets, turns).
%    
%    \label{Lfuncgroup}
%    New amino acid groups are defined with the
%    |\funcgroup| command. This command needs six parameters:
%    \medskip
%    
%    \quad|\funcgroup{|\meta{descr}|}{|\meta{residues}|}{|\meta{res.col.}|}{|\meta{shad.col.}|}|
%    
%    \hfill|{|\meta{case}|}{|\meta{style}|}|
%    \medskip
%    
%    \meta{descr} contains descriptive text which is displayed in the legend.
%    The second parameter \meta{residues} holds the amino acids to be
%    grouped. The colors for the foreground and background are set 
%    with the following two parameters, the case and style is declared by the 
%    last parameters. The example below defines a 
%    funcional group named `acidic ($-$)' containing the amino acids
%    aspartic and glutamic acid with white letters on a red background:
%    \bigskip
%    
%    \quad|\funcgroup{acidic ($-$)}{DE}{White}{Red}{upper}{up}|
%    \bigskip
%    
%    For the usage of colors see section \ref{colors}. Up to nine
%    individual groups can be defined. New groups are simply added to the
%    already existing groups, i.\,e.\ if an extension of the group
%    definitions of an existing shading mode is desired there is
%    no need to clear these groups und re-define them again. Just 
%    add the new groups with the |\funcgroup| command. To create 
%    completely new modes use the command
%    |\shadingmode{functional}| without an option 
%    \emph{before} setting the new groups. The new definitions are active
%    only in the functional shading mode---so be sure to
%    have it switched on before setting the new groups.
%    Remember, |\shadingmode{functional}| without an optional parameter
%    clears all groups defined before, see above. The following example
%    shows the definitions needed to produce an output which is identical 
%    to the functional mode `charge':
%    \bigskip
%    
%    \quad|\begin{texshade}{|\meta{alignmentfile}|}|
%    \medskip
%    
%    \quad\quad |\shadingmode{functional}|
%    
%    \quad\quad |\funcgroup{acidic ($-$)}{DE}{White}{Red}{upper}{up}|
%    
%    \quad\quad |\funcgroup{basic ($+$)}{HKR}{White}{Blue}{upper}{up}|
%    \medskip
%    
%    \quad|\end{texshade}|
%    
%    
%    \subsection{Appearance of the consensus line}
%    
%    \label{Lthreshold}
%    An important parameter for the calculation of the consensus is the
%    threshold percentage. Default setting is 50\%, i.\,e.\ to become
%    the consensus residue more than half of the residues at this
%    position must be identical or similar, depending on the shading
%    mode. Any percentage between 0 and 100 is allowed and can be 
%    set with 
%    |\threshold{|\meta{percentage}|}|, e.\,g.\ |\threshold{50}|.
%
%    Additionally, an optional parameter can be set, e.\,g.\,
%    |\threshold[90]{50}|, to label residues that are highly conserved
%    in a special color (see example on page \pageref{ident}).
%    
%    \label{Lconstosingleseq}
%    Another possibility is to set one sequence of the alignment
%    as consensus and 
%    compare the other sequences to this one. Therefore, the
%    command 
%    |\constosingleseq{|\meta{seqref}|}| is provided. The
%    \meta{seqref} selects the sequence to be used as consensus
%    (numbering according to the appearance in the alignment file;
%    top sequence is number~1, or use the sequence name). 
%    Nevertheless, the threshold percentage is also taken into 
%    account, i.\,e.\ with a threshold of 50\% half
%    of the sequences must be identical or similar compared to the 
%    specified consensus sequence in order to be shaded.
%    \label{Lconstoallseqs} With |\constoallseqs| the 
%    consensus is calculated considering all sequences (the case 
%    described in the paragraph above).
%    
%    \label{Lshowconsensus}\label{Lhideconsensus}
%    \label{Lnameconsensus}
%    Consensus lines are displayed either on the top or at the bottom
%    of the alignment by calling 
%    \medskip
%
%    |\showconsensus[|\meta{color/scale}|[,|\meta{color/scale}|]]{|\meta{position}|}| 
%    \medskip
%
%    with 
%    \meta{scale} |Gray|, |BlueRed|, |RedBlue|, |GreenRed|, |RedGreen|,
%    |ColdHot| (recommended), |HotCold|, or |T-Coffee| \ref{TCoffee}, and \meta{position} |top| 
%    or |bottom|. 
%
%    The first color defines the foreground, i.e. the letters, the
%    second color---if specified---defines the background. 
%    If a color scale is named the consensus will be shaded according
%    to the level of sequence conservation (see section \ref{resweight} on residue 
%    weight tables below). For an example output see page
%    \pageref{shadecons}. You can find more information on color scales 
%    on page \pageref{Lgraphs}. The calculated consensus colors can be exported
%    as a Pymol [8] \label{Lexportconsensus}
%    script by |\exportconsensus[|\meta{filename}|]{|\meta{seqref}|}|.
%    If no \meta{filename} is specified |export.txt| will be used. The 
%    generated file can be opened in Pymol in order to shade a 3D model
%    of the sequence \meta{seqref}.
%
%    To hide the consensus use 
%    |\hideconsensus|. The consensus
%    line is named `consensus' in english texts, `consenso' in spanish
%    or `Konsensus' if the |german.sty| is used. With 
%    |\nameconsensus{|\meta{name}|}| any name can be set.
%    
%    \label{Ldefconsensus} 
%    You can tell \TeXshade{} which symbols or letters to use in
%    the consensus line for different matching qualities by 
%    \bigskip
%    
%    \quad|\defconsensus{|\meta{symbol1}|}{|\meta{symbol2}|}{|\meta{symbol3}|}|.
%    \bigskip
%    
%    The following parameters are allowed for symobols 1--3:
%    
%    \begin{enumerate}
%    
%    \item \meta{symbol1} = no match symbol (if below threshold)
%    
%          \begin{itemize}
%             \item any character or letter
%             \item |{}| (empty braces) for blank space
%          \end{itemize}
%          
%    \item \meta{symbol2} = conserved symbol (if threshold is exceeded)
%    
%          \begin{itemize}
%             \item |upper| (prints the consensus residue in uppercase)
%             \item |lower| (prints the consensus residue in lowercase)
%             \item any character or letter
%             \item |{}| (empty braces) for blank space
%          \end{itemize}
%          
%    \item \meta{symbol3} = highly conserved symbol (if
%
%                                 \hfill|\allmatchspecial| is active)
%    
%          \begin{itemize}
%             \item see \meta{symbol2}
%          \end{itemize}
%    
%    \end{enumerate}
%    
%    Example: |\defconsensus{{}}{*}{upper}| does not show non matching
%    residues in the consensus line, marks conserved residues 
%    with `|*|', and displays the uppercase letter of the consensus
%    residue at positions with high conservation.
%    
%    
%    \label{Lconsensuscolors}
%    Finally, the colors of the above defined symbols are adjustable
%    by the command: 
%    
%    \begin{tabbing}
%    \quad|\consensuscolors|\=|{|\meta{res.col.1}|}{|\meta{shad.col.1}|}|\\
%    
%    \>|{|\meta{res.col.2}|}{|\meta{shad.col.2}|}|\\
%    
%    \>|{|\meta{res.col.3}|}{|\meta{shad.col.3}|}|\\
%    \end{tabbing}
%    
%    The color definitions are in the same order as in the
%    |\defconsensus| command:
%    
%    \begin{enumerate}
%    
%    \item \meta{res.col.1} = no match residue color (if below threshold)
%    
%          \meta{shad.col.1} = no match background color 
%          
%    \item \meta{res.col.2} = conserved residue color (if threshold is exceeded)
%    
%          \meta{shad.col.2} = conserved background color 
%          
%    \item \meta{res.col.3} = highly conserved residue color (if 
%
%                                 \hfill|\allmatchspecial| is active)
%    
%          \meta{shad.col.3} = highly conserved background color 
%    
%    \end{enumerate}
%    
%    For colors which are not to be changed empty braces can be used.
%    
%    Example:\medskip
%    
%    \quad|\consensuscolors{}{}{Blue}{White}{Red}{Green}|
%    \medskip
%    
%    Non matching symbol colors are not changed, 
%    conserved residues are displayed blue on white and highly conserved residues
%    appear as red symbols on a green background in the 
%    consensus line.
%    
%    
%    \subsubsection{Residue weight tables} \label{resweight}
%    
%    The degree of similarity between two amino acid residues is defined using
%    so-called \emph{residue weight tables}. The values usually range roughly
%    from $-10$ to 10, with positive values denoting similarity and negative
%    values dissimilarity. The most simple table sets pairs
%    of identical residues to a value of 10 and all others to 0, i.e. the
%    |identity| matrix. Several more matrices based on extensive protein alignments
%    exist and can be used, e.g. |PAM250| (|Point Accepted Mutations|), |PAM100|,
%    or |BLOSUM62| (|BLOcks of amino acid SUbstitution Matrix|); for details see
%    respective sources and section \ref{weightmatrix}. \TeXshade{} further contains a 
%    |structural| matrix where 
%    similarity is defined on simple comparisons of the sidechain properties with
%    respect to volume and hydropathy.
%    
%    For calculation of the consensus color shading or for bar graphs or color scales
%    in the |\feature| lines (\ref{Lfeature}), a residue weight table can be selected
%    by \label{Lweighttable} |\weighttable{|\meta{table}|}| with \meta{table}
%    being |identity|, |structural|, |PAM250|, |PAM100|, or |BLOSUM62| (default is
%    |identity|). Which matrix suits the analysis best needs to be decided case by case.
%    Due to the all-positive values of the |structural| matrix (section \ref{weightmatrix}) the similarity level
%    appears usually very high; the |identity| matrix simply represents the number
%    of identical residues at each position. The |PAM| and |BLOSUM| matrices provide
%    more differentiated results.
%    One can change individual values or even define his own weight table
%    using the command \label{Lsetweight} 
%    |\setweight{|\meta{res.1}|}{|\meta{res.2}|}{|\meta{value}|}|, e.g.
%    |\setweight{E}{Q}{2}| or |\setweight{K}{C}{-5}|. A full table, thus, needs 200
%    entries ($20 * 20 / 2$). \label{Lgappenalty}
%    A value for the gap penalty is set with |\gappenalty{|\meta{value}|}|.
%    
%    \subsection{Display of logos}
%    
%    \subsubsection{Sequence logos}
%    
%    \label{Lshowsequencelogo}\label{Lhidesequencelogo}
%    In a sequence logo [7], the information content $I(P_i)$ of 
%    each alignment position $i$ is defined as 
%
%    \[ 
%    I(P_i) = \log_2 \vert\Sigma\vert + \sum P_{ij} \cdot \log_2 P_{ij}
%    \]
%
%    \noindent
%    with $\vert\Sigma\vert$ being the cardinality of the used alphabet, 
%    i.\,e. 4 for DNA and 20 for protein sequences, and $P_{ij}$ 
%    being the frequency of residue $j$ at this position. Each position 
%    is displayed as a stack of residue symbols whose heights 
%    represent their proportion of the information content (example on
%    p.\pageref{logo}).
%
%    The display of sequence logos can be either on the top or at the bottom
%    of a nucleotide or protein alignment. Logos will be shown after the
%    command: |\showsequencelogo[|\meta{colorset}|]{|\meta{top/bottom}|}|. If no optional
%    \meta{colorset} is selected the residues will be shaded as follows:\medskip
%
%    \begin{itemize}
%    \item Nucleotide sequences
%
%    \begin{itemize}
%      \item[G]: Black
%      \item[A]: Green
%      \item[T,U]: Red
%      \item[C]: Blue
%    \end{itemize}
%    
%
%    \item Protein sequences (similar to rasmol)
%
%    \begin{itemize}
%      \item[D,E]: Red
%      \item[C,M]: Yellow
%      \item[K,R]: Blue
%      \item[S,T]: Orange
%      \item[F,Y]: MidnightBlue
%      \item[N,Q]: Cyan
%      \item[G]:   LightGray
%      \item[L,V,I]: Green
%      \item[A]: DarkGray
%      \item[W]: CarnationPink
%      \item[H]: CornflowerBlue
%      \item[P]: Apricot
%      \item[B,Z]: LightMagenta
%    \end{itemize}
%    \end{itemize}
%    
%    Optional color sets correspond to the functional shading modes
%    |chemical|, |rasmol|, |hydropathy|, |structure|, |standard area|, 
%    |accessible area| (see p.\pageref{funcdef}). The |\showsequencelogo| 
%    command can be reversed by |\hidesequencelogo|.
%
%    \label{Llogocolor}\label{Lclearlogocolors}
%    Logo colors can be turned to `Black' with the command
%    |\clearlogocolors[|\meta{color}|]| with the optional parameter
%    not set. The optional parameter can be used to set all 
%    residue colors to \meta{color}, e.g.\ |\clearlogocolors[Blue]|.
%    User specific logo color sets are defined by using
%    |\logocolor{|\meta{residues}|}{|\meta{color}|}|, e.g.\ 
%    |\logocolor{DE}{Red} \logocolor{CM}{Yellow}| etc.
%
%    \label{Ldofrequencycorrection}\label{Lundofrequencycorrection}
%    It is common practice for protein sequence logos to correct
%    amino acid frequencies to the background frequency in the
%    alignment, which usually differs from the equal distribution
%    of 5\% for each residue. Frequency correction can be turned on
%    by |\dofrequencycorrection| and off by |\undofrequencycorrection|.
%
%    \label{Llogostretch}The vertical extent of the logo can be changed by 
%    |\logostretch{|\meta{factor}|}|, e.g.\ |\logostretch{1.5}|.
%    The width of the logo characters is dependent on the character
%    width set for the alignment, see |\charstretch| on p.\pageref{Lcharstretch}.
%
%    \label{Lshowlogoscale}\label{Lhidelogoscale}Finally, the bit-scale 
%    can be turned off and on using |\hidelogoscale| and
%    |\showlogoscale[|\meta{color}|]{|\meta{position}|}|, respectively, with
%    \meta{position} |left|, |right|, or |leftright| and an optional
%    \meta{color}. 
%    \label{Lnamesequencelogo}
%    A name for the sequence logo can be set, which is displayed
%    next to the scale by |\namesequencelogo{|\meta{name}|}|.
%
%
%    \subsubsection{Subfamily logos}
%
%    Subfamily logos provide a novel tool to visualize
%    subfamily-specific sequence deviations at alignment positions with
%    a high information content in an intuitive way [14].
%
%    This is achieved by subtracting from the frequency of a residue within 
%    a pre-defined subset of sequences, i.\,e. a subfamily, the frequency of 
%    this residue in the remaining set of sequences. The difference is then 
%    weighted by the information content, see above section on sequence logos.
%    An example is shown on p.\pageref{sublogo}.
%
%    Subtraction of frequencies produces values from $-1$ to $1$. Positive 
%    values correspond to residues which are characteristic for the subfamily 
%    (shown upright in the output), negative values to those that are typical 
%    for the remaining sequences (shown upside-down). Positions with an equal 
%    distribution of the residue result in a zero value. 
%
%    \label{Lshowsubfamilylogo}\label{Lhidesubfamilylogo}\label{Lsetsubfamily}
%    Subfamily logos are displayed analogous to sequence logos by the command 
%    |\showsubfamilylogo[|\meta{colorset}|]{|\meta{top/bottom}|}| and hidden by
%    |\hidesubfamilylogo|. To calculate a subfamily logo, it is further required
%    to define a subfamily within the alignment by 
%    |\setsubfamily{|\meta{seqrefs}|}|, e.g. |\setsubfamily{1-10,20,AQP3}|. 
%
%    For coloring residues, display/stretching of the scales, and frequency 
%    correction the same commands as for sequence logos apply with two exceptions.
%    \label{Lshownegatives}\label{Lhidenegatives}
%    First, subfamily logos contain negative values, which can be displayed
%    |\shownegatives[|\meta{weak, medium, strong}|]| or hidden 
%    |\hidenegatives|. Without the optional parameter negative residues will
%    be tinted by 50\%, i.e. |medium|. This greatly improves readability. 
%    \label{Lnamesubfamilylogo}
%    Second, a name for the subfamily logo is set by 
%    |\namesubfamilylogo[|\meta{neg.name}|]{|\meta{pos.name}|}| with a required
%    name for the positive part of the logo and an optional name for the negative
%    part.
%
%    \label{Lrelevance}\label{Lshowrelevance}\label{Lhiderelevance}
%    In order to better recognize relevant positions in the subfamily logo, a 
%    bit-value can be set above which the deviation is considered relevant
%    by the command |\relevance{|\meta{bit-value}|}|. If this command is
%    not given 2.321\,bit is assumed for proteins, i.\,e.
%    $\log_2 5$, and 1\,bit for DNA, i.\,e. $\log_2 2$. Such positions will 
%    be labeled by 
%    |\showrelevance[|\meta{color}|]{|\meta{symbol}|}|, e.\,g.
%    |\showrelevance[Blue]{$\nabla$}|. The symbol will be hidden with
%    |\hiderelevance|. 
%
%
%
%    \subsubsection{Structure memes}
%
%    \label{Lmemes}
%
%    Structure memes combine sequence logo and subfamily logo
%    information with 3D protein structure data. \TeXshade{} generates
%    a `command file' output that can be loaded into the molecular
%    structure display software Chimera.
%    
%    A sequence logo or subfamily logo must be calculated as described in
%    the two subsections above. Then, the command
%    |\structurememe[|\meta{seqref}|]{|\meta{filename}|}| generates an
%    Chimera output file (use the ending |.cmd| or |.com|) with 
%    \meta{seqref} as the reference sequence for which a protein structure
%    file must be available for display. Two example files are distributed with
%    this \TeXshade{} version to load into Chimera; the PDB file 1FX8 must
%    be used for the structure meme.
%    
%    The |\memeStandardcolors| are: |DE| (Red), |CGNQSTY| (Yellow), 
%    |HKR| (Blue), |APV| (White), |FILMW| (Black).
%    
%    The color grouping can be changed using |\memeRed{|\meta{residues}|}|,
%    and accordingly |\memeYellow|, |\memeBlue|, 
%    |\memeWhite|, and |\memeBlack|. If a 3D structure file contains
%    more than one chain, the structure meme display can be 
%    restricted to one chain using |\chimerachain{|\meta{chainlabel}|}| with
%    \meta{chainlabel} typically being a letter e.g. |{A}|. For the radius of the 
%    $\alpha$-carbon spheres the logo information bit-value is used, which can
%    be scaled using |\chimeraballScale{|\meta{factor}|}| with \meta{factor}
%    from 0--1. The residue composition at a sequence position is stored in
%    a name label of the command file and will be visible in Chimera when 
%    pointing to it. To eliminate irrelevant residues a cut-off bit-value can be 
%    set |\memelabelcutoff{|\meta{bit-value}|}|. The symbols |>| |:| and |.|  
%    indicate the 3, 2, and 1 bit levels. 
%    
%    Structure memes from subfamily logos will duplicate the structure model
%    and shift it by 40 \AA{} along the x-axis to achieve a side-by-side arrangement.
%    The translational direction can be changed using
%     |\chimeraxisdistance{|\meta{x-dist}|}{|\meta{y-dist}|}{|\meta{z-dist}|}| in 
%     \AA{}ngstr\o{}m.
%    
%    Eventually, with |\echostructurefile{|\meta{strucref}|}| a reference to a
%    suitable structure file, e.g. from the PDB, can be named, which will be put
%    into the command file and which will appear when Chimera asks for a
%    structure file in a load dialog window.
%
%    \subsection{Appearance of the sequence lines}
%    
%    \label{seqlines}
%    
%    \subsubsection{Names, numbers and gaps}
%     \label{Lshownames}\label{Lshownumbering}
%    Many parameters that influence the appearance of the actual sequence
%    lines can be changed for customization.
%    Thus, the sequence names can be shown colored via \meta{color}
%    either left or right by
%    \medskip
%    
%    \quad|\shownames[|\meta{color}|]{|\meta{position}|}|
%    \medskip
%
%    with \meta{position} set to |left| or |right|. The numbering can be
%    displayed either left or right and even on both sides by
%    \medskip
%    
%    \quad|\shownumbering[|\meta{color}|]{|\meta{position}|}|
%    \medskip
%    
%    with \meta{position} |left|, |right| or |leftright|. Both,  
%    names and numbering can be displayed on the same side.     
%    \label{Lnamescolor}\label{Lnumberingcolor}    
%    The colors can also be set with |\namescolor{|\meta{color}|}| and 
%    |\numberingcolor{|\meta{color}|}|, respectively. 
%
%    \label{Lnameseq}
%    \TeXshade{} uses the sequence names from the 
%    alignment input file. This can cause some
%    problems during the \TeX-run when special characters are present
%    in those names! \TeXshade{} does not accept the following characters
%    in sequence names: |\ { } @| spaces and the tilde. Those have to be replaced in 
%    the input file. The characters |#| and |%| can only be used with a
%    leading backslash, e.\,g. |\#|. This must also be changed in the
%    input file. All other special characters should be displayed 
%    properly.    
%    
%    Sequence names that are accepted by \TeXshade{} can further be
%    changed in the |texshade| environment: 
%    \medskip
%    
%    \quad|\nameseq{|\meta{seqref}|}{|\meta{name}|}| 
%    \medskip
%    
%    \meta{seqref} selects the sequence whose name is to be changed. 
%    The basis for the \meta{seqref} is the appearance in
%    the alignment input file with the top sequence = 1, or the old 
%    name. 
%    \label{Lnamecolor}\label{Lnumbercolor}    
%    In order to change the colors only of some sequence names or numbers 
%    the commands 
%    |\namecolor{|\meta{seq1}|, ... ,|\meta{seq n}|}{|\meta{color}|}| and 
%    |\numbercolor{|\meta{seq1}|, ... ,|\meta{seq n}|}{|\meta{color}|}| 
%    are provided.
%    
%    \label{Lhidenames}\label{Lhidename}
%    \label{Lhidenumbering}\label{Lhidenumber}
%     In order to hide all names or the numbering use the command
%     |\hidenames| or |\hidenumbering|. If only the names or numbers of
%     some sequences should be hidden apply 
%     
%     |\hidename{|\meta{seq1}|, ... ,|\meta{seq n}|}| or
%     
%     |\hidenumber{|\meta{seq1}|, ... ,|\meta{seq n}|}|, respectively.
%    
%    \label{Lstartnumber} \label{Lallowzero} \label{Ldisallowzero}
%    In some situations, e.\,g.\ when only sections of sequences are 
%    displayed, one
%    may not want to have the residue numbering start out with number~1.
%    The command 
%    |\startnumber[|\meta{start..stop}|]{|\meta{seqref}|}{|\meta{startnumber}|}|
%    allows one to set the starting number of any sequence to any value
%    incl.\ negative values but except `0' which is not usually used in
%    sequence numbering (the transition from negative to positive
%    values is like this: \ldots\ $-2$, $-1$, 1, 2 \ldots). If, however,
%    the use of the number `0' is wanted as sometimes in sequence logos
%    this can be turned on by |\allowzero| and off with |\disallowzero|.
%    The optional parameter can be used to truncate the sequence display
%    to a certain section (see also |\setends| below). 
%    
%    \label{Lseqlength}
%    \TeXshade{} needs to know the correct length of the sequences
%    to be able to break them right after the last residue. If
%    MSF files are used as an input the length is already given
%    but the calculation is usually wrong because the gaps are
%    also counted. Thus, \TeXshade{} counts the number of residues
%    during each run by itself and stores the values in the |.aux| file. That
%    means that it needs two runs to get the numbers right. Again,
%    this is only important if the gap symbol after the sequence end
%    should be suppressed, see below (|\hideleadinggaps|).
%    
%    If you know the correct length of the sequences you can use the
%    command 
%    \medskip
%    
%    \quad|\seqlength{|\meta{seqref}|}{|\meta{length}|}| 
%    \medskip
%    
%    in order to set the values by hand and have the gaps break 
%    properly already in the first \TeX{} run.
%    \medskip
%    
%    Example: |\seqlength{1}{346}| means that sequence no.~1 is 346 
%    residues long.
%    
%    
%     \label{Lshowruler}\label{Lhideruler}
%     \label{Lrulersteps}\label{Lrulercolor}
%     \label{Lrotateruler}\label{Lunrotateruler}
%     \label{Lnamerulerpos}\label{Lrulerspace}
%     Another possibility to label sequence positions is to switch
%    on a ruler on the top or at the bottom of the sequence block
%    using \label{ruler}
%    |\showruler[|\meta{color}|]{|\meta{position}|}{|\meta{seqref}|}|. 
%    The residue ruler of one sequence \meta{seqref} or the consensus 
%    (declare `|consensus|' as \meta{seqref}) can be 
%    displayed at \meta{position} |top| or |bottom|.
%    The ruler is hidden with  |\hideruler|. The steps between two
%    numbers are set by |\rulersteps{|\meta{number}|}|. If the steps
%    are set to be very close ($< 4$) or when every position is numbered, the
%    numbering is automatically rotated by 90$^\circ$. Using |\rotateruler|
%    and |\unrotateruler| this can be done and undone manually. 
%    In order to change the
%    ruler color use the optional parameter or the command 
%    |\rulercolor{|\meta{color}|}|. Also, the label and its color at individual 
%    ruler positions can be changed by the user to a string using 
%    |\namerulerpos{|\meta{number}|}{|\meta{text}|[|\meta{color}|]}| 
%    (see example on p.\ \pageref{ras}). Finally, to adjust the distance
%    between the ruler and the top or bottom sequence row use
%    |\rulerspace{|\meta{length}|}|, e.\,g. |\rulerspace{1mm}|.
%    
%    \label{Lgapchar}\label{Lgaprule}
%    \label{Lgapcolors}\label{gapchar}
%    Further, the symbol which is displayed in sequence gaps is freely 
%    selectable with  
%    |\gapchar{|\meta{symbol}|}|. \meta{symbol} can be any character
%    or symbol. If math symbols are to be used math mode must be
%    activated by |$| characters, i.\,e. |\gapchar{{$\triangle$}}|.
%    Note the double curly braces in the last command. Everytime a
%    `complex' character is used, i.\,e. a character definition consisting 
%    of more than one letter, it must be braced in order to be interpreted as one 
%    character. One exception is |\gapchar{rule}|; with this 
%    parameter lines are drawn in the sequence gaps with a certain 
%    thickness defined by |\gaprule{|\meta{thickness}|}|, e.\,g.
%    |\gaprule{1.5pt}|. The colors of the gaps and gap symbols are set by 
%    |\gapcolors{|\meta{symbol color}|}{|\meta{background color}|}|.
%    
%    There are some discussions whether or not to display gap symbols before
%    and after the actual sequence. Since v1.3a one can control the
%    appearance of those gap symbols by the commands
%    \label{Lshowleadinggaps} \label{Lhideleadinggaps}
%    |\showleadinggaps| and |\hideleadinggaps|. By default, leading 
%    gaps are indicated by symbols despite my personal
%    thinking that it could suggest that 
%    there are some not displayed residues upstream resp.\ downstream of the 
%    gap.
%    
%    At certain instances a protein alignment input file may contain stop 
%    positions, e.\,g. due to frame shifts in the underlying DNA sequence.
%    If such positions are labeled in the input with an |*| this will be 
%    shown in the output as well as an asterisk, i.\,e. distinguishable from
%    a normal gap symbol. The character shown at stop positions can be
%    changed by |\stopchar{|\meta{symbol}|}|. \label{Lstopchar}
%    
%    \subsubsection{Displaying selected residues in the alignment}
%    
%    \label{Lsetends}
%    \TeXshade{} can display a section of the complete alignment
%    without the need to edit the alignment input file or even 
%    to re-calculate
%    the entire alignment. This allows one to use one single 
%    alignment of the full length proteins or open reading frames for
%    multiple visualizations of different sections in a document as
%    done in this manual. Thus, the file |AQPpro.MSF| contains
%    the full-length multiple protein alignment of five aquaporins but
%    only sections are displayed as examples in 
%    \ref{ident} through \ref{accarea}. The definition of a section
%    is done by  
%    \medskip
%    
%    \quad|\setends[|\meta{startnumber}|]{|\meta{seqref}|}{|\meta{start..stop}|}|.
%    \medskip
%    
%    Again, \meta{seqref} is the sequence number based on the 
%    appearance in the alignment file, or the name; further, in order to use
%    the consensus as a measure for the sequence section the
%    string `|consensus|' as \meta{seqref} is accepted. The 
%    specified sequence is truncated at
%    positions \meta{start} and \meta{stop}. All other 
%    sequences are cut accordingly. If the number of the first
%    residue in the sequence is set to a new value with the 
%    |\startnumber| command (s.\,a.) this is taken into account. The 
%    \meta{startnumber} can be set as an optional parameter directly
%    in the |\setends| command as well.
%    \medskip
%
%    Some examples:
%    \medskip
%    
%    \quad a) |\setends{1}{21..100}|
%    \medskip
%    
%    \quad b) |\startnumber{1}{101} \setends{1}{121..200}|
%    \medskip
%    
%    Both commands select the same sequence section from the alignment but
%    numbering for sequence 1 starts at position~21 in the first example and at 
%    position~121 in the latter.
%    \medskip
%    
%    \quad c) |\setends[101]{1}{121..200}| equals example b.
%    \medskip
%    
%    \medskip
%    
%    \quad d) |\startnumber[121..200]{1}{101}| also equals example b.
%    \medskip
%    
%    \medskip
%    
%    \quad e) |\setends{consensus}{21..100}|
%    \medskip
%    
%    This may describe a very different section of the multiple
%    sequence alignment because the consensus counts every position
%    including gaps.
%    
%    \label{Lsetdomain}
%    The output can be even further restricted to
%    individually selected residues, e.g.\ to eliminate uninteresting alignment stretches or
%    to condense the output, by:
%    \medskip
%    
%    \quad|\setdomain{|\meta{seqref}|}{|\meta{selection}|}|
%    \medskip
%    
%    Here, \meta{seqref} denotes the reference sequence by its number or name.
%    This sequence is used to define the alignment positions \meta{selection} 
%    to be shown. 
%    \meta{selection} can have two different formats
%    depending on whether (a) the user wants to select the positions manually
%    or (b) \TeXshade{} is supposed to select the residues based on 3D
%    coordinates provided by a PDB file.
%    
%    To select the residues manually, the user provides a position list of
%    the following format:
%    \medskip
%    
%    \quad|{|\meta{start1}..\meta{stop1}|,|\meta{start2}..\meta{stop2}|,|\ldots|,|\meta{start n}..\meta{stop n}|}|
%    \medskip
%
%    For how to select positions by 3D coordinates
%    provided by a PDB file, see \ref{Lshaderegion}. 
%    \medskip
%    
%    Examples (see also p.\,\pageref{ident}ff):
%
%    \quad|\setdomain{1}{20..80}|
%    
%    \quad|\setdomain{consensus}{20..80,100..150,200..220}|
%    
%    \quad|\setdomain{AQP1}{point[6]:1FX8.pdb,173[side]}|
%    
%    \quad|\setdomain{3}{plane[0.5]:1JN4.pdb,66[CA],73[side],199[CA]}|
%    \medskip
%
%    
%    It is helpful to show a ruler (e.g. single-stepped, see p.\,\pageref{Lshowruler}) to 
%    label the residue positions. 
%    
%    The resulting gaps between sequence stretches are marked by a vertical rule, which
%    can be changed in thickness by 
%    \medskip
%
%    \quad|\domaingaprule{|\meta{thickness}|}|, e.\,g.\
%    |\domaingaprule{1pt}|. \label{Ldomaingaprule}
%    \medskip
%
%    Also, the colors can be set by 
%    \medskip
%
%    \quad|\domaingapcolors{|\meta{foreground}|}{|\meta{background}|}| \label{Ldomaingapcolors}
%    \medskip
%
%     e.\,g.\ |\domaingapcolors{Blue}{Yellow}|.
%    
%    
%    \subsubsection{Hiding, killing, separating and ordering}
%    
%    \label{kill}
%    
%    \label{Lhideseq}\label{Lhideseqs}\label{Lshowseqs}\label{Lkillseq}
%    If one or more sequences from the alignment input file should be used for
%    the calculation of the consensus but it is desired not to
%    display these sequences in the final output use the command 
%    |\hideseq{|\meta{seq1}|,|\meta{seq2}|,|\ldots|,|\meta{seq n}|}|. 
%    For consecutive sequence numbers a dash can be used, e.\,g. 
%    |\hideseq{1-3}| instead of |\hideseq{1,2,3}|. Decending series 
%    are also permitted, e.\,g. |\hideseq{3-1}|.
%    This command allows one for example to hide 
%    the sequence which has been defined as the consensus sequence
%    with |\constosingleseq|. When all sequences should be hidden, e.g. to
%    show a sequence logo alone, one can simply say |\hideseqs|. This
%    command is reversed by |\showseqs|.
%
%    In order to completely exclude sequences the command 
%    |\killseq{|\meta{seq1}|,|\meta{seq2}|,|\ldots|,|\meta{seq n}|}| is 
%    provided. Again, for number series the dash can be used (s.\,a.). The 
%    designated sequences are neither displayed nor
%    considered for the calculation of the consensus. This is
%    another possibility to comment out sequences in addition
%    to the use of an exclamation point infront of the |Name:|
%    definition in an MSF-file (see figure on page \pageref{commout}).
%    
%    \label{Ldonotshade}
%    The command 
%    |\donotshade{|\meta{seq1}|,|\meta{seq2}\ldots|,|\meta{seq n}|}| 
%    makes 
%    one or more sequences (remember the dash, s.\,a.) appear unshaded 
%    in black letters on white background.
%    This does not influence any other sequences or the consensus
%    calculation.
%    
%    \label{Lhideresidues}\label{Lshowresidues}
%    If a very graphical output of the sequences is desired, the 
%    residue symbols or letters can be blanked out by
%    |\hideresidues|. Now, only the shaded boxes are printed. 
%    In combination with |\gapchar{rule}| one obtains alignments
%    in a style \`a la Mondrian.
%    The residues reappear with |\showresidues|.
%    
%    \label{Lseparationline}\label{Lsmallsep}
%    \label{Lmedsep}\label{Lbigsep}
%    \label{Lvsepspace}
%    If an alignment contains members of several subgroups of a
%    protein or a gene family it may be rather helpful to visualize the group
%    divisions by a separation line. Therefore, the command
%    |\separationline{|\meta{seqref}|}| is applicable. This 
%    command inserts vertical space after the sequence which is
%    refered to by \meta{seqref}. How much space is inserted
%    is defined by one of the following commands: 
%    |\smallsep|, |\medsep| (default) or |\bigsep|. These lengths
%    correspond to the known |\small|-, |\med|- and |\bigskip| commands. 
%    With |\vsepspace{|\meta{length}|}| any length with any
%    \TeX{} unit can be assigned, e.\,g. |\vsepspace{2mm}|.
%    
%    \label{Lorderseqs}
%    The sequence order given by the alignment input file is easily
%    reorganized by 
%    |\orderseqs{|\meta{seq1}|,|\meta{seq2}|,|\ldots|,|\meta{seq n}|}|
%    without the need for editing the alignment input file (which 
%    would be a big copy'n'paste job).
%    Make sure that all sequences are assigned in this
%    command. If there are more sequences present than numbers or names in the
%    command an error message will occur. Here also, the dash can be 
%    used for sequence number series. Example: |\orderseqs{1-3,6-4,7}| 
%    is equivalent to |\orderseqs{1,2,3,6,5,4,7}|.
%    Reordering of sequences only changes the output; all commands using
%    the parameter \meta{seqref} are not influenced, because \meta{seqref} 
%    always corresponds to the appearance in the alignment file. Thus, 
%    to completely reverse the order of a five sequence alignment simply type 
%    |\orderseqs{5-1}|.
%    
%    
%    \subsubsection{Residues per line and further settings}
%    
%    \label{Lresiduesperline}\label{Lresiduesperline*}
%    By default \TeXshade{} puts the highest possible by five 
%    divisible number of residues in one line depending on the 
%    |\linewidth|. With |\residuesperline{|\meta{number}|}| a new 
%    value can be set. If this value exceeds the highest possible 
%    number of residues per line it is ignored; lower values are
%    accepted of course. But also in the latter case the number
%    of residues printed per line is rounded such to be divisible by five. 
%    To force \TeXshade{}
%    to set lines with exactly the desired number of residues use
%    the asterisk-extended command |\residuesperline*{|\meta{number}|}|. 
%    You have to take care yourself of the alignment width after this command, 
%    because in this mode \TeXshade{} does not check the length of the 
%    lines any more.
%    
%    \label{Lcharstretch}\label{Llinestretch}
%    \TeXshade{} calculates the dimensions of a shaded box from
%    the width and height of the uppercase letter `M' and the depth of
%    the lowercase `g'. Depending on the font used for the
%    sequence residues the box dimensions might not be fully
%    satisfactory. With |\charstretch{|\meta{factor}|}| and
%    |\linestretch{|\meta{factor}|}| the width and height/depth,
%    respectively, of the boxes can be multiplied individually by a 
%    \meta{factor} to stretch ($>1$) or shrink ($<1$) the dimensions.
%    
%    \label{Lnumberingwidth}
%    The reserved space for the sequence numbering is set by the
%    command |\numberingwidth{|\meta{n digits}|}|. Here, the default setting
%    is four-digit numbering, i.\,e.\ $-999$ through 9999. If this range
%    is to be changed assign the desired number as parameter
%    \meta{n digits}, e.\,g.\ |\numberingwidth{111111}| reserves
%    space for 6 digit numbering.
%    
%    The vertical space between the sequence blocks can be controlled
%    by the commands |\smallblockskip|, |\medblockskip| (default 
%    setting), 
%    \label{Lsmallblockskip}\label{Lmedblockskip}
%    \label{Lbigblockskip}\label{Lnoblockskip}
%    \label{Lvblockspace}
%    |\bigblockskip| or |\noblockskip|. Further, the command 
%    |\vblockspace{|\meta{length}|}| allows one to set a defined space
%    length using any \TeX{} unit, e.\,g.\ |\vblockspace{0.4in}|.
%    
%    Two more commands set the space between the sequence blocks to be
%    \label{Lflexblockspace}\label{Lfixblockspace}
%    flexible (|\flexblockspace|) (default) or fixed (|\fixblockspace|).
%    Flexible means, that only the vertical white space between the
%    blocks is kept to the settings by
%    e.\,g. |\medblockskip|. This results in flexible space between 
%    the actual blocks depending on the presence of feature lines. When 
%    switching to fixed space the distance of the blocks is kept constant
%    by using more white space between blocks without feature lines.
%    Thus, a difference between flexible and fixed space will only be 
%    noticeable when features are used.
%    
%    \label{Lalignment}
%    The position of the output can be aligned left, right
%    or centered on the page by |\alignment{|\meta{position}|}| 
%    with the \meta{position} parameter |left|, |center| or 
%    |right|.
%    
%    
%    
%    \subsubsection{Fingerprinting}
%    
%    \label{fingerprint}
%    
%    \label{Lfingerprint}
%    An easy way to gain an overview on complete alignments is 
%    provided by displaying a so called alignment `fingerprint'.
%    In this style the whole sequence can be shown in one line. Due to
%    the lacking space the residue names are hidden and the shaded
%    boxes are reduced to thin vertical colored lines. The command
%    |\fingerprint{|\meta{res. per line}|}| takes one argument stating
%    the desired number of residues per line, e.\,g. |\fingerprint{1000}|.
%    All \TeXshade{} commands are compatible with |\fingerprint|, 
%    i.\,e. all shading modes are applicable for displaying overviews
%    on similarity or every functional aspect. Also, all kinds of
%    labeling---as described in the following---work with this
%    command.
%    
%    
%    \subsection{Individual shading and labeling of sequence stretches}
%    
%    Computer calculated conservation shading is informative---but
%    even more information can be visualized by additional labeling
%    of positions and regions of interest with different colors,
%    text styles or graphical marks and descriptive text. All this 
%    is provided by easy to handle \TeXshade{} commands.
%    
%    
%    \subsubsection{Shading of regions and blocks}
%    \label{shaderegion}
%    
%    \label{Lshaderegion}
%    Besides the shading calculated by \TeXshade{} any region can be
%    additionally shaded with a color specified by the user. This is very
%    useful to highlight secondary protein modification
%    sites such as phosphorylation or glycosylation sites, or longer
%    motifs for example protein/protein interaction sites or protein domains.
%    This is done with the following command:
%    \medskip
%    
%    \quad|\shaderegion{|\meta{seqref}|}{|\meta{selection}|}{|\meta{res.col.}|}{|\meta{shad.col.}|}|
%    \medskip
%    
%    Here, \meta{seqref} refers to the sequence by its name or number
%    within the alignment. The 
%    foreground and background colors can be set with the last two 
%    parameters. 
%    \meta{selection} can have three different formats
%    depending on whether (a) the user wants to select the positions manually,
%    (b) the selection should be based on a sequence motif to be found by \TeXshade{}, 
%    or (c) \TeXshade{} is supposed to select the residues based on 3D
%    coordinates provided by a PDB file.
%    \medskip
%    
%    (a) To select the residues manually, the user provides a position list of
%    the following format:
%    \medskip
%    
%    \quad|{|\meta{start1}..\meta{stop1}|,|\meta{start2}..\meta{stop2}|,|\ldots|,|\meta{start n}..\meta{stop n}|}|
%    \medskip
%    
%    In order to shade residue number 13 and the region
%    20--30 of sequence number 1 in red letters on a green background 
%    type the following command:
%    \medskip
%    
%    \quad|\shaderegion{1}{13..13,20..30}{Red}{Green}|
%    \medskip
%    
%    If the consensus is to be shaded use |consensus| as 
%    \meta{seqref}.
%    \medskip
%    
%    (b) A sequence motif can be given, which will be found and labeled. A simple 
%    example selecting all sequence motifs `NNAD' in sequence `1' would be:
%    \medskip
%
%    \quad|\shaderegion{1}{NNAD}{Red}{Green}|
%    \medskip
%
%    The definition can further include `X' for any residue and groups of residues in brackets
%    at uncertain positions. The example below will find all motifs in sequence `1'
%    that start with an asparagine, followed by any two residues, an acidic, a
%    basic residue, again any two residues and finally a glutamine.
%    \medskip
%    
%    \quad|\shaderegion{1}{NXX[DE][KR]XXQ}{Red}{Green}|
%    \medskip
%    
%    (c) In order to select positions based on the 3D structure a PDB
%    structure file is required. \TeXshade{} can select residues within
%    a given distance in \AA{} around a point, along a line, or above and
%    below a plane, which are defined by one to three residues. The points 
%    can be further specified to be the $\alpha$-carbon atom, i.\,e.\ the
%    protein backbone, or the most distant atom of the sidechain. Accordingly,
%    \meta{selection} has one of the formats:
%    \medskip
%
%    \quad|{point[|\meta{dist}|]:|\meta{file}|,|\meta{num}|[CA/side]}|
%    \medskip
%    
%    \quad|{line[|\meta{dist}|]:|\meta{file}|,|\meta{num1}|[CA/side],|\meta{num2}|[CA/side]}|
%    \medskip
%
%    \vbox{%    
%    \quad|{plane[|\meta{dist}|]:|\meta{file}|,|\meta{num1}|[CA/side],|\meta{num2}|[CA/side],|
%    \medskip
%
%    \hfill\meta{num3}|[CA/side]}|
%    }
%    \medskip
%    
%
%    Example: in order to select and shade as above all residues that are within an 8\ \AA{} 
%    sphere around the $\alpha$-carbon of residue 81 and the data are
%    provided in the PDB file |1J4N.pdb|, type:
%    \medskip
%    
%    \quad|\shaderegion{1}{point[8]:1J4N.pdb,81[CA]}{Red}{Green}|
%    \medskip
%    
%    Example: two points denote a line, hence, give two residues to select 
%    everything within 1\ \AA{} along the line between the $\alpha$-carbon of
%    residue 81 and the sidechain of residue 168 with:
%    \medskip
%    
%    \quad|\shaderegion{1}{line[1]:1J4N.pdb,81[CA],168[side]}|
%    \medskip
%
%    \hfill|{Red}{Green}|
%    \medskip
%    
%    Definition of a plane follows the same format but requires three points. 
%    If the optional parameters |[\meta{dist}]| and |[CA/side]| are not given,
%    \TeXshade{} assumes |[1]| and |[side]|, respectively.
%
%    In case one needs the position numbers of the selected residues for usage
%    in other applications, those can be either printed in the \TeX{} document with
%    |\printPDBlist{|\meta{selection}|}| or shown during the \TeX{} run with
%    |\messagePDBlist{|\meta{selection}|}|. The commands can be used outside of the 
%    \TeXshade{} environment. \label{LprintPDBlist}\label{LmessagePDBlist}
%    
%    Both selection formats, i.\,e. a manually given list and the 3D selection,
%    can be used with |\shadeblock| (see below), 
%    |\tintregion|, |\tintblock|, |\emphregion|, |\emphblock|, |\lowerregion|, |\lowerblock|,
%    |\frameblock| (all in \ref{Lframeblock}), and |\feature| (\ref{Lfeature}).
%
%
%    \label{Lshadeblock}
%    In analogy to |\shaderegion| which is restricted to a single
%    sequence, |\shadeblock| shades the corresponding region in all 
%    other sequences as well
%    except the consensus. If also the consensus is to be shaded 
%    define the region using |consensus| as \meta{seqref}.
%    \medskip
%    
%    \quad|\shadeblock{|\meta{seqref}|}{|\meta{selection}|}{|\meta{res.col.}|}{|\meta{shad.col.}|}|
%    \medskip
%    
%    
%    Another option is to change the whole set of shading colors for certain
%    sequence blocks, e.\,g.\ from |blues| to |reds| or self-defined color sets (see \ref{Lshadingmode}).
%    Therefore, the following command was implemented: 
%    \medskip
%    
%    \quad |\changeshadingcolors{|\meta{seqref}|}{|\meta{selection}|}{|\meta{name}|}|
%    \medskip \label{Lchangeshadingcolors}
%
%    Examples:
%    \medskip
%
%    \quad |\changeshadingcolors{1}{10..50}{reds}|
%    \medskip
%
%    \quad |\changeshadingcolors{AQP1}{[AS]NKD}{my_set}|
%    \medskip
%
%    etc.
%    
%    \subsubsection{Emphasizing, tinting, lowercasing, and framing}
%    
%    \label{Lemphregion}\label{Lemphblock}
%    If it is preferred to keep the calculated shading colors 
%    but distinct regions or blocks are yet to be emphasized one
%    can use the following commands to change the font style of 
%    such regions:
%    \medskip
%    
%    \quad|\emphregion{|\meta{seqref}|}{|\meta{selection}|}|
%    \medskip   
%    
%    and
%    \medskip
%    
%    \quad|\emphblock{|\meta{seqref}|}{|\meta{selection}|}|
%    \medskip   
%    
%    For the format possibilities of the \meta{selection} parameter please see \ref{Lshaderegion}.
%
%    \label{Lemphdefault}
%    Which style \TeXshade{} uses for emphasizing regions is defined by 
%    |\emphdefault{|\meta{style}|}|. Default setting is the
%    \emph{italics} font shape (set by |\emphdefault{it}|). In order to change
%    this setting choose one of the styles |bf, md, up, it, sl, rm, sf, tt|.
%    
%    Example: |\emphdefault{bf}|
%    \medskip
%    
%    \label{Ltintregion}\label{Ltintblock}
%    Further, it is possible to tint the region or block in question
%    or to switch the characters to lowercase
%    by using the commands (for example see hydropathy-figure on page 
%    \pageref{hydro}):
%    \medskip
%    
%    \quad|\tintregion{|\meta{seqref}|}{|\meta{selection}|}|
%    \medskip   
%    
%    and
%    \medskip
%    
%    \quad|\tintblock{|\meta{seqref}|}{|\meta{selection}|}|
%    \medskip   
%    
%    \label{Ltintdefault}
%    The level of tinting in the region in question can be set by 
%    |\tintdefault{|\meta{level}|}| with |weak|, |normal|, and
%    |strong| as possible \meta{level}s.
%    \medskip
%    
%    \label{Llowerregion}\label{Llowerblock}
%    \quad|\lowerregion{|\meta{seqref}|}{|\meta{selection}|}|
%    \medskip   
%    
%    and
%    \medskip
%    
%    \quad|\lowerblock{|\meta{seqref}|}{|\meta{selection}|}|
%    \medskip   
%    
%    
%    Another option is to draw a bounding box around the sequence block
%    in question (for an example see diversity mode-figure on page 
%    \pageref{frame}) with the 
%    command:\footnote{Thanks to Alan Robinson for inspiration.}
%    \medskip\label{Lframeblock}
%
%    \quad|\frameblock{|\meta{seqref}|}{|\meta{selection}|}{|\meta{color}|[|\meta{length}|]}|
%    \medskip
%
%    With the optional parameter the default line thickness of the frame can
%    be changed, example: |\frameblock{1}{10..20,50..70}{Red[2pt]}|
%
%
%
%
%    \subsubsection{Graphical labeling of sequence features}
%    
%    \label{feature}
%    
%    \label{Lfeature}
%    The |\feature| command is designed to fulfill most needs for the
%    graphical labeling of sequence stretches and the setting of descriptive 
%    text. It needs five parameters:
%    \medskip
%    
%    \quad|\feature{|\meta{position}|}{|\meta{seqref}|}{|\meta{selection}|}{|\meta{labelstyle}|}{|\meta{text}|}|
%    \medskip
%    
%    In the following paragraphs all possible parameter settings of 
%    this rather complex but mighty command are discussed in detail.
%    The parameter \meta{position} tells \TeXshade{} where to display
%    the feature label, i.\,e. on the top of the alignment (|top|),
%    or at the bottom (|bottom|). Further, there can be three more feature lines
%    ontop of the top feature line (|ttop|, |tttop|, and |ttttop|) or below the bottom
%    feature line (|bbottom|, |bbbottom|, |bbbbottom|). Thus, up to eight features
%    overlapping in eight different lines may be displayed. 
%    Depending on the content of the feature lines the gaps between 
%    them might be not satisfactory. 
%    \label{Ltopspace}\label{Lttopspace}
%    \label{Ltttopspace}\label{Lttttopspace}
%    \label{Lbottomspace}\label{Lbbottomspace}
%    \label{Lbbbottomspace}\label{Lbbbbottomspace}
%    Therefore, eight separate commands can be employed to change the 
%    space below |ttttop|, |tttop|, |ttop|, or |top| 
%    (|\topspace{|\meta{length}|}| etc.), and above 
%    |bottom|, |bbottom|, |bbbottom|, or |bbbbottom| 
%    (|\bottomspace{|\meta{length}|}| etc.). Use positive
%    values to further separate the lines, e.\,g. 
%    |\ttopspace{3mm}| or negative values to reduce the space, e.\,g.
%    |\bottomspace{-0.1in}|.
%    
%    The argument \meta{seqref} and the third 
%    parameter containing the \meta{selection} of the specified residues
%    are identical to the ones described before in several commands, e.\,g.
%    |\shaderegion| (\ref{shaderegion}). 
%    
%    New is the fourth parameter for the definition of the label style.
%    There are many possibilities like braces, helices, boxes, arrows, bars, any
%    fill character, bar graphs, color scales or even translations of the 
%    specified regions. 
%    \medskip
%
%    \textbf{Braces:}\\
%    In order to display an over- or underbrace as
%    a label use the parameter |{brace}|. Depending on the 
%    \meta{position} (|ttttop|, |tttop|, |ttop|, |top|, |bottom|,
%    |bbottom|, |bbbottom|, or |bbbbottom|) the respective brace is
%    displayed. The standard color of braces is 
%    black. It can be changed by an optional parameter directly after 
%    the definition of the symbol, e.\,g. |{brace[Red]}|.
%    \medskip
%
%    \textbf{Protein $\alpha$-Helices:}\\
%    The parameter |{helix}| will plot a symbolized $\alpha$-helix
%    as a label. The standard color of the helix spiral is 
%    black. It can be changed by an optional parameter directly after 
%    the definition of the symbol, e.\,g. |{helix[Red]}|.
%    \medskip
%
%    \textbf{Filling a stretch with a symbol:}\\
%    A region can be filled with any character for
%    labeling purposes using the parameter |{fill:|\meta{symbol}|}|.
%    The \meta{symbol} is freely selectable; the usage is like
%    in |\gapchar| (\ref{gapchar}). Do not use spaces before or after 
%    the expression \meta{symbol}; this will shift the symbols to the
%    respective direction. The standard color of the fill symbol is 
%    black. It can be changed by an optional parameter directly after 
%    the definition of the symbol, e.\,g. |{fill:$\bullet$[Red]}|.
%    
%    The |\feature| command does not like special characters in
%    text mode, e.\,g. |\dag|. One has to use the math version of
%    those symbols between |$|-signs. The following quite common
%    text symbols have also a math equivalent\footnote{Thanks to 
%    Darrell Conklin for reporting this problem}:
%    
%    \begin{center}
%    \begin{tabular}{cll}
%    \emph{symbol} & \emph{command} & \emph{description} \\ \hline
%    $\dagger$ & |$\dagger$| & dagger\\
%    $\ddagger$ & |$\ddagger$| & double dagger\\
%    $\mathparagraph$ & |$\mathparagraph$| & paragraph mark\\
%    $\mathsection$ & |$\mathsection$| & section mark\\
%    $\mathdollar$ & |$\mathdollar$| & dollar\\
%    $\lbrace$ & |$\lbrace$| & left brace\\
%    $\rbrace$ & |$\rbrace$| & right brace\\
%    \end{tabular}
%    \end{center}
%    \medskip
%
%    \textbf{Labeling restriction or protease cutting sites:}\\
%    If a label is needed that points between two residues, e.\,g.
%    for showing restriction sites, simply use the feature style
%    |{restriction[|\meta{color}|]}|. This will show a filled 
%    triangle with the tip right between the residues to be labeled,
%    e.\,g. |\feature{top}{1}{25..26}{restriction[Blue]}{EcoR I}|.
%
%    \medskip
%    
%    \textbf{Boxes:}\\
%    Boxed text is printed using the parameter |{box:|\meta{text}|}|.
%    By default black letters in a white framed box are displayed. In 
%    order to change these colors optional parameters can be included
%    in the argument: 
%    \medskip
%
%    \quad|{box[|\meta{framecolor,boxcolor}|][|\meta{length}|]:|\meta{text}|[|\meta{textcolor}|]}|.
%    \medskip
%
%    If the box frame and fill colors are the same it is sufficient to
%    use only this one color as an argument in the command. The optional
%    parameter \meta{length} defines the thickness of the box frame. If
%    this parameter is not set in the command the value from the 
%    |\featurerule{|\meta{length}|}| command (see below) is used.
%    \medskip
%    
%    Examples: 
%    \medskip
%
%    \quad|{box[Blue]:$\alpha$-helix[Yellow]}|
%    \smallskip
%
%    \quad|{box[Blue,Red]:$\alpha$-helix[Yellow]}|
%    \smallskip
%
%    \quad|{box[Blue,Red][2pt]:$\alpha$-helix[Yellow]}|
%    \medskip
%
%    \medskip
%    
%    \textbf{Horizontal bars and arrows:}\\
%    For displaying bars and arrows a simple selection scheme 
%    consisting of three consecutive characters is
%    used as the \meta{labelstyle} parameter. Each bar or arrow is
%    defined by its left end, the middle part, and the right end.
%    The following table gives some examples for the construction
%    of arrows and bars.
%    
%    \begin{center}
%    \begin{tabular}{cl}
%    middle & \\
%    \hbox to 1.6cm{\hss left end} \raisebox{1mm}{$\downarrow$} \hbox to 1.6cm{right end} & \\ \hline
%    |---|& plain bar \\
%    |===|& double bar \\
%    |-->|& right arrow \\
%    |'->|& right arrow with up hook \\
%    |<-|$\vert$ & left \emph{maps to} arrow \\
%    |<-o|& left arrow with ball at right end\\
%    |<=>|& double arrow, two heads \\
%    |,-,|& plain bar with down hooks\\
%    $\vert$|=|$\vert$ & double bar with vertical ends\\
%    |S-S|& labels disulfide bridges\\
%    \end{tabular}
%    \end{center}
%    
%    Combinations of the left-end-characters 
%    (|-=<',|$\vert$o), the middle-characters (|-=|),
%    and the right-end-characters (|-=>',|$\vert$o) are
%    allowed and produce the desired arrow or bar.
%    Adding an underscore in the definition for top feature hooks,
%    e.g. |{_,-,}| makes the hooks 
%    fully reach down to the alignment.
%    The color is changed as described above.
%    \label{Lfeaturerule} The thickness can be generally 
%    set by the separate command |\featurerule{|\meta{length}|}|
%    with any \TeX{} measure as \meta{length}, e.\,g.\ |\featurerule{3pt}|. 
%    This value is then used for all arrows, bars, and boxes (see above)
%    throughout the alignment. If an individual thickness for a 
%    particular arrow should be set one can add an optional
%    parameter to the \meta{labelstyle} parameter, e.g.
%    |{o->[Red][1mm]}|. Similar to the boxes described above, a text can be put on
%    the arrow or bar, e.\,g.\ |{<->[Red][1mm]:$\beta$-sheet[Blue]}|.
%
%
%    In \TeXshade{} versions before v1.9, the original \LaTeX{}-arrows 
%    were used. These have now been replaced by more modern looking
%    arrows with scalable line thickness. If the classical look is
%    requested, use |v| instead of |<| or |>| in the arrow definition,
%    e.\,g.\ |{--v}|, to get them back. The new arrow style makes use of
%    of the AMS math symbol font (amssymb.sty). Thus, in order to 
%    display the arrow heads correctly make sure that this style is    
%    present on your system (usually it is in a common \LaTeX{} installation).
%    \medskip
%
%    \textbf{Sequence translations:}\\
%    With the option |{translate}|, sequence stretches can be 
%    translated from nucleotide to peptide sequences as well as 
%    backtranslations from peptide to nucleotide sequences are
%    possible. Default setting for the translations is the standard
%    genetic code. Of course, the codons can be re-defined by the
%    user. The command \label{Lcodon}
%    |\codon{|\meta{amino acid}|}{|\meta{triplet1, \ldots, triplet n}|}|
%    has been implemented for this issue. The usage is simple. Replace
%    \meta{amino acid} by the single letter code of the amino acid
%    to be defined and add a list of triplets for this residue.    
%    Example definition for the amino acid \emph{alanine}:
%    \medskip
%    
%    \quad |\codon{A}{GCA,GCG,GCC,GCT,GCU,GCN}|
%    \medskip
%    
%    Note the last triplet in the list. It contains an ambiguity code
%    |N| which stands for \emph{any} nucleotide. This triplet has been 
%    added at the last position because the last triplet is used
%    for the generation of the backtranslated nucleotide sequence from
%    a peptide. Two files are included in the \TeXshade{}
%    distribution as examples (|standard.cod, ciliate.cod|). If you 
%    want to define a new genetic code store your commands in a file 
%    like the examples. Such files with the suffix |.cod| can be 
%    loaded in the \TeXshade{} environment by \label{Lgeneticcode}
%    |\geneticcode{|\meta{filename}|}|, e.\,g. |\geneticcode{ciliate}|.
%    Do not designate the suffix |.cod| in \meta{filename}. Please 
%    note, when inspecting the example files, that only the exchanges
%    compared to the standard code need to be defined in a new genetic code file. 
%    
%    When DNA sequences are translated to protein the resulting amino
%    acids are aligned to the second nucleotide of each triplet. 
%    It is more difficult to produce a satisfactory display of
%    backtranslated nucleotide sequences due to the lack of space. 
%    You need thrice as much space than the original peptide sequence,
%    because single letter amino acid code is translated to a triplet
%    code. Therefore, the user can choose from five display styles
%    for backtranslations depending on personal preferences:
%    \medskip\label{Lbacktranslabel}
%    
%    \quad |\backtranslabel[|\meta{size}|]{|\meta{style}|}|, with
%    \medskip
%    
%    \begin{tabbing}
%    \qquad\qquad|{|\meta{style}|}|\ \= = |{horizontal}|\\
%       \> = |{alternating}|\\
%       \> = |{zigzag}|\\
%       \> = |{oblique}|\\
%       \> = |{vertical}|
%    \end{tabbing}
%    
%    \meta{size} can be any \TeX{} size from |tiny| up to |Huge|, but 
%    |tiny| is recommended (and default setting). Translations
%    can be colored as all other labels, see above.
%    \medskip
%
%    \textbf{Bar graphs and color scales:}\label{Lgraphs}\\
%    Sequence related numeral data, such as hydropathy or solvent
%    accessibility data etc., can be shown in a feature line as bar graphs
%    or color scales. The data are (a) pre-defined or calculated by 
%    \TeXshade{} due to amino acid properties or conservation, (b) are 
%    provided in a separate file or (c) may be entered by hand in the 
%    |\feature| command. 
%
%    (a) Currently, three different
%    properties can be plotted, i.e. |hydrophobicity|, |molweight|, and 
%    |charge|. Further, the level of sequence conservation at the given
%    protein sequence stretch can be shown (|conservation|). See 
%    \ref{resweight} for selecting an appropriate \emph{residue weight table}.
%
%    (b) The format of a data file is simple: every value must 
%    appear in a separate line. Numbers and the Java-typical `NaN' for
%    `Not a Number' are permitted. Comments are allowed, because \TeXshade{}
%    ignores all lines starting with a letter except `NaN' lines (avoid 
%    `|-|' as the first
%    character of a comment line as this is interpreted as a negative number).
%    Make sure that there are as many values as positions defined as the 
%    sequence stretch in the feature command.
%    \TeXshade{} will read this file and determine the minimal and maximal 
%    values. These data are then normalized for plotting.
%    Due to \TeX's limited calculation capabilities no values above 10\,737
%    are allowed and the difference between minimum and maximum must not
%    exceed this very number. Values below 0.001 may be susceptible to major
%    rounding errors. Thus, try to provide your data already normalized to
%    moderate scales, e.g. 0.0\,--\,1.0 or -100\,--\,100. 

%    (c) Data which is  
%    directly entered in the |\feature| command must be normalized to integer
%    values with a maximal difference of 100 between the highest and lowest
%    value, e.g. -50\,--\,50 or 0\,--\,100.
%
%    For (b) and (c), the range to be plotted can be set by hand as an optional parameter
%    in the |\feature| command. This can be necessary when the data file
%    contains values between e.g. $-0.44$ and $0.87$. Without help \TeXshade{} 
%    will assume $-0.44$ as minimum and $0.87$ as maximum. But if the actual 
%    range to be plotted should be $-1.0$\,--\,$1.0$ this needs to be set
%    manually, see examples below. Be aware of the fact, that if you 
%    define a scale by hand, which is more narrow than the values of the
%    input, this will stretch the bars accordingly. It is NOT recommended
%    to use this method for stretching bars vertically. Instead another
%    command has been introduced. 
%    \label{Lbargraphstretch}\label{Lcolorscalestretch}
%    The plotted bars can be stretched by a factor if the appearance is
%    not as desired: |\bargraphstretch{|\meta{factor}|}|. Here, the factor 
%    is multiplied with the bar length, e.g |\bargraphstretch{2}| will double
%    the bar height, |\bargraphstretch{0.5}| will make them half as high.
%    Similarly, color scales can be stretched vertically with 
%    |\colorscalestretch{|\meta{factor}|}|.
%
%    The default color of bar graphs is gray and can be changed by an 
%    optional parameter at the end of the |label| definition. Further, an optional 
%    background color can be chosen for the bars. Doing so will visualize
%    the maximal bar extension.
%
%    Default for
%    color scales is a 5\% gray scale from very light gray to black (|Gray|).
%    More colorful scales have been implemented, i.e. |BlueRed|, |RedBlue|, |GreenRed|,
%    |RedGreen|, |ColdHot| and |HotCold|, the latter two being particularly 
%    useful for ranges from negative to positive values. Further, a scale called |T-Coffee|
%    is available if |T-Coffee| shading information has been imported as the 
%    |\shadingmode| \ref{Lshadingmode}.
%
%
%
%    The general format of this feature label definition for bar graphs is:
%    \medskip
%
%    \quad |{bar[|\meta{min}|,|\meta{max}|]:|\meta{properties/file/data}|[|\meta{color(,bgcolor)}|]}|
%    \medskip
%
%    and for color scales:
%    \medskip
%
%    \quad|{color[|\meta{min}|,|\meta{max}|]:|\meta{properties/file/data}|[|\meta{scale}|]}|
%    \medskip
%
%    Some examples:
%    \medskip
%
%    \qquad |{bar:conservation}|
%    \medskip
%
%    \qquad |{bar:conservation[T-Coffee]}|
%    \medskip
%
%    \qquad |{bar:hydrophobicity}|
%    \medskip
%
%    \qquad |{bar:charge[Red]}|
%    \medskip
%
%    \qquad |{bar:molweight[Red,Gray10]}|
%    \medskip
%
%    \qquad |{bar:10,20,30,40,50[Red]}|
%    \medskip
%
%    \qquad |{bar[-20,40]:-10,0,10,20,30[Red,Gray10]}|
%    \medskip
%
%    \qquad |{bar:data.txt}|
%    \medskip
%
%    \qquad |{bar[-10,10]:data.txt[Red,Gray10]}|
%    \medskip
%
%    \qquad |{color:conservation[BlueRed]}|
%    \medskip
%
%    \qquad |{color:conservation[T-Coffee]}|
%    \medskip
%
%    \qquad |{color:hydrophobicity[GreenRed]}|
%    \medskip
%
%    \qquad |{color:charge}|
%    \medskip
%
%    \qquad |{color:molweight}|
%    \medskip
%
%    \qquad |{color[-10,10]:data.txt[ColdHot]}|
%    \medskip
%
%    \qquad |{color[-0.1,0.1]:otherdata.txt[ColdHot]}|
%    \medskip
%
%    See also the example output in section \ref{graphs} on page
%    \pageref{graphs}.
%
%    \medskip
%
%    \textbf{No graphical label, only text:}\\
%    If no graphical label is
%    wanted the fourth parameter of |\feature| can be empty
%    braces.
%    \medskip
%    
%    Finally, the fifth parameter of the |\feature| command contains 
%    the descriptive text
%    for the labeled region. Type whatever you want incl. symbols and
%    math chars. The text field can also contain sequence translations.
%    In this case just set \meta{text} = |{translate}|. There is a 
%    command for setting the size and style of backtranslated sequences
%    in the feature \meta{text} which corresponds to the one
%    described above: \label{Lbacktranstext}
%    \medskip
%    
%    \quad |\backtranstext[|\meta{size}|]{|\meta{style}|}|
%    \medskip    
%    
%    Again, the color can be set by an
%    optional parameter appended to the text. For how to change the 
%    font size of text or symbols in the feature style line 
%    (|featurestyles|) or the in descriptive text line (|features|)
%    see section \ref{Lsetsize}, page \pageref{Lsetsize}.
%    
%    Another set of commands can be used to set a name for a feature
%    line, which is printed together with the sequence names at the
%    left or right side of the alignment, i.\,e.\ 
%    \medskip
%
%                    \label{Lshowfeaturename} \label{Lshowfeaturestylename}
%                    \label{Lhidefeaturename} \label{Lhidefeaturestylename}
%                    \label{Lhidefeaturenames} \label{Lhidefeaturestylenames}
%    \quad |\showfeaturename{|\meta{ttttop...bbbbottom}|}{|\meta{name}|}|,
%    \medskip
%
%    \quad |\showfeaturestylename{|\meta{ttttop...bbbbottom}|}{|\meta{name}|}|
%    \medskip
%
%    \quad |\hidefeaturename{|\meta{ttttop...bbbbottom}|}|
%    \medskip
%
%    \quad |\hidefeaturestylename{|\meta{ttttop...bbbbottom}|}|
%    \medskip
%
%    \quad |\hidefeaturenames|, and |\hidefeaturestylenames|.
%    \medskip
%
%    Using |\showfeaturename| will print the name in the same line as
%    the descriptive text of the feature whereas | \showfeaturestylename|
%    will put the name in the same line as the feature symbols.
%
%    The color of such names can be generally changed with \label{Lfeaturenamescolor}
%    \medskip
%
%    \quad |\featurenamescolor{|\meta{color}|}| and \label{Lfeaturestylenamescolor}
%    \medskip
%
%    \quad |\featurestylenamescolor{|\meta{color}|}| 
%    \medskip
%
%    or individually with \label{Lfeaturenamecolor}
%    \medskip
%
%    \quad |\featurenamecolor{|\meta{ttttop...bbbbottom}|}{|\meta{color}|}| and \label{Lfeaturestylenamecolor}
%    \medskip
%
%    \quad |\featurestylenamecolor{|\meta{ttttop...bbbbottom}|}{|\meta{color}|}| 
%    \medskip
%
%    See section \ref{colors} for how to select colors in \TeX{}shade.
%
%    Font styles can be set as usual (see section \ref{Lsetfamily}), e.\,g.\ 
%    \medskip
%    
%    |\setsize{featurenames}{large}| or |\featurestylenamesrm| etc. 
%    \medskip
%
%
%    
%    Examples for the appearance of features are given in the 
%    overview section (\ref{over}), see:
%    \medskip
%    
%    \emph{similarity mode} (\ref{similar}): fill-character; here, only 
%         one position is labeled. It is also possible to label a longer 
%         stretch, then, the character is printed several times to
%         fill the specified region.
%         
%    \quad|\feature{top}{1}{93..93}{fill:$\downarrow$}{first...}|   
%     
%    \quad|\feature{bottom}{1}{98..98}{fill:$\uparrow$}{second...}|    
%    \medskip
%    
%    \emph{T-Coffee mode} (\ref{TCoffee}): |T-Coffee| color scale
%    
%    \quad|\feature{top}{1}{30..63}{color:conservation[T-Coffee]}{}|
%
%    \quad|\showfeaturestylename{bottom}{cons}|
%    \medskip
%    
%    \emph{diversity mode} (\ref{diverse}): frames, text only
%    
%    \quad|\feature{top}{1}{77..109}{}{AQP2 species variants}|
%    
%    \quad|\frameblock{1}{82..82,106..106}{Red[1pt]}|
%    \medskip
%    
%    \emph{functional mode} (\ref{func}): bar graph, color scale, tinting, box, arrow, 
%                                         translation, brace, helix
%    
%    \quad|\feature{top}{3}{153..165}|
%
%    \quad\quad\quad\quad\quad|{bar[-50,50]:-50,-45,-40,...,40,45,50}{}|
%    \medskip
%
%    \quad|\feature{top}{3}{167..186}|
%
%    \quad\quad\quad\quad\quad|{color:5,10,15,...,90,95,100[ColdHot]}{}|
%    \medskip
%
%    \quad |\feature{top}{1}{158..163}{brace}{tinted}|
%
%    \quad|\tintblock{1}{158..163}|
%    \medskip
%
%    \quad|\feature{top}{1}{138..157}|
%    
%    \quad\quad\quad\quad\quad|{box[Blue,Red][0.5pt]:$\alpha$-helix[Yellow]}|
%    
%    \quad\quad\quad\quad\quad|{transmembrane domain 4}|
%    
%    \quad|\feature{top}{1}{164..170}{o->[Red]}{trans. dom. 5}|
%    
%    \quad|\feature{top}{1}{158..163}{translate[Blue]}{}|
%    
%    \quad|\backtranslabel{oblique}|
%    
%    \quad|\feature{bottom}{1}{158..163}|
%    
%    \quad\quad\quad\quad\quad|{brace[Blue]}{loop D[Blue]}|
%    \medskip
%
%    \quad|\feature{top}{1}{138..157,164..170}{helix}{membr.}|
%
%    \quad|\feature{top}{1}{158..163}{---}{loop}|
%
%    \quad|\featurerule{1mm}|
%    \medskip
%    
%    \emph{bar graphs and color scales} (\ref{graphs}): sequence conservation, 
%                                         charge, molecular weight, hydrophobicity
%    
%    \quad|\feature{ttop}{1}{138..170}{bar:conservation}{}|
%
%    \quad|\feature{top}{1}{138..170}{color:charge}{}|
%
%    \quad|\feature{bottom}{1}{138..170}|
%
%    \quad\quad\quad\quad\quad|{color:molweight[ColdHot]}{}|
%
%    \quad|\feature{bbottom}{1}{138..170}|
%
%    \quad\quad\quad\quad\quad|{bar:hydrophobicity[Red,Gray10]}{}|
%    \medskip 
%    
%    \subsubsection{Including secondary structure information}
%    
%    \label{structure}
%    
%    \label{LincludeDSSP}
%    \label{LincludeSTRIDE}
%    \label{LincludePHDsec}
%    \label{LincludePHDtopo}
%    \label{LincludeHMMTOP}
%    The DSSP [9], STRIDE [10], PHD [11] and HMMTOP [12] algorithms produce 
%    secondary protein structure predictions. PHD files contain both,
%    secondary structure information and topology data. This information can be
%    displayed in an alignment by one of the commands:
%    \bigskip
%    
%    \begin{tabular}{ll}
%    |\includeDSSP| & sec. structure calculated by DSSP\\
%    
%    |\includeSTRIDE| & sec. structure calculated by STRIDE \\
%    
%    |\includePHDsec| & sec. structure calculated by PHD \\
%    
%    |\includePHDtopo| & topology data calculated by PHD \\
%    
%    |\includeHMMTOP| & topology data calculated by HMMTOP \\
%    \end{tabular}
%    \bigskip
%    
%    The syntax is |\includeDSSP{|\meta{seqref}|}{|\meta{filename}|}|,
%    with |seqref| indicating the number or name of the sequence for which 
%    the secondary structure data is calculated and |filename| designating the 
%    corresponding structure file to be included.
%    
%    Several types of secondary structures are predicted by these 
%    programs; in order to designate them in \TeXshade{} use the names
%    from the right column:
%    
%    \begin{center}
%    \begin{tabular}{ll}
%    secondary structure & designation\\[3mm]
%    \emph{DSSP and STRIDE} & \\[2mm]
%    4-helix ($\alpha$-helix) & |alpha| \\
%    isolated $\beta$-bridge  & |bridge| \\
%    extended strand ($\beta$-strand) & |beta| \\
%    3-helix (3$_{10}$-helix) & |3-10| \\
%    5-helix ($\pi$-helix)    & |pi| \\
%    H-bonded turn            & |turn| \\[3mm]
%    \emph{PHDsec} & \\[2mm]
%    helix                    & |alpha| \\
%    sheet                    & |beta| \\[3mm]
%    \emph{PHDtopo and HMMTOP} & \\[2mm]
%    internal region          & |internal| \\
%    external region          & |external| \\
%    transmembrane domain     & |TM| \\
%    \end{tabular}
%    \end{center}
%    
%    \label{LshowonDSSP}
%    \label{LshowonSTRIDE}
%    \label{LshowonPHDsec}
%    \label{LshowonPHDtopo}
%    \label{LshowonHMMTOP}
%    \label{LhideonDSSP}
%    \label{LhideonSTRIDE}
%    \label{LhideonPHDsec}
%    \label{LhideonPHDtopo}
%    \label{LhideonHMMTOP}
%    By default all three types of helices and the strands are 
%    displayed whereas turns and bridges are skipped. If it is 
%    desired to shown them as well, call for example |\shownonDSSP{bridge,turn}|.
%    In analogy to this example all structure features can be activated
%    in DSSP, STRIDE, PHDsec, PHDtopo and HMMTOP. In order to hide
%    certain structure types use for example |\hideonDSSP{3-10,pi}|.
%    
%    The DSSP format has two columns of sequence numberings. The first 
%    column is consecutive, whereas the second column contains the 
%    actual sequence numbering. This can be different from the first
%    column when sequence parts are missing in the DSSP file. One can
%    choose which column will be read by \TeXshade{} by 
%    \label{LfirstcolumnDSSP} \label{LsecondcolumnDSSP}
%    `|\firstcolumnDSSP|' and |\secondcolumnDSSP|'. The second column 
%    is still default.
%    
%    The HMMTOP algorithm can present its results as plain text or 
%    as HTML---plain text needs to be selected here. Further, the 
%    output can be formatted in a single line or in an extended form 
%    (see the HMMTOP documentation). Both can be read and interpreted 
%    by \TeXshade{}. Importantly, HMMTOP files can contain topology
%    predictions of multiple sequences. \TeXshade{} tries to find
%    the correct data based on the respective sequence name. If the
%    sequence name is not found in the file, the first topology data
%    is used. Using an optional parameter (number of the prediction
%    in the file or name) one can define which data from the file is
%    to be used:
%    \medskip
%    
%    |\includeHMMTOP{|\meta{seqref in texshade}|[|\meta{seqref in file}|]}{|\meta{filename}|}|
%    \medskip
%    
%    PHD predictions: when starting the PHD software do not
%    restrict the prediction to secondary structure or topology alone.
%    This leads to changes in the PHD output file which are not
%    correctly interpretable by \TeXshade{} due to ambiguities. There
%    is no way around it---thus, run the full prediction.
%
%    Now, some information on how \TeXshade{} extracts and displays
%    secondary structure features. In short, it is a two step process.
%    First, \TeXshade{} analyzes the secondary structure file and 
%    extracts all necessary data. This data is converted into a 
%    format which is readable and processable by \TeXshade{} using the 
%    |feature| command (see \ref{feature}). This command allows one to 
%    label sequence stretches graphically. For a detailed explanation
%    see the indicated reference. A list of feature commands is saved
%    in a file with the ending `|.sec|' for DSSP, STRIDE and PHDsec
%    or `|.top|' for PHDtopo. Then, in a second step, this file is loaded 
%    again and executed. When \TeXshade{} encouters this file a
%    second time, i.\,e. in a second \TeX{} run, it uses the already
%    existing file for the output. The great advantage of this method 
%    is its flexibility. Due to the simple reason that the feature 
%    file can be edited in the meantime. Thus, the user has the 
%    ability to change the computer-generated file according to his
%    personal needs. On the other hand, one can force \TeXshade{} to
%    write a new file every time by the optional argument |[make new]| in the 
%    include command, e.\,g. |\includePHDsec[make new]{1}{AQP.phd}|.
%    
%    \label{Lappearance}
%    Finally, the appearance of the feature labels can be assigned by 
%    the command 
%    
%    |\appearance{|\meta{filetype}|}{|\meta{type}|}{|\meta{position}|}{|\meta{labelstyle}|}{|\meta{text}|}|.
%    
%    Here, \meta{filetype} stands for one of the following secondary structure 
%    file types: |DSSP|, |STRIDE|, |PHDsec|, |PHDtopo| or |HMMTOP| and 
%    \meta{type} designates the secondary structure type as shown in 
%    the right column of the table above. The other
%    arguments \meta{position}, \meta{labelstyle} and \meta{text} 
%    are almost as described in \ref{feature}. 
%    \label{Lnumcount}
%    \label{Lalphacount}
%    \label{LAlphacount}
%    \label{Lromancount}
%    \label{LRomancount}
%    One further possibility
%    is to include internal counters for each secondary structure type.
%    Just add one of the following commands 
%    to the text in the feature description. 
%    
%    \begin{center}
%    \begin{tabular}{ll}
%    \emph{counter} & \emph{display} \\[2mm]
%    |\numcount|     & 1, 2, 3 \ldots \\
%    |\alphacount|   & a, b, c \ldots \\
%    |\Alphacount|   & A, B, C \ldots \\
%    |\romancount|   & i, ii, iii \ldots \\
%    |\Romancount|   & I, II, III \ldots \\
%    \end{tabular}
%    \end{center}
%    
%    Examples: 
%    
%    \quad|\appearance{DSSP}{alpha}{ttop}|
%    
%    \quad\quad\quad\quad\quad\quad\quad|{-->}{$\alpha$-helix~\Alphacount}|
%    
%    \quad|\appearance{PHDtopo}{TM}{bottom}|
%    
%    \quad\quad\quad\quad\quad\quad\quad|{box[Blue]:TM\numcount[Yellow]}{}|
%    
%    
%    \subsection{Displaying and building legends}
%    
%    \label{Lshowlegend}\label{Lhidelegend}\label{Lmovelegend}
%    \label{Lgermanlanguage}\label{Lenglishlanguage}\label{Llegendcolor}
%    \label{Lspanishlanguage}
%    For each predefined shading mode \TeXshade{} can print an appropriate
%    legend to explain the used
%    shading colors. The commands |\showlegend| and |\hidelegend|
%    display or clear the legend at the end of the alignment. 
%    The legend is displayed by default beneath the first residue 
%    of the last alignment line. The location can be changed by 
%    |\movelegend{|\meta{x-offset}|}{|\meta{y-offset}|}|. Both 
%    parameters require a \TeX{} length, e.\,g. |\movelegend{5cm}{-2cm}|
%    moves the legend 5\,cm to the right and 2\,cm up.
%    
%    The language for the descriptions is english by default;
%    if the |\german.sty| package is active legend texts are in 
%    german. So far, german, spanish and english are implemented. With the 
%    commands |\germanlanguage|, |\spanishlanguage| and |\englishlanguage| 
%    switching between the languages 
%    is made possible. For the addition of other languages contact me.
%    Finally, the color of the describing legend texts can be set 
%    with the command |\legendcolor{|\meta{color}|}|.
%    
%    User defined legends are easily built with the following command
%    \label{Lshadebox}|\shadebox{|\meta{color}|}|. Use this command outside 
%    the \TeXshade{} environment, e.\,g. in the text or in the caption. As 
%    \meta{color} any color can be designated (see section \ref{colors}) or 
%    one of the following parameters:
%    
%    \begin{itemize}
%    \item |nomatch| = the color used for nonmatching residues
%    
%    \item |similar| = the color used for similar residues
%    
%    \item |conserved| = the color used for conserved residues
%    
%    \item |allmatch| = the color used for highly conserved residues 
%                       (if |\allmatchspecial| is active)
%    
%    \end{itemize}
%    
%    The command simply prints a shaded box in the specified color 
%    then a describing text can be appended. Examples:
%    \medskip
%    
%    \quad|\shadebox{nomatch}---nonmatching residues| 
%    
%    \quad|\shadebox{similar}: similar residues| 
%    
%    \quad|\shadebox{conserved}~conserved residues| 
%    
%    \quad|\shadebox{Yellow}\quad PKA phosphorylation sites|
%    
%    
%    
%    \subsection{Adding captions to the alignment}
%    
%    Since \TeXshade{} v1.5 captions can be added to the alignment. 
%    So far, captions were difficult to use when the alignment was
%    bigger than one page and therefore did not fit into a 
%    figure environment. The \TeXshade{} captions behave exactly as 
%    normal figure captions. They
%    adopt their style, use the figure counter number and appear in
%    the list of figures as any other figure.
%    
%    The usage is slightly different from normal captions but 
%    intuitive: \label{Lshowcaption}
%    \medskip
%    
%    \quad |\showcaption[|\meta{position}|]{|\meta{text}|}|
%    \medskip
%    
%    The optional \meta{position} tells \TeXshade{} to put the caption on
%    |top| or at the |bottom| of the alignment. If nothing is stated here
%    the caption will appear at the bottom. The parameter
%    \meta{text} just holds the caption text as in the normal |\caption|.
%    The command can be used at any position within the |texshade|
%    environment. A simple example would be: 
%    \medskip
%    
%    \quad |\showcaption{A beautiful \TeXshade{} alignment.}|
%    \medskip
%
%    \label{Lshortcaption}
%    In order to show a short version of the caption in the 
%    "List of Figures" the |\shortcaption{|\meta{short caption text}|}| 
%    command can be used.
%    
%    \subsection{Font handling}
%    
%    \subsubsection{Changing font styles}
%    
%    \label{Lsetfamily}\label{Lsetseries}
%    \label{Lsetshape}\label{Lsetsize}
%    The font styles for the numbering, the sequence names,
%    the sequence residues, the descriptive feature texts
%    and the legends can be changed by several commands.
%    \medskip
%    
%    \quad|\setfamily{|\meta{text}|}{|\meta{family}|}|
%    
%    \quad|\setseries{|\meta{text}|}{|\meta{series}|}|
%    
%    \quad|\setshape{|\meta{text}|}{|\meta{shape}|}|
%    
%    \quad|\setsize{|\meta{text}|}{|\meta{size}|}|
%    \medskip
%    
%    The first parameter selects the text whose style is to be
%    changed. Possible first parameters are
%    |numbering|, |names|, |residues|, |features|, |featurestyles|,
%    |hideblock|, |ruler|, and |legend|.
%    \medskip
%    
%    The style is set by the second parameter:
%    
%    \begin{center}
%    \begin{tabular}{lll}
%    command & \meta{2. parameter} & \\
%    \hline
%    |\setfamily| & |rm| & modern roman font family \\
%           & |sf| & sans serif font family \\
%           & |tt| & typewriter font family \\ \hline
%    |\setseries| & |bf| & bold face series \\
%           & |md| & normal series \\ \hline
%    |\setshape|  & |it| & italics shape \\
%           & |sl| & slanted shape \\
%           & |sc| & small capitals shape \\
%           & |up| & upright shape \\ \hline
%    |\setsize|   & |tiny| & the known \TeX{} sizes \\
%           & |scriptsize| & \\
%           & |footnotesize| & \\
%           & |small| & \\
%           & |normalsize| & \\
%           & |large| & \\
%           & |Large| & \\
%           & |LARGE| & \\
%           & |huge| & \\
%           & |Huge| & \\ \hline
%    \end{tabular}
%    \end{center}
%    
%    Example: |\setfamily{features}{it} \setseries{features}{bf}|
%    \medskip
%    
%    \label{Lsetfont}
%    With the command
%    \medskip 
%    
%    \quad|\setfont{|\meta{text}|}{|\meta{family}|}{|\meta{series}|}{|\meta{shape}|}{|\meta{size}|}|
%    \medskip
%    
%    all four font attributes of one \meta{text} can be changed 
%    simultaneously. The order of the parameters is as indicated.
%    \medskip
%    
%    Example: |\setfont{features}{rm}{it}{bf}{normalsize}|
%    \medskip
%    
%    Further, short commands are provided to change single font
%    attributes quickly. The following commands set attributes
%    of feature texts. 
%    \medskip
%    \enlargethispage{\baselineskip}
%    
%    \quad |\featuresrm| \quad |\featurestiny| \label{Lfeaturesrm}
%    
%    \quad |\featuressf| \quad |\featuresscriptsize|
%    
%    \quad |\featurestt| \quad |\featuresfootnotesize|
%    
%    \quad |\featuresbf| \quad |\featuressmall|
%    
%    \quad |\featuresmd| \quad |\featuresnormalsize|
%    
%    \quad |\featuresit| \quad |\featureslarge|
%    
%    \quad |\featuressl| \quad |\featuresLarge|
%    
%    \quad |\featuressc| \quad |\featuresLARGE|
%    
%    \quad |\featuresup| \quad |\featureshuge|
%    
%    \quad |           | \quad |\featuresHuge|
%    \medskip
%    
%    Corresponding sets are provided for the
%    numbering (|\numberingrm| etc.), 
%    featurestyles (|featurestylesrm| etc.), names (|\namesrm| etc.), 
%    featurenames (|\featurenamesrm| etc.),
%    featurestylenames (|\featurestylenames| etc.),
%    residues (|\residuesrm| etc.),
%    hideblock labels (|hideblockrm| etc.),
%    rulers (|\rulerrm| etc.), and
%    legend texts (|legendrm| etc.).
%     
%    
%    \subsubsection{Using PostScript fonts}
%    
%    As already mentioned \TeXshade{} makes intensive use of
%    \textsc{PostScript} for shading. Now, that 
%    \textsc{PostScript} output is active anyway, including \textsc{PostScript}
%    fonts is very easy. Just declare in the document header
%    \medskip
%    
%    \quad |\usepackage{|\meta{PS-font}|}|.
%    \medskip
%    
%    
%    The typewriter font of \TeX{} is always a topic of discussions.
%    By including the package |\usepackage{courier}| \TeX's 
%    typewriter font is replaced by the widely accepted \textsc{Courier}.
%    Have a look into the directory |..texinputs:latex:psnfss|; there, 
%    some styles are located which exchange the common \TeX{} fonts by 
%    \textsc{PostScript} fonts, e.\,g.\ |avant.sty|, |bookman.sty|,
%    |chancery.sty|, |courier.sty|, |helvet.sty| or |utopia.sty|.
%    Depending on the style used the |\rmdefault|-, |\sfdefault|-,
%    and |\ttdefault| fonts are substituted partly or completely.
%    Thus, |courier.sty| for instance exchanges only the typewriter font, 
%    whereas |bookman.sty| sets \textsc{Bookman} as |\rmdefault|,
%    \textsc{Avantgarde} as |\sfdefault| and \textsc{Courier} as
%    |\ttdefault|.
%    
%    For further information see \textsc{Tomas Rokicki}'s 
%    |dvips| manual [13].
%    
%    
%    
%    
%    \subsection{Goodies}
%    
%    The following commands give information on sequence properties,
%    such as molecular weight, charge or similarity data. They can
%    be used outside the |texshade| environment directly in the 
%    document or in a caption text.
%
%    \subsubsection{Molweight and charge}
%    
%    \label{molcharge}
%    
%    \label{Lmolweight}\label{Lcharge}
%    During the process of sequence setting \TeXshade{}
%    sums up the molecular weight and charge of the
%    aligned proteins. This data can be accessed by the 
%    following commands.
%    \medskip
%    
%    \quad|\molweight{|\meta{seqref}|}{|\meta{Da/kDa}|}|
%    
%    \quad|\charge{|\meta{seqref}|}{|\meta{i/o/N/C}|}|
%    \medskip
%    
%    The first parameter \meta{seqref} selects the sequence. The
%    second parameter in the |\molweight| command allows one to
%    switch the units between Dalton (|Da|) and kilo-Dalton 
%    (|kDa|). The |\charge| command needs the second parameter
%    for the correct consideration of the charged protein termini.
%    Thus, `|i|' refers to internal sequences, `|o|' to the 
%    overall charge, `|N|' to N-terminal sequence parts, and
%    `|C|' to the C-terminal end of a protein.
%    \medskip
%    
%    Example: \quad Charge: |\charge{1}{o}|; Weight: |\molweight{1}{Da}|
%    
%    \subsubsection{Similarity/identity data and tables}
%    
%    \label{simtable}
%    The degree of similarity and identity in percent for any two
%    sequences in the displayed alignment section can be read out 
%    with the commands
%    \label{Lpercentsimilarity} \label{Lpercentidentity} 
%    |\percentsimilarity{|\meta{seqref1}|}{|\meta{seqref2}|}| and
%    |\percentidentity{|\meta{seqref1}|}{|\meta{seqref2}|}|.
%    
%    Using the example alignment on page \pageref{simtableEx} and
%    typing outside the |texshade| environment in the document text the
%    following phrase:
%    \bigskip
% 
%    \quad\quad |AQP1 and AQP2 share a sequence similarity|
%
%    \quad\quad |of \percentsimilarity{AQP1.pro}{AQP2.pro}\%| 
%    \bigskip
% 
%    will result in the text:
%    \bigskip
%
%    \quad\quad AQP1 and AQP2 share a sequence similarity of 69.6\%
%    \bigskip
%    
%    Likewise, |\percentidentity{1}{2}| will give the value |48.4|; note that 
%    sequences can be referred to by their number or their assigned name. The
%    percent value is calculated by dividing the number of identical or similar
%    positions, respectively, by the total of non-gap positions shared by both 
%    sequences. Here, only the part of the alignment is taken into account that
%    is actually displayed. Two residues are considered similar when this is 
%    defined by the command |\pepsims| (see page \pageref{Lpepsims}).
%
%    A full similarity/identity table showing values for all sequences of the
%    alignment can be set using |\similaritytable|.\label{Lsimilaritytable} 
%    The labels and number format will be adjusted according to the language 
%    settings (\ref{Lgermanlanguage}).
%    \medskip
% 
%    Example (see section \ref{simtableEx}):
%    \medskip
%
%    \vbox{%
%    \quad |\begin{center}|
%    \medskip
%
%    \quad\quad|\similaritytable|
%    \medskip
%
%    \quad |\end{center}|
%    }
%    \bigskip
%
%    The command generates a valid \LaTeX{} |tabular|
%    environment, which can be embedded into a |table| environment, e.g.
%    \medskip
%    
%    \vbox{%
%    \quad |\begin{table}[htdp]|
%
%    \quad |\caption{Text ...}|
%
%    \quad |\begin{center}|
%    \medskip
%
%    \quad\quad|\similaritytable|
%    \medskip
%
%    \quad |\end{center}|
%
%    \quad |\end{table}|
%    }
%    
%    \newpage
%    \section{The PostScript color selection scheme}
%    
%    \label{colors}
%    
%    \textsc{PostScript} provides 64 standard colors. All these 
%    colors are predefined in the |color.sty|. Each color
%    has a pictorial name such as |Bittersweet| and a distinct 
%    composition, e.\,g.\ 0\% cyan + 75\% magenta + 100\% yellow + 
%    24\% black---the so-called CMYK scheme. \TeXshade{} enhances this
%    color scheme by gray scales in 5\% steps. 
%    The following colors and grays can be used in \TeXshade{} by 
%    simply declaring the name of the color in the respective 
%    command, e.\,g.\ |\consensuscolors|:
%    
%    
%    \begin{footnotesize}
%    \begin{tabbing}
%    \emph{name}\hspace{2.7cm}\= \emph{CMYK}\hspace{1.6cm}
%    \=\emph{name}\hspace{2.5cm}\= \emph{CMYK}\\
%    
%     \textcolor{GreenYellow}{$\bullet$}GreenYellow   \>{0.15,0,0.69,0}\>\textcolor{Yellow}{$\bullet$}Yellow        \>{0,0,1,0}\\
%     \textcolor{Goldenrod}{$\bullet$}Goldenrod     \>{0,0.10,0.84,0}\>\textcolor{Dandelion}{$\bullet$}Dandelion     \>{0,0.29,0.84,0}\\
%     \textcolor{Apricot}{$\bullet$}Apricot       \>{0,0.32,0.52,0}\>\textcolor{Peach}{$\bullet$}Peach         \>{0,0.50,0.70,0}\\
%     \textcolor{Melon}{$\bullet$}Melon         \>{0,0.46,0.50,0}\>\textcolor{YellowOrange}{$\bullet$}YellowOrange  \>{0,0.42,1,0}\\
%     \textcolor{Orange}{$\bullet$}Orange        \>{0,0.61,0.87,0}\>\textcolor{BurntOrange}{$\bullet$}BurntOrange   \>{0,0.51,1,0}\\
%     \textcolor{Bittersweet}{$\bullet$}Bittersweet   \>{0,0.75,1,0.24}\>\textcolor{RedOrange}{$\bullet$}RedOrange     \>{0,0.77,0.87,0}\\
%     \textcolor{Mahagony}{$\bullet$}Mahagony      \>{0,0.85,0.87,0.35}\>\textcolor{Maroon}{$\bullet$}Maroon        \>{0,0.87,0.68,0.32}\\
%     \textcolor{BrickRed}{$\bullet$}BrickRed      \>{0,0.89,0.94,0.28}\>\textcolor{Red}{$\bullet$}Red           \>{0,1,1,0}\\
%     \textcolor{OrangeRed}{$\bullet$}OrangeRed     \>{0,1,0.50,0}\>\textcolor{RubineRed}{$\bullet$}RubineRed     \>{0,1,0.13,0}\\
%     \textcolor{WildStrawberry}{$\bullet$}WildStrawberry\>{0,0.96,0.39,0}\>\textcolor{Salmon}{$\bullet$}Salmon        \>{0,0.53,0.38,0}\\
%     \textcolor{CarnationPink}{$\bullet$}CarnationPink \>{0,0.63,0,0}\>\textcolor{Magenta}{$\bullet$}Magenta       \>{0,1,0,0}\\
%     \textcolor{VioletRed}{$\bullet$}VioletRed     \>{0,0.81,0,0}\>\textcolor{Rhodamine}{$\bullet$}Rhodamine     \>{0,0.82,0,0}\\
%     \textcolor{Mulberry}{$\bullet$}Mulberry      \>{0.34,0.90,0,0.02}\>\textcolor{RedViolet}{$\bullet$}RedViolet     \>{0.07,0.90,0,0.34}\\
%     \textcolor{Fuchsia}{$\bullet$}Fuchsia       \>{0.47,0.91,0,0.08}\>\textcolor{Lavender}{$\bullet$}Lavender      \>{0,0.48,0,0}\\
%     \textcolor{Thistle}{$\bullet$}Thistle       \>{0.12,0.59,0,0}\>\textcolor{Orchid}{$\bullet$}Orchid        \>{0.32,0.64,0,0}\\
%     \textcolor{DarkOrchid}{$\bullet$}DarkOrchid    \>{0.40,0.80,0.20,0}\>\textcolor{Purple}{$\bullet$}Purple        \>{0.45,0.86,0,0}\\
%     \textcolor{Plum}{$\bullet$}Plum          \>{0.50,1,0,0}\>\textcolor{Violet}{$\bullet$}Violet        \>{0.79,0.88,0,0}\\
%     \textcolor{RoyalPurple}{$\bullet$}RoyalPurple   \>{0.75,0.90,0,0}\>\textcolor{BlueViolet}{$\bullet$}BlueViolet    \>{0.86,0.91,0,0.04}\\
%     \textcolor{Periwinkle}{$\bullet$}Periwinkle    \>{0.57,0.55,0,0}\>\textcolor{CadetBlue}{$\bullet$}CadetBlue     \>{0.62,0.57,0.23,0}\\
%     \textcolor{CornflowerBlue}{$\bullet$}CornflowerBlue\>{0.65,0.13,0,0}\>\textcolor{MidnightBlue}{$\bullet$}MidnightBlue  \>{0.98,0.13,0,0.43}\\
%     \textcolor{NavyBlue}{$\bullet$}NavyBlue      \>{0.94,0.54,0,0}\>\textcolor{RoyalBlue}{$\bullet$}RoyalBlue     \>{1,0.50,0,0}\\
%     \textcolor{Blue}{$\bullet$}Blue          \>{1,1,0,0}\>\textcolor{Cerulean}{$\bullet$}Cerulean      \>{0.94,0.11,0,0}\\
%     \textcolor{Cyan}{$\bullet$}Cyan          \>{1,0,0,0}\>\textcolor{ProcessBlue}{$\bullet$}ProcessBlue   \>{0.96,0,0,0}\\
%     \textcolor{SkyBlue}{$\bullet$}SkyBlue       \>{0.62,0,0.12,0}\>\textcolor{Turquoise}{$\bullet$}Turquoise     \>{0.85,0,0.20,0}\\
%     \textcolor{TealBlue}{$\bullet$}TealBlue      \>{0.86,0,0.34,0.02}\>\textcolor{Aquamarine}{$\bullet$}Aquamarine    \>{0.82,0,0.30,0}\\
%     \textcolor{BlueGreen}{$\bullet$}BlueGreen     \>{0.85,0,0.33,0}\>\textcolor{Emerald}{$\bullet$}Emerald       \>{1,0,0.50,0}\\
%     \textcolor{JungleGreen}{$\bullet$}JungleGreen   \>{0.99,0,0.52,0}\>\textcolor{SeaGreen}{$\bullet$}SeaGreen      \>{0.69,0,0.50,0}\\
%     \textcolor{Green}{$\bullet$}Green         \>{1,0,1,0}\>\textcolor{ForestGreen}{$\bullet$}ForestGreen   \>{0.91,0,0.88,0.12}\\
%     \textcolor{PineGreen}{$\bullet$}PineGreen     \>{0.92,0,0.59,0.25}\>\textcolor{LimeGreen}{$\bullet$}LimeGreen     \>{0.50,0,1,0}\\
%     \textcolor{YellowGreen}{$\bullet$}YellowGreen   \>{0.44,0,0.74,0}\>\textcolor{SpringGreen}{$\bullet$}SpringGreen   \>{0.26,0,0.76,0}\\
%     \textcolor{OliveGreen}{$\bullet$}OliveGreen    \>{0.64,0,0.95,0.40}\>\textcolor{RawSienna}{$\bullet$}RawSienna     \>{0,0.72,1,0.45}\\
%     \textcolor{Sepia}{$\bullet$}Sepia         \>{0,0.83,1,0.70}\>\textcolor{Brown}{$\bullet$}Brown         \>{0,0.81,1,0.60}\\
%     \textcolor{Tan}{$\bullet$}Tan           \>{0.14,0.42,0.56,0}\>\>\\
%     \textcolor{White}{$\bullet$}White (Gray0)  \>{0,0,0,0}\>\textcolor{Black}{$\bullet$}Black (Gray100)    \>{0,0,0,1}\\
%     \textcolor{Gray5}{$\bullet$}Gray5          \>{0,0,0,0.05}\>\textcolor{Gray10}{$\bullet$}Gray10          \>{0,0,0,0.10}\\
%     \textcolor{Gray15}{$\bullet$}Gray15          \>{0,0,0,0.15}\>\textcolor{Gray20}{$\bullet$}Gray20          \>{0,0,0,0.20}\\
%     \textcolor{Gray25}{$\bullet$}Gray25          \>{0,0,0,0.25}\>\textcolor{Gray30}{$\bullet$}Gray30          \>{0,0,0,0.30}\\
%     \textcolor{LightGray}{$\bullet$}LightGray       \>{0,0,0,0.33}\>\textcolor{Gray35}{$\bullet$}Gray35          \>{0,0,0,0.35}\\
%     \textcolor{Gray40}{$\bullet$}Gray40          \>{0,0,0,0.40}\>\textcolor{Gray45}{$\bullet$}Gray45          \>{0,0,0,0.45}\\
%     \textcolor{Gray50}{$\bullet$}Gray50          \>{0,0,0,0.50}\>\textcolor{Gray}{$\bullet$}Gray          \>{0,0,0,0.50}\\
%     \textcolor{Gray55}{$\bullet$}Gray55          \>{0,0,0,0.55}\>\textcolor{Gray60}{$\bullet$}Gray60          \>{0,0,0,0.60}\\
%     \textcolor{Gray65}{$\bullet$}Gray65          \>{0,0,0,0.65}\>\textcolor{DarkGray}{$\bullet$}DarkGray        \>{0,0,0,0.66}\\
%     \textcolor{Gray70}{$\bullet$}Gray70          \>{0,0,0,0.70}\>\textcolor{Gray75}{$\bullet$}Gray75          \>{0,0,0,0.75}\\
%     \textcolor{Gray80}{$\bullet$}Gray80          \>{0,0,0,0.80}\>\textcolor{Gray85}{$\bullet$}Gray85          \>{0,0,0,0.85}\\
%     \textcolor{Gray90}{$\bullet$}Gray90          \>{0,0,0,0.90}\>\textcolor{Gray95}{$\bullet$}Gray95          \>{0,0,0,0.95}\\
%     \textcolor{LightGreenYellow}{$\bullet$}LightGreenYellow\>{0.08,0,0.35,0}\>\textcolor{LightYellow}{$\bullet$}LightYellow \>{0,0,0.50,0}\\
%     \textcolor{LightGoldenrod}{$\bullet$}LightGoldenrod  \>{0,0.05,0.42,0}\>\textcolor{LightDandelion}{$\bullet$}LightDandelion\> {0,0.15,0.42,0}\\
%     \textcolor{LightApricot}{$\bullet$}LightApricot    \>{0,0.16,0.26,0}\>\textcolor{LightPeach}{$\bullet$}LightPeach   \>{0,0.25,0.35,0}\\
%     \textcolor{LightMelon}{$\bullet$}LightMelon      \>{0,0.23,0.25,0}\>\textcolor{LightYellowOrange}{$\bullet$}LightYellowOrange \>{0,0.21,0.50,0}\\
%     \textcolor{LightOrange}{$\bullet$}LightOrange     \>{0,0.31,0.44,0}\>\textcolor{LightBurntOrange}{$\bullet$}LightBurntOrange \>{0,0.26,0.50,0}\\
%     \textcolor{LightBittersweet}{$\bullet$}LightBittersweet\>{0,0.38,0.50,0.12}\>\textcolor{LightRedOrange}{$\bullet$}LightRedOrange\>{0,0.39,0.44,0}\\
%     \textcolor{LightMahagony}{$\bullet$}LightMahagony   \>{0,0.43,0.44,0.18}\>\textcolor{LightMaroon}{$\bullet$}LightMaroon \>{0,0.44,0.34,0.16}\\
%     \textcolor{LightBrickRed}{$\bullet$}LightBrickRed   \>{0,0.45,0.47,0.14}\>\textcolor{LightRed}{$\bullet$}LightRed    \>{0,0.50,0.50,0}\\
%     \textcolor{LightOrangeRed}{$\bullet$}LightOrangeRed  \>{0,0.50,0.25,0}\>\textcolor{LightRubineRed}{$\bullet$}LightRubineRed  \>{0,0.50,0.07,0}\\
%     \textcolor{LightWildStrawberry}{$\bullet$}LightWildStrawberry\>{0,0.48,0.20,0}\>\textcolor{LightSalmon}{$\bullet$}LightSalmon  \>{0,0.27,0.19,0}\\
%     \textcolor{LightCarnationPink}{$\bullet$}LightCarnationPink \>{0,0.32,0,0} \>\textcolor{LightMagenta}{$\bullet$}LightMagenta   \>{0,0.50,0,0}\\
%     \textcolor{LightVioletRed}{$\bullet$}LightVioletRed  \>{0,0.40,0,0} \>\textcolor{LightRhodamine}{$\bullet$}LightRhodamine   \>{0,0.41,0,0}\\
%     \textcolor{LightMulberry}{$\bullet$}LightMulberry   \>{0.17,0.45,0,0.01}\>\textcolor{LightRedViolet}{$\bullet$}LightRedViolet \>{0.04,0.45,0,0.17}\\
%     \textcolor{LightFuchsia}{$\bullet$}LightFuchsia    \>{0.24,0.46,0,0.04}\>\textcolor{LightLavender}{$\bullet$}LightLavender \> {0,0.24,0,0}\\
%     \textcolor{LightThistle}{$\bullet$}LightThistle    \>{0.06,0.30,0,0} \>\textcolor{LightOrchid}{$\bullet$}LightOrchid    \>{0.16,0.32,0,0}\\
%     \textcolor{LightDarkOrchid}{$\bullet$}LightDarkOrchid \>{0.20,0.40,0.10,0}\>\textcolor{LightPurple}{$\bullet$}LightPurple  \>{0.23,0.43,0,0}\\
%     \textcolor{LightPlum}{$\bullet$}LightPlum       \>{0.25,0.50,0,0} \>\textcolor{LightViolet}{$\bullet$}LightViolet \>{0.40,0.44,0,0}\\
%     \textcolor{LightRoyalPurple}{$\bullet$}LightRoyalPurple\>{0.38,0.45,0,0} \>\textcolor{LightBlueViolet}{$\bullet$}LightBlueViolet \>{0.43,0.46,0,0.02}\\
%     \textcolor{LightPeriwinkle}{$\bullet$}LightPeriwinkle \>{0.29,0.28,0,0} \>\textcolor{LightCadetBlue}{$\bullet$}LightCadetBlue \> {0.31,0.29,0.12,0}\\
%     \textcolor{LightCornflowerBlue}{$\bullet$}LightCornflowerBlue\>{0.33,0.07,0,0}\>\textcolor{LightMidnightBlue}{$\bullet$}LightMidnightBlue\>{0.49,0.07,0,0.22}\\
%     \textcolor{LightNavyBlue}{$\bullet$}LightNavyBlue   \>{0.47,0.27,0,0} \>\textcolor{LightRoyalBlue}{$\bullet$}LightRoyalBlue \> {0.50,0.25,0,0}\\
%     \textcolor{LightBlue}{$\bullet$}LightBlue       \>{0.50,0.50,0,0} \>\textcolor{LightCerulean}{$\bullet$}LightCerulean \> {0.47,0.06,0,0}\\
%     \textcolor{LightCyan}{$\bullet$}LightCyan       \>{0.50,0,0,0} \>\textcolor{LightProcessBlue}{$\bullet$}LightProcessBlue \> {0.48,0,0,0}\\
%     \textcolor{LightSkyBlue}{$\bullet$}LightSkyBlue    \>{0.31,0,0.06,0} \>\textcolor{LightTurquoise}{$\bullet$}LightTurquoise \>{0.43,0,0.10,0}\\
%     \textcolor{LightTealBlue}{$\bullet$}LightTealBlue   \>{0.43,0,0.17,0.01}\>\textcolor{LightAquamarine}{$\bullet$}LightAquamarine \>{0.41,0,0.15,0}\\
%     \textcolor{LightBlueGreen}{$\bullet$}LightBlueGreen  \>{0.43,0,0.17,0}\>\textcolor{LightEmerald}{$\bullet$}LightEmerald   \>{0.50,0,0.25,0}\\
%     \textcolor{LightJungleGreen}{$\bullet$}LightJungleGreen\>{0.50,0,0.26,0} \>\textcolor{LightSeaGreen}{$\bullet$}LightSeaGreen  \>{0.35,0,0.25,0}\\
%     \textcolor{LightGreen}{$\bullet$}LightGreen      \>{0.50,0,0.50,0} \>\textcolor{LightForestGreen}{$\bullet$}LightForestGreen\>{0.46,0,0.44,0.06}\\
%     \textcolor{LightPineGreen}{$\bullet$}LightPineGreen  \>{0.46,0,0.30,0.13}\>\textcolor{LightLimeGreen}{$\bullet$}LightLimeGreen\>{0.25,0,0.50,0}\\
%     \textcolor{LightYellowGreen}{$\bullet$}LightYellowGreen\>{0.22,0,0.37,0} \>\textcolor{LightSpringGreen}{$\bullet$}LightSpringGreen \>{0.13,0,0.38,0}\\
%     \textcolor{LightOliveGreen}{$\bullet$}LightOliveGreen \>{0.32,0,0.48,0.20} \>\textcolor{LightRawSienna}{$\bullet$}LightRawSienna\>{0,0.36,0.50,0.23}\\
%     \textcolor{LightSepia}{$\bullet$}LightSepia      \>{0,0.44,0.50,0.35} \>\textcolor{LightBrown}{$\bullet$}LightBrown    \>{0,0.41,0.50,0.30}\\
%     \textcolor{LightTan}{$\bullet$}LightTan        \>{0.07,0.21,0.28,0}\\ 
%     LightLight- and LightLightLight-versions were derived by dividing all values\\ 
%     from Light-color definitions by 2 and 4, respectively. 
%    \end{tabbing}
%
%    \begin{tabbing}
%    \emph{name}\hspace{2.5cm}\= \emph{RGB\quad}\hspace{1.8cm}
%    \=\emph{name}\hspace{2.5cm}\= \emph{RGB\quad}\\
%    
%     \textcolor{BlueRed5}{$\bullet$}BlueRed5        \>{0.15,0.17,0.55} \>\textcolor{BlueRed10}{$\bullet$}BlueRed10 \> {0.20,0.23,0.57}\\
%     \textcolor{BlueRed15}{$\bullet$}BlueRed15       \> {0.24,0.29,0.60} \>\textcolor{BlueRed20}{$\bullet$}BlueRed20 \> {0.33,0.35,0.64}\\
%     \textcolor{BlueRed25}{$\bullet$}BlueRed25       \> {0.43,0.43,0.68} \>\textcolor{BlueRed30}{$\bullet$}BlueRed30 \> {0.52,0.52,0.73}\\
%     \textcolor{BlueRed35}{$\bullet$}BlueRed35       \> {0.60,0.60,0.78} \>\textcolor{BlueRed40}{$\bullet$}BlueRed40 \> {0.70,0.70,0.84}\\
%     \textcolor{BlueRed45}{$\bullet$}BlueRed45       \> {0.80,0.80,0.85} \>\textcolor{BlueRed50}{$\bullet$}BlueRed50 \> {0.86,0.82,0.82}\\
%     \textcolor{BlueRed55}{$\bullet$}BlueRed55       \> {0.87,0.73,0.73} \>\textcolor{BlueRed60}{$\bullet$}BlueRed60 \> {0.89,0.64,0.64}\\
%     \textcolor{BlueRed65}{$\bullet$}BlueRed65       \> {0.90,0.55,0.55} \>\textcolor{BlueRed70}{$\bullet$}BlueRed70 \> {0.91,0.47,0.46}\\
%     \textcolor{BlueRed75}{$\bullet$}BlueRed75       \> {0.91,0.39,0.37} \>\textcolor{BlueRed80}{$\bullet$}BlueRed80 \> {0.90,0.33,0.28}\\
%     \textcolor{BlueRed85}{$\bullet$}BlueRed85       \> {0.89,0.25,0.20} \>\textcolor{BlueRed90}{$\bullet$}BlueRed90 \> {0.88,0.23,0.14}\\
%     \textcolor{BlueRed95}{$\bullet$}BlueRed95       \> {0.87,0.21,0.09} \>\textcolor{BlueRed100}{$\bullet$}BlueRed100\> {0.87,0.16,0.04}\\
%     \textcolor{GreenRed5}{$\bullet$}GreenRed5       \> {0,1,0}          \>\textcolor{GreenRed10}{$\bullet$}GreenRed10\> {0.05,0.95,0}\\
%     \textcolor{GreenRed15}{$\bullet$}GreenRed15      \> {0.10,0.90,0}    \>\textcolor{GreenRed20}{$\bullet$}GreenRed20\> {0.15,0.85,0}\\
%     \textcolor{GreenRed25}{$\bullet$}GreenRed25      \> {0.20,0.80,0}    \>\textcolor{GreenRed30}{$\bullet$}GreenRed30\> {0.25,0.75,0}\\
%     \textcolor{GreenRed35}{$\bullet$}GreenRed35      \> {0.30,0.70,0}    \>\textcolor{GreenRed40}{$\bullet$}GreenRed40\> {0.35,0.65,0}\\
%     \textcolor{GreenRed45}{$\bullet$}GreenRed45      \> {0.40,0.60,0}    \>\textcolor{GreenRed50}{$\bullet$}GreenRed50\> {0.45,0.55,0}\\
%     \textcolor{GreenRed55}{$\bullet$}GreenRed55      \> {0.50,0.50,0}    \>\textcolor{GreenRed60}{$\bullet$}GreenRed60\> {0.55,0.45,0}\\
%     \textcolor{GreenRed65}{$\bullet$}GreenRed65      \> {0.60,0.40,0}    \>\textcolor{GreenRed70}{$\bullet$}GreenRed70\> {0.65,0.35,0}\\
%     \textcolor{GreenRed75}{$\bullet$}GreenRed75      \> {0.70,0.30,0}    \>\textcolor{GreenRed80}{$\bullet$}GreenRed80\> {0.75,0.25,0}\\
%     \textcolor{GreenRed85}{$\bullet$}GreenRed85      \> {0.80,0.20,0}    \>\textcolor{GreenRed90}{$\bullet$}GreenRed90\> {0.85,0.15,0}\\
%     \textcolor{GreenRed95}{$\bullet$}GreenRed95      \> {0.90,0.10,0}    \>\textcolor{GreenRed100}{$\bullet$}GreenRed100\> {0.95,0.05,0}\\
%     \textcolor{ColdHot5}{$\bullet$}ColdHot5        \> {0,0.08,1}        \>\textcolor{ColdHot10}{$\bullet$}ColdHot10 \> {0,0.29,1}\\
%     \textcolor{ColdHot15}{$\bullet$}ColdHot15       \> {0,0.49,1}       \>\textcolor{ColdHot20}{$\bullet$}ColdHot20 \> {0,0.70,1}\\
%     \textcolor{ColdHot25}{$\bullet$}ColdHot25       \> {0,0.90,1}       \>\textcolor{ColdHot30}{$\bullet$}ColdHot30 \> {0,1,0.87}\\
%     \textcolor{ColdHot35}{$\bullet$}ColdHot35       \> {0,1,0.68}       \>\textcolor{ColdHot40}{$\bullet$}ColdHot40 \> {0,1,0.46}\\
%     \textcolor{ColdHot45}{$\bullet$}ColdHot45       \> {0,1,0.25}       \>\textcolor{ColdHot50}{$\bullet$}ColdHot50 \> {0,1,0.04}\\
%     \textcolor{ColdHot55}{$\bullet$}ColdHot55       \> {0.16,1,0}       \>\textcolor{ColdHot60}{$\bullet$}ColdHot60 \> {0.35,1,0}\\
%     \textcolor{ColdHot65}{$\bullet$}ColdHot65       \> {0.56,1,0}       \>\textcolor{ColdHot70}{$\bullet$}ColdHot70 \> {0.79,1,0}\\
%     \textcolor{ColdHot75}{$\bullet$}ColdHot75       \> {0.98,1,0}       \>\textcolor{ColdHot80}{$\bullet$}ColdHot80 \> {1,0.82,0}\\
%     \textcolor{ColdHot85}{$\bullet$}ColdHot85       \> {1,0.60,0}       \>\textcolor{ColdHot90}{$\bullet$}ColdHot90 \> {1,0.40,0}\\
%     \textcolor{ColdHot95}{$\bullet$}ColdHot95       \> {1,0.20,0}       \>\textcolor{ColdHot100}{$\bullet$}ColdHot100\> {0.91,0,0}\\
%     and reverse definitions: |RedBlue|, |RedGreen|, |HotCold|.\\
%    \end{tabbing}
%    \end{footnotesize}
%    
%    Type the color names with the upper case letters exactly as described above.
%    For the definition of new colors use one of the |color.sty| commands:
%    \medskip
%    
%    \quad|\definecolor{|\meta{name}|}{cmyk}{|\meta{C,M,Y,K}|}|
%    \medskip
%    
%    \quad|\definecolor{|\meta{name}|}{rgb}{|\meta{R,G,B}|}|
%    \medskip
%    
%    The \meta{name} can be chosen freely, the values for the color
%    composition must be in the range 0--1, i\,e.\ 0--100\% of the 
%    respective component (`C' -- cyan, `M' -- magenta, `Y' -- yellow,
%    `K' -- black; or `R' -- red, `G' -- green, `Blue' -- blue) separated by 
%    commas.
%    \medskip
%    
%    Examples: 
%    \medskip
%    
%    |\definecolor{Salmon}{cmyk}{0,0.53,0.38,0}|
%    \medskip
%
%    |\definecolor{ColdHot15}{rgb}{0,0.49,1}|
%    \medskip
%
%    \newpage
%    \section{Listing of the \texttt{texshade} default settings}
%    
%    \subsection{Standard definitions}
%    
%    The file |texshade.def| mirrors all commands which are
%    carried out at the beginning of the |texshade| environment.
%    Short comments are also included, thus, it is refered to
%    this file for further information. 
%    
%    \subsection{Colors used in the different shading modes}
%    
%    \vspace{5mm}
%    
%    Color scheme \emph{blues}:
%    \medskip
%    
%    \begin{tabbing}
%    \hspace{1cm}\=\emph{res.color}\hspace{1.5cm}\=\emph{shad.color}\hspace{1.5cm} \= \emph{residues} \\
%    \>Black \> White   \> no match \\
%    \>Black \> Magenta \> similar  \\
%    \>White \> RoyalBlue \> identical \\
%    \>Goldenrod \> RoyalPurple \> all match\\
%    \end{tabbing}
%    \medskip
%    
%    Color scheme \emph{greens}:
%    \medskip
%    
%    \begin{tabbing}
%    \hspace{1cm}\=\emph{res.color}\hspace{1.5cm}\=\emph{shad.color}\hspace{1.5cm} \= \emph{residues} \\
%    \>Black \> White   \> no match \\
%    \>Black \> GreenYellow \> similar  \\
%    \>White \> PineGreen \> identical \\
%    \>YellowOrange \> OliveGreen \> all match\\
%    \end{tabbing}
%    \medskip
%    
%    Color scheme \emph{reds}:
%    \medskip
%    
%    \begin{tabbing}
%    \hspace{1cm}\=\emph{res.color}\hspace{1.5cm}\=\emph{shad.color}\hspace{1.5cm} \= \emph{residues} \\
%    \>Black \> White   \> no match \\
%    \>Black \> YellowOrange \> similar  \\
%    \>White \> BrickRed \> identical \\
%    \>YellowGreen \> Mahagony \> all match\\
%    \end{tabbing}
%    \medskip
%    
%    \newpage
%    Color scheme \emph{grays}:
%    \medskip
%    
%    \begin{tabbing}
%    \hspace{1cm}\=\emph{res.color}\hspace{1.5cm}\=\emph{shad.color}\hspace{1.5cm} \= \emph{residues} \\
%    \>Black \> White   \> no match \\
%    \>Black \> LightGray \> similar  \\
%    \>White \> DarkGray \> identical \\
%    \>White \> Black \> all match\\
%    \end{tabbing}
%    \medskip
%    
%    Color scheme \emph{black}:
%    \medskip
%    
%    \begin{tabbing}
%    \hspace{1cm}\=\emph{res.color}\hspace{1.5cm}\=\emph{shad.color}\hspace{1.5cm} \= \emph{residues} \\
%    \>Black \> White   \> no match \\
%    \>Black \> White \> similar  \\
%    \>White \> Black \> identical \\
%    \>White \> Black \> all match\\
%    \end{tabbing}
%    \medskip
%    
%    Functional mode \emph{charge}:
%    \medskip
%    
%    \begin{tabbing}
%    \hspace{1cm}\=\emph{res.color}\hspace{1.5cm}\=\emph{shad.color}\hspace{1.5cm} \= \emph{residues} \\
%    \>Black \> White   \> no match \\
%    \>White \> Red \> acidic  \\
%    \>White \> Blue \> basic \\
%    \end{tabbing}
%    \medskip
%    
%    Functional mode \emph{hydropathy}:
%    \medskip
%    
%    \begin{tabbing}
%    \hspace{1cm}\=\emph{res.color}\hspace{1.5cm}\=\emph{shad.color}\hspace{1.5cm} \= \emph{residues} \\
%    \>Black \> White   \> no match \\
%    \>White \> Red \> acidic  \\
%    \>White \> Blue \> basic \\
%    \>Black \> Yellow   \> polar uncharged \\
%    \>White \> Green   \> hydrophobic nonpolar \\
%    \end{tabbing}
%    \medskip
%    
%    \newpage
%    Functional mode \emph{chemical}:
%    \medskip
%    
%    \begin{tabbing}
%    \hspace{1cm}\=\emph{res.color}\hspace{1.5cm}\=\emph{shad.color}\hspace{1.5cm} \= \emph{residues} \\
%    \>Black \> White   \> no match \\
%    \>White \> Red \> acidic  \\
%    \>White \> Black \> aliphatic  \\
%    \>White \> Gray \> aliphatic (small) \\
%    \>White \> Green \> amide  \\
%    \>White \> Brown \> aromatic  \\
%    \>White \> Blue \> basic \\
%    \>Black \> Magenta   \> hydroxyl \\
%    \>Black \> Orange   \> imino \\
%    \>Black \> Yellow   \> sulfur \\
%    \end{tabbing}
%    \medskip
%    
%    Functional mode \emph{rasmol}:
%    \medskip
%    
%    \begin{tabbing}
%    \hspace{1cm}\=\emph{res.color}\hspace{1.5cm}\=\emph{shad.color}\hspace{1.5cm} \= \emph{residues} \\
%    \>Black \> White   \> no match \\
%    \>Red \> White   \> Asp, Glu \\
%    \>Blue \> White \> Arg, Lys, His  \\
%    \>MidnightBlue \> White \> Phe, Tyr, Trp  \\
%    \>Gray \> White \> Ala, Gly  \\
%    \>Yellow \> White \> Cys, Met  \\
%    \>Orange \> White \> Ser, Thr \\
%    \>Cyan \> White   \> Asn, Gln \\
%    \>Gree \> White   \> Leu, Val, Ile \\
%    \>Apricot \> White   \> Pro \\
%    \end{tabbing}
%    \medskip
%    
%    Functional mode \emph{structure}:
%    \medskip
%    
%    \begin{tabbing}
%    \hspace{1cm}\=\emph{res.color}\hspace{1.5cm}\=\emph{shad.color}\hspace{1.5cm} \= \emph{residues} \\
%    \>Black \> White   \> no match \\
%    \>Black \> Orange \> external  \\
%    \>Black \> Yellow   \> ambivalent \\
%    \>White \> Green \> internal  \\
%    \end{tabbing}
%    \medskip
%    
%    \newpage
%    Functional mode \emph{standard area}:
%    \medskip
%    
%    \begin{tabbing}
%    \hspace{1cm}\=\emph{res.color}\hspace{1.5cm}\=\emph{shad.color}\hspace{1.5cm} \= \emph{residues} \\
%    \>Black \> White   \> no match \\
%    \>Black \> BrickRed \> G\\
%    \>Black \> Orange  \> A, S\\
%    \>Black \> Yellow   \> C, P  \\
%    \>Black \> YellowGreen   \> T, D, V, N \\
%    \>White \> PineGreen   \> I, E \\
%    \>Black \> SkyBlue   \> L, Q, H, M \\
%    \>White \> RoyalPurple \> F, K \\
%    \>White \> RedViolet   \> Y \\
%    \>White \> Black   \> R, W \\
%    \end{tabbing}
%    \medskip
%    
%    Functional mode \emph{accessible area}:
%    \medskip
%    
%    \begin{tabbing}
%    \hspace{1cm}\=\emph{res.color}\hspace{1.5cm}\=\emph{shad.color}\hspace{1.5cm} \= \emph{residues} \\
%    \>Black \> White   \> no match \\
%    \>Black \> BrickRed \> C \\
%    \>Black \> Orange  \> I, V, G \\
%    \>Black \> Yellow   \> F, L, M, A \\
%    \>Black \> YellowGreen   \> W, S, T, H \\
%    \>White \> PineGreen   \> P \\
%    \>Black \> SkyBlue   \> Y, D, N \\
%    \>White \> RoyalPurple \> E, Q  \\
%    \>White \> RedViolet   \> R \\
%    \>White \> Black   \> K \\
%    \end{tabbing}
%    \medskip
%    
%    
%    \newpage
%    \subsection{Residue weight tables}\label{weightmatrix}
%    \bigskip
%
%    |identity|
%    \bigskip
%    \bigskip
%    
%{\footnotesize\tt
%
%\weighttable{identity}
%\hspace*{-2cm}
%\begin{tabular}{lrrrrrrrrrrrrrrrrrrrr}
%
%C & \consCC &&&&&&&&&&&&&&&&&&& \\
%
%S & \consSC & \consSS &&&&&&&&&&&&&&&&&& \\
%
%T & \consTC & \consTS & \consTT &&&&&&&&&&&&&&&&& \\
%
%P & \consPC & \consPS & \consPT & \consPP &&&&&&&&&&&&&&&& \\
%
%A & \consAC & \consAS & \consAT & \consAP & \consAA
%&&&&&&&&&&&&&&& \\
%
%G & \consGC & \consGS & \consGT & \consGP & \consGA
%& \consGG &&&&&&&&&&&&&& \\
%
%N & \consNC & \consNS & \consNT & \consNP & \consNA
%& \consNG & \consNN &&&&&&&&&&&&& \\
%
%D & \consDC & \consDS & \consDT & \consDP & \consDA
%& \consDG & \consDN & \consDD &&&&&&&&&&&& \\
%
%E & \consEC & \consES & \consET & \consEP & \consEA
%& \consEG & \consEN & \consED & \consEE &&&&&&&&&&& \\
%
%Q & \consQC & \consQS & \consQT & \consQP & \consQA
%& \consQG & \consQN & \consQD & \consQE & \consQQ
%&&&&&&&&&& \\
%
%H & \consHC & \consHS & \consHT & \consHP & \consHA
%& \consHG & \consHN & \consHD & \consHE & \consHQ
%& \consHH &&&&&&&&& \\
%
%R & \consRC & \consRS & \consRT & \consRP & \consRA
%& \consRG & \consRN & \consRD & \consRE & \consRQ
%& \consRH & \consRR &&&&&&&& \\
%
%K & \consKC & \consKS & \consKT & \consKP & \consKA
%& \consKG & \consKN & \consKD & \consKE & \consKQ
%& \consKH & \consKR & \consKK &&&&&&& \\
%
%M & \consMC & \consMS & \consMT & \consMP & \consMA
%& \consMG & \consMN & \consMD & \consME & \consMQ
%& \consMH & \consMR & \consMK & \consMM &&&&&& \\
%
%I & \consIC & \consIS & \consIT & \consIP & \consIA
%& \consIG & \consIN & \consID & \consIE & \consIQ
%& \consIH & \consIR & \consIK & \consIM & \consII
%&&&&& \\
%
%L & \consLC & \consLS & \consLT & \consLP & \consLA
%& \consLG & \consLN & \consLD & \consLE & \consLQ
%& \consLH & \consLR & \consLK & \consLM & \consLI
%& \consLL &&&& \\
%
%V & \consVC & \consVS & \consVT & \consVP & \consVA
%& \consVG & \consVN & \consVD & \consVE & \consVQ
%& \consVH & \consVR & \consVK & \consVM & \consVI
%& \consVL & \consVV &&& \\
%
%F & \consFC & \consFS & \consFT & \consFP & \consFA
%& \consFG & \consFN & \consFD & \consFE & \consFQ
%& \consFH & \consFR & \consFK & \consFM & \consFI
%& \consFL & \consFV & \consFF && \\
%
%Y & \consYC & \consYS & \consYT & \consYP & \consYA
%& \consYG & \consYN & \consYD & \consYE & \consYQ
%& \consYH & \consYR & \consYK & \consYM & \consYI
%& \consYL & \consYV & \consYF & \consYY & \\
%
%W & \consWC & \consWS & \consWT & \consWP & \consWA
%& \consWG & \consWN & \consWD & \consWE & \consWQ
%& \consWH & \consWR & \consWK & \consWM & \consWI
%& \consWL & \consWV & \consWF & \consWY & \consWW \\[1.5ex]
%
%& C & S & T & P & A & G & N & D & E & Q & H & R & K & M & I & L & V & F & Y & W \\
%
%\end{tabular}
%}
%    \newpage
%
%    |structural|
%    \bigskip
%    \bigskip
%    
%{\footnotesize\tt
%
%\weighttable{structural}
%\hspace*{-2cm}
%\begin{tabular}{lrrrrrrrrrrrrrrrrrrrr}
%
%C & \consCC &&&&&&&&&&&&&&&&&&& \\
%
%S & \consSC & \consSS &&&&&&&&&&&&&&&&&& \\
%
%T & \consTC & \consTS & \consTT &&&&&&&&&&&&&&&&& \\
%
%P & \consPC & \consPS & \consPT & \consPP &&&&&&&&&&&&&&&& \\
%
%A & \consAC & \consAS & \consAT & \consAP & \consAA
%&&&&&&&&&&&&&&& \\
%
%G & \consGC & \consGS & \consGT & \consGP & \consGA
%& \consGG &&&&&&&&&&&&&& \\
%
%N & \consNC & \consNS & \consNT & \consNP & \consNA
%& \consNG & \consNN &&&&&&&&&&&&& \\
%
%D & \consDC & \consDS & \consDT & \consDP & \consDA
%& \consDG & \consDN & \consDD &&&&&&&&&&&& \\
%
%E & \consEC & \consES & \consET & \consEP & \consEA
%& \consEG & \consEN & \consED & \consEE &&&&&&&&&&& \\
%
%Q & \consQC & \consQS & \consQT & \consQP & \consQA
%& \consQG & \consQN & \consQD & \consQE & \consQQ
%&&&&&&&&&& \\
%
%H & \consHC & \consHS & \consHT & \consHP & \consHA
%& \consHG & \consHN & \consHD & \consHE & \consHQ
%& \consHH &&&&&&&&& \\
%
%R & \consRC & \consRS & \consRT & \consRP & \consRA
%& \consRG & \consRN & \consRD & \consRE & \consRQ
%& \consRH & \consRR &&&&&&&& \\
%
%K & \consKC & \consKS & \consKT & \consKP & \consKA
%& \consKG & \consKN & \consKD & \consKE & \consKQ
%& \consKH & \consKR & \consKK &&&&&&& \\
%
%M & \consMC & \consMS & \consMT & \consMP & \consMA
%& \consMG & \consMN & \consMD & \consME & \consMQ
%& \consMH & \consMR & \consMK & \consMM &&&&&& \\
%
%I & \consIC & \consIS & \consIT & \consIP & \consIA
%& \consIG & \consIN & \consID & \consIE & \consIQ
%& \consIH & \consIR & \consIK & \consIM & \consII
%&&&&& \\
%
%L & \consLC & \consLS & \consLT & \consLP & \consLA
%& \consLG & \consLN & \consLD & \consLE & \consLQ
%& \consLH & \consLR & \consLK & \consLM & \consLI
%& \consLL &&&& \\
%
%V & \consVC & \consVS & \consVT & \consVP & \consVA
%& \consVG & \consVN & \consVD & \consVE & \consVQ
%& \consVH & \consVR & \consVK & \consVM & \consVI
%& \consVL & \consVV &&& \\
%
%F & \consFC & \consFS & \consFT & \consFP & \consFA
%& \consFG & \consFN & \consFD & \consFE & \consFQ
%& \consFH & \consFR & \consFK & \consFM & \consFI
%& \consFL & \consFV & \consFF && \\
%
%Y & \consYC & \consYS & \consYT & \consYP & \consYA
%& \consYG & \consYN & \consYD & \consYE & \consYQ
%& \consYH & \consYR & \consYK & \consYM & \consYI
%& \consYL & \consYV & \consYF & \consYY & \\
%
%W & \consWC & \consWS & \consWT & \consWP & \consWA
%& \consWG & \consWN & \consWD & \consWE & \consWQ
%& \consWH & \consWR & \consWK & \consWM & \consWI
%& \consWL & \consWV & \consWF & \consWY & \consWW \\[1.5ex]
%
%& C & S & T & P & A & G & N & D & E & Q & H & R & K & M & I & L & V & F & Y & W \\
%
%\end{tabular}
%}
%    \newpage
%
%    |PAM250|
%    \bigskip
%    \bigskip
%
%{\footnotesize\tt
%
%\weighttable{PAM250}
%\hspace*{-2cm}
%\begin{tabular}{lrrrrrrrrrrrrrrrrrrrr}
%
%C & \consCC &&&&&&&&&&&&&&&&&&& \\
%
%S & \consSC & \consSS &&&&&&&&&&&&&&&&&& \\
%
%T & \consTC & \consTS & \consTT &&&&&&&&&&&&&&&&& \\
%
%P & \consPC & \consPS & \consPT & \consPP &&&&&&&&&&&&&&&& \\
%
%A & \consAC & \consAS & \consAT & \consAP & \consAA
%&&&&&&&&&&&&&&& \\
%
%G & \consGC & \consGS & \consGT & \consGP & \consGA
%& \consGG &&&&&&&&&&&&&& \\
%
%N & \consNC & \consNS & \consNT & \consNP & \consNA
%& \consNG & \consNN &&&&&&&&&&&&& \\
%
%D & \consDC & \consDS & \consDT & \consDP & \consDA
%& \consDG & \consDN & \consDD &&&&&&&&&&&& \\
%
%E & \consEC & \consES & \consET & \consEP & \consEA
%& \consEG & \consEN & \consED & \consEE &&&&&&&&&&& \\
%
%Q & \consQC & \consQS & \consQT & \consQP & \consQA
%& \consQG & \consQN & \consQD & \consQE & \consQQ
%&&&&&&&&&& \\
%
%H & \consHC & \consHS & \consHT & \consHP & \consHA
%& \consHG & \consHN & \consHD & \consHE & \consHQ
%& \consHH &&&&&&&&& \\
%
%R & \consRC & \consRS & \consRT & \consRP & \consRA
%& \consRG & \consRN & \consRD & \consRE & \consRQ
%& \consRH & \consRR &&&&&&&& \\
%
%K & \consKC & \consKS & \consKT & \consKP & \consKA
%& \consKG & \consKN & \consKD & \consKE & \consKQ
%& \consKH & \consKR & \consKK &&&&&&& \\
%
%M & \consMC & \consMS & \consMT & \consMP & \consMA
%& \consMG & \consMN & \consMD & \consME & \consMQ
%& \consMH & \consMR & \consMK & \consMM &&&&&& \\
%
%I & \consIC & \consIS & \consIT & \consIP & \consIA
%& \consIG & \consIN & \consID & \consIE & \consIQ
%& \consIH & \consIR & \consIK & \consIM & \consII
%&&&&& \\
%
%L & \consLC & \consLS & \consLT & \consLP & \consLA
%& \consLG & \consLN & \consLD & \consLE & \consLQ
%& \consLH & \consLR & \consLK & \consLM & \consLI
%& \consLL &&&& \\
%
%V & \consVC & \consVS & \consVT & \consVP & \consVA
%& \consVG & \consVN & \consVD & \consVE & \consVQ
%& \consVH & \consVR & \consVK & \consVM & \consVI
%& \consVL & \consVV &&& \\
%
%F & \consFC & \consFS & \consFT & \consFP & \consFA
%& \consFG & \consFN & \consFD & \consFE & \consFQ
%& \consFH & \consFR & \consFK & \consFM & \consFI
%& \consFL & \consFV & \consFF && \\
%
%Y & \consYC & \consYS & \consYT & \consYP & \consYA
%& \consYG & \consYN & \consYD & \consYE & \consYQ
%& \consYH & \consYR & \consYK & \consYM & \consYI
%& \consYL & \consYV & \consYF & \consYY & \\
%
%W & \consWC & \consWS & \consWT & \consWP & \consWA
%& \consWG & \consWN & \consWD & \consWE & \consWQ
%& \consWH & \consWR & \consWK & \consWM & \consWI
%& \consWL & \consWV & \consWF & \consWY & \consWW \\[1.5ex]
%
%& C & S & T & P & A & G & N & D & E & Q & H & R & K & M & I & L & V & F & Y & W \\
%
%\end{tabular}
%}
%    \newpage
%
%    |PAM100|
%    \bigskip
%    \bigskip
%    
%{\footnotesize\tt
%
%\weighttable{PAM100}
%\hspace*{-3cm}
%\begin{tabular}{lrrrrrrrrrrrrrrrrrrrr}
%
%C & \consCC &&&&&&&&&&&&&&&&&&& \\
%
%S & \consSC & \consSS &&&&&&&&&&&&&&&&&& \\
%
%T & \consTC & \consTS & \consTT &&&&&&&&&&&&&&&&& \\
%
%P & \consPC & \consPS & \consPT & \consPP &&&&&&&&&&&&&&&& \\
%
%A & \consAC & \consAS & \consAT & \consAP & \consAA
%&&&&&&&&&&&&&&& \\
%
%G & \consGC & \consGS & \consGT & \consGP & \consGA
%& \consGG &&&&&&&&&&&&&& \\
%
%N & \consNC & \consNS & \consNT & \consNP & \consNA
%& \consNG & \consNN &&&&&&&&&&&&& \\
%
%D & \consDC & \consDS & \consDT & \consDP & \consDA
%& \consDG & \consDN & \consDD &&&&&&&&&&&& \\
%
%E & \consEC & \consES & \consET & \consEP & \consEA
%& \consEG & \consEN & \consED & \consEE &&&&&&&&&&& \\
%
%Q & \consQC & \consQS & \consQT & \consQP & \consQA
%& \consQG & \consQN & \consQD & \consQE & \consQQ
%&&&&&&&&&& \\
%
%H & \consHC & \consHS & \consHT & \consHP & \consHA
%& \consHG & \consHN & \consHD & \consHE & \consHQ
%& \consHH &&&&&&&&& \\
%
%R & \consRC & \consRS & \consRT & \consRP & \consRA
%& \consRG & \consRN & \consRD & \consRE & \consRQ
%& \consRH & \consRR &&&&&&&& \\
%
%K & \consKC & \consKS & \consKT & \consKP & \consKA
%& \consKG & \consKN & \consKD & \consKE & \consKQ
%& \consKH & \consKR & \consKK &&&&&&& \\
%
%M & \consMC & \consMS & \consMT & \consMP & \consMA
%& \consMG & \consMN & \consMD & \consME & \consMQ
%& \consMH & \consMR & \consMK & \consMM &&&&&& \\
%
%I & \consIC & \consIS & \consIT & \consIP & \consIA
%& \consIG & \consIN & \consID & \consIE & \consIQ
%& \consIH & \consIR & \consIK & \consIM & \consII
%&&&&& \\
%
%L & \consLC & \consLS & \consLT & \consLP & \consLA
%& \consLG & \consLN & \consLD & \consLE & \consLQ
%& \consLH & \consLR & \consLK & \consLM & \consLI
%& \consLL &&&& \\
%
%V & \consVC & \consVS & \consVT & \consVP & \consVA
%& \consVG & \consVN & \consVD & \consVE & \consVQ
%& \consVH & \consVR & \consVK & \consVM & \consVI
%& \consVL & \consVV &&& \\
%
%F & \consFC & \consFS & \consFT & \consFP & \consFA
%& \consFG & \consFN & \consFD & \consFE & \consFQ
%& \consFH & \consFR & \consFK & \consFM & \consFI
%& \consFL & \consFV & \consFF && \\
%
%Y & \consYC & \consYS & \consYT & \consYP & \consYA
%& \consYG & \consYN & \consYD & \consYE & \consYQ
%& \consYH & \consYR & \consYK & \consYM & \consYI
%& \consYL & \consYV & \consYF & \consYY & \\
%
%W & \consWC & \consWS & \consWT & \consWP & \consWA
%& \consWG & \consWN & \consWD & \consWE & \consWQ
%& \consWH & \consWR & \consWK & \consWM & \consWI
%& \consWL & \consWV & \consWF & \consWY & \consWW \\[1.5ex]
%
%& C & S & T & P & A & G & N & D & E & Q & H & R & K & M & I & L & V & F & Y & W \\
%
%\end{tabular}
%}
%    \newpage
%
%    |BLOSUM62|
%\bigskip
%\bigskip
%    
%{\footnotesize\tt
%
%\weighttable{BLOSUM62}
%\hspace*{-2cm}
%\begin{tabular}{lrrrrrrrrrrrrrrrrrrrr}
%
%C & \consCC &&&&&&&&&&&&&&&&&&& \\
%
%S & \consSC & \consSS &&&&&&&&&&&&&&&&&& \\
%
%T & \consTC & \consTS & \consTT &&&&&&&&&&&&&&&&& \\
%
%P & \consPC & \consPS & \consPT & \consPP &&&&&&&&&&&&&&&& \\
%
%A & \consAC & \consAS & \consAT & \consAP & \consAA
%&&&&&&&&&&&&&&& \\
%
%G & \consGC & \consGS & \consGT & \consGP & \consGA
%& \consGG &&&&&&&&&&&&&& \\
%
%N & \consNC & \consNS & \consNT & \consNP & \consNA
%& \consNG & \consNN &&&&&&&&&&&&& \\
%
%D & \consDC & \consDS & \consDT & \consDP & \consDA
%& \consDG & \consDN & \consDD &&&&&&&&&&&& \\
%
%E & \consEC & \consES & \consET & \consEP & \consEA
%& \consEG & \consEN & \consED & \consEE &&&&&&&&&&& \\
%
%Q & \consQC & \consQS & \consQT & \consQP & \consQA
%& \consQG & \consQN & \consQD & \consQE & \consQQ
%&&&&&&&&&& \\
%
%H & \consHC & \consHS & \consHT & \consHP & \consHA
%& \consHG & \consHN & \consHD & \consHE & \consHQ
%& \consHH &&&&&&&&& \\
%
%R & \consRC & \consRS & \consRT & \consRP & \consRA
%& \consRG & \consRN & \consRD & \consRE & \consRQ
%& \consRH & \consRR &&&&&&&& \\
%
%K & \consKC & \consKS & \consKT & \consKP & \consKA
%& \consKG & \consKN & \consKD & \consKE & \consKQ
%& \consKH & \consKR & \consKK &&&&&&& \\
%
%M & \consMC & \consMS & \consMT & \consMP & \consMA
%& \consMG & \consMN & \consMD & \consME & \consMQ
%& \consMH & \consMR & \consMK & \consMM &&&&&& \\
%
%I & \consIC & \consIS & \consIT & \consIP & \consIA
%& \consIG & \consIN & \consID & \consIE & \consIQ
%& \consIH & \consIR & \consIK & \consIM & \consII
%&&&&& \\
%
%L & \consLC & \consLS & \consLT & \consLP & \consLA
%& \consLG & \consLN & \consLD & \consLE & \consLQ
%& \consLH & \consLR & \consLK & \consLM & \consLI
%& \consLL &&&& \\
%
%V & \consVC & \consVS & \consVT & \consVP & \consVA
%& \consVG & \consVN & \consVD & \consVE & \consVQ
%& \consVH & \consVR & \consVK & \consVM & \consVI
%& \consVL & \consVV &&& \\
%
%F & \consFC & \consFS & \consFT & \consFP & \consFA
%& \consFG & \consFN & \consFD & \consFE & \consFQ
%& \consFH & \consFR & \consFK & \consFM & \consFI
%& \consFL & \consFV & \consFF && \\
%
%Y & \consYC & \consYS & \consYT & \consYP & \consYA
%& \consYG & \consYN & \consYD & \consYE & \consYQ
%& \consYH & \consYR & \consYK & \consYM & \consYI
%& \consYL & \consYV & \consYF & \consYY & \\
%
%W & \consWC & \consWS & \consWT & \consWP & \consWA
%& \consWG & \consWN & \consWD & \consWE & \consWQ
%& \consWH & \consWR & \consWK & \consWM & \consWI
%& \consWL & \consWV & \consWF & \consWY & \consWW \\[1.5ex]
%
%& C & S & T & P & A & G & N & D & E & Q & H & R & K & M & I & L & V & F & Y & W \\
%
%\end{tabular}
%}
%    \bigskip
%
%
%    \newpage
%    \section{Quick Reference}
%    
%    \textbf{The \TeXshade{} logo}
%    \medskip
%    
%    \quad |\TeXshade|
%    
%    \vspace{1.5\baselineskip}
%    
%    \textbf{The \TeXshade{} environment} (\pageref{tsenvironment}\,ff.)
%    \medskip
%    
%    \begin{quote}
%    |\begin{texshade}[|\meta{parameterfile}|]|
%         |{|\meta{alignmentfile}|}|
%         
%    \quad\emph{further \emph{\TeXshade} commands, if needed}
%    
%    |\end{texshade}|
%    \end{quote}
%    \bigskip
%    
%    \textbf{Predefined shading modes}
%    \medskip
%    
%    \quad|\seqtype{|\meta{type}|}| 
%        \hfill(|P| -- peptide, |N| -- nucleotide) \hfill[\pageref{Lseqtype}]
%        
%    \medskip
%    
%    \quad|\shadingmode[|\meta{option}|]{|\meta{mode}|}|
%        \hfill[\pageref{Lshadingmode}]
%        
%    \medskip
%        
%    \begin{center}    
%    \begin{tabular}{lll}
%    \meta{mode} & \meta{option} &\\ \hline
%    |identical| & |allmatchspecial/number| &\\
%    |similar|   & |allmatchspecial/number| &\\
%    |T-Coffee|  & \meta{filename} &\\
%    |diverse|   & \meta{seqref} &\\
%    |functional|& \meta{type} & |charge| \\
%                &             & |hydropathy| \\
%                &             & |structure| \\
%                &             & |chemical| \\
%                &             & |rasmol| \\
%                &             & |standard area| \\
%                &             & |accessible area| \\ \hline
%    \end{tabular}
%    \end{center}
%    \medskip
%    
%    \quad|\allmatchspecial[|\meta{percentage}|]|
%        \hfill[\pageref{Lallmatchspecial}]
%    
%    \quad|\hideallmatchpositions|
%        \hfill[\pageref{Lhideallmatchpositions}]
%    
%    \quad|\shadeallresidues|
%        \hfill[\pageref{Lshadeallresidues}]
%        
%    \vspace{1.5\baselineskip}
%                
%    \textbf{Shading colors} (\pageref{Lshadingcolors}\,ff.)
%    \medskip
%    
%    \quad|\shadingcolors{|\meta{scheme}|}| \,\, (|blues|, |reds|, 
%                                       |greens|, |grays|, |black|)
%    
%    \quad|\nomatchresidues{|\meta{res.col.}|}{|\meta{shad.col.}|}{|\meta{case}|}{|\meta{style}|}|
%       
%    \quad|\similarresidues{|\meta{res.col.}|}{|\meta{shad.col.}|}{|\meta{case}|}{|\meta{style}|}|
%    
%    \quad|\conservedresidues{|\meta{res.col.}|}{|\meta{shad.col.}|}{|\meta{case}|}{|\meta{style}|}|
%       
%    \quad|\allmatchresidues{|\meta{res.col.}|}{|\meta{shad.col.}|}{|\meta{case}|}{|\meta{style}|}|
%       
%    \quad|\defshadingcolors{|\meta{name}|}|
%    \newpage
%    
%    \quad|\funcshadingstyle{|\meta{residue}|}{|\meta{res.col.}|}{|\meta{shad.color}|}|
%    
%    \hfill|{|\meta{case}|}{|\meta{style}|}| [\pageref{Lfuncshadingstyle}]
%    
%    \vspace{1.5\baselineskip}            
%    
%    \textbf{Residue grouping}
%    \medskip
%    
%    \quad|\pepsims{|\meta{residue}|}{|\meta{similars}|}|
%       \hfill[\pageref{Lpepsims}]
%    
%    \quad|\pepgroups{|\meta{group1}|,|\meta{group2}|, ... , |\meta{groupn}|}|
%       \hfill[\pageref{Lpepgroups}]
%    
%    \quad|\DNAsims{|\meta{residue}|}{|\meta{similars}|}|
%       \hfill[\pageref{LDNAsims}]
%    
%    \quad|\DNAgroups{|\meta{group1}|,|\meta{group2}|, ... , |\meta{groupn}|}|
%       \hfill[\pageref{LDNAgroups}]
%       
%    \vspace{1.5\baselineskip}
%    
%    \textbf{Definition of new functional shading modes}
%    \medskip
%    
%    \quad|\clearfuncgroups| \hfill [\pageref{Lclearfuncgroups}]
%    
%    \quad|\funcgroup{|\meta{descr}|}{|\meta{residues}|}{|\meta{res.col.}|}{|\meta{shad.col.}|}|
%
%    \hfill|{|\meta{case}|}{|\meta{style}|}|
%       \hfill[\pageref{Lfuncgroup}]
%       
%    \vspace{1.5\baselineskip}
%    
%    \textbf{Appearance of the consensus line}
%    \medskip
%    
%    \quad|\threshold[|\meta{percentage}|]{|\meta{percentage}|}|
%       \hfill[\pageref{Lthreshold}]
%    
%    \quad|\constosingleseq{|\meta{seqref}|}| 
%       \hfill[\pageref{Lconstosingleseq}]
%    
%    \quad|\showconsensus[|\meta{color/scale}|[,|\meta{color/scale}|]]{|\meta{top/bot.}|}|
%       \hfill[\pageref{Lshowconsensus}]
%
%    \quad|\exportconsensus[|\meta{filename}|]{|\meta{seqref}|}|
%       \hfill[\pageref{Lexportconsensus}]
%
%    \quad|\hideconsensus|
%       \hfill[\pageref{Lhideconsensus}]
%    
%    \quad|\nameconsensus{|\meta{name}|}|
%       \hfill[\pageref{Lnameconsensus}]
%    
%    \quad|\defconsensus{|\meta{symbol1}|}{|\meta{symbol2}|}{|\meta{symbol3}|}|
%       \hfill[\pageref{Ldefconsensus}]
%       
%    \vspace*{-0.5\baselineskip}
%    
%    \begin{tabbing}
%    \quad|\consensuscolors|\=|{|\meta{res.col.1}|}{|\meta{shad.col.1}|}|\\
%    
%    \>|{|\meta{res.col.2}|}{|\meta{shad.col.2}|}|\\
%    
%    \>|{|\meta{res.col.3}|}{|\meta{shad.col.3}|}|
%        \hspace{1.2in}[\pageref{Lconsensuscolors}]\\
%    \end{tabbing}
%
%    \quad|\weighttable{|\meta{table}|}| 
%       \hfill[\pageref{Lweighttable}]
%    
%    \quad|\setweight{|\meta{res.1}|}{|\meta{res.2}|}{|\meta{value}|}|
%       \hfill[\pageref{Lsetweight}]
%    
%    \quad|\gappenalty{|\meta{value}|}|
%       \hfill[\pageref{Lgappenalty}]
%    
%    \vspace{1.5\baselineskip}
%    
%    \textbf{Sequence logos}
%    \medskip
%    
%    \quad|\showsequencelogo[|\meta{colorset}|]{|\meta{top/bottom}|}|
%       \hfill[\pageref{Lshowsequencelogo}]
%    
%    \quad|\hidesequencelogo|
%       \hfill[\pageref{Lhidesequencelogo}]
%    
%    \quad|\clearlogocolors[|\meta{color}|]|
%       \hfill[\pageref{Lclearlogocolors}]
%    
%    \quad|\logocolor{|\meta{residues}|}{|\meta{color}|}|
%       \hfill[\pageref{Llogocolor}]
%    
%    \quad|\showlogoscale[|\meta{color}|]{|\meta{left/right/leftright}|}|
%       \hfill[\pageref{Lshowlogoscale}]
%    
%    \quad|\hidelogoscale|
%       \hfill[\pageref{Lhidelogoscale}]
%    
%    \quad|\logostretch{|\meta{factor}|}|
%       \hfill[\pageref{Llogostretch}]
%    
%    \quad|\namesequencelogo{|\meta{name}|}|
%       \hfill[\pageref{Lnamesequencelogo}]
%    
%    \quad|\dofrequencycorrection|
%       \hfill[\pageref{Ldofrequencycorrection}]
%    
%    \quad|\undofrequencycorrection|
%       \hfill[\pageref{Lundofrequencycorrection}]
%    
%    \vspace{1.5\baselineskip}
%    
%    \textbf{Subfamily logos}
%    \medskip
%    
%    \quad|\showsubfamilylogo[|\meta{colorset}|]{|\meta{top/bottom}|}|
%       \hfill[\pageref{Lshowsubfamilylogo}]
%    
%    \quad|\hidesubfamilylogo|
%       \hfill[\pageref{Lhidesubfamilylogo}]
%    
%    \quad|\setsubfamily{|\meta{seqrefs}|}|
%       \hfill[\pageref{Lsetsubfamily}]
%    
%    \quad|\shownegatives[|\meta{weak, medium, strong}|]|
%       \hfill[\pageref{Lshownegatives}]
%    
%    \quad|\hidenegatives|
%       \hfill[\pageref{Lhidenegatives}]
%    
%    \quad|\namesubfamilylogo[|\meta{neg.name}|]{|\meta{name}|}|
%       \hfill[\pageref{Lnamesubfamilylogo}]
%    
%    \quad|\relevance{|\meta{bit-value}|}|
%       \hfill[\pageref{Lrelevance}]
%    
%    \quad|\showrelevance[|\meta{color}|]{|\meta{symbol}|}|
%       \hfill[\pageref{Lshowrelevance}]
%    
%    \quad|\hiderelevance|
%       \hfill[\pageref{Lhiderelevance}]
%    
%
%
%    \vspace{1.5\baselineskip}
%    
%    \textbf{Structure memes}
%    \medskip
%    
%    \quad|\structurememe[|\meta{seqref}|]{|\meta{filename}|}|
%       \hfill[\pageref{Lmemes}]
%    
%    \quad|\memeStandardcolors| 
%       \hfill[\pageref{Lmemes}]
%    
%    \quad|\memeRed{|\meta{residues}|}|
%       \hfill[\pageref{Lmemes}]
%    
%    \quad|\memeYellow{|\meta{residues}|}|
%       \hfill[\pageref{Lmemes}]
%    
%    \quad|\memeBlue{|\meta{residues}|}|
%       \hfill[\pageref{Lmemes}]
%    
%    \quad|\memeWhite{|\meta{residues}|}|
%       \hfill[\pageref{Lmemes}]
%    
%    \quad|\memeBlack{|\meta{residues}|}|
%       \hfill[\pageref{Lmemes}]
%    
%    \quad|\chimerachain{|\meta{chainlabel}|}|
%       \hfill[\pageref{Lmemes}]
%    
%    \quad|\chimeraballScale{|\meta{factor}|}| 
%       \hfill[\pageref{Lmemes}]
%    
%    \quad|\memelabelcutoff{|\meta{bit-value}|}|
%       \hfill[\pageref{Lmemes}]
%    
%    \quad|\chimeraxisdistance{|\meta{x-dist}|}{|\meta{y-dist}|}{|\meta{z-dist}|}|
%       \hfill[\pageref{Lmemes}]
%    
%    \quad|\echostructurefile{|\meta{strucref}|}|
%       \hfill[\pageref{Lmemes}]
%    
%
%    
%    \vspace{1.5\baselineskip}
%    
%    \textbf{Appearance of the sequence lines}
%    \medskip
%    
%    \quad|\shownames[|\meta{color}|]{|\meta{left/right}|}|
%       \hfill[\pageref{Lshownames}]
%    
%    \quad|\shownumbering[|\meta{color}|]{|\meta{left/right/leftright}|}|
%       \hfill[\pageref{Lshownumbering}]
%    
%    \quad|\nameseq{|\meta{seqref}|}{|\meta{name}|}| 
%       \hfill[\pageref{Lnameseq}]
%    
%    \quad|\namescolor{|\meta{color}|}|
%       \hfill[\pageref{Lnamescolor}]
%    
%    \quad|\namecolor{|\meta{seq1}|, ... ,|\meta{seq n}|}{|\meta{color}|}|
%       \hfill[\pageref{Lnamecolor}]
%    
%    \quad|\hidenames|
%       \hfill[\pageref{Lhidenames}]
%    
%    \quad|\hidename{|\meta{seq1}|, ... ,|\meta{seq n}|}|
%       \hfill[\pageref{Lhidename}]
%    
%    \quad|\numberingcolor{|\meta{color}|}|
%       \hfill[\pageref{Lnumberingcolor}]
%    
%    \quad|\numbercolor{|\meta{seq1}|, ... ,|\meta{seq n}|}{|\meta{color}|}|
%       \hfill[\pageref{Lnumbercolor}]
%    
%    \quad|\hidenumbering|
%       \hfill[\pageref{Lhidenumbering}]
%    
%    \quad|\hidenumber{|\meta{seq1}|, ... ,|\meta{seq n}|}|
%       \hfill[\pageref{Lhidenumber}]
%    
%    \quad|\hideresidues|
%       \hfill[\pageref{Lhideresidues}]
%    
%    \quad|\showresidues|
%       \hfill[\pageref{Lshowresidues}]
%    
%    \quad|\startnumber[|\meta{start..stop}|]{|\meta{seqref}|}{|\meta{startnumber}|}|
%       \hfill[\pageref{Lstartnumber}]
%    
%    \quad|\allowzero|
%       \hfill[\pageref{Lallowzero}]
%    
%    \quad|\disallowzero|
%       \hfill[\pageref{Lallowzero}]
%    
%    \quad|\seqlength{|\meta{seqref}|}{|\meta{length}|}|
%       \hfill[\pageref{Lseqlength}]
%    
%    \quad|\showruler[|\meta{color}|]{|\meta{top/bottom}|}{|\meta{seqref}|}|
%       \hfill[\pageref{Lshowruler}]
%    
%    \quad|\rulersteps{|\meta{number}|}|
%       \hfill[\pageref{Lrulersteps}]
%    
%    \quad|\rulercolor{|\meta{color}|}|
%       \hfill[\pageref{Lrulercolor}]
%    
%    \quad|\hideruler|
%       \hfill[\pageref{Lhideruler}]
%    
%    \quad|\rotateruler|
%       \hfill[\pageref{Lrotateruler}]
%    
%    \quad|\unrotateruler|
%       \hfill[\pageref{Lunrotateruler}]
%
%    \quad|\namerulerpos{|\meta{number}|}{|\meta{text}|[|\meta{color}|]}|
%       \hfill[\pageref{Lnamerulerpos}]
%    
%    \quad|\rulerspace{|\meta{length}|}|
%       \hfill[\pageref{Lrulerspace}]
%    
%    \quad|\gapchar{|\meta{symbol}|}| 
%       \qquad (incl. |rule|) \hfill [\pageref{Lgapchar}]
%    
%    \quad|\gapcolors{|\meta{symbol color}|}{|\meta{background color}|}|
%       \hfill[\pageref{Lgapcolors}]
%    
%    \quad|\showleadinggaps|
%       \hfill[\pageref{Lshowleadinggaps}]
%    
%    \quad|\hideleadinggaps|
%       \hfill[\pageref{Lhideleadinggaps}]
%    
%    \quad|\stopchar{|\meta{symbol}|}| 
%       \hfill [\pageref{Lstopchar}]
%    
%    \quad|\fingerprint{|\meta{res. per line}|}|
%       \hfill[\pageref{Lfingerprint}]
%       
%    \vspace{1.5\baselineskip}
%    
%    \textbf{Displaying selected residues in the alignment}
%    \medskip
%    
%    \quad|\setends[|\meta{startnumber}|]{|\meta{seqref}|}{|\meta{start..stop}|}|
%       \hfill[\pageref{Lsetends}]
%    
%    \quad|\setdomain{|\meta{seqref}|}{|\meta{selection}|}| (see |\shaderegion| p.\,\pageref{Lshaderegion})
%       \hfill[\pageref{Lsetdomain}]
%
%    \quad|\domaingaprule{|\meta{thickness}|}|
%       \hfill[\pageref{Lsetdomain}]
%
%    \quad|\domaingapcolors{|\meta{foreground}|}{|\meta{background}|}|
%       \hfill[\pageref{Lsetdomain}]
%
%    \vspace{1.5\baselineskip}
%    
%    \textbf{Hiding, killing, separating and ordering}
%    \medskip
%    
%    \quad|\hideseq{|\meta{seq1}|,|\meta{seq2}|,|\ldots|,|\meta{seq n}|}|
%       \hfill[\pageref{Lhideseq}]
%    
%    \quad|\hideseqs|
%       \hfill[\pageref{Lhideseqs}]
%    
%    \quad|\showseqs|
%       \hfill[\pageref{Lshowseqs}]
%    
%    \quad|\killseq{|\meta{seq1}|,|\meta{seq2}|,|\ldots|,|\meta{seq n}|}|
%       \hfill[\pageref{Lkillseq}]
%    
%    \quad|\donotshade{|\meta{seq1}|,|\meta{seq2},\ldots|,|\meta{seq n}|}|
%       \hfill[\pageref{Ldonotshade}]
%    
%    \quad|\separationline{|\meta{seqref}|}|
%       \hfill[\pageref{Lseparationline}]
%    
%    \quad|\smallsep|
%       \hfill[\pageref{Lsmallsep}]
%    
%    \quad|\medsep|
%       \hfill[\pageref{Lmedsep}]
%    
%    \quad|\bigsep|
%       \hfill[\pageref{Lbigsep}]
%    
%    \quad|\vsepspace{|\meta{length}|}|
%       \hfill[\pageref{Lvsepspace}]
%    
%    \quad|\orderseqs{|\meta{seq1}|,|\meta{seq2}|,|\ldots|,|\meta{seq n}|}|
%       \hfill[\pageref{Lorderseqs}]
%
%    \vspace{1.5\baselineskip}
%    
%    \textbf{Residues per line and further settings}
%    \medskip
%    
%    \quad|\residuesperline{|\meta{number}|}|
%       \hfill[\pageref{Lresiduesperline}]
%    
%    \quad|\residuesperline*{|\meta{number}|}|
%       \hfill[\pageref{Lresiduesperline*}]
%    
%    \quad|\charstretch{|\meta{factor}|}|
%       \hfill[\pageref{Lcharstretch}]
%    
%    \quad|\linestretch{|\meta{factor}|}|
%       \hfill[\pageref{Llinestretch}]
%    
%    \quad|\numberingwidth{|\meta{n digits}|}|
%       \hfill[\pageref{Lnumberingwidth}]
%    
%    \quad|\smallblockskip|
%       \hfill[\pageref{Lsmallblockskip}]
%    
%    \quad|\medblockskip|
%       \hfill[\pageref{Lmedblockskip}]
%    
%    \quad|\bigblockskip|
%       \hfill[\pageref{Lbigblockskip}]
%    
%    \quad|\noblockskip|
%       \hfill[\pageref{Lnoblockskip}]
%    
%    \quad|\vblockspace{|\meta{length}|}|
%       \hfill[\pageref{Lvblockspace}]
%    
%    \quad|\flexblockspace|
%       \hfill[\pageref{Lflexblockspace}]
%    
%    \quad|\fixblockspace|
%       \hfill[\pageref{Lfixblockspace}]
%    
%    \quad|\alignment{|\meta{left/center/right}|}|
%       \hfill[\pageref{Lalignment}]
%       
%    \vspace{1.5\baselineskip}
%    
%    
%
%    
%    \textbf{Individual shading and labeling of sequence stretches}
%    \medskip
%    
%    \quad|\shaderegion{|\meta{seqref}|}{|\meta{selection}|}{|\meta{res.col.}|}{|\meta{shad.col.}|}|
%       \hfill[\pageref{Lshaderegion}]
%
%    \medskip
%    
%    \quad |{|\meta{selection}|}| =
%    \medskip%    
%    
%    \quad\quad|{|\meta{start1}..\meta{stop1}|,|\meta{start2}..\meta{stop2}|,|\ldots|,|\meta{start n}..\meta{stop n}|}|
%    \medskip
%    
%    \quad\quad|{point[|\meta{dist}|]:|\meta{file}|,|\meta{num}|[CA/side]}|
%    \medskip
%    
%    \quad\quad|{line[|\meta{dist}|]:|\meta{file}|,|\meta{num1}|[CA/side],|\meta{num2}|[CA/side]}|
%    \medskip
%    
%    \quad\quad|{plane[|\meta{dist}|]:|\meta{file}|,|\meta{num1}|[CA/side],|\meta{num2}|[CA/side],|
%    \medskip
%
%    \hfill\meta{num3}|[CA/side]}|
%    \medskip
%    
%    
%    \quad|\printPDBlist{|\meta{selection}|}| \quad|\messagePDBlist{|\meta{selection}|}|
%       \hfill[\pageref{LprintPDBlist}]
%    
%    \quad|\shadeblock{|\meta{seqref}|}{|\meta{selection}|}{|\meta{res.col.}|}{|\meta{shad.col.}|}|
%       \hfill[\pageref{Lshadeblock}]
%    
%    \quad|\changeshadingcolors{|\meta{seqref}|}{|\meta{selection}|}{|\meta{name}|}|
%       \hfill[\pageref{Lchangeshadingcolors}]
%    
%    \quad|\emphregion{|\meta{seqref}|}{|\meta{selection}|}|
%       \hfill[\pageref{Lemphregion}]
%    
%    \quad|\emphblock{|\meta{seqref}|}{|\meta{selection}|}|
%       \hfill[\pageref{Lemphblock}]
%    
%    \quad|\emphdefault{|\meta{style}|}|
%       \hfill[\pageref{Lemphdefault}]
%    
%    \quad|\tintregion{|\meta{seqref}|}{|\meta{selection}|}|
%       \hfill[\pageref{Ltintregion}]
%    
%    \quad|\tintblock{|\meta{seqref}|}{|\meta{selection}|}|
%       \hfill[\pageref{Ltintblock}]
%    
%    \quad|\tintdefault{|\meta{effect}|}| \qquad\qquad|weak, normal, strong|
%       \hfill[\pageref{Ltintdefault}]
%    
%    \quad|\lowerregion{|\meta{seqref}|}{|\meta{selection}|}|
%       \hfill[\pageref{Llowerregion}]
%    
%    \quad|\lowerblock{|\meta{seqref}|}{|\meta{selection}|}|
%       \hfill[\pageref{Llowerblock}]
%    
%    \quad|\frameblock{|\meta{seqref}|}{|\meta{selection}|}{|\meta{color}|[|\meta{length}|]}|
%       \hfill[\pageref{Lframeblock}]
%    
%    \quad|\feature{|\meta{position}|}{|\meta{seqref}|}{|\meta{selection}|}|
%
%       \hfill |{|\meta{labelstyle}|}{|\meta{text}|}| [\pageref{Lfeature}]
%    
%    \begin{tabbing}
%    \quad\quad\quad|{|\meta{labelstyle}|}|\ \= = |{brace[|\meta{color}|]}|\\
%       \> = |{fill:|\meta{symbol}|[|\meta{textcolor}|]}|\\
%       \> = |{restriction[|\meta{color}|]}|\\
%       \> = |{helix[|\meta{helixcolor}|]}|\\
%       \> = |{box[|\meta{framecolor,boxcolor}|][|\meta{length}|]:|\\
%               \hspace{8.7cm}\meta{text}|[|\meta{textcolor}|]}|\\
%       \> = |{S-S}|\\
%       \> = arrows and bars (|-=<',|$\vert$|o|)(|-=|)(|-=>',|$\vert$|o|)\\
%       \> = |{translate[|\meta{color}|]}|\\
%       \> = |{bar[|\meta{min}|,|\meta{max}|]:|\\
%               \hspace{5cm}\meta{properties/file/data}|[|\meta{color(,bgcolor)}|]}|\\
%       \> = |{color[|\meta{min}|,|\meta{max}|]:|\\
%               \hspace{5cm}\meta{properties/file/data}|[|\meta{scale}|]}|\\
%               \hspace{5cm}\meta{properties}: |hydrophobicity|, |charge|,\\
%               \hspace{7.4cm}|molweight|, |conservation|\\
%               \hspace{5cm}\meta{scale}: |Gray|, |BlueRed|, |RedBlue|,\\ 
%               \hspace{6.5cm}|GreenRed|, |RedGreen|, |ColdHot|,\\
%               \hspace{6.5cm}|HotCold|, |T-Coffee|\\
%    \end{tabbing}
%    
%    \quad|\ttopspace{|\meta{length}|}|
%       \hfill[\pageref{Lttopspace}]
%    
%    \quad|\topspace{|\meta{length}|}|
%       \hfill[\pageref{Ltopspace}]
%    
%    \quad|\bottomspace{|\meta{length}|}|
%       \hfill[\pageref{Lbottomspace}]
%    
%    \quad|\bbottomspace{|\meta{length}|}|
%       \hfill[\pageref{Lbottomspace}]
%    
%    \quad|\featurerule{|\meta{length}|}|
%       \hfill[\pageref{Lfeaturerule}]
%       
%    \quad|\bargraphstretch{|\meta{factor}|}|
%       \hfill[\pageref{Lbargraphstretch}]
%       
%    \quad|\colorscalestretch{|\meta{factor}|}|
%       \hfill[\pageref{Lcolorscalestretch}]
%       
%    \quad|\codon{|\meta{amino acid}|}{|\meta{triplet1,\ldots, triplet n}|}|
%       \hfill[\pageref{Lcodon}]
%    
%    \quad|\geneticcode{|\meta{filename}|}|
%       \hfill[\pageref{Lgeneticcode}]
%    
%    \quad|\backtranslabel[|\meta{size}|]{|\meta{style}|}|
%       \hfill[\pageref{Lbacktranslabel}]
%       
%    \quad|\backtranstext[|\meta{size}|]{|\meta{style}|}|
%       \hfill[\pageref{Lbacktranstext}]
%       
%    \begin{tabbing}
%    \quad\quad\quad|{|\meta{style}|}|\ \= = |{horizontal}|\\
%       \> = |{alternating}|\\
%       \> = |{zigzag}|\\
%       \> = |{oblique}|\\
%       \> = |{vertical}|
%    \end{tabbing}
%       
%    \quad |\showfeaturename{|\meta{ttttop...bbbbottom}|}{|\meta{name}|}|
%       \hfill[\pageref{Lshowfeaturename}]
%
%    \quad |\showfeaturestylename{|\meta{ttttop...bbbbottom}|}{|\meta{name}|}|
%       \hfill[\pageref{Lshowfeaturestylename}]
%
%    \quad |\hidefeaturename{|\meta{ttttop...bbbbottom}|}|
%       \hfill[\pageref{Lhidefeaturename}]
%
%    \quad |\hidefeaturestylename{|\meta{ttttop...bbbbottom}|}|
%       \hfill[\pageref{Lhidefeaturestylename}]
%
%    \quad |\hidefeaturenames|
%       \hfill[\pageref{Lhidefeaturenames}]
%
%    \quad|\hidefeaturestylenames|
%       \hfill[\pageref{Lhidefeaturestylenames}]
%
%    \quad |\featurenamescolor{|\meta{color}|}|
%       \hfill[\pageref{Lfeaturenamescolor}]
%
%    \quad |\featurestylenamescolor{|\meta{color}|}| 
%       \hfill[\pageref{Lfeaturestylenamescolor}]
%
%    \quad |\featurenamecolor{|\meta{ttttop...bbbbottom}|}{|\meta{color}|}|
%       \hfill[\pageref{Lfeaturenamecolor}]
%
%    \quad |\featurestylenamecolor{|\meta{ttttop...bbbbottom}|}{|\meta{color}|}| 
%       \hfill[\pageref{Lfeaturestylenamecolor}]
%
%    \vspace{1.5\baselineskip}
%    
%
%
%    \textbf{Including secondary structure information}
%    \medskip
%    
%    \quad|\includeDSSP[make new]{|\meta{seqref}|}{|\meta{filename}|}|
%       \hfill[\pageref{LincludeDSSP}]
%    
%    \quad|\includeSTRIDE[make new]{|\meta{seqref}|}{|\meta{filename}|}|
%       \hfill[\pageref{LincludeSTRIDE}]
%    
%    \quad|\includePHDsec[make new]{|\meta{seqref}|}{|\meta{filename}|}|
%       \hfill[\pageref{LincludePHDsec}]
%    
%    \quad|\includePHDtopo[make new]{|\meta{seqref}|}{|\meta{filename}|}|
%       \hfill[\pageref{LincludePHDtopo}]
%    
%    \quad|\includeHMMTOP[make new]{|\meta{seqref}|[|\meta{seqref}|]}{|\meta{filename}|}|
%       \hfill[\pageref{LincludeHMMTOP}]
%    
%    \quad|\showonDSSP{|\meta{structures}|}|
%       \hfill[\pageref{LshowonDSSP}]
%    
%    \quad|\showonSTRIDE{|\meta{structures}|}|
%       \hfill[\pageref{LshowonSTRIDE}]
%    
%    \quad|\showonPHDsec{|\meta{structures}|}|
%       \hfill[\pageref{LshowonPHDsec}]
%    
%    \quad|\showonPHDtopo{|\meta{structures}|}|
%       \hfill[\pageref{LshowonPHDtopo}]
%    
%    \quad|\showonHMMTOP{|\meta{structures}|}|
%       \hfill[\pageref{LshowonHMMTOP}]
%    
%    \quad|\hideonDSSP{|\meta{structures}|}|
%       \hfill[\pageref{LhideonDSSP}]
%    
%    \quad|\hideonSTRIDE{|\meta{structures}|}|
%       \hfill[\pageref{LhideonSTRIDE}]
%    
%    \quad|\hideonPHDsec{|\meta{structures}|}|
%       \hfill[\pageref{LhideonPHDsec}]
%    
%    \quad|\hideonPHDtopo{|\meta{structures}|}|
%       \hfill[\pageref{LhideonPHDtopo}]
%    
%    \quad|\hideonHMMTOP{|\meta{structures}|}|
%       \hfill[\pageref{LhideonHMMTOP}]
%    
%    \quad|\appearance{|\meta{type}|}{|\meta{position}|}{|\meta{labelstyle}|}{|\meta{text}|}|
%       \hfill[\pageref{Lappearance}]
%    
%    \quad|\numcount|
%       \hfill[\pageref{Lnumcount}]
%    
%    \quad|\alphacount|
%       \hfill[\pageref{Lalphacount}]
%    
%    \quad|\Alphacount|
%       \hfill[\pageref{LAlphacount}]
%       
%    \quad|\romancount|
%       \hfill[\pageref{Lromancount}]
%    
%    \quad|\Romancount|
%       \hfill[\pageref{LRomancount}]
%    
%    \quad|\firstcolumnDSSP|
%       \hfill[\pageref{LfirstcolumnDSSP}]   
%             
%    \quad|\secondcolumnDSSP|
%       \hfill[\pageref{LsecondcolumnDSSP}]   
%             
%    \vspace{1.5\baselineskip}
%    
%
%    \textbf{Displaying and building legends}
%    \medskip
%    
%    \quad|\showlegend|
%       \hfill[\pageref{Lshowlegend}]
%    
%    \quad|\hidelegend|
%       \hfill[\pageref{Lhidelegend}]
%    
%    \quad|\movelegend{|\meta{x-offset}|}{|\meta{y-offset}|}|
%       \hfill[\pageref{Lmovelegend}]
%    
%    \quad|\germanlanguage|, |\spanishlanguage|, |\englishlanguage|
%       \hfill[\pageref{Lgermanlanguage}]
%    
%    \quad|\legendcolor{|\meta{color}|}|
%       \hfill[\pageref{Llegendcolor}]
%    
%    \quad|\shadebox{|\meta{color}|}|
%       \hfill[\pageref{Lshadebox}]
%       
%    \vspace{1.5\baselineskip}
%    
%    
%    \textbf{Adding captions to the alignment}
%    \medskip
%    
%    \quad|\showcaption[|\meta{top/bottom}|]{|\meta{text}|}|
%       \hfill[\pageref{Lshowcaption}]
%    
%    \quad|\shortcaption{|\meta{text}|}|
%       \hfill[\pageref{Lshortcaption}]
%    
%    \vspace{1.5\baselineskip}
%    
%    \newpage    
%    
%    \textbf{Font handling}
%    \medskip
%    
%    \quad|\setfamily{|\meta{text}|}{|\meta{family}|}|
%       \hfill[\pageref{Lsetfamily}]
%    
%    \quad|\setseries{|\meta{text}|}{|\meta{series}|}|
%       \hfill[\pageref{Lsetseries}]
%    
%    \quad|\setshape{|\meta{text}|}{|\meta{shape}|}|
%       \hfill[\pageref{Lsetshape}]
%    
%    \quad|\setsize{|\meta{text}|}{|\meta{size}|}|
%       \hfill[\pageref{Lsetsize}]
%    
%    \quad|\setfont{|\meta{text}|}{|\meta{family}|}{|\meta{series}|}{|\meta{shape}|}{|\meta{size}|}|
%       \hfill[\pageref{Lsetfont}]
%       
%    \medskip
%    
%    \quad |\featuresrm| \quad |\featurestiny| \hfill[\pageref{Lfeaturesrm}]
%    
%    \quad |\featuressf| \quad |\featuresscriptsize|
%    
%    \quad |\featurestt| \quad |\featuresfootnotesize|
%    
%    \quad |\featuresbf| \quad |\featuressmall|
%    
%    \quad |\featuresmd| \quad |\featuresnormalsize|
%    
%    \quad |\featuresit| \quad |\featureslarge|
%    
%    \quad |\featuressl| \quad |\featuresLarge|
%    
%    \quad |\featuressc| \quad |\featuresLARGE|
%    
%    \quad |\featuresup| \quad |\featureshuge|
%    
%    \quad |           | \quad |\featuresHuge|
%    \medskip
%    
%    Corresponding sets are provided for the
%    numbering (|\numberingrm| etc.), 
%    featurestyles (|featurestylesrm| etc.), names (|\namesrm| etc.), 
%    featurenames (|\featurenamesrm| etc.),
%    featurestylenames (|\featurestylenames| etc.),
%    residues (|\residuesrm| etc.),
%    hideblock labels (|hideblockrm| etc.), 
%    rulers (|\rulerrm| etc.), and
%    legend texts (|legendrm| etc.).
%    \bigskip
%    
%   
%    \textbf{Goodies---molweight and charge}
%    \medskip
%    
%    \quad|\molweight{|\meta{seqref}|}{|\meta{Da/kDa}|}|
%       \hfill[\pageref{Lmolweight}]
%    
%    \quad|\charge{|\meta{seqref}|}{|\meta{i/o/N/C}|}|    
%       \hfill[\pageref{Lcharge}]
%    
%    \quad|\percentsimilarity{|\meta{seqref1}|}{|\meta{seqref2}|}| 
%       \hfill[\pageref{Lpercentsimilarity}]
%
%    \quad|\percentidentity{|\meta{seqref1}|}{|\meta{seqref2}|}|
%       \hfill[\pageref{Lpercentidentity}]
%
%    \quad|\similaritytable|
%       \hfill[\pageref{Lsimilaritytable}]
%    
%    \StopEventually{%
%    \newpage
%    \section{References}
%      [1] \textsc{Carlisle, D.} The Standard \LaTeX{} `Graphics 
%          Bundle', |color.sty|.
%    
%      [2] \textsc{Karlin, S.; Ghandour, G.} (1985) Multiple-alphabet
%          amino acid sequence comparisons of the immunoglobulin
%          $\kappa$-chain constant domain.
%          \newblock \textit{Proc. Natl. Acad. Sci. USA}: \textbf{82},
%          8597--8601.
%          
%      [3] \textsc{Kyte, J.; Doolittle, R. F.} (1982) A simple
%          method for displaying the hydropathic character of a
%          protein.
%          \newblock \textit{J. Mol. Biol.}: \textbf{157}, 105--132.
%          
%      [4] \textsc{Rose, G. D.; Geselowitz, A. R.; Lesser, G. J.; 
%          Lee, R. H.; Zehfus, M. H.} (1985) Hydrophobicity of amino 
%          acid residues in globular proteins.
%          \newblock \textit{Science}: \textbf{229}, 835--838.
%          
%      [5] \textsc{Lesser, G. J.; Rose, G. D.} (1990) Hydrophobicity
%          of amino acid subgroups in proteins.
%          \newblock \textit{Proteins: structure, function and 
%          genetics}: \textbf{8}, 6--13.
%        
%      [6] \textsc{Fr\"ohlich, K.-U.} (1994) Sequence similarity 
%          presenter: a tool for the graphic display of similarities
%          of long sequences for use in presentations.
%          \newblock \textit{Comput. Applic. Biosci.}:
%          \textbf{10}, 179--183.
%          
%      [7] \textsc{Schneider, T.D.; Stephens, R.M.} (1990) Sequence logos:
%          a new way to display consensus
%          \newblock \textit{Nucleic Acid Res.}: \textbf{18}, 6097--6100.
% 
%      [8] DeLano Scientific LLC `www.pymol.org'
%         
%      [9] \textsc{Kabsch, W.; Sander, C.} (1983) Dictionary of
%          protein secondary structure: pattern recognition of
%          hydrogen-bonded and geometrical features.
%          \newblock \textit{Biopolymers}: \textbf{22}, 2577--2637.
%          
%      [10] \textsc{Frishman, D.; Argos, P.} (1995) Knowledge-based
%          protein secondary structure assignment.
%          \newblock \textit{Proteins: structure, function and 
%          genetics}: \textbf{23}, 566--579.    
%          
%      [11] \textsc{Rost, B.; Sander, C.} (1994)
%          Combining evolutionary information and neural networks to predict
%          protein secondary structure.
%          \newblock \textit{Proteins: structure, function and 
%          genetics}: \textbf{19}, 55--72.
%
%      [12] \textsc{Tusnady, G.E.; Simon, I.} (2001) 
%           The HMMTOP transmembrane topology prediction server.
%           \newblock \textit{Bioinformatics}: \textbf{17}, 849-850.
%              
%      [13] \textsc{Rokicki, T.} DVIPS: A \TeX{} driver.
%         
%      [14] \textsc{Beitz, E.} (2006) Subfamily logos: visualization of sequence 
%           deviations at alignment positions with high information content.
%           \newblock \textit{BMC Bioinformatics}: \textbf{7}:313.
%         
%    }
%    \section{Implementation}
%    \subsection{Documentation Driver}
%    \begin{macrocode}
%<*driver>
\documentclass[12pt,a4paper]{ltxdoc} 
\usepackage[dvips]{texshade}
\usepackage{graphicx}
\openin\structurefile = hyperref.sty
\ifeof\structurefile
 \else
%     \usepackage[dvips,colorlinks]{hyperref}
     \usepackage[pdftex,colorlinks]{hyperref}
 \fi
\closein\structurefile
\DisableCrossrefs
\sloppy
\def\BioTeX{\textsc{Bio}\kern-0.5ex\TeX}
\def\TeXtopo{\mbox{\TeX\textsf{topo}}}
\begin{document}
  \OnlyDescription
  \DocInput{texshade.dtx}
\end{document}
%</driver>
%    \end{macrocode}
%    \subsection{\texttt{texshade.sty}---no comments}
%    \begin{macrocode}
%<*texshade>
\NeedsTeXFormat{LaTeX2e}
\ProvidesPackage{texshade}[2022/08/29 LaTeX TeXshade (v1.26a)]
\message{Package `texshade', Version 1.26a of 2022/08/29.}

\PassOptionsToPackage{dvips}{color}
\PassOptionsToPackage{dvips}{graphicx}
\DeclareOption*{%
   \PassOptionsToPackage{\CurrentOption}{color}%
   \PassOptionsToPackage{\CurrentOption}{graphicx}%
}
\ProcessOptions
\RequirePackage{color,graphics}

\expandafter\ifx\csname TeXshade\endcsname\relax \else \endinput \fi

\expandafter\ifx\csname TeXtopo\endcsname\relax \else
 \PackageError{TeXtopo}
  {TeXtopo loaded before TeXshade}
  {\MessageBreak
   For the proper function of the TeXtopo/TeXshade combo the \MessageBreak
   TeXshade package must be loaded before the TeXtopo package.\MessageBreak
   Please change the order of the \noexpand\usepackage commands in your
   \MessageBreak
   document header section or use the `biotex.sty'.\MessageBreak\MessageBreak
   Quit here by typing \space X <return>. \MessageBreak
}
\fi

\catcode`\@11

\def\rotopo#1{%
  \Grot@setangle{#1}%
  \setbox\z@\hbox\bgroup\ignorespaces}
\def\endrotopo{%
  \unskip\egroup
  \Grot@x\z@
  \Grot@y\z@
  \wd0\z@\dp0\z@\ht0\z@
  \Grot@box
}

\newread\structurefile  \newwrite\featurefile
\newread\alignfile      \newread\sublogofile
\newwrite\exp@rtfile
\newwrite\exp@rt@chimerafile

\expandafter\ifx\csname blacktriangleright\endcsname\relax
 \openin\structurefile = amssymb.sty
 \ifeof\structurefile
  \message{<AMS symbol style `amssymb.sty' not installed - using round heads>}
  \def\blacktriangleright{%
    \rule[\width@tmp]{0.65ex}{\temp@@length}\kern-0.55ex\ensuremath{\bullet}%
  }
  \def\blacktriangleleft{%
    \ensuremath{\bullet}\kern-0.55ex\rule[\width@tmp]{0.65ex}{\temp@@length}%
  }
 \else \RequirePackage[]{amssymb} \fi
 \closein\structurefile
\fi

\DeclareSymbolFont{alphahelix}{OML}{cmm}{m}{it}
\DeclareMathSymbol{\helixhook}{\mathrel}{alphahelix}{"5E}

\newcount\loopcount     \newcount\innerloopcount   \newcount\outerloopcount
\newcount\seq@count     \newcount\killseq@count
\newcount\seq@percent   \newcount\res@count
\newcount\seq@pointer   \newcount\pos@count
\newcount\res@perline   \newcount\end@count
\newcount\cons@count    \newcount\total@count
\newcount\temp@count    \newcount\triple@count
\newcount\temp@@count   \newcount\pos@sum

\newlength\box@width    \newlength\name@width
\newlength\box@depth    \newlength\width@tmp
\newlength\box@height   \newlength\number@width
\newlength\line@stretch
\newlength\center@fill  \newlength\arrow@width
\newlength\arrow@height \newlength\rule@thick
\newlength\arrow@thick  \newlength\logo@height
\newlength\equal@width  \newlength\equal@tmp
\newlength\equal@height \newlength\temp@@length
\newlength\vspace@legend
\newlength\hspace@legend
\newlength\bar@length

\newif\ifletter       \newif\ifnumber
\newif\ifnewres       \newif\ifall@shade
\newif\ifnames@right
\newif\ifnumbers@left \newif\ifnumbers@right
\newif\ifhide@cons    \newif\ifshow@cons
\newif\iffuncmode     \newif\iflegend@
\newif\ifT@coffee
\newif\ifnumbers@     \newif\ifnames@
\newif\ifgerm@n       \newif\ifsp@nish
\newif\ifrpl@fix
\newif\ifnosh@de      \newif\ifregionalshade
\newif\ifstart@       \newif\ifstop@
\newif\iftopfeature   \newif\ifbottomfeature
\newif\ifttopfeature  \newif\ifbbottomfeature
\newif\iftttopfeature \newif\ifbbbottomfeature
\newif\ifttttopfeature\newif\ifbbbbottomfeature
\newif\ifall@fshade   \newif\ifregionalemph 
\newif\ifregionallower    
\newif\ifframe@       \newif\ifregionaltint
\newif\ifshading@
\newif\ifshow@logo    \newif\ifshow@sublogo
\newif\ifhidechar     \newif\ifsh@wg@ps
\newif\ifsimmode
\newif\ifregionaltintnow
\newif\ifregionalemphnow
\newif\ifregionallowernow
\newif\ifregionalshadenow
\newif\iftopfeaturenow
\newif\ifttopfeaturenow
\newif\iftttopfeaturenow
\newif\ifttttopfeaturenow
\newif\ifbottomfeaturenow
\newif\ifbbottomfeaturenow
\newif\ifbbbottomfeaturenow
\newif\ifbbbbottomfeaturenow
\newif\ifframenow
\newif\ifshadingnow
\newif\iffix@

\expandafter\ifx\csname mdqon\endcsname\relax 
     \germ@nfalse \sp@nishfalse \def\cons@name{consensus}
\else \germ@ntrue \sp@nishfalse \def\cons@name{Konsensus} \fi

\def\n@me{Name:} \def\@msf{MSF:} \def\he@derend{//} \def\ampers@nd{&}
\def\comm@{,} \def\@loc{LOC} \def\@asg{ASG} \def\@t{@} \def\@HP{>HP:}
\def\gre@ter{>} \def\sm@ller{<} \def\N@{N} \def\equ@l{=} \def\H@{H}
\def\gap@char{.} \def\dom@char{{\dom@rule}} \def\yes{yes} \def\y@{y} \def\n@{n} 
\def\o@{o} \def\d@t{.} \def\questi@n{?} \def\st@p@char{*} \def\semic@n{;}
\def\br@cket{[} \def\p@r@gr@ph{€}
\def\@TOM{ATOM} \def\C@lpha{CA} \def\@point@{point} \def\@line@{line} \def\@plane@{plane}
\def\gap@rule{\rule[0.3\box@height]{\box@width}{\gap@rulethick}}
\def\dom@rule{\vrule depth\box@depth height\box@height width\domgap@rulethick}
\def\fgroup@num{0} \def\max@seqnumber{0}  \def\@lign@count{0}
\def\resn@m@tch{upper} \def\ressimm@tch{upper}
\def\resm@tch{upper}   \def\res@llm@tch{upper}
\def\tr@ns{translate}  \def\gr@ydef@ult{GrayDefault}
\xdef\par@{\string\par}
\expandafter\def\csname fg@textcolor/\endcsname{White}
\expandafter\def\csname fg@color/\endcsname{White}
\expandafter\def\csname func@style/\endcsname{\csname textup\endcsname}
\expandafter\def\csname func@style*\endcsname{\csname textup\endcsname}
\expandafter\def\csname funcm@tch/\endcsname{upper}
\expandafter\def\csname funcm@tch*\endcsname{upper}

\setlength\hspace@legend{0pt} \setlength\vspace@legend{0pt}

\expandafter\xdef\csname log2@1\endcsname{0}
\expandafter\xdef\csname log2@2\endcsname{1000}
\expandafter\xdef\csname log2@3\endcsname{1585}
\expandafter\xdef\csname log2@4\endcsname{2000}
\expandafter\xdef\csname log2@5\endcsname{2322}
\expandafter\xdef\csname log2@6\endcsname{2585}
\expandafter\xdef\csname log2@7\endcsname{2807}
\expandafter\xdef\csname log2@8\endcsname{3000}
\expandafter\xdef\csname log2@9\endcsname{3170}
\expandafter\xdef\csname log2@10\endcsname{3322}
\expandafter\xdef\csname log2@11\endcsname{3459}
\expandafter\xdef\csname log2@12\endcsname{3585}
\expandafter\xdef\csname log2@13\endcsname{3700}
\expandafter\xdef\csname log2@14\endcsname{3807}
\expandafter\xdef\csname log2@15\endcsname{3907}
\expandafter\xdef\csname log2@16\endcsname{4000}
\expandafter\xdef\csname log2@17\endcsname{4087}
\expandafter\xdef\csname log2@18\endcsname{4170}
\expandafter\xdef\csname log2@19\endcsname{4248}
\expandafter\xdef\csname log2@20\endcsname{4322}
\expandafter\xdef\csname log2@21\endcsname{4392}
\expandafter\xdef\csname log2@22\endcsname{4459}
\expandafter\xdef\csname log2@23\endcsname{4524}
\expandafter\xdef\csname log2@24\endcsname{4585}
\expandafter\xdef\csname log2@25\endcsname{4644}
\expandafter\xdef\csname log2@26\endcsname{4700}
\expandafter\xdef\csname log2@27\endcsname{4755}
\expandafter\xdef\csname log2@28\endcsname{4807}
\expandafter\xdef\csname log2@29\endcsname{4858}
\expandafter\xdef\csname log2@30\endcsname{4907}
\expandafter\xdef\csname log2@31\endcsname{4954}
\expandafter\xdef\csname log2@32\endcsname{5000}
\expandafter\xdef\csname log2@33\endcsname{5044}
\expandafter\xdef\csname log2@34\endcsname{5087}
\expandafter\xdef\csname log2@35\endcsname{5129}
\expandafter\xdef\csname log2@36\endcsname{5170}
\expandafter\xdef\csname log2@37\endcsname{5209}
\expandafter\xdef\csname log2@38\endcsname{5248}
\expandafter\xdef\csname log2@39\endcsname{5285}
\expandafter\xdef\csname log2@40\endcsname{5322}
\expandafter\xdef\csname log2@41\endcsname{5358}
\expandafter\xdef\csname log2@42\endcsname{5392}
\expandafter\xdef\csname log2@43\endcsname{5426}
\expandafter\xdef\csname log2@44\endcsname{5459}
\expandafter\xdef\csname log2@45\endcsname{5492}
\expandafter\xdef\csname log2@46\endcsname{5524}
\expandafter\xdef\csname log2@47\endcsname{5555}
\expandafter\xdef\csname log2@48\endcsname{5585}
\expandafter\xdef\csname log2@49\endcsname{5615}
\expandafter\xdef\csname log2@50\endcsname{5644}
\expandafter\xdef\csname log2@51\endcsname{5672}
\expandafter\xdef\csname log2@52\endcsname{5700}
\expandafter\xdef\csname log2@53\endcsname{5728}
\expandafter\xdef\csname log2@54\endcsname{5755}
\expandafter\xdef\csname log2@55\endcsname{5781}
\expandafter\xdef\csname log2@56\endcsname{5807}
\expandafter\xdef\csname log2@57\endcsname{5833}
\expandafter\xdef\csname log2@58\endcsname{5858}
\expandafter\xdef\csname log2@59\endcsname{5883}
\expandafter\xdef\csname log2@60\endcsname{5907}
\expandafter\xdef\csname log2@61\endcsname{5931}
\expandafter\xdef\csname log2@62\endcsname{5954}
\expandafter\xdef\csname log2@63\endcsname{5977}
\expandafter\xdef\csname log2@64\endcsname{6000}
\expandafter\xdef\csname log2@65\endcsname{6022}
\expandafter\xdef\csname log2@66\endcsname{6044}
\expandafter\xdef\csname log2@67\endcsname{6066}
\expandafter\xdef\csname log2@68\endcsname{6087}
\expandafter\xdef\csname log2@69\endcsname{6109}
\expandafter\xdef\csname log2@70\endcsname{6129}
\expandafter\xdef\csname log2@71\endcsname{6149}
\expandafter\xdef\csname log2@72\endcsname{6170}
\expandafter\xdef\csname log2@73\endcsname{6190}
\expandafter\xdef\csname log2@74\endcsname{6209}
\expandafter\xdef\csname log2@75\endcsname{6229}
\expandafter\xdef\csname log2@76\endcsname{6248}
\expandafter\xdef\csname log2@77\endcsname{6267}
\expandafter\xdef\csname log2@78\endcsname{6285}
\expandafter\xdef\csname log2@79\endcsname{6304}
\expandafter\xdef\csname log2@80\endcsname{6322}
\expandafter\xdef\csname log2@81\endcsname{6340}
\expandafter\xdef\csname log2@82\endcsname{6358}
\expandafter\xdef\csname log2@83\endcsname{6375}
\expandafter\xdef\csname log2@84\endcsname{6392}
\expandafter\xdef\csname log2@85\endcsname{6409}
\expandafter\xdef\csname log2@86\endcsname{6426}
\expandafter\xdef\csname log2@87\endcsname{6443}
\expandafter\xdef\csname log2@88\endcsname{6459}
\expandafter\xdef\csname log2@89\endcsname{6476}
\expandafter\xdef\csname log2@90\endcsname{6492}
\expandafter\xdef\csname log2@91\endcsname{6508}
\expandafter\xdef\csname log2@92\endcsname{6524}
\expandafter\xdef\csname log2@93\endcsname{6539}
\expandafter\xdef\csname log2@94\endcsname{6555}
\expandafter\xdef\csname log2@95\endcsname{6570}
\expandafter\xdef\csname log2@96\endcsname{6585}
\expandafter\xdef\csname log2@97\endcsname{6600}
\expandafter\xdef\csname log2@98\endcsname{6615}
\expandafter\xdef\csname log2@99\endcsname{6629}
\expandafter\xdef\csname log2@100\endcsname{6644}
\expandafter\xdef\csname log2@101\endcsname{6658}
\expandafter\xdef\csname log2@102\endcsname{6672}
\expandafter\xdef\csname log2@103\endcsname{6687}
\expandafter\xdef\csname log2@104\endcsname{6700}
\expandafter\xdef\csname log2@105\endcsname{6714}
\expandafter\xdef\csname log2@106\endcsname{6728}
\expandafter\xdef\csname log2@107\endcsname{6741}
\expandafter\xdef\csname log2@108\endcsname{6755}
\expandafter\xdef\csname log2@109\endcsname{6768}
\expandafter\xdef\csname log2@110\endcsname{6781}
\expandafter\xdef\csname log2@111\endcsname{6794}
\expandafter\xdef\csname log2@112\endcsname{6807}
\expandafter\xdef\csname log2@113\endcsname{6820}
\expandafter\xdef\csname log2@114\endcsname{6833}
\expandafter\xdef\csname log2@115\endcsname{6845}
\expandafter\xdef\csname log2@116\endcsname{6858}
\expandafter\xdef\csname log2@117\endcsname{6870}
\expandafter\xdef\csname log2@118\endcsname{6883}
\expandafter\xdef\csname log2@119\endcsname{6895}
\expandafter\xdef\csname log2@120\endcsname{6907}
\expandafter\xdef\csname log2@121\endcsname{6919}
\expandafter\xdef\csname log2@122\endcsname{6931}
\expandafter\xdef\csname log2@123\endcsname{6943}
\expandafter\xdef\csname log2@124\endcsname{6954}
\expandafter\xdef\csname log2@125\endcsname{6966}
\expandafter\xdef\csname log2@126\endcsname{6977}
\expandafter\xdef\csname log2@127\endcsname{6989}
\expandafter\xdef\csname log2@128\endcsname{7000}
\expandafter\xdef\csname log2@129\endcsname{7011}
\expandafter\xdef\csname log2@130\endcsname{7022}
\expandafter\xdef\csname log2@131\endcsname{7033}
\expandafter\xdef\csname log2@132\endcsname{7044}
\expandafter\xdef\csname log2@133\endcsname{7055}
\expandafter\xdef\csname log2@134\endcsname{7066}
\expandafter\xdef\csname log2@135\endcsname{7077}
\expandafter\xdef\csname log2@136\endcsname{7088}
\expandafter\xdef\csname log2@137\endcsname{7098}
\expandafter\xdef\csname log2@138\endcsname{7108}
\expandafter\xdef\csname log2@139\endcsname{7119}
\expandafter\xdef\csname log2@140\endcsname{7129}
\expandafter\xdef\csname log2@141\endcsname{7140}
\expandafter\xdef\csname log2@142\endcsname{7150}
\expandafter\xdef\csname log2@143\endcsname{7160}
\expandafter\xdef\csname log2@144\endcsname{7170}
\expandafter\xdef\csname log2@145\endcsname{7180}
\expandafter\xdef\csname log2@146\endcsname{7190}
\expandafter\xdef\csname log2@147\endcsname{7200}
\expandafter\xdef\csname log2@148\endcsname{7209}
\expandafter\xdef\csname log2@149\endcsname{7219}
\expandafter\xdef\csname log2@150\endcsname{7229}
\expandafter\xdef\csname log2@151\endcsname{7238}
\expandafter\xdef\csname log2@152\endcsname{7248}
\expandafter\xdef\csname log2@153\endcsname{7257}
\expandafter\xdef\csname log2@154\endcsname{7267}
\expandafter\xdef\csname log2@155\endcsname{7276}
\expandafter\xdef\csname log2@156\endcsname{7285}
\expandafter\xdef\csname log2@157\endcsname{7295}
\expandafter\xdef\csname log2@158\endcsname{7304}
\expandafter\xdef\csname log2@159\endcsname{7313}
\expandafter\xdef\csname log2@160\endcsname{7323}
\expandafter\xdef\csname log2@161\endcsname{7331}
\expandafter\xdef\csname log2@162\endcsname{7340}
\expandafter\xdef\csname log2@163\endcsname{7349}
\expandafter\xdef\csname log2@164\endcsname{7358}
\expandafter\xdef\csname log2@165\endcsname{7366}
\expandafter\xdef\csname log2@166\endcsname{7375}
\expandafter\xdef\csname log2@167\endcsname{7374}
\expandafter\xdef\csname log2@168\endcsname{7392}
\expandafter\xdef\csname log2@169\endcsname{7401}
\expandafter\xdef\csname log2@170\endcsname{7409}
\expandafter\xdef\csname log2@171\endcsname{7418}
\expandafter\xdef\csname log2@172\endcsname{7463}
\expandafter\xdef\csname log2@173\endcsname{7435}
\expandafter\xdef\csname log2@174\endcsname{7443}
\expandafter\xdef\csname log2@175\endcsname{7451}
\expandafter\xdef\csname log2@176\endcsname{7460}
\expandafter\xdef\csname log2@177\endcsname{7468}
\expandafter\xdef\csname log2@178\endcsname{7476}
\expandafter\xdef\csname log2@179\endcsname{7484}
\expandafter\xdef\csname log2@180\endcsname{7492}
\expandafter\xdef\csname log2@181\endcsname{7500}
\expandafter\xdef\csname log2@182\endcsname{7508}
\expandafter\xdef\csname log2@183\endcsname{7516}
\expandafter\xdef\csname log2@184\endcsname{7524}
\expandafter\xdef\csname log2@185\endcsname{7531}
\expandafter\xdef\csname log2@186\endcsname{7539}
\expandafter\xdef\csname log2@187\endcsname{7547}
\expandafter\xdef\csname log2@188\endcsname{7555}
\expandafter\xdef\csname log2@189\endcsname{7562}
\expandafter\xdef\csname log2@190\endcsname{7570}
\expandafter\xdef\csname log2@191\endcsname{7577}
\expandafter\xdef\csname log2@192\endcsname{7585}
\expandafter\xdef\csname log2@193\endcsname{7592}
\expandafter\xdef\csname log2@194\endcsname{7600}
\expandafter\xdef\csname log2@195\endcsname{7607}
\expandafter\xdef\csname log2@196\endcsname{7615}
\expandafter\xdef\csname log2@197\endcsname{7622}
\expandafter\xdef\csname log2@198\endcsname{7629}
\expandafter\xdef\csname log2@199\endcsname{7637}
\expandafter\xdef\csname log2@200\endcsname{7644}
\expandafter\xdef\csname log2@201\endcsname{7651}
\expandafter\xdef\csname log2@202\endcsname{7658}
\expandafter\xdef\csname log2@203\endcsname{7665}
\expandafter\xdef\csname log2@204\endcsname{7672}
\expandafter\xdef\csname log2@205\endcsname{7679}
\expandafter\xdef\csname log2@206\endcsname{7687}
\expandafter\xdef\csname log2@207\endcsname{7693}
\expandafter\xdef\csname log2@208\endcsname{7700}
\expandafter\xdef\csname log2@209\endcsname{7707}
\expandafter\xdef\csname log2@210\endcsname{7714}
\expandafter\xdef\csname log2@211\endcsname{7721}
\expandafter\xdef\csname log2@212\endcsname{7728}
\expandafter\xdef\csname log2@213\endcsname{7735}
\expandafter\xdef\csname log2@214\endcsname{7741}
\expandafter\xdef\csname log2@215\endcsname{7748}
\expandafter\xdef\csname log2@216\endcsname{7755}
\expandafter\xdef\csname log2@217\endcsname{7761}
\expandafter\xdef\csname log2@218\endcsname{7768}
\expandafter\xdef\csname log2@219\endcsname{7775}
\expandafter\xdef\csname log2@220\endcsname{7781}


\expandafter\xdef\csname ch@r@65\endcsname{A}
\expandafter\xdef\csname ch@r@66\endcsname{B}
\expandafter\xdef\csname ch@r@67\endcsname{C}
\expandafter\xdef\csname ch@r@68\endcsname{D}
\expandafter\xdef\csname ch@r@69\endcsname{E}
\expandafter\xdef\csname ch@r@70\endcsname{F}
\expandafter\xdef\csname ch@r@71\endcsname{G}
\expandafter\xdef\csname ch@r@72\endcsname{H}
\expandafter\xdef\csname ch@r@73\endcsname{I}
\expandafter\xdef\csname ch@r@74\endcsname{J}
\expandafter\xdef\csname ch@r@75\endcsname{K}
\expandafter\xdef\csname ch@r@76\endcsname{L}
\expandafter\xdef\csname ch@r@77\endcsname{M}
\expandafter\xdef\csname ch@r@78\endcsname{N}
\expandafter\xdef\csname ch@r@79\endcsname{O}
\expandafter\xdef\csname ch@r@80\endcsname{P}
\expandafter\xdef\csname ch@r@81\endcsname{Q}
\expandafter\xdef\csname ch@r@82\endcsname{R}
\expandafter\xdef\csname ch@r@83\endcsname{S}
\expandafter\xdef\csname ch@r@84\endcsname{T}
\expandafter\xdef\csname ch@r@85\endcsname{U}
\expandafter\xdef\csname ch@r@86\endcsname{V}
\expandafter\xdef\csname ch@r@87\endcsname{W}
\expandafter\xdef\csname ch@r@88\endcsname{X}
\expandafter\xdef\csname ch@r@89\endcsname{Y}
\expandafter\xdef\csname ch@r@90\endcsname{Z}

\def\clear@sims{%
    \expandafter\xdef\csname \prfx simA\endcsname{(1)A}
    \expandafter\xdef\csname \prfx simB\endcsname{(1)B}
    \expandafter\xdef\csname \prfx simC\endcsname{(1)C}
    \expandafter\xdef\csname \prfx simD\endcsname{(1)D}
    \expandafter\xdef\csname \prfx simE\endcsname{(1)E}
    \expandafter\xdef\csname \prfx simF\endcsname{(1)F}
    \expandafter\xdef\csname \prfx simG\endcsname{(1)G}
    \expandafter\xdef\csname \prfx simH\endcsname{(1)H}
    \expandafter\xdef\csname \prfx simI\endcsname{(1)I}
    \expandafter\xdef\csname \prfx simJ\endcsname{(1)J}
    \expandafter\xdef\csname \prfx simK\endcsname{(1)K}
    \expandafter\xdef\csname \prfx simL\endcsname{(1)L}
    \expandafter\xdef\csname \prfx simM\endcsname{(1)M}
    \expandafter\xdef\csname \prfx simN\endcsname{(1)N}
    \expandafter\xdef\csname \prfx simO\endcsname{(1)O}
    \expandafter\xdef\csname \prfx simP\endcsname{(1)P}
    \expandafter\xdef\csname \prfx simQ\endcsname{(1)Q}
    \expandafter\xdef\csname \prfx simR\endcsname{(1)R}
    \expandafter\xdef\csname \prfx simS\endcsname{(1)S}
    \expandafter\xdef\csname \prfx simT\endcsname{(1)T}
    \expandafter\xdef\csname \prfx simU\endcsname{(1)U}
    \expandafter\xdef\csname \prfx simV\endcsname{(1)V}
    \expandafter\xdef\csname \prfx simW\endcsname{(1)W}
    \expandafter\xdef\csname \prfx simX\endcsname{(1)X}
    \expandafter\xdef\csname \prfx simY\endcsname{(1)Y}
    \expandafter\xdef\csname \prfx simZ\endcsname{(1)Z}
}

\def\clear@simpairs{%
\xdef\simpairCC{1} \xdef\simpairCS{0} \xdef\simpairCT{0} \xdef\simpairCP{0} \xdef\simpairCA{0}
\xdef\simpairCG{0} \xdef\simpairCN{0} \xdef\simpairCD{0} \xdef\simpairCE{0} \xdef\simpairCQ{0}
\xdef\simpairCH{0} \xdef\simpairCR{0} \xdef\simpairCK{0} \xdef\simpairCM{0} \xdef\simpairCI{0}
\xdef\simpairCL{0} \xdef\simpairCV{0} \xdef\simpairCF{0} \xdef\simpairCY{0} \xdef\simpairCW{0}
\xdef\simpairCB{0}  \xdef\simpairCJ{0}  \xdef\simpairCO{0}  \xdef\simpairCU{0}  \xdef\simpairCX{0}  \xdef\simpairCZ{0}

\xdef\simpairSC{0} \xdef\simpairSS{1} \xdef\simpairST{0} \xdef\simpairSP{0} \xdef\simpairSA{0}
\xdef\simpairSG{0} \xdef\simpairSN{0} \xdef\simpairSD{0} \xdef\simpairSE{0} \xdef\simpairSQ{0}
\xdef\simpairSH{0} \xdef\simpairSR{0} \xdef\simpairSK{0} \xdef\simpairSM{0} \xdef\simpairSI{0}
\xdef\simpairSL{0} \xdef\simpairSV{0} \xdef\simpairSF{0} \xdef\simpairSY{0} \xdef\simpairSW{0}
\xdef\simpairSB{0}  \xdef\simpairSJ{0}  \xdef\simpairSO{0}  \xdef\simpairSU{0}  \xdef\simpairSX{0}  \xdef\simpairSZ{0}

\xdef\simpairTC{0} \xdef\simpairTS{0} \xdef\simpairTT{1} \xdef\simpairTP{0} \xdef\simpairTA{0}
\xdef\simpairTG{0} \xdef\simpairTN{0} \xdef\simpairTD{0} \xdef\simpairTE{0} \xdef\simpairTQ{0}
\xdef\simpairTH{0} \xdef\simpairTR{0} \xdef\simpairTK{0} \xdef\simpairTM{0} \xdef\simpairTI{0}
\xdef\simpairTL{0} \xdef\simpairTV{0} \xdef\simpairTF{0} \xdef\simpairTY{0} \xdef\simpairTW{0}
\xdef\simpairTB{0}  \xdef\simpairTJ{0}  \xdef\simpairTO{0}  \xdef\simpairTU{0}  \xdef\simpairTX{0}  \xdef\simpairTZ{0}

\xdef\simpairPC{0} \xdef\simpairPS{0} \xdef\simpairPT{0} \xdef\simpairPP{1} \xdef\simpairPA{0}
\xdef\simpairPG{0} \xdef\simpairPN{0} \xdef\simpairPD{0} \xdef\simpairPE{0} \xdef\simpairPQ{0}
\xdef\simpairPH{0} \xdef\simpairPR{0} \xdef\simpairPK{0} \xdef\simpairPM{0} \xdef\simpairPI{0}
\xdef\simpairPL{0} \xdef\simpairPV{0} \xdef\simpairPF{0} \xdef\simpairPY{0} \xdef\simpairPW{0}
\xdef\simpairPB{0}  \xdef\simpairPJ{0}  \xdef\simpairPO{0}  \xdef\simpairPU{0}  \xdef\simpairPX{0}  \xdef\simpairPZ{0}

\xdef\simpairAC{0} \xdef\simpairAS{0} \xdef\simpairAT{0} \xdef\simpairAP{0} \xdef\simpairAA{1}
\xdef\simpairAG{0} \xdef\simpairAN{0} \xdef\simpairAD{0} \xdef\simpairAE{0} \xdef\simpairAQ{0}
\xdef\simpairAH{0} \xdef\simpairAR{0} \xdef\simpairAK{0} \xdef\simpairAM{0} \xdef\simpairAI{0}
\xdef\simpairAL{0} \xdef\simpairAV{0} \xdef\simpairAF{0} \xdef\simpairAY{0} \xdef\simpairAW{0} 
\xdef\simpairAB{0}  \xdef\simpairAJ{0}  \xdef\simpairAO{0}  \xdef\simpairAU{0}  \xdef\simpairAX{0}  \xdef\simpairAZ{0}

\xdef\simpairGC{0} \xdef\simpairGS{0} \xdef\simpairGT{0} \xdef\simpairGP{0} \xdef\simpairGA{0}
\xdef\simpairGG{1} \xdef\simpairGN{0} \xdef\simpairGD{0} \xdef\simpairGE{0} \xdef\simpairGQ{0}
\xdef\simpairGH{0} \xdef\simpairGR{0} \xdef\simpairGK{0} \xdef\simpairGM{0} \xdef\simpairGI{0}
\xdef\simpairGL{0} \xdef\simpairGV{0} \xdef\simpairGF{0} \xdef\simpairGY{0} \xdef\simpairGW{0}
\xdef\simpairGB{0}  \xdef\simpairGJ{0}  \xdef\simpairGO{0}  \xdef\simpairGU{0}  \xdef\simpairGX{0}  \xdef\simpairGZ{0}

\xdef\simpairNC{0} \xdef\simpairNS{0} \xdef\simpairNT{0} \xdef\simpairNP{0} \xdef\simpairNA{0}
\xdef\simpairNG{0} \xdef\simpairNN{1} \xdef\simpairND{0} \xdef\simpairNE{0} \xdef\simpairNQ{0}
\xdef\simpairNH{0} \xdef\simpairNR{0} \xdef\simpairNK{0} \xdef\simpairNM{0} \xdef\simpairNI{0}
\xdef\simpairNL{0} \xdef\simpairNV{0} \xdef\simpairNF{0} \xdef\simpairNY{0} \xdef\simpairNW{0}
\xdef\simpairNB{0}  \xdef\simpairNJ{0}  \xdef\simpairNO{0}  \xdef\simpairNU{0}  \xdef\simpairNX{0}  \xdef\simpairNZ{0}

\xdef\simpairDC{0} \xdef\simpairDS{0} \xdef\simpairDT{0} \xdef\simpairDP{0} \xdef\simpairDA{0}
\xdef\simpairDG{0} \xdef\simpairDN{0} \xdef\simpairDD{1} \xdef\simpairDE{0} \xdef\simpairDQ{0}
\xdef\simpairDH{0} \xdef\simpairDR{0} \xdef\simpairDK{0} \xdef\simpairDM{0} \xdef\simpairDI{0}
\xdef\simpairDL{0} \xdef\simpairDV{0} \xdef\simpairDF{0} \xdef\simpairDY{0} \xdef\simpairDW{0}
\xdef\simpairDB{0}  \xdef\simpairDJ{0}  \xdef\simpairDO{0}  \xdef\simpairDU{0}  \xdef\simpairDX{0}  \xdef\simpairDZ{0}

\xdef\simpairEC{0} \xdef\simpairES{0} \xdef\simpairET{0} \xdef\simpairEP{0} \xdef\simpairEA{0}
\xdef\simpairEG{0} \xdef\simpairEN{0} \xdef\simpairED{0} \xdef\simpairEE{1} \xdef\simpairEQ{0}
\xdef\simpairEH{0} \xdef\simpairER{0} \xdef\simpairEK{0} \xdef\simpairEM{0} \xdef\simpairEI{0}
\xdef\simpairEL{0} \xdef\simpairEV{0} \xdef\simpairEF{0} \xdef\simpairEY{0} \xdef\simpairEW{0}
\xdef\simpairEB{0}  \xdef\simpairEJ{0}  \xdef\simpairEO{0}  \xdef\simpairEU{0}  \xdef\simpairEX{0}  \xdef\simpairEZ{0}

\xdef\simpairQC{0} \xdef\simpairQS{0} \xdef\simpairQT{0} \xdef\simpairQP{0} \xdef\simpairQA{0}
\xdef\simpairQG{0} \xdef\simpairQN{0} \xdef\simpairQD{0} \xdef\simpairQE{0} \xdef\simpairQQ{1}
\xdef\simpairQH{0} \xdef\simpairQR{0} \xdef\simpairQK{0} \xdef\simpairQM{0} \xdef\simpairQI{0}
\xdef\simpairQL{0} \xdef\simpairQV{0} \xdef\simpairQF{0} \xdef\simpairQY{0} \xdef\simpairQW{0}
\xdef\simpairQB{0}  \xdef\simpairQJ{0}  \xdef\simpairQO{0}  \xdef\simpairQU{0}  \xdef\simpairQX{0}  \xdef\simpairQZ{0}

\xdef\simpairHC{0} \xdef\simpairHS{0} \xdef\simpairHT{0} \xdef\simpairHP{0} \xdef\simpairHA{0}
\xdef\simpairHG{0} \xdef\simpairHN{0} \xdef\simpairHD{0} \xdef\simpairHE{0} \xdef\simpairHQ{0}
\xdef\simpairHH{1} \xdef\simpairHR{0} \xdef\simpairHK{0} \xdef\simpairHM{0} \xdef\simpairHI{0}
\xdef\simpairHL{0} \xdef\simpairHV{0} \xdef\simpairHF{0} \xdef\simpairHY{0} \xdef\simpairHW{0}
\xdef\simpairHB{0}  \xdef\simpairHJ{0}  \xdef\simpairHO{0}  \xdef\simpairHU{0}  \xdef\simpairHX{0}  \xdef\simpairHZ{0}

\xdef\simpairRC{0} \xdef\simpairRS{0} \xdef\simpairRT{0} \xdef\simpairRP{0} \xdef\simpairRA{0}
\xdef\simpairRG{0} \xdef\simpairRN{0} \xdef\simpairRD{0} \xdef\simpairRE{0} \xdef\simpairRQ{0}
\xdef\simpairRH{0} \xdef\simpairRR{1} \xdef\simpairRK{0} \xdef\simpairRM{0} \xdef\simpairRI{0}
\xdef\simpairRL{0} \xdef\simpairRV{0} \xdef\simpairRF{0} \xdef\simpairRY{0} \xdef\simpairRW{0}
\xdef\simpairRB{0}  \xdef\simpairRJ{0}  \xdef\simpairRO{0}  \xdef\simpairRU{0}  \xdef\simpairRX{0}  \xdef\simpairRZ{0}

\xdef\simpairKC{0} \xdef\simpairKS{0} \xdef\simpairKT{0} \xdef\simpairKP{0} \xdef\simpairKA{0}
\xdef\simpairKG{0} \xdef\simpairKN{0} \xdef\simpairKD{0} \xdef\simpairKE{0} \xdef\simpairKQ{0}
\xdef\simpairKH{0} \xdef\simpairKR{0} \xdef\simpairKK{1} \xdef\simpairKM{0} \xdef\simpairKI{0}
\xdef\simpairKL{0} \xdef\simpairKV{0} \xdef\simpairKF{0} \xdef\simpairKY{0} \xdef\simpairKW{0}
\xdef\simpairKB{0}  \xdef\simpairKJ{0}  \xdef\simpairKO{0}  \xdef\simpairKU{0}  \xdef\simpairKX{0}  \xdef\simpairKZ{0}

\xdef\simpairMC{0} \xdef\simpairMS{0} \xdef\simpairMT{0} \xdef\simpairMP{0} \xdef\simpairMA{0}
\xdef\simpairMG{0} \xdef\simpairMN{0} \xdef\simpairMD{0} \xdef\simpairME{0} \xdef\simpairMQ{0}
\xdef\simpairMH{0} \xdef\simpairMR{0} \xdef\simpairMK{0} \xdef\simpairMM{1} \xdef\simpairMI{0}
\xdef\simpairML{0} \xdef\simpairMV{0} \xdef\simpairMF{0} \xdef\simpairMY{0} \xdef\simpairMW{0}
\xdef\simpairMB{0}  \xdef\simpairMJ{0}  \xdef\simpairMO{0}  \xdef\simpairMU{0}  \xdef\simpairMX{0}  \xdef\simpairMZ{0}

\xdef\simpairIC{0} \xdef\simpairIS{0} \xdef\simpairIT{0} \xdef\simpairIP{0} \xdef\simpairIA{0}
\xdef\simpairIG{0} \xdef\simpairIN{0} \xdef\simpairID{0} \xdef\simpairIE{0} \xdef\simpairIQ{0}
\xdef\simpairIH{0} \xdef\simpairIR{0} \xdef\simpairIK{0} \xdef\simpairIM{0} \xdef\simpairII{1}
\xdef\simpairIL{0} \xdef\simpairIV{0} \xdef\simpairIF{0} \xdef\simpairIY{0} \xdef\simpairIW{0}
\xdef\simpairIB{0}  \xdef\simpairIJ{0}  \xdef\simpairIO{0}  \xdef\simpairIU{0}  \xdef\simpairIX{0}  \xdef\simpairIZ{0}

\xdef\simpairLC{0} \xdef\simpairLS{0} \xdef\simpairLT{0} \xdef\simpairLP{0} \xdef\simpairLA{0}
\xdef\simpairLG{0} \xdef\simpairLN{0} \xdef\simpairLD{0} \xdef\simpairLE{0} \xdef\simpairLQ{0}
\xdef\simpairLH{0} \xdef\simpairLR{0} \xdef\simpairLK{0} \xdef\simpairLM{0} \xdef\simpairLI{0}
\xdef\simpairLL{1} \xdef\simpairLV{0} \xdef\simpairLF{0} \xdef\simpairLY{0} \xdef\simpairLW{0}
\xdef\simpairLB{0}  \xdef\simpairLJ{0}  \xdef\simpairLO{0}  \xdef\simpairLU{0}  \xdef\simpairLX{0}  \xdef\simpairLZ{0}

\xdef\simpairVC{0} \xdef\simpairVS{0} \xdef\simpairVT{0} \xdef\simpairVP{0} \xdef\simpairVA{0}
\xdef\simpairVG{0} \xdef\simpairVN{0} \xdef\simpairVD{0} \xdef\simpairVE{0} \xdef\simpairVQ{0}
\xdef\simpairVH{0} \xdef\simpairVR{0} \xdef\simpairVK{0} \xdef\simpairVM{0} \xdef\simpairVI{0}
\xdef\simpairVL{0} \xdef\simpairVV{1} \xdef\simpairVF{0} \xdef\simpairVY{0} \xdef\simpairVW{0}
\xdef\simpairVB{0}  \xdef\simpairVJ{0}  \xdef\simpairVO{0}  \xdef\simpairVU{0}  \xdef\simpairVX{0}  \xdef\simpairVZ{0}

\xdef\simpairFC{0} \xdef\simpairFS{0} \xdef\simpairFT{0} \xdef\simpairFP{0} \xdef\simpairFA{0}
\xdef\simpairFG{0} \xdef\simpairFN{0} \xdef\simpairFD{0} \xdef\simpairFE{0} \xdef\simpairFQ{0}
\xdef\simpairFH{0} \xdef\simpairFR{0} \xdef\simpairFK{0} \xdef\simpairFM{0} \xdef\simpairFI{0}
\xdef\simpairFL{0} \xdef\simpairFV{0} \xdef\simpairFF{1} \xdef\simpairFY{0} \xdef\simpairFW{0}
\xdef\simpairFB{0}  \xdef\simpairFJ{0}  \xdef\simpairFO{0}  \xdef\simpairFU{0}  \xdef\simpairFX{0}  \xdef\simpairFZ{0}

\xdef\simpairYC{0} \xdef\simpairYS{0} \xdef\simpairYT{0} \xdef\simpairYP{0} \xdef\simpairYA{0}
\xdef\simpairYG{0} \xdef\simpairYN{0} \xdef\simpairYD{0} \xdef\simpairYE{0} \xdef\simpairYQ{0}
\xdef\simpairYH{0} \xdef\simpairYR{0} \xdef\simpairYK{0} \xdef\simpairYM{0} \xdef\simpairYI{0}
\xdef\simpairYL{0} \xdef\simpairYV{0} \xdef\simpairYF{0} \xdef\simpairYY{1} \xdef\simpairYW{0}
\xdef\simpairYB{0}  \xdef\simpairYJ{0}  \xdef\simpairYO{0}  \xdef\simpairYU{0}  \xdef\simpairYX{0}  \xdef\simpairYZ{0}

\xdef\simpairWC{0} \xdef\simpairWS{0} \xdef\simpairWT{0} \xdef\simpairWP{0} \xdef\simpairWA{0}
\xdef\simpairWG{0} \xdef\simpairWN{0} \xdef\simpairWD{0} \xdef\simpairWE{0} \xdef\simpairWQ{0}
\xdef\simpairWH{0} \xdef\simpairWR{0} \xdef\simpairWK{0} \xdef\simpairWM{0} \xdef\simpairWI{0}
\xdef\simpairWL{0} \xdef\simpairWV{0} \xdef\simpairWF{0} \xdef\simpairWY{0} \xdef\simpairWW{1}
\xdef\simpairWB{0}  \xdef\simpairWJ{0}  \xdef\simpairWO{0}  \xdef\simpairWU{0}  \xdef\simpairWX{0}  \xdef\simpairWZ{0}

\xdef\simpairBC{0} \xdef\simpairBS{0} \xdef\simpairBT{0} \xdef\simpairBP{0} \xdef\simpairBA{0}
\xdef\simpairBG{0} \xdef\simpairBN{0} \xdef\simpairBD{0} \xdef\simpairBE{0} \xdef\simpairBQ{0}
\xdef\simpairBH{0} \xdef\simpairBR{0} \xdef\simpairBK{0} \xdef\simpairBM{0} \xdef\simpairBI{0}
\xdef\simpairBL{0} \xdef\simpairBV{0} \xdef\simpairBF{0} \xdef\simpairBY{0} \xdef\simpairBW{0}
\xdef\simpairBB{1} \xdef\simpairBJ{0} \xdef\simpairBO{0} \xdef\simpairBU{0}  \xdef\simpairBX{0}  \xdef\simpairBZ{0}

\xdef\simpairJC{0} \xdef\simpairJS{0} \xdef\simpairJT{0} \xdef\simpairJP{0} \xdef\simpairJA{0}
\xdef\simpairJG{0} \xdef\simpairJN{0} \xdef\simpairJD{0} \xdef\simpairJE{0} \xdef\simpairJQ{0}
\xdef\simpairJH{0} \xdef\simpairJR{0} \xdef\simpairJK{0} \xdef\simpairJM{0} \xdef\simpairJI{0}
\xdef\simpairJL{0} \xdef\simpairJV{0} \xdef\simpairJF{0} \xdef\simpairJY{0} \xdef\simpairJW{0}
\xdef\simpairJB{0} \xdef\simpairJJ{1} \xdef\simpairJO{0} \xdef\simpairJU{0}  \xdef\simpairJX{0}  \xdef\simpairJZ{0}

\xdef\simpairOC{0} \xdef\simpairOS{0} \xdef\simpairOT{0} \xdef\simpairOP{0} \xdef\simpairOA{0}
\xdef\simpairOG{0} \xdef\simpairON{0} \xdef\simpairOD{0} \xdef\simpairOE{0} \xdef\simpairOQ{0}
\xdef\simpairOH{0} \xdef\simpairOR{0} \xdef\simpairOK{0} \xdef\simpairOM{0} \xdef\simpairOI{0}
\xdef\simpairOL{0} \xdef\simpairOV{0} \xdef\simpairOF{0} \xdef\simpairOY{0} \xdef\simpairOW{0}
\xdef\simpairOB{0} \xdef\simpairOJ{0} \xdef\simpairOO{1} \xdef\simpairOU{0}  \xdef\simpairOX{0}  \xdef\simpairOZ{0}

\xdef\simpairUC{0} \xdef\simpairUS{0} \xdef\simpairUT{0} \xdef\simpairUP{0} \xdef\simpairUA{0}
\xdef\simpairUG{0} \xdef\simpairUN{0} \xdef\simpairUD{0} \xdef\simpairUE{0} \xdef\simpairUQ{0}
\xdef\simpairUH{0} \xdef\simpairUR{0} \xdef\simpairUK{0} \xdef\simpairUM{0} \xdef\simpairUI{0}
\xdef\simpairUL{0} \xdef\simpairUV{0} \xdef\simpairUF{0} \xdef\simpairUY{0} \xdef\simpairUW{0}
\xdef\simpairUB{0} \xdef\simpairUJ{0} \xdef\simpairUO{0} \xdef\simpairUU{1}  \xdef\simpairUX{0}  \xdef\simpairUZ{0}

\xdef\simpairXC{0} \xdef\simpairXS{0} \xdef\simpairXT{0} \xdef\simpairXP{0} \xdef\simpairXA{0}
\xdef\simpairXG{0} \xdef\simpairXN{0} \xdef\simpairXD{0} \xdef\simpairXE{0} \xdef\simpairXQ{0}
\xdef\simpairXH{0} \xdef\simpairXR{0} \xdef\simpairXK{0} \xdef\simpairXM{0} \xdef\simpairXI{0}
\xdef\simpairXL{0} \xdef\simpairXV{0} \xdef\simpairXF{0} \xdef\simpairXY{0} \xdef\simpairXW{0}
\xdef\simpairXB{0} \xdef\simpairXJ{0} \xdef\simpairXO{0} \xdef\simpairXU{0}  \xdef\simpairXX{1}  \xdef\simpairXZ{0}

\xdef\simpairZC{0} \xdef\simpairZS{0} \xdef\simpairZT{0} \xdef\simpairZP{0} \xdef\simpairZA{0}
\xdef\simpairZG{0} \xdef\simpairZN{0} \xdef\simpairZD{0} \xdef\simpairZE{0} \xdef\simpairZQ{0}
\xdef\simpairZH{0} \xdef\simpairZR{0} \xdef\simpairZK{0} \xdef\simpairZM{0} \xdef\simpairZI{0}
\xdef\simpairZL{0} \xdef\simpairZV{0} \xdef\simpairZF{0} \xdef\simpairZY{0} \xdef\simpairZW{0}
\xdef\simpairZB{0} \xdef\simpairZJ{0} \xdef\simpairZO{0} \xdef\simpairZU{0}  \xdef\simpairZX{0}  \xdef\simpairZZ{1}

}

\def\clear@sim@count{
   \def\clear@sim@nums{%
      \advance\innerloopcount by 1
      \ifnum\innerloopcount>\seq@count
      \else
         \expandafter\xdef\csname poscount\the\outerloopcount @\the\innerloopcount\endcsname{0}
         \expandafter\xdef\csname identcount\the\outerloopcount @\the\innerloopcount\endcsname{0}
         \expandafter\xdef\csname simcount\the\outerloopcount @\the\innerloopcount\endcsname{0}
         \clear@sim@nums
      \fi
   }

   \outerloopcount=1
   \loop
      \innerloopcount=\outerloopcount
      \clear@sim@nums
      \advance\outerloopcount by 1
   \ifnum\outerloopcount<\seq@count \repeat
}

\xdef\pepmwA{711}  \xdef\pepmwB{1146} \xdef\pepmwC{1032}
\xdef\pepmwD{1151} \xdef\pepmwE{1291} \xdef\pepmwF{1472}
\xdef\pepmwG{571}  \xdef\pepmwH{1372} \xdef\pepmwI{1132}
\xdef\pepmwJ{0}    \xdef\pepmwK{1282} \xdef\pepmwL{1132}
\xdef\pepmwM{1312} \xdef\pepmwN{1141} \xdef\pepmwO{0}
\xdef\pepmwP{971}  \xdef\pepmwQ{1281} \xdef\pepmwR{1562}
\xdef\pepmwS{871}  \xdef\pepmwT{1011} \xdef\pepmwU{0}
\xdef\pepmwV{991}  \xdef\pepmwW{1862} \xdef\pepmwX{1282}
\xdef\pepmwY{1632} \xdef\pepmwZ{1286}

\xdef\DNAmwA{3462}  \xdef\DNAmwB{0}     \xdef\DNAmwC{3222}
\xdef\DNAmwD{0}     \xdef\DNAmwE{0}     \xdef\DNAmwF{0}
\xdef\DNAmwG{3622}  \xdef\DNAmwH{0}     \xdef\DNAmwI{0}
\xdef\DNAmwJ{0}     \xdef\DNAmwK{0}     \xdef\DNAmwL{0}
\xdef\DNAmwM{0}     \xdef\DNAmwN{0}     \xdef\DNAmwO{0}
\xdef\DNAmwP{0}     \xdef\DNAmwQ{0}     \xdef\DNAmwR{0}
\xdef\DNAmwS{0}     \xdef\DNAmwT{3372}  \xdef\DNAmwU{3232}
\xdef\DNAmwV{0}     \xdef\DNAmwW{0}     \xdef\DNAmwX{0}
\xdef\DNAmwY{0}     \xdef\DNAmwZ{0}

\xdef\pepchargeA{0}     \xdef\pepchargeB{0}     \xdef\pepchargeC{-30}
\xdef\pepchargeD{-1000} \xdef\pepchargeE{-1000} \xdef\pepchargeF{0}
\xdef\pepchargeG{0}     \xdef\pepchargeH{165}   \xdef\pepchargeI{0}
\xdef\pepchargeJ{0}     \xdef\pepchargeK{1000}  \xdef\pepchargeL{0}
\xdef\pepchargeM{0}     \xdef\pepchargeN{0}     \xdef\pepchargeO{0}
\xdef\pepchargeP{0}     \xdef\pepchargeQ{0}     \xdef\pepchargeR{1000}
\xdef\pepchargeS{0}     \xdef\pepchargeT{0}     \xdef\pepchargeU{0}
\xdef\pepchargeV{0}     \xdef\pepchargeW{0}     \xdef\pepchargeX{0}
\xdef\pepchargeY{0}     \xdef\pepchargeZ{0}
\xdef\chargeNterm{910}  \xdef\chargeCterm{-1000}

\xdef\chargeA{0}     \xdef\chargeB{0}     \xdef\chargeC{0}
\xdef\chargeD{-50}   \xdef\chargeE{-50}   \xdef\chargeF{0}
\xdef\chargeG{0}     \xdef\chargeH{30}    \xdef\chargeI{0}
\xdef\chargeJ{0}     \xdef\chargeK{50}    \xdef\chargeL{0}
\xdef\chargeM{0}     \xdef\chargeN{0}     \xdef\chargeO{0}
\xdef\chargeP{0}     \xdef\chargeQ{0}     \xdef\chargeR{50}
\xdef\chargeS{0}     \xdef\chargeT{0}     \xdef\chargeU{0}
\xdef\chargeV{0}     \xdef\chargeW{0}     \xdef\chargeX{0}
\xdef\chargeY{0}     \xdef\chargeZ{0}

\xdef\molwA{11}   \xdef\molwB{45}   \xdef\molwC{36}
\xdef\molwD{45}   \xdef\molwE{66}   \xdef\molwF{70}
\xdef\molwG{1}    \xdef\molwH{62}   \xdef\molwI{44}
\xdef\molwJ{N}    \xdef\molwK{55}   \xdef\molwL{44}
\xdef\molwM{58}   \xdef\molwN{44}   \xdef\molwO{N}
\xdef\molwP{31}   \xdef\molwQ{55}   \xdef\molwR{77}
\xdef\molwS{19}   \xdef\molwT{34}   \xdef\molwU{N}
\xdef\molwV{33}   \xdef\molwW{100}  \xdef\molwX{55}
\xdef\molwY{82}   \xdef\molwZ{66}

\xdef\HydroA{21}  \xdef\HydroB{N}   \xdef\HydroC{10}
\xdef\HydroD{-31} \xdef\HydroE{-25} \xdef\HydroF{41}
\xdef\HydroG{16}  \xdef\HydroH{-14} \xdef\HydroI{47}
\xdef\HydroJ{N}   \xdef\HydroK{-52} \xdef\HydroL{36}
\xdef\HydroM{22}  \xdef\HydroN{-27} \xdef\HydroO{N}
\xdef\HydroP{4}   \xdef\HydroQ{-29} \xdef\HydroR{-53}
\xdef\HydroS{-6}  \xdef\HydroT{-2}  \xdef\HydroU{N}
\xdef\HydroV{37}  \xdef\HydroW{28}  \xdef\HydroX{N}
\xdef\HydroY{9}   \xdef\HydroZ{N}

\xdef\hardnessA{98}  \xdef\hardnessB{0}   \xdef\hardnessC{37}
\xdef\hardnessD{100} \xdef\hardnessE{96} \xdef\hardnessF{48}
\xdef\hardnessG{92}  \xdef\hardnessH{50} \xdef\hardnessI{98}
\xdef\hardnessJ{86}   \xdef\hardnessK{73} \xdef\hardnessL{98}
\xdef\hardnessM{25}  \xdef\hardnessN{86} \xdef\hardnessO{0}
\xdef\hardnessP{91}   \xdef\hardnessQ{95} \xdef\hardnessR{69}
\xdef\hardnessS{94}  \xdef\hardnessT{98}  \xdef\hardnessU{-34}
\xdef\hardnessV{98}  \xdef\hardnessW{0}  \xdef\hardnessX{0}
\xdef\hardnessY{28}   \xdef\hardnessZ{0}



\xdef\consCB{0}  \xdef\consCJ{0}  \xdef\consCO{0}  \xdef\consCU{0}  \xdef\consCX{0}  \xdef\consCZ{0}
\expandafter\xdef\csname consC?\endcsname{0}

\xdef\consSB{0}  \xdef\consSJ{0}  \xdef\consSO{0}  \xdef\consSU{0}  \xdef\consSX{0}  \xdef\consSZ{0}
\expandafter\xdef\csname consS?\endcsname{0}

\xdef\consTB{0}  \xdef\consTJ{0}  \xdef\consTO{0}  \xdef\consTU{0}  \xdef\consTX{0}  \xdef\consTZ{0}
\expandafter\xdef\csname consT?\endcsname{0}

\xdef\consPB{0}  \xdef\consPJ{0}  \xdef\consPO{0}  \xdef\consPU{0}  \xdef\consPX{0}  \xdef\consPZ{0}
\expandafter\xdef\csname consP?\endcsname{0}

\xdef\consAB{0}  \xdef\consAJ{0}  \xdef\consAO{0}  \xdef\consAU{0}  \xdef\consAX{0}  \xdef\consAZ{0}
\expandafter\xdef\csname consA?\endcsname{0}

\xdef\consGB{0}  \xdef\consGJ{0}  \xdef\consGO{0}  \xdef\consGU{0}  \xdef\consGX{0}  \xdef\consGZ{0}
\expandafter\xdef\csname consG?\endcsname{0}

\xdef\consNB{0}  \xdef\consNJ{0}  \xdef\consNO{0}  \xdef\consNU{0}  \xdef\consNX{0}  \xdef\consNZ{0}
\expandafter\xdef\csname consN?\endcsname{0}

\xdef\consDB{0}  \xdef\consDJ{0}  \xdef\consDO{0}  \xdef\consDU{0}  \xdef\consDX{0}  \xdef\consDZ{0}
\expandafter\xdef\csname consD?\endcsname{0}

\xdef\consEB{0}  \xdef\consEJ{0}  \xdef\consEO{0}  \xdef\consEU{0}  \xdef\consEX{0}  \xdef\consEZ{0}
\expandafter\xdef\csname consE?\endcsname{0}

\xdef\consQB{0}  \xdef\consQJ{0}  \xdef\consQO{0}  \xdef\consQU{0}  \xdef\consQX{0}  \xdef\consQZ{0}
\expandafter\xdef\csname consQ?\endcsname{0}

\xdef\consHB{0}  \xdef\consHJ{0}  \xdef\consHO{0}  \xdef\consHU{0}  \xdef\consHX{0}  \xdef\consHZ{0}
\expandafter\xdef\csname consH?\endcsname{0}

\xdef\consRB{0}  \xdef\consRJ{0}  \xdef\consRO{0}  \xdef\consRU{0}  \xdef\consRX{0}  \xdef\consRZ{0}
\expandafter\xdef\csname consR?\endcsname{0}

\xdef\consKB{0}  \xdef\consKJ{0}  \xdef\consKO{0}  \xdef\consKU{0}  \xdef\consKX{0}  \xdef\consKZ{0}
\expandafter\xdef\csname consK?\endcsname{0}

\xdef\consMB{0}  \xdef\consMJ{0}  \xdef\consMO{0}  \xdef\consMU{0}  \xdef\consMX{0}  \xdef\consMZ{0}
\expandafter\xdef\csname consM?\endcsname{0}

\xdef\consIB{0}  \xdef\consIJ{0}  \xdef\consIO{0}  \xdef\consIU{0}  \xdef\consIX{0}  \xdef\consIZ{0}
\expandafter\xdef\csname consI?\endcsname{0}

\xdef\consLB{0}  \xdef\consLJ{0}  \xdef\consLO{0}  \xdef\consLU{0}  \xdef\consLX{0}  \xdef\consLZ{0}
\expandafter\xdef\csname consL?\endcsname{0}

\xdef\consVB{0}  \xdef\consVJ{0}  \xdef\consVO{0}  \xdef\consVU{0}  \xdef\consVX{0}  \xdef\consVZ{0}
\expandafter\xdef\csname consV?\endcsname{0}

\xdef\consFB{0}  \xdef\consFJ{0}  \xdef\consFO{0}  \xdef\consFU{0}  \xdef\consFX{0}  \xdef\consFZ{0}
\expandafter\xdef\csname consF?\endcsname{0}

\xdef\consYB{0}  \xdef\consYJ{0}  \xdef\consYO{0}  \xdef\consYU{0}  \xdef\consYX{0}  \xdef\consYZ{0}
\expandafter\xdef\csname consY?\endcsname{0}

\xdef\consWB{0}  \xdef\consWJ{0}  \xdef\consWO{0}  \xdef\consWU{0}  \xdef\consWX{0}  \xdef\consWZ{0}
\expandafter\xdef\csname consW?\endcsname{0}

\xdef\consBC{0} \xdef\consBS{0} \xdef\consBT{0} \xdef\consBP{0} \xdef\consBA{0}
\xdef\consBG{0} \xdef\consBN{0} \xdef\consBD{0} \xdef\consBE{0} \xdef\consBQ{0}
\xdef\consBH{0} \xdef\consBR{0} \xdef\consBK{0} \xdef\consBM{0} \xdef\consBI{0}
\xdef\consBL{0} \xdef\consBV{0} \xdef\consBF{0} \xdef\consBY{0} \xdef\consBW{0}
\xdef\consBB{10} \xdef\consBJ{0} \xdef\consBO{0} \xdef\consBU{0}  \xdef\consBX{0}  \xdef\consBZ{0}
\expandafter\xdef\csname consB.\endcsname{0}
\expandafter\xdef\csname consB?\endcsname{0}

\xdef\consJC{0} \xdef\consJS{0} \xdef\consJT{0} \xdef\consJP{0} \xdef\consJA{0}
\xdef\consJG{0} \xdef\consJN{0} \xdef\consJD{0} \xdef\consJE{0} \xdef\consJQ{0}
\xdef\consJH{0} \xdef\consJR{0} \xdef\consJK{0} \xdef\consJM{0} \xdef\consJI{0}
\xdef\consJL{0} \xdef\consJV{0} \xdef\consJF{0} \xdef\consJY{0} \xdef\consJW{0}
\xdef\consJB{0} \xdef\consJJ{10} \xdef\consJO{0} \xdef\consJU{0}  \xdef\consJX{0}  \xdef\consJZ{0}
\expandafter\xdef\csname consJ.\endcsname{0}
\expandafter\xdef\csname consJ?\endcsname{0}

\xdef\consOC{0} \xdef\consOS{0} \xdef\consOT{0} \xdef\consOP{0} \xdef\consOA{0}
\xdef\consOG{0} \xdef\consON{0} \xdef\consOD{0} \xdef\consOE{0} \xdef\consOQ{0}
\xdef\consOH{0} \xdef\consOR{0} \xdef\consOK{0} \xdef\consOM{0} \xdef\consOI{0}
\xdef\consOL{0} \xdef\consOV{0} \xdef\consOF{0} \xdef\consOY{0} \xdef\consOW{0}
\xdef\consOB{0} \xdef\consOJ{0} \xdef\consOO{10} \xdef\consOU{0}  \xdef\consOX{0}  \xdef\consOZ{0}
\expandafter\xdef\csname consO.\endcsname{0}
\expandafter\xdef\csname consO?\endcsname{0}

\xdef\consUC{0} \xdef\consUS{0} \xdef\consUT{0} \xdef\consUP{0} \xdef\consUA{0}
\xdef\consUG{0} \xdef\consUN{0} \xdef\consUD{0} \xdef\consUE{0} \xdef\consUQ{0}
\xdef\consUH{0} \xdef\consUR{0} \xdef\consUK{0} \xdef\consUM{0} \xdef\consUI{0}
\xdef\consUL{0} \xdef\consUV{0} \xdef\consUF{0} \xdef\consUY{0} \xdef\consUW{0}
\xdef\consUB{0} \xdef\consUJ{0} \xdef\consUO{0} \xdef\consUU{10}  \xdef\consUX{0}  \xdef\consUZ{0}
\expandafter\xdef\csname consU.\endcsname{0}
\expandafter\xdef\csname consU?\endcsname{0}

\xdef\consXC{0} \xdef\consXS{0} \xdef\consXT{0} \xdef\consXP{0} \xdef\consXA{0}
\xdef\consXG{0} \xdef\consXN{0} \xdef\consXD{0} \xdef\consXE{0} \xdef\consXQ{0}
\xdef\consXH{0} \xdef\consXR{0} \xdef\consXK{0} \xdef\consXM{0} \xdef\consXI{0}
\xdef\consXL{0} \xdef\consXV{0} \xdef\consXF{0} \xdef\consXY{0} \xdef\consXW{0}
\xdef\consXB{0} \xdef\consXJ{0} \xdef\consXO{0} \xdef\consXU{0}  \xdef\consXX{10}  \xdef\consXZ{0}
\expandafter\xdef\csname consX.\endcsname{0}
\expandafter\xdef\csname consX?\endcsname{0}

\xdef\consZC{0} \xdef\consZS{0} \xdef\consZT{0} \xdef\consZP{0} \xdef\consZA{0}
\xdef\consZG{0} \xdef\consZN{0} \xdef\consZD{0} \xdef\consZE{0} \xdef\consZQ{0}
\xdef\consZH{0} \xdef\consZR{0} \xdef\consZK{0} \xdef\consZM{0} \xdef\consZI{0}
\xdef\consZL{0} \xdef\consZV{0} \xdef\consZF{0} \xdef\consZY{0} \xdef\consZW{0}
\xdef\consZB{0} \xdef\consZJ{0} \xdef\consZO{0} \xdef\consZU{0}  \xdef\consZX{0}  \xdef\consZZ{10}
\expandafter\xdef\csname consZ.\endcsname{0}
\expandafter\xdef\csname consZ?\endcsname{0}

\expandafter\xdef\csname cons..\endcsname{0}
\expandafter\xdef\csname cons.?\endcsname{0}

\expandafter\xdef\csname cons?C\endcsname{0}
\expandafter\xdef\csname cons?S\endcsname{0}
\expandafter\xdef\csname cons?T\endcsname{0}
\expandafter\xdef\csname cons?P\endcsname{0}
\expandafter\xdef\csname cons?A\endcsname{0}
\expandafter\xdef\csname cons?G\endcsname{0}
\expandafter\xdef\csname cons?N\endcsname{0}
\expandafter\xdef\csname cons?D\endcsname{0}
\expandafter\xdef\csname cons?E\endcsname{0}
\expandafter\xdef\csname cons?Q\endcsname{0}
\expandafter\xdef\csname cons?H\endcsname{0}
\expandafter\xdef\csname cons?R\endcsname{0}
\expandafter\xdef\csname cons?K\endcsname{0}
\expandafter\xdef\csname cons?M\endcsname{0}
\expandafter\xdef\csname cons?I\endcsname{0}
\expandafter\xdef\csname cons?L\endcsname{0}
\expandafter\xdef\csname cons?V\endcsname{0}
\expandafter\xdef\csname cons?F\endcsname{0}
\expandafter\xdef\csname cons?Y\endcsname{0}
\expandafter\xdef\csname cons?W\endcsname{0}
\expandafter\xdef\csname cons?B\endcsname{0}
\expandafter\xdef\csname cons?J\endcsname{0}
\expandafter\xdef\csname cons?O\endcsname{0}
\expandafter\xdef\csname cons?U\endcsname{0}
\expandafter\xdef\csname cons?X\endcsname{0}
\expandafter\xdef\csname cons?Z\endcsname{0}
\expandafter\xdef\csname cons?.\endcsname{0}
\expandafter\xdef\csname cons??\endcsname{0}

\expandafter\xdef\csname cons==\endcsname{0}


\def\gappenalty#1{%
\expandafter\xdef\csname cons.C\endcsname{#1} \expandafter\xdef\csname consC.\endcsname{#1}
\expandafter\xdef\csname cons.S\endcsname{#1} \expandafter\xdef\csname consS.\endcsname{#1}
\expandafter\xdef\csname cons.T\endcsname{#1} \expandafter\xdef\csname consT.\endcsname{#1}
\expandafter\xdef\csname cons.P\endcsname{#1} \expandafter\xdef\csname consP.\endcsname{#1}
\expandafter\xdef\csname cons.A\endcsname{#1} \expandafter\xdef\csname consA.\endcsname{#1}
\expandafter\xdef\csname cons.G\endcsname{#1} \expandafter\xdef\csname consG.\endcsname{#1}
\expandafter\xdef\csname cons.N\endcsname{#1} \expandafter\xdef\csname consN.\endcsname{#1}
\expandafter\xdef\csname cons.D\endcsname{#1} \expandafter\xdef\csname consD.\endcsname{#1}
\expandafter\xdef\csname cons.E\endcsname{#1} \expandafter\xdef\csname consE.\endcsname{#1}
\expandafter\xdef\csname cons.Q\endcsname{#1} \expandafter\xdef\csname consQ.\endcsname{#1}
\expandafter\xdef\csname cons.H\endcsname{#1} \expandafter\xdef\csname consH.\endcsname{#1}
\expandafter\xdef\csname cons.R\endcsname{#1} \expandafter\xdef\csname consR.\endcsname{#1}
\expandafter\xdef\csname cons.K\endcsname{#1} \expandafter\xdef\csname consK.\endcsname{#1}
\expandafter\xdef\csname cons.M\endcsname{#1} \expandafter\xdef\csname consM.\endcsname{#1}
\expandafter\xdef\csname cons.I\endcsname{#1} \expandafter\xdef\csname consI.\endcsname{#1}
\expandafter\xdef\csname cons.L\endcsname{#1} \expandafter\xdef\csname consL.\endcsname{#1}
\expandafter\xdef\csname cons.V\endcsname{#1} \expandafter\xdef\csname consV.\endcsname{#1}
\expandafter\xdef\csname cons.F\endcsname{#1} \expandafter\xdef\csname consF.\endcsname{#1}
\expandafter\xdef\csname cons.Y\endcsname{#1} \expandafter\xdef\csname consY.\endcsname{#1}
\expandafter\xdef\csname cons.W\endcsname{#1} \expandafter\xdef\csname consW.\endcsname{#1}
\expandafter\xdef\csname cons.B\endcsname{#1} \expandafter\xdef\csname consB.\endcsname{#1}
\expandafter\xdef\csname cons.J\endcsname{#1} \expandafter\xdef\csname consJ.\endcsname{#1}
\expandafter\xdef\csname cons.O\endcsname{#1} \expandafter\xdef\csname consO.\endcsname{#1}
\expandafter\xdef\csname cons.U\endcsname{#1} \expandafter\xdef\csname consU.\endcsname{#1}
\expandafter\xdef\csname cons.X\endcsname{#1} \expandafter\xdef\csname consX.\endcsname{#1}
\expandafter\xdef\csname cons.Z\endcsname{#1} \expandafter\xdef\csname consZ.\endcsname{#1}
}

\def\weighttable#1{%
\xdef\first@{#1}

\xdef\temp@{identity}
\ifx\first@\temp@


%%%% identity matrix

\xdef\consCC{10} \xdef\consCS{0} \xdef\consCT{0} \xdef\consCP{0} \xdef\consCA{0}
\xdef\consCG{0} \xdef\consCN{0} \xdef\consCD{0} \xdef\consCE{0} \xdef\consCQ{0}
\xdef\consCH{0} \xdef\consCR{0} \xdef\consCK{0} \xdef\consCM{0} \xdef\consCI{0}
\xdef\consCL{0} \xdef\consCV{0} \xdef\consCF{0} \xdef\consCY{0} \xdef\consCW{0}

\xdef\consSC{0} \xdef\consSS{10} \xdef\consST{0} \xdef\consSP{0} \xdef\consSA{0}
\xdef\consSG{0} \xdef\consSN{0} \xdef\consSD{0} \xdef\consSE{0} \xdef\consSQ{0}
\xdef\consSH{0} \xdef\consSR{0} \xdef\consSK{0} \xdef\consSM{0} \xdef\consSI{0}
\xdef\consSL{0} \xdef\consSV{0} \xdef\consSF{0} \xdef\consSY{0} \xdef\consSW{0}

\xdef\consTC{0} \xdef\consTS{0} \xdef\consTT{10} \xdef\consTP{0} \xdef\consTA{0}
\xdef\consTG{0} \xdef\consTN{0} \xdef\consTD{0} \xdef\consTE{0} \xdef\consTQ{0}
\xdef\consTH{0} \xdef\consTR{0} \xdef\consTK{0} \xdef\consTM{0} \xdef\consTI{0}
\xdef\consTL{0} \xdef\consTV{0} \xdef\consTF{0} \xdef\consTY{0} \xdef\consTW{0}

\xdef\consPC{0} \xdef\consPS{0} \xdef\consPT{0} \xdef\consPP{10} \xdef\consPA{0}
\xdef\consPG{0} \xdef\consPN{0} \xdef\consPD{0} \xdef\consPE{0} \xdef\consPQ{0}
\xdef\consPH{0} \xdef\consPR{0} \xdef\consPK{0} \xdef\consPM{0} \xdef\consPI{0}
\xdef\consPL{0} \xdef\consPV{0} \xdef\consPF{0} \xdef\consPY{0} \xdef\consPW{0}

\xdef\consAC{0} \xdef\consAS{0} \xdef\consAT{0} \xdef\consAP{0} \xdef\consAA{10}
\xdef\consAG{0} \xdef\consAN{0} \xdef\consAD{0} \xdef\consAE{0} \xdef\consAQ{0}
\xdef\consAH{0} \xdef\consAR{0} \xdef\consAK{0} \xdef\consAM{0} \xdef\consAI{0}
\xdef\consAL{0} \xdef\consAV{0} \xdef\consAF{0} \xdef\consAY{0} \xdef\consAW{0} 

\xdef\consGC{0} \xdef\consGS{0} \xdef\consGT{0} \xdef\consGP{0} \xdef\consGA{0}
\xdef\consGG{10} \xdef\consGN{0} \xdef\consGD{0} \xdef\consGE{0} \xdef\consGQ{0}
\xdef\consGH{0} \xdef\consGR{0} \xdef\consGK{0} \xdef\consGM{0} \xdef\consGI{0}
\xdef\consGL{0} \xdef\consGV{0} \xdef\consGF{0} \xdef\consGY{0} \xdef\consGW{0}

\xdef\consNC{0} \xdef\consNS{0} \xdef\consNT{0} \xdef\consNP{0} \xdef\consNA{0}
\xdef\consNG{0} \xdef\consNN{10} \xdef\consND{0} \xdef\consNE{0} \xdef\consNQ{0}
\xdef\consNH{0} \xdef\consNR{0} \xdef\consNK{0} \xdef\consNM{0} \xdef\consNI{0}
\xdef\consNL{0} \xdef\consNV{0} \xdef\consNF{0} \xdef\consNY{0} \xdef\consNW{0}

\xdef\consDC{0} \xdef\consDS{0} \xdef\consDT{0} \xdef\consDP{0} \xdef\consDA{0}
\xdef\consDG{0} \xdef\consDN{0} \xdef\consDD{10} \xdef\consDE{0} \xdef\consDQ{0}
\xdef\consDH{0} \xdef\consDR{0} \xdef\consDK{0} \xdef\consDM{0} \xdef\consDI{0}
\xdef\consDL{0} \xdef\consDV{0} \xdef\consDF{0} \xdef\consDY{0} \xdef\consDW{0}

\xdef\consEC{0} \xdef\consES{0} \xdef\consET{0} \xdef\consEP{0} \xdef\consEA{0}
\xdef\consEG{0} \xdef\consEN{0} \xdef\consED{0} \xdef\consEE{10} \xdef\consEQ{0}
\xdef\consEH{0} \xdef\consER{0} \xdef\consEK{0} \xdef\consEM{0} \xdef\consEI{0}
\xdef\consEL{0} \xdef\consEV{0} \xdef\consEF{0} \xdef\consEY{0} \xdef\consEW{0}

\xdef\consQC{0} \xdef\consQS{0} \xdef\consQT{0} \xdef\consQP{0} \xdef\consQA{0}
\xdef\consQG{0} \xdef\consQN{0} \xdef\consQD{0} \xdef\consQE{0} \xdef\consQQ{10}
\xdef\consQH{0} \xdef\consQR{0} \xdef\consQK{0} \xdef\consQM{0} \xdef\consQI{0}
\xdef\consQL{0} \xdef\consQV{0} \xdef\consQF{0} \xdef\consQY{0} \xdef\consQW{0}

\xdef\consHC{0} \xdef\consHS{0} \xdef\consHT{0} \xdef\consHP{0} \xdef\consHA{0}
\xdef\consHG{0} \xdef\consHN{0} \xdef\consHD{0} \xdef\consHE{0} \xdef\consHQ{0}
\xdef\consHH{10} \xdef\consHR{0} \xdef\consHK{0} \xdef\consHM{0} \xdef\consHI{0}
\xdef\consHL{0} \xdef\consHV{0} \xdef\consHF{0} \xdef\consHY{0} \xdef\consHW{0}

\xdef\consRC{0} \xdef\consRS{0} \xdef\consRT{0} \xdef\consRP{0} \xdef\consRA{0}
\xdef\consRG{0} \xdef\consRN{0} \xdef\consRD{0} \xdef\consRE{0} \xdef\consRQ{0}
\xdef\consRH{0} \xdef\consRR{10} \xdef\consRK{0} \xdef\consRM{0} \xdef\consRI{0}
\xdef\consRL{0} \xdef\consRV{0} \xdef\consRF{0} \xdef\consRY{0} \xdef\consRW{0}

\xdef\consKC{0} \xdef\consKS{0} \xdef\consKT{0} \xdef\consKP{0} \xdef\consKA{0}
\xdef\consKG{0} \xdef\consKN{0} \xdef\consKD{0} \xdef\consKE{0} \xdef\consKQ{0}
\xdef\consKH{0} \xdef\consKR{0} \xdef\consKK{10} \xdef\consKM{0} \xdef\consKI{0}
\xdef\consKL{0} \xdef\consKV{0} \xdef\consKF{0} \xdef\consKY{0} \xdef\consKW{0}

\xdef\consMC{0} \xdef\consMS{0} \xdef\consMT{0} \xdef\consMP{0} \xdef\consMA{0}
\xdef\consMG{0} \xdef\consMN{0} \xdef\consMD{0} \xdef\consME{0} \xdef\consMQ{0}
\xdef\consMH{0} \xdef\consMR{0} \xdef\consMK{0} \xdef\consMM{10} \xdef\consMI{0}
\xdef\consML{0} \xdef\consMV{0} \xdef\consMF{0} \xdef\consMY{0} \xdef\consMW{0}

\xdef\consIC{0} \xdef\consIS{0} \xdef\consIT{0} \xdef\consIP{0} \xdef\consIA{0}
\xdef\consIG{0} \xdef\consIN{0} \xdef\consID{0} \xdef\consIE{0} \xdef\consIQ{0}
\xdef\consIH{0} \xdef\consIR{0} \xdef\consIK{0} \xdef\consIM{0} \xdef\consII{10}
\xdef\consIL{0} \xdef\consIV{0} \xdef\consIF{0} \xdef\consIY{0} \xdef\consIW{0}

\xdef\consLC{0} \xdef\consLS{0} \xdef\consLT{0} \xdef\consLP{0} \xdef\consLA{0}
\xdef\consLG{0} \xdef\consLN{0} \xdef\consLD{0} \xdef\consLE{0} \xdef\consLQ{0}
\xdef\consLH{0} \xdef\consLR{0} \xdef\consLK{0} \xdef\consLM{0} \xdef\consLI{0}
\xdef\consLL{10} \xdef\consLV{0} \xdef\consLF{0} \xdef\consLY{0} \xdef\consLW{0}

\xdef\consVC{0} \xdef\consVS{0} \xdef\consVT{0} \xdef\consVP{0} \xdef\consVA{0}
\xdef\consVG{0} \xdef\consVN{0} \xdef\consVD{0} \xdef\consVE{0} \xdef\consVQ{0}
\xdef\consVH{0} \xdef\consVR{0} \xdef\consVK{0} \xdef\consVM{0} \xdef\consVI{0}
\xdef\consVL{0} \xdef\consVV{10} \xdef\consVF{0} \xdef\consVY{0} \xdef\consVW{0}

\xdef\consFC{0} \xdef\consFS{0} \xdef\consFT{0} \xdef\consFP{0} \xdef\consFA{0}
\xdef\consFG{0} \xdef\consFN{0} \xdef\consFD{0} \xdef\consFE{0} \xdef\consFQ{0}
\xdef\consFH{0} \xdef\consFR{0} \xdef\consFK{0} \xdef\consFM{0} \xdef\consFI{0}
\xdef\consFL{0} \xdef\consFV{0} \xdef\consFF{10} \xdef\consFY{0} \xdef\consFW{0}

\xdef\consYC{0} \xdef\consYS{0} \xdef\consYT{0} \xdef\consYP{0} \xdef\consYA{0}
\xdef\consYG{0} \xdef\consYN{0} \xdef\consYD{0} \xdef\consYE{0} \xdef\consYQ{0}
\xdef\consYH{0} \xdef\consYR{0} \xdef\consYK{0} \xdef\consYM{0} \xdef\consYI{0}
\xdef\consYL{0} \xdef\consYV{0} \xdef\consYF{0} \xdef\consYY{10} \xdef\consYW{0}

\xdef\consWC{0} \xdef\consWS{0} \xdef\consWT{0} \xdef\consWP{0} \xdef\consWA{0}
\xdef\consWG{0} \xdef\consWN{0} \xdef\consWD{0} \xdef\consWE{0} \xdef\consWQ{0}
\xdef\consWH{0} \xdef\consWR{0} \xdef\consWK{0} \xdef\consWM{0} \xdef\consWI{0}
\xdef\consWL{0} \xdef\consWV{0} \xdef\consWF{0} \xdef\consWY{0} \xdef\consWW{10}

\gappenalty{0} \xdef\m@trixf@ctor{10}

\else

\xdef\temp@{structural}
\ifx\first@\temp@


%%%% structural

\xdef\consCC{10} \xdef\consCS{6} \xdef\consCT{3} \xdef\consCP{3} \xdef\consCA{3}
\xdef\consCG{5} \xdef\consCN{3} \xdef\consCD{1} \xdef\consCE{0} \xdef\consCQ{1}
\xdef\consCH{3} \xdef\consCR{3} \xdef\consCK{0} \xdef\consCM{3} \xdef\consCI{3}
\xdef\consCL{3} \xdef\consCV{3} \xdef\consCF{5} \xdef\consCY{5} \xdef\consCW{5}

\xdef\consSC{6} \xdef\consSS{10} \xdef\consST{8} \xdef\consSP{6} \xdef\consSA{8}
\xdef\consSG{8} \xdef\consSN{8} \xdef\consSD{6} \xdef\consSE{5} \xdef\consSQ{5}
\xdef\consSH{5} \xdef\consSR{5} \xdef\consSK{5} \xdef\consSM{3} \xdef\consSI{3}
\xdef\consSL{3} \xdef\consSV{6} \xdef\consSF{5} \xdef\consSY{5} \xdef\consSW{3}

\xdef\consTC{3} \xdef\consTS{8} \xdef\consTT{10} \xdef\consTP{6} \xdef\consTA{8}
\xdef\consTG{6} \xdef\consTN{6} \xdef\consTD{5} \xdef\consTE{5} \xdef\consTQ{5}
\xdef\consTH{3} \xdef\consTR{5} \xdef\consTK{6} \xdef\consTM{5} \xdef\consTI{5}
\xdef\consTL{3} \xdef\consTV{6} \xdef\consTF{3} \xdef\consTY{3} \xdef\consTW{1}

\xdef\consPC{3} \xdef\consPS{6} \xdef\consPT{6} \xdef\consPP{10} \xdef\consPA{8}
\xdef\consPG{6} \xdef\consPN{3} \xdef\consPD{5} \xdef\consPE{5} \xdef\consPQ{5}
\xdef\consPH{5} \xdef\consPR{5} \xdef\consPK{3} \xdef\consPM{3} \xdef\consPI{3}
\xdef\consPL{5} \xdef\consPV{6} \xdef\consPF{5} \xdef\consPY{3} \xdef\consPW{3}

\xdef\consAC{3} \xdef\consAS{8} \xdef\consAT{8} \xdef\consAP{8} \xdef\consAA{10}
\xdef\consAG{8} \xdef\consAN{5} \xdef\consAD{6} \xdef\consAE{6} \xdef\consAQ{5}
\xdef\consAH{3} \xdef\consAR{3} \xdef\consAK{5} \xdef\consAM{5} \xdef\consAI{3}
\xdef\consAL{3} \xdef\consAV{8} \xdef\consAF{5} \xdef\consAY{3} \xdef\consAW{3} 

\xdef\consGC{5} \xdef\consGS{8} \xdef\consGT{6} \xdef\consGP{6} \xdef\consGA{8}
\xdef\consGG{10} \xdef\consGN{5} \xdef\consGD{6} \xdef\consGE{6} \xdef\consGQ{3}
\xdef\consGH{1} \xdef\consGR{5} \xdef\consGK{3} \xdef\consGM{1} \xdef\consGI{3}
\xdef\consGL{3} \xdef\consGV{6} \xdef\consGF{3} \xdef\consGY{3} \xdef\consGW{5}

\xdef\consNC{3} \xdef\consNS{8} \xdef\consNT{6} \xdef\consNP{3} \xdef\consNA{5}
\xdef\consNG{5} \xdef\consNN{10} \xdef\consND{8} \xdef\consNE{6} \xdef\consNQ{5}
\xdef\consNH{6} \xdef\consNR{5} \xdef\consNK{6} \xdef\consNM{1} \xdef\consNI{3}
\xdef\consNL{1} \xdef\consNV{3} \xdef\consNF{3} \xdef\consNY{5} \xdef\consNW{0}

\xdef\consDC{1} \xdef\consDS{6} \xdef\consDT{5} \xdef\consDP{5} \xdef\consDA{6}
\xdef\consDG{6} \xdef\consDN{8} \xdef\consDD{10} \xdef\consDE{8} \xdef\consDQ{6}
\xdef\consDH{5} \xdef\consDR{3} \xdef\consDK{5} \xdef\consDM{3} \xdef\consDI{1}
\xdef\consDL{1} \xdef\consDV{5} \xdef\consDF{1} \xdef\consDY{3} \xdef\consDW{0}

\xdef\consEC{0} \xdef\consES{5} \xdef\consET{5} \xdef\consEP{5} \xdef\consEA{6}
\xdef\consEG{6} \xdef\consEN{6} \xdef\consED{8} \xdef\consEE{10} \xdef\consEQ{6}
\xdef\consEH{3} \xdef\consER{5} \xdef\consEK{6} \xdef\consEM{3} \xdef\consEI{1}
\xdef\consEL{1} \xdef\consEV{6} \xdef\consEF{3} \xdef\consEY{1} \xdef\consEW{1}

\xdef\consQC{1} \xdef\consQS{5} \xdef\consQT{5} \xdef\consQP{5} \xdef\consQA{5}
\xdef\consQG{3} \xdef\consQN{5} \xdef\consQD{6} \xdef\consQE{6} \xdef\consQQ{10}
\xdef\consQH{6} \xdef\consQR{5} \xdef\consQK{6} \xdef\consQM{3} \xdef\consQI{1}
\xdef\consQL{3} \xdef\consQV{3} \xdef\consQF{1} \xdef\consQY{3} \xdef\consQW{1}

\xdef\consHC{3} \xdef\consHS{5} \xdef\consHT{3} \xdef\consHP{5} \xdef\consHA{3}
\xdef\consHG{1} \xdef\consHN{6} \xdef\consHD{5} \xdef\consHE{3} \xdef\consHQ{6}
\xdef\consHH{10} \xdef\consHR{6} \xdef\consHK{5} \xdef\consHM{3} \xdef\consHI{3}
\xdef\consHL{5} \xdef\consHV{1} \xdef\consHF{3} \xdef\consHY{5} \xdef\consHW{1}

\xdef\consRC{3} \xdef\consRS{5} \xdef\consRT{5} \xdef\consRP{5} \xdef\consRA{3}
\xdef\consRG{5} \xdef\consRN{5} \xdef\consRD{3} \xdef\consRE{5} \xdef\consRQ{5}
\xdef\consRH{6} \xdef\consRR{10} \xdef\consRK{8} \xdef\consRM{3} \xdef\consRI{3}
\xdef\consRL{3} \xdef\consRV{3} \xdef\consRF{1} \xdef\consRY{1} \xdef\consRW{3}

\xdef\consKC{0} \xdef\consKS{5} \xdef\consKT{6} \xdef\consKP{3} \xdef\consKA{5}
\xdef\consKG{3} \xdef\consKN{6} \xdef\consKD{5} \xdef\consKE{6} \xdef\consKQ{6}
\xdef\consKH{5} \xdef\consKR{8} \xdef\consKK{10} \xdef\consKM{3} \xdef\consKI{3}
\xdef\consKL{3} \xdef\consKV{5} \xdef\consKF{1} \xdef\consKY{1} \xdef\consKW{1}

\xdef\consMC{3} \xdef\consMS{5} \xdef\consMT{5} \xdef\consMP{3} \xdef\consMA{5}
\xdef\consMG{1} \xdef\consMN{1} \xdef\consMD{3} \xdef\consME{3} \xdef\consMQ{3}
\xdef\consMH{3} \xdef\consMR{3} \xdef\consMK{3} \xdef\consMM{10} \xdef\consMI{6}
\xdef\consML{8} \xdef\consMV{6} \xdef\consMF{5} \xdef\consMY{3} \xdef\consMW{5}

\xdef\consIC{3} \xdef\consIS{3} \xdef\consIT{5} \xdef\consIP{3} \xdef\consIA{3}
\xdef\consIG{3} \xdef\consIN{3} \xdef\consID{1} \xdef\consIE{1} \xdef\consIQ{1}
\xdef\consIH{3} \xdef\consIR{3} \xdef\consIK{3} \xdef\consIM{6} \xdef\consII{10}
\xdef\consIL{8} \xdef\consIV{3} \xdef\consIF{6} \xdef\consIY{5} \xdef\consIW{5}

\xdef\consLC{3} \xdef\consLS{3} \xdef\consLT{3} \xdef\consLP{5} \xdef\consLA{3}
\xdef\consLG{3} \xdef\consLN{1} \xdef\consLD{1} \xdef\consLE{1} \xdef\consLQ{3}
\xdef\consLH{5} \xdef\consLR{3} \xdef\consLK{3} \xdef\consLM{8} \xdef\consLI{8}
\xdef\consLL{10} \xdef\consLV{3} \xdef\consLF{6} \xdef\consLY{5} \xdef\consLW{6}

\xdef\consVC{3} \xdef\consVS{6} \xdef\consVT{6} \xdef\consVP{6} \xdef\consVA{8}
\xdef\consVG{6} \xdef\consVN{3} \xdef\consVD{5} \xdef\consVE{6} \xdef\consVQ{3}
\xdef\consVH{1} \xdef\consVR{3} \xdef\consVK{5} \xdef\consVM{6} \xdef\consVI{3}
\xdef\consVL{3} \xdef\consVV{10} \xdef\consVF{6} \xdef\consVY{5} \xdef\consVW{5}

\xdef\consFC{5} \xdef\consFS{5} \xdef\consFT{3} \xdef\consFP{5} \xdef\consFA{5}
\xdef\consFG{3} \xdef\consFN{3} \xdef\consFD{1} \xdef\consFE{3} \xdef\consFQ{1}
\xdef\consFH{3} \xdef\consFR{1} \xdef\consFK{1} \xdef\consFM{5} \xdef\consFI{6}
\xdef\consFL{6} \xdef\consFV{6} \xdef\consFF{10} \xdef\consFY{8} \xdef\consFW{5}

\xdef\consYC{5} \xdef\consYS{5} \xdef\consYT{3} \xdef\consYP{3} \xdef\consYA{3}
\xdef\consYG{3} \xdef\consYN{5} \xdef\consYD{3} \xdef\consYE{1} \xdef\consYQ{3}
\xdef\consYH{5} \xdef\consYR{1} \xdef\consYK{1} \xdef\consYM{3} \xdef\consYI{5}
\xdef\consYL{5} \xdef\consYV{5} \xdef\consYF{8} \xdef\consYY{10} \xdef\consYW{5}

\xdef\consWC{5} \xdef\consWS{3} \xdef\consWT{1} \xdef\consWP{3} \xdef\consWA{3}
\xdef\consWG{5} \xdef\consWN{0} \xdef\consWD{0} \xdef\consWE{1} \xdef\consWQ{1}
\xdef\consWH{1} \xdef\consWR{3} \xdef\consWK{1} \xdef\consWM{5} \xdef\consWI{5}
\xdef\consWL{6} \xdef\consWV{5} \xdef\consWF{5} \xdef\consWY{5} \xdef\consWW{10}

\gappenalty{0} \xdef\m@trixf@ctor{10}

\else

\xdef\temp@{PAM250}
\ifx\first@\temp@

%%%% PAM250

\xdef\consCC{4} \xdef\consCS{0} \xdef\consCT{-2} \xdef\consCP{-3} \xdef\consCA{-2}
\xdef\consCG{-3} \xdef\consCN{-4} \xdef\consCD{2} \xdef\consCE{-5} \xdef\consCQ{-5}
\xdef\consCH{-3} \xdef\consCR{-4} \xdef\consCK{-5} \xdef\consCM{-5} \xdef\consCI{-2}
\xdef\consCL{-6} \xdef\consCV{-2} \xdef\consCF{-4} \xdef\consCY{0} \xdef\consCW{-8}
\xdef\consCB{-4} \xdef\consCZ{-5} \xdef\consCX{-3}

\xdef\consSC{0} \xdef\consSS{3} \xdef\consST{1} \xdef\consSP{1} \xdef\consSA{1}
\xdef\consSG{1} \xdef\consSN{1} \xdef\consSD{0} \xdef\consSE{0} \xdef\consSQ{-1}
\xdef\consSH{-1} \xdef\consSR{0} \xdef\consSK{0} \xdef\consSM{-2} \xdef\consSI{-1}
\xdef\consSL{-3} \xdef\consSV{-1} \xdef\consSF{-3} \xdef\consSY{-3} \xdef\consSW{-2}
\xdef\consSB{0} \xdef\consSZ{0} \xdef\consSX{0}

\xdef\consTC{-2} \xdef\consTS{1} \xdef\consTT{3} \xdef\consTP{0} \xdef\consTA{1}
\xdef\consTG{0} \xdef\consTN{0} \xdef\consTD{0} \xdef\consTE{0} \xdef\consTQ{-1}
\xdef\consTH{-1} \xdef\consTR{-1} \xdef\consTK{0} \xdef\consTM{-1} \xdef\consTI{0}
\xdef\consTL{-2} \xdef\consTV{0} \xdef\consTF{-2} \xdef\consTY{-3} \xdef\consTW{-5}
\xdef\consTB{0} \xdef\consTZ{-1} \xdef\consTX{0}

\xdef\consPC{-3} \xdef\consPS{1} \xdef\consPT{0} \xdef\consPP{6} \xdef\consPA{1}
\xdef\consPG{-1} \xdef\consPN{-1} \xdef\consPD{-1} \xdef\consPE{-1} \xdef\consPQ{0}
\xdef\consPH{0} \xdef\consPR{0} \xdef\consPK{-1} \xdef\consPM{-2} \xdef\consPI{-2}
\xdef\consPL{-3} \xdef\consPV{-1} \xdef\consPF{-5} \xdef\consPY{-5} \xdef\consPW{-6}
\xdef\consPB{-1} \xdef\consPZ{0} \xdef\consPX{-1}

\xdef\consAC{-2} \xdef\consAS{1} \xdef\consAT{1} \xdef\consAP{1} \xdef\consAA{2}
\xdef\consAG{1} \xdef\consAN{0} \xdef\consAD{0} \xdef\consAE{0} \xdef\consAQ{0}
\xdef\consAH{-1} \xdef\consAR{-2} \xdef\consAK{-1} \xdef\consAM{-1} \xdef\consAI{-1}
\xdef\consAL{-2} \xdef\consAV{0} \xdef\consAF{-4} \xdef\consAY{-3} \xdef\consAW{-6} 
\xdef\consAB{0} \xdef\consAZ{0} \xdef\consAX{0}

\xdef\consGC{-3} \xdef\consGS{1} \xdef\consGT{0} \xdef\consGP{-1} \xdef\consGA{1}
\xdef\consGG{5} \xdef\consGN{0} \xdef\consGD{1} \xdef\consGE{0} \xdef\consGQ{-1}
\xdef\consGH{-2} \xdef\consGR{-3} \xdef\consGK{-2} \xdef\consGM{-3} \xdef\consGI{-3}
\xdef\consGL{-4} \xdef\consGV{-1} \xdef\consGF{-5} \xdef\consGY{-5} \xdef\consGW{-7}
\xdef\consGB{0} \xdef\consGZ{0} \xdef\consGX{-1}

\xdef\consNC{-4} \xdef\consNS{1} \xdef\consNT{0} \xdef\consNP{-1} \xdef\consNA{0}
\xdef\consNG{0} \xdef\consNN{2} \xdef\consND{2} \xdef\consNE{1} \xdef\consNQ{1}
\xdef\consNH{2} \xdef\consNR{0} \xdef\consNK{1} \xdef\consNM{-2} \xdef\consNI{-2}
\xdef\consNL{-3} \xdef\consNV{-2} \xdef\consNF{-4} \xdef\consNY{-2} \xdef\consNW{-4}
\xdef\consNB{2} \xdef\consNZ{1} \xdef\consNX{0}

\xdef\consDC{-5} \xdef\consDS{0} \xdef\consDT{0} \xdef\consDP{-1} \xdef\consDA{0}
\xdef\consDG{1} \xdef\consDN{2} \xdef\consDD{4} \xdef\consDE{3} \xdef\consDQ{2}
\xdef\consDH{1} \xdef\consDR{-1} \xdef\consDK{0} \xdef\consDM{-3} \xdef\consDI{-2}
\xdef\consDL{-4} \xdef\consDV{-2} \xdef\consDF{-6} \xdef\consDY{-4} \xdef\consDW{-7}
\xdef\consDB{3} \xdef\consDZ{3} \xdef\consDX{-1}

\xdef\consEC{-5} \xdef\consES{0} \xdef\consET{0} \xdef\consEP{-1} \xdef\consEA{0}
\xdef\consEG{0} \xdef\consEN{1} \xdef\consED{3} \xdef\consEE{4} \xdef\consEQ{2}
\xdef\consEH{1} \xdef\consER{-1} \xdef\consEK{0} \xdef\consEM{-2} \xdef\consEI{-2}
\xdef\consEL{-3} \xdef\consEV{-2} \xdef\consEF{-5} \xdef\consEY{-4} \xdef\consEW{-7}
\xdef\consEB{3} \xdef\consEZ{3} \xdef\consEX{-1}

\xdef\consQC{-5} \xdef\consQS{-1} \xdef\consQT{-1} \xdef\consQP{0} \xdef\consQA{0}
\xdef\consQG{-1} \xdef\consQN{1} \xdef\consQD{2} \xdef\consQE{2} \xdef\consQQ{4}
\xdef\consQH{3} \xdef\consQR{1} \xdef\consQK{1} \xdef\consQM{-1} \xdef\consQI{-2}
\xdef\consQL{-2} \xdef\consQV{-2} \xdef\consQF{-5} \xdef\consQY{-4} \xdef\consQW{-5}
\xdef\consQB{1} \xdef\consQZ{3} \xdef\consQX{-1}

\xdef\consHC{-3} \xdef\consHS{-1} \xdef\consHT{-1} \xdef\consHP{0} \xdef\consHA{-1}
\xdef\consHG{-2} \xdef\consHN{2} \xdef\consHD{1} \xdef\consHE{1} \xdef\consHQ{3}
\xdef\consHH{6} \xdef\consHR{2} \xdef\consHK{0} \xdef\consHM{-2} \xdef\consHI{-2}
\xdef\consHL{-2} \xdef\consHV{-2} \xdef\consHF{-2} \xdef\consHY{-5} \xdef\consHW{-7}
\xdef\consHB{1} \xdef\consHZ{2} \xdef\consHX{-1}

\xdef\consRC{-4} \xdef\consRS{0} \xdef\consRT{-1} \xdef\consRP{0} \xdef\consRA{-2}
\xdef\consRG{-3} \xdef\consRN{0} \xdef\consRD{-1} \xdef\consRE{-1} \xdef\consRQ{1}
\xdef\consRH{2} \xdef\consRR{6} \xdef\consRK{3} \xdef\consRM{0} \xdef\consRI{-2}
\xdef\consRL{-3} \xdef\consRV{-2} \xdef\consRF{-4} \xdef\consRY{-4} \xdef\consRW{2}
\xdef\consRB{-1} \xdef\consRZ{0} \xdef\consRX{-1}

\xdef\consKC{-5} \xdef\consKS{0} \xdef\consKT{0} \xdef\consKP{-1} \xdef\consKA{-1}
\xdef\consKG{-2} \xdef\consKN{1} \xdef\consKD{0} \xdef\consKE{0} \xdef\consKQ{1}
\xdef\consKH{0} \xdef\consKR{3} \xdef\consKK{5} \xdef\consKM{0} \xdef\consKI{-2}
\xdef\consKL{-3} \xdef\consKV{-2} \xdef\consKF{-5} \xdef\consKY{-4} \xdef\consKW{-3}
\xdef\consKB{1} \xdef\consKZ{-2} \xdef\consKX{-1}

\xdef\consMC{-5} \xdef\consMS{-2} \xdef\consMT{-1} \xdef\consMP{-2} \xdef\consMA{-1}
\xdef\consMG{-3} \xdef\consMN{-2} \xdef\consMD{-3} \xdef\consME{-2} \xdef\consMQ{-1}
\xdef\consMH{-2} \xdef\consMR{0} \xdef\consMK{0} \xdef\consMM{6} \xdef\consMI{2}
\xdef\consML{4} \xdef\consMV{2} \xdef\consMF{0} \xdef\consMY{-2} \xdef\consMW{-4}
\xdef\consMB{-2} \xdef\consMZ{-2} \xdef\consMX{-1}

\xdef\consIC{-2} \xdef\consIS{-1} \xdef\consIT{0} \xdef\consIP{-2} \xdef\consIA{-1}
\xdef\consIG{-3} \xdef\consIN{-2} \xdef\consID{-2} \xdef\consIE{-2} \xdef\consIQ{-2}
\xdef\consIH{-2} \xdef\consIR{-2} \xdef\consIK{-2} \xdef\consIM{2} \xdef\consII{5}
\xdef\consIL{2} \xdef\consIV{4} \xdef\consIF{1} \xdef\consIY{-1} \xdef\consIW{-5}
\xdef\consIB{-2} \xdef\consIZ{-2} \xdef\consIX{-1}

\xdef\consLC{-6} \xdef\consLS{-3} \xdef\consLT{-2} \xdef\consLP{-3} \xdef\consLA{-2}
\xdef\consLG{-4} \xdef\consLN{-3} \xdef\consLD{-4} \xdef\consLE{-3} \xdef\consLQ{-2}
\xdef\consLH{-2} \xdef\consLR{-3} \xdef\consLK{-3} \xdef\consLM{4} \xdef\consLI{2}
\xdef\consLL{6} \xdef\consLV{2} \xdef\consLF{2} \xdef\consLY{-1} \xdef\consLW{-2}
\xdef\consLB{-3} \xdef\consLZ{-3} \xdef\consLX{-1}

\xdef\consVC{-2} \xdef\consVS{-1} \xdef\consVT{0} \xdef\consVP{-1} \xdef\consVA{0}
\xdef\consVG{-1} \xdef\consVN{-2} \xdef\consVD{-2} \xdef\consVE{-2} \xdef\consVQ{-2}
\xdef\consVH{-2} \xdef\consVR{-2} \xdef\consVK{-2} \xdef\consVM{2} \xdef\consVI{4}
\xdef\consVL{2} \xdef\consVV{4} \xdef\consVF{-1} \xdef\consVY{-2} \xdef\consVW{-6}
\xdef\consVB{-2} \xdef\consVZ{-2} \xdef\consVX{-1}

\xdef\consFC{-4} \xdef\consFS{-3} \xdef\consFT{-2} \xdef\consFP{-5} \xdef\consFA{-4}
\xdef\consFG{-5} \xdef\consFN{-4} \xdef\consFD{-6} \xdef\consFE{-5} \xdef\consFQ{-5}
\xdef\consFH{-2} \xdef\consFR{-4} \xdef\consFK{-5} \xdef\consFM{0} \xdef\consFI{1}
\xdef\consFL{2} \xdef\consFV{-1} \xdef\consFF{9} \xdef\consFY{7} \xdef\consFW{0}
\xdef\consFB{-4} \xdef\consFZ{-5} \xdef\consFX{-2}

\xdef\consYC{0} \xdef\consYS{-3} \xdef\consYT{-3} \xdef\consYP{-5} \xdef\consYA{-3}
\xdef\consYG{-5} \xdef\consYN{-2} \xdef\consYD{-4} \xdef\consYE{-4} \xdef\consYQ{-4}
\xdef\consYH{0} \xdef\consYR{-4} \xdef\consYK{-4} \xdef\consYM{-2} \xdef\consYI{-1}
\xdef\consYL{-1} \xdef\consYV{-2} \xdef\consYF{7} \xdef\consYY{10} \xdef\consYW{0}
\xdef\consYB{-3} \xdef\consYZ{-4} \xdef\consYX{-2}

\xdef\consWC{-8} \xdef\consWS{-2} \xdef\consWT{-5} \xdef\consWP{-6} \xdef\consWA{-6}
\xdef\consWG{-7} \xdef\consWN{-4} \xdef\consWD{-7} \xdef\consWE{-7} \xdef\consWQ{-5}
\xdef\consWH{-3} \xdef\consWR{2} \xdef\consWK{-3} \xdef\consWM{-4} \xdef\consWI{-5}
\xdef\consWL{-2} \xdef\consWV{-6} \xdef\consWF{0} \xdef\consWY{0} \xdef\consWW{17}
\xdef\consWB{-5} \xdef\consWZ{-6} \xdef\consWX{-4}

\xdef\consBC{-4} \xdef\consBS{0} \xdef\consBT{0} \xdef\consBP{-1} \xdef\consBA{0}
\xdef\consBG{0} \xdef\consBN{2} \xdef\consBD{3} \xdef\consBE{3} \xdef\consBQ{1}
\xdef\consBH{1} \xdef\consBR{-1} \xdef\consBK{1} \xdef\consBM{-2} \xdef\consBI{-2}
\xdef\consBL{-3} \xdef\consBV{-2} \xdef\consBF{-4} \xdef\consBY{-3} \xdef\consBW{-5}
\xdef\consBB{3} \xdef\consBZ{2}  \xdef\consBX{-1}

\xdef\consZC{-5} \xdef\consZS{0} \xdef\consZT{-1} \xdef\consZP{0} \xdef\consZA{0}
\xdef\consZG{0} \xdef\consZN{1} \xdef\consZD{3} \xdef\consZE{3} \xdef\consZQ{3}
\xdef\consZH{2} \xdef\consZR{0} \xdef\consZK{0} \xdef\consZM{-2} \xdef\consZI{-2}
\xdef\consZL{-3} \xdef\consZV{-2} \xdef\consZF{-5} \xdef\consZY{-4} \xdef\consZW{-6}
\xdef\consZB{2} \xdef\consZZ{3}  \xdef\consZX{-1}

\xdef\consXC{-3} \xdef\consXS{0} \xdef\consXT{0} \xdef\consXP{-1} \xdef\consXA{0}
\xdef\consXG{-1} \xdef\consXN{0} \xdef\consXD{-1} \xdef\consXE{-1} \xdef\consXQ{-1}
\xdef\consXH{-1} \xdef\consXR{-1} \xdef\consXK{-1} \xdef\consXM{-1} \xdef\consXI{-1}
\xdef\consXL{-1} \xdef\consXV{-1} \xdef\consXF{-2} \xdef\consXY{-2} \xdef\consXW{-4}
\xdef\consXB{-1} \xdef\consXZ{-1}  \xdef\consXX{-1}

\expandafter\xdef\csname cons..\endcsname{1}
\gappenalty{-8} \xdef\m@trixf@ctor{10}

\else
\xdef\temp@{PAM100} \xdef\m@trixf@ctor{10}
\ifx\first@\temp@

%%%% PAM100

\xdef\consCC{14} \xdef\consCS{-1} \xdef\consCT{-5} \xdef\consCP{-6} \xdef\consCA{-5}
\xdef\consCG{-8} \xdef\consCN{-8} \xdef\consCD{-11} \xdef\consCE{-11} \xdef\consCQ{-11}
\xdef\consCH{-6} \xdef\consCR{-6} \xdef\consCK{-11} \xdef\consCM{-11} \xdef\consCI{-5}
\xdef\consCL{-12} \xdef\consCV{-4} \xdef\consCF{-10} \xdef\consCY{-2} \xdef\consCW{-13}
\xdef\consCB{-6} \xdef\consCZ{-8}  \xdef\consCX{-5}

\xdef\consSC{-1} \xdef\consSS{6} \xdef\consST{2} \xdef\consSP{1} \xdef\consSA{2}
\xdef\consSG{1} \xdef\consSN{2} \xdef\consSD{-1} \xdef\consSE{-2} \xdef\consSQ{-3}
\xdef\consSH{-4} \xdef\consSR{-1} \xdef\consSK{-2} \xdef\consSM{-4} \xdef\consSI{-4}
\xdef\consSL{-7} \xdef\consSV{-4} \xdef\consSF{-5} \xdef\consSY{-6} \xdef\consSW{-4}
\xdef\consSB{0} \xdef\consSZ{-2}  \xdef\consSX{1}

\xdef\consTC{-5} \xdef\consTS{2} \xdef\consTT{7} \xdef\consTP{-1} \xdef\consTA{2}
\xdef\consTG{-3} \xdef\consTN{0} \xdef\consTD{-2} \xdef\consTE{-3} \xdef\consTQ{-3}
\xdef\consTH{-5} \xdef\consTR{-4} \xdef\consTK{-1} \xdef\consTM{-2} \xdef\consTI{-1}
\xdef\consTL{-5} \xdef\consTV{-1} \xdef\consTF{-6} \xdef\consTY{-6} \xdef\consTW{-10}
\xdef\consTB{-1} \xdef\consTZ{-2}  \xdef\consTX{-1}

\xdef\consPC{-6} \xdef\consPS{1} \xdef\consPT{-1} \xdef\consPP{10} \xdef\consPA{1}
\xdef\consPG{-3} \xdef\consPN{-3} \xdef\consPD{-4} \xdef\consPE{-3} \xdef\consPQ{-1}
\xdef\consPH{-2} \xdef\consPR{-2} \xdef\consPK{-4} \xdef\consPM{-6} \xdef\consPI{-6}
\xdef\consPL{-5} \xdef\consPV{-4} \xdef\consPF{-9} \xdef\consPY{-11} \xdef\consPW{-11}
\xdef\consPB{-3} \xdef\consPZ{-1}  \xdef\consPX{-2}

\xdef\consAC{-5} \xdef\consAS{2} \xdef\consAT{2} \xdef\consAP{1} \xdef\consAA{6}
\xdef\consAG{1} \xdef\consAN{-1} \xdef\consAD{-1} \xdef\consAE{0} \xdef\consAQ{-2}
\xdef\consAH{-5} \xdef\consAR{-5} \xdef\consAK{-4} \xdef\consAM{-3} \xdef\consAI{-3}
\xdef\consAL{-5} \xdef\consAV{0} \xdef\consAF{-7} \xdef\consAY{-6} \xdef\consAW{-11} 
\xdef\consAB{-1} \xdef\consAZ{-1}  \xdef\consAX{-1}

\xdef\consGC{-8} \xdef\consGS{1} \xdef\consGT{-3} \xdef\consGP{-3} \xdef\consGA{1}
\xdef\consGG{8} \xdef\consGN{-1} \xdef\consGD{-1} \xdef\consGE{-2} \xdef\consGQ{-5}
\xdef\consGH{-7} \xdef\consGR{-8} \xdef\consGK{-5} \xdef\consGM{-8} \xdef\consGI{-7}
\xdef\consGL{-8} \xdef\consGV{-4} \xdef\consGF{-8} \xdef\consGY{-11} \xdef\consGW{-13}
\xdef\consGB{-1} \xdef\consGZ{-2}  \xdef\consGX{-2}

\xdef\consNC{-8} \xdef\consNS{2} \xdef\consNT{0} \xdef\consNP{-3} \xdef\consNA{-1}
\xdef\consNG{-1} \xdef\consNN{7} \xdef\consND{4} \xdef\consNE{1} \xdef\consNQ{-1}
\xdef\consNH{2} \xdef\consNR{-3} \xdef\consNK{1} \xdef\consNM{-5} \xdef\consNI{-4}
\xdef\consNL{-6} \xdef\consNV{-5} \xdef\consNF{-6} \xdef\consNY{-3} \xdef\consNW{-8}
\xdef\consNB{4} \xdef\consNZ{0}  \xdef\consNX{-1}

\xdef\consDC{-11} \xdef\consDS{-1} \xdef\consDT{-2} \xdef\consDP{-4} \xdef\consDA{-1}
\xdef\consDG{-1} \xdef\consDN{4} \xdef\consDD{8} \xdef\consDE{5} \xdef\consDQ{1}
\xdef\consDH{-1} \xdef\consDR{-6} \xdef\consDK{-2} \xdef\consDM{-8} \xdef\consDI{-6}
\xdef\consDL{-9} \xdef\consDV{-6} \xdef\consDF{-11} \xdef\consDY{-9} \xdef\consDW{-13}
\xdef\consDB{4} \xdef\consDZ{3}  \xdef\consDX{-2}

\xdef\consEC{-11} \xdef\consES{-2} \xdef\consET{-3} \xdef\consEP{-3} \xdef\consEA{0}
\xdef\consEG{-2} \xdef\consEN{1} \xdef\consED{5} \xdef\consEE{8} \xdef\consEQ{4}
\xdef\consEH{-2} \xdef\consER{-5} \xdef\consEK{-2} \xdef\consEM{-6} \xdef\consEI{-5}
\xdef\consEL{-7} \xdef\consEV{-5} \xdef\consEF{-11} \xdef\consEY{-7} \xdef\consEW{-14}
\xdef\consEB{3} \xdef\consEZ{4}  \xdef\consEX{-2}

\xdef\consQC{-11} \xdef\consQS{-3} \xdef\consQT{-3} \xdef\consQP{-1} \xdef\consQA{-2}
\xdef\consQG{-5} \xdef\consQN{-1} \xdef\consQD{1} \xdef\consQE{4} \xdef\consQQ{9}
\xdef\consQH{4} \xdef\consQR{1} \xdef\consQK{-1} \xdef\consQM{-2} \xdef\consQI{-5}
\xdef\consQL{-3} \xdef\consQV{-5} \xdef\consQF{-10} \xdef\consQY{-9} \xdef\consQW{-11}
\xdef\consQB{0} \xdef\consQZ{5}  \xdef\consQX{-2}

\xdef\consHC{-6} \xdef\consHS{-4} \xdef\consHT{-5} \xdef\consHP{-2} \xdef\consHA{-5}
\xdef\consHG{-7} \xdef\consHN{2} \xdef\consHD{-1} \xdef\consHE{-2} \xdef\consHQ{4}
\xdef\consHH{11} \xdef\consHR{1} \xdef\consHK{-3} \xdef\consHM{-7} \xdef\consHI{-7}
\xdef\consHL{-5} \xdef\consHV{-6} \xdef\consHF{-4} \xdef\consHY{-1} \xdef\consHW{-7}
\xdef\consHB{1} \xdef\consHZ{1}  \xdef\consHX{-2}

\xdef\consRC{-6} \xdef\consRS{-1} \xdef\consRT{-4} \xdef\consRP{-2} \xdef\consRA{-5}
\xdef\consRG{-8} \xdef\consRN{-3} \xdef\consRD{-6} \xdef\consRE{-5} \xdef\consRQ{1}
\xdef\consRH{1} \xdef\consRR{10} \xdef\consRK{3} \xdef\consRM{-2} \xdef\consRI{-4}
\xdef\consRL{-7} \xdef\consRV{-6} \xdef\consRF{-7} \xdef\consRY{-10} \xdef\consRW{1}
\xdef\consRB{-3} \xdef\consRZ{-1}  \xdef\consRX{-2}

\xdef\consKC{-11} \xdef\consKS{-2} \xdef\consKT{-1} \xdef\consKP{-4} \xdef\consKA{-4}
\xdef\consKG{-5} \xdef\consKN{1} \xdef\consKD{-2} \xdef\consKE{-2} \xdef\consKQ{-1}
\xdef\consKH{-3} \xdef\consKR{3} \xdef\consKK{8} \xdef\consKM{1} \xdef\consKI{-4}
\xdef\consKL{-6} \xdef\consKV{-6} \xdef\consKF{-11} \xdef\consKY{-10} \xdef\consKW{-9}
\xdef\consKB{0} \xdef\consKZ{-1}  \xdef\consKX{-2}

\xdef\consMC{-11} \xdef\consMS{-4} \xdef\consMT{-2} \xdef\consMP{-6} \xdef\consMA{-3}
\xdef\consMG{-8} \xdef\consMN{-5} \xdef\consMD{-8} \xdef\consME{-6} \xdef\consMQ{-2}
\xdef\consMH{-7} \xdef\consMR{-2} \xdef\consMK{1} \xdef\consMM{13} \xdef\consMI{2}
\xdef\consML{4} \xdef\consMV{1} \xdef\consMF{-2} \xdef\consMY{-8} \xdef\consMW{-11}
\xdef\consMB{-4} \xdef\consMZ{-2}  \xdef\consMX{-2}

\xdef\consIC{-5} \xdef\consIS{-4} \xdef\consIT{-1} \xdef\consIP{-6} \xdef\consIA{-3}
\xdef\consIG{-7} \xdef\consIN{-4} \xdef\consID{-6} \xdef\consIE{-5} \xdef\consIQ{-5}
\xdef\consIH{-7} \xdef\consIR{-4} \xdef\consIK{-4} \xdef\consIM{2} \xdef\consII{9}
\xdef\consIL{2} \xdef\consIV{5} \xdef\consIF{0} \xdef\consIY{-4} \xdef\consIW{-12}
\xdef\consIB{-3} \xdef\consIZ{-3}  \xdef\consIX{-2}

\xdef\consLC{-12} \xdef\consLS{-7} \xdef\consLT{-5} \xdef\consLP{-5} \xdef\consLA{-5}
\xdef\consLG{-8} \xdef\consLN{-6} \xdef\consLD{-9} \xdef\consLE{-7} \xdef\consLQ{-3}
\xdef\consLH{-5} \xdef\consLR{-7} \xdef\consLK{-6} \xdef\consLM{4} \xdef\consLI{2}
\xdef\consLL{9} \xdef\consLV{1} \xdef\consLF{0} \xdef\consLY{-5} \xdef\consLW{-7}
\xdef\consLB{-5} \xdef\consLZ{-4}  \xdef\consLX{-3}

\xdef\consVC{-4} \xdef\consVS{-4} \xdef\consVT{-1} \xdef\consVP{-4} \xdef\consVA{0}
\xdef\consVG{-4} \xdef\consVN{-5} \xdef\consVD{-6} \xdef\consVE{-5} \xdef\consVQ{-5}
\xdef\consVH{-6} \xdef\consVR{-6} \xdef\consVK{-6} \xdef\consVM{1} \xdef\consVI{5}
\xdef\consVL{1} \xdef\consVV{8} \xdef\consVF{-5} \xdef\consVY{-6} \xdef\consVW{-14}
\xdef\consVB{-4} \xdef\consVZ{-3}  \xdef\consVX{-2}

\xdef\consFC{-10} \xdef\consFS{-5} \xdef\consFT{-6} \xdef\consFP{-9} \xdef\consFA{-7}
\xdef\consFG{-8} \xdef\consFN{-6} \xdef\consFD{-11} \xdef\consFE{-11} \xdef\consFQ{-10}
\xdef\consFH{-4} \xdef\consFR{-7} \xdef\consFK{-11} \xdef\consFM{-2} \xdef\consFI{0}
\xdef\consFL{0} \xdef\consFV{-5} \xdef\consFF{12} \xdef\consFY{6} \xdef\consFW{-2}
\xdef\consFB{-6} \xdef\consFZ{-7}  \xdef\consFX{-4}

\xdef\consYC{-2} \xdef\consYS{-6} \xdef\consYT{-6} \xdef\consYP{-11} \xdef\consYA{-6}
\xdef\consYG{-11} \xdef\consYN{-3} \xdef\consYD{-9} \xdef\consYE{-7} \xdef\consYQ{-9}
\xdef\consYH{-1} \xdef\consYR{-10} \xdef\consYK{-10} \xdef\consYM{-8} \xdef\consYI{-4}
\xdef\consYL{-5} \xdef\consYV{-6} \xdef\consYF{6} \xdef\consYY{13} \xdef\consYW{-2}
\xdef\consYB{-4} \xdef\consYZ{-6}  \xdef\consYX{-4}

\xdef\consWC{-13} \xdef\consWS{-4} \xdef\consWT{-10} \xdef\consWP{-11} \xdef\consWA{-11}
\xdef\consWG{-13} \xdef\consWN{-8} \xdef\consWD{-13} \xdef\consWE{-14} \xdef\consWQ{-11}
\xdef\consWH{-7} \xdef\consWR{1} \xdef\consWK{-9} \xdef\consWM{-11} \xdef\consWI{-12}
\xdef\consWL{-7} \xdef\consWV{-14} \xdef\consWF{-2} \xdef\consWY{-2} \xdef\consWW{19}
\xdef\consWB{-6} \xdef\consWZ{-8}  \xdef\consWX{-6}

\xdef\consBC{-6} \xdef\consBS{0} \xdef\consBT{-1} \xdef\consBP{-3} \xdef\consBA{-1}
\xdef\consBG{-1} \xdef\consBN{4} \xdef\consBD{4} \xdef\consBE{3} \xdef\consBQ{0}
\xdef\consBH{1} \xdef\consBR{-3} \xdef\consBK{0} \xdef\consBM{-4} \xdef\consBI{-3}
\xdef\consBL{-5} \xdef\consBV{-4} \xdef\consBF{-6} \xdef\consBY{-4} \xdef\consBW{-6}
\xdef\consBB{4} \xdef\consBZ{2}  \xdef\consBX{-2}

\xdef\consZC{-8} \xdef\consZS{-2} \xdef\consZT{-2} \xdef\consZP{-1} \xdef\consZA{-1}
\xdef\consZG{-2} \xdef\consZN{0} \xdef\consZD{3} \xdef\consZE{4} \xdef\consZQ{5}
\xdef\consZH{1} \xdef\consZR{-1} \xdef\consZK{-1} \xdef\consZM{-2} \xdef\consZI{-3}
\xdef\consZL{-4} \xdef\consZV{-3} \xdef\consZF{-7} \xdef\consZY{-6} \xdef\consZW{-8}
\xdef\consZB{2} \xdef\consZZ{5}  \xdef\consZX{-2}

\xdef\consXC{-5} \xdef\consXS{-1} \xdef\consXT{-1} \xdef\consXP{-2} \xdef\consXA{-1}
\xdef\consXG{-2} \xdef\consXN{-1} \xdef\consXD{-2} \xdef\consXE{-2} \xdef\consXQ{-2}
\xdef\consXH{-2} \xdef\consXR{-2} \xdef\consXK{-2} \xdef\consXM{-2} \xdef\consXI{-2}
\xdef\consXL{-3} \xdef\consXV{-2} \xdef\consXF{-4} \xdef\consXY{-4} \xdef\consXW{-6}
\xdef\consXB{-2} \xdef\consXZ{-2}  \xdef\consXX{-2}

\expandafter\xdef\csname cons..\endcsname{1}
\gappenalty{-9} \xdef\m@trixf@ctor{10}

\else
\xdef\temp@{BLOSUM62}
\ifx\first@\temp@

%%%% BLOSUM62

\xdef\consCC{9} \xdef\consCS{-1} \xdef\consCT{-1} \xdef\consCP{-3} \xdef\consCA{0}
\xdef\consCG{-3} \xdef\consCN{-3} \xdef\consCD{-3} \xdef\consCE{-4} \xdef\consCQ{-3}
\xdef\consCH{-3} \xdef\consCR{-3} \xdef\consCK{-3} \xdef\consCM{-1} \xdef\consCI{-1}
\xdef\consCL{-1} \xdef\consCV{-1} \xdef\consCF{-2} \xdef\consCY{-2} \xdef\consCW{-2}

\xdef\consSC{-1} \xdef\consSS{4} \xdef\consST{1} \xdef\consSP{-1} \xdef\consSA{1}
\xdef\consSG{0} \xdef\consSN{1} \xdef\consSD{0} \xdef\consSE{0} \xdef\consSQ{0}
\xdef\consSH{-1} \xdef\consSR{-1} \xdef\consSK{0} \xdef\consSM{-1} \xdef\consSI{-2}
\xdef\consSL{-2} \xdef\consSV{-2} \xdef\consSF{-2} \xdef\consSY{-2} \xdef\consSW{-3}

\xdef\consTC{-1} \xdef\consTS{1} \xdef\consTT{4} \xdef\consTP{1} \xdef\consTA{-1}
\xdef\consTG{1} \xdef\consTN{0} \xdef\consTD{1} \xdef\consTE{0} \xdef\consTQ{0}
\xdef\consTH{0} \xdef\consTR{-1} \xdef\consTK{0} \xdef\consTM{-1} \xdef\consTI{-2}
\xdef\consTL{-2} \xdef\consTV{-2} \xdef\consTF{-2} \xdef\consTY{-2} \xdef\consTW{-3}

\xdef\consPC{-3} \xdef\consPS{-1} \xdef\consPT{1} \xdef\consPP{7} \xdef\consPA{-1}
\xdef\consPG{-2} \xdef\consPN{-1} \xdef\consPD{-1} \xdef\consPE{-1} \xdef\consPQ{-1}
\xdef\consPH{-2} \xdef\consPR{-2} \xdef\consPK{-1} \xdef\consPM{-2} \xdef\consPI{-3}
\xdef\consPL{-3} \xdef\consPV{-2} \xdef\consPF{-4} \xdef\consPY{-3} \xdef\consPW{-4}

\xdef\consAC{0} \xdef\consAS{1} \xdef\consAT{-1} \xdef\consAP{-1} \xdef\consAA{4}
\xdef\consAG{0} \xdef\consAN{-1} \xdef\consAD{-2} \xdef\consAE{-1} \xdef\consAQ{-1}
\xdef\consAH{-2} \xdef\consAR{-1} \xdef\consAK{-1} \xdef\consAM{-1} \xdef\consAI{-1}
\xdef\consAL{-1} \xdef\consAV{-2} \xdef\consAF{-2} \xdef\consAY{-2} \xdef\consAW{-3} 

\xdef\consGC{-3} \xdef\consGS{0} \xdef\consGT{1} \xdef\consGP{-2} \xdef\consGA{0}
\xdef\consGG{6} \xdef\consGN{-2} \xdef\consGD{-1} \xdef\consGE{-2} \xdef\consGQ{-2}
\xdef\consGH{-2} \xdef\consGR{-2} \xdef\consGK{-2} \xdef\consGM{-3} \xdef\consGI{-4}
\xdef\consGL{-4} \xdef\consGV{0} \xdef\consGF{-3} \xdef\consGY{-3} \xdef\consGW{-2}

\xdef\consNC{-3} \xdef\consNS{1} \xdef\consNT{0} \xdef\consNP{-2} \xdef\consNA{-2}
\xdef\consNG{0} \xdef\consNN{6} \xdef\consND{1} \xdef\consNE{0} \xdef\consNQ{0}
\xdef\consNH{-1} \xdef\consNR{0} \xdef\consNK{0} \xdef\consNM{-2} \xdef\consNI{-3}
\xdef\consNL{-3} \xdef\consNV{-3} \xdef\consNF{-3} \xdef\consNY{-2} \xdef\consNW{-4}

\xdef\consDC{-3} \xdef\consDS{0} \xdef\consDT{1} \xdef\consDP{-1} \xdef\consDA{-2}
\xdef\consDG{-1} \xdef\consDN{1} \xdef\consDD{6} \xdef\consDE{2} \xdef\consDQ{0}
\xdef\consDH{-1} \xdef\consDR{-2} \xdef\consDK{-1} \xdef\consDM{-3} \xdef\consDI{-3}
\xdef\consDL{-4} \xdef\consDV{-3} \xdef\consDF{-3} \xdef\consDY{-3} \xdef\consDW{-4}

\xdef\consEC{-4} \xdef\consES{0} \xdef\consET{0} \xdef\consEP{-1} \xdef\consEA{-1}
\xdef\consEG{-2} \xdef\consEN{0} \xdef\consED{2} \xdef\consEE{5} \xdef\consEQ{2}
\xdef\consEH{0} \xdef\consER{0} \xdef\consEK{1} \xdef\consEM{-2} \xdef\consEI{-3}
\xdef\consEL{-3} \xdef\consEV{-3} \xdef\consEF{-3} \xdef\consEY{-2} \xdef\consEW{-3}

\xdef\consQC{-3} \xdef\consQS{0} \xdef\consQT{0} \xdef\consQP{-1} \xdef\consQA{-1}
\xdef\consQG{-2} \xdef\consQN{0} \xdef\consQD{0} \xdef\consQE{2} \xdef\consQQ{5}
\xdef\consQH{0} \xdef\consQR{1} \xdef\consQK{1} \xdef\consQM{0} \xdef\consQI{-3}
\xdef\consQL{-2} \xdef\consQV{-2} \xdef\consQF{-3} \xdef\consQY{-1} \xdef\consQW{-2}

\xdef\consHC{-3} \xdef\consHS{-1} \xdef\consHT{0} \xdef\consHP{-2} \xdef\consHA{-2}
\xdef\consHG{-2} \xdef\consHN{1} \xdef\consHD{1} \xdef\consHE{0} \xdef\consHQ{0}
\xdef\consHH{8} \xdef\consHR{0} \xdef\consHK{-1} \xdef\consHM{-2} \xdef\consHI{-3}
\xdef\consHL{-3} \xdef\consHV{-2} \xdef\consHF{-1} \xdef\consHY{2} \xdef\consHW{-2}

\xdef\consRC{-3} \xdef\consRS{-1} \xdef\consRT{-1} \xdef\consRP{-2} \xdef\consRA{-1}
\xdef\consRG{-2} \xdef\consRN{0} \xdef\consRD{-2} \xdef\consRE{0} \xdef\consRQ{1}
\xdef\consRH{0} \xdef\consRR{5} \xdef\consRK{2} \xdef\consRM{-1} \xdef\consRI{-3}
\xdef\consRL{-2} \xdef\consRV{-3} \xdef\consRF{-3} \xdef\consRY{-2} \xdef\consRW{-3}

\xdef\consKC{-3} \xdef\consKS{0} \xdef\consKT{0} \xdef\consKP{-1} \xdef\consKA{-1}
\xdef\consKG{-2} \xdef\consKN{0} \xdef\consKD{-1} \xdef\consKE{1} \xdef\consKQ{1}
\xdef\consKH{-1} \xdef\consKR{2} \xdef\consKK{5} \xdef\consKM{-1} \xdef\consKI{-3}
\xdef\consKL{-2} \xdef\consKV{-3} \xdef\consKF{-3} \xdef\consKY{-2} \xdef\consKW{-3}

\xdef\consMC{-1} \xdef\consMS{-1} \xdef\consMT{-1} \xdef\consMP{-2} \xdef\consMA{-1}
\xdef\consMG{-3} \xdef\consMN{-2} \xdef\consMD{-3} \xdef\consME{-2} \xdef\consMQ{0}
\xdef\consMH{-2} \xdef\consMR{-1} \xdef\consMK{-1} \xdef\consMM{5} \xdef\consMI{1}
\xdef\consML{2} \xdef\consMV{-2} \xdef\consMF{0} \xdef\consMY{-1} \xdef\consMW{-1}

\xdef\consIC{-1} \xdef\consIS{-2} \xdef\consIT{-2} \xdef\consIP{-3} \xdef\consIA{-1}
\xdef\consIG{-4} \xdef\consIN{-3} \xdef\consID{-3} \xdef\consIE{-3} \xdef\consIQ{-3}
\xdef\consIH{-3} \xdef\consIR{-3} \xdef\consIK{-3} \xdef\consIM{1} \xdef\consII{4}
\xdef\consIL{2} \xdef\consIV{1} \xdef\consIF{0} \xdef\consIY{-1} \xdef\consIW{-3}

\xdef\consLC{-1} \xdef\consLS{-2} \xdef\consLT{-2} \xdef\consLP{-3} \xdef\consLA{-1}
\xdef\consLG{-4} \xdef\consLN{-3} \xdef\consLD{-4} \xdef\consLE{-3} \xdef\consLQ{-2}
\xdef\consLH{-3} \xdef\consLR{-2} \xdef\consLK{-2} \xdef\consLM{2} \xdef\consLI{2}
\xdef\consLL{4} \xdef\consLV{3} \xdef\consLF{0} \xdef\consLY{-1} \xdef\consLW{-2}

\xdef\consVC{-1} \xdef\consVS{-2} \xdef\consVT{-2} \xdef\consVP{-2} \xdef\consVA{0}
\xdef\consVG{-3} \xdef\consVN{-3} \xdef\consVD{-3} \xdef\consVE{-2} \xdef\consVQ{-2}
\xdef\consVH{-3} \xdef\consVR{-3} \xdef\consVK{-2} \xdef\consVM{1} \xdef\consVI{3}
\xdef\consVL{1} \xdef\consVV{4} \xdef\consVF{-1} \xdef\consVY{-1} \xdef\consVW{-3}

\xdef\consFC{-2} \xdef\consFS{-2} \xdef\consFT{-2} \xdef\consFP{-4} \xdef\consFA{-2}
\xdef\consFG{-3} \xdef\consFN{-3} \xdef\consFD{-3} \xdef\consFE{-3} \xdef\consFQ{-3}
\xdef\consFH{-1} \xdef\consFR{-3} \xdef\consFK{-3} \xdef\consFM{0} \xdef\consFI{0}
\xdef\consFL{0} \xdef\consFV{-1} \xdef\consFF{6} \xdef\consFY{3} \xdef\consFW{1}

\xdef\consYC{-2} \xdef\consYS{-2} \xdef\consYT{-2} \xdef\consYP{-3} \xdef\consYA{-2}
\xdef\consYG{-3} \xdef\consYN{-2} \xdef\consYD{-3} \xdef\consYE{-2} \xdef\consYQ{-1}
\xdef\consYH{2} \xdef\consYR{-2} \xdef\consYK{-2} \xdef\consYM{-1} \xdef\consYI{-1}
\xdef\consYL{-1} \xdef\consYV{-1} \xdef\consYF{3} \xdef\consYY{7} \xdef\consYW{2}

\xdef\consWC{-2} \xdef\consWS{-3} \xdef\consWT{-3} \xdef\consWP{-4} \xdef\consWA{-3}
\xdef\consWG{-2} \xdef\consWN{-4} \xdef\consWD{-4} \xdef\consWE{-3} \xdef\consWQ{-2}
\xdef\consWH{-2} \xdef\consWR{-3} \xdef\consWK{-3} \xdef\consWM{-1} \xdef\consWI{-3}
\xdef\consWL{-2} \xdef\consWV{-3} \xdef\consWF{1} \xdef\consWY{2} \xdef\consWW{11}

\gappenalty{0} \xdef\m@trixf@ctor{10}

\else

\message{<Unknown residue weight matrix - using `identity'>} \weighttable{identity}

\fi
\fi
\fi
\fi
\fi
}


\def\setweight#1#2#3{%
   \expandafter\xdef\csname cons#1#2\endcsname{#3}
   \expandafter\xdef\csname cons#2#1\endcsname{#3}
}

\weighttable{identity}
\gappenalty{0}


\def\c@d@ns{%
\codon{A}{GCA,GCG,GCC,GCT,GCU,GCN}
\codon{B}{---}
\codon{C}{TGC,TGT,UGC,UGU,TGY}
\codon{D}{GAC,GAT,GAU,GAY}
\codon{E}{GAA,GAG,GAR}
\codon{F}{TTC,TTT,UUC,UUU,TTY}
\codon{G}{GGA,GGG,GGC,GGT,GGU,GGN}
\codon{H}{CAC,CAT,CAY}
\codon{I}{ATA,ATC,ATT,AUA,AUC,AUU,ATH}
\codon{J}{---}
\codon{K}{AAA,AAG,AAG,AAR}
\codon{L}{CTA,CTG,CTC,CTT,TTA,TTG,CUG,CUG,CUC,CUU,UUA,UUG,YTN}
\codon{M}{ATG,AUG,ATG}
\codon{N}{AAC,AAT,AAU,AAY}
\codon{O}{---}
\codon{P}{CCA,CCG,CCC,CCT,CCU,CCN}
\codon{Q}{CAA,CAG,CAR}
\codon{R}{AGA,AGG,CGA,CGG,CGC,CGT,CGU,MGN}
\codon{S}{TCT,TCC,TCG,TCA,AGT,AGC,UCU,UCC,UCG,UCA,AGU,WSN}
\codon{T}{ACT,ACC,ACG,ACA,ACU,ACN}
\codon{U}{---}
\codon{V}{GTA,GTG,GTC,GTT,GUA,GUG,GUC,GUU,GTN}
\codon{W}{TGG,UGG,TGG}
\codon{X}{---}
\codon{Y}{TAC,TAT,UAC,UAU,TAY}
\codon{Z}{---}
\codon{.}{TAA,TAG,TGA,UAA,UAG,UGA,TRR}
}

\definecolor{GreenYellow}   {cmyk}{0.15,0,0.69,0}
\definecolor{Yellow}        {cmyk}{0,0,1,0}
\definecolor{Goldenrod}     {cmyk}{0,0.10,0.84,0}
\definecolor{Dandelion}     {cmyk}{0,0.29,0.84,0}
\definecolor{Apricot}       {cmyk}{0,0.32,0.52,0}
\definecolor{Peach}         {cmyk}{0,0.50,0.70,0}
\definecolor{Melon}         {cmyk}{0,0.46,0.50,0}
\definecolor{YellowOrange}  {cmyk}{0,0.42,1,0}
\definecolor{Orange}        {cmyk}{0,0.61,0.87,0}
\definecolor{BurntOrange}   {cmyk}{0,0.51,1,0}
\definecolor{Bittersweet}   {cmyk}{0,0.75,1,0.24}
\definecolor{RedOrange}     {cmyk}{0,0.77,0.87,0}
\definecolor{Mahagony}      {cmyk}{0,0.85,0.87,0.35}
\definecolor{Maroon}        {cmyk}{0,0.87,0.68,0.32}
\definecolor{BrickRed}      {cmyk}{0,0.89,0.94,0.28}
\definecolor{Red}           {cmyk}{0,1,1,0}
\definecolor{OrangeRed}     {cmyk}{0,1,0.50,0}
\definecolor{RubineRed}     {cmyk}{0,1,0.13,0}
\definecolor{WildStrawberry}{cmyk}{0,0.96,0.39,0}
\definecolor{Salmon}        {cmyk}{0,0.53,0.38,0}
\definecolor{CarnationPink} {cmyk}{0,0.63,0,0}
\definecolor{Magenta}       {cmyk}{0,1,0,0}
\definecolor{VioletRed}     {cmyk}{0,0.81,0,0}
\definecolor{Rhodamine}     {cmyk}{0,0.82,0,0}
\definecolor{Mulberry}      {cmyk}{0.34,0.90,0,0.02}
\definecolor{RedViolet}     {cmyk}{0.07,0.90,0,0.34}
\definecolor{Fuchsia}       {cmyk}{0.47,0.91,0,0.08}
\definecolor{Lavender}      {cmyk}{0,0.48,0,0}
\definecolor{Thistle}       {cmyk}{0.12,0.59,0,0}
\definecolor{Orchid}        {cmyk}{0.32,0.64,0,0}
\definecolor{DarkOrchid}    {cmyk}{0.40,0.80,0.20,0}
\definecolor{Purple}        {cmyk}{0.45,0.86,0,0}
\definecolor{Plum}          {cmyk}{0.50,1,0,0}
\definecolor{Violet}        {cmyk}{0.79,0.88,0,0}
\definecolor{RoyalPurple}   {cmyk}{0.75,0.90,0,0}
\definecolor{BlueViolet}    {cmyk}{0.86,0.91,0,0.04}
\definecolor{Periwinkle}    {cmyk}{0.57,0.55,0,0}
\definecolor{CadetBlue}     {cmyk}{0.62,0.57,0.23,0}
\definecolor{CornflowerBlue}{cmyk}{0.65,0.13,0,0}
\definecolor{MidnightBlue}  {cmyk}{0.98,0.13,0,0.43}
\definecolor{NavyBlue}      {cmyk}{0.94,0.54,0,0}
\definecolor{RoyalBlue}     {cmyk}{1,0.50,0,0}
\definecolor{Blue}          {cmyk}{1,1,0,0}
\definecolor{Cerulean}      {cmyk}{0.94,0.11,0,0}
\definecolor{Cyan}          {cmyk}{1,0,0,0}
\definecolor{ProcessBlue}   {cmyk}{0.96,0,0,0}
\definecolor{SkyBlue}       {cmyk}{0.62,0,0.12,0}
\definecolor{Turquoise}     {cmyk}{0.85,0,0.20,0}
\definecolor{TealBlue}      {cmyk}{0.86,0,0.34,0.02}
\definecolor{Aquamarine}    {cmyk}{0.82,0,0.30,0}
\definecolor{BlueGreen}     {cmyk}{0.85,0,0.33,0}
\definecolor{Emerald}       {cmyk}{1,0,0.50,0}
\definecolor{JungleGreen}   {cmyk}{0.99,0,0.52,0}
\definecolor{SeaGreen}      {cmyk}{0.69,0,0.50,0}
\definecolor{Green}         {cmyk}{1,0,1,0}
\definecolor{ForestGreen}   {cmyk}{0.91,0,0.88,0.12}
\definecolor{PineGreen}     {cmyk}{0.92,0,0.59,0.25}
\definecolor{LimeGreen}     {cmyk}{0.50,0,1,0}
\definecolor{YellowGreen}   {cmyk}{0.44,0,0.74,0}
\definecolor{SpringGreen}   {cmyk}{0.26,0,0.76,0}
\definecolor{OliveGreen}    {cmyk}{0.64,0,0.95,0.40}
\definecolor{RawSienna}     {cmyk}{0,0.72,1,0.45}
\definecolor{Sepia}         {cmyk}{0,0.83,1,0.70}
\definecolor{Brown}         {cmyk}{0,0.81,1,0.60}
\definecolor{Tan}           {cmyk}{0.14,0.42,0.56,0}
\definecolor{White}         {cmyk}{0,0,0,0}
\definecolor{Gray0}         {cmyk}{0,0,0,0}
\definecolor{Gray5}         {cmyk}{0,0,0,0.05}
\definecolor{Gray10}        {cmyk}{0,0,0,0.10}
\definecolor{Gray15}        {cmyk}{0,0,0,0.15}
\definecolor{Gray20}        {cmyk}{0,0,0,0.20}
\definecolor{Gray25}        {cmyk}{0,0,0,0.25}
\definecolor{Gray30}        {cmyk}{0,0,0,0.30}
\definecolor{LightGray}     {cmyk}{0,0,0,0.33}
\definecolor{Gray35}        {cmyk}{0,0,0,0.35}
\definecolor{Gray40}        {cmyk}{0,0,0,0.40}
\definecolor{Gray45}        {cmyk}{0,0,0,0.45}
\definecolor{Gray50}        {cmyk}{0,0,0,0.50}
\definecolor{Gray}          {cmyk}{0,0,0,0.50}
\definecolor{GrayDefault}   {cmyk}{0,0,0,0.50}
\definecolor{Gray55}        {cmyk}{0,0,0,0.55}
\definecolor{Gray60}        {cmyk}{0,0,0,0.60}
\definecolor{Gray65}        {cmyk}{0,0,0,0.65}
\definecolor{DarkGray}      {cmyk}{0,0,0,0.66}
\definecolor{Gray70}        {cmyk}{0,0,0,0.70}
\definecolor{Gray75}        {cmyk}{0,0,0,0.75}
\definecolor{Gray80}        {cmyk}{0,0,0,0.80}
\definecolor{Gray85}        {cmyk}{0,0,0,0.85}
\definecolor{Gray90}        {cmyk}{0,0,0,0.90}
\definecolor{Gray95}        {cmyk}{0,0,0,0.95}
\definecolor{Black}         {cmyk}{0,0,0,1}
\definecolor{Gray100}       {cmyk}{0,0,0,1}
\definecolor{TC0}           {cmyk}{0.4,0.4,0,0}
\definecolor{TC1}           {cmyk}{0.6,0,0.7,0}
\definecolor{TC2}           {cmyk}{0.4,0,0.7,0}
\definecolor{TC3}           {cmyk}{0.2,0,1,0}
\definecolor{TC4}           {cmyk}{0,0,1,0}
\definecolor{TC5}           {cmyk}{0,0.2,1,0}
\definecolor{TC6}           {cmyk}{0,0.4,1,0}
\definecolor{TC7}           {cmyk}{0,0.6,1,0}
\definecolor{TC8}           {cmyk}{0,0.8,1,0}
\definecolor{TC9}           {cmyk}{0,0.875,1,0}
\definecolor{TC99}          {cmyk}{0,0,0,0}
\definecolor{LightGreenYellow}   {cmyk}{0.08,0,0.35,0}
\definecolor{LightYellow}        {cmyk}{0,0,0.50,0}
\definecolor{LightGoldenrod}     {cmyk}{0,0.05,0.42,0}
\definecolor{LightDandelion}     {cmyk}{0,0.15,0.42,0}
\definecolor{LightApricot}       {cmyk}{0,0.16,0.26,0}
\definecolor{LightPeach}         {cmyk}{0,0.25,0.35,0}
\definecolor{LightMelon}         {cmyk}{0,0.23,0.25,0}
\definecolor{LightYellowOrange}  {cmyk}{0,0.21,0.50,0}
\definecolor{LightOrange}        {cmyk}{0,0.31,0.44,0}
\definecolor{LightBurntOrange}   {cmyk}{0,0.26,0.50,0}
\definecolor{LightBittersweet}   {cmyk}{0,0.38,0.50,0.12}
\definecolor{LightRedOrange}     {cmyk}{0,0.39,0.44,0}
\definecolor{LightMahagony}      {cmyk}{0,0.43,0.44,0.18}
\definecolor{LightMaroon}        {cmyk}{0,0.44,0.34,0.16}
\definecolor{LightBrickRed}      {cmyk}{0,0.45,0.47,0.14}
\definecolor{LightRed}           {cmyk}{0,0.50,0.50,0}
\definecolor{LightOrangeRed}     {cmyk}{0,0.50,0.25,0}
\definecolor{LightRubineRed}     {cmyk}{0,0.50,0.07,0}
\definecolor{LightWildStrawberry}{cmyk}{0,0.48,0.20,0}
\definecolor{LightSalmon}        {cmyk}{0,0.27,0.19,0}
\definecolor{LightCarnationPink} {cmyk}{0,0.32,0,0}
\definecolor{LightMagenta}       {cmyk}{0,0.50,0,0}
\definecolor{LightVioletRed}     {cmyk}{0,0.40,0,0}
\definecolor{LightRhodamine}     {cmyk}{0,0.41,0,0}
\definecolor{LightMulberry}      {cmyk}{0.17,0.45,0,0.01}
\definecolor{LightRedViolet}     {cmyk}{0.04,0.45,0,0.17}
\definecolor{LightFuchsia}       {cmyk}{0.24,0.46,0,0.04}
\definecolor{LightLavender}      {cmyk}{0,0.24,0,0}
\definecolor{LightThistle}       {cmyk}{0.06,0.30,0,0}
\definecolor{LightOrchid}        {cmyk}{0.16,0.32,0,0}
\definecolor{LightDarkOrchid}    {cmyk}{0.20,0.40,0.10,0}
\definecolor{LightPurple}        {cmyk}{0.23,0.43,0,0}
\definecolor{LightPlum}          {cmyk}{0.25,0.50,0,0}
\definecolor{LightViolet}        {cmyk}{0.40,0.44,0,0}
\definecolor{LightRoyalPurple}   {cmyk}{0.38,0.45,0,0}
\definecolor{LightBlueViolet}    {cmyk}{0.43,0.46,0,0.02}
\definecolor{LightPeriwinkle}    {cmyk}{0.29,0.28,0,0}
\definecolor{LightCadetBlue}     {cmyk}{0.31,0.29,0.12,0}
\definecolor{LightCornflowerBlue}{cmyk}{0.33,0.07,0,0}
\definecolor{LightMidnightBlue}  {cmyk}{0.49,0.07,0,0.22}
\definecolor{LightNavyBlue}      {cmyk}{0.47,0.27,0,0}
\definecolor{LightRoyalBlue}     {cmyk}{0.50,0.25,0,0}
\definecolor{LightBlue}          {cmyk}{0.50,0.50,0,0}
\definecolor{LightCerulean}      {cmyk}{0.47,0.06,0,0}
\definecolor{LightCyan}          {cmyk}{0.50,0,0,0}
\definecolor{LightProcessBlue}   {cmyk}{0.48,0,0,0}
\definecolor{LightSkyBlue}       {cmyk}{0.31,0,0.06,0}
\definecolor{LightTurquoise}     {cmyk}{0.43,0,0.10,0}
\definecolor{LightTealBlue}      {cmyk}{0.43,0,0.17,0.01}
\definecolor{LightAquamarine}    {cmyk}{0.41,0,0.15,0}
\definecolor{LightBlueGreen}     {cmyk}{0.43,0,0.17,0}
\definecolor{LightEmerald}       {cmyk}{0.50,0,0.25,0}
\definecolor{LightJungleGreen}   {cmyk}{0.50,0,0.26,0}
\definecolor{LightSeaGreen}      {cmyk}{0.35,0,0.25,0}
\definecolor{LightGreen}         {cmyk}{0.50,0,0.50,0}
\definecolor{LightForestGreen}   {cmyk}{0.46,0,0.44,0.06}
\definecolor{LightPineGreen}     {cmyk}{0.46,0,0.30,0.13}
\definecolor{LightLimeGreen}     {cmyk}{0.25,0,0.50,0}
\definecolor{LightYellowGreen}   {cmyk}{0.22,0,0.37,0}
\definecolor{LightSpringGreen}   {cmyk}{0.13,0,0.38,0}
\definecolor{LightOliveGreen}    {cmyk}{0.32,0,0.48,0.20}
\definecolor{LightRawSienna}     {cmyk}{0,0.36,0.50,0.23}
\definecolor{LightSepia}         {cmyk}{0,0.44,0.50,0.35}
\definecolor{LightBrown}         {cmyk}{0,0.41,0.50,0.30}
\definecolor{LightTan}           {cmyk}{0.07,0.21,0.28,0}
\definecolor{LightWhite}         {cmyk}{0,0,0,0}
\definecolor{LightGray0}         {cmyk}{0,0,0,0}
\definecolor{LightGray5}         {cmyk}{0,0,0,0.02}
\definecolor{LightGray10}        {cmyk}{0,0,0,0.05}
\definecolor{LightGray15}        {cmyk}{0,0,0,0.07}
\definecolor{LightGray20}        {cmyk}{0,0,0,0.10}
\definecolor{LightGray25}        {cmyk}{0,0,0,0.12}
\definecolor{LightGray30}        {cmyk}{0,0,0,0.15}
\definecolor{LightLightGray}     {cmyk}{0,0,0,0.16}
\definecolor{LightGray35}        {cmyk}{0,0,0,0.17}
\definecolor{LightGray40}        {cmyk}{0,0,0,0.20}
\definecolor{LightGray45}        {cmyk}{0,0,0,0.22}
\definecolor{LightGray50}        {cmyk}{0,0,0,0.25}
\definecolor{LightGray}          {cmyk}{0,0,0,0.25}
\definecolor{LightGray55}        {cmyk}{0,0,0,0.27}
\definecolor{LightGray60}        {cmyk}{0,0,0,0.30}
\definecolor{LightGray65}        {cmyk}{0,0,0,0.32}
\definecolor{LightDarkGray}      {cmyk}{0,0,0,0.33}
\definecolor{LightGray70}        {cmyk}{0,0,0,0.35}
\definecolor{LightGray75}        {cmyk}{0,0,0,0.37}
\definecolor{LightGray80}        {cmyk}{0,0,0,0.40}
\definecolor{LightGray85}        {cmyk}{0,0,0,0.42}
\definecolor{LightGray90}        {cmyk}{0,0,0,0.45}
\definecolor{LightGray95}        {cmyk}{0,0,0,0.47}
\definecolor{LightBlack}         {cmyk}{0,0,0,0.50}
\definecolor{LightGray100}       {cmyk}{0,0,0,0.50}
\definecolor{LightTC0}           {cmyk}{0.2,0.2,0,0}
\definecolor{LightTC1}           {cmyk}{0.3,0,0.35,0}
\definecolor{LightTC2}           {cmyk}{0.2,0,0.35,0}
\definecolor{LightTC3}           {cmyk}{0.1,0,0.5,0}
\definecolor{LightTC4}           {cmyk}{0,0,0.5,0}
\definecolor{LightTC5}           {cmyk}{0,0.1,0.5,0}
\definecolor{LightTC6}           {cmyk}{0,0.2,0.5,0}
\definecolor{LightTC7}           {cmyk}{0,0.3,0.5,0}
\definecolor{LightTC8}           {cmyk}{0,0.4,0.5,0}
\definecolor{LightTC99}          {cmyk}{0,0,0,0}
\definecolor{LightLightGreenYellow}   {cmyk}{0.04,0,0.17,0}
\definecolor{LightLightYellow}        {cmyk}{0,0,0.25,0}
\definecolor{LightLightGoldenrod}     {cmyk}{0,0.02,0.21,0}
\definecolor{LightLightDandelion}     {cmyk}{0,0.07,0.21,0}
\definecolor{LightLightApricot}       {cmyk}{0,0.08,0.13,0}
\definecolor{LightLightPeach}         {cmyk}{0,0.12,0.17,0}
\definecolor{LightLightMelon}         {cmyk}{0,0.11,0.12,0}
\definecolor{LightLightYellowOrange}  {cmyk}{0,0.10,0.25,0}
\definecolor{LightLightOrange}        {cmyk}{0,0.15,0.22,0}
\definecolor{LightLightBurntOrange}   {cmyk}{0,0.13,0.25,0}
\definecolor{LightLightBittersweet}   {cmyk}{0,0.19,0.25,0.06}
\definecolor{LightLightRedOrange}     {cmyk}{0,0.14,0.22,0}
\definecolor{LightLightMahagony}      {cmyk}{0,0.21,0.22,0.09}
\definecolor{LightLightMaroon}        {cmyk}{0,0.22,0.17,0.08}
\definecolor{LightLightBrickRed}      {cmyk}{0,0.22,0.23,0.07}
\definecolor{LightLightRed}           {cmyk}{0,0.25,0.25,0}
\definecolor{LightLightOrangeRed}     {cmyk}{0,0.25,0.12,0}
\definecolor{LightLightRubineRed}     {cmyk}{0,0.25,0.03,0}
\definecolor{LightLightWildStrawberry}{cmyk}{0,0.24,0.10,0}
\definecolor{LightLightSalmon}        {cmyk}{0,0.13,0.09,0}
\definecolor{LightLightCarnationPink} {cmyk}{0,0.16,0,0}
\definecolor{LightLightMagenta}       {cmyk}{0,0.25,0,0}
\definecolor{LightLightVioletRed}     {cmyk}{0,0.20,0,0}
\definecolor{LightLightRhodamine}     {cmyk}{0,0.20,0,0}
\definecolor{LightLightMulberry}      {cmyk}{0.08,0.22,0,0.005}
\definecolor{LightLightRedViolet}     {cmyk}{0.02,0.22,0,0.08}
\definecolor{LightLightFuchsia}       {cmyk}{0.12,0.23,0,0.02}
\definecolor{LightLightLavender}      {cmyk}{0,0.12,0,0}
\definecolor{LightLightThistle}       {cmyk}{0.03,0.15,0,0}
\definecolor{LightLightOrchid}        {cmyk}{0.08,0.16,0,0}
\definecolor{LightLightDarkOrchid}    {cmyk}{0.10,0.20,0.05,0}
\definecolor{LightLightPurple}        {cmyk}{0.11,0.21,0,0}
\definecolor{LightLightPlum}          {cmyk}{0.12,0.25,0,0}
\definecolor{LightLightViolet}        {cmyk}{0.20,0.22,0,0}
\definecolor{LightLightRoyalPurple}   {cmyk}{0.19,0.22,0,0}
\definecolor{LightLightBlueViolet}    {cmyk}{0.21,0.23,0,0.01}
\definecolor{LightLightPeriwinkle}    {cmyk}{0.14,0.14,0,0}
\definecolor{LightLightCadetBlue}     {cmyk}{0.15,0.14,0.06,0}
\definecolor{LightLightCornflowerBlue}{cmyk}{0.16,0.03,0,0}
\definecolor{LightLightMidnightBlue}  {cmyk}{0.24,0.03,0,0.11}
\definecolor{LightLightNavyBlue}      {cmyk}{0.23,0.13,0,0}
\definecolor{LightLightRoyalBlue}     {cmyk}{0.25,0.12,0,0}
\definecolor{LightLightBlue}          {cmyk}{0.25,0.25,0,0}
\definecolor{LightLightCerulean}      {cmyk}{0.23,0.03,0,0}
\definecolor{LightLightCyan}          {cmyk}{0.25,0,0,0}
\definecolor{LightLightProcessBlue}   {cmyk}{0.24,0,0,0}
\definecolor{LightLightSkyBlue}       {cmyk}{0.15,0,0.03,0}
\definecolor{LightLightTurquoise}     {cmyk}{0.21,0,0.05,0}
\definecolor{LightLightTealBlue}      {cmyk}{0.21,0,0.08,0.005}
\definecolor{LightLightAquamarine}    {cmyk}{0.20,0,0.07,0}
\definecolor{LightLightBlueGreen}     {cmyk}{0.21,0,0.08,0}
\definecolor{LightLightEmerald}       {cmyk}{0.25,0,0.12,0}
\definecolor{LightLightJungleGreen}   {cmyk}{0.25,0,0.13,0}
\definecolor{LightLightSeaGreen}      {cmyk}{0.17,0,0.12,0}
\definecolor{LightLightGreen}         {cmyk}{0.25,0,0.25,0}
\definecolor{LightLightForestGreen}   {cmyk}{0.23,0,0.22,0.03}
\definecolor{LightLightPineGreen}     {cmyk}{0.23,0,0.15,0.06}
\definecolor{LightLightLimeGreen}     {cmyk}{0.12,0,0.25,0}
\definecolor{LightLightYellowGreen}   {cmyk}{0.11,0,0.18,0}
\definecolor{LightLightSpringGreen}   {cmyk}{0.06,0,0.19,0}
\definecolor{LightLightOliveGreen}    {cmyk}{0.16,0,0.24,0.10}
\definecolor{LightLightRawSienna}     {cmyk}{0,0.18,0.25,0.11}
\definecolor{LightLightSepia}         {cmyk}{0,0.22,0.25,0.17}
\definecolor{LightLightBrown}         {cmyk}{0,0.20,0.25,0.15}
\definecolor{LightLightTan}           {cmyk}{0.03,0.10,0.14,0}
\definecolor{LightLightWhite}         {cmyk}{0,0,0,0}
\definecolor{LightLightGray0}         {cmyk}{0,0,0,0}
\definecolor{LightLightGray5}         {cmyk}{0,0,0,0.01}
\definecolor{LightLightGray10}        {cmyk}{0,0,0,0.02}
\definecolor{LightLightGray15}        {cmyk}{0,0,0,0.03}
\definecolor{LightLightGray20}        {cmyk}{0,0,0,0.05}
\definecolor{LightLightGray25}        {cmyk}{0,0,0,0.06}
\definecolor{LightLightGray30}        {cmyk}{0,0,0,0.07}
\definecolor{LightLightLightGray}     {cmyk}{0,0,0,0.08}
\definecolor{LightLightGray35}        {cmyk}{0,0,0,0.09}
\definecolor{LightLightGray40}        {cmyk}{0,0,0,0.10}
\definecolor{LightLightGray45}        {cmyk}{0,0,0,0.11}
\definecolor{LightLightGray50}        {cmyk}{0,0,0,0.12}
\definecolor{LightLightGray}          {cmyk}{0,0,0,0.13}
\definecolor{LightLightGray55}        {cmyk}{0,0,0,0.14}
\definecolor{LightLightGray60}        {cmyk}{0,0,0,0.15}
\definecolor{LightLightGray65}        {cmyk}{0,0,0,0.16}
\definecolor{LightLightDarkGray}      {cmyk}{0,0,0,0.17}
\definecolor{LightLightGray70}        {cmyk}{0,0,0,0.18}
\definecolor{LightLightGray75}        {cmyk}{0,0,0,0.19}
\definecolor{LightLightGray80}        {cmyk}{0,0,0,0.20}
\definecolor{LightLightGray85}        {cmyk}{0,0,0,0.21}
\definecolor{LightLightGray90}        {cmyk}{0,0,0,0.22}
\definecolor{LightLightGray95}        {cmyk}{0,0,0,0.23}
\definecolor{LightLightBlack}         {cmyk}{0,0,0,0.25}
\definecolor{LightLightGray100}       {cmyk}{0,0,0,0.25}
\definecolor{LightLightTC0}           {cmyk}{0.1,0.1,0,0}
\definecolor{LightLightTC1}           {cmyk}{0.15,0,0.175,0}
\definecolor{LightLightTC2}           {cmyk}{0.1,0,0.175,0}
\definecolor{LightLightTC3}           {cmyk}{0.05,0,0.25,0}
\definecolor{LightLightTC4}           {cmyk}{0,0,0.25,0}
\definecolor{LightLightTC5}           {cmyk}{0,0.05,0.25,0}
\definecolor{LightLightTC6}           {cmyk}{0,0.1,0.25,0}
\definecolor{LightLightTC7}           {cmyk}{0,0.15,0.25,0}
\definecolor{LightLightTC8}           {cmyk}{0,0.2,0.25,0}
\definecolor{LightLightTC99}          {cmyk}{0,0,0,0}
\definecolor{LightLightLightGreenYellow}   {cmyk}{0.02,0,0.08,0}
\definecolor{LightLightLightYellow}        {cmyk}{0,0,0.12,0}
\definecolor{LightLightLightGoldenrod}     {cmyk}{0,0.01,0.10,0}
\definecolor{LightLightLightDandelion}     {cmyk}{0,0.03,0.10,0}
\definecolor{LightLightLightApricot}       {cmyk}{0,0.04,0.06,0}
\definecolor{LightLightLightPeach}         {cmyk}{0,0.06,0.08,0}
\definecolor{LightLightLightMelon}         {cmyk}{0,0.05,0.06,0}
\definecolor{LightLightLightYellowOrange}  {cmyk}{0,0.05,0.12,0}
\definecolor{LightLightLightOrange}        {cmyk}{0,0.07,0.11,0}
\definecolor{LightLightLightBurntOrange}   {cmyk}{0,0.06,0.12,0}
\definecolor{LightLightLightBittersweet}   {cmyk}{0,0.09,0.12,0.03}
\definecolor{LightLightLightRedOrange}     {cmyk}{0,0.07,0.11,0}
\definecolor{LightLightLightMahagony}      {cmyk}{0,0.10,0.11,0.04}
\definecolor{LightLightLightMaroon}        {cmyk}{0,0.11,0.08,0.04}
\definecolor{LightLightLightBrickRed}      {cmyk}{0,0.11,0.11,0.03}
\definecolor{LightLightLightRed}           {cmyk}{0,0.12,0.12,0}
\definecolor{LightLightLightOrangeRed}     {cmyk}{0,0.12,0.06,0}
\definecolor{LightLightLightRubineRed}     {cmyk}{0,0.12,0.01,0}
\definecolor{LightLightLightWildStrawberry}{cmyk}{0,0.12,0.05,0}
\definecolor{LightLightLightSalmon}        {cmyk}{0,0.06,0.04,0}
\definecolor{LightLightLightCarnationPink} {cmyk}{0,0.08,0,0}
\definecolor{LightLightLightMagenta}       {cmyk}{0,0.12,0,0}
\definecolor{LightLightLightLightMagenta}  {cmyk}{0,0.06,0,0}
\definecolor{LightLightLightVioletRed}     {cmyk}{0,0.10,0,0}
\definecolor{LightLightLightRhodamine}     {cmyk}{0,0.10,0,0}
\definecolor{LightLightLightMulberry}      {cmyk}{0.04,0.11,0,0.002}
\definecolor{LightLightLightRedViolet}     {cmyk}{0.01,0.11,0,0.04}
\definecolor{LightLightLightFuchsia}       {cmyk}{0.06,0.11,0,0.01}
\definecolor{LightLightLightLavender}      {cmyk}{0,0.06,0,0}
\definecolor{LightLightLightThistle}       {cmyk}{0.01,0.07,0,0}
\definecolor{LightLightLightOrchid}        {cmyk}{0.04,0.08,0,0}
\definecolor{LightLightLightDarkOrchid}    {cmyk}{0.05,0.10,0.02,0}
\definecolor{LightLightLightPurple}        {cmyk}{0.05,0.10,0,0}
\definecolor{LightLightLightPlum}          {cmyk}{0.06,0.12,0,0}
\definecolor{LightLightLightViolet}        {cmyk}{0.10,0.11,0,0}
\definecolor{LightLightLightRoyalPurple}   {cmyk}{0.09,0.11,0,0}
\definecolor{LightLightLightBlueViolet}    {cmyk}{0.10,0.11,0,0.005}
\definecolor{LightLightLightPeriwinkle}    {cmyk}{0.07,0.07,0,0}
\definecolor{LightLightLightCadetBlue}     {cmyk}{0.07,0.07,0.03,0}
\definecolor{LightLightLightCornflowerBlue}{cmyk}{0.08,0.01,0,0}
\definecolor{LightLightLightMidnightBlue}  {cmyk}{0.12,0.01,0,0.05}
\definecolor{LightLightLightNavyBlue}      {cmyk}{0.11,0.06,0,0}
\definecolor{LightLightLightRoyalBlue}     {cmyk}{0.12,0.06,0,0}
\definecolor{LightLightLightBlue}          {cmyk}{0.12,0.12,0,0}
\definecolor{LightLightLightCerulean}      {cmyk}{0.11,0.01,0,0}
\definecolor{LightLightLightCyan}          {cmyk}{0.12,0,0,0}
\definecolor{LightLightLightProcessBlue}   {cmyk}{0.12,0,0,0}
\definecolor{LightLightLightSkyBlue}       {cmyk}{0.07,0,0.01,0}
\definecolor{LightLightLightTurquoise}     {cmyk}{0.10,0,0.02,0}
\definecolor{LightLightLightTealBlue}      {cmyk}{0.10,0,0.04,0.002}
\definecolor{LightLightLightAquamarine}    {cmyk}{0.10,0,0.03,0}
\definecolor{LightLightLightBlueGreen}     {cmyk}{0.10,0,0.04,0}
\definecolor{LightLightLightEmerald}       {cmyk}{0.12,0,0.06,0}
\definecolor{LightLightLightJungleGreen}   {cmyk}{0.12,0,0.06,0}
\definecolor{LightLightLightSeaGreen}      {cmyk}{0.08,0,0.06,0}
\definecolor{LightLightLightGreen}         {cmyk}{0.12,0,0.12,0}
\definecolor{LightLightLightForestGreen}   {cmyk}{0.11,0,0.11,0.01}
\definecolor{LightLightLightPineGreen}     {cmyk}{0.11,0,0.07,0.03}
\definecolor{LightLightLightLimeGreen}     {cmyk}{0.06,0,0.12,0}
\definecolor{LightLightLightYellowGreen}   {cmyk}{0.05,0,0.09,0}
\definecolor{LightLightLightSpringGreen}   {cmyk}{0.03,0,0.09,0}
\definecolor{LightLightLightOliveGreen}    {cmyk}{0.08,0,0.12,0.05}
\definecolor{LightLightLightRawSienna}     {cmyk}{0,0.09,0.12,0.05}
\definecolor{LightLightLightSepia}         {cmyk}{0,0.11,0.12,0.06}
\definecolor{LightLightLightBrown}         {cmyk}{0,0.10,0.12,0.07}
\definecolor{LightLightLightTan}           {cmyk}{0.01,0.05,0.07,0}
\definecolor{LightLightLightWhite}         {cmyk}{0,0,0,0}
\definecolor{LightLightLightGray0}         {cmyk}{0,0,0,0}
\definecolor{LightLightLightGray5}         {cmyk}{0,0,0,0.005}
\definecolor{LightLightLightGray10}        {cmyk}{0,0,0,0.01}
\definecolor{LightLightLightGray15}        {cmyk}{0,0,0,0.015}
\definecolor{LightLightLightGray20}        {cmyk}{0,0,0,0.025}
\definecolor{LightLightLightGray25}        {cmyk}{0,0,0,0.03}
\definecolor{LightLightLightGray30}        {cmyk}{0,0,0,0.035}
\definecolor{LightLightLightLightGray}     {cmyk}{0,0,0,0.04}
\definecolor{LightLightLightGray35}        {cmyk}{0,0,0,0.045}
\definecolor{LightLightLightGray40}        {cmyk}{0,0,0,0.05}
\definecolor{LightLightLightGray45}        {cmyk}{0,0,0,0.055}
\definecolor{LightLightLightGray50}        {cmyk}{0,0,0,0.06}
\definecolor{LightLightLightGray}          {cmyk}{0,0,0,0.065}
\definecolor{LightLightLightGray55}        {cmyk}{0,0,0,0.07}
\definecolor{LightLightLightGray60}        {cmyk}{0,0,0,0.075}
\definecolor{LightLightLightGray65}        {cmyk}{0,0,0,0.08}
\definecolor{LightLightLightDarkGray}      {cmyk}{0,0,0,0.085}
\definecolor{LightLightLightGray70}        {cmyk}{0,0,0,0.09}
\definecolor{LightLightLightGray75}        {cmyk}{0,0,0,0.095}
\definecolor{LightLightLightGray80}        {cmyk}{0,0,0,0.10}
\definecolor{LightLightLightGray85}        {cmyk}{0,0,0,0.105}
\definecolor{LightLightLightGray90}        {cmyk}{0,0,0,0.11}
\definecolor{LightLightLightGray95}        {cmyk}{0,0,0,0.115}
\definecolor{LightLightLightBlack}         {cmyk}{0,0,0,0.12}
\definecolor{LightLightLightGray100}       {cmyk}{0,0,0,0.125}
\definecolor{LightLightLightTC0}           {cmyk}{0.05,0.05,0,0}
\definecolor{LightLightLightTC1}           {cmyk}{0.075,0,0.086,0}
\definecolor{LightLightLightTC2}           {cmyk}{0.05,0,0.086,0}
\definecolor{LightLightLightTC3}           {cmyk}{0.025,0,0.125,0}
\definecolor{LightLightLightTC4}           {cmyk}{0,0,0.125,0}
\definecolor{LightLightLightTC5}           {cmyk}{0,0.025,0.125,0}
\definecolor{LightLightLightTC6}           {cmyk}{0,0.05,0.125,0}
\definecolor{LightLightLightTC7}           {cmyk}{0,0.0175,0.125,0}
\definecolor{LightLightLightTC8}           {cmyk}{0,0.1,0.125,0}
\definecolor{LightLightLightTC99}          {cmyk}{0,0,0,0}
\definecolor{BlueRed5}           {rgb} {0.15,0.17,0.55}
\definecolor{BlueRed10}          {rgb} {0.20,0.23,0.57}
\definecolor{BlueRed15}          {rgb} {0.24,0.29,0.60}
\definecolor{BlueRed20}          {rgb} {0.33,0.35,0.64}
\definecolor{BlueRed25}          {rgb} {0.43,0.43,0.68}
\definecolor{BlueRed30}          {rgb} {0.52,0.52,0.73}
\definecolor{BlueRed35}          {rgb} {0.60,0.60,0.78}
\definecolor{BlueRed40}          {rgb} {0.70,0.70,0.84}
\definecolor{BlueRed45}          {rgb} {0.80,0.80,0.85}
\definecolor{BlueRed50}          {rgb} {0.86,0.82,0.82}
\definecolor{BlueRed55}          {rgb} {0.87,0.73,0.73}
\definecolor{BlueRed60}          {rgb} {0.89,0.64,0.64}
\definecolor{BlueRed65}          {rgb} {0.90,0.55,0.55}
\definecolor{BlueRed70}          {rgb} {0.91,0.47,0.46}
\definecolor{BlueRed75}          {rgb} {0.91,0.39,0.37}
\definecolor{BlueRed80}          {rgb} {0.90,0.33,0.28}
\definecolor{BlueRed85}          {rgb} {0.89,0.25,0.20}
\definecolor{BlueRed90}          {rgb} {0.88,0.23,0.14}
\definecolor{BlueRed95}          {rgb} {0.87,0.21,0.09}
\definecolor{BlueRed100}         {rgb} {0.87,0.16,0.04}
\definecolor{RedBlue100}         {rgb} {0.15,0.17,0.55}
\definecolor{RedBlue95}          {rgb} {0.15,0.17,0.55}
\definecolor{RedBlue90}          {rgb} {0.20,0.23,0.57}
\definecolor{RedBlue85}          {rgb} {0.24,0.29,0.60}
\definecolor{RedBlue80}          {rgb} {0.33,0.35,0.64}
\definecolor{RedBlue75}          {rgb} {0.43,0.43,0.68}
\definecolor{RedBlue70}          {rgb} {0.52,0.52,0.73}
\definecolor{RedBlue65}          {rgb} {0.60,0.60,0.78}
\definecolor{RedBlue60}          {rgb} {0.70,0.70,0.84}
\definecolor{RedBlue55}          {rgb} {0.80,0.80,0.85}
\definecolor{RedBlue50}          {rgb} {0.86,0.82,0.82}
\definecolor{RedBlue45}          {rgb} {0.87,0.73,0.73}
\definecolor{RedBlue40}          {rgb} {0.89,0.64,0.64}
\definecolor{RedBlue35}          {rgb} {0.90,0.55,0.55}
\definecolor{RedBlue30}          {rgb} {0.91,0.47,0.46}
\definecolor{RedBlue25}          {rgb} {0.91,0.39,0.37}
\definecolor{RedBlue20}          {rgb} {0.90,0.33,0.28}
\definecolor{RedBlue15}          {rgb} {0.89,0.25,0.20}
\definecolor{RedBlue10}          {rgb} {0.88,0.23,0.14}
\definecolor{RedBlue5}           {rgb} {0.87,0.21,0.09}
\definecolor{GreenRed5}          {rgb} {0,1,0}
\definecolor{GreenRed10}         {rgb} {0.05,0.95,0}
\definecolor{GreenRed15}         {rgb} {0.10,0.90,0}
\definecolor{GreenRed20}         {rgb} {0.15,0.85,0}
\definecolor{GreenRed25}         {rgb} {0.20,0.80,0}
\definecolor{GreenRed30}         {rgb} {0.25,0.75,0}
\definecolor{GreenRed35}         {rgb} {0.30,0.70,0}
\definecolor{GreenRed40}         {rgb} {0.35,0.65,0}
\definecolor{GreenRed45}         {rgb} {0.40,0.60,0}
\definecolor{GreenRed50}         {rgb} {0.45,0.55,0}
\definecolor{GreenRed55}         {rgb} {0.50,0.50,0}
\definecolor{GreenRed60}         {rgb} {0.55,0.45,0}
\definecolor{GreenRed65}         {rgb} {0.60,0.40,0}
\definecolor{GreenRed70}         {rgb} {0.65,0.35,0}
\definecolor{GreenRed75}         {rgb} {0.70,0.30,0}
\definecolor{GreenRed80}         {rgb} {0.75,0.25,0}
\definecolor{GreenRed85}         {rgb} {0.80,0.20,0}
\definecolor{GreenRed90}         {rgb} {0.85,0.15,0}
\definecolor{GreenRed95}         {rgb} {0.90,0.10,0}
\definecolor{GreenRed100}        {rgb} {0.95,0.05,0}
\definecolor{RedGreen100}        {rgb} {0.05,0.95,0}
\definecolor{RedGreen95}         {rgb} {0.10,0.90,0}
\definecolor{RedGreen90}         {rgb} {0.15,0.85,0}
\definecolor{RedGreen85}         {rgb} {0.20,0.80,0}
\definecolor{RedGreen80}         {rgb} {0.25,0.75,0}
\definecolor{RedGreen75}         {rgb} {0.30,0.70,0}
\definecolor{RedGreen70}         {rgb} {0.35,0.65,0}
\definecolor{RedGreen65}         {rgb} {0.40,0.60,0}
\definecolor{RedGreen60}         {rgb} {0.45,0.55,0}
\definecolor{RedGreen55}         {rgb} {0.50,0.50,0}
\definecolor{RedGreen50}         {rgb} {0.55,0.45,0}
\definecolor{RedGreen45}         {rgb} {0.60,0.40,0}
\definecolor{RedGreen40}         {rgb} {0.65,0.35,0}
\definecolor{RedGreen35}         {rgb} {0.70,0.30,0}
\definecolor{RedGreen30}         {rgb} {0.75,0.25,0}
\definecolor{RedGreen25}         {rgb} {0.80,0.20,0}
\definecolor{RedGreen20}         {rgb} {0.85,0.15,0}
\definecolor{RedGreen15}         {rgb} {0.90,0.10,0}
\definecolor{RedGreen10}         {rgb} {0.95,0.05,0}
\definecolor{RedGreen5}          {rgb} {1,0,0}
\definecolor{ColdHot5}           {rgb} {0,0.08,1}
\definecolor{ColdHot10}          {rgb} {0,0.29,1}
\definecolor{ColdHot15}          {rgb} {0,0.49,1}
\definecolor{ColdHot20}          {rgb} {0,0.70,1}
\definecolor{ColdHot25}          {rgb} {0,0.90,1}
\definecolor{ColdHot30}          {rgb} {0,1,0.87}
\definecolor{ColdHot35}          {rgb} {0,1,0.68}
\definecolor{ColdHot40}          {rgb} {0,1,0.46}
\definecolor{ColdHot45}          {rgb} {0,1,0.25}
\definecolor{ColdHot50}          {rgb} {0,1,0.04}
\definecolor{ColdHot55}          {rgb} {0.16,1,0}
\definecolor{ColdHot60}          {rgb} {0.35,1,0}
\definecolor{ColdHot65}          {rgb} {0.56,1,0}
\definecolor{ColdHot70}          {rgb} {0.79,1,0}
\definecolor{ColdHot75}          {rgb} {0.98,1,0}
\definecolor{ColdHot80}          {rgb} {1,0.82,0}
\definecolor{ColdHot85}          {rgb} {1,0.60,0}
\definecolor{ColdHot90}          {rgb} {1,0.40,0}
\definecolor{ColdHot95}          {rgb} {1,0.20,0}
\definecolor{ColdHot100}         {rgb} {0.91,0,0}
\definecolor{HotCold100}         {rgb} {0,0.08,1}
\definecolor{HotCold95}          {rgb} {0,0.29,1}
\definecolor{HotCold90}          {rgb} {0,0.49,1}
\definecolor{HotCold85}          {rgb} {0,0.70,1}
\definecolor{HotCold80}          {rgb} {0,0.90,1}
\definecolor{HotCold75}          {rgb} {0,1,0.87}
\definecolor{HotCold70}          {rgb} {0,1,0.68}
\definecolor{HotCold65}          {rgb} {0,1,0.46}
\definecolor{HotCold60}          {rgb} {0,1,0.25}
\definecolor{HotCold55}          {rgb} {0,1,0.04}
\definecolor{HotCold50}          {rgb} {0.16,1,0}
\definecolor{HotCold45}          {rgb} {0.35,1,0}
\definecolor{HotCold40}          {rgb} {0.56,1,0}
\definecolor{HotCold35}          {rgb} {0.79,1,0}
\definecolor{HotCold30}          {rgb} {0.98,1,0}
\definecolor{HotCold25}          {rgb} {1,0.82,0}
\definecolor{HotCold20}          {rgb} {1,0.60,0}
\definecolor{HotCold15}          {rgb} {1,0.40,0}
\definecolor{HotCold10}          {rgb} {1,0.20,0}
\definecolor{HotCold5}           {rgb} {0.91,0,0}
\expandafter\def\csname BlueRed5\endcsname{[0.15,0.17,0.55]}
\expandafter\def\csname BlueRed10\endcsname{[0.20,0.23,0.57]}
\expandafter\def\csname BlueRed15\endcsname{[0.24,0.29,0.60]}
\expandafter\def\csname BlueRed20\endcsname{[0.33,0.35,0.64]}
\expandafter\def\csname BlueRed25\endcsname{[0.43,0.43,0.68]}
\expandafter\def\csname BlueRed30\endcsname{[0.52,0.52,0.73]}
\expandafter\def\csname BlueRed35\endcsname{[0.60,0.60,0.78]}
\expandafter\def\csname BlueRed40\endcsname{[0.70,0.70,0.84]}
\expandafter\def\csname BlueRed45\endcsname{[0.80,0.80,0.85]}
\expandafter\def\csname BlueRed50\endcsname{[0.86,0.82,0.82]}
\expandafter\def\csname BlueRed55\endcsname{[0.87,0.73,0.73]}
\expandafter\def\csname BlueRed60\endcsname{[0.89,0.64,0.64]}
\expandafter\def\csname BlueRed65\endcsname{[0.90,0.55,0.55]}
\expandafter\def\csname BlueRed70\endcsname{[0.91,0.47,0.46]}
\expandafter\def\csname BlueRed75\endcsname{[0.91,0.39,0.37]}
\expandafter\def\csname BlueRed80\endcsname{[0.90,0.33,0.28]}
\expandafter\def\csname BlueRed85\endcsname{[0.89,0.25,0.20]}
\expandafter\def\csname BlueRed90\endcsname{[0.88,0.23,0.14]}
\expandafter\def\csname BlueRed95\endcsname{[0.87,0.21,0.09]}
\expandafter\def\csname BlueRed100\endcsname{[0.87,0.16,0.04]}
\expandafter\def\csname RedBlue100\endcsname{[0.15,0.17,0.55]}
\expandafter\def\csname RedBlue95\endcsname{[0.15,0.17,0.55]}
\expandafter\def\csname RedBlue90\endcsname{[0.20,0.23,0.57]}
\expandafter\def\csname RedBlue85\endcsname{[0.24,0.29,0.60]}
\expandafter\def\csname RedBlue80\endcsname{[0.33,0.35,0.64]}
\expandafter\def\csname RedBlue75\endcsname{[0.43,0.43,0.68]}
\expandafter\def\csname RedBlue70\endcsname{[0.52,0.52,0.73]}
\expandafter\def\csname RedBlue65\endcsname{[0.60,0.60,0.78]}
\expandafter\def\csname RedBlue60\endcsname{[0.70,0.70,0.84]}
\expandafter\def\csname RedBlue55\endcsname{[0.80,0.80,0.85]}
\expandafter\def\csname RedBlue50\endcsname{[0.86,0.82,0.82]}
\expandafter\def\csname RedBlue45\endcsname{[0.87,0.73,0.73]}
\expandafter\def\csname RedBlue40\endcsname{[0.89,0.64,0.64]}
\expandafter\def\csname RedBlue35\endcsname{[0.90,0.55,0.55]}
\expandafter\def\csname RedBlue30\endcsname{[0.91,0.47,0.46]}
\expandafter\def\csname RedBlue25\endcsname{[0.91,0.39,0.37]}
\expandafter\def\csname RedBlue20\endcsname{[0.90,0.33,0.28]}
\expandafter\def\csname RedBlue15\endcsname{[0.89,0.25,0.20]}
\expandafter\def\csname RedBlue10\endcsname{[0.88,0.23,0.14]}
\expandafter\def\csname RedBlue5\endcsname{[0.87,0.21,0.09]}
\expandafter\def\csname GreenRed5\endcsname{[0,1,0]}
\expandafter\def\csname GreenRed10\endcsname{[0.05,0.95,0]}
\expandafter\def\csname GreenRed15\endcsname{[0.10,0.90,0]}
\expandafter\def\csname GreenRed20\endcsname{[0.15,0.85,0]}
\expandafter\def\csname GreenRed25\endcsname{[0.20,0.80,0]}
\expandafter\def\csname GreenRed30\endcsname{[0.25,0.75,0]}
\expandafter\def\csname GreenRed35\endcsname{[0.30,0.70,0]}
\expandafter\def\csname GreenRed40\endcsname{[0.35,0.65,0]}
\expandafter\def\csname GreenRed45\endcsname{[0.40,0.60,0]}
\expandafter\def\csname GreenRed50\endcsname{[0.45,0.55,0]}
\expandafter\def\csname GreenRed55\endcsname{[0.50,0.50,0]}
\expandafter\def\csname GreenRed60\endcsname{[0.55,0.45,0]}
\expandafter\def\csname GreenRed65\endcsname{[0.60,0.40,0]}
\expandafter\def\csname GreenRed70\endcsname{[0.65,0.35,0]}
\expandafter\def\csname GreenRed75\endcsname{[0.70,0.30,0]}
\expandafter\def\csname GreenRed80\endcsname{[0.75,0.25,0]}
\expandafter\def\csname GreenRed85\endcsname{[0.80,0.20,0]}
\expandafter\def\csname GreenRed90\endcsname{[0.85,0.15,0]}
\expandafter\def\csname GreenRed95\endcsname{[0.90,0.10,0]}
\expandafter\def\csname GreenRed100\endcsname{[0.95,0.05,0]}
\expandafter\def\csname RedGreen100\endcsname{[0.05,0.95,0]}
\expandafter\def\csname RedGreen95\endcsname{[0.10,0.90,0]}
\expandafter\def\csname RedGreen90\endcsname{[0.15,0.85,0]}
\expandafter\def\csname RedGreen85\endcsname{[0.20,0.80,0]}
\expandafter\def\csname RedGreen80\endcsname{[0.25,0.75,0]}
\expandafter\def\csname RedGreen75\endcsname{[0.30,0.70,0]}
\expandafter\def\csname RedGreen70\endcsname{[0.35,0.65,0]}
\expandafter\def\csname RedGreen65\endcsname{[0.40,0.60,0]}
\expandafter\def\csname RedGreen60\endcsname{[0.45,0.55,0]}
\expandafter\def\csname RedGreen55\endcsname{[0.50,0.50,0]}
\expandafter\def\csname RedGreen50\endcsname{[0.55,0.45,0]}
\expandafter\def\csname RedGreen45\endcsname{[0.60,0.40,0]}
\expandafter\def\csname RedGreen40\endcsname{[0.65,0.35,0]}
\expandafter\def\csname RedGreen35\endcsname{[0.70,0.30,0]}
\expandafter\def\csname RedGreen30\endcsname{[0.75,0.25,0]}
\expandafter\def\csname RedGreen25\endcsname{[0.80,0.20,0]}
\expandafter\def\csname RedGreen20\endcsname{[0.85,0.15,0]}
\expandafter\def\csname RedGreen15\endcsname{[0.90,0.10,0]}
\expandafter\def\csname RedGreen10\endcsname{[0.95,0.05,0]}
\expandafter\def\csname RedGreen5\endcsname{[1,0,0]}
\expandafter\def\csname ColdHot5\endcsname{[0,0.08,1]}
\expandafter\def\csname ColdHot10\endcsname{[0,0.29,1]}
\expandafter\def\csname ColdHot15\endcsname{[0,0.49,1]}
\expandafter\def\csname ColdHot20\endcsname{[0,0.70,1]}
\expandafter\def\csname ColdHot25\endcsname{[0,0.90,1]}
\expandafter\def\csname ColdHot30\endcsname{[0,1,0.87]}
\expandafter\def\csname ColdHot35\endcsname{[0,1,0.68]}
\expandafter\def\csname ColdHot40\endcsname{[0,1,0.46]}
\expandafter\def\csname ColdHot45\endcsname{[0,1,0.25]}
\expandafter\def\csname ColdHot50\endcsname{[0,1,0.04]}
\expandafter\def\csname ColdHot55\endcsname{[0.16,1,0]}
\expandafter\def\csname ColdHot60\endcsname{[0.35,1,0]}
\expandafter\def\csname ColdHot65\endcsname{[0.56,1,0]}
\expandafter\def\csname ColdHot70\endcsname{[0.79,1,0]}
\expandafter\def\csname ColdHot75\endcsname{[0.98,1,0]}
\expandafter\def\csname ColdHot80\endcsname{[1,0.82,0]}
\expandafter\def\csname ColdHot85\endcsname{[1,0.60,0]}
\expandafter\def\csname ColdHot90\endcsname{[1,0.40,0]}
\expandafter\def\csname ColdHot95\endcsname{[1,0.20,0]}
\expandafter\def\csname ColdHot100\endcsname{[0.91,0,0]}
\expandafter\def\csname HotCold100\endcsname{[0,0.08,1]}
\expandafter\def\csname HotCold95\endcsname{[0,0.29,1]}
\expandafter\def\csname HotCold90\endcsname{[0,0.49,1]}
\expandafter\def\csname HotCold85\endcsname{[0,0.70,1]}
\expandafter\def\csname HotCold80\endcsname{[0,0.90,1]}
\expandafter\def\csname HotCold75\endcsname{[0,1,0.87]}
\expandafter\def\csname HotCold70\endcsname{[0,1,0.68]}
\expandafter\def\csname HotCold65\endcsname{[0,1,0.46]}
\expandafter\def\csname HotCold60\endcsname{[0,1,0.25]}
\expandafter\def\csname HotCold55\endcsname{[0,1,0.04]}
\expandafter\def\csname HotCold50\endcsname{[0.16,1,0]}
\expandafter\def\csname HotCold45\endcsname{[0.35,1,0]}
\expandafter\def\csname HotCold40\endcsname{[0.56,1,0]}
\expandafter\def\csname HotCold35\endcsname{[0.79,1,0]}
\expandafter\def\csname HotCold30\endcsname{[0.98,1,0]}
\expandafter\def\csname HotCold25\endcsname{[1,0.82,0]}
\expandafter\def\csname HotCold20\endcsname{[1,0.60,0]}
\expandafter\def\csname HotCold15\endcsname{[1,0.40,0]}
\expandafter\def\csname HotCold10\endcsname{[1,0.20,0]}
\expandafter\def\csname HotCold5\endcsname{[0.91,0,0]}


\def\make@lower{%
\if\first@ A\xdef\first@{a}\else \if\first@ B\xdef\first@{b}\else
\if\first@ C\xdef\first@{c}\else \if\first@ D\xdef\first@{d}\else
\if\first@ E\xdef\first@{e}\else \if\first@ F\xdef\first@{f}\else
\if\first@ G\xdef\first@{g}\else \if\first@ H\xdef\first@{h}\else
\if\first@ I\xdef\first@{i}\else \if\first@ J\xdef\first@{j}\else
\if\first@ K\xdef\first@{k}\else \if\first@ L\xdef\first@{l}\else
\if\first@ M\xdef\first@{m}\else \if\first@ N\xdef\first@{n}\else
\if\first@ O\xdef\first@{o}\else \if\first@ P\xdef\first@{p}\else
\if\first@ Q\xdef\first@{q}\else \if\first@ R\xdef\first@{r}\else
\if\first@ S\xdef\first@{s}\else \if\first@ T\xdef\first@{t}\else
\if\first@ U\xdef\first@{u}\else \if\first@ V\xdef\first@{v}\else
\if\first@ W\xdef\first@{w}\else \if\first@ X\xdef\first@{x}\else
\if\first@ Y\xdef\first@{y}\else \if\first@ Z\xdef\first@{z}\else
\fi\fi\fi\fi\fi\fi\fi\fi\fi\fi\fi\fi\fi\fi\fi\fi\fi\fi\fi\fi\fi\fi\fi\fi\fi\fi}

\def\make@upper{%
\if\first@ a\xdef\first@{A}\else \if\first@ b\xdef\first@{B}\else
\if\first@ c\xdef\first@{C}\else \if\first@ d\xdef\first@{D}\else
\if\first@ e\xdef\first@{E}\else \if\first@ f\xdef\first@{F}\else
\if\first@ g\xdef\first@{G}\else \if\first@ h\xdef\first@{H}\else
\if\first@ i\xdef\first@{I}\else \if\first@ j\xdef\first@{J}\else
\if\first@ k\xdef\first@{K}\else \if\first@ l\xdef\first@{L}\else
\if\first@ m\xdef\first@{M}\else \if\first@ n\xdef\first@{N}\else
\if\first@ o\xdef\first@{O}\else \if\first@ p\xdef\first@{P}\else
\if\first@ q\xdef\first@{Q}\else \if\first@ r\xdef\first@{R}\else
\if\first@ s\xdef\first@{S}\else \if\first@ t\xdef\first@{T}\else
\if\first@ u\xdef\first@{U}\else \if\first@ v\xdef\first@{V}\else
\if\first@ w\xdef\first@{W}\else \if\first@ x\xdef\first@{X}\else
\if\first@ y\xdef\first@{Y}\else \if\first@ z\xdef\first@{Z}\else
\fi\fi\fi\fi\fi\fi\fi\fi\fi\fi\fi\fi\fi\fi\fi\fi\fi\fi\fi\fi\fi\fi\fi\fi\fi\fi}

\def\type@get#1 MSF: #2 Type: #3 #4@{\seqtype{#3}}
\def\inf@@get#1 #2 #3 #4 #5 #6@{%
   \def\first@{#1}
   \def\second@{#2}
   \def\third@{#3}
   \xdef\fourth@{#4 @}
   \expandafter\check@letter\fourth@
   \ifnumber 
      \def\fourth@{#4} 
   \else 
      \xdef\fourth@{#5 @}
      \expandafter\check@letter\fourth@
      \ifnumber 
         \def\fourth@{#5} 
      \else
         \def\fourth@{99999999}
      \fi
   \fi
   \def\fifth@{#5}
   \def\last@{#6}}
\def\check@char#1{\letterfalse \ifnum\catcode`#1=11 \lettertrue \fi
                  \numberfalse \ifnum\catcode`#1=12 \numbertrue \fi
                  \xdef\code@num{\the\catcode`#1}
                  \xdef\char@num{`#1}}
\def\check@letter#1#2@{\letterfalse\ifnum\catcode`#1=11 \lettertrue\fi%
                       \numberfalse\ifnum\catcode`#1=12 \numbertrue\fi}
\def\seq@get#1 #2@{\def\first@{#1} \def\seq@line{#2}}
\def\res@get#1#2@{\def\first@{#1} \def\seq@line{#2&@}}
\def\tot@get#1#2@{%
   \xdef\first@{#1}
   \expandafter\xdef\csname seq\the\triple@count\endcsname{#2&@}}
\def\dis@get#1#2@{%
   \xdef\first@{#1}
   \expandafter\xdef\csname res\the\triple@count\endcsname{#1}
   \expandafter\xdef\csname seq\the\triple@count\endcsname{#2&@}}
\def\firstchar@get#1#2@{\def\first@{#1}\def\third@{#2@}}
\def\residue@get#1#2@{\xdef\first@{#1}
                      \xdef\char@num{`#1}
                      \expandafter\xdef\csname res\the\loopcount\endcsname{\first@}
                      \expandafter\xdef\csname sequence\the\loopcount\endcsname{#2@}}
\def\remove@fromseq#1@{\expandafter\xdef\csname sequence\the\loopcount\endcsname{#1}}
\def\get@temp#1@{\xdef\temp@{#1}}
\def\TC@get#1#2@{\xdef\char@num{`#1}
                 \ifnum\char@num=45 \xdef\TC@first@{*}
                 \else
                       \ifnum\char@num>57 \xdef\TC@first@{*}
                       \else
                          \xdef\TC@first@{#1}
                       \fi
                 \fi
                 \expandafter\xdef\csname TC@num\the\loopcount\endcsname{\TC@first@}
                 \expandafter\xdef\csname T@coffee\the\loopcount\endcsname{#2@}}
\def\sublogo@get#1 @{\xdef\first@{#1}}
\def\re@write#1,#2@{%
   \xdef\third@{#2,@}
   \xdef\first@{\csname hide@@@seq#1\endcsname}
   \ifx\first@\last@
   \else
      \xdef\first@{#1&}
      \ifx\first@\ampers@nd \else \advance\innerloopcount by 1 \fi
      \if\second@ e \xdef\second@{\csname @rd#1\endcsname}
      \else \xdef\second@{\second@,\csname @rd#1\endcsname}\fi
   \fi
   }
\def\order@set#1,#2@{%
   \xdef\second@{#2,@}
      \expandafter\xdef\csname @rd#1\endcsname{\the\loopcount}
      \expandafter\xdef\csname res@count\the\loopcount\endcsname{%
                    \csname pos#1\endcsname}
      \expandafter\xdef\csname hide@seq\the\loopcount\endcsname{%
                    \csname hide@@seq#1\endcsname}
      \expandafter\xdef\csname hide@name\the\loopcount\endcsname{%
                    \csname hide@@name#1\endcsname}
      \expandafter\xdef\csname hide@number\the\loopcount\endcsname{%
                    \csname hide@@number#1\endcsname}
      \expandafter\xdef\csname seq@start\the\loopcount\endcsname{%
                    \csname seq@@start#1\endcsname}
      \expandafter\xdef\csname seq@len\the\loopcount\endcsname{%
                    \csname seq@@len#1\endcsname}
      \expandafter\xdef\csname seqname\the\loopcount\endcsname{%
                    \csname seq@name#1\endcsname}
      \expandafter\xdef\csname newseqname\the\loopcount\endcsname{%
                    \csname newseq@name#1\endcsname}
      \expandafter\xdef\csname seq@gap\the\loopcount\endcsname{%
                    \csname seq@@gap#1\endcsname}
      \expandafter\xdef\csname name@col\the\loopcount\endcsname{%
                    \csname name@@col#1\endcsname}
      \expandafter\xdef\csname number@col\the\loopcount\endcsname{%
                    \csname number@@col#1\endcsname}
      \expandafter\xdef\csname stack@reg\the\loopcount\endcsname{%
                    \csname stack@@reg#1\endcsname}
      \expandafter\xdef\csname stack@tintreg\the\loopcount\endcsname{%
                    \csname stack@@tintreg#1\endcsname}
      \expandafter\xdef\csname stack@emphreg\the\loopcount\endcsname{%
                    \csname stack@@emphreg#1\endcsname}
      \expandafter\xdef\csname stack@lowerreg\the\loopcount\endcsname{%
                    \csname stack@@lowerreg#1\endcsname}
      \expandafter\xdef\csname stack@framereg\the\loopcount\endcsname{%
                    \csname stack@@framereg#1\endcsname}
      \expandafter\xdef\csname stack@shadingreg\the\loopcount\endcsname{%
                    \csname stack@@shadingreg#1\endcsname}
      \expandafter\xdef\csname stack@top\the\loopcount\endcsname{%
                    \csname stack@@top#1\endcsname}
      \expandafter\xdef\csname stack@ttop\the\loopcount\endcsname{%
                    \csname stack@@ttop#1\endcsname}
      \expandafter\xdef\csname stack@tttop\the\loopcount\endcsname{%
                    \csname stack@@tttop#1\endcsname}
      \expandafter\xdef\csname stack@ttttop\the\loopcount\endcsname{%
                    \csname stack@@ttttop#1\endcsname}
      \expandafter\xdef\csname stack@bottom\the\loopcount\endcsname{%
                    \csname stack@@bottom#1\endcsname}
      \expandafter\xdef\csname stack@bbottom\the\loopcount\endcsname{%
                    \csname stack@@bbottom#1\endcsname}
      \expandafter\xdef\csname stack@bbbottom\the\loopcount\endcsname{%
                    \csname stack@@bbbottom#1\endcsname}
      \expandafter\xdef\csname stack@bbbbottom\the\loopcount\endcsname{%
                    \csname stack@@bbbbottom#1\endcsname}
}
\def\reorder@seqs#1{%
   \loopcount=0
   \loop
      \advance\loopcount by 1
         \expandafter\xdef\csname pos\the\loopcount\endcsname{%
                 \csname res@count\the\loopcount\endcsname}
         \expandafter\xdef\csname hide@@seq\the\loopcount\endcsname{%
            \csname hide@seq\the\loopcount\endcsname}
         \expandafter\xdef\csname hide@@name\the\loopcount\endcsname{%
            \csname hide@name\the\loopcount\endcsname}
         \expandafter\xdef\csname hide@@number\the\loopcount\endcsname{%
            \csname hide@number\the\loopcount\endcsname}
         \expandafter\xdef\csname seq@@start\the\loopcount\endcsname{%
            \csname seq@start\the\loopcount\endcsname}
         \expandafter\xdef\csname seq@@len\the\loopcount\endcsname{%
            \csname seq@len\the\loopcount\endcsname}
         \expandafter\xdef\csname seq@name\the\loopcount\endcsname{%
            \csname seqname\the\loopcount\endcsname}
         \expandafter\xdef\csname newseq@name\the\loopcount\endcsname{%
            \csname newseqname\the\loopcount\endcsname}
         \expandafter\xdef\csname seq@@gap\the\loopcount\endcsname{%
            \csname seq@gap\the\loopcount\endcsname}
         \expandafter\xdef\csname name@@col\the\loopcount\endcsname{%
            \csname name@col\the\loopcount\endcsname}
         \expandafter\xdef\csname number@@col\the\loopcount\endcsname{%
            \csname number@col\the\loopcount\endcsname}
         \expandafter\xdef\csname stack@@reg\the\loopcount\endcsname{%
            \csname stack@reg\the\loopcount\endcsname}
         \expandafter\xdef\csname stack@@tintreg\the\loopcount\endcsname{%
            \csname stack@tintreg\the\loopcount\endcsname}
         \expandafter\xdef\csname stack@@emphreg\the\loopcount\endcsname{%
            \csname stack@emphreg\the\loopcount\endcsname}
         \expandafter\xdef\csname stack@@lowerreg\the\loopcount\endcsname{%
            \csname stack@lowerreg\the\loopcount\endcsname}
         \expandafter\xdef\csname stack@@framereg\the\loopcount\endcsname{%
            \csname stack@framereg\the\loopcount\endcsname}
         \expandafter\xdef\csname stack@@shadingreg\the\loopcount\endcsname{%
            \csname stack@shadingreg\the\loopcount\endcsname}
         \expandafter\xdef\csname stack@@top\the\loopcount\endcsname{%
            \csname stack@top\the\loopcount\endcsname}
         \expandafter\xdef\csname stack@@ttop\the\loopcount\endcsname{%
            \csname stack@ttop\the\loopcount\endcsname}
         \expandafter\xdef\csname stack@@tttop\the\loopcount\endcsname{%
            \csname stack@tttop\the\loopcount\endcsname}
         \expandafter\xdef\csname stack@@ttttop\the\loopcount\endcsname{%
            \csname stack@ttttop\the\loopcount\endcsname}
         \expandafter\xdef\csname stack@@bottom\the\loopcount\endcsname{%
            \csname stack@bottom\the\loopcount\endcsname}
         \expandafter\xdef\csname stack@@bbottom\the\loopcount\endcsname{%
            \csname stack@bbottom\the\loopcount\endcsname}
         \expandafter\xdef\csname stack@@bbbottom\the\loopcount\endcsname{%
            \csname stack@bbbottom\the\loopcount\endcsname}
         \expandafter\xdef\csname stack@@bbbbottom\the\loopcount\endcsname{%
            \csname stack@bbbbottom\the\loopcount\endcsname}
   \ifnum\loopcount<\killseq@count \repeat
   \xdef\third@{#1}
   \loopcount=0 \innerloopcount=0 \xdef\last@{kill} \xdef\second@{e}
   \loop
      \advance\loopcount by 1
      \expandafter\re@write\third@
   \ifnum\loopcount<\seq@count \repeat
   \ifnum\innerloopcount<\killseq@count
       \@latex@error{Not enough sequences specified in `orderseqs'
       (\the\innerloopcount/\the\killseq@count)}\@ehc
   \fi
   \loopcount=0 \xdef\second@{\second@,@}
   \loop
      \advance\loopcount by 1
      \expandafter\order@set\second@
   \ifnum\loopcount<\killseq@count \repeat
   \ifnum\cons@num>0 \xdef\cons@num{\csname @rd\cons@num\endcsname} \fi
   \ifnum\rule@num>0 \xdef\rule@num{\csname @rd\rule@num\endcsname} \fi
   \ifnum\divref@>0  \xdef\divref@{\csname @rd\divref@\endcsname} \fi
   \loopcount=0
   \loop
      \advance\loopcount by 1
      \expandafter\xdef\csname pos\the\loopcount\endcsname{0}
   \ifnum\loopcount<\killseq@count \repeat
}
\def\group@get#1,#2@{%
   \def\group@set{\expandafter\residue@get\second@
                 \ifnum\char@num>96 \make@upper \fi
                 \ifx\first@\ampers@nd
                 \else \expandafter\xdef\csname \prfx grp\first@\endcsname{\the\loopcount}
                 \xdef\second@{\csname sequence\the\loopcount\endcsname} \group@set
                 \fi}
   \xdef\second@{#1 &@} \xdef\third@{#2&,@} \group@set}
\def\get@item#1,#2@{\xdef\first@{#2@}\xdef\first@@{#2}\xdef\fourth@{#1}}
\def\get@first@@#1-#2@{\xdef\first@@{#1}}
\def\get@digit#1,#2@{%
   \def\check@series##1-##2##3@{%
                    \xdef\first@@{##2}\xdef\fourth@{##1}\xdef\fourth@@{##3}}
   \xdef\first@{#2@}
   \xdef\fourth@{#1-&@}
   \expandafter\check@series\fourth@
   \ifx\first@@\ampers@nd
   \else
      \xdef\first@@{\first@@\fourth@@-&@}
      \expandafter\get@first@@\first@@
      \loopcount=\fourth@ 
      \ifnum\first@@>\fourth@
         \advance\loopcount by 1
         \xdef\first@{\the\loopcount-\first@@,#2@}
      \else
         \ifnum\first@@<\fourth@
            \advance\loopcount by -1
            \xdef\first@{\the\loopcount-\first@@,#2@}
         \fi
      \fi
   \fi
}
\def\donot@shade{%
   \expandafter\check@letter\first@
   \ifletter
      \expandafter\get@item\first@
      \xdef\first@{\fourth@}
      \get@name@number
      \xdef\third@{\csname hide@seq\first@\endcsname} \xdef\second@{kill}
      \ifx\third@\second@
      \else
         \expandafter\xdef\csname hide@seq\first@\endcsname{noshade}
      \fi
      \xdef\first@{\first@@ @}
      \donot@shade
   \else
      \expandafter\get@digit\first@
      \ifx\fourth@\ampers@nd
      \else
         \xdef\third@{\csname hide@seq\fourth@\endcsname} \xdef\second@{kill}
         \ifx\third@\second@
         \else
            \expandafter\xdef\csname hide@seq\fourth@\endcsname{noshade}
         \fi
        \donot@shade
      \fi
   \fi}
\def\clear@res@nums#1{%
   \temp@count=65
   \loop
      \xdef\first@@{\csname ch@r@\the\temp@count\endcsname}
      \expandafter\xdef\csname res@num\first@@ #1\endcsname{0}
      \advance\temp@count by 1
   \ifnum\temp@count>90\else\repeat
   \expandafter\xdef\csname res@num\d@t #1\endcsname{0}   
   \expandafter\xdef\csname res@num\questi@n #1\endcsname{0}   
}
\def\setsubfamily@{%
   \expandafter\check@letter\first@
   \ifletter
      \expandafter\get@item\first@
      \xdef\first@{\fourth@}
      \get@name@number
      \advance\res@count by 1 
      \expandafter\xdef\csname subfamily@num\first@\endcsname{\subfamily@count}
      \xdef\first@{\first@@ @}
      \setsubfamily@
   \else
      \expandafter\get@digit\first@
      \ifx\fourth@\ampers@nd
      \else
         \advance\res@count by 1 
         \expandafter\xdef\csname subfamily@num\fourth@\endcsname{\subfamily@count}
         \setsubfamily@
      \fi
   \fi
}
\def\hideseq@{%
   \expandafter\check@letter\first@
   \ifletter
      \expandafter\get@item\first@
      \xdef\first@{\fourth@}
      \get@name@number
      \xdef\third@{\csname hide@seq\first@\endcsname} \xdef\second@{kill}
      \ifx\third@\second@
      \else
         \expandafter\xdef\csname hide@seq\first@\endcsname{true}
      \fi
      \xdef\first@{\first@@ @}
      \hideseq@
   \else
      \expandafter\get@digit\first@
      \ifx\fourth@\ampers@nd
      \else
         \xdef\third@{\csname hide@seq\fourth@\endcsname} \xdef\second@{kill}
         \ifx\third@\second@
         \else
            \expandafter\xdef\csname hide@seq\fourth@\endcsname{true}
         \fi
        \hideseq@
      \fi
   \fi}
\def\hidename@{%
   \expandafter\check@letter\first@
   \ifletter
      \expandafter\get@item\first@
      \xdef\first@{\fourth@}
      \get@name@number
      \expandafter\xdef\csname hide@name\first@\endcsname{yes}
      \xdef\first@{\first@@ @}
      \hidename@
   \else
      \expandafter\get@digit\first@
      \ifx\fourth@\ampers@nd
      \else
         \expandafter\xdef\csname hide@name\fourth@\endcsname{yes}
         \hidename@
      \fi
   \fi}
\def\hidenumber@{%
   \expandafter\check@letter\first@
   \ifletter
      \expandafter\get@item\first@
      \xdef\first@{\fourth@}
      \get@name@number
      \expandafter\xdef\csname hide@number\first@\endcsname{yes}
      \xdef\first@{\first@@ @}
      \hidenumber@
   \else
      \expandafter\get@digit\first@
      \ifx\fourth@\ampers@nd
      \else
         \expandafter\xdef\csname hide@number\fourth@\endcsname{yes}
         \hidenumber@
      \fi
   \fi}
\def\namecolor@{%
   \expandafter\check@letter\first@
   \ifletter
      \expandafter\get@item\first@
      \xdef\first@{\fourth@}
      \get@name@number
      \expandafter\xdef\csname name@col\first@\endcsname{\third@}
      \xdef\first@{\first@@ @}
      \namecolor@
   \else
      \expandafter\get@digit\first@
      \ifx\fourth@\ampers@nd
      \else
         \expandafter\xdef\csname name@col\fourth@\endcsname{\third@}
         \namecolor@
      \fi
   \fi}
\def\numbercolor@{%
   \expandafter\check@letter\first@
   \ifletter
      \expandafter\get@item\first@
      \xdef\first@{\fourth@}
      \get@name@number
      \expandafter\xdef\csname number@col\first@\endcsname{\third@}
      \xdef\first@{\first@@ @}
      \numbercolor@
   \else
      \expandafter\get@digit\first@
      \ifx\fourth@\ampers@nd
      \else
         \expandafter\xdef\csname number@col\fourth@\endcsname{\third@}
         \numbercolor@
      \fi
   \fi}
\def\killseq@{%
   \expandafter\check@letter\first@
   \ifletter
      \expandafter\get@item\first@
      \xdef\first@{\fourth@}
      \get@name@number
      \ifnum\first@>\seq@count
      \else
      \ifnum\killseq@count>1
         \xdef\third@{\csname hide@seq\first@\endcsname} \xdef\second@{kill}
         \ifx\third@\second@
         \else
            \expandafter\xdef\csname hide@seq\first@\endcsname{kill}
            \expandafter\xdef\csname hide@@@seq\first@\endcsname{kill}
            \ifnum\first@=\cons@num \xdef\cons@num{0} \fi
            \advance\killseq@count by -1
%            \seq@percent=100 \divide\seq@percent by \killseq@count
         \fi
      \fi\fi
      \xdef\first@{\first@@ @}
      \killseq@
   \else
      \expandafter\get@digit\first@
      \ifx\fourth@\ampers@nd
      \else
      \ifnum\fourth@>\seq@count
      \else
      \ifnum\killseq@count>1
         \xdef\third@{\csname hide@seq\fourth@\endcsname} \xdef\second@{kill}
         \ifx\third@\second@
         \else
            \expandafter\xdef\csname hide@seq\fourth@\endcsname{kill}
            \expandafter\xdef\csname hide@@@seq\fourth@\endcsname{kill}
            \ifnum\fourth@=\cons@num \xdef\cons@num{0} \fi
            \advance\killseq@count by -1
%            \seq@percent=100 \divide\seq@percent by \killseq@count
         \fi
      \fi\fi
      \killseq@
      \fi
   \fi}
\def\kill@loop{%
   \advance\innerloopcount by 1
   \xdef\first@{\csname hide@seq\the\innerloopcount\endcsname}
   \ifx\first@\second@ \kill@loop \fi}
\def\kill@seqnow{%
   \xdef\first@{\csname hide@seq\the\loopcount\endcsname} \xdef\second@{kill}
   \innerloopcount=\loopcount
   \ifx\first@\second@
      \kill@loop
      \expandafter\xdef\csname hide@seq\the\loopcount\endcsname{%
            \csname hide@seq\the\innerloopcount\endcsname}
      \expandafter\xdef\csname seqname\the\loopcount\endcsname{%
            \csname seqname\the\innerloopcount\endcsname}
      \expandafter\xdef\csname newseqname\the\loopcount\endcsname{%
            \csname newseqname\the\innerloopcount\endcsname}
      \expandafter\xdef\csname seq@gap\the\loopcount\endcsname{%
            \csname seq@gap\the\innerloopcount\endcsname}
      \expandafter\xdef\csname stack@reg\the\loopcount\endcsname{%
            \csname stack@reg\the\innerloopcount\endcsname}
      \expandafter\xdef\csname stack@tintreg\the\loopcount\endcsname{%
            \csname stack@tintreg\the\innerloopcount\endcsname}
      \expandafter\xdef\csname stack@emphreg\the\loopcount\endcsname{%
            \csname stack@emphreg\the\innerloopcount\endcsname}
      \expandafter\xdef\csname stack@lowerreg\the\loopcount\endcsname{%
            \csname stack@lowerreg\the\innerloopcount\endcsname}
      \expandafter\xdef\csname stack@framereg\the\loopcount\endcsname{%
            \csname stack@framereg\the\innerloopcount\endcsname}
      \expandafter\xdef\csname stack@shadingreg\the\loopcount\endcsname{%
            \csname stack@shadingreg\the\innerloopcount\endcsname}
      \expandafter\xdef\csname stack@top\the\loopcount\endcsname{%
            \csname stack@top\the\innerloopcount\endcsname}
      \expandafter\xdef\csname stack@ttop\the\loopcount\endcsname{%
            \csname stack@ttop\the\innerloopcount\endcsname}
      \expandafter\xdef\csname stack@tttop\the\loopcount\endcsname{%
            \csname stack@tttop\the\innerloopcount\endcsname}
      \expandafter\xdef\csname stack@ttttop\the\loopcount\endcsname{%
            \csname stack@ttttop\the\innerloopcount\endcsname}
      \expandafter\xdef\csname stack@bottom\the\loopcount\endcsname{%
            \csname stack@bottom\the\innerloopcount\endcsname}
      \expandafter\xdef\csname stack@bbottom\the\loopcount\endcsname{%
            \csname stack@bbottom\the\innerloopcount\endcsname}
      \expandafter\xdef\csname stack@bbbottom\the\loopcount\endcsname{%
            \csname stack@bbbottom\the\innerloopcount\endcsname}
      \expandafter\xdef\csname stack@bbbbottom\the\loopcount\endcsname{%
            \csname stack@bbbbottom\the\innerloopcount\endcsname}
      \expandafter\xdef\csname res@count\the\loopcount\endcsname{%
            \csname res@count\the\innerloopcount\endcsname}
      \expandafter\xdef\csname hide@seq\the\innerloopcount\endcsname{kill}
      \ifnum\loopcount=\rule@num \hideruler \fi
   \fi
   \expandafter\xdef\csname @rd\the\innerloopcount\endcsname{\the\loopcount}
   \advance\loopcount by 1
   \ifnum\loopcount>\killseq@count 
      \ifnum\rule@num>0\xdef\rule@num{\csname @rd\rule@num\endcsname}\fi
      \ifnum\cons@num>0\xdef\cons@num{\csname @rd\cons@num\endcsname}\fi
   \else 
      \kill@seqnow
   \fi}
\def\get@sim#1#2@{\xdef\sim@char{#1} \xdef\last@{#2 &@}}
\def\getsim@char{%
   \advance\innerloopcount by 1
   \ifnum\innerloopcount>\m@x
   \else
      \expandafter\get@sim\last@
      \ifx\second@\sim@char \xdef\third@{2} \innerloopcount=\m@x \fi
      \getsim@char
   \fi}
\def\get@count(#1)#2@{\xdef\last@{#2 &@} \xdef\m@x{#1}}
\def\get@nums#1..#2@{\xdef\first@{#1} \xdef\second@{#2}}
\def\func@shading#1{%
   \clearfuncgroups
   \xdef\temp@{#1}
   \xdef\second@{charge}
   \ifx\temp@\second@
      \ifgerm@n
      \funcgroup{sauer ($-$)}{DE}{White}{Red}{upper}{up}
      \funcgroup{basisch ($+$)}{KRH}{White}{Blue}{upper}{up}
      \else
      \ifsp@nish
      \funcgroup{\'acido ($-$)}{DE}{White}{Red}{upper}{up}
      \funcgroup{b\'asico ($+$)}{KRH}{White}{Blue}{upper}{up}
      \else
      \funcgroup{acidic ($-$)}{DE}{White}{Red}{upper}{up}
      \funcgroup{basic ($+$)}{KRH}{White}{Blue}{upper}{up}
      \fi
      \fi
   \else
   \xdef\second@{hydropathy}
   \ifx\temp@\second@
      \ifgerm@n
      \funcgroup{sauer ($-$)}{DE}{White}{Red}{upper}{up}
      \funcgroup{basisch ($+$)}{KRH}{White}{Blue}{upper}{up}
      \funcgroup{polar ungeladen}{YSTGNQC}{Black}{Yellow}{upper}{up}
      \funcgroup{hydrophob unpolar}{AFPMWVIL}{White}{Green}{upper}{up}
      \else
      \ifsp@nish
      \funcgroup{\'acido ($-$)}{DE}{White}{Red}{upper}{up}
      \funcgroup{b\'asico ($+$)}{KRH}{White}{Blue}{upper}{up}
      \funcgroup{polar sin carga}{YSTGNQC}{Black}{Yellow}{upper}{up}
      \funcgroup{hidrof\'obico no polar}{AFPMWVIL}{White}{Green}{upper}{up}
      \else
      \funcgroup{acidic ($-$)}{DE}{White}{Red}{upper}{up}
      \funcgroup{basic ($+$)}{KRH}{White}{Blue}{upper}{up}
      \funcgroup{polar uncharged}{YSTGNQC}{Black}{Yellow}{upper}{up}
      \funcgroup{hydrophobic nonpolar}{AFPMWVIL}{White}{Green}{upper}{up}
      \fi
      \fi
   \else
   \xdef\second@{chemical}
   \ifx\temp@\second@
      \ifgerm@n
      \funcgroup{sauer ($-$)}{DE}{White}{Red}{upper}{up}
      \funcgroup{aliphatisch}{AGVIL}{White}{Black}{upper}{up}
      \funcgroup{aliphatisch (klein)}{AG}{White}{Gray}{upper}{up}
      \funcgroup{Amid}{NQ}{White}{Green}{upper}{up}
      \funcgroup{aromatisch}{FYW}{White}{Brown}{upper}{up}
      \funcgroup{basisch (+)}{KRH}{White}{Blue}{upper}{up}
      \funcgroup{Hydroxyl}{ST}{Black}{Magenta}{upper}{up}
      \funcgroup{Imin}{P}{Black}{Orange}{upper}{up}
      \funcgroup{Schwefel}{CM}{Black}{Yellow}{upper}{up}
      \else
      \ifsp@nish
      \funcgroup{\'acido ($-$)}{DE}{White}{Red}{upper}{up}
      \funcgroup{alif\'atico}{AGVIL}{White}{Black}{upper}{up}
      \funcgroup{alif\'atico (parvo)}{AG}{White}{Gray}{upper}{up}
      \funcgroup{amida}{NQ}{White}{Green}{upper}{up}
      \funcgroup{arom\'atico}{FYW}{White}{Brown}{upper}{up}
      \funcgroup{b\'asico ($+$)}{KRH}{White}{Blue}{upper}{up}
      \funcgroup{hidr\'oxido}{ST}{Black}{Magenta}{upper}{up}
      \funcgroup{imino}{P}{Black}{Orange}{upper}{up}
      \funcgroup{azufre}{CM}{Black}{Yellow}{upper}{up}
      \else
      \funcgroup{acidic ($-$)}{DE}{White}{Red}{upper}{up}
      \funcgroup{aliphatic}{AGVIL}{White}{Black}{upper}{up}
      \funcgroup{aliphatic (small)}{AG}{White}{Gray}{upper}{up}
      \funcgroup{amide}{NQ}{White}{Green}{upper}{up}
      \funcgroup{aromatic}{FYW}{White}{Brown}{upper}{up}
      \funcgroup{basic ($+$)}{KRH}{White}{Blue}{upper}{up}
      \funcgroup{hydroxyl}{ST}{Black}{Magenta}{upper}{up}
      \funcgroup{imino}{P}{Black}{Orange}{upper}{up}
      \funcgroup{sulfur}{CM}{Black}{Yellow}{upper}{up}
      \fi
      \fi
   \else
   \xdef\second@{structure}
   \ifx\temp@\second@
      \ifgerm@n
      \funcgroup{extern}{DEHKNQR}{Black}{Orange}{upper}{up}
      \funcgroup{ambivalent}{ACGPSTWY}{Black}{Yellow}{upper}{up}
      \funcgroup{intern}{FILMV}{White}{Green}{upper}{up}
      \else
      \ifsp@nish
      \funcgroup{externo}{DEHKNQR}{Black}{Orange}{upper}{up}
      \funcgroup{ambivalente}{ACGPSTWY}{Black}{Yellow}{upper}{up}
      \funcgroup{interno}{FILMV}{White}{Green}{upper}{up}
      \else
      \funcgroup{external}{DEHKNQR}{Black}{Orange}{upper}{up}
      \funcgroup{ambivalent}{ACGPSTWY}{Black}{Yellow}{upper}{up}
      \funcgroup{internal}{FILMV}{White}{Green}{upper}{up}
      \fi
      \fi
   \else
   \xdef\second@{standard area}
   \ifx\temp@\second@
      \ifgerm@n
      \funcgroup{\ 88,1 (G); Standard Seitenkettenfl\"ache %
                             (\AA$^2$)}%
                {G}{Black}{BrickRed}{upper}{up}
      \funcgroup{118,2 (A); 129,8 (S)}{AS}{Black}{Orange}{upper}{up}
      \funcgroup{146,1 (C); 146,8 (P)}%
                {CP}{Black}{Yellow}{upper}{up}
      \funcgroup{152,5 (T); 158,7 (D); 164,5 (V); 165,5 (N)}%
                {TDVN}{Black}{YellowGreen}{upper}{up}
      \funcgroup{181,0 (I); 186,2 (E)}{IE}{White}{PineGreen}{upper}{up}
      \funcgroup{193,1 (L); 193,2 (Q); 202,5 (H); 203,3 (M)}%
                {LQHM}{Black}{SkyBlue}{upper}{up}
      \funcgroup{222,8 (F); 225,8 (K)}{FK}{White}{RoyalPurple}{upper}{up}
      \funcgroup{238,8 (Y)}{Y}{White}{RedViolet}{upper}{up}
      \funcgroup{256,0 (R); 266,2 (W)}{RW}{White}{Black}{upper}{up}
      \else
      \ifsp@nish
      \funcgroup{\ 88.1 (G); Superficie est\'andar de la cadena lateral (\AA$^2$)}%
                {G}{Black}{BrickRed}{upper}{up}
      \funcgroup{118.2 (A); 129.8 (S)}{AS}{Black}{Orange}{upper}{up}
      \funcgroup{146.1 (C); 146.8 (P)}%
                {CP}{Black}{Yellow}{upper}{up}
      \funcgroup{152.5 (T); 158.7 (D); 164.5 (V); 165.5 (N)}%
                {TDVN}{Black}{YellowGreen}{upper}{up}
      \funcgroup{181.0 (I); 186.2 (E)}{IE}{White}{PineGreen}{upper}{up}
      \funcgroup{193.1 (L); 193.2 (Q); 202.5 (H); 203.3 (M)}%
                {LQHM}{Black}{SkyBlue}{upper}{up}
      \funcgroup{222.8 (F); 225.8 (K)}{FK}{White}{RoyalPurple}{upper}{up}
      \funcgroup{238.8 (Y)}{Y}{White}{RedViolet}{upper}{up}
      \funcgroup{256.0 (R); 266.2 (W)}{RW}{White}{Black}{upper}{up}
      \else
      \funcgroup{\ 88.1 (G); Standard sidechain area (\AA$^2$)}%
                {G}{Black}{BrickRed}{upper}{up}
      \funcgroup{118.2 (A); 129.8 (S)}{AS}{Black}{Orange}{upper}{up}
      \funcgroup{146.1 (C); 146.8 (P)}%
                {CP}{Black}{Yellow}{upper}{up}
      \funcgroup{152.5 (T); 158.7 (D); 164.5 (V); 165.5 (N)}%
                {TDVN}{Black}{YellowGreen}{upper}{up}
      \funcgroup{181.0 (I); 186.2 (E)}{IE}{White}{PineGreen}{upper}{up}
      \funcgroup{193.1 (L); 193.2 (Q); 202.5 (H); 203.3 (M)}%
                {LQHM}{Black}{SkyBlue}{upper}{up}
      \funcgroup{222.8 (F); 225.8 (K)}{FK}{White}{RoyalPurple}{upper}{up}
      \funcgroup{238.8 (Y)}{Y}{White}{RedViolet}{upper}{up}
      \funcgroup{256.0 (R); 266.2 (W)}{RW}{White}{Black}{upper}{up}
      \fi
      \fi
   \else
   \xdef\second@{accessible area}
   \ifx\temp@\second@
      \ifgerm@n
      \funcgroup{\ 13,9 (C); Zug\"angliche Seitenkettenfl\"ache %
                            (\AA$^2$)}%
                {CIV}{Black}{BrickRed}{upper}{up}
      \funcgroup{\ 23,0 (I); 23,5 (V); 25,2 (G)}%
                {IVG}{Black}{Orange}{upper}{up}
      \funcgroup{\ 28,7 (F); 29,0 (L); 30,5 (M); 31,5 (A)}%
                {FLMA}{Black}{Yellow}{upper}{up}
      \funcgroup{\ 41,7 (W); 44,2 (S); 46,0 (T); 46,7 (H)}%
                {WSTH}{Black}{YellowGreen}{upper}{up}
      \funcgroup{\ 53,7 (P)}{P}{White}{PineGreen}{upper}{up}
      \funcgroup{\ 59,1 (Y); 60,9 (D); 62,2 (N)}%
                {YDN}{Black}{SkyBlue}{upper}{up}
      \funcgroup{\ 72,3 (E); 74,0 (Q)}{EQ}{White}{RoyalPurple}{upper}{up}
      \funcgroup{\ 93,8 (R)}{R}{White}{RedViolet}{upper}{up}
      \funcgroup{110,3 (K)}{K}{White}{Black}{upper}{up}
      \else
      \ifsp@nish
      \funcgroup{\ 13.9 (C); Superficie accesible de la cadena lateral (\AA$^2$)}%
                {CIV}{Black}{BrickRed}{upper}{up}
      \funcgroup{\ 23.0 (I); 23.5 (V); 25.2 (G)}%
                {IVG}{Black}{Orange}{upper}{up}
      \funcgroup{\ 28.7 (F); 29.0 (L); 30.5 (M); 31.5 (A)}%
                {FLMA}{Black}{Yellow}{upper}{up}
      \funcgroup{\ 41.7 (W); 44.2 (S); 46.0 (T); 46.7 (H)}%
                {WSTH}{Black}{YellowGreen}{upper}{up}
      \funcgroup{\ 53.7 (P)}{P}{White}{PineGreen}{upper}{up}
      \funcgroup{\ 59.1 (Y); 60.9 (D); 62.2 (N)}%
                {YDN}{Black}{SkyBlue}{upper}{up}
      \funcgroup{\ 72.3 (E); 74.0 (Q)}{EQ}{White}{RoyalPurple}{upper}{up}
      \funcgroup{\ 93.8 (R)}{R}{White}{RedViolet}{upper}{up}
      \funcgroup{110.3 (K)}{K}{White}{Black}{upper}{up}
      \else
      \funcgroup{\ 13.9 (C); Accessible sidechain area (\AA$^2$)}%
                {CIV}{Black}{BrickRed}{upper}{up}
      \funcgroup{\ 23.0 (I); 23.5 (V); 25.2 (G)}%
                {IVG}{Black}{Orange}{upper}{up}
      \funcgroup{\ 28.7 (F); 29.0 (L); 30.5 (M); 31.5 (A)}%
                {FLMA}{Black}{Yellow}{upper}{up}
      \funcgroup{\ 41.7 (W); 44.2 (S); 46.0 (T); 46.7 (H)}%
                {WSTH}{Black}{YellowGreen}{upper}{up}
      \funcgroup{\ 53.7 (P)}{P}{White}{PineGreen}{upper}{up}
      \funcgroup{\ 59.1 (Y); 60.9 (D); 62.2 (N)}%
                {YDN}{Black}{SkyBlue}{upper}{up}
      \funcgroup{\ 72.3 (E); 74.0 (Q)}{EQ}{White}{RoyalPurple}{upper}{up}
      \funcgroup{\ 93.8 (R)}{R}{White}{RedViolet}{upper}{up}
      \funcgroup{110.3 (K)}{K}{White}{Black}{upper}{up}
      \fi
      \fi
   \else
   \xdef\second@{rasmol}
   \ifx\temp@\second@
      \funcgroup{Asp, Glu}{DE}{Red}{White}{upper}{up}
      \funcgroup{Arg, Lys, His}{KRH}{Blue}{White}{upper}{up}
      \funcgroup{Phe, Tyr, Trp}{FYW}{MidnightBlue}{White}{upper}{up}
      \funcgroup{Ala, Gly}{AG}{Gray}{White}{upper}{up}
      \funcgroup{Cys, Met}{CM}{Yellow}{White}{upper}{up}
      \funcgroup{Ser, Thr}{ST}{Orange}{White}{upper}{up}
      \funcgroup{Asn, Gln}{NQ}{Cyan}{White}{upper}{up}
      \funcgroup{Leu, Val, Ile}{LVI}{Green}{White}{upper}{up}
      \funcgroup{Pro}{P}{Apricot}{White}{upper}{up}
   \else
   \xdef\second@{hardness}
   \ifx\temp@\second@
      \funcgroup{22-23 eV (ADEGILPQSTV) | 2$\times$ HOMO-LUMO gap}{ADEGILPQSTV}{White}{BlueRed5}{upper}{up}
      \funcgroup{21-22 eV (KN)}{KN}{White}{BlueRed20}{upper}{up}
      \funcgroup{20-21 eV (R)}{R}{Black}{BlueRed40}{upper}{up}
      \funcgroup{19-20 eV (CFH)}{CFH}{Black}{BlueRed60}{upper}{up}
      \funcgroup{18-19 eV (MY)}{MY}{White}{BlueRed80}{upper}{up}
      \funcgroup{17-18 eV (W)}{W}{White}{BlueRed100}{upper}{up}
   \else
   \message{<Unknown shading mode - clearing `funcgroups'>}
   \fi\fi\fi\fi\fi\fi\fi\fi
}
\def\shadeallresidues{\all@fshadetrue}
\def\get@fromstack#1;#2;#3;#4;#5@{%
   \xdef\first@{#1} \xdef\second@{#2}
   \xdef\third@{#3} \xdef\fourth@{#4} \xdef\last@{#5@}
}
\def\getregion@fromstack#1{%
   \xdef\temp@{#1}
   \xdef\first@{\csname stack@reg\temp@\endcsname}
   \expandafter\get@fromstack\first@
   \expandafter\xdef\csname style\temp@\endcsname{\first@}
   \expandafter\xdef\csname start\temp@\endcsname{\second@}
   \expandafter\xdef\csname stop\temp@\endcsname{\third@}
   \expandafter\xdef\csname all\temp@\endcsname{\fourth@}
   \expandafter\xdef\csname stack@reg\temp@\endcsname{\last@}
}
\def\sort@stack{%
   \expandafter\get@fromstack\last@ 
   \ifx\first@\ampers@nd  
      \xdef\tmpstack{\tmpstack\the\loopcount;\st@rt;\st@p;\@ll;&;&;&;&;@}
   \else 
   \ifnum\st@rt<\second@
      \xdef\tmpstack{\tmpstack\the\loopcount;\st@rt;\st@p;\@ll;%
                     \first@;\second@;\third@;\fourth@;\last@} 
   \else
      \xdef\tmpstack{\tmpstack\first@;\second@;\third@;\fourth@;} 
      \sort@stack
   \fi\fi
}
\def\get@regions#1..#2,#3@{%
   \xdef\st@rt{#1} \xdef\st@p{#2} \xdef\list@{#3}
   \xdef\last@{\csname stack@reg\seq@\endcsname}
   \xdef\tmpstack{}
   \sort@stack
   \expandafter\xdef\csname stack@reg\seq@\endcsname{\tmpstack}
}
\def\get@fromemphstack#1;#2;#3;#4@{%
   \xdef\first@{#1} \xdef\second@{#2}
   \xdef\third@{#3} \xdef\last@{#4@}
}
\def\getregion@fromemphstack#1{%
   \xdef\temp@{#1}
   \xdef\first@{\csname stack@emphreg\temp@\endcsname}
   \expandafter\get@fromemphstack\first@
   \expandafter\xdef\csname emphstart\temp@\endcsname{\first@}
   \expandafter\xdef\csname emphstop\temp@\endcsname{\second@}
   \expandafter\xdef\csname emphall\temp@\endcsname{\third@}
   \expandafter\xdef\csname stack@emphreg\temp@\endcsname{\last@}
}
\def\sort@emphstack{%
   \expandafter\get@fromemphstack\last@
   \ifx\first@\ampers@nd
      \xdef\tmpstack{\tmpstack\st@rt;\st@p;\@ll;&;&;&;@}
   \else
   \ifnum\st@rt<\second@
      \xdef\tmpstack{\tmpstack\st@rt;\st@p;\@ll;\first@;\second@;\third@;\last@}
   \else
      \xdef\tmpstack{\tmpstack\first@;\second@;\third@;} \sort@emphstack
   \fi\fi
}
\def\get@emphregions#1..#2,#3@{%
   \regionalemphtrue
   \xdef\st@rt{#1} \xdef\st@p{#2} \xdef\list@{#3}
   \xdef\last@{\csname stack@emphreg\seq@\endcsname}
   \xdef\tmpstack{}
   \sort@emphstack
   \expandafter\xdef\csname stack@emphreg\seq@\endcsname{\tmpstack}
}
\def\getregion@fromlowerstack#1{%
   \xdef\temp@{#1}
   \xdef\first@{\csname stack@lowerreg\temp@\endcsname}
   \expandafter\get@fromemphstack\first@
   \expandafter\xdef\csname lowerstart\temp@\endcsname{\first@}
   \expandafter\xdef\csname lowerstop\temp@\endcsname{\second@}
   \expandafter\xdef\csname lowerall\temp@\endcsname{\third@}
   \expandafter\xdef\csname stack@lowerreg\temp@\endcsname{\last@}
}
\def\sort@lowerstack{%
   \expandafter\get@fromemphstack\last@
   \ifx\first@\ampers@nd
      \xdef\tmpstack{\tmpstack\st@rt;\st@p;\@ll;&;&;&;@}
   \else
   \ifnum\st@rt<\second@
      \xdef\tmpstack{\tmpstack\st@rt;\st@p;\@ll;\first@;\second@;\third@;\last@}
   \else
      \xdef\tmpstack{\tmpstack\first@;\second@;\third@;} \sort@lowerstack
   \fi\fi
}
\def\get@lowerregions#1..#2,#3@{%
   \regionallowertrue
   \xdef\st@rt{#1} \xdef\st@p{#2} \xdef\list@{#3}
   \xdef\last@{\csname stack@lowerreg\seq@\endcsname}
   \xdef\tmpstack{}
   \sort@lowerstack
   \expandafter\xdef\csname stack@lowerreg\seq@\endcsname{\tmpstack}
}
\def\get@fromdomstack#1;#2;#3@{%
   \xdef\first@{#1} \xdef\second@{#2}
   \xdef\last@{#3@}
}
\def\sort@domstack{%
   \expandafter\get@fromdomstack\last@
   \ifx\first@\ampers@nd
      \xdef\tmpstack{\tmpstack\st@rt;\st@p;&;&;@}
   \else
   \ifnum\st@rt<\second@
      \xdef\tmpstack{\tmpstack\st@rt;\st@p;\first@;\second@;\last@}
   \else
      \xdef\tmpstack{\tmpstack\first@;\second@;} \sort@domstack
   \fi\fi
}
\def\get@domainregions#1..#2,#3@{%
   \xdef\st@rt{#1} \xdef\st@p{#2} \xdef\list@{#3}
   \xdef\last@{\stack@dom}
   \xdef\tmpstack{}
   \sort@domstack
   \xdef\stack@dom{\tmpstack}
}
\def\get@domain{%
   \expandafter\get@fromdomstack\last@
   \ifx\first@\ampers@nd
   \else
      \loopcount=\first@
      \loop
         \xdef\domain@num@list{\domain@num@list \the\loopcount,}
         \advance\loopcount by 1
      \ifnum\loopcount>\second@\else\repeat
      \get@domain
   \fi
}
\def\getregion@fromtintstack#1{%
   \xdef\temp@{#1}
   \xdef\first@{\csname stack@tintreg\temp@\endcsname}
   \expandafter\get@fromemphstack\first@
   \expandafter\xdef\csname tintstart\temp@\endcsname{\first@}
   \expandafter\xdef\csname tintstop\temp@\endcsname{\second@}
   \expandafter\xdef\csname tintall\temp@\endcsname{\third@}
   \expandafter\xdef\csname stack@tintreg\temp@\endcsname{\last@}
}
\def\get@tintregions#1..#2,#3@{%
   \regionaltinttrue
   \xdef\st@rt{#1} \xdef\st@p{#2} \xdef\list@{#3}
   \xdef\last@{\csname stack@tintreg\seq@\endcsname}
   \xdef\tmpstack{}
   \sort@emphstack
   \expandafter\xdef\csname stack@tintreg\seq@\endcsname{\tmpstack}
}
\def\getregion@fromframestack#1{%
   \xdef\temp@{#1}
   \xdef\first@{\csname stack@framereg\temp@\endcsname}
   \expandafter\get@fromemphstack\first@
   \expandafter\xdef\csname framestart\temp@\endcsname{\first@}
   \expandafter\xdef\csname framestop\temp@\endcsname{\second@}
   \expandafter\xdef\csname framestyle\temp@\endcsname{\third@}
   \expandafter\xdef\csname stack@framereg\temp@\endcsname{\last@}
}
\def\get@frameregions#1..#2,#3@{%
   \frame@true
   \xdef\st@rt{#1} \xdef\st@p{#2} \xdef\list@{#3}
   \xdef\last@{\csname stack@framereg\seq@\endcsname}
   \xdef\tmpstack{}
   \sort@emphstack
   \expandafter\xdef\csname stack@framereg\seq@\endcsname{\tmpstack}
}
\def\getregion@fromshadingstack#1{%
   \xdef\temp@{#1}
   \xdef\first@{\csname stack@shadingreg\temp@\endcsname}
   \expandafter\get@fromemphstack\first@
   \expandafter\xdef\csname shadingstart\temp@\endcsname{\first@}
   \expandafter\xdef\csname shadingstop\temp@\endcsname{\second@}
   \expandafter\xdef\csname shadingstyle\temp@\endcsname{\third@}
   \expandafter\xdef\csname stack@shadingreg\temp@\endcsname{\last@}
}
\def\get@shadingregions#1..#2,#3@{%
   \shading@true
   \xdef\st@rt{#1} \xdef\st@p{#2} \xdef\list@{#3}
   \xdef\last@{\csname stack@shadingreg\seq@\endcsname}
   \xdef\tmpstack{}
   \sort@emphstack
   \expandafter\xdef\csname stack@shadingreg\seq@\endcsname{\tmpstack}
}
\def\getregion@fromfstack#1{%
   \xdef\temp@{#1}
   \xdef\first@{\csname stack@\bottop@\temp@\endcsname}
   \expandafter\get@fromstack\first@
   \expandafter\xdef\csname text\bottop@\temp@\endcsname{\first@}
   \expandafter\xdef\csname start\bottop@\temp@\endcsname{\second@}
   \expandafter\xdef\csname stop\bottop@\temp@\endcsname{\third@}
   \ifx\fourth@\ampers@nd \xdef\fourth@{///} \fi
   \expandafter\xdef\csname style\bottop@\temp@\endcsname{\fourth@}
   \expandafter\xdef\csname stack@\bottop@\temp@\endcsname{\last@}
}
\def\sort@fstack{%
   \expandafter\get@fromstack\last@
   \ifx\first@\ampers@nd
      \xdef\tmpstack{\tmpstack\f@text@;\st@rt;\st@p;\style@;&;&;&;&;@}
   \else
   \ifnum\st@rt<\second@
      \xdef\tmpstack{\tmpstack\f@text@;\st@rt;\st@p;\style@;%
                     \first@;\second@;\third@;\fourth@;\last@}
   \else
      \xdef\tmpstack{\tmpstack\first@;\second@;\third@;\fourth@;}
      \sort@fstack
   \fi\fi
}
\def\get@fregions#1..#2,#3@{%
   \xdef\st@rt{#1}
   \ifx\temp@\y@
      \xdef\st@p{#1}
   \else
      \xdef\st@p{#2}
   \fi
   \xdef\list@{#3}
   \xdef\last@{\csname stack@\bottop@\seq@\endcsname}
   \xdef\tmpstack{}
   \sort@fstack
   \expandafter\xdef\csname stack@\bottop@\seq@\endcsname{\tmpstack}
}
\def\getarrow@shape#1#2#3#4&{%
   \xdef\first@@{#1}\xdef\second@@{#2}\xdef\third@@{#3}\xdef\fourth@@{#4}
   \ifx\temp@@\y@
      \if\first@@ ` \xdef\first@@{a} \else
      \if\first@@ ' \xdef\first@@{a} \else
      \if\first@@ , \xdef\first@@{b} \else
      \if\first@@ S \xdef\first@@{c} \xdef\second@@{-}
      \fi\fi\fi\fi
      \if\third@@ ` \xdef\third@@{a} \else
      \if\third@@ ' \xdef\third@@{a} \else
      \if\third@@ , \xdef\third@@{b} \else
      \if\third@@ S \xdef\third@@{c} \xdef\second@@{-}
      \fi\fi\fi\fi
      \xdef\style@{\first@@ -\third@@\fourth@@}
   \fi
   \if\first@@ S
      \xdef\style@{\first@@ -\third@@\fourth@@}
   \else
   \if\first@@ v
      \if\second@@ = 
      \else \xdef\style@{\first@@ v\third@@\fourth@@} \fi
   \else
   \if\third@@ S
      \xdef\style@{\first@@ -\third@@\fourth@@}
   \else
   \if\third@@ v
      \if\second@@ =
      \else \xdef\style@{\first@@ v\third@@\fourth@@} \fi
   \fi\fi\fi\fi
}
\def\get@shape#1#2#3{%
   \xdef\first@@{#1}\xdef\second@@{#2}\xdef\third@@{#3}%
   \if\second@@ v \xdef\second@@{arrow}%
   \fi%
   \if\second@@ = \xdef\second@@{doublearrow}%
   \fi%
}
\def\getstyle@left#1#2#3#4@{%
   \ifstop@
      \xdef\style@@{\csname fstyle\bottop@\the\loopcount\endcsname}
   \else
   \xdef\style@@{#2}
   \xdef\temp@{-}
   \ifx\style@@\temp@ \xdef\style@@{#1#2-#4}
   \else
   \xdef\temp@{v}
   \ifx\style@@\temp@ \xdef\style@@{#1#2-#4}
   \else   
   \xdef\temp@{=}
   \ifx\style@@\temp@ \xdef\style@@{#1#2=#4}
   \else   
      \xdef\style@@{\csname fstyle\bottop@\the\loopcount\endcsname}
   \fi\fi\fi\fi
}
\def\getstyle@right#1#2#3#4@{%
   \xdef\style@@{#2}
   \xdef\temp@{-}
   \ifx\style@@\temp@
      \expandafter\xdef\csname fstyle\bottop@\the\loopcount\endcsname{-#2#3#4}
   \fi
   \xdef\temp@{v}
   \ifx\style@@\temp@
      \expandafter\xdef\csname fstyle\bottop@\the\loopcount\endcsname{-#2#3#4}
   \fi
   \xdef\temp@{=}
   \ifx\style@@\temp@
      \expandafter\xdef\csname fstyle\bottop@\the\loopcount\endcsname{=#2#3#4}
   \fi
}
\def\opt@color#1[#2]#3&{\xdef\fourth@{#1}\xdef\f@color{#2}}
\def\graph@opt@color#1[#2]:#3[#4][#5]#6&{%
   \xdef\fourth@{#1}
   \xdef\f@color{#2}
   \xdef\ffourth@{#3}
   \xdef\ff@color{#4}
   \xdef\fffourth@{#5}}
\def\arrow@col@width#1[#2][#3]#4&{\xdef\fourth@{#1}\xdef\f@color{#2}\xdef\rule@@thick{#3}}
\def\second@color#1,&{\xdef\back@color{#1}}
\def\two@opt@color#1,#2@{%
   \xdef\sixth@{#1}%
   \xdef\seventh@{#2}%
   \ifx\first@\ampers@nd%
   \else%
      \xdef\frame@color{#1}%
      \ifx\seventh@\ampers@nd%
         \xdef\back@color{#1}%
         \xdef\backtext@color{White}%
      \else%
         \expandafter\second@color\seventh@%
      \fi%
   \fi%
}
\def\test@get#1#2@{%
   \xdef\temp@{#1} \xdef\third@@{#2@}
   \ifx\temp@\ampers@nd
   \else
      \if \temp@ X
         \xdef\temp@{\second@@}
      \fi
      \ifx\temp@\second@@
         \xdef\third@@{\second@@}
      \else
         \expandafter\test@get\third@@
      \fi
   \fi
}
\def\get@@second@#1#2@{\xdef\second@@{#1}\xdef\first@@{#2@}}
\def\get@@third@#1]#2@{\xdef\third@@{#1&@}\xdef\fifth@@{#2@}}
\def\get@third@#1#2@{\xdef\third@@{#1}\xdef\fifth@@{#2@}}
\def\find@motif{%
   \expandafter\get@@second@\first@@
   \ifx\second@@\ampers@nd
   \else
      \ifx\second@@\d@t
      \else
         \advance\temp@count by 1
         \if \third@@ X
            \xdef\third@@{\second@@}
         \fi
         \ifx\third@@\br@cket
            \expandafter\get@@third@\fifth@@
            \expandafter\test@get\third@@
         \fi
         \ifx\second@@\third@@
            \ifx\nineth@@\n@
               \xdef\st@rt{\the\temp@count}
               \xdef\st@p{n}
            \fi
            \xdef\nineth@@{y}
            \expandafter\get@third@\fifth@@
            \ifx\third@@\ampers@nd
               \xdef\st@p{\the\temp@count}
               \xdef\list@{}
               \xdef\tmpstack{}
               \xdef\temp@{shade}
               \ifx\temp@\temp@@
                  \xdef\last@{\csname stack@reg\seq@\endcsname} 
                  \loopcount=\style@
                  \sort@stack
                  \expandafter\xdef\csname stack@reg\seq@\endcsname{\tmpstack}
               \else
               \xdef\temp@{emph}
               \ifx\temp@\temp@@
                  \regionalemphtrue
                  \xdef\last@{\csname stack@emphreg\seq@\endcsname} 
                  \sort@emphstack
                  \expandafter\xdef\csname stack@emphreg\seq@\endcsname{\tmpstack}
               \else
               \xdef\temp@{lower}
               \ifx\temp@\temp@@
                  \regionallowertrue
                  \xdef\last@{\csname stack@lowerreg\seq@\endcsname} 
                  \sort@lowerstack
                  \expandafter\xdef\csname stack@lowerreg\seq@\endcsname{\tmpstack}
               \else
               \xdef\temp@{tint}
               \ifx\temp@\temp@@
                  \regionaltinttrue
                  \xdef\last@{\csname stack@tintreg\seq@\endcsname}
                  \sort@emphstack
                  \expandafter\xdef\csname stack@tintreg\seq@\endcsname{\tmpstack}
               \else
               \xdef\temp@{frame}
               \ifx\temp@\temp@@
                  \frame@true
                  \xdef\last@{\csname stack@framereg\seq@\endcsname} 
                  \sort@emphstack
                  \expandafter\xdef\csname stack@framereg\seq@\endcsname{\tmpstack}
               \else
               \xdef\temp@{shading}
               \ifx\temp@\temp@@
                  \shading@true
                  \xdef\last@{\csname stack@shadingreg\seq@\endcsname} 
                  \sort@emphstack
                  \expandafter\xdef\csname stack@shadingreg\seq@\endcsname{\tmpstack}
               \else
               \xdef\temp@{ttttop}
               \ifx\temp@\temp@@
                  \xdef\last@{\csname stack@ttttop\seq@\endcsname} 
                  \sort@fstack
                  \expandafter\xdef\csname stack@ttttop\seq@\endcsname{\tmpstack}
               \else
               \xdef\temp@{tttop}
               \ifx\temp@\temp@@
                  \xdef\last@{\csname stack@tttop\seq@\endcsname} 
                  \sort@fstack
                  \expandafter\xdef\csname stack@tttop\seq@\endcsname{\tmpstack}
               \else
               \xdef\temp@{ttop}
               \ifx\temp@\temp@@
                  \xdef\last@{\csname stack@ttop\seq@\endcsname} 
                  \sort@fstack
                  \expandafter\xdef\csname stack@ttop\seq@\endcsname{\tmpstack}
               \else
               \xdef\temp@{top}
               \ifx\temp@\temp@@
                  \xdef\last@{\csname stack@top\seq@\endcsname}
                  \sort@fstack
                  \expandafter\xdef\csname stack@top\seq@\endcsname{\tmpstack}
               \else
               \xdef\temp@{bottom}
               \ifx\temp@\temp@@
                  \xdef\last@{\csname stack@bottom\seq@\endcsname} 
                  \sort@fstack
                  \expandafter\xdef\csname stack@bottom\seq@\endcsname{\tmpstack}
               \else
               \xdef\temp@{bbottom}
               \ifx\temp@\temp@@
                  \xdef\last@{\csname stack@bbottom\seq@\endcsname} 
                  \sort@fstack
                  \expandafter\xdef\csname stack@bbottom\seq@\endcsname{\tmpstack}
               \else
               \xdef\temp@{bbbottom}
               \ifx\temp@\temp@@
                  \xdef\last@{\csname stack@bbbottom\seq@\endcsname} 
                  \sort@fstack
                  \expandafter\xdef\csname stack@bbbottom\seq@\endcsname{\tmpstack}
               \else
               \xdef\temp@{bbbbottom}
               \ifx\temp@\temp@@
                  \xdef\last@{\csname stack@bbbbottom\seq@\endcsname} 
                  \sort@fstack
                  \expandafter\xdef\csname stack@bbbbottom\seq@\endcsname{\tmpstack}
               \fi
               \fi
               \fi
               \fi
               \fi
               \fi
               \fi
               \fi
               \fi
               \fi
               \fi
               \fi
               \fi
               \fi
               \xdef\fifth@@{\m@tif} \expandafter\get@third@\fifth@@
               \xdef\nineth@@{n}
            \fi
         \else
            \ifx\nineth@@\y@
               \xdef\first@@{\second@@\first@@}
               \advance\temp@count by -1
            \fi
            \xdef\nineth@@{n} \xdef\st@p{n}
            \xdef\fifth@@{\m@tif} \expandafter\get@third@\fifth@@
         \fi
      \fi
   \find@motif
   \fi
}
\def\prep@motif{%

   \def\get@motif@specs##1,##2,##3,##4,##5,&{%
      \xdef\seq@{##1} \xdef\temp@@{##2} \xdef\@ll{##3} \def\fourth@@{##4} \xdef\fifth@@{##5&@}
   }

   \loopcount=0
   \loop
      \advance\loopcount by 1
      \expandafter\xdef\csname stack@reg\the\loopcount\endcsname{\csname style\the\loopcount\endcsname;%
                                                                 \csname start\the\loopcount\endcsname;%
                                                                 \csname stop\the\loopcount\endcsname;%
                                                                 \csname all\the\loopcount\endcsname;%
                                                                 \csname stack@reg\the\loopcount\endcsname}
      \expandafter\xdef\csname stack@emphreg\the\loopcount\endcsname{\csname emphstart\the\loopcount\endcsname;%
                                                                 \csname emphstop\the\loopcount\endcsname;%
                                                                 \csname emphall\the\loopcount\endcsname;%
                                                                 \csname stack@emphreg\the\loopcount\endcsname}
      \expandafter\xdef\csname stack@lowerreg\the\loopcount\endcsname{\csname lowerstart\the\loopcount\endcsname;%
                                                                 \csname lowerstop\the\loopcount\endcsname;%
                                                                 \csname lowerall\the\loopcount\endcsname;%
                                                                 \csname stack@lowerreg\the\loopcount\endcsname}
      \expandafter\xdef\csname stack@tintreg\the\loopcount\endcsname{\csname tintstart\the\loopcount\endcsname;%
                                                                 \csname tintstop\the\loopcount\endcsname;%
                                                                 \csname tintall\the\loopcount\endcsname;%
                                                                 \csname stack@tintreg\the\loopcount\endcsname}
      \expandafter\xdef\csname stack@framereg\the\loopcount\endcsname{\csname framestart\the\loopcount\endcsname;%
                                                                 \csname framestop\the\loopcount\endcsname;%
                                                                 \csname framestyle\the\loopcount\endcsname;%
                                                                 \csname stack@framereg\the\loopcount\endcsname}
      \expandafter\xdef\csname stack@shadingreg\the\loopcount\endcsname{\csname shadingstart\the\loopcount\endcsname;%
                                                                 \csname shadingstop\the\loopcount\endcsname;%
                                                                 \csname shadingstyle\the\loopcount\endcsname;%
                                                                 \csname stack@shadingreg\the\loopcount\endcsname}
      \expandafter\xdef\csname stack@ttttop\the\loopcount\endcsname{\csname textttttop\the\loopcount\endcsname;%
                                                                 \csname startttttop\the\loopcount\endcsname;%
                                                                 \csname stopttttop\the\loopcount\endcsname;%
                                                                 \csname stylettttop\the\loopcount\endcsname;%
                                                                 \csname stack@ttttop\the\loopcount\endcsname}
      \expandafter\xdef\csname stack@tttop\the\loopcount\endcsname{\csname texttttop\the\loopcount\endcsname;%
                                                                 \csname starttttop\the\loopcount\endcsname;%
                                                                 \csname stoptttop\the\loopcount\endcsname;%
                                                                 \csname styletttop\the\loopcount\endcsname;%
                                                                 \csname stack@tttop\the\loopcount\endcsname}
      \expandafter\xdef\csname stack@ttop\the\loopcount\endcsname{\csname textttop\the\loopcount\endcsname;%
                                                                 \csname startttop\the\loopcount\endcsname;%
                                                                 \csname stopttop\the\loopcount\endcsname;%
                                                                 \csname stylettop\the\loopcount\endcsname;%
                                                                 \csname stack@ttop\the\loopcount\endcsname}
      \expandafter\xdef\csname stack@top\the\loopcount\endcsname{\csname texttop\the\loopcount\endcsname;%
                                                                 \csname starttop\the\loopcount\endcsname;%
                                                                 \csname stoptop\the\loopcount\endcsname;%
                                                                 \csname styletop\the\loopcount\endcsname;%
                                                                 \csname stack@top\the\loopcount\endcsname}
      \expandafter\xdef\csname stack@bottom\the\loopcount\endcsname{\csname textbottom\the\loopcount\endcsname;%
                                                                 \csname startbottom\the\loopcount\endcsname;%
                                                                 \csname stopbottom\the\loopcount\endcsname;%
                                                                 \csname stylebottom\the\loopcount\endcsname;%
                                                                 \csname stack@bottom\the\loopcount\endcsname}
      \expandafter\xdef\csname stack@bbottom\the\loopcount\endcsname{\csname textbbottom\the\loopcount\endcsname;%
                                                                 \csname startbbottom\the\loopcount\endcsname;%
                                                                 \csname stopbbottom\the\loopcount\endcsname;%
                                                                 \csname stylebbottom\the\loopcount\endcsname;%
                                                                 \csname stack@bbottom\the\loopcount\endcsname}
      \expandafter\xdef\csname stack@bbbottom\the\loopcount\endcsname{\csname textbbbottom\the\loopcount\endcsname;%
                                                                 \csname startbbbottom\the\loopcount\endcsname;%
                                                                 \csname stopbbbottom\the\loopcount\endcsname;%
                                                                 \csname stylebbbottom\the\loopcount\endcsname;%
                                                                 \csname stack@bbbottom\the\loopcount\endcsname}
      \expandafter\xdef\csname stack@bbbbottom\the\loopcount\endcsname{\csname textbbbbottom\the\loopcount\endcsname;%
                                                                 \csname startbbbbottom\the\loopcount\endcsname;%
                                                                 \csname stopbbbbottom\the\loopcount\endcsname;%
                                                                 \csname stylebbbbottom\the\loopcount\endcsname;%
                                                                 \csname stack@bbbbottom\the\loopcount\endcsname}
   \ifnum\loopcount<\seq@num\repeat
   \temp@@count=0
   \loop
      \advance\temp@@count by 1
      \xdef\first@@{\csname motif\the\temp@@count\endcsname}
      \expandafter\get@motif@specs\first@@
      \xdef\first@@{\csname sequence\seq@\endcsname &@}
      \temp@count=\csname res@count\seq@\endcsname
      \expandafter\ifnum\csname seq@start\seq@\endcsname<0 \ifx\allow@zero\n@ \advance\temp@count by 1 \fi\fi
      \xdef\f@text@{\@ll}
      \xdef\style@{\fourth@@}
      \xdef\m@tif{\fifth@@} \expandafter\get@third@\fifth@@
      \xdef\nineth@@{n} \xdef\st@p{n}
      \find@motif
   \ifnum\temp@@count<\motif@num\repeat
   \loopcount=0
   \loop
      \advance\loopcount by 1
      \expandafter\getregion@fromstack{\the\loopcount}
      \expandafter\getregion@fromtintstack{\the\loopcount}
      \expandafter\getregion@fromemphstack{\the\loopcount}
      \expandafter\getregion@fromlowerstack{\the\loopcount}
      \expandafter\getregion@fromframestack{\the\loopcount}
      \expandafter\getregion@fromshadingstack{\the\loopcount}
      \xdef\bottop@{top}       \expandafter\getregion@fromfstack{\the\loopcount}
      \xdef\bottop@{ttop}      \expandafter\getregion@fromfstack{\the\loopcount}
      \xdef\bottop@{tttop}     \expandafter\getregion@fromfstack{\the\loopcount}
      \xdef\bottop@{ttttop}    \expandafter\getregion@fromfstack{\the\loopcount}
      \xdef\bottop@{bottom}    \expandafter\getregion@fromfstack{\the\loopcount}
      \xdef\bottop@{bbottom}   \expandafter\getregion@fromfstack{\the\loopcount}
      \xdef\bottop@{bbbottom}  \expandafter\getregion@fromfstack{\the\loopcount}
      \xdef\bottop@{bbbbottom} \expandafter\getregion@fromfstack{\the\loopcount}
   \ifnum\loopcount<\seq@num\repeat
}

\def\test@motif#1#2@{%
   \xdef\temp@{shade}
   \ifx\temp@\temp@@ \xdef\style@{\the\loopcount} \fi
   \xdef\temp@{emph}
   \ifx\temp@\temp@@ \xdef\style@{0} \fi
   \xdef\temp@{lower}
   \ifx\temp@\temp@@ \xdef\style@{0} \fi
   \xdef\temp@{tint}
   \ifx\temp@\temp@@ \xdef\style@{0} \fi
   \xdef\temp@{frame}
   \ifx\temp@\temp@@ \xdef\style@{0} \fi
   \xdef\temp@{shading}
   \ifx\temp@\temp@@ \xdef\style@{0} \fi
   \xdef\temp@{#1#2}
   \xdef\first@{#1}
   \ifx\first@\br@cket
      \xdef\label@motif{y}
      \temp@count=\motif@num
      \advance\temp@count by 1
      \xdef\motif@num{\the\temp@count} 
      \expandafter\xdef\csname motif\the\temp@count\endcsname{\seq@,\temp@@,\@ll,\style@,\temp@}
      \xdef\list@{&}
   \else
      \expandafter\check@char\first@
      \ifletter
         \xdef\label@motif{y}
         \temp@count=\motif@num
         \advance\temp@count by 1
         \xdef\motif@num{\the\temp@count} 
         \expandafter\xdef\csname motif\the\temp@count\endcsname{\seq@,\temp@@,\@ll,\style@,\temp@}
         \xdef\list@{&}
      \fi
   \fi   
}

\def\test@PDB#1:#2,#3,#4,#5,#6@{%
   \xdef\last@{#1[,]&}\expandafter\opt@color\last@
   \xdef\PDB@distance{no}
   \ifx\f@color\comm@
   \else
      \xdef\f@color{\f@color .&&@}
      \expandafter\coord@get\f@color
      \xdef\temp@{\c@@rd @}
      \expandafter\check@letter\temp@
      \ifletter
      \else
         \temp@count=\c@@rd
         \multiply\temp@count by \temp@count
         \xdef\PDB@distance{\the\temp@count}
      \fi
   \fi
   \xdef\PDB@hitpos{0}
   \xdef\PDB@stack{}
   \xdef\PDB@name{#2}
   \xdef\PDB@type{point}
   \ifx\fourth@\PDB@type
      \xdef\PDB@{y}
      \xdef\last@{#3[,]&}\expandafter\opt@color\last@
      \xdef\PDB@refnum@a{\fourth@}
      \ifx\f@color\C@lpha
         \xdef\PDB@back@side@a{CA}
      \else
         \xdef\PDB@back@side@a{side}
      \fi
      \xdef\PDB@refnum@b{0} \xdef\PDB@back@side@b{x}
      \xdef\PDB@refnum@c{0} \xdef\PDB@back@side@c{x}
      \ifx\PDB@distance\n@ \xdef\PDB@distance{3600}\fi
      \load@PDB 
      \expandafter\get@item\PDB@stack \xdef\list@{\first@@}
   \else
      \xdef\PDB@type{line}
      \ifx\fourth@\PDB@type
         \xdef\PDB@{y}
         \xdef\last@{#3[,]&}\expandafter\opt@color\last@
         \xdef\PDB@refnum{\fourth@}
         \ifx\f@color\C@lpha
            \xdef\PDB@back@side{CA}
         \else
            \xdef\PDB@back@side{side}
         \fi
         \xdef\last@{#4[,]&}\expandafter\opt@color\last@
         \ifnum\PDB@refnum<\fourth@ 
            \xdef\PDB@refnum@a{\PDB@refnum} 
            \xdef\PDB@refnum@b{\fourth@}
            \xdef\PDB@back@side@a{\PDB@back@side}
            \ifx\f@color\C@lpha
               \xdef\PDB@back@side@b{CA}
            \else
               \xdef\PDB@back@side@b{side}
            \fi
         \else
            \xdef\PDB@refnum@b{\PDB@refnum} 
            \xdef\PDB@refnum@a{\fourth@}
            \xdef\PDB@back@side@b{\PDB@back@side}
            \ifx\f@color\C@lpha
               \xdef\PDB@back@side@a{CA}
            \else
               \xdef\PDB@back@side@a{side}
            \fi
         \fi
         \xdef\PDB@refnum@c{0} \xdef\PDB@back@side@c{x}
         \ifx\PDB@distance\n@ \xdef\PDB@distance{100}\fi
         \load@PDB 
         \expandafter\get@item\PDB@stack \xdef\list@{\first@@}
      \else
         \xdef\PDB@type{plane}
         \ifx\fourth@\PDB@type
            \xdef\PDB@{y}
            \xdef\last@{#3[,]&}\expandafter\opt@color\last@
            \xdef\PDB@refnum{\fourth@}
            \ifx\f@color\C@lpha
               \xdef\PDB@back@side{CA}
            \else
               \xdef\PDB@back@side{side}
            \fi
            \xdef\last@{#4[,]&}\expandafter\opt@color\last@
            \ifnum\PDB@refnum<\fourth@ 
               \xdef\PDB@refnum@a{\PDB@refnum} 
               \xdef\PDB@refnum@b{\fourth@}
               \xdef\PDB@back@side@a{\PDB@back@side}
               \ifx\f@color\C@lpha
                  \xdef\PDB@back@side@b{CA}
               \else
                  \xdef\PDB@back@side@b{side}
               \fi
            \else
               \xdef\PDB@refnum@b{\PDB@refnum} 
               \xdef\PDB@refnum@a{\fourth@}
               \xdef\PDB@back@side@b{\PDB@back@side}
               \ifx\f@color\C@lpha
                  \xdef\PDB@back@side@a{CA}
               \else
                  \xdef\PDB@back@side@a{side}
               \fi
            \fi
            \xdef\last@{#5[,]&}\expandafter\opt@color\last@
            \xdef\PDB@refnum{\fourth@}
            \ifx\f@color\C@lpha
               \xdef\PDB@back@side{CA}
            \else
               \xdef\PDB@back@side{side}
            \fi
            \ifnum\PDB@refnum>\PDB@refnum@b
               \xdef\PDB@refnum@c{\PDB@refnum}
               \xdef\PDB@back@side@c{\PDB@back@side}
            \else
               \ifnum\PDB@refnum>\PDB@refnum@a
                  \xdef\PDB@refnum@c{\PDB@refnum@b}
                  \xdef\PDB@back@side@c{\PDB@back@side@b}
                  \xdef\PDB@refnum@b{\PDB@refnum}
                  \xdef\PDB@back@side@b{\PDB@back@side}
               \else
                  \xdef\PDB@refnum@c{\PDB@refnum@b}
                  \xdef\PDB@back@side@c{\PDB@back@side@b}
                  \xdef\PDB@refnum@b{\PDB@refnum@a}
                  \xdef\PDB@back@side@b{\PDB@back@side@a}
                  \xdef\PDB@refnum@a{\PDB@refnum}
                  \xdef\PDB@back@side@a{\PDB@back@side}
               \fi
            \fi
            \ifx\PDB@distance\n@ \xdef\PDB@distance{100}\fi
            \load@PDB 
            \expandafter\get@item\PDB@stack \xdef\list@{\first@@}
         \else
            \xdef\temp@{\list@ @}
            \expandafter\test@motif\temp@
         \fi
      \fi
   \fi
}

\def\test@fill#1:#2:#3&{%
   \xdef\last@{#1[,][,]&}\expandafter\arrow@col@width\last@%
   \xdef\second@@{\fourth@}%
   \xdef\last@{///}%
   \ifx\fourth@\last@%
     \xdef\second@@{empty}%
   \else
   \xdef\last@{translate}%
   \ifx\fourth@\last@%
     \xdef\second@@{translate}%
     \xdef\fill@char{#2[,]&}%
     \expandafter\opt@color\fill@char%
     \xdef\fill@char{\fourth@}%
   \else
   \xdef\last@{fill}%
   \ifx\fourth@\last@
     \xdef\second@@{fill}%
     \xdef\fill@char{#2[,]&}%
     \expandafter\opt@color\fill@char%
     \xdef\fill@char{\fourth@}%
   \else
  \xdef\last@{bar}%
   \ifx\fourth@\last@
     \xdef\second@@{bar}%
     \xdef\b@r{bar}%
     \ifx\f@color\comm@
        \xdef\g@min{,}\xdef\g@max{,}%
     \else
        \xdef\f@color{\f@color @}%
        \expandafter\get@item\f@color%
        \xdef\g@min{\fourth@} \xdef\g@max{\first@@}%
        \xdef\pm@shift{\fourth@}%
     \fi%
     \xdef\fill@char{#2[,]&}%
     \expandafter\opt@color\fill@char%
     \ifx\f@color\comm@\xdef\f@color{GrayDefault}\fi%
     \xdef\fill@char{\fourth@}%
     \xdef\box@color{\f@color,&@}%
     \expandafter\two@opt@color\box@color%
   \else
   \xdef\last@{color}%
   \ifx\fourth@\last@
     \xdef\second@@{color}%
     \ifx\f@color\comm@
        \xdef\g@min{,}\xdef\g@max{,}%
     \else
        \xdef\f@color{\f@color @}%
        \expandafter\get@item\f@color%
        \xdef\g@min{\fourth@} \xdef\g@max{\first@@}%
     \fi%
     \xdef\fill@char{#2[,]&}%
     \expandafter\opt@color\fill@char%
     \ifx\f@color\comm@\xdef\f@color{GrayDefault}\fi%
     \xdef\fill@char{\fourth@}%
   \else
     \xdef\last@{box}%
     \ifx\fourth@\last@%
        \xdef\second@@{box}%
        \ifx\f@color\comm@\xdef\f@color{White}\fi%
        \xdef\seventh@{\f@color @}\expandafter\check@letter\seventh@%
        \ifnumber%
           \ifx\rule@@thick\comm@%
              \xdef\rule@@thick{\f@color}%
              \xdef\box@color{White,&@}%
           \else%
              \xdef\seventh@{\rule@@thick}%
              \xdef\rule@@thick{\f@color}%
              \xdef\box@color{\seventh@,&@}%
           \fi%
        \else%
           \xdef\box@color{\f@color,&@}%
        \fi%
        \expandafter\two@opt@color\box@color%
        \xdef\fill@char{#2[,]&}%
        \expandafter\opt@color\fill@char%
        \xdef\fill@char{\fourth@}%
   \else
     \xdef\last@{plotcolor}%
     \ifx\fourth@\last@%
        \xdef\second@@{plotcolor}%
        \ifx\f@color\comm@%
           \xdef\pm@shift{0}%
        \else%
           \xdef\pm@shift{\f@color}%
        \fi%
        \xdef\fill@char{#2[,]&}%
        \expandafter\opt@color\fill@char%
        \xdef\fill@char{\fourth@}%
   \else
     \xdef\last@{plotbar}%
     \ifx\fourth@\last@%
        \xdef\second@@{plotbar}%
        \xdef\b@r{bar}%
        \ifx\f@color\comm@
           \xdef\pm@shift{0}%
        \else
           \xdef\pm@shift{\f@color}%
        \fi%
        \xdef\fill@char{#2[,]&}%
        \expandafter\opt@color\fill@char%
        \ifx\f@color\comm@\xdef\f@color{GrayDefault}\fi%
        \xdef\fill@char{\fourth@}%
        \xdef\box@color{\f@color,&@}%
         \expandafter\two@opt@color\box@color%
     \else
        \xdef\arrow@color{\f@color}%
        \ifx\arrow@color\comm@\xdef\arrow@color{Black}\fi%
        \xdef\seventh@{\f@color @}\expandafter\check@letter\seventh@%
        \ifnumber%
           \ifx\rule@@thick\comm@%
              \xdef\rule@@thick{\f@color}%
              \xdef\arrow@color{Black}%
           \else%
              \xdef\seventh@{\rule@@thick}%
              \xdef\rule@@thick{\f@color}%
              \xdef\arrow@color{\seventh@}%
           \fi%
        \fi%
        \xdef\fill@char{#2[,]&}%
        \expandafter\opt@color\fill@char%
        \xdef\fill@char{\fourth@}%
        \ifx\f@color\comm@\xdef\f@color{Black,White}\fi%
        \xdef\backtext@color{White}%
        \xdef\f@color{\f@color,&@}%
        \expandafter\two@opt@color\f@color%
        \xdef\f@color{\arrow@color}%
     \fi
   \fi
   \fi
   \fi
   \fi
   \fi
   \fi
   \fi
   \xdef\seventh@{\f@color @}\expandafter\check@letter\seventh@%
   \ifnumber%
      \ifx\rule@@thick\comm@%
         \xdef\rule@@thick{\f@color}%
         \xdef\f@color{Black}%
      \else%
         \xdef\seventh@{\rule@@thick}%
         \xdef\rule@@thick{\f@color}%
         \xdef\f@color{\seventh@}%
      \fi%
   \fi%
   \ifx\rule@@thick\comm@\xdef\rule@@thick{\the\rule@thick}\fi%
   \ifx\f@color\comm@\xdef\f@color{Black}\fi}
\def\clear@groups{%
    \expandafter\xdef\csname \prfx grpA\endcsname{ -1}
    \expandafter\xdef\csname \prfx grpB\endcsname{ -2}
    \expandafter\xdef\csname \prfx grpC\endcsname{ -3}
    \expandafter\xdef\csname \prfx grpD\endcsname{ -4}
    \expandafter\xdef\csname \prfx grpE\endcsname{ -5}
    \expandafter\xdef\csname \prfx grpF\endcsname{ -6}
    \expandafter\xdef\csname \prfx grpG\endcsname{ -7}
    \expandafter\xdef\csname \prfx grpH\endcsname{ -8}
    \expandafter\xdef\csname \prfx grpI\endcsname{ -9}
    \expandafter\xdef\csname \prfx grpJ\endcsname{-10}
    \expandafter\xdef\csname \prfx grpK\endcsname{-11}
    \expandafter\xdef\csname \prfx grpL\endcsname{-12}
    \expandafter\xdef\csname \prfx grpM\endcsname{-13}
    \expandafter\xdef\csname \prfx grpN\endcsname{-14}
    \expandafter\xdef\csname \prfx grpO\endcsname{-15}
    \expandafter\xdef\csname \prfx grpP\endcsname{-16}
    \expandafter\xdef\csname \prfx grpQ\endcsname{-17}
    \expandafter\xdef\csname \prfx grpR\endcsname{-18}
    \expandafter\xdef\csname \prfx grpS\endcsname{-19}
    \expandafter\xdef\csname \prfx grpT\endcsname{-20}
    \expandafter\xdef\csname \prfx grpU\endcsname{-21}
    \expandafter\xdef\csname \prfx grpV\endcsname{-22}
    \expandafter\xdef\csname \prfx grpW\endcsname{-23}
    \expandafter\xdef\csname \prfx grpX\endcsname{-24}
    \expandafter\xdef\csname \prfx grpY\endcsname{-25}
    \expandafter\xdef\csname \prfx grpZ\endcsname{-26}
    \expandafter\xdef\csname \prfx grp-\endcsname{-999}
    \expandafter\xdef\csname \prfx grp.\endcsname{-999}
    \expandafter\xdef\csname \prfx grp=\endcsname{-999}
}
\def\inactivate@chars{%
   \catcode`\#=12
   \catcode`\"=12
   \catcode`\~=12
   \catcode`\^=12
   \catcode`\_=12
   }
\def\numcount{\the\loopcount}
\def\Alphacount{\@Alph\loopcount}
\def\alphacount{\@alph\loopcount}
\def\romancount{\@roman\loopcount}
\def\Romancount{\@Roman\loopcount}
\def\cut@name#1.#2@{\global\xdef\file@n@me{#1}}
\def\coord@get#1.#2#3@{%
   \xdef\temp@{#2}
   \ifx\temp@\ampers@nd \xdef\temp@{0} \fi
   \xdef\c@@rd{#1\temp@}
}
\def\struc@get#1 #2 #3 #4 #5 #6 #7 #8 #9@{%
    \xdef\first@{#1} \xdef\second@{#2} \xdef\third@{#3}  \xdef\fourth@{#4}
    \xdef\fifth@{#5} \xdef\sixth@{#6}  \xdef\seventh@{#7}\xdef\eighth@{#8}
    \xdef\nineth@{#9}}
\def\get@PHD#1|#2|#3@{\xdef\PHD@line{\PHD@line #2}}
\def\write@PHDsec{%
 \expandafter\get@sim\last@
 \ifx\sim@char\c@mp
      \def\end@{\the\innerloopcount}
      \advance\innerloopcount by 1
 \else
      \ifnum\end@<\begin@ \xdef\end@{\begin@}\fi
      \if\c@mp .
      \else
      \if\c@mp L
      \else
      \if\c@mp H
        \ifx\show@Hsec\yes
         \loopcount=\first@
         \advance\loopcount by 1
         \xdef\first@{\the\loopcount}
         \ifx\m@p\yes
         \immediate\write\featurefile{%
         \string\feature{\bottop@Hsec}{\begin@..\end@}%
               {\label@Hsec}{\text@Hsec}}
         \else
         \immediate\write\featurefile{%
         \string\feature{\bottop@Hsec}{\st@rt}{\begin@..\end@}%
               {\label@Hsec}{\text@Hsec}}\fi
        \fi
      \else
      \if\c@mp E
        \ifx\show@Esec\yes
         \loopcount=\second@
         \advance\loopcount by 1
         \xdef\second@{\the\loopcount}
         \ifx\m@p\yes
         \immediate\write\featurefile{%
         \string\feature{\bottop@Esec}{\begin@..\end@}%
               {\label@Esec}{\text@Esec}}
         \else
         \immediate\write\featurefile{%
         \string\feature{\bottop@Esec}{\st@rt}{\begin@..\end@}%
               {\label@Esec}{\text@Esec}}\fi
        \fi
      \fi\fi\fi\fi
      \xdef\c@mp{\sim@char}
      \advance\innerloopcount by 1
      \xdef\begin@{\the\innerloopcount}
   \fi
 \ifx\sim@char\ampers@nd\else\write@PHDsec\fi
}
\def\write@PHDtopo{%
 \expandafter\get@sim\last@
 \ifx\sim@char\c@mp
      \def\end@{\the\innerloopcount}
      \advance\innerloopcount by 1
 \else
      \ifnum\end@<\begin@ \xdef\end@{\begin@}\fi
      \if\c@mp .
      \else
      \if\c@mp L
      \else
      \if\c@mp T
        \ifx\show@TMtop\yes
         \loopcount=\first@
         \advance\loopcount by 1
         \xdef\first@{\the\loopcount}
         \ifx\m@p\yes
         \immediate\write\featurefile{%
         \string\feature{\bottop@TMtop}{\begin@..\end@}%
               {\label@TMtop}{\text@TMtop}}
         \else
         \immediate\write\featurefile{%
         \string\feature{\bottop@TMtop}{\st@rt}{\begin@..\end@}%
               {\label@TMtop}{\text@TMtop}}\fi
        \fi
      \else
      \if\c@mp i
        \ifx\show@itop\yes
         \loopcount=\second@
         \advance\loopcount by 1
         \xdef\second@{\the\loopcount}
         \ifx\m@p\yes
         \immediate\write\featurefile{%
         \string\feature{\bottop@itop}{\begin@..\end@}%
               {\label@itop}{\text@itop}}
         \else
         \immediate\write\featurefile{%
         \string\feature{\bottop@itop}{\st@rt}{\begin@..\end@}%
               {\label@itop}{\text@itop}}\fi
        \fi
      \else
      \if\c@mp o
        \ifx\show@etop\yes
         \loopcount=\second@
         \advance\loopcount by 1
         \xdef\second@{\the\loopcount}
         \ifx\m@p\yes
         \immediate\write\featurefile{%
         \string\feature{\bottop@etop}{\begin@..\end@}%
               {\label@etop}{\text@etop}}
         \else
         \immediate\write\featurefile{%
         \string\feature{\bottop@etop}{\st@rt}{\begin@..\end@}%
               {\label@etop}{\text@etop}}\fi
        \fi
      \fi\fi\fi\fi\fi
      \xdef\c@mp{\sim@char}
      \advance\innerloopcount by 1
      \xdef\begin@{\the\innerloopcount}
 \fi
 \ifx\sim@char\ampers@nd\else\write@PHDtopo\fi
}
\def\include@T@coffee{%
    \bgroup
    \immediate\openin\structurefile = \TC@first@\relax
    \ifeof\structurefile
      \PackageError{TeXshade}
      {File `\TC@first@' not found}
      {\MessageBreak
       The file you specified is missing or you have \MessageBreak
       misspelled it. \MessageBreak\MessageBreak
       No T-Coffee shading will be displayed (using "similar"). \MessageBreak
       Type <return> to proceed. \MessageBreak
       Type X <return> to quit. \MessageBreak
}
       \immediate\closein\structurefile\egroup
       \simmodetrue  \funcmodefalse \T@coffeefalse
    \else
    \message{<T-Coffee shading>}
    \xdef\first@{} \xdef\temp@{*}
    \loop
       \read\structurefile to \readline
       \xdef\test@{\expandafter\string\readline}
       \ifx\test@\par@
       \else
          \xdef\structureline{\readline & & & & & & & & &@}
          \expandafter\struc@get\structureline
          \ifx\first@\temp@ \xdef\temp@{y}
             \read\structurefile to \readline
             \read\structurefile to \readline
          \fi
       \fi
    \ifx\temp@\y@ \else\repeat
    \loopcount=0
    \xdef\temp@{cons}
    \loop
       \read\structurefile to \readline
       \xdef\test@{\expandafter\string\readline}
       \ifx\test@\par@
       \else
          \xdef\structureline{\readline & & & & & & & & &@}
          \expandafter\struc@get\structureline     
          \if\second@ : 
             \advance\loopcount by 1\relax
             \expandafter\xdef\csname T@coffee\the\loopcount\endcsname{}
          \fi
       \fi
    \ifx\first@\temp@\expandafter\xdef\csname T@coffee0\endcsname{} \else\repeat
    \xdef\temp@{\the\loopcount}
    \loopcount=0
    \loop
       \read\structurefile to \readline
       \xdef\test@{\expandafter\string\readline}
       \ifx\test@\par@
       \else
          \xdef\structureline{\readline & & & & & & & & &@}
          \expandafter\struc@get\structureline
          \advance\loopcount by 1
          \ifnum\loopcount=\temp@ \loopcount=0 \fi
          \expandafter\xdef\csname T@coffee\the\loopcount\endcsname{\csname T@coffee\the\loopcount\endcsname \third@}
       \fi
    \ifeof\structurefile \else\repeat
    \closein\structurefile
    \loopcount=0
    \loop
       \expandafter\xdef\csname T@coffee\the\loopcount\endcsname{\csname T@coffee\the\loopcount\endcsname @}
       \advance\loopcount by 1
    \ifnum\loopcount=\temp@ \else\repeat
    \egroup
    \def\n@m@tch{{$\bullet$}}
    \expandafter\def\csname fg@color0\endcsname{TC0}
    \expandafter\def\csname fg@color1\endcsname{TC1}
    \expandafter\def\csname fg@color2\endcsname{TC2}
    \expandafter\def\csname fg@color3\endcsname{TC3}
    \expandafter\def\csname fg@color4\endcsname{TC4}
    \expandafter\def\csname fg@color5\endcsname{TC5}
    \expandafter\def\csname fg@color6\endcsname{TC6}
    \expandafter\def\csname fg@color7\endcsname{TC7}
    \expandafter\def\csname fg@color8\endcsname{TC8}
    \expandafter\def\csname fg@color9\endcsname{TC9}
    \expandafter\def\csname fg@textcolor0\endcsname{Black}
    \expandafter\def\csname fg@textcolor1\endcsname{Black}
    \expandafter\def\csname fg@textcolor2\endcsname{Black}
    \expandafter\def\csname fg@textcolor3\endcsname{Black}
    \expandafter\def\csname fg@textcolor4\endcsname{Black}
    \expandafter\def\csname fg@textcolor5\endcsname{Black}
    \expandafter\def\csname fg@textcolor6\endcsname{Black}
    \expandafter\def\csname fg@textcolor7\endcsname{Black}
    \expandafter\def\csname fg@textcolor8\endcsname{Black}
    \expandafter\def\csname fg@textcolor9\endcsname{Black}
   \fi
 }

\def\include@DSSP{%
  \xdef\first@{\csname optiondssp\the\loopcount\endcsname}
  \xdef\bottop@{\csname bottopdssp\the\loopcount\endcsname}
  \xdef\st@rt{\csname doseqdssp\the\loopcount\endcsname}
  \xdef\structurefilename{\csname filenamedssp\the\loopcount\endcsname}
  \bgroup
   \xdef\file@n@me{\structurefilename .@}
   \expandafter\cut@name\file@n@me
   \xdef\file@n@me{\file@n@me.sec}
   \ifx\first@\file@n@me
   \else
      \immediate\openin\alignfile = \file@n@me\relax
      \ifeof\alignfile \xdef\first@{make new} \fi
      \immediate\closein\alignfile
   \fi
   \xdef\temp@{make new}
   \ifx\first@\temp@
    \def\par{}
    \inactivate@chars
    \immediate\openin\structurefile = \structurefilename\relax
    \ifeof\structurefile
      \PackageError{TeXshade}
      {File `\structurefilename' not found}
      {\MessageBreak
       The `DSSP' file you specified is missing or you have \MessageBreak
       misspelled it. \MessageBreak\MessageBreak
       No labels for secondary structures will be displayed. \MessageBreak
       Type <return> to proceed. \MessageBreak
       Type X <return> to quit. \MessageBreak
}
       \immediate\closein\structurefile\egroup
    \else
    \message{[\structurefilename] ->}
    \xdef\second@{} \xdef\temp@{RESIDUE}
    \loop
       \read\structurefile to \readline
       \xdef\test@{\expandafter\string\readline}
       \xdef\structureline{\readline & & & & & & & & &@}
       \expandafter\struc@get\structureline
       \ifx\second@\temp@ \xdef\temp@{AA}\fi
    \ifx\third@\temp@ \else\repeat
    \immediate\openout\featurefile = \file@n@me
    \xdef\c@mp{+}
    \xdef\begin@{\csname seq@start\st@rt\endcsname}
    \xdef\end@{\csname seq@start\st@rt\endcsname}
    \xdef\st@rt@{\begin@}
    \expandafter\innerloopcount=\csname seq@start\st@rt\endcsname
    \advance\innerloopcount by -1
    \xdef\first@@{0} \xdef\second@@{0} \xdef\third@@{0}
    \xdef\fourth@@{0} \xdef\fifth@@{0} \xdef\sixth@@{0}
    \xdef\seventh@@{0}
    \loop
       \read\structurefile to \readline
       \xdef\test@{\expandafter\string\readline}
       \ifx\test@\par@
       \else
             \xdef\structureline{\readline & & & & & & & & &@}
             \expandafter\struc@get\structureline
             \if\c@mp + \xdef\c@mp{\fifth@}\fi
             \ifx\fifth@\c@mp
                \ifx\fc@DSSP\y@
                   \temp@count=\first@
                \else
                   \temp@count=\second@
                \fi
                \advance\temp@count by \st@rt@
                \advance\temp@count by -1
                \xdef\end@{\the\temp@count}
             \else
                   \ifnum\begin@>0
                      \ifnum\end@<\begin@ \xdef\end@{\begin@}\fi
                      \if\c@mp C
                      \else
                      \if\c@mp H
                        \ifx\show@Hdssp\yes
                         \loopcount=\first@@
                         \advance\loopcount by 1
                         \xdef\first@@{\the\loopcount}
                         \ifx\m@p\yes
                         \immediate\write\featurefile{%
                            \string\feature{\bottop@Hdssp}{\begin@..\end@}%
                            {\label@Hdssp}{\text@Hdssp}}
                         \else
                         \immediate\write\featurefile{%
                            \string\feature{\bottop@Hdssp}{\st@rt}{\begin@..\end@}%
                            {\label@Hdssp}{\text@Hdssp}}\fi
                        \fi
                      \else
                      \if\c@mp G
                        \ifx\show@Gdssp\yes
                         \loopcount=\second@@
                         \advance\loopcount by 1
                         \xdef\second@@{\the\loopcount}
                         \ifx\m@p\yes
                         \immediate\write\featurefile{%
                            \string\feature{\bottop@Gdssp}{\begin@..\end@}%
                            {\label@Gdssp}{\text@Gdssp}}
                         \else
                         \immediate\write\featurefile{%
                            \string\feature{\bottop@Gdssp}{\st@rt}{\begin@..\end@}%
                            {\label@Gdssp}{\text@Gdssp}}\fi
                        \fi
                      \else
                      \if\c@mp I
                        \ifx\show@Idssp\yes
                         \loopcount=\third@@
                         \advance\loopcount by 1
                         \xdef\third@@{\the\loopcount}
                         \ifx\m@p\yes
                         \immediate\write\featurefile{%
                            \string\feature{\bottop@Idssp}{\begin@..\end@}%
                            {\label@Idssp}{\text@Idssp}}
                         \else
                         \immediate\write\featurefile{%
                            \string\feature{\bottop@Idssp}{\st@rt}{\begin@..\end@}%
                            {\label@Idssp}{\text@Idssp}}\fi
                        \fi
                      \else
                      \if\c@mp E
                        \ifx\show@Edssp\yes
                         \loopcount=\fourth@@
                         \advance\loopcount by 1
                         \xdef\fourth@@{\the\loopcount}
                         \advance\innerloopcount by 1
                         \ifx\m@p\yes
                         \immediate\write\featurefile{%
                            \string\feature{\bottop@Edssp}{\begin@..\end@}%
                            {\label@Edssp}{\text@Edssp}}
                         \else
                         \immediate\write\featurefile{%
                            \string\feature{\bottop@Edssp}{\st@rt}{\begin@..\end@}%
                            {\label@Edssp}{\text@Edssp}}\fi
                        \fi
                      \else
                      \if\c@mp B
                        \ifx\show@Bdssp\yes
                         \loopcount=\fifth@@
                         \advance\loopcount by 1
                         \xdef\fifth@@{\the\loopcount}
                         \ifx\m@p\yes
                         \immediate\write\featurefile{%
                            \string\feature{\bottop@Bdssp}{\begin@..\end@}%
                            {\label@Bdssp}{\text@Bdssp}}
                         \else
                         \immediate\write\featurefile{%
                            \string\feature{\bottop@Bdssp}{\st@rt}{\begin@..\end@}%
                            {\label@Bdssp}{\text@Bdssp}}\fi
                        \fi
                      \else
                      \if\c@mp T
                        \ifx\show@Tdssp\yes
                         \loopcount=\sixth@@
                         \advance\loopcount by 1
                         \xdef\sixth@@{\the\loopcount}
                         \ifx\m@p\yes
                         \immediate\write\featurefile{%
                            \string\feature{\bottop@Tdssp}{\begin@..\end@}%
                            {\label@Tdssp}{\text@Tdssp}}
                         \else
                         \immediate\write\featurefile{%
                            \string\feature{\bottop@Tdssp}{\st@rt}{\begin@..\end@}%
                            {\label@Tdssp}{\text@Tdssp}}\fi
                        \fi
                      \else
                      \if\c@mp S
                        \ifx\show@Sdssp\yes
                         \loopcount=\seventh@@
                         \advance\loopcount by 1
                         \xdef\seventh@@{\the\loopcount}
                         \ifx\m@p\yes
                         \immediate\write\featurefile{%
                            \string\feature{\bottop@Sdssp}{\begin@..\end@}%
                            {\label@Sdssp}{\text@Sdssp}}
                         \else
                         \immediate\write\featurefile{%
                            \string\feature{\bottop@Sdssp}{\st@rt}{\begin@..\end@}%
                            {\label@Sdssp}{\text@Sdssp}}\fi
                        \fi
                      \fi\fi\fi\fi\fi\fi\fi\fi
                   \fi
                   \xdef\c@mp{\fifth@}
                   \ifx\fc@DSSP\y@
                      \temp@count=\first@
                   \else
                      \temp@count=\second@
                   \fi
                   \advance\temp@count by \st@rt@
                   \advance\temp@count by -1
                   \xdef\begin@{\the\temp@count}
             \fi
       \fi
    \ifeof\structurefile \else\repeat
    \closein\structurefile
    \immediate\closeout\featurefile
    \egroup
    \input{\file@n@me}
   \fi
   \else
    \egroup
    \message{using existing file:}
    \input{\file@n@me}
   \fi
 }

\def\include@HMMTOP{%
 \def\get@HMMTOP@TMs##1-##2 ##3@{%
     \xdef\temp@@{\temp@@\fourth@@##1\fourth@@##2}
     \xdef\structureline{##3 @}
  }
  \def\get@HMMTOP{%
     \ifnum\temp@count<\fifth@
        \advance\temp@count by 1
        \expandafter\get@HMMTOP@TMs\structureline
        \get@HMMTOP
     \fi
  }  
  \def\rem@ve@TM@info Transmembrane helices: ##1@{%
     \xdef\structureline{##1 @}
     \temp@count=0
     \get@HMMTOP
  }
 \xdef\bottop@{\csname bottopHMMTOP\the\loopcount\endcsname}
  \xdef\st@rt{\csname doseqHMMTOP\the\loopcount\endcsname}
  \xdef\first@{\csname optionHMMTOP\the\loopcount\endcsname}
  \xdef\structurefilename{\csname filenameHMMTOP\the\loopcount\endcsname}
  \bgroup
   \xdef\file@n@me{\structurefilename .@}
   \expandafter\cut@name\file@n@me
   \xdef\file@n@me{\file@n@me.top}
   \ifx\first@\file@n@me
   \else
      \immediate\openin\alignfile = \file@n@me\relax
      \ifeof\alignfile \xdef\first@{make new} \fi
      \immediate\closein\alignfile
   \fi
   \xdef\temp@{make new}
   \ifx\first@\temp@
    \def\par{}
    \inactivate@chars
    \immediate\openin\structurefile = \structurefilename\relax
    \ifeof\structurefile
      \PackageError{TeXshade}
      {File `\structurefilename' not found}
      {\MessageBreak
       The `HMMTOP' file you specified is missing or you have \MessageBreak
       misspelled it. \MessageBreak\MessageBreak
       No labels for secondary structures will be displayed. \MessageBreak
       Type <return> to proceed. \MessageBreak
       Type X <return> to quit. \MessageBreak
}
       \immediate\closein\structurefile\egroup
    \else
    \message{[\structurefilename] ->}
    \immediate\openout\featurefile = \file@n@me
    
    \xdef\first@{\csname fileseqHMMTOP\the\loopcount\endcsname @}
    \expandafter\check@letter\first@
    \ifletter
       \xdef\st@p{\csname fileseqHMMTOP\the\loopcount\endcsname}
    \else
       \xdef\first@{\csname fileseqHMMTOP\the\loopcount\endcsname}
       \ifnum\first@=0
          \xdef\st@p{0}
       \else
          \xdef\st@p{\csname fileseqHMMTOP\the\loopcount\endcsname}
       \fi
    \fi
    \xdef\temp@{yes} \innerloopcount=0
    \loop
       \read\structurefile to \readline
       \xdef\test@{\expandafter\string\readline}
       \ifx\test@\par@
       \else
          \xdef\structureline{\readline & & & & & & & & &@}
          \expandafter\struc@get\structureline
          \ifx\first@\@HP
             \ifx\temp@\yes \xdef\temp@{\readline} \fi
             \ifletter
                \ifx\st@p\third@ \xdef\temp@{\readline} \fi
             \else
             \ifnum\st@p=0
                \expandafter\ifx\csname seqname\st@rt\endcsname\third@
                   \xdef\temp@{\readline} \xdef\st@p{-1}
                \else
                   \expandafter\ifx\csname newseqname\st@rt\endcsname\third@
                      \xdef\temp@{\readline} \xdef\st@p{-1}
	               \fi
	            \fi
	         \else
	            \advance\innerloopcount by 1
	            \ifnum\st@p=\innerloopcount \xdef\temp@{\readline} \fi
	         \fi
	         \fi
	      \else
	      \xdef\first@@{Protein:}
	      \ifx\first@\first@@
	         \xdef\second@@{\second@}
	         \xdef\temp@@{>HP:} \xdef\fourth@@{ }
             \read\structurefile to \readline
             \xdef\structureline{\readline & & & & & & & & &@}
             \expandafter\struc@get\structureline
             \xdef\temp@@{\temp@@\fourth@@\second@\fourth@@\second@@}
             \read\structurefile to \readline
             \xdef\structureline{\readline & & & & & & & & &@}
             \expandafter\struc@get\structureline
             \xdef\temp@@{\temp@@\fourth@@\second@}
             \read\structurefile to \readline
             \xdef\structureline{\readline & & & & & & & & &@}
             \expandafter\struc@get\structureline
             \xdef\temp@@{\temp@@\fourth@@\fifth@}
             \read\structurefile to \readline
             \xdef\structureline{\readline @}
             \expandafter\rem@ve@TM@info\structureline
             \xdef\temp@@{\temp@@\fourth@@}
             \ifx\temp@\yes \xdef\temp@{\temp@@} \fi
             \ifletter
                \ifx\st@p\third@ \xdef\temp@{\temp@@} \fi
             \else
             \ifnum\st@p=0
                \expandafter\ifx\csname seqname\st@rt\endcsname\second@@
                   \xdef\temp@{\temp@@}
                \else
                   \expandafter\ifx\csname newseqname\st@rt\endcsname\second@@
                      \xdef\temp@{\temp@@}
	               \fi
	            \fi
	         \else
	            \advance\innerloopcount by 1
	            \ifnum\st@p=\innerloopcount \xdef\temp@{\temp@@} \fi
	         \fi
	         \fi
	      \fi
          \fi
       \fi
    \ifeof\structurefile \else\repeat
    \xdef\seq@line{\temp@ @}
    \innerloopcount=0
    \loop
       \advance\innerloopcount by 1
       \expandafter\seq@get\seq@line
       \xdef\seq@line{\seq@line @}
       \ifnum\innerloopcount=4 \xdef\c@mp{\first@} \fi
    \ifnum\innerloopcount=5 \xdef\st@p{\first@} \else \repeat
    \xdef\first@{IN}
    \ifx\c@mp\first@ \xdef\c@mp{i} \else \xdef\c@mp{e} \fi
    \xdef\begin@{1} \xdef\first@@{0} \xdef\second@@{0}
    \innerloopcount=0
    \loop
       \advance\innerloopcount by 1
       \expandafter\seq@get\seq@line
       \xdef\seq@line{\seq@line @}
       \temp@count=\first@
       \advance\temp@count by -1
       \xdef\end@{\the\temp@count}
       \loopcount=\second@@
       \advance\loopcount by 1
       \xdef\second@@{\the\loopcount}
       \if\c@mp i
          \ifx\show@i@HMMTOP\yes
          \immediate\write\featurefile{%
             \string\feature{\bottop@i@HMMTOP}{\st@rt}{\begin@..\end@}%
             {\label@i@HMMTOP}{\text@i@HMMTOP}}\fi
          \xdef\c@mp{e}
       \else
          \ifx\show@e@HMMTOP\yes
          \immediate\write\featurefile{%
             \string\feature{\bottop@e@HMMTOP}{\st@rt}{\begin@..\end@}%
             {\label@e@HMMTOP}{\text@e@HMMTOP}}\fi
          \xdef\c@mp{i}
       \fi
       \advance\temp@count by 1
       \xdef\begin@{\the\temp@count}
       \expandafter\seq@get\seq@line
       \xdef\seq@line{\seq@line @}
       \xdef\end@{\first@}
       \loopcount=\first@@
       \advance\loopcount by 1
       \xdef\first@@{\the\loopcount}
       \ifx\show@TM@HMMTOP\yes
          \immediate\write\featurefile{%
          \string\feature{\bottop@TM@HMMTOP}{\st@rt}{\begin@..\end@}%
          {\label@TM@HMMTOP}{\text@TM@HMMTOP}}\fi
       \temp@count=\end@ \advance\temp@count by 1 \xdef\begin@{\the\temp@count}
    \ifnum\innerloopcount=\st@p\else\repeat
    \closein\structurefile
    \immediate\closeout\featurefile
    \egroup
    
    \input{\file@n@me}
   \fi
   \else
    \egroup
    \message{using existing file:}
    \input{\file@n@me}
   \fi
 }

\def\include@stride{%
  \xdef\first@{\csname optionstride\the\loopcount\endcsname}
  \xdef\bottop@{\csname bottopstride\the\loopcount\endcsname}
  \xdef\st@rt{\csname doseqstride\the\loopcount\endcsname}
  \xdef\structurefilename{\csname filenamestride\the\loopcount\endcsname}
  \bgroup
   \xdef\file@n@me{\structurefilename .@}
   \expandafter\cut@name\file@n@me
   \xdef\file@n@me{\file@n@me.sec}
   \ifx\first@\file@n@me
   \else
      \immediate\openin\alignfile = \file@n@me\relax
      \ifeof\alignfile \xdef\first@{make new} \fi
      \immediate\closein\alignfile
   \fi
   \xdef\temp@{make new}
   \ifx\first@\temp@
    \def\par{}
    \inactivate@chars
    \immediate\openin\structurefile = \structurefilename\relax
    \ifeof\structurefile
      \PackageError{TeXshade}
      {File `\structurefilename' not found}
      {\MessageBreak
       The `STRIDE' file you specified is missing or you have \MessageBreak
       misspelled it. \MessageBreak\MessageBreak
       No labels for secondary structures will be displayed. \MessageBreak
       Type <return> to proceed. \MessageBreak
       Type X <return> to quit. \MessageBreak
}
       \immediate\closein\structurefile\egroup
    \else
    \message{[\structurefilename] ->}
    \immediate\openout\featurefile = \file@n@me
    \xdef\c@mp{+}
    \xdef\begin@{\csname seq@start\st@rt\endcsname}
    \xdef\end@{\csname seq@start\st@rt\endcsname}
    \xdef\st@rt@{\begin@}
    \expandafter\innerloopcount=\csname seq@start\st@rt\endcsname
    \advance\innerloopcount by -1
    \xdef\first@@{0} \xdef\second@@{0} \xdef\third@@{0}
    \xdef\fourth@@{0} \xdef\fifth@@{0} \xdef\sixth@@{0}
    \loop
       \read\structurefile to \readline
       \xdef\test@{\expandafter\string\readline}
       \ifx\test@\par@
       \else
          \xdef\structureline{\readline & & & & & & & & &@}
          \expandafter\struc@get\structureline
          \ifx\first@\@asg
             \if\c@mp + \xdef\c@mp{\sixth@}\fi
             \ifx\sixth@\c@mp
                \temp@count=\fifth@
                \advance\temp@count by \st@rt@
                \advance\temp@count by -1
                \xdef\end@{\the\temp@count}
             \else
                   \ifnum\begin@>0
                      \ifnum\end@<\begin@ \xdef\end@{\begin@}\fi
                      \if\c@mp C
                      \else
                      \if\c@mp H
                        \ifx\show@Hstride\yes
                         \loopcount=\first@@
                         \advance\loopcount by 1
                         \xdef\first@@{\the\loopcount}
                         \ifx\m@p\yes
                         \immediate\write\featurefile{%
                            \string\feature{\bottop@Hstride}{\begin@..\end@}%
                            {\label@Hstride}{\text@Hstride}}
                         \else
                         \immediate\write\featurefile{%
                            \string\feature{\bottop@Hstride}{\st@rt}{\begin@..\end@}%
                            {\label@Hstride}{\text@Hstride}}\fi
                        \fi
                      \else
                      \if\c@mp G
                        \ifx\show@Gstride\yes
                         \loopcount=\second@@
                         \advance\loopcount by 1
                         \xdef\second@@{\the\loopcount}
                         \ifx\m@p\yes
                         \immediate\write\featurefile{%
                            \string\feature{\bottop@Gstride}{\begin@..\end@}%
                            {\label@Gstride}{\text@Gstride}}
                         \else
                         \immediate\write\featurefile{%
                            \string\feature{\bottop@Gstride}{\st@rt}{\begin@..\end@}%
                            {\label@Gstride}{\text@Gstride}}\fi
                        \fi
                      \else
                      \if\c@mp I
                        \ifx\show@Istride\yes
                         \loopcount=\third@@
                         \advance\loopcount by 1
                         \xdef\third@@{\the\loopcount}
                         \ifx\m@p\yes
                         \immediate\write\featurefile{%
                            \string\feature{\bottop@Istride}{\begin@..\end@}%
                            {\label@Istride}{\text@Istride}}
                         \else
                         \immediate\write\featurefile{%
                            \string\feature{\bottop@Istride}{\st@rt}{\begin@..\end@}%
                            {\label@Istride}{\text@Istride}}\fi
                        \fi
                      \else
                      \if\c@mp E
                        \ifx\show@Estride\yes
                         \loopcount=\fourth@@
                         \advance\loopcount by 1
                         \xdef\fourth@@{\the\loopcount}
                         \advance\innerloopcount by 1
                         \ifx\m@p\yes
                         \immediate\write\featurefile{%
                            \string\feature{\bottop@Estride}{\begin@..\end@}%
                            {\label@Estride}{\text@Estride}}
                         \else
                         \immediate\write\featurefile{%
                            \string\feature{\bottop@Estride}{\st@rt}{\begin@..\end@}%
                            {\label@Estride}{\text@Estride}}\fi
                        \fi
                      \else
                      \if\c@mp B
                        \ifx\show@Bstride\yes
                         \loopcount=\fifth@@
                         \advance\loopcount by 1
                         \xdef\fifth@@{\the\loopcount}
                         \ifx\m@p\yes
                         \immediate\write\featurefile{%
                            \string\feature{\bottop@Bstride}{\begin@..\end@}%
                            {\label@Bstride}{\text@Bstride}}
                         \else
                         \immediate\write\featurefile{%
                            \string\feature{\bottop@Bstride}{\st@rt}{\begin@..\end@}%
                            {\label@Bstride}{\text@Bstride}}\fi
                        \fi
                      \else
                      \if\c@mp T
                        \ifx\show@Tstride\yes
                         \loopcount=\sixth@@
                         \advance\loopcount by 1
                         \xdef\sixth@@{\the\loopcount}
                         \ifx\m@p\yes
                         \immediate\write\featurefile{%
                            \string\feature{\bottop@Tstride}{\begin@..\end@}%
                            {\label@Tstride}{\text@Tstride}}
                         \else
                         \immediate\write\featurefile{%
                            \string\feature{\bottop@Tstride}{\st@rt}{\begin@..\end@}%
                            {\label@Tstride}{\text@Tstride}}\fi
                        \fi
                      \fi\fi\fi\fi\fi\fi\fi
                   \fi
                   \xdef\c@mp{\sixth@}
                   \temp@count=\fifth@
                   \advance\temp@count by \st@rt@
                   \advance\temp@count by -1
                   \xdef\begin@{\the\temp@count}
             \fi
          \fi
       \fi
    \ifeof\structurefile \else\repeat
    \closein\structurefile
    \immediate\closeout\featurefile
    \egroup
    \input{\file@n@me}
   \fi
   \else
    \egroup
    \message{using existing file:}
    \input{\file@n@me}
   \fi
 }
\def\include@PHD{%
  \xdef\first@{\csname optionphd\the\loopcount\endcsname}
  \xdef\bottop@{\csname bottopphd\the\loopcount\endcsname}
  \xdef\st@rt{\csname doseqphd\the\loopcount\endcsname}
  \xdef\m@de{\csname modephd\the\loopcount\endcsname}
  \xdef\structurefilename{\csname filenamephd\the\loopcount\endcsname}
  \bgroup
   \xdef\file@n@me{\structurefilename .@}
   \expandafter\cut@name\file@n@me
   \xdef\temp@{structure}
   \ifx\m@de\temp@
      \xdef\temp@{\file@n@me .sec}
      \immediate\openin\alignfile = \temp@\relax
      \ifeof\alignfile \xdef\first@{make new} \fi
      \immediate\closein\alignfile
   \else
   \xdef\temp@{topology}
   \ifx\m@de\temp@
      \xdef\temp@{\file@n@me .top}
      \immediate\openin\alignfile = \temp@\relax
      \ifeof\alignfile \xdef\first@{make new} \fi
      \immediate\closein\alignfile
   \else
      \message{<Unknown type - ignoring \noexpand\includePHD>}
      \xdef\first@{ignore}
   \fi\fi
   \xdef\temp@{make new}
   \ifx\first@\temp@
    \def\par{}
    \xdef\PHD@line{}
    \inactivate@chars
    \immediate\openin\structurefile=\structurefilename\relax
    \ifeof\structurefile
      \PackageError{TeXshade}%
      {File `\structurefilename' not found}%
      {\MessageBreak
       The `PHD' file you specified is missing or you have \MessageBreak
       misspelled it. \MessageBreak\MessageBreak
       No labels for secondary structures will be displayed. \MessageBreak
       Type <return> to proceed. \MessageBreak
       Type X <return> to quit. \MessageBreak
}
       \immediate\closein\structurefile\egroup
    \else
    \message{[\structurefilename] ->}
    \loop
       \read\structurefile to \readline
       \xdef\structureline{\readline & & & & & & & & &@}
       \expandafter\struc@get\structureline
       \xdef\temp@{structure}
       \ifx\m@de\temp@
          \xdef\temp@{SUB}
          \ifx\temp@\first@
             \xdef\temp@{sec}
             \ifx\temp@\second@
                \xdef\third@{\third@ @}
                \expandafter\get@PHD\third@
             \fi
          \else
             \ifx\temp@\second@
                \xdef\temp@{sec}
                \ifx\temp@\third@
                   \xdef\fourth@{\fourth@ @}
                   \expandafter\get@PHD\fourth@
                \fi
             \fi
          \fi
       \else
       \xdef\temp@{topology}
       \ifx\m@de\temp@
          \xdef\temp@{PHDThtm}
          \ifx\temp@\first@
             \xdef\second@{\second@ @}
             \expandafter\get@PHD\second@
          \fi
       \fi\fi
    \ifeof\structurefile \else\repeat
    \closein\structurefile
    \xdef\c@mp{+}
    \xdef\begin@{\csname seq@start\st@rt\endcsname}
    \xdef\end@{\csname seq@start\st@rt\endcsname}
    \expandafter\innerloopcount=\csname seq@start\st@rt\endcsname
    \advance\innerloopcount by -1
    \xdef\first@{0} \xdef\second@{0} \xdef\third@{0}
    \xdef\fourth@{0} \xdef\fifth@{0} \xdef\sixth@{0}
    \xdef\last@{\PHD@line &@}
    \xdef\temp@{structure}
    \ifx\temp@\m@de
       \xdef\file@n@me{\file@n@me .sec}
       \immediate\openout\featurefile = \file@n@me\relax
       \write@PHDsec
    \fi
    \xdef\temp@{topology}
    \ifx\temp@\m@de
       \xdef\file@n@me{\file@n@me .top}
       \immediate\openout\featurefile = \file@n@me\relax
       \write@PHDtopo
    \fi
    \immediate\closeout\featurefile
    \egroup
    \input{\file@n@me}
   \fi
   \else
    \egroup
    \xdef\temp@{ignore}
    \ifx\temp@\first@
    \else
       \message{using existing file:}
       \xdef\temp@{structure}
       \ifx\temp@\m@de \xdef\file@n@me{\file@n@me .sec} \fi
       \xdef\temp@{topology}
       \ifx\temp@\m@de \xdef\file@n@me{\file@n@me .top} \fi
       \input{\file@n@me}
    \fi
   \fi}
\def\show@DSSP{%
   \expandafter\get@item\first@
   \ifx\fourth@\ampers@nd
   \else
      \xdef\temp@{alpha}
      \ifx\fourth@\temp@ \xdef\show@Hdssp{\second@}\fi
      \xdef\temp@{3-10}
      \ifx\fourth@\temp@ \xdef\show@Gdssp{\second@}\fi
      \xdef\temp@{pi}
      \ifx\fourth@\temp@ \xdef\show@Idssp{\second@}\fi
      \xdef\temp@{beta}
      \ifx\fourth@\temp@ \xdef\show@Edssp{\second@}\fi
      \xdef\temp@{bridge}
      \ifx\fourth@\temp@ \xdef\show@Bdssp{\second@}\fi
      \xdef\temp@{turn}
      \ifx\fourth@\temp@ \xdef\show@Tdssp{\second@}\fi
      \xdef\temp@{bend}
      \ifx\fourth@\temp@ \xdef\show@Sdssp{\second@}\fi
      \show@DSSP
   \fi}
\def\show@HMMTOP{%
   \expandafter\get@item\first@
   \ifx\fourth@\ampers@nd
   \else
      \xdef\temp@{internal}
      \ifx\fourth@\temp@ \xdef\show@i@HMMTOP{\second@}\fi
      \xdef\temp@{external}
      \ifx\fourth@\temp@ \xdef\show@e@HMMTOP{\second@}\fi
      \xdef\temp@{TM}
      \ifx\fourth@\temp@ \xdef\show@TM@HMMTOP{\second@}\fi
      \show@HMMTOP
   \fi}
\def\show@STRIDE{%
   \expandafter\get@item\first@
   \ifx\fourth@\ampers@nd
   \else
      \xdef\temp@{alpha}
      \ifx\fourth@\temp@ \xdef\show@Hstride{\second@}\fi
      \xdef\temp@{3-10}
      \ifx\fourth@\temp@ \xdef\show@Gstride{\second@}\fi
      \xdef\temp@{pi}
      \ifx\fourth@\temp@ \xdef\show@Istride{\second@}\fi
      \xdef\temp@{beta}
      \ifx\fourth@\temp@ \xdef\show@Estride{\second@}\fi
      \xdef\temp@{bridge}
      \ifx\fourth@\temp@ \xdef\show@Bstride{\second@}\fi
      \xdef\temp@{turn}
      \ifx\fourth@\temp@ \xdef\show@Tstride{\second@}\fi
      \show@STRIDE
   \fi}
\def\show@PHDtopo{%
   \expandafter\get@item\first@
   \ifx\fourth@\ampers@nd
   \else
      \xdef\temp@{internal}
      \ifx\fourth@\temp@ \xdef\show@itop{\second@}\fi
      \xdef\temp@{external}
      \ifx\fourth@\temp@ \xdef\show@etop{\second@}\fi
      \xdef\temp@{TM}
      \ifx\fourth@\temp@ \xdef\show@TMtop{\second@}\fi
      \show@PHDtopo
   \fi}
\def\show@PHDsec{%
   \expandafter\get@item\first@
   \ifx\fourth@\ampers@nd
   \else
      \xdef\temp@{alpha}
      \ifx\fourth@\temp@ \xdef\show@Hsec{\second@}\fi
      \xdef\temp@{beta}
      \ifx\fourth@\temp@ \xdef\show@Esec{\second@}\fi
      \show@PHDsec
   \fi}
\def\get@triplet#1,#2@{%
   \xdef\third@{#1}
   \ifx\third@\ampers@nd
   \else
      \expandafter\xdef\csname @\third@\endcsname{\first@}
      \expandafter\xdef\csname rev@\first@\endcsname{\third@}
  \xdef\fourth@{#2,&,@}
  \expandafter\get@triplet\fourth@
   \fi}
\def\get@dom@count#1,#2@{\res@count=#1 \xdef\temp@{#2 @}}   
\def\get@name@number{%
   \xdef\second@{n}
   \loopcount=0
   \loop
      \advance\loopcount by 1
      \expandafter\ifx\csname newseqname\the\loopcount\endcsname\first@
         \xdef\first@{\the\loopcount} \loopcount=\seq@count 
         \xdef\second@{y}
      \fi
   \ifnum\loopcount=\seq@count \else \repeat
   \ifx\second@\n@
      \message{<Sequence name `\first@' was not found - using first sequence.>}
      \xdef\first@{1}
   \fi
}
\def\get@name@number@table{%
   \xdef\second@{n}%
   \loopcount=0%
   \loop%
      \advance\loopcount by 1%
      \expandafter\ifx\csname newseqname\the\loopcount\endcsname\first@%
         \xdef\first@{\the\loopcount}\loopcount=\seq@num%
         \xdef\second@{y}%
      \fi%
   \ifnum\loopcount=\seq@num\else\repeat%
   \ifx\second@\n@%
      \message{<Sequence name `\first@' was not found - using first sequence.>}%
      \xdef\first@{1}%
   \fi%
}
 
\def\calc@distance@point{%
   \loopcount=\xref@coord@a \advance\loopcount by -\x@coord\relax 
   \multiply\loopcount by \loopcount \xdef\x@distance{\the\loopcount}
   \loopcount=\yref@coord@a \advance\loopcount by -\y@coord\relax 
   \multiply\loopcount by \loopcount \xdef\y@distance{\the\loopcount}
   \loopcount=\zref@coord@a \advance\loopcount by -\z@coord\relax 
   \multiply\loopcount by \loopcount \xdef\z@distance{\the\loopcount}
   \advance\loopcount by \y@distance\relax
   \advance\loopcount by \x@distance\relax
   \ifnum\loopcount > \PDB@distance \else 
      \loopcount=\PDB@hitpos \advance\loopcount by 1 
      \ifnum\loopcount=\PDB@pos
      \else
         \xdef\PDB@stack{\PDB@stack\PDB@hitpos,\PDB@pos..}
      \fi
      \xdef\PDB@hitpos{\PDB@pos}
   \fi
}

\def\calc@distance@line{%
   \loopcount=\x@coord\relax \multiply\loopcount by \x@orient
   \temp@count=\y@coord \relax \multiply\temp@count by \y@orient
   \advance\loopcount by \temp@count 
   \temp@count=\z@coord \relax \multiply\temp@count by \z@orient
   \advance\loopcount by \temp@count 
   \advance\loopcount by -\c@nst@ 
   \xdef\z@hler{\the\loopcount}
   \multiply\loopcount by \x@orient \divide\loopcount by \f@ct@r \advance\loopcount by \xref@coord@a
   \advance\loopcount by -\x@coord
   \multiply\loopcount by \loopcount \xdef\x@distance{\the\loopcount}
   \loopcount=\z@hler   
   \multiply\loopcount by \y@orient \divide\loopcount by \f@ct@r \advance\loopcount by \yref@coord@a
   \advance\loopcount by -\y@coord
   \multiply\loopcount by \loopcount \xdef\y@distance{\the\loopcount}
   \loopcount=\z@hler   
   \multiply\loopcount by \z@orient \divide\loopcount by \f@ct@r \advance\loopcount by \zref@coord@a
   \advance\loopcount by -\z@coord
   \multiply\loopcount by \loopcount \xdef\z@distance{\the\loopcount}
   \advance\loopcount by \y@distance\relax
   \advance\loopcount by \x@distance\relax
   \ifnum\loopcount > \PDB@distance \else 
      \loopcount=\PDB@hitpos \advance\loopcount by 1 
      \ifnum\loopcount=\PDB@pos
      \else
         \xdef\PDB@stack{\PDB@stack\PDB@hitpos,\PDB@pos..}
      \fi
      \xdef\PDB@hitpos{\PDB@pos}
   \fi
}


\def\intr@@t{%
   \loopcount=\r@@t
   \temp@count=1
   \loop
      \advance\loopcount by \temp@count
      \divide\loopcount by 2
      \temp@count=\r@@t
      \divide\temp@count by \loopcount
   \ifnum\loopcount>\temp@count\repeat
}


\def\calc@distance@plane{%
   \loopcount=\x@coord\relax
   \multiply\loopcount by \x@orient
   \temp@count=\y@coord \relax
   \multiply\temp@count by \y@orient
   \advance\loopcount by \temp@count 
   \temp@count=\z@coord \relax
   \multiply\temp@count by \z@orient
   \advance\loopcount by \temp@count
   \advance\loopcount by \c@nst@
   \divide\loopcount by \f@ct@r  
   \multiply\loopcount by \loopcount
   \ifnum\loopcount > \PDB@distance \else 
      \loopcount=\PDB@hitpos \advance\loopcount by 1 
      \ifnum\loopcount=\PDB@pos
      \else
         \xdef\PDB@stack{\PDB@stack\PDB@hitpos,\PDB@pos..}
      \fi
      \xdef\PDB@hitpos{\PDB@pos}
   \fi
}

   
\def\load@PDB{%
    \bgroup
    \immediate\openin\structurefile = \PDB@name\relax
    \ifeof\structurefile
      \PackageError{TeXshade}
      {File `\PDB@name' not found}
      {\MessageBreak
       The file you specified is missing or you have \MessageBreak
       misspelled it. \MessageBreak\MessageBreak
       No PDB domain labeling will be displayed. \MessageBreak
       Type <return> to proceed. \MessageBreak
       Type X <return> to quit. \MessageBreak
}
       \immediate\closein\structurefile\egroup
    \else
    \message{<reading PDB: \PDB@name>}
    \xdef\first@{}
    \loop
       \read\structurefile to \readline
       \xdef\test@{\expandafter\string\readline}
       \ifx\test@\par@
       \else
          \xdef\structureline{\readline & & & & & & & & &@}
          \expandafter\struc@get\structureline
          \ifx\first@\@TOM 
             \ifx\PDB@back@side@a\C@lpha
                \ifx\third@\C@lpha
                   \xdef\temp@@@{\fifth@ &&@}
                   \expandafter\check@letter\temp@@@
                   \ifnumber
                      \ifnum\fifth@=\PDB@refnum@a
                         \xdef\sixth@{\sixth@ &.&&@} \expandafter\coord@get\sixth@ \xdef\xref@coord@a{\c@@rd}
                         \xdef\seventh@{\seventh@ &.&&@} \expandafter\coord@get\seventh@ \xdef\yref@coord@a{\c@@rd}
                         \xdef\eighth@{\eighth@ &.&&@} \expandafter\coord@get\eighth@ \xdef\zref@coord@a{\c@@rd}
                      \fi
                   \else
                      \ifnum\sixth@=\PDB@refnum@a
                         \xdef\seventh@{\seventh@ &.&&@} \expandafter\coord@get\seventh@ \xdef\xref@coord@a{\c@@rd}
                         \xdef\eighth@{\eighth@ &.&&@} \expandafter\coord@get\eighth@ \xdef\yref@coord@a{\c@@rd}
                         \xdef\nineth@{\nineth@ &.&&@} \expandafter\coord@get\nineth@ \xdef\zref@coord@a{\c@@rd}
                      \fi
                   \fi
                \fi
             \else
                \xdef\temp@@@{\fifth@ &&@}
                \expandafter\check@letter\temp@@@
                \ifnumber
                   \ifnum\fifth@=\PDB@refnum@a
                      \xdef\sixth@{\sixth@ &.&&@} \expandafter\coord@get\sixth@ \xdef\xref@coord@a{\c@@rd}
                      \xdef\seventh@{\seventh@ &.&&@} \expandafter\coord@get\seventh@ \xdef\yref@coord@a{\c@@rd}
                      \xdef\eighth@{\eighth@ &.&&@} \expandafter\coord@get\eighth@ \xdef\zref@coord@a{\c@@rd}
                   \fi
                \else
                   \ifnum\sixth@=\PDB@refnum@a
                      \xdef\seventh@{\seventh@ &.&&@} \expandafter\coord@get\seventh@ \xdef\xref@coord@a{\c@@rd}
                      \xdef\eighth@{\eighth@ &.&&@} \expandafter\coord@get\eighth@ \xdef\yref@coord@a{\c@@rd}
                      \xdef\nineth@{\nineth@ &.&&@} \expandafter\coord@get\nineth@ \xdef\zref@coord@a{\c@@rd}
                   \fi
                \fi
             \fi
             \ifx\PDB@back@side@b\C@lpha
                \ifx\third@\C@lpha
                   \xdef\temp@@@{\fifth@ &&@}
                   \expandafter\check@letter\temp@@@
                   \ifnumber
                      \ifnum\fifth@=\PDB@refnum@b
                         \xdef\sixth@{\sixth@ &.&&@} \expandafter\coord@get\sixth@ \xdef\xref@coord@b{\c@@rd}
                         \xdef\seventh@{\seventh@ &.&&@} \expandafter\coord@get\seventh@ \xdef\yref@coord@b{\c@@rd}
                         \xdef\eighth@{\eighth@ &.&&@} \expandafter\coord@get\eighth@ \xdef\zref@coord@b{\c@@rd}
                      \fi
                   \else
                      \ifnum\sixth@=\PDB@refnum@b
                         \xdef\seventh@{\seventh@ &.&&@} \expandafter\coord@get\seventh@ \xdef\xref@coord@b{\c@@rd}
                         \xdef\eighth@{\eighth@ &.&&@} \expandafter\coord@get\eighth@ \xdef\yref@coord@b{\c@@rd}
                         \xdef\nineth@{\nineth@ &.&&@} \expandafter\coord@get\nineth@ \xdef\zref@coord@b{\c@@rd}
                      \fi
                   \fi
                \fi
             \else
                \xdef\temp@@@{\fifth@ &&@}
                \expandafter\check@letter\temp@@@
                \ifnumber
                   \ifnum\fifth@=\PDB@refnum@b
                      \xdef\sixth@{\sixth@ &.&&@} \expandafter\coord@get\sixth@ \xdef\xref@coord@b{\c@@rd}
                      \xdef\seventh@{\seventh@ &.&&@} \expandafter\coord@get\seventh@ \xdef\yref@coord@b{\c@@rd}
                      \xdef\eighth@{\eighth@ &.&&@} \expandafter\coord@get\eighth@ \xdef\zref@coord@b{\c@@rd}
                   \fi
                \else
                   \ifnum\sixth@=\PDB@refnum@b
                      \xdef\seventh@{\seventh@ &.&&@} \expandafter\coord@get\seventh@ \xdef\xref@coord@b{\c@@rd}
                      \xdef\eighth@{\eighth@ &.&&@} \expandafter\coord@get\eighth@ \xdef\yref@coord@b{\c@@rd}
                      \xdef\nineth@{\nineth@ &.&&@} \expandafter\coord@get\nineth@ \xdef\zref@coord@b{\c@@rd}
                   \fi
                \fi
             \fi
             \ifx\PDB@back@side@c\C@lpha
                \ifx\third@\C@lpha
                   \xdef\temp@@@{\fifth@ &&@}
                   \expandafter\check@letter\temp@@@
                   \ifnumber
                      \ifnum\fifth@=\PDB@refnum@c
                         \xdef\sixth@{\sixth@ &.&&@} \expandafter\coord@get\sixth@ \xdef\xref@coord@c{\c@@rd}
                         \xdef\seventh@{\seventh@ &.&&@} \expandafter\coord@get\seventh@ \xdef\yref@coord@c{\c@@rd}
                         \xdef\eighth@{\eighth@ &.&&@} \expandafter\coord@get\eighth@ \xdef\zref@coord@c{\c@@rd}
                      \fi
                   \else
                      \ifnum\sixth@=\PDB@refnum@c
                         \xdef\seventh@{\seventh@ &.&&@} \expandafter\coord@get\seventh@ \xdef\xref@coord@c{\c@@rd}
                         \xdef\eighth@{\eighth@ &.&&@} \expandafter\coord@get\eighth@ \xdef\yref@coord@c{\c@@rd}
                         \xdef\nineth@{\nineth@ &.&&@} \expandafter\coord@get\nineth@ \xdef\zref@coord@c{\c@@rd}
                      \fi
                   \fi
                \fi
             \else
                \xdef\temp@@@{\fifth@ &&@}
                \expandafter\check@letter\temp@@@
                \ifnumber
                   \ifnum\fifth@=\PDB@refnum@c
                      \xdef\sixth@{\sixth@ &.&&@} \expandafter\coord@get\sixth@ \xdef\xref@coord@c{\c@@rd}
                      \xdef\seventh@{\seventh@ &.&&@} \expandafter\coord@get\seventh@ \xdef\yref@coord@c{\c@@rd}
                      \xdef\eighth@{\eighth@ &.&&@} \expandafter\coord@get\eighth@ \xdef\zref@coord@c{\c@@rd}
                   \fi
                \else
                   \ifnum\sixth@=\PDB@refnum@c
                      \xdef\seventh@{\seventh@ &.&&@} \expandafter\coord@get\seventh@ \xdef\xref@coord@c{\c@@rd}
                      \xdef\eighth@{\eighth@ &.&&@} \expandafter\coord@get\eighth@ \xdef\yref@coord@c{\c@@rd}
                      \xdef\nineth@{\nineth@ &.&&@} \expandafter\coord@get\nineth@ \xdef\zref@coord@c{\c@@rd}
                   \fi
                \fi
             \fi
          \fi
       \fi
    \ifeof\structurefile \else\repeat
    \closein\structurefile
    \ifx\PDB@type\@line@
       \loopcount=\xref@coord@a \advance\loopcount by -\xref@coord@b \xdef\x@orient{\the\loopcount}
       \loopcount=\yref@coord@a \advance\loopcount by -\yref@coord@b \xdef\y@orient{\the\loopcount}
       \loopcount=\zref@coord@a \advance\loopcount by -\zref@coord@b \xdef\z@orient{\the\loopcount}
       \loopcount=\xref@coord@a \multiply\loopcount by \x@orient
       \temp@count=\yref@coord@a \multiply\temp@count by \y@orient
       \advance\loopcount by \temp@count
       \temp@count=\zref@coord@a \multiply\temp@count by \z@orient
       \advance\loopcount by \temp@count 
       \xdef\c@nst@{\the\loopcount}
       \loopcount=\x@orient \multiply\loopcount by \x@orient
       \temp@count=\y@orient \multiply\temp@count by \y@orient
       \advance\loopcount by \temp@count
       \temp@count=\z@orient \multiply\temp@count by \z@orient
       \advance\loopcount by \temp@count
       \xdef\f@ct@r{\the\loopcount}
    \fi
    \ifx\PDB@type\@plane@
       \loopcount=\xref@coord@a \advance\loopcount by -\xref@coord@b \xdef\x@orient@r{\the\loopcount}
       \loopcount=\yref@coord@a \advance\loopcount by -\yref@coord@b \xdef\y@orient@r{\the\loopcount}
       \loopcount=\zref@coord@a \advance\loopcount by -\zref@coord@b \xdef\z@orient@r{\the\loopcount}
       \loopcount=\xref@coord@a \advance\loopcount by -\xref@coord@c \xdef\x@orient@s{\the\loopcount}
       \loopcount=\yref@coord@a \advance\loopcount by -\yref@coord@c \xdef\y@orient@s{\the\loopcount}
       \loopcount=\zref@coord@a \advance\loopcount by -\zref@coord@c \xdef\z@orient@s{\the\loopcount}
       \loopcount=\y@orient@r \multiply\loopcount by \z@orient@s
       \temp@count=\z@orient@r \multiply\temp@count by \y@orient@s
       \advance\loopcount by -\temp@count \divide\loopcount by 100 \xdef\x@orient{\the\loopcount}
       \loopcount=\z@orient@r \multiply\loopcount by \x@orient@s
       \temp@count=\x@orient@r \multiply\temp@count by \z@orient@s
       \advance\loopcount by -\temp@count \divide\loopcount by 100 \xdef\y@orient{\the\loopcount}
       \loopcount=\x@orient@r \multiply\loopcount by \y@orient@s
       \temp@count=\y@orient@r \multiply\temp@count by \x@orient@s
       \advance\loopcount by -\temp@count \divide\loopcount by 100 \xdef\z@orient{\the\loopcount}
       \loopcount=\x@orient \multiply\loopcount by \xref@coord@a
       \temp@count=\y@orient \multiply\temp@count by \yref@coord@a
       \advance\loopcount by \temp@count
       \temp@count=\z@orient \multiply\temp@count by \zref@coord@a
       \advance\loopcount by \temp@count 
       \multiply\loopcount by -1\relax 
       \xdef\c@nst@{\the\loopcount}
       \loopcount=\x@orient \multiply\loopcount by \loopcount
       \temp@count=\y@orient \multiply\temp@count by \temp@count
       \advance\loopcount by \temp@count
       \temp@count=\z@orient \multiply\temp@count by \temp@count
       \advance\loopcount by \temp@count
       \xdef\r@@t{\the\loopcount} \intr@@t
       \xdef\f@ct@r{\the\loopcount}
    \fi
    \immediate\openin\structurefile = \PDB@name\relax
    \xdef\first@{}
    \loop
       \read\structurefile to \readline
       \xdef\test@{\expandafter\string\readline}
       \ifx\test@\par@
       \else
          \xdef\structureline{\readline & & & & & & & & &@}
          \expandafter\struc@get\structureline
          \ifx\first@\@TOM
                \xdef\temp@@@{\fifth@ &&@}
                \expandafter\check@letter\temp@@@
                \ifnumber
                      \xdef\PDB@pos{\fifth@}
                      \xdef\sixth@{\sixth@ &.&&@} \expandafter\coord@get\sixth@ \xdef\x@coord{\c@@rd}
                      \xdef\seventh@{\seventh@ &.&&@} \expandafter\coord@get\seventh@ \xdef\y@coord{\c@@rd}
                      \xdef\eighth@{\eighth@ &.&&@} \expandafter\coord@get\eighth@ \xdef\z@coord{\c@@rd}
                \else
                      \xdef\PDB@pos{\sixth@}
                      \xdef\seventh@{\seventh@ &.&&@} \expandafter\coord@get\seventh@ \xdef\x@coord{\c@@rd}
                      \xdef\eighth@{\eighth@ &.&&@} \expandafter\coord@get\eighth@ \xdef\y@coord{\c@@rd}
                      \xdef\nineth@{\nineth@ &.&&@} \expandafter\coord@get\nineth@ \xdef\z@coord{\c@@rd}
                \fi
                \ifnum\PDB@hitpos=\PDB@pos
                \else
                   \ifx\PDB@type\@point@
                      \calc@distance@point
                   \else
                      \ifx\PDB@type\@line@
                         \calc@distance@line
                      \else
                         \ifx\PDB@type\@plane@
                            \calc@distance@plane
                         \fi
                      \fi
                   \fi
                \fi
          \fi
       \fi
    \ifeof\structurefile \xdef\PDB@stack{\PDB@stack\PDB@hitpos,&@}
    \else\repeat
    \closein\structurefile
    
   \fi
   \egroup
 }
 


%%%%% Definition of user commands

\def\clearfuncgroups{\xdef\prfx{func} \clear@groups \xdef\fgroup@num{0}}
\clearfuncgroups
\def\germanlanguage{\germ@ntrue \sp@nishfalse \def\cons@name{Konsensus}}
\def\spanishlanguage{\germ@nfalse \sp@nishtrue \def\cons@name{consenso}}
\def\englishlanguage{\germ@nfalse \sp@nishfalse \def\cons@name{consensus}}
\def\showlegend{\legend@true}
\def\hidelegend{\legend@false}
\def\movelegend#1#2{%
   \setlength\hspace@legend{#1}
   \setlength\vspace@legend{#2}
}
\newcommand{\showcaption}[2][bottom]{\def\cap@pos{#1}\def\c@p{#2}}
\def\shortcaption#1{\def\c@pshort{#1}}
\def\funcgroup#1#2#3#4#5#6{%
   \xdef\first@{#1}
   \loopcount=0 \innerloopcount=0
   \loop
      \advance\loopcount by 1
      \ifx\csname fgroup@name\the\loopcount\endcsname\first@
          \innerloopcount=\loopcount \loopcount=\fgroup@num
      \fi
   \ifnum\loopcount<\fgroup@num \repeat
   \ifnum\innerloopcount=0
     \ifnum\fgroup@num<9
      \innerloopcount=\fgroup@num
      \advance\innerloopcount by 1 \xdef\fgroup@num{\the\innerloopcount}
     \else \message{<Too many \noexpand\funcgroups>}
     \fi
   \fi
   \ifnum\innerloopcount>0
      \expandafter\xdef\csname fgroup@name\the\innerloopcount\endcsname{\first@}
      \expandafter\xdef\csname fg@textcolor\the\innerloopcount\endcsname{#3}
      \expandafter\xdef\csname fg@color\the\innerloopcount\endcsname{#4}
      \expandafter\xdef\csname funcm@tch\the\innerloopcount\endcsname{#5}
      \expandafter\def\csname func@style\the\innerloopcount\endcsname{%
          \csname text#6\endcsname}
      \xdef\prfx{func}
      \xdef\third@{#2&,@} \loopcount=\innerloopcount
      \expandafter\group@get\third@
  \fi}
\def\pepgroups#1{%
   \xdef\prfx{pep}
    \clear@groups
    \xdef\third@{#1&,@} \loopcount=0
    \loop \expandafter\group@get\third@ \advance\loopcount by 1
    \ifnum\loopcount<10 \repeat}
\def\DNAgroups#1{%
   \xdef\prfx{DNA}
    \clear@groups
    \xdef\third@{#1&,@} \loopcount=0
    \loop \expandafter\group@get\third@ \advance\loopcount by 1
    \ifnum\loopcount<10 \repeat}
\def\pepsims#1#2{\xdef\prfx{pep}
   \def\sim@set{\expandafter\residue@get\second@
                 \ifx\first@\ampers@nd
                 \else \advance\innerloopcount by 1
                 
                 \expandafter\xdef\csname simpair\third@\first@\endcsname{1}
                 \expandafter\xdef\csname simpair\first@\third@\endcsname{1}
                 
                 \xdef\second@{\csname sequence\the\loopcount\endcsname} \sim@set
                 \fi}
   \xdef\first@{#1} \make@upper \xdef\third@{\first@}
   \xdef\last@{#2} \xdef\second@{#2 &@} \innerloopcount=0 \sim@set
   \expandafter\xdef\csname \prfx sim\third@\endcsname{%
               (\the\innerloopcount)\last@}}
\def\DNAsims#1#2{\xdef\prfx{DNA}
   \def\sim@set{\expandafter\residue@get\second@
                 \ifx\first@\ampers@nd
                 \else \advance\innerloopcount by 1
                 \xdef\second@{\csname sequence\the\loopcount\endcsname} \sim@set
                 \fi}
   \xdef\first@{#1} \make@upper \xdef\third@{\first@}
   \xdef\last@{#2} \xdef\second@{#2 &@} \innerloopcount=0 \sim@set
   \expandafter\xdef\csname \prfx sim\third@\endcsname{%
               (\the\innerloopcount)\last@}}
\def\fingerprint#1{%
   \ifnum #1 >0
   \residuesperline*{#1}
   \def\finger@linenum{#1}
   \shownames{left}
   \hidenumbering
   \rulersteps{100}
   \nomatchresidues{}{Gray10}{}{}
   \loopcount=0
   \loop
      \advance\loopcount by 1
      \separationline{\the\loopcount}
   \ifnum\loopcount<\seq@count\repeat
   \fi}
\def\printPDBlist#1{%
     \xdef\list@{}
     \xdef\temp@{#1,,,:,,,,@} \expandafter\test@PDB\temp@
     \xdef\first@{\list@ &}
     \ifx\first@\ampers@nd
     \else
        \xdef\first@{\list@ @} 
        \loop
           \expandafter\get@item\first@
           \fourth@
        \ifx\first@@\ampers@nd\else{,\ }\repeat
     \fi
}
\def\messagePDBlist#1{%
     \xdef\list@{}
     \xdef\temp@{#1,,,:,,,,@} \expandafter\test@PDB\temp@
     \xdef\first@{\list@ &}
     \ifx\first@\ampers@nd
     \else
        \message{(#1:}
        \xdef\first@{\list@ @}%
        \loop%
           \expandafter\get@item\first@%
           \message{\fourth@}%
        \ifx\first@@\ampers@nd\else\repeat%
        \message{)}
     \fi
}
\def\shaderegion#1#2#3#4{%
   \regionalshadetrue
   \xdef\seq@{#1}
   \xdef\temp@{consensus} \ifx\seq@\temp@ \xdef\seq@{0} \fi
   \xdef\first@{\seq@ @} \expandafter\check@letter\first@
   \xdef\first@{\seq@}
   \ifletter \get@name@number \xdef\seq@{\first@} \fi
   \ifnum\seq@>\seq@count
     \message{<Ignored `\seq@' in \noexpand\shaderegion or \noexpand\shadeblock>}
   \else
     \ifx\@ll\yes \xdef\@ll{y}\else \xdef\@ll{n}\fi
     \loopcount=\seq@regions
     \advance\loopcount by 1
     \xdef\seq@regions{\the\loopcount}
     \expandafter\xdef\csname fgseqregion\the\loopcount\endcsname{#3}
     \expandafter\xdef\csname bgseqregion\the\loopcount\endcsname{#4}
     \xdef\list@{#2,&}
     \xdef\temp@@{shade}
     \xdef\temp@{#2,,,:,,,,@} \expandafter\test@PDB\temp@
     \ifx\list@\ampers@nd
     \else
        \loop
           \xdef\list@{\list@ @}
           \expandafter\get@regions\list@
        \ifx\list@\ampers@nd\else\repeat
     \fi
   \fi
}
\def\shadeblock#1#2#3#4{%
   \xdef\seq@{#1}
      \xdef\@ll{yes}
      \shaderegion{#1}{#2}{#3}{#4}
      \xdef\@ll{}
}
\def\tintregion#1#2{%
   \xdef\seq@{#1}
   \xdef\temp@{consensus} \ifx\seq@\temp@ \xdef\seq@{0} \fi
   \xdef\first@{\seq@ @} \expandafter\check@letter\first@
   \xdef\first@{\seq@}
   \ifletter \get@name@number \xdef\seq@{\first@} \fi
   \ifnum\seq@>\seq@count
     \message{<Ignored seq `\seq@' in \noexpand\tintregion or \noexpand\tintblock>}
   \else
     \ifx\@ll\yes \xdef\@ll{y}\else \xdef\@ll{n}\fi
     \xdef\list@{#2,&}
     \xdef\temp@@{tint} \loopcount=0
     \xdef\temp@{#2,,,:,,,,@} \expandafter\test@PDB\temp@
     \ifx\list@\ampers@nd
     \else
        \loop
           \xdef\list@{\list@ @}
           \expandafter\get@tintregions\list@
        \ifx\list@\ampers@nd\else\repeat
     \fi
   \fi
}
\def\tintblock#1#2{%
   \xdef\seq@{#1}
      \xdef\@ll{yes}
      \tintregion{#1}{#2}
      \xdef\@ll{}
}
\def\tintdefault#1{%
   \xdef\first@{#1}
   \xdef\second@{strong}
   \ifx\first@\second@
      \xdef\light@{LightLightLight}
   \else
   \xdef\second@{medium}
   \ifx\first@\second@
      \xdef\light@{LightLight}
   \else
   \xdef\second@{weak}
   \ifx\first@\second@
      \xdef\light@{Light}
   \else
      \xdef\light@{LightLight}
   \fi\fi\fi
}
\def\lowerregion#1#2{%
   \xdef\seq@{#1}
   \xdef\temp@{consensus} \ifx\seq@\temp@ \xdef\seq@{0} \fi
   \xdef\first@{\seq@ @} \expandafter\check@letter\first@
   \xdef\first@{\seq@}
   \ifletter \get@name@number \xdef\seq@{\first@} \fi
   \ifnum\seq@>\seq@count
     \message{<Ignored seq `\seq@' in \noexpand\lowerregion or \noexpand\lowerblock>}
   \else
     \ifx\@ll\yes \xdef\@ll{y}\else \xdef\@ll{n}\fi
     \xdef\list@{#2,&}
     \xdef\temp@@{lower} \loopcount=0
     \xdef\temp@{#2,,,:,,,,@} \expandafter\test@PDB\temp@
     \ifx\list@\ampers@nd
     \else
        \loop
           \xdef\list@{\list@ @}
           \expandafter\get@lowerregions\list@
        \ifx\list@\ampers@nd\else\repeat
     \fi
   \fi
}
\def\lowerblock#1#2{%
   \xdef\seq@{#1}
      \xdef\@ll{yes}
      \lowerregion{#1}{#2}
      \xdef\@ll{}
}
\def\emphregion#1#2{%
   \xdef\seq@{#1}
   \xdef\temp@{consensus} \ifx\seq@\temp@ \xdef\seq@{0} \fi
   \xdef\first@{\seq@ @} \expandafter\check@letter\first@
   \xdef\first@{\seq@}
   \ifletter \get@name@number \xdef\seq@{\first@} \fi
   \ifnum\seq@>\seq@count
     \message{<Ignored seq `\seq@' in \noexpand\emphregion or \noexpand\emphblock>}
   \else
     \ifx\@ll\yes \xdef\@ll{y}\else \xdef\@ll{n}\fi
     \xdef\list@{#2,&}
     \xdef\temp@@{emph} \loopcount=0
     \xdef\temp@{#2,,,:,,,,@} \expandafter\test@PDB\temp@
     \ifx\list@\ampers@nd
     \else
        \loop
           \xdef\list@{\list@ @}
           \expandafter\get@emphregions\list@
        \ifx\list@\ampers@nd\else\repeat
     \fi
   \fi
}
\def\emphblock#1#2{%
   \xdef\seq@{#1}
      \xdef\@ll{yes}
      \emphregion{#1}{#2}
      \xdef\@ll{}
}
\def\emphdefault#1{\def\res@style{\csname text#1\endcsname}}
\def\frameblock#1#2#3{%
   \xdef\seq@{#1}
   \xdef\temp@{consensus} \ifx\seq@\temp@ \xdef\seq@{0} \fi
   \xdef\first@{\seq@ @} \expandafter\check@letter\first@
   \xdef\first@{\seq@}
   \ifletter \get@name@number \xdef\seq@{\first@} \fi
   \ifnum\seq@>\seq@count
     \message{<Ignored seq `\seq@' in \noexpand\frameblock>}
   \else
     \xdef\@ll{#3}
     \xdef\list@{#2,&}
     \xdef\temp@@{frame} \loopcount=0
     \xdef\temp@{#2,,,:,,,,@} \expandafter\test@PDB\temp@
     \ifx\list@\ampers@nd
     \else
        \loop
           \xdef\list@{\list@ @}
           \expandafter\get@frameregions\list@
        \ifx\list@\ampers@nd\else\repeat
     \fi
   \fi
}
\def\changeshadingcolors#1#2#3{%
   \xdef\seq@{#1}
   \xdef\temp@{consensus} \ifx\seq@\temp@ \xdef\seq@{0} \fi
   \xdef\first@{\seq@ @} \expandafter\check@letter\first@
   \xdef\first@{\seq@}
   \ifletter \get@name@number \xdef\seq@{\first@} \fi
   \ifnum\seq@>\seq@count
     \message{<Ignored seq `\seq@' in \noexpand\changeshadingcolors>}
   \else
     \xdef\@ll{#3}
     \xdef\list@{#2,&}
     \xdef\temp@@{shading} \loopcount=0
     \xdef\temp@{#2,,,:,,,,@}    \expandafter\test@PDB\temp@
     \ifx\list@\ampers@nd
     \else
        \loop
           \xdef\list@{\list@ @}
           \expandafter\get@shadingregions\list@
        \ifx\list@\ampers@nd\else\repeat
     \fi
   \fi
}
\def\bargraphstretch#1{\def\b@r@stretch{#1}}
\def\colorscalestretch#1{\def\sc@le@stretch{#1}}
\def\rm@@measure#1#2@{%
   \xdef\test@{#1}
   \ifx\test@\ampers@nd
   \else
      \expandafter\check@char\test@
      \ifnumber\xdef\first@{\first@ #1}
         \xdef\second@{#2 &@}
         \expandafter\rm@@measure\second@
      \fi
   \fi
}
\def\rm@measure#1.#2@{\xdef\first@{#1.}\xdef\second@{#2 @}\xdef\third@{#1}\expandafter\rm@@measure\second@}
\def\pm@calc{%
    \temp@@length=100000sp
    \temp@@length=\g@min\temp@@length
    \innerloopcount=\temp@@length
    \xdef\min@{\the\innerloopcount}
    \arrow@height=\temp@@length
    \temp@@length=100000sp
    \temp@@length=\g@max\temp@@length
    \advance\temp@@length by -\arrow@height
    \innerloopcount=\temp@@length
    \divide\innerloopcount by 100
    \ifnum\innerloopcount=0 \innerloopcount=1 \fi
    \xdef\m@x{\the\innerloopcount}
    \xdef\test@{\g@min pt}
    \setlength\arrow@width{\test@}
    \xdef\test@{\g@max pt}
    \setlength\arrow@height{\test@}
    \advance\arrow@height by -\arrow@width
    \ifdim\arrow@width<0pt\temp@@length=-\arrow@width\xdef\test@{y}\else\temp@@length=\arrow@width\xdef\test@{n}\fi
    \ifdim\arrow@height>0pt
       \ifx\test@\n@ \xdef\test@{y} \else \xdef\test@{n} \fi
       \ifdim\temp@@length<\arrow@height\temp@@length=\arrow@height\fi
    \else
       \ifdim\temp@@length>-\arrow@height\temp@@length=\arrow@height\fi
    \fi
    \ifdim\temp@@length<100pt\arrow@width=100\arrow@width\arrow@height=100\arrow@height\else
    \ifdim\temp@@length<10pt\arrow@width=1000\arrow@width\arrow@height=1000\arrow@height\else
    \ifdim\temp@@length<1pt\arrow@width=10000\arrow@width\arrow@height=10000\arrow@height\else
    \ifdim\temp@@length<0.1pt\arrow@width=100000\arrow@width\arrow@height=100000\arrow@height\else
    \ifdim\temp@@length<0.01pt\arrow@width=1000000\arrow@width \arrow@height=1000000\arrow@height\else
    \ifdim\temp@@length<0.001pt\arrow@width=10000000\arrow@width \arrow@height=10000000\arrow@height\else
    \ifdim\temp@@length<0.0001pt\arrow@width=100000000\arrow@width \arrow@height=100000000\arrow@height\else
    \ifdim\temp@@length<0.00001pt\arrow@width=1000000000\arrow@width \arrow@height=1000000000\arrow@height
    \fi\fi\fi\fi\fi\fi\fi\fi
    \ifx\test@\y@
       \xdef\pm@{0}
    \else
       \xdef\test@{-\the\arrow@height @}
       \expandafter\rm@measure\test@
       \divide\arrow@width by \third@
       \xdef\pm@{\the\arrow@width @}
       \expandafter\rm@measure\pm@
       \xdef\pm@{\first@}
    \fi
}
\def\read@graph{%
    \bgroup
    \immediate\openin\structurefile = \fill@char\relax
    \ifeof\structurefile
      \PackageError{TeXshade}
      {File `\fill@char' not found}
      {\MessageBreak
       The file you specified is missing or you have \MessageBreak
       misspelled it. \MessageBreak\MessageBreak
       No feature graph will be displayed. \MessageBreak
       Type <return> to proceed. \MessageBreak
       Type X <return> to quit. \MessageBreak
}
       \immediate\closein\structurefile\egroup
    \else
    \ifx\g@min\comm@
    \def\par@{}
    \xdef\g@min{,} \xdef\g@max{,}
    \loop
       \read\structurefile to \readline
       \xdef\test@{\expandafter\string\readline}
       \ifx\test@\par@
       \else
       \xdef\second@{\test@ @}
       \expandafter\check@letter\second@
       \ifletter
          \xdef\second@{\expandafter\string\readline}
          \xdef\second@{\second@ @}
          \expandafter\firstchar@get\second@
          \if\first@ - \numbertrue\fi
       \fi
       \ifnumber
          \temp@@length=1pt
          \temp@@length=\test@\temp@@length
          \innerloopcount=\temp@@length
          \ifx\g@min\comm@ \xdef\min@{\test@ pt} \xdef\g@min{\test@}
          \else
             \ifdim\temp@@length<\min@\relax \xdef\min@{\the\temp@@length} \xdef\g@min{\test@} \fi\fi%
          \ifx\g@max\comm@ \xdef\m@x{\test@ pt} \xdef\g@max{\test@}
          \else
             \ifdim\temp@@length>\m@x\relax \xdef\m@x{\the\temp@@length} \xdef\g@max{\test@} \fi\fi%
       \fi
       \fi
    \ifeof\structurefile\else\repeat
    \fi
    \immediate\closein\structurefile
    \pm@calc
    \expandafter\temp@count=\csname seq@start\seq@\endcsname
    \advance\temp@count by -1
    \xdef\temp@@@{n}
    \immediate\openin\structurefile = \fill@char\relax
    \loop
       \read\structurefile to \readline
       \xdef\test@{\expandafter\string\readline}
       \ifx\test@\par@
       \else
       \xdef\second@{\test@ @}
       \expandafter\check@letter\second@
       \ifletter
          \xdef\second@{\expandafter\string\readline}
          \xdef\second@{\second@ @}
          \expandafter\firstchar@get\second@
          \if\first@ - \numbertrue
          \else
             \if\first@ N
                \expandafter\firstchar@get\third@
                \if\first@ a
                   \expandafter\firstchar@get\third@
                   \if\first@ N
                     \advance\temp@count by 1 
                     \ifnum\temp@count=0 \temp@count=1 \fi
                     \ifnum\temp@count<\st@rt
                     \else
                        \ifnum\temp@count>\st@p
                        \else
                           \ifx\temp@@@\n@
                              \xdef\temp@@@{N}
                           \else
                              \xdef\temp@@@{\temp@@@,N}
                           \fi
                        \fi
                     \fi
                   \fi
                \fi   
             \fi
          \fi
       \fi
       \ifnumber
          \advance\temp@count by 1 
          \ifnum\temp@count=0 \temp@count=1 \fi
          \ifnum\temp@count<\st@rt
          \else
             \ifnum\temp@count>\st@p
             \else
                \temp@@length=100000sp
                \temp@@length=\test@\temp@@length
                \innerloopcount=\temp@@length
                \advance\innerloopcount by 1
                \xdef\test@{\pm@ pt}
                \ifdim\test@=0pt
                   \advance\innerloopcount by -\min@
                \else
                   \ifx\b@r\n@
                      \advance\innerloopcount by -\min@
                   \fi
                \fi
                \divide\innerloopcount by \m@x
                \ifx\temp@@@\n@
                   \xdef\temp@@@{\the\innerloopcount}
                \else
                   \xdef\temp@@@{\temp@@@,\the\innerloopcount}
                \fi
             \fi
          \fi
       \fi
       \fi
    \ifeof\structurefile\else\repeat
    \immediate\closein\structurefile
    \xdef\temp@@@{\temp@@@,@} \xdef\temp@@{y}
    \fi
    \egroup
}
\def\sort@gstack{%
   \expandafter\get@fromstack\last@
   \ifx\first@\ampers@nd
      \xdef\tmpstack{\tmpstack\f@text@&;&;&;&;@}
   \else
   \ifnum\loopcount<\second@
      \xdef\tmpstack{\tmpstack\f@text@\first@;\second@;\third@;\fourth@;\last@}
   \else
      \xdef\tmpstack{\tmpstack\first@;\second@;\third@;\fourth@;}
      \sort@gstack
   \fi\fi
}
\def\get@gregion#1..#2,#3&{%
   \xdef\st@rt{#1} \xdef\st@p{#2} \xdef\list@{#3}
}
\def\do@bargraph{%
   \expandafter\get@gregion\list@
   \expandafter\ifnum\csname seq@start\seq@\endcsname>\st@rt
   \else
   \xdef\temp@@{n}
   \xdef\b@r{y}
   \xdef\pm@{0}
   \xdef\temp@@@{\fill@char,@}
   \expandafter\check@letter\temp@@@
   \ifletter 
      \read@graph
   \else
      \ifx\g@min\comm@ \xdef\g@min{0} \fi
      \ifx\g@max\comm@ \xdef\g@max{100} \fi
      \pm@calc
      \xdef\temp@@@{\fill@char,@}
      \xdef\temp@@{y}
   \fi
   \ifx\temp@@\y@
   \loopcount=\st@rt
   \xdef\tmpstack{}
   \loop
      \expandafter\get@item\temp@@@ 
      \xdef\temp@@@{\first@}
      \xdef\style@{bar[\pm@,0]:\fourth@[\f@color]}
      \xdef\tmpstack{\tmpstack\f@text@;\the\loopcount;\the\loopcount;\style@;}
      \advance\loopcount by 1
      \ifnum\loopcount=0 \loopcount=1 \fi
   \ifnum\loopcount>\st@p
   \else\repeat
      \xdef\f@text@{\tmpstack} \xdef\tmpstack{}
      \xdef\last@{\csname stack@\bottop@\seq@\endcsname}
      \sort@gstack
      \expandafter\xdef\csname stack@\bottop@\seq@\endcsname{\tmpstack}
   \fi
   \fi
   \xdef\list@{\list@ &}
   \ifx\list@\ampers@nd\else\do@bargraph\fi
}
\def\do@colorgraph{%
   \expandafter\get@gregion\list@
   \expandafter\ifnum\csname seq@start\seq@\endcsname>\st@rt
   \else
   \xdef\temp@@{n}
   \xdef\b@r{n}
   \xdef\pm@{0}
   \xdef\temp@@@{\fill@char,@}
   \expandafter\check@letter\temp@@@
   \ifletter 
      \read@graph
   \else
      \ifx\g@min\comm@ \xdef\g@min{0} \fi
      \ifx\g@max\comm@ \xdef\g@max{100} \fi
      \pm@calc
      \xdef\temp@@@{\fill@char,@}
      \xdef\temp@@{y}
   \fi
   \ifx\temp@@\y@
   \loopcount=\st@rt
   \xdef\tmpstack{}
   \loop
      \xdef\last@{\csname stack@\bottop@\seq@\endcsname}
      \expandafter\get@item\temp@@@ 
      \xdef\temp@@@{\first@}
      \ifx\fourth@\N@
         \xdef\style@{color:50[White]}
      \else
         \ifnum\fourth@<1 \xdef\fourth@{1} \fi
         \innerloopcount=\fourth@
         \advance\innerloopcount by 4
         \divide\innerloopcount by 5
         \multiply\innerloopcount by 5
         \ifnum\innerloopcount>100 \innerloopcount=100 \fi
         \xdef\style@{color:50[\f@color\the\innerloopcount]}
      \fi    
      \xdef\tmpstack{\tmpstack\f@text@;\the\loopcount;\the\loopcount;\style@;}
      \advance\loopcount by 1
      \ifnum\loopcount=0 \loopcount=1 \fi
   \ifnum\loopcount>\st@p
   \else\repeat
      \xdef\f@text@{\tmpstack} \xdef\tmpstack{}
      \xdef\last@{\csname stack@\bottop@\seq@\endcsname}
      \sort@gstack
      \expandafter\xdef\csname stack@\bottop@\seq@\endcsname{\tmpstack}
   \fi
   \fi
   \xdef\list@{\list@ &}
   \ifx\list@\ampers@nd\else\do@colorgraph\fi
}
\def\feature#1#2#3#4#5{%
   \xdef\bottop@{#1}
   \xdef\temp@{top}
   \ifx\bottop@\temp@ \topfeaturetrue\fi
   \xdef\temp@{ttop}
   \ifx\bottop@\temp@ \ttopfeaturetrue\fi
   \xdef\temp@{tttop}
   \ifx\bottop@\temp@ \tttopfeaturetrue\fi
   \xdef\temp@{ttttop}
   \ifx\bottop@\temp@ \ttttopfeaturetrue\fi
   \xdef\temp@{bottom}
   \ifx\bottop@\temp@ \bottomfeaturetrue\fi
   \xdef\temp@{bbottom}
   \ifx\bottop@\temp@ \bbottomfeaturetrue\fi
   \xdef\temp@{bbbottom}
   \ifx\bottop@\temp@ \bbbottomfeaturetrue\fi
   \xdef\temp@{bbbbottom}
   \ifx\bottop@\temp@ \bbbbottomfeaturetrue\fi
   \xdef\seq@{#2}
   \xdef\temp@{consensus} \ifx\seq@\temp@ \xdef\seq@{0} \fi
   \xdef\first@{\seq@ @} \expandafter\check@letter\first@
   \xdef\first@{\seq@}
   \ifletter \get@name@number \xdef\seq@{\first@} \fi
   \ifnum\seq@>\seq@count
      \message{<Ignored seq `\seq@' in \noexpand\feature>}
   \else
   \ifnum\seq@>-1
   \xdef\temp@{#4::&}\expandafter\test@fill\temp@
   \xdef\last@{bar}
   \ifx\second@@\last@
      \xdef\last@{hydrophobicity}
      \ifx\last@\fourth@
         \xdef\second@@{bh}
      \else
         \xdef\last@{molweight}
         \ifx\last@\fourth@
            \xdef\second@@{bm}
         \else
            \xdef\last@{charge}
            \ifx\last@\fourth@
               \xdef\second@@{bc}
            \else
               \xdef\last@{conservation}
               \ifx\last@\fourth@
                  \xdef\second@@{bcons}
               \fi
            \fi
         \fi
      \fi
   \fi
   \xdef\last@{color}
   \ifx\second@@\last@
      \xdef\last@{hydrophobicity}
      \ifx\last@\fourth@
         \xdef\second@@{ch}
      \else
         \xdef\last@{molweight}
         \ifx\last@\fourth@
            \xdef\second@@{cm}
         \else
            \xdef\last@{charge}
            \ifx\last@\fourth@
               \xdef\second@@{cc}
            \else
               \xdef\last@{conservation}
               \ifx\last@\fourth@
                  \xdef\second@@{ccons}
               \fi
            \fi
         \fi
      \fi
   \fi
   \xdef\last@{bar}
   \ifx\second@@\last@
      \xdef\list@{#3,&}
      \xdef\style@{#4}
      \def\f@text@{#5}
      \do@bargraph
      \xdef\temp@{bottom}
      \ifx\bottop@\temp@ \xdef\bottom@stretch{y}\fi
      \xdef\temp@{bbottom}
      \ifx\bottop@\temp@ \xdef\bbottom@stretch{y}\fi
      \xdef\temp@{bbbottom}
      \ifx\bottop@\temp@ \xdef\bbbottom@stretch{y}\fi
      \xdef\temp@{bbbbottom}
      \ifx\bottop@\temp@ \xdef\bbbbottom@stretch{y}\fi
   \else
   \xdef\last@{color}
   \ifx\second@@\last@
      \xdef\list@{#3,&}
      \xdef\style@{#4}
      \def\f@text@{#5}
      \do@colorgraph
   \else
   \xdef\f@@color{\f@color}
   \def\f@text@{#5}
   \xdef\f@color{\f@@color}
   \xdef\temp@{#4&}
   \ifx\temp@\ampers@nd
      \xdef\list@{#3,&}
      \xdef\style@{&}
      \def\f@text@{#5}
   \else
      \xdef\last@{restriction}
      \ifx\second@@\last@
         \xdef\temp@{\bottop@ @}
         \expandafter\firstchar@get\temp@
         \xdef\temp@{y}
         \if\first@ t
            \xdef\style@{fill:\kern0.9\box@width$\blacktriangledown$[\f@color]}
         \else
            \xdef\style@{fill:\kern0.9\box@width$\blacktriangle$[\f@color]}
         \fi
         \xdef\f@text@{\kern0.9\box@width#5}
      \else
      \xdef\last@{bh}
      \ifx\second@@\last@
         \ifx\f@color\gr@ydef@ult\xdef\f@color{Gray50}\fi
         \xdef\style@{plot[bar]:Hydro[\f@color][-53]}
      \else
      \xdef\last@{bm}
      \ifx\second@@\last@
         \ifx\f@color\gr@ydef@ult\xdef\f@color{Gray50}\fi
         \xdef\style@{plot[bar]:molw[\f@color][0]}
      \else
      \xdef\last@{bc}
      \ifx\second@@\last@
         \ifx\f@color\gr@ydef@ult\xdef\f@color{Gray50}\fi
         \xdef\style@{plot[bar]:charge[\f@color][-50]}
      \else
      \xdef\last@{bcons}
      \ifx\second@@\last@
         \ifx\f@color\gr@ydef@ult\xdef\f@color{Gray50}\fi
         \xdef\last@{T-Coffee}
         \ifx\f@color\last@
            \xdef\T@coffee@bcons{y}
            \xdef\f@color{Gray50}
         \fi
         \xdef\style@{cons[bar]:cons[\f@color][0]}
      \else
      \xdef\last@{ch}
      \ifx\second@@\last@
         \ifx\f@color\gr@ydef@ult\xdef\f@color{RedGreen}\fi
         \xdef\style@{plot[color]:Hydro[\f@color][-53]}
      \else
      \xdef\last@{cm}
      \ifx\second@@\last@
         \ifx\f@color\gr@ydef@ult\xdef\f@color{Gray}\fi
         \xdef\style@{plot[color]:molw[\f@color][0]}
      \else
      \xdef\last@{cc}
      \ifx\second@@\last@
         \ifx\f@color\gr@ydef@ult\xdef\f@color{RedBlue}\fi
         \xdef\style@{plot[color]:charge[\f@color][-50]}
      \else
      \xdef\last@{ccons}
      \ifx\second@@\last@
         \ifx\f@color\gr@ydef@ult\xdef\f@color{ColdHot}\fi
         \xdef\last@{T-Coffee}
         \ifx\f@color\last@
            \xdef\T@coffee@ccons{y}
            \xdef\f@color{TC}
         \fi
         \xdef\style@{cons[color]:cons[\f@color][0]}
      \else
         \xdef\style@{#4} \xdef\temp@@{n}
         \expandafter\get@firstfill\temp@
         \if\second@@ ^ \xdef\second@@{_} \fi
         \if\second@@ _ \xdef\temp@{\fill@char} \xdef\temp@@{y}\fi
         \expandafter\getarrow@shape\temp@
      \fi
      \fi
      \fi
      \fi
      \fi
      \fi
      \fi
      \fi
      \fi
   \xdef\list@{#3,&}
   \xdef\@ll{\f@text@}
   \xdef\temp@@{\bottop@}
   \xdef\temp@{#3,,,:,,,,@} \expandafter\test@PDB\temp@
   \fi
   \loopcount=0 
   \ifx\list@\ampers@nd
   \else
      \loop
         \advance\loopcount by 1\relax
         \xdef\list@{\list@ @}
         \expandafter\get@fregions\list@
      \ifx\list@\ampers@nd\else\repeat
   \fi
   \fi\fi\fi
   \fi
}
\def\showfeaturename#1#2{\expandafter\xdef\csname featuretextn@me#1\endcsname{#2}}
\def\showfeaturestylename#1#2{\expandafter\xdef\csname featurestylesn@me#1\endcsname{#2}}
\def\hidefeaturename#1{\expandafter\xdef\csname featuretextn@me#1\endcsname{}}
\def\hidefeaturenames{\xdef\featuretextn@mettop{}\xdef\featuretextn@metop{}\xdef\featuretextn@mebottom{}\xdef\featuretextn@mebbottom{}
    \xdef\featuretextn@metttop{}\xdef\featuretextn@mettttop{}\xdef\featuretextn@mebbbbottom{}\xdef\featuretextn@mebbbottom{}}
\def\hidefeaturestylename#1{\expandafter\xdef\csname featurestylesn@me#1\endcsname{}}
\def\hidefeaturestylenames{\xdef\featurestylesn@mettop{}\xdef\featurestylesn@metop{}\xdef\featurestylesn@mebottom{}\xdef\featurestylesn@mebbottom{}
    \xdef\featurestylesn@metttop{}\xdef\featurestylesn@mettttop{}\xdef\featurestylesn@mebbbbottom{}\xdef\featurestylesn@mebbbottom{}}
\def\seqtype#1{\xdef\seq@type{#1}
               \if\seq@type P \xdef\prefix@{pep}
               \else \if\seq@type p \xdef\seq@type{P} \xdef\prefix@{pep}
               \else \xdef\seq@type{N} \xdef\prefix@{DNA} \fi\fi}
\def\nameseq#1#2{%
   \xdef\first@{#1 @} \expandafter\check@letter\first@
   \xdef\first@{#1}   \ifletter \get@name@number \fi
   \expandafter\xdef\csname newseqname\first@\endcsname{#2}
}
\newcommand\threshold[2][n]{%
   \xdef\first@{#1}
   \ifx\first@\n@
   \xdef\thresh@ld{#2}
   \else
      \all@shadetrue
      \ifnum\first@>#2
         \xdef\thresh@ld{#2}
         \xdef\all@thresh@ld{#1}
      \else
         \xdef\thresh@ld{#1}
         \xdef\all@thresh@ld{#2}
      \fi
   \fi
}
\def\constosingleseq#1{%
   \xdef\first@{#1 @} \expandafter\check@letter\first@
   \xdef\first@{#1}   \ifletter \get@name@number \fi
   \ifnum\first@>\seq@count
   \message{<Ignored seq `#1' in \noexpand\constosingleseq>}
   \else
   \ifnum\first@>0 \xdef\cons@num{\first@} \hideconsensus\fi\fi
}
\def\constoallseqs{\xdef\cons@num{0}}
\def\residuesperline{%
   \def\@rplfix*##1{%
      \res@perline=##1
      \ifnum\res@perline<1 \res@perline=1\fi \rpl@fixtrue}
   \def\@rplvar ##1{%
      \res@perline=##1
      \ifnum\res@perline<5 \res@perline=5\fi \rpl@fixfalse}
   \def\decide@{\ifx\l@@k * \let\next\@rplfix \else \let\next\@rplvar \fi \next}
   \futurelet\l@@k\decide@}
\def\numberingwidth#1{\def\num@width{#1}}
\def\charstretch#1{%
   \def\char@stretch{#1}
   \xdef\temp@{\char@stretch .00@} \expandafter\coord@get\temp@ \loopcount=\c@@rd
   \multiply \loopcount by 7
   \divide \loopcount by 10
   \ifnum\loopcount<10
      \xdef\temp@{0\the\loopcount @}
   \else
      \xdef\temp@{\the\loopcount @}
   \fi
   \expandafter\decimal@B\temp@
   \ifnum\loopcount>99
      \xdef\char@stretch@W{#1}
   \else
      \xdef\char@stretch@W{\temp@}
   \fi
}
\def\linestretch#1{\def\line@stretch{#1}}
\def\logostretch#1{%
   \def\logo@stretch{#1}
   \setlength\temp@@length{1000sp}
   \setlength\temp@@length{\logo@stretch\temp@@length}
   \loopcount=\temp@@length
   \xdef\logo@stretch@IOOO{\the\loopcount}
}
\def\noblockskip{\def\block@skip{\vspace{0pt}}}
\def\smallblockskip{\def\block@skip{\vspace{\baselineskip}}}
\def\medblockskip{\def\block@skip{\vspace{1.5\baselineskip}}}
\def\bigblockskip{\def\block@skip{\vspace{2\baselineskip}}}
\def\vblockspace#1{\def\block@skip{\vspace{#1}}}
\def\topspace#1{\def\t@sp@ce{#1}}
\def\ttopspace#1{\def\tt@sp@ce{#1}}
\def\tttopspace#1{\def\ttt@sp@ce{#1}}
\def\ttttopspace#1{\def\tttt@sp@ce{#1}}
\def\bottomspace#1{\def\b@sp@ce{#1}}
\def\bbottomspace#1{\def\bb@sp@ce{#1}}
\def\bbbottomspace#1{\def\bbb@sp@ce{#1}}
\def\bbbbottomspace#1{\def\bbbb@sp@ce{#1}}
\def\rulerspace#1{\def\ruler@sp@ce{#1}}
\def\fixblockspace{\fix@true}
\def\flexblockspace{\fix@false}
\def\nosepline{\def\seq@skip{\relax}}
\def\smallsepline{\def\seq@skip{\vspace{3pt}}\def\sep@space{3pt}}
\def\medsepline{\def\seq@skip{\vspace{6pt}}\def\sep@space{6pt}}
\def\bigsepline{\def\seq@skip{\vspace{12pt}}\def\sep@space{12pt}}
\def\vsepspace#1{\def\seq@skip{\vspace{#1}}\xdef\sep@space{#1}}
\def\separationline#1{%
   \xdef\first@@{#1}
   \xdef\first@{#1 @} \expandafter\check@letter\first@
   \xdef\first@{#1}
   \ifletter \get@name@number \xdef\first@@{\first@} \fi
   \ifnum\first@@>\seq@count \xdef\first@@{1}
   \else
      \ifnum\first@@<0 \xdef\first@@{1}
      \else
         \expandafter\def\csname seq@gap\first@@\endcsname{yes}
         \loopcount=\seq@gap@num
         \advance\loopcount by 1
         \xdef\seq@gap@num{\the\loopcount}
      \fi
   \fi}
\newcommand{\shadingmode}[2][-1]{%
   \T@coffeefalse
   \xdef\last@{#2}
   \xdef\first@{identical}
   \ifx\first@\last@
      \simmodefalse \funcmodefalse
      \xdef\second@{#1 @} \expandafter\check@letter\second@
      \ifnumber \all@shadetrue 
         \ifnum#1<0 \all@shadefalse
         \else
         \ifnum#1>100 \xdef\all@thresh@ld{100}
         \else 
            \xdef\all@thresh@ld{#1}
         \fi\fi
      \else
         \xdef\last@{#1}
         \xdef\first@{allmatchspecial}
         \ifx\last@\first@
            \xdef\all@thresh@ld{100} \all@shadetrue
         \fi
      \fi
   \else
   \xdef\first@{similar}
   \ifx\first@\last@
      \simmodetrue  \funcmodefalse
      \xdef\second@{#1 @} \expandafter\check@letter\second@
      \ifnumber \all@shadetrue 
         \ifnum#1<0 \all@shadefalse
         \else
         \ifnum#1>100 \xdef\all@thresh@ld{100}
         \else 
            \xdef\all@thresh@ld{#1}
         \fi\fi
      \else
         \xdef\last@{#1}
         \xdef\first@{allmatchspecial}
         \ifx\last@\first@
            \xdef\all@thresh@ld{100} \all@shadetrue
         \fi
      \fi
   \else
   \xdef\first@{functional}
   \ifx\first@\last@
      \if\seq@type N \message{<No functional shading on DNA sequences>}
      \else \simmodefalse \funcmodetrue \func@shading{#1} 
            \xdef\seq@type{P} \xdef\prefix@{pep} \fi
   \else
   \xdef\first@{T-Coffee}
   \ifx\first@\last@
      \simmodefalse \funcmodefalse \T@coffeetrue 
      \xdef\second@{#1 @} \expandafter\check@letter\second@
      \ifletter \xdef\TC@first@{#1}\include@T@coffee\fi
   \else
   \xdef\first@{diverse}
   \ifx\first@\last@
      \xdef\last@{#1}
      \ifnum\last@>\seq@count \xdef\last@{1}\fi
      \ifnum\last@<1 \xdef\last@{1}\fi
      \simmodetrue  \funcmodefalse
      \threshold{0}
      \xdef\divref@{\last@}
      \constosingleseq{\last@}
      \nomatchresidues{Black}{White}{lower}{up}
      \similarresidues{Black}{White}{lower}{up}
      \conservedresidues{Black}{White}{{.}}{up}
      \allmatchresidues{Black}{White}{{.}}{up}
      \gapchar{-}
      \hideconsensus
   \else
   \message{<Unknown shading mode - using `similar'>}
   \simmodetrue  \funcmodefalse
   \fi\fi\fi\fi\fi}
\def\hideallmatchpositions{\xdef\all@out{y}}
\def\showallmatchpositions{\xdef\all@out{n}}
\newcommand\allmatchspecial[1][100]{%
   \ifnum#1<0
      \xdef\all@thresh@ld{0}
   \else
   \ifnum#1>100
      \xdef\all@thresh@ld{100}
   \else
      \xdef\all@thresh@ld{#1}
   \fi\fi
   \all@shadetrue}
\def\allmatchspecialoff{\all@shadefalse}
\def\stopchar#1{\def\st@p@char{#1}}
\def\gapchar#1{%
   \xdef\first@{rule}\xdef\second@{#1}
   \ifx\first@\second@\def\gap@char{o}
   \else\def\gap@char{#1}\fi}
\def\gaprule#1{\def\gap@rulethick{#1}}
\def\domaingaprule#1{\def\domgap@rulethick{#1}}
\newcommand{\setends}[3][&]{%
  \xdef\start@seq{#2}
  \xdef\temp@{consensus}
  \ifx\start@seq\temp@ 
%%%     \message{<\noexpand\setends does not accept `consensus'>}
     \xdef\start@seq{0}
%%%%%%%%%%
            \xdef\second@{#3@} \expandafter\get@nums\second@ 
            \xdef\start@num{\first@} \xdef\end@num{\second@}
            \ifnum\start@num<1 \xdef\start@num{1}\fi
            \ifnum\end@num<\start@num \xdef\end@num{\start@num}\fi
            \loopcount=\end@num
            \advance\loopcount by 1
            \advance\loopcount by -\start@num
            \xdef\end@num{\the\loopcount}
            \start@false
%%%%%%%%%%
  \else
     \xdef\first@{\start@seq @} \expandafter\check@letter\first@
     \xdef\first@{\start@seq}
     \ifletter \get@name@number \xdef\start@seq{\first@} \fi
     \ifnum\start@seq>\seq@count 
         \message{<\noexpand\setends{} error: sequence `#2' not defined>}
         \xdef\start@seq{0}
     \else
         \ifnum\start@seq<1
            \message{<\noexpand\setends{} error: sequence `#2' not defined>}
            \xdef\start@seq{0}
         \else
            \xdef\second@{#3@} \expandafter\get@nums\second@ 
            \xdef\start@num{\first@} \xdef\end@num{\second@}
            \ifnum\start@num=0 \xdef\allow@zero{y}\fi
            \ifnum\end@num=0 \xdef\allow@zero{y}\fi
            \start@false
         \fi
     \fi
  \fi
  \xdef\first@{#1}\ifx\first@\ampers@nd\else\startnumber{#2}{#1}\fi
}
\newcommand{\startnumber}[3][&]{%
   \xdef\first@{#2} \xdef\second@{consensus}
   \ifx\first@\second@
      \ifnum#3=0 \xdef\allow@zero{y} \fi
      \expandafter\xdef\csname seq@start0\endcsname{#2}
      \cons@count=#3 \advance\cons@count by -1\relax 
      \expandafter\xdef\csname res@count0\endcsname{\the\cons@count}
   \else
      \xdef\first@{#2 @} \expandafter\check@letter\first@
      \xdef\first@{#2}   \ifletter \get@name@number \fi
      \ifnum\first@>\seq@count \message{<Ignored seq `#2' in \noexpand\startnumber>}
      \else
         \xdef\second@{#3} 
         \ifnum\second@=0 \xdef\allow@zero{y} \fi
         \expandafter\xdef\csname seq@start\first@\endcsname{\second@}
         \res@count=\second@
         \advance\res@count by -1
         \expandafter\xdef\csname res@count\first@\endcsname{\the\res@count}
      \fi
   \fi
   \xdef\first@{#1}\ifx\first@\ampers@nd\else\setends{#2}{#1}\fi
   }
\def\seqlength#1#2{%
   \xdef\first@{#1 @} \expandafter\check@letter\first@
   \xdef\first@{#1}   \ifletter \get@name@number \fi
   \ifnum\first@>\seq@count \message{<Ignored seq `#1' in \noexpand\seqlength>}
   \else
      \xdef\second@{#2} \ifnum\second@<0 \xdef\second@{1} \fi
      \expandafter\xdef\csname seq@len\first@\endcsname{\second@}
   \fi}
\newcommand\shownumbering[2][n]{%
   \xdef\first@{#1}\ifx\first@\n@\else\xdef\numbering@fg{#1}\fi
   \xdef\first@{#2}
   \xdef\second@{left}
   \ifx\first@\second@ \numbers@lefttrue \numbers@rightfalse \numbers@true \fi
   \xdef\second@{right}
   \ifx\first@\second@ \numbers@leftfalse \numbers@righttrue \numbers@true \fi
   \xdef\second@{leftright}
   \ifx\first@\second@ \numbers@lefttrue \numbers@righttrue \numbers@true \fi
   }
\def\hidenumbering{\numbers@false}
\def\hidenumber#1{\xdef\first@{#1,&,@} \hidenumber@}
\newcommand\shownames[2][n]{%
   \xdef\first@{#1}\ifx\first@\n@\else\xdef\names@fg{#1}\fi
   \xdef\first@{#2} \xdef\second@{left}
   \ifx\first@\second@ \names@rightfalse \else \names@righttrue \fi
   \names@true}
\def\hidenames{\names@false}
\def\hidename#1{\xdef\first@{#1,&,@} \hidename@}
\def\hideresidues{\hidechartrue}
\def\showresidues{\hidecharfalse}
\def\alignment#1{%
   \xdef\first@{#1}
   \xdef\temp@{left}
      \ifx\first@\temp@ \xdef\c@factor{0}
      \else
         \xdef\temp@{center}
         \ifx\first@\temp@ \xdef\c@factor{0.5}
         \else
            \xdef\temp@{right}
            \ifx\first@\temp@ \xdef\c@factor{1}
   \fi\fi\fi}
\def\donotshade#1{%
   \xdef\temp@{consensus}
   \xdef\first@{#1}
   \ifx\first@\temp@
      \consensuscolors{Black}{White}{Black}{White}{Black}{White}
   \else
      \xdef\first@{#1,&,@} \donot@shade
   \fi}
\def\hideseqs{\xdef\hide@seqs{y}}
\def\showseqs{\xdef\hide@seqs{n}}
\def\hideseq#1{\xdef\first@{#1,&,@} \hideseq@}
\def\killseq#1{\xdef\first@{#1,&,@} \killseq@}
\def\hidesequencelogo{\show@logofalse}
\def\hidesubfamilylogo{\show@sublogofalse}
\def\logo@group@get#1#2@{%
   \xdef\first@{#1}\xdef\third@{#2@}
   \ifx\first@\ampers@nd
   \else
      \ifnum`#1>96 \make@upper\fi
      \expandafter\xdef\csname logo@col\first@\endcsname{\second@}
      \expandafter\logo@group@get\third@
   \fi
}
\def\logocolor#1#2{%
   \xdef\logo@colors@set{yes}
   \xdef\second@{#2}
   \xdef\third@{#1&@}
   \expandafter\logo@group@get\third@
}
\newcommand\clearlogocolors[1][Black]{%
   \logocolor{ABCDEFGHIJKLMNOPQRSTUVWXYZ}{#1}
}
\def\dofrequencycorrection{\xdef\do@freq@correction{y}}
\def\undofrequencycorrection{\xdef\do@freq@correction{n}}
\newcommand\showlogoscale[2][Black]{\xdef\logo@scalecol{#1}\xdef\show@logoscale{#2}}
\def\hidelogoscale{\xdef\show@logoscale{n}}
\def\hidenegatives{\xdef\hide@negatives{y}\xdef\sublogo@tint{}}
\newcommand\shownegatives[1][medium]{%
   \xdef\hide@negatives{n}
   \xdef\first@{#1}
   \xdef\second@{full}
   \ifx\first@\second@
      \xdef\sublogo@tint{}
   \else
   \xdef\second@{strong}
   \ifx\first@\second@
      \xdef\sublogo@tint{Light}
   \else
   \xdef\second@{medium}
   \ifx\first@\second@
      \xdef\sublogo@tint{LightLight}
   \else
   \xdef\second@{weak}
   \ifx\first@\second@
      \xdef\sublogo@tint{LightLightLight}
   \else
      \xdef\sublogo@tint{LightLight}
   \fi\fi\fi\fi
}
\def\set@logocolors{%
   \xdef\second@{undefined}
   \ifx\first@\second@
      \if\logo@colors@set\n@
         \xdef\first@{standard}
      \fi
   \fi
   \xdef\second@{standard}
   \ifx\first@\second@
    \if\seq@type A
      \xdef\logo@colors@set{n}
    \else
      \if\seq@type P
         \xdef\first@{rasmol}
      \else
         \xdef\first@{nucleotide}
      \fi
    \fi
   \fi
   \xdef\second@{nucleotide}
   \ifx\first@\second@
      \clearlogocolors
      \logocolor{G}{Black}
      \logocolor{A}{Green}
      \logocolor{TU}{Red}
      \logocolor{C}{Blue}
   \else
   \xdef\second@{rasmol}
   \ifx\first@\second@
      \clearlogocolors
      \logocolor{DE}{Red}
      \logocolor{CM}{Yellow}
      \logocolor{KR}{Blue}
      \logocolor{ST}{Orange}
      \logocolor{FY}{MidnightBlue}
      \logocolor{NQ}{Cyan}
      \logocolor{G}{LightGray}
      \logocolor{LVI}{Green}
      \logocolor{A}{DarkGray}
      \logocolor{W}{CarnationPink}
      \logocolor{H}{CornflowerBlue}
      \logocolor{P}{Apricot}
      \logocolor{BZ}{LightMagenta}
   \else
   \xdef\second@{chemical}
   \ifx\first@\second@
      \clearlogocolors
      \logocolor{DE}{Red}
      \logocolor{VIL}{Black}
      \logocolor{AG}{Gray}
      \logocolor{NQ}{Green}
      \logocolor{FYW}{Brown}
      \logocolor{KRH}{Blue}
      \logocolor{ST}{Magenta}
      \logocolor{P}{Orange}
      \logocolor{CM}{Yellow}
   \else
   \xdef\second@{hydropathy}
   \ifx\first@\second@
      \clearlogocolors
      \logocolor{DE}{Red}
      \logocolor{KRH}{Blue}
      \logocolor{YSTGNQC}{Yellow}
      \logocolor{AFPMWVIL}{Green}
   \else
   \xdef\second@{structure}
   \ifx\first@\second@
      \clearlogocolors
      \logocolor{DEHKNQR}{Orange}
      \logocolor{ACGPSTWY}{Yellow}
      \logocolor{FILMV}{Green}
   \else
   \xdef\second@{standard area}
   \ifx\first@\second@
      \clearlogocolors
      \logocolor{G}{BrickRed}
      \logocolor{AS}{Orange}
      \logocolor{CP}{Yellow}
      \logocolor{TDVN}{YellowGreen}
      \logocolor{IE}{PineGreen}
      \logocolor{LQHM}{SkyBlue}
      \logocolor{FK}{RoyalPurple}
      \logocolor{Y}{RedViolet}
      \logocolor{RW}{Black}
   \else
   \xdef\second@{accessible area}
   \ifx\first@\second@
      \clearlogocolors
      \logocolor{C}{BrickRed}
      \logocolor{IVG}{Orange}
      \logocolor{FLMA}{Yellow}
      \logocolor{WSTH}{YellowGreen}
      \logocolor{P}{PineGreen}
      \logocolor{YDN}{SkyBlue}
      \logocolor{EQ}{RoyalPurple}
      \logocolor{R}{RedViolet}
      \logocolor{K}{Black}
   \else
   \xdef\second@{hardness}
   \ifx\first@\second@
      \clearlogocolors
      \logocolor{ADEGILPQSTV}{BlueRed5}
      \logocolor{KN}{BlueRed20}
      \logocolor{R}{BlueRed40}
      \logocolor{CFH}{BlueRed60}
      \logocolor{MY}{BlueRed80}
      \logocolor{W}{BlueRed100}
   \fi\fi\fi\fi\fi\fi\fi\fi
}
\newcommand\findsubfamily[2][n]{%
   \xdef\first@{#1}
   \ifx\first@\n@
   \else
      \xdef\subfamily@threshold{#1}
   \fi
   \xdef\subfamily@seq{#2}
}
\def\setsubfamily#1{%
   \xdef\sub@family@setting{#1}
   \loopcount=1
   \loop
      \expandafter\xdef\csname subfamily@num\the\loopcount\endcsname{1}
      \advance\loopcount by 1
   \ifnum\loopcount>\seq@count\else\repeat
   \xdef\subfamily@count{2}
   \res@count=0 
   \xdef\first@{#1,&,@}
   \setsubfamily@
   \clear@res@nums{1}
   \clear@res@nums{2}
   \expandafter\xdef\csname group@num2\endcsname{\the\res@count}
   \loopcount=\seq@count
   \advance\loopcount by -\res@count
   \expandafter\xdef\csname group@num1\endcsname{\the\loopcount}
}
\def\subfamilythreshold#1{\xdef\subfamily@threshold{#1}}
\newcommand\showsubfamilylogo[2][undefined]{%
   \xdef\first@{#1} \set@logocolors
   \xdef\first@{#2} \xdef\last@{top}
   \ifx\first@\last@\xdef\sublogo@top{0}\else\xdef\sublogo@top{1}\fi
   \show@sublogotrue
}
\def\relevance#1{%
   \def\sig@max{#1}
   \setlength\temp@@length{1000sp}
   \setlength\temp@@length{\sig@max\temp@@length}
   \loopcount=\temp@@length
   \xdef\sig@max{\the\loopcount}
}
\newcommand\showrelevance[2][Black]{\def\sig@color{#1}\def\sig@char{#2}\xdef\hide@sig{n}}
\def\hiderelevance{\xdef\hide@sig{y}}
\newcommand\showsequencelogo[2][undefined]{%
   \xdef\first@{#1} \set@logocolors
   \xdef\first@{#2} \xdef\last@{top}
   \ifx\first@\last@\xdef\logo@top{0}\else\xdef\logo@top{1}\fi
   \show@logotrue   
}
\newcommand\showconsensus[2][n]{%
   \xdef\text@scale{n}
   \xdef\box@scale{n}
   \xdef\first@{#1}
   \ifx\first@\n@
      \xdef\collect@cons@colors{no}
   \else
      \xdef\first@{#1,&,@}
      \expandafter\get@item\first@
      \xdef\c@nsc@l{\fourth@}
      \xdef\first@@{Gray}
      \ifx\c@nsc@l\first@@\xdef\text@scale{y}\else
      \xdef\first@@{RedBlue}
      \ifx\c@nsc@l\first@@\xdef\text@scale{y}\else
      \xdef\first@@{BlueRed}
      \ifx\c@nsc@l\first@@\xdef\text@scale{y}\else
      \xdef\first@@{RedGreen}
      \ifx\c@nsc@l\first@@\xdef\text@scale{y}\else
      \xdef\first@@{GreenRed}
      \ifx\c@nsc@l\first@@\xdef\text@scale{y}\else
      \xdef\first@@{ColdHot}
      \ifx\c@nsc@l\first@@\xdef\text@scale{y}\else
      \xdef\first@@{HotCold}
      \ifx\c@nsc@l\first@@\xdef\text@scale{y}
      \fi\fi\fi\fi\fi\fi\fi
      \expandafter\get@item\first@
      \ifx\fourth@\ampers@nd
         \xdef\c@nssc@le{White}
      \else
         \xdef\c@nssc@le{\fourth@}
         \xdef\first@{Gray}
         \ifx\c@nssc@le\first@\xdef\box@scale{y}\else
         \xdef\first@{RedBlue}
         \ifx\c@nssc@le\first@\xdef\box@scale{y}\else
         \xdef\first@{BlueRed}
         \ifx\c@nssc@le\first@\xdef\box@scale{y}\else
         \xdef\first@{RedGreen}
         \ifx\c@nssc@le\first@\xdef\box@scale{y}\else
         \xdef\first@{GreenRed}
         \ifx\c@nssc@le\first@\xdef\box@scale{y}\else
         \xdef\first@{ColdHot}
         \ifx\c@nssc@le\first@\xdef\box@scale{y}\else
         \xdef\first@{HotCold}
         \ifx\c@nssc@le\first@\xdef\box@scale{y}\else
         \fi\fi\fi\fi\fi\fi\fi
      \fi
      \xdef\collect@cons@colors{y}
   \fi
   \xdef\first@{#2} \xdef\last@{top}
   \ifx\first@\last@\xdef\cons@top{0}\else\xdef\cons@top{1}\fi
   \show@construe}
\def\consensuscolors#1#2#3#4#5#6{%
   \xdef\last@{\ampers@nd}
   \xdef\first@{#1&}\xdef\second@{#2&}
   \ifx\first@\last@\else\def\ConsTextNomatch{#1}\fi
   \ifx\second@\last@\else\def\ConsNomatch{#2}\fi
   \xdef\first@{#3&}\xdef\second@{#4&}
   \ifx\first@\last@\else\def\ConsTextMatch{#3}\fi
   \ifx\second@\last@\else\def\ConsMatch{#4}\fi
   \xdef\first@{#5&}\xdef\second@{#6&}
   \ifx\first@\last@\else\def\ConsTextAllmatch{#5}\fi
   \ifx\second@\last@\else\def\ConsAllmatch{#6}\fi
}
\def\defconsensus#1#2#3{%
   \xdef\second@{#1&}
   \ifx\second@\ampers@nd   \else \def\n@m@tch{#1}\fi
   \xdef\second@{#2&}
   \ifx\second@\ampers@nd   \else \def\m@tch{#2}\fi
   \xdef\second@{#3&}
   \ifx\second@\ampers@nd   \else \def\@llm@tch{#3}\fi}
\def\hideconsensus{\show@consfalse}
\def\nameconsensus#1{\def\cons@name{#1}}
\def\namesequencelogo#1{\def\logo@name@user{#1}}
\newcommand\namesubfamilylogo[2][]{\def\sublogo@name@neg{#1}\def\sublogo@name@user{#2}}
\def\hideleadinggaps{\sh@wg@psfalse}
\def\showleadinggaps{\sh@wg@pstrue}
\newcommand\showruler[3][n]{%
   \xdef\first@{#1}
   \ifx\first@\n@\else\xdef\ruler@fg{#1}\fi
   \xdef\first@{consensus} \xdef\second@{#3}
   \xdef\third@{bottom}    \xdef\fourth@{#2}
   \ifx\third@\fourth@ \xdef\rule@top{1}\else\xdef\rule@top{0}\fi
   \ifx\first@\second@ \xdef\rule@num{0}
   \else
     \xdef\first@{#3 @} \expandafter\check@letter\first@
     \xdef\first@{#3}   \ifletter \get@name@number \fi
     \ifnum\first@>\seq@count
     \else
        \ifnum\first@>0
           \xdef\rule@num{\first@}
        \fi
      \fi
   \fi
   \xdef\ruler@{}}
\def\hideruler{\xdef\rule@num{-1}}
\def\allowzero{\xdef\allow@zero{y}}
\def\disallowzero{\xdef\allow@zero{n}}
\def\rulersteps#1{%
   \xdef\ruler@step{#1}
   \ifnum#1<4 \xdef\ruler@rot{90}\fi
}
\def\rotateruler{\xdef\ruler@rot{90}}
\def\unrotateruler{\xdef\ruler@rot{0}}
\def\namerulerpos#1#2{%
   \expandafter\xdef\csname alt@ruler#1\endcsname{#2}
}
\def\featurerule#1{\setlength\rule@thick{#1}}
\def\orderseqs#1{%
   \def\order@loop{%
      \expandafter\check@letter\first@
      \ifletter
         \expandafter\get@item\first@
         \xdef\first@{\fourth@}
         \get@name@number
         \xdef\seq@order{\seq@order,\first@}
         \xdef\first@{\first@@ @}
         \order@loop
      \else
         \expandafter\get@digit\first@
         \ifx\fourth@\ampers@nd
         \else
            \xdef\seq@order{\seq@order,\fourth@}
            \order@loop
         \fi
      \fi}
   \xdef\first@{#1,&,@}
   \xdef\seq@order{}
   \order@loop
   \xdef\seq@order{\seq@order @}
   \expandafter\get@item\seq@order
   \xdef\seq@order{\first@@,@}
}
\def\setfamily#1#2{%
   \xdef\second@{#2&}
   \ifx\second@\ampers@nd
   \else
   \xdef\first@{#1}
   \xdef\second@{#2}
   \xdef\temp@{rm}
   \ifx\second@\temp@
      \xdef\third@{\rmdefault}
   \else
      \xdef\temp@{sf}
      \ifx\second@\temp@
         \xdef\third@{\sfdefault}
      \else
         \xdef\temp@{tt}
         \ifx\second@\temp@
            \xdef\third@{\ttdefault}
         \else
            \xdef\third@{\second@}
   \fi\fi\fi
   \xdef\temp@{featurenames}
   \ifx\first@\temp@ \xdef\ftext@family{\third@}
   \else
   \xdef\temp@{featurestylenames}
   \ifx\first@\temp@ \xdef\fstyles@family{\third@}
   \else
   \xdef\temp@{features}
   \ifx\first@\temp@ \xdef\featuretext@family{\third@}
   \else
   \xdef\temp@{featurestyles}
   \ifx\first@\temp@ \xdef\featurestyles@family{\third@}
   \else
   \xdef\temp@{numbering}
   \ifx\first@\temp@ \xdef\numbertext@family{\third@}
   \else
   \xdef\temp@{names}
   \ifx\first@\temp@ \xdef\namestext@family{\third@}
   \else
   \xdef\temp@{residues}
   \ifx\first@\temp@ \xdef\residues@family{\third@}
   \else
   \xdef\temp@{legend}
   \ifx\first@\temp@ \xdef\legend@family{\third@}
   \else
   \xdef\temp@{labels}
   \ifx\first@\temp@ \xdef\label@family{\third@}
   \else
   \xdef\temp@{ruler}
   \ifx\first@\temp@ \xdef\ruler@family{\second@} 
   \else
   \xdef\temp@{all}
   \ifx\first@\temp@
      \xdef\ftext@family{\third@}
      \xdef\fstyles@family{\third@}
      \xdef\featuretext@family{\third@}
      \xdef\featurestyles@family{\third@}
      \xdef\numbertext@family{\third@}
      \xdef\namestext@family{\third@}
      \xdef\residues@family{\third@}
      \xdef\legend@family{\third@}
      \xdef\label@family{\third@}
      \xdef\ruler@family{\second@}
   \fi\fi\fi\fi\fi\fi\fi\fi\fi\fi\fi
   \fi
}
\def\setseries#1#2{%
   \xdef\second@{#2&}
   \ifx\second@\ampers@nd
   \else
   \xdef\first@{#1}
   \xdef\second@{#2}
   \xdef\temp@{bf}
   \ifx\second@\temp@
      \xdef\third@{\bfdefault}
   \else
      \xdef\temp@{md}
      \ifx\second@\temp@
         \xdef\third@{\mddefault}
      \else
         \xdef\third@{\second@}
   \fi\fi
   \xdef\temp@{featurenames}
   \ifx\first@\temp@ \xdef\ftext@series{\third@}
   \else
   \xdef\temp@{featurestylenames}
   \ifx\first@\temp@ \xdef\fstyles@series{\third@}
   \else
   \xdef\temp@{features}
   \ifx\first@\temp@ \xdef\featuretext@series{\third@}
   \else
   \xdef\temp@{featurestyles}
   \ifx\first@\temp@ \xdef\featurestyles@series{\third@}
   \else
   \xdef\temp@{numbering}
   \ifx\first@\temp@ \xdef\numbertext@series{\third@}
   \else
   \xdef\temp@{names}
   \ifx\first@\temp@ \xdef\namestext@series{\third@}
   \else
   \xdef\temp@{residues}
   \ifx\first@\temp@ \xdef\residues@series{\third@}
   \else
   \xdef\temp@{legend}
   \ifx\first@\temp@ \xdef\legend@series{\third@}
   \else
   \xdef\temp@{labels}
   \ifx\first@\temp@ \xdef\label@series{\third@}
   \else
   \xdef\temp@{all}
   \ifx\first@\temp@
      \xdef\ftext@series{\third@}
      \xdef\fstyles@series{\third@}
      \xdef\featuretext@series{\third@}
      \xdef\featurestyles@series{\third@}
      \xdef\numbertext@series{\third@}
      \xdef\namestext@series{\third@}
      \xdef\residues@series{\third@}
      \xdef\legend@series{\third@}
      \xdef\label@series{\third@}
   \fi\fi\fi\fi\fi\fi\fi\fi\fi\fi
   \fi
}
\def\setshape#1#2{%
   \xdef\second@{#2&}
   \ifx\second@\ampers@nd
   \else
   \xdef\first@{#1}
   \xdef\second@{#2}
   \xdef\temp@{it}
   \ifx\second@\temp@
      \xdef\third@{\itdefault}
   \else
      \xdef\temp@{sl}
      \ifx\second@\temp@
         \xdef\third@{\sldefault}
      \else
         \xdef\temp@{sc}
         \ifx\second@\temp@
            \xdef\third@{\scdefault}
         \else
            \xdef\temp@{up}
            \ifx\second@\temp@
               \xdef\third@{\updefault}
            \else
               \xdef\third@{\second@}
   \fi\fi\fi\fi
   \xdef\temp@{featurenames}
   \ifx\first@\temp@ \xdef\ftext@shape{\third@}
   \else
   \xdef\temp@{featurestylenames}
   \ifx\first@\temp@ \xdef\fstyles@shape{\third@}
   \else
   \xdef\temp@{features}
   \ifx\first@\temp@ \xdef\featuretext@shape{\third@}
   \else
   \xdef\temp@{featurestyles}
   \ifx\first@\temp@ \xdef\featurestyles@shape{\third@}
   \else
   \xdef\temp@{numbering}
   \ifx\first@\temp@ \xdef\numbertext@shape{\third@}
   \else
   \xdef\temp@{names}
   \ifx\first@\temp@ \xdef\namestext@shape{\third@}
   \else
   \xdef\temp@{residues}
   \ifx\first@\temp@ \xdef\residues@shape{\third@}
   \else
   \xdef\temp@{legend}
   \ifx\first@\temp@ \xdef\legend@shape{\third@}
   \else
   \xdef\temp@{labels}
   \ifx\first@\temp@ \xdef\label@shape{\third@}
   \else
   \xdef\temp@{all}
   \ifx\first@\temp@
      \xdef\ftext@shape{\third@}
      \xdef\fstyles@shape{\third@}
      \xdef\featuretext@shape{\third@}
      \xdef\featurestyles@shape{\third@}
      \xdef\numbertext@shape{\third@}
      \xdef\namestext@shape{\third@}
      \xdef\residues@shape{\third@}
      \xdef\legend@shape{\third@}
      \xdef\label@shape{\third@}
   \fi\fi\fi\fi\fi\fi\fi\fi\fi\fi
   \fi
}
\def\setsize#1#2{%
   \xdef\second@{#2&}
   \ifx\second@\ampers@nd
   \else
   \xdef\first@{#1}
   \xdef\temp@{features}
   \ifx\first@\temp@
      \def\featuretext@size{\csname #2\endcsname}
   \else
   \xdef\temp@{featurestyles}
   \ifx\first@\temp@
      \def\featurestyles@size{\csname #2\endcsname}
   \else
   \xdef\temp@{featurenames}
   \ifx\first@\temp@
      \def\ftext@size{\csname #2\endcsname}
   \else
   \xdef\temp@{featurestylenames}
   \ifx\first@\temp@
      \def\fstyles@size{\csname #2\endcsname}
   \else
   \xdef\temp@{numbering}
   \ifx\first@\temp@
      \def\numbertext@size{\csname #2\endcsname}
   \else
   \xdef\temp@{names}
   \ifx\first@\temp@
      \def\namestext@size{\csname #2\endcsname}
   \else
   \xdef\temp@{legend}
   \ifx\first@\temp@
      \def\legend@size{\csname #2\endcsname}
   \else
   \xdef\temp@{labels}
   \ifx\first@\temp@
      \def\label@size{\csname #2\endcsname}
   \else
   \xdef\temp@{residues}
   \ifx\first@\temp@
      \def\residues@size{\csname #2\endcsname}
      \xdef\res@size{#2}
   \else
   \xdef\temp@{all}
   \ifx\first@\temp@
      \def\featuretext@size{\csname #2\endcsname}
      \def\featurestyles@size{\csname #2\endcsname}
      \def\ftext@size{\csname #2\endcsname}
      \def\fstyles@size{\csname #2\endcsname}
      \def\numbertext@size{\csname #2\endcsname}
      \def\namestext@size{\csname #2\endcsname}
      \def\legend@size{\csname #2\endcsname}
      \def\label@size{\csname #2\endcsname}
      \def\residues@size{\csname #2\endcsname}
      \xdef\res@size{#2}
   \fi\fi\fi\fi\fi\fi\fi\fi\fi\fi
   \xdef\temp@{Huge}
   \ifx\temp@\res@size
      \def\bottomruler@size{\csname Large\endcsname}
   \else
   \xdef\temp@{huge}
   \ifx\temp@\res@size
      \def\bottomruler@size{\csname large\endcsname}
   \else
   \xdef\temp@{LARGE}
   \ifx\temp@\res@size
      \def\bottomruler@size{\csname normalsize\endcsname}
   \else
   \xdef\temp@{Large}
   \ifx\temp@\res@size
      \def\bottomruler@size{\csname small\endcsname}
   \else
   \xdef\temp@{large}
   \ifx\temp@\res@size
      \def\bottomruler@size{\csname footnotesize\endcsname}
   \else
   \xdef\temp@{normalsize}
   \ifx\temp@\res@size
      \def\bottomruler@size{\csname scriptsize\endcsname}
   \else
      \def\bottomruler@size{\csname tiny\endcsname}
   \fi\fi\fi\fi\fi\fi
   \fi
   \xdef\temp@{ruler}
   \ifx\first@\temp@
      \def\bottomruler@size{\csname #2\endcsname}
   \fi
}
\def\setfont#1#2#3#4#5{%
   \setfamily{#1}{#2}\setseries{#1}{#3}
   \setshape{#1}{#4}\setsize{#1}{#5}}
\def\featurenamesrm{\setfamily{featurenames}{rm}}
\def\featurenamessf{\setfamily{featurenames}{sf}}
\def\featurenamestt{\setfamily{featurenames}{tt}}
\def\featurenamesmd{\setseries{featurenames}{md}}
\def\featurenamesbf{\setseries{featurenames}{bf}}
\def\featurenamesup{\setshape {featurenames}{up}}
\def\featurenamesit{\setshape {featurenames}{it}}
\def\featurenamessl{\setshape {featurenames}{sl}}
\def\featurenamessc{\setshape {featurenames}{sc}}
\def\featurenamestiny        {\setsize{featurenames}{tiny}}
\def\featurenamesscriptsize  {\setsize{featurenames}{scriptsize}}
\def\featurenamesfootnotesize{\setsize{featurenames}{footnotesize}}
\def\featurenamessmall       {\setsize{featurenames}{small}}
\def\featurenamesnormalsize  {\setsize{featurenames}{normalsize}}
\def\featurenameslarge       {\setsize{featurenames}{large}}
\def\featurenamesLarge       {\setsize{featurenames}{Large}}
\def\featurenamesLARGE       {\setsize{featurenames}{LARGE}}
\def\featurenameshuge        {\setsize{featurenames}{huge}}
\def\featurenamesHuge        {\setsize{featurenames}{Huge}}
\def\featurestylenamesrm{\setfamily{featurestylenames}{rm}}
\def\featurestylenamessf{\setfamily{featurestylenames}{sf}}
\def\featurestylenamestt{\setfamily{featurestylenames}{tt}}
\def\featurestylenamesmd{\setseries{featurestylenames}{md}}
\def\featurestylenamesbf{\setseries{featurestylenames}{bf}}
\def\featurestylenamesup{\setshape {featurestylenames}{up}}
\def\featurestylenamesit{\setshape {featurestylenames}{it}}
\def\featurestylenamessl{\setshape {featurestylenames}{sl}}
\def\featurestylenamessc{\setshape {featurestylenames}{sc}}
\def\featurestylenamestiny        {\setsize{featurestylenames}{tiny}}
\def\featurestylenamesscriptsize  {\setsize{featurestylenames}{scriptsize}}
\def\featurestylenamesfootnotesize{\setsize{featurestylenames}{footnotesize}}
\def\featurestylenamessmall       {\setsize{featurestylenames}{small}}
\def\featurestylenamesnormalsize  {\setsize{featurestylenames}{normalsize}}
\def\featurestylenameslarge       {\setsize{featurestylenames}{large}}
\def\featurestylenamesLarge       {\setsize{featurestylenames}{Large}}
\def\featurestylenamesLARGE       {\setsize{featurestylenames}{LARGE}}
\def\featurestylenameshuge        {\setsize{featurestylenames}{huge}}
\def\featurestylenamesHuge        {\setsize{featurestylenames}{Huge}}
\def\featuresrm{\setfamily{features}{rm}}
\def\featuressf{\setfamily{features}{sf}}
\def\featurestt{\setfamily{features}{tt}}
\def\featuresmd{\setseries{features}{md}}
\def\featuresbf{\setseries{features}{bf}}
\def\featuresup{\setshape {features}{up}}
\def\featuresit{\setshape {features}{it}}
\def\featuressl{\setshape {features}{sl}}
\def\featuressc{\setshape {features}{sc}}
\def\featurestiny        {\setsize{features}{tiny}}
\def\featuresscriptsize  {\setsize{features}{scriptsize}}
\def\featuresfootnotesize{\setsize{features}{footnotesize}}
\def\featuressmall       {\setsize{features}{small}}
\def\featuresnormalsize  {\setsize{features}{normalsize}}
\def\featureslarge       {\setsize{features}{large}}
\def\featuresLarge       {\setsize{features}{Large}}
\def\featuresLARGE       {\setsize{features}{LARGE}}
\def\featureshuge        {\setsize{features}{huge}}
\def\featuresHuge        {\setsize{features}{Huge}}
\def\featurestylesrm{\setfamily{featurestyles}{rm}}
\def\featurestylessf{\setfamily{featurestyles}{sf}}
\def\featurestylestt{\setfamily{featurestyles}{tt}}
\def\featurestylesmd{\setseries{featurestyles}{md}}
\def\featurestylesbf{\setseries{featurestyles}{bf}}
\def\featurestylesup{\setshape {featurestyles}{up}}
\def\featurestylesit{\setshape {featurestyles}{it}}
\def\featurestylessl{\setshape {featurestyles}{sl}}
\def\featurestylessc{\setshape {featurestyles}{sc}}
\def\featurestylestiny        {\setsize{featurestyles}{tiny}}
\def\featurestylesscriptsize  {\setsize{featurestyles}{scriptsize}}
\def\featurestylesfootnotesize{\setsize{featurestyles}{footnotesize}}
\def\featurestylessmall       {\setsize{featurestyles}{small}}
\def\featurestylesnormalsize  {\setsize{featurestyles}{normalsize}}
\def\featurestyleslarge       {\setsize{featurestyles}{large}}
\def\featurestylesLarge       {\setsize{featurestyles}{Large}}
\def\featurestylesLARGE       {\setsize{featurestyles}{LARGE}}
\def\featurestyleshuge        {\setsize{featurestyles}{huge}}
\def\featurestylesHuge        {\setsize{featurestyles}{Huge}}
\def\numberingrm{\setfamily{numbering}{rm}}
\def\numberingsf{\setfamily{numbering}{sf}}
\def\numberingtt{\setfamily{numbering}{tt}}
\def\numberingmd{\setseries{numbering}{md}}
\def\numberingbf{\setseries{numbering}{bf}}
\def\numberingup{\setshape {numbering}{up}}
\def\numberingit{\setshape {numbering}{it}}
\def\numberingsl{\setshape {numbering}{sl}}
\def\numberingsc{\setshape {numbering}{sc}}
\def\numberingtiny        {\setsize{numbering}{tiny}}
\def\numberingscriptsize  {\setsize{numbering}{scriptsize}}
\def\numberingfootnotesize{\setsize{numbering}{footnotesize}}
\def\numberingsmall       {\setsize{numbering}{small}}
\def\numberingnormalsize  {\setsize{numbering}{normalsize}}
\def\numberinglarge       {\setsize{numbering}{large}}
\def\numberingLarge       {\setsize{numbering}{Large}}
\def\numberingLARGE       {\setsize{numbering}{LARGE}}
\def\numberinghuge        {\setsize{numbering}{huge}}
\def\numberingHuge        {\setsize{numbering}{Huge}}
\def\namesrm{\setfamily{names}{rm}}
\def\namessf{\setfamily{names}{sf}}
\def\namestt{\setfamily{names}{tt}}
\def\namesmd{\setseries{names}{md}}
\def\namesbf{\setseries{names}{bf}}
\def\namesup{\setshape {names}{up}}
\def\namesit{\setshape {names}{it}}
\def\namessl{\setshape {names}{sl}}
\def\namessc{\setshape {names}{sc}}
\def\namestiny        {\setsize{names}{tiny}}
\def\namesscriptsize  {\setsize{names}{scriptsize}}
\def\namesfootnotesize{\setsize{names}{footnotesize}}
\def\namessmall       {\setsize{names}{small}}
\def\namesnormalsize  {\setsize{names}{normalsize}}
\def\nameslarge       {\setsize{names}{large}}
\def\namesLarge       {\setsize{names}{Large}}
\def\namesLARGE       {\setsize{names}{LARGE}}
\def\nameshuge        {\setsize{names}{huge}}
\def\namesHuge        {\setsize{names}{Huge}}
\def\residuesrm{\setfamily{residues}{rm}}
\def\residuessf{\setfamily{residues}{sf}}
\def\residuestt{\setfamily{residues}{tt}}
\def\residuesmd{\setseries{residues}{md}}
\def\residuesbf{\setseries{residues}{bf}}
\def\residuesup{\setshape {residues}{up}}
\def\residuesit{\setshape {residues}{it}}
\def\residuessl{\setshape {residues}{sl}}
\def\residuessc{\setshape {residues}{sc}}
\def\residuestiny        {\setsize{residues}{tiny}}
\def\residuesscriptsize  {\setsize{residues}{scriptsize}}
\def\residuesfootnotesize{\setsize{residues}{footnotesize}}
\def\residuessmall       {\setsize{residues}{small}}
\def\residuesnormalsize  {\setsize{residues}{normalsize}}
\def\residueslarge       {\setsize{residues}{large}}
\def\residuesLarge       {\setsize{residues}{Large}}
\def\residuesLARGE       {\setsize{residues}{LARGE}}
\def\residueshuge        {\setsize{residues}{huge}}
\def\residuesHuge        {\setsize{residues}{Huge}}
\def\legendrm{\setfamily{legend}{rm}}
\def\legendsf{\setfamily{legend}{sf}}
\def\legendtt{\setfamily{legend}{tt}}
\def\legendmd{\setseries{legend}{md}}
\def\legendbf{\setseries{legend}{bf}}
\def\legendup{\setshape {legend}{up}}
\def\legendit{\setshape {legend}{it}}
\def\legendsl{\setshape {legend}{sl}}
\def\legendsc{\setshape {legend}{sc}}
\def\legendtiny        {\setsize{legend}{tiny}}
\def\legendscriptsize  {\setsize{legend}{scriptsize}}
\def\legendfootnotesize{\setsize{legend}{footnotesize}}
\def\legendsmall       {\setsize{legend}{small}}
\def\legendnormalsize  {\setsize{legend}{normalsize}}
\def\legendlarge       {\setsize{legend}{large}}
\def\legendLarge       {\setsize{legend}{Large}}
\def\legendLARGE       {\setsize{legend}{LARGE}}
\def\legendhuge        {\setsize{legend}{huge}}
\def\legendHuge        {\setsize{legend}{Huge}}
\def\rulerrm{\setfamily{ruler}{rm}}
\def\rulersf{\setfamily{ruler}{sf}}
\def\rulertt{\setfamily{ruler}{tt}}
\def\rulertiny        {\setsize{ruler}{tiny}}
\def\rulerscriptsize  {\setsize{ruler}{scriptsize}}
\def\rulerfootnotesize{\setsize{ruler}{footnotesize}}
\def\rulersmall       {\setsize{ruler}{small}}
\def\rulernormalsize  {\setsize{ruler}{normalsize}}
\def\rulerlarge       {\setsize{ruler}{large}}
\def\rulerLarge       {\setsize{ruler}{Large}}
\def\rulerLARGE       {\setsize{ruler}{LARGE}}
\def\rulerhuge        {\setsize{ruler}{huge}}
\def\rulerHuge        {\setsize{ruler}{Huge}}
\def\funcshadingstyle#1#2#3#4#5{%
  \xdef\temp@{nomatch} \xdef\first@{#1}
  \ifx\temp@\first@
     \xdef\first@{0}
  \else
     \xdef\first@{\csname funcgrp#1\endcsname}
  \fi
  \ifnum\first@>-1
     \expandafter\xdef\csname fg@textcolor\first@\endcsname{#2}
     \expandafter\xdef\csname fg@color\first@\endcsname{#3}
     \expandafter\xdef\csname funcm@tch\first@\endcsname{#4}
     \expandafter\def\csname func@style\first@\endcsname{\csname text#5\endcsname}
  \fi}
\def\defshadingcolors#1{%  
   \expandafter\xdef\csname TextNomatch@#1\endcsname{\TextNomatch}
   \expandafter\xdef\csname Nomatch@#1\endcsname{\Nomatch}
   \expandafter\xdef\csname resn@m@tch@#1\endcsname{\resn@m@tch}
      
   \expandafter\xdef\csname TextSimilar@#1\endcsname{\TextSimilar}
   \expandafter\xdef\csname Similar@#1\endcsname{\Similar}
   \expandafter\xdef\csname ressimm@tch@#1\endcsname{\ressimm@tch}
   
   \expandafter\xdef\csname TextIdentical@#1\endcsname{\TextIdentical}
   \expandafter\xdef\csname Identical@#1\endcsname{\Identical}
   \expandafter\xdef\csname resm@tch@#1\endcsname{\resm@tch}
   
   \expandafter\xdef\csname TextAllmatch@#1\endcsname{\TextAllmatch}
   \expandafter\xdef\csname Allmatch@#1\endcsname{\Allmatch}
   \expandafter\xdef\csname res@llm@tch@#1\endcsname{\res@llm@tch}
   
   \expandafter\xdef\csname gap@fg@#1\endcsname{\gap@fg}
   \expandafter\xdef\csname gap@bg@#1\endcsname{\gap@bg}
   
}
\def\shadingcolors#1{%
  \gapcolors{Black}{White}
  \nomatchresidues{Black}{White}{upper}{up}
  \xdef\first@{#1} \xdef\second@{blues}
  \ifx\first@\second@
    \similarresidues{Black}{Magenta}{upper}{up}
    \conservedresidues{White}{RoyalBlue}{upper}{up}
    \allmatchresidues{Goldenrod}{RoyalPurple}{upper}{up}
  \else \xdef\second@{greens}
  \ifx\first@\second@
    \similarresidues{Black}{GreenYellow}{upper}{up}
    \conservedresidues{White}{PineGreen}{upper}{up}
    \allmatchresidues{YellowOrange}{OliveGreen}{upper}{up}
  \else \xdef\second@{reds}
  \ifx\first@\second@
    \similarresidues{Black}{YellowOrange}{upper}{up}
    \conservedresidues{White}{BrickRed}{upper}{up}
    \allmatchresidues{YellowGreen}{Mahagony}{upper}{up}
  \else \xdef\second@{black}
  \ifx\first@\second@
    \similarresidues{Black}{White}{upper}{sl}
    \conservedresidues{White}{Black}{upper}{up}
    \allmatchresidues{White}{Black}{upper}{sl}
  \else \xdef\second@{grays}
  \ifx\first@\second@
    \similarresidues{Black}{LightGray}{upper}{up}
    \conservedresidues{White}{DarkGray}{upper}{up}
    \allmatchresidues{White}{Black}{upper}{up}
  \else 
    \xdef\TextNomatch{\csname TextNomatch@#1\endcsname}
    \xdef\Nomatch    {\csname Nomatch@#1\endcsname}
    \xdef\resn@m@tch {\csname resn@m@tch@#1\endcsname}
    \xdef\TextSimilar{\csname TextSimilar@#1\endcsname}
    \xdef\Similar    {\csname Similar@#1\endcsname}
    \xdef\ressimm@tch{\csname ressimm@tch@#1\endcsname}
    \xdef\TextIdentical{\csname TextIdentical@#1\endcsname}
    \xdef\Identical    {\csname Identical@#1\endcsname}
    \xdef\resm@tch     {\csname resm@tch@#1\endcsname}
    \xdef\TextAllmatch{\csname TextAllmatch@#1\endcsname}
    \xdef\Allmatch    {\csname Allmatch@#1\endcsname}
    \xdef\res@llm@tch {\csname res@llm@tch@#1\endcsname}
    \xdef\gap@fg    {\csname gap@fg@#1\endcsname}
    \xdef\gap@bg    {\csname gap@bg@#1\endcsname}
    \xdef\domgap@fg {\csname domgap@fg@#1\endcsname}
    \xdef\domgap@bg {\csname domgap@bg@#1\endcsname}
  \fi\fi\fi\fi\fi
  \xdef\first@{#1}
  \expandafter\defshadingcolors{\first@}
}
\def\nomatchresidues#1#2#3#4  {\xdef\first@{#1&}\xdef\second@{#2&}\xdef\third@{#3&}
                             \xdef\last@{\ampers@nd}
                             \ifx\first@\last@\else\def\TextNomatch{#1}
                               \expandafter\def\csname fg@textcolor0\endcsname{#1}
                             \fi
                             \ifx\second@\last@\else\gdef\Nomatch{#2}
                               \expandafter\def\csname fg@color0\endcsname{#2}
                             \fi
                             \ifx\third@\last@\else\def\resn@m@tch{#3}
                             \fi
                             \xdef\first@{#4&}
                             \ifx\first@\last@\else
                               \def\no@style{\csname text#4\endcsname}
                               \expandafter\def\csname func@style0\endcsname%
                                  {\csname text#4\endcsname}\fi}
\def\similarresidues#1#2#3#4  {\xdef\first@{#1&}\xdef\second@{#2&}\xdef\third@{#3&}
                             \xdef\last@{\ampers@nd}
                             \ifx\first@\last@\else\def\TextSimilar{#1}\fi
                             \ifx\second@\last@\else\gdef\Similar{#2}\fi
                             \ifx\third@\last@\else\def\ressimm@tch{#3}\fi
                             \xdef\first@{#4&}
                             \ifx\first@\last@\else
                               \def\sim@style{\csname text#4\endcsname}\fi}
\def\conservedresidues#1#2#3#4{\xdef\first@{#1&}\xdef\second@{#2&}\xdef\third@{#3&}
                             \xdef\last@{\ampers@nd}
                             \ifx\first@\last@\else\def\TextIdentical{#1}\fi
                             \ifx\second@\last@\else\gdef\Identical{#2}\fi
                             \ifx\third@\last@\else\def\resm@tch{#3}\fi
                             \xdef\first@{#4&}
                             \ifx\first@\last@\else
                               \def\id@style{\csname text#4\endcsname}\fi}
\def\allmatchresidues#1#2#3#4 {\xdef\first@{#1&}\xdef\second@{#2&}\xdef\third@{#3&}
                             \xdef\last@{\ampers@nd}
                             \ifx\first@\last@\else\def\TextAllmatch{#1}\fi
                             \ifx\second@\last@\else\gdef\Allmatch{#2}\fi
                             \ifx\third@\last@\else\def\res@llm@tch{#3}\fi
                             \xdef\first@{#4&}
                             \ifx\first@\last@\else
                               \def\all@style{\csname text#4\endcsname}\fi}
\def\gapcolors#1#2           {\xdef\first@{#1&}\xdef\second@{#2&}
                             \xdef\last@{\ampers@nd}
                             \ifx\first@\last@\else\def\gap@fg{#1}
                                \expandafter\def\csname fg@textcolor*\endcsname{#1}\fi
                             \ifx\second@\last@\else\def\gap@bg{#2}
                                \expandafter\def\csname fg@color*\endcsname{#2}\fi}
\def\domaingapcolors#1#2     {\xdef\first@{#1&}\xdef\second@{#2&}
                             \xdef\last@{\ampers@nd}
                             \ifx\first@\last@\else\def\domgap@fg{#1}
                                \expandafter\def\csname fg@textcolor!\endcsname{#1}\fi
                             \ifx\second@\last@\else\def\domgap@bg{#2}
                                \expandafter\def\csname fg@color!\endcsname{#2}\fi}
\def\shadebox#1{%
   \xdef\first@{White}%
   \xdef\third@{#1}%
   \xdef\second@{nomatch}%
   \ifx\second@\third@
      \ifx\Nomatch\first@\white@box\else\textcolor{\Nomatch}{\box@rule}\fi%
   \else
   \xdef\second@{similar}%
   \ifx\second@\third@
      \ifx\Similar\first@\white@box\else\textcolor{\Similar}{\box@rule}\fi%
   \else
   \xdef\second@{conserved}%
   \ifx\second@\third@
      \ifx\Identical\first@\white@box\else\textcolor{\Identical}{\box@rule}\fi%
   \else
   \xdef\second@{allmatch}%
   \ifx\second@\third@
      \ifx\Allmatch\first@\white@box\else\textcolor{\Allmatch}{\box@rule}\fi%
   \else
   \ifx\third@\first@\white@box\else\textcolor{\third@}{\box@rule}\fi
   \fi\fi\fi\fi}
\def\featurenamescolor#1{%
   \expandafter\xdef\csname ftext@fg@ttttop\endcsname{#1}
   \expandafter\xdef\csname ftext@fg@tttop\endcsname{#1}
   \expandafter\xdef\csname ftext@fg@ttop\endcsname{#1}
   \expandafter\xdef\csname ftext@fg@top\endcsname{#1}
   \expandafter\xdef\csname ftext@fg@bottom\endcsname{#1}
   \expandafter\xdef\csname ftext@fg@bbottom\endcsname{#1}
   \expandafter\xdef\csname ftext@fg@bbbottom\endcsname{#1}
   \expandafter\xdef\csname ftext@fg@bbbbottom\endcsname{#1}
}
\def\featurestylenamescolor#1{%
   \expandafter\xdef\csname fstyles@fg@ttttop\endcsname{#1}
   \expandafter\xdef\csname fstyles@fg@tttop\endcsname{#1}
   \expandafter\xdef\csname fstyles@fg@ttop\endcsname{#1}
   \expandafter\xdef\csname fstyles@fg@top\endcsname{#1}
   \expandafter\xdef\csname fstyles@fg@bottom\endcsname{#1}
   \expandafter\xdef\csname fstyles@fg@bbottom\endcsname{#1}
   \expandafter\xdef\csname fstyles@fg@bbbottom\endcsname{#1}
   \expandafter\xdef\csname fstyles@fg@bbbbottom\endcsname{#1}
}
\def\featurenamecolor#1#2{\expandafter\xdef\csname ftext@fg@#1\endcsname{#2}}
\def\featurestylenamecolor#1#2{\expandafter\xdef\csname fstyles@fg@#1\endcsname{#2}}
\def\namescolor#1{\xdef\names@fg{#1}}
\def\namecolor#1#2{%
   \xdef\first@{consensus} \xdef\second@{#1}
   \ifx\first@\second@
      \expandafter\xdef\csname name@col0\endcsname{#2}
   \else
   \xdef\first@{featurenames} \xdef\second@{#1}
   \ifx\first@\second@
     \xdef\ftext@fg{#2}
   \else
   \xdef\first@{featurestylenames} \xdef\second@{#1}
   \ifx\first@\second@
     \xdef\fstyles@fg{#2}
   \else
      \xdef\first@{#1,&,@} \xdef\third@{#2} \namecolor@
   \fi\fi\fi
}
\def\numberingcolor#1{\xdef\numbering@fg{#1}}
\def\numbercolor#1#2{%
   \xdef\first@{consensus} \xdef\second@{#1}
   \ifx\first@\second@
      \expandafter\xdef\csname number@col0\endcsname{#2}
   \else
      \xdef\first@{#1,&,@} \xdef\third@{#2} \numbercolor@
   \fi
}
\def\legendcolor#1{\xdef\legend@fg{#1}}
\def\rulercolor#1{\xdef\ruler@fg{#1}}
\def\molweight#1#2{%
   \xdef\temp@{Da}%
   \xdef\second@{#2}%
   \ifx\second@\temp@\xdef\third@{Da}\else\xdef\third@{kDa}\fi%
   \xdef\first@{#1 @} \expandafter\check@letter\first@
   \xdef\first@{#1}   \ifletter \get@name@number \fi
   \xdef\first@{\csname @rd\first@\endcsname}%
   \loopcount=\csname mol@weight\first@\endcsname%
   \divide\loopcount by 10\relax%
   \innerloopcount=\loopcount%
   \hbox{%
   \ifnum\loopcount>1000%
      \divide\loopcount by 1000\relax%
      \pos@count=\loopcount%
      \multiply\loopcount by 1000\relax%
      \advance\innerloopcount by -\loopcount%
      \loopcount=\innerloopcount%
      \ifnum\loopcount>949\advance\pos@count by 1\relax\fi%
      \the\pos@count%
      \ifx\temp@\second@\ifgerm@n .\else {,}\fi\fi%
   \else%
      \ifx\second@\temp@ \else 0\fi%
   \fi%
   \ifx\second@\temp@%
      \the\loopcount%
      \loopcount=\csname mol@weight\first@\endcsname%
      \innerloopcount=\loopcount%
      \divide\loopcount by 10\relax%
      \multiply\loopcount by 10\relax%
      \advance\innerloopcount by -\loopcount\relax%
   \else%
      \divide\innerloopcount by 10\relax%
      \advance\innerloopcount by 5\relax%
      \divide\innerloopcount by 10\relax%
   \fi%
   \ifnum\innerloopcount>9\relax\innerloopcount=0\relax\fi%
   \ifgerm@n {,}\else .\fi%
   \the\innerloopcount~\third@}}
\newcommand\charge[2][o]{%
 \xdef\temp@{pep}%
 \ifx\prefix@\temp@%
   \xdef\first@{#2 @} \expandafter\check@letter\first@
   \xdef\first@{#2}   \ifletter \get@name@number \fi
   \xdef\second@{\csname @rd\first@\endcsname}%
   \loopcount=\csname ch@rge\second@\endcsname%
   \xdef\first@{#1}\make@lower%
   \if\first@ i\fi%
   \if\first@ o\advance\loopcount by \chargeNterm%
               \advance\loopcount by \chargeCterm\fi%
   \if\first@ n\advance\loopcount by \chargeNterm\fi%
   \if\first@ c\advance\loopcount by \chargeCterm\fi%
   \hbox{\ensuremath{%
   \ifnum\loopcount>0 +%
   \else\ifnum\loopcount=0 \pm%
   \else -\multiply\loopcount by -1\relax%
   \fi\fi%
   \innerloopcount=\loopcount%
   \divide\loopcount by 1000\relax%
   \the\loopcount%
   \multiply\loopcount by 1000\relax%
   \advance\innerloopcount by -\loopcount\relax%
   \divide\innerloopcount by 10\relax%
   \ifnum\innerloopcount=0%
   \else%
      \ifgerm@n {,}\else .\fi%
      \ifnum\innerloopcount<10 0\fi%
      \the\innerloopcount%
   \fi}}%
 \fi}
\def\TeXshade{%
   \setbox1=\hbox{\texttt{H}}%
   \def\logo@rule{\vrule depth0.25\ht1 height1.25\ht1 width\wd1}%
   \TeX%
   \logo@rule\kern-\wd1\textcolor{White}{\texttt{s}}%
   \logo@rule\kern-\wd1\textcolor{White}{\texttt{h}}%
   \texttt{a}%
   \logo@rule\kern-\wd1\textcolor{White}{\texttt{d}}%
   \texttt{e}}

\def\includeTCoffee#1{%
   \xdef\TC@first@{#1}\include@T@coffee}   
   
\def\firstcolumnDSSP{\xdef\fc@DSSP{y}}
\def\secondcolumnDSSP{\xdef\fc@DSSP{n}}
   
\newcommand{\includeDSSP}[3][existing]{%
  \temp@count=\dssp@num
  \advance\temp@count by 1
  \xdef\dssp@num{\the\temp@count}
  \xdef\first@{#1} \xdef\temp@{existing}
  \ifx\first@\temp@ \else\xdef\first@{make new}\fi
  \expandafter\xdef\csname optiondssp\the\temp@count\endcsname{\first@}
  \xdef\first@{#2 @} \expandafter\check@letter\first@
  \xdef\first@{#2}   \ifletter \get@name@number \fi
  \expandafter\xdef\csname doseqdssp\the\temp@count\endcsname{\first@}
  \expandafter\xdef\csname filenamedssp\the\temp@count\endcsname{#3}
  \expandafter\ifnum\csname doseqdssp\the\temp@count\endcsname>\seq@count
     \message{<Ignoring `#2' in \noexpand\includeDSSP>}
     \advance\temp@count by -1
     \xdef\dssp@num{\the\temp@count}
  \fi
}

\newcommand{\includeHMMTOP}[3][existing]{%
  \temp@count=\HMMTOP@num
  \advance\temp@count by 1
  \xdef\HMMTOP@num{\the\temp@count}
  \xdef\first@{#1} \xdef\temp@{existing}
  \ifx\first@\temp@ \else\xdef\first@{make new}\fi
  \expandafter\xdef\csname optionHMMTOP\the\temp@count\endcsname{\first@}
  \xdef\first@{#2[,]&}\expandafter\opt@color\first@
  \ifx\f@color\comm@ 
     \expandafter\xdef\csname fileseqHMMTOP\the\temp@count\endcsname{0} 
  \else
     \expandafter\xdef\csname fileseqHMMTOP\the\temp@count\endcsname{\f@color} 
  \fi
  \xdef\first@{\fourth@ @} \expandafter\check@letter\first@
  \xdef\first@{\fourth@}   \ifletter \get@name@number \fi       
  \expandafter\xdef\csname doseqHMMTOP\the\temp@count\endcsname{\first@}
  \expandafter\xdef\csname filenameHMMTOP\the\temp@count\endcsname{#3}
  \expandafter\ifnum\csname doseqHMMTOP\the\temp@count\endcsname>\seq@count
     \message{<Ignoring `#2' in \noexpand\includeHMMTOP>}
     \advance\temp@count by -1
     \xdef\HMMTOP@num{\the\temp@count}
  \fi
}

\newcommand{\includeSTRIDE}[3][existing]{%
  \temp@count=\stride@num
  \advance\temp@count by 1
  \xdef\stride@num{\the\temp@count}
  \xdef\first@{#1} \xdef\temp@{existing}
  \ifx\first@\temp@ \else\xdef\first@{make new}\fi
  \expandafter\xdef\csname optionstride\the\temp@count\endcsname{\first@}
  \xdef\first@{#2 @} \expandafter\check@letter\first@
  \xdef\first@{#2}   \ifletter \get@name@number \fi
  \expandafter\xdef\csname doseqstride\the\temp@count\endcsname{\first@}
  \expandafter\xdef\csname filenamestride\the\temp@count\endcsname{#3}
  \expandafter\ifnum\csname doseqstride\the\temp@count\endcsname>\seq@count
     \message{<Ignoring `#2' in \noexpand\includeSTRIDE>}
     \advance\temp@count by -1
     \xdef\stride@num{\the\temp@count}
  \fi
}
\newcommand{\includePHDsec}[3][existing]{%
  \temp@count=\PHD@num
  \advance\temp@count by 1
  \xdef\PHD@num{\the\temp@count}
  \xdef\first@{#1} \xdef\temp@{existing}
  \ifx\first@\temp@ \else\xdef\first@{make new}\fi
  \expandafter\xdef\csname optionphd\the\temp@count\endcsname{\first@}
  \xdef\first@{#2 @} \expandafter\check@letter\first@
  \xdef\first@{#2}   \ifletter \get@name@number \fi
  \expandafter\xdef\csname doseqphd\the\temp@count\endcsname{\first@}
  \expandafter\xdef\csname modephd\the\temp@count\endcsname{structure}
  \expandafter\xdef\csname filenamephd\the\temp@count\endcsname{#3}
  \expandafter\ifnum\csname doseqphd\the\temp@count\endcsname>\seq@count
     \message{<Ignoring `#2' in \noexpand\includePHDsec>}
     \advance\temp@count by -1
     \xdef\PHD@num{\the\temp@count}
  \fi
}
\newcommand{\includePHDtopo}[3][existing]{%
  \temp@count=\PHD@num
  \advance\temp@count by 1
  \xdef\PHD@num{\the\temp@count}
  \xdef\first@{#1} \xdef\temp@{existing}
  \ifx\first@\temp@ \else\xdef\first@{make new}\fi
  \expandafter\xdef\csname optionphd\the\temp@count\endcsname{\first@}
  \xdef\first@{#2 @} \expandafter\check@letter\first@
  \xdef\first@{#2}   \ifletter \get@name@number \fi
  \expandafter\xdef\csname doseqphd\the\temp@count\endcsname{\first@}
  \expandafter\xdef\csname modephd\the\temp@count\endcsname{topology}
  \expandafter\xdef\csname filenamephd\the\temp@count\endcsname{#3}
  \expandafter\ifnum\csname doseqphd\the\temp@count\endcsname>\seq@count
     \message{<Ignoring `#2' in \noexpand\includePHDtopo>}
     \advance\temp@count by -1
     \xdef\PHD@num{\the\temp@count}
  \fi
}
\def\appearance#1#2#3#4#5{%
   \xdef\first@{#1} \xdef\second@{#2}
   \xdef\temp@{PHDsec}
   \ifx\temp@\first@
     \xdef\temp@{alpha}
     \ifx\second@\temp@
      \def\bottop@Hsec{#3}
      \def\label@Hsec{#4}
      \def\text@Hsec{#5}
     \else
     \xdef\temp@{beta}
     \ifx\second@\temp@
      \def\bottop@Esec{#3}
      \def\label@Esec{#4}
      \def\text@Esec{#5}
     \fi\fi
   \else
   \xdef\temp@{PHDtopo}
   \ifx\temp@\first@
     \xdef\temp@{internal}
     \ifx\second@\temp@
      \def\bottop@itop{#3}
      \def\label@itop{#4}
      \def\text@itop{#5}
     \else
     \xdef\temp@{external}
     \ifx\second@\temp@
      \def\bottop@etop{#3}
      \def\label@etop{#4}
      \def\text@etop{#5}
     \else
     \xdef\temp@{TM}
     \ifx\second@\temp@
      \def\bottop@TMtop{#3}
      \def\label@TMtop{#4}
      \def\text@TMtop{#5}
     \fi\fi\fi
   \else
   \xdef\temp@{STRIDE}
   \ifx\temp@\first@
     \xdef\temp@{alpha}
     \ifx\second@\temp@
      \def\bottop@Hstride{#3}
      \def\label@Hstride{#4}
      \def\text@Hstride{#5}
     \else
     \xdef\temp@{3-10}
     \ifx\second@\temp@
      \def\bottop@Gstride{#3}
      \def\label@Gstride{#4}
      \def\text@Gstride{#5}
     \else
     \xdef\temp@{pi}
     \ifx\second@\temp@
      \def\bottop@Istride{#3}
      \def\label@Istride{#4}
      \def\text@Istride{#5}
     \else
     \xdef\temp@{beta}
     \ifx\second@\temp@
      \def\bottop@Estride{#3}
      \def\label@Estride{#4}
      \def\text@Estride{#5}
     \else
     \xdef\temp@{bridge}
     \ifx\second@\temp@
      \def\bottop@Bstride{#3}
      \def\label@Bstride{#4}
      \def\text@Bstride{#5}
     \else
     \xdef\temp@{turn}
     \ifx\second@\temp@
      \def\bottop@Tstride{#3}
      \def\label@Tstride{#4}
      \def\text@Tstride{#5}
     \fi\fi\fi\fi\fi\fi
   \else
   \xdef\temp@{DSSP}
   \ifx\temp@\first@
     \xdef\temp@{alpha}
     \ifx\second@\temp@
      \def\bottop@Hdssp{#3}
      \def\label@Hdssp{#4}
      \def\text@Hdssp{#5}
     \else
     \xdef\temp@{3-10}
     \ifx\second@\temp@
      \def\bottop@Gdssp{#3}
      \def\label@Gdssp{#4}
      \def\text@Gdssp{#5}
     \else
     \xdef\temp@{pi}
     \ifx\second@\temp@
      \def\bottop@Idssp{#3}
      \def\label@Idssp{#4}
      \def\text@Idssp{#5}
     \else
     \xdef\temp@{beta}
     \ifx\second@\temp@
      \def\bottop@Edssp{#3}
      \def\label@Edssp{#4}
      \def\text@Edssp{#5}
     \else
     \xdef\temp@{bridge}
     \ifx\second@\temp@
      \def\bottop@Bdssp{#3}
      \def\label@Bdssp{#4}
      \def\text@Bdssp{#5}
     \else
     \xdef\temp@{turn}
     \ifx\second@\temp@
      \def\bottop@Tdssp{#3}
      \def\label@Tdssp{#4}
      \def\text@Tdssp{#5}
     \else
     \xdef\temp@{bend}
     \ifx\second@\temp@
      \def\bottop@Sdssp{#3}
      \def\label@Sdssp{#4}
      \def\text@Sdssp{#5}
     \fi\fi\fi\fi\fi\fi\fi
   \else
   \xdef\temp@{HMMTOP}
   \ifx\temp@\first@
     \xdef\temp@{internal}
     \ifx\second@\temp@
      \def\bottop@i@HMMTOP{#3}
      \def\label@i@HMMTOP{#4}
      \def\text@i@HMMTOP{#5}
     \else
     \xdef\temp@{external}
     \ifx\second@\temp@
      \def\bottop@e@HMMTOP{#3}
      \def\label@e@HMMTOP{#4}
      \def\text@e@HMMTOP{#5}
     \else
     \xdef\temp@{TM}
     \ifx\second@\temp@
      \def\bottop@TM@HMMTOP{#3}
      \def\label@TM@HMMTOP{#4}
      \def\text@TM@HMMTOP{#5}
     \fi\fi\fi
   \fi\fi\fi\fi\fi
}

\def\showonDSSP#1{%
   \xdef\first@{#1,&,@} \xdef\second@{yes} \show@DSSP}
\def\hideonDSSP#1{%
   \xdef\first@{#1,&,@} \xdef\second@{no} \show@DSSP}

\def\showonSTRIDE#1{%
   \xdef\first@{#1,&,@} \xdef\second@{yes} \show@STRIDE}
\def\hideonSTRIDE#1{%
   \xdef\first@{#1,&,@} \xdef\second@{no} \show@STRIDE}

\def\showonPHDtopo#1{%
   \xdef\first@{#1,&,@} \xdef\second@{yes} \show@PHDtopo}
\def\hideonPHDtopo#1{%
   \xdef\first@{#1,&,@} \xdef\second@{no} \show@PHDtopo}

\def\showonPHDsec#1{%
   \xdef\first@{#1,&,@} \xdef\second@{yes} \show@PHDsec}
\def\hideonPHDsec#1{%
   \xdef\first@{#1,&,@} \xdef\second@{no} \show@PHDsec}

\def\showonHMMTOP#1{%
   \xdef\first@{#1,&,@} \xdef\second@{yes} \show@HMMTOP}
\def\hideonHMMTOP#1{%
   \xdef\first@{#1,&,@} \xdef\second@{no} \show@HMMTOP}

\def\codon#1#2{%
   \xdef\first@{#1}
   \xdef\second@{#2,&,@}
   \expandafter\get@triplet\second@}
\def\geneticcode#1{%
   \xdef\first@{#1}
   \xdef\temp@{standard}
   \ifx\first@\temp@
      \c@d@ns
   \else
  \input{#1.cod}
   \fi}
\newcommand{\backtranslabel}[2][tiny]{%
   \def\trans@size{\csname #1\endcsname}
   \xdef\first@{#2}
   \xdef\temp@{horizontal}
   \ifx\temp@\first@ \xdef\tr@nsstyle{0}\fi
   \xdef\temp@{zigzag}
   \ifx\temp@\first@ \xdef\tr@nsstyle{1}\fi
   \xdef\temp@{alternating}
   \ifx\temp@\first@ \xdef\tr@nsstyle{2}\fi
   \xdef\temp@{oblique}
   \ifx\temp@\first@ \xdef\tr@nsstyle{3}\fi
   \xdef\temp@{vertical}
   \ifx\temp@\first@ \xdef\tr@nsstyle{4}\fi
}
\newcommand{\backtranstext}[2][tiny]{%
   \def\transtext@size{\csname #1\endcsname}
   \xdef\first@{#2}
   \xdef\temp@{horizontal}
   \ifx\temp@\first@ \xdef\tr@nstextstyle{0}\fi
   \xdef\temp@{zigzag}
   \ifx\temp@\first@ \xdef\tr@nstextstyle{1}\fi
   \xdef\temp@{alternating}
   \ifx\temp@\first@ \xdef\tr@nstextstyle{2}\fi
   \xdef\temp@{oblique}
   \ifx\temp@\first@ \xdef\tr@nstextstyle{3}\fi
   \xdef\temp@{vertical}
   \ifx\temp@\first@ \xdef\tr@nstextstyle{4}\fi
}

\newcommand\exportconsensus[2][export.txt]{%
   \ifx\exp@rt\n@
      \xdef\first@{#2 @} \expandafter\check@letter\first@
      \xdef\first@{#2}   \ifletter \get@name@number \fi
      \xdef\exp@rt@num{\first@}
      \xdef\exp@rt{y}
      \immediate\openout\exp@rtfile = #1
   \fi
}

\def\setdomain#1#2{%
   \xdef\first@{#1}
   \xdef\temp@{consensus}
   \ifx\first@\temp@ \xdef\domain@seq{0}
   \else
      \xdef\first@{#1 @} \expandafter\check@letter\first@
      \xdef\first@{#1}   \ifletter \get@name@number \fi
      \xdef\domain@seq{\csname @rd\first@\endcsname}%
   \fi
   \xdef\list@{#2,&}
   \xdef\temp@{#2,,,:,,,,@} \expandafter\test@PDB\temp@
   \loop
        \xdef\list@{\list@ @}
        \expandafter\get@domainregions\list@
   \ifx\list@\ampers@nd\else\repeat
   \xdef\dom@in{y}
}

\def\decimal@IOO#1#2#3{#1#2\temp@#3}
\def\decimal@IO#1#2{#1\temp@#2}
\def\decimal@I#1{0\temp@#1}

\def\percentsimilarity#1#2{%
   \xdef\first@{#2@}\expandafter\check@letter\first@%
   \xdef\first@{#2}\xdef\second@@{#2}%
   \ifletter\get@name@number@table\xdef\second@@{\first@}\fi%
   \xdef\first@{#1@}\expandafter\check@letter\first@%
   \xdef\first@{#1}\xdef\first@@{#1}%
   \ifletter\get@name@number@table\xdef\first@@{\first@}\fi%
   \ifnum\first@@<\second@@%
      \loopcount=\csname simcount\first@@ @\second@@\endcsname%
      \temp@count=\csname poscount\first@@ @\second@@\endcsname%
   \else%
   \ifnum\first@@=\second@@ \loopcount=1 \temp@count=1%
   \else%
      \loopcount=\csname simcount\second@@ @\first@@\endcsname%
      \temp@count=\csname poscount\second@@ @\first@@\endcsname%
   \fi%
   \fi%
   \multiply\loopcount by 1000%
   \divide\loopcount by \temp@count%
   \xdef\first@{\the\loopcount}%
   \ifgerm@n\xdef\temp@{,}\else\xdef\temp@{.}\fi%
   \ifnum\first@=1000 100\temp@0%
   \else%
   \ifnum\first@>99 \expandafter\decimal@IOO\first@%
   \else%
   \ifnum\first@>9 \expandafter\decimal@IO\first@%
   \else%
   \ifnum\first@>0 \expandafter\decimal@I\first@%
   \else 0\temp@0%
   \fi%
   \fi%
   \fi%
   \fi%  
}

\def\percentidentity#1#2{%
   \xdef\first@{#2@}\expandafter\check@letter\first@%
   \xdef\first@{#2}\xdef\second@@{#2}%
   \ifletter\get@name@number@table\xdef\second@@{\first@}\fi%
   \xdef\first@{#1@}\expandafter\check@letter\first@%
   \xdef\first@{#1}\xdef\first@@{#1}%
   \ifletter\get@name@number@table\xdef\first@@{\first@}\fi%
   \ifnum\first@@<\second@@%
      \loopcount=\csname identcount\first@@ @\second@@\endcsname%
      \temp@count=\csname poscount\first@@ @\second@@\endcsname%
   \else%
   \ifnum\first@@=\second@@ \loopcount=1 \temp@count=1%
   \else%
      \loopcount=\csname identcount\second@@ @\first@@\endcsname%
      \temp@count=\csname poscount\second@@ @\first@@\endcsname%
   \fi%
   \fi%
   \multiply\loopcount by 1000%
   \divide\loopcount by \temp@count%
   \xdef\first@{\the\loopcount}%
   \ifgerm@n\xdef\temp@{,}\else\xdef\temp@{.}\fi%
   \ifnum\first@=1000 100\temp@0%
   \else%
   \ifnum\first@>99 \expandafter\decimal@IOO\first@%
   \else%
   \ifnum\first@>9 \expandafter\decimal@IO\first@%
   \else%
   \ifnum\first@>0 \expandafter\decimal@I\first@%
   \else 0\temp@0%
   \fi%
   \fi%
   \fi%
   \fi%  
}

\def\similaritytable{%
   \def\name@loop{%
      \advance\outerloopcount by 1
      \ifnum\outerloopcount>\seq@num\relax
      \else
      \immediate\write\exp@rtfile{&\string\multicolumn{1}{c}{\string\kern1ex\string\begin{rotopo}{90}\csname newseqname\the\outerloopcount\endcsname\string\end{rotopo}}}
      \name@loop\fi
   }
   \def\inner@loop{% 
      \advance\innerloopcount by 1
      \ifnum\innerloopcount>\seq@num\relax
      \else
      \ifnum\innerloopcount=\outerloopcount\relax \immediate\write\exp@rtfile{& {---\hss}}\fi
      \ifnum\innerloopcount>\outerloopcount\relax \immediate\write\exp@rtfile{& \string\percentsimilarity{\the\outerloopcount}{\the\innerloopcount}}\fi
      \ifnum\innerloopcount<\outerloopcount\relax \immediate\write\exp@rtfile{& \string\percentidentity{\the\outerloopcount}{\the\innerloopcount}}\fi
      \inner@loop\fi
   }
   \def\outer@loop{%
      \advance\outerloopcount by 1
      \ifnum\outerloopcount>\seq@num\relax
      \else
      \immediate\write\exp@rtfile{\csname newseqname\the\outerloopcount\endcsname}
      \innerloopcount=0 \inner@loop \immediate\write\exp@rtfile{&\string\\}
      \outer@loop\fi
   }
   \xdef\temp@{l|}
   \outerloopcount=0
   \loop
      \advance\outerloopcount by 1
      \xdef\temp@{\temp@ r}
   \ifnum\outerloopcount<\seq@num\repeat
   \advance\outerloopcount by 1 \xdef\seq@num@plus{\the\outerloopcount}
   \advance\outerloopcount by 1 \xdef\seq@num@plus@plus{\the\outerloopcount}
   \ifnum\seq@num<3 \xdef\temp@@{$\approx$} \else
   \ifnum\seq@num<5 \xdef\temp@@{simil.}\ifgerm@n\def\temp@@{\string\"{A}hnl.}\fi\ifsp@nish\xdef\temp@@{simil.}\fi\else
   \xdef\temp@@{similarity} 
      \ifgerm@n\ifnum\seq@num=5\def\temp@@{\string\"{A}hnlichk.}\else\def\temp@@{\string\"{A}hnlichkeit}\fi\fi\ifsp@nish\xdef\temp@@{similitud}\fi\fi
   \fi
   \xdef\temp@@@{identity}\ifgerm@n\def\temp@@@{Identit\string\"{a}t}\fi\ifsp@nish\xdef\temp@@@{identidad}\fi
   
   \immediate\openout\exp@rtfile = simtable.tmp
   \immediate\write\exp@rtfile{\string\begin{tabular}{\temp@ |r}}
   \immediate\write\exp@rtfile{\string\multicolumn{\seq@num@plus@plus}{c}{}\string\\[\name@@width]}
   \immediate\write\exp@rtfile{\string\multicolumn{1}{c}{}}
   \outerloopcount=0 \name@loop \immediate\write\exp@rtfile{&\string\\ \string\cline{2-\seq@num@plus}}
   \immediate\write\exp@rtfile{&\string\multicolumn{\seq@num}{|c|}{}&\string\\[-2ex]}
   \outerloopcount=0 \outer@loop 
   \immediate\write\exp@rtfile{\string\cline{2-\seq@num@plus}}
   \immediate\write\exp@rtfile{\string\multicolumn{\seq@num@plus}{c}{}&}
   \immediate\write\exp@rtfile{\string\multicolumn{1}{r}{\string\kern1.5ex\string\begin{rotopo}{90}\kern1em\% \temp@@\string\end{rotopo}}\string\\[-2ex]}
   \immediate\write\exp@rtfile{\string\multicolumn{1}{c}{}&\string\multicolumn{\seq@num}{l}{\% \temp@@@}&\string\\}
   \immediate\write\exp@rtfile{\string\end{tabular}}
   \immediate\closeout\exp@rtfile

   \input{simtable.tmp}
}

\def\identitytable{\similaritytable}

%%%%%  Calculate consensus

\def\check@sim{%
 \xdef\first@{\csname res\the\loopcount\endcsname}
 \xdef\first@{\csname \prefix@ grp\first@\endcsname}
 \newrestrue
 \ifnum\first@<0 \newresfalse
 \else
 \innerloopcount=\loopcount
 \ifnum\loopcount=\cons@num \innerloopcount=\seq@count \fi
 \ifnum\innerloopcount<\seq@count
 \loop
    \advance\innerloopcount by 1
    \xdef\second@{\csname res\the\innerloopcount\endcsname}
    \expandafter\ifx\csname \prefix@ grp\second@\endcsname\first@
       \newresfalse \innerloopcount=\seq@count \fi
 \ifnum\innerloopcount<\seq@count \repeat
 \fi
 \fi

 \ifnewres
  \pos@sum=0
  \innerloopcount=0
  \loop
    \advance\innerloopcount by 1
    \xdef\second@{\csname res\the\innerloopcount\endcsname}
    \expandafter\ifx\csname \prefix@ grp\second@\endcsname\first@
       \advance\pos@sum by 1 \fi
  \ifnum\innerloopcount<\seq@count \repeat

%  \multiply\pos@sum by \seq@percent
  \expandafter\xdef\csname pos\the\loopcount\endcsname{\the\pos@sum}
%  \expandafter\ifnum\csname pos\the\loopcount\endcsname<\thresh@ld
  \expandafter\ifnum\csname pos\the\loopcount\endcsname<\thresh@ld@
  \else
    \expandafter\ifnum\csname pos\the\loopcount\endcsname>\m@x
      \xdef\m@x{\csname pos\the\loopcount\endcsname}
      \xdef\cons@seq{\the\loopcount} \xdef\match@case{\c@se}
      \xdef\simgroup@{\first@}
    \else
     \expandafter\ifnum\csname pos\the\loopcount\endcsname=\m@x
      \xdef\match@case{0}
     \fi
    \fi
  \fi
 \fi

 \ifnum\loopcount=\cons@num \loopcount=1 \fi
 \advance\loopcount by -1
 \ifnum\loopcount>0 \check@sim \fi}

\def\check@ident{%
 \xdef\first@{\csname res\the\loopcount\endcsname}
 \newrestrue \expandafter\check@char\first@
 \ifletter
   \innerloopcount=\loopcount
   \ifnum\loopcount=\cons@num \innerloopcount=\seq@count \fi
   \ifnum\innerloopcount<\seq@count
     \loop
       \advance\innerloopcount by 1
       \expandafter\ifx\csname res\the\innerloopcount\endcsname\first@
          \newresfalse \innerloopcount=\seq@count  \fi
     \ifnum\innerloopcount<\seq@count \repeat
   \fi
 \else
    \newresfalse
%    \expandafter\xdef\csname res\the\loopcount\endcsname{\d@t}  %%%%% or \equ@l for domains!!!
 \fi

 \ifnewres
  \pos@sum=0
  \innerloopcount=0
  \loop
    \advance\innerloopcount by 1
    \expandafter\ifx\csname res\the\innerloopcount\endcsname\first@
       \advance\pos@sum by 1 \fi
  \ifnum\innerloopcount<\seq@count \repeat

  \expandafter\xdef\csname pos\the\loopcount\endcsname{\the\pos@sum}
  \expandafter\ifnum\csname pos\the\loopcount\endcsname=\seq@count
    \xdef\cons@seq{\the\loopcount} \xdef\match@case{2} \loopcount=1
  \else
%   \multiply\pos@sum by \seq@percent
   \expandafter\xdef\csname pos\the\loopcount\endcsname{\the\pos@sum}
%   \expandafter\ifnum\csname pos\the\loopcount\endcsname<\thresh@ld
   \expandafter\ifnum\csname pos\the\loopcount\endcsname<\thresh@ld@
   \else
    \expandafter\ifnum\csname pos\the\loopcount\endcsname>\m@x
      \xdef\m@x{\csname pos\the\loopcount\endcsname}
%      \expandafter\ifnum\csname pos\the\loopcount\endcsname<\all@thresh@ld
      \expandafter\ifnum\csname pos\the\loopcount\endcsname<\all@thresh@ld@
         \xdef\cons@seq{\the\loopcount} \xdef\match@case{1}
      \else
         \xdef\cons@seq{\the\loopcount} \xdef\match@case{2}
      \fi      
    \else
     \expandafter\ifnum\csname pos\the\loopcount\endcsname=\m@x
      \xdef\match@case{0}
     \fi
    \fi
   \fi
  \fi
 \fi

 \ifnum\loopcount=\cons@num \loopcount=1 \fi
 \advance\loopcount by -1
 \ifnum\loopcount>0 \check@ident \fi}

\def\get@simchar{%
 \xdef\first@{\csname res\the\loopcount\endcsname}
 \newrestrue \expandafter\check@char\first@
 \ifletter
   \innerloopcount=\loopcount
   \ifnum\loopcount=\cons@num \innerloopcount=\seq@count \fi
   \ifnum\innerloopcount<\seq@count
     \loop
       \advance\innerloopcount by 1
       \expandafter\ifx\csname res\the\innerloopcount\endcsname\first@
          \newresfalse \innerloopcount=\seq@count  \fi
     \ifnum\innerloopcount<\seq@count \repeat
   \fi
 \else
    \newresfalse
 \fi

 \ifnewres
  \pos@sum=0
  \innerloopcount=0
  \loop
    \advance\innerloopcount by 1
    \expandafter\ifx\csname res\the\innerloopcount\endcsname\first@
       \xdef\second@{\csname res\the\innerloopcount\endcsname}
       \expandafter\ifx\csname \prefix@ grp\second@\endcsname\simgroup@
          \advance\pos@sum by 1 \fi
    \fi
  \ifnum\innerloopcount<\seq@count \repeat

  \expandafter\xdef\csname pos\the\loopcount\endcsname{\the\pos@sum}
  \expandafter\ifnum\csname pos\the\loopcount\endcsname>\m@x
      \xdef\m@x{\csname pos\the\loopcount\endcsname}
      \xdef\cons@seq{\the\loopcount}
  \fi
 \fi

 \ifnum\loopcount=\cons@num \loopcount=1 \fi
 \advance\loopcount by -1
 \ifnum\loopcount>0 \get@simchar \fi}

\def\unc@nserved{%
 \ifsimmode
    \ifnum\cons@num>0 \loopcount=\cons@num \else \loopcount=\seq@count \fi
    \xdef\match@case{0} \xdef\m@x{1} \check@sim
    \ifnum\match@case=0
       \loopcount=0
       \loop
          \advance\loopcount by 1
          \xdef\first@{noshade}
          \expandafter\ifx\csname hide@seq\the\loopcount\endcsname\first@
            \expandafter\ifx\csname res\the\loopcount\endcsname\d@t
             \expandafter\xdef\csname seq\the\loopcount\endcsname{%
                \csname seq\the\loopcount\endcsname3\gap@char}
            \else
            \expandafter\ifx\csname res\the\loopcount\endcsname\questi@n
             \expandafter\xdef\csname seq\the\loopcount\endcsname{%
                \csname seq\the\loopcount\endcsname3\st@p@char}
            \else
             \expandafter\ifx\csname res\the\loopcount\endcsname\equ@l
              \expandafter\xdef\csname seq\the\loopcount\endcsname{%
                 \csname seq\the\loopcount\endcsname!\dom@char} \xdef\g@p{y}
             \else
             \expandafter\xdef\csname seq\the\loopcount\endcsname{%
                \csname seq\the\loopcount\endcsname%
                3\csname res\the\loopcount\endcsname}
             \fi
             \fi
             \fi
          \else
           \xdef\first@{\csname res\the\loopcount\endcsname}
           \ifx\first@\equ@l
             \def\first@{\dom@char} \def\third@{9}
             \expandafter\ifx\csname lower@seq\the\loopcount\endcsname\y@
                   \expandafter\xdef\csname lower@seq\the\loopcount\endcsname{n}
             \xdef\first@{;\first@}\fi
             \expandafter\ifx\csname tint@seq\the\loopcount\endcsname\y@
                   \expandafter\xdef\csname tint@seq\the\loopcount\endcsname{n}
             \xdef\first@{=\first@}\fi
             \expandafter\ifx\csname emph@seq\the\loopcount\endcsname\y@
                   \expandafter\xdef\csname emph@seq\the\loopcount\endcsname{n}
             \xdef\first@{,\first@}\fi
             \expandafter\xdef\csname seq\the\loopcount\endcsname{%
                              \csname seq\the\loopcount\endcsname\third@\first@} \xdef\g@p{y}
           \else
             \ifx\first@\d@t
                 \expandafter\ifnum\csname res@count\the\loopcount\endcsname%
                      <\csname seq@start\the\loopcount\endcsname
                    \ifsh@wg@ps
                       \def\first@{\gap@char}\def\third@{7}
                    \else
                       \def\first@{{}} \def\third@{8}
                    \fi
                 \else
                    \expandafter\ifnum\csname res@count\the\loopcount\endcsname%
                        <\csname seq@len\the\loopcount\endcsname
                       \def\third@{7}\def\first@{\gap@char}
                    \else
                    \ifsh@wg@ps
                       \def\first@{\gap@char}\def\third@{7}
                    \else
                       \def\first@{{}} \def\third@{8}
                    \fi
                    \fi
                 \fi
                \expandafter\ifx\csname lower@seq\the\loopcount\endcsname\y@
                   \expandafter\xdef\csname lower@seq\the\loopcount\endcsname{n}
                \xdef\first@{;\first@}\fi
                \expandafter\ifx\csname tint@seq\the\loopcount\endcsname\y@
                   \expandafter\xdef\csname tint@seq\the\loopcount\endcsname{n}
                \xdef\first@{=\first@}\fi
                \expandafter\ifx\csname emph@seq\the\loopcount\endcsname\y@
                   \expandafter\xdef\csname emph@seq\the\loopcount\endcsname{n}
                \xdef\first@{,\first@}\fi
                \expandafter\xdef\csname seq\the\loopcount\endcsname{%
                                 \csname seq\the\loopcount\endcsname\third@\first@}
             \else
             \ifx\first@\questi@n
                 \expandafter\ifnum\csname res@count\the\loopcount\endcsname%
                      <\csname seq@start\the\loopcount\endcsname
                    \def\first@{{}} \def\third@{8}
                 \else
                    \expandafter\ifnum\csname res@count\the\loopcount\endcsname%
                        <\csname seq@len\the\loopcount\endcsname
                       \def\third@{7}\def\first@{\st@p@char}
                    \else
                       \def\first@{{}} \def\third@{8}
                    \fi
                 \fi
                \expandafter\ifx\csname lower@seq\the\loopcount\endcsname\y@
                   \expandafter\xdef\csname lower@seq\the\loopcount\endcsname{n}
                \xdef\first@{;\first@}\fi
                \expandafter\ifx\csname tint@seq\the\loopcount\endcsname\y@
                   \expandafter\xdef\csname tint@seq\the\loopcount\endcsname{n}
                \xdef\first@{=\first@}\fi
                \expandafter\ifx\csname emph@seq\the\loopcount\endcsname\y@
                   \expandafter\xdef\csname emph@seq\the\loopcount\endcsname{n}
                \xdef\first@{,\first@}\fi
                \expandafter\xdef\csname seq\the\loopcount\endcsname{%
                                 \csname seq\the\loopcount\endcsname\third@\first@}
             \else
               \ifnum\divref@=\loopcount\else 
               \xdef\low@up{lower} \ifx\resn@m@tch\low@up \make@lower \else
               \xdef\low@up{upper} \ifx\resn@m@tch\low@up
               \else \xdef\first@{\resn@m@tch} \fi\fi\fi
               \expandafter\ifx\csname lower@seq\the\loopcount\endcsname\y@
                  \expandafter\xdef\csname lower@seq\the\loopcount\endcsname{n}
                  \xdef\first@{;\first@}\fi
               \expandafter\ifx\csname tint@seq\the\loopcount\endcsname\y@
                  \expandafter\xdef\csname tint@seq\the\loopcount\endcsname{n}
                  \xdef\first@{=\first@}\fi
               \expandafter\ifx\csname emph@seq\the\loopcount\endcsname\y@
                  \expandafter\xdef\csname emph@seq\the\loopcount\endcsname{n}
                  \xdef\first@{,\first@}\fi
               \expandafter\xdef\csname seq\the\loopcount\endcsname{%
                  \csname seq\the\loopcount\endcsname3\first@}
             \fi
             \fi
           \fi
          \fi
       \ifnum\loopcount<\seq@count \repeat
       \xdef\low@up{lower} \ifx\n@m@tch\low@up \xdef\first@{{ }} \else
       \xdef\low@up{upper} \ifx\n@m@tch\low@up \xdef\first@{{ }}
       \else \xdef\first@{\n@m@tch} \fi\fi
       \expandafter\ifx\csname lower@seq0\endcsname\y@
          \expandafter\xdef\csname lower@seq0\endcsname{n}
          \xdef\first@{;\first@}\fi
       \expandafter\ifx\csname tint@seq0\endcsname\y@
          \expandafter\xdef\csname tint@seq0\endcsname{n}
          \xdef\first@{=\first@}\fi
       \expandafter\ifx\csname emph@seq0\endcsname\y@
          \expandafter\xdef\csname emph@seq0\endcsname{n}
          \xdef\first@{,\first@}\fi
       \ifx\g@p\y@
          \xdef\consensus{\consensus 8{}} 
       \else
          \xdef\consensus{\consensus 4\first@} 
          \expandafter\ifx\csname res\cons@num\endcsname\d@t
          \else\xdef\constopo{\constopo 0}\fi
       \fi
    \else
     \ifnum\cons@num>0
       \xdef\tmp@{\csname res\cons@num\endcsname}
     \else
       \xdef\m@x{0} \loopcount=\seq@count \get@simchar
       \xdef\tmp@{\csname res\cons@seq\endcsname}
     \fi
     \xdef\second@{\csname \prefix@ grp\tmp@\endcsname}
     \loopcount=0
     \loop
        \advance\loopcount by 1
        \xdef\first@{noshade}
        \expandafter\ifx\csname hide@seq\the\loopcount\endcsname\first@
             \xdef\first@{\csname res\the\loopcount\endcsname}
             \ifx\first@\d@t\def\first@{\gap@char}\fi
             \ifx\first@\questi@n\def\first@{\st@p@char}\fi
             \ifx\first@\equ@l\def\first@{\dom@char} \xdef\g@p{y} \fi
             \expandafter\ifx\csname lower@seq\the\loopcount\endcsname\y@
                \expandafter\xdef\csname lower@seq\the\loopcount\endcsname{n}
                \xdef\first@{;\first@}\fi
             \expandafter\ifx\csname tint@seq\the\loopcount\endcsname\y@
                \expandafter\xdef\csname tint@seq\the\loopcount\endcsname{n}
                \xdef\first@{=\first@}\fi
             \expandafter\ifx\csname emph@seq\the\loopcount\endcsname\y@
                \expandafter\xdef\csname emph@seq\the\loopcount\endcsname{n}
                \xdef\first@{,\first@}\fi
             \expandafter\xdef\csname seq\the\loopcount\endcsname{%
                \csname seq\the\loopcount\endcsname3\first@} 
        \else
          \xdef\first@{\csname res\the\loopcount\endcsname}
          \xdef\last@{\csname res\the\loopcount\endcsname}
          \expandafter\ifnum\csname \prefix@ grp\last@\endcsname=\second@
              \xdef\third@{2}
              \ifnum\divref@=\loopcount\else 
              \xdef\low@up{lower} \ifx\ressimm@tch\low@up \make@lower \else
              \xdef\low@up{upper} \ifx\ressimm@tch\low@up
              \else \xdef\first@{\ressimm@tch} \fi\fi\fi
          \else
            \ifx\first@\equ@l
              \def\first@{\dom@char} \def\third@{9}  \xdef\g@p{y}
            \else
              \ifx\first@\d@t
                 \expandafter\ifnum\csname res@count\the\loopcount\endcsname%
                      <\csname seq@start\the\loopcount\endcsname
                    \ifsh@wg@ps
                       \def\first@{\gap@char}\def\third@{7}
                    \else
                       \def\first@{{}} \def\third@{8}
                    \fi
                 \else
                    \expandafter\ifnum\csname res@count\the\loopcount\endcsname%
                         <\csname seq@len\the\loopcount\endcsname
                       \def\third@{7}\def\first@{\gap@char}
                    \else
                    \ifsh@wg@ps
                       \def\first@{\gap@char}\def\third@{7}
                    \else
                       \def\first@{{}} \def\third@{8}
                    \fi
                    \fi
                 \fi
              \else
              \ifx\first@\questi@n
                 \expandafter\ifnum\csname res@count\the\loopcount\endcsname%
                      <\csname seq@start\the\loopcount\endcsname
                       \def\first@{{}} \def\third@{8}
                 \else
                    \expandafter\ifnum\csname res@count\the\loopcount\endcsname%
                         <\csname seq@len\the\loopcount\endcsname
                       \def\third@{7}\def\first@{\st@p@char}
                    \else
                       \def\first@{{}} \def\third@{8}
                    \fi
                 \fi
              \else
                 \xdef\third@{3}
                 \ifnum\divref@=\loopcount\else 
                 \xdef\low@up{lower} \ifx\resn@m@tch\low@up \make@lower \else
                 \xdef\low@up{upper} \ifx\resn@m@tch\low@up
                 \else \xdef\first@{\resn@m@tch} \fi\fi\fi
              \fi
              \fi
            \fi
          \fi
          \expandafter\ifx\csname lower@seq\the\loopcount\endcsname\y@
                \expandafter\xdef\csname lower@seq\the\loopcount\endcsname{n}
                \xdef\first@{;\first@}\fi
          \expandafter\ifx\csname tint@seq\the\loopcount\endcsname\y@
                \expandafter\xdef\csname tint@seq\the\loopcount\endcsname{n}
                \xdef\first@{=\first@}\fi
          \expandafter\ifx\csname emph@seq\the\loopcount\endcsname\y@
                \expandafter\xdef\csname emph@seq\the\loopcount\endcsname{n}
                \xdef\first@{,\first@}\fi
          \expandafter\xdef\csname seq\the\loopcount\endcsname{%
              \csname seq\the\loopcount\endcsname\third@\first@}
        \fi
     \ifnum\loopcount<\seq@count \repeat
     \xdef\first@{\tmp@}
     \xdef\low@up{lower} \ifx\m@tch\low@up \make@lower \else
     \xdef\low@up{upper} \ifx\m@tch\low@up
     \else \xdef\first@{\m@tch} \fi\fi
     \expandafter\ifx\csname lower@seq0\endcsname\y@
                \expandafter\xdef\csname lower@seq0\endcsname{n}
                \xdef\first@{;\first@}\fi
     \expandafter\ifx\csname tint@seq0\endcsname\y@
                \expandafter\xdef\csname tint@seq0\endcsname{n}
                \xdef\first@{=\first@}\fi
     \expandafter\ifx\csname emph@seq0\endcsname\y@
                \expandafter\xdef\csname emph@seq0\endcsname{n}
                \xdef\first@{,\first@}\fi
     \ifx\g@p\y@
        \xdef\consensus{\consensus 8{}}
     \else
        \xdef\consensus{\consensus5\first@}
        \xdef\constopo{\constopo 1}
     \fi
    \fi
 \else
 \iffuncmode
    \loopcount=0
    \loop
       \advance\loopcount by 1
       \xdef\first@{\csname res\the\loopcount\endcsname}
       \ifx\first@\equ@l
          \def\first@{\dom@char} \def\third@{!}  \xdef\g@p{y}
          \expandafter\ifx\csname lower@seq\the\loopcount\endcsname\y@
                   \expandafter\xdef\csname lower@seq\the\loopcount\endcsname{n}
                   \xdef\first@{;\first@}\fi
          \expandafter\ifx\csname tint@seq\the\loopcount\endcsname\y@
                   \expandafter\xdef\csname tint@seq\the\loopcount\endcsname{n}
                   \xdef\first@{=\first@}\fi
          \expandafter\ifx\csname emph@seq\the\loopcount\endcsname\y@
                   \expandafter\xdef\csname emph@seq\the\loopcount\endcsname{n}
                   \xdef\first@{,\first@}\fi
          \expandafter\xdef\csname seq\the\loopcount\endcsname{%
             \csname seq\the\loopcount\endcsname\third@\first@}
       \else
       \ifx\first@\d@t
            \expandafter\ifnum\csname res@count\the\loopcount\endcsname%
                 <\csname seq@start\the\loopcount\endcsname
                    \ifsh@wg@ps
                       \def\first@{\gap@char}\def\third@{*}
                    \else
                       \def\first@{{}} \def\third@{/}
                    \fi
            \else
               \expandafter\ifnum\csname res@count\the\loopcount\endcsname%
                    <\csname seq@len\the\loopcount\endcsname
                  \def\first@{\gap@char}\def\third@{*}
               \else
                    \ifsh@wg@ps
                       \def\first@{\gap@char}\def\third@{*}
                    \else
                       \def\first@{{}} \def\third@{/}
                    \fi
               \fi
             \fi
          \expandafter\ifx\csname lower@seq\the\loopcount\endcsname\y@
                   \expandafter\xdef\csname lower@seq\the\loopcount\endcsname{n}
                   \xdef\first@{;\first@}\fi
          \expandafter\ifx\csname tint@seq\the\loopcount\endcsname\y@
                   \expandafter\xdef\csname tint@seq\the\loopcount\endcsname{n}
                   \xdef\first@{=\first@}\fi
          \expandafter\ifx\csname emph@seq\the\loopcount\endcsname\y@
                   \expandafter\xdef\csname emph@seq\the\loopcount\endcsname{n}
                   \xdef\first@{,\first@}\fi
          \expandafter\xdef\csname seq\the\loopcount\endcsname{%
             \csname seq\the\loopcount\endcsname\third@\first@}
       \else
       \ifx\first@\questi@n
            \expandafter\ifnum\csname res@count\the\loopcount\endcsname%
                 <\csname seq@start\the\loopcount\endcsname
                  \def\first@{{}} \def\third@{/}
            \else
               \expandafter\ifnum\csname res@count\the\loopcount\endcsname%
                    <\csname seq@len\the\loopcount\endcsname
                  \def\first@{\st@p@char}\def\third@{*}
               \else
                  \def\first@{{}} \def\third@{/}
               \fi
             \fi
          \expandafter\ifx\csname lower@seq\the\loopcount\endcsname\y@
                   \expandafter\xdef\csname lower@seq\the\loopcount\endcsname{n}
                   \xdef\first@{;\st@p@char}\fi
          \expandafter\ifx\csname tint@seq\the\loopcount\endcsname\y@
                   \expandafter\xdef\csname tint@seq\the\loopcount\endcsname{n}
                   \xdef\first@{=\st@p@char}\fi
          \expandafter\ifx\csname emph@seq\the\loopcount\endcsname\y@
                   \expandafter\xdef\csname emph@seq\the\loopcount\endcsname{n}
                   \xdef\first@{,\st@p@char}\fi
          \expandafter\xdef\csname seq\the\loopcount\endcsname{%
             \csname seq\the\loopcount\endcsname\third@\first@}
       \else
          \xdef\low@up{lower} \ifx\resn@m@tch\low@up \make@lower \else
          \xdef\low@up{upper} \ifx\resn@m@tch\low@up
          \else \xdef\first@{\resn@m@tch} \fi\fi
          \expandafter\ifx\csname lower@seq\the\loopcount\endcsname\y@
                   \expandafter\xdef\csname lower@seq\the\loopcount\endcsname{n}
                   \xdef\first@{;\first@}\fi
          \expandafter\ifx\csname tint@seq\the\loopcount\endcsname\y@
                   \expandafter\xdef\csname tint@seq\the\loopcount\endcsname{n}
                   \xdef\first@{=\first@}\fi
          \expandafter\ifx\csname emph@seq\the\loopcount\endcsname\y@
                   \expandafter\xdef\csname emph@seq\the\loopcount\endcsname{n}
                   \xdef\first@{,\first@}\fi
          \expandafter\xdef\csname seq\the\loopcount\endcsname{%
             \csname seq\the\loopcount\endcsname0\first@}
       \fi
       \fi
       \fi
    \ifnum\loopcount<\seq@count \repeat
 \else
    \loopcount=0
    \loop
        \advance\loopcount by 1
        \xdef\first@{noshade}
        \expandafter\ifx\csname hide@seq\the\loopcount\endcsname\first@
             \xdef\first@{\csname res\the\loopcount\endcsname}
             \ifx\first@\d@t\def\first@{\gap@char}\fi
             \ifx\first@\questi@n\def\first@{\st@p@char}\fi
             \ifx\first@\equ@l\def\first@{\dom@char}  \xdef\g@p{y} \fi
             \expandafter\ifx\csname lower@seq\the\loopcount\endcsname\y@
                \expandafter\xdef\csname lower@seq\the\loopcount\endcsname{n}
                \xdef\first@{;\first@}\fi
             \expandafter\ifx\csname tint@seq\the\loopcount\endcsname\y@
                \expandafter\xdef\csname tint@seq\the\loopcount\endcsname{n}
                \xdef\first@{=\first@}\fi
             \expandafter\ifx\csname emph@seq\the\loopcount\endcsname\y@
                \expandafter\xdef\csname emph@seq\the\loopcount\endcsname{n}
                \xdef\first@{,\first@}\fi
             \expandafter\xdef\csname seq\the\loopcount\endcsname{%
                \csname seq\the\loopcount\endcsname%
                3\first@}
        \else
         \xdef\first@{\csname res\the\loopcount\endcsname}
         \ifx\first@\equ@l
            \def\first@{\dom@char} \def\third@{9}  \xdef\g@p{y}
           \expandafter\ifx\csname lower@seq\the\loopcount\endcsname\y@
                   \expandafter\xdef\csname lower@seq\the\loopcount\endcsname{n}
                   \xdef\first@{;\first@}\fi
           \expandafter\ifx\csname tint@seq\the\loopcount\endcsname\y@
                   \expandafter\xdef\csname tint@seq\the\loopcount\endcsname{n}
                   \xdef\first@{=\first@}\fi
           \expandafter\ifx\csname emph@seq\the\loopcount\endcsname\y@
                   \expandafter\xdef\csname emph@seq\the\loopcount\endcsname{n}
                   \xdef\first@{,\first@}\fi
           \expandafter\xdef\csname seq\the\loopcount\endcsname{%
              \csname seq\the\loopcount\endcsname\third@\first@}
         \else
         \ifx\first@\d@t
            \expandafter\ifnum\csname res@count\the\loopcount\endcsname%
                 <\csname seq@start\the\loopcount\endcsname
                    \ifsh@wg@ps
                       \def\first@{\gap@char}\def\third@{7}
                    \else
                       \def\first@{{}} \def\third@{8}
                    \fi
            \else
               \expandafter\ifnum\csname res@count\the\loopcount\endcsname%
                    <\csname seq@len\the\loopcount\endcsname
                       \def\third@{7}\def\first@{\gap@char}
               \else
                    \ifsh@wg@ps
                       \def\first@{\gap@char}\def\third@{7}
                    \else
                       \def\first@{{}} \def\third@{8}
                    \fi
               \fi
            \fi
           \expandafter\ifx\csname lower@seq\the\loopcount\endcsname\y@
                   \expandafter\xdef\csname lower@seq\the\loopcount\endcsname{n}
                   \xdef\first@{;\first@}\fi
           \expandafter\ifx\csname tint@seq\the\loopcount\endcsname\y@
                   \expandafter\xdef\csname tint@seq\the\loopcount\endcsname{n}
                   \xdef\first@{=\first@}\fi
           \expandafter\ifx\csname emph@seq\the\loopcount\endcsname\y@
                   \expandafter\xdef\csname emph@seq\the\loopcount\endcsname{n}
                   \xdef\first@{,\first@}\fi
           \expandafter\xdef\csname seq\the\loopcount\endcsname{%
              \csname seq\the\loopcount\endcsname\third@\first@}
         \else
         \ifx\first@\questi@n
            \expandafter\ifnum\csname res@count\the\loopcount\endcsname%
                 <\csname seq@start\the\loopcount\endcsname
                       \def\first@{{}} \def\third@{8}
            \else
               \expandafter\ifnum\csname res@count\the\loopcount\endcsname%
                    <\csname seq@len\the\loopcount\endcsname
                       \def\third@{7}\def\first@{\st@p@char}
               \else
                       \def\first@{{}} \def\third@{8}
               \fi
            \fi
           \expandafter\ifx\csname lower@seq\the\loopcount\endcsname\y@
                   \expandafter\xdef\csname lower@seq\the\loopcount\endcsname{n}
                   \xdef\first@{;\first@}\fi
           \expandafter\ifx\csname tint@seq\the\loopcount\endcsname\y@
                   \expandafter\xdef\csname tint@seq\the\loopcount\endcsname{n}
                   \xdef\first@{=\first@}\fi
           \expandafter\ifx\csname emph@seq\the\loopcount\endcsname\y@
                   \expandafter\xdef\csname emph@seq\the\loopcount\endcsname{n}
                   \xdef\first@{,\first@}\fi
           \expandafter\xdef\csname seq\the\loopcount\endcsname{%
              \csname seq\the\loopcount\endcsname\third@\first@}
         \else
           \ifnum\divref@=\loopcount\else 
           \xdef\low@up{lower} \ifx\resn@m@tch\low@up \make@lower \else
           \xdef\low@up{upper} \ifx\resn@m@tch\low@up
           \else \xdef\first@{\resn@m@tch} \fi\fi\fi
           \expandafter\ifx\csname lower@seq\the\loopcount\endcsname\y@
                   \expandafter\xdef\csname lower@seq\the\loopcount\endcsname{n}
                   \xdef\first@{;\first@}\fi
           \expandafter\ifx\csname tint@seq\the\loopcount\endcsname\y@
                   \expandafter\xdef\csname tint@seq\the\loopcount\endcsname{n}
                   \xdef\first@{=\first@}\fi
           \expandafter\ifx\csname emph@seq\the\loopcount\endcsname\y@
                   \expandafter\xdef\csname emph@seq\the\loopcount\endcsname{n}
                   \xdef\first@{,\first@}\fi
           \expandafter\xdef\csname seq\the\loopcount\endcsname{%
             \csname seq\the\loopcount\endcsname3\first@}
         \fi
         \fi
         \fi
        \fi
    \ifnum\loopcount<\seq@count \repeat
    \xdef\low@up{lower} \ifx\n@m@tch\low@up \xdef\first@{{ }} \else
    \xdef\low@up{upper} \ifx\n@m@tch\low@up \xdef\first@{{ }}
    \else \xdef\first@{\n@m@tch} \fi\fi
    \expandafter\ifx\csname lower@seq0\endcsname\y@
                \expandafter\xdef\csname lower@seq0\endcsname{n}
                \xdef\first@{;\first@}\fi
    \expandafter\ifx\csname tint@seq0\endcsname\y@
                \expandafter\xdef\csname tint@seq0\endcsname{n}
                \xdef\first@{=\first@}\fi
    \expandafter\ifx\csname emph@seq0\endcsname\y@
                \expandafter\xdef\csname emph@seq0\endcsname{n}
                \xdef\first@{,\first@}\fi
    \ifx\g@p\y@
       \xdef\consensus{\consensus 8{}}
    \else
       \xdef\consensus{\consensus 4\first@}
       \xdef\constopo{\constopo 0}
    \fi
 \fi\fi}

\def\c@nserved{%
 \xdef\tmp@{\csname res\cons@seq\endcsname}
 \loopcount=0
 \loop
    \advance\loopcount by 1
    \xdef\first@{noshade}
    \expandafter\ifx\csname hide@seq\the\loopcount\endcsname\first@
             \xdef\first@{\csname res\the\loopcount\endcsname}
             \ifx\first@\d@t\def\first@{\gap@char}\fi
             \ifx\first@\questi@n\def\first@{\st@p@char}\fi
             \ifx\first@\equ@l\def\first@{\dom@char} \xdef\g@p{y} \fi
             \expandafter\ifx\csname lower@seq\the\loopcount\endcsname\y@
                \expandafter\xdef\csname lower@seq\the\loopcount\endcsname{n}
                \xdef\first@{;\first@}\fi
             \expandafter\ifx\csname tint@seq\the\loopcount\endcsname\y@
                \expandafter\xdef\csname tint@seq\the\loopcount\endcsname{n}
                \xdef\first@{=\first@}\fi
             \expandafter\ifx\csname emph@seq\the\loopcount\endcsname\y@
                \expandafter\xdef\csname emph@seq\the\loopcount\endcsname{n}
                \xdef\first@{,\first@}\fi
             \expandafter\xdef\csname seq\the\loopcount\endcsname{%
                \csname seq\the\loopcount\endcsname%
                3\first@}
    \else
    \xdef\second@{\csname res\the\loopcount\endcsname}
    \ifx\tmp@\second@
       \xdef\third@{1}
    \else
       \xdef\third@{3}
       \ifsimmode
          \xdef\last@{\csname \prefix@ sim\tmp@\endcsname &@}
          \expandafter\get@count\last@
          \innerloopcount=0 \getsim@char
       \fi
    \fi
    \xdef\first@{\csname res\the\loopcount\endcsname}
    \ifcase\third@ \or
    \ifnum\divref@=\loopcount\else 
    \xdef\low@up{lower} \ifx\resm@tch\low@up \make@lower \else
    \xdef\low@up{upper} \ifx\resm@tch\low@up
    \else \xdef\first@{\resm@tch} \fi\fi\fi
    \or
    \ifnum\divref@=\loopcount\else 
    \xdef\low@up{lower} \ifx\ressimm@tch\low@up \make@lower \else
    \xdef\low@up{upper} \ifx\ressimm@tch\low@up
    \else \xdef\first@{\ressimm@tch} \fi\fi\fi
    \else
    \ifnum\divref@=\loopcount\else 
    \xdef\low@up{lower} \ifx\resn@m@tch\low@up \make@lower \else
    \xdef\low@up{upper} \ifx\resn@m@tch\low@up
    \else \xdef\first@{\resn@m@tch} \fi\fi\fi
    \fi
    \ifx\first@\equ@l \def\first@{\dom@char} \def\third@{9}  \xdef\g@p{y} \fi
    \ifx\first@\d@t
            \expandafter\ifnum\csname res@count\the\loopcount\endcsname%
                 <\csname seq@start\the\loopcount\endcsname
		    \ifsh@wg@ps
		       \def\first@{\gap@char}\def\third@{7}
		    \else
                       \def\first@{{}} \def\third@{8}
		    \fi
            \else
               \expandafter\ifnum\csname res@count\the\loopcount\endcsname%
                    <\csname seq@len\the\loopcount\endcsname
		       \def\third@{7}\def\first@{\gap@char}
               \else
                    \ifsh@wg@ps
                       \def\first@{\gap@char}\def\third@{7}
                    \else
                       \def\first@{{}} \def\third@{8}
                    \fi
               \fi
            \fi
    \fi
    \ifx\first@\questi@n
            \expandafter\ifnum\csname res@count\the\loopcount\endcsname%
                 <\csname seq@start\the\loopcount\endcsname
                      \def\first@{{}} \def\third@{8}
            \else
               \expandafter\ifnum\csname res@count\the\loopcount\endcsname%
                    <\csname seq@len\the\loopcount\endcsname
		               \def\third@{7}\def\first@{\st@p@char}
               \else
                       \def\first@{{}} \def\third@{8}
               \fi
            \fi
    \fi
    \expandafter\ifx\csname lower@seq\the\loopcount\endcsname\y@
                \expandafter\xdef\csname lower@seq\the\loopcount\endcsname{n}
                \xdef\first@{;\first@}\fi
    \expandafter\ifx\csname tint@seq\the\loopcount\endcsname\y@
                \expandafter\xdef\csname tint@seq\the\loopcount\endcsname{n}
                \xdef\first@{=\first@}\fi
    \expandafter\ifx\csname emph@seq\the\loopcount\endcsname\y@
                \expandafter\xdef\csname emph@seq\the\loopcount\endcsname{n}
                \xdef\first@{,\first@}\fi
    \expandafter\xdef\csname seq\the\loopcount\endcsname{%
       \csname seq\the\loopcount\endcsname\third@\first@}
    \fi
 \ifnum\loopcount<\seq@count \repeat
 \xdef\first@{\csname res\cons@seq\endcsname}
 \xdef\second@{lower} \ifx\m@tch\second@ \make@lower \else
 \xdef\second@{upper} \ifx\m@tch\second@
 \else \xdef\first@{\m@tch} \fi\fi
 \expandafter\ifx\csname lower@seq0\endcsname\y@
                \expandafter\xdef\csname lower@seq0\endcsname{n}
                \xdef\first@{;\first@}\fi
 \expandafter\ifx\csname tint@seq0\endcsname\y@
                \expandafter\xdef\csname tint@seq0\endcsname{n}
                \xdef\first@{=\first@}\fi
 \expandafter\ifx\csname emph@seq0\endcsname\y@
                \expandafter\xdef\csname emph@seq0\endcsname{n}
                \xdef\first@{,\first@}\fi
 \ifx\g@p\y@
    \xdef\consensus{\consensus 8{}}
 \else
    \xdef\consensus{\consensus 5\first@}
    \xdef\constopo{\constopo 2}
 \fi}
 
\def\allm@tch{%
 \xdef\tmp@{\csname res\cons@seq\endcsname}
 \loopcount=0
 \loop
    \advance\loopcount by 1
    \xdef\first@{noshade}
    \expandafter\ifx\csname hide@seq\the\loopcount\endcsname\first@
             \xdef\first@{\csname res\the\loopcount\endcsname}
             \ifx\first@\d@t\def\first@{\gap@char}\fi
             \ifx\first@\questi@n\def\first@{\st@p@char}\fi
             \ifx\first@\equ@l\def\first@{\dom@char} \xdef\g@p{y} \fi
             \expandafter\ifx\csname lower@seq\the\loopcount\endcsname\y@
                \expandafter\xdef\csname lower@seq\the\loopcount\endcsname{n}
                \xdef\first@{;\first@}\fi
             \expandafter\ifx\csname tint@seq\the\loopcount\endcsname\y@
                \expandafter\xdef\csname tint@seq\the\loopcount\endcsname{n}
                \xdef\first@{=\first@}\fi
             \expandafter\ifx\csname emph@seq\the\loopcount\endcsname\y@
                \expandafter\xdef\csname emph@seq\the\loopcount\endcsname{n}
                \xdef\first@{,\first@}\fi
             \expandafter\xdef\csname seq\the\loopcount\endcsname{%
                \csname seq\the\loopcount\endcsname%
                3\first@}
    \else
    \xdef\second@{\csname res\the\loopcount\endcsname}
    \ifx\tmp@\second@
       \ifall@shade \xdef\third@{0} \else \xdef\third@{1} \fi
    \else
       \xdef\third@{3}
       \ifsimmode
          \xdef\last@{\csname \prefix@ sim\tmp@\endcsname &@}
          \expandafter\get@count\last@
          \innerloopcount=0 \getsim@char
       \fi
    \fi
    \xdef\first@{\csname res\the\loopcount\endcsname}
    \ifcase\third@ 
    \ifnum\divref@=\loopcount\else 
    \xdef\low@up{lower} \ifx\res@llm@tch\low@up \make@lower \else
    \xdef\low@up{upper} \ifx\res@llm@tch\low@up
    \else \xdef\first@{\res@llm@tch} \fi\fi\fi
    \or
    \ifnum\divref@=\loopcount\else 
    \xdef\low@up{lower} \ifx\res@llm@tch\low@up \make@lower \else
    \xdef\low@up{upper} \ifx\res@llm@tch\low@up
    \else \xdef\first@{\res@llm@tch} \fi\fi\fi
    \or
    \ifnum\divref@=\loopcount\else 
    \xdef\low@up{lower} \ifx\ressimm@tch\low@up \make@lower \else
    \xdef\low@up{upper} \ifx\ressimm@tch\low@up
    \else \xdef\first@{\ressimm@tch} \fi\fi\fi
    \else
    \ifnum\divref@=\loopcount\else 
    \xdef\low@up{lower} \ifx\resn@m@tch\low@up \make@lower \else
    \xdef\low@up{upper} \ifx\resn@m@tch\low@up
    \else \xdef\first@{\resn@m@tch} \fi\fi\fi
    \fi
    \ifx\first@\equ@l \def\first@{\dom@char} \def\third@{9}  \xdef\g@p{y} \fi
    \ifx\first@\d@t
            \expandafter\ifnum\csname res@count\the\loopcount\endcsname%
                 <\csname seq@start\the\loopcount\endcsname
		    \ifsh@wg@ps
		       \def\first@{\gap@char}\def\third@{7}
		    \else
                       \def\first@{{}} \def\third@{8}
		    \fi
            \else
               \expandafter\ifnum\csname res@count\the\loopcount\endcsname%
                    <\csname seq@len\the\loopcount\endcsname
		       \def\third@{7}\def\first@{\gap@char}
               \else
                    \ifsh@wg@ps
                       \def\first@{\gap@char}\def\third@{7}
                    \else
                       \def\first@{{}} \def\third@{8}
                    \fi
               \fi
            \fi
    \fi
    \ifx\first@\questi@n
            \expandafter\ifnum\csname res@count\the\loopcount\endcsname%
                 <\csname seq@start\the\loopcount\endcsname
                       \def\first@{{}} \def\third@{8}
            \else
               \expandafter\ifnum\csname res@count\the\loopcount\endcsname%
                    <\csname seq@len\the\loopcount\endcsname
		                \def\third@{7}\def\first@{\st@p@char}
               \else
                       \def\first@{{}} \def\third@{8}
               \fi
            \fi
    \fi
    \expandafter\ifx\csname lower@seq\the\loopcount\endcsname\y@
                \expandafter\xdef\csname lower@seq\the\loopcount\endcsname{n}
                \xdef\first@{;\first@}\fi
    \expandafter\ifx\csname tint@seq\the\loopcount\endcsname\y@
                \expandafter\xdef\csname tint@seq\the\loopcount\endcsname{n}
                \xdef\first@{=\first@}\fi
    \expandafter\ifx\csname emph@seq\the\loopcount\endcsname\y@
                \expandafter\xdef\csname emph@seq\the\loopcount\endcsname{n}
                \xdef\first@{,\first@}\fi
    \expandafter\xdef\csname seq\the\loopcount\endcsname{%
       \csname seq\the\loopcount\endcsname\third@\first@}
    \fi
 \ifnum\loopcount<\seq@count \repeat
 \xdef\first@{\csname res\cons@seq\endcsname}
 \xdef\second@{lower} \ifx\@llm@tch\second@ \make@lower \else
 \xdef\second@{upper} \ifx\@llm@tch\second@
 \else \xdef\first@{\@llm@tch} \fi\fi
 \expandafter\ifx\csname lower@seq0\endcsname\y@
                \expandafter\xdef\csname lower@seq0\endcsname{n}
                \xdef\first@{;\first@}\fi
 \expandafter\ifx\csname tint@seq0\endcsname\y@
                \expandafter\xdef\csname tint@seq0\endcsname{n}
                \xdef\first@{=\first@}\fi
 \expandafter\ifx\csname emph@seq0\endcsname\y@
                \expandafter\xdef\csname emph@seq0\endcsname{n}
                \xdef\first@{,\first@}\fi
 \ifx\g@p\y@
    \xdef\consensus{\consensus 8{}}
 \else
    \xdef\consensus{\consensus 6\first@}
    \xdef\constopo{\constopo 3}
 \fi}

\def\functi@nal{%
     \ifnum\cons@num>0
       \xdef\first@{\csname res\cons@num\endcsname}
     \else
       \xdef\first@{\csname res\cons@seq\endcsname}
     \fi
     \xdef\second@{\csname funcgrp\first@\endcsname}
     \loopcount=0
     \loop
        \advance\loopcount by 1
        \xdef\third@{noshade}
        \expandafter\ifx\csname hide@seq\the\loopcount\endcsname\third@
             \xdef\first@{\csname res\the\loopcount\endcsname}
             \expandafter\ifx\csname lower@seq\the\loopcount\endcsname\y@
                \expandafter\xdef\csname lower@seq\the\loopcount\endcsname{n}
                \xdef\first@{;\first@}\fi
             \expandafter\ifx\csname tint@seq\the\loopcount\endcsname\y@
                \expandafter\xdef\csname tint@seq\the\loopcount\endcsname{n}
                \xdef\first@{=\first@}\fi
             \expandafter\ifx\csname emph@seq\the\loopcount\endcsname\y@
                \expandafter\xdef\csname emph@seq\the\loopcount\endcsname{n}
                \xdef\first@{,\first@}\fi
             \expandafter\xdef\csname seq\the\loopcount\endcsname{%
                \csname seq\the\loopcount\endcsname%
                0\first@}
        \else
        \xdef\first@{\csname res\the\loopcount\endcsname}
         \ifx\first@\equ@l
            \def\first@{\dom@char} \def\third@{!} \xdef\g@p{y}
         \else
         \ifx\first@\d@t
            \expandafter\ifnum\csname res@count\the\loopcount\endcsname%
                 <\csname seq@start\the\loopcount\endcsname
                    \ifsh@wg@ps
                       \def\first@{\gap@char}\def\third@{*}
                    \else
                       \def\first@{{}} \def\third@{/}
                    \fi
            \else
               \expandafter\ifnum\csname res@count\the\loopcount\endcsname%
                    <\csname seq@len\the\loopcount\endcsname
                  \def\first@{\gap@char}\def\third@{*}
               \else
                    \ifsh@wg@ps
                       \def\first@{\gap@char}\def\third@{*}
                    \else
                       \def\first@{{}} \def\third@{/}
                    \fi
               \fi
            \fi
         \else
         \ifx\first@\questi@n
            \expandafter\ifnum\csname res@count\the\loopcount\endcsname%
                 <\csname seq@start\the\loopcount\endcsname
                    \def\first@{{}} \def\third@{/}
            \else
               \expandafter\ifnum\csname res@count\the\loopcount\endcsname%
                    <\csname seq@len\the\loopcount\endcsname
                    \def\first@{\st@p@char}\def\third@{*}
               \else
                    \def\first@{{}} \def\third@{/}
               \fi
            \fi
         \else
           \expandafter\ifnum\csname funcgrp\first@\endcsname=\second@
              \xdef\low@up{lower}
           \expandafter\ifx\csname funcm@tch\second@\endcsname\low@up
                       \make@lower \fi
              \xdef\third@{\second@}
           \else \xdef\third@{0}
              \xdef\low@up{lower} \ifx\resn@m@tch\low@up \make@lower \else
              \xdef\low@up{upper} \ifx\resn@m@tch\low@up
              \else \xdef\first@{\resn@m@tch} \fi\fi
         \fi
         \fi
         \fi
        \fi
        \expandafter\ifx\csname lower@seq\the\loopcount\endcsname\y@
                \expandafter\xdef\csname lower@seq\the\loopcount\endcsname{n}
                \xdef\first@{;\first@}\fi
        \expandafter\ifx\csname tint@seq\the\loopcount\endcsname\y@
                \expandafter\xdef\csname tint@seq\the\loopcount\endcsname{n}
                \xdef\first@{=\first@}\fi
        \expandafter\ifx\csname emph@seq\the\loopcount\endcsname\y@
                \expandafter\xdef\csname emph@seq\the\loopcount\endcsname{n}
                \xdef\first@{,\first@}\fi
        \expandafter\xdef\csname seq\the\loopcount\endcsname{%
          \csname seq\the\loopcount\endcsname\third@\first@}
        \fi
     \ifnum\loopcount<\seq@count \repeat}

\def\all@funcshade{%
   \loopcount=0
   \loop
     \advance\loopcount by 1
     \xdef\first@{\csname res\the\loopcount\endcsname}
     \xdef\second@{\csname funcgrp\first@\endcsname}
     \ifnum\second@<0 \xdef\second@{0} \fi
     \ifx\first@\equ@l \def\first@{\dom@char} \def\second@{!}  \xdef\g@p{y} \fi
     \ifx\first@\d@t
            \expandafter\ifnum\csname res@count\the\loopcount\endcsname%
                 <\csname seq@start\the\loopcount\endcsname
                    \ifsh@wg@ps
                       \def\first@{\gap@char}\def\second@{*}
                    \else
                       \def\first@{{}} \def\second@{/}
                    \fi
            \else
               \expandafter\ifnum\csname res@count\the\loopcount\endcsname%
                    <\csname seq@len\the\loopcount\endcsname
                  \def\first@{\gap@char}\def\second@{*}
               \else
                    \ifsh@wg@ps
                       \def\first@{\gap@char}\def\second@{*}
                    \else
                       \def\first@{{}} \def\second@{/}
                    \fi
               \fi
            \fi
     \fi
     \ifx\first@\questi@n
            \expandafter\ifnum\csname res@count\the\loopcount\endcsname%
                 <\csname seq@start\the\loopcount\endcsname
                    \def\first@{{}} \def\second@{/}
            \else
               \expandafter\ifnum\csname res@count\the\loopcount\endcsname%
                    <\csname seq@len\the\loopcount\endcsname
                    \def\first@{\st@p@char}\def\second@{*}
               \else
                    \def\first@{{}} \def\second@{/}
               \fi
            \fi
     \fi
     \xdef\low@up{lower}
     \expandafter\ifx\csname funcm@tch\second@\endcsname\low@up
                 \make@lower \fi
     \xdef\third@{noshade}
     \expandafter\ifx\csname hide@seq\the\loopcount\endcsname\third@
             \expandafter\ifx\csname lower@seq\the\loopcount\endcsname\y@
                \expandafter\xdef\csname lower@seq\the\loopcount\endcsname{n}
                \xdef\first@{;\first@}\fi
             \expandafter\ifx\csname tint@seq\the\loopcount\endcsname\y@
                \expandafter\xdef\csname tint@seq\the\loopcount\endcsname{n}
                \xdef\first@{=\first@}\fi
             \expandafter\ifx\csname emph@seq\the\loopcount\endcsname\y@
                \expandafter\xdef\csname emph@seq\the\loopcount\endcsname{n}
                \xdef\first@{,\first@}\fi
             \expandafter\xdef\csname seq\the\loopcount\endcsname{%
                \csname seq\the\loopcount\endcsname%
                0\first@}
     \else
             \expandafter\ifx\csname lower@seq\the\loopcount\endcsname\y@
                \expandafter\xdef\csname lower@seq\the\loopcount\endcsname{n}
                \xdef\first@{;\first@}\fi
             \expandafter\ifx\csname tint@seq\the\loopcount\endcsname\y@
                \expandafter\xdef\csname tint@seq\the\loopcount\endcsname{n}
                \xdef\first@{=\first@}\fi
             \expandafter\ifx\csname emph@seq\the\loopcount\endcsname\y@
                \expandafter\xdef\csname emph@seq\the\loopcount\endcsname{n}
                \xdef\first@{,\first@}\fi
             \expandafter\xdef\csname seq\the\loopcount\endcsname{%
                \csname seq\the\loopcount\endcsname%
                \second@\first@}
     \fi
   \ifnum\loopcount<\seq@count \repeat}

\def\T@coffee@shade{%
   \loopcount=0
   \loop
     \advance\loopcount by 1
     \xdef\first@{\csname res\the\loopcount\endcsname}
     \xdef\second@{\csname TC@num\the\loopcount\endcsname}
     \ifx\first@\equ@l \def\first@{\dom@char} \def\second@{!}  \xdef\g@p{y} \fi
     \ifx\first@\d@t
            \expandafter\ifnum\csname res@count\the\loopcount\endcsname%
                 <\csname seq@start\the\loopcount\endcsname
                    \ifsh@wg@ps
                       \def\first@{\gap@char}\def\second@{*}
                    \else
                       \def\first@{{}} \def\second@{/}
                    \fi
            \else
               \expandafter\ifnum\csname res@count\the\loopcount\endcsname%
                    <\csname seq@len\the\loopcount\endcsname
                  \def\first@{\gap@char}\def\second@{*}
               \else
                    \ifsh@wg@ps
                       \def\first@{\gap@char}\def\second@{*}
                    \else
                       \def\first@{{}} \def\second@{/}
                    \fi
               \fi
            \fi
     \fi
     \ifx\first@\questi@n
            \expandafter\ifnum\csname res@count\the\loopcount\endcsname%
                 <\csname seq@start\the\loopcount\endcsname
                       \def\first@{{}} \def\second@{/}
            \else
               \expandafter\ifnum\csname res@count\the\loopcount\endcsname%
                    <\csname seq@len\the\loopcount\endcsname
                       \def\first@{\st@p@char}\def\second@{*}
               \else
                       \def\first@{{}} \def\second@{/}
               \fi
            \fi
     \fi
     \xdef\third@{noshade}
     \expandafter\ifx\csname hide@seq\the\loopcount\endcsname\third@
             \expandafter\ifx\csname lower@seq\the\loopcount\endcsname\y@
                \expandafter\xdef\csname lower@seq\the\loopcount\endcsname{n}
                \xdef\first@{;\first@}\fi
             \expandafter\ifx\csname tint@seq\the\loopcount\endcsname\y@
                \expandafter\xdef\csname tint@seq\the\loopcount\endcsname{n}
                \xdef\first@{=\first@}\fi
             \expandafter\ifx\csname emph@seq\the\loopcount\endcsname\y@
                \expandafter\xdef\csname emph@seq\the\loopcount\endcsname{n}
                \xdef\first@{,\first@}\fi
             \expandafter\xdef\csname seq\the\loopcount\endcsname{%
                \csname seq\the\loopcount\endcsname%
                *\first@}
     \else
             \expandafter\ifx\csname lower@seq\the\loopcount\endcsname\y@
                \expandafter\xdef\csname lower@seq\the\loopcount\endcsname{n}
                \xdef\first@{;\first@}\fi
             \expandafter\ifx\csname tint@seq\the\loopcount\endcsname\y@
                \expandafter\xdef\csname tint@seq\the\loopcount\endcsname{n}
                \xdef\first@{=\first@}\fi
             \expandafter\ifx\csname emph@seq\the\loopcount\endcsname\y@
                \expandafter\xdef\csname emph@seq\the\loopcount\endcsname{n}
                \xdef\first@{,\first@}\fi
             \expandafter\xdef\csname seq\the\loopcount\endcsname{%
                \csname seq\the\loopcount\endcsname%
                \second@\first@}
     \fi
   \ifnum\loopcount<\seq@count \repeat
   \xdef\first@{\csname res0\endcsname}
   \expandafter\ifx\csname lower@seq0\endcsname\y@
                \expandafter\xdef\csname lower@seq0\endcsname{n}
                \xdef\first@{;\first@}\fi
   \expandafter\ifx\csname tint@seq0\endcsname\y@
                \expandafter\xdef\csname tint@seq0\endcsname{n}
                \xdef\first@{=\first@}\fi
   \expandafter\ifx\csname emph@seq0\endcsname\y@
                \expandafter\xdef\csname emph@seq0\endcsname{n}
                \xdef\first@{,\first@}\fi
   \xdef\second@{\csname TC@num0\endcsname}
   \ifx\g@p\y@
      \xdef\consensus{\consensus 8{}}
   \else
      \xdef\consensus{\consensus\second@\n@m@tch}
      \xdef\constopo{\constopo\second@}
   \fi
   }

\def\getregion@fromstack@first{%
      \expandafter\getregion@fromstack{\the\loopcount}
      \expandafter\ifx\csname start\the\loopcount\endcsname\ampers@nd
      \else
      \expandafter\ifnum\csname stop\the\loopcount\endcsname<\first@@
         \getregion@fromstack@first
      \else
         \expandafter\ifx\csname all\the\loopcount\endcsname\y@
            \innerloopcount=\csname style\the\loopcount\endcsname
         \fi
         \expandafter\xdef\csname shade@style\the\loopcount\endcsname{%
                          \csname style\the\loopcount\endcsname}
      \fi
      \fi
}

\def\calc@regshade{%
      \loopcount=-1 \innerloopcount=0
      \loop
         \advance\loopcount by 1
         \expandafter\xdef\csname shade@style\the\loopcount\endcsname{y}
         \ifnum\loopcount=0 \xdef\first@@{\the\cons@count}
         \else \xdef\first@@{\csname res@count\the\loopcount\endcsname} \fi
         \expandafter\ifx\csname start\the\loopcount\endcsname\ampers@nd
         \else
         \expandafter\ifnum\csname start\the\loopcount\endcsname>\first@@
         \else
            \expandafter\ifnum\csname stop\the\loopcount\endcsname<\first@@
               \getregion@fromstack@first
            \else
               \expandafter\ifx\csname all\the\loopcount\endcsname\y@
                  \innerloopcount=\csname style\the\loopcount\endcsname
               \fi
               \expandafter\xdef\csname shade@style\the\loopcount\endcsname{%
                                \csname style\the\loopcount\endcsname}
               \expandafter\ifnum\csname stop\the\loopcount\endcsname=\first@@
                  \expandafter\getregion@fromstack{\the\loopcount}
               \fi
            \fi
         \fi\fi
      \ifnum\loopcount<\seq@count \repeat
      \loopcount=0
      \expandafter\ifx\csname shade@style\the\loopcount\endcsname\y@
      \else
         \xdef\consensus{\consensus&\csname shade@style\the\loopcount\endcsname)}
      \fi
      \loop
         \advance\loopcount by 1
         \expandafter\ifx\csname shade@style\the\loopcount\endcsname\y@
            \ifnum\innerloopcount>0
               \expandafter\xdef\csname seq\the\loopcount\endcsname{%
                                \csname seq\the\loopcount\endcsname%
                                &\the\innerloopcount)}
            \fi
         \else
             \expandafter\xdef\csname seq\the\loopcount\endcsname{%
                              \csname seq\the\loopcount\endcsname%
                              &\csname shade@style\the\loopcount\endcsname)}
         \fi
      \ifnum\loopcount<\seq@count \repeat
}

\def\getregion@fromshadingstack@first{%
      \expandafter\getregion@fromshadingstack{\the\loopcount}
      \expandafter\ifx\csname shadingstart\the\loopcount\endcsname\ampers@nd
      \else
      \expandafter\ifnum\csname shadingstop\the\loopcount\endcsname<\first@@
         \getregion@fromshadingstack@first
      \fi
      \fi
}

\def\calc@shading{%
      \xdef\shading@style{&}
      \loopcount=-1
      \loop
         \advance\loopcount by 1
         \ifnum\loopcount=0 \xdef\first@@{\the\cons@count}
         \else \xdef\first@@{\csname res@count\the\loopcount\endcsname} \fi
         \expandafter\ifx\csname shadingstart\the\loopcount\endcsname\ampers@nd
         \else
         \expandafter\ifnum\csname shadingstart\the\loopcount\endcsname>\first@@
         \else
            \expandafter\ifnum\csname shadingstop\the\loopcount\endcsname<\first@@
               \getregion@fromshadingstack@first
            \else
               \xdef\shading@style{\csname shadingstyle\the\loopcount\endcsname}
               \expandafter\ifnum\csname shadingstop\the\loopcount\endcsname=\first@@
                  \expandafter\getregion@fromshadingstack{\the\loopcount}
               \fi
            \fi
         \fi\fi
      \ifnum\loopcount<\seq@count \repeat
      \loopcount=0
      \loop
         \advance\loopcount by 1
         \ifx\shading@style\ampers@nd
         \else
             \expandafter\xdef\csname seq\the\loopcount\endcsname{%
                              \csname seq\the\loopcount\endcsname€\shading@style)}
         \fi
      \ifnum\loopcount<\seq@count \repeat
}

\def\getregion@fromemphstack@first{%
      \expandafter\getregion@fromemphstack{\the\loopcount}
      \expandafter\ifx\csname emphstart\the\loopcount\endcsname\ampers@nd
      \else
      \expandafter\ifnum\csname emphstop\the\loopcount\endcsname<\first@@
         \getregion@fromemphstack@first
      \else
         \expandafter\ifx\csname emphall\the\loopcount\endcsname\y@
            \innerloopcount=1
         \fi
         \expandafter\xdef\csname emph@seq\the\loopcount\endcsname{y}
      \fi
      \fi
}

\def\calc@regemph{%
      \loopcount=-1 \innerloopcount=0
      \loop
         \advance\loopcount by 1
         \ifnum\loopcount=0 \xdef\first@@{\the\cons@count}
         \else \xdef\first@@{\csname res@count\the\loopcount\endcsname} \fi
         \expandafter\ifx\csname emphstart\the\loopcount\endcsname\ampers@nd
         \else
         \expandafter\ifnum\csname emphstart\the\loopcount\endcsname>\first@@
         \else
            \expandafter\ifnum\csname emphstop\the\loopcount\endcsname<\first@@
              \getregion@fromemphstack@first
            \else
               \expandafter\ifx\csname emphall\the\loopcount\endcsname\y@
                  \innerloopcount=1
               \fi
               \expandafter\xdef\csname emph@seq\the\loopcount\endcsname{y}
               \expandafter\ifnum\csname emphstop\the\loopcount\endcsname=\first@@
                  \expandafter\getregion@fromemphstack{\the\loopcount}
               \fi
            \fi
         \fi\fi
      \ifnum\loopcount<\seq@count \repeat
      \loopcount=0
      \expandafter\ifx\csname emph@seq\the\loopcount\endcsname\y@
         \expandafter\xdef\csname emph@seq\the\loopcount\endcsname{y}
      \fi
      \ifnum\innerloopcount>0
         \loop
            \advance\loopcount by 1
            \expandafter\xdef\csname emph@seq\the\loopcount\endcsname{y}
         \ifnum\loopcount<\seq@count \repeat
      \fi
}

\def\getregion@fromlowerstack@first{%
      \expandafter\getregion@fromlowerstack{\the\loopcount}
      \expandafter\ifx\csname lowerstart\the\loopcount\endcsname\ampers@nd
      \else
      \expandafter\ifnum\csname lowerstop\the\loopcount\endcsname<\first@@
         \getregion@fromlowerstack@first
      \else
         \expandafter\ifx\csname lowerall\the\loopcount\endcsname\y@
            \innerloopcount=1
         \fi
         \expandafter\xdef\csname lower@seq\the\loopcount\endcsname{y}
      \fi
      \fi
}

\def\calc@reglower{%
      \loopcount=-1 \innerloopcount=0
      \loop
         \advance\loopcount by 1
         \ifnum\loopcount=0 \xdef\first@@{\the\cons@count}
         \else \xdef\first@@{\csname res@count\the\loopcount\endcsname} \fi
         \expandafter\ifx\csname lowerstart\the\loopcount\endcsname\ampers@nd
         \else
         \expandafter\ifnum\csname lowerstart\the\loopcount\endcsname>\first@@
         \else
            \expandafter\ifnum\csname lowerstop\the\loopcount\endcsname<\first@@
              \getregion@fromlowerstack@first
            \else
               \expandafter\ifx\csname lowerall\the\loopcount\endcsname\y@
                  \innerloopcount=1
               \fi
               \expandafter\xdef\csname lower@seq\the\loopcount\endcsname{y}
               \expandafter\ifnum\csname lowerstop\the\loopcount\endcsname=\first@@
                  \expandafter\getregion@fromlowerstack{\the\loopcount}
               \fi
            \fi
         \fi\fi
      \ifnum\loopcount<\seq@count \repeat
      \loopcount=0
      \expandafter\ifx\csname lower@seq\the\loopcount\endcsname\y@
         \expandafter\xdef\csname lower@seq\the\loopcount\endcsname{y}
      \fi
      \ifnum\innerloopcount>0
         \loop
            \advance\loopcount by 1
            \expandafter\xdef\csname lower@seq\the\loopcount\endcsname{y}
         \ifnum\loopcount<\seq@count \repeat
      \fi
}

\def\getregion@fromtintstack@first{%
      \expandafter\getregion@fromtintstack{\the\loopcount}
      \expandafter\ifx\csname tintstart\the\loopcount\endcsname\ampers@nd
      \else
      \expandafter\ifnum\csname tintstop\the\loopcount\endcsname<\first@@
         \getregion@fromtintstack@first
      \else
         \expandafter\ifx\csname tintall\the\loopcount\endcsname\y@
            \innerloopcount=1
         \fi
         \expandafter\xdef\csname tint@seq\the\loopcount\endcsname{y}
      \fi
      \fi
}

\def\calc@regtint{%
      \loopcount=-1 \innerloopcount=0
      \loop
         \advance\loopcount by 1
         \ifnum\loopcount=0 \xdef\first@@{\the\cons@count}
         \else \xdef\first@@{\csname res@count\the\loopcount\endcsname} \fi
         \expandafter\ifx\csname tintstart\the\loopcount\endcsname\ampers@nd
         \else
         \expandafter\ifnum\csname tintstart\the\loopcount\endcsname>\first@@
         \else
            \expandafter\ifnum\csname tintstop\the\loopcount\endcsname<\first@@
              \getregion@fromtintstack@first
            \else
               \expandafter\ifx\csname tintall\the\loopcount\endcsname\y@
                  \innerloopcount=1
               \fi
               \expandafter\xdef\csname tint@seq\the\loopcount\endcsname{y}
               \expandafter\ifnum\csname tintstop\the\loopcount\endcsname=\first@@
                  \expandafter\getregion@fromtintstack{\the\loopcount}
               \fi
            \fi
         \fi\fi
      \ifnum\loopcount<\seq@count \repeat
      \loopcount=0
      \expandafter\ifx\csname tint@seq\the\loopcount\endcsname\y@
         \expandafter\xdef\csname tint@seq\the\loopcount\endcsname{y}
      \fi
      \ifnum\innerloopcount>0
         \loop
            \advance\loopcount by 1
            \expandafter\xdef\csname tint@seq\the\loopcount\endcsname{y}
         \ifnum\loopcount<\seq@count \repeat
      \fi
}

\def\getregion@fromframestack@first{%
      \expandafter\getregion@fromframestack{\the\loopcount}
      \expandafter\ifx\csname framestart\the\loopcount\endcsname\ampers@nd
      \else
      \expandafter\ifnum\csname framestop\the\loopcount\endcsname<\first@@
         \getregion@fromframestack@first
      \else
         \ifnum\frame@on=0
            \xdef\frame@on{1}
            \xdef\frame@{1}
            \expandafter\xdef\csname fr@style\the\loopcount\endcsname{%
                             \csname framestyle\the\loopcount\endcsname}
            \innerloopcount=\pos@count
            \advance\innerloopcount by -\frame@pos
            \xdef\styleframe{\styleframe&\the\innerloopcount;&;}
            \xdef\frame@pos{\the\pos@count}
         \fi
         \expandafter\ifnum\csname framestop\the\loopcount\endcsname=\first@@
            \expandafter\getregion@fromframestack{\the\loopcount}
            \ifnum\frame@on=1
               \xdef\frame@on{0}
               \innerloopcount=\pos@count \advance\innerloopcount by 1
               \advance\innerloopcount by -\frame@pos
               \xdef\styleframe{\styleframe&\the\innerloopcount;%
                                \csname fr@style\the\loopcount\endcsname;}
               \innerloopcount=\pos@count \advance\innerloopcount by 1
               \xdef\frame@pos{\the\innerloopcount}
            \fi
         \fi
         \ifnum\pos@count=\res@perline
            \ifnum\frame@on=1
               \innerloopcount=\pos@count
               \advance\innerloopcount by 1
               \advance\innerloopcount by -\frame@pos
               \xdef\styleframe{\styleframe&\the\innerloopcount;%
                                \csname fr@style\the\loopcount\endcsname;}
               \innerloopcount=\pos@count \advance\innerloopcount by 1
               \xdef\frame@pos{\the\innerloopcount}
            \fi
         \fi
         \expandafter\ifnum\csname res@count\the\loopcount\endcsname=\end@num\relax
            \ifnum\frame@on=1
               \innerloopcount=\pos@count
               \advance\innerloopcount by 1
               \advance\innerloopcount by -\frame@pos
               \xdef\styleframe{\styleframe&\the\innerloopcount;%
                                \csname fr@style\the\loopcount\endcsname;}
               \innerloopcount=\pos@count \advance\innerloopcount by 1
               \xdef\frame@pos{\the\innerloopcount}
            \fi
         \fi
      \fi
      \fi
}

\def\calc@frame{%
      \loopcount=-1 \innerloopcount=0
      \loop
         \advance\loopcount by 1
         \ifnum\loopcount=0 \xdef\first@@{\the\cons@count}
         \else \xdef\first@@{\csname res@count\the\loopcount\endcsname} \fi
         \expandafter\ifx\csname framestart\the\loopcount\endcsname\ampers@nd
         \else
         \expandafter\ifnum\csname framestart\the\loopcount\endcsname>\first@@
         \else
            \expandafter\ifnum\csname framestop\the\loopcount\endcsname<\first@@
              \getregion@fromframestack@first
            \else
               \ifnum\frame@on=0
                  \xdef\frame@on{1}
                  \xdef\frame@{1}
                  \expandafter\xdef\csname fr@style\the\loopcount\endcsname{%
                                   \csname framestyle\the\loopcount\endcsname}
                  \innerloopcount=\pos@count 
                  \advance\innerloopcount by -\frame@pos
                  \xdef\styleframe{\styleframe&\the\innerloopcount;&;}
                  \xdef\frame@pos{\the\pos@count}
               \fi
               \expandafter\ifnum\csname framestop\the\loopcount\endcsname=\first@@
                  \expandafter\getregion@fromframestack{\the\loopcount}
                  \ifnum\frame@on=1
                     \xdef\frame@on{0}
                     \innerloopcount=\pos@count \advance\innerloopcount by 1
                     \advance\innerloopcount by -\frame@pos
                     \xdef\styleframe{\styleframe&\the\innerloopcount;%
                                      \csname fr@style\the\loopcount\endcsname;}
                     \innerloopcount=\pos@count \advance\innerloopcount by 1
                     \xdef\frame@pos{\the\innerloopcount}
                  \fi
               \fi
               \ifnum\pos@count=\res@perline
                  \ifnum\frame@on=1
                     \innerloopcount=\pos@count 
                     \advance\innerloopcount by 1
                     \advance\innerloopcount by -\frame@pos
                     \xdef\styleframe{\styleframe&-\the\innerloopcount;%
                                      \csname fr@style\the\loopcount\endcsname;}
                     \innerloopcount=\pos@count \advance\innerloopcount by 1
                     \xdef\frame@pos{\the\innerloopcount}
                  \fi
               \fi
               \expandafter\ifnum\csname res@count\the\loopcount\endcsname=\end@num\relax
                  \ifnum\frame@on=1
                     \innerloopcount=\pos@count
                     \advance\innerloopcount by 1
                     \advance\innerloopcount by -\frame@pos
                     \xdef\styleframe{\styleframe&\the\innerloopcount;%
                                      \csname fr@style\the\loopcount\endcsname;}
                     \innerloopcount=\pos@count \advance\innerloopcount by 1
                     \xdef\frame@pos{\the\innerloopcount}
                  \fi
               \fi
            \fi
         \fi\fi
      \ifnum\loopcount<\seq@count \repeat
%      \advance\pos@count by 1 
}

\def\get@nextres#1#2:{%
    \xdef\first@{#1}
    \xdef\temp@{#2:}
    \ifx\first@\gap@char \expandafter\get@nextres\temp@
    \else
       \if\first@ @
          \xdef\temp@{}
       \else
          \xdef\temp@{+\first@}
          \expandafter\xdef\csname last@res\bottop@\endcsname{}
       \fi
    \fi}
\def\get@@nextres#1#2:{%
    \xdef\first@{#1}
    \xdef\temp@@{#2:}
    \ifx\first@\gap@char \expandafter\get@@nextres\temp@@
    \else
       \if\first@ @
          \xdef\temp@@{}
       \else
          \xdef\temp@@{+\first@}
          \expandafter\xdef\csname last@res\bottop@\endcsname{}
       \fi
    \fi}

\def\char@get#1#2@{\xdef\first@{#1} \xdef\tr@nsl@ted{#2@}} 
\def\trans@now#1#2@{%
   \xdef\first@{#1}
   \xdef\tr@nsl@ted{#2@}
   \ifx\first@\ampers@nd
   \else
      \expandafter\check@char\first@
      \ifletter
         \xdef\triplet@{\triplet@\first@}
         \advance\triple@count by 1
         \ifnum\triple@count=1
            \xdef\out@{\out@{-}}
         \fi
         \ifnum\triple@count=3
            \expandafter\ifx\csname @\triplet@\endcsname\relax
            \expandafter\xdef\csname @\triplet@\endcsname{?} \fi
            \xdef\out@{\out@\csname @\triplet@\endcsname\out@@{-}}
            \triple@count=0
            \xdef\triplet@{}
            \xdef\out@@{}
         \fi
      \fi
      \if\first@ -
         \ifnum\triple@count<2
            \xdef\out@{\out@{-}}
         \else
            \xdef\out@@{\out@@{-}}
         \fi
      \fi
      \if\first@ +
         \expandafter\char@get\tr@nsl@ted
         \xdef\triplet@{\triplet@\first@}
         \advance\triple@count by 1
         \ifnum\triple@count=3
            \expandafter\ifx\csname @\triplet@\endcsname\relax
            \expandafter\xdef\csname @\triplet@\endcsname{?} \fi
            \xdef\out@{\out@\csname @\triplet@\endcsname\out@@}
            \triple@count=0
            \xdef\triplet@{}
            \xdef\out@@{}
         \fi
      \fi
      \if\first@ 2
         \loop
            \expandafter\char@get\tr@nsl@ted
            \ifx\first@\ampers@nd
               \lettertrue
               \xdef\tr@nsl@ted{&@}
            \else
               \expandafter\check@char\first@
            \fi
         \ifletter\else\xdef\out@{\out@{-}}\repeat
         \xdef\out@{\out@{-}}
      \fi
   \expandafter\trans@now\tr@nsl@ted
   \fi
}

\def\do@translation{%
   \xdef\triplet@{}
   \xdef\out@{}
   \xdef\out@@{}
   \xdef\tr@nsl@ted{\tr@nsl@ted &@}
   \triple@count=0
   \expandafter\trans@now\tr@nsl@ted
   \xdef\tr@nsl@ted{\out@}
}

\def\trans@pep#1#2@{%
   \xdef\first@{#1}
   \xdef\tr@nsl@ted{#2@}
   \ifx\first@\ampers@nd
   \else
      \expandafter\check@char\first@
  \ifletter
     \xdef\out@{\out@\csname rev@\first@\endcsname}
  \else
     \xdef\out@{\out@{-}{-}{-}}
  \fi
   \expandafter\trans@pep\tr@nsl@ted
   \fi
}

\def\rev@translation{%
   \xdef\out@{}
   \xdef\tr@nsl@ted{\tr@nsl@ted &@}
   \expandafter\trans@pep\tr@nsl@ted
   \xdef\tr@nsl@ted{\out@}
}


\def\sum@up{%
   \advance\innerloopcount by 1
   \xdef\second@@@{\csname res\the\innerloopcount\endcsname}
   \xdef\third@@@{\csname cons\first@@@\second@@@\endcsname}   
   \advance\temp@count by \third@@@   
   \ifnum\innerloopcount<\seq@count\sum@up\fi
}

\def\sum@up@cons{%
   \innerloopcount=\outerloopcount
   \xdef\first@@@{\csname res\the\outerloopcount\endcsname}
   \sum@up
   \advance\outerloopcount by 1\relax
   \ifnum\outerloopcount<\seq@count
      \sum@up@cons
   \else
      \innerloopcount=\seq@count
      \advance\innerloopcount by -1
      \multiply\innerloopcount by \seq@count
      \multiply\temp@count by 2
      \multiply\temp@count by \m@trixf@ctor
      \divide\temp@count by \innerloopcount
      
      \ifnum\temp@count>100 \temp@count=100 \fi
      \ifnum\temp@count<0   \temp@count=0 \fi
      \xdef\cons@val{\the\temp@count}
      \ifx\first@\equ@l \xdef\cons@val{N} \fi
      
      \ifx\T@coffee@bcons\y@
         \xdef\cons@val{\csname TC@num0\endcsname}
         \if\cons@val * 
            \xdef\cons@val{99}
         \fi
      \fi
      \ifx\T@coffee@ccons\y@
         \xdef\cons@val{\csname TC@num0\endcsname}
         \if\cons@val * 
            \xdef\cons@val{99}
         \fi
      \fi
   \fi
}

\def\collect@cons@res{%
   \xdef\temp@{\temp@\csname res\the\innerloopcount\endcsname}
   \advance\innerloopcount by 1
   \ifnum\innerloopcount>\seq@count\relax
      \outerloopcount=1\relax
      \temp@count=0\relax
      \sum@up@cons
   \else
      \collect@cons@res
   \fi
}

\def\sum@@up{%
   \advance\innerloopcount by 1
   \xdef\second@@@{\csname res\the\innerloopcount\endcsname}
   
   \ifx\first@@@\d@t
   \else
    \ifx\first@@@\questi@n
    \else
     \ifx\first@@@\equ@l
     \else       
         \ifx\second@@@\d@t
         \else
          \ifx\second@@@\questi@n
          \else
           \ifx\second@@@\equ@l
           \else
               \xdef\third@@@{\csname poscount\the\outerloopcount @\the\innerloopcount\endcsname}
               \temp@@count=\third@@@
               \advance\temp@@count by 1
               \expandafter\xdef\csname poscount\the\outerloopcount @\the\innerloopcount\endcsname{\the\temp@@count}
               
               \xdef\third@@@{\csname simcount\the\outerloopcount @\the\innerloopcount\endcsname}
               \temp@@count=\third@@@
               \xdef\third@@@{\csname simpair\first@@@\second@@@\endcsname}
               \advance\temp@@count by \third@@@
               \expandafter\xdef\csname simcount\the\outerloopcount @\the\innerloopcount\endcsname{\the\temp@@count}
               
               \ifx\first@@@\second@@@
                  \xdef\third@@@{\csname identcount\the\outerloopcount @\the\innerloopcount\endcsname}
                  \temp@@count=\third@@@
                  \advance\temp@@count by 1
                  \expandafter\xdef\csname identcount\the\outerloopcount @\the\innerloopcount\endcsname{\the\temp@@count}
               \fi
           \fi
          \fi
         \fi
     \fi
    \fi
   \fi
   
   \ifnum\innerloopcount<\seq@count\sum@@up\fi
}

\def\sum@up@sim{%
   \innerloopcount=\outerloopcount
   \xdef\first@@@{\csname res\the\outerloopcount\endcsname}
   \sum@@up
   \advance\outerloopcount by 1\relax
   \ifnum\outerloopcount<\seq@count
      \sum@up@sim
   \fi
}

\def\collect@similarity{%
   \xdef\temp@{\temp@\csname res\the\innerloopcount\endcsname}
   \advance\innerloopcount by 1
   \ifnum\innerloopcount>\seq@count\relax
      \outerloopcount=1\relax
      \temp@count=0\relax
      \sum@up@sim
   \else
      \collect@similarity
   \fi
}


\def\sum@up@bits{%
   \ifx\first@@@\second@@@
      \ifx\first@@@\d@t
         \temp@count=0\relax
      \else
         \temp@count=100\relax
      \fi
   \else
      \temp@count=0
%      \xdef\third@@@{\csname cons\first@@@\second@@@\endcsname}\temp@count=\third@@@   %%% or if only identical =0
   \fi
   \multiply\temp@count by \last@ 
   \divide\temp@count by 100  
   \expandafter\ifx\csname info@\subfamily@seq @\the\outerloopcount\endcsname\relax
      \expandafter\xdef\csname info@\subfamily@seq @\the\outerloopcount\endcsname{0}
   \fi
   \xdef\third@@@{\csname info@\subfamily@seq @\the\outerloopcount\endcsname} 
   \advance\temp@count by \third@@@ 
   \expandafter\xdef\csname info@\subfamily@seq @\the\outerloopcount\endcsname{\the\temp@count}
}
\def\sum@up@info{%
   \ifnum\outerloopcount=\subfamily@seq \advance\outerloopcount by 1 \fi
   \ifnum\outerloopcount>\seq@count
   \else
      \xdef\first@@@{\csname res\the\outerloopcount\endcsname}
      \sum@up@bits
      \advance\outerloopcount by 1\relax
   \fi
   \ifnum\outerloopcount>\seq@count\else\sum@up@info\fi
}
\def\collect@info{%
   \outerloopcount=1\relax
   \temp@count=0\relax
   \xdef\second@@@{\csname res\subfamily@seq\endcsname}   
   \sum@up@info
}

\def\calc@grouping{%
   \xdef\second@{\csname res@num\d@t\endcsname}
   \innerloopcount=\seq@count
   \advance\innerloopcount by -\second@
   \xdef\second@{\the\innerloopcount}
   \xdef\seventh@{0}
   \innerloopcount=0
   \outerloopcount=65
   \loop
      \xdef\first@@{\csname ch@r@\the\outerloopcount\endcsname}
      \xdef\first@{\csname res@num\first@@\endcsname}
      \ifnum\first@>0
         \temp@count=\csname log2@\first@\endcsname
         \advance\temp@count by -\csname log2@\second@\endcsname
         \advance\temp@count by \csname res@corr\first@@\endcsname
         \multiply\temp@count by -\first@
         \divide\temp@count by \second@
         \advance\innerloopcount by \temp@count\relax
         \temp@count=\second@              %%%
         \advance\temp@count by -\first@   %%%
         \advance\temp@count by \seventh@  %%%
         \xdef\seventh@{\the\temp@count}   %%%
      \fi
      \expandafter\xdef\csname res@num\first@@\endcsname{0}
      \advance\outerloopcount by 1
   \ifnum\outerloopcount>90\else\repeat
   \expandafter\xdef\csname res@num\d@t\endcsname{0}
   \if\seq@type P
      \temp@count=\csname log2@20\endcsname 
      \ifnum\sig@max=100000 \xdef\sig@max{2321} \fi
   \else
      \temp@count=\csname log2@4\endcsname
      \ifnum\sig@max=100000 \xdef\sig@max{1000} \fi
   \fi
   \xdef\bit@max{\the\temp@count}
   \advance\temp@count by -\innerloopcount
   \multiply\temp@count by \second@
   \divide\temp@count by \seq@count
   \xdef\last@{\the\temp@count}
   \expandafter\xdef\csname bit@pos\the\loopcount\endcsname{\last@}
   \temp@count=\bit@total
   \advance\temp@count by \last@
   \xdef\bit@total{\the\temp@count}
   \collect@info
}
\def\do@grouping{%
   \ifnum\loopcount>\total@pos
   \else
      \triple@count=1
      \loop
         \xdef\first@{\csname seq\the\triple@count\endcsname}
         \expandafter\dis@get\first@
         \ifx\first@\ampers@nd
         \else
            \innerloopcount=\csname res@num\first@\endcsname
            \advance\innerloopcount by 1 
            \expandafter\xdef\csname res@num\first@\endcsname{\the\innerloopcount}
         \fi
         \advance\triple@count by 1\relax 
      \ifnum\triple@count>\seq@count
         \calc@grouping
         \advance\loopcount by 1
         \do@grouping
      \else
      \repeat 
   \fi
}
\def\total@frequency@correction{%
   \innerloopcount=0 
   \outerloopcount=65
   \loop
      \xdef\first@@{\csname ch@r@\the\outerloopcount\endcsname}
      \xdef\first@{\csname res@corr\first@@\endcsname}
      \advance\innerloopcount by \first@\relax  
      \advance\outerloopcount by 1
   \ifnum\outerloopcount>90\else\repeat
   \xdef\res@num@total{\the\innerloopcount}
   \outerloopcount=65
   \loop
      \xdef\first@@{\csname ch@r@\the\outerloopcount\endcsname}
      \xdef\first@{\csname res@corr\first@@\endcsname}
      \innerloopcount=\first@\relax  
      \ifnum\innerloopcount>0
         \multiply\innerloopcount by 1000\relax   
         \divide\innerloopcount by \res@num@total\relax  
         \innerloopcount=\csname log2@\the\innerloopcount\endcsname  
         \advance\innerloopcount by -5644\relax  
         \multiply\innerloopcount by -1\relax
      \fi 
      \expandafter\xdef\csname res@corr\first@@\endcsname{\the\innerloopcount}
      \ifnum\the\innerloopcount>\corr@max \xdef\corr@max{\the\innerloopcount}\fi
      \advance\outerloopcount by 1  
   \ifnum\outerloopcount>90\else\repeat
}
\def\calc@total@frequency{%
   \ifnum\loopcount>\total@pos
   \else
      \triple@count=1
      \loop
         \xdef\first@{\csname seq\the\triple@count\endcsname}
         \expandafter\tot@get\first@ 
         \ifx\first@\ampers@nd
         \else
            \innerloopcount=\csname res@corr\first@\endcsname 
            \advance\innerloopcount by 1 
            \expandafter\xdef\csname res@corr\first@\endcsname{\the\innerloopcount}
         \fi
         \advance\triple@count by 1\relax 
      \ifnum\triple@count>\seq@count
         \advance\loopcount by 1
         \calc@total@frequency
      \else
      \repeat 
   \fi
}
\def\prep@logo{%
    \if\seq@type P
     \ifx\do@freq@correction\y@
       \xdef\do@freq@correction{n}
       \loopcount=1
       \loop
          \expandafter\xdef\csname seq\the\loopcount\endcsname{\csname sequence\the\loopcount\endcsname &@}
          \advance\loopcount by 1
       \ifnum\loopcount>\seq@count\else\repeat
       \xdef\corr@max{0}
       \loopcount=1
       \calc@total@frequency   
       \total@frequency@correction 
     \fi
    \fi
}
\def\define@subfamilies{%
  \xdef\first@{\csname group@num1\endcsname}
  \advance\loopcount by -1 
  \xdef\last@{\the\loopcount}
  \outerloopcount=\bit@total
  \divide\outerloopcount by \last@
  \xdef\bit@mean{\the\outerloopcount}
  \multiply\outerloopcount by \subfamily@threshold\relax
  \divide\outerloopcount by 100\relax
  \xdef\sub@threshold{\the\outerloopcount}
  \ifnum\first@=0
    \clear@res@nums{1}
    \clear@res@nums{2}
    \immediate\write\featurefile{TeXshade subfamily logo data file for \alignfilename}
    \immediate\write\featurefile{--}
    \immediate\write\featurefile{Average information content [1000*bits per position]: \bit@mean}
    \immediate\write\featurefile{Subfamily threshold setting [percent]: \subfamily@threshold}
    \immediate\write\featurefile{=> \bit@mean\space * 0.\subfamily@threshold\space = \sub@threshold\space [1000*bits per position]}
    \immediate\write\featurefile{}
    \expandafter\xdef\csname subfamily@num\subfamily@seq\endcsname{2}
    \immediate\write\featurefile{Automatic subfamily assignment around sequence no. \subfamily@seq:} 
    \immediate\write\featurefile{[sequence pair: average shared information *1000, subfamily number]}
    \temp@count=1 
    \outerloopcount=1
    \loop
       \ifnum\outerloopcount=\subfamily@seq \advance\outerloopcount by 1 \fi
       \ifnum\outerloopcount>\seq@count
       \else
          \xdef\first@{\csname info@\subfamily@seq @\the\outerloopcount\endcsname}
          \loopcount=\first@\relax 
          \divide\loopcount by \last@\relax 
          \ifnum\loopcount>\sub@threshold
             \expandafter\xdef\csname subfamily@num\the\outerloopcount\endcsname{2}
             \advance\temp@count by 1\relax 
          \else
             \expandafter\xdef\csname subfamily@num\the\outerloopcount\endcsname{1}
          \fi    
          \xdef\first@@{\csname subfamily@num\the\outerloopcount\endcsname}          
          \immediate\write\featurefile{\subfamily@seq-\the\outerloopcount:\space\the\loopcount,\space subfamily: \first@@}
          \advance\outerloopcount by 1 
       \fi
    \ifnum\outerloopcount>\seq@count\else\repeat
    \immediate\write\featurefile{}
    \immediate\write\featurefile{//}
    \ifnum\temp@count>0
       \expandafter\xdef\csname group@num2\endcsname{\the\temp@count}
    \else
       \message{<No subfamilies found (check threshold) - hiding subfamily logo>}
       \show@sublogofalse
       \immediate\write\featurefile{<No subfamilies found (check threshold) - hiding subfamily logo>}
    \fi
    \loopcount=\seq@count 
    \advance\loopcount by -\temp@count
    \ifnum\loopcount>0
       \expandafter\xdef\csname group@num1\endcsname{\the\loopcount}
    \else
       \message{<No subfamilies found (check threshold) - hiding subfamily logo>}
       \immediate\write\featurefile{<No subfamilies found (check threshold) - hiding subfamily logo>}
       \show@sublogofalse
    \fi
  \else
    \immediate\write\featurefile{TeXshade subfamily logo data file for \alignfilename}
    \immediate\write\featurefile{--}
    \immediate\write\featurefile{Average information content [1000*bits per position]: \bit@mean}
    \immediate\write\featurefile{Subfamily threshold setting [percent]: \subfamily@threshold}
    \immediate\write\featurefile{=> \bit@mean\space * 0.\subfamily@threshold\space = \sub@threshold\space [1000*bits per position]}
    \immediate\write\featurefile{}
    \immediate\write\featurefile{User defined subfamily: \sub@family@setting}
    \immediate\write\featurefile{}
    \immediate\write\featurefile{//}
  \fi
}
\def\prep@sublogo{%
    \message{<Calculating subfamily logo...>}
    \immediate\openout\featurefile = sublogo.txt 
    \if\seq@type P
     \ifx\do@freq@correction\y@
       \xdef\do@freq@correction{n}
       \loopcount=1
       \loop
          \xdef\first@{\csname @rd\the\loopcount\endcsname} %%%
          \expandafter\xdef\csname seq\the\loopcount\endcsname{\csname sequence\first@\endcsname &@}
          \advance\loopcount by 1
       \ifnum\loopcount>\seq@count\else\repeat
       \xdef\corr@max{0}
       \loopcount=1
       \calc@total@frequency   
       \total@frequency@correction 
     \fi
    \fi
    \loopcount=1
    \loop
          \xdef\first@{\csname @rd\the\loopcount\endcsname} %%%
          \expandafter\xdef\csname seq\the\loopcount\endcsname{\csname sequence\first@\endcsname &@}
       \advance\loopcount by 1
    \ifnum\loopcount>\seq@count\else\repeat
    \loopcount=1
    \do@grouping 
    \define@subfamilies
    \ifshow@sublogo
       \loopcount=1
       \loop
          \xdef\first@{\csname @rd\the\loopcount\endcsname} %%%
          \expandafter\xdef\csname seq\the\loopcount\endcsname{\csname sequence\first@\endcsname &@}
          \advance\loopcount by 1
       \ifnum\loopcount>\seq@count\else\repeat
       \xdef\first@{\csname group@num1\endcsname}
       \xdef\first@@{\csname log2@\first@\endcsname}
       \xdef\second@{\csname group@num2\endcsname}
       \xdef\second@@{\csname log2@\second@\endcsname}
       \loopcount=\first@ 
       \advance\loopcount by \second@ 
       \xdef\third@{\the\loopcount}
       \xdef\third@@{\csname log2@\third@\endcsname}
       \loopcount=\first@@ 
       \advance\loopcount by -\third@@ 
       \multiply\loopcount by \first@ 
       \divide\loopcount by \third@
       \xdef\first@@@{\the\loopcount}
       \loopcount=\second@@ 
       \advance\loopcount by -\third@@ 
       \multiply\loopcount by \second@ 
       \divide\loopcount by \third@
       \xdef\second@@@{\the\loopcount}
       \loopcount=\bit@max 
       \advance\loopcount by \first@@@
       \advance\loopcount by \second@@@
       \xdef\group@correction{\the\loopcount}
       \loopcount=1 \total@count=0
       \do@sublogo 
       \immediate\closeout\featurefile  
       \openin\sublogofile = sublogo.txt\relax
       \read@header     
    \else
       \immediate\closeout\featurefile  
    \fi 
}
\def\read@header{%
   \read\sublogofile to \first@\relax
   \xdef\first@{\expandafter\string\first@}
   \ifx\first@\par@ \read@header
   \else
      \xdef\first@{\first@ @}
      \expandafter\seq@get\first@
      \ifx\first@\he@derend
      \else\read@header
      \fi
   \fi
}
\def\read@sublogo{%
   \xdef\stack@sublogo{}%
   \loopcount=1 %
   \loop%
      \ifeof\sublogofile%
         \advance\loopcount by \res@perline%
      \else%
         \read\sublogofile to \first@\relax%
         \xdef\first@{\expandafter\string\first@}%
         \ifx\first@\par@%
            \advance\loopcount by \res@perline%
         \else
            \xdef\first@{\first@ @}%
            \expandafter\sublogo@get\first@%
            \xdef\stack@sublogo{\stack@sublogo\first@}%
            \advance\loopcount by 1%
         \fi%
      \fi%
   \ifnum\loopcount=\res@perline\else\repeat%
}
\def\calc@sublogo{%
   \xdef\seventh@{\csname res@num.2\endcsname}
   \xdef\first@{\csname group@num2\endcsname}
   \innerloopcount=\first@
   \advance\innerloopcount by -\seventh@
   \xdef\seventh@{\the\innerloopcount}
   \xdef\eighth@{\csname res@num.1\endcsname}
   \xdef\first@{\csname group@num1\endcsname}
   \innerloopcount=\first@
   \advance\innerloopcount by -\eighth@
   \xdef\eighth@{\the\innerloopcount}
   \xdef\nineth@{n}
   \xdef\pos@max{0}
   \xdef\pos@min{0}
   \outerloopcount=65
   \loop
      \xdef\first@@{\csname ch@r@\the\outerloopcount\endcsname}
      \xdef\first@{\csname res@num\first@@ 1\endcsname} \xdef\temp@{\first@}
      \temp@count=\first@ 
      \ifnum\seventh@>0 \multiply\temp@count by \seventh@ \fi 
      \xdef\first@{\the\temp@count}
      \xdef\second@{\csname res@num\first@@ 2\endcsname}
      \temp@count=\temp@
      \advance\temp@count by \second@\relax 
      \ifnum\temp@count>0
         \temp@count=\second@ 
         \ifnum\eighth@>0 \multiply\temp@count by \eighth@ \fi 
         \advance\temp@count by -\first@ 
         \xdef\last@{\csname bit@pos\the\loopcount\endcsname}
         \multiply\temp@count by \last@ 
         \ifnum\seventh@>0 \divide\temp@count by \seventh@ \fi 
         \ifnum\eighth@>0 \divide\temp@count by \eighth@ \fi 
         \multiply\temp@count by \bit@max  
         \divide\temp@count by \group@correction   
         \ifx\hide@negatives\y@
            \ifnum\temp@count>0 \else \temp@count=1 \fi%
         \else
            \ifnum\temp@count=0 \temp@count=1 \fi%
         \fi 
         \expandafter\xdef\csname res@val\first@@\endcsname{\the\temp@count}
         \expandafter\xdef\csname res@num\first@@ 1\endcsname{0}
         \expandafter\xdef\csname res@num\first@@ 2\endcsname{0}
         \xdef\nineth@{y}
      \else 
         \expandafter\xdef\csname res@val\first@@\endcsname{0}
         \expandafter\xdef\csname res@num\first@@ 1\endcsname{0}
         \expandafter\xdef\csname res@num\first@@ 2\endcsname{0}
      \fi
      \ifnum\temp@count>\pos@max \xdef\pos@max{\the\temp@count}\fi
      \ifnum\temp@count<\pos@min \xdef\pos@min{\the\temp@count}\fi
      \advance\outerloopcount by 1
   \ifnum\outerloopcount>90\else\repeat
   \expandafter\xdef\csname res@val\d@t\endcsname{0}
   \expandafter\xdef\csname res@num\d@t 1\endcsname{0}
   \expandafter\xdef\csname res@num\d@t 2\endcsname{0}
   \temp@count=\pos@min\relax 
   \multiply\temp@count by -1\relax 
   \xdef\pos@min{\the\temp@count}
   \ifnum\pos@min>\pos@max \xdef\pos@max{\pos@min}\fi
   \temp@count=\pos@max
   \multiply\temp@count by 100 \relax
   \divide\temp@count by \bit@max \relax
   \xdef\sublogo@num{\sublogo@num\the\temp@count,}
   \ifnum\pos@max<\sig@max \xdef\sublogo@sig{\sublogo@sig n} \else  \xdef\sublogo@sig{\sublogo@sig y} \fi
   \ifx\nineth@\y@ 
      \xdef\tmpstack{&:&,}
      \outerloopcount=65
      \loop
         \xdef\first@@{\csname ch@r@\the\outerloopcount\endcsname}
         \xdef\second@@{\csname res@val\first@@\endcsname}
         \ifnum\second@@=0 
         \else
            \xdef\third@{\tmpstack @} \xdef\tmpstack{} 
            \sort@logostack
         \fi
         \advance\outerloopcount by 1
      \ifnum\outerloopcount>90 
         \immediate\write\featurefile{\expandafter\string\tmpstack&:&,} 
      \else\repeat
   \else
      \immediate\write\featurefile{\expandafter\string O:0,&:&,} 
   \fi
}
\def\do@sublogo{%
   \ifnum\loopcount>\total@pos 
   \xdef\sublogo@sig{\sublogo@sig &@}
   \else 
      \triple@count=1
      \loop
         \xdef\first@{\csname seq\the\triple@count\endcsname}
         \expandafter\dis@get\first@
         \ifnum\triple@count=\start@seq
            \ifx\first@\d@t\else \advance\total@count by 1\relax\fi  
         \fi
         \ifnum\loopcount<\start@number
         \else
         \ifnum\total@count>\end@num
         \else
            \ifx\first@\ampers@nd
            \else
               \xdef\second@{\csname subfamily@num\the\triple@count\endcsname}
               \innerloopcount=\csname res@num\first@\second@\endcsname
               \advance\innerloopcount by 1 
               \expandafter\xdef\csname res@num\first@\second@\endcsname{\the\innerloopcount}
            \fi
         \fi
         \fi
         \advance\triple@count by 1\relax 
      \ifnum\triple@count>\seq@count
         \ifnum\loopcount<\start@number
         \else
            \ifnum\total@count>\end@num
            \else
               \calc@sublogo 
            \fi
         \fi
         \advance\loopcount by 1   
         \do@sublogo
      \else
      \repeat
   \fi
}
\def\calc@logo{%
   \xdef\second@{\csname res@num\d@t\endcsname}
   \innerloopcount=\seq@count
   \advance\innerloopcount by -\second@
   \xdef\second@{\the\innerloopcount}
   \xdef\seventh@{0}
   \innerloopcount=0
   \outerloopcount=65
   \loop
      \xdef\first@@{\csname ch@r@\the\outerloopcount\endcsname}
      \xdef\first@{\csname res@num\first@@\endcsname}
      \ifnum\first@>0
         \temp@count=\csname log2@\first@\endcsname
         \advance\temp@count by -\csname log2@\second@\endcsname
         \advance\temp@count by \csname res@corr\first@@\endcsname     %%%% correct for background freq.
         \multiply\temp@count by -\first@
         \divide\temp@count by \second@
         \advance\innerloopcount by \temp@count\relax
      \fi
      \advance\outerloopcount by 1
   \ifnum\outerloopcount>90\else\repeat
   \if\seq@type P
      \temp@count=\csname log2@20\endcsname
   \else
      \temp@count=\csname log2@4\endcsname
   \fi
   \advance\temp@count by -\innerloopcount
   \multiply\temp@count by \second@
   \divide\temp@count by \seq@count
   \xdef\last@{\the\temp@count}
   \xdef\nineth@{n}
   \outerloopcount=65
   \loop
      \xdef\first@@{\csname ch@r@\the\outerloopcount\endcsname}
      \xdef\first@{\csname res@num\first@@\endcsname}
      \ifnum\first@>0
         \temp@count=\last@
         \multiply\temp@count by \first@\relax 
         \divide\temp@count by \second@\relax 
         \ifnum\temp@count=0 %
            \temp@count=1 %
         \fi%
         \expandafter\xdef\csname res@val\first@@\endcsname{\the\temp@count}
         \expandafter\xdef\csname res@num\first@@\endcsname{0}
         \xdef\nineth@{y}
      \else
         \expandafter\xdef\csname res@val\first@@\endcsname{0}
      \fi
      \advance\outerloopcount by 1
   \ifnum\outerloopcount>90\else\repeat
   \expandafter\xdef\csname res@val\d@t\endcsname{0}
   \expandafter\xdef\csname res@num\d@t\endcsname{0}
   \ifx\nineth@\y@
      \xdef\tmpstack{&:&,}
      \outerloopcount=65
      \loop
         \xdef\first@@{\csname ch@r@\the\outerloopcount\endcsname}
         \xdef\second@@{\csname res@val\first@@\endcsname}
         \ifnum\second@@>0
            \xdef\third@{\tmpstack @} \xdef\tmpstack{}
            \sort@logostack
         \fi
         \advance\outerloopcount by 1
      \ifnum\outerloopcount>90
         \xdef\stack@sequencelogo{\stack@sequencelogo\tmpstack&:&,}
      \else\repeat
   \else
      \xdef\stack@sequencelogo{\stack@sequencelogo O:0,&:&,}
   \fi
}
\def\get@fromlogostack#1:#2,#3@{%
   \xdef\first@{#1} \xdef\second@{#2} \xdef\third@{#3}
}
\def\sort@logostack{%
   \expandafter\get@fromlogostack\third@
   \ifx\first@\ampers@nd 
      \xdef\tmpstack{\tmpstack\first@@:\second@@,}
   \else
   \ifnum\second@@<0 
      \ifnum\second@>0 
         \xdef\tmpstack{\tmpstack\first@@:\second@@,\first@:\second@,\third@}
      \else       
         \temp@count=\second@@\relax \multiply\temp@count by -1\relax  
         \triple@count=\second@\relax \multiply\triple@count by -1\relax  
         \ifnum\temp@count<\triple@count
            \xdef\tmpstack{\tmpstack\first@@:\second@@,\first@:\second@,\third@}
         \else
            \xdef\tmpstack{\tmpstack\first@:\second@,}
            \xdef\third@@{\third@ .}
            \ifx\third@@\d@t
               \xdef\third@{&:&,@}
            \else 
               \xdef\third@{\third@ @}
            \fi
            \sort@logostack
         \fi
      \fi
   \else
      \ifnum\second@@<\second@ 
         \xdef\tmpstack{\tmpstack\first@@:\second@@,\first@:\second@,\third@}
      \else
         \xdef\tmpstack{\tmpstack\first@:\second@,}
         \xdef\third@@{\third@ .}
         \ifx\third@@\d@t
            \xdef\third@{&:&,@}
         \else 
            \xdef\third@{\third@ @}
         \fi
         \sort@logostack
      \fi
   \fi
   \fi
}


\def\calc@feature{%
%      \advance\pos@count by -1
      \loopcount=-1 \innerloopcount=0
      \loop
         \advance\loopcount by 1
         \ifnum\loopcount=0 \xdef\first@@{\the\cons@count}
         \else \xdef\first@@{\csname res@count\the\loopcount\endcsname} \fi
         \expandafter\ifx\csname start\bottop@\the\loopcount\endcsname\ampers@nd
         \else
         \expandafter\ifnum\csname start\bottop@\the\loopcount\endcsname>\first@@
         \else
            \expandafter\ifnum\csname stop\bottop@\the\loopcount\endcsname<\first@@
               \expandafter\getregion@fromfstack{\the\loopcount}
            \else
               \innerloopcount=\loopcount
               \expandafter\ifnum\csname featureon\bottop@\endcsname=0
                  \expandafter\xdef\csname featureon\bottop@\endcsname{1}
                  \expandafter\xdef\csname feature@\bottop@\endcsname{1}
                  \expandafter\xdef\csname ftext\bottop@\the\loopcount\endcsname{%
                                   \csname text\bottop@\the\loopcount\endcsname}
                  \expandafter\xdef\csname fstyle\bottop@\the\loopcount\endcsname{%
                                   \csname style\bottop@\the\loopcount\endcsname}
                  \innerloopcount=\pos@count
                  \advance\innerloopcount by -\csname featurepos\bottop@\endcsname
                  \expandafter\xdef\csname textfeature\bottop@\endcsname{%
                                   \csname textfeature\bottop@\endcsname%
                                   &\the\innerloopcount;{};}
                  \expandafter\xdef\csname stylefeature\bottop@\endcsname{%
                                   \csname stylefeature\bottop@\endcsname%
                                   &\the\innerloopcount;&;}
                  \expandafter\xdef\csname featurepos\bottop@\endcsname{\the\pos@count}
                  \xdef\temp@@@{n}
                  \xdef\fourth@{}
                  \xdef\temp@{\csname ftext\bottop@\the\loopcount\endcsname}
                  \xdef\temp@{\temp@[,]:[,][]&}\expandafter\graph@opt@color\temp@
                  \ifx\fourth@\tr@ns \xdef\temp@@@{y} \fi
                  \xdef\fourth@{}
                  \xdef\temp@{\csname fstyle\bottop@\the\loopcount\endcsname}
                  \xdef\temp@{\temp@[,]:[,][]&}\expandafter\graph@opt@color\temp@
                  \ifx\fourth@\tr@ns \xdef\temp@@@{y} \fi
                  \ifx\temp@@@\y@
                    \ifnum\loopcount=0
                     \message{<No translations of the consensus sequence>}
                     \expandafter\xdef\csname ftext\bottop@\the\loopcount\endcsname{%
                                  No consensus translations!}
                     \expandafter\xdef\csname fstyle\bottop@\the\loopcount\endcsname{%
                                          ///}
                    \else
                     \if\seq@type P
                        \expandafter\xdef\csname collect@res\bottop@\endcsname{yes}
                        \expandafter\xdef\csname tr@nsseq\bottop@\endcsname{\the\loopcount}
                        \expandafter\xdef\csname tr@nslate\bottop@\endcsname{%
                        \csname res\the\loopcount\endcsname}
                     \else
                        \expandafter\xdef\csname collect@res\bottop@\endcsname{yes}
                        \expandafter\xdef\csname tr@nsseq\bottop@\endcsname{\the\loopcount}
                        \expandafter\ifx\csname res\the\loopcount\endcsname\gap@char
                           \expandafter\xdef\csname triple@count\bottop@\endcsname{0}
                           \expandafter\xdef\csname tr@nslate\bottop@\endcsname{-}
                        \else
                           \expandafter\xdef\csname triple@count\bottop@\endcsname{1}
                           \expandafter\xdef\csname tr@nslate\bottop@\endcsname{%
                           \csname res\the\loopcount\endcsname}
                           \expandafter\xdef\csname last@@res\bottop@\endcsname{%
                           \csname res\the\loopcount\endcsname}
                        \fi
                     \fi
                    \fi
                  \fi
                  \xdef\temp@{plot}
                  \ifx\temp@\fourth@
                    \ifnum\loopcount=0
                     \message{<No bar graphs/scales with the consensus sequence>}
                     \expandafter\xdef\csname ftext\bottop@\the\loopcount\endcsname{}
                     \expandafter\xdef\csname fstyle\bottop@\the\loopcount\endcsname{%
                                          ///}
                    \else
                     \expandafter\xdef\csname collect@val\bottop@\endcsname{yes}
                     \expandafter\xdef\csname v@lseq\bottop@\endcsname{\the\loopcount}
                     \expandafter\xdef\csname ffourth@\bottop@\endcsname{\ffourth@}
                     \xdef\temp@{\ffourth@\csname res\the\loopcount\endcsname}
                     \expandafter\xdef\csname v@l\bottop@\endcsname{\csname \temp@\endcsname}
                    \fi
                  \fi
                  \xdef\temp@{cons}
                  \ifx\temp@\fourth@
                     \expandafter\xdef\csname collect@cons@graph\bottop@\endcsname{yes}
                     \expandafter\xdef\csname v@lseq\bottop@\endcsname{\the\loopcount}
                     \expandafter\xdef\csname ffourth@\bottop@\endcsname{\ffourth@}
                     \innerloopcount=1
                     \collect@cons@res
                     \expandafter\xdef\csname v@l\bottop@\endcsname{\cons@val}
                  \fi
               \else
                  \ifnum\pos@count=1\relax
                     \xdef\temp@{\csname fstyle\bottop@\the\loopcount\endcsname @}
                     \expandafter\getstyle@right\temp@
                     \xdef\temp@@@{n}
                     \xdef\fourth@{}
                     \xdef\temp@{\csname ftext\bottop@\the\loopcount\endcsname}
                     \xdef\temp@{\temp@[,]:[,][]&}\expandafter\graph@opt@color\temp@
                     \ifx\fourth@\tr@ns \xdef\temp@@@{y} \fi
                     \xdef\fourth@{}
                     \xdef\temp@{\csname fstyle\bottop@\the\loopcount\endcsname}
                     \xdef\temp@{\temp@[,]:[,][]&}\expandafter\graph@opt@color\temp@
                     \ifx\fourth@\tr@ns \xdef\temp@@@{y} \fi
                     \ifx\temp@@@\y@
                        \if\seq@type P
                           \expandafter\xdef\csname tr@nslate\bottop@\endcsname{%
                           \csname res\the\loopcount\endcsname}
                        \else
                           \expandafter\ifnum\csname triple@count\bottop@\endcsname=2
                              \expandafter\ifx\csname res\the\loopcount\endcsname\gap@char
                              \else
                                 \expandafter\xdef\csname tr@nslate\bottop@\endcsname{%
                                 +\csname last@res\bottop@\endcsname%
                                 \csname tr@nslate\bottop@\endcsname}
                              \fi
                           \fi
                           \expandafter\ifnum\csname triple@count\bottop@\endcsname=1
                              \expandafter\xdef\csname last@res\bottop@\endcsname{%
                              \csname last@@res\bottop@\endcsname}
                              \expandafter\ifx\csname res\the\loopcount\endcsname\gap@char
                                 \expandafter\xdef\csname tr@nslate\bottop@\endcsname{%
                                  +\csname last@res\bottop@\endcsname%
                                  \csname tr@nslate\bottop@\endcsname}
                              \fi
                           \fi
                        \fi
                     \fi
                     \xdef\temp@{plot}
                     \ifx\temp@\fourth@
                        \expandafter\ifx\csname res\the\loopcount\endcsname\gap@char
                           \expandafter\xdef\csname v@l\bottop@\endcsname{N}
                        \else
                           \xdef\temp@{\ffourth@\csname res\the\loopcount\endcsname}
                           \expandafter\xdef\csname v@l\bottop@\endcsname{%
                           \csname \temp@\endcsname}
                        \fi
                     \fi 
                     \xdef\temp@{cons}
                     \ifx\temp@\fourth@
                        \innerloopcount=1
                        \collect@cons@res
                        \expandafter\xdef\csname v@l\bottop@\endcsname{\cons@val}
                     \fi
                  \fi
               \fi
               \expandafter\ifnum\csname stop\bottop@\the\loopcount\endcsname=\first@@
                  \expandafter\getregion@fromfstack{\the\loopcount}
                  \xdef\temp@@@{n}
                  \xdef\fourth@{}
                  \xdef\temp@{\csname ftext\bottop@\the\loopcount\endcsname}
                  \xdef\temp@{\temp@[,]:[,][]&}\expandafter\graph@opt@color\temp@
                  \ifx\fourth@\tr@ns
                     \ifx\f@color\comm@
                        \xdef\f@color{} \else \xdef\f@color{[\f@color]}
                     \fi
                     \xdef\tr@nsl@ted{\csname tr@nslate\bottop@\endcsname}
                     \if\seq@type N \do@translation \else \rev@translation \fi
                     \xdef\temp@{translate:\tr@nsl@ted\f@color}
                     \xdef\temp@@@{y}
                  \else
                     \xdef\temp@{\csname ftext\bottop@\the\loopcount\endcsname}
                  \fi
                  \xdef\fourth@{}
                  \xdef\temp@@{\csname fstyle\bottop@\the\loopcount\endcsname}
                  \xdef\temp@@{\temp@@[,]:[,][]&}\expandafter\graph@opt@color\temp@@
                  \ifx\fourth@\tr@ns
                     \ifx\f@color\comm@
                         \xdef\f@color{} \else \xdef\f@color{[\f@color]}
                     \fi
                     \xdef\tr@nsl@ted{\csname tr@nslate\bottop@\endcsname}
                     \if\seq@type N \do@translation \else \rev@translation \fi
                     \xdef\temp@@{translate:\tr@nsl@ted\f@color}
                     \xdef\temp@@@{y}
                  \else
                     \xdef\temp@@{plot}
                     \ifx\temp@@\fourth@
                        \xdef\temp@@{plot\f@color[\fffourth@]:\csname v@l\bottop@\endcsname[\ff@color]}
                        \expandafter\xdef\csname collect@val\bottop@\endcsname{no}
                        \expandafter\xdef\csname v@l\bottop@\endcsname{}
                     \else
                     \xdef\temp@@{cons}
                     \ifx\temp@@\fourth@
                        \xdef\temp@@{plot\f@color[\fffourth@]:\csname v@l\bottop@\endcsname[\ff@color]}
                        \expandafter\xdef\csname collect@cons@graph\bottop@\endcsname{no}
                        \expandafter\xdef\csname v@l\bottop@\endcsname{}
                     \else
                        \xdef\temp@@{\csname fstyle\bottop@\the\loopcount\endcsname}
                     \fi
                     \fi
                  \fi
                  \ifx\temp@@@\y@
                     \expandafter\xdef\csname collect@res\bottop@\endcsname{no}
                     \expandafter\xdef\csname tr@nslate\bottop@\endcsname{}
                  \fi
                  \expandafter\ifnum\csname featureon\bottop@\endcsname=1
                     \expandafter\xdef\csname featureon\bottop@\endcsname{0}
                     \innerloopcount=\pos@count \advance\innerloopcount by 1
                     \advance\innerloopcount by -\csname featurepos\bottop@\endcsname
                     \expandafter\xdef\csname textfeature\bottop@\endcsname{%
                                      \csname textfeature\bottop@\endcsname%
                                      &\the\innerloopcount;\temp@;}
                     \expandafter\xdef\csname stylefeature\bottop@\endcsname{%
                                      \csname stylefeature\bottop@\endcsname%
                                      &\the\innerloopcount;\temp@@;}
                     \innerloopcount=\pos@count \advance\innerloopcount by 1
                     \expandafter\xdef\csname %
                     featurepos\bottop@\endcsname{\the\innerloopcount}
                  \fi
               \fi
               \ifnum\pos@count=\res@perline
                  \expandafter\ifnum\csname featureon\bottop@\endcsname=1
                     \innerloopcount=\pos@count
                     \advance\innerloopcount by 1
                     \advance\innerloopcount by -\csname featurepos\bottop@\endcsname
                     \xdef\temp@{\csname fstyle\bottop@\the\loopcount\endcsname @}
                     \expandafter\getstyle@left\temp@
                     \xdef\temp@@@{n}
                     \xdef\fourth@{}
                     \xdef\temp@{\csname ftext\bottop@\the\loopcount\endcsname}
                     \xdef\temp@{\temp@[,]:[,][]&}\expandafter\graph@opt@color\temp@
                     \ifx\fourth@\tr@ns
                        \ifx\f@color\comm@
                           \xdef\f@color{} \else \xdef\f@color{[\f@color]}
                        \fi
                        \expandafter\ifnum\csname triple@count\bottop@\endcsname=2
                           \if\seq@type N
                              \xdef\temp@{\csname sequence\the\loopcount\endcsname:}
                              \expandafter\get@nextres\temp@
                              \xdef\tr@nsl@ted{\csname tr@nslate\bottop@\endcsname\temp@}
                              \do@translation
                              \xdef\temp@@@{2}
                           \else
                              \xdef\tr@nsl@ted{\csname tr@nslate\bottop@\endcsname}
                              \rev@translation
                             \xdef\temp@@@{}
                           \fi
                           \xdef\temp@{translate:\tr@nsl@ted\f@color}
                        \else
                           \xdef\tr@nsl@ted{\csname tr@nslate\bottop@\endcsname}
                           \if\seq@type N \do@translation \else \rev@translation \fi
                           \xdef\temp@{translate:\tr@nsl@ted\f@color}
                           \xdef\temp@@@{}
                        \fi
                     \else \xdef\temp@{\csname ftext\bottop@\the\loopcount\endcsname}
                     \fi
                     \xdef\fourth@{}
                     \xdef\temp@@{\csname fstyle\bottop@\the\loopcount\endcsname}
                     \xdef\temp@@{\temp@@[,]:[,][]&}\expandafter\graph@opt@color\temp@@
                     \ifx\fourth@\tr@ns
                        \ifx\f@color\comm@
                           \xdef\f@color{} \else \xdef\f@color{[\f@color]}
                        \fi
                        \expandafter\ifnum\csname triple@count\bottop@\endcsname=2
                           \if\seq@type N
                              \xdef\temp@@{\csname sequence\the\loopcount\endcsname:}
                              \expandafter\get@@nextres\temp@@
                              \xdef\tr@nsl@ted{\csname tr@nslate\bottop@\endcsname\temp@@}
                              \do@translation
                              \xdef\temp@@@{2}
                           \else
                              \xdef\tr@nsl@ted{\csname tr@nslate\bottop@\endcsname}
                              \rev@translation
                              \xdef\temp@@@{}
                           \fi
                           \xdef\temp@@{translate:\tr@nsl@ted\f@color}
                        \else
                           \xdef\tr@nsl@ted{\csname tr@nslate\bottop@\endcsname}
                           \if\seq@type N \do@translation \else \rev@translation \fi
                           \xdef\temp@@{translate:\tr@nsl@ted\f@color}
                           \xdef\temp@@@{}
                        \fi
                     \else
                        \xdef\temp@@{plot}
                        \ifx\temp@@\fourth@
                          \xdef\temp@@{plot\f@color[\fffourth@]:\csname v@l\bottop@\endcsname[\ff@color]}
                        \else
                        \xdef\temp@@{cons}
                        \ifx\temp@@\fourth@
                           \xdef\temp@@{plot\f@color[\fffourth@]:\csname v@l\bottop@\endcsname[\ff@color]}
                        \else
                           \xdef\temp@@{\style@@}
                        \fi
                        \fi
                     \fi
                     \ifx\temp@@@\n@
                     \else
                        \expandafter\xdef\csname tr@nslate\bottop@\endcsname{\temp@@@}
                     \fi
                     \expandafter\xdef\csname textfeature\bottop@\endcsname{%
                                      \csname textfeature\bottop@\endcsname%
                                      &\the\innerloopcount;\temp@;}
                     \expandafter\xdef\csname stylefeature\bottop@\endcsname{%
                                      \csname stylefeature\bottop@\endcsname%
                                      &\the\innerloopcount;\temp@@;}
                     \innerloopcount=\pos@count \advance\innerloopcount by 1
                     \expandafter\xdef\csname %
                     featurepos\bottop@\endcsname{\the\innerloopcount}
                  \fi
               \fi
            \fi
         \fi
         \fi
      \ifnum\loopcount<\seq@count \repeat
%      \advance\pos@count by 1
}

\def\add@to@rule@tens{%
    \advance\innerloopcount by \ruler@step\relax
    \ifnum\innerloopcount=0 
          \ifx\allow@zero\n@ \innerloopcount=\ruler@step \fi
    \fi
    \expandafter\ifnum\csname res@count\rule@num\endcsname>\innerloopcount \add@to@rule@tens
    \else
       \xdef\rule@tens{\the\innerloopcount}
    \fi
}


\def\c@nsensus{%
 \ifnum\pos@count>\res@perline
    \else
    \loopcount=0
    \ifx\g@p\n@
       \global\advance\cons@count by 1\relax
    \fi
    \ifnum\cons@count=0\relax
       \ifx\allow@zero\n@ \global\advance\cons@count by 1 \fi
    \fi
    \ifT@coffee 
            \xdef\TC@line{\csname T@coffee0\endcsname}
            \expandafter\TC@get\TC@line
    \fi
    \expandafter\xdef\csname res@count0\endcsname{\the\cons@count}
    \loop
         \advance\loopcount by 1
         \ifT@coffee
            \xdef\TC@line{\csname T@coffee\the\loopcount\endcsname}
            \expandafter\TC@get\TC@line
            \xdef\TC@first@{\csname TC@num\the\loopcount\endcsname}
         \fi
         \xdef\seq@line{\csname sequence\the\loopcount\endcsname}
         \expandafter\residue@get\seq@line
         \xdef\first@{\csname res\the\loopcount\endcsname}
         \innerloopcount=\csname res@num\first@\endcsname
         \advance\innerloopcount by 1
         \expandafter\xdef\csname res@num\first@\endcsname{\the\innerloopcount}
         \expandafter\check@char\first@
         \ifx\dom@in\y@ %%%%%%%%%%%*****
                \xdef\temp@{\csname dom@num\the\loopcount\endcsname}
                \expandafter\get@dom@count\temp@
                \expandafter\xdef\csname dom@num\the\loopcount\endcsname{\temp@}    
         \fi            %%%%%%%%%%%*****
         \ifletter
             \ifx\dom@in\y@
             \else
                \res@count=\csname res@count\the\loopcount\endcsname
                \advance\res@count by 1
                \ifnum\res@count=0\relax
                   \ifx\allow@zero\n@ \advance\res@count by 1 \fi
             \fi
             \fi
             \expandafter\xdef\csname res@count\the\loopcount\endcsname{\the\res@count}
             \ifnum\loopcount=\chimera@seqnum \xdef\chimera@stack{\chimera@stack \the\res@count,} \fi
             \ifnum\loopcount=\exp@rt@num \xdef\sixth@{\the\res@count}\xdef\seventh@{y}\fi
             \ifnum\loopcount=\rule@num\relax 
              \expandafter\ifnum\csname res@count\rule@num\endcsname>\rule@tens 
                 \innerloopcount=\rule@tens \add@to@rule@tens
              \fi  
              \expandafter\ifnum\csname res@count\rule@num\endcsname=\rule@tens
                \expandafter\ifx\csname alt@ruler\rule@tens\endcsname\relax
                   \xdef\temp@{\rule@tens}
                \else
                   \xdef\temp@{\csname alt@ruler\rule@tens\endcsname}
                \fi
                \xdef\ruler@{\ruler@ !<\temp@>}
              \else
                \xdef\ruler@{\ruler@ -}
              \fi
              \xdef\temp@{\csname res@count\rule@num\endcsname}
              \expandafter\ifx\csname alt@ruler\temp@\endcsname\relax
              \else
                \expandafter\xdef\csname alt@ruler\temp@\endcsname{\temp@}
              \fi
             \fi
             \ifx\collect@valtop\yes
             \ifnum\v@lseqtop=\loopcount
                \xdef\v@ltop{\v@ltop,\csname \ffourth@top\first@\endcsname}
             \fi\fi
             \ifx\collect@valttop\yes
             \ifnum\v@lseqttop=\loopcount
                \xdef\v@lttop{\v@lttop,\csname \ffourth@ttop\first@\endcsname}
             \fi\fi
             \ifx\collect@valtttop\yes
             \ifnum\v@lseqtttop=\loopcount
                \xdef\v@ltttop{\v@ltttop,\csname \ffourth@tttop\first@\endcsname}
             \fi\fi
             \ifx\collect@valttttop\yes
             \ifnum\v@lseqttttop=\loopcount
                \xdef\v@lttttop{\v@lttttop,\csname \ffourth@ttttop\first@\endcsname}
             \fi\fi
             \ifx\collect@valbottom\yes
             \ifnum\v@lseqbottom=\loopcount
                \xdef\v@lbottom{\v@lbottom,\csname \ffourth@bottom\first@\endcsname}
             \fi\fi
             \ifx\collect@valbbottom\yes
             \ifnum\v@lseqbbottom=\loopcount
                \xdef\v@lbbottom{\v@lbbottom,\csname \ffourth@bbottom\first@\endcsname}
             \fi\fi
             \ifx\collect@valbbbottom\yes
             \ifnum\v@lseqbbbottom=\loopcount
                \xdef\v@lbbbottom{\v@lbbbottom,\csname \ffourth@bbbottom\first@\endcsname}
             \fi\fi
             \ifx\collect@valbbbbottom\yes
             \ifnum\v@lseqbbbbottom=\loopcount
                \xdef\v@lbbbbottom{\v@lbbbbottom,\csname \ffourth@bbbbottom\first@\endcsname}
             \fi\fi
             \ifx\collect@restop\yes
             \ifnum\tr@nsseqtop=\loopcount
                \xdef\last@restop{\last@@restop}
                \xdef\last@@restop{\first@}
                \xdef\tr@nslatetop{\tr@nslatetop\first@}
                \innerloopcount=\triple@counttop
                \advance\innerloopcount by 1
                \ifnum\innerloopcount>3 \innerloopcount=1 \fi
                \xdef\triple@counttop{\the\innerloopcount}
             \fi\fi
             \ifx\collect@resttop\yes
             \ifnum\tr@nsseqttop=\loopcount
                \xdef\last@resttop{\last@@resttop}
                \xdef\last@@resttop{\first@}
                \xdef\tr@nslatettop{\tr@nslatettop\first@}
                \innerloopcount=\triple@countttop
                \advance\innerloopcount by 1
                \ifnum\innerloopcount>3 \innerloopcount=1 \fi
                \xdef\triple@countttop{\the\innerloopcount}
             \fi\fi
             \ifx\collect@restttop\yes
             \ifnum\tr@nsseqtttop=\loopcount
                \xdef\last@restttop{\last@@restttop}
                \last@@restttop{\first@}
                \xdef\tr@nslatetttop{\tr@nslatetttop\first@}
                \innerloopcount=\triple@counttttop
                \advance\innerloopcount by 1
                \ifnum\innerloopcount>3 \innerloopcount=1 \fi
                \xdef\triple@counttttop{\the\innerloopcount}
             \fi\fi
             \ifx\collect@resttttop\yes
             \ifnum\tr@nsseqttttop=\loopcount
                \xdef\last@resttttop{\last@@resttttop}
                \xdef\last@@resttttop{\first@}
                \xdef\tr@nslatettttop{\tr@nslatettttop\first@}
                \innerloopcount=\triple@countttttop
                \advance\innerloopcount by 1
                \ifnum\innerloopcount>3 \innerloopcount=1 \fi
                \xdef\triple@countttttop{\the\innerloopcount}
             \fi\fi
             \ifx\collect@resbottom\yes
             \ifnum\tr@nsseqbottom=\loopcount
                 \xdef\last@resbottom{\last@@resbottom}
                 \xdef\last@@resbottom{\first@}
                 \xdef\tr@nslatebottom{\tr@nslatebottom\first@}
                 \innerloopcount=\triple@countbottom
                 \advance\innerloopcount by 1
                 \ifnum\innerloopcount>3 \innerloopcount=1 \fi
                 \xdef\triple@countbottom{\the\innerloopcount}
             \fi\fi
             \ifx\collect@resbbottom\yes
             \ifnum\tr@nsseqbbottom=\loopcount
                 \xdef\last@resbbottom{\last@@resbbottom}
                 \xdef\last@@resbbottom{\first@}
                 \xdef\tr@nslatebbottom{\tr@nslatebbottom\first@}
                 \innerloopcount=\triple@countbbottom
                 \advance\innerloopcount by 1
                 \ifnum\innerloopcount>3 \innerloopcount=1 \fi
                 \xdef\triple@countbbottom{\the\innerloopcount}
             \fi\fi
             \ifx\collect@resbbbottom\yes
             \ifnum\tr@nsseqbbbottom=\loopcount
                 \xdef\last@resbbbottom{\last@@resbbbottom}
                 \xdef\last@@resbbbottom{\first@}
                 \xdef\tr@nslatebbbottom{\tr@nslatebbbottom\first@}
                 \innerloopcount=\triple@countbbbottom
                 \advance\innerloopcount by 1
                 \ifnum\innerloopcount>3 \innerloopcount=1 \fi
                 \xdef\triple@countbbbottom{\the\innerloopcount}
             \fi\fi
             \ifx\collect@resbbbbottom\yes
             \ifnum\tr@nsseqbbbbottom=\loopcount
                 \xdef\last@resbbbbottom{\last@@resbbbbottom}
                 \xdef\last@@resbbbbottom{\first@}
                 \xdef\tr@nslatebbbbottom{\tr@nslatebbbbottom\first@}
                 \innerloopcount=\triple@countbbbbottom
                 \advance\innerloopcount by 1
                 \ifnum\innerloopcount>3 \innerloopcount=1 \fi
                 \xdef\triple@countbbbbottom{\the\innerloopcount}
             \fi\fi
          \innerloopcount=\csname mol@weight\the\loopcount\endcsname
          \advance\innerloopcount by \csname \prefix@ mw\first@\endcsname
          \expandafter\xdef\csname mol@weight\the\loopcount\endcsname{%
                           \the\innerloopcount}
          \innerloopcount=\csname ch@rge\the\loopcount\endcsname
          \advance\innerloopcount by \csname pepcharge\first@\endcsname
          \expandafter\xdef\csname ch@rge\the\loopcount\endcsname{%
                           \the\innerloopcount}
         \else
          \ifnum\loopcount=\rule@num \xdef\ruler@{\ruler@ -} \fi
          \ifnum\loopcount=\chimera@seqnum \xdef\chimera@stack{\chimera@stack &,} \fi
          \ifx\collect@valtop\yes
          \ifnum\v@lseqtop=\loopcount
             \xdef\v@ltop{\v@ltop,N}
          \fi\fi
          \ifx\collect@valttop\yes
          \ifnum\v@lseqttop=\loopcount
             \xdef\v@lttop{\v@lttop,N}
          \fi\fi
          \ifx\collect@valtttop\yes
          \ifnum\v@lseqtttop=\loopcount
             \xdef\v@ltttop{\v@ltttop,N}
          \fi\fi
          \ifx\collect@valttttop\yes
          \ifnum\v@lseqttttop=\loopcount
             \xdef\v@lttttop{\v@lttttop,N}
          \fi\fi
          \ifx\collect@valbottom\yes
          \ifnum\v@lseqbottom=\loopcount
             \xdef\v@lbottom{\v@lbottom,N}
          \fi\fi
          \ifx\collect@valbbottom\yes
          \ifnum\v@lseqbbottom=\loopcount
             \xdef\v@lbbottom{\v@lbbottom,N}
          \fi\fi
          \ifx\collect@valbbbottom\yes
          \ifnum\v@lseqbbbottom=\loopcount
             \xdef\v@lbbbottom{\v@lbbbottom,N}
          \fi\fi
          \ifx\collect@valbbbbottom\yes
          \ifnum\v@lseqbbbbottom=\loopcount
             \xdef\v@lbbbbottom{\v@lbbbbottom,N}
          \fi\fi
          \ifx\collect@restop\yes
          \ifnum\tr@nsseqtop=\loopcount
              \xdef\tr@nslatetop{\tr@nslatetop -}
          \fi\fi
          \ifx\collect@resttop\yes
          \ifnum\tr@nsseqttop=\loopcount
              \xdef\tr@nslatettop{\tr@nslatettop -}
          \fi\fi
          \ifx\collect@restttop\yes
          \ifnum\tr@nsseqtttop=\loopcount
              \xdef\tr@nslatetttop{\tr@nslatetttop -}
          \fi\fi
          \ifx\collect@resttttop\yes
          \ifnum\tr@nsseqttttop=\loopcount
              \xdef\tr@nslatettttop{\tr@nslatettttop -}
          \fi\fi
          \ifx\collect@resbottom\yes
          \ifnum\tr@nsseqbottom=\loopcount
              \xdef\tr@nslatebottom{\tr@nslatebottom -}
          \fi\fi
          \ifx\collect@resbbottom\yes
          \ifnum\tr@nsseqbbottom=\loopcount
              \xdef\tr@nslatebbottom{\tr@nslatebbottom -}
          \fi\fi
          \ifx\collect@resbbbottom\yes
          \ifnum\tr@nsseqbbbottom=\loopcount
              \xdef\tr@nslatebbbottom{\tr@nslatebbbottom -}
          \fi\fi
          \ifx\collect@resbbbbottom\yes
          \ifnum\tr@nsseqbbbbottom=\loopcount
              \xdef\tr@nslatebbbbottom{\tr@nslatebbbbottom -}
          \fi\fi
         \fi
    \ifnum\loopcount<\seq@count \repeat
         \ifx\collect@cons@graphtop\yes
                \outerloopcount=1\relax
                \temp@count=0\relax
                \sum@up@cons
                \xdef\v@ltop{\v@ltop,\cons@val}
         \fi
         \ifx\collect@cons@graphttop\yes
                \outerloopcount=1\relax
                \temp@count=0\relax
                \sum@up@cons
                \xdef\v@lttop{\v@lttop,\cons@val}
         \fi
         \ifx\collect@cons@graphtttop\yes
                \outerloopcount=1\relax
                \temp@count=0\relax
                \sum@up@cons
                \xdef\v@ltttop{\v@ltttop,\cons@val}
         \fi
         \ifx\collect@cons@graphttttop\yes
                \outerloopcount=1\relax
                \temp@count=0\relax
                \sum@up@cons
                \xdef\v@lttttop{\v@lttttop,\cons@val}
         \fi
         \ifx\collect@cons@graphbottom\yes
                \outerloopcount=1\relax
                \temp@count=0\relax
                \sum@up@cons
                \xdef\v@lbottom{\v@lbottom,\cons@val}
         \fi
         \ifx\collect@cons@graphbbottom\yes
                \outerloopcount=1\relax
                \temp@count=0\relax
                \sum@up@cons
                \xdef\v@lbbottom{\v@lbbottom,\cons@val}
         \fi
         \ifx\collect@cons@graphbbbottom\yes
                \outerloopcount=1\relax
                \temp@count=0\relax
                \sum@up@cons
                \xdef\v@lbbbottom{\v@lbbbottom,\cons@val}
         \fi
         \ifx\collect@cons@graphbbbbottom\yes
                \outerloopcount=1\relax
                \temp@count=0\relax
                \sum@up@cons
                \xdef\v@lbbbbottom{\v@lbbbbottom,\cons@val}
         \fi
         \ifx\collect@cons@colors\y@
                \outerloopcount=1\relax
                \temp@count=0\relax
                \sum@up@cons
                \outerloopcount=\cons@val
                \advance\outerloopcount by 4
                \divide\outerloopcount by 5
                \multiply\outerloopcount by 5
                \ifnum\outerloopcount<5\relax\outerloopcount=5\fi
                \xdef\c@nscol{\c@nscol\the\outerloopcount,}
         \fi
         \ifx\exp@rt\y@
             \ifx\seventh@\y@
                \outerloopcount=1\relax
                \temp@count=0\relax
                \sum@up@cons
                \outerloopcount=\cons@val
                \advance\outerloopcount by 4
                \divide\outerloopcount by 5
                \multiply\outerloopcount by 5
                \ifnum\outerloopcount<5\relax\outerloopcount=5\fi
                \xdef\seventh@{n}
                \immediate\write\exp@rtfile{\string color col\the\outerloopcount, resi \sixth@}
             \fi
         \fi

    \xdef\g@p{n}
    \ifregionalshadenow \calc@regshade \fi
    \ifregionaltintnow  \calc@regtint  \fi
    \ifregionalemphnow  \calc@regemph  \fi
    \ifregionallowernow \calc@reglower \fi
    \ifframenow         \calc@frame    \fi
    \ifshadingnow\iffuncmode\else\ifT@coffee\else   \calc@shading \fi\fi\fi
    \iftopfeaturenow       \xdef\bottop@{top}       \calc@feature \fi
    \ifttopfeaturenow      \xdef\bottop@{ttop}      \calc@feature \fi
    \iftttopfeaturenow     \xdef\bottop@{tttop}     \calc@feature \fi
    \ifttttopfeaturenow    \xdef\bottop@{ttttop}    \calc@feature \fi
    \ifbottomfeaturenow    \xdef\bottop@{bottom}    \calc@feature \fi
    \ifbbottomfeaturenow   \xdef\bottop@{bbottom}   \calc@feature \fi
    \ifbbbottomfeaturenow  \xdef\bottop@{bbbottom}  \calc@feature \fi
    \ifbbbbottomfeaturenow \xdef\bottop@{bbbbottom} \calc@feature \fi
    \ifnum\seq@count>1 \innerloopcount=1 \collect@similarity \fi
    \advance\pos@count by 1 
     \ifshow@logo         \calc@logo     \fi
     \ifall@fshade \all@funcshade
     \else
      \ifT@coffee \T@coffee@shade
      \else
       \ifnum\cons@num>0 \loopcount=\cons@num \else \loopcount=\seq@count \fi
       \xdef\match@case{0} \xdef\m@x{1}
       \iffuncmode
          \xdef\prfx{\prefix@} \xdef\prefix@{func} \xdef\c@se{3} \check@sim
          \xdef\prefix@{\prfx}
       \else \xdef\c@se{1} \check@ident \fi
       \ifcase\match@case \unc@nserved \or \c@nserved \or \allm@tch \else \functi@nal \fi
      \fi
     \fi
     \ifnum\rule@num=0 
         \ifx\g@p\y@  
           \xdef\ruler@{\ruler@ -}
         \else
           \ifnum\cons@count=\rule@tens
              \expandafter\ifx\csname alt@ruler\rule@tens\endcsname\relax
                 \xdef\temp@{\rule@tens}
              \else
                 \xdef\temp@{\csname alt@ruler\rule@tens\endcsname}
              \fi
              \xdef\ruler@{\ruler@ !<\temp@>}
              \innerloopcount=\rule@tens \advance\innerloopcount by \ruler@step\relax
              \ifnum\innerloopcount=0 
                 \ifx\allow@zero\n@ \innerloopcount=\ruler@step \fi
              \fi
              \xdef\rule@tens{\the\innerloopcount}
           \else
              \xdef\ruler@{\ruler@ -}
           \fi 
           \expandafter\ifx\csname alt@ruler\the\cons@count\endcsname\relax
           \else
             \expandafter\xdef\csname alt@ruler\the\cons@count\endcsname{\the\cons@count}
           \fi
         \fi
     \fi 
    \expandafter\ifnum\csname res@count\start@seq\endcsname<\end@num\relax 
      \c@nsensus 
    \else 
      \global\stop@true
      \loopcount=\pos@count \advance\loopcount by -1 \relax
      \res@perline=\loopcount
      \iftopfeature       \xdef\bottop@{top}       \calc@feature \fi
      \ifttopfeature      \xdef\bottop@{ttop}      \calc@feature \fi
      \iftttopfeature     \xdef\bottop@{tttop}     \calc@feature \fi
      \ifttttopfeature    \xdef\bottop@{ttttop}    \calc@feature \fi
      \ifbottomfeature    \xdef\bottop@{bottom}    \calc@feature \fi
      \ifbbottomfeature   \xdef\bottop@{bbottom}   \calc@feature \fi
      \ifbbbottomfeature  \xdef\bottop@{bbbottom}  \calc@feature \fi
      \ifbbbbottomfeature \xdef\bottop@{bbbbottom} \calc@feature \fi
      \pos@count=0
    \fi
 \fi}
 
\def\prep@reexp@rtfile{
   \loopcount=0
   \loop
      \advance\loopcount by 5
      \immediate\write\exp@rtfile{\string set_color col\the\loopcount, \csname\c@nsc@l\the\loopcount\endcsname}
   \ifnum\loopcount>95\else\repeat
}

\def\c@unt{%
       \advance\loopcount by 1
       \ifT@coffee
          \xdef\seq@line{\csname T@coffee\the\loopcount\endcsname}
          \expandafter\TC@get\seq@line
       \fi
       \xdef\seq@line{\csname sequence\the\loopcount\endcsname}
       \expandafter\residue@get\seq@line
       \xdef\first@{\csname res\the\loopcount\endcsname}
       \expandafter\check@char\first@
       \ifletter
           \res@count=\csname res@count\the\loopcount\endcsname
           \advance\res@count by 1
           \ifnum\res@count=0
              \ifx\allow@zero\n@ \advance\res@count by 1 \fi
           \fi
           \ifx\dom@in\y@ 
              \xdef\temp@{\csname dom@num\the\loopcount\endcsname}
              \expandafter\get@dom@count\temp@ 
              \expandafter\xdef\csname dom@num\the\loopcount\endcsname{\temp@}
           \fi %%%%%%%%%%%*****
           \expandafter\xdef\csname res@count\the\loopcount\endcsname{\the\res@count}
           \ifnum\rule@num=\loopcount
              \temp@count=\csname res@count\the\loopcount\endcsname
              \advance\temp@count by \ruler@step  
              \divide\temp@count by \ruler@step
              \multiply\temp@count by \ruler@step
              \xdef\rule@tens{\the\temp@count}  
           \fi
       \fi
       \ifnum\loopcount<\seq@count \c@unt\fi}
\def\count@first{%
    \advance\end@count by 1
    \ifnum\end@count<\start@number
       \loopcount=0
       \ifT@coffee
          \xdef\seq@line{\csname T@coffee\the\loopcount\endcsname}
          \expandafter\TC@get\seq@line
       \fi
       \c@unt
       \count@first
    \fi}
\def\findc@nsensus{%
 \loopcount=0
 \loop
 \advance\loopcount by 1
    \expandafter\xdef\csname sequence\the\loopcount\endcsname{%
                     \csname sequence\the\loopcount\endcsname @}
 \ifnum\loopcount<\seq@count \repeat
 \end@count=0 \count@first \end@count=0 \xdef\start@number{0}
 \topfeaturenowfalse     \bottomfeaturenowfalse
 \ttopfeaturenowfalse    \bbottomfeaturenowfalse
 \tttopfeaturenowfalse   \bbbottomfeaturenowfalse
 \ttttopfeaturenowfalse  \bbbbottomfeaturenowfalse
 \innerloopcount=\cons@count
 \advance\innerloopcount by \res@perline \advance\innerloopcount by 1
 \loopcount=0
    \ifregionalshade
         \expandafter\ifx\csname start\the\loopcount\endcsname\ampers@nd
         \else
            \expandafter\ifnum\csname start\the\loopcount\endcsname>%
            \innerloopcount
            \else
                  \regionalshadenowtrue
            \fi
         \fi
    \fi
    \ifregionaltint
         \expandafter\ifx\csname tintstart\the\loopcount\endcsname\ampers@nd
         \else
            \expandafter\ifnum\csname tintstart\the\loopcount\endcsname>%
            \innerloopcount
            \else
                  \regionaltintnowtrue
            \fi
         \fi
    \fi
    \ifregionalemph
         \expandafter\ifx\csname emphstart\the\loopcount\endcsname\ampers@nd
         \else
            \expandafter\ifnum\csname emphstart\the\loopcount\endcsname>%
            \innerloopcount
            \else
                  \regionalemphnowtrue
            \fi
         \fi
    \fi
    \ifregionallower
         \expandafter\ifx\csname lowerstart\the\loopcount\endcsname\ampers@nd
         \else
            \expandafter\ifnum\csname lowerstart\the\loopcount\endcsname>%
            \innerloopcount
            \else
                  \regionallowernowtrue
            \fi
         \fi
    \fi
    \ifframe@
         \expandafter\ifx\csname framestart\the\loopcount\endcsname\ampers@nd
         \else
            \expandafter\ifnum\csname framestart\the\loopcount\endcsname>%
            \innerloopcount
            \else
               \framenowtrue
            \fi
         \fi
    \fi
    \ifshading@
         \expandafter\ifx\csname shadingstart\the\loopcount\endcsname\ampers@nd
         \else
            \expandafter\ifnum\csname shadingstart\the\loopcount\endcsname>%
            \innerloopcount
            \else
               \shadingnowtrue
            \fi
         \fi
    \fi
    \iftopfeature
       \xdef\bottop@{top}
         \expandafter\ifx\csname start\bottop@\the\loopcount\endcsname\ampers@nd
         \else
            \expandafter\ifnum\csname start\bottop@\the\loopcount\endcsname>%
            \innerloopcount
            \else
               \topfeaturenowtrue
            \fi
         \fi
    \fi
    \ifttopfeature
       \xdef\bottop@{ttop}
         \expandafter\ifx\csname start\bottop@\the\loopcount\endcsname\ampers@nd
         \else
            \expandafter\ifnum\csname start\bottop@\the\loopcount\endcsname>%
            \innerloopcount
            \else
               \ttopfeaturenowtrue
            \fi
         \fi
    \fi
    \iftttopfeature
       \xdef\bottop@{tttop}
         \expandafter\ifx\csname start\bottop@\the\loopcount\endcsname\ampers@nd
         \else
            \expandafter\ifnum\csname start\bottop@\the\loopcount\endcsname>%
            \innerloopcount
            \else
               \tttopfeaturenowtrue
            \fi
         \fi
    \fi
    \ifttttopfeature
       \xdef\bottop@{ttttop}
         \expandafter\ifx\csname start\bottop@\the\loopcount\endcsname\ampers@nd
         \else
            \expandafter\ifnum\csname start\bottop@\the\loopcount\endcsname>%
            \innerloopcount
            \else
               \ttttopfeaturenowtrue
            \fi
         \fi
    \fi
    \ifbottomfeature
       \xdef\bottop@{bottom}
         \expandafter\ifx\csname start\bottop@\the\loopcount\endcsname\ampers@nd
         \else
            \expandafter\ifnum\csname start\bottop@\the\loopcount\endcsname>%
            \innerloopcount
            \else
               \bottomfeaturenowtrue
            \fi
         \fi
    \fi
    \ifbbottomfeature
       \xdef\bottop@{bbottom}
         \expandafter\ifx\csname start\bottop@\the\loopcount\endcsname\ampers@nd
         \else
            \expandafter\ifnum\csname start\bottop@\the\loopcount\endcsname>%
            \innerloopcount
            \else
               \bbottomfeaturenowtrue
            \fi
         \fi
    \fi
    \ifbbbottomfeature
       \xdef\bottop@{bbbottom}
         \expandafter\ifx\csname start\bottop@\the\loopcount\endcsname\ampers@nd
         \else
            \expandafter\ifnum\csname start\bottop@\the\loopcount\endcsname>%
            \innerloopcount
            \else
               \bbbottomfeaturenowtrue
            \fi
         \fi
    \fi
    \ifbbbbottomfeature
       \xdef\bottop@{bbbbottom}
         \expandafter\ifx\csname start\bottop@\the\loopcount\endcsname\ampers@nd
         \else
            \expandafter\ifnum\csname start\bottop@\the\loopcount\endcsname>%
            \innerloopcount
            \else
               \bbbbottomfeaturenowtrue
            \fi
         \fi
    \fi
\loop
 \advance\loopcount by 1
    \ifnumbers@left
       \innerloopcount=\csname seq@len\the\loopcount\endcsname\relax 
       \expandafter\ifnum\csname res@count\the\loopcount\endcsname=%
                   \innerloopcount
          \res@count=\innerloopcount
       \else
          \res@count=\csname res@count\the\loopcount\endcsname
          \advance\res@count by 1
          \ifnum\res@count=0
             \ifx\allow@zero\n@ \advance\res@count by 1 \fi
          \fi
       \fi
       \ifx\dom@in\y@
            \xdef\temp@{\csname dom@num\the\loopcount\endcsname}
            \expandafter\get@dom@count\temp@
            \expandafter\xdef\csname dom@num\the\loopcount\endcsname{\the\res@count,\temp@}
       \fi         %%%%%%%%%%*****
       \expandafter\xdef\csname res@count\the\loopcount\endcsname{\the\res@count}
       \ifnames@right
          \expandafter\xdef\csname seq\the\loopcount\endcsname{%
              \csname res@count\the\loopcount\endcsname)}
       \else
          \expandafter\xdef\csname seq\the\loopcount\endcsname{%
              <\csname newseqname\the\loopcount\endcsname>
              \csname res@count\the\loopcount\endcsname)}
       \fi
       \ifx\dom@in\y@
       \else
          \expandafter\ifnum\csname res@count\the\loopcount\endcsname=%
                      \innerloopcount
          \else
             \res@count=\csname res@count\the\loopcount\endcsname
             \advance\res@count by -1
             \ifnum\res@count=0
                \ifx\allow@zero\n@ \advance\res@count by -1 \fi
             \fi
	         \expandafter\xdef\csname res@count\the\loopcount\endcsname{\the\res@count}
          \fi
       \fi
    \else
       \ifnames@right
          \expandafter\xdef\csname seq\the\loopcount\endcsname{}
       \else
          \expandafter\xdef\csname seq\the\loopcount\endcsname{%
              <\csname newseqname\the\loopcount\endcsname>}
       \fi
    \fi
    \ifx\dom@in\y@
         \xdef\temp@{\csname dom@num@break\the\loopcount\endcsname}
         \expandafter\get@dom@count\temp@
         \expandafter\xdef\csname dom@num@break\the\loopcount\endcsname{\temp@}
         \innerloopcount=\res@count
    \else        
         \innerloopcount=\csname res@count\the\loopcount\endcsname
         \advance\innerloopcount by \res@perline
    \fi
    \advance\innerloopcount by 1
    \ifregionalshade
         \expandafter\ifx\csname start\the\loopcount\endcsname\ampers@nd
         \else
            \expandafter\ifnum\csname start\the\loopcount\endcsname>%
            \innerloopcount
            \else
               \regionalshadenowtrue
            \fi
         \fi
    \fi
    \ifregionaltint
         \expandafter\ifx\csname tintstart\the\loopcount\endcsname\ampers@nd
         \else
            \expandafter\ifnum\csname tintstart\the\loopcount\endcsname>%
            \innerloopcount
            \else
                  \regionaltintnowtrue
            \fi
         \fi
    \fi
    \ifregionalemph
         \expandafter\ifx\csname emphstart\the\loopcount\endcsname\ampers@nd
         \else
            \expandafter\ifnum\csname emphstart\the\loopcount\endcsname>%
            \innerloopcount
            \else
               \regionalemphnowtrue
            \fi
         \fi
    \fi
    \ifregionallower
         \expandafter\ifx\csname lowerstart\the\loopcount\endcsname\ampers@nd
         \else
            \expandafter\ifnum\csname lowerstart\the\loopcount\endcsname>%
            \innerloopcount
            \else
               \regionallowernowtrue
            \fi
         \fi
    \fi
    \ifframe@
         \expandafter\ifx\csname framestart\the\loopcount\endcsname\ampers@nd
         \else
            \expandafter\ifnum\csname framestart\the\loopcount\endcsname>%
            \innerloopcount
            \else
               \framenowtrue
            \fi
         \fi
    \fi
    \ifshading@
         \expandafter\ifx\csname shadingstart\the\loopcount\endcsname\ampers@nd
         \else
            \expandafter\ifnum\csname shadingstart\the\loopcount\endcsname>%
            \innerloopcount
            \else
               \shadingnowtrue
            \fi
         \fi
    \fi
    \iftopfeature
       \xdef\bottop@{top}
         \expandafter\ifx\csname start\bottop@\the\loopcount\endcsname\ampers@nd
         \else
            \expandafter\ifnum\csname start\bottop@\the\loopcount\endcsname>%
            \innerloopcount
            \else
               \topfeaturenowtrue
            \fi
         \fi
    \fi
    \ifttopfeature
       \xdef\bottop@{ttop}
         \expandafter\ifx\csname start\bottop@\the\loopcount\endcsname\ampers@nd
         \else
            \expandafter\ifnum\csname start\bottop@\the\loopcount\endcsname>%
            \innerloopcount
            \else
               \ttopfeaturenowtrue
            \fi
         \fi
    \fi
    \iftttopfeature
       \xdef\bottop@{tttop}
         \expandafter\ifx\csname start\bottop@\the\loopcount\endcsname\ampers@nd
         \else
            \expandafter\ifnum\csname start\bottop@\the\loopcount\endcsname>%
            \innerloopcount
            \else
               \tttopfeaturenowtrue
            \fi
         \fi
    \fi
    \ifttttopfeature
       \xdef\bottop@{ttttop}
         \expandafter\ifx\csname start\bottop@\the\loopcount\endcsname\ampers@nd
         \else
            \expandafter\ifnum\csname start\bottop@\the\loopcount\endcsname>%
            \innerloopcount
            \else
               \ttttopfeaturenowtrue
            \fi
         \fi
    \fi
    \ifbottomfeature
       \xdef\bottop@{bottom}
         \expandafter\ifx\csname start\bottop@\the\loopcount\endcsname\ampers@nd
         \else
            \expandafter\ifnum\csname start\bottop@\the\loopcount\endcsname>%
            \innerloopcount
            \else
               \bottomfeaturenowtrue
            \fi
         \fi
    \fi
    \ifbbottomfeature
       \xdef\bottop@{bbottom}
         \expandafter\ifx\csname start\bottop@\the\loopcount\endcsname\ampers@nd
         \else
            \expandafter\ifnum\csname start\bottop@\the\loopcount\endcsname>%
            \innerloopcount
            \else
               \bbottomfeaturenowtrue
            \fi
         \fi
    \fi
    \ifbbbottomfeature
       \xdef\bottop@{bbbottom}
         \expandafter\ifx\csname start\bottop@\the\loopcount\endcsname\ampers@nd
         \else
            \expandafter\ifnum\csname start\bottop@\the\loopcount\endcsname>%
            \innerloopcount
            \else
               \bbbottomfeaturenowtrue
            \fi
         \fi
    \fi
    \ifbbbbottomfeature
       \xdef\bottop@{bbbbottom}
         \expandafter\ifx\csname start\bottop@\the\loopcount\endcsname\ampers@nd
         \else
            \expandafter\ifnum\csname start\bottop@\the\loopcount\endcsname>%
            \innerloopcount
            \else
               \bbbbottomfeaturenowtrue
            \fi
         \fi
    \fi
\ifnum\loopcount<\seq@count \repeat
 \c@nsensus}

%%%%%  Output routines

\def\white@box{%
   \bgroup
   \fboxsep-0.5pt\fboxrule0.5pt
   \fcolorbox{Black}{White}{\box@hstrut\box@wstrut}\egroup}
\def\box@rule{\vrule depth\box@depth height\box@height width\box@width}
\def\box@hstrut{\vrule depth\box@depth height\box@height width 0pt}
\def\box@wstrut{\vrule depth 0pt height 0pt width\box@width}
\def\do@legend{%
      \baselineskip=1.2\baselineskip
      \xdef\first@{White}
      \fontfamily{\legend@family}%
      \fontseries{\legend@series}%
      \fontshape{\legend@shape}%
      \ifT@coffee
         \noindent%
         \ifnames@\ifnames@right\else\hbox to \name@width{\hss}\fi\fi%
         \ifnumbers@\ifnumbers@left\hbox to \number@width{\hss}\fi\fi%
	     \hbox to \hspace@legend{\hss}%
         \def\res@@style{\csname no@style\endcsname}%
         \def\temp@{X}\xdef\low@up{lower}%
         \expandafter\ifx\csname n@m@tch\the\loopcount\endcsname\low@up%
         \def\temp@{x}\fi%
         \ifhidechar\xdef\temp@{}\fi%
         \textcolor{TC0}{\box@rule}%
         \kern-\box@width\textcolor{Black}{\hbox to \box@width{%
               \legend@size{\res@@style{\hss\temp@\hss}}}}%
         \kern2ex\textcolor{\legend@fg}{%
         \ifgerm@n\kern2ex\legend@size{schlecht}%
         \else\ifsp@nish\kern2ex\legend@size{mala}%
         \else\kern2ex\legend@size{bad}\fi\fi}
         \newline\hbox{}%
         \noindent%
         \ifnames@\ifnames@right\else\hbox to \name@width{\hss}\fi\fi%
         \ifnumbers@\ifnumbers@left\hbox to \number@width{\hss}\fi\fi%
	     \hbox to \hspace@legend{\hss}%
         \textcolor{TC1}{\box@rule}%
         \kern-\box@width\textcolor{Black}{\hbox to \box@width{%
               \legend@size{\res@@style{\hss\temp@\hss}}}}%
         \kern2ex\textcolor{\legend@fg}{\legend@size{}}
         \newline\hbox{}%
         \noindent%
         \ifnames@\ifnames@right\else\hbox to \name@width{\hss}\fi\fi%
         \ifnumbers@\ifnumbers@left\hbox to \number@width{\hss}\fi\fi%
	     \hbox to \hspace@legend{\hss}%
         \textcolor{TC2}{\box@rule}%
         \kern-\box@width\textcolor{Black}{\hbox to \box@width{%
               \legend@size{\res@@style{\hss\temp@\hss}}}}%
         \kern2ex\textcolor{\legend@fg}{\legend@size{}}
         \newline\hbox{}%
         \noindent%
         \ifnames@\ifnames@right\else\hbox to \name@width{\hss}\fi\fi%
         \ifnumbers@\ifnumbers@left\hbox to \number@width{\hss}\fi\fi%
	     \hbox to \hspace@legend{\hss}%
         \textcolor{TC3}{\box@rule}%
         \kern-\box@width\textcolor{Black}{\hbox to \box@width{%
               \legend@size{\res@@style{\hss\temp@\hss}}}}%
         \kern2ex\textcolor{\legend@fg}{\legend@size{}}
         \newline\hbox{}%
         \noindent%
         \ifnames@\ifnames@right\else\hbox to \name@width{\hss}\fi\fi%
         \ifnumbers@\ifnumbers@left\hbox to \number@width{\hss}\fi\fi%
	     \hbox to \hspace@legend{\hss}%
         \textcolor{TC4}{\box@rule}%
         \kern-\box@width\textcolor{Black}{\hbox to \box@width{%
               \legend@size{\res@@style{\hss\temp@\hss}}}}%
         \kern2ex\textcolor{\legend@fg}{%
         \ifgerm@n\kern2ex\legend@size{Grad der}%
         \else\ifsp@nish\kern2ex\legend@size{conservaci\'on}%
         \else\kern2ex\legend@size{level of}\fi\fi}
         \newline\hbox{}%
         \noindent%
         \ifnames@\ifnames@right\else\hbox to \name@width{\hss}\fi\fi%
         \ifnumbers@\ifnumbers@left\hbox to \number@width{\hss}\fi\fi%
	     \hbox to \hspace@legend{\hss}%
         \textcolor{TC5}{\box@rule}%
         \kern-\box@width\textcolor{Black}{\hbox to \box@width{%
               \legend@size{\res@@style{\hss\temp@\hss}}}}%
         \kern2ex\textcolor{\legend@fg}{%
         \ifgerm@n\kern2ex\legend@size{Konservierung}%
         \else\ifsp@nish\kern2ex\legend@size{}%
         \else\kern2ex\legend@size{conservation}\fi\fi}
         \newline\hbox{}%
         \noindent%
         \ifnames@\ifnames@right\else\hbox to \name@width{\hss}\fi\fi%
         \ifnumbers@\ifnumbers@left\hbox to \number@width{\hss}\fi\fi%
	     \hbox to \hspace@legend{\hss}%
         \textcolor{TC6}{\box@rule}%
         \kern-\box@width\textcolor{Black}{\hbox to \box@width{%
               \legend@size{\res@@style{\hss\temp@\hss}}}}%
         \kern2ex\textcolor{\legend@fg}{\legend@size{}}
         \newline\hbox{}%
         \noindent%
         \ifnames@\ifnames@right\else\hbox to \name@width{\hss}\fi\fi%
         \ifnumbers@\ifnumbers@left\hbox to \number@width{\hss}\fi\fi%
	     \hbox to \hspace@legend{\hss}%
         \textcolor{TC7}{\box@rule}%
         \kern-\box@width\textcolor{Black}{\hbox to \box@width{%
               \legend@size{\res@@style{\hss\temp@\hss}}}}%
         \kern2ex\textcolor{\legend@fg}{\legend@size{}}
         \newline\hbox{}%
         \noindent%
         \ifnames@\ifnames@right\else\hbox to \name@width{\hss}\fi\fi%
         \ifnumbers@\ifnumbers@left\hbox to \number@width{\hss}\fi\fi%
	     \hbox to \hspace@legend{\hss}%
         \textcolor{TC8}{\box@rule}%
         \kern-\box@width\textcolor{Black}{\hbox to \box@width{%
               \legend@size{\res@@style{\hss\temp@\hss}}}}%
         \kern2ex\textcolor{\legend@fg}{\legend@size{}}
         \newline\hbox{}%
         \noindent%
         \ifnames@\ifnames@right\else\hbox to \name@width{\hss}\fi\fi%
         \ifnumbers@\ifnumbers@left\hbox to \number@width{\hss}\fi\fi%
	     \hbox to \hspace@legend{\hss}%
         \textcolor{TC9}{\box@rule}%
         \kern-\box@width\textcolor{Black}{\hbox to \box@width{%
               \legend@size{\res@@style{\hss\temp@\hss}}}}%
         \kern2ex\textcolor{\legend@fg}{%
         \ifgerm@n\kern2ex\legend@size{gut}%
         \else\ifsp@nish\kern2ex\legend@size{buena}%
         \else\kern2ex\legend@size{good}\fi\fi}
         \newline\hbox{}%
      \else
      \iffuncmode
        \ifnum\fgroup@num>0
         \loopcount=0
         \loop
         \advance\loopcount by 1
         \xdef\third@{\csname fg@color\the\loopcount\endcsname}%
         \noindent%
         \ifnames@\ifnames@right\else\hbox to \name@width{\hss}\fi\fi%
         \ifnumbers@\ifnumbers@left\hbox to \number@width{\hss}\fi\fi%
	 \hbox to \hspace@legend{\hss}%
         \ifx\third@\first@\white@box\else\textcolor{\third@}{\box@rule}\fi%
         \xdef\third@{\csname fg@textcolor\the\loopcount\endcsname}%
         \def\res@@style{\csname func@style\the\loopcount\endcsname}%
         \def\temp@{X}\xdef\low@up{lower}%
         \expandafter\ifx\csname funcm@tch\the\loopcount\endcsname\low@up%
         \def\temp@{x}\fi%
         \ifhidechar\xdef\temp@{}\fi%
         \kern-\box@width\textcolor{\third@}{\hbox to \box@width{%
               \legend@size{\res@@style{\hss\temp@\hss}}}}%
         \kern2ex\textcolor{\legend@fg}{%
         \legend@size{\csname fgroup@name\the\loopcount\endcsname}}
         \newline\hbox{}%
         \ifnum\loopcount<\fgroup@num \repeat
        \fi
      \else
         \noindent
         \ifnames@\ifnames@right\else\hbox to \name@width{\hss}\fi\fi%
         \ifnumbers@\ifnumbers@left \hbox to \number@width{\hss}\fi\fi%
	 \hbox to \hspace@legend{\hss}%
         \ifx\Nomatch\first@\white@box\else\textcolor{\Nomatch}{\box@rule}\fi%
         \def\res@@style{\csname no@style\endcsname}%
         \def\temp@{X}\xdef\low@up{lower}\ifx\resn@m@tch\low@up\def\temp@{x}\fi%
         \ifhidechar\xdef\temp@{}\fi%
         \kern-\box@width\textcolor{\TextNomatch}{\hbox to \box@width{%
               \legend@size{\res@@style{\hss\temp@\hss}}}}%
         \textcolor{\legend@fg}{%
         \ifgerm@n\kern2ex\legend@size{nicht konserviert}%
         \else\ifsp@nish\kern2ex\legend@size{no conservado}%
         \else\kern2ex\legend@size{non-conserved}\fi\fi}
         \newline\hbox{}\noindent%
         \ifsimmode%
         \ifnames@\ifnames@right\else\hbox to \name@width{\hss}\fi\fi
         \ifnumbers@\ifnumbers@left \hbox to \number@width{\hss}\fi\fi
	 \hbox to \hspace@legend{\hss}%
         \ifx\Similar\first@\white@box\else\textcolor{\Similar}{\box@rule}\fi
         \def\res@@style{\csname sim@style\endcsname}%
         \def\temp@{X}\xdef\low@up{lower}\ifx\ressimm@tch\low@up\def\temp@{x}\fi%
         \ifhidechar\xdef\temp@{}\fi%
         \kern-\box@width\textcolor{\TextSimilar}{\hbox to \box@width{%
               \legend@size{\res@@style{\hss\temp@\hss}}}}%
         \textcolor{\legend@fg}{%
         \ifgerm@n\kern2ex\legend@size{\"ahnlich}%
         \else\ifsp@nish\kern2ex\legend@size{similar}%
         \else\kern2ex\legend@size{similar}\fi\fi}
         \newline\hbox{}\noindent%
         \fi%
         \ifnames@\ifnames@right\else\hbox to \name@width{\hss}\fi\fi
         \ifnumbers@\ifnumbers@left \hbox to \number@width{\hss}\fi\fi
	 \hbox to \hspace@legend{\hss}%
         \ifx\Identical\first@\white@box\else\textcolor{\Identical}{\box@rule}\fi
         \def\res@@style{\csname id@style\endcsname}%
         \def\temp@{X}\xdef\low@up{lower}\ifx\resm@tch\low@up\def\temp@{x}\fi%
         \ifhidechar\xdef\temp@{}\fi%
         \kern-\box@width\textcolor{\TextIdentical}{\hbox to \box@width{%
               \legend@size{\res@@style{\hss\temp@\hss}}}}%
         \textcolor{\legend@fg}{%
         \ifgerm@n\kern2ex\legend@size{$\ge$\,\thresh@ld\% konserviert}%
         \else\ifsp@nish\kern2ex\legend@size{$\ge$\,\thresh@ld\% conservado}%
         \else\kern2ex\legend@size{$\ge$\,\thresh@ld\% conserved}\fi\fi}
         \newline\hbox{}\noindent%
         \ifall@shade%
         \ifnames@\ifnames@right\else\hbox to \name@width{\hss}\fi\fi
         \ifnumbers@\ifnumbers@left \hbox to \number@width{\hss}\fi\fi
	 \hbox to \hspace@legend{\hss}%
         \ifx\Allmatch\first@\white@box\else\textcolor{\Allmatch}{\box@rule}\fi
         \def\res@@style{\csname all@style\endcsname}%
         \def\temp@{X}\xdef\low@up{lower}\ifx\res@llm@tch\low@up\def\temp@{x}\fi%
         \ifhidechar\xdef\temp@{}\fi%
         \kern-\box@width\textcolor{\TextAllmatch}{\hbox to \box@width{%
               \legend@size{\res@@style{\hss\temp@\hss}}}}%
         \ifnum\all@thresh@ld=100
         \textcolor{\legend@fg}{%
         \ifgerm@n\kern2ex\legend@size{alle identisch}%
         \else\ifsp@nish\kern2ex\legend@size{todos id\'enticos}%
         \else\kern2ex\legend@size{all match}\fi\fi}
         \else
         \textcolor{\legend@fg}{%
         \ifgerm@n\kern2ex\legend@size{$\ge$\,\all@thresh@ld\% konserviert}%
         \else\ifsp@nish\kern2ex\legend@size{$\ge$\,\all@thresh@ld\% conservado}%
         \else\kern2ex\legend@size{$\ge$\,\all@thresh@ld\% conserved}\fi\fi}
         \fi
         \newline\hbox{}\noindent
         \fi
      \fi
      \fi
}
\def\put@name<#1>#2@{%
   \ifnames@%
   \expandafter\ifx\csname hide@name\the\loopcount\endcsname\yes%
      \xdef\temp@{}\else\xdef\temp@{#1}\fi%
   \expandafter\ifx\csname name@col\the\loopcount\endcsname\yes%
           \def\second@{\names@fg}%
      \else\def\second@{\csname name@col\the\loopcount\endcsname}\fi%
   \fontfamily{\namestext@family}%
   \fontseries{\namestext@series}%
   \fontshape{\namestext@shape}%
   \selectfont%
   \textcolor{\second@}{%
   \hbox to \name@width{\@kern\namestext@size{\temp@}\hss}}\fi%
   \xdef\first@{#2@}%
}
\def\put@number#1)#2@{%
   \ifnumbers@%
   \expandafter\ifx\csname hide@number\the\loopcount\endcsname\yes%
      \xdef\temp@{}\else\xdef\temp@{#1}\fi%
   \expandafter\ifx\csname number@col\the\loopcount\endcsname\yes%
           \def\second@{\numbering@fg}%
      \else\def\second@{\csname number@col\the\loopcount\endcsname}\fi%
   \fontfamily{\numbertext@family}%
   \fontseries{\numbertext@series}%
   \fontshape{\numbertext@shape}%
   \selectfont%
   \textcolor{\second@}{%
   \hbox to \number@width{\hss\numbertext@size{\temp@}\@kern}}\fi%
   \xdef\first@{#2@}%
}
\def\special@shading#1)#2#3#4@{%
   \xdef\second@{@\second@#1}%
   \xdef\third@{\second@}%
   \xdef\first@{#4@}%
   \xdef\second@{#3}%
   \xdef\last@{#2}%
    \def\res@@style{\csname relax\endcsname}%
}
\def\special@shade#1)#2#3#4@{%
   \xdef\second@{\second@#1}%
   \xdef\boxc@l@r{\csname bgseqregion\second@\endcsname}%
   \xdef\textc@l@r{\csname fgseqregion\second@\endcsname}%
   \xdef\first@{#4@}%
   \xdef\second@{#3}%
    \def\res@@style{\csname relax\endcsname}%
}
\def\get@second@#1#2@{\xdef\second@{#1}\xdef\first@{#2@}}
\def\next@char#1#2#3@{%
   \xdef\first@{#3@}%
   \xdef\second@{#2}%
   \xdef\last@{#1}%
   \xdef\temp@@{}%
   \xdef\third@{}%
   \ifx\last@\p@r@gr@ph\expandafter\special@shading\first@\fi%
   \ifx\last@\ampers@nd\def\last@{0}\expandafter\special@shade\first@%
   \else%
   \ifT@coffee%
         \xdef\boxc@l@r{\csname fg@color#1\endcsname}%
         \xdef\textc@l@r{\csname fg@textcolor#1\endcsname}%
         \def\res@@style{\csname func@style#1\endcsname}%
         \if\last@ *\def\last@{0}\fi
         \if\last@ /\def\last@{10}\fi
         \if\last@ !\def\last@{11}\fi
   \else%
   \iffuncmode%
         \xdef\boxc@l@r{\csname fg@color#1\endcsname}%
         \xdef\textc@l@r{\csname fg@textcolor#1\endcsname}%
         \def\res@@style{\csname func@style#1\endcsname}%
         \if\last@ *\def\last@{0}\fi
         \if\last@ /\def\last@{10}\fi
         \if\last@ !\def\last@{11}\fi
   \else%
   \ifcase\last@\xdef\boxc@l@r{\csname Allmatch\third@\endcsname}%
             \xdef\textc@l@r{\csname TextAllmatch\third@\endcsname}%
             \def\res@@style{\all@style}%
         \or\xdef\boxc@l@r{\csname Identical\third@\endcsname}%
            \xdef\textc@l@r{\csname TextIdentical\third@\endcsname}%
             \def\res@@style{\id@style}%
         \or\xdef\boxc@l@r{\csname Similar\third@\endcsname}%
            \xdef\textc@l@r{\csname TextSimilar\third@\endcsname}%
             \def\res@@style{\sim@style}%
         \or\xdef\boxc@l@r{\csname Nomatch\third@\endcsname}%
            \xdef\textc@l@r{\csname TextNomatch\third@\endcsname}%
             \def\res@@style{\no@style}%
         \or\xdef\boxc@l@r{\csname ConsNomatch\third@\endcsname}%
            \xdef\textc@l@r{\csname ConsTextNomatch\third@\endcsname}%
             \def\res@@style{\csname relax\endcsname}%
         \or\xdef\boxc@l@r{\csname ConsMatch\third@\endcsname}%
            \xdef\textc@l@r{\csnam ConsTextMatch\third@\endcsname}%
             \def\res@@style{\csname relax\endcsname}%
         \or\xdef\boxc@l@r{\csname ConsAllmatch\third@\endcsname}%
            \xdef\textc@l@r{\csname ConsTextAllmatch\third@\endcsname}%
             \def\res@@style{\csname relax\endcsname}%
         \or\xdef\boxc@l@r{\csname gap@bg\third@\endcsname}%
            \xdef\textc@l@r{\csname gap@fg\third@\endcsname}%
             \def\res@@style{\csname relax\endcsname}%
         \or\xdef\boxc@l@r{White}\xdef\textc@l@r{White}%
             \def\res@@style{\csname relax\endcsname}%
         \else\xdef\boxc@l@r{\csname domgap@bg\endcsname}%
              \xdef\textc@l@r{\csname domgap@fg\endcsname}%
             \def\res@@style{\csname relax\endcsname}%
   \fi\fi\fi\fi%
   \ifx\second@\comm@%
         \def\temp@{\res@style}%
         \expandafter\get@second@\first@%
   \else%
         \def\temp@{\csname relax\endcsname}%
   \fi%
   \ifx\second@\equ@l%
         \xdef\temp@@{\light@}%
         \expandafter\get@second@\first@%
   \fi%
   \ifx\second@\semic@n%
         \expandafter\get@second@\first@%
         \xdef\first@@{\first@}%
         \xdef\first@{\second@}%
         \make@lower%
         \xdef\second@{\first@}%
         \xdef\first@{\first@@}%
   \fi%
   \textcolor{\temp@@\boxc@l@r}{\box@rule}%
   \ifhidechar%
      \ifx\second@\o@%
         \def\second@{\gap@rule}%
         \hbox to -\box@width{\hss\textcolor{\temp@@\textc@l@r}%
                             {\residues@size{\res@@style{\temp@{\second@}}}}\hss}%
                              \kern\box@width%
      \fi%
   \else%
   \ifx\second@\o@\def\second@{\gap@rule}\fi%
   \hbox to -\box@width{\hss\textcolor{\temp@@\textc@l@r}%
                       {\residues@size{\res@@style{\temp@{\second@}}}}\hss}%
                        \kern\box@width\fi%
}
\def\put@char{%
   \ifnum\innerloopcount>\res@perline
   \else
      \expandafter\next@char\first@
      \advance\innerloopcount by 1
      \put@char%
   \fi}
\def\next@cons#1#2#3@{%
   \xdef\last@{#1}%
   \xdef\third@{}%
   \ifx\last@\p@r@gr@ph\expandafter\special@shading\first@\fi%
   \ifx\last@\ampers@nd\def\last@{0}\expandafter\special@shade\first@%
   \else%
   \ifx\collect@cons@colors\y@%
      \if\last@ *%
         \xdef\boxc@l@r{\gap@bg}\xdef\textc@l@r{\gap@bg}%
      \else%
      \if\last@ !%
         \xdef\boxc@l@r{\domgap@bg}\xdef\textc@l@r{\domgap@bg}%
      \else%
      \if\last@ /%
         \xdef\boxc@l@r{White}\xdef\textc@l@r{White}%
      \else%
      \expandafter\get@item\first@@@%
      \xdef\first@@@{\first@}%
      \ifx\box@scale\y@%
         \xdef\boxc@l@r{\c@nssc@le\fourth@}%
      \else%
         \xdef\first@{T-Coffee}%
         \ifx\first@\c@nssc@le%
            \xdef\boxc@l@r{TC\last@}%         
         \else
            \xdef\boxc@l@r{\c@nssc@le}%
         \fi%
      \fi%
      \ifx\text@scale\y@%
         \xdef\textc@l@r{\c@nsc@l\fourth@}%
      \else%
         \xdef\first@{T-Coffee}%
         \ifx\first@\c@nsc@l%
            \xdef\textc@l@r{TC\last@}%         
         \else
            \xdef\textc@l@r{\c@nsc@l}%
         \fi%
      \fi\fi\fi\fi%
      \def\res@@style{\csname relax\endcsname}%
   \else%
   \ifT@coffee%
      \if\last@ *%
         \xdef\boxc@l@r{\gap@bg}\xdef\textc@l@r{\gap@bg}%
      \else%
         \if\last@ !%
            \xdef\boxc@l@r{\domgap@bg}\xdef\textc@l@r{\domgap@bg}%
         \else%
            \if\last@ /%
               \xdef\boxc@l@r{White}\xdef\textc@l@r{White}%
            \else%
               \xdef\boxc@l@r{White}\xdef\textc@l@r{TC\last@}%
               \def\res@@style{\csname relax\endcsname}%
            \fi
         \fi%
      \fi%
   \else%
   \ifcase#1\xdef\boxc@l@r{\csname Allmatch\third@\endcsname}%
            \xdef\textc@l@r{\csname TextAllmatch\third@\endcsname}%
             \def\res@@style{\all@style}%
         \or\xdef\boxc@l@r{\csname Identical\third@\endcsname}%
            \xdef\textc@l@r{\csname TextIdentical\third@\endcsname}%
             \def\res@@style{\id@style}%
         \or\xdef\boxc@l@r{\csname Similar\third@\endcsname}%
            \xdef\textc@l@r{\csname TextSimilar\third@\endcsname}%
             \def\res@@style{\sim@style}%
         \or\xdef\boxc@l@r{\csname Nomatch\third@\endcsname}%
            \xdef\textc@l@r{\csname TextNomatch\third@\endcsname}%
             \def\res@@style{\no@style}%
         \or\xdef\boxc@l@r{\ConsNomatch}\xdef\textc@l@r{\ConsTextNomatch}%
             \def\res@@style{\csname relax\endcsname}%
         \or\xdef\boxc@l@r{\ConsMatch}\xdef\textc@l@r{\ConsTextMatch}%
             \def\res@@style{\csname relax\endcsname}%
         \or\xdef\boxc@l@r{\ConsAllmatch}\xdef\textc@l@r{\ConsTextAllmatch}%
             \def\res@@style{\csname relax\endcsname}%
         \or\xdef\boxc@l@r{\gap@bg}\xdef\textc@l@r{\gap@fg}%
             \def\res@@style{\csname relax\endcsname}%
         \or\xdef\boxc@l@r{White}\xdef\textc@l@r{White}%
             \def\res@@style{\csname relax\endcsname}%
         \else\xdef\boxc@l@r{\domgap@bg}\xdef\textc@l@r{\domgap@fg}%
             \def\res@@style{\csname relax\endcsname}%
   \fi\fi\fi\fi%
   \xdef\first@{#3@}%
   \xdef\second@{#2}%
   \ifx\second@\comm@%
         \def\temp@{\res@style}%
         \expandafter\get@second@\first@%
   \else%
         \def\temp@{\csname relax\endcsname}%
   \fi%
   \ifx\second@\semic@n%
         \expandafter\get@second@\first@%
         \xdef\first@@{\first@}%
         \xdef\first@{\second@}%
         \make@lower%
         \xdef\second@{\first@}%
         \xdef\first@{\first@@}%
   \fi%
   \textcolor{\boxc@l@r}{\box@rule}%
   \ifhidechar%
      \ifx\second@\o@%
         \def\second@{\gap@rule}%
         \hbox to -\box@width{\hss\textcolor{\textc@l@r}%
                             {\residues@size{\res@@style{\temp@{\second@}}}}\hss}%
                              \kern\box@width%
      \fi%
   \else%
   \ifx\second@\o@\def\second@{\gap@rule}\fi%
   \hbox to -\box@width{\hss\textcolor{\textc@l@r}%
                       {\residues@size{\res@@style{\temp@{\second@}}}}\hss}%
                        \kern\box@width\fi%
}
\def\put@cons{%
   \ifnum\innerloopcount>\res@perline
   \else
      \expandafter\next@cons\first@
      \advance\innerloopcount by 1
      \put@cons%
   \fi}
\def\put@line{%
   \ifnames@right\else\def\@kern{\kern0em}\expandafter\put@name\first@\fi
   \ifnumbers@left\def\@kern{\kern1em}\expandafter\put@number\first@\fi
   \fontfamily{\residues@family}%
   \fontseries{\residues@series}%
   \fontshape{\residues@shape}%
   \selectfont%
   \ifx\cons@now\y@%
      \innerloopcount=1\relax\put@cons%
   \else
      \innerloopcount=1\relax\put@char%
   \fi
   \ifnumbers@right\def\@kern{\kern0em}\expandafter\put@number\first@\fi
   \ifnames@right\def\@kern{\kern1em}\expandafter\put@name\first@\fi
   \newline\hbox{}%
}

\def\set@consensus{%
      \ifnames@right
        \ifnumbers@left
          \ifnumbers@right
           \xdef\consensus{)\consensus)<\cons@name>}%
          \else
           \xdef\consensus{)\consensus<\cons@name>}%
          \fi
        \else
           \xdef\consensus{\consensus)<\cons@name>}%
        \fi
      \else
        \ifnumbers@left
          \ifnumbers@right
           \xdef\consensus{<\cons@name>)\consensus)}%
          \else
           \xdef\consensus{<\cons@name>)\consensus}%
          \fi
        \else
           \xdef\consensus{<\cons@name>\consensus)}%
        \fi
      \fi}

\def\get@rulenum<#1>#2@{%
   \xdef\first@{#2@}%
   \xdef\fill@char{#1[,]&}%
   \expandafter\opt@color\fill@char%
   \ifx\f@color\comm@\xdef\f@color{\ruler@fg}\fi%
   \ifcase\rule@top
     \ifnum\ruler@rot=0 %
      \xdef\temp@{tt}%
      \ifx\ruler@family\temp@%
         \def\third@{\bottomruler@size\ensuremath{\,\stackrel{\mathtt{\fourth@}}{\textcolor{\ruler@fg}{.}}}}%
      \else
      \xdef\temp@{sf}%
      \ifx\ruler@family\temp@%
         \def\third@{\bottomruler@size\ensuremath{\,\stackrel{\mathsf{\fourth@}}{\textcolor{\ruler@fg}{.}}}}%
      \else
      \xdef\temp@{rm}%
      \ifx\ruler@family\temp@%
         \def\third@{\bottomruler@size\ensuremath{\,\stackrel{\mathrm{\fourth@}}{\textcolor{\ruler@fg}{.}}}}%
      \fi\fi\fi
     \else
      \xdef\temp@{tt}%
      \ifx\ruler@family\temp@%
         \def\third@{\tt\bottomruler@size\,\,\,\,%
           \begin{rotopo}{90}\hbox to \ruler@width{\,\textcolor{\ruler@fg}{\ensuremath{\cdot}}\,\fourth@\hss}\end{rotopo}}%
      \else
      \xdef\temp@{sf}%
      \ifx\ruler@family\temp@%
         \def\third@{\sf\bottomruler@size\,\,\,\,%
           \begin{rotopo}{90}\hbox to \ruler@width{\,\textcolor{\ruler@fg}{\ensuremath{\cdot}}\,\fourth@\hss}\end{rotopo}}%
      \else
      \xdef\temp@{rm}%
      \ifx\ruler@family\temp@%
         \def\third@{\rm\bottomruler@size\,\,\,\,%
           \begin{rotopo}{90}\hbox to \ruler@width{\,\textcolor{\ruler@fg}{\ensuremath{\cdot}}\,\fourth@\hss}\end{rotopo}}%
      \fi\fi\fi
     \fi
   \else
     \ifnum\ruler@rot=0 %
      \xdef\temp@{tt}%
      \ifx\ruler@family\temp@%
         \def\third@{\bottomruler@size\ensuremath{\,\stackrel{\textcolor{\ruler@fg}{.}}{\mathtt{\fourth@}}}}%
      \else
      \xdef\temp@{sf}%
      \ifx\ruler@family\temp@%
         \def\third@{\bottomruler@size\ensuremath{\,\stackrel{\textcolor{\ruler@fg}{.}}{\mathsf{\fourth@}}}}%
      \else
      \xdef\temp@{rm}%
      \ifx\ruler@family\temp@%
         \def\third@{\bottomruler@size\ensuremath{\,\stackrel{\textcolor{\ruler@fg}{.}}{\mathrm{\fourth@}}}}%
      \fi\fi\fi
     \else
      \xdef\temp@{tt}%
      \ifx\ruler@family\temp@%
         \def\third@{\tt\bottomruler@size\,\,\,\,%
           \begin{rotopo}{90}\hbox to \ruler@width{\hss\fourth@\,\textcolor{\ruler@fg}{\ensuremath{\cdot}}}\end{rotopo}}%
      \else
      \xdef\temp@{sf}%
      \ifx\ruler@family\temp@%
         \def\third@{\sf\bottomruler@size\,\,\,\,%
           \begin{rotopo}{90}\hbox to \ruler@width{\hss\fourth@\,\textcolor{\ruler@fg}{\ensuremath{\cdot}}}\end{rotopo}}%
      \else
      \xdef\temp@{rm}%
      \ifx\ruler@family\temp@%
         \def\third@{\rm\bottomruler@size\,\,\,\,%
           \begin{rotopo}{90}\hbox to \ruler@width{\hss\fourth@\,\textcolor{\ruler@fg}{\ensuremath{\cdot}}}\end{rotopo}}%
      \fi\fi\fi
     \fi
   \fi}

\def\next@rulechar#1#2@{%
   \xdef\third@{#1}%
   \xdef\first@{#2@}%
   \xdef\second@{!}%
   \ifx\third@\second@ \expandafter\get@rulenum\first@%
   \else \xdef\third@{}\xdef\f@color{Black}\fi
   \textcolor{\f@color}{\hbox to \box@width{\hss\third@\hss}}%
   }

\def\put@rulechar{%
   \ifnum\innerloopcount>\res@perline
   \else
      \expandafter\next@rulechar\first@
      \advance\innerloopcount by 1
      \put@rulechar%
   \fi}

\def\put@ruler{%
   \ifnames@right%
        \ifnumbers@left%
          \ifnumbers@right%
            \xdef\ruler@{)\ruler@)<>}%
          \else%
            \xdef\ruler@{)\ruler@<>}%
          \fi%
        \else%
            \xdef\ruler@{\ruler@)<>}%
        \fi%
   \else%
        \ifnumbers@left%
          \ifnumbers@right%
            \xdef\ruler@{<>)\ruler@)}%
          \else%
            \xdef\ruler@{<>)\ruler@}%
          \fi%
        \else%
            \xdef\ruler@{<>\ruler@)}%
        \fi%
   \fi%
   \xdef\first@{\ruler@ @}%
   \ifnum\rule@top=1%
      \vspace{\ruler@sp@ce}%
      \vspace{-\baselineskip}%
      \newline\hbox{}%
   \fi%
   \ifnames@right\else\def\@kern{\kern0em}\expandafter\put@name\first@\fi%
   \ifnumbers@left\def\@kern{\kern1em}\expandafter\put@number\first@\fi%
   \vspace{-0.25\baselineskip}%
   \ifnum\rule@top=0%
      \vspace{\ruler@sp@ce}%
   \fi%
   \fontfamily{\ruler@family}%
   \fontseries{m}%
   \fontshape{n}%
   \selectfont%
   \innerloopcount=1\relax\put@rulechar%
   \newline\hbox{}%
}
\def\get@firstfill#1#2&{\xdef\second@@{#1}\xdef\fill@char{#2&}}
\def\get@firstv@l#1,#2&{\xdef\second@@{#1}\xdef\fill@char{#2&}}
\def\get@tripletfill#1#2#3#4&{%
   \multiply\temp@count by -1%
   \def\second@@{#1}\def\second@@@{#2}\def\second@@@@{#3}\def\fill@char{#4&}}


\def\putfeature@style#1{%
   \residues@size%
   \setbox1=\hbox{\ensuremath{\overrightarrow{\hbox{}}}}%
   \arrow@height=\ht1%
   \arrow@width=\wd1%
   \xdef\second@@{#1}%
   \xdef\last@{\second@@::&}\expandafter\test@fill\last@%
   \xdef\last@{empty}%
   \ifx\second@@\last@%
      \hbox to \second@\box@width{\hss}%
   \else%
    \xdef\last@{translate}%
    \ifx\second@@\last@%
     \fontfamily{\featurestyles@family}%
     \fontseries{\featurestyles@series}%
     \fontshape{\featurestyles@shape}%
     \selectfont%
     \xdef\fill@char{\fill@char &}%
     \if\seq@type N%
        \loop%
            \expandafter\get@firstfill\fill@char%
            \if\second@@ -\def\second@@{\hss}\fi%
            \hbox to \box@width{\hss\textcolor{\f@color}{\featurestyles@size{\second@@}}\hss}%
        \ifx\fill@char\ampers@nd\else\repeat%
     \else%
     \temp@count=1%
     \loop%
      \expandafter\get@tripletfill\fill@char%
      \if\second@@ -\ifnum\tr@nsstyle>0\hbox to \box@width{\hss}\fi%
      \else%
        \ifcase\tr@nsstyle%
           \hbox{\trans@size\textcolor{\f@color}%
                {\hss\second@@\second@@@\second@@@@\,\hss}}%
        \or%
           \ifnum\temp@count=1%
              \vbox{\trans@size%
              \hbox to \box@width{\textcolor{\f@color}{\second@@\hss}}%
              \hbox to \box@width{\textcolor{\f@color}{\hss\second@@@\hss}}%
              \hbox to \box@width{\textcolor{\f@color}{\hss\second@@@@}}}%
           \else%
              \vbox{\trans@size%
              \hbox to \box@width{\textcolor{\f@color}{\hss\second@@@@}}%
              \hbox to \box@width{\textcolor{\f@color}{\hss\second@@@\hss}}%
              \hbox to \box@width{\textcolor{\f@color}{\second@@\hss}}}%
           \fi%
        \or%
           \ifnum\temp@count=1%
              \vbox{\trans@size%
              \hbox to \box@width{\textcolor{\f@color}%
                                 {\hss\second@@\second@@@\second@@@@\hss}}%
              \hbox to \box@width{\textcolor{\f@color}{\hss}}}%
           \else%
              \vbox{\trans@size%
              \hbox to \box@width{\textcolor{\f@color}{\hss}}%
              \hbox to \box@width{\textcolor{\f@color}%
                                 {\hss\second@@\second@@@\second@@@@\hss}}}%
           \fi%
        \or%
           \vbox{\trans@size%
           \hbox to \box@width{\textcolor{\f@color}{\second@@\hss}}%
           \hbox to \box@width{\textcolor{\f@color}{\hss\second@@@\hss}}%
           \hbox to \box@width{\textcolor{\f@color}{\hss\second@@@@}}}%
        \or%
           \vbox{\trans@size%
           \hbox to \box@width{\textcolor{\f@color}{\hss\second@@\hss}}%
           \hbox to \box@width{\textcolor{\f@color}{\hss\second@@@\hss}}%
           \hbox to \box@width{\textcolor{\f@color}{\hss\second@@@@\hss}}}%
        \fi%
      \fi%
     \ifx\fill@char\ampers@nd\else\repeat%
     \fi%
    \else%
     \xdef\last@{brace}%
     \ifx\second@@\last@%
      \xdef\last@{top}%
      \textcolor{\f@color}{%
      \ifx\bottop@\last@%
         \ensuremath{\overbrace{\hbox to \second@\box@width{\hss%
             \rule[0.2\arrow@height]{0pt}{0pt}\hss}}}%
      \else%
      \xdef\last@{ttop}%
      \ifx\bottop@\last@%
         \ensuremath{\overbrace{\hbox to \second@\box@width{\hss%
             \rule[0.2\arrow@height]{0pt}{0pt}\hss}}}%
      \else%
      \xdef\last@{tttop}%
      \ifx\bottop@\last@%
         \ensuremath{\overbrace{\hbox to \second@\box@width{\hss%
             \rule[0.2\arrow@height]{0pt}{0pt}\hss}}}%
      \else%
      \xdef\last@{ttttop}%
      \ifx\bottop@\last@%
         \ensuremath{\overbrace{\hbox to \second@\box@width{\hss%
             \rule[0.2\arrow@height]{0pt}{0pt}\hss}}}%
      \else%
         \raisebox{1.75\arrow@height}{\ensuremath{\underbrace{\hbox to %
             \second@\box@width{}}}}%
      \fi\fi\fi\fi}%
     \else%
      \xdef\last@{fill}%
      \ifx\second@@\last@%
       \fontfamily{\featurestyles@family}%
       \fontseries{\featurestyles@series}%
       \fontshape{\featurestyles@shape}%
       \selectfont%
       \xdef\last@{bottom}%
       \ifx\bottop@\temp@\xdef\last@{0.35}\else\xdef\last@{0}\fi%
       \kern0.15\box@width%
       \loopcount=0\relax%
       \loop%
         \advance\loopcount by 1\relax%
         \raisebox{\last@\arrow@height}{%
            \hbox to \box@width{\hss{\textcolor{\f@color}{\featurestyles@size\fill@char}}\hss}}%
       \ifnum\loopcount<\second@ \repeat%
      \else%
      \xdef\last@{bar}%
      \ifx\second@@\last@%
       \setlength\arrow@width{\pm@shift\box@height}%
       \setlength\arrow@width{\b@r@stretch\arrow@width}%
       \ifx\fill@char\N@%
          \raisebox{\pm@shift\box@height}{\vrule width\box@width}%
       \else%
          \kern-0.6\box@width%
          \ifx\frame@color\back@color%
          \else%
             \setlength\arrow@height{\box@height}%
             \advance\arrow@height by -\pm@shift\box@height%
             \setlength\arrow@height{\b@r@stretch\arrow@height}%
             \raisebox{\arrow@width}{%
                \hbox to \box@width{\hss{\textcolor{\back@color}{\vrule depth\arrow@width%
                height\arrow@height width\box@width}}\hss}}%
             \kern-\box@width%
          \fi%
          \setlength\arrow@height{\fill@char\box@height}%
          \setlength\arrow@height{\b@r@stretch\arrow@height}%
          \divide\arrow@height by 100\relax%
          \ifdim\arrow@height<0pt%
             \arrow@height=-\arrow@height%
             \raisebox{\arrow@width}{%
                \hbox to \box@width{\hss{\textcolor{\frame@color}{\vrule depth\arrow@height width0.8\box@width}}\hss}}%
          \else%
             \raisebox{\arrow@width}{%
                \hbox to \box@width{\hss{\textcolor{\frame@color}{\vrule height\arrow@height width0.8\box@width}}\hss}}\fi%
          \kern-\box@width%
          \raisebox{\arrow@width}{\textcolor{Black}{\vrule height0.25pt depth0.25pt width\box@width}}%
       \fi%
      \else%
      \xdef\last@{color}%
      \ifx\second@@\last@%
       \setlength\arrow@height{\fill@char\box@height}%
       \setlength\arrow@height{\sc@le@stretch\arrow@height}%
       \divide\arrow@height by 100\relax%
       \raisebox{0.2\box@height}{%
         \hbox to \box@width{\hss{\textcolor{\f@color}{\vrule height\arrow@height width\box@width}}\hss}}%
      \else%
      \xdef\last@{plotcolor}%
      \ifx\second@@\last@%
        \xdef\fill@char{\fill@char,&}%
        \loop%
            \expandafter\get@firstv@l\fill@char%
            \ifx\second@@\N@\hbox to \box@width{\hss}%
            \else
              \loopcount=\second@@%
              \ifx\T@coffee@ccons\y@%
              \else%
                 \advance\loopcount by -\pm@shift%
                 \advance\loopcount by 4%
                 \divide\loopcount by 5%
                 \multiply\loopcount by 5%
                 \ifnum\loopcount>100\loopcount=100\fi%
                 \ifnum\loopcount<5\loopcount=5\fi%
              \fi%
              \setlength\arrow@height{50\box@height}%
              \divide\arrow@height by 100\relax%
              \setlength\arrow@height{\sc@le@stretch\arrow@height}%
              \raisebox{0.2\box@height}{%
                 \hbox to \box@width{\hss{\textcolor{\f@color\the\loopcount}{\vrule height\arrow@height width\box@width}}\hss}}%
            \fi
        \ifx\fill@char\ampers@nd\else\repeat%
      \else%
      \xdef\last@{plotbar}%
      \ifx\second@@\last@%
        \xdef\fill@char{\fill@char,&}%
        \ifnum\pm@shift>0%
           \setlength\arrow@width{0pt}%
        \else
           \setlength\arrow@width{-\pm@shift\box@height}%
           \divide\arrow@width by 100%
           \setlength\arrow@width{\b@r@stretch\arrow@width}%
        \fi%
        \loop%
            \expandafter\get@firstv@l\fill@char%
            \ifx\second@@\N@\hbox to \box@width{\hss}%
            \else\relax%
               \ifx\frame@color\back@color%
               \else%
                  \setlength\arrow@height{\b@r@stretch\box@height}%
                  \hbox to \box@width{\hss{\textcolor{\back@color}%
                      {\vrule height\arrow@height width\box@width}}\hss}%
                  \kern-\box@width%
               \fi%
               \ifx\T@coffee@bcons\y@%
                  \ifnum\second@@=99 \xdef\second@@{0}\fi%
               \fi
               \setlength\arrow@height{\second@@\box@height}%
               \ifx\T@coffee@bcons\y@%
                  \setlength\arrow@height{11\arrow@height}%
               \fi
               \divide\arrow@height by 100\relax%
               \setlength\arrow@height{\b@r@stretch\arrow@height}%
               \ifdim\arrow@height<0pt%
                  \arrow@height=-\arrow@height%
                  \raisebox{\arrow@width}{%
                     \hbox to \box@width{\hss{\textcolor{\frame@color}%
                        {\vrule depth\arrow@height width0.8\box@width}}\hss}}%
               \else%
                  \raisebox{\arrow@width}{%
                     \hbox to \box@width{\hss{\textcolor{\frame@color}%
                        {\vrule height\arrow@height width0.8\box@width}}\hss}}\fi%
               \kern-\box@width%
               \raisebox{\arrow@width}{\textcolor{Black}%
                        {\vrule height0.25pt depth0.25pt width\box@width}}%
            \fi
        \ifx\fill@char\ampers@nd\else\repeat%
      \else%
      \xdef\last@{helix}%
      \ifx\second@@\last@%
       \fontfamily{cmr}%
       \fontseries{m}%
       \fontshape{it}%
       \selectfont%
       \ifx\bottop@\temp@\xdef\last@{0.35}\else\xdef\last@{0}\fi%
       \kern0.15\box@width%
       \setbox1=\hbox{\ensuremath{\helixhook}\kern-1.13exo\kern-1.02ex}%
       \arrow@width=\second@\box@width%
       \divide\arrow@width by \wd1%
       \arrow@width=2\wd1%
       \loop%
         \textcolor{\f@color}{\raisebox{-0.25ex}{\ensuremath{\helixhook}}%
                              \kern-1.13ex\raisebox{0.3ex}{o}}\kern-1.02ex%
         \advance\arrow@width by \wd1\relax%
       \ifdim\arrow@width<\second@\box@width \repeat%
       \textcolor{\f@color}{\raisebox{-0.25ex}{\ensuremath{\helixhook}}}%
      \else%
      \xdef\last@{box}%
      \ifx\second@@\last@%
         \fontfamily{\featurestyles@family}%
         \fontseries{\featurestyles@series}%
         \fontshape{\featurestyles@shape}%
         \selectfont%
         \kern-\second@\box@width%
         \bgroup%
             \textcolor{\back@color}{%
                \vrule width\second@\box@width height\box@height depth\box@depth}%
             \kern-\second@\box@width%
             \fboxsep-\rule@@thick\fboxrule\rule@@thick%
             \textcolor{\frame@color}{%
                \fbox{\makebox[\second@\box@width]%
                     {\vrule\@height\box@height\@depth\box@depth \@width\z@}}}%
         \egroup%
         \setbox1=\hbox{\residues@size{\fill@char}}%
         \temp@count=\wd1 \xdef\wd@{\the\temp@count}%
         \width@tmp=\second@\box@width%
         \temp@count=\width@tmp%
         \xdef\sb@{\the\temp@count}%
         \ifnum\wd@>\sb@ \xdef\fill@char{}\fi%
         \hbox to -\second@\box@width{\hss\textcolor{\f@color}%
                       {\residues@size{\fill@char}}\hss}%
     \else%
      \expandafter\get@shape\second@@%
      \xdef\last@{arrow}%
      \ifx\second@@\last@%
        \kern-0.75\box@width%
        \ifx\bottop@\temp@ \xdef\last@{0.35}\else\xdef\last@{-0.55}\fi%
        \textcolor{\f@color}{%
        \raisebox{\last@\arrow@height}{%
         \if\first@@ b \xdef\first@@{,}\fi%
         \if\first@@ ,%
            \rule{0.1\arrow@height}{\arrow@height}\kern-0.35\arrow@height%
         \else%
            \if\first@@ |%
               \rule{0.1\arrow@height}{2\arrow@height}\kern-0.35\arrow@height%
            \else%
               \if\first@@ `\xdef\first@@{'}\fi%
               \if\first@@ '%
                  \rule[\arrow@height]%
                     {0.1\arrow@height}{\arrow@height}\kern-0.35\arrow@height%
               \else%
               \if\first@@ '%
                  \rule[\arrow@height]%
                     {0.1\arrow@height}{\arrow@height}\kern-0.35\arrow@height%
               \else%
                  \if\first@@ -%
                     \rule{0pt}{0pt}\kern-0.35\arrow@height%
                  \fi%
               \fi%
               \fi%
            \fi%
         \fi%
         \if\third@@ v%
            \if\first@@ v%
               \ensuremath{\overleftarrow{\hbox to %
               \second@\box@width{\rule[0.4\arrow@height]{0pt}{0pt}\hss}}}%
               \kern-\arrow@width%
               \ensuremath{\overrightarrow{\hbox%
                          {\rule[0.4\arrow@height]{0pt}{0pt}\hss}}}%
            \else%
               \ensuremath{\overrightarrow{\hbox to %
               \second@\box@width{\rule[0.4\arrow@height]{0pt}{0pt}\hss}}}%
            \fi%
         \else%
            \if\first@@ v%
               \ensuremath{\overleftarrow{\hbox to %
                   \second@\box@width{\rule[0.4\arrow@height]{0pt}{0pt}\hss}}}%
            \else
               \kern0.35\arrow@height%
               \rule[0.9\arrow@height]{\second@\box@width}{0.1\arrow@height}%
               \kern0.35\arrow@height%
            \fi
            \if\third@@ ,%
               \kern-0.4\arrow@height\rule{0.1\arrow@height}{\arrow@height}%
            \else%
               \if\third@@ |%
                  \kern-0.4\arrow@height\rule{0.1\arrow@height}{2\arrow@height}%
               \else%
                  \if\third@@ `\xdef\third@@{'}\fi%
                  \if\third@@ '%
                     \kern-0.4\arrow@height%
                     \rule[\arrow@height]{0.1\arrow@height}{\arrow@height}%
                  \fi%
               \fi%
            \fi%
         \fi}}%
      \else%
       \xdef\last@{doublearrow}%
       \ifx\second@@\last@%
          \setbox1=\hbox{\ensuremath{\Rightarrow}}%
          \arrow@height=\ht1%
          \arrow@width=\wd1%
          \setbox1=\hbox to \second@\box@width{}%
          \width@tmp=\wd1%
          \kern-0.75\box@width%
        \xdef\temp@{top}%
        \ifx\bottop@\temp@ \xdef\last@{0.25}%
        \else%
           \xdef\temp@{ttop}%
           \ifx\bottop@\temp@ \xdef\last@{0.25}%
           \else%
             \xdef\temp@{tttop}%
             \ifx\bottop@\temp@ \xdef\last@{0.25}%
             \else%
               \xdef\temp@{ttttop}%
               \ifx\bottop@\temp@ \xdef\last@{0.25}%
               \else%
                 \xdef\last@{-0.25}\fi\fi\fi\fi%
        \textcolor{\f@color}{%
        \raisebox{\last@\arrow@height}{%
         \if\first@@ b \xdef\first@@{,}\fi%
         \if\first@@ ,%
            \rule[-0.5\arrow@height]{0.1\arrow@height}{1.5\arrow@height}%
            \kern-0.1\arrow@height%
         \else%
            \if\first@@ |%
               \rule[-0.5\arrow@height]{0.1\arrow@height}{2.25\arrow@height}%
               \kern-0.1\arrow@height%
            \else%
               \if\first@@ a \xdef\first@@{'}\fi%
               \if\first@@ ` \xdef\first@@{'}\fi%
               \if\first@@ '%
                  \rule[0.4\arrow@height]%
                     {0.1\arrow@height}{1.5\arrow@height}%
                  \kern-0.1\arrow@height%
               \else%
                  \if\first@@ <%
                     \ensuremath{\Leftarrow}\kern-0.5\arrow@width%
                     \advance\width@tmp by -0.5\arrow@width
                  \else
                     \rule{0pt}{0pt}%
                  \fi%
               \fi%
            \fi%
         \fi%
         \if\third@@ >%
            \advance\width@tmp by -0.5\arrow@width%
            \rule[0.37\arrow@height]{\width@tmp}{0.1\arrow@height}%
            \kern-\width@tmp%
            \rule[0.9\arrow@height]{\width@tmp}{0.1\arrow@height}%
            \kern-0.5\arrow@width\ensuremath{\Rightarrow}%
         \else%
            \rule[0.37\arrow@height]{\width@tmp}{0.1\arrow@height}%
            \kern-\width@tmp%
            \rule[0.9\arrow@height]{\width@tmp}{0.1\arrow@height}%
            \if\first@@ b \xdef\first@@{,}\fi%
            \if\third@@ ,%
               \kern-0.05\arrow@height%
               \rule[-0.5\arrow@height]{0.1\arrow@height}{1.5\arrow@height}%
            \else%
               \if\third@@ |%
                  \kern-0.05\arrow@height%
                  \rule[-0.5\arrow@height]{0.1\arrow@height}{2.25\arrow@height}%
               \else%
                  \if\first@@ a \xdef\first@@{'}\fi%
                  \if\third@@ ` \xdef\third@@{'}\fi%
                  \if\third@@ '%
                     \kern-0.05\arrow@height%
                     \rule[0.4\arrow@height]{0.1\arrow@height}{1.5\arrow@height}%
                  \else%
                  \if\third@@ a%
                     \kern-0.05\arrow@height%
                     \rule[0.4\arrow@height]{0.1\arrow@height}{1.5\arrow@height}%
                  \fi%
                  \fi%
               \fi%
            \fi%
         \fi}}%
       \else
        \loopcount=0\relax%
        \width@tmp=\arrow@height%
        \temp@@length=\rule@@thick%
        \advance\width@tmp by -0.5\temp@@length%
        \if\first@@ o\xdef\first@@{O}\fi%
        \if\third@@ o\xdef\third@@{O}\fi%
        \xdef\last@{bottom}%
        \ifx\fe@turep@s\last@%
           \vspace{-20pt}\message{(((-20pt)))}%
        \fi%
        \textcolor{\f@color}{%
         \if\first@@ ,%
           \rule{\temp@@length}{\arrow@height}\kern-\temp@@length%
         \else%
         \if\first@@ b%
           \xdef\last@{bottom}%
           \ifx\fe@turep@s\last@%
            \rule{\temp@@length}{\arrow@height}\kern-\temp@@length%
           \else
            \xdef\shift@feature{y}%
            \kern-\temp@@length%
            \rule[-\bar@length]{\temp@@length}{\bar@length}%
            \kern-\temp@@length%
            \rule{\temp@@length}{\arrow@height}\kern-\temp@@length%
           \fi
         \else%
            \if\first@@ |%
               \rule{\temp@@length}{2\arrow@height}\kern-\temp@@length%
            \else%
              \if\first@@ O%
                \raisebox{0.06ex}{\ensuremath{\bullet}}\kern-0.55ex%
                \rule[\width@tmp]{0.65ex}{\temp@@length}%
              \else%
              \if\first@@ S%
                \xdef\last@{bottom}%
                \ifx\fe@turep@s\last@%
                   \advance\width@tmp by 0.5\temp@@length%
                   \kern0.5\box@width%
                   \kern-2\temp@@length%
                   \rule[\width@tmp]{\temp@@length}{\arrow@height}%
                   \advance\width@tmp by -0.5\temp@@length%
                   \kern-\temp@@length%
                   \rule[\width@tmp]{0.5\temp@@length}{\temp@@length}%
                   \rule[\width@tmp]{0.5\box@width}{\temp@@length}%
                \else%
                   \kern0.5\box@width%
                   \kern-2\temp@@length%
                   \rule{\temp@@length}{\arrow@height}%
                   \kern-\temp@@length%
                   \rule[\width@tmp]{0.5\temp@@length}{\temp@@length}%
                   \rule[\width@tmp]{0.5\box@width}{\temp@@length}%
                \fi%
              \else%
              \if\first@@ c%
                \xdef\last@{bottom}%
                \ifx\fe@turep@s\last@%
                   \advance\width@tmp by 0.5\temp@@length%
                   \raisebox{\arrow@height}{\rule[\box@height]{0pt}{\bar@length}}%%%%%%%%%%%%%%%%%%%%%%%
                   \kern0.5\box@width%
                   \kern-2\temp@@length%
                   \rule[2\arrow@height]{\temp@@length}{\bar@length}%%%%%%%%%%%%%%%%%%%%%%
                   \kern-\temp@@length%%%%%%%%%%%%%%%%%%%%%%%%%
                   \rule[\width@tmp]{\temp@@length}{\arrow@height}%
                   \advance\width@tmp by -0.5\temp@@length%
                   \kern-\temp@@length%
                   \rule[\width@tmp]{0.5\temp@@length}{\temp@@length}%
                   \rule[\width@tmp]{0.5\box@width}{\temp@@length}%
                \else%
                   \xdef\shift@feature{y}%
                   \kern0.5\box@width%
                   \kern-2\temp@@length%
                   \rule[-\bar@length]{\temp@@length}{\bar@length}%
                   \kern-\temp@@length%
                   \rule{\temp@@length}{\arrow@height}%
                   \kern-\temp@@length%
                   \rule[\width@tmp]{0.5\temp@@length}{\temp@@length}%
                   \rule[\width@tmp]{0.5\box@width}{\temp@@length}%
                \fi%
              \else%
               \if\first@@ <%
                 \raisebox{0.06ex}{\ensuremath{\blacktriangleleft}}\kern-0.35ex%
               \else
                \if\first@@ `\xdef\first@@{'}\fi%
                \if\first@@ '%
                  \advance\width@tmp by 0.5\temp@@length%
                  \rule[\width@tmp]{\temp@@length}{\arrow@height}\kern-\temp@@length%
                  \advance\width@tmp by -0.5\temp@@length%
                \fi%
               \fi%
              \fi%
              \fi%
              \fi%
            \fi%
         \fi%
         \fi%
         \loopcount=\second@%
         \if\first@@ <\advance\loopcount by -1\relax\ifnum\loopcount<0 \loopcount=0\fi\fi%
         \if\third@@ >\advance\loopcount by -1\relax\ifnum\loopcount<0 \loopcount=0\fi\fi%
         \if\first@@ O\advance\loopcount by -1\relax\ifnum\loopcount<0 \loopcount=0\fi\fi%
         \if\third@@ O\advance\loopcount by -1\relax\ifnum\loopcount<0 \loopcount=0\fi\fi%
         \if\first@@ S\advance\loopcount by -1\relax\ifnum\loopcount<0 \loopcount=0\fi\fi%
         \if\third@@ S\advance\loopcount by -1\relax\ifnum\loopcount<0 \loopcount=0\fi\fi%
         \if\first@@ c\advance\loopcount by -1\relax\ifnum\loopcount<0 \loopcount=0\fi\fi%
         \if\third@@ c\advance\loopcount by -1\relax\ifnum\loopcount<0 \loopcount=0\fi\fi%
         \xdef\second@{\the\loopcount}%
         \rule[\width@tmp]{\second@\box@width}{\temp@@length}%
         \setbox1=\hbox{\residues@size{\fill@char}}%
         \kern-\second@\box@width%
         \hbox to \second@\box@width{\textcolor{\backtext@color}{\hss\rule[\width@tmp]{1.2\wd1}{\temp@@length}\hss}}%
         \kern-\second@\box@width%
         \hbox to \second@\box@width{\textcolor{\frame@color}{\residues@size{\hss\fill@char\hss}}}%
         \if\third@@ ,%
         \kern-\temp@@length\rule{\temp@@length}{\arrow@height}%
         \else%
         \if\third@@ b%
           \xdef\last@{bottom}%
           \ifx\fe@turep@s\last@%
            \kern-\temp@@length\rule{\temp@@length}{\arrow@height}%
           \else
            \xdef\shift@feature{y}%
            \kern-\temp@@length%
            \rule[-\bar@length]{\temp@@length}{\bar@length}%
            \kern-\temp@@length\rule{\temp@@length}{\arrow@height}%
           \fi
         \else%
            \if\third@@ |%
               \kern-\temp@@length\rule{\temp@@length}{2\arrow@height}%
            \else%
              \if\third@@ O%
                \rule[\width@tmp]{0.65ex}{\temp@@length}%
                \kern-0.55ex\raisebox{0.06ex}{\ensuremath{\bullet}}%
              \else%
              \if\third@@ S%
                \rule[\width@tmp]{0.5\box@width}{\temp@@length}%
                \rule[\width@tmp]{0.5\temp@@length}{\temp@@length}%
                \xdef\last@{bottom}%
                \ifx\fe@turep@s\last@%
                  \advance\width@tmp by 0.5\temp@@length%
                  \kern-\temp@@length\rule[\width@tmp]{\temp@@length}{\arrow@height}%
                \else
                  \kern-\temp@@length\rule{\temp@@length}{\arrow@height}%
                \fi%
              \else%
              \if\third@@ c%
                \xdef\shift@feature{y}%
                \rule[\width@tmp]{0.5\box@width}{\temp@@length}%
                \xdef\last@{bottom}%
                \ifx\fe@turep@s\last@%
                  \kern-\temp@@length%%%%%%%%%%%%%%%%%%%%%%%%%
                  \rule[2\arrow@height]{\temp@@length}{\bar@length}%%%%%%%%%%%%%%%%%%%%%%
                  \advance\width@tmp by 0.5\temp@@length%
                  \kern-\temp@@length\rule[\width@tmp]{\temp@@length}{\arrow@height}%
                \else
                  \kern-\temp@@length\rule[-\bar@length]{\temp@@length}{\bar@length}%
                  \kern-\temp@@length\rule{\temp@@length}{\arrow@height}%
                \fi%
              \else%
               \if\third@@ >%
                \kern-0.35ex\raisebox{0.06ex}{\ensuremath{\blacktriangleright}}%
               \else
                \if\third@@ `\xdef\third@@{'}\fi%
                \if\third@@ '%
                  \advance\width@tmp by 0.5\temp@@length%
                  \kern-\temp@@length\rule[\width@tmp]{\temp@@length}{\arrow@height}%
                \fi%
               \fi%
              \fi%
              \fi%
              \fi%
            \fi%
         \fi%
         \fi}%
       \fi
      \fi
      \fi
      \fi
      \fi
     \fi
     \fi
     \fi
    \fi
   \fi
   \fi
   \fi
}
\def\next@featuretext&#1;#2;#3@{%
   \xdef\first@{#3@}%
   \xdef\last@{fill:#2:&}\expandafter\test@fill\last@%
   \ifx\f@color\comm@\xdef\f@color{Black}\fi%
   \fontfamily{\featuretext@family}%
   \fontseries{\featuretext@series}%
   \fontshape{\featuretext@shape}%
   \selectfont%
   \xdef\last@{translate}%
   \ifx\last@\second@@%
     \xdef\fill@char{\fill@char &}%
     \if\seq@type N
  \loop%
  \expandafter\get@firstfill\fill@char%
  \if\second@@ -\def\second@@{\hss}\fi%
      \hbox to \box@width{\hss\textcolor{\f@color}{%
                      \featuretext@size{\strut\second@@}}\hss}%
  \ifx\fill@char\ampers@nd\else\repeat%
 \else
   \hbox to #1\box@width{\hss%
    \temp@count=1%
    \loop%
      \expandafter\get@tripletfill\fill@char%
      \if\second@@ -\ifnum\tr@nstextstyle>0\hbox to \box@width{\hss}\fi%
  \else%
   \ifcase\tr@nstextstyle%
     \hbox{\transtext@size\textcolor{\f@color}%
      {\hss\second@@\second@@@\second@@@@\,\hss}}%
   \or%
     \ifnum\temp@count=1%
        \vbox{\transtext@size%
          \hbox to \box@width{\textcolor{\f@color}{\second@@\hss}}%
          \hbox to \box@width{\textcolor{\f@color}{\hss\second@@@\hss}}%
          \hbox to \box@width{\textcolor{\f@color}{\hss\second@@@@}}}%
     \else%
        \vbox{\transtext@size%
          \hbox to \box@width{\textcolor{\f@color}{\hss\second@@@@}}%
          \hbox to \box@width{\textcolor{\f@color}{\hss\second@@@\hss}}%
          \hbox to \box@width{\textcolor{\f@color}{\second@@\hss}}}%
     \fi%
   \or%
     \ifnum\temp@count=1%
        \vbox{\transtext@size%
          \hbox to \box@width{\textcolor{\f@color}%
                             {\hss\second@@\second@@@\second@@@@\hss}}%
          \hbox to \box@width{\textcolor{\f@color}{\hss}}}%
     \else%
        \vbox{\transtext@size%
          \hbox to \box@width{\textcolor{\f@color}{\hss}}%
          \hbox to \box@width{\textcolor{\f@color}%
                             {\hss\second@@\second@@@\second@@@@\hss}}}%
     \fi%
   \or%
        \vbox{\transtext@size%
          \hbox to \box@width{\textcolor{\f@color}{\second@@\hss}}%
          \hbox to \box@width{\textcolor{\f@color}{\hss\second@@@\hss}}%
          \hbox to \box@width{\textcolor{\f@color}{\hss\second@@@@}}}%
   \or%
        \vbox{\transtext@size%
          \hbox to \box@width{\textcolor{\f@color}{\hss\second@@\hss}}%
          \hbox to \box@width{\textcolor{\f@color}{\hss\second@@@\hss}}%
          \hbox to \box@width{\textcolor{\f@color}{\hss\second@@@@\hss}}}%
   \fi%
  \fi%
    \ifx\fill@char\ampers@nd\else\repeat%
   \hss}
 \fi%
   \else%
      \textcolor{\f@color}{%
      \hbox to #1\box@width{\hss\featuretext@size{\strut\fourth@}\hss}}%
   \fi%
   \advance\loopcount by -#1%
}
\def\put@featuretext{%
   \if\first@ @%
   \else
      \expandafter\next@featuretext\first@%
      \put@featuretext%
   \fi}
\def\next@featurestyle&#1;#2;#3@{%
   \xdef\first@{#2}%
   \xdef\second@{#1}%
   \ifx\first@\ampers@nd \hbox to \second@\box@width{\hss}%
   \else%
   \hbox to \second@\box@width%
           {\hss\expandafter\putfeature@style{\first@}\hss}\fi%
   \xdef\first@{#3@}%
   \advance\loopcount by -\second@%
}
\def\put@featurestyle{%
   \if\first@ @%
   \else
      \expandafter\next@featurestyle\first@%
      \put@featurestyle%
   \fi
}
\def\put@feature{%
   \vspace{-\baselineskip}%
   \newline\hbox{}%
   \xdef\temp@{ttttop}%
   \ifx\temp@\bottop@%
      \ifnames@right\else\ifnames@%
      \fontfamily{\ftext@family}\fontseries{\ftext@series}\fontshape{\ftext@shape}%
      \selectfont%
      \textcolor{\csname ftext@fg@\bottop@\endcsname}{\hbox to \name@width{\ftext@size\csname featuretextn@me\bottop@\endcsname\hss}}\fi\fi%
      \ifnumbers@left\ifnumbers@\hbox to \number@width{\hss}\fi\fi%
      \xdef\first@{\csname textfeature\bottop@\endcsname @}%
      \loopcount=\res@perline%
      \put@featuretext%
      \ifnumbers@right\ifnumbers@\hbox to \number@width{\hss}\fi\fi%
      \ifnames@right\ifnames@%
      \fontfamily{\ftext@family}\fontseries{\ftext@series}\fontshape{\ftext@shape}%
      \selectfont%
      \kern1em\hbox to \loopcount\box@width{\hss}%
      \textcolor{\csname ftext@fg@\bottop@\endcsname}{\hbox to \name@width{\ftext@size\csname featuretextn@me\bottop@\endcsname\hss}}\fi\fi%
   \else%
   \xdef\temp@{tttop}%
   \ifx\temp@\bottop@%
      \ifnames@right\else\ifnames@%
      \fontfamily{\ftext@family}\fontseries{\ftext@series}\fontshape{\ftext@shape}%
      \selectfont%
      \textcolor{\csname ftext@fg@\bottop@\endcsname}{\hbox to \name@width{\ftext@size\csname featuretextn@me\bottop@\endcsname\hss}}\fi\fi%
      \ifnumbers@left\ifnumbers@\hbox to \number@width{\hss}\fi\fi%
      \xdef\first@{\csname textfeature\bottop@\endcsname @}%
      \loopcount=\res@perline%
      \put@featuretext%
      \ifnumbers@right\ifnumbers@\hbox to \number@width{\hss}\fi\fi%
      \ifnames@right\ifnames@%
      \fontfamily{\ftext@family}\fontseries{\ftext@series}\fontshape{\ftext@shape}%
      \selectfont%
      \kern1em\hbox to \loopcount\box@width{\hss}%
      \textcolor{\csname ftext@fg@\bottop@\endcsname}{\hbox to \name@width{\ftext@size\csname featuretextn@me\bottop@\endcsname\hss}}\fi\fi%
   \else%
   \xdef\temp@{ttop}%
   \ifx\temp@\bottop@%
      \ifnames@right\else\ifnames@%
      \fontfamily{\ftext@family}\fontseries{\ftext@series}\fontshape{\ftext@shape}%
      \selectfont%
      \textcolor{\csname ftext@fg@\bottop@\endcsname}{\hbox to \name@width{\ftext@size\csname featuretextn@me\bottop@\endcsname\hss}}\fi\fi%
      \ifnumbers@left\ifnumbers@\hbox to \number@width{\hss}\fi\fi%
      \xdef\first@{\csname textfeature\bottop@\endcsname @}%
      \loopcount=\res@perline%
      \put@featuretext%
      \ifnumbers@right\ifnumbers@\hbox to \number@width{\hss}\fi\fi%
      \ifnames@right\ifnames@%
      \fontfamily{\ftext@family}\fontseries{\ftext@series}\fontshape{\ftext@shape}%
      \selectfont%
      \kern1em\hbox to \loopcount\box@width{\hss}%
      \textcolor{\csname ftext@fg@\bottop@\endcsname}{\hbox to \name@width{\ftext@size\csname featuretextn@me\bottop@\endcsname\hss}}\fi\fi%
   \else%
      \xdef\temp@{top}%
      \ifx\temp@\bottop@%
         \ifnames@right\else\ifnames@%
         \fontfamily{\ftext@family}\fontseries{\ftext@series}\fontshape{\ftext@shape}%
         \selectfont%
         \textcolor{\csname ftext@fg@\bottop@\endcsname}{\hbox to \name@width{\ftext@size\csname featuretextn@me\bottop@\endcsname\hss}}\fi\fi%
         \ifnumbers@left\ifnumbers@\hbox to \number@width{\hss}\fi\fi%
         \xdef\first@{\csname textfeature\bottop@\endcsname @}%
         \loopcount=\res@perline%
         \put@featuretext%
         \ifnumbers@right\ifnumbers@\hbox to \number@width{\hss}\fi\fi%
         \ifnames@right\ifnames@%
         \fontfamily{\ftext@family}\fontseries{\ftext@series}\fontshape{\ftext@shape}%
         \selectfont%
         \kern1em\hbox to \loopcount\box@width{\hss}%
         \textcolor{\csname ftext@fg@\bottop@\endcsname}{\hbox to \name@width{\ftext@size\csname featuretextn@me\bottop@\endcsname\hss}}\fi\fi%
      \else
         \ifnames@right\else\ifnames@%
         \fontfamily{\fstyles@family}\fontseries{\fstyles@series}\fontshape{\fstyles@shape}%
         \selectfont%
         \textcolor{\csname fstyles@fg@\bottop@\endcsname}{\hbox to \name@width{\fstyles@size\csname featurestylesn@me\bottop@\endcsname\hss}}\fi\fi%
         \ifnumbers@left\ifnumbers@\hbox to \number@width{\hss}\fi\fi%
         \xdef\first@{\csname stylefeature\bottop@\endcsname @}%
         \loopcount=\res@perline%
         \put@featurestyle%
         \ifnumbers@right\ifnumbers@\hbox to \number@width{\hss}\fi\fi%
         \ifnames@right\ifnames@%
         \fontfamily{\fstyles@family}\fontseries{\fstyles@series}\fontshape{\fstyles@shape}%
         \selectfont%
         \kern1em\hbox to \loopcount\box@width{\hss}%
         \textcolor{\csname fstyles@fg@\bottop@\endcsname}{\hbox to \name@width{\fstyles@size\csname featurestylesn@me\bottop@\endcsname\hss}}\fi\fi%
   \fi\fi\fi\fi%
   \newline\hbox{}%
   \ifx\temp@\bottop@%
      \ifnames@right\else\ifnames@%
      \fontfamily{\fstyles@family}\fontseries{\fstyles@series}\fontshape{\fstyles@shape}%
      \selectfont%
      \textcolor{\csname fstyles@fg@\bottop@\endcsname}{\hbox to \name@width{\fstyles@size\csname featurestylesn@me\bottop@\endcsname\hss}}\fi\fi%
      \ifnumbers@left\ifnumbers@\hbox to \number@width{\hss}\fi\fi%
      \xdef\first@{\csname stylefeature\bottop@\endcsname @}%
      \loopcount=\res@perline%
      \put@featurestyle%
      \ifnumbers@right\ifnumbers@\hbox to \number@width{\hss}\fi\fi%
      \ifnames@right\ifnames@%
      \fontfamily{\fstyles@family}\fontseries{\fstyles@series}\fontshape{\fstyles@shape}%
      \selectfont%
      \kern1em\hbox to \loopcount\box@width{\hss}%
      \textcolor{\csname fstyles@fg@\bottop@\endcsname}{\hbox to \name@width{\fstyles@size\csname featurestylesn@me\bottop@\endcsname\hss}}\fi\fi%
   \else%
      \xdef\temp@{ttop}%
      \ifx\temp@\bottop@%
         \ifnames@right\else\ifnames@%
         \fontfamily{\fstyles@family}\fontseries{\fstyles@series}\fontshape{\fstyles@shape}%
         \selectfont%
         \textcolor{\csname fstyles@fg@\bottop@\endcsname}{\hbox to \name@width{\fstyles@size\csname featurestylesn@me\bottop@\endcsname\hss}}\fi\fi%
         \ifnumbers@left\ifnumbers@\hbox to \number@width{\hss}\fi\fi%
         \xdef\first@{\csname stylefeature\bottop@\endcsname @}%
         \loopcount=\res@perline%
         \put@featurestyle%
         \ifnumbers@right\ifnumbers@\hbox to \number@width{\hss}\fi\fi%
         \ifnames@right\ifnames@%
         \fontfamily{\fstyles@family}\fontseries{\fstyles@series}\fontshape{\fstyles@shape}%
         \selectfont%
         \kern1em\hbox to \loopcount\box@width{\hss}%
         \textcolor{\csname fstyles@fg@\bottop@\endcsname}{\hbox to \name@width{\fstyles@size\csname featurestylesn@me\bottop@\endcsname\hss}}\fi\fi%
      \else
      \xdef\temp@{tttop}%
      \ifx\temp@\bottop@%
         \ifnames@right\else\ifnames@%
         \fontfamily{\fstyles@family}\fontseries{\fstyles@series}\fontshape{\fstyles@shape}%
         \selectfont%
         \textcolor{\csname fstyles@fg@\bottop@\endcsname}{\hbox to \name@width{\fstyles@size\csname featurestylesn@me\bottop@\endcsname\hss}}\fi\fi%
         \ifnumbers@left\ifnumbers@\hbox to \number@width{\hss}\fi\fi%
         \xdef\first@{\csname stylefeature\bottop@\endcsname @}%
         \loopcount=\res@perline%
         \put@featurestyle%
         \ifnumbers@right\ifnumbers@\hbox to \number@width{\hss}\fi\fi%
         \ifnames@right\ifnames@%
         \fontfamily{\fstyles@family}\fontseries{\fstyles@series}\fontshape{\fstyles@shape}%
         \selectfont%
         \kern1em\hbox to \loopcount\box@width{\hss}%
         \textcolor{\csname fstyles@fg@\bottop@\endcsname}{\hbox to \name@width{\fstyles@size\csname featurestylesn@me\bottop@\endcsname\hss}}\fi\fi%
      \else
      \xdef\temp@{ttttop}%
      \ifx\temp@\bottop@%
         \ifnames@right\else\ifnames@%
         \fontfamily{\fstyles@family}\fontseries{\fstyles@series}\fontshape{\fstyles@shape}%
         \selectfont%
         \textcolor{\csname fstyles@fg@\bottop@\endcsname}{\hbox to \name@width{\fstyles@size\csname featurestylesn@me\bottop@\endcsname\hss}}\fi\fi%
         \ifnumbers@left\ifnumbers@\hbox to \number@width{\hss}\fi\fi%
         \xdef\first@{\csname stylefeature\bottop@\endcsname @}%
         \loopcount=\res@perline%
         \put@featurestyle%
         \ifnumbers@right\ifnumbers@\hbox to \number@width{\hss}\fi\fi%
         \ifnames@right\ifnames@%
         \fontfamily{\fstyles@family}\fontseries{\fstyles@series}\fontshape{\fstyles@shape}%
         \selectfont%
         \kern1em\hbox to \loopcount\box@width{\hss}%
         \textcolor{\csname fstyles@fg@\bottop@\endcsname}{\hbox to \name@width{\fstyles@size\csname featurestylesn@me\bottop@\endcsname\hss}}\fi\fi%
      \else
         \ifnames@right\else\ifnames@%
         \fontfamily{\ftext@family}\fontseries{\ftext@series}\fontshape{\ftext@shape}%
         \selectfont%
         \textcolor{\csname ftext@fg@\bottop@\endcsname}{\hbox to \name@width{\ftext@size\csname featuretextn@me\bottop@\endcsname\hss}}\fi\fi%
         \ifnumbers@left\ifnumbers@\hbox to \number@width{\hss}\fi\fi%
         \xdef\first@{\csname textfeature\bottop@\endcsname @}%
         \xdef\first@{\csname textfeature\bottop@\endcsname @}%
         \loopcount=\res@perline%
         \put@featuretext%
         \ifnumbers@right\ifnumbers@\hbox to \number@width{\hss}\fi\fi%
         \ifnames@right\ifnames@%
         \fontfamily{\ftext@family}\fontseries{\ftext@series}\fontshape{\ftext@shape}%
         \selectfont%
         \kern1em\hbox to \loopcount\box@width{\hss}%
         \textcolor{\csname ftext@fg@\bottop@\endcsname}{\hbox to \name@width{\ftext@size\csname featuretextn@me\bottop@\endcsname\hss}}\fi\fi%
   \fi\fi\fi\fi%
   \newline\hbox{}%
}
\def\put@@@frame#1{%
   \xdef\last@{#1[,]&}\expandafter\opt@color\last@%
   \xdef\second@@{\fourth@}%
   \ifx\f@color\comm@%
      \xdef\third@@{0.2\box@width}%
   \else%
      \xdef\third@@{\f@color}%
   \fi%
   \ifnum\second@<0%
      \loopcount=-\second@\relax%
      \xdef\second@{\the\loopcount}%
      \xdef\right@frame@open{y}%
   \fi%
   \setlength\arrow@width{\temp@@length}%
   \advance\arrow@width by -\third@@%
   \textcolor{\second@@}{%
      \rule{\second@\box@width}{\third@@}%
      \kern-\second@\box@width%
      \ifx\left@frame@open\n@%
         \rule{\third@@}{\arrow@width}%
         \kern-\third@@%
      \else
         \xdef\left@frame@open{n}%
      \fi      
      \rule[\arrow@width]{\second@\box@width}{\third@@}%
      \kern-\third@@%
      \ifx\right@frame@open\n@%
         \rule{\third@@}{\arrow@width}%
      \else
         \xdef\left@frame@open{y}%
         \xdef\right@frame@open{n}%
      \fi%
      }%
}
\def\next@frame&#1;#2;#3@{%
   \xdef\first@{#2}%
   \xdef\second@{#1}%
   \ifx\first@\ampers@nd \hbox to \second@\box@width{\hss}%
   \else%
      \expandafter\put@@@frame{\first@}%
   \fi%
   \xdef\first@{#3@}%
}
\def\put@@frame{%
   \if\first@ @%
   \else
      \expandafter\next@frame\first@%
      \put@@frame%
   \fi%
}
\def\put@frame{%
   \ifx\hide@seqs\n@%
   \ifshow@cons\ifnum\cons@top=1 \vspace{-\baselineskip}\fi\fi%
   \temp@count=\seq@count%
   \ifx\hide@seqs\y@
      \advance\temp@count by -1
   \else
   \loopcount=1%
   \xdef\first@{true}%
   \loop%
      \expandafter\ifx\csname hide@seq\the\loopcount\endcsname\first@%
         \advance\temp@count by -1\fi%
      \advance\loopcount by 1%
   \ifnum\loopcount>\seq@count\else\repeat%
   \fi
   \ifnum\temp@count>0%
   \setlength\temp@@length{\box@height}%
   \advance\temp@@length by \box@depth%
   \setlength\arrow@height{0.5\temp@@length}%
   \setlength\temp@@length{\temp@count\temp@@length}%
   \advance\temp@@length by \arrow@height%
   \setlength\arrow@width{\sep@space}%
   \setlength\arrow@width{\seq@gap@num\arrow@width}%
   \advance\temp@@length by \arrow@width%
   \vspace{\arrow@height}%
   \vspace{-\temp@@length}%
   \vspace{-\baselineskip}%
   \newline\hbox{}%
   \ifnames@right\else\ifnames@\hbox to \name@width{\hss}\fi\fi%
   \ifnumbers@left\ifnumbers@\hbox to \number@width{\hss}\fi\fi%
   \xdef\first@{\styleframe @}%
   \put@@frame%
   \vspace{-\arrow@height}%
   \ifshow@cons\ifnum\cons@top=1 \vspace{\baselineskip}\fi\fi%
   \newline\hbox{}%
   \fi%
   \fi%
}
\def\decimal@A#1#2@{%
   \def\decimal@AB##1##2@{\xdef\temp@{##1.##2}}%
   \xdef\temp@@{#1}%
   \xdef\temp@{#2@}%
   \expandafter\decimal@AB\temp@%
   \xdef\temp@{\temp@@\temp@}%
}
\def\decimal@B#1#2@{\xdef\temp@{#1.#2}}
\def\decimal@C#1#2@{\xdef\temp@{0.#1#2}}
\def\decimal@D#1#2@{\xdef\temp@{0.0#1#2}}
\def\decimal@E#1#2@{\xdef\temp@{0.00#1#2}}
\def\correct@CGSO{%
   \setlength\temp@@length{0.12\logo@height}%
   \multiply\temp@@length by \third@ %
   \if\seq@type P %
      \divide\temp@@length by 4322 % 
   \else%
      \divide\temp@@length by 4000 %
   \fi%
   \loopcount=\third@%
   \multiply\loopcount by 94 %
   \divide\loopcount by 100 %
   \xdef\third@{\the\loopcount}%
}
\def\correct@Q{%
   \setlength\temp@@length{0.063\logo@height}%
   \multiply\temp@@length by \third@%
   \if\seq@type P %
      \divide\temp@@length by 4322 % 
   \else%
      \divide\temp@@length by 4000 %
   \fi%
   \multiply\temp@@length by 10 %
   \loopcount=\third@%
   \multiply\loopcount by 83 %
   \divide\loopcount by 100 %
   \xdef\third@{\the\loopcount}%
}
\def\correct@JUV{%
   \setlength\temp@@length{0.12\logo@height}%
   \multiply\temp@@length by \third@ %
   \if\seq@type P %
      \divide\temp@@length by 4322 % 
   \else%
      \divide\temp@@length by 4000 %
   \fi%
   \loopcount=\third@%
   \multiply\loopcount by 97 %
   \divide\loopcount by 100 %
   \xdef\third@{\the\loopcount}%
}
%%%%%%%%%%%%%%%%%%%%%%%%%%%%%%%%%%%
%%%%                         Chimera export routines                      %%%%
%%%%%%%%%%%%%%%%%%%%%%%%%%%%%%%%%%%
\def\chimeraballScale#1{\xdef\chimerab@llSc@le{#1}}
\def\chimeraballscale#1{\xdef\chimerab@llSc@le{#1}}
\def\chimeraaxisdistance#1#2#3{\xdef\chimera@xdist@nce{#1}\xdef\chimera@ydist@nce{#2}\xdef\chimera@zdist@nce{#3}}
\def\memelabelcutoff#1{%
   \xdef\label@threshold@lower{#1}
   \setlength\temp@@length{1000sp}
   \setlength\temp@@length{\label@threshold@lower\temp@@length}
   \loopcount=\temp@@length
   \xdef\label@threshold@lowerIOOO{\the\loopcount}
   \xdef\label@threshold@lowerremainingIOOO{-\the\loopcount}
}
\def\chimera@color@get#1#2@{%
   \xdef\first@{#1}\xdef\third@{#2@}
   \ifx\first@\ampers@nd
   \else
      \ifnum`#1>96 \make@upper\fi
      \if\second@ O \else \expandafter\xdef\csname chimera@col@\first@\endcsname{\second@} \fi
      \expandafter\chimera@color@get\third@
   \fi
}
\def\memered#1{%
      \xdef\third@{#1&@}
      \xdef\second@{red}
      \expandafter\chimera@color@get\third@
}
\def\memeblue#1{%
      \xdef\third@{#1&@}
      \xdef\second@{blue}
      \expandafter\chimera@color@get\third@
}
\def\memeyellow#1{%
      \xdef\third@{#1&@}
      \xdef\second@{yellow}
      \expandafter\chimera@color@get\third@
}
\def\memeblack#1{%
      \xdef\third@{#1&@}
      \xdef\second@{black}
      \expandafter\chimera@color@get\third@
}
\def\memewhite#1{%
      \xdef\third@{#1&@}
      \xdef\second@{white}
      \expandafter\chimera@color@get\third@
}
\def\memenone#1{%
      \xdef\third@{#1&@}
      \xdef\second@{none}
      \expandafter\chimera@color@get\third@
}
\def\memeRed#1{\memered{#1}}
\def\memeBlue#1{\memeblue{#1}}
\def\memeYellow#1{\memeyellow{#1}}
\def\memeBlack#1{\memeblack{#1}}
\def\memeWhite#1{\memewhite{#1}}
\def\memeNone#1{\memenone{#1}}
\def\chimera@colors#1#2@{%
   \xdef\first@{#1}\xdef\third@{#2@}
   \ifx\first@\ampers@nd
   \else
      \xdef\second@{\csname chimera@col@\first@\endcsname}
      \expandafter\xdef\csname chimera@list@\second@\endcsname{\csname chimera@list@\second@\endcsname \first@}
      \expandafter\chimera@colors\third@
   \fi
}
\def\memeStandardcolors{%
   \xdef\chimera@col@A{white}
   \xdef\chimera@col@B{none}
   \xdef\chimera@col@C{yellow}
   \xdef\chimera@col@D{red}
   \xdef\chimera@col@E{red}
   \xdef\chimera@col@F{black}
   \xdef\chimera@col@G{yellow}
   \xdef\chimera@col@H{blue}
   \xdef\chimera@col@I{black}
   \xdef\chimera@col@J{none}
   \xdef\chimera@col@K{blue}
   \xdef\chimera@col@L{black}
   \xdef\chimera@col@M{black}
   \xdef\chimera@col@N{yellow}
   \xdef\chimera@col@O{none}
   \xdef\chimera@col@P{white}
   \xdef\chimera@col@Q{yellow}
   \xdef\chimera@col@R{blue}
   \xdef\chimera@col@S{yellow}
   \xdef\chimera@col@T{yellow}
   \xdef\chimera@col@U{none}
   \xdef\chimera@col@V{white}
   \xdef\chimera@col@W{black}
   \xdef\chimera@col@X{none}
   \xdef\chimera@col@Y{yellow}
   \xdef\chimera@col@Z{none}
}
\def\echostructurefile#1{\xdef\echo@filename{, e.g. #1}}
\def\chimerachain#1{\xdef\chimerach@in{.#1}}
\catcode`\"=6 \catcode`\#=12 \catcode`\~=12
\newcommand{\structurememe}[2][chimera.cmd]{%
   \ifx\exp@rt@chimera\n@
      \xdef\first@{"2 @} \expandafter\check@letter\first@
      \xdef\first@{"2}   \ifletter \get@name@number \fi
      \xdef\chimera@seqnum{\first@}
      \xdef\exp@rt@chimera{y}
      \xdef\first@{"1}
      \immediate\openout\exp@rt@chimerafile = \first@
   \fi
}
\def\init@chimera@file{%
      \xdef\chimera@list@red{} \xdef\chimera@list@blue{} \xdef\chimera@list@yellow{} \xdef\chimera@list@black{} \xdef\chimera@list@white{} 
      \xdef\chimera@list@none{}
      \xdef\first@{ABCDEFGHIJKLMNOPQRSTUVWXYZ&@}
      \expandafter\chimera@colors\first@
      \immediate\write\exp@rt@chimerafile{############################################################}
      \immediate\write\exp@rt@chimerafile{###}
      \immediate\write\exp@rt@chimerafile{###\unexpanded{ }\unexpanded{ }\unexpanded{ }. O ¡ o STRUCTURE MEME ¡ . O o}
      \immediate\write\exp@rt@chimerafile{###}
      \immediate\write\exp@rt@chimerafile{###\unexpanded{ }\unexpanded{ }\unexpanded{ }Chimera command file generated from TeXshade logo data}
      \immediate\write\exp@rt@chimerafile{###\unexpanded{ }\unexpanded{ }\unexpanded{ }use '.cmd' or '.com' suffix}
      \immediate\write\exp@rt@chimerafile{###}
      \immediate\write\exp@rt@chimerafile{###\unexpanded{ }\unexpanded{ }\unexpanded{ }Reference sequence: '\csname newseqname\chimera@seqnum\endcsname'\echo@filename}
      \immediate\write\exp@rt@chimerafile{###}
      \immediate\write\exp@rt@chimerafile{############################################################}
      \immediate\write\exp@rt@chimerafile{}
      \immediate\write\exp@rt@chimerafile{############################################################}
      \immediate\write\exp@rt@chimerafile{###\unexpanded{ }\unexpanded{ }\unexpanded{ }residue colors: (set with \string\memeRed{residues}\unexpanded{ }etc.)}
      \immediate\write\exp@rt@chimerafile{###\unexpanded{ }\unexpanded{ }\unexpanded{ }\chimera@list@red \unexpanded{ }(Red)}
      \immediate\write\exp@rt@chimerafile{###\unexpanded{ }\unexpanded{ }\unexpanded{ }\chimera@list@blue \unexpanded{ }(Blue)}
      \immediate\write\exp@rt@chimerafile{###\unexpanded{ }\unexpanded{ }\unexpanded{ }\chimera@list@yellow \unexpanded{ }(Yellow)}
      \immediate\write\exp@rt@chimerafile{###\unexpanded{ }\unexpanded{ }\unexpanded{ }\chimera@list@black \unexpanded{ }(Black)}
      \immediate\write\exp@rt@chimerafile{###\unexpanded{ }\unexpanded{ }\unexpanded{ }\chimera@list@white \unexpanded{ }(White)}
      \immediate\write\exp@rt@chimerafile{############################################################}
      \immediate\write\exp@rt@chimerafile{}
      \immediate\write\exp@rt@chimerafile{############################################################}
      \immediate\write\exp@rt@chimerafile{###\unexpanded{ }\unexpanded{ }\unexpanded{ }residues in labels: (\string\memelabelcutoff{bit-value})}
      \immediate\write\exp@rt@chimerafile{###\unexpanded{ }\unexpanded{ }\unexpanded{ }\label@threshold@lower\unexpanded{ }bit required to appear}
      \immediate\write\exp@rt@chimerafile{###\unexpanded{ }\unexpanded{ }\unexpanded{ }> : . indicate 3, 2, and 1 bit levels}
      \immediate\write\exp@rt@chimerafile{############################################################}
      \immediate\write\exp@rt@chimerafile{}
      \immediate\write\exp@rt@chimerafile{############################################################}
      \immediate\write\exp@rt@chimerafile{###\unexpanded{ }\unexpanded{ }\unexpanded{ }initialize & load structure}
      \immediate\write\exp@rt@chimerafile{############################################################}
      \immediate\write\exp@rt@chimerafile{close session; background solid white;}
      \immediate\write\exp@rt@chimerafile{echo load structure file for '\csname newseqname\chimera@seqnum\endcsname'\echo@filename}
      \immediate\write\exp@rt@chimerafile{open #0; ~ribbon; ~surface}
      \immediate\write\exp@rt@chimerafile{}
      \immediate\write\exp@rt@chimerafile{############################################################}
      \immediate\write\exp@rt@chimerafile{###\unexpanded{ }\unexpanded{ }\unexpanded{ }to set ball scale, change factor ]0.0 - 1.0]}
      \immediate\write\exp@rt@chimerafile{###\unexpanded{ }\unexpanded{ }\unexpanded{ }in 'setattr m ballScale <factor> #0' below,}
      \immediate\write\exp@rt@chimerafile{###\unexpanded{ }\unexpanded{ }\unexpanded{ }or use the \string\chimeraballScale{factor}\unexpanded{ }command}
      \immediate\write\exp@rt@chimerafile{############################################################}
      \immediate\write\exp@rt@chimerafile{setattr m ballScale \chimerab@llSc@le \unexpanded{ }#0}
      \immediate\write\exp@rt@chimerafile{}
      \immediate\write\exp@rt@chimerafile{############################################################}
      \immediate\write\exp@rt@chimerafile{###\unexpanded{ }\unexpanded{ }\unexpanded{ }display alpha-carbons only}
      \immediate\write\exp@rt@chimerafile{############################################################}
      \immediate\write\exp@rt@chimerafile{represent bs @ca; show @ca}
      \immediate\write\exp@rt@chimerafile{}
      \ifshow@sublogo
            \immediate\write\exp@rt@chimerafile{############################################################}
            \immediate\write\exp@rt@chimerafile{###\unexpanded{ }\unexpanded{ }\unexpanded{ }duplicate structure;}
            \immediate\write\exp@rt@chimerafile{###\unexpanded{ }\unexpanded{ }\unexpanded{ }to alter direction or Angstrom distance, set}
            \immediate\write\exp@rt@chimerafile{###\unexpanded{ }\unexpanded{ }\unexpanded{ }'move <x, y, or z> 1 <distance> models #1' below,}
            \immediate\write\exp@rt@chimerafile{###\unexpanded{ }\unexpanded{ }\unexpanded{ }or use the \string\chimeraaxisdistance{x}{y}{z}\unexpanded{ }command}
            \immediate\write\exp@rt@chimerafile{############################################################}
            \immediate\write\exp@rt@chimerafile{combine #0}
            \ifnum\chimera@xdist@nce>0
               \immediate\write\exp@rt@chimerafile{move x 1 \chimera@xdist@nce \unexpanded{ }models #1}
            \fi
            \ifnum\chimera@ydist@nce>0
               \immediate\write\exp@rt@chimerafile{move y 1 \chimera@ydist@nce \unexpanded{ }models #1}
            \fi
            \ifnum\chimera@zdist@nce>0
               \immediate\write\exp@rt@chimerafile{move z 1 \chimera@zdist@nce \unexpanded{ }models #1}
            \fi
            \immediate\write\exp@rt@chimerafile{}
            \immediate\write\exp@rt@chimerafile{############################################################}
            \immediate\write\exp@rt@chimerafile{###\unexpanded{ }\unexpanded{ }\unexpanded{ }focus view and name structure models}
            \immediate\write\exp@rt@chimerafile{############################################################}
            \immediate\write\exp@rt@chimerafile{wait; focus #0#1}
            \ifx\sublogo@name@user\ampers@nd
                \immediate\write\exp@rt@chimerafile{setattr m name subfamily #0}
            \else
                \immediate\write\exp@rt@chimerafile{setattr m name '\sublogo@name@user' \unexpanded{ }#0}
            \fi
            \xdef\first@{\sublogo@name@neg &}
            \ifx\first@\ampers@nd
                \immediate\write\exp@rt@chimerafile{setattr m name remaining #1}
            \else
                \immediate\write\exp@rt@chimerafile{setattr m name '\sublogo@name@neg' \unexpanded{ }#1}
            \fi
     \else
            \immediate\write\exp@rt@chimerafile{\string ############################################################}
            \immediate\write\exp@rt@chimerafile{\string ###\unexpanded{ }\unexpanded{ }\unexpanded{ }name structure model}
            \immediate\write\exp@rt@chimerafile{\string ############################################################}
            \ifx\logo@name@user\ampers@nd
                \immediate\write\exp@rt@chimerafile{setattr m name logo #0}
            \else
                \immediate\write\exp@rt@chimerafile{setattr m name '\logo@name@user' \unexpanded{ }#0}
            \fi
      \fi
      \immediate\write\exp@rt@chimerafile{}
      \immediate\write\exp@rt@chimerafile{############################################################}
      \immediate\write\exp@rt@chimerafile{###\unexpanded{ }\unexpanded{ }\unexpanded{ }logo data}
      \immediate\write\exp@rt@chimerafile{############################################################}
}
\def\make@HEX{%
     \xdef\sixth@{\the\loopcount}%
     \divide\loopcount by 4096\relax%
     \ifnum\loopcount=10\relax%
        \xdef\seventh@{A}%
     \else%
     \ifnum\loopcount=11\relax%
        \xdef\seventh@{B}%
     \else%
     \ifnum\loopcount=12\relax%
        \xdef\seventh@{C}%
     \else%
     \ifnum\loopcount=13\relax%
        \xdef\seventh@{D}%
     \else%
     \ifnum\loopcount=14\relax%
        \xdef\seventh@{E}%
     \else%
     \ifnum\loopcount>14\relax%
        \xdef\seventh@{F}%
     \else%
        \xdef\seventh@{\the\loopcount}%
     \fi\fi\fi\fi\fi\fi%
     \multiply\loopcount by -4096\relax%
     \advance\loopcount by \sixth@%
     \xdef\sixth@{\the\loopcount}%
     \divide\loopcount by 256\relax%
     \ifnum\loopcount=10 \relax%
        \xdef\eighth@{A}%
     \else%
     \ifnum\loopcount=11 \relax%
        \xdef\eighth@{B}%
     \else%
     \ifnum\loopcount=12 \relax%
        \xdef\eighth@{C}%
     \else%
     \ifnum\loopcount=13 \relax%
        \xdef\eighth@{D}%
     \else%
     \ifnum\loopcount=14 \relax%
        \xdef\eighth@{E}%
     \else%
     \ifnum\loopcount>14 \relax%
        \xdef\eighth@{F}%
     \else%
        \xdef\eighth@{\the\loopcount}%
     \fi\fi\fi\fi\fi\fi%
     \xdef\seventh@{\seventh@\eighth@}%
     \multiply\loopcount by -256\relax%
     \advance\loopcount by \sixth@%
     \xdef\sixth@{\the\loopcount}%
     \divide\loopcount by 16\relax%
     \ifnum\loopcount=10\relax%
        \xdef\eighth@{A}%
     \else%
     \ifnum\loopcount=11\relax%
        \xdef\eighth@{B}%
     \else%
     \ifnum\loopcount=12\relax%
        \xdef\eighth@{C}%
     \else%
     \ifnum\loopcount=13\relax%
        \xdef\eighth@{D}%
     \else%
     \ifnum\loopcount=14\relax%
        \xdef\eighth@{E}%
     \else%
     \ifnum\loopcount>14\relax%
        \xdef\eighth@{F}%
     \else%
        \xdef\eighth@{\the\loopcount}%
     \fi\fi\fi\fi\fi\fi%
     \xdef\seventh@{\seventh@\eighth@}%
     \multiply\loopcount by -16\relax%
     \advance\loopcount by \sixth@\relax%
     \ifnum\loopcount=10\relax%
        \xdef\eighth@{A}%
     \else%
     \ifnum\loopcount=11\relax%
        \xdef\eighth@{B}%
     \else%
     \ifnum\loopcount=12\relax%
        \xdef\eighth@{C}%
     \else%
     \ifnum\loopcount=13\relax%
        \xdef\eighth@{D}%
     \else%
     \ifnum\loopcount=14\relax%
        \xdef\eighth@{E}%
     \else%
     \ifnum\loopcount>14\relax%
        \xdef\eighth@{F}%
     \else%
        \xdef\eighth@{\the\loopcount}%
     \fi\fi\fi\fi\fi\fi%
     \xdef\seventh@{\seventh@\eighth@}%
}
\def\write@chimera@data{%
      \ifnum\temp@count>0 \else\temp@count=1\relax\temp@@count=0\fi%
      \loopcount=\temp@count%
      \xdef\fourth@{\the\loopcount @}%
      \ifnum\loopcount>9999 \expandafter\decimal@A\fourth@ \else%
      \ifnum\loopcount>999  \expandafter\decimal@B\fourth@ \else%
      \ifnum\loopcount>99   \expandafter\decimal@C\fourth@ \else%
      \ifnum\loopcount>9    \expandafter\decimal@D\fourth@ \else%
                            \expandafter\decimal@E\fourth@%
      \fi\fi\fi\fi%
      \loopcount=\chimera@subfamily@col@black%
      \multiply\loopcount by 65535\relax%
      \divide\loopcount by \temp@count%
      \xdef\chimera@subfamily@col@black{\the\loopcount}%
      \loopcount=65535\relax%
      \advance\loopcount by -\chimera@subfamily@col@black%
      \xdef\temp@@{\the\loopcount}%
%%%%%%%%%%%%%%
%	 *  This code was taken from:
%	 * 
%      *  Color Space - A RGB to RYB converter.
%      *  
%      *  Copyright (C) 2013 Joshua Gentry
%      *
%      *  This program is free software: you can redistribute it and/or modify
%      *  it under the terms of the GNU Affero General Public License as
%      *  published by the Free Software Foundation, either version 3 of the
%      *  License, or (at your option) any later version.
%.     *
%      *  This program is distributed in the hope that it will be useful,
%      *  but WITHOUT ANY WARRANTY; without even the implied warranty of
%      *  MERCHANTABILITY or FITNESS FOR A PARTICULAR PURPOSE.  See the
%      *  GNU Affero General Public License for more details.
%.     *
%      *  You should have received a copy of the GNU Affero General Public License
%      *  along with this program.  If not, see <http://www.gnu.org/licenses/>.
%      */
%	RgbRyb.prototype.setRyb = function(iRed, iYellow, iBlue)
%	{
%		//-----------------------------------------------------------------------------------------
%		// Validate the arguments.
%		VALIDATOR.validate(E10002, "setRyb", arguments, [
%		    {type: "number", max: 255, min: 0, arrayType: "int"},
%		    {type: "number", max: 255, min: 0, arrayType: "int"},
% 		    {type: "number", max: 255, min: 0, arrayType: "int"}
% 		]);
%
%		//-----------------------------------------------------------------------------------------
%		// Save the RYB
%		this.pRyb = [iRed, iYellow, iBlue];
%		
%		// Remove the whiteness from the color.
%		var iWhite = Math.min(iRed, iYellow, iBlue);
%		
%		iRed    -= iWhite;
%		iYellow -= iWhite;
%		iBlue   -= iWhite;
%
%		var iMaxYellow = Math.max(iRed, iYellow, iBlue);
%
%		// Get the green out of the yellow and blue
%		var iGreen = Math.min(iYellow, iBlue);
%		
%		iYellow -= iGreen;
%		iBlue   -= iGreen;
%
%		if (iBlue > 0 && iGreen > 0)
%		{
%			iBlue  *= 2.0;
%			iGreen *= 2.0;
%		}
%		
%		// Redistribute the remaining yellow.
%		iRed   += iYellow;
%		iGreen += iYellow;
%
%		// Normalize to values.
%		var iMaxGreen = Math.max(iRed, iGreen, iBlue);
%		
%		if (iMaxGreen > 0)
%		{
%			var iN = iMaxYellow / iMaxGreen;
%			
%			iRed   *= iN;
%			iGreen *= iN;
%			iBlue  *= iN;
%		}
%		
%		// Add the white back in.
%		iRed   += iWhite;
%		iGreen += iWhite;
%		iBlue  += iWhite;
%
%		// Save the RGB
%		this.pRgb = [Math.floor(iRed), Math.floor(iGreen), Math.floor(iBlue)];
%	}
%%%%%%%%%%%%%%
      \loopcount=\chimera@subfamily@col@blue\relax%
      \ifnum\chimera@subfamily@col@yellow>\loopcount\relax%
         \loopcount=\chimera@subfamily@col@yellow\relax%
      \else%
         \ifnum\chimera@subfamily@col@red>\loopcount\relax%
            \loopcount=\chimera@subfamily@col@red\relax%
         \fi%
      \fi%
      \xdef\max@yellow{\the\loopcount}%
      \ifnum\chimera@subfamily@col@blue<\chimera@subfamily@col@yellow\relax%
         \xdef\chimera@subfamily@col@green{\chimera@subfamily@col@blue}%
      \else%
         \xdef\chimera@subfamily@col@green{\chimera@subfamily@col@yellow}%
      \fi%
      \loopcount=\chimera@subfamily@col@yellow\relax%
      \advance\loopcount by -\chimera@subfamily@col@green%
      \xdef\chimera@subfamily@col@yellow{\the\loopcount}%
      \loopcount=\chimera@subfamily@col@blue\relax%
      \advance\loopcount by -\chimera@subfamily@col@green%
      \xdef\chimera@subfamily@col@blue{\the\loopcount}%
      \ifnum\chimera@subfamily@col@blue>0\relax%
         \ifnum\chimera@subfamily@col@green>0\relax%
            \loopcount=\chimera@subfamily@col@blue\relax%
            \multiply\loopcount by 2\relax \xdef\chimera@subfamily@col@blue{\the\loopcount}%
            \loopcount=\chimera@subfamily@col@green\relax%
            \multiply\loopcount by 2\relax \xdef\chimera@subfamily@col@green{\the\loopcount}%
         \fi%
      \fi%
      \loopcount=\chimera@subfamily@col@green\relax%
      \advance\loopcount by \chimera@subfamily@col@yellow%
      \xdef\chimera@subfamily@col@green{\the\loopcount}%
      \loopcount=\chimera@subfamily@col@red\relax%
      \advance\loopcount by \chimera@subfamily@col@yellow%
      \xdef\chimera@subfamily@col@red{\the\loopcount}%
      \loopcount=\chimera@subfamily@col@blue\relax%
      \ifnum\chimera@subfamily@col@green>\loopcount\relax%
         \loopcount=\chimera@subfamily@col@green\relax%
      \else%
         \ifnum\chimera@subfamily@col@red>\loopcount\relax%
            \loopcount=\chimera@subfamily@col@red\relax%
         \fi%
      \fi%
      \xdef\max@green{\the\loopcount}%
      \loopcount=\chimera@subfamily@col@blue%
      \ifnum\max@green>0\relax%
         \multiply\loopcount by \max@yellow%
         \divide\loopcount by \max@green%
      \fi%
      \advance\loopcount by \chimera@subfamily@col@white%
      \multiply\loopcount by \temp@@%
      \divide\loopcount by \temp@count%
      \make@HEX%
      \xdef\nineth@{\seventh@}%
      \loopcount=\chimera@subfamily@col@green%
      \ifnum\max@green>0\relax%
         \multiply\loopcount by \max@yellow%
         \divide\loopcount by \max@green%
      \fi%
      \advance\loopcount by \chimera@subfamily@col@white%
      \multiply\loopcount by \temp@@%
      \divide\loopcount by \temp@count%
      \make@HEX%
      \xdef\nineth@{\seventh@\nineth@}%
      \loopcount=\chimera@subfamily@col@red%
      \ifnum\max@green>0\relax%
         \multiply\loopcount by \max@yellow%
         \divide\loopcount by \max@green%
      \fi%
      \advance\loopcount by \chimera@subfamily@col@white%
      \multiply\loopcount by \temp@@%
      \divide\loopcount by \temp@count%
      \make@HEX%
      \xdef\nineth@{\seventh@\nineth@}%
                     \ifx\chimera@labelOne@set\n@%
                        \xdef\chimera@labelOne@set{y}%
                        \xdef\chimera@label@subfamily{.\chimera@label@subfamily}%
                     \fi%
      \immediate\write\exp@rt@chimerafile{setattr a radius \temp@ \unexpanded{} #0:\chimera@res@num\chimerach@in @ca}%
      \immediate\write\exp@rt@chimerafile{setattr a color #\nineth@ \unexpanded{} #0:\chimera@res@num\chimerach@in @ca}%
      \immediate\write\exp@rt@chimerafile{setattr a name '|| \chimera@label@subfamily' #0:\chimera@res@num\chimerach@in @ca}%
      \ifshow@sublogo
         \ifnum\temp@@count>0 \else\temp@@count=1\fi%
         \loopcount=\temp@@count%
         \xdef\fourth@{\the\loopcount @}%
         \ifnum\loopcount>9999 \expandafter\decimal@A\fourth@ \else%
         \ifnum\loopcount>999  \expandafter\decimal@B\fourth@ \else%
         \ifnum\loopcount>99   \expandafter\decimal@C\fourth@ \else%
         \ifnum\loopcount>9    \expandafter\decimal@D\fourth@ \else%
                               \expandafter\decimal@E\fourth@%
         \fi\fi\fi\fi%
         \loopcount=-\chimera@remaining@col@red\relax%
         \xdef\chimera@remaining@col@red{\the\loopcount}%
         \loopcount=-\chimera@remaining@col@yellow\relax%
         \xdef\chimera@remaining@col@yellow{\the\loopcount}%
         \loopcount=-\chimera@remaining@col@blue\relax%
         \xdef\chimera@remaining@col@blue{\the\loopcount}%
         \loopcount=-\chimera@remaining@col@white\relax%
         \xdef\chimera@remaining@col@white{\the\loopcount}%
         \loopcount=-\chimera@remaining@col@black%
         \multiply\loopcount by 65535\relax%
         \divide\loopcount by \temp@count%
         \xdef\chimera@remaining@col@black{\the\loopcount}%
         \loopcount=65535\relax%
         \advance\loopcount by -\chimera@remaining@col@black%
         \xdef\temp@@{\the\loopcount}%
         \loopcount=\chimera@remaining@col@blue\relax%
         \ifnum\chimera@remaining@col@yellow>\loopcount\relax%
            \loopcount=\chimera@remaining@col@yellow\relax%
         \else%
            \ifnum\chimera@remaining@col@red>\loopcount\relax%
               \loopcount=\chimera@remaining@col@red\relax%
            \fi%
         \fi%
         \xdef\max@yellow{\the\loopcount}%
         \ifnum\chimera@remaining@col@blue<\chimera@remaining@col@yellow\relax%
            \xdef\chimera@remaining@col@green{\chimera@remaining@col@blue}%
         \else%
            \xdef\chimera@remaining@col@green{\chimera@remaining@col@yellow}%
         \fi%
         \loopcount=\chimera@remaining@col@yellow\relax%
         \advance\loopcount by -\chimera@remaining@col@green%
         \xdef\chimera@remaining@col@yellow{\the\loopcount}%
         \loopcount=\chimera@remaining@col@blue\relax%
         \advance\loopcount by -\chimera@remaining@col@green%
         \xdef\chimera@remaining@col@blue{\the\loopcount}%
         \ifnum\chimera@remaining@col@blue>0\relax%
            \ifnum\chimera@remaining@col@green>0\relax%
               \loopcount=\chimera@remaining@col@blue\relax%
               \multiply\loopcount by 2\relax \xdef\chimera@remaining@col@blue{\the\loopcount}%
               \loopcount=\chimera@remaining@col@green\relax%
               \multiply\loopcount by 2\relax \xdef\chimera@remaining@col@green{\the\loopcount}%
            \fi%
         \fi%
         \loopcount=\chimera@remaining@col@green\relax%
         \advance\loopcount by \chimera@remaining@col@yellow%
         \xdef\chimera@remaining@col@green{\the\loopcount}%
         \loopcount=\chimera@remaining@col@red\relax%
         \advance\loopcount by \chimera@remaining@col@yellow%
         \xdef\chimera@remaining@col@red{\the\loopcount}%
         \loopcount=\chimera@remaining@col@blue\relax%
         \ifnum\chimera@remaining@col@green>\loopcount\relax%
            \loopcount=\chimera@remaining@col@green\relax%
         \else%
            \ifnum\chimera@remaining@col@red>\loopcount\relax%
               \loopcount=\chimera@remaining@col@red\relax%
            \fi%
         \fi%
         \xdef\max@green{\the\loopcount}%
         \loopcount=\chimera@remaining@col@blue%
         \ifnum\max@green>0\relax%
            \multiply\loopcount by \max@yellow%
            \divide\loopcount by \max@green%
         \fi%
         \advance\loopcount by \chimera@remaining@col@white%
         \multiply\loopcount by \temp@@%
         \divide\loopcount by \temp@count%
         \make@HEX%
         \xdef\nineth@{\seventh@}%
         \loopcount=\chimera@remaining@col@green%
         \ifnum\max@green>0\relax%
            \multiply\loopcount by \max@yellow%
            \divide\loopcount by \max@green%
         \fi%
         \advance\loopcount by \chimera@remaining@col@white%
         \multiply\loopcount by \temp@@%
         \divide\loopcount by \temp@count%
         \make@HEX%
         \xdef\nineth@{\seventh@\nineth@}%
         \loopcount=\chimera@remaining@col@red%
         \ifnum\max@green>0\relax%
            \multiply\loopcount by \max@yellow%
            \divide\loopcount by \max@green%
         \fi%
         \advance\loopcount by \chimera@remaining@col@white%
         \multiply\loopcount by \temp@@%
         \divide\loopcount by \temp@count%
         \make@HEX%
         \xdef\nineth@{\seventh@\nineth@}%
                     \ifx\chimera@label@remainingOne@set\n@%
                        \xdef\chimera@label@remainingOne@set{y}%
                        \xdef\chimera@label@remaining{.\chimera@label@remaining}%
                     \fi%
         \immediate\write\exp@rt@chimerafile{setattr a radius \temp@ \unexpanded{} #1:\chimera@res@num\chimerach@in @ca}%
         \immediate\write\exp@rt@chimerafile{setattr a color #\nineth@ \unexpanded{} #1:\chimera@res@num\chimerach@in @ca}%
         \immediate\write\exp@rt@chimerafile{setattr a name '|| \chimera@label@remaining' #1:\chimera@res@num\chimerach@in @ca}%
      \fi%
}
\catcode`\"=12 \catcode`\#=6 \catcode`\~=13
\def\get@chimera@res@num#1,#2@{\xdef\chimera@res@num{#1}\xdef\chimera@stack{#2}}
%%%%%%%%%%%%%%%%%%%%%%%%%%%%%%%%%%%
%%%%                     Chimera export routines   end                %%%%
%%%%%%%%%%%%%%%%%%%%%%%%%%%%%%%%%%%
\def\next@logo#1:#2,#3@{%
   \xdef\first@{#1}\xdef\second@{#2}\xdef\last@{#3}%
   \ifx\first@\ampers@nd%
      \ifx\exp@rt@chimera\y@%
        \xdef\chimera@stack{\chimera@stack @}%
        \expandafter\get@chimera@res@num\chimera@stack%
        \ifx\chimera@res@num\ampers@nd%
        \else%
            \write@chimera@data%
        \fi%
        \xdef\chimera@label@subfamily{}%
        \xdef\chimera@label@remaining{}%
        \xdef\chimera@subfamily@col@red{0}%
        \xdef\chimera@subfamily@col@blue{0}%
        \xdef\chimera@subfamily@col@yellow{0}%
        \xdef\chimera@subfamily@col@black{0}%
        \xdef\chimera@subfamily@col@white{0}%
        \xdef\chimera@remaining@col@red{0}%
        \xdef\chimera@remaining@col@blue{0}%
        \xdef\chimera@remaining@col@yellow{0}%
        \xdef\chimera@remaining@col@black{0}%
        \xdef\chimera@remaining@col@white{0}%
        \xdef\chimera@labelOne@set{n}%
        \xdef\chimera@labelTwo@set{n}%
        \xdef\chimera@labelThree@set{n}%
        \xdef\chimera@label@remainingOne@set{n}%
        \xdef\chimera@label@remainingTwo@set{n}%
        \xdef\chimera@label@remainingThree@set{n}%
      \fi%
     \ifx\clear@logo\n@%
     \ifx\hide@sig\n@%
      \expandafter\firstchar@get\sublogo@sig%
      \xdef\sublogo@sig{\third@}%
      \ifx\first@\y@%
         \loopcount=\temp@count%
         \ifnum\loopcount<0 \loopcount=0\fi%
         \xdef\fourth@{\the\loopcount @}%
         \ifnum\loopcount>9999 \expandafter\decimal@A\fourth@ \else%
         \ifnum\loopcount>999  \expandafter\decimal@B\fourth@ \else%
         \ifnum\loopcount>99   \expandafter\decimal@C\fourth@ \else%
         \ifnum\loopcount>9    \expandafter\decimal@D\fourth@ \else%
                               \expandafter\decimal@E\fourth@%
         \fi\fi\fi\fi%
         \raisebox{\temp@\logo@height}{%
         \hbox to \box@width{\textcolor{\sig@color}{\hss\residues@size{\sig@char}\hss}}}%
         \kern-\box@width%
      \fi%
     \fi%
     \fi%
     \temp@count=0%
     \temp@@count=0%
     \kern\box@width%
     \advance\outerloopcount by 1 %
   \else%
      \if\seq@type N%
         \loopcount=\second@%
         \multiply\loopcount by 2%
         \xdef\second@{\the\loopcount}%
      \fi%
      \xdef\third@{\second@}%
      \if\first@ C\correct@CGSO \else%
      \if\first@ G\correct@CGSO \else%
      \if\first@ S\correct@CGSO \else%
      \if\first@ O\correct@CGSO \else%
      \if\first@ Q\correct@Q    \else%
      \if\first@ J\correct@JUV  \else%
      \if\first@ U\correct@JUV  \else%
      \if\first@ V\correct@JUV  \else%
                  \setlength\temp@@length{0pt}%
      \fi\fi\fi\fi\fi\fi\fi\fi%
      \loopcount=\third@
      \multiply\loopcount by \logo@stretch@IOOO%
      \ifnum\loopcount>0 %
         \divide\loopcount by 1000 \xdef\fl@g{}%
         \xdef\tint@{}%
      \else%
         \divide\loopcount by -1000 \xdef\fl@g{-}%
         \xdef\tint@{\sublogo@tint}%
      \fi%
      \ifnum\loopcount=0\relax\loopcount=1\relax\fi
      \xdef\third@{\the\loopcount @}%
      \ifnum\loopcount>9999 \expandafter\decimal@A\third@ \else%
      \ifnum\loopcount>999  \expandafter\decimal@B\third@ \else%
      \ifnum\loopcount>99   \expandafter\decimal@C\third@ \else%
      \ifnum\loopcount>9    \expandafter\decimal@D\third@ \else%
                            \expandafter\decimal@E\third@%
      \fi\fi\fi\fi%
      \xdef\third@{\fl@g\temp@}%
      \ifnum\temp@count<0 %
         \ifnum\second@<0 \else\temp@@count=-\temp@count\relax\temp@count=0 \fi%
      \fi%
      \ifnum\temp@count>0 %
         \loopcount=\temp@count \xdef\fl@g{}%
      \else
         \loopcount=-\temp@count \xdef\fl@g{-}%
      \fi%
      \xdef\fourth@{\the\loopcount @}%
      \ifnum\loopcount>9999 \expandafter\decimal@A\fourth@ \else%
      \ifnum\loopcount>999  \expandafter\decimal@B\fourth@ \else%
      \ifnum\loopcount>99   \expandafter\decimal@C\fourth@ \else%
      \ifnum\loopcount>9    \expandafter\decimal@D\fourth@ \else%
                             \expandafter\decimal@E\fourth@%
      \fi\fi\fi\fi%
      \xdef\fourth@{\fl@g\temp@}%
      \if\first@ W%  
         \hbox to \box@width{\hss%
         \raisebox{\temp@@length}{\raisebox{\fourth@\logo@height}{\scalebox{\char@stretch@W}[\third@]%
         {\hbox to \box@width{\textcolor{\tint@\csname logo@col\first@\endcsname}{\hss\residues@size{W}\hss}}}}}\hss}%
         \kern-\box@width%
      \else%
         \hbox to \box@width{\hss%
         \raisebox{\temp@@length}{\raisebox{\fourth@\logo@height}{\scalebox{\char@stretch}[\third@]%
         {\hbox to \box@width{\textcolor{\tint@\csname logo@col\first@\endcsname}{\hss\residues@size{\first@}\hss}}}}}\hss}%
         \kern-\box@width%
      \fi%
      \advance\temp@count by \second@%
      \xdef\temp@{\the\temp@count}%
      \ifx\exp@rt@chimera\y@%
         \xdef\temp@@{\csname chimera@col@\first@\endcsname}%
         \ifnum\second@<0 %
            \xdef\temp@@@{\csname chimera@remaining@col@\temp@@\endcsname}%
            \temp@count=\temp@@@%
            \advance\temp@count by \second@%
            \expandafter\xdef\csname chimera@remaining@col@\temp@@\endcsname{\the\temp@count}%
            \ifnum\second@<\label@threshold@lowerremainingIOOO%
                  \ifnum\second@<-999%
                     \ifx\chimera@label@remainingOne@set\n@%
                        \xdef\chimera@label@remainingOne@set{y}%
                        \xdef\chimera@label@remaining{.\chimera@label@remaining}%
                     \fi%
                  \fi%
                  \ifnum\second@<-1999%
                     \ifx\chimera@label@remainingTwo@set\n@%
                        \xdef\chimera@label@remainingTwo@set{y}%
                        \xdef\chimera@label@remaining{:\chimera@label@remaining}%
                     \fi%
                  \fi%
                  \ifnum\second@<-2999%
                     \ifx\chimera@label@remainingThree@set\n@%
                        \xdef\chimera@label@remainingThree@set{y}%
                        \xdef\chimera@label@remaining{>\chimera@label@remaining}%
                     \fi%
                  \fi%
                  \xdef\chimera@label@remaining{\first@\chimera@label@remaining}%
            \fi%
         \else%
            \xdef\temp@@@{\csname chimera@subfamily@col@\temp@@\endcsname}%
            \temp@count=\temp@@@%
            \advance\temp@count by \second@%
            \expandafter\xdef\csname chimera@subfamily@col@\temp@@\endcsname{\the\temp@count}%
            \ifnum\second@>\label@threshold@lowerIOOO%
                  \ifnum\second@>999%
                     \ifx\chimera@labelOne@set\n@%
                        \xdef\chimera@labelOne@set{y}%
                        \xdef\chimera@label@subfamily{.\chimera@label@subfamily}%
                     \fi%
                  \fi%
                  \ifnum\second@>1999%
                     \ifx\chimera@labelTwo@set\n@%
                        \xdef\chimera@labelTwo@set{y}%
                        \xdef\chimera@label@subfamily{:\chimera@label@subfamily}%
                     \fi%
                  \fi%
                  \ifnum\second@>2999%
                     \ifx\chimera@labelThree@set\n@%
                        \xdef\chimera@labelThree@set{y}%
                        \xdef\chimera@label@subfamily{>\chimera@label@subfamily}%
                     \fi%
                  \fi%
                  \xdef\chimera@label@subfamily{\first@\chimera@label@subfamily}%
            \fi%
         \fi%
      \fi%
      \temp@count=\temp@%
   \fi%
   \xdef\temp@{\last@.}%
   \ifx\temp@\d@t%
   \else%
      \xdef\last@{#3@}%
      \ifnum\outerloopcount>\res@perline%
      \else
         \expandafter\next@logo\last@%
      \fi%
   \fi%
}

\def\put@logo{%
   \ifx\clear@logo\y@%
   \ifnum\logo@top=1 %
      \vspace{-0.5\baselineskip}%
   \else
      \vspace{-0.75\baselineskip}%
   \fi%
   \else%
   \ifnum\sublogo@top=1 %
      \vspace{-0.5\baselineskip}%
   \else
      \vspace{-0.75\baselineskip}%
   \fi%
   \fi%
   \newline\hbox{}%
   \ifnames@right%
   \else%
      \fontfamily{\namestext@family}\fontseries{\namestext@series}\fontshape{\namestext@shape}%
      \selectfont%
      \ifnames@
         \textcolor{\names@fg}{%
            \raisebox{2\logo@height}{\hbox to \name@width{\namestext@size{\logo@name}\hss}}%
            \ifx\clear@logo\n@
            \ifx\hide@negatives\n@%
               \kern-\name@width\raisebox{-3\logo@height}{\hbox to \name@width{\namestext@size{\sublogo@name@neg}\hss}}%
            \fi%
            \fi}%
      \fi%
   \fi%
   \ifnumbers@left\ifnumbers@\hbox to \number@width{\hss}\fi\fi%
   \fontfamily{cmss}\fontseries{m}\fontshape{n}\selectfont%
   \ifx\show@logoscale\n@%
      \if\seq@type P%
         \ifnum\corr@max>600\relax%
            \rule[5.4\logo@height]{0pt}{0pt}%
         \else%
            \rule[4.4\logo@height]{0pt}{0pt}%
         \fi%
      \else%
         \rule[4.0688\logo@height]{0pt}{0pt}%
      \fi%
   \else%
   \xdef\first@{right}%
   \ifx\show@logoscale\first@%
   \else%
      \textcolor{\logo@scalecol}{%
      \kern-0.7\box@width%
      \rule[-0.05\logo@height]{0.4\box@width}{0.1\box@width}\kern-0.4\box@width%
      \rule[2\logo@height]{0.4\box@width}{0.1\box@width}\kern-0.4\box@width%
      \rule[4\logo@height]{0.4\box@width}{0.1\box@width}%
      \if\seq@type P%
         \kern-0.4\box@width%
         \rule[1\logo@height]{0.4\box@width}{0.1\box@width}\kern-0.4\box@width%
         \rule[3\logo@height]{0.4\box@width}{0.1\box@width}%
         \ifnum\corr@max>600\relax%
            \kern-0.4\box@width\rule[5\logo@height]{0.4\box@width}{0.1\box@width}\fi%
         \setbox1=\hbox{\bottomruler@size 2}%
         \kern-1.2\box@width%
         \ifdim\logo@height>1.25\ht1 \relax%
            \raisebox{-0.5\ht1}{\raisebox{1\logo@height}{\hbox to 0.6\box@width{\bottomruler@size\hss 1}}}%
            \kern-0.6\box@width%
            \raisebox{-0.5\ht1}{\raisebox{3\logo@height}{\hbox to 0.6\box@width{\bottomruler@size\hss 3}}}%
            \kern-0.6\box@width%
            \ifnum\corr@max>600\relax%
               \raisebox{-0.5\ht1}{\raisebox{5\logo@height}{\hbox to 0.6\box@width{\bottomruler@size\hss 5}}}%
               \kern-0.6\box@width%
            \fi
         \fi
         \raisebox{-0.5\ht1}{\raisebox{2\logo@height}{\hbox to 0.6\box@width{\bottomruler@size\hss 2}}}%
         \kern-0.6\box@width%
         \raisebox{-0.5\ht1}{\raisebox{4\logo@height}{\hbox to 0.6\box@width{\bottomruler@size\hss 4}}}%
         \kern0.6\box@width%
         \ifnum\corr@max>600\relax%
            \rule[5.4\logo@height]{0pt}{0pt}%
         \else
            \rule[4.4\logo@height]{0pt}{0pt}%
         \fi
      \else%
         \setbox1=\hbox{\bottomruler@size 2}%
         \kern-1.2\box@width%
         \raisebox{-0.5\ht1}{\raisebox{2\logo@height}{\hbox to 0.6\box@width{\bottomruler@size\hss 1}}}%
         \kern-0.6\box@width%
         \raisebox{-0.5\ht1}{\raisebox{4\logo@height}{\hbox to 0.6\box@width{\bottomruler@size\hss 2}}}%
         \kern0.6\box@width%
         \rule[4.0688\logo@height]{0pt}{0pt}%
      \fi%
      \kern-0.1\box@width%
      \ifnum\corr@max>600\relax%
         \rule{0.1\box@width}{5\logo@height}%
      \else%
         \rule{0.1\box@width}{4\logo@height}%
      \fi%
      \ifx\clear@logo\n@%
      \ifx\hide@negatives\n@%
      \kern-0.4\box@width%
      \rule[-2\logo@height]{0.4\box@width}{0.1\box@width}\kern-0.4\box@width%
      \rule[-4\logo@height]{0.4\box@width}{0.1\box@width}%
      \if\seq@type P%
         \kern-0.4\box@width%
         \rule[-\logo@height]{0.4\box@width}{0.1\box@width}\kern-0.4\box@width%
         \rule[-3\logo@height]{0.4\box@width}{0.1\box@width}%
         \ifnum\corr@max>600\relax%
            \kern-0.4\box@width\rule[-5\logo@height]{0.4\box@width}{0.1\box@width}\fi%
         \setbox1=\hbox{\bottomruler@size 2}%
         \kern-1.2\box@width%
         \ifdim\logo@height>1.25\ht1 \relax%
            \raisebox{-0.5\ht1}{\raisebox{-\logo@height}{\hbox to 0.6\box@width{\bottomruler@size\hss -1}}}%
            \kern-0.6\box@width%
            \raisebox{-0.5\ht1}{\raisebox{-3\logo@height}{\hbox to 0.6\box@width{\bottomruler@size\hss -3}}}%
            \kern-0.6\box@width%
            \ifnum\corr@max>600\relax%
               \raisebox{-0.5\ht1}{\raisebox{-5\logo@height}{\hbox to 0.6\box@width{\bottomruler@size\hss -5}}}%
               \kern-0.6\box@width%
            \fi
         \fi
         \raisebox{-0.5\ht1}{\raisebox{-2\logo@height}{\hbox to 0.6\box@width{\bottomruler@size\hss -2}}}%
         \kern-0.6\box@width%
         \raisebox{-0.5\ht1}{\raisebox{-4\logo@height}{\hbox to 0.6\box@width{\bottomruler@size\hss -4}}}%
         \kern0.6\box@width%
         \ifnum\corr@max>600\relax%
            \rule[-5.4\logo@height]{0pt}{0pt}%
         \else
            \rule[-4.4\logo@height]{0pt}{0pt}%
         \fi
      \else%
         \setbox1=\hbox{\bottomruler@size 2}%
         \kern-1.2\box@width%
         \raisebox{-0.5\ht1}{\raisebox{-2\logo@height}{\hbox to 0.6\box@width{\bottomruler@size\hss -1}}}%
         \kern-0.6\box@width%
         \raisebox{-0.5\ht1}{\raisebox{-4\logo@height}{\hbox to 0.6\box@width{\bottomruler@size\hss -2}}}%
         \kern0.6\box@width%
         \rule[-4.0688\logo@height]{0pt}{0pt}%
      \fi%
      \kern-0.1\box@width%
      \ifnum\corr@max>600\relax%
         \rule[-5\logo@height]{0.1\box@width}{5\logo@height}%
      \else%
         \rule[-4\logo@height]{0.1\box@width}{4\logo@height}%
      \fi%
      \fi%
      \fi
      \kern0.3\box@width}%
   \fi%
   \fi%
   \temp@count=0%
   \outerloopcount=1 \expandafter\next@logo\last@%
   \ifx\clear@logo\n@%
      \ifx\hide@negatives\n@%
         \textcolor{\logo@scalecol}{%
         \kern-\res@perline\box@width\rule[-0.03\box@width]{\res@perline\box@width}{0.06\box@width}}%
      \fi%
   \fi%
   \def\@kern{\kern1em}%
   \ifx\show@logoscale\n@%
   \else%
   \xdef\first@{left}%
   \ifx\show@logoscale\first@%
   \kern0.1\box@width%
   \else%
      \def\@kern{\kern2em}%
      \textcolor{\logo@scalecol}{%
      \kern0.3\box@width%
      \rule[-0.05\logo@height]{0.4\box@width}{0.1\box@width}\kern-0.4\box@width%
      \rule[2\logo@height]{0.4\box@width}{0.1\box@width}\kern-0.4\box@width%
      \rule[4\logo@height]{0.4\box@width}{0.1\box@width}%
      \if\seq@type P%
         \kern-0.4\box@width%
         \rule[1\logo@height]{0.4\box@width}{0.1\box@width}\kern-0.4\box@width%
         \rule[3\logo@height]{0.4\box@width}{0.1\box@width}%
         \ifnum\corr@max>600\relax%
            \kern-0.4\box@width\rule[5\logo@height]{0.4\box@width}{0.1\box@width}\fi%
         \setbox1=\hbox{\bottomruler@size 2}%
         \kern0.2\box@width%
         \ifdim\logo@height>1.25\ht1 \relax%
            \raisebox{-0.5\ht1}{\raisebox{1\logo@height}{\hbox to 0.6\box@width{\bottomruler@size 1 \hss}}}%
            \kern-0.6\box@width%
            \raisebox{-0.5\ht1}{\raisebox{3\logo@height}{\hbox to 0.6\box@width{\bottomruler@size 3 \hss}}}%
            \kern-0.6\box@width%
            \ifnum\corr@max>600\relax%
               \raisebox{-0.5\ht1}{\raisebox{5\logo@height}{\hbox to 0.6\box@width{\bottomruler@size 5 \hss}}}%
               \kern-0.6\box@width%
            \fi
         \fi
         \raisebox{-0.5\ht1}{\raisebox{2\logo@height}{\hbox to 0.6\box@width{\bottomruler@size 2 \hss}}}%
         \kern-0.6\box@width%
         \raisebox{-0.5\ht1}{\raisebox{4\logo@height}{\hbox to 0.6\box@width{\bottomruler@size 4 \hss}}}%
         \kern-0.8\box@width%
         \ifnum\corr@max>600\relax%
            \rule[5.4\logo@height]{0pt}{0pt}%
         \else
            \rule[4.4\logo@height]{0pt}{0pt}%
         \fi
      \else%
         \setbox1=\hbox{\bottomruler@size 2}%
         \kern0.2\box@width%
         \raisebox{-0.5\ht1}{\raisebox{2\logo@height}{\hbox to 0.6\box@width{\bottomruler@size 1 \hss}}}%
         \kern-0.6\box@width%
         \raisebox{-0.5\ht1}{\raisebox{4\logo@height}{\hbox to 0.6\box@width{\bottomruler@size 2 \hss}}}%
         \kern-0.8\box@width%
         \rule[4.0688\logo@height]{0pt}{0pt}%
      \fi%
      \kern-0.4\box@width%
      \ifnum\corr@max>600\relax%
         \rule{0.1\box@width}{5\logo@height}%
      \else%
         \rule{0.1\box@width}{4\logo@height}%
      \fi%
      \ifx\clear@logo\n@%
      \ifx\hide@negatives\n@%
      \kern-0.1\box@width%
      \rule[-2\logo@height]{0.4\box@width}{0.1\box@width}\kern-0.4\box@width%
      \rule[-4\logo@height]{0.4\box@width}{0.1\box@width}%
      \if\seq@type P%
         \kern-0.4\box@width%
         \rule[-\logo@height]{0.4\box@width}{0.1\box@width}\kern-0.4\box@width%
         \rule[-3\logo@height]{0.4\box@width}{0.1\box@width}%
         \ifnum\corr@max>600\relax%
            \kern-0.4\box@width\rule[-5\logo@height]{0.4\box@width}{0.1\box@width}\fi%
         \setbox1=\hbox{\bottomruler@size 2}%
         \kern0.2\box@width%
         \ifdim\logo@height>1.25\ht1 \relax%
            \raisebox{-0.5\ht1}{\raisebox{-\logo@height}{\hbox to 0.6\box@width{\bottomruler@size -1 \hss}}}%
            \kern-0.6\box@width%
            \raisebox{-0.5\ht1}{\raisebox{-3\logo@height}{\hbox to 0.6\box@width{\bottomruler@size -3 \hss}}}%
            \kern-0.6\box@width%
            \ifnum\corr@max>600\relax%
               \raisebox{-0.5\ht1}{\raisebox{-5\logo@height}{\hbox to 0.6\box@width{\bottomruler@size -5 \hss}}}%
               \kern-0.6\box@width%
            \fi
         \fi
         \raisebox{-0.5\ht1}{\raisebox{-2\logo@height}{\hbox to 0.6\box@width{\bottomruler@size -2 \hss}}}%
         \kern-0.6\box@width%
         \raisebox{-0.5\ht1}{\raisebox{-4\logo@height}{\hbox to 0.6\box@width{\bottomruler@size -4 \hss}}}%
         \kern-0.8\box@width%
         \ifnum\corr@max>600\relax%
            \rule[-5.4\logo@height]{0pt}{0pt}%
         \else
            \rule[-4.4\logo@height]{0pt}{0pt}%
         \fi
      \else%
         \setbox1=\hbox{\bottomruler@size 2}%
         \kern0.2\box@width%
         \raisebox{-0.5\ht1}{\raisebox{-2\logo@height}{\hbox to 0.6\box@width{\bottomruler@size -1 \hss}}}%
         \kern-0.6\box@width%
         \raisebox{-0.5\ht1}{\raisebox{-4\logo@height}{\hbox to 0.6\box@width{\bottomruler@size -2 \hss}}}%
         \kern-0.8\box@width%
         \rule[-4.0688\logo@height]{0pt}{0pt}%
      \fi%
      \kern-0.4\box@width%
      \ifnum\corr@max>600\relax%
         \rule[-5\logo@height]{0.1\box@width}{5\logo@height}%
      \else%
         \rule[-4\logo@height]{0.1\box@width}{4\logo@height}%
      \fi%
      \fi%
      \fi%
      \kern-0.3\box@width}%
   \fi%
   \fi%
   \ifnumbers@right\ifnumbers@\hbox to \number@width{\hss}\def\@kern{\kern1em}\fi\fi%
   \ifnames@right%
      \fontfamily{\namestext@family}\fontseries{\namestext@series}\fontshape{\namestext@shape}%
      \selectfont%
      \ifnames@
         \textcolor{\names@fg}{%
            \raisebox{2\logo@height}{\hbox to \name@width{\namestext@size{\@kern\logo@name}\hss}}%
            \ifx\clear@logo\n@
            \ifx\hide@negatives\n@%
               \kern-\name@width\raisebox{-3\logo@height}{\hbox to \name@width{\namestext@size{\@kern\sublogo@name@neg}\hss}}%
            \fi%
            \fi}%
      \fi%
   \fi%
   \ifx\clear@logo\y@\xdef\stack@sequencelogo{}%
   \else%
      \ifx\hide@negatives\n@ \newline\hbox{}\fi%
   \fi%
   \newline\hbox{}%
}

\def\set@lines{%
   \pos@count=1
   \xdef\frame@pos{1}\xdef\bar@pos{0}%
   \xdef\featurepostop{1}    \xdef\featureposbottom{1}%
   \xdef\featureposttop{1}   \xdef\featureposbbottom{1}%
   \xdef\featurepostttop{1}  \xdef\featureposbbbottom{1}%
   \xdef\featureposttttop{1} \xdef\featureposbbbbottom{1}%
   \findc@nsensus%
   \noindent%
   \setlength{\bar@length}{0pt}%
   \xdef\fe@turep@s{top}%
   \ifnum\feature@ttttop=1
      \advance\bar@length by \feature@tttop\baselineskip%
      \advance\bar@length by \feature@ttop\baselineskip%
      \advance\bar@length by \feature@top\baselineskip%
      \multiply\bar@length by 2%
      \advance\bar@length by \tttt@sp@ce%
      \advance\bar@length by \ttt@sp@ce%
      \advance\bar@length by \tt@sp@ce%
      \ifnum\featureonttttop=0 \xdef\feature@ttttop{0} \fi
      \xdef\bottop@{ttttop}%
      \xdef\shift@feature{n}%
      \put@feature%
      \advance\bar@length by -\baselineskip%
      \ifx\shift@feature\y@ \vspace{-\bar@length}\fi%
      \setlength{\bar@length}{0pt}%
      \vspace{\tttt@sp@ce}%
      \vspace{-\baselineskip}%
      \newline\hbox{}%
   \else
      \iffix@\ifttttopfeature \vspace{\tttt@sp@ce}\newline\hbox{}\newline\hbox{}\fi\fi
   \fi
   \ifnum\feature@tttop=1
      \advance\bar@length by \feature@ttop\baselineskip%
      \advance\bar@length by \feature@top\baselineskip%
      \multiply\bar@length by 2%
      \advance\bar@length by \ttt@sp@ce%
      \advance\bar@length by \tt@sp@ce%
      \ifnum\featureontttop=0 \xdef\feature@tttop{0} \fi
      \xdef\bottop@{tttop}%
      \xdef\shift@feature{n}%
      \put@feature%
      \advance\bar@length by -\baselineskip%
      \ifx\shift@feature\y@ \vspace{-\bar@length}\fi%
      \setlength{\bar@length}{0pt}%
      \vspace{\ttt@sp@ce}%
      \vspace{-\baselineskip}%
      \newline\hbox{}%
   \else
      \iffix@\iftttopfeature \vspace{\ttt@sp@ce}\newline\hbox{}\newline\hbox{}\fi\fi
   \fi
   \ifnum\feature@ttop=1
      \advance\bar@length by \feature@top\baselineskip%
      \multiply\bar@length by 2%
      \advance\bar@length by \tt@sp@ce\message{tt\the\bar@length tt}%
      \ifnum\featureonttop=0 \xdef\feature@ttop{0} \fi
      \xdef\bottop@{ttop}%
      \xdef\shift@feature{n}%
      \put@feature%
      \advance\bar@length by -\baselineskip%
      \ifx\shift@feature\y@ \vspace{-\bar@length}\fi%
      \setlength{\bar@length}{0pt}%
      \vspace{\tt@sp@ce}%
      \vspace{-\baselineskip}%
      \newline\hbox{}%
   \else
      \iffix@\ifttopfeature \vspace{\tt@sp@ce}\newline\hbox{}\newline\hbox{}\fi\fi
   \fi
   \ifnum\feature@top=1
      \ifnum\featureontop=0 \xdef\feature@top{0} \fi
      \xdef\bottop@{top}%
      \put@feature%
      \vspace{\t@sp@ce}%
      \vspace{-\baselineskip}%
      \newline\hbox{}%
   \else
      \iffix@\iftopfeature \vspace{\t@sp@ce}\newline\hbox{}\newline\hbox{}\fi\fi
   \fi
   \ifnum\rule@num<0 \else \ifnum\rule@top=0 \loopcount=0 \put@ruler \fi\fi
   \ifshow@logo\ifnum\logo@top=0
      \xdef\last@{\stack@sequencelogo @}%
      \ifx\logo@name@user\ampers@nd\xdef\logo@name{logo}\else\xdef\logo@name{\logo@name@user}\fi
      \xdef\clear@logo{y}%
      \put@logo%
      \xdef\clear@logo{n}%
   \fi\fi
   \ifshow@sublogo\ifnum\sublogo@top=0
      \read@sublogo%
      \xdef\last@{\stack@sublogo @}%
      \ifx\sublogo@name@user\ampers@nd\xdef\logo@name{subfamily}\else\xdef\logo@name{\sublogo@name@user}\fi
      \put@logo%
   \fi\fi
   \set@consensus
   \ifshow@cons\ifnum\cons@top=0
      \xdef\cons@now{y}%
      \xdef\first@{\consensus @}%
      \xdef\first@@@{\c@nscol @}%
      \loopcount=0\relax%
      \put@line%
      \xdef\c@nscol{}%
      \xdef\cons@now{no}%
   \fi\fi
   \ifx\hide@seqs\n@%
   \loopcount=0 %
   \loop
     \advance\loopcount by 1 %
     \xdef\second@{\csname hide@seq\the\loopcount\endcsname}%
     \xdef\first@{noshade}\nosh@defalse%
     \ifx\second@\first@ \nosh@detrue \xdef\second@{false}\fi
     \xdef\first@{false}%
     \ifx\second@\first@
       \ifnames@right
         \ifnumbers@right
            \expandafter\xdef\csname seq\the\loopcount\endcsname{%
               \csname seq\the\loopcount\endcsname%
               \csname res@count\the\loopcount\endcsname)%
               <\csname newseqname\the\loopcount\endcsname>}%
         \else
            \expandafter\xdef\csname seq\the\loopcount\endcsname{%
               \csname seq\the\loopcount\endcsname%
               <\csname newseqname\the\loopcount\endcsname>}%
         \fi
        \else
         \ifnumbers@right
            \expandafter\xdef\csname seq\the\loopcount\endcsname{%
               \csname seq\the\loopcount\endcsname%
               \csname res@count\the\loopcount\endcsname)}%
         \fi
       \fi
       \xdef\first@{\csname seq\the\loopcount\endcsname @}%
       \expandafter\ifx\csname seq@gap\the\loopcount\endcsname\yes%
       \ifnum\finger@linenum=0\seq@skip%
       \else%
           \ifnum\loopcount<\seq@count\seq@skip\fi%
       \fi%
     \fi%
     \put@line%
     \fi%
   \ifnum\loopcount<\seq@count\repeat%
   \fi%
   \ifshow@cons\ifnum\cons@top=1 %
      \xdef\cons@now{y}%
      \xdef\first@{\consensus @}%
      \xdef\first@@@{\c@nscol @}%
      \loopcount=0\relax%
      \put@line%
      \xdef\c@nscol{}%
      \xdef\cons@now{no}%
   \fi\fi%
   \ifnum\frame@=1 %
      \ifnum\frame@on=0 \xdef\frame@{0}\fi%
      \put@frame%
   \fi%
   \ifshow@logo\ifnum\logo@top=1 %
      \xdef\last@{\stack@sequencelogo @}%
      \ifx\logo@name@user\ampers@nd\xdef\logo@name{logo}\else\xdef\logo@name{\logo@name@user}\fi
      \xdef\clear@logo{y}%
      \put@logo%
      \xdef\clear@logo{n}%
   \fi\fi%
   \ifshow@sublogo\ifnum\sublogo@top=1
      \read@sublogo%
      \xdef\last@{\stack@sublogo @}
      \ifx\sublogo@name@user\ampers@nd\xdef\logo@name{subfamily}\else\xdef\logo@name{\sublogo@name@user}\fi
      \put@logo%
   \fi\fi%
   \ifnum\rule@num<0 %
   \else%
      \ifnum\rule@top=1 %
         \loopcount=0\relax%
         \ifnum\ruler@rot=0 %
         \else \vspace{\ruler@width}\vspace{-1.75\baselineskip}\newline\hbox{}%
         \fi
         \put@ruler%
         \ifnum\ruler@rot=0 \vspace{0.25\baselineskip}\fi%
      \fi%
   \fi%
   \xdef\fe@turep@s{bottom}%
   \setlength{\bar@length}{0pt}%
   \xdef\b@feature@count{0}
   \ifnum\feature@bottom=1 %
      \ifnum\featureonbottom=0 \xdef\feature@bottom{0}\fi%
      \xdef\bottop@{bottom}%
      \vspace{\b@sp@ce}%
      \if\bottom@stretch y%
         \vspace{-\box@height}%
         \vspace{\b@r@stretch\box@height}%
      \fi%
      \put@feature%
      \temp@@count=\b@feature@count%
      \advance\temp@@count by 1%
      \xdef\b@feature@count{\the\temp@@count}%
   \else
      \iffix@
      \if\bottom@stretch y%
         \vspace{-\box@height}%
         \vspace{\b@r@stretch\box@height}%
      \fi%
      \ifbottomfeature%
         \vspace{\b@sp@ce}\newline\hbox{}\newline\hbox{}%
      \fi%
      \fi%
   \fi%
   \ifnum\feature@bbottom=1 %
      \advance\bar@length by \b@feature@count \baselineskip%
      \multiply\bar@length by \b@r@stretch%
      \multiply\bar@length by 2%
      \advance\bar@length by \bb@sp@ce\message{bb\the\bar@length bb}%
      \ifnum\featureonbbottom=0 \xdef\feature@bbottom{0}\fi%
      \xdef\bottop@{bbottom}%
      \vspace{\bb@sp@ce}%
      \if\bbottom@stretch y%
         \advance\bar@length by -\box@height%
         \advance\bar@length by \b@r@stretch\box@height%
         \vspace{-\box@height}%
         \vspace{\b@r@stretch\box@height}%
      \fi%
      \setlength{\bar@length}{0pt}%
      \put@feature%
      \temp@@count=\b@feature@count%
      \advance\temp@@count by 1%
      \xdef\b@feature@count{\the\temp@@count}%
   \else
      \iffix@
      \if\bbottom@stretch y%
         \vspace{-\box@height}%
         \vspace{\b@r@stretch\box@height}%
      \fi%
      \ifbbottomfeature
         \vspace{\bb@sp@ce}\newline\hbox{}\newline\hbox{}%
      \fi
      \fi
   \fi
   \ifnum\feature@bbbottom=1 %
      \advance\bar@length by \b@feature@count\baselineskip%
      \multiply\bar@length by \b@r@stretch%
      \multiply\bar@length by 2%
      \advance\bar@length by \bbb@sp@ce\message{bbb\the\bar@length bbb}%
      \ifnum\featureonbbbottom=0 \xdef\feature@bbbottom{0}\fi%
      \xdef\bottop@{bbbottom}%
      \vspace{\bbb@sp@ce}%
      \if\bbbottom@stretch y%
         \vspace{-\box@height}%
         \vspace{\b@r@stretch\box@height}%
      \fi%
      \put@feature%
      \temp@@count=\b@feature@count%
      \advance\temp@@count by 1%
      \xdef\b@feature@count{\the\temp@@count}%
   \else
      \iffix@
      \if\bbbottom@stretch y%
         \vspace{-\box@height}%
         \vspace{\b@r@stretch\box@height}%
      \fi%
      \ifbbbottomfeature
         \vspace{\bbb@sp@ce}\newline\hbox{}\newline\hbox{}%
      \fi
      \fi
   \fi
   \ifnum\feature@bbbbottom=1 %
      \advance\bar@length by \b@feature@count\baselineskip%
      \multiply\bar@length by \b@r@stretch%
      \multiply\bar@length by 2%
      \advance\bar@length by \bbbb@sp@ce\message{bbbb\the\bar@length bbbb}%
      \ifnum\featureonbbbbottom=0 \xdef\feature@bbbbottom{0}\fi%
      \xdef\bottop@{bbbbottom}%
      \vspace{\bbbb@sp@ce}%
      \if\bbbbottom@stretch y%
         \vspace{-\box@height}%
         \vspace{\b@r@stretch\box@height}%
      \fi%
      \put@feature%
   \else
      \iffix@
      \if\bbbbottom@stretch y%
         \vspace{-\box@height}%
         \vspace{\b@r@stretch\box@height}%
      \fi%
      \ifbbbbottomfeature
         \vspace{\bbbb@sp@ce}\newline\hbox{}\newline\hbox{}%
      \fi
      \fi
   \fi
   \xdef\consensus{} \xdef\ruler@{}
   \xdef\styleframe{}         
   \xdef\textfeaturetop{}     \xdef\textfeaturebottom{}
   \xdef\textfeaturettop{}    \xdef\textfeaturebbottom{}
   \xdef\textfeaturetttop{}   \xdef\textfeaturebbbottom{}
   \xdef\textfeaturettttop{}  \xdef\textfeaturebbbbottom{}
   \xdef\stylefeaturetop{}    \xdef\stylefeaturebottom{}
   \xdef\stylefeaturettop{}   \xdef\stylefeaturebbottom{}
   \xdef\stylefeaturetttop{}  \xdef\stylefeaturebbbottom{}
   \xdef\stylefeaturettttop{} \xdef\stylefeaturebbbbottom{}
   \loopcount=0
   \loop
         \advance\loopcount by 1
         \xdef\seq@line{\csname sequence\the\loopcount\endcsname}
         \expandafter\remove@fromseq\seq@line
   \ifnum\loopcount<\seq@count \repeat}

\def\block@output{%
   \expandafter\ifnum\csname res@count\start@seq\endcsname<\end@num\relax
     \message{.}
     \ifx\out@put\y@
        \vbox{\set@lines}\par
        \block@skip
     \else
        \pos@count=1
        \findc@nsensus
        \loopcount=0
        \loop
           \advance\loopcount by 1\relax
           \xdef\seq@line{\csname sequence\the\loopcount\endcsname}
           \expandafter\remove@fromseq\seq@line
           \ifT@coffee
              \xdef\TC@cons{\csname TC0\endcsname}
              \expandafter\remove@fromTC\TC@cons
           \fi
        \ifnum\loopcount<\seq@count \repeat
     \fi
     \ifstop@
     \else
        \advance\seq@pointer by -\res@perline
        \ifnum\seq@pointer>\res@perline \block@output \fi
     \fi
   \fi}

%%%%% Basic input routines

\def\savedseqlength#1#2#3{%
   \expandafter\xdef\csname savelength#1seq#2\endcsname{#3}}

\def\set@savedseqlength{%
   \loopcount=0
   \loop
      \advance\loopcount by 1
      \expandafter\ifx%
         \csname savelength\@lign@count seq\the\loopcount\endcsname\relax 
      \else
         \expandafter\xdef\csname seq@len\the\loopcount\endcsname{%
            \csname savelength\@lign@count seq\the\loopcount\endcsname}
      \fi
   \ifnum\loopcount<\seq@count\repeat
}
   
\def\save@lengths{%
   \loopcount=0
   \loop
       \advance\loopcount by 1\relax
       \immediate\write\@auxout{%
       \string\savedseqlength{\@lign@count}{\the\loopcount}% 
           {\csname res@count\the\loopcount\endcsname}}
   \ifnum\loopcount<\seq@count \repeat
}
   
\def\do@cleanup{%
   \expandafter\firstchar@get\third@
   \expandafter\check@char\first@
   \ifletter
      \xdef\second@{\second@\first@}
   \else
      \ifnumber
         \ifx\first@\gre@ter
	    \xdef\second@{\second@{$>$}}
	 \else
         \ifx\first@\sm@ller
	    \xdef\second@{\second@{$<$}}
	 \else
         \xdef\second@{\second@\first@}
	 \fi\fi
      \else
         \ifnum\code@num=6
            \xdef\second@{\second@\#}	 
	 \else
         \ifnum\code@num=14
            \xdef\second@{\second@\%}	 
	 \else
            \xdef\second@{\second@\noexpand\string\first@}
	 \fi\fi
      \fi
   \fi
   \ifx\third@\@t \else \do@cleanup \fi
}
    
\def\cleanup@name{%
   \xdef\third@{\csname seqname\the\loopcount\endcsname @}
   \xdef\second@{}
   \do@cleanup
   \expandafter\xdef\csname newseqname\the\loopcount\endcsname{\second@}
}
    
\def\clear@seq{%
    \loopcount=0
    \loop
       \advance\loopcount by 1
       \expandafter\def\csname sequence\the\loopcount\endcsname{}
    \ifnum\loopcount<\seq@count \repeat
    \xdef\consensus{}
    \xdef\constopo{}
    \xdef\frame@on{0}
    \xdef\textfeaturetop{}       \xdef\featureontop{0}
    \xdef\textfeaturettop{}      \xdef\featureonttop{0}
    \xdef\textfeaturetttop{}     \xdef\featureontttop{0}
    \xdef\textfeaturettttop{}    \xdef\featureonttttop{0}
    \xdef\textfeaturebottom{}    \xdef\featureonbottom{0}
    \xdef\textfeaturebbottom{}   \xdef\featureonbbottom{0}
    \xdef\textfeaturebbbottom{}  \xdef\featureonbbbottom{0}
    \xdef\textfeaturebbbbottom{} \xdef\featureonbbbbottom{0}
    \xdef\styleframe{} 
    \xdef\stylefeaturetop{}
    \xdef\stylefeaturettop{}
    \xdef\stylefeaturetttop{}
    \xdef\stylefeaturettttop{}
    \xdef\stylefeaturebottom{}
    \xdef\stylefeaturebbottom{}
    \xdef\stylefeaturebbbottom{}
    \xdef\stylefeaturebbbbottom{}}
\def\guess@protein{\seqtype{P}\message{<Seqtype guess: protein>}}
\def\guess@DNA{\seqtype{N}\message{<Seqtype guess: nucleotide>}}
\def\allmatch@out{%
    \xdef\leave@in{n}
    \loopcount=1
    \xdef\seq@line{\csname seq\the\loopcount\endcsname}
    \expandafter\res@get\seq@line
    \expandafter\xdef\csname seq\the\loopcount\endcsname{\seq@line}
    \ifnum\loopcount=\start@seq\relax
       \expandafter\check@char\first@
       \ifletter
          \advance\temp@@count by 1
          \ifnum\temp@@count=\end@num\relax 
             \advance\temp@count by 1
             \xdef\second@{\the\temp@count}
             \advance\temp@count by -1
          \fi
       \fi
    \fi
    \ifx\first@\ampers@nd
       \ifx\domain@active\y@
          \xdef\stack@dom{\stack@dom\the\temp@count;}
       \fi   
       \xdef\stack@dom{\stack@dom &;&;@}
    \else
       \xdef\first@@{\first@}
       \advance\temp@count by 1
       \loop
          \advance\loopcount by 1
          \xdef\seq@line{\csname seq\the\loopcount\endcsname}
          \expandafter\res@get\seq@line
          \expandafter\xdef\csname seq\the\loopcount\endcsname{\seq@line}
          \ifx\first@@\first@\else\xdef\leave@in{y}\fi
          \ifnum\loopcount=\start@seq\relax
             \expandafter\check@char\first@
             \ifletter
                \advance\temp@@count by 1
                \ifnum\temp@@count=\end@num\relax \xdef\second@{\the\temp@count}\fi
             \fi
          \fi
       \ifnum\loopcount<\seq@count\repeat
       \ifx\leave@in\y@
          \ifx\domain@active\n@
             \ifnum\start@seq=0
                \xdef\domain@active{y}
                \xdef\stack@dom{\stack@dom\the\temp@count;}
             \else
                \ifnum\temp@@count<\start@num\relax
                \else
                   \ifnum\temp@@count>\end@num\relax
                   \else
                      \xdef\domain@active{y}
                      \xdef\stack@dom{\stack@dom\the\temp@count;}
                   \fi
                \fi
             \fi
          \else
             \ifnum\start@seq>0
                \ifnum\temp@@count<\end@num\relax
                \else
                   \xdef\domain@active{n}
                   \xdef\stack@dom{\stack@dom\second@;}
                \fi
             \fi
          \fi
       \else
          \ifx\domain@active\y@
             \xdef\domain@active{n}
             \advance\temp@count by -1
             \xdef\stack@dom{\stack@dom\the\temp@count;}
             \advance\temp@count by 1
          \fi
       \fi
       \allmatch@out
    \fi
}
\def\domain@count{%
    \xdef\@ddg@p{n}
    \loopcount=0
    \ifT@coffee
       \xdef\seq@line{\csname T@cof\the\loopcount\endcsname}
       \expandafter\res@get\seq@line
       \expandafter\xdef\csname TC@num\the\loopcount\endcsname{\first@}
       \expandafter\xdef\csname T@cof\the\loopcount\endcsname{\seq@line}
    \fi
    \loop
       \advance\loopcount by 1
       \xdef\seq@line{\csname seq\the\loopcount\endcsname}
       \expandafter\res@get\seq@line
       \expandafter\xdef\csname res\the\loopcount\endcsname{\first@}
       \expandafter\xdef\csname seq\the\loopcount\endcsname{\seq@line}
       \ifT@coffee
          \xdef\seq@line{\csname T@cof\the\loopcount\endcsname}
          \expandafter\res@get\seq@line
          \expandafter\xdef\csname TC@num\the\loopcount\endcsname{\first@}
          \expandafter\xdef\csname T@cof\the\loopcount\endcsname{\seq@line}
       \fi
    \ifnum\loopcount<\seq@count\repeat
    \ifx\first@\ampers@nd
    \else
     \ifnum\next@domain@num>-99999\relax
         \ifnum\domain@seq=0
             \temp@count=\csname dom@count0\endcsname 
             \advance\temp@count by 1
             \expandafter\xdef\csname dom@count0\endcsname{\the\temp@count}
         \fi
         \loopcount=0
         \loop
           \advance\loopcount by 1
            \xdef\first@{\csname res\the\loopcount\endcsname}
            \expandafter\check@char\first@
            \ifletter
               \temp@count=\csname dom@count\the\loopcount\endcsname 
               \advance\temp@count by 1
               \expandafter\xdef\csname dom@count\the\loopcount\endcsname{\the\temp@count}
               \if\seq@type A \relax
                  \temp@count=\char@num
                  \ifnum\temp@count>96 \advance\temp@count by -32 \fi
                  \ifnum\temp@count=69 \guess@protein
                  \else
                  \ifnum\temp@count=70 \guess@protein
                  \else
                  \ifnum\temp@count=73 \guess@protein
                  \else
                  \ifnum\temp@count=76 \guess@protein
                  \else
                  \ifnum\temp@count=80 \guess@protein
                  \else
                  \ifnum\temp@count=81 \guess@protein
                  \fi\fi\fi\fi\fi\fi
               \fi
            \else
               \ifx\d@m@in\y@
                  \ifnum\loopcount=\domain@seq \xdef\@ddg@p{y} \fi
               \fi
            \fi
         \ifnum\loopcount<\seq@count\repeat 
         \ifx\@ddg@p\y@ 
           \ifx\g@p\n@
            \loopcount=0
            \loop
            \advance\loopcount by 1
               \ifT@coffee
                  \expandafter\xdef\csname T@coffee\the\loopcount\endcsname{%
                                   \csname T@coffee\the\loopcount\endcsname \csname TC@num\the\loopcount\endcsname}
               \fi
               \expandafter\xdef\csname sequence\the\loopcount\endcsname{%
                                \csname sequence\the\loopcount\endcsname \csname res\the\loopcount\endcsname}
               \expandafter\xdef\csname dom@num\the\loopcount\endcsname{%
                                \csname dom@num\the\loopcount\endcsname \csname dom@count\the\loopcount\endcsname,}
            \ifnum\loopcount<\seq@count\repeat
            \advance\total@count by 1
            \advance\outerloopcount by 1 \ifnum\outerloopcount=\res@perline \outerloopcount=0 \fi
           \fi
         \fi 
         \expandafter\ifnum\csname dom@count\domain@seq\endcsname=\next@domain@num
            \xdef\d@m@in{y}
            \advance\total@count by 1
            \advance\outerloopcount by 1 \ifnum\outerloopcount=\res@perline \outerloopcount=0 \fi
            \loopcount=0
            \ifx\g@p\y@
               \ifT@coffee
                  \expandafter\xdef\csname T@coffee\the\loopcount\endcsname{%
                                   \csname T@coffee\the\loopcount\endcsname \csname TC@num\the\loopcount\endcsname}
               \fi
            \fi
            \ifT@coffee
               \expandafter\xdef\csname T@coffee\the\loopcount\endcsname{%
                                \csname T@coffee\the\loopcount\endcsname \csname TC@num\the\loopcount\endcsname}
            \fi
            \loop
               \advance\loopcount by 1
               \ifx\g@p\y@
                  \ifT@coffee
                     \expandafter\xdef\csname T@coffee\the\loopcount\endcsname{%
                                      \csname T@coffee\the\loopcount\endcsname \csname TC@num\the\loopcount\endcsname}
                  \fi
                  \expandafter\xdef\csname sequence\the\loopcount\endcsname{%
                                   \csname sequence\the\loopcount\endcsname\equ@l}
                  \expandafter\xdef\csname dom@num\the\loopcount\endcsname{%
                                   \csname dom@num\the\loopcount\endcsname \csname dom@count\the\loopcount\endcsname,}
                  \ifnum\outerloopcount=0
                     \expandafter\xdef\csname dom@num@break\the\loopcount\endcsname{%
                                      \csname dom@num@break\the\loopcount\endcsname \csname dom@count\the\loopcount\endcsname,}
                  \fi
               \fi
               \ifT@coffee
                  \expandafter\xdef\csname T@coffee\the\loopcount\endcsname{%
                                   \csname T@coffee\the\loopcount\endcsname \csname TC@num\the\loopcount\endcsname}
               \fi
               \expandafter\xdef\csname sequence\the\loopcount\endcsname{%
                                \csname sequence\the\loopcount\endcsname \csname res\the\loopcount\endcsname}
               \expandafter\xdef\csname dom@num\the\loopcount\endcsname{%
                                \csname dom@num\the\loopcount\endcsname \csname dom@count\the\loopcount\endcsname,}
               \ifnum\outerloopcount=0
                  \expandafter\xdef\csname dom@num@break\the\loopcount\endcsname{%
                                   \csname dom@num@break\the\loopcount\endcsname \csname dom@count\the\loopcount\endcsname,}
               \fi
            \ifnum\loopcount<\seq@count\repeat
            \ifx\@ddg@p\n@
               \expandafter\get@item\domain@num@list
               \xdef\domain@num@list{\first@}
            \fi
            \ifx\fourth@\ampers@nd \xdef\next@domain@num{-99999}
            \else
               \loopcount=\next@domain@num
               \advance\loopcount by 1
               \ifnum\loopcount=\fourth@
                  \xdef\g@p{n}
               \else
                  \xdef\g@p{y}
                     \advance\total@count by 1
                     \advance\outerloopcount by 1
                     \ifnum\outerloopcount=\res@perline
                        \outerloopcount=0
                        \loopcount=0
                        \loop
                           \advance\loopcount by 1
                           \expandafter\xdef\csname dom@num@break\the\loopcount\endcsname{%
                                            \csname dom@num@break\the\loopcount\endcsname \csname dom@count\the\loopcount\endcsname,}
                        \ifnum\loopcount<\seq@count\repeat
                     \fi
               \fi
               \xdef\next@domain@num{\fourth@} \xdef\end@num{\fourth@}
            \fi
         \fi 
         \domain@count
     \fi
    \fi
}
\def\residue@count{%
    \expandafter\res@get\seq@line
    \if\first@\ampers@nd
    \else
         \advance\total@count by 1
         \advance\innerloopcount by 1
         \ifstart@ \advance\res@count by 1 \fi
         \ifnum\start@seq=0
            \advance\end@count by 1
%%%%%%%%%%            
            \ifnum\end@count=\start@num\relax
               \xdef\start@number{\the\innerloopcount}
               \start@true \res@count=1
            \fi
%%%%%%%%%%
         \else
            \expandafter\check@char\first@
            \ifletter 
               \advance\end@count by 1
               \if\seq@type A \relax
                  \temp@count=\char@num
                  \ifnum\temp@count>96 \advance\temp@count by -32 \fi
                  \ifnum\temp@count=69 \guess@protein
                  \else
                  \ifnum\temp@count=70 \guess@protein
                  \else
                  \ifnum\temp@count=73 \guess@protein
                  \else
                  \ifnum\temp@count=76 \guess@protein
                  \else
                  \ifnum\temp@count=80 \guess@protein
                  \else
                  \ifnum\temp@count=81 \guess@protein
                  \fi\fi\fi\fi\fi\fi
               \fi
            \fi
            \ifnum\end@count=\start@num\relax
               \xdef\start@number{\the\innerloopcount}
               \start@true \res@count=1
            \fi
      \fi
    \residue@count
    \fi}
\def\clean@seq{%
    \expandafter\res@get\seq@line
    \if\first@\ampers@nd \xdef\seq@line{\seq@@line}
    \else
       \expandafter\check@char\first@
       \ifnum\char@num>64
          \ifnum\char@num>96 \make@upper \fi
          \xdef\seq@@line{\seq@@line\first@}
       \else
       \ifnum\char@num=46 
          \xdef\seq@@line{\seq@@line\d@t}
       \else
       \ifnum\char@num=45 
          \xdef\seq@@line{\seq@@line\d@t}
       \else
       \ifnum\char@num=42 
          \xdef\seq@@line{\seq@@line\questi@n}
       \fi
       \fi
       \fi
       \fi
       \clean@seq
    \fi}
\def\read@loop{%
    \read\alignfile to \inline
    \xdef\last@{\expandafter\string\inline}
    \ifx\last@\par@
      \else
         \xdef\inline{\inline @}
         \expandafter\seq@get\inline
         \ifstop@
         \else
         \innerloopcount=\csname @rd\the\loopcount\endcsname\relax
         \expandafter\ifx\csname seq@name\the\loopcount\endcsname\first@
             \xdef\seq@@line{} \xdef\seq@line{\seq@line &@} \clean@seq
             \expandafter\xdef\csname sequence\the\innerloopcount\endcsname{%
                  \csname sequence\the\innerloopcount\endcsname\seq@line}
             \advance\loopcount by 1
             \ifnum\loopcount>\seq@count
                \loopcount=1
             \fi
         \fi
         \fi
       \fi
    \ifeof\alignfile
     \ifx\all@out\y@
       \xdef\domain@seq{0}
       \loopcount=0
       \loop
          \advance\loopcount by 1
          \expandafter\xdef\csname seq\the\loopcount\endcsname{\csname sequence\the\loopcount\endcsname &@}
       \ifnum\loopcount<\seq@count\repeat
       \xdef\domain@active{n}
       \xdef\stack@dom{}
       \temp@count=0
       \temp@@count=0
       \allmatch@out
       \xdef\dom@in{y}
     \fi
     \ifx\dom@in\y@
       \loopcount=0
       \ifT@coffee
          \xdef\temp@{\csname T@coffee\the\loopcount\endcsname}
          \expandafter\get@temp\temp@
          \expandafter\xdef\csname T@cof\the\loopcount\endcsname{\temp@ &@}
          \expandafter\xdef\csname T@coffee\the\loopcount\endcsname{}
       \fi
       \expandafter\xdef\csname dom@count\the\loopcount\endcsname{0}
       \loop
          \advance\loopcount by 1
          \ifT@coffee
             \xdef\temp@{\csname T@coffee\the\loopcount\endcsname}
             \expandafter\get@temp\temp@
             \expandafter\xdef\csname T@cof\the\loopcount\endcsname{\temp@ &@}
             \expandafter\xdef\csname T@coffee\the\loopcount\endcsname{}
          \fi
          \expandafter\xdef\csname seq\the\loopcount\endcsname{\csname sequence\the\loopcount\endcsname &@}
          \expandafter\xdef\csname sequence\the\loopcount\endcsname{}
          \expandafter\xdef\csname dom@num\the\loopcount\endcsname{}
          \expandafter\xdef\csname dom@num@break\the\loopcount\endcsname{}
          \expandafter\xdef\csname dom@count\the\loopcount\endcsname{\csname res@count\the\loopcount\endcsname}
       \ifnum\loopcount<\seq@count\repeat
       \xdef\domain@num@list{}
       \xdef\last@{\stack@dom}
       \get@domain
       \xdef\domain@num@list{\domain@num@list &,@}
       \expandafter\get@item\domain@num@list
       \xdef\next@domain@num{\fourth@} \xdef\domain@num@list{\first@}
       \xdef\d@m@in{n}
       \domain@count
       \loopcount=0
       \ifT@coffee
          \expandafter\xdef\csname T@coffee\the\loopcount\endcsname{\csname T@coffee\the\loopcount\endcsname @}
       \fi
       \loop
          \advance\loopcount by 1
          \ifT@coffee
             \expandafter\xdef\csname T@coffee\the\loopcount\endcsname{\csname T@coffee\the\loopcount\endcsname @}
          \fi
          \expandafter\xdef\csname dom@num\the\loopcount\endcsname{\csname dom@num\the\loopcount\endcsname @}
          \expandafter\xdef\csname dom@num@break\the\loopcount\endcsname{%
                           \csname dom@num@break\the\loopcount\endcsname \csname dom@count\the\loopcount\endcsname,@}
       \ifnum\loopcount<\seq@count\repeat
       \if\seq@type A \guess@DNA \fi
       \advance\seq@pointer by \total@count\relax
       \xdef\total@pos{\the\total@count}
    
     \else   

       \ifnum\start@seq=0
          \xdef\seq@line{\csname sequence1\endcsname &@} 
       \else
          \xdef\seq@line{\csname sequence\start@seq\endcsname &@} 
       \fi
       \res@count=0 
       \innerloopcount=0
       \residue@count 
       \if\seq@type A \guess@DNA \fi
       \advance\seq@pointer by \res@count\relax 
       \xdef\total@pos{\the\total@count}
       \ifx\end@num\n@\xdef\end@num{\the\end@count}\fi
       \total@count=\end@count\relax
       \xdef\first@{\csname res@count\start@seq\endcsname}
       \advance\total@count by \first@\relax
       \xdef\first@{\the\total@count}
       \ifnum\end@num>\total@count\xdef\end@num{\the\total@count}\fi
     \fi       
     \ifshow@sublogo \prep@sublogo \fi
     \ifshow@logo  \prep@logo  \fi
     \ifx\label@motif\y@  \prep@motif \fi
     \ifnum\seq@pointer>\res@perline \block@output \fi 
    \else \read@loop
    \fi}
\def\read@fasta{%
    \read\alignfile to \inline
    \xdef\last@{\expandafter\string\inline}
    \ifx\last@\par@
      \else
         \xdef\seq@line{\inline}
         \xdef\inline{\inline @}
	 \expandafter\firstchar@get\inline
         \ifstop@
         \else
	 \ifx\first@\gre@ter
	    \advance\loopcount by 1\relax
	    \expandafter\seq@get\third@
	    \xdef\seq@name{\first@}
	    \xdef\second@{\first@ &}
	    \ifx\second@\ampers@nd \xdef\seq@name{seq\the\loopcount}\fi
	    \innerloopcount=\csname @rd\the\loopcount\endcsname\relax
	 \else
	    \expandafter\ifx\csname seq@name\the\loopcount\endcsname\seq@name
              \xdef\seq@@line{} \xdef\seq@line{\seq@line &@} \clean@seq
              \expandafter\xdef\csname sequence\the\innerloopcount\endcsname{%
                  \csname sequence\the\innerloopcount\endcsname\seq@line}
            \fi
	 \fi
         \fi
       \fi
    \ifeof\alignfile
     \ifx\all@out\y@
       \xdef\domain@seq{0}
       \loopcount=0
       \loop
          \advance\loopcount by 1
          \expandafter\xdef\csname seq\the\loopcount\endcsname{\csname sequence\the\loopcount\endcsname &@}
       \ifnum\loopcount<\seq@count\repeat
       \xdef\domain@active{n}
       \xdef\stack@dom{}
       \temp@count=0
       \temp@@count=0
       \allmatch@out
       \xdef\dom@in{y}
     \fi
     \ifx\dom@in\y@
       \loopcount=0
       \ifT@coffee
          \xdef\temp@{\csname T@coffee\the\loopcount\endcsname}
          \expandafter\get@temp\temp@
          \expandafter\xdef\csname T@cof\the\loopcount\endcsname{\temp@ &@}
          \expandafter\xdef\csname T@coffee\the\loopcount\endcsname{}
       \fi
       \expandafter\xdef\csname dom@count\the\loopcount\endcsname{0}
       \loop
          \advance\loopcount by 1
          \ifT@coffee
             \xdef\temp@{\csname T@coffee\the\loopcount\endcsname}
             \expandafter\get@temp\temp@
             \expandafter\xdef\csname T@cof\the\loopcount\endcsname{\temp@ &@} 
             \expandafter\xdef\csname T@coffee\the\loopcount\endcsname{}
          \fi
          \expandafter\xdef\csname seq\the\loopcount\endcsname{\csname sequence\the\loopcount\endcsname &@}
          \expandafter\xdef\csname sequence\the\loopcount\endcsname{}
          \expandafter\xdef\csname dom@num\the\loopcount\endcsname{}
          \expandafter\xdef\csname dom@num@break\the\loopcount\endcsname{}
          \expandafter\xdef\csname dom@count\the\loopcount\endcsname{\csname res@count\the\loopcount\endcsname}
       \ifnum\loopcount<\seq@count\repeat
       \xdef\domain@num@list{}
       \xdef\last@{\stack@dom}
       \get@domain
       \xdef\domain@num@list{\domain@num@list &,@}
       \expandafter\get@item\domain@num@list
       \xdef\next@domain@num{\fourth@} \xdef\domain@num@list{\first@}
       \xdef\d@m@in{n}
       \domain@count
       \loopcount=0
       \ifT@coffee
          \expandafter\xdef\csname T@coffee\the\loopcount\endcsname{\csname T@coffee\the\loopcount\endcsname @}
       \fi
       \loop
          \advance\loopcount by 1
          \ifT@coffee
             \expandafter\xdef\csname T@coffee\the\loopcount\endcsname{\csname T@coffee\the\loopcount\endcsname @}
          \fi
          \expandafter\xdef\csname dom@num\the\loopcount\endcsname{\csname dom@num\the\loopcount\endcsname @}
          \expandafter\xdef\csname dom@num@break\the\loopcount\endcsname{%
                           \csname dom@num@break\the\loopcount\endcsname \csname dom@count\the\loopcount\endcsname,@}
       \ifnum\loopcount<\seq@count\repeat
       \if\seq@type A \guess@DNA \fi
       \advance\seq@pointer by \total@count\relax 
       \xdef\total@pos{\the\total@count}
    
     \else   
    
       \ifnum\start@seq=0
          \xdef\seq@line{\csname sequence1\endcsname &@} 
       \else
          \xdef\seq@line{\csname sequence\start@seq\endcsname &@} 
       \fi
       \res@count=0 
       \innerloopcount=0
       \residue@count
       \if\seq@type A \guess@DNA \fi
       \advance\seq@pointer by \res@count
       \xdef\total@pos{\the\total@count}
       \ifx\end@num\n@\xdef\end@num{\the\end@count}\fi
       \total@count=\end@count\relax
       \xdef\first@{\csname res@count\start@seq\endcsname}
       \advance\total@count by \first@\relax
       \xdef\first@{\the\total@count}
       \ifnum\end@num>\total@count\xdef\end@num{\the\total@count}\fi
     \fi
     \ifshow@sublogo \prep@sublogo \fi
     \ifshow@logo  \prep@logo  \fi
     \ifx\label@motif\y@  \prep@motif \fi
     \ifnum\seq@pointer>\res@perline \block@output \fi
    \else 
       \read@fasta 
    \fi}
\def\read@lines{%
    \openin\alignfile = \alignfilename
    \clear@seq
    \ifnum\start@seq>0
       \xdef\start@seq{\csname @rd\start@seq\endcsname} \fi
    \loopcount=\start@num
    \advance\loopcount by -\csname res@count\start@seq\endcsname\relax
    \expandafter\ifnum\csname res@count\start@seq\endcsname<\start@num\relax
       \xdef\start@num{\the\loopcount}
    \else \start@true \fi
    \res@count=0 \seq@pointer=0 \end@count=0 \total@count=0
    \xdef\start@number{0}
    \ifx\exp@rt\y@ \prep@reexp@rtfile \fi
    \ifx\f@st@\y@
       \loopcount=0
       \read@fasta
    \else
       \loopcount=1
       \read@loop
    \fi
    \ifnum\seq@pointer>0
       \ifstop@
       \else
          \res@count=\res@perline
          \res@perline=\seq@pointer
          \block@output
       \fi
    \fi
    \ifshow@sublogo \closein\sublogofile \fi
    \ifx\exp@rt\y@ \closeout\exp@rtfile \fi
    \ifx\exp@rt@chimera\y@ \closeout\exp@rt@chimerafile \fi%
    \closein\alignfile
    \message{)}
}


%%%%%  Read alignment, decide whether MSF or ALN, interpret
    
\def\readalignfile#1{%
 \def\alignfilename{#1}
 \xdef\first@{byhand}
 \ifx\alignfilename\first@
 \else
 \openin\alignfile = #1
 \ifeof\alignfile
    \PackageError{TeXshade}
    {File `#1' not found}
    {\MessageBreak
     The alignment file you specified is missing or you have \MessageBreak
     misspelled it. \MessageBreak\MessageBreak
     Stop here, otherwise you're likely getting in trouble. \MessageBreak
     Type X <return> to quit. \MessageBreak
}
 \else
    \message{(\alignfilename :}
    \xdef\seq@type{A} \xdef\he@der{no} \xdef\f@st@{no} 
    \xdef\first@line{y}
    \seq@count=0 \loopcount=0 \innerloopcount=0 \temp@count=0
    \loop
       \read\alignfile to \inline 
       \xdef\test@{\expandafter\string\inline}
       \ifx\test@\par@ \innerloopcount=0
       \else
          \xdef\msfline{\inline & & & & @}
          \expandafter\inf@@get\msfline
          \ifx\first@\@msf \expandafter\type@get\msfline \fi
          \ifx\second@\@msf   \expandafter\type@get\msfline \fi
          \ifx\third@\@msf \expandafter\type@get\msfline \fi
          \ifx\first@\n@me \advance\loopcount by 1\relax
             \expandafter\xdef\csname seqname\the\loopcount\endcsname{\second@}
             \expandafter\xdef\csname seq@len\the\loopcount\endcsname{\fourth@}
          \fi
          \ifx\first@\he@derend
             \ifnum\loopcount>0 \xdef\he@der{yes} \fi
       \fi
       \xdef\alnline{\inline @} 
       \expandafter\check@letter\alnline
       \ifletter
          \expandafter\seq@get\alnline
          \advance\innerloopcount by 1\relax
          \seq@count=\innerloopcount
          \expandafter\xdef\csname newseqname\the\seq@count\endcsname{\first@}
	   \else
	     \expandafter\firstchar@get\alnline
	     \ifx\first@\gre@ter
	        \ifx\first@line\y@ \xdef\f@st@{y}\fi
                \expandafter\seq@get\third@
                \advance\temp@count by 1\relax
	        \xdef\second@{\first@ &}
	        \ifx\second@\ampers@nd \xdef\first@{seq\the\temp@count}\fi
                \expandafter\xdef\csname seq@name\the\temp@count\endcsname{\first@}
             \fi
          \fi
       \fi
      \xdef\first@line{n}
    \ifeof\alignfile \else\repeat
    \closein\alignfile
    \xdef\first@{no}
    \ifx\he@der\first@
      \loopcount=0
      \ifx\f@st@\y@
        \seq@count=\temp@count \loopcount=0
        \loop
          \advance\loopcount by 1
          \expandafter\xdef\csname seqname\the\loopcount\endcsname{%
                              \csname seq@name\the\loopcount\endcsname}
          \expandafter\xdef\csname newseqname\the\loopcount\endcsname{%
                              \csname seq@name\the\loopcount\endcsname}
          \expandafter\xdef\csname seq@len\the\loopcount\endcsname{99999999}
        \ifnum\loopcount<\seq@count \repeat
      \else
        \loop
          \advance\loopcount by 1
          \expandafter\xdef\csname seqname\the\loopcount\endcsname{%
                              \csname newseqname\the\loopcount\endcsname}
          \expandafter\xdef\csname seq@len\the\loopcount\endcsname{99999999}
        \ifnum\loopcount<\seq@count \repeat
      \fi
    \else
       \seq@count=\loopcount \loopcount=0 \box@width=0pt
       \loop
          \advance\loopcount by 1
          \expandafter\xdef\csname newseqname\the\loopcount\endcsname{%
                              \csname seqname\the\loopcount\endcsname}
       \ifnum\loopcount<\seq@count \repeat
    \fi
    \set@savedseqlength
    \loopcount=0 \xdef\seq@order{}
    \expandafter\xdef\csname stack@dom\endcsname{&;&;&;&;@}
    \expandafter\xdef\csname stack@reg\the\loopcount\endcsname{&;&;&;&;@}
    \expandafter\xdef\csname stack@tintreg\the\loopcount\endcsname{&;&;&;@}
    \expandafter\xdef\csname stack@emphreg\the\loopcount\endcsname{&;&;&;@}
    \expandafter\xdef\csname stack@lowerreg\the\loopcount\endcsname{&;&;&;@}
    \expandafter\xdef\csname stack@framereg\the\loopcount\endcsname{&;&;&;@}
    \expandafter\xdef\csname stack@shadingreg\the\loopcount\endcsname{&;&;&;@}
    \expandafter\xdef\csname stack@top\the\loopcount\endcsname{&;&;&;&;@}
    \expandafter\xdef\csname stack@ttop\the\loopcount\endcsname{&;&;&;&;@}
    \expandafter\xdef\csname stack@tttop\the\loopcount\endcsname{&;&;&;&;@}
    \expandafter\xdef\csname stack@ttttop\the\loopcount\endcsname{&;&;&;&;@}
    \expandafter\xdef\csname stack@bottom\the\loopcount\endcsname{&;&;&;&;@}
    \expandafter\xdef\csname stack@bbottom\the\loopcount\endcsname{&;&;&;&;@}
    \expandafter\xdef\csname stack@bbbottom\the\loopcount\endcsname{&;&;&;&;@}
    \expandafter\xdef\csname stack@bbbbottom\the\loopcount\endcsname{&;&;&;&;@}
    \expandafter\xdef\csname seq@start\the\loopcount\endcsname{1}
    \expandafter\xdef\csname name@col\the\loopcount\endcsname{yes}
    \expandafter\xdef\csname number@col\the\loopcount\endcsname{yes}
    \loop
       \advance\loopcount by 1
       \expandafter\xdef\csname @rd\the\loopcount\endcsname{\the\loopcount}
       \expandafter\xdef\csname res@count\the\loopcount\endcsname{0}
       \cleanup@name
       \expandafter\xdef\csname tint@seq\the\loopcount\endcsname{n}
       \expandafter\xdef\csname emph@seq\the\loopcount\endcsname{n}
       \expandafter\xdef\csname lower@seq\the\loopcount\endcsname{n}
       \expandafter\xdef\csname hide@seq\the\loopcount\endcsname{false}
       \expandafter\xdef\csname hide@@@seq\the\loopcount\endcsname{false}
       \expandafter\xdef\csname hide@name\the\loopcount\endcsname{no}
       \expandafter\xdef\csname name@col\the\loopcount\endcsname{yes}
       \expandafter\xdef\csname hide@number\the\loopcount\endcsname{no}
       \expandafter\xdef\csname number@col\the\loopcount\endcsname{yes}
       \expandafter\xdef\csname stack@reg\the\loopcount\endcsname{&;&;&;&;@}
       \expandafter\xdef\csname stack@tintreg\the\loopcount\endcsname{&;&;&;@}
       \expandafter\xdef\csname stack@emphreg\the\loopcount\endcsname{&;&;&;@}
       \expandafter\xdef\csname stack@lowerreg\the\loopcount\endcsname{&;&;&;@}
       \expandafter\xdef\csname stack@framereg\the\loopcount\endcsname{&;&;&;@}
       \expandafter\xdef\csname stack@shadingreg\the\loopcount\endcsname{&;&;&;@}
       \expandafter\xdef\csname stack@top\the\loopcount\endcsname{&;&;&;&;@}
       \expandafter\xdef\csname stack@ttop\the\loopcount\endcsname{&;&;&;&;@}
       \expandafter\xdef\csname stack@tttop\the\loopcount\endcsname{&;&;&;&;@}
       \expandafter\xdef\csname stack@ttttop\the\loopcount\endcsname{&;&;&;&;@}
       \expandafter\xdef\csname stack@bottom\the\loopcount\endcsname{&;&;&;&;@}
       \expandafter\xdef\csname stack@bbottom\the\loopcount\endcsname{&;&;&;&;@}
       \expandafter\xdef\csname stack@bbbottom\the\loopcount\endcsname{&;&;&;&;@}
       \expandafter\xdef\csname stack@bbbbottom\the\loopcount\endcsname{&;&;&;&;@}
       \expandafter\xdef\csname seq@gap\the\loopcount\endcsname{no}
       \expandafter\xdef\csname seq@start\the\loopcount\endcsname{1}
       \expandafter\xdef\csname mol@weight\the\loopcount\endcsname{0}
       \expandafter\xdef\csname ch@rge\the\loopcount\endcsname{0}
       \ifnum\loopcount=1 \xdef\seq@order{\the\loopcount}
       \else              \xdef\seq@order{\seq@order,\the\loopcount}\fi
    \ifnum\loopcount<\seq@count \repeat
    \xdef\seq@order{\seq@order,@}
 \killseq@count=\seq@count
 \fi
 \fi
}

%%%%%  TeXshade

\def\calc@widths{%
   \fontfamily{cmss}\fontseries{m}\fontshape{n}
   \selectfont
   \setbox1=\hbox{\residues@size{A}}\logo@height=\ht1
   \setlength\logo@height{\logo@stretch\logo@height}
   \fontfamily{\residues@family}
   \fontseries{\residues@series}
   \fontshape{\residues@shape}
   \selectfont
   \setbox1=\hbox{\residues@size{W}}\box@width=1.15\wd1
   \global\setlength\box@width{\char@stretch\box@width}
   \box@height=1.2\ht1
   \setbox1=\hbox{\residues@size{g}}\box@depth=1.1\dp1
   \global\setlength\box@height{\line@stretch\box@height}
   \global\setlength\box@depth{\line@stretch\box@depth}
   \baselineskip=\box@height \advance\baselineskip by \box@depth
   \lineskip=0pt
   \ifshow@cons\setbox1=\hbox{\cons@name}\name@width=\wd1\else\name@width=0pt\fi
   \loopcount=0
   \expandafter\getregion@fromstack{\the\loopcount}
   \expandafter\getregion@fromtintstack{\the\loopcount}
   \expandafter\getregion@fromemphstack{\the\loopcount}
   \expandafter\getregion@fromlowerstack{\the\loopcount}
   \expandafter\getregion@fromframestack{\the\loopcount}
   \expandafter\getregion@fromshadingstack{\the\loopcount}
   \xdef\bottop@{top}       \expandafter\getregion@fromfstack{\the\loopcount}
   \xdef\bottop@{ttop}      \expandafter\getregion@fromfstack{\the\loopcount}
   \xdef\bottop@{tttop}     \expandafter\getregion@fromfstack{\the\loopcount}
   \xdef\bottop@{ttttop}    \expandafter\getregion@fromfstack{\the\loopcount}
   \xdef\bottop@{bottom}    \expandafter\getregion@fromfstack{\the\loopcount}
   \xdef\bottop@{bbottom}   \expandafter\getregion@fromfstack{\the\loopcount}
   \xdef\bottop@{bbbottom}  \expandafter\getregion@fromfstack{\the\loopcount}
   \xdef\bottop@{bbbbottom} \expandafter\getregion@fromfstack{\the\loopcount}
   \fontfamily{\namestext@family}
   \fontseries{\namestext@series}
   \fontshape{\namestext@shape}
   \selectfont
   \ifshow@logo
      \ifx\logo@name@user\ampers@nd
         \setbox1=\hbox{\namestext@size logo\kern1em}
      \else
         \setbox1=\hbox{\namestext@size\logo@name@user\kern1em}
      \fi
      \ifnum\wd1>\name@width \name@width=\wd1 \fi
   \fi
   \ifshow@sublogo
      \ifx\sublogo@name@user\ampers@nd
         \setbox1=\hbox{\namestext@size family\kern1em}
      \else
         \setbox1=\hbox{\namestext@size\sublogo@name@user\kern1em}
      \fi
      \ifnum\wd1>\name@width \name@width=\wd1 \fi
      \ifx\hide@negatives\n@
         \setbox1=\hbox{\namestext@size\sublogo@name@neg\kern1em}
      \fi
      \ifnum\wd1>\name@width \name@width=\wd1 \fi
   \fi
   \loopcount=1
   \loop
       \setbox1=\hbox{\namestext@size\csname newseqname\the\loopcount\endcsname}
       \ifnum\wd1>\name@width \name@width=\wd1 \fi
       \ifnum\ht1>\box@height
          \box@height=1.1\ht1
          \global\setlength\box@height{\line@stretch\box@height}
          \baselineskip=\box@height \advance\baselineskip by \box@depth
       \fi
       \ifnum\dp1>\box@depth
          \box@depth=1.1\dp1
          \global\setlength\box@depth{\line@stretch\box@depth}
          \baselineskip=\box@height \advance\baselineskip by \box@depth
       \fi
       \expandafter\getregion@fromstack{\the\loopcount}
       \expandafter\getregion@fromtintstack{\the\loopcount}
       \expandafter\getregion@fromemphstack{\the\loopcount}
       \expandafter\getregion@fromlowerstack{\the\loopcount}
       \expandafter\getregion@fromframestack{\the\loopcount}
       \expandafter\getregion@fromshadingstack{\the\loopcount}
       \xdef\bottop@{top}       \expandafter\getregion@fromfstack{\the\loopcount}
       \xdef\bottop@{ttop}      \expandafter\getregion@fromfstack{\the\loopcount}
       \xdef\bottop@{tttop}     \expandafter\getregion@fromfstack{\the\loopcount}
       \xdef\bottop@{ttttop}    \expandafter\getregion@fromfstack{\the\loopcount}
       \xdef\bottop@{bottom}    \expandafter\getregion@fromfstack{\the\loopcount}
       \xdef\bottop@{bbottom}   \expandafter\getregion@fromfstack{\the\loopcount}
       \xdef\bottop@{bbbottom}  \expandafter\getregion@fromfstack{\the\loopcount}
       \xdef\bottop@{bbbbottom} \expandafter\getregion@fromfstack{\the\loopcount}
       \innerloopcount = \csname seq@len\the\loopcount\endcsname
       \advance\innerloopcount by \csname seq@start\the\loopcount\endcsname
       \advance\innerloopcount by -1
       \expandafter\xdef\csname seq@len\the\loopcount\endcsname{\the\innerloopcount}
       \advance\loopcount by 1
   \ifnum\loopcount>\seq@count \else \repeat
   \advance\name@width by -2em
   \global\xdef\name@@width{\the\name@width}
   \advance\name@width by 3em
   \fontfamily{\numbertext@family}
   \fontseries{\numbertext@series}
   \fontshape{\numbertext@shape}
   \selectfont
   \setbox1=\hbox{\bottomruler@size\num@width\,-}
   \xdef\ruler@width{\the\wd1}
   \setbox1=\hbox{\numbertext@size\num@width}
   \number@width=\wd1
   \advance\number@width by 1em
   \loopcount=\linewidth
   \ifnames@   \advance\loopcount by -\name@width \fi
   \width@tmp=1.6\box@width
   \ifnumbers@
      \ifnumbers@left  
         \advance\loopcount by -\number@width
      \else
         \xdef\first@{left}
         \ifx\show@logoscale\first@
            \advance\loopcount by -\width@tmp
         \fi
      \fi
      \ifnumbers@right  
         \advance\loopcount by -\number@width
      \else
         \xdef\first@{right}
         \ifx\show@logoscale\first@
            \advance\loopcount by -\width@tmp
         \fi
      \fi
   \else
      \xdef\first@{left}
      \ifx\show@logoscale\first@
         \advance\loopcount by -\width@tmp
      \fi
      \xdef\first@{right}
      \ifx\show@logoscale\first@
         \advance\loopcount by -\width@tmp
      \fi
      \xdef\first@{leftright}
      \ifx\show@logoscale\first@
         \advance\loopcount by -\width@tmp
         \advance\loopcount by -\width@tmp
      \fi
   \fi
   \divide\loopcount by \box@width
   \divide\loopcount by 5 \multiply\loopcount by 5
   \ifrpl@fix\else
   \ifnum\res@perline>\loopcount \res@perline=\loopcount\fi\fi
   \ifnum\finger@linenum>0
      \width@tmp=\linewidth
      \ifnames@   \advance\width@tmp by -\name@width \fi
      \ifnumbers@ \advance\width@tmp by -\number@width \fi
      \divide\width@tmp by \finger@linenum
      \global\setlength\box@width{\width@tmp}
   \fi
   \center@fill=\linewidth
   \ifnames@   \advance\center@fill by -\name@width \fi
   \ifnumbers@ \advance\center@fill by -\number@width \fi
   \width@tmp=\box@width \multiply\width@tmp by \res@perline
   \advance\center@fill by -\width@tmp
   \ifx\out@put\y@\leftskip\c@factor\center@fill\fi
   \ifx\logo@colors@set\n@
     \if\seq@type N
      \clearlogocolors[Yellow]
      \logocolor{G}{Black}
      \logocolor{A}{Green}
      \logocolor{TU}{Red}
      \logocolor{C}{Blue}
     \else
      \clearlogocolors
      \logocolor{DE}{Red}
      \logocolor{CM}{Yellow}
      \logocolor{KR}{Blue}
      \logocolor{ST}{Orange}
      \logocolor{FY}{MidnightBlue}
      \logocolor{NQ}{Cyan}
      \logocolor{G}{LightGray}
      \logocolor{LVI}{Green}
      \logocolor{A}{DarkGray}
      \logocolor{W}{CarnationPink}
      \logocolor{H}{CornflowerBlue}
      \logocolor{P}{Apricot}
      \logocolor{BZ}{LightMagenta}
     \fi   
   \fi
}

\def\multiple@dssp{%
   \advance\loopcount by 1
   \include@DSSP
   \ifnum\loopcount<\dssp@num\multiple@dssp\fi}

\def\multiple@stride{%
   \advance\loopcount by 1
   \include@stride
   \ifnum\loopcount<\stride@num\multiple@stride\fi}

\def\multiple@PHD{%
   \advance\loopcount by 1
   \include@PHD
   \ifnum\loopcount<\PHD@num\multiple@PHD\fi}

\def\multiple@HMMTOP{%
   \advance\loopcount by 1
   \include@HMMTOP
   \ifnum\loopcount<\HMMTOP@num\multiple@HMMTOP\fi}

\newenvironment{texshade}[2][&]%
  {\standarddefinitions
   \c@d@ns
   \xdef\first@{#1}\ifx\first@\ampers@nd\else\input{#1}\fi
   \readalignfile{#2}
   }%
  {\ifnum\seq@count>0
      \loopcount=0
      \ifnum\loopcount<\dssp@num \multiple@dssp\fi
      \loopcount=0
      \ifnum\loopcount<\stride@num \multiple@stride\fi
      \loopcount=0
      \ifnum\loopcount<\PHD@num \multiple@PHD\fi
      \loopcount=0
      \ifnum\loopcount<\HMMTOP@num \multiple@HMMTOP\fi
      \loopcount=1 \kill@seqnow
      \reorder@seqs\seq@order \seq@count=\killseq@count
      \global\xdef\seq@num{\the\seq@count}
      \clear@sim@count
      \ifnum\csname res@count\start@seq\endcsname<0
       \ifnum\start@num>0
         \loopcount=\start@num
         \advance\loopcount by -1
         \xdef\start@num{\the\loopcount}
       \fi
      \fi
      \ifnum\rule@num>0  
         \loopcount=\csname res@count\rule@num\endcsname
         \divide\loopcount by \ruler@step
         \multiply\loopcount by \ruler@step
         \ifnum\loopcount<0 
            \ifnum\ruler@step<3
               \advance\loopcount by \ruler@step 
            \fi
         \else 
            \advance\loopcount by \ruler@step 
         \fi
         \xdef\rule@tens{\the\loopcount} 
      \else
         \loopcount=\cons@count  
         \divide\loopcount by \ruler@step
         \multiply\loopcount by \ruler@step
         \ifnum\loopcount<0 
            \ifnum\ruler@step<3
               \advance\loopcount by \ruler@step 
            \fi
         \else 
            \advance\loopcount by \ruler@step 
         \fi
         \xdef\rule@tens{\the\loopcount}  
      \fi
      \xdef\first@{top}
      \ifx\cap@pos\first@
         \xdef\@captype{figure}
         \ifx\c@pshort\n@
           \caption{\c@p}
         \else
           \caption[\c@pshort]{\c@p}
         \fi
      \fi
      \loopcount = \thresh@ld \multiply\loopcount by \seq@count \divide\loopcount by 100
      \xdef\thresh@ld@{\the\loopcount}
      \loopcount = \all@thresh@ld \multiply\loopcount by \seq@count \divide\loopcount by 100
      \xdef\all@thresh@ld@{\the\loopcount}
      \if\exp@rt@chimera\y@%
         \ifshow@logo%
             \init@chimera@file%
         \else%
         \ifshow@sublogo%
             \init@chimera@file%
         \fi\fi%
      \fi%
      \bgroup
         \ifx\out@put\y@\bigskip\fi
         \iffuncmode \show@consfalse \fi
         \ifall@fshade \iffuncmode \else \all@fshadefalse \fi\fi
         \ifnum\finger@linenum>0
             \show@consfalse
             \hidechartrue
             \message{<Fingerprinting---please wait>}
         \fi
         \calc@widths
         \read@lines
	     \save@lengths
         \iflegend@ 
            \vspace{\vspace@legend}
            \setbox1=\vbox{\do@legend}
            \vbox{\do@legend}\par
            \ifnum\ht1<-\vspace@legend
            \vspace{-\ht1}\vspace{-\vspace@legend}
            \fi
         \fi
      \egroup
      \xdef\first@{bottom}
      \ifx\cap@pos\first@ 
         \vspace{-\baselineskip}
         \xdef\@captype{figure}
         \ifx\c@pshort\n@
           \caption{\c@p}
         \else
           \caption[\c@pshort]{\c@p}
         \fi
      \fi
   \fi}

\def\standarddefinitions{%
\xdef\prfx{pep}\clear@groups\clear@sims
\xdef\prfx{DNA}\clear@groups\clear@sims
\clearfuncgroups
\clear@simpairs
\shadingcolors{blues}
\shadingcolors{reds}
\shadingcolors{greens}
\shadingcolors{grays}
\shadingcolors{black}
\loopcount=\@lign@count
\advance\loopcount by 1\relax
\xdef\@lign@count{\the\loopcount}
\start@true \xdef\start@num{1} \xdef\start@seq{0}
\stop@false \xdef\end@num{n} \xdef\seq@regions{0}
\cons@count=0
\expandafter\xdef\csname res@count0\endcsname{0}
\xdef\allow@zero{n} \xdef\c@ns@shift{0}
\regionalshadefalse\regionalemphfalse\regionallowerfalse\regionaltintfalse
\frame@false\shading@false
\xdef\ruler@rot{0}
\topfeaturefalse    \bottomfeaturefalse
\ttopfeaturefalse   \bbottomfeaturefalse
\tttopfeaturefalse  \bbbottomfeaturefalse
\ttttopfeaturefalse \bbbbottomfeaturefalse
\all@fshadefalse\hidecharfalse
\xdef\finger@linenum{0}
\xdef\hide@seqs{n}
\xdef\dssp@num{0} \xdef\stride@num{0} \xdef\PHD@num{0} \xdef\HMMTOP@num{0}
\xdef\bottop@{top}      \expandafter\xdef\csname feature@\bottop@\endcsname{0}
\xdef\bottop@{ttop}     \expandafter\xdef\csname feature@\bottop@\endcsname{0}
\xdef\bottop@{tttop}    \expandafter\xdef\csname feature@\bottop@\endcsname{0}
\xdef\bottop@{ttttop}   \expandafter\xdef\csname feature@\bottop@\endcsname{0}
\xdef\bottop@{bottom}   \expandafter\xdef\csname feature@\bottop@\endcsname{0}
\xdef\bottop@{bbottom}  \expandafter\xdef\csname feature@\bottop@\endcsname{0}
\xdef\bottop@{bbbottom} \expandafter\xdef\csname feature@\bottop@\endcsname{0}
\xdef\bottop@{bbbbottom}\expandafter\xdef\csname feature@\bottop@\endcsname{0}
\xdef\frame@{0}
\xdef\show@Hdssp{no} \xdef\show@Gdssp{no} \xdef\show@Idssp{no}
\xdef\show@Edssp{no} \xdef\show@Bdssp{no} \xdef\show@Tdssp{no}
\xdef\show@Sdssp{no}
\xdef\show@Hstride{no} \xdef\show@Gstride{no} \xdef\show@Istride{no}
\xdef\show@Estride{no} \xdef\show@Bstride{no} \xdef\show@Tstride{no}
\xdef\show@itop{no} \xdef\show@etop{no} \xdef\show@TMtop{no}
\xdef\show@i@HMMTOP{no} \xdef\show@e@HMMTOP{no} \xdef\show@TM@HMMTOP{no}
\xdef\show@Hsec{no} \xdef\show@Esec{no}
\xdef\collect@restop{no}     \xdef\collect@resttop{no}
\xdef\collect@restttop{no}   \xdef\collect@resttttop{no}
\xdef\collect@resbottom{no}  \xdef\collect@resbbottom{no}
\xdef\collect@resbbbottom{no}\xdef\collect@resbbbbottom{no}
\xdef\tr@nslatetop{}         \xdef\tr@nslatettop{}
\xdef\tr@nslatetttop{}       \xdef\tr@nslatettttop{}
\xdef\tr@nslatebottom{}      \xdef\tr@nslatebbottom{}
\xdef\tr@nslatebbbottom{}    \xdef\tr@nslatebbbbottom{}
\xdef\tr@nsseqtop{0}         \xdef\tr@nsseqttop{0}
\xdef\tr@nsseqtttop{0}       \xdef\tr@nsseqttttop{0}
\xdef\tr@nsseqbottom{0}      \xdef\tr@nsseqbbottom{0}
\xdef\tr@nsseqbbbottom{0}    \xdef\tr@nsseqbbbbottom{0}
\xdef\triple@counttop{0}     \xdef\triple@countttop{0}
\xdef\triple@counttttop{0}   \xdef\triple@countttttop{0}
\xdef\triple@countbottom{0}  \xdef\triple@countbbottom{0}
\xdef\triple@countbbbottom{0}\xdef\triple@countbbbbottom{0}
\xdef\last@@restop{}         \xdef\last@@resttop{}
\xdef\last@@restttop{}       \xdef\last@@resttttop{}
\xdef\last@@resbottom{}      \xdef\last@@resbbottom{}
\xdef\last@@resbbbottom{}    \xdef\last@@resbbbbottom{}
\xdef\out@put{y}             \xdef\m@p{no}
\xdef\t@sp@ce{0mm}           \xdef\tt@sp@ce{0mm}
\xdef\ttt@sp@ce{0mm}         \xdef\tttt@sp@ce{0mm}
\xdef\b@sp@ce{0mm}           \xdef\bb@sp@ce{0mm}
\xdef\bbb@sp@ce{0mm}         \xdef\bbbb@sp@ce{0mm}
\xdef\ruler@sp@ce{0mm}
\xdef\seq@gap@num{0}        
\xdef\h@ndalign{no}         \xdef\sep@space{0pt}
\xdef\c@pshort{n}
\xdef\bottom@stretch{n}     \xdef\bbottom@stretch{n}
\xdef\bbbottom@stretch{n}   \xdef\bbbbottom@stretch{n}
\xdef\c@nscol{}             \xdef\c@nssc@le{ColdHot}
\xdef\collect@cons@colors{no} \xdef\cons@now{no}
\xdef\res@numA{0}         \xdef\res@numB{0}
\xdef\res@numC{0}         \xdef\res@numD{0}
\xdef\res@numE{0}         \xdef\res@numF{0}
\xdef\res@numG{0}         \xdef\res@numH{0}
\xdef\res@numI{0}         \xdef\res@numJ{0}
\xdef\res@numK{0}         \xdef\res@numL{0}
\xdef\res@numM{0}         \xdef\res@numN{0}
\xdef\res@numO{0}         \xdef\res@numP{0}
\xdef\res@numQ{0}         \xdef\res@numR{0}
\xdef\res@numS{0}         \xdef\res@numT{0}
\xdef\res@numU{0}         \xdef\res@numV{0}
\xdef\res@numW{0}         \xdef\res@numX{0}
\xdef\res@numY{0}         \xdef\res@numZ{0}
\expandafter\xdef\csname res@num\questi@n\endcsname{0}
\expandafter\xdef\csname res@num\d@t\endcsname{0}
\expandafter\xdef\csname res@num\equ@l\endcsname{0}
\xdef\res@corrA{0}        \xdef\res@corrB{0}
\xdef\res@corrC{0}        \xdef\res@corrD{0}
\xdef\res@corrE{0}        \xdef\res@corrF{0}
\xdef\res@corrG{0}        \xdef\res@corrH{0}
\xdef\res@corrI{0}        \xdef\res@corrJ{0}
\xdef\res@corrK{0}        \xdef\res@corrL{0}
\xdef\res@corrM{0}        \xdef\res@corrN{0}
\xdef\res@corrO{0}        \xdef\res@corrP{0}
\xdef\res@corrQ{0}        \xdef\res@corrR{0}
\xdef\res@corrS{0}        \xdef\res@corrT{0}
\xdef\res@corrU{0}        \xdef\res@corrV{0}
\xdef\res@corrW{0}        \xdef\res@corrX{0}
\xdef\res@corrY{0}        \xdef\res@corrZ{0}
\expandafter\xdef\csname res@corr\questi@n\endcsname{0}
\expandafter\xdef\csname res@corr\d@t\endcsname{0}
\expandafter\xdef\csname res@corr\equ@l\endcsname{0}
\xdef\corr@max{0}         \xdef\do@freq@correction{n}
\xdef\res@total{0}
\xdef\stack@sequencelogo{} \xdef\bit@total{0}
\clearlogocolors \xdef\logo@colors@set{n}
\show@sublogofalse    \xdef\stack@sublogo{}
\xdef\clear@logo{n}       \xdef\subfamily@threshold{50}
\xdef\subfamily@seq{1}    \xdef\sig@max{100000}
\xdef\sublogo@num{}       \xdef\sublogo@sig{}
\xdef\subfamily@count{0}  \xdef\sublogo@name@neg{}
\xdef\logo@name@user{&}   \xdef\sublogo@name@user{&}
\expandafter\xdef\csname group@num1\endcsname{0}
\expandafter\xdef\csname group@num2\endcsname{0}
\xdef\exp@rt@num{0}       \xdef\exp@rt{n}
\xdef\exp@rt@chimera@num{1}   \xdef\exp@rt@chimera{n} \xdef\chimera@stack{}
\xdef\chimerab@llSc@le{1.0}
\memelabelcutoff{0.1}
\xdef\chimera@labelOne@set{n}
\xdef\chimera@labelTwo@set{n}
\xdef\chimera@labelThree@set{n}
\xdef\chimera@labelOne@remaining@set{n}
\xdef\chimera@labelTwo@remaining@set{n}
\xdef\chimera@labelThree@remaining@set{n}
\xdef\chimera@xdist@nce{0} \xdef\chimera@ydist@nce{40} \xdef\chimera@zdist@nce{0}
\xdef\chimerach@in{} \xdef\chimera@seqnum{0}
\xdef\chimera@label@subfamily{} \xdef\chimera@label@remaining{}
\memeStandardcolors
\xdef\chimera@subfamily@col@red{0}%
\xdef\chimera@subfamily@col@blue{0}%
\xdef\chimera@subfamily@col@yellow{0}%
\xdef\chimera@subfamily@col@black{0}%
\xdef\chimera@subfamily@col@white{0}%
\xdef\chimera@subfamily@col@none{0}%
\xdef\chimera@remaining@col@red{0}%
\xdef\chimera@remaining@col@blue{0}%
\xdef\chimera@remaining@col@yellow{0}%
\xdef\chimera@remaining@col@black{0}%
\xdef\chimera@remaining@col@white{0}%
\xdef\chimera@remaining@col@none{0}%
\xdef\echo@filename{}
\xdef\divref@{0}
\xdef\all@thresh@ld{100}  \all@shadefalse
\hidefeaturenames         \hidefeaturestylenames
\xdef\T@coffee@ccons{n}   \xdef\T@coffee@bcons{n}
\xdef\dom@in{no}          \xdef\label@motif{no}
\xdef\motif@num{0}        \xdef\all@out{n}
\hfuzz9999pt
\xdef\left@frame@open{n}
\xdef\right@frame@open{n}

%%%%%%%%%%%%%%%%%%%%%%%%%%%%%%%%%%%%%%%%%%%%%%%%%%%%%%%%%%%%%%%%%%%%%%%%%%%%%%%%
%%%%%                                                                      %%%%%
%%%%%   Default parameter settings for the  LaTeX  ``TeXshade''  package   %%%%%
%%%%%                                                                      %%%%%
%%%%%%%%%%%%%%%%%%%%%%%%%%%%%%%%%%%%%%%%%%%%%%%%%%%%%%%%%%%%%%%%%%%%%%%%%%%%%%%%
%%%%%                                                                      %%%%%
%%%%%   Under any circumstances:                                           %%%%%
%%%%%                                                                      %%%%%
%%%%%                   DO NOT CHANGE ANY SETTINGS HERE !!!                %%%%%
%%%%%                                                                      %%%%%
%%%%%   Please define your personal parameter file!  Store your new file   %%%%%
%%%%%   together with this style-file in the same directory and load the   %%%%%
%%%%%   file by naming it as an optional parameter in the `texshade' en-   %%%%%
%%%%%   vironment. The file `texshade.def' can be used as a template for   %%%%%
%%%%%   the new creation. See the manual for further help.                 %%%%%
%%%%%                                                                      %%%%%
%%%%%                              THANK YOU !!!                           %%%%%
%%%%%                                                                      %%%%%
%%%%%%%%%%%%%%%%%%%%%%%%%%%%%%%%%%%%%%%%%%%%%%%%%%%%%%%%%%%%%%%%%%%%%%%%%%%%%%%%

\shadingmode{identical}             %  Shade identical residues only
\shadingcolors{blues}               %  Select the blue color scheme for shading
\constoallseqs                      %  Calculate consensus considering all seqs
\threshold{50}                      %  Consensus threshold percentage is 50
\residuesperline{999}               %  As many residues as possible per line
\numberingwidth{9999}               %  Assign space for 4 digit numbering
\charstretch{1}                     %  Do not stretch character width
\linestretch{1}                     %  Do not stretch lines vertically
\gapchar{.}                         %  . is printed in sequence gaps
\gaprule{0.3pt}                     %  If a rule is printed in gaps use 0.3 pt
\gapcolors{Black}{White}            %  Gap symbols appear `Black on White'
\domaingaprule{1pt}                 %  If a rule is set between domains use 1 pt
\domaingapcolors{Red}{White}        %  Domain separation line appears `Red on White'
\numberingcolor{Black}              %  Numbering color is `Black'
\shownumbering{right}               %  Show sequence numbering on the left
\namescolor{Black}                  %  Names' color is `Black'
\featurenamescolor{Black}           %  Feature names' color is `Black'
\featurestylenamescolor{Black}      %  Feature names' color is `Black'
\shownames{left}                    %  Show sequence names on the right
\consensuscolors{Black}{White}      %  All consensus symbols/letters
                {Black}{White}      %  appear `Black on White'
                {Black}{White}      %
\showconsensus{bottom}              %  Show consensus line at bottom with
\hidesequencelogo                   %  Do not show a sequence logo as consensus
\showlogoscale{leftright}           %  Show vertical scale bar if logo is on
\logostretch{1}                     %  Do not stretch sequence logo vertically
\hidesubfamilylogo                  %  Do not show a subfamily logo
\subfamilythreshold{50}             %  Set subfamily threshold to 50%
\shownegatives                      %  Show negative values in subfamily logo
\showrelevance[Black]{*}            %  Label relevant subfamily deviations
\defconsensus{{}}{*}{!}             %  Blank =no match; * =match; ! =all match
\showleadinggaps                    %  Show gap symbols before sequence start
\rulercolor{Black}                  %  Ruler's color is `Black'
\hideruler                          %  Do not show the ruler
\rulersteps{10}                     %  Ruler ticks every 10 residues
\legendcolor{Black}                 %  Legend text color is `Black'
\hidelegend                         %  Do not show the legend
\alignment{center}                  %  Center alignment on page
\medsepline                         %  Medium height if separation line is on
\medblockskip                       %  Medium skip between sequence blocks
\flexblockspace                     %  Use optimized space between blocks
\featurerule{0.5ex}                 %  Set feature rule thickness to 1/2 ex
\bargraphstretch{1}                 %  Do not stretch bars in feature graphs
\colorscalestretch{1}               %  Do not stretch color scales in features
\backtranstext{horizontal}          %  Horizontal triplets in feature texts
\backtranslabel{alternating}        %  Alternating triplets in feature styles
\setfamily{residues}{tt}            %  Use typewriter family for residues
\setseries{residues}{md}            %  Use normal series for residues
\setshape {residues}{up}            %  Use upright shape for residues
\setsize  {residues}{normalsize}    %  Use normal size for residues
\setfamily{numbering}{tt}           %  Use typewriter family for numbering
\setseries{numbering}{md}           %  Use normal series for numbering
\setshape {numbering}{up}           %  Use upright shape for numbering
\setsize  {numbering}{normalsize}   %  Use normal size for numbering
\setfamily{names}{tt}               %  Use typewriter family for names
\setseries{names}{md}               %  Use normal series for names
\setshape {names}{up}               %  Use upright shape for names
\setsize  {names}{normalsize}       %  Use normal size for names
\setfamily{featurenames}{tt}        %  Use typewriter family for feature names
\setseries{featurenames}{md}        %  Use normal series for feature names
\setshape {featurenames}{up}        %  Use upright shape for feature names
\setsize  {featurenames}{normalsize}%  Use normal size for feature  names
\setfamily{featurestylenames}{tt}   %  Use typewriter family for feature style names
\setseries{featurestylenames}{md}   %  Use normal series for feature style names
\setshape {featurestylenames}{up}   %  Use upright shape for feature style names
\setsize  {featurestylenames}{normalsize}%  Use normal size for feature style names
\setfamily{features}{rm}            %  Use roman family for feature texts
\setseries{features}{md}            %  Use normal series for feature texts
\setshape {features}{it}            %  Use italics shape for feature texts
\setsize  {features}{normalsize}    %  Use normal size for feature texts
\setfamily{featurestyles}{tt}       %  Use typewriter family for feature styles
\setseries{featurestyles}{md}       %  Use normal series for feature styles
\setshape {featurestyles}{up}       %  Use upright shape for feature styles
\setsize {featurestyles}{normalsize}%  Use normal size for feature styles
\setfamily{legend}{tt}              %  Use typewriter family for legend texts
\setseries{legend}{md}              %  Use normal series for legend texts
\setshape {legend}{up}              %  Use upright shape for legend texts
\setsize  {legend}{normalsize}      %  Use normal size for legend texts
\setfamily{ruler}{sf}               %  Use sans serif font for ruler numbers
\tintdefault{medium}                %  Use medium tint intensity
\emphdefault{it}                    %  Use italics to emphasize regions
\showonPHDsec{alpha,beta}           %  Show helices and strands (PHD input)
\showonPHDtopo{internal,external,TM}%  Show int., ext. and TM's (PHD input)
\showonHMMTOP{TM}                   %  Show TM's (not int/ext on HMMTOP input)
\showonSTRIDE{alpha,3-10,pi,beta}   %  Show helices and strands (STRIDE input)
\showonDSSP{alpha,3-10,pi,beta}     %  Show helices and strands (DSSP input)
\secondcolumnDSSP                   %  Use numbering from 2. column in DSSP
\appearance{PHDtopo}{internal}      % \
           {bottom}{'-'}            %  \
           {int.\ \Alphacount}      %   |
\appearance{PHDtopo}{external}      %   |
           {top}{,-,}               %   |
           {ext.\ \Alphacount}      %   |
\appearance{PHDtopo}{TM}{top}       %   |
     {box[LightGray]:TM\numcount}{} %   |
\appearance{HMMTOP}{internal}       %   |
           {bottom}{'-'}            %   |
           {int.\ \Alphacount}      %   |
\appearance{HMMTOP}{external}       %   |
           {top}{,-,}               %   |
           {ext.\ \Alphacount}      %   |
\appearance{HMMTOP}{TM}{top}        %   |
           {helix}{TM\numcount}     %   |
\appearance{PHDsec}{alpha}{top}     %   |
          {box:$\alpha$\numcount}{} %   |
\appearance{PHDsec}{beta}{top}      %   |
          {-->}{$\beta$\numcount}   %   |
\appearance{STRIDE}{alpha}{top}     %   |
          {box:$\alpha$\numcount}{} %   |    Definitions for the appearance of
\appearance{STRIDE}{3-10}{top}      %    \
           {fill:$\circ$}{3$_{10}$} %     >  secondary structures included from|
\appearance{STRIDE}{pi}             %    /
           {top}{---}{$\pi$}        %   |    PHD-, STRIDE-, or DSSP-files.
\appearance{STRIDE}{beta}{top}      %   |
          {-->}{$\beta$\numcount}   %   |
\appearance{STRIDE}{bridge}         %   |
           {top}{fill:$\uparrow$}{} %   |
\appearance{STRIDE}{turn}           %   |
           {top}{,-,}{turn}         %   |
\appearance{DSSP}{alpha}{top}       %   |
          {box:$\alpha$\numcount}{} %   |
\appearance{DSSP}{3-10}{top}        %   |
           {fill:$\circ$}{3$_{10}$} %   |
\appearance{DSSP}{pi}               %   |
           {top}{---}{$\pi$}        %   |
\appearance{DSSP}{beta}{top}        %   |
          {-->}{$\beta$\numcount}   %   |
\appearance{DSSP}{bridge}           %   |
           {top}{fill:$\uparrow$}{} %   |
\appearance{DSSP}{turn}             %   |
           {top}{,-,}{turn}         %   |
\appearance{DSSP}{bend}{top}        %  /
           {fill:$\diamond$}{}      % /

\pepgroups{FYW,ILVM,RK,DE,GA,ST,NQ} %  Amino acid grouping due to similarity

\pepsims{F}{YW}                     %  Y and W are similar to F
\pepsims{Y}{WF}                     %  W and F are similar to Y
\pepsims{W}{YF}                     %  Y and F are similar to W

\pepsims{I}{LVM}                    %  L, V and M are similar to I
\pepsims{L}{VMI}                    %  V, M and I are similar to L
\pepsims{V}{MIL}                    %  M, I and L are similar to V

\pepsims{R}{KH}                     %  K and H are similar to R
\pepsims{K}{HR}                     %  H and R are similar to K
\pepsims{H}{RK}                     %  R and K are similar to H

\pepsims{A}{GS}                     %  G and S are similar to A
\pepsims{G}{A}                      %  A (but not S) is similar to G

\pepsims{S}{TA}                     %  T and A are similar to S
\pepsims{T}{S}                      %  S (but not A) is similar to T

\pepsims{D}{EN}                     %  E and N (but not Q) are similar to D
\pepsims{E}{DQ}                     %  D and Q (but not N) are similar to E
\pepsims{N}{QD}                     %  Q and D (but not E) are similar to N
\pepsims{Q}{NE}                     %  N and E (but not D) are similar to Q

\DNAgroups{GAR,CTY}                 %  Nucleotide grouping due to similarity

\DNAsims{A}{GR}                     %  G and R are similar to A
\DNAsims{G}{AR}                     %  A and R are similar to G
\DNAsims{R}{AG}                     %  A and G are similar to R

\DNAsims{C}{TY}                     %  T and Y are similar to C
\DNAsims{T}{CY}                     %  C and Y are similar to T
\DNAsims{Y}{CT}                     %  C and T are similar to Y

}
\catcode`\@=12
%</texshade>
%    \end{macrocode}
%    \subsection{\file{texshade.def}}
%    \begin{macrocode}
%<*definitions>
%%%%%%%%%%%%%%%%%%%%%%%%%%%%%%%%%%%%%%%%%%%%%%%%%%%%%%%%%%%%%%%%%%%%%%%%%%%%%%%%
%%%%%                                                                      %%%%%
%%%%%   Default parameter settings for the  LaTeX  ``TeXshade''  package   %%%%%
%%%%%                                                                      %%%%%
%%%%%%%%%%%%%%%%%%%%%%%%%%%%%%%%%%%%%%%%%%%%%%%%%%%%%%%%%%%%%%%%%%%%%%%%%%%%%%%%
%%%%%                                                                      %%%%%
%%%%%   This example file contains all standard settings of the TeXshade   %%%%%
%%%%%   package. It can be used as a template for the creation of perso-   %%%%%
%%%%%   nal parameter files.  All TeXshade user commands are allowed and   %%%%%
%%%%%   functional when specified here.                                    %%%%%
%%%%%                                                                      %%%%%
%%%%%   To activate these settings for your alignment load  this file by   %%%%%
%%%%%   naming it as optional parameter at the beginning of the texshade   %%%%%
%%%%%   environment, e.g.                                                  %%%%%
%%%%%                                                                      %%%%%
%%%%%          \begin{texshade}[myparameterfile]{alignmentfile}            %%%%%
%%%%%             .                                                        %%%%%
%%%%%             .                                                        %%%%%
%%%%%          \end{texshade}                                              %%%%%
%%%%%                                                                      %%%%%
%%%%%%%%%%%%%%%%%%%%%%%%%%%%%%%%%%%%%%%%%%%%%%%%%%%%%%%%%%%%%%%%%%%%%%%%%%%%%%%%


\shadingmode{identical}             %  Shade identical residues only
\shadingcolors{blues}               %  Select the blue color scheme for shading
\constoallseqs                      %  Calculate consensus considering all seqs
\threshold{50}                      %  Consensus threshold percentage is 50
\residuesperline{999}               %  As many residues as possible per line
\numberingwidth{9999}               %  Assign space for 4 digit numbering
\charstretch{1}                     %  Do not stretch character width
\linestretch{1}                     %  Do not stretch lines vertically
\gapchar{.}                         %  . is printed in sequence gaps
\gaprule{0.3pt}                     %  If a rule is printed in gaps use 0.3 pt
\gapcolors{Black}{White}            %  Gap symbols appear `Black on White'
\domaingaprule{1pt}                 %  If a rule is set between domains use 1 pt
\domaingapcolors{Red}{White}        %  Domain separation line appears `Red on White'
\numberingcolor{Black}              %  Numbering color is `Black'
\shownumbering{right}               %  Show sequence numbering on the left
\namescolor{Black}                  %  Names' color is `Black'
\featurenamescolor{Black}           %  Feature names' color is `Black'
\featurestylenamescolor{Black}      %  Feature names' color is `Black'
\shownames{left}                    %  Show sequence names on the right
\consensuscolors{Black}{White}      %  All consensus symbols/letters
                {Black}{White}      %  appear `Black on White'
                {Black}{White}      %
\showconsensus{bottom}              %  Show consensus line at bottom with
\hidesequencelogo                   %  Do not show a sequence logo as consensus
\showlogoscale{leftright}           %  Show vertical scale bar if logo is on
\logostretch{1}                     %  Do not stretch sequence logo vertically
\hidesubfamilylogo                  %  Do not show a subfamily logo
\subfamilythreshold{50}             %  Set subfamily threshold to 50%
\shownegatives                      %  Show negative values in subfamily logo
\showrelevance[Black]{*}            %  Label relevant subfamily deviations
\defconsensus{{}}{*}{!}             %  Blank =no match; * =match; ! =all match
\showleadinggaps                    %  Show gap symbols before sequence start
\rulercolor{Black}                  %  Ruler's color is `Black'
\hideruler                          %  Do not show the ruler
\rulersteps{10}                     %  Ruler ticks every 10 residues
\legendcolor{Black}                 %  Legend text color is `Black'
\hidelegend                         %  Do not show the legend
\alignment{center}                  %  Center alignment on page
\medsepline                         %  Medium height if separation line is on
\medblockskip                       %  Medium skip between sequence blocks
\flexblockspace                     %  Use optimized space between blocks
\featurerule{0.5ex}                 %  Set feature rule thickness to 1/2 ex
\bargraphstretch{1}                 %  Do not stretch bars in feature graphs
\colorscalestretch{1}               %  Do not stretch color scales in features
\backtranstext{horizontal}          %  Horizontal triplets in feature texts
\backtranslabel{alternating}        %  Alternating triplets in feature styles
\setfamily{residues}{tt}            %  Use typewriter family for residues
\setseries{residues}{md}            %  Use normal series for residues
\setshape {residues}{up}            %  Use upright shape for residues
\setsize  {residues}{normalsize}    %  Use normal size for residues
\setfamily{numbering}{tt}           %  Use typewriter family for numbering
\setseries{numbering}{md}           %  Use normal series for numbering
\setshape {numbering}{up}           %  Use upright shape for numbering
\setsize  {numbering}{normalsize}   %  Use normal size for numbering
\setfamily{names}{tt}               %  Use typewriter family for names
\setseries{names}{md}               %  Use normal series for names
\setshape {names}{up}               %  Use upright shape for names
\setsize  {names}{normalsize}       %  Use normal size for names
\setfamily{featurenames}{tt}        %  Use typewriter family for feature names
\setseries{featurenames}{md}        %  Use normal series for feature names
\setshape {featurenames}{up}        %  Use upright shape for feature names
\setsize  {featurenames}{normalsize}%  Use normal size for feature  names
\setfamily{featurestylenames}{tt}   %  Use typewriter family for feature style names
\setseries{featurestylenames}{md}   %  Use normal series for feature style names
\setshape {featurestylenames}{up}   %  Use upright shape for feature style names
\setsize  {featurestylenames}{normalsize}%  Use normal size for feature style names
\setfamily{features}{rm}            %  Use roman family for feature texts
\setseries{features}{md}            %  Use normal series for feature texts
\setshape {features}{it}            %  Use italics shape for feature texts
\setsize  {features}{normalsize}    %  Use normal size for feature texts
\setfamily{featurestyles}{tt}       %  Use typewriter family for feature styles
\setseries{featurestyles}{md}       %  Use normal series for feature styles
\setshape {featurestyles}{up}       %  Use upright shape for feature styles
\setsize {featurestyles}{normalsize}%  Use normal size for feature styles
\setfamily{legend}{tt}              %  Use typewriter family for legend texts
\setseries{legend}{md}              %  Use normal series for legend texts
\setshape {legend}{up}              %  Use upright shape for legend texts
\setsize  {legend}{normalsize}      %  Use normal size for legend texts
\setfamily{ruler}{sf}               %  Use sans serif font for ruler numbers
\tintdefault{medium}                %  Use medium tint intensity
\emphdefault{it}                    %  Use italics to emphasize regions
\showonPHDsec{alpha,beta}           %  Show helices and strands (PHD input)
\showonPHDtopo{internal,external,TM}%  Show int., ext. and TM's (PHD input)
\showonHMMTOP{TM}                   %  Show TM's (not int/ext on HMMTOP input)
\showonSTRIDE{alpha,3-10,pi,beta}   %  Show helices and strands (STRIDE input)
\showonDSSP{alpha,3-10,pi,beta}     %  Show helices and strands (DSSP input)
\secondcolumnDSSP                   %  Use numbering from 2. column in DSSP
\appearance{PHDtopo}{internal}      % \
           {bottom}{'-'}            %  \
           {int.\ \Alphacount}      %   |
\appearance{PHDtopo}{external}      %   |
           {top}{,-,}               %   |
           {ext.\ \Alphacount}      %   |
\appearance{PHDtopo}{TM}{top}       %   |
     {box[LightGray]:TM\numcount}{} %   |
\appearance{HMMTOP}{internal}       %   |
           {bottom}{'-'}            %   |
           {int.\ \Alphacount}      %   |
\appearance{HMMTOP}{external}       %   |
           {top}{,-,}               %   |
           {ext.\ \Alphacount}      %   |
\appearance{HMMTOP}{TM}{top}        %   |
           {helix}{TM\numcount}     %   |
\appearance{PHDsec}{alpha}{top}     %   |
          {box:$\alpha$\numcount}{} %   |
\appearance{PHDsec}{beta}{top}      %   |
          {-->}{$\beta$\numcount}   %   |
\appearance{STRIDE}{alpha}{top}     %   |
          {box:$\alpha$\numcount}{} %   |    Definitions for the appearance of
\appearance{STRIDE}{3-10}{top}      %    \
           {fill:$\circ$}{3$_{10}$} %     >  secondary structures included from|
\appearance{STRIDE}{pi}             %    /
           {top}{---}{$\pi$}        %   |    PHD-, STRIDE-, or DSSP-files.
\appearance{STRIDE}{beta}{top}      %   |
          {-->}{$\beta$\numcount}   %   |
\appearance{STRIDE}{bridge}         %   |
           {top}{fill:$\uparrow$}{} %   |
\appearance{STRIDE}{turn}           %   |
           {top}{,-,}{turn}         %   |
\appearance{DSSP}{alpha}{top}       %   |
          {box:$\alpha$\numcount}{} %   |
\appearance{DSSP}{3-10}{top}        %   |
           {fill:$\circ$}{3$_{10}$} %   |
\appearance{DSSP}{pi}               %   |
           {top}{---}{$\pi$}        %   |
\appearance{DSSP}{beta}{top}        %   |
          {-->}{$\beta$\numcount}   %   |
\appearance{DSSP}{bridge}           %   |
           {top}{fill:$\uparrow$}{} %   |
\appearance{DSSP}{turn}             %   |
           {top}{,-,}{turn}         %   |
\appearance{DSSP}{bend}{top}        %  /
           {fill:$\diamond$}{}      % /

\pepgroups{FYW,ILVM,RK,DE,GA,ST,NQ} %  Amino acid grouping due to similarity

\pepsims{F}{YW}                     %  Y and W are similar to F
\pepsims{Y}{WF}                     %  W and F are similar to Y
\pepsims{W}{YF}                     %  Y and F are similar to W

\pepsims{I}{LVM}                    %  L, V and M are similar to I
\pepsims{L}{VMI}                    %  V, M and I are similar to L
\pepsims{V}{MIL}                    %  M, I and L are similar to V

\pepsims{R}{KH}                     %  K and H are similar to R
\pepsims{K}{HR}                     %  H and R are similar to K
\pepsims{H}{RK}                     %  R and K are similar to H

\pepsims{A}{GS}                     %  G and S are similar to A
\pepsims{G}{A}                      %  A (but not S) is similar to G

\pepsims{S}{TA}                     %  T and A are similar to S
\pepsims{T}{S}                      %  S (but not A) is similar to T

\pepsims{D}{EN}                     %  E and N (but not Q) are similar to D
\pepsims{E}{DQ}                     %  D and Q (but not N) are similar to E
\pepsims{N}{QD}                     %  Q and D (but not E) are similar to N
\pepsims{Q}{NE}                     %  N and E (but not D) are similar to Q

\DNAgroups{GAR,CTY}                 %  Nucleotide grouping due to similarity

\DNAsims{A}{GR}                     %  G and R are similar to A
\DNAsims{G}{AR}                     %  A and R are similar to G
\DNAsims{R}{AG}                     %  A and G are similar to R

\DNAsims{C}{TY}                     %  T and Y are similar to C
\DNAsims{T}{CY}                     %  C and Y are similar to T
\DNAsims{Y}{CT}                     %  C and T are similar to Y
%</definitions>
%    \end{macrocode}
%    \begin{macrocode}
%<*AQPDNA>
AQPDNA.MSF  MSF: 979   Type: N  Freitag, 12. Februar 1999 Check: 2594 ..
Name: AQP1nuc.SEQ  Len: 807  Check:  8330    Weight: 1.00
Name: AQP2nuc.SEQ  Len: 813  Check:  7220    Weight: 1.00
Name: AQP3nuc.SEQ  Len: 855  Check:  7590    Weight: 1.00
Name: AQP4nuc.SEQ  Len: 960  Check:  8696    Weight: 1.00
Name: AQP5nuc.SEQ  Len: 795  Check:  758    Weight: 1.00
//
              1                                                      60
AQP1nuc.SEQ  ATGGCCAGCGAAATCAAGAAGAAGC.................TCTTCT........GGAG
AQP2nuc.SEQ  ATGTG....GGAACTCAG.........................ATCCAT...........
AQP3nuc.SEQ  ATG........AACC........GTTGCGGGG.AGATG.....CTCC.............
AQP4nuc.SEQ  ATGAGTGACGGAGCTGCAGCGAGGCGGTGGGGTAAGTGTGGACCTCCCTGCAGCAGAGAG
AQP5nuc.SEQ  ATGAAAAA.GGAGGTGTG.........................CTCCCT...........

              61                                                      120
AQP1nuc.SEQ  GGC..TGTGGTGGCT.....GAGTTCCTGGCCATGA.CCCTCTTCG..............
AQP2nuc.SEQ  ...................................AGCCTTCTCCCGAGCAGTGCTGGCT
AQP3nuc.SEQ  .ACATCC.....GCTACCGG......CTG.........CTTCGCCA....GGCTCTGGCG
AQP4nuc.SEQ  AGCATCATGGTGGCTTTCAAAGGCGTCTGGACTCAAGCCTTCTGGAAGGCGGTCACAGCA
AQP5nuc.SEQ  ...................................TGCCTTCTTCAAGGCGGTGTTCGCA

              121                                                      180
AQP1nuc.SEQ  ....TCTTCATCAGCATCGGTTCTGCCCTA...GGCTT.....CAATTACCCACTGGAGA
AQP2nuc.SEQ  GAGTTCTTGGCCACGCTCCTTTTTGTCTTCTTTGGCCTTGGCTCAGCCCTCCA.....GT
AQP3nuc.SEQ  GAGTGCCTGGGGACCCTCATCCTTGTGATGTTCGGCTGTGGTTCCGTGGCTCAA.GTGGT
AQP4nuc.SEQ  GAGTTCCTGGCCATGCTCATCTTTGTTCTGCTCAGCGTGGGATCCACCATTAACTGGGGT
AQP5nuc.SEQ  GAGTTCCTGGCCACCCTCATCTTCGTCTTCTTTGGCCTGGGCTCAGCACTCAA.....GT

              181                                                      240
AQP1nuc.SEQ  GA...AACCAGACGCTGGTCCA.GGACAATGTGAAGGTGTCACTGGCCTTTGGTCTGAGC
AQP2nuc.SEQ  GGGCCAGCT....CCCCACCCTC...TGTGCTCCAGATCGCCGTGGCCTTTGGTCTGGGC
AQP3nuc.SEQ  GCTCAGCCGAGGGACCCATG.GTGG.CTTCCTCACCATCAACTTGGCTTTTGGCTTCGCT
AQP4nuc.SEQ  GGCTCAGAGAACCCCCTACCTGTGGACATGGTCCTCATCTCCCTCTGCTTTGGACTCAGC
AQP5nuc.SEQ  GGCCCTCGG....CTCTGCCCAC...CATTCTGCAAATCTCAATTGCCTTTGGCCTGGCC

              241                                                      300
AQP1nuc.SEQ  ATCGCTACTCTGGCCCAAAGTGTGGGTCACATCAGTGGTGCTCACTCCAACCCAGCGGTC
AQP2nuc.SEQ  ATCGGCATCCTGGTTCAGGCTCTGGGCCATGTCAGCGGGGCACACATCAACCCCGCCGTG
AQP3nuc.SEQ  GTCACCCTTGCCATCTTGGTGGCTGGCCAAGTGTCTGGAGCCCACTTGAACCCTGCTGTG
AQP4nuc.SEQ  ATTGCCACCATGGTTCAGTGCTTCGGCCACATCAGCGGTGGCCACATCAACCCAGCGGTG
AQP5nuc.SEQ  ATAGGTACCTTAGCCCAAGCTCTGGGACCTGTGAGTGGTGGCCACATCAATCCAGCCATT

              301                                                      360
AQP1nuc.SEQ  ACACTGGGGCTTCTGCTCAGCTGTCAGATCAGCATCCTCCGGGCTGTCA.TGTATATCAT
AQP2nuc.SEQ  ACTGTGGCATGCCTGGTGGGTTGCCATGTCTCCTTCCTTCGAGCTGCCT.TCTATGTGGC
AQP3nuc.SEQ  ACCTTTGCAATG.TGCTTCCTGGCACGAGAGCCCTGGATCAAGCTGCCCATCTACACACT
AQP4nuc.SEQ  ACAGTGGCCATGGTGTGCACACGAAAGATCAGCATCGCCAAGTCTGTCT.TCTACATCAC
AQP5nuc.SEQ  ACTCTGGCCCTCTTAATAGGAAACCAGATCTCGCTGCTCCGAGCTGTCT.TCTACGTGGC

              361                                                      420
AQP1nuc.SEQ  CGCCCAGTGTGTGGGAGCCATCGTTGCCTCCGCCATCCTCTCCGGCATCACCTCCTCCCT
AQP2nuc.SEQ  TGCCCAGCTGCTGGGCGCCGTGGCTGGGGCTGCCATCCTCCATGAGATTAC.TCCAGTAG
AQP3nuc.SEQ  GGCACAGACCCTCGGGGCCTTCTTGGGTGCTGGGATTGTTTTTGGGCT..CTACTA..TG
AQP4nuc.SEQ  TGCGCAGTGCCTGGGGGCCATCATCGGAGCTGGGATCCTCTACCTGGTCAC.ACCCCCCA
AQP5nuc.SEQ  AGCCCAGCTGGTGGGCGCCATTGCTGGGGCAGGCATCCTGTACTGGCTGGC.GCCACTCA

              421                                                      480
AQP1nuc.SEQ  GCTCGAGAACTCACTTGGCCGA.AATGACCTGGCTCGAGGTGTGAACTCCGGCCAGGGCC
AQP2nuc.SEQ  AAATCCGTGGGGACCTGGCTGTCAATGCTCTCCACAACAACGCCACAGCTGGCCAGGCTG
AQP3nuc.SEQ  ATGCAATCTGGGCCTTTGCTGGCAATGAGCT.........TGTTGTCTCCGGCC.....C
AQP4nuc.SEQ  GCGTGGTGGGAGGATTGGGAGTCACCACGGTTCATGGAAACCTCACTGCTGGCCATGGGC
AQP5nuc.SEQ  ATGCCCGGGGTAACCTGGCCGTCAATGCGCTGAACAACAACACAACGCCTGGCAAGGCCA

              481                                                      540
AQP1nuc.SEQ  TGGGCATTGAGATCATTGGCACCCTGCAGCTGGTGCTGTGCGT.TCTGGCTACCACTGAC
AQP2nuc.SEQ  TGACTGTAGAGCTCTTCCTGACCATGCAGCTGGTGCTGTGCAT.CTTTGCCTCCACCGAC
AQP3nuc.SEQ  CAATGGCACAGCTGGTATC..TTTGCCACCTATCCCTCTGGACACTTGGATATGGTCAAT
AQP4nuc.SEQ  TCCTGGTGGAGCTAATAATCACTTTCCAGCTGGTATTCACCAT.TTTTGCCAGCTGTGAT
AQP5nuc.SEQ  TGGTGGTGGAGTTAATCTTGACTTTCCAGCTAGCCCTCTGCAT.CTTCTCCTCCACCGAC

              541                                                      600
AQP1nuc.SEQ  CGGAGGCGCCGAGACTTAGGTGGCTCAGCCCCACTTGCCATTGGCTTGTCTGTGGCTCTT
AQP2nuc.SEQ  GAGCGCCGCGGTGACAACCTGGGTAGCCCTGCCCTCTCCATTGGTTTCTCTGTTACCCTG
AQP3nuc.SEQ  GGCTTCTTTGATCAGTTCATAGGCACAGCAGCCCTTATTGTGTGTGTGCTGGCCATTGTT
AQP4nuc.SEQ  TCCAAACGGACTGATGTTACTGGTTCCGTTGCTTTAGCAATTGGGTTTTCCGTTGCAATT
AQP5nuc.SEQ  TCTCGCCGAACCAGCCCTGTGGGCTCCCCAGCCTTATCCATTGGCTTGTCTGTCACACTG

              601                                                      660
AQP1nuc.SEQ  GGACACCTGCTGGCCATTGACTACACTGGCTGTGGGATCAACCCTGCCCGGTCATT.TGG
AQP2nuc.SEQ  GGCCACCTCCTTGGGATCTATTTCACCGGTTGCTCCATGAATCCAGCCCGCTCCCT.GGC
AQP3nuc.SEQ  GACC..CTTATAACAACCCTGTGCCCCGGGGCCTGGAGGCCTTCACTGTGGGCCTTGTGG
AQP4nuc.SEQ  GGACATTTGTTTGCAATCAATTATACCGGAGCCAGCATGAATCCAGCTCGATCCTT.TGG
AQP5nuc.SEQ  GGCCATCTTGTGGGGATCTACTTCACCGGCTGTTCCATGAACCCAGCCCGATCTTT.CGG

              661                                                      720
AQP1nuc.SEQ  CTCTGCTGTGCTCACCCGCAACTTCTCAAAC...CACTGGATTTTCTGGGTGGGACCATT
AQP2nuc.SEQ  TCCAGCAGTTGTCACTGGCAAGTTTGATGA...TCACTGGGTCTTCTGGATCGGACCCCT
AQP3nuc.SEQ  TCCTG.....GTCATTGGGACCTCCATGGGCTTCAATTCTGGCTATGCCGTCAACCCAGC
AQP4nuc.SEQ  CCCTGCAGTTATCATGGGAAACTGGGAAAAC...CACTGGATATATTGGGTTGGACCAAT
AQP5nuc.SEQ  CCCTGCGGTGGTCATGAACCGGTTCAGCCCCTCTCACTGGGTCTTCTGGGTAGGGCCTAT

              721                                                      780
AQP1nuc.SEQ  CATTGGGAGTGCCCTGGCAGTGCTGATCTATGACTTCATC..CTGGCCCCACGC..AGC.
AQP2nuc.SEQ  GGTGGGCGCCATCATCGGCTCCCTCCTCTACAACTAC..CTGCTGTTC..........CC
AQP3nuc.SEQ  T.....CGTGACTTTGG..ACCTCGCCTTTTCACTGCCCTGGCTGGC......TGGGGTT
AQP4nuc.SEQ  CATAGGCGCTGTGCTGGCAGGTGCACTTTACGAGTATGTCTTCTGTCCTGACGTGGAGCT
AQP5nuc.SEQ  TGTGGGGGCCATGCTGGCGGCCATCCTCTATTTCTAC..CTGCTCTTC..........CC

              781                                                      840
AQP1nuc.SEQ  ..AGCG.........................ACTTTACAG.............ACCGCAT
AQP2nuc.SEQ  C.....TCGGCAAAG...AGCCTGCAGGAGCGCTTGGCAGTGCTCAAGGG.......CCT
AQP3nuc.SEQ  CAGAAGTC.TTTACGACTGGCC...AGAACTGGTGGTGGGTACCCATCGTCTCTCCACTC
AQP4nuc.SEQ  CAAACGTCGCCTAAAGGAAGCCTTCAGCAAAGCTGCACAGCAGACGAAAGGGAGCTACAT
AQP5nuc.SEQ  C.....TCCTCTCTG...AGCCTCCATGATCGCGTGGCTGTCGTCAAAGG.......CAC

              841                                                      900
AQP1nuc.SEQ  GAAGGTGTGGACCAGT...GGCCAAGTGGA.....GGAGTATGACCTGGATGC.......
AQP2nuc.SEQ  GGAGCCCGACACCGACTGGGA.......GGAACGTGAAGTGCGG..CGGCGGCAGTCGGT
AQP3nuc.SEQ  CTGGGTTC.CATTGGTGGTGTCTTCGTGT.ACCAGCT..CATGAT.TGGCTGCCACC..T
AQP4nuc.SEQ  GGAGGTGGAGGACAACCGGAGCCAAGTGGAGACAGAAGACTTGATCCTGAAGCCCGGGGT
AQP5nuc.SEQ  ATA...TGA.GCCGG..AGGA.......GGACTGGGAAGATCAT..CGAGAGGAGAGGAA

              901                                                      960
AQP1nuc.SEQ  ........TGAT.GATATCAACTCCAGGGTGGAGATGAAG....................
AQP2nuc.SEQ  GGAGC......TC..CACTCTCCTCAGAG...................CCTGCCTCGCG.
AQP3nuc.SEQ  GGAGCA.GCCCCCGCCTTCCACT..GAGGCAGAGAATGTGAAGCTGG.CCCACATGAAGC
AQP4nuc.SEQ  GGTGCATGTGATCGACATTGACCGTGGAGACGAGAAGAAGGGGAAGGACTCGTCTGGAGA
AQP5nuc.SEQ  GAAG............ACCATC....GAG........................CTGACG.

              961                    979
AQP1nuc.SEQ  ..........CCCAAATAG
AQP2nuc.SEQ  .GCAGCAAGGCCTG....A
AQP3nuc.SEQ  ACAAGGA..GCAGATCTGA
AQP4nuc.SEQ  GGTATTATCTTCTGTATGA
AQP5nuc.SEQ  .GCA.CA....CTG....A

%</AQPDNA>
%    \end{macrocode}
%    \begin{macrocode}
%<*AQPpro>
AQPpro.MSF  MSF: 356   Type: P  Freitag, 12. Februar 1999 Check: 2586 ..
Name: AQP1.PRO  Len: 269  Check:  5367    Weight: 1.00
Name: AQP2.PRO  Len: 271  Check:  6176    Weight: 1.00
Name: AQP3.PRO  Len: 285  Check:  2893    Weight: 1.00
Name: AQP4.PRO  Len: 323  Check:  9737    Weight: 1.00
Name: AQP5.PRO  Len: 265  Check:  8413    Weight: 1.00
//
          1                                                      60
AQP1.PRO  MAS........................EIKKKLFWRAVVAEFLAMTLFVFISIGSALGFN
AQP2.PRO  MW.........................ELRSIAFSRAVLAEFLATLLFVFFGLGSALQWA
AQP3.PRO  M.........NRCG.....EMLHIRYR......LLRQALAECLGTLILVMFGCGSVAQVV
AQP4.PRO  MSDGAAARRWGKCGPPCSRESIMVAFKGVWTQAFWKAVTAEFLAMLIFVLLSVGSTINWG
AQP5.PRO  MK........................KEVCSLAFFKAVFAEFLATLIFVFFGLGSALKWP

          61                                                      120
AQP1.PRO  YPLERNQTLVQDNVKVSLAFGLSIATLAQSVGHISGAHSNPAVTLGLLLSCQISILRAVM
AQP2.PRO  ...SS....PPSVLQIAVAFGLGIGILVQALGHVSGAHINPAVTVACLVGCHVSFLRAAF
AQP3.PRO  LSRGTH....GGFLTINLAFGFAVTLAILVAGQVSGAHLNPAVTFAMCFLAREPWIKLPI
AQP4.PRO  ...GSENPLPVDMVLISLCFGLSIATMVQCFGHISGGHINPAVTVAMVCTRKISIAKSVF
AQP5.PRO  ...SA....LPTILQISIAFGLAIGTLAQALGPVSGGHINPAITLALLIGNQISLLRAVF

          121                                                      180
AQP1.PRO  YIIAQCVGAIVASAILSGI..........TSSLLENSLGRNDLARGVNSGQ.....GLGI
AQP2.PRO  YVAAQLLGAVAGAAILHEI..........TPVEIRGDLAVNALHNNATAGQ.....AVTV
AQP3.PRO  YTLAQTLGAFLGAGIVFGLYYDAIWAFAGNELVVSGPNGTAGIFATYPSGHLDMVNGFFD
AQP4.PRO  YITAQCLGAIIGAGILYLV..........TPPSVVGGLGVTTVHGNLTAGH.....GLLV
AQP5.PRO  YVAAQLVGAIAGAGILYWL..........APLNARGNLAVNALNNNTTPGK.....AMVV

          181                                                      240
AQP1.PRO  EIIGTLQLVLCVLATTDR.RRRDLGGSAPLAIGLSVALGHLLAIDYTGCGINPARSFGSA
AQP2.PRO  ELFLTMQLVLCIFASTDE.RRGDNLGSPALSIGFSVTLGHLLGIYFTGCSMNPARSLAPA
AQP3.PRO  QFIGTAALIVCVLAIVDPYNNPVPRGLEAFTVGLVVLVIGTSMGFNSGYAVNPARDFGPR
AQP4.PRO  ELIITFQLVFTIFASCDS.KRTDVTGSVALAIGFSVAIGHLFAINYTGASMNPARSFGPA
AQP5.PRO  ELILTFQLALCIFSSTDS.RRTSPVGSPALSIGLSVTLGHLVGIYFTGCSMNPARSFGPA

          241                                                      300
AQP1.PRO  VLTR..NFS.N......HWIFWVGPFIGSALAVL..IYDFILAPRSSDFTDRMK......
AQP2.PRO  VVTG..KFD.D......HWVFWIGPLVGAIIGSL..LYNYLLFPSAKSLQERL..AVLKG
AQP3.PRO  LFTALAGWGSEVFTTGQNW..WWVPIVSPLLGSIGGVFVYQL..................
AQP4.PRO  VIMG..NWE.N......HWIYWVGPIIGAVLAGA..LYEYV.FCPDVELKRRLKEAFSKA
AQP5.PRO  VVMN..RFSPS......HWVFWVGPIVGAMLAAI..LYFYLLFPSSLSLHDRV..AVVKG

          301                                                      354
AQP1.PRO  ......VWTS.....GQVEEYDLDAD.......DINSRVEMKPK..........
AQP2.PRO  .LEPDTDWEEREVRRRQ..SVELHSPQSLPRG...................SKA
AQP3.PRO  .................MIGCHLEQPPPSTEAENV.KLAHMKHKE.......QI
AQP4.PRO  AQQTKGSYMEVEDNRSQVETEDLILKPGVVHVIDIDRGDEKKGKDSSGEVLSSV
AQP5.PRO  TYEPEEDWEDHREERKK..TIELTAH............................

%</AQPpro>
%    \end{macrocode}
%    \begin{macrocode}
%<*AQP2spec>
AQP2bt  SIAFSRAVLAEFLATLLFVFFGLGSALNWPQALPSVLQIAMAFGLAIGTLVQALGHVSGA
AQP2cf  SVAFSRAVFAEFLATLLFVFFGLGSALNWPQALPSVLQIAMAFGLGIGTLVQALGHVSGA
AQP2dd  SIAFSRAVFSEFLATLLFVFFGLGSALNWPQALPSVLQIAMAFGLAIGTLVQALGHISGA
AQP2ec  SIAFSRAVLAEFLATLLFVFFGLGSALNWPQAMPSVLQIAMAFGLAIGTLVQALGHVSGA
AQP2em  SIAFSRAVFSEFLATLLFVFFGLGSALNWPQALPSVLQIAMAFGLAIGTLVQTLGHISGA


AQP2bt  HINPAVTVACLVGCHVSFLRAVFYVAAQLLGAVAGAALLHEITPPAIRG
AQP2cf  HINPAVTVACLVGCHVSFLRAAFYVAAQLLGAVAGAALLHEITPPHVRG
AQP2dd  HINPAVTVACLVGCHVSFLRATFYLAAQLLGAVAGAAILHEITPPDIRG
AQP2ec  HINPAVTVACLVGCHVSFLRAAFYVAAQLLGAVAGAALLHEITPPDIRR
AQP2em  HINPAVTVACLVGCHVSFLRATFYLAAQLLGAVAGAALLHELTPPDIRG

%</AQP2spec>
%    \end{macrocode}
%    \begin{macrocode}
%<*AQP1topo>    
\feature{bottom}{1}{1..14}{'-'}{int.\ A}
\feature{top}{1}{15..32}{box[LightGray]:TM1}{}
\feature{top}{1}{33..49}{,-,}{ext.\ B}
\feature{top}{1}{50..68}{box[LightGray]:TM2}{}
\feature{bottom}{1}{69..81}{'-'}{int.\ C}
\feature{top}{1}{82..106}{box[LightGray]:TM3}{}
\feature{top}{1}{107..136}{,-,}{ext.\ D}
\feature{top}{1}{137..154}{box[LightGray]:TM4}{}
\feature{bottom}{1}{155..168}{'-'}{int.\ E}
\feature{top}{1}{169..186}{box[LightGray]:TM5}{}
\feature{top}{1}{187..211}{,-,}{ext.\ F}
\feature{top}{1}{212..230}{box[LightGray]:TM6}{}
\feature{bottom}{1}{231..269}{'-'}{int.\ G}
%</AQP1topo>
%    \end{macrocode}
%    \begin{macrocode}
%<*AQP1PHD>

From phd@EMBL-Heidelberg.de Wed Nov 25 10:24:25 1998
Date: Tue, 24 Nov 1998 17:45:25 +0100
From: Protein Prediction <phd@EMBL-Heidelberg.de>
To: eric.beitz@uni-tuebingen.de
Subject: PredictProtein




The following information has been received by the server:
~~~~~~~~~~~~~~~~~~~~~~~~~~~~~~~~~~~~~~~~~~~~~~~~~~~~~~~~~

________________________________________________________________________________

reference predict_h25873 (Tue Nov 24 17:43:21 MET 1998)
from eric.beitz@uni-tuebingen.de
password(###)
resp MAIL
orig HTML
prediction of: -secondary structure   (PHDsec)-solvent accessibility (PHDacc)-
return msf format
# no description
MASEIKKKLFWRAVVAEFLAMTLFVFISIGSALGFNYPLERNQTLVQDNVKVSLAFGLSIATLAQSVGHISGAHSNPAVT
LGLLLSCQISILRAVMYIIAQCVGAIVASAILSGITSSLLENSLGRNDLARGVNSGQGLGIEIIGTLQLVLCVLATTDRR
RRDLGGSAPLAIGLSVALGHLLAIDYTGCGINPARSFGSAVLTRNFSNHWIFWVGPFIGSALAVLIYDFILAPRSSDFTD
RMKVWTSGQVEEYDLDADDINSRVEMKPK

________________________________________________________________________________





Result of PROSITE search (Amos Bairoch):    
~~~~~~~~~~~~~~~~~~~~~~~~~~~~~~~~~~~~~~~

please quote: A Bairoch, P Bucher & K Hofmann: The PROSITE database,
its status in 1997. Nucl. Acids Res., 1997, 25, 217-221.

________________________________________________________________________________


--------------------------------------------------------

--------------------------------------------------------

Pattern-ID: ASN_GLYCOSYLATION PS00001 PDOC00001
Pattern-DE: N-glycosylation site
Pattern:    N[^P][ST][^P]
   42       NQTL
   250      NFSN

Pattern-ID: GLYCOSAMINOGLYCAN PS00002 PDOC00002
Pattern-DE: Glycosaminoglycan attachment site
Pattern:    SG.G
   135      SGQG

Pattern-ID: PKC_PHOSPHO_SITE PS00005 PDOC00005
Pattern-DE: Protein kinase C phosphorylation site
Pattern:    [ST].[RK]
   157      TDR
   398      TDR

Pattern-ID: CK2_PHOSPHO_SITE PS00006 PDOC00006
Pattern-DE: Casein kinase II phosphorylation site
Pattern:    [ST].{2}[DE]
   118      SLLE
   383      SRVE

Pattern-ID: MYRISTYL PS00008 PDOC00008
Pattern-DE: N-myristoylation site
Pattern:    G[^EDRKHPFYW].{2}[STAGCN][^P]
   30       GSALGF
   92       GLSIAT
   179      GLLLSC
   288      GAIVAS
   407      GITSSL
   544      GVNSGQ
   722      GLSVAL
   917      GINPAR
   1141     GSALAV

Pattern-ID: PROKAR_LIPOPROTEIN PS00013 PDOC00013
Pattern-DE: Prokaryotic membrane lipoprotein lipid attachment site
Pattern:    [^DERK]{6}[LIVMFWSTAG]{2}[LIVMFYSTAGCQ][AGS]C
   77       PAVTLGLLLSC

Pattern-ID: MIP PS00221 PDOC00193
Pattern-DE: MIP family signature
Pattern:    [HNQA].NP[STA][LIVMF][ST][LIVMF][GSTAFY]
   74       HSNPAVTLG




________________________________________________________________________________





Result of ProDom domain search (Corpet, Gouzy, Kahn):
~~~~~~~~~~~~~~~~~~~~~~~~~~~~~~~~~~~~~~~~~~~~~~~~~~~~

- please quote: ELL Sonnhammer & D Kahn, Prot. Sci., 1994, 3, 482-492

________________________________________________________________________________


--- ------------------------------------------------------------
--- Results from running BLAST against PRODOM domains
---
--- PLEASE quote:
---       F Corpet, J Gouzy, D Kahn (1998).  The ProDom database
---       of protein domain families. Nucleic Ac Res 26:323-326.
---
--- BEGIN of BLASTP output
BLASTP 1.4.7 [16-Oct-94] [Build 17:06:52 Oct 31 1994]

Reference:  Altschul, Stephen F., Warren Gish, Webb Miller, Eugene W. Myers,
and David J. Lipman (1990).  Basic local alignment search tool.  J. Mol. Biol.
215:403-10.

Query=  prot (#) ppOld, no description /home/phd/server/work/predict_h25873
        (269 letters)

Database:  /home/phd/ut/prodom/prodom_34_2
           53,597 sequences; 6,740,067 total letters.
Searching..................................................done

                                                                     Smallest
                                                                       Sum
                                                              High  Probability
Sequences producing High-scoring Segment Pairs:              Score  P(N)      N

 390 p34.2 (45) MIP(6) AQP1(4) GLPF(4)  // PROTEIN INTRIN...   270  2.0e-32   1
 45663 p34.2 (1) AQPZ_ECOLI // AQUAPORIN Z.                     90  3.2e-13   2
 45611 p34.2 (1) AQP2_HUMAN // AQUAPORIN-CD (AQP-CD) (WAT...   136  6.0e-13   1
 304 p34.2 (61) AQP2(10) GLPF(6) MIP(5)  // PROTEIN CHANN...   121  9.2e-11   1
 45607 p34.2 (1) PMIP_NICAL // POLLEN-SPECIFIC MEMBRANE I...    80  1.2e-07   2
 45606 p34.2 (1) BIB_DROME // NEUROGENIC PROTEIN BIG BRAIN.     80  1.2e-05   2
 2027 p34.2 (15) GLPF(9) AQP3(2)  // PROTEIN FACILITATOR ...    60  3.4e-05   2
 45615 p34.2 (1) GLPF_STRPN // GLYCEROL UPTAKE FACILITATO...    63  0.024     1
 45638 p34.2 (1) AQP5_HUMAN // AQUAPORIN 5.                     61  0.044     1



>390 p34.2 (45) MIP(6) AQP1(4) GLPF(4)  // PROTEIN INTRINSIC CHANNEL WATER
  AQUAPORIN TONOPLAST MEMBRANE FOR PLASMA LENS
  Length = 88

 Score = 270 (125.3 bits), Expect = 2.0e-32, P = 2.0e-32
 Identities = 47/67 (70%), Positives = 56/67 (83%)

Query:   156 TTDRRRRDLGGSAPLAIGLSVALGHLLAIDYTGCGINPARSFGSAVLTRNFSNHWIFWVG 215
             T D+RR  +GGSAPL IG SVALGHL+ I YTGCG+NPARSFG AV+T NF+NHW++WVG
Sbjct:    22 TDDKRRGSVGGSAPLPIGFSVALGHLIGIPYTGCGMNPARSFGPAVVTGNFTNHWVYWVG 81

Query:   216 PFIGSAL 222
             P IG+ L
Sbjct:    82 PIIGAVL 88

 Score = 95 (44.1 bits), Expect = 2.3e-06, P = 2.3e-06
 Identities = 20/33 (60%), Positives = 23/33 (69%)

Query:   136 GQGLGIEIIGTLQLVLCVLATTDRRRRDLGGSA 168
             GQ L +EIIGT QLV CV ATTD +RR   G +
Sbjct:     1 GQNLVVEIIGTFQLVYCVFATTDDKRRGSVGGS 33


>45663 p34.2 (1) AQPZ_ECOLI // AQUAPORIN Z.
  Length = 96

 Score = 90 (41.8 bits), Expect = 3.2e-13, Sum P(2) = 3.2e-13
 Identities = 18/36 (50%), Positives = 25/36 (69%)

Query:   166 GSAPLAIGLSVALGHLLAIDYTGCGINPARSFGSAV 201
             G AP+AIGL++ L HL++I  T   +NPARS   A+
Sbjct:    25 GFAPIAIGLALTLIHLISIPVTNTSVNPARSTAVAI 60

 Score = 63 (29.2 bits), Expect = 3.2e-13, Sum P(2) = 3.2e-13
 Identities = 11/25 (44%), Positives = 14/25 (56%)

Query:   210 WIFWVGPFIGSALAVLIYDFILAPR 234
             W FWV P +G  +  LIY  +L  R
Sbjct:    71 WFFWVVPIVGGIIGGLIYRTLLEKR 95


>45611 p34.2 (1) AQP2_HUMAN // AQUAPORIN-CD (AQP-CD) (WATER CHANNEL PROTEIN FOR
  RENAL COLLECTING DUCT) (ADH WATER CHANNEL) (AQUAPORIN 2) (COLLECTING DUCT
  WATER CHANNEL PROTEIN) (WCH-CD).
  Length = 49

 Score = 136 (63.1 bits), Expect = 6.0e-13, P = 6.0e-13
 Identities = 23/42 (54%), Positives = 34/42 (80%)

Query:    50 VKVSLAFGLSIATLAQSVGHISGAHSNPAVTLGLLLSCQISI 91
             +++++AFGL I TL Q++GHISGAH NPAVT+  L+ C +S+
Sbjct:     8 LQIAMAFGLGIGTLVQALGHISGAHINPAVTVACLVGCHVSV 49


>304 p34.2 (61) AQP2(10) GLPF(6) MIP(5)  // PROTEIN CHANNEL WATER AQUAPORIN
  INTRINSIC DUCT COLLECTING FOR TONOPLAST WCH-CD
  Length = 43

 Score = 121 (56.1 bits), Expect = 9.2e-11, P = 9.2e-11
 Identities = 24/43 (55%), Positives = 31/43 (72%)

Query:    70 ISGAHSNPAVTLGLLLSCQISILRAVMYIIAQCVGAIVASAIL 112
             ISG H NPAVT+GLL+  +   LRAV YI AQ +GA+  +A+L
Sbjct:     1 ISGGHINPAVTIGLLIGGRFPFLRAVFYIAAQLLGAVAGAALL 43


>45607 p34.2 (1) PMIP_NICAL // POLLEN-SPECIFIC MEMBRANE INTEGRAL PROTEIN.
  Length = 69

 Score = 80 (37.1 bits), Expect = 1.2e-07, Sum P(2) = 1.2e-07
 Identities = 17/54 (31%), Positives = 32/54 (59%)

Query:   149 LVLCVLATTDRRRRDLGGSAPLAIGLSVALGHLLAIDYTGCGINPARSFGSAVL 202
             L++ V++      R +G  A +A+G+++ L   +A   +G  +NPARS G A++
Sbjct:    13 LLMFVISGVATDDRAIGQVAGIAVGMTITLNVFVAGPISGASMNPARSIGPAIV 66

 Score = 34 (15.8 bits), Expect = 1.2e-07, Sum P(2) = 1.2e-07
 Identities = 8/18 (44%), Positives = 11/18 (61%)

Query:   136 GQGLGIEIIGTLQLVLCV 153
             GQ L IEII +  L+  +
Sbjct:     1 GQSLAIEIIISFLLMFVI 18


>45606 p34.2 (1) BIB_DROME // NEUROGENIC PROTEIN BIG BRAIN.
  Length = 119

 Score = 80 (37.1 bits), Expect = 1.2e-05, Sum P(2) = 1.2e-05
 Identities = 15/34 (44%), Positives = 24/34 (70%)

Query:     1 MASEIKKKLFWRAVVAEFLAMTLFVFISIGSALG 34
             M +EI+   FWR++++E LA  ++VFI  G+A G
Sbjct:    55 MQAEIRTLEFWRSIISECLASFMYVFIVCGAAAG 88

 Score = 39 (18.1 bits), Expect = 1.2e-05, Sum P(2) = 1.2e-05
 Identities = 9/17 (52%), Positives = 12/17 (70%)

Query:    53 SLAFGLSIATLAQSVGH 69
             +LA GL++ATL Q   H
Sbjct:   103 ALASGLAMATLTQCFLH 119


>2027 p34.2 (15) GLPF(9) AQP3(2)  // PROTEIN FACILITATOR GLYCEROL UPTAKE
  AQUAPORIN DIFFUSION UPTAKE/EFFLUX PEPX 5'REGION ORF1
  Length = 55

 Score = 60 (27.8 bits), Expect = 3.4e-05, Sum P(2) = 3.4e-05
 Identities = 17/46 (36%), Positives = 20/46 (43%)

Query:   156 TTDRRRRDLGGSAPLAIGLSVALGHLLAIDYTGCGINPARSFGSAV 201
             T D      GG  PL +G  V    +     TG  INPAR FG  +
Sbjct:    10 TDDGNNVPSGGLHPLMVGFLVMGIGMSLGGTTGYAINPARDFGPRI 55

 Score = 37 (17.2 bits), Expect = 3.4e-05, Sum P(2) = 3.4e-05
 Identities = 7/10 (70%), Positives = 8/10 (80%)

Query:   149 LVLCVLATTD 158
             L+ CVLA TD
Sbjct:     2 LIACVLALTD 11


>45615 p34.2 (1) GLPF_STRPN // GLYCEROL UPTAKE FACILITATOR PROTEIN.
  Length = 26

 Score = 63 (29.2 bits), Expect = 0.025, P = 0.024
 Identities = 13/23 (56%), Positives = 18/23 (78%)

Query:   205 NFSNHWIFWVGPFIGSALAVLIY 227
             ++S  WI  VGP IG+ALAVL++
Sbjct:     1 DWSYAWIPVVGPVIGAALAVLVF 23


>45638 p34.2 (1) AQP5_HUMAN // AQUAPORIN 5.
  Length = 27

 Score = 61 (28.3 bits), Expect = 0.045, P = 0.044
 Identities = 11/19 (57%), Positives = 18/19 (94%)

Query:    50 VKVSLAFGLSIATLAQSVG 68
             ++++LAFGL+I TLAQ++G
Sbjct:     8 LQIALAFGLAIGTLAQALG 26


Parameters:
  E=0.1
  B=500

  V=500
  -ctxfactor=1.00

  Query                        -----  As Used  -----    -----  Computed  ----
  Frame  MatID Matrix name     Lambda    K       H      Lambda    K       H
   +0      0   BLOSUM62        0.322   0.138   0.394    same    same    same

  Query
  Frame  MatID  Length  Eff.Length   E    S W   T  X     E2  S2
   +0      0      269       269     0.10 69 3  11 22    0.22 33


Statistics:
  Query          Expected         Observed           HSPs       HSPs
  Frame  MatID  High Score       High Score       Reportable  Reported
   +0      0    59 (27.4 bits)  270 (125.3 bits)       14         14

  Query         Neighborhd  Word      Excluded    Failed   Successful  Overlaps
  Frame  MatID   Words      Hits        Hits    Extensions Extensions  Excluded
   +0      0      5349     3124825      609708     2510548     4569         2

  Database:  /home/phd/ut/prodom/prodom_34_2
    Release date:  unknown
    Posted date:  12:24 PM MET DST May 06, 1998
  # of letters in database:  6,740,067
  # of sequences in database:  53,597
  # of database sequences satisfying E:  9
  No. of states in DFA:  564 (111 KB)
  Total size of DFA:  226 KB (256 KB)
  Time to generate neighborhood:  0.03u 0.00s 0.03t  Real: 00:00:00
  Time to search database:  9.80u 0.03s 9.83t  Real: 00:00:10
  Total cpu time:  9.90u 0.06s 9.96t  Real: 00:00:10
--- END of BLASTP output
--- ------------------------------------------------------------
---
--- Again: these results were obtained based on the domain data-
--- base collected by Daniel Kahn and his coworkers in Toulouse.
---
--- PLEASE quote:
---       F Corpet, J Gouzy, D Kahn (1998).  The ProDom database
---       of protein domain families. Nucleic Ac Res 26:323-326.
---
--- The general WWW page is on:
----      ---------------------------------------
---       http://www.toulouse.inra.fr/prodom.html
----      ---------------------------------------
---
--- For WWW graphic interfaces to PRODOM, in particular for your
--- protein family, follow the following links (each line is ONE
--- single link for your protein!!):
---
http://www.toulouse.inra.fr/prodom/cgi-bin/ReqProdomII.pl?id_dom1=390 ==> multiple alignment, consensus, PDB and PROSITE links of domain 390
http://www.toulouse.inra.fr/prodom/cgi-bin/ReqProdomII.pl?id_dom2=390 ==> graphical output of all proteins having domain 390
http://www.toulouse.inra.fr/prodom/cgi-bin/ReqProdomII.pl?id_dom1=45663 ==> multiple alignment, consensus, PDB and PROSITE links of domain 45663
http://www.toulouse.inra.fr/prodom/cgi-bin/ReqProdomII.pl?id_dom2=45663 ==> graphical output of all proteins having domain 45663
http://www.toulouse.inra.fr/prodom/cgi-bin/ReqProdomII.pl?id_dom1=45611 ==> multiple alignment, consensus, PDB and PROSITE links of domain 45611
http://www.toulouse.inra.fr/prodom/cgi-bin/ReqProdomII.pl?id_dom2=45611 ==> graphical output of all proteins having domain 45611
http://www.toulouse.inra.fr/prodom/cgi-bin/ReqProdomII.pl?id_dom1=304 ==> multiple alignment, consensus, PDB and PROSITE links of domain 304
http://www.toulouse.inra.fr/prodom/cgi-bin/ReqProdomII.pl?id_dom2=304 ==> graphical output of all proteins having domain 304
http://www.toulouse.inra.fr/prodom/cgi-bin/ReqProdomII.pl?id_dom1=45607 ==> multiple alignment, consensus, PDB and PROSITE links of domain 45607
http://www.toulouse.inra.fr/prodom/cgi-bin/ReqProdomII.pl?id_dom2=45607 ==> graphical output of all proteins having domain 45607
http://www.toulouse.inra.fr/prodom/cgi-bin/ReqProdomII.pl?id_dom1=45606 ==> multiple alignment, consensus, PDB and PROSITE links of domain 45606
http://www.toulouse.inra.fr/prodom/cgi-bin/ReqProdomII.pl?id_dom2=45606 ==> graphical output of all proteins having domain 45606
http://www.toulouse.inra.fr/prodom/cgi-bin/ReqProdomII.pl?id_dom1=2027 ==> multiple alignment, consensus, PDB and PROSITE links of domain 2027
http://www.toulouse.inra.fr/prodom/cgi-bin/ReqProdomII.pl?id_dom2=2027 ==> graphical output of all proteins having domain 2027
http://www.toulouse.inra.fr/prodom/cgi-bin/ReqProdomII.pl?id_dom1=45615 ==> multiple alignment, consensus, PDB and PROSITE links of domain 45615
http://www.toulouse.inra.fr/prodom/cgi-bin/ReqProdomII.pl?id_dom2=45615 ==> graphical output of all proteins having domain 45615
http://www.toulouse.inra.fr/prodom/cgi-bin/ReqProdomII.pl?id_dom1=45638 ==> multiple alignment, consensus, PDB and PROSITE links of domain 45638
http://www.toulouse.inra.fr/prodom/cgi-bin/ReqProdomII.pl?id_dom2=45638 ==> graphical output of all proteins having domain 45638
---
--- NOTE: if you want to use the link, make sure the entire line
---       is pasted as URL into your browser!
---
--- END of PRODOM
--- ------------------------------------------------------------

________________________________________________________________________________


--- Database used for sequence comparison:
--- SEQBASE    RELEASE 34.0 OF EMBL/SWISS-PROT WITH  59021 SEQUENCES




The alignment that has been used as input to the network is:
~~~~~~~~~~~~~~~~~~~~~~~~~~~~~~~~~~~~~~~~~~~~~~~~~~~~~~~~~~~

________________________________________________________________________________

--- ------------------------------------------------------------
--- MAXHOM multiple sequence alignment
--- ------------------------------------------------------------
---
--- MAXHOM ALIGNMENT HEADER: ABBREVIATIONS FOR SUMMARY
--- ID           : identifier of aligned (homologous) protein
--- STRID        : PDB identifier (only for known structures)
--- PIDE         : percentage of pairwise sequence identity
--- WSIM         : percentage of weighted similarity
--- LALI         : number of residues aligned
--- NGAP         : number of insertions and deletions (indels)
--- LGAP         : number of residues in all indels
--- LSEQ2        : length of aligned sequence
--- ACCNUM       : SwissProt accession number
--- NAME         : one-line description of aligned protein
---
--- MAXHOM ALIGNMENT HEADER: SUMMARY
ID         STRID  IDE WSIM LALI NGAP LGAP LEN2 ACCNUM NAME
aqp1_rat          100  100  269    0    0  269 P29975 PROXIMAL TUBULE) (AQUAPOR
aqp1_mouse         98   99  269    0    0  269 Q02013 PROXIMAL TUBULE) (AQUAPOR
aqp1_human         93   97  269    0    0  269 P29972 PROXIMAL TUBULE) (AQUAPOR
aqp1_bovin         90   95  269    1    2  271 P47865 PROXIMAL TUBULE) (AQUAPOR
aqp1_sheep         90   94  269    2    3  272 P56401 PROXIMAL TUBULE) (AQUAPOR
aqpa_ranes         78   89  268    2    5  272 P50501 AQUAPORIN FA-CHIP.
aqp2_dasno         49   73  109    1    7  109 P79164 PROTEIN) (WCH-CD) (FRAGME
aqp2_bovin         49   73  109    1    7  109 P79099 PROTEIN) (WCH-CD) (FRAGME
aqp2_canfa         48   72  109    1    7  109 P79144 PROTEIN) (WCH-CD) (FRAGME
aqp2_rabit         48   73  109    1    7  109 P79213 PROTEIN) (WCH-CD) (FRAGME
aqp2_elema         47   72  109    1    7  109 P79168 PROTEIN) (WCH-CD) (FRAGME
aqp2_horse         47   72  109    1    7  109 P79165 PROTEIN) (WCH-CD) (FRAGME
aqp2_proha         47   73  109    1    7  109 P79229 PROTEIN) (WCH-CD) (FRAGME
mip_rat            46   73  259    1    7  261 P09011 LENS FIBER MAJOR INTRINSI
aqp2_oryaf         46   72  109    1    7  109 P79200 PROTEIN) (WCH-CD) (FRAGME
mip_mouse          46   73  261    1    7  263 P51180 LENS FIBER MAJOR INTRINSI
mip_ranpi          45   73  261    1    7  263 Q06019 LENS FIBER MAJOR INTRINSI
mip_bovin          45   73  261    1    7  263 P06624 LENS FIBER MAJOR INTRINSI
mip_human          45   73  261    1    7  263 P30301 LENS FIBER MAJOR INTRINSI
mip_chick          45   72  110    1    1  112 P28238 LENS FIBER MAJOR INTRINSI
aqp5_rat           44   71  262    2    8  265 P47864 AQUAPORIN 5.
aqp5_human         44   71  262    2    8  265 P55064 AQUAPORIN 5.
aqp2_human         44   72  261    2    8  271 P41181 PROTEIN) (WCH-CD).
aqp4_human         43   70  266    2    5  323 P55087 AQUAPORIN 4 (WCH4) (MERCU
aqp4_rat           43   70  266    2    5  323 P47863 AQUAPORIN 4 (WCH4) (MERCU
aqp4_mouse         43   69  265    3    6  322 P55088 AQUAPORIN 4 (WCH4) (MERCU
aqp2_rat           42   71  261    2    8  271 P34080 PROTEIN) (WCH-CD).
aqp2_mouse         42   71  261    2    8  271 P56402 PROTEIN) (WCH-CD).
wc2a_arath         42   67  248    4   12  287 P43286 PLASMA MEMBRANE INTRINSIC
aqp6_human         42   68  260    2    9  282 Q13520 AQUAPORIN 6 (AQUAPORIN-2
wc2c_arath         41   66  248    4   12  285 P30302 INTRINSIC PROTEIN) (WSI-T
wc2b_arath         41   66  248    4   12  285 P43287 PLASMA MEMBRANE INTRINSIC
wc1c_arath         41   65  238    4   10  286 Q08733 (TMP-B).
wc1b_arath         41   65  238    4   10  286 Q06611 (TMP-A).
tipw_lyces         40   65  237    4   10  286 Q08451 (RIPENING-ASSOCIATED MEMB
wc1a_arath         40   64  238    4   10  286 P43285 PLASMA MEMBRANE INTRINSIC
tipw_pea           40   64  237    4   11  289 P25794 RESPONSIVE PROTEIN 7A).
tipa_arath         38   64  250    3    9  268 P26587 TONOPLAST INTRINSIC PROTE
aqua_atrca         38   64  246    4   10  282 P42767 AQUAPORIN.
dip_antma          38   65  242    2    4  250 P33560 PROBABLE TONOPLAST INTRIN
aqpz_ecoli         37   59  220    4   17  231 P48838 AQUAPORIN Z (BACTERIAL NO
tip2_tobac         37   64  242    2    4  250 P24422 TONOPLAST INTRINSIC PROTE
tip1_tobac         37   64  242    2    4  250 P21653 TONOPLAST INTRINSIC PROTE
tipg_arath         33   62  241    2    4  251 P25818 TONOPLAST INTRINSIC PROTE
bib_drome          33   60  260    4   10  700 P23645 NEUROGENIC PROTEIN BIG BR
tipr_arath         33   62  243    2    4  253 P21652 TONOPLAST INTRINSIC PROTE
tipa_phavu         33   62  246    2    4  256 P23958 TONOPLAST INTRINSIC PROTE
tipg_orysa         32   62  240    2    5  250 P50156 TONOPLAST INTRINSIC PROTE
---
--- MAXHOM ALIGNMENT: IN MSF FORMAT
MSF of: /home/phd/server/work/predict_h25873-22040.hssp from:    1 to:  269
 /home/phd/server/work/predict_h25873-22040.msfRet  MSF:  269  Type: P 24-Nov-98  17:44:5  Check: 3448  ..


 Name: predict_h258    Len:   269  Check: 8331  Weight:  1.00
 Name: aqp1_rat        Len:   269  Check: 8331  Weight:  1.00
 Name: aqp1_mouse      Len:   269  Check: 7552  Weight:  1.00
 Name: aqp1_human      Len:   269  Check: 6501  Weight:  1.00
 Name: aqp1_bovin      Len:   269  Check: 7067  Weight:  1.00
 Name: aqp1_sheep      Len:   269  Check: 7582  Weight:  1.00
 Name: aqpa_ranes      Len:   269  Check: 4844  Weight:  1.00
 Name: aqp2_dasno      Len:   269  Check: 8933  Weight:  1.00
 Name: aqp2_bovin      Len:   269  Check: 9649  Weight:  1.00
 Name: aqp2_canfa      Len:   269  Check: 8990  Weight:  1.00
 Name: aqp2_rabit      Len:   269  Check: 8787  Weight:  1.00
 Name: aqp2_elema      Len:   269  Check: 9381  Weight:  1.00
 Name: aqp2_horse      Len:   269  Check: 8993  Weight:  1.00
 Name: aqp2_proha      Len:   269  Check: 8855  Weight:  1.00
 Name: mip_rat         Len:   269  Check: 9773  Weight:  1.00
 Name: aqp2_oryaf      Len:   269  Check: 8554  Weight:  1.00
 Name: mip_mouse       Len:   269  Check: 9723  Weight:  1.00
 Name: mip_ranpi       Len:   269  Check: 5937  Weight:  1.00
 Name: mip_bovin       Len:   269  Check: 1430  Weight:  1.00
 Name: mip_human       Len:   269  Check:  372  Weight:  1.00
 Name: mip_chick       Len:   269  Check: 4658  Weight:  1.00
 Name: aqp5_rat        Len:   269  Check: 9033  Weight:  1.00
 Name: aqp5_human      Len:   269  Check: 6547  Weight:  1.00
 Name: aqp2_human      Len:   269  Check: 6209  Weight:  1.00
 Name: aqp4_human      Len:   269  Check: 2589  Weight:  1.00
 Name: aqp4_rat        Len:   269  Check: 4412  Weight:  1.00
 Name: aqp4_mouse      Len:   269  Check: 2845  Weight:  1.00
 Name: aqp2_rat        Len:   269  Check: 5748  Weight:  1.00
 Name: aqp2_mouse      Len:   269  Check: 6526  Weight:  1.00
 Name: wc2a_arath      Len:   269  Check: 4866  Weight:  1.00
 Name: aqp6_human      Len:   269  Check: 9404  Weight:  1.00
 Name: wc2c_arath      Len:   269  Check: 6187  Weight:  1.00
 Name: wc2b_arath      Len:   269  Check: 7328  Weight:  1.00
 Name: wc1c_arath      Len:   269  Check: 8575  Weight:  1.00
 Name: wc1b_arath      Len:   269  Check: 9544  Weight:  1.00
 Name: tipw_lyces      Len:   269  Check: 9283  Weight:  1.00
 Name: wc1a_arath      Len:   269  Check:  598  Weight:  1.00
 Name: tipw_pea        Len:   269  Check: 9253  Weight:  1.00
 Name: tipa_arath      Len:   269  Check: 6544  Weight:  1.00
 Name: aqua_atrca      Len:   269  Check: 2848  Weight:  1.00
 Name: dip_antma       Len:   269  Check: 9619  Weight:  1.00
 Name: aqpz_ecoli      Len:   269  Check: 5641  Weight:  1.00
 Name: tip2_tobac      Len:   269  Check:  490  Weight:  1.00
 Name: tip1_tobac      Len:   269  Check:  622  Weight:  1.00
 Name: tipg_arath      Len:   269  Check: 3231  Weight:  1.00
 Name: bib_drome       Len:   269  Check: 7687  Weight:  1.00
 Name: tipr_arath      Len:   269  Check: 4476  Weight:  1.00
 Name: tipa_phavu      Len:   269  Check: 5563  Weight:  1.00
 Name: tipg_orysa      Len:   269  Check: 3537  Weight:  1.00

//


              1                                                   50
predict_h258  MASEIKKKLF WRAVVAEFLA MTLFVFISIG SALGFNYPLE RNQTLVQDNV
aqp1_rat      MASEIKKKLF WRAVVAEFLA MTLFVFISIG SALGFNYPLE RNQTLVQDNV
aqp1_mouse    MASEIKKKLF WRAVVAEFLA MTLFVFISIG SALGFNYPLE RNQTLVQDNV
aqp1_human    MASEFKKKLF WRAVVAEFLA TTLFVFISIG SALGFKYPVG NNQTAVQDNV
aqp1_bovin    MASEFKKKLF WRAVVAEFLA MILFIFISIG SALGFHYPIK SNQTtvQDNV
aqp1_sheep    MASEFKKKLF WRAVVAEFLA MILFIFISIG SALGFHYPIK SNQTtvQDNV
aqpa_ranes    MASEFKKKAF WRAVIAEFLA MILFVFISIG AALGFNFPIE EKANQtqDIV
aqp2_dasno    ......SVAF SRAVLAEFLA TLIFVFFGLG SALSWPQALP S.......VL
aqp2_bovin    ......SIAF SRAVLAEFLA TLLFVFFGLG SALNWPQALP S.......VL
aqp2_canfa    ......SVAF SRAVFAEFLA TLLFVFFGLG SALNWPQALP S.......VL
aqp2_rabit    ......SIAF SRAVFAEFLA TLLFVFFGLG SALNWPSALP S.......TL
aqp2_elema    ......SIAF SRAVFSEFLA TLLFVFFGLG SALNWPQALP S.......VL
aqp2_horse    ......SIAF SRAVLAEFLA TLLFVFFGLG SALNWPQAMP S.......VL
aqp2_proha    ......SIAF SRAVLSEFLA TLLFVFFGLG SALNWPQALP S.......VL
mip_rat       ...ELRSASF WRAIFAEFFA TLFYVFFGLG SSLRWA.... ...PGPLHVL
aqp2_oryaf    ......SIAF SKAVFSEFLA TLLFVFFGLG SALNWPQALP S.......GL
mip_mouse     .MWELRSASF WRAIFAEFFA TLFYVFFGLG ASLRWA.... ...PGPLHVL
mip_ranpi     .MWEFRSFSF WRAVFAEFFG TMFYVFFGLG ASLKWAAGPA .......NVL
mip_bovin     .MWELRSASF WRAICAEFFA SLFYVFFGLG ASLRWA.... ...PGPLHVL
mip_human     .MWELRSASF WRAIFAEFFA TLFYVFFGLG SSLRWA.... ...PGPLHVL
mip_chick     .......... .......... .......... .......... ..........
aqp5_rat      MKKEVCSLAF FKAVFAEFLA TLIFVFFGLG SALKWPSALP T.......IL
aqp5_human    MKKEVCSVAF LKAVFAEFLA TLIFVFFGLG SALKWPSALP T.......IL
aqp2_human    .MWELRSIAF SRAVFAEFLA TLLFVFFGLG SALNWPQALP S.......VL
aqp4_human    AFKGVWTQAF WKAVTAEFLA MLIFVLLSLG STINWG...G TEKPLPVDMV
aqp4_rat      AFKGVWTQAF WKAVTAEFLA MLIFVLLSVG STINWG...G SENPLPVDMV
aqp4_mouse    AFKGVWTQAF WKAVSAEFLA TLIFVL.GVG STINWG...G SENPLPVDMV
aqp2_rat      .MWELRSIAF SRAVLAEFLA TLLFVFFGLG SALQWASSPP S.......VL
aqp2_mouse    .MWELRSIAY CRAVLAEFLA TLLFVFFGLG SALQWASSPP S.......VL
wc2a_arath    DGAELKKWSF YRAVIAEFVA TLLFLYITVL TVIGYKIQSD TDAGGVdgIL
aqp6_human    MLACRLWKAI SRALFAEFLA TGLYVFFGVG SVMRWPTALP S.......VL
wc2c_arath    DAEELTKWSL YRAVIAEFVA TLLFLYVTVL TVIGYKIQSD TKAGGVdgIL
wc2b_arath    DADELTKWSL YRAVIAEFVA TLLFLYITVL TVIGYKIQSD TKAGGVdgIL
wc1c_arath    EPGELSSWSF YRAGIAEFIA TFLFLYITVL TVMGVKRA.. PNMCASVGIQ
wc1b_arath    EPGELASWSF WRAGIAEFIA TFLFLYITVL TVMGVKR..S PNMCASVGIQ
tipw_lyces    EPGELSSWSF YRAGIAEFMA TFLFLYITIL TVMGLKRSDS LCSSV..GIQ
wc1a_arath    EPGELSSWSF WRAGIAEFIA TFLFLYITVL TVMGVKR..S PNMCASVGIQ
tipw_pea      EPSELTSWSF YRAGIAEFIA TFLFLYITVL TVMGVVRESS KCKTV..GIQ
tipa_arath    RADEATHPDS IRATLAEFLS TFVFVFAAEG SILSLDKLYW EHAAHAGTni
aqua_atrca    DMGELKLWSF WRAAIAEFIA TLLFLYITVA TVIGYKKETD PCASVGL..L
dip_antma     SIGDSFSVAS IKAYVAEFIA TLLFVFAGVG SAIAYNKLTS DAALDPAGLV
aqpz_ecoli    .........M FRKLAAECFG TFWLVFGGCG SAVLAAGFPE ....LGIGFA
tip2_tobac    SIGDSFSVGS LKAYVAEFIA TLLFVFAGVG SAIAYNKLTA DAALDPAGLV
tip1_tobac    SIGDSFSVGS LKAYVAEFIA TLLFVFAGVG SAIAYNKLTA DAALDPAGLV
tipg_arath    RPDEATRPDA LKAALAEFIS TLIFVVAGSG SGMAFNKLTE NGATTPSGLV
bib_drome     MQAEIRTLEF WRSIISECLA SFMYVFIVCG AAAGVGVGAS VSSVL....L
tipr_arath    RPDEATRPDA LKAALAEFIS TLIFVVAGSG SGMAFNKLTE NGATTPSGLV
tipa_phavu    RTDEATHPDS MRASLAEFAS TFIFVFAGEG SGLALVKIYQ DSAFSAGELL
tipg_orysa    SHQEVYHPGA LKAALAEFIS TLIFVFAGQG SGMAFSKLTG GGATTPAGLI

              51                                                 100
predict_h258  KVSLAFGLSI ATLAQSVGHI SGAHSNPAVT LGLLLSCQIS ILRAVMYIIA
aqp1_rat      KVSLAFGLSI ATLAQSVGHI SGAHSNPAVT LGLLLSCQIS ILRAVMYIIA
aqp1_mouse    KVSLAFGLSI ATLAQSVGHI SGAHLNPAVT LGLLLSCQIS ILRAVMYIIA
aqp1_human    KVSLAFGLSI ATLAQSVGHI SGAHLNPAVT LGLLLSCQIS IFRALMYIIA
aqp1_bovin    KVSLAFGLSI ATLAQSVGHI SGAHLNPAVT LGLLLSCQIS VLRAIMYIIA
aqp1_sheep    KVSLAFGLSI ATLAQSVGHI SGAHLNPAVT LGLLLSCQIS ILRAIMYIIA
aqpa_ranes    KVSLAFGISI ATMAQSVGHV SGAHLNPAVT LGCLLSCQIS ILKAVMYIIA
aqp2_dasno    QIALAFGLAI GTLVQALGHV SGAHINPAVT VACLVGCHVS FLRAAFYVAA
aqp2_bovin    QIAMAFGLAI GTLVQALGHV SGAHINPAVT VACLVGCHVS FLRAVFYVAA
aqp2_canfa    QIAMAFGLGI GTLVQALGHV SGAHINPAVT VACLVGCHVS FLRAAFYVAA
aqp2_rabit    QIAMAFGLGI GTLVQALGHV SGAHINPAVT VACLVGCHVS FLRAAFYVAA
aqp2_elema    QIAMAFGLAI GTLVQTLGHI SGAHINPAVT VACLVGCHVS FLRATFYLAA
aqp2_horse    QIAMAFGLAI GTLVQALGHV SGAHINPAVT VACLVGCHVS FLRAAFYVAA
aqp2_proha    QIAMAFGLAI GTLVQTLGHI SGAHINPAVT IACLVGCHVS FLRALFYLAA
mip_rat       QVALAFGLAL ATLVQTVGHI SGAHVNPAVT FAFLVGSQMS LLRAFCYIAA
aqp2_oryaf    QIAMAFGLAI GTLVQTLGHI SGAHINPAVT VACLVGCHVS FLRAIFYVAA
mip_mouse     QVALAFGLAL ATLVQTVGHI SGAHVNPAVT FAFLVGSQMS LLRAFCYIAA
mip_ranpi     VIALAFGLVL ATMVQSIGHV SGAHINPAVT FAFLIGSQMS LFRAIFYIAA
mip_bovin     QVALAFGLAL ATLVQAVGHI SGAHVNPAVT FAFLVGSQMS LLRAICYMVA
mip_human     QVAMAFGLAL ATLVQSVGHI SGAHVNPAVT FAFLVGSQMS LLRAFCYMAA
mip_chick     .......... .......... .......... .......... ..........
aqp5_rat      QISIAFGLAI GTLAQALGPV SGGHINPAIT LALLIGNQIS LLRAVFYVAA
aqp5_human    QIALAFGLAI GTLAQALGPV SGGHINPAIT LALLVGNQIS LLRAFFYVAA
aqp2_human    QIAMAFGLGI GTLVQALGHI SGAHINPAVT VACLVGCHVS VLRAAFYVAA
aqp4_human    LISLCFGLSI ATMVQCFGHI SGGHINPAVT VAMVCTRKIS IAKSVFYIAA
aqp4_rat      LISLCFGLSI ATMVQCFGHI SGGHINPAVT VAMVCTRKIS IAKSVFYITA
aqp4_mouse    LISLCFGLSI ATMVQCLGHI SGGHINPAVT VAMVCTRKIS IAKSVFYIIA
aqp2_rat      QIAVAFGLGI GILVQALGHV SGAHINPAVT VACLVGCHVS FLRAAFYVAA
aqp2_mouse    QIAVAFGLGI GTLVQALGHV SGAHINPAVT VACLVGCHVS FLRAAFYVAA
wc2a_arath    GIAWAFGGMI FILVYCTAGI SGGHINPAVT FGLFLARKVS LPRALLYIIA
aqp6_human    QIAITFNLVT AMAVQVTWKT SGAHANPAVT LAFLVGSHIS LPRAVAYVAA
wc2c_arath    GIAWAFGGMI FILVYCTAGI SGGHINPAVT FGLFLARKVS LIRAVLYMVA
wc2b_arath    GIAWAFGGMI FILVYCTAGI SGGHINPAVT FGLFLARKVS LIRAVLYMVA
wc1c_arath    GIAWAFGGMI FALVYCTAGI SGGHINPAVT FGLFLARKLS LTRAVFYIVM
wc1b_arath    GIAWAFGGMI FALVYCTAGI SGGHINPAVT FGLFLARKLS LTRAVYYIVM
tipw_lyces    GVAWAFGGMI FALVYCTAGI SGGHINPAVT FGLFLARKLS LTRAVFYMVM
wc1a_arath    GIAWAFGGMI FALVYCTAGI SGGHINPAVT FGLFLARKLS LTRALYYIVM
tipw_pea      GIAWAFGGMI FALVYCTAGI SGGHINPAVT FGLFLARKLS LTRAIFYMVM
tipa_arath    LVALAHAFAL FAAVSAAINV SGGHVNPAVT FGALVGGRVT AIRAIYYWIA
aqua_atrca    GIAWSFGGMI FVLVYCTAGI SGGHINPAVT FGLFLARKVS LLRALVYMIA
dip_antma     AVAVAHAFAL FVGVSMAANV SGGHLNPAVT LGLAVGGNIT ILTGLFYWIA
aqpz_ecoli    GVALAFGLTV LTMAFAVGHI SGGHFNPAVT IGLWAGGRFP AKEVVGYVIA
tip2_tobac    AVAVAHAFAL FVGVSIAANI SGGHLNPAVT LGLAVGGNIT ILTGFFYWIA
tip1_tobac    AVAVAHAFAL FVGVSIAANI SGGHLNPAVT LGLAVGGNIT ILTGFFYWIA
tipg_arath    AAAVAHAFGL FVAVSVGANI SGGHVNPAVT FGAFIGGNIT LLRGILYWIA
bib_drome     ATALASGLAM ATLTQCFLHI SGAHINPAVT LALCVVRSIS PIRAAMYITA
tipr_arath    AAAVAHAFGL FVAVSVGANI SGGHVNPAVT FGAFIGGNIT LLRGILYWIA
tipa_phavu    ALALAHAFAL FAAVSASMHV SGGHVNPAVS FGALIGGRIS VIRAVYYWIA
tipg_orysa    AAAVAHAFAL FVAVSVGANI SGGHVNPAVT FGAFVGGNIT LFRGLLYWIA

              101                                                150
predict_h258  QCVGAIVASA ILSGITSSLL ENSLGRNDLA RGVNSGQGLG IEIIGTLQLV
aqp1_rat      QCVGAIVASA ILSGITSSLL ENSLGRNDLA RGVNSGQGLG IEIIGTLQLV
aqp1_mouse    QCVGAIVATA ILSGITSSLV DNSLGRNDLA HGVNSGQGLG IEIIGTLQLV
aqp1_human    QCVGAIVATA ILSGITSSLT GNSLGRNDLA DGVNSGQGLG IEIIGTLQLV
aqp1_bovin    QCVGAIVATA ILSGITSSLP DNSLGLNALA PGVNSGQGLG IEIIGTLQLV
aqp1_sheep    QCVGAIVATV ILSGITSSLP DNSLGLNALA PGVNSGQGLG IEIIGTLQLV
aqpa_ranes    QCLGAVVATA ILSGITSGLE NNSLGLNGLS PGVSAGQGLG VEILVTFQLV
aqp2_dasno    QLLGAVAGAA ILHEITPPDV RG........ .......... ..........
aqp2_bovin    QLLGAVAGAA LLHEITPPAI RG........ .......... ..........
aqp2_canfa    QLLGAVAGAA LLHEITPPHV RG........ .......... ..........
aqp2_rabit    QLLGAVAGAA LLHEITPAEV RG........ .......... ..........
aqp2_elema    QLLGAVAGAA LLHELTPPDI RG........ .......... ..........
aqp2_horse    QLLGAVAGAA LLHEITPPDI RR........ .......... ..........
aqp2_proha    QLLGAVAGAA LLHELTPPDI RG........ .......... ..........
mip_rat       QLLGAVAGAA VLYSVTPPAV RGNLALNTLH AGVSVGQATT VEIFLTLQFV
aqp2_oryaf    QLLGAVAGAA LLHELTPPDI RG........ .......... ..........
mip_mouse     QLLGAVAGAA VLYSVTPPAV RGNLALNTLH TGVSVGQATT VEIFLTLQFV
mip_ranpi     QLLGAVAGAA VLYGVTPAAI RGNLALNTLH PGVSLGQATT VEIFLTLQFV
mip_bovin     QLLGAVAGAA VLYSVTPPAV RGNLALNTLH PGVSVGQATI VEIFLTLQFV
mip_human     QLLGAVAGAA VLYSVTPPAV RGNLALNTLH PAVSVGQATT VEIFLTLQFV
mip_chick     .......... .......... .......... .......... ..........
aqp5_rat      QLVGAIAGAG ILYWLAPLNA RGNLAVNALN NNTTPGKAMV VELILTFQLA
aqp5_human    QLVGAIAGAG ILYGVAPLNA RGNLAVNALN NNTTQGQAMV VELILTFQLA
aqp2_human    QLLGAVAGAA LLHEITPADI RGDLAVNALS NSTTAGQAVT VELFLTLQLV
aqp4_human    QCLGAIIGAG ILYLVTPPSV VGGLGVTMVH GNLTAGHGLL VELIITFQLV
aqp4_rat      QCLGAIIGAG ILYLVTPPSV VGGLGVTTVH GNLTAGHGLL VELIITFQLV
aqp4_mouse    QCLGAIIGAG ILYLVTPPSV VGGLGVTTVH GNLTAGHGLL VELIITFQLV
aqp2_rat      QLLGAVAGAA ILHEITPVEI RGDLAVNALH NNATAGQAVT VELFLTMQLV
aqp2_mouse    QLLGAVAGAA ILHEITPVEI RGDLAVNALH NNATAGQAVT VELFLTMQLV
wc2a_arath    QCLGAICGVG FVKAFQSSYY TRYGGgnSLA DGYSTGTGLA AEIIGTFVLV
aqp6_human    QLVGATVGAA LLYGVMPGDI RETLGINVVR NSVSTGQAVA VELLLTLQLV
wc2c_arath    QCLGAICGVG FVKAFQSSHY VNYGGgnFLA DGYNTGTGLA AEIIGTFVLV
wc2b_arath    QCLGAICGVG FRQSFQSSYY DRYGGgnSLA DGYNTGTGLA AEIIGTFVLV
wc1c_arath    QCLGAICGAG VVKGFQPNPY QtgGGANTVA HGYTKGSGLG AEIIGTFVLV
wc1b_arath    QCLGAICGAG VVKGFQPKQY QagGGANTIA HGYTKGSGLG AEIIGTFVLV
tipw_lyces    QCLGAICGAG VVKGFMVGPY QrgGGANVVN PGYTKGDGLG AEIIGTFVLV
wc1a_arath    QCLGAICGAG VVKGFQPKQY QagGGANTVA HGYTKGSGLG AEIIGTFVLV
tipw_pea      QVLGAICGAG VVKGFEGKQR FGDLNgnFVA PGYTKGDGLG AEIVGTFILV
tipa_arath    QLLGAILACL LLRLTTNGMR PVGFR...LA SGVGAVNGLV LEIILTFGLV
aqua_atrca    QCAGAICGVG LVKAFMKGPY NqgGGANSVA LGYNKGTAFG AELIGTFVLV
dip_antma     QCLGSTVACL LLKFVTNGL. ..SVPTHGVA AGMDAIQGVV MEIIITFALV
aqpz_ecoli    QVVGGIVAAA LLYLIASGKT GFDAAASGFA sgYSMLSALV VELVLSAGFL
tip2_tobac    QLLGSTVACL LLKYVTNGL. ..AVPTHGVA AGLNGFQGVV MEIIITFALV
tip1_tobac    QLLGSTVACL LLKYVTNGL. ..AVPTHGVA AGLNGLQGVV MEIIITFALV
tipg_arath    QLLGSVVACL ILKFATGGLA VPAFG...LS AGVGVLNAFV FEIVMTFGLV
bib_drome     QCGGGIAGAA LLYGVTVPGY QGNLQAasHS AALAAWERFG VEFILTSLVV
tipr_arath    QLLGSVVACL ILKFATGGLA VPPFG...LS AGVGVLNAFV FEIVMTFGLV
tipa_phavu    QLLGSIVAAL VLRLVTNNMR PSGF...HVS PGVGVGHMFI LEVVMTFGLM
tipg_orysa    QLLGSTVACF LLRFSTGGLA TGTFGL.... TGVSVWEALV LEIVMTFGLV

              151                                                200
predict_h258  LCVLATTDRR RRDLGGSAPL AIGLSVALGH LLAIDYTGCG INPARSFGSA
aqp1_rat      LCVLATTDRR RRDLGGSAPL AIGLSVALGH LLAIDYTGCG INPARSFGSA
aqp1_mouse    LCVLATTDRR RRDLGGSAPL AIGLSVALGH LLAIDYTGCG INPARSFGSA
aqp1_human    LCVLATTDRR RRDLGGSAPL AIGLSVALGH LLAIDYTGCG INPARSFGSA
aqp1_bovin    LCVLATTDRR RRDLGGSGPL AIGFSVALGH LLAIDYTGCG INPARSFGSS
aqp1_sheep    LCVLATTDRR RrdLGDSGPL AIGFSVALGH LLAIDYTGCG INPARSFGSS
aqpa_ranes    LCVVAVTDRR RHDVSGSVPL AIGLSVALGH LIAIDYTGCG MNPARSFGSA
aqp2_dasno    .......... .......... .......... .......... ..........
aqp2_bovin    .......... .......... .......... .......... ..........
aqp2_canfa    .......... .......... .......... .......... ..........
aqp2_rabit    .......... .......... .......... .......... ..........
aqp2_elema    .......... .......... .......... .......... ..........
aqp2_horse    .......... .......... .......... .......... ..........
aqp2_proha    .......... .......... .......... .......... ..........
mip_rat       LCIFATYDER RNGRMGSVAL AVGFSLTLGH LFGMYYTGAG MNPARSFAPA
aqp2_oryaf    .......... .......... .......... .......... ..........
mip_mouse     LCIFATYDER RNGRMGSVAL AVGFSLTLGH LFGMYYTGAG MNPARSFAPA
mip_ranpi     LCIFATYDER RNGRLGSVSL AIGFSLTLGH LFGLYYTGAS MNPARSFAPA
mip_bovin     LCIFATYDER RNGRLGSVAL AVGFSLTLGH LFGMYYTGAG MNPARSFAPA
mip_human     LCIFATYDER RNGQLGSVAL AVGFSLALGH LFGMYYTGAG MNPARSFAPA
mip_chick     ........DR HDGRPGSAAL PVGFSLALGH LFGIPFTGAG MNPARSFAPA
aqp5_rat      LCIFSSTDSR RTSPVGSPAL SIGLSVTLGH LVGIYFTGCS MNPARSFGPA
aqp5_human    LCIFASTDSR RTSPVGSPAL SIGLSVTLGH LVGIYFTGCS MNPARSFGPA
aqp2_human    LCIFASTDER RGENPGTPAL SIGFSVALGH LLGIHYTGCS MNPARSLAPA
aqp4_human    FTIFASCDSK RTDVTGSIAL AIGFSVAIGH LFAINYTGAS MNPARSFGPA
aqp4_rat      FTIFASCDSK RTDVTGSVAL AIGFSVAIGH LFAINYTGAS MNPARSFGPA
aqp4_mouse    FTVFASCDSK RTDVTGSIAL AIGFSVAIGH LFAINYTGAS MNPARSFGPA
aqp2_rat      LCIFASTDER RGDNLGSPAL SIGFSVTLGH LLGIYFTGCS MNPARSLAPA
aqp2_mouse    LCIFASTDER RSDNLGSPAL SIGFSVTLGH LLGIYFTGCS MNPARSLAPA
wc2a_arath    YTVFSATDPK RSavPVLAPL PIGFAVFMVH LATIPITGTG INPARSFGAA
aqp6_human    LCVFASTDSR QTS..GSPAT MIGISWALGH LIGILFTGCS MNPARSFGPA
wc2c_arath    YTVFSATDPK RNavPVLAPL PIGFAVFMVH LATIPITGTG INPARSFGAA
wc2b_arath    YTVFSATDPK RNavPVLAPL PIGFAVFMVH LATIPITGTG INPARSFGAS
wc1c_arath    YTVFSATDAK RSavPILAPL PIGFAVFLVH LATIPITGTG INPARSLGAA
wc1b_arath    YTVFSATDAK RNavPILAPL PIGFAVFLVH LATIPITGTG INPARSLGAA
tipw_lyces    YTVFSATDAK RNavPILAPL PIGFAVFLVH LATIPITGTG INPARSLGAA
wc1a_arath    YTVFSATDAK RNavPILAPL PIGFAVFLVH LATIPITATG INPARSLGAA
tipw_pea      YTVFSATDAK RSavPILAPL PIGFAVFLVH LATIPITGTG INPARSLGAA
tipa_arath    YVVYStiDPK RGSLGIIAPL AIGLIVGANI LVGGPFSGAS MNPARAFGPA
aqua_atrca    YTVFSATDPK RSavPILAPL PIGFAVFMVH LATIPITGTG INPARSFGAA
dip_antma     YTVYAtaDPK KGSLGVIAPI AIGFIVGANI LAAGPFSGGS MNPARSFGPA
aqpz_ecoli    LVIHGATDKF APA..GFAPI AIGLALTLIH LISIPVTNTS VNPARSTAVA
tip2_tobac    YTVYAtaDPK KGSLGTIAPI AIGFIVGANI LAAGPFSGGS MNPARSFGPA
tip1_tobac    YTVYAtaDPK KGSLGTIAPI AIGFIVGANI LAAGPFSGGS MNPARSFGPA
tipg_arath    YTVYAtiDPK NGSLGTIAPI AIGFIVGANI LAGGAFSGAS MNPAVAFGPA
bib_drome     LCYFVSTDPM KKFMGNS.AA SIGCAYSACC FVSMPYLN.. ..PARSLGPS
tipr_arath    YTVYAtiDPK NGSLGTIAPI AIGFIVGANI LAGGAFSGAS MNPAVAFGPA
tipa_phavu    YTVYGtiDPK RGAVSYIAPL AIGLIVGANI LVGGPFDGAC MNPALAFGPS
tipg_orysa    YTVYAtvDPK KGSLGTIAPI AIGFIVGANI LVGGAFDGAS MNPAVSFGPA

              201                                                250
predict_h258  VLTRNFSNHW IFWVGPFIGS ALAVLIYDFI LAPRSSDFTD RMKVWTSGQV
aqp1_rat      VLTRNFSNHW IFWVGPFIGS ALAVLIYDFI LAPRSSDFTD RMKVWTSGQV
aqp1_mouse    VLTRNFSNHW IFWVGPFIGG ALAVLIYDFI LAPRSSDFTD RMKVWTSGQV
aqp1_human    VITHNFSNHW IFWVGPFIGG ALAVLIYDFI LAPRSSDLTD RVKVWTSGQV
aqp1_bovin    VITHNFQDHW IFWVGPFIGA ALAVLIYDFI LAPRSSDLTD RVKVWTSGQV
aqp1_sheep    VITHNFQDHW IFWVGPFIGA ALAVLIYDFI LAPRSSDLTD RVKVWTSGQV
aqpa_ranes    VLTKNFTYHW IFWVGPMIGG AAAAIIYDFI LAPRTSDLTD RMKVWTNGQV
aqp2_dasno    .......... .......... .......... .......... ..........
aqp2_bovin    .......... .......... .......... .......... ..........
aqp2_canfa    .......... .......... .......... .......... ..........
aqp2_rabit    .......... .......... .......... .......... ..........
aqp2_elema    .......... .......... .......... .......... ..........
aqp2_horse    .......... .......... .......... .......... ..........
aqp2_proha    .......... .......... .......... .......... ..........
mip_rat       ILTRNFSNHW VYWVGPIIGG GLGSLLYDFL LFPRLKSVSE RLSILKGARP
aqp2_oryaf    .......... .......... .......... .......... ..........
mip_mouse     ILTRNFSNHW VYWVGPIIGG GLGSLLYDFL LFPRLKSVSE RLSILKGARP
mip_ranpi     VLTRNFTNHW VYWVGPIIGG ALGGLVYDFI LFPRMRGLSE RLSILKGARP
mip_bovin     ILTRNFTNHW VYWVGPVIGA GLGSLLYDFL LFPRLKSVSE RLSILKGSRP
mip_human     ILTGNFTNHW VYWVGPIIGG GLGSLLYDFL LFPRLKSISE RLSVLKGAKP
mip_chick     VITRNFTNHW VFWAGPLLGA ALAALLYELA LCPRARSMAE RLAV.LRGEP
aqp5_rat      VVMNRFssHW VFWVGPIVGA MLAAILYFYL LFPSSLSLHD RVAVVKGTYE
aqp5_human    VVMNRFsaHW VFWVGPIVGA VLAAILYFYL LFPNSLSLSE RVAIIKGTYE
aqp2_human    VVTGKFDDHW VFWIGPLVGA ILGSLLYNYV LFPPAKSLSE RLAVLKGLEp
aqp4_human    VIMGNWENHW IYWVGPIIGA VLAGGLYEYV FCPDVEFKRR FKEAFSKaqT
aqp4_rat      VIMGNWENHW IYWVGPIIGA VLAGALYEYV FCPDVELKRR LKEAFSKaqT
aqp4_mouse    VIMGNWANHW IYWVGPIMGA VLAGALYEYV FCPDVELKRR LKEAFSKaqT
aqp2_rat      VVTGKFDDHW VFWIGPLVGA IIGSLLYNYL LFPSAKSLQE RLAVLKGLEp
aqp2_mouse    VVTGKFDDHW VFWIGPLVGA IIGSLLYNYL LFPSTKSLQE RLAVLKGLEp
wc2a_arath    VIYnpWDDHW IFWVGPFIGA AIAAFYHQFV LRASGSKSLG SFRSAANV..
aqp6_human    IIIGKFTVHW VFWVGPLMGA LLASLIYNFV LFPDTKTLAQ RLAILTGTVE
wc2c_arath    VIFnpWDDHW IFWVGPFIGA TIAAFYHQFV LRASGSKSLG SFRSAANV..
wc2b_arath    VIYnpWDDHW IFWVGPFIGA AIAAFYHQFV LRASGSKSLG SFRSAANV..
wc1c_arath    IIYnaWDDHW IFWVGPFIGA ALAALYHQLV IRAIPFKSRS ..........
wc1b_arath    IIFnaWDDHW VFWVGPFIGA ALAALYHVIV IRAIPFKSRS ..........
tipw_lyces    IIYnaWNDHW IFWVGPMIGA ALAAIYHQII IRAMPFHRS. ..........
wc1a_arath    IIYnsWDDHW VFWVGPFIGA ALAALYHVVV IRAIPFKSRS ..........
tipw_pea      IVFngWNDHW IFWVGPFIGA ALAALYHQVV IRAIPFKSK. ..........
tipa_arath    LVGWRWHDHW IYWVGPFIGS ALAALIYEYM VIPTEPPTHH AHGVHQPLAP
aqua_atrca    VIyrVWDDHW IFWVGPFVGA LAAAAYHQYV LRAAAIKALG SFRSNPTN..
dip_antma     VASGDFSQNW IYWAGPLIGG ALAGFIYGDV FITAHAPLPT SEDYA.....
aqpz_ecoli    IFQgaLEQLW FFWVVPIVGG IIGGLIYRTL LEKRD..... ..........
tip2_tobac    VVAGDFSQNW IYWAGPLIGG GLAGFIYGDV FIGCHTPLPT SEDYA.....
tip1_tobac    VVAGDFSQNW IYWAGPLIGG GLAGFIYGDV FIGCHTPLPT SEDYA.....
tipg_arath    VVSWTWTNHW VYWAGPLVGG GIAGLIYEVF FINTTHEQLP TTDY......
bib_drome     FVLNKWDSHW VYWFGPLVGG MASGLVYEYI FNSRNRNLRH NKGSIDNDSS
tipr_arath    VVSWTWTNHW VYWAGPLVGG GIAGLIYEVF FINTTHTSSS NHRLLN....
tipa_phavu    LVGWQWHQHW IFWVGPLLGA ALAALVYEYA VIPIEPPPHH HQPLATEDY.
tipg_orysa    LVSWSWESQW VYWVGPLIGG GLAGVIYEVL FISHTHEQLP TTDY......

              251              269
predict_h258  EEYDLDADDI NSRVEMKPK
aqp1_rat      EEYDLDADDI NSRVEMKPK
aqp1_mouse    EEYDLDADDI NSRVEMKPK
aqp1_human    EEYDLDADDI NSRVEMKPK
aqp1_bovin    EEYDLDADDI NSRVEMKPK
aqp1_sheep    EEYDLDADDI NSRVEMKPK
aqpa_ranes    EEYELDGDD. NTRVEMKPK
aqp2_dasno    .......... .........
aqp2_bovin    .......... .........
aqp2_canfa    .......... .........
aqp2_rabit    .......... .........
aqp2_elema    .......... .........
aqp2_horse    .......... .........
aqp2_proha    .......... .........
mip_rat       SDSNGQPEGT GEPVELKTQ
aqp2_oryaf    .......... .........
mip_mouse     SDSNGQPEGT GEPVELKTQ
mip_ranpi     AEPEGQQEAT GEPIELKTQ
mip_bovin     SESNGQPEVT GEPVELKTQ
mip_human     DVSNGQPEVT GEPVELNTQ
mip_chick     PAAAPPPEPP AEPLELKTQ
aqp5_rat      PEEDWEDHRE ERKKTIELT
aqp5_human    PDEDWEEQRE ERKKTMELT
aqp2_human    tDWEEREVRR RQSVELHSP
aqp4_human    KGSYMEVEDN RSQVETDDL
aqp4_rat      KGSYMEVEDN RSQVETEDL
aqp4_mouse    KGSYMEVEDN RSQVETEDL
aqp2_rat      tDWEEREVRR RQSVELHSP
aqp2_mouse    tDWEEREVRR RQSVELHSP
wc2a_arath    .......... .........
aqp6_human    VGTGARAGAE PLKKESQPG
wc2c_arath    .......... .........
wc2b_arath    .......... .........
wc1c_arath    .......... .........
wc1b_arath    .......... .........
tipw_lyces    .......... .........
wc1a_arath    .......... .........
tipw_pea      .......... .........
tipa_arath    EDY....... .........
aqua_atrca    .......... .........
dip_antma     .......... .........
aqpz_ecoli    .......... .........
tip2_tobac    .......... .........
tip1_tobac    .......... .........
tipg_arath    .......... .........
bib_drome     SIHSEDELNY DMDMEKPNK
tipr_arath    .......... .........
tipa_phavu    .......... .........
tipg_orysa    .......... .........


________________________________________________________________________________




   Prediction of:           

    - secondary structure,          by PHDsec       
    - solvent accessibility,        by PHDacc       
    - and helical transmembrane regions,    by PHDhtm       

   PHD: Profile fed neural network systems from HeiDelberg
   ~~~~~~~~~~~~~~~~~~~~~~~~~~~~~~~~~~~~~~~~~~~~~~~~~~~~~~~

   Author:             Burkhard Rost        
                       EMBL, Heidelberg, FRG
                       Meyerhofstrasse 1, 69 117 Heidelberg
                       Internet: Predict-Help@EMBL-Heidelberg.DE

   All rights reserved.





   ~~~~~~~~~~~~~~~~~~~~~~~~~~~~~~~~~~~~~~~~~    
   Secondary structure prediction by PHDsec:
   ~~~~~~~~~~~~~~~~~~~~~~~~~~~~~~~~~~~~~~~~~    

   Author:             Burkhard Rost        
                       EMBL, Heidelberg, FRG
                       Meyerhofstrasse 1, 69 117 Heidelberg
                       Internet: Rost@EMBL-Heidelberg.DE        

   All rights reserved.




About the network method
~~~~~~~~~~~~~~~~~~~~~~~

The network procedure is described in detail in:
1) Rost, Burkhard; Sander, Chris:
  Prediction of protein structure at better than 70% accuracy.
  J. Mol. Biol., 1993, 232, 584-599.            

A brief description is given in:
  Rost, Burkhard; Sander, Chris:
  Improved prediction of protein secondary structure by use of se-
  quence profiles and neural networks.
  Proc. Natl. Acad. Sci. U.S.A., 1993, 90, 7558-7562.           

The PHD mail server is described in:
2) Rost, Burkhard; Sander, Chris; Schneider, Reinhard:
  PHD - an automatic mail server for protein secondary structure
  prediction.
  CABIOS, 1994, 10, 53-60.

The latest improvement steps (up to 72%) are explained in:
3) Rost, Burkhard; Sander, Chris:
  Combining evolutionary information and neural networks to predict
  protein secondary structure.
  Proteins, 1994,  19, 55-72.

To be quoted for publications of PHD output:
  Papers 1-3 for the prediction of secondary structure and the pre-
  diction server.



About the input to the network
~~~~~~~~~~~~~~~~~~~~~~~~~~~~~~

The prediction is performed by a system of neural networks.
The input is a multiple sequence alignment. It is taken from an HSSP
file (produced by the program MaxHom:
  Sander, Chris & Schneider, Reinhard: Database of Homology-Derived
  Structures and the Structural Meaning of Sequence Alignment.
  Proteins, 1991, 9, 56-68.

For optimal results the alignment should contain sequences with varying
degrees of sequence similarity relative to the input protein.
The following is an ideal situation:

+-----------------+----------------------+
|   sequence:     |  sequence identity   |
+-----------------+----------------------+
| target sequence |  100 %               |
| aligned seq. 1  |   90 %               |
| aligned seq. 2  |   80 %               |
|      ...        |   ...                |
| aligned seq. 7  |   30 %               |
+-----------------+----------------------+



Estimated Accuracy of Prediction
~~~~~~~~~~~~~~~~~~~~~~~~~~~~~~~~

A careful cross validation test on some 250 protein chains (in total
about 55,000 residues) with less than 25% pairwise sequence identity
gave the following results:

++================++-----------------------------------------+
|| Qtotal = 72.1% ||      ("overall three state accuracy")   |
++================++-----------------------------------------+

+----------------------------+-----------------------------+
| Qhelix (% of observed)=70% | Qhelix (% of predicted)=77% |
| Qstrand(% of observed)=62% | Qstrand(% of predicted)=64% |
| Qloop  (% of observed)=79% | Qloop  (% of predicted)=72% |
+----------------------------+-----------------------------+
..........................................................................

These percentages are defined by:
~~~~~~~~~~~~~~~~~~~~~~~~~~~~~~~~

|                    number of correctly predicted residues
|Qtotal =            ---------------------------------------      (*100)
|                          number of all residues
|
|                    no of res correctly predicted to be in helix
|Qhelix (% of obs) = -------------------------------------------- (*100)
|                    no of all res observed to be in helix
|
|
|                    no of res correctly predicted to be in helix
|Qhelix (% of pred)= -------------------------------------------- (*100)
|                    no of all residues predicted to be in helix

..........................................................................

Averaging over single chains
~~~~~~~~~~~~~~~~~~~~~~~~~~~

The most reasonable way to compute the overall accuracies is the above
quoted percentage of correctly predicted residues.  However, since the
user is mainly interested in the expected performance of the prediction
for a particular protein, the mean value when averaging over protein
chains might be of help as well.  Computing first the three state
accuracy for each protein chain, and then averaging over 250 chains
yields the following average:

+-------------------------------====--+
| Qtotal/averaged over chains = 72.2% |
+-------------------------------====--+
| standard deviation          =  9.3% |
+-------------------------------------+

..........................................................................

Further measures of performance
~~~~~~~~~~~~~~~~~~~~~~~~~~~~~~

Matthews correlation coefficient:

+---------------------------------------------+
| Chelix = 0.63, Cstrand = 0.53, Cloop = 0.52 |
+---------------------------------------------+
..........................................................................

Average length of predicted secondary structure segments:

.           +------------+----------+
.           |  predicted | observed |
+-----------+------------+----------+
| Lhelix  = |    10.3    |    9.3   |
| Lstrand = |     5.0    |    5.3   |
| Lloop   = |     7.2    |    5.9   |
+-----------+------------+----------+
..........................................................................

The accuracy matrix in detail:

+---------------------------------------+
|    number of residues with H, E, L    |
+---------+------+------+------+--------+
|         |net H |net E |net L |sum obs |
+---------+------+------+------+--------+
| obs H   |12447 | 1255 | 3990 |  17692 |
| obs E   |  949 | 7493 | 3750 |  12192 |
| obs L   | 2604 | 2875 |19962 |  25441 |
+---------+------+------+------+--------+
| sum Net |16000 |11623 |27702 |  55325 |
+---------+------+------+------+--------+

Note: This table is to be read in the following manner:
     12447 of all residues predicted to be in helix, were observed to
     be in helix, 949 however belong to observed strands, 2604 to
     observed loop regions.  The term "observed" refers to the DSSP
     assignment of secondary structure calculated from 3D coordinates
     of experimentally determined structures (Dictionary of Secondary
     Structure  of Proteins: Kabsch & Sander (1983) Biopolymers, 22,
     2577-2637).



Position-specific reliability index
~~~~~~~~~~~~~~~~~~~~~~~~~~~~~~~~~~

The network predicts the three secondary structure types using real
numbers from the output units. The prediction is assigned by choosing
the maximal unit ("winner takes all").  However, the real numbers
contain additional information.
E.g. the difference between the maximal and the second largest output
unit can be used to derive a "reliability index".  This index is given
for each residue along with the prediction.  The index is scaled to
have values between 0 (lowest reliability), and 9 (highest).
The accuracies (Qtot) to be expected for residues with values above a
particular value of the index are given below as well as the fraction
of such residues (%res).:

+------+-----+-----+-----+-----+-----+-----+-----+-----+-----+-----+
| index|  0  |  1  |  2  |  3  |  4  |  5  |  6  |  7  |  8  |  9  |
| %res |100.0| 99.2| 90.4| 80.9| 71.6| 62.5| 52.8| 42.3| 29.8| 14.1|
+------+-----+-----+-----+-----+-----+-----+-----+-----+-----+-----+
|      |     |     |     |     |     |     |     |     |     |     |
| Qtot | 72.1| 72.3| 74.8| 77.7| 80.3| 82.9| 85.7| 88.5| 91.1| 94.2|
|      |     |     |     |     |     |     |     |     |     |     |
+------+-----+-----+-----+-----+-----+-----+-----+-----+-----+-----+
| H%obs| 70.4| 70.6| 73.7| 77.1| 80.1| 83.1| 86.0| 89.3| 92.5| 96.4|
| E%obs| 61.5| 61.7| 63.7| 66.6| 69.1| 71.7| 74.6| 77.0| 77.8| 68.1|
|      |     |     |     |     |     |     |     |     |     |     |
| H%prd| 77.8| 78.0| 80.0| 82.6| 84.7| 86.9| 89.2| 91.3| 93.1| 95.4|
| E%prd| 64.5| 64.7| 67.8| 71.0| 74.2| 77.6| 81.4| 85.1| 89.8| 93.5|
+------+-----+-----+-----+-----+-----+-----+-----+-----+-----+-----+

The above table gives the cumulative results, e.g. 62.5% of all
residues have a reliability of at least 5.  The overall three-state
accuracy for this subset of almost two thirds of all residues is 82.9%.
For this subset, e.g., 83.1% of the observed helices are correctly
predicted, and 86.9% of all residues predicted to be in helix are
correct.

..........................................................................

The following table gives the non-cumulative quantities, i.e. the
values per reliability index range.  These numbers answer the question:
how reliable is the prediction for all residues labeled with the
particular index i.

+------+-----+-----+-----+-----+-----+-----+-----+-----+-----+
| index|  1  |  2  |  3  |  4  |  5  |  6  |  7  |  8  |  9  |
| %res |  8.8|  9.5|  9.3|  9.1|  9.7| 10.5| 12.5| 15.7| 14.1|
+------+-----+-----+-----+-----+-----+-----+-----+-----+-----+
|      |     |     |     |     |     |     |     |     |     |
| Qtot | 46.6| 50.6| 57.7| 62.6| 67.9| 74.2| 82.2| 88.3| 94.2|
|      |     |     |     |     |     |     |     |     |     |
+------+-----+-----+-----+-----+-----+-----+-----+-----+-----+
| H%obs| 36.8| 42.3| 49.5| 55.2| 61.7| 69.9| 78.8| 87.4| 96.4|
| E%obs| 44.7| 44.5| 52.1| 55.4| 60.9| 68.0| 75.9| 81.0| 68.1|
|      |     |     |     |     |     |     |     |     |     |
| H%prd| 49.9| 52.5| 60.3| 64.2| 69.2| 77.5| 85.4| 89.9| 95.4|
| E%prd| 41.7| 47.1| 53.6| 57.0| 64.0| 71.6| 78.8| 88.8| 93.5|
+------+-----+-----+-----+-----+-----+-----+-----+-----+-----+

For example, for residues with Relindex = 5 64% of all predicted betha-
strand residues are correctly identified.





   ~~~~~~~~~~~~~~~~~~~~~~~~~~~~~~~~~~~~~~~~~~~      
   Solvent accessibility prediction by PHDacc:
   ~~~~~~~~~~~~~~~~~~~~~~~~~~~~~~~~~~~~~~~~~~~      

   Author:             Burkhard Rost        
                       EMBL, Heidelberg, FRG
                       Meyerhofstrasse 1, 69 117 Heidelberg
                       Internet: Rost@EMBL-Heidelberg.DE        

   All rights reserved.




About the network method
~~~~~~~~~~~~~~~~~~~~~~~

The network for prediction of secondary structure is described in
detail in:
  Rost, Burkhard; Sander, Chris:
  Prediction of protein structure at better than 70% accuracy.
  J. Mol. Biol., 1993, 232, 584-599.

The analysis of the prediction of solvent exposure is given in:
  Rost, Burkhard; Sander, Chris:
  Conservation and prediction of solvent accessibility in protein
  families.  Proteins, 1994, 20, 216-226.

To be quoted for publications of PHD exposure prediction:
  Both papers quoted above.



Definition of accessibility
~~~~~~~~~~~~~~~~~~~~~~~~~~

For training the residue solvent accessibility the DSSP (Dictionary of
Secondary Structure of Proteins; Kabsch & Sander (1983) Biopolymers, 22,
2577-2637) values of accessible surface area have been used.  The
prediction provides values for the relative solvent accessibility.  The
normalisation is the following:

|                           ACCESSIBILITY (from DSSP in Angstrom)
|RELATIVE_ACCESSIBILITY =   ------------------------------------- * 100
|                               MAXIMAL_ACC (amino acid type i)

where MAXIMAL_ACC (i) is the maximal accessibility of amino acid type i.
The maximal values are:

+----+----+----+----+----+----+----+----+----+----+----+----+
|  A |  B |  C |  D |  E |  F |  G |  H |  I |  K |  L |  M |
| 106| 160| 135| 163| 194| 197|  84| 184| 169| 205| 164| 188|
+----+----+----+----+----+----+----+----+----+----+----+----+
|  N |  P |  Q |  R |  S |  T |  V |  W |  X |  Y |  Z |
| 157| 136| 198| 248| 130| 142| 142| 227| 180| 222| 196|
+----+----+----+----+----+----+----+----+----+----+----+

Notation: one letter code for amino acid, B stands for D or N; Z stands
  for E or Q; and X stands for undetermined.

The relative solvent accessibility can be used to estimate the number
of water molecules (W) in contact with the residue:

W = ACCESSIBILITY /10

The prediction is given in 10 states for relative accessibility, with

RELATIVE_ACCESSIBILITY = (PREDICTED_ACC * PREDICTED_ACC)

where PREDICTED_ACC = 0 - 9.



Estimated Accuracy of Prediction
~~~~~~~~~~~~~~~~~~~~~~~~~~~~~~~~

A careful cross validation test on some 238 protein chains (in total
about 62,000 residues) with less than 25% pairwise sequence identity
gave the following results:


Correlation
...........

The correlation between observed and predicted solvent accessibility
is:

-----------
corr = 0.53
-----------

This value ought to be compared to the worst and best case prediction
scenario: random prediction (corr = 0.0) and homology modelling
(corr = 0.66).  (Note: homology modelling yields a relative accurate
prediction in 3D if, and only if, a significantly identical sequence
has a known 3D structure.)


3-state accuracy
................

Often the relative accessibility is projected onto, e.g., 3 states:
  b  = buried       (here defined as < 9% relative accessibility),
  i  = intermediate ( 9% <= rel. acc. < 36% ),
  e  = exposed      ( rel. acc. >= 36% ).

A projection onto 3 states or 2 states (buried/exposed) enables the
compilation of a 3- and 2-state prediction accuracy.  PHD reaches an
overall 3-state accuracy of:
  Q3 = 57.5%
(compared to 35% for random prediction and 70% for homology modelling).

In detail:

+-----------------------------------+-------------------------+
| Qburied       (% of observed)=77% | Qb (% of predicted)=60% |
| Qintermediate (% of observed)= 9% | Qi (% of predicted)=44% |
| Qexposed      (% of observed)=78% | Qe (% of predicted)=56% |
+-----------------------------------+-------------------------+


10-state accuracy
.................

The network predicts relative solvent accessibility in 10 states, with
state i (i = 0-9) corresponding to a relative solvent accessibility of
i*i %.  The 10-state accuracy of the network is:

  Q10 = 24.5%

..........................................................................

These percentages are defined by:
~~~~~~~~~~~~~~~~~~~~~~~~~~~~~~~~

|                     number of correctly predicted residues
|Q3               = ---------------------------------------      (*100)
|                           number of all residues
|
|                     no of res. correctly predicted to be buried
|Qburied (% of obs) = ------------------------------------------- (*100)
|                     no of all res. observed to be buried
|
|
|                     no of res. correctly predicted to be buried
|Qburied (% of pred)= ------------------------------------------- (*100)
|                     no of all residues predicted to be buried

..........................................................................

Averaging over single chains
~~~~~~~~~~~~~~~~~~~~~~~~~~~

The most reasonable way to compute the overall accuracies is the above
quoted percentage of correctly predicted residues.  However, since the
user is mainly interested in the expected performance of the prediction
for a particular protein, the mean value when averaging over protein
chains might be of help as well.  Computing first the correlation
between observed and predicted accessibility for each protein chan, and
then averaging over all 238 chains yields the following average:

+-------------------------------====--+
| corr/averaged over chains   = 0.53  |
+-------------------------------====--+
| standard deviation          = 0.11  |
+-------------------------------------+

..........................................................................

Further details of performance accuracy
~~~~~~~~~~~~~~~~~~~~~~~~~~~~~~~~~~~~~~

The accuracy matrix in detail:
..............................

-------+----------------------------------------------------+-----------
\ PHD |    0    1   2   3    4    5     6     7    8    9  |  SUM  %obs
-------+----------------------------------------------------+-----------
OBS  0 | 8611  140   8  44   82  169   772   334   27    0  | 10187 16.6
OBS  1 | 4367  164   0  50  106  231   738   346   44    3  |  6049  9.8
OBS  2 | 3194  168   1  68  125  303   951   513   42    7  |  5372  8.7
OBS  3 | 2760  159   8  80  136  327  1246   746   58   19  |  5539  9.0
OBS  4 | 2312  144   2  72  166  396  1615  1245  124   19  |  6095  9.9
OBS  5 | 1873   96   3  84  138  425  1979  1834  187   27  |  6646 10.8
OBS  6 | 1387   67   1  60   80  278  2237  2627  231   51  |  7019 11.4
OBS  7 | 1082   35   0  32   56  225  1871  3107  302   60  |  6770 11.0
OBS  8 |  660   25   0  27   43  136  1206  2374  325   87  |  4883  7.9
OBS  9 |  325   20   2  27   29   74   648  1159  366  214  |  2864  4.7
-------+----------------------------------------------------+-----------
SUM    |26571 1018  25 544  961 2564 13263 14285 1706  487  |
%pred  | 43.3  1.7 0.0 0.9  1.6  4.2  21.6  23.3  2.8  0.8  |
-------+----------------------------------------------------+-----------

Note: This table is to be read in the following manner:
     8611 of all residues predicted to be in exposed by 0%, were
     observed with 0% relative accessibility.  However, 325 of all
     residues predicted to have 0% are observed as completely exposed
     (obs = 9 -> rel. acc. >= 81%).  The term "observed" refers to the
     DSSP compilation of area of solvent accessibility calculated from
     3D coordinates of experimentally determined structures (Diction-
     ary of Secondary Structure  of Proteins: Kabsch & Sander (1983)
     Biopolymers, 22, 2577-2637).


Accuracy for each amino acid:
.............................

+---+------------------------------+-----+-------+------+
|AA |   Q3 b%o b%p i%o i%p e%o e%p | Q10 |  corr |    N |
+---+------------------------------+-----+-------+------+
| A | 59.0  87  60   2  38  66  57 |  31 | 0.530 | 5054 |
| C | 62.0  91  67   5  39  25  21 |  34 | 0.244 |  893 |
| D | 56.5  21  45   6  49  94  57 |  20 | 0.321 | 3536 |
| E | 60.8   9  40   3  41  98  61 |  21 | 0.347 | 3743 |
| F | 63.3  94  67   9  46  29  37 |  27 | 0.366 | 2436 |
| G | 52.1  75  51   1  31  67  53 |  22 | 0.405 | 4787 |
| H | 50.9  63  53  23  45  71  50 |  18 | 0.442 | 1366 |
| I | 64.9  95  68   6  41  30  38 |  34 | 0.360 | 3437 |
| K | 66.6   2  11   2  37  98  67 |  23 | 0.267 | 3652 |
| L | 61.6  93  65   8  44  31  40 |  31 | 0.368 | 5016 |
| M | 60.1  92  64   5  39  45  44 |  29 | 0.452 | 1371 |
| N | 55.5  45  45   8  38  87  59 |  17 | 0.410 | 2923 |
| P | 53.0  48  48   9  39  83  56 |  18 | 0.364 | 2920 |
| Q | 54.3  27  44   7  44  92  56 |  20 | 0.344 | 2225 |
| R | 49.9  15  47  36  47  76  51 |  18 | 0.372 | 2765 |
| S | 55.6  69  53   3  51  81  56 |  22 | 0.464 | 3981 |
| T | 51.8  61  51   8  38  78  53 |  21 | 0.432 | 3740 |
| V | 61.1  93  65   5  40  39  42 |  34 | 0.418 | 4156 |
| W | 56.2  85  62  20  49  29  27 |  21 | 0.318 |  891 |
| Y | 49.7  73  52  33  49  36  38 |  19 | 0.359 | 2301 |
+---+------------------------------+-----+-------+------+

Abbreviations:

AA:   amino acid in one-letter code
b%o, i%o, e%o:   = Qburied, Qintermediate, Qexposed (% of observed),
     i.e. percentage of correct prediction in each state, see above
b%p, i%p, e%p:   = Qburied, Qintermediate, Qexposed (% of predicted),
     i.e. probability of correct prediction in each state, see above
b%o:  = Qburied (% of observed), see above
Q10:  percentage of correctly predicted residues in each of the 10
     states of predicted relative accessibility.
corr: correlation between predicted and observed rel. acc.
N:    number of residues in data set


Accuracy for different secondary structure:
...........................................

+--------+------------------------------+----+-------+-------+
| type   |   Q3 b%o b%p i%o i%p e%o e%p |Q10 |  corr |     N |
+--------+------------------------------+----+-------+-------+
| helix  | 59.5  79  64   8  44  80  56 | 27 | 0.574 | 20100 |
| strand | 61.3  84  73   9  46  69  37 | 35 | 0.524 | 13356 |
| loop   | 54.4  64  43  11  44  78  61 | 18 | 0.442 | 27968 |
+--------+------------------------------+----+-------+-------+

Abbreviations as before.



Position-specific reliability index
~~~~~~~~~~~~~~~~~~~~~~~~~~~~~~~~~~

The network predicts the 10 states for relative accessibility using real
numbers from the output units. The prediction is assigned by choosing
the maximal unit ("winner takes all").  However, the real numbers
contain additional information.
E.g. the difference between the maximal and the second largest output
unit (with the constraint that the second largest output is compiled
among all units at least 2 positions off the maximal unit) can be used
to derive a "reliability index".  This index is given for each residue
along with the prediction.  The index is scaled to have values between
0 (lowest reliability), and 9 (highest).
The accuracies (Q3, corr, asf.) to be expected for residues with values
above a particular value of the index are given below as well as the
fraction of such residues (%res).:

+---+------------------------------+----+-------+-------+
|RI |   Q3 b%o b%p i%o i%p e%o e%p |Q10 |  corr |  %res |
+---+------------------------------+----+-------+-------+
| 0 | 57.5  77  60   9  44  78  56 | 24 | 0.535 | 100.0 |
| 1 | 59.1  76  63   9  45  82  57 | 25 | 0.560 |  91.2 |
| 2 | 61.7  79  66   4  47  87  58 | 27 | 0.594 |  77.1 |
| 3 | 66.6  87  70   1  51  89  63 | 30 | 0.650 |  57.1 |
| 4 | 70.0  89  72   0  83  91  67 | 32 | 0.686 |  45.8 |
| 5 | 72.9  92  75   0   0  93  70 | 34 | 0.722 |  35.6 |
| 6 | 76.3  95  77   0   0  93  75 | 36 | 0.769 |  24.7 |
| 7 | 79.0  97  79   0   0  93  78 | 39 | 0.803 |  16.0 |
| 8 | 80.9  98  80   0   0  91  81 | 43 | 0.824 |   9.6 |
| 9 | 81.2  99  80   0   0  88  83 | 45 | 0.828 |   5.9 |
+---+------------------------------+----+-------+-------+

Abbreviations as before.

The above table gives the cumulative results, e.g. 45.8% of all
residues have a reliability of at least 4.  The correlation for this
most reliably predicted half of the residues is 0.686, i.e. a value
comparable to what could be expected if homology modelling were
possible.  For this subset of 45.8% of all residues, 89% of the buried
residues are correctly predicted, and 72% of all residues predicted to
be buried are correct.

..........................................................................

The following table gives the non-cumulative quantities, i.e. the
values per reliability index range.  These numbers answer the question:
how reliable is the prediction for all residues labeled with the
particular index i.

+---+------------------------------+----+-------+-------+
|RI |   Q3 b%o b%p i%o i%p e%o e%p |Q10 |  corr |  %res |
+---+------------------------------+----+-------+-------+
| 0 | 40.9  79  40  16  41  21  40 | 14 | 0.175 |   8.8 |
| 1 | 45.4  61  46  28  44  48  44 | 17 | 0.278 |  14.1 |
| 2 | 47.4  53  52  10  46  80  44 | 19 | 0.343 |  19.9 |
| 3 | 52.9  75  59   4  50  77  47 | 23 | 0.439 |  11.4 |
| 4 | 60.0  81  63   0  83  84  56 | 25 | 0.547 |  10.1 |
| 5 | 65.2  82  70   0   0  93  62 | 28 | 0.607 |  10.9 |
| 6 | 71.3  90  72   0   0  94  70 | 31 | 0.692 |   8.8 |
| 7 | 76.0  94  76   0   0  95  75 | 34 | 0.762 |   6.3 |
| 8 | 80.5  97  81   0   0  94  79 | 39 | 0.808 |   3.8 |
| 9 | 81.2  99  80   0   0  88  83 | 45 | 0.828 |   5.9 |
+---+------------------------------+----+-------+-------+

For example, for residues with RI = 4 83% of all predicted intermediate
residues are correctly predicted as such.





   ~~~~~~~~~~~~~~~~~~~~~~~~~~~~~~~~~~~~~~~~~~~~~~~~~~~~~~~
   Prediction of helical transmembrane segments by PHDhtm:
   ~~~~~~~~~~~~~~~~~~~~~~~~~~~~~~~~~~~~~~~~~~~~~~~~~~~~~~~

   Author:             Burkhard Rost        
                       EMBL, Heidelberg, FRG
                       Meyerhofstrasse 1, 69 117 Heidelberg
                       Internet: Rost@EMBL-Heidelberg.DE        

   All rights reserved.




About the network method
~~~~~~~~~~~~~~~~~~~~~~~

The PHD mail server is described in:
  Rost, Burkhard; Sander, Chris; Schneider, Reinhard:
  PHD - an automatic mail server for protein secondary structure
  prediction.
  CABIOS, 1994, 10, 53-60.      

To be quoted for publications of PHDhtm output:
  Rost, Burkhard; Casadio, Rita; Fariselli, Piero; Sander, Chris:
  Prediction of helical transmembrane segments at 95% accuracy.
  Protein Science, 1995, 4, 521-533.                



Estimated Accuracy of Prediction
~~~~~~~~~~~~~~~~~~~~~~~~~~~~~~~~

A cross validation test on 69 helical trans-membrane  proteins (in total
about 30,000 residues) with less than 25% pairwise sequence identity
gave the following results:

++================++-----------------------------------------+
|| Qtotal = 94.7% ||      ("overall two state accuracy")     |
++================++-----------------------------------------+

+----------------------------+-----------------------------+
| Qhelix (% of observed)=92% | Qhelix (% of predicted)=83% |
| Qloop  (% of observed)=96% | Qloop  (% of predicted)=97% |
+----------------------------+-----------------------------+

..........................................................................

These percentages are defined by:
~~~~~~~~~~~~~~~~~~~~~~~~~~~~~~~~

|                    number of correctly predicted residues
|Qtotal =            ---------------------------------------      (*100)
|                          number of all residues
|
|                    no of res correctly predicted to be in helix
|Qhelix (% of obs) = -------------------------------------------- (*100)
|                    no of all res observed to be in helix
|
|
|                    no of res correctly predicted to be in helix
|Qhelix (% of pred)= -------------------------------------------- (*100)
|                    no of all residues predicted to be in helix

..........................................................................

Further measures of performance
~~~~~~~~~~~~~~~~~~~~~~~~~~~~~~

Matthews correlation coefficient:

+---------------------------------------------+
| Chelix = 0.84, Cloop = 0.84                 |
+---------------------------------------------+
..........................................................................

Average length of predicted secondary structure segments:

|           +------------+----------+
|           |  predicted | observed |
+-----------+------------+----------+
| Lhelix  = |    24.6    |   22.2   |
+-----------+------------+----------+
..........................................................................

The accuracy matrix in detail:

+---------------------------------+
|    number of residues with H, L |
+---------+------+-------+--------+
|         |net H | net L |sum obs |
+---------+------+-------+--------+
| obs H   | 5214 |   492 |   5706 |
| obs L   | 1050 | 22423 |  23473 |
+---------+------+-------+--------+
| sum Net | 6264 | 22915 |  29179 |
+---------+------+-------+--------+

Note: This table is to be read in the following manner:
     5214 of all residues predicted to be in a helical trans-membrane
     region, were observed to be in the lipid bilayer, 1050 however
     were observed either inside or outside of the protein, i.e. in
     loop (or non-membrane) regions. The term "observed" refers to DSSP
     assignment of secondary structure calculated from 3D coordinates
     of experimentally determined structures (Dictionary of Secondary
     Structure  of Proteins: Kabsch & Sander (1983) Biopolymers, 22,
     2577-2637) where these were available.  For all other proteins,
     the assignment of trans-membrane segments has been taken from the
     Swissprot data bank (Bairoch, A.; Boeckmann, B.: The SWISS-PROT
     protein sequence data bank. Nucl. Acids Res. 20: 2019-2022, 1992).

..........................................................................

Overlap between predicted and observed segments:

+-----------------+---------------+----------------+
| segment overlap | % of observed | % of predicted |
|   Sov helix     |      95.6%    |      95.5%     |
|   Sov loop      |      83.6%    |      97.2%     |
+-----------------+---------------+----------------+
|   Sov total     |      86.0%    |      96.8%     |
+-----------------+---------------+----------------+

     Definition of Sov in: Rost et al., JMB, 1994, 235, 13-26.

     As helical trans-membrane segments are longer than globular heli-
     ces, correctly predicted segments can easily be made out.  PHDhtm
     misses 5 out of 258 observed segments, predicts 6 where non is
     observed and 3 times the predicted helical segment overlaps two
     observed regions.  Thus, in total more than 95% of all segments
     are correctly predicted.

..........................................................................

Entropy of prediction (information measure):

+-----------------+
| I = 0.64        |
+-----------------+

     (For comparison: homology modelling of globular proteins in three
     states: I=0.62.)
     Definition of Sov in: Rost et al., JMB, 1994, 235, 13-26.



Position-specific reliability index
~~~~~~~~~~~~~~~~~~~~~~~~~~~~~~~~~~

The network predicts two states: helical trans-membrane region and rest
using two output units.  The prediction is assigned by choosing the ma-
ximal unit ("winner takes all").  However, the real numbers of the out-
put units contain additional information.
E.g. the difference between the two output units can be used to derive
a "reliability index".  This index is given for each residue along with
the prediction.  The index is scaled to have values between 0 (lowest
reliability), and 9 (highest).
The accuracies (Qtot) to be expected for residues with values above a
particular value of the index are given below as well as the fraction
of such residues (%res).:

+------+-----+-----+-----+-----+-----+-----+-----+-----+-----+-----+
| index|  0  |  1  |  2  |  3  |  4  |  5  |  6  |  7  |  8  |  9  |
| %res |100.0| 98.8| 97.3| 95.9| 94.1| 92.3| 89.9| 86.2| 75.0| 66.8|
+------+-----+-----+-----+-----+-----+-----+-----+-----+-----+-----+
|      |     |     |     |     |     |     |     |     |     |     |
| Qtot | 94.7| 95.2| 95.6| 96.2| 96.7| 97.2| 97.7| 98.4| 99.4| 99.8|
|      |     |     |     |     |     |     |     |     |     |     |
+------+-----+-----+-----+-----+-----+-----+-----+-----+-----+-----+
| H%obs| 91.8| 92.9| 93.8| 94.4| 95.0| 95.7| 96.2| 96.8| 95.5| 78.7|
| L%obs| 95.3| 95.7| 96.1| 96.6| 97.0| 97.5| 98.1| 98.8| 99.7|100.0|
|      |     |     |     |     |     |     |     |     |     |     |
| H%prd| 82.7| 83.8| 85.0| 86.7| 88.1| 89.7| 91.4| 93.8| 96.3| 97.1|
| L%prd| 97.9| 98.3| 98.5| 98.7| 98.8| 99.0| 99.2| 99.4| 99.7| 99.9|
+------+-----+-----+-----+-----+-----+-----+-----+-----+-----+-----+

The above table gives the cumulative results, e.g. 92.3% of all
residues have a reliability of at least 5.  The overall two-state
accuracy for this subset is 97.2%.  For this subset, e.g., 95.7% of
the observed helical trans-membrane residues are correctly predicted,
and 89.7% of all residues predicted to be in helical trans-membrane
segment are correct.







The resulting network (PHD) prediction is:
~~~~~~~~~~~~~~~~~~~~~~~~~~~~~~~~~~~~~~~~~

________________________________________________________________________________



 PHD: Profile fed neural network systems from HeiDelberg
 ~~~~~~~~~~~~~~~~~~~~~~~~~~~~~~~~~~~~~~~~~~~~~~~~~~~~~~~

 Prediction of:         
    secondary structure,               by PHDsec        
    solvent accessibility,             by PHDacc        
    and helical transmembrane regions,     by PHDhtm        

 Author:                                    
    Burkhard Rost                           
    EMBL, 69012 Heidelberg, Germany                 
    Internet: Rost@EMBL-Heidelberg.DE               

 All rights reserved.



 The network systems are described in:                              

 PHDsec:    B Rost & C Sander: JMB, 1993, 232, 584-599.     
        B Rost & C Sander: Proteins, 1994, 19, 55-72.       
 PHDacc:    B Rost & C Sander: Proteins, 1994, 20, 216-226.     
 PHDhtm:    B Rost et al.:     Prot. Science, 1995, 4, 521-533. 



 Some statistics
 ~~~~~~~~~~~~~~~

 Percentage of amino acids:
 +--------------+--------+--------+--------+--------+--------+
 | AA:          |    L   |    A   |    S   |    G   |    I   |
 | % of AA:     |   13.0 |   10.0 |    9.7 |    8.9 |    8.6 |
 +--------------+--------+--------+--------+--------+--------+
 | AA:          |    V   |    R   |    T   |    F   |    D   |
 | % of AA:     |    7.8 |    5.2 |    4.5 |    4.5 |    4.5 |
 +--------------+--------+--------+--------+--------+--------+
 | AA:          |    N   |    Q   |    E   |    P   |    K   |
 | % of AA:     |    4.1 |    3.0 |    3.0 |    2.6 |    2.6 |
 +--------------+--------+--------+--------+--------+--------+
 | AA:          |    Y   |    M   |    W   |    H   |    C   |
 | % of AA:     |    1.9 |    1.9 |    1.5 |    1.5 |    1.5 |
 +--------------+--------+--------+--------+--------+--------+

 Percentage of secondary structure predicted:
 +--------------+--------+--------+--------+
 | SecStr:      |    H   |    E   |    L   |
 | % Predicted: |   43.9 |   16.7 |   39.4 |
 +--------------+--------+--------+--------+

 According to the following classes:
    all-alpha:   %H>45 and %E< 5; all-beta : %H<5 and %E>45
    alpha-beta : %H>30 and %E>20; mixed:    rest,
 this means that the predicted class is:           mixed class



 PHD output for your protein
 ~~~~~~~~~~~~~~~~~~~~~~~~~~~

 Tue Nov 24 17:44:57 1998
 Jury on:       10    different architectures (version   5.94_317 ).
 Note: differently trained architectures, i.e., different versions can
 result in different predictions.



 About the protein
 ~~~~~~~~~~~~~~~~~

 HEADER     /home/phd/server/work/predict_h25873-220
 COMPND
 SOURCE
 AUTHOR
 SEQLENGTH   269
 NCHAIN        1 chain(s) in predict_h25873-22040 data set
 NALIGN       48
 (=number of aligned sequences in HSSP file)



 Abbreviations: PHDsec
 ~~~~~~~~~~~~~~~~~~~~~

 sequence:
    AA : amino acid sequence
 secondary structure:
    HEL: H=helix, E=extended (sheet), blank=other (loop)
    PHD: Profile network prediction HeiDelberg
    Rel: Reliability index of prediction (0-9)
 detail:
    prH: 'probability' for assigning helix
    prE: 'probability' for assigning strand
    prL: 'probability' for assigning loop
         note: the 'probabilites' are scaled to the interval 0-9, e.g.,
               prH=5 means, that the first output node is 0.5-0.6
 subset:
    SUB: a subset of the prediction, for all residues with an expected
         average accuracy > 82% (tables in header)
         note: for this subset the following symbols are used:
      L: is loop (for which above " " is used)
    ".": means that no prediction is made for this residue, as the
         reliability is:  Rel < 5

 Abbreviations: PHDacc
 ~~~~~~~~~~~~~~~~~~~~~

    SS : secondary structure
    HEL: H=helix, E=extended (sheet), blank=other (loop)
 solvent accessibility:
    3st: relative solvent accessibility (acc) in 3 states:
         b = 0-9%, i = 9-36%, e = 36-100%.
    PHD: Profile network prediction HeiDelberg
    Rel: Reliability index of prediction (0-9)
    O_3: observed relative acc. in 3 states: B, I, E
         note: for convenience a blank is used intermediate (i).
    P_3: predicted relative accessibility in 3 states
    10st:relative accessibility in 10 states:
         = n corresponds to a relative acc. of n*n %
 subset:
    SUB: a subset of the prediction, for all residues with an expected
         average correlation > 0.69 (tables in header)
         note: for this subset the following symbols are used:
    "I": is intermediate (for which above " " is used)
    ".": means that no prediction is made for this residue, as the
         reliability is: Rel < 4


 Abbreviations: PHDhtm
 ~~~~~~~~~~~~~~~~~~~~~

 secondary structure:
    HL:  T=helical transmembrane region, blank=other (loop)
    PHD: Profile network prediction HeiDelberg
    PHDF:filtered prediction, i.e., too long transmembrane segments
         are split, too short ones are deleted
    Rel: Reliability index of prediction (0-9)
 detail:
    prH: 'probability' for assigning helical transmembrane region
    prL: 'probability' for assigning loop
         note: the 'probabilites' are scaled to the interval 0-9, e.g.,
               prH=5 means, that the first output node is 0.5-0.6
 subset:
    SUB: a subset of the prediction, for all residues with an expected
         average accuracy > 82% (tables in header)
         note: for this subset the following symbols are used:
      L: is loop (for which above " " is used)
    ".": means that no prediction is made for this residue, as the
         reliability is:  Rel < 5



 protein:       predict        length      269

                  ....,....1....,....2....,....3....,....4....,....5....,....6
         AA      |MASEIKKKLFWRAVVAEFLAMTLFVFISIGSALGFNYPLERNQTLVQDNVKVSLAFGLSI|
         PHD sec |       HHHHHHHHHHHHHHHHHHHHHHHHHHEE            HHHHHHHHHHHHH|
         Rel sec |998443148899999999999998997676530312469989998623353579999999|
 detail:
         prH sec |001223468899999999999998888777653112210000000145566788999999|
         prE sec |000011000000000000000001001111233542100000000000323211000000|
         prL sec |998665420100000000000000000011112244578988998753100000000000|
 subset: SUB sec |LLL.....HHHHHHHHHHHHHHHHHHHHHHH......LLLLLLLLL...H.HHHHHHHHH|

 ACCESSIBILITY
 3st:    P_3 acc |eeeebee bbb bbbbbbbbbbbbbbbbbbbbbebeee eeeeeeeeebbbbbbbbbbbb|
 10st:   PHD acc |997706650005000000000000000000000607775779776677000000000000|
         Rel acc |735421110541467608662789996343122133420454330023453975664547|
 subset: SUB acc |e.ee.....bb.bbbb.bbb.bbbbbb.b.......e..eee......bb.bbbbbbbbb|
                  ....,....7....,....8....,....9....,....10...,....11...,....12
         AA      |ATLAQSVGHISGAHSNPAVTLGLLLSCQISILRAVMYIIAQCVGAIVASAILSGITSSLL|
         PHD sec |HHHHHHHHHE      HHHHEHHHHHHHHHHHHHHHHHHHHHHHHHHHHHHHHHH     |
         Rel sec |999996412122653167703135552356779999999999999999999998467213|
 detail:
         prH sec |998986544334223477843456665567779999999999999999999998611343|
         prE sec |001001123420010000145432101221110000000000000000000000000000|
         prL sec |000001232245765521000000123210000000000000000000000000278555|
 subset: SUB sec |HHHHHH......LL..HHH....HHH..HHHHHHHHHHHHHHHHHHHHHHHHHH.LL...|

 ACCESSIBILITY
 3st:    P_3 acc |bbbbebbbebbbbbb bbbbbbbbbbbebbbbbbbbbbbbbbbbbbbbbbbbeebbeeeb|
 10st:   PHD acc |000060006000000500000000000600000000000000000000000067006760|
         Rel acc |456515321655013144869663400154551757478936465465467713401400|
 subset: SUB acc |bbbb.b...bbb....bbbbbbb.b...bbbb.bbbbbbb.bbbbbbbbbbb..b..e..|
                  ....,....13...,....14...,....15...,....16...,....17...,....18
         AA      |ENSLGRNDLARGVNSGQGLGIEIIGTLQLVLCVLATTDRRRRDLGGSAPLAIGLSVALGH|
         PHD sec |       HHH       EEEEEEEEEEEEEEEEEEE            E E   HHHHHH|
         Rel sec |359985212134223651899898866789799875436658889963211351457756|
 detail:
         prH sec |320002345432332111000000000000100000221120000000001113567767|
         prE sec |100000000000011014899888877789789886100000000013544222221111|
         prL sec |568986543466545763100000011100000112567768889975454564210111|
 subset: SUB sec |.LLLLL.........LL.EEEEEEEEEEEEEEEEEE..LLLLLLLLL.....L..HHHHH|

 ACCESSIBILITY
 3st:    P_3 acc |eeebbbebbbeebeebeebbbbbbbbbbbbbbbbbbbeeeeeeeebbbbbbbbbbbbbbb|
 10st:   PHD acc |677000600077076077000000000000000000077767767000000000000000|
         Rel acc |133100124043040233247198656399879530035414413123255869586654|
 subset: SUB acc |........b.e..e.....bb.bbbbb.bbbbbb....ee.ee......bbbbbbbbbbb|
                  ....,....19...,....20...,....21...,....22...,....23...,....24
         AA      |LLAIDYTGCGINPARSFGSAVLTRNFSNHWIFWVGPFIGSALAVLIYDFILAPRSSDFTD|
         PHD sec |HEEEE E          HHHEEEE    EEEEEE HHHHHHHHHHHHHEEEEE       |
         Rel sec |321341126989622145152653534229996251699999999973147525556642|
 detail:
         prH sec |521100000000145432463121122000000114789999999875421111121124|
         prE sec |244564431000000000015765121358997510000000000013467642110000|
         prL sec |233234457889754567411012655530002364200000000010010136667765|
 subset: SUB sec |........LLLLL....H.H.EE.L....EEEE.L.HHHHHHHHHHH...EE.LLLLL..|

 ACCESSIBILITY
 3st:    P_3 acc |bbbbebbbbbbebb bbbbbbbbeebeebbbbbbbbbbbbbbbbbbbbbbbbeeeee ee|
 10st:   PHD acc |000060000006005000000007606600000000000000000000000076777577|
         Rel acc |754424240102242141047612131118967874356346635751777031345044|
 subset: SUB acc |bbbb.b.b.....b..b..bbb.......bbbbbbb.bb.bbb.bbb.bbb....ee.ee|
                  ....,....25...,....26...,....27...,....28...,....29...,....30
         AA      |RMKVWTSGQVEEYDLDADDINSRVEMKPK|
         PHD sec |HHHHHH                       |
         Rel sec |66775259975467555457776422699|
 detail:
         prH sec |77887520012221222221111100000|
         prE sec |00000000000000000000001233200|
         prL sec |11112379987678777678887655799|
 subset: SUB sec |HHHHH.LLLLL.LLLLL.LLLLL...LLL|

 ACCESSIBILITY
 3st:    P_3 acc |ebebbeeeeeeeeeeeeeeeeeebeeeee|
 10st:   PHD acc |60700787677777677777767067789|
         Rel acc |10411563134335144444514212559|
 subset: SUB acc |..e..ee...e..e.eeeeee.e...eee|


 PHDhtm Helical transmembrane prediction
        note: PHDacc and PHDsec are reliable for water-
              soluble globular proteins, only.  Thus,
              please take the  predictions above with
              particular caution wherever transmembrane
              helices are predicted by PHDhtm!


 PHDhtm
---
--- PhdTopology REFINEMENT AND TOPOLOGY PREDICTION: SYMBOLS
--- AA           : amino acid in one-letter code
--- PHD htm      : HTM's predicted by the PHD neural network
---                system (T=HTM, ' '=not HTM)
--- Rel htm      : Reliability index of prediction (0-9, 0 is low)
--- detail       : Neural network output in detail
--- prH htm      : 'Probability' for assigning a helical trans-
---                membrane region (HTM)
--- prL htm      : 'Probability' for assigning a non-HTM region
---          note: 'Probabilites' are scaled to the interval
---                0-9, e.g., prH=5 means, that the first
---                output node is 0.5-0.6
--- subset       : Subset of more reliable predictions
--- SUB htm      : All residues for which the expected average
---                accuracy is > 82% (tables in header).
---          note: for this subset the following symbols are used:
---             L: is loop (for which above ' ' is used)
---           '.': means that no prediction is made for this,
---                residue as the reliability is:  Rel < 5
--- other        : predictions derived based on PHDhtm
--- PHDFhtm      : filtered prediction, i.e., too long HTM's are
---                split, too short ones are deleted
--- PHDRhtm      : refinement of neural network output
--- PHDThtm      : topology prediction based on refined model
---                symbols used:
---             i: intra-cytoplasmic
---             T: transmembrane region
---             o: extra-cytoplasmic
---
--- PhdTopology REFINEMENT AND TOPOLOGY PREDICTION
                  ....,....1....,....2....,....3....,....4....,....5....,....6
         AA      |MASEIKKKLFWRAVVAEFLAMTLFVFISIGSALGFNYPLERNQTLVQDNVKVSLAFGLSI|
         PHD htm |              TTTTTTTTTTTTTTTTTTT               TTTTTTTTTTTT|
 detail:         |                                                            |
         prH htm |000000000001136788999999999988875321110000000123678889999988|
         prL htm |999999999998863211000000000011124678889999999876321110000011|
  other:         |                                                            |
         PHDFhtm |              TTTTTTTTTTTTTTTTTTT                TTTTTTTTTTT|
         PHDRhtm |              TTTTTTTTTTTTTTTTTT                 TTTTTTTTTTT|
         PHDThtm |iiiiiiiiiiiiiiTTTTTTTTTTTTTTTTTToooooooooooooooooTTTTTTTTTTT|
 subset:         |                                                            |
         SUB htm |............................................................|
                  ....,....7....,....8....,....9....,....10...,....11...,....12
         AA      |ATLAQSVGHISGAHSNPAVTLGLLLSCQISILRAVMYIIAQCVGAIVASAILSGITSSLL|
         PHD htm |TTTTTTTTTTTTTTTTTTTTTTTTTTTTTTTTTTTTTTTTTTTTTTTTTTTTTTT     |
 detail:         |                                                            |
         prH htm |888888877777666677788888888888888888888888888888888876543211|
         prL htm |111111122222333322211111111111111111111111111111111123456788|
  other:         |                                                            |
         PHDFhtm |TTTTTTTTTTTTTTTT    TTTTTTTTTTTTTTTTTTTTTTTTTTTTTTTTT       |
         PHDRhtm |TTTTTTTT             TTTTTTTTTTTTTTTTTTTTTTTTT              |
         PHDThtm |TTTTTTTTiiiiiiiiiiiiiTTTTTTTTTTTTTTTTTTTTTTTTToooooooooooooo|
 subset:         |                                                            |
         SUB htm |............................................................|
                  ....,....13...,....14...,....15...,....16...,....17...,....18
         AA      |ENSLGRNDLARGVNSGQGLGIEIIGTLQLVLCVLATTDRRRRDLGGSAPLAIGLSVALGH|
         PHD htm |               TTTTTTTTTTTTTTTTTTT             TTTTTTTTTTTTT|
 detail:         |                                                            |
         prH htm |000000000001234567788888999988887643211111111235788899998888|
         prL htm |999999999998765432211111000011112356788888888764211100001111|
  other:         |                                                            |
         PHDFhtm |               TTTTTTTTTTTTTTTTTTT             TTTTTTTTTTTTT|
         PHDRhtm |                TTTTTTTTTTTTTTTTTT              TTTTTTTTTTTT|
         PHDThtm |ooooooooooooooooTTTTTTTTTTTTTTTTTTiiiiiiiiiiiiiiTTTTTTTTTTTT|
 subset:         |                                                            |
         SUB htm |............................................................|
                  ....,....19...,....20...,....21...,....22...,....23...,....24
         AA      |LLAIDYTGCGINPARSFGSAVLTRNFSNHWIFWVGPFIGSALAVLIYDFILAPRSSDFTD|
         PHD htm |TTTTTTTTT            TTTTTTTTTTTTTTTTTTTTTTTTTTTTTTT        |
 detail:         |                                                            |
         prH htm |888887765443432233334566777777788888888888888888887542100000|
         prL htm |111112234556567766665433222222211111111111111111112457899999|
  other:         |                                                            |
         PHDFhtm |TTTTTTTTT            TTTTTTTTTTTTTTTTTTTTTTTTTTTTTTT        |
         PHDRhtm |TTTTTT                         TTTTTTTTTTTTTTTTTTT          |
         PHDThtm |TTTTTToooooooooooooooooooooooooTTTTTTTTTTTTTTTTTTTiiiiiiiiii|
 subset:         |                                                            |
         SUB htm |............................................................|
                  ....,....25...,....26...,....27...,....28...,....29...,....30
         AA      |RMKVWTSGQVEEYDLDADDINSRVEMKPK|
         PHD htm |                             |
 detail:         |                             |
         prH htm |00000000000000000000000000000|
         prL htm |99999999999999999999999999999|
  other:         |                             |
         PHDFhtm |                             |
         PHDRhtm |                             |
         PHDThtm |iiiiiiiiiiiiiiiiiiiiiiiiiiiii|
 subset:         |                             |
         SUB htm |.............................|
---
--- PhdTopology REFINEMENT AND TOPOLOGY PREDICTION END
---

________________________________________________________________________________



________________________________________________________________________________


-----------------------------------------------------------------------------
---   PredictProtein: NEWS from January, 1997                             ---
---                                                                       ---
---   Dear user,                                                          ---
---                                                                       ---
---      as of  January 1, 1997,  EMBL has effectively decided to not     ---
---   support the PredictProtein service by personal resources.  I do     ---
---   maintain the program, so to speak, in my private time. However,     ---
---   my contract obliges me to do science, instead.   Unfortunately,     ---
---   the computer environment at  EMBL  is at the same time starting     ---
---   to become increasingly unstable.   Consequence of these two re-     ---
---   cent developments is that the  PredictProtein service is not as     ---
---   stable as it was.                                                   ---
---                                                                       ---
---      I apologise for the problems this may cause.  In particular,     ---
---   I apologise for my inability to reply to the 20-30 daily,  per-     ---
---   sonal mails, and suggest to  re-submit requests after 24 hours!     ---
---                                                                       ---
---   Hoping  that I shall  find  a more convenient  solution for the     ---
---   future of the PredictProtein I remain with my best regards,         ---
---                                                                       ---
---   Burkhard Rost                                                       ---
-----------------------------------------------------------------------------
---   PredictProtein: NEWS from April, 1998                               ---
---                                                                       ---
--------------------------------                                          ---
---   MOVING PredictProtein                                               ---
---   There appears to be light on the horizon! PP will may be having     ---
---   many hickups over the next months (as I shall leave EMBL). How-     ---
---   ever, the server seems to have a fair chance of survival thanks     ---
---   to a major support that is being raised by Columbia University,     ---
---   New York, U.S.A.).   I hope that this will settle the issue for     ---
---   the years to come ...                                               ---
--------------------------------                                          ---
---   WARNING                                                             ---
---   After a  major  rewriting of most of the PP code over the last,     ---
---   I am afraid  that  not all errors have been traced by me,  yet.     ---
---   Thus, please have mercy and report any bug you'll encounter!        ---
---                                             THANKS, Burkhard Rost     ---
--------------------------------                                          ---
---   NEW PREDICTION DEFAULTS                                             ---
---   * Coiled-coil regions: now by default the program COILS written by  ---
---     Andrei Lupas is run on your sequence. An output is returned if a  ---
---     coiled-coil region has been detected.                             ---
---   * Functional sequence motifs: now by default the PROSITE database   ---
---     written by Amos Bairoch, Philip Bucher and Kay Hofmann is scanned ---
---     for sequence motifs. An output is returned if any motif has been  ---
---     detected.                                                         ---
--------------------------------                                          ---
---   see http://www.embl-heidelberg.de/predictprotein/ppNews.html        ---
---     for a description of the following new options.                   ---
---   NEW INPUT OPTION                                                    ---
---   * Your input sequence(s) in FASTA-list format ("# FASTA list ")     ---
---   NEW OUTPUT OPTIONS                                                  ---
---   * Return also BLASTP output ("return blast")                        ---
---   * Return prediction additionally in RDB format ("return phd rdb")   ---
---   * Return topits hssp  ("return topits hssp")                        ---
---   * Return topits strip ("return topits strip")                       ---
---   * Return topits own   ("return topits own")                         ---
---   * Return no coils     ("return no coils")                           ---
---   * Return no prosite   ("return no prosite")                         ---
-----------------------------------------------------------------------------
%</AQP1PHD>
%    \end{macrocode}
%    \begin{macrocode}
%<*AQPHMMsgl>
>HP: 269 AQP1  IN    6   14    33    54    73    94   112   139   156   165   184   211   230 
>HP: 271 AQP2  IN    6   17    35    44    65    86   104   131   148   157   176   203   224 
>HP: 285 AQP3  IN    6   22    41    50    72   103   122   153   172   185   207   238   260 
>HP: 323 AQP4  IN    6   37    57    70    92   123   147   160   177   186   205   232   254 
>HP: 265 AQP5  IN    6   13    32    59    78    87   110   131   149   158   177   204   228 
%</AQPHMMsgl>
%    \end{macrocode}
%    \begin{macrocode}
%<*AQPHMMext>
Protein: AQP1
Length:  269
N-terminus:  IN 
Number of transmembrane helices: 6
Transmembrane helices: 14-33 54-73 94-112 139-156 165-184 211-230 

Total entropy of the model:  17.0025
Entropy of the best path:  17.0049

The best path:

     seq  MASEIKKKLF WRAVVAEFLA MTLFVFISIG SALGFNYPLE RNQTLVQDNV    50
     pred IIIIiiiiii iiiHHHHHHH HHHHHHHHHH HHHooooooo oooooooooo 

     seq  KVSLAFGLSI ATLAQSVGHI SGAHSNPAVT LGLLLSCQIS ILRAVMYIIA   100
     pred oooHHHHHHH HHHHHHHHHH HHHiiiiiii iiiiiiiiii iiiHHHHHHH 

     seq  QCVGAIVASA ILSGITSSLL ENSLGRNDLA RGVNSGQGLG IEIIGTLQLV   150
     pred HHHHHHHHHH HHoooooooo oooooooooo ooooooooHH HHHHHHHHHH 

     seq  LCVLATTDRR RRDLGGSAPL AIGLSVALGH LLAIDYTGCG INPARSFGSA   200
     pred HHHHHHiiii iiiiHHHHHH HHHHHHHHHH HHHHoooooo oooooooooo 

     seq  VLTRNFSNHW IFWVGPFIGS ALAVLIYDFI LAPRSSDFTD RMKVWTSGQV   250
     pred oooooooooo HHHHHHHHHH HHHHHHHHHH iiiiiiiiii iiiiiIIIII 

     seq  EEYDLDADDI NSRVEMKPK  269
     pred IIIIIIIIII IIIIIIIII

Protein: AQP2
Length:  271
N-terminus:  IN 
Number of transmembrane helices: 6
Transmembrane helices: 17-35 44-65 86-104 131-148 157-176 203-224 

Total entropy of the model:  17.0017
Entropy of the best path:  17.0046

The best path:

     seq  MWELRSIAFS RAVLAEFLAT LLFVFFGLGS ALQWASSPPS VLQIAVAFGL    50
     pred IIIIIIiiii iiiiiiHHHH HHHHHHHHHH HHHHHooooo oooHHHHHHH 

     seq  GIGILVQALG HVSGAHINPA VTVACLVGCH VSFLRAAFYV AAQLLGAVAG   100
     pred HHHHHHHHHH HHHHHiiiii iiiiiiiiii iiiiiHHHHH HHHHHHHHHH 

     seq  AAILHEITPV EIRGDLAVNA LHNNATAGQA VTVELFLTMQ LVLCIFASTD   150
     pred HHHHoooooo oooooooooo oooooooooo HHHHHHHHHH HHHHHHHHii 

     seq  ERRGDNLGSP ALSIGFSVTL GHLLGIYFTG CSMNPARSLA PAVVTGKFDD   200
     pred iiiiiiHHHH HHHHHHHHHH HHHHHHoooo oooooooooo oooooooooo 

     seq  HWVFWIGPLV GAIIGSLLYN YLLFPSAKSL QERLAVLKGL EPDTDWEERE   250
     pred ooHHHHHHHH HHHHHHHHHH HHHHiiiiii iiiiiiiiiI IIIIIIIIII 

     seq  VRRRQSVELH SPQSLPRGSK A  271
     pred IIIIIIIIII IIIIIIIIII I

Protein: AQP3
Length:  285
N-terminus:  IN 
Number of transmembrane helices: 6
Transmembrane helices: 22-41 50-72 103-122 153-172 185-207 238-260 

Total entropy of the model:  17.0059
Entropy of the best path:  17.0075

The best path:

     seq  MNRCGEMLHI RYRLLRQALA ECLGTLILVM FGCGSVAQVV LSRGTHGGFL    50
     pred IIIIIIiiii iiiiiiiiii iHHHHHHHHH HHHHHHHHHH HooooooooH 

     seq  TINLAFGFAV TLAILVAGQV SGAHLNPAVT FAMCFLAREP WIKLPIYTLA   100
     pred HHHHHHHHHH HHHHHHHHHH HHiiiiiiii iiiiiiiiii iiiiiiiiii 

     seq  QTLGAFLGAG IVFGLYYDAI WAFAGNELVV SGPNGTAGIF ATYPSGHLDM   150
     pred iiHHHHHHHH HHHHHHHHHH HHoooooooo oooooooooo oooooooooo 

     seq  VNGFFDQFIG TAALIVCVLA IVDPYNNPVP RGLEAFTVGL VVLVIGTSMG   200
     pred ooHHHHHHHH HHHHHHHHHH HHiiiiiiii iiiiHHHHHH HHHHHHHHHH 

     seq  FNSGYAVNPA RDFGPRLFTA LAGWGSEVFT TGQNWWWVPI VSPLLGSIGG   250
     pred HHHHHHHooo oooooooooo oooooooooo oooooooHHH HHHHHHHHHH 

     seq  VFVYQLMIGC HLEQPPPSTE AENVKLAHMK HKEQI  285
     pred HHHHHHHHHH iiiiiiiiii iiiiiIIIII IIIII

Protein: AQP4
Length:  323
N-terminus:  IN 
Number of transmembrane helices: 6
Transmembrane helices: 37-57 70-92 123-147 160-177 186-205 232-254 

Total entropy of the model:  17.0058
Entropy of the best path:  17.0091

The best path:

     seq  MSDGAAARRW GKCGPPCSRE SIMVAFKGVW TQAFWKAVTA EFLAMLIFVL    50
     pred IIIIIIIIII IIIIIIIIII Iiiiiiiiii iiiiiiHHHH HHHHHHHHHH 

     seq  LSVGSTINWG GSENPLPVDM VLISLCFGLS IATMVQCFGH ISGGHINPAV   100
     pred HHHHHHHooo oooooooooH HHHHHHHHHH HHHHHHHHHH HHiiiiiiii 

     seq  TVAMVCTRKI SIAKSVFYIT AQCLGAIIGA GILYLVTPPS VVGGLGVTTV   150
     pred iiiiiiiiii iiiiiiiiii iiHHHHHHHH HHHHHHHHHH HHHHHHHooo 

     seq  HGNLTAGHGL LVELIITFQL VFTIFASCDS KRTDVTGSVA LAIGFSVAIG   200
     pred oooooooooH HHHHHHHHHH HHHHHHHiii iiiiiHHHHH HHHHHHHHHH 

     seq  HLFAINYTGA SMNPARSFGP AVIMGNWENH WIYWVGPIIG AVLAGALYEY   250
     pred HHHHHooooo oooooooooo oooooooooo oHHHHHHHHH HHHHHHHHHH 

     seq  VFCPDVELKR RLKEAFSKAA QQTKGSYMEV EDNRSQVETE DLILKPGVVH   300
     pred HHHHiiiiii iiiiiiiiiI IIIIIIIIII IIIIIIIIII IIIIIIIIII 

     seq  VIDIDRGDEK KGKDSSGEVL SSV  323
     pred IIIIIIIIII IIIIIIIIII III

Protein: AQP5
Length:  265
N-terminus:  IN 
Number of transmembrane helices: 6
Transmembrane helices: 13-32 59-78 87-110 131-149 158-177 204-228 

Total entropy of the model:  17.0020
Entropy of the best path:  17.0052

The best path:

     seq  MKKEVCSLAF FKAVFAEFLA TLIFVFFGLG SALKWPSALP TILQISIAFG    50
     pred IIIIIIIIii iiHHHHHHHH HHHHHHHHHH HHoooooooo oooooooooo 

     seq  LAIGTLAQAL GPVSGGHINP AITLALLIGN QISLLRAVFY VAAQLVGAIA   100
     pred ooooooooHH HHHHHHHHHH HHHHHHHHii iiiiiiHHHH HHHHHHHHHH 

     seq  GAGILYWLAP LNARGNLAVN ALNNNTTPGK AMVVELILTF QLALCIFSST   150
     pred HHHHHHHHHH oooooooooo oooooooooo HHHHHHHHHH HHHHHHHHHi 

     seq  DSRRTSPVGS PALSIGLSVT LGHLVGIYFT GCSMNPARSF GPAVVMNRFS   200
     pred iiiiiiiHHH HHHHHHHHHH HHHHHHHooo oooooooooo oooooooooo 

     seq  PSHWVFWVGP IVGAMLAAIL YFYLLFPSSL SLHDRVAVVK GTYEPEEDWE   250
     pred oooHHHHHHH HHHHHHHHHH HHHHHHHHii iiiiiiiiii iiiIIIIIII 

     seq  DHREERKKTI ELTAH  265
     pred IIIIIIIIII IIIII
%</AQPHMMext>
%    \end{macrocode}
%    \begin{macrocode}
%<*TCoffee>
T-COFFEE, Version_5.31(Fri Oct 26 17:01:36 2007)
Cedric Notredame 
CPU TIME:1 sec.
SCORE=54
*
 BAD AVG GOOD
*
AQP1.PRO   :  55
AQP2.PRO   :  58
AQP3.PRO   :  46
AQP4.PRO   :  55
AQP5.PRO   :  58
cons       :  54

AQP1.PRO       1 65------------------------3777666677777766666666556666666554   36 
AQP2.PRO       1 -2------------------------2776667777777777777667666666666654   35 
AQP3.PRO       1 5---------9644-----33322111------457777665555556555657555544   40 
AQP4.PRO       1 63DGAAARRW9644PPCSR33322113667667677777666666666666767666665   60 
AQP5.PRO       1 65------------------------4877667667777777776666666667666665   36 
cons           1 54--------9644-----33322113777667667777766666666666666666554   60 


AQP1.PRO      37 111221111223454555555555444444444554455555555544544445555444   96 
AQP2.PRO      36 ---33----323444444444444444454444455556666555555555455565555   88 
AQP3.PRO      41 111211----11444444444433333232333344455565555544444445555444   96 
AQP4.PRO      61 ---331111223554555555544444454455555666666666555555455555555  117 
AQP5.PRO      37 ---43----323555555555554444454455555555666555555555555565555   89 
cons          61 111331111323454545444444444444444455555665555555555455555555  120 


AQP1.PRO      97 5555555556555655555----------4444445555555555555555-----5555  141 
AQP2.PRO      89 6556666666666665555----------5555555566666666666666-----5556  133 
AQP3.PRO      97 4445555555555554443YYDAIWAFAG2322222333333344443454LDMVN4445  156 
AQP4.PRO     118 5556666666666665544----------4445545555555555665666-----6666  162 
AQP5.PRO      90 5555555556666665554----------4455555666666666666677-----6676  134 
cons         121 5555556666666655554----------4444444555555555555566-----5566  180 


AQP1.PRO     142 6677777777777766666666-5666656666666555666556677777778777766  200 
AQP2.PRO     134 5566777777777766666555-5565556666666555556556788888888888777  192 
AQP3.PRO     157 7666777777777777765444V3454455554444322222224577777777777665  216 
AQP4.PRO     163 7777777777777777776666-5555555555666555555555677777777777766  221 
AQP5.PRO     135 7777888788887777776666-5666666666666555555556677777777777776  193 
cons         181 7677777777777777666555-5565555565556544555555677777777777766  240 


AQP1.PRO     201 7766--666---------66666666665555666666555566544455555220111-  248 
AQP2.PRO     193 7766--666---------6666666665555555555555556664456666633-222-  239 
AQP3.PRO     217 5544LA2220EVFTTGQN223333434333334433332322------------------  258 
AQP4.PRO     222 6555--555---------55655666656666655555555454-33345544100111-  268 
AQP5.PRO     194 6666--5440--------5555556655555555555544445664444555533-222T  242 
cons         241 6665--5550--------55555555555555555555444455544455555220221-  300 


AQP1.PRO     249 22222--2111111111----11000----------------------------3  269 
AQP2.PRO     240 33333--42222222223--3333333233322-------------------223  271 
AQP3.PRO     259 --------------------2111222122211122-111112221-------12  285 
AQP4.PRO     269 11111KG22111111123VE3222333233322122D111112221SSGEVL224  323 
AQP5.PRO     243 23333--33222223334--3333333----------------------------  265 
cons         301 23322--32222222223--3222223233322122-111112221------223  355 
%</TCoffee>
%    \end{macrocode}
%    \begin{macrocode}
%<*AQPlogo>
############################################################
###
###   . O ¡ o STRUCTURE MEME ¡ . O o
###
###   Chimera command file generated from TeXshade logo data
###   use '.cmd' or '.com' suffix
###
###   Reference sequence: 'E_coli_GlpF', e.g. PDB 1FX8
###
############################################################

############################################################
###   residue colors: (set with \memeRed{residues} etc.)
###   DE (Red)
###   HKR (Blue)
###   CGNQSTY (Yellow)
###   FILMW (Black)
###   APV (White)
############################################################

############################################################
###   residues in labels: (\memelabelcutoff{bit-value})
###   0.1 bit required to appear
###   > : . indicate 3, 2, and 1 bit levels
############################################################

############################################################
###   initialize & load structure
############################################################
close session; background solid white;
echo load structure file for 'E_coli_GlpF', e.g. PDB 1FX8
open #0; ~ribbon; ~surface

############################################################
###   to set ball scale, change factor ]0.0 - 1.0]
###   in 'setattr m ballScale <factor> #0' below,
###   or use the \chimeraballScale{factor} command
############################################################
setattr m ballScale 1.0 #0

############################################################
###   display alpha-carbons only
############################################################
represent bs @ca; show @ca

############################################################
###   name structure model
############################################################
setattr m name logo #0

############################################################
###   logo data
############################################################
setattr a radius 1.473 #0:7@ca
setattr a color #0CC010570517 #0:7@ca
setattr a name '|| .FLS' #0:7@ca
setattr a radius 1.061 #0:8@ca
setattr a color #22962541099F #0:8@ca
setattr a name '|| .WYLF' #0:8@ca
setattr a radius 2.096 #0:9@ca
setattr a color #09032F7DC2EF #0:9@ca
setattr a name '|| R.KQ' #0:9@ca
setattr a radius 2.178 #0:10@ca
setattr a color #CA40C889AC09 #0:10@ca
setattr a name '|| A.EKQ' #0:10@ca
setattr a radius 1.232 #0:11@ca
setattr a color #8A718A714538 #0:11@ca
setattr a name '|| .VGAYLF' #0:11@ca
setattr a radius 1.913 #0:12@ca
setattr a color #09F109F105D4 #0:12@ca
setattr a name '|| .ILFV' #0:12@ca
setattr a radius 3.112 #0:13@ca
setattr a color #FC8DFC8DCC2F #0:13@ca
setattr a name '|| A:.GSV' #0:13@ca
setattr a radius 4.322 #0:14@ca
setattr a color #FFFF00000000 #0:14@ca
setattr a name '|| E>:.' #0:14@ca
setattr a radius 3.440 #0:15@ca
setattr a color #016B016B0083 #0:15@ca
setattr a name '|| F:.LC' #0:15@ca
setattr a radius 2.230 #0:16@ca
setattr a color #075807A505D3 #0:16@ca
setattr a name '|| .LIVF' #0:16@ca
setattr a radius 2.657 #0:17@ca
setattr a color #FFFFFFFF783D #0:17@ca
setattr a name '|| A.GS' #0:17@ca
setattr a radius 3.773 #0:18@ca
setattr a color #D7E8D7E80000 #0:18@ca
setattr a name '|| T>:.M' #0:18@ca
setattr a radius 2.097 #0:19@ca
setattr a color #088A088A0420 #0:19@ca
setattr a name '|| .LFA' #0:19@ca
setattr a radius 2.481 #0:20@ca
setattr a color #00F200F200BD #0:20@ca
setattr a name '|| L.IFVW' #0:20@ca
setattr a radius 3.057 #0:21@ca
setattr a color #004A004A000E #0:21@ca
setattr a name '|| F.L' #0:21@ca
setattr a radius 2.701 #0:22@ca
setattr a color #3BC23BC23BC2 #0:22@ca
setattr a name '|| V.LI' #0:22@ca
setattr a radius 2.024 #0:23@ca
setattr a color #1ECB1ECB069F #0:23@ca
setattr a name '|| .FYLV' #0:23@ca
setattr a radius 1.768 #0:24@ca
setattr a color #1B441B441476 #0:24@ca
setattr a name '|| .IFALV' #0:24@ca
setattr a radius 2.841 #0:25@ca
setattr a color #FFFFFFFF08A6 #0:25@ca
setattr a name '|| G.TS' #0:25@ca
setattr a radius 1.326 #0:26@ca
setattr a color #879582864490 #0:26@ca
setattr a name '|| .VILC' #0:26@ca
setattr a radius 2.916 #0:27@ca
setattr a color #9F8A9F8A175B #0:27@ca
setattr a name '|| G.LA' #0:27@ca
setattr a radius 2.371 #0:28@ca
setattr a color #EF03EF032197 #0:28@ca
setattr a name '|| S.TAVC' #0:28@ca
setattr a radius 2.180 #0:29@ca
setattr a color #E860E860B33C #0:29@ca
setattr a name '|| V.AGC' #0:29@ca
setattr a radius 1.761 #0:30@ca
setattr a color #1CBC1CBC1779 #0:30@ca
setattr a name '|| .IMALV' #0:30@ca
setattr a radius 1.470 #0:31@ca
setattr a color #C42ECDFE4527 #0:31@ca
setattr a name '|| .GAQN' #0:31@ca
setattr a radius 1.181 #0:32@ca
setattr a color #5DEA64CE3980 #0:32@ca
setattr a name '|| .VYFWA' #0:32@ca
setattr a radius 0.902 #0:33@ca
setattr a color #75B6A8AB46D8 #0:33@ca
setattr a name '|| .KNV' #0:33@ca
setattr a radius 0.889 #0:34@ca
setattr a color #10BD5883373D #0:34@ca
setattr a name '|| .KLR' #0:34@ca
setattr a radius 0.906 #0:35@ca
setattr a color #74E177B12C96 #0:35@ca
setattr a name '|| .SLA' #0:35@ca
setattr a radius 0.623 #0:36@ca
setattr a color #83EF89E63287 #0:36@ca
setattr a name '|| .T' #0:36@ca
setattr a radius 0.820 #0:37@ca
setattr a color #D292B0543AE0 #0:37@ca
setattr a name '|| .SA' #0:37@ca
setattr a radius 0.623 #0:38@ca
setattr a color #925787241E9C #0:38@ca
setattr a name '|| .S' #0:38@ca
setattr a radius 0.436 #0:39@ca
setattr a color #4EE56E191F33 #0:39@ca
setattr a name '|| .' #0:39@ca
setattr a radius 0.550 #0:40@ca
setattr a color #B86FC42D53F3 #0:40@ca
setattr a name '|| .GA' #0:40@ca
setattr a radius 0.420 #0:41@ca
setattr a color #19491949024C #0:41@ca
setattr a name '|| .W' #0:41@ca
setattr a radius 0.498 #0:42@ca
setattr a color #01D201360000 #0:42@ca
setattr a name '|| .LI' #0:42@ca
setattr a radius 0.312 #0:43@ca
setattr a color #CCB5A5D24091 #0:43@ca
setattr a name '|| .' #0:43@ca
setattr a radius 2.906 #0:44@ca
setattr a color #1B971B971B0A #0:44@ca
setattr a name '|| I.VA' #0:44@ca
setattr a radius 2.647 #0:45@ca
setattr a color #FC87FC87A050 #0:45@ca
setattr a name '|| A.SN' #0:45@ca
setattr a radius 1.815 #0:46@ca
setattr a color #04BD04BD0447 #0:46@ca
setattr a name '|| .LWIVF' #0:46@ca
setattr a radius 2.362 #0:47@ca
setattr a color #E816E81695BC #0:47@ca
setattr a name '|| A.GCST' #0:47@ca
setattr a radius 2.672 #0:48@ca
setattr a color #006602D30BD3 #0:48@ca
setattr a name '|| F.HW' #0:48@ca
setattr a radius 3.512 #0:49@ca
setattr a color #FFFFFFFF284F #0:49@ca
setattr a name '|| G:.A' #0:49@ca
setattr a radius 2.212 #0:50@ca
setattr a color #142D142D00F9 #0:50@ca
setattr a name '|| .LGFI' #0:50@ca
setattr a radius 2.009 #0:51@ca
setattr a color #8C0F8C0F44B5 #0:51@ca
setattr a name '|| .AMGS' #0:51@ca
setattr a radius 2.611 #0:52@ca
setattr a color #1CAF1CAF1AFA #0:52@ca
setattr a name '|| I.VL' #0:52@ca
setattr a radius 1.862 #0:53@ca
setattr a color #329F329F1E62 #0:53@ca
setattr a name '|| .FAGMTL' #0:53@ca
setattr a radius 1.664 #0:54@ca
setattr a color #760E760E49F7 #0:54@ca
setattr a name '|| .VTAMILF' #0:54@ca
setattr a radius 1.938 #0:55@ca
setattr a color #3AB33AB324F0 #0:55@ca
setattr a name '|| .LAGM' #0:55@ca
setattr a radius 2.923 #0:56@ca
setattr a color #A74DA74DA74D #0:56@ca
setattr a name '|| V.IA' #0:56@ca
setattr a radius 1.937 #0:57@ca
setattr a color #BBE9C06E0F4D #0:57@ca
setattr a name '|| .YQSAL' #0:57@ca
setattr a radius 1.734 #0:58@ca
setattr a color #BF91BF91517A #0:58@ca
setattr a name '|| .CAVSIT' #0:58@ca
setattr a radius 1.519 #0:59@ca
setattr a color #94C794C73679 #0:59@ca
setattr a name '|| .TAVLGFS' #0:59@ca
setattr a radius 2.729 #0:60@ca
setattr a color #E1A0E1A074F3 #0:60@ca
setattr a name '|| AG.' #0:60@ca
setattr a radius 1.939 #0:61@ca
setattr a color #7125B8920DDC #0:61@ca
setattr a name '|| .GHN' #0:61@ca
setattr a radius 2.932 #0:62@ca
setattr a color #102110210EDC #0:62@ca
setattr a name '|| I.VL' #0:62@ca
setattr a radius 4.044 #0:63@ca
setattr a color #F914F91401AF #0:63@ca
setattr a name '|| S>:.' #0:63@ca
setattr a radius 0.728 #0:64@ca
setattr a color #FFFFFFFF20B3 #0:64@ca
setattr a name '|| .G' #0:64@ca
setattr a radius 3.132 #0:65@ca
setattr a color #FFFFFFFF5555 #0:65@ca
setattr a name '|| GA.' #0:65@ca
setattr a radius 3.981 #0:66@ca
setattr a color #06F10C37FAC9 #0:66@ca
setattr a name '|| H>:.A' #0:66@ca
setattr a radius 2.420 #0:67@ca
setattr a color #06C906C906C9 #0:67@ca
setattr a name '|| I.LVF' #0:67@ca
setattr a radius 4.322 #0:68@ca
setattr a color #FFFFFFFF0000 #0:68@ca
setattr a name '|| N>:.' #0:68@ca
setattr a radius 4.218 #0:69@ca
setattr a color #FFFFFFFFFFFF #0:69@ca
setattr a name '|| P>:.' #0:69@ca
setattr a radius 4.264 #0:70@ca
setattr a color #FFFFFFFFFE41 #0:70@ca
setattr a name '|| A>:.' #0:70@ca
setattr a radius 3.811 #0:71@ca
setattr a color #DB83DB83DB83 #0:71@ca
setattr a name '|| V>:.I' #0:71@ca
setattr a radius 3.595 #0:72@ca
setattr a color #FFFFFFFF01B5 #0:72@ca
setattr a name '|| T:.S' #0:72@ca
setattr a radius 2.524 #0:73@ca
setattr a color #03D603D603A0 #0:73@ca
setattr a name '|| F.LIV' #0:73@ca
setattr a radius 2.925 #0:74@ca
setattr a color #EF21EF214BCA #0:74@ca
setattr a name '|| G.A' #0:74@ca
setattr a radius 2.254 #0:75@ca
setattr a color #0D920D920731 #0:75@ca
setattr a name '|| L.MAFC' #0:75@ca
setattr a radius 1.789 #0:76@ca
setattr a color #2518251817DB #0:76@ca
setattr a name '|| .FLACW' #0:76@ca
setattr a radius 1.947 #0:77@ca
setattr a color #226A230D1A05 #0:77@ca
setattr a name '|| .LVI' #0:77@ca
setattr a radius 1.627 #0:78@ca
setattr a color #B633B9314FD0 #0:78@ca
setattr a name '|| .AGSF' #0:78@ca
setattr a radius 2.199 #0:79@ca
setattr a color #50ACA8551176 #0:79@ca
setattr a name '|| .GRCAS' #0:79@ca
setattr a radius 1.325 #0:80@ca
setattr a color #17719FE59351 #0:80@ca
setattr a name '|| .KRQHN' #0:80@ca
setattr a radius 2.039 #0:81@ca
setattr a color #0B5509E10989 #0:81@ca
setattr a name '|| .ILVFM' #0:81@ca
setattr a radius 2.139 #0:82@ca
setattr a color #EBA3E17B31D8 #0:82@ca
setattr a name '|| S.PT' #0:82@ca
setattr a radius 1.836 #0:83@ca
setattr a color #0885099206BB #0:83@ca
setattr a name '|| .LIWV' #0:83@ca
setattr a radius 1.108 #0:84@ca
setattr a color #257E35D21A7C #0:84@ca
setattr a name '|| .LTIA' #0:84@ca
setattr a radius 2.072 #0:85@ca
setattr a color #09DB2D14C16B #0:85@ca
setattr a name '|| R.KT' #0:85@ca
setattr a radius 1.938 #0:86@ca
setattr a color #B352B3527D7B #0:86@ca
setattr a name '|| .AGVLSF' #0:86@ca
setattr a radius 1.921 #0:87@ca
setattr a color #40DE40DE40DE #0:87@ca
setattr a name '|| .VILPA' #0:87@ca
setattr a radius 1.176 #0:88@ca
setattr a color #2C092D871C5C #0:88@ca
setattr a name '|| .FLPV' #0:88@ca
setattr a radius 3.969 #0:89@ca
setattr a color #EEE0EEE001AE #0:89@ca
setattr a name '|| Y>:.M' #0:89@ca
setattr a radius 2.150 #0:90@ca
setattr a color #093809380797 #0:90@ca
setattr a name '|| .IWMV' #0:90@ca
setattr a radius 2.112 #0:91@ca
setattr a color #463346333E14 #0:91@ca
setattr a name '|| .IVAL' #0:91@ca
setattr a radius 2.760 #0:92@ca
setattr a color #BC20B92EA154 #0:92@ca
setattr a name '|| A.MS' #0:92@ca
setattr a radius 4.322 #0:93@ca
setattr a color #FFFFFFFF0000 #0:93@ca
setattr a name '|| Q>:.' #0:93@ca
setattr a radius 1.979 #0:94@ca
setattr a color #4BD14BD1101F #0:94@ca
setattr a name '|| .CLVIF' #0:94@ca
setattr a radius 2.917 #0:95@ca
setattr a color #09F209F209F2 #0:95@ca
setattr a name '|| L:.VIA' #0:95@ca
setattr a radius 4.045 #0:96@ca
setattr a color #FFFFFFFF0C26 #0:96@ca
setattr a name '|| G>:.A' #0:96@ca
setattr a radius 3.214 #0:97@ca
setattr a color #FFFFFFFFBD68 #0:97@ca
setattr a name '|| A:.SG' #0:97@ca
setattr a radius 2.067 #0:98@ca
setattr a color #106310630AA7 #0:98@ca
setattr a name '|| I.VFMT' #0:98@ca
setattr a radius 1.755 #0:99@ca
setattr a color #89AB89AB53A6 #0:99@ca
setattr a name '|| .CVAIL' #0:99@ca
setattr a radius 3.132 #0:100@ca
setattr a color #FC93FC93634B #0:100@ca
setattr a name '|| GA.' #0:100@ca
setattr a radius 2.092 #0:101@ca
setattr a color #FFFFFFFF9C24 #0:101@ca
setattr a name '|| A.CVTS' #0:101@ca
setattr a radius 2.127 #0:102@ca
setattr a color #9C4D9C4D48B2 #0:102@ca
setattr a name '|| .GALF' #0:102@ca
setattr a radius 2.026 #0:103@ca
setattr a color #16D316D311A8 #0:103@ca
setattr a name '|| .LIVTF' #0:103@ca
setattr a radius 2.680 #0:104@ca
setattr a color #3D723D723930 #0:104@ca
setattr a name '|| LV.AI' #0:104@ca
setattr a radius 1.577 #0:105@ca
setattr a color #0B586EDB1869 #0:105@ca
setattr a name '|| .KYS' #0:105@ca
setattr a radius 1.426 #0:106@ca
setattr a color #74436E742DDA #0:106@ca
setattr a name '|| .GALFV' #0:106@ca
setattr a radius 1.409 #0:107@ca
setattr a color #1E7B205718D6 #0:107@ca
setattr a name '|| .FVIL' #0:107@ca
setattr a radius 1.493 #0:108@ca
setattr a color #B93CAEE1162B #0:108@ca
setattr a name '|| .TYQFA' #0:108@ca
setattr a radius 1.141 #0:109@ca
setattr a color #AB61C5D44348 #0:109@ca
setattr a name '|| .PSGY' #0:109@ca
setattr a radius 1.168 #0:110@ca
setattr a color #B827B85D369C #0:110@ca
setattr a name '|| .GPKA' #0:110@ca
setattr a radius 0.603 #0:111@ca
setattr a color #61F26B8B1B77 #0:111@ca
setattr a name '|| .L' #0:111@ca
setattr a radius 0.534 #0:112@ca
setattr a color #33354166273D #0:112@ca
setattr a name '|| .A' #0:112@ca
setattr a radius 0.728 #0:113@ca
setattr a color #791975332C70 #0:113@ca
setattr a name '|| .YV' #0:113@ca
setattr a radius 0.443 #0:114@ca
setattr a color #7D67851D497E #0:114@ca
setattr a name '|| .' #0:114@ca
setattr a radius 0.645 #0:115@ca
setattr a color #809B8FD62A4E #0:115@ca
setattr a name '|| .G' #0:115@ca
setattr a radius 0.848 #0:116@ca
setattr a color #3774413D1985 #0:116@ca
setattr a name '|| .LAYHF' #0:116@ca
setattr a radius 1.620 #0:117@ca
setattr a color #E965E3C01878 #0:117@ca
setattr a name '|| .GSNT' #0:117@ca
setattr a radius 1.568 #0:118@ca
setattr a color #67CA6B080B2C #0:118@ca
setattr a name '|| .GL' #0:118@ca
setattr a radius 0.157 #0:119@ca
setattr a color #5BDC80002424 #0:119@ca
setattr a name '|| .' #0:119@ca
setattr a radius 0.277 #0:120@ca
setattr a color #0E4B1D010000 #0:120@ca
setattr a name '|| .L' #0:120@ca
setattr a radius 0.250 #0:121@ca
setattr a color #22AB181D0A2C #0:121@ca
setattr a name '|| .' #0:121@ca
setattr a radius 0.396 #0:122@ca
setattr a color #5CF9625257A0 #0:122@ca
setattr a name '|| .V' #0:122@ca
setattr a radius 0.343 #0:123@ca
setattr a color #8894B31F5515 #0:123@ca
setattr a name '|| .' #0:123@ca
setattr a radius 0.490 #0:124@ca
setattr a color #FFFFF50615F1 #0:124@ca
setattr a name '|| .G' #0:124@ca
setattr a radius 0.400 #0:125@ca
setattr a color #F0DDC24E3E14 #0:125@ca
setattr a name '|| .T' #0:125@ca
setattr a radius 0.267 #0:126@ca
setattr a color #A41789990BF7 #0:126@ca
setattr a name '|| .' #0:126@ca
setattr a radius 0.286 #0:127@ca
setattr a color #EAF7540A1494 #0:127@ca
setattr a name '|| .' #0:127@ca
setattr a radius 0.331 #0:128@ca
setattr a color #E10FE8CB5CCE #0:128@ca
setattr a name '|| .' #0:128@ca
setattr a radius 0.314 #0:129@ca
setattr a color #96CE89AA19A7 #0:129@ca
setattr a name '|| .' #0:129@ca
setattr a radius 1.806 #0:130@ca
setattr a color #E28BDE5F3880 #0:130@ca
setattr a name '|| G.A' #0:130@ca
setattr a radius 1.235 #0:131@ca
setattr a color #884B94166CF0 #0:131@ca
setattr a name '|| .AVTL' #0:131@ca
setattr a radius 1.615 #0:132@ca
setattr a color #B60EB60E2063 #0:132@ca
setattr a name '|| .NALT' #0:132@ca
setattr a radius 0.847 #0:133@ca
setattr a color #CBE0BC7D3F78 #0:133@ca
setattr a name '|| .GATS' #0:133@ca
setattr a radius 0.494 #0:134@ca
setattr a color #39A84CAA1FFF #0:134@ca
setattr a name '|| .I' #0:134@ca
setattr a radius 2.132 #0:135@ca
setattr a color #292A2B572577 #0:135@ca
setattr a name '|| .VLFS' #0:135@ca
setattr a radius 1.717 #0:136@ca
setattr a color #D340E65C7367 #0:136@ca
setattr a name '|| .ASHPN' #0:136@ca
setattr a radius 1.055 #0:137@ca
setattr a color #BFCBBC894344 #0:137@ca
setattr a name '|| .TPAG' #0:137@ca
setattr a radius 2.127 #0:138@ca
setattr a color #DE03DC6B04E4 #0:138@ca
setattr a name '|| G.YN' #0:138@ca
setattr a radius 1.372 #0:139@ca
setattr a color #BEF8BA9E5CAA #0:139@ca
setattr a name '|| .YVPL' #0:139@ca
setattr a radius 1.504 #0:140@ca
setattr a color #E7D3EB661D9D #0:140@ca
setattr a name '|| .STNGA' #0:140@ca
setattr a radius 0.797 #0:141@ca
setattr a color #55ED82B25517 #0:141@ca
setattr a name '|| .KVA' #0:141@ca
setattr a radius 1.416 #0:142@ca
setattr a color #885D96A80FA1 #0:142@ca
setattr a name '|| .GL' #0:142@ca
setattr a radius 0.600 #0:143@ca
setattr a color #9522810D182C #0:143@ca
setattr a name '|| .Q' #0:143@ca
setattr a radius 2.204 #0:144@ca
setattr a color #F158EBEE4AFF #0:144@ca
setattr a name '|| G.AS' #0:144@ca
setattr a radius 0.303 #0:145@ca
setattr a color #282728270B60 #0:145@ca
setattr a name '|| .' #0:145@ca
setattr a radius 0.233 #0:146@ca
setattr a color #43BB474C3871 #0:146@ca
setattr a name '|| .' #0:146@ca
setattr a radius 0.290 #0:147@ca
setattr a color #8E7EC0482A0F #0:147@ca
setattr a name '|| .' #0:147@ca
setattr a radius 0.584 #0:148@ca
setattr a color #FFFFF8FB5037 #0:148@ca
setattr a name '|| .GNA' #0:148@ca
setattr a radius 2.455 #0:149@ca
setattr a color #0A9E0A9E0937 #0:149@ca
setattr a name '|| L.FV' #0:149@ca
setattr a radius 1.604 #0:150@ca
setattr a color #744574453878 #0:150@ca
setattr a name '|| .GVLAF' #0:150@ca
setattr a radius 1.186 #0:151@ca
setattr a color #7DAE68F55464 #0:151@ca
setattr a name '|| .AVDM' #0:151@ca
setattr a radius 3.778 #0:152@ca
setattr a color #F922148501AB #0:152@ca
setattr a name '|| E>:.Q' #0:152@ca
setattr a radius 2.401 #0:153@ca
setattr a color #0531053104B5 #0:153@ca
setattr a name '|| I.LFV' #0:153@ca
setattr a radius 2.400 #0:154@ca
setattr a color #214121412006 #0:154@ca
setattr a name '|| I.VFL' #0:154@ca
setattr a radius 2.040 #0:155@ca
setattr a color #53A153A10AE6 #0:155@ca
setattr a name '|| .GLIM' #0:155@ca
setattr a radius 3.786 #0:156@ca
setattr a color #F5D2F5D20361 #0:156@ca
setattr a name '|| T>:.S' #0:156@ca
setattr a radius 2.596 #0:157@ca
setattr a color #093B093B059E #0:157@ca
setattr a name '|| F.ALMG' #0:157@ca
setattr a radius 1.530 #0:158@ca
setattr a color #8F9F8F9F5639 #0:158@ca
setattr a name '|| .VQIAGL' #0:158@ca
setattr a radius 3.733 #0:159@ca
setattr a color #000200020002 #0:159@ca
setattr a name '|| L>:.F' #0:159@ca
setattr a radius 2.438 #0:160@ca
setattr a color #92229222851F #0:160@ca
setattr a name '|| V.LIM' #0:160@ca
setattr a radius 2.187 #0:161@ca
setattr a color #4BCB4BCB0B56 #0:161@ca
setattr a name '|| .YLFV' #0:161@ca
setattr a radius 2.103 #0:162@ca
setattr a color #C540C540185C #0:162@ca
setattr a name '|| T.CVIG' #0:162@ca
setattr a radius 3.048 #0:163@ca
setattr a color #7B427B4278EE #0:163@ca
setattr a name '|| V:.I' #0:163@ca
setattr a radius 1.928 #0:164@ca
setattr a color #16A918E9051A #0:164@ca
setattr a name '|| .FLY' #0:164@ca
setattr a radius 2.345 #0:165@ca
setattr a color #E4D0E3298B54 #0:165@ca
setattr a name '|| A.SGVM' #0:165@ca
setattr a radius 1.698 #0:166@ca
setattr a color #A81CA81C5372 #0:166@ca
setattr a name '|| .ATSIVL' #0:166@ca
setattr a radius 2.140 #0:167@ca
setattr a color #E572E5723D0A #0:167@ca
setattr a name '|| T.AGC' #0:167@ca
setattr a radius 2.367 #0:168@ca
setattr a color #D7F7436B34AF #0:168@ca
setattr a name '|| D.VA' #0:168@ca
setattr a radius 1.108 #0:169@ca
setattr a color #A9827F5A49E1 #0:169@ca
setattr a name '|| .DPAS' #0:169@ca
setattr a radius 1.558 #0:170@ca
setattr a color #352B5625BC11 #0:170@ca
setattr a name '|| .KRP' #0:170@ca
setattr a radius 1.995 #0:171@ca
setattr a color #113C4342B925 #0:171@ca
setattr a name '|| .RKNA' #0:171@ca
setattr a radius 1.077 #0:172@ca
setattr a color #79A6B0F7151D #0:172@ca
setattr a name '|| .NSKT' #0:172@ca
setattr a radius 1.117 #0:173@ca
setattr a color #CD2DCD2D2C1C #0:173@ca
setattr a name '|| .SV' #0:173@ca
setattr a radius 1.221 #0:174@ca
setattr a color #3DB56D3C278A #0:174@ca
setattr a name '|| .HSDP' #0:174@ca
setattr a radius 1.262 #0:175@ca
setattr a color #4AC15B844421 #0:175@ca
setattr a name '|| .VLRI' #0:175@ca
setattr a radius 1.507 #0:176@ca
setattr a color #D5BAD946595C #0:176@ca
setattr a name '|| .GP' #0:176@ca
setattr a radius 1.273 #0:177@ca
setattr a color #88E986162BA8 #0:177@ca
setattr a name '|| .GITV' #0:177@ca
setattr a radius 1.967 #0:178@ca
setattr a color #10FF1171032A #0:178@ca
setattr a name '|| .LSIF' #0:178@ca
setattr a radius 2.318 #0:179@ca
setattr a color #DE99D13DB6D6 #0:179@ca
setattr a name '|| A.PEGI' #0:179@ca
setattr a radius 3.108 #0:180@ca
setattr a color #F913FC89ECF8 #0:180@ca
setattr a name '|| P:.AG' #0:180@ca
setattr a radius 2.738 #0:181@ca
setattr a color #008F008F0066 #0:181@ca
setattr a name '|| L.IF' #0:181@ca
setattr a radius 1.791 #0:182@ca
setattr a color #B391B3919978 #0:182@ca
setattr a name '|| .APSL' #0:182@ca
setattr a radius 3.459 #0:183@ca
setattr a color #063D063D04E3 #0:183@ca
setattr a name '|| I:.V' #0:183@ca
setattr a radius 4.160 #0:184@ca
setattr a color #F5C3F5C30000 #0:184@ca
setattr a name '|| G>:.' #0:184@ca
setattr a radius 2.710 #0:185@ca
setattr a color #00BE00BE005F #0:185@ca
setattr a name '|| F.LI' #0:185@ca
setattr a radius 1.819 #0:186@ca
setattr a color #78E478E44447 #0:186@ca
setattr a name '|| .ASLIV' #0:186@ca
setattr a radius 3.278 #0:187@ca
setattr a color #9F729F729F72 #0:187@ca
setattr a name '|| V:.IL' #0:187@ca
setattr a radius 1.622 #0:188@ca
setattr a color #57A257A227B5 #0:188@ca
setattr a name '|| .FAGTV' #0:188@ca
setattr a radius 1.771 #0:189@ca
setattr a color #339633962976 #0:189@ca
setattr a name '|| .LAMVI' #0:189@ca
setattr a radius 1.839 #0:190@ca
setattr a color #80527CAC33D7 #0:190@ca
setattr a name '|| .VIGN' #0:190@ca
setattr a radius 2.490 #0:191@ca
setattr a color #0AC7457E76A2 #0:191@ca
setattr a name '|| H.IGV' #0:191@ca
setattr a radius 2.864 #0:192@ca
setattr a color #0531053101BB #0:192@ca
setattr a name '|| L:.TM' #0:192@ca
setattr a radius 2.174 #0:193@ca
setattr a color #9F7F9F7F7A81 #0:193@ca
setattr a name '|| A.SVLIF' #0:193@ca
setattr a radius 1.675 #0:194@ca
setattr a color #880886CE2770 #0:194@ca
setattr a name '|| .TGALM' #0:194@ca
setattr a radius 2.890 #0:195@ca
setattr a color #3BA53BA501AA #0:195@ca
setattr a name '|| IG.' #0:195@ca
setattr a radius 1.530 #0:196@ca
setattr a color #C85DC1DB8318 #0:196@ca
setattr a name '|| .PGAY' #0:196@ca
setattr a radius 1.528 #0:197@ca
setattr a color #2A8E2A8E0CB4 #0:197@ca
setattr a name '|| .IFYNV' #0:197@ca
setattr a radius 3.127 #0:198@ca
setattr a color #FFFFF9090370 #0:198@ca
setattr a name '|| T:.S' #0:198@ca
setattr a radius 3.941 #0:199@ca
setattr a color #FFFFFFFF0000 #0:199@ca
setattr a name '|| G>:.N' #0:199@ca
setattr a radius 1.755 #0:200@ca
setattr a color #E8A2E8A23372 #0:200@ca
setattr a name '|| .TACYG' #0:200@ca
setattr a radius 2.477 #0:201@ca
setattr a color #F96DF96D36FC #0:201@ca
setattr a name '|| .SGA' #0:201@ca
setattr a radius 2.727 #0:202@ca
setattr a color #02A902A902A9 #0:202@ca
setattr a name '|| MI.VL' #0:202@ca
setattr a radius 4.264 #0:203@ca
setattr a color #FFFFFFFF0000 #0:203@ca
setattr a name '|| N>:.' #0:203@ca
setattr a radius 4.264 #0:204@ca
setattr a color #FFFFFE41FE41 #0:204@ca
setattr a name '|| P>:.' #0:204@ca
setattr a radius 4.017 #0:205@ca
setattr a color #FFFFFFFFF90C #0:205@ca
setattr a name '|| A>:.' #0:205@ca
setattr a radius 3.906 #0:206@ca
setattr a color #115E1312FE4A #0:206@ca
setattr a name '|| R>:.V' #0:206@ca
setattr a radius 3.077 #0:207@ca
setattr a color #FFFFC4C214CC #0:207@ca
setattr a name '|| S:.DA' #0:207@ca
setattr a radius 2.971 #0:208@ca
setattr a color #004B004B0000 #0:208@ca
setattr a name '|| FL.T' #0:208@ca
setattr a radius 3.624 #0:209@ca
setattr a color #FFFFFE4C169A #0:209@ca
setattr a name '|| G>:.A' #0:209@ca
setattr a radius 2.802 #0:210@ca
setattr a color #FC89FC89E7D9 #0:210@ca
setattr a name '|| P.AS' #0:210@ca
setattr a radius 3.028 #0:211@ca
setattr a color #B86EC26DF40F #0:211@ca
setattr a name '|| A:.RK' #0:211@ca
setattr a radius 2.586 #0:212@ca
setattr a color #3F353F353F35 #0:212@ca
setattr a name '|| V.IL' #0:212@ca
setattr a radius 1.898 #0:213@ca
setattr a color #28F028F026D5 #0:213@ca
setattr a name '|| .IVFAL' #0:213@ca
setattr a radius 1.162 #0:214@ca
setattr a color #75D481051BA5 #0:214@ca
setattr a name '|| .TYFSA' #0:214@ca
setattr a radius 1.263 #0:215@ca
setattr a color #7EBC93D410C6 #0:215@ca
setattr a name '|| .NGWA' #0:215@ca
setattr a radius 0.680 #0:216@ca
setattr a color #263344180922 #0:216@ca
setattr a name '|| .KL' #0:216@ca
setattr a radius 0.989 #0:217@ca
setattr a color #382C25100F84 #0:217@ca
setattr a name '|| .FWDA' #0:217@ca
setattr a radius 0.854 #0:218@ca
setattr a color #7292989422BA #0:218@ca
setattr a name '|| .GSK' #0:218@ca
setattr a radius 0.528 #0:219@ca
setattr a color #2BFF271C138E #0:219@ca
setattr a name '|| .W' #0:219@ca
setattr a radius 0.541 #0:220@ca
setattr a color #BAE9DD740000 #0:220@ca
setattr a name '|| .G' #0:220@ca
setattr a radius 0.224 #0:221@ca
setattr a color #77206FBD1802 #0:221@ca
setattr a name '|| .' #0:221@ca
setattr a radius 0.218 #0:222@ca
setattr a color #6FF9650C3FB9 #0:222@ca
setattr a name '|| .' #0:222@ca
setattr a radius 0.362 #0:223@ca
setattr a color #B5B9B5B99C1A #0:223@ca
setattr a name '|| .VA' #0:223@ca
setattr a radius 0.503 #0:224@ca
setattr a color #024503560000 #0:224@ca
setattr a name '|| .F' #0:224@ca
setattr a radius 0.309 #0:225@ca
setattr a color #A983B04B333B #0:225@ca
setattr a name '|| .T' #0:225@ca
setattr a radius 0.357 #0:226@ca
setattr a color #B502CB36347C #0:226@ca
setattr a name '|| .G' #0:226@ca
setattr a radius 0.692 #0:227@ca
setattr a color #934FAD710688 #0:227@ca
setattr a name '|| .G' #0:227@ca
setattr a radius 0.400 #0:228@ca
setattr a color #9C38955F0000 #0:228@ca
setattr a name '|| .N' #0:228@ca
setattr a radius 0.777 #0:229@ca
setattr a color #D807CC3C6F1E #0:229@ca
setattr a name '|| .AS' #0:229@ca
setattr a radius 0.857 #0:230@ca
setattr a color #624829810000 #0:230@ca
setattr a name '|| .DW' #0:230@ca
setattr a radius 0.253 #0:231@ca
setattr a color #5E8E817C33F6 #0:231@ca
setattr a name '|| .' #0:231@ca
setattr a radius 1.421 #0:232@ca
setattr a color #BF4F7782047F #0:232@ca
setattr a name '|| .DYN' #0:232@ca
setattr a radius 1.940 #0:233@ca
setattr a color #109D39D6842F #0:233@ca
setattr a name '|| H.WQA' #0:233@ca
setattr a radius 3.685 #0:234@ca
setattr a color #001B001B0012 #0:234@ca
setattr a name '|| W>:.' #0:234@ca
setattr a radius 2.966 #0:235@ca
setattr a color #2AAD2AAD29F8 #0:235@ca
setattr a name '|| IV.' #0:235@ca
setattr a radius 2.740 #0:236@ca
setattr a color #4F9C4F9C20FB #0:236@ca
setattr a name '|| F.YP' #0:236@ca
setattr a radius 3.026 #0:237@ca
setattr a color #0440044003D1 #0:237@ca
setattr a name '|| W:.VI' #0:237@ca
setattr a radius 2.881 #0:238@ca
setattr a color #9FA19FA19E4C #0:238@ca
setattr a name '|| V:.IAL' #0:238@ca
setattr a radius 3.342 #0:239@ca
setattr a color #FCA3FCA325F8 #0:239@ca
setattr a name '|| G:.AS' #0:239@ca
setattr a radius 4.160 #0:240@ca
setattr a color #FFFFFE45FAD3 #0:240@ca
setattr a name '|| P>:.' #0:240@ca
setattr a radius 2.345 #0:241@ca
setattr a color #00460046001C #0:241@ca
setattr a name '|| .LFIM' #0:241@ca
setattr a radius 2.392 #0:242@ca
setattr a color #19081908176D #0:242@ca
setattr a name '|| I.VL' #0:242@ca
setattr a radius 4.017 #0:243@ca
setattr a color #FFFFFFFF0DE4 #0:243@ca
setattr a name '|| G>:.A' #0:243@ca
setattr a radius 3.083 #0:244@ca
setattr a color #FFFFFFFF924E #0:244@ca
setattr a name '|| AG.S' #0:244@ca
setattr a radius 1.637 #0:245@ca
setattr a color #82C582C559C0 #0:245@ca
setattr a name '|| .AGVIL' #0:245@ca
setattr a radius 2.327 #0:246@ca
setattr a color #0B420B42097B #0:246@ca
setattr a name '|| L.IAV' #0:246@ca
setattr a radius 3.218 #0:247@ca
setattr a color #F5EAF5EAB6C2 #0:247@ca
setattr a name '|| A:.G' #0:247@ca
setattr a radius 2.479 #0:248@ca
setattr a color #F258F25887A4 #0:248@ca
setattr a name '|| A.GVS' #0:248@ca
setattr a radius 1.773 #0:249@ca
setattr a color #0F440F4409C4 #0:249@ca
setattr a name '|| .LIFAV' #0:249@ca
setattr a radius 1.935 #0:250@ca
setattr a color #445F445F1DAC #0:250@ca
setattr a name '|| .IYVL' #0:250@ca
setattr a radius 2.590 #0:251@ca
setattr a color #501E8C4C048C #0:251@ca
setattr a name '|| Y.HWF' #0:251@ca
setattr a radius 0.903 #0:252@ca
setattr a color #92D47EB2186D #0:252@ca
setattr a name '|| .QE' #0:252@ca
setattr a radius 1.003 #0:253@ca
setattr a color #46EF47CA1A99 #0:253@ca
setattr a name '|| .FYVL' #0:253@ca
setattr a radius 1.957 #0:254@ca
setattr a color #2F9A30592D5E #0:254@ca
setattr a name '|| .VLIFM' #0:254@ca
setattr a radius 1.931 #0:255@ca
setattr a color #04C505050243 #0:255@ca
setattr a name '|| .LIF' #0:255@ca
setattr a radius 0.984 #0:256@ca
setattr a color #192B52294732 #0:256@ca
setattr a name '|| .RI' #0:256@ca
setattr a radius 1.032 #0:257@ca
setattr a color #BC94BC948612 #0:257@ca
setattr a name '|| .AP' #0:257@ca
setattr a radius 0.726 #0:258@ca
setattr a color #501D60C71220 #0:258@ca
setattr a name '|| .SI' #0:258@ca
setattr a radius 0.319 #0:259@ca
setattr a color #115914170EEA #0:259@ca
setattr a name '|| .L' #0:259@ca
setattr a radius 0.267 #0:260@ca
setattr a color #D5CFB5356B62 #0:260@ca
setattr a name '|| .P' #0:260@ca
setattr a radius 0.155 #0:261@ca
setattr a color #BD3DB2331614 #0:261@ca
setattr a name '|| .' #0:261@ca
setattr a radius 0.220 #0:262@ca
setattr a color #A69CA69C854A #0:262@ca
setattr a name '|| .' #0:262@ca
setattr a radius 0.231 #0:263@ca
setattr a color #B9C28208662A #0:263@ca
setattr a name '|| .' #0:263@ca
setattr a radius 0.215 #0:264@ca
setattr a color #A565AD6A4F30 #0:264@ca
setattr a name '|| .' #0:264@ca
setattr a radius 0.120 #0:265@ca
setattr a color #8C0082AA2555 #0:265@ca
setattr a name '|| .' #0:265@ca
setattr a radius 0.133 #0:266@ca
setattr a color #6CD88C653BF1 #0:266@ca
setattr a name '|| .' #0:266@ca
setattr a radius 0.684 #0:267@ca
setattr a color #CCC3A7885B01 #0:267@ca
setattr a name '|| .PA' #0:267@ca
setattr a radius 0.598 #0:268@ca
setattr a color #215145631FCA #0:268@ca
setattr a name '|| .F' #0:268@ca
setattr a radius 0.897 #0:269@ca
setattr a color #3ED96AAF3F60 #0:269@ca
setattr a name '|| .KE' #0:269@ca
setattr a radius 0.771 #0:270@ca
setattr a color #7DFB8C3642F1 #0:270@ca
setattr a name '|| .SLPA' #0:270@ca
setattr a radius 0.412 #0:271@ca
setattr a color #62D273CCB108 #0:271@ca
setattr a name '|| .K' #0:271@ca
setattr a radius 0.237 #0:272@ca
setattr a color #EEB6E6125BD0 #0:272@ca
setattr a name '|| .' #0:272@ca
setattr a radius 0.254 #0:273@ca
setattr a color #CAABA0853237 #0:273@ca
setattr a name '|| .' #0:273@ca
setattr a radius 0.262 #0:274@ca
setattr a color #4A0258CF3B35 #0:274@ca
setattr a name '|| .' #0:274@ca
setattr a radius 0.253 #0:275@ca
setattr a color #841384135601 #0:275@ca
setattr a name '|| .' #0:275@ca
setattr a radius 0.200 #0:276@ca
setattr a color #49027CD3A992 #0:276@ca
setattr a name '|| .' #0:276@ca
setattr a radius 0.300 #0:277@ca
setattr a color #C28EAE1351EB #0:277@ca
setattr a name '|| .' #0:277@ca
setattr a radius 0.264 #0:278@ca
setattr a color #22614CB10DE2 #0:278@ca
setattr a name '|| .' #0:278@ca
setattr a radius 0.288 #0:279@ca
setattr a color #434B36849955 #0:279@ca
setattr a name '|| .K' #0:279@ca
setattr a radius 0.192 #0:280@ca
setattr a color #558E6CE37C71 #0:280@ca
setattr a name '|| .' #0:280@ca
setattr a radius 0.232 #0:281@ca
setattr a color #2E471A71358C #0:281@ca
setattr a name '|| .' #0:281@ca
%</AQPlogo>
%    \end{macrocode}
%    \begin{macrocode}
%<*AQPsublogo>
############################################################
###
###   . O ¡ o STRUCTURE MEME ¡ . O o
###
###   Chimera command file generated from TeXshade logo data
###   use '.cmd' or '.com' suffix
###
###   Reference sequence: 'E_coli_GlpF', e.g. PDB 1FX8
###
############################################################

############################################################
###   residue colors: (set with \memeRed{residues} etc.)
###   DE (Red)
###   HKR (Blue)
###   CGNQSTY (Yellow)
###   FILMW (Black)
###   APV (White)
############################################################

############################################################
###   residues in labels: (\memelabelcutoff{bit-value})
###   0.1 bit required to appear
###   > : . indicate 3, 2, and 1 bit levels
############################################################

############################################################
###   initialize & load structure
############################################################
close session; background solid white;
echo load structure file for 'E_coli_GlpF', e.g. PDB 1FX8
open #0; ~ribbon; ~surface

############################################################
###   to set ball scale, change factor ]0.0 - 1.0]
###   in 'setattr m ballScale <factor> #0' below,
###   or use the \chimeraballScale{factor} command
############################################################
setattr m ballScale 1.0 #0

############################################################
###   display alpha-carbons only
############################################################
represent bs @ca; show @ca

############################################################
###   duplicate structure;
###   to alter direction or Angstrom distance, set
###   'move <x, y, or z> 1 <distance> models #1' below,
###   or use the \chimeraaxisdistance{x}{y}{z} command
############################################################
combine #0
move x 1 40 models #1

############################################################
###   focus view and name structure models
############################################################
wait; focus #0#1
setattr m name GlpFs #0
setattr m name AQPs #1

############################################################
###   logo data
############################################################
setattr a radius 1.413 #0:7@ca
setattr a color #017A017A046F #0:7@ca
setattr a name '|| .LWMK' #0:7@ca
setattr a radius 1.408 #1:7@ca
setattr a color #17361752081A #1:7@ca
setattr a name '|| FSA' #1:7@ca
setattr a radius 0.795 #0:8@ca
setattr a color #368759C23733 #0:8@ca
setattr a name '|| .MKVA' #0:8@ca
setattr a radius 0.799 #1:8@ca
setattr a color #26AC26AC0000 #1:8@ca
setattr a name '|| .WY' #1:8@ca
setattr a radius 1.377 #0:9@ca
setattr a color #BB08AF463B97 #0:9@ca
setattr a name '|| .VSGTAN' #0:9@ca
setattr a radius 1.386 #1:9@ca
setattr a color #00000137DFF9 #1:9@ca
setattr a name '|| .RK' #1:9@ca
setattr a radius 2.223 #0:10@ca
setattr a color #A97E52390E6C #0:10@ca
setattr a name '|| .EQKWND' #0:10@ca
setattr a radius 2.225 #1:10@ca
setattr a color #F078F5B6E5FC #1:10@ca
setattr a name '|| A:.' #1:10@ca
setattr a radius 0.931 #0:11@ca
setattr a color #33DA33DA03B6 #0:11@ca
setattr a name '|| .FCT' #0:11@ca
setattr a radius 0.930 #1:11@ca
setattr a color #ABB8ABB867E3 #1:11@ca
setattr a name '|| .VGL' #1:11@ca
setattr a radius 0.906 #0:12@ca
setattr a color #02C102C10000 #0:12@ca
setattr a name '|| .L' #0:12@ca
setattr a radius 0.907 #1:12@ca
setattr a color #51EA51EA37E3 #1:12@ca
setattr a name '|| .IVT' #1:12@ca
setattr a radius 1.929 #0:13@ca
setattr a color #E361E3610000 #0:13@ca
setattr a name '|| .GSTM' #0:13@ca
setattr a radius 1.931 #1:13@ca
setattr a color #F042F042F042 #1:13@ca
setattr a name '|| A.V' #1:13@ca
setattr a radius 0.001 #0:14@ca
setattr a color #FFFF00000000 #0:14@ca
setattr a name '|| .' #0:14@ca
setattr a radius 0.001 #1:14@ca
setattr a color #000000000000 #1:14@ca
setattr a name '|| .' #1:14@ca
setattr a radius 0.947 #0:15@ca
setattr a color #DA0DDA0D2AE9 #0:15@ca
setattr a name '|| .CVTG' #0:15@ca
setattr a radius 0.949 #1:15@ca
setattr a color #000000000000 #1:15@ca
setattr a name '|| .F' #1:15@ca
setattr a radius 1.205 #0:16@ca
setattr a color #008300830000 #0:16@ca
setattr a name '|| L.' #0:16@ca
setattr a radius 1.201 #1:16@ca
setattr a color #23C025AE1D9F #1:16@ca
setattr a name '|| .IV' #1:16@ca
setattr a radius 2.268 #0:17@ca
setattr a color #FFFFFFFF0000 #0:17@ca
setattr a name '|| G:.' #0:17@ca
setattr a radius 2.265 #1:17@ca
setattr a color #FFA8FFA8D161 #1:17@ca
setattr a name '|| A.S' #1:17@ca
setattr a radius 0.610 #0:18@ca
setattr a color #FFFFFFFF0000 #0:18@ca
setattr a name '|| .T' #0:18@ca
setattr a radius 0.610 #1:18@ca
setattr a color #0A3D0A3D0000 #1:18@ca
setattr a name '|| .MS' #1:18@ca
setattr a radius 1.223 #0:19@ca
setattr a color #B26AB26A77B3 #0:19@ca
setattr a name '|| .AMYG' #0:19@ca
setattr a radius 1.224 #1:19@ca
setattr a color #00C100C10000 #1:19@ca
setattr a name '|| .LFI' #1:19@ca
setattr a radius 1.021 #0:20@ca
setattr a color #3F1F3F1F3F1F #0:20@ca
setattr a name '|| .IV' #0:20@ca
setattr a radius 1.016 #1:20@ca
setattr a color #00BA00BA0000 #1:20@ca
setattr a name '|| .LWM' #1:20@ca
setattr a radius 3.234 #0:21@ca
setattr a color #004900490049 #0:21@ca
setattr a name '|| L:.MV' #0:21@ca
setattr a radius 3.235 #1:21@ca
setattr a color #006A006A0000 #1:21@ca
setattr a name '|| F>:.Y' #1:21@ca
setattr a radius 2.123 #0:22@ca
setattr a color #000000000000 #0:22@ca
setattr a name '|| I.M' #0:22@ca
setattr a radius 2.122 #1:22@ca
setattr a color #503B503B503B #1:22@ca
setattr a name '|| V.L' #1:22@ca
setattr a radius 1.631 #0:23@ca
setattr a color #198619861986 #0:23@ca
setattr a name '|| .IVLMA' #0:23@ca
setattr a radius 1.630 #1:23@ca
setattr a color #494249420000 #1:23@ca
setattr a name '|| .YF' #1:23@ca
setattr a radius 1.693 #0:24@ca
setattr a color #005A005A0000 #0:24@ca
setattr a name '|| .FL' #0:24@ca
setattr a radius 1.691 #1:24@ca
setattr a color #47B447B4371E #1:24@ca
setattr a name '|| .IAVG' #1:24@ca
setattr a radius 1.749 #0:25@ca
setattr a color #FFFFFFFF0000 #0:25@ca
setattr a name '|| G.' #0:25@ca
setattr a radius 1.746 #1:25@ca
setattr a color #FF8EFF8E163F #1:25@ca
setattr a name '|| T.SA' #1:25@ca
setattr a radius 0.938 #0:26@ca
setattr a color #FFFFFFFF26C1 #0:26@ca
setattr a name '|| .NCAGT' #0:26@ca
setattr a radius 0.936 #1:26@ca
setattr a color #849179BF529B #1:26@ca
setattr a name '|| .VIS' #1:26@ca
setattr a radius 1.511 #0:27@ca
setattr a color #FFFFFFFF0000 #0:27@ca
setattr a name '|| G.S' #0:27@ca
setattr a radius 1.507 #1:27@ca
setattr a color #201720172017 #1:27@ca
setattr a name '|| .LAV' #1:27@ca
setattr a radius 1.494 #0:28@ca
setattr a color #FFFFFFFF8BA6 #0:28@ca
setattr a name '|| .VCA' #0:28@ca
setattr a radius 1.496 #1:28@ca
setattr a color #FA8BFA8B0000 #1:28@ca
setattr a name '|| .TS' #1:28@ca
setattr a radius 1.126 #0:29@ca
setattr a color #9E249E24761D #0:29@ca
setattr a name '|| .VIC' #0:29@ca
setattr a radius 1.125 #1:29@ca
setattr a color #FFC4FFC4C29C #1:29@ca
setattr a name '|| .ATG' #1:29@ca
setattr a radius 1.985 #0:30@ca
setattr a color #FFFFFFFFF274 #0:30@ca
setattr a name '|| A.' #0:30@ca
setattr a radius 1.980 #1:30@ca
setattr a color #055E055E033D #1:30@ca
setattr a name '|| .IMLV' #1:30@ca
setattr a radius 0.860 #0:31@ca
setattr a color #73229C7B0000 #0:31@ca
setattr a name '|| .QHM' #0:31@ca
setattr a radius 0.853 #1:31@ca
setattr a color #90879CBE1EEE #1:31@ca
setattr a name '|| .G' #1:31@ca
setattr a radius 0.902 #0:32@ca
setattr a color #6306630650D8 #0:32@ca
setattr a name '|| .VLMT' #0:32@ca
setattr a radius 0.905 #1:32@ca
setattr a color #2A49355B0000 #1:32@ca
setattr a name '|| .YWF' #1:32@ca
setattr a radius 0.680 #0:33@ca
setattr a color #E7C0CD29A547 #0:33@ca
setattr a name '|| .VT' #0:33@ca
setattr a radius 0.675 #1:33@ca
setattr a color #B909CA8F530E #1:33@ca
setattr a name '|| .ASG' #1:33@ca
setattr a radius 1.038 #0:34@ca
setattr a color #07C807C807C8 #0:34@ca
setattr a name '|| .LV' #0:34@ca
setattr a radius 1.025 #1:34@ca
setattr a color #038A7C8944CB #1:34@ca
setattr a name '|| .KRS' #1:34@ca
setattr a radius 0.578 #0:35@ca
setattr a color #0D49A28B664F #0:35@ca
setattr a name '|| .K' #0:35@ca
setattr a radius 0.576 #1:35@ca
setattr a color #292A184A076A #1:35@ca
setattr a name '|| .L' #1:35@ca
setattr a radius 0.646 #0:36@ca
setattr a color #00008D9294F1 #0:36@ca
setattr a name '|| .KGR' #0:36@ca
setattr a radius 0.640 #1:36@ca
setattr a color #B433A77E4332 #1:36@ca
setattr a name '|| .TP' #1:36@ca
setattr a radius 0.452 #0:37@ca
setattr a color #9ECE9ECE1FAB #0:37@ca
setattr a name '|| .TG' #0:37@ca
setattr a radius 0.450 #1:37@ca
setattr a color #CD8F816D2FDB #1:37@ca
setattr a name '|| .DN' #1:37@ca
setattr a radius 0.422 #0:38@ca
setattr a color #0D429CCA926A #0:38@ca
setattr a name '|| .KS' #0:38@ca
setattr a radius 0.418 #1:38@ca
setattr a color #D4E6630C3B33 #1:38@ca
setattr a name '|| .DP' #1:38@ca
setattr a radius 0.356 #0:39@ca
setattr a color #31ED4EDB2A76 #0:39@ca
setattr a name '|| .F' #0:39@ca
setattr a radius 0.356 #1:39@ca
setattr a color #1C6865342C63 #1:39@ca
setattr a name '|| .KT' #1:39@ca
setattr a radius 0.518 #0:40@ca
setattr a color #A11CBD470000 #0:40@ca
setattr a name '|| .G' #0:40@ca
setattr a radius 0.521 #1:40@ca
setattr a color #CCF6CCF6B3D6 #1:40@ca
setattr a name '|| .A' #1:40@ca
setattr a radius 0.325 #0:41@ca
setattr a color #406540650587 #0:41@ca
setattr a name '|| .' #0:41@ca
setattr a radius 0.330 #1:41@ca
setattr a color #000000000000 #1:41@ca
setattr a name '|| .W' #1:41@ca
setattr a radius 0.588 #0:42@ca
setattr a color #020001590000 #0:42@ca
setattr a name '|| .LI' #0:42@ca
setattr a radius 0.593 #1:42@ca
setattr a color #000000000000 #1:42@ca
setattr a name '|| .F' #1:42@ca
setattr a radius 0.355 #0:43@ca
setattr a color #D53DA87C33FE #0:43@ca
setattr a name '|| .T' #0:43@ca
setattr a radius 0.364 #1:43@ca
setattr a color #379937993799 #1:43@ca
setattr a name '|| .FA' #1:43@ca
setattr a radius 0.566 #0:44@ca
setattr a color #000000000000 #0:44@ca
setattr a name '|| .IL' #0:44@ca
setattr a radius 0.564 #1:44@ca
setattr a color #FF17FF17F185 #1:44@ca
setattr a name '|| .AVP' #1:44@ca
setattr a radius 1.724 #0:45@ca
setattr a color #FFFFFFFF2A2B #0:45@ca
setattr a name '|| N.TV' #0:45@ca
setattr a radius 1.728 #1:45@ca
setattr a color #F855F855F855 #1:45@ca
setattr a name '|| A.' #1:45@ca
setattr a radius 0.698 #0:46@ca
setattr a color #22F922F9119F #0:46@ca
setattr a name '|| .FA' #0:46@ca
setattr a radius 0.703 #1:46@ca
setattr a color #01CB01CB01CB #1:46@ca
setattr a name '|| .WL' #1:46@ca
setattr a radius 2.055 #0:47@ca
setattr a color #932593250000 #0:47@ca
setattr a name '|| G.IM' #0:47@ca
setattr a radius 2.056 #1:47@ca
setattr a color #F01EF01EBADB #1:47@ca
setattr a name '|| A.CS' #1:47@ca
setattr a radius 2.214 #0:48@ca
setattr a color #079F079F01E5 #0:48@ca
setattr a name '|| W.Y' #0:48@ca
setattr a radius 2.212 #1:48@ca
setattr a color #000002BD253F #1:48@ca
setattr a name '|| F.H' #1:48@ca
setattr a radius 0.626 #0:49@ca
setattr a color #FFFFFFFF0000 #0:49@ca
setattr a name '|| .GS' #0:49@ca
setattr a radius 0.626 #1:49@ca
setattr a color #FFFFFFFFE153 #1:49@ca
setattr a name '|| .A' #1:49@ca
setattr a radius 0.995 #0:50@ca
setattr a color #19D219D20CE9 #0:50@ca
setattr a name '|| .IFVM' #0:50@ca
setattr a radius 0.998 #1:50@ca
setattr a color #CC96CC960000 #1:50@ca
setattr a name '|| .GL' #1:50@ca
setattr a radius 1.489 #0:51@ca
setattr a color #FFFFFFFF6E34 #0:51@ca
setattr a name '|| .GA' #0:51@ca
setattr a radius 1.485 #1:51@ca
setattr a color #31B031B00834 #1:51@ca
setattr a name '|| .MSTV' #1:51@ca
setattr a radius 2.601 #0:52@ca
setattr a color #FFFFFFFFF6F0 #0:52@ca
setattr a name '|| V:.' #0:52@ca
setattr a radius 2.601 #1:52@ca
setattr a color #004400440016 #1:52@ca
setattr a name '|| I.L' #1:52@ca
setattr a radius 1.498 #0:53@ca
setattr a color #8A498A491FA6 #0:53@ca
setattr a name '|| .TMLSV' #0:53@ca
setattr a radius 1.502 #1:53@ca
setattr a color #06AE06AE0000 #1:53@ca
setattr a name '|| F.G' #1:53@ca
setattr a radius 1.328 #0:54@ca
setattr a color #000000000000 #0:54@ca
setattr a name '|| .MLF' #0:54@ca
setattr a radius 1.328 #1:54@ca
setattr a color #E272E2728541 #1:54@ca
setattr a name '|| .TAV' #1:54@ca
setattr a radius 1.775 #0:55@ca
setattr a color #FFFFFFFF7892 #0:55@ca
setattr a name '|| .GAP' #0:55@ca
setattr a radius 1.769 #1:55@ca
setattr a color #024A024A0091 #1:55@ca
setattr a name '|| L.MS' #1:55@ca
setattr a radius 1.021 #0:56@ca
setattr a color #000000000000 #0:56@ca
setattr a name '|| .IL' #0:56@ca
setattr a radius 1.021 #1:56@ca
setattr a color #FFFFFFFFFFFF #1:56@ca
setattr a name '|| .VA' #1:56@ca
setattr a radius 1.326 #0:57@ca
setattr a color #075516E60529 #0:57@ca
setattr a name '|| .LYHF' #0:57@ca
setattr a radius 1.329 #1:57@ca
setattr a color #D9F2D9F221CD #1:57@ca
setattr a name '|| .QSAW' #1:57@ca
setattr a radius 1.095 #0:58@ca
setattr a color #5D915D914A58 #0:58@ca
setattr a name '|| .VLITA' #0:58@ca
setattr a radius 1.100 #1:58@ca
setattr a color #F12AF12A0000 #1:58@ca
setattr a name '|| .CS' #1:58@ca
setattr a radius 1.021 #0:59@ca
setattr a color #FF7EFF7EA2D0 #0:59@ca
setattr a name '|| .AS' #0:59@ca
setattr a radius 1.015 #1:59@ca
setattr a color #82DD82DD07BD #1:59@ca
setattr a name '|| .TLG' #1:59@ca
setattr a radius 1.235 #0:60@ca
setattr a color #9E929E920000 #0:60@ca
setattr a name '|| .GFNS' #0:60@ca
setattr a radius 1.236 #1:60@ca
setattr a color #C6DDC6DDC6DD #1:60@ca
setattr a name '|| A.' #1:60@ca
setattr a radius 1.099 #0:61@ca
setattr a color #6465B2321012 #0:61@ca
setattr a name '|| .QSRK' #0:61@ca
setattr a radius 1.098 #1:61@ca
setattr a color #0D0BA4737082 #1:61@ca
setattr a name '|| .HN' #1:61@ca
setattr a radius 1.551 #0:62@ca
setattr a color #2B512B511F4C #0:62@ca
setattr a name '|| .LVMYF' #0:62@ca
setattr a radius 1.557 #1:62@ca
setattr a color #001500150015 #1:62@ca
setattr a name '|| I.' #1:62@ca
setattr a radius 0.725 #0:63@ca
setattr a color #5C285C281EB8 #0:63@ca
setattr a name '|| .FTNA' #0:63@ca
setattr a radius 0.726 #1:63@ca
setattr a color #F059F0590000 #1:63@ca
setattr a name '|| .S' #1:63@ca
setattr a radius 0.887 #0:64@ca
setattr a color #FFFFFFFF209C #0:64@ca
setattr a name '|| .G' #0:64@ca
setattr a radius 0.001 #1:64@ca
setattr a color #000000000000 #1:64@ca
setattr a name '|| .' #1:64@ca
setattr a radius 2.076 #0:65@ca
setattr a color #FFFFFFFFE45F #0:65@ca
setattr a name '|| A.SN' #0:65@ca
setattr a radius 2.076 #1:65@ca
setattr a color #FFFFFFFF0000 #1:65@ca
setattr a name '|| G.C' #1:65@ca
setattr a radius 0.240 #0:66@ca
setattr a color #FFFFFFFF0000 #0:66@ca
setattr a name '|| .Q' #0:66@ca
setattr a radius 0.240 #1:66@ca
setattr a color #B554B554FFFF #1:66@ca
setattr a name '|| .A' #1:66@ca
setattr a radius 1.657 #0:67@ca
setattr a color #000000000000 #0:67@ca
setattr a name '|| L.M' #0:67@ca
setattr a radius 1.655 #1:67@ca
setattr a color #248924892489 #1:67@ca
setattr a name '|| .IVF' #1:67@ca
setattr a radius 0.001 #0:68@ca
setattr a color #FFFFFFFF0000 #0:68@ca
setattr a name '|| .' #0:68@ca
setattr a radius 0.001 #1:68@ca
setattr a color #000000000000 #1:68@ca
setattr a name '|| .' #1:68@ca
setattr a radius 0.302 #0:69@ca
setattr a color #FFFFFFFFFFFF #0:69@ca
setattr a name '|| .A' #0:69@ca
setattr a radius 0.302 #1:69@ca
setattr a color #FFFFFFFFFFFF #1:69@ca
setattr a name '|| .P' #1:69@ca
setattr a radius 0.045 #0:70@ca
setattr a color #FFFFFFFFFFFF #0:70@ca
setattr a name '|| .' #0:70@ca
setattr a radius 0.045 #1:70@ca
setattr a color #FFFFFFFF0000 #1:70@ca
setattr a name '|| .' #1:70@ca
setattr a radius 0.832 #0:71@ca
setattr a color #000000000000 #0:71@ca
setattr a name '|| .ILF' #0:71@ca
setattr a radius 0.833 #1:71@ca
setattr a color #E844E844E844 #1:71@ca
setattr a name '|| .V' #1:71@ca
setattr a radius 1.245 #0:72@ca
setattr a color #FFFFFFFF0000 #0:72@ca
setattr a name '|| S.' #0:72@ca
setattr a radius 1.245 #1:72@ca
setattr a color #FFFFFFFF0804 #1:72@ca
setattr a name '|| T.' #1:72@ca
setattr a radius 1.159 #0:73@ca
setattr a color #000000000000 #0:73@ca
setattr a name '|| I.' #0:73@ca
setattr a radius 1.155 #1:73@ca
setattr a color #0393039302E3 #1:73@ca
setattr a name '|| .FLV' #1:73@ca
setattr a radius 1.632 #0:74@ca
setattr a color #763676365FEF #0:74@ca
setattr a name '|| .AFTM' #0:74@ca
setattr a radius 1.635 #1:74@ca
setattr a color #F077F0770000 #1:74@ca
setattr a name '|| G.' #1:74@ca
setattr a radius 0.982 #0:75@ca
setattr a color #1CE81CE80000 #0:75@ca
setattr a name '|| .MNFY' #0:75@ca
setattr a radius 0.983 #1:75@ca
setattr a color #2EB72EB71923 #1:75@ca
setattr a name '|| .LAC' #1:75@ca
setattr a radius 1.654 #0:76@ca
setattr a color #BF5ABF5A654C #0:76@ca
setattr a name '|| .ACW' #0:76@ca
setattr a radius 1.653 #1:76@ca
setattr a color #02C802C801DC #1:76@ca
setattr a name '|| .FLV' #1:76@ca
setattr a radius 1.133 #0:77@ca
setattr a color #562E603A2CB1 #0:77@ca
setattr a name '|| .AFSTI' #0:77@ca
setattr a radius 1.138 #1:77@ca
setattr a color #252125211C3B #1:77@ca
setattr a name '|| .LVC' #1:77@ca
setattr a radius 1.622 #0:78@ca
setattr a color #10FC18230C19 #0:78@ca
setattr a name '|| .FLVSK' #0:78@ca
setattr a radius 1.622 #1:78@ca
setattr a color #EFD9EFD966F7 #1:78@ca
setattr a name '|| .GAT' #1:78@ca
setattr a radius 1.465 #0:79@ca
setattr a color #FFFFFFFF7BB7 #0:79@ca
setattr a name '|| .GA' #0:79@ca
setattr a radius 1.460 #1:79@ca
setattr a color #0000A8A0AA5F #1:79@ca
setattr a name '|| .RC' #1:79@ca
setattr a radius 1.045 #0:80@ca
setattr a color #5F5167C712B5 #0:80@ca
setattr a name '|| .RDC' #0:80@ca
setattr a radius 1.048 #1:80@ca
setattr a color #01B689EB1A74 #1:80@ca
setattr a name '|| .KQHN' #1:80@ca
setattr a radius 1.682 #0:81@ca
setattr a color #0C0E02690000 #0:81@ca
setattr a name '|| F.ME' #0:81@ca
setattr a radius 1.682 #1:81@ca
setattr a color #149614961496 #1:81@ca
setattr a name '|| I.VL' #1:81@ca
setattr a radius 2.209 #0:82@ca
setattr a color #DC788A5884A0 #0:82@ca
setattr a name '|| .PEDA' #0:82@ca
setattr a radius 2.213 #1:82@ca
setattr a color #F075F0750000 #1:82@ca
setattr a name '|| S.T' #1:82@ca
setattr a radius 1.917 #0:83@ca
setattr a color #00C206860000 #0:83@ca
setattr a name '|| W.GK' #0:83@ca
setattr a radius 1.914 #1:83@ca
setattr a color #042104210376 #1:83@ca
setattr a name '|| L.IA' #1:83@ca
setattr a radius 0.816 #0:84@ca
setattr a color #0A7E6C7585C1 #0:84@ca
setattr a name '|| .RSK' #0:84@ca
setattr a radius 0.816 #1:84@ca
setattr a color #289C289C095F #1:84@ca
setattr a name '|| .LT' #1:84@ca
setattr a radius 1.922 #0:85@ca
setattr a color #0BCE277CDA6C #0:85@ca
setattr a name '|| K.H' #0:85@ca
setattr a radius 1.924 #1:85@ca
setattr a color #00EC1B66E059 #1:85@ca
setattr a name '|| R.T' #1:85@ca
setattr a radius 1.981 #0:86@ca
setattr a color #3A413A41375B #0:86@ca
setattr a name '|| .VFL' #0:86@ca
setattr a radius 1.981 #1:86@ca
setattr a color #FFFFFFFFA1A9 #1:86@ca
setattr a name '|| A.GS' #1:86@ca
setattr a radius 1.349 #0:87@ca
setattr a color #CEFCCEFCCEFC #0:87@ca
setattr a name '|| P.L' #0:87@ca
setattr a radius 1.349 #1:87@ca
setattr a color #54B654B654B6 #1:87@ca
setattr a name '|| .VI' #1:87@ca
setattr a radius 0.824 #0:88@ca
setattr a color #6C82805E413F #0:88@ca
setattr a name '|| .IVQPG' #0:88@ca
setattr a radius 0.823 #1:88@ca
setattr a color #0A530A5302E1 #1:88@ca
setattr a name '|| .FLM' #1:88@ca
setattr a radius 0.423 #0:89@ca
setattr a color #71C671C638E3 #0:89@ca
setattr a name '|| .GFA' #0:89@ca
setattr a radius 0.425 #1:89@ca
setattr a color #5C4E5C4E0000 #1:89@ca
setattr a name '|| .YM' #1:89@ca
setattr a radius 1.094 #0:90@ca
setattr a color #5AA65AA6252D #0:90@ca
setattr a name '|| .ISVTF' #0:90@ca
setattr a radius 1.100 #1:90@ca
setattr a color #000000000000 #1:90@ca
setattr a name '|| .WM' #1:90@ca
setattr a radius 1.272 #0:91@ca
setattr a color #06FA06FA0000 #0:91@ca
setattr a name '|| L.G' #0:91@ca
setattr a radius 1.269 #1:91@ca
setattr a color #2572257223D3 #1:91@ca
setattr a name '|| .IVA' #1:91@ca
setattr a radius 0.741 #0:92@ca
setattr a color #FFFFFFFF6D83 #0:92@ca
setattr a name '|| .GA' #0:92@ca
setattr a radius 0.741 #1:92@ca
setattr a color #064603230323 #1:92@ca
setattr a name '|| .M' #1:92@ca
setattr a radius 0.001 #0:93@ca
setattr a color #FFFFFFFF0000 #0:93@ca
setattr a name '|| .' #0:93@ca
setattr a radius 0.001 #1:93@ca
setattr a color #000000000000 #1:93@ca
setattr a name '|| .' #1:93@ca
setattr a radius 1.560 #0:94@ca
setattr a color #31CA31CA147B #0:94@ca
setattr a name '|| .FTVM' #0:94@ca
setattr a radius 1.561 #1:94@ca
setattr a color #969896980000 #1:94@ca
setattr a name '|| C.L' #1:94@ca
setattr a radius 0.472 #0:95@ca
setattr a color #FEEAFEEAFEEA #0:95@ca
setattr a name '|| .A' #0:95@ca
setattr a radius 0.467 #1:95@ca
setattr a color #650265026502 #1:95@ca
setattr a name '|| .VL' #1:95@ca
setattr a radius 0.305 #0:96@ca
setattr a color #FFFFFFFF0000 #0:96@ca
setattr a name '|| .G' #0:96@ca
setattr a radius 0.305 #1:96@ca
setattr a color #FFFFFFFFFFFF #1:96@ca
setattr a name '|| .A' #1:96@ca
setattr a radius 0.728 #0:97@ca
setattr a color #FFFFFFFFFFFF #0:97@ca
setattr a name '|| .A' #0:97@ca
setattr a radius 0.726 #1:97@ca
setattr a color #FF4AFF4A0000 #1:97@ca
setattr a name '|| .GS' #1:97@ca
setattr a radius 2.017 #0:98@ca
setattr a color #005500550000 #0:98@ca
setattr a name '|| F.M' #0:98@ca
setattr a radius 2.014 #1:98@ca
setattr a color #1FF11FF113DB #1:98@ca
setattr a name '|| I.VT' #1:98@ca
setattr a radius 1.160 #0:99@ca
setattr a color #06A706A70000 #0:99@ca
setattr a name '|| .LFIT' #0:99@ca
setattr a radius 1.165 #1:99@ca
setattr a color #F119F119A512 #1:99@ca
setattr a name '|| .ACV' #1:99@ca
setattr a radius 0.659 #0:100@ca
setattr a color #B060B0600000 #0:100@ca
setattr a name '|| .GSF' #0:100@ca
setattr a radius 0.657 #1:100@ca
setattr a color #FF38FF38F2C9 #1:100@ca
setattr a name '|| .A' #1:100@ca
setattr a radius 1.111 #0:101@ca
setattr a color #FFFFFFFF8706 #0:101@ca
setattr a name '|| .AQ' #0:101@ca
setattr a radius 1.108 #1:101@ca
setattr a color #FF4EFF4E53DF #1:101@ca
setattr a name '|| .CVS' #1:101@ca
setattr a radius 1.173 #0:102@ca
setattr a color #AD9DAD9D76CC #0:102@ca
setattr a name '|| .ATV' #0:102@ca
setattr a radius 1.179 #1:102@ca
setattr a color #9DD19DD10000 #1:102@ca
setattr a name '|| .GF' #1:102@ca
setattr a radius 0.846 #0:103@ca
setattr a color #4DAC4DAC0000 #0:103@ca
setattr a name '|| .TML' #0:103@ca
setattr a radius 0.842 #1:103@ca
setattr a color #A090A0909BC1 #1:103@ca
setattr a name '|| .VFA' #1:103@ca
setattr a radius 1.630 #0:104@ca
setattr a color #C4C1C4C182AA #0:104@ca
setattr a name '|| .VTI' #0:104@ca
setattr a radius 1.630 #1:104@ca
setattr a color #000100010001 #1:104@ca
setattr a name '|| L.' #1:104@ca
setattr a radius 1.448 #0:105@ca
setattr a color #4A794A790B40 #0:105@ca
setattr a name '|| .YFWV' #0:105@ca
setattr a radius 1.450 #1:105@ca
setattr a color #04F34DF7C711 #1:105@ca
setattr a name '|| .KS' #1:105@ca
setattr a radius 0.532 #0:106@ca
setattr a color #24E824E8028E #0:106@ca
setattr a name '|| .LTM' #0:106@ca
setattr a radius 0.532 #1:106@ca
setattr a color #8C6971544C9F #1:106@ca
setattr a name '|| .AF' #1:106@ca
setattr a radius 0.943 #0:107@ca
setattr a color #01DA0FD90000 #0:107@ca
setattr a name '|| .LHQ' #0:107@ca
setattr a radius 0.940 #1:107@ca
setattr a color #242D242D1B0E #1:107@ca
setattr a name '|| .FV' #1:107@ca
setattr a radius 1.773 #0:108@ca
setattr a color #835983590000 #0:108@ca
setattr a name '|| Y.F' #0:108@ca
setattr a radius 1.766 #1:108@ca
setattr a color #D37BC4C01F4D #1:108@ca
setattr a name '|| .TQAME' #1:108@ca
setattr a radius 1.159 #0:109@ca
setattr a color #471970900000 #0:109@ca
setattr a name '|| .YWR' #0:109@ca
setattr a radius 1.149 #1:109@ca
setattr a color #F2E9F0516CEB #1:109@ca
setattr a name '|| .PGS' #1:109@ca
setattr a radius 0.808 #0:110@ca
setattr a color #E5797E494CCA #0:110@ca
setattr a name '|| .DPN' #0:110@ca
setattr a radius 0.804 #1:110@ca
setattr a color #CEDECE913B71 #1:110@ca
setattr a name '|| .GA' #1:110@ca
setattr a radius 0.338 #0:111@ca
setattr a color #7E7B92EEEE93 #0:111@ca
setattr a name '|| .AH' #0:111@ca
setattr a radius 0.328 #1:111@ca
setattr a color #38444C451901 #1:111@ca
setattr a name '|| .' #1:111@ca
setattr a radius 0.567 #0:112@ca
setattr a color #01D501D50000 #0:112@ca
setattr a name '|| .IF' #0:112@ca
setattr a radius 0.562 #1:112@ca
setattr a color #8CDDBA777D83 #1:112@ca
setattr a name '|| .AR' #1:112@ca
setattr a radius 0.612 #0:113@ca
setattr a color #034F1EF810C0 #0:113@ca
setattr a name '|| .KWM' #0:113@ca
setattr a radius 0.612 #1:113@ca
setattr a color #F1F9F1F95E59 #1:113@ca
setattr a name '|| .YV' #1:113@ca
setattr a radius 0.416 #0:114@ca
setattr a color #8B139313B13A #0:114@ca
setattr a name '|| .A' #0:114@ca
setattr a radius 0.409 #1:114@ca
setattr a color #7322989E585B #1:114@ca
setattr a name '|| .P' #1:114@ca
setattr a radius 0.668 #0:115@ca
setattr a color #088E0B5900CF #0:115@ca
setattr a name '|| .FY' #0:115@ca
setattr a radius 0.657 #1:115@ca
setattr a color #909AA34E11A0 #1:115@ca
setattr a name '|| .GR' #1:115@ca
setattr a radius 0.718 #0:116@ca
setattr a color #E03F76216F74 #0:116@ca
setattr a name '|| .AE' #0:116@ca
setattr a radius 0.713 #1:116@ca
setattr a color #014B23BF067B #1:116@ca
setattr a name '|| .HFL' #1:116@ca
setattr a radius 0.777 #0:117@ca
setattr a color #94C899283483 #0:117@ca
setattr a name '|| .QP' #0:117@ca
setattr a radius 0.781 #1:117@ca
setattr a color #D3D8EA941C55 #1:117@ca
setattr a name '|| .SG' #1:117@ca
setattr a radius 1.279 #0:118@ca
setattr a color #A921A0853CA5 #0:118@ca
setattr a name '|| .NKADT' #0:118@ca
setattr a radius 1.280 #1:118@ca
setattr a color #2A8F2A8F01B6 #1:118@ca
setattr a name '|| .LG' #1:118@ca
setattr a radius 0.180 #0:119@ca
setattr a color #3761647415E4 #0:119@ca
setattr a name '|| .' #0:119@ca
setattr a radius 0.190 #1:119@ca
setattr a color #FE37FE3738E3 #1:119@ca
setattr a name '|| .S' #1:119@ca
setattr a radius 0.288 #0:120@ca
setattr a color #0417105E12E3 #0:120@ca
setattr a name '|| .L' #0:120@ca
setattr a radius 0.296 #1:120@ca
setattr a color #2C652C650000 #1:120@ca
setattr a name '|| .FN' #1:120@ca
setattr a radius 0.242 #0:121@ca
setattr a color #1D640EB20A09 #0:121@ca
setattr a name '|| .' #0:121@ca
setattr a radius 0.248 #1:121@ca
setattr a color #F657F6574195 #1:121@ca
setattr a name '|| .G' #1:121@ca
setattr a radius 0.480 #0:122@ca
setattr a color #9A18A34A9A18 #0:122@ca
setattr a name '|| .V' #0:122@ca
setattr a radius 0.488 #1:122@ca
setattr a color #039703970000 #1:122@ca
setattr a name '|| .FL' #1:122@ca
setattr a radius 0.398 #0:123@ca
setattr a color #54E792A624A9 #0:123@ca
setattr a name '|| .RT' #0:123@ca
setattr a radius 0.399 #1:123@ca
setattr a color #F0A3F0A3EC0E #1:123@ca
setattr a name '|| .A' #1:123@ca
setattr a radius 0.131 #0:124@ca
setattr a color #FFFFCD306792 #0:124@ca
setattr a name '|| .' #0:124@ca
setattr a radius 0.136 #1:124@ca
setattr a color #F9C4F9C40000 #1:124@ca
setattr a name '|| .G' #1:124@ca
setattr a radius 0.489 #0:125@ca
setattr a color #FFFFCE4341F6 #0:125@ca
setattr a name '|| .T' #0:125@ca
setattr a radius 0.492 #1:125@ca
setattr a color #000000000000 #1:125@ca
setattr a name '|| .F' #1:125@ca
setattr a radius 0.325 #0:126@ca
setattr a color #AC0990AF0CF4 #0:126@ca
setattr a name '|| .' #0:126@ca
setattr a radius 0.335 #1:126@ca
setattr a color #000000000000 #1:126@ca
setattr a name '|| .W' #1:126@ca
setattr a radius 0.347 #0:127@ca
setattr a color #EA584C3B0A96 #0:127@ca
setattr a name '|| .ED' #0:127@ca
setattr a radius 0.355 #1:127@ca
setattr a color #F5E5F5E5F5E5 #1:127@ca
setattr a name '|| .P' #1:127@ca
setattr a radius 0.185 #0:128@ca
setattr a color #FFFFE7169999 #0:128@ca
setattr a name '|| .P' #0:128@ca
setattr a radius 0.181 #1:128@ca
setattr a color #4AB9AFBC6D51 #1:128@ca
setattr a name '|| .' #1:128@ca
setattr a radius 0.305 #0:129@ca
setattr a color #8668866830A7 #0:129@ca
setattr a name '|| .G' #0:129@ca
setattr a radius 0.311 #1:129@ca
setattr a color #F508CA470000 #1:129@ca
setattr a name '|| .N' #1:129@ca
setattr a radius 1.450 #0:130@ca
setattr a color #CE64B5BE5329 #0:130@ca
setattr a name '|| .ASND' #0:130@ca
setattr a radius 1.456 #1:130@ca
setattr a color #D616D6160000 #1:130@ca
setattr a name '|| G.F' #1:130@ca
setattr a radius 1.096 #0:131@ca
setattr a color #4E416E140000 #0:131@ca
setattr a name '|| .TKLI' #0:131@ca
setattr a radius 1.094 #1:131@ca
setattr a color #C295CDC6BB0C #1:131@ca
setattr a name '|| .AV' #1:131@ca
setattr a radius 1.891 #0:132@ca
setattr a color #60556055520C #0:132@ca
setattr a name '|| .ALQ' #0:132@ca
setattr a radius 1.885 #1:132@ca
setattr a color #E288E28802B5 #1:132@ca
setattr a name '|| N.T' #1:132@ca
setattr a radius 0.560 #0:133@ca
setattr a color #CA26E49D0000 #0:133@ca
setattr a name '|| .G' #0:133@ca
setattr a radius 0.555 #1:133@ca
setattr a color #D667A5A34AD4 #1:133@ca
setattr a name '|| .TV' #1:133@ca
setattr a radius 0.404 #0:134@ca
setattr a color #4B394B39369D #0:134@ca
setattr a name '|| .IV' #0:134@ca
setattr a radius 0.405 #1:134@ca
setattr a color #180781C1A967 #1:134@ca
setattr a name '|| .H' #1:134@ca
setattr a radius 2.457 #0:135@ca
setattr a color #0000000000FA #0:135@ca
setattr a name '|| F:.MH' #0:135@ca
setattr a radius 2.453 #1:135@ca
setattr a color #4E694FA444DA #1:135@ca
setattr a name '|| VL.S' #1:135@ca
setattr a radius 0.822 #0:136@ca
setattr a color #DB69DB6915E9 #0:136@ca
setattr a name '|| .SCT' #0:136@ca
setattr a radius 0.821 #1:136@ca
setattr a color #7429B4513A3C #1:136@ca
setattr a name '|| .HPNQG' #1:136@ca
setattr a radius 1.179 #0:137@ca
setattr a color #FFFFFFFF056D #0:137@ca
setattr a name '|| T.' #0:137@ca
setattr a radius 1.168 #1:137@ca
setattr a color #A4EAA0014F2F #1:137@ca
setattr a name '|| .PGANHD' #1:137@ca
setattr a radius 1.997 #0:138@ca
setattr a color #8DF07F6E04F6 #0:138@ca
setattr a name '|| Y.FDTI' #0:138@ca
setattr a radius 1.998 #1:138@ca
setattr a color #F413F72602F2 #1:138@ca
setattr a name '|| G.N' #1:138@ca
setattr a radius 1.462 #0:139@ca
setattr a color #FBD0F37AE804 #0:139@ca
setattr a name '|| P.' #0:139@ca
setattr a radius 1.454 #1:139@ca
setattr a color #D7B2D30551BB #1:139@ca
setattr a name '|| .YV' #1:139@ca
setattr a radius 1.190 #0:140@ca
setattr a color #BABAB66CFBB1 #0:140@ca
setattr a name '|| .APHK' #0:140@ca
setattr a radius 1.187 #1:140@ca
setattr a color #F8E5F5740000 #1:140@ca
setattr a name '|| .TSG' #1:140@ca
setattr a radius 0.613 #0:141@ca
setattr a color #75554E993B17 #0:141@ca
setattr a name '|| .IP' #0:141@ca
setattr a radius 0.613 #1:141@ca
setattr a color #31C9A8B9A983 #1:141@ca
setattr a name '|| .KA' #1:141@ca
setattr a radius 1.519 #0:142@ca
setattr a color #328A7D2C7DC3 #0:142@ca
setattr a name '|| .HYPFKDR' #0:142@ca
setattr a radius 1.524 #1:142@ca
setattr a color #B2CEB2CE0437 #1:142@ca
setattr a name '|| G.L' #1:142@ca
setattr a radius 0.500 #0:143@ca
setattr a color #202B252F0AA7 #0:143@ca
setattr a name '|| .L' #0:143@ca
setattr a radius 0.502 #1:143@ca
setattr a color #DB21C9B923D6 #1:143@ca
setattr a name '|| .Q' #1:143@ca
setattr a radius 2.341 #0:144@ca
setattr a color #FFFFD99C0000 #0:144@ca
setattr a name '|| S.TND' #0:144@ca
setattr a radius 2.338 #1:144@ca
setattr a color #EBBFEBBF6698 #1:144@ca
setattr a name '|| G.A' #1:144@ca
setattr a radius 0.288 #0:145@ca
setattr a color #448544851350 #0:145@ca
setattr a name '|| .' #0:145@ca
setattr a radius 0.296 #1:145@ca
setattr a color #000000000000 #1:145@ca
setattr a name '|| .L' #1:145@ca
setattr a radius 0.251 #0:146@ca
setattr a color #5C63618F4C3D #0:146@ca
setattr a name '|| .' #0:146@ca
setattr a radius 0.258 #1:146@ca
setattr a color #000000000000 #1:146@ca
setattr a name '|| .W' #1:146@ca
setattr a radius 0.348 #0:147@ca
setattr a color #9B80C7372BB6 #0:147@ca
setattr a name '|| .' #0:147@ca
setattr a radius 0.362 #1:147@ca
setattr a color #00000000FA4B #1:147@ca
setattr a name '|| .K' #1:147@ca
setattr a radius 0.458 #0:148@ca
setattr a color #FFFFF4D17FFF #0:148@ca
setattr a name '|| .NA' #0:148@ca
setattr a radius 0.462 #1:148@ca
setattr a color #F23BF23B0000 #1:148@ca
setattr a name '|| .G' #1:148@ca
setattr a radius 2.142 #0:149@ca
setattr a color #000000000000 #0:149@ca
setattr a name '|| F:.' #0:149@ca
setattr a radius 2.135 #1:149@ca
setattr a color #19C219C215CC #1:149@ca
setattr a name '|| L.V' #1:149@ca
setattr a radius 1.282 #0:150@ca
setattr a color #0581058103D6 #0:150@ca
setattr a name '|| .LIFA' #0:150@ca
setattr a radius 1.283 #1:150@ca
setattr a color #EC39EC3945CB #1:150@ca
setattr a name '|| .GVT' #1:150@ca
setattr a radius 1.153 #0:151@ca
setattr a color #FFFF7B720000 #0:151@ca
setattr a name '|| .DSNT' #0:151@ca
setattr a radius 1.150 #1:151@ca
setattr a color #4FFC4B624B62 #1:151@ca
setattr a name '|| .AMLF' #1:151@ca
setattr a radius 1.491 #0:152@ca
setattr a color #FFFFFFFF0000 #0:152@ca
setattr a name '|| Q.' #0:152@ca
setattr a radius 1.489 #1:152@ca
setattr a color #E4F0067F067F #1:152@ca
setattr a name '|| E.' #1:152@ca
setattr a radius 1.807 #0:153@ca
setattr a color #355C355C2FDD #0:153@ca
setattr a name '|| .FV' #0:153@ca
setattr a radius 1.806 #1:153@ca
setattr a color #000D000D0000 #1:153@ca
setattr a name '|| I.L' #1:153@ca
setattr a radius 0.615 #0:154@ca
setattr a color #327232723272 #0:154@ca
setattr a name '|| .VF' #0:154@ca
setattr a radius 0.612 #1:154@ca
setattr a color #03F503F50151 #1:154@ca
setattr a name '|| .I' #1:154@ca
setattr a radius 0.880 #0:155@ca
setattr a color #FFFFFFFF5F20 #0:155@ca
setattr a name '|| .GAS' #0:155@ca
setattr a radius 0.881 #1:155@ca
setattr a color #016001600000 #1:155@ca
setattr a name '|| .LM' #1:155@ca
setattr a radius 0.553 #0:156@ca
setattr a color #92439243208A #0:156@ca
setattr a name '|| .SM' #0:156@ca
setattr a radius 0.555 #1:156@ca
setattr a color #BB49BB490000 #1:156@ca
setattr a name '|| .T' #1:156@ca
setattr a radius 1.589 #0:157@ca
setattr a color #B359B359631C #0:157@ca
setattr a name '|| .AGMS' #0:157@ca
setattr a radius 1.593 #1:157@ca
setattr a color #004B004B0000 #1:157@ca
setattr a name '|| F.L' #1:157@ca
setattr a radius 0.671 #0:158@ca
setattr a color #02DD02DD02B4 #0:158@ca
setattr a name '|| .IMF' #0:158@ca
setattr a radius 0.675 #1:158@ca
setattr a color #F185F18505B9 #1:158@ca
setattr a name '|| .QG' #1:158@ca
setattr a radius 0.166 #0:159@ca
setattr a color #AE5EAE5EAE5E #0:159@ca
setattr a name '|| .V' #0:159@ca
setattr a radius 0.166 #1:159@ca
setattr a color #000000000000 #1:159@ca
setattr a name '|| .M' #1:159@ca
setattr a radius 2.102 #0:160@ca
setattr a color #0AE90AE90000 #0:160@ca
setattr a name '|| .IMLQF' #0:160@ca
setattr a radius 2.110 #1:160@ca
setattr a color #F0F8F0F8F104 #1:160@ca
setattr a name '|| V.C' #1:160@ca
setattr a radius 1.524 #0:161@ca
setattr a color #5E1A5E1A362E #0:161@ca
setattr a name '|| .LGAVF' #0:161@ca
setattr a radius 1.528 #1:161@ca
setattr a color #F3B7F3B70000 #1:161@ca
setattr a name '|| Y.' #1:161@ca
setattr a radius 1.555 #0:162@ca
setattr a color #52EC52EC0BB6 #0:162@ca
setattr a name '|| .GLI' #0:162@ca
setattr a radius 1.563 #1:162@ca
setattr a color #F150F1500000 #1:162@ca
setattr a name '|| T.' #1:162@ca
setattr a radius 1.901 #0:163@ca
setattr a color #035D035D035D #0:163@ca
setattr a name '|| I.LA' #0:163@ca
setattr a radius 1.900 #1:163@ca
setattr a color #F751F751EED8 #1:163@ca
setattr a name '|| V.' #1:163@ca
setattr a radius 1.295 #0:164@ca
setattr a color #07A807A802C5 #0:164@ca
setattr a name '|| .LMS' #0:164@ca
setattr a radius 1.296 #1:164@ca
setattr a color #36373EC40000 #1:164@ca
setattr a name '|| .FY' #1:164@ca
setattr a radius 1.422 #0:165@ca
setattr a color #E2A5D46ACA97 #0:165@ca
setattr a name '|| .AV' #0:165@ca
setattr a radius 1.421 #1:165@ca
setattr a color #C3C3C3C307DF #1:165@ca
setattr a name '|| .SGM' #1:165@ca
setattr a radius 1.771 #0:166@ca
setattr a color #004D004D0000 #0:166@ca
setattr a name '|| .ILF' #0:166@ca
setattr a radius 1.769 #1:166@ca
setattr a color #FFB4FFB46329 #1:166@ca
setattr a name '|| .TAS' #1:166@ca
setattr a radius 1.649 #0:167@ca
setattr a color #A24DA24D2C9F #0:167@ca
setattr a name '|| .GVFS' #0:167@ca
setattr a radius 1.654 #1:167@ca
setattr a color #F9ABF9AB6DE8 #1:167@ca
setattr a name '|| .TAC' #1:167@ca
setattr a radius 0.766 #0:168@ca
setattr a color #0CB53F510B25 #0:168@ca
setattr a name '|| .LNK' #0:168@ca
setattr a radius 0.768 #1:168@ca
setattr a color #B3447DEE773B #1:168@ca
setattr a name '|| .VDI' #1:168@ca
setattr a radius 0.623 #0:169@ca
setattr a color #69C873DE0000 #0:169@ca
setattr a name '|| .NF' #0:169@ca
setattr a radius 0.627 #1:169@ca
setattr a color #E0C48D5AB183 #1:169@ca
setattr a name '|| .AE' #1:169@ca
setattr a radius 1.303 #0:170@ca
setattr a color #BAC897340000 #0:170@ca
setattr a name '|| .GYNDQ' #0:170@ca
setattr a radius 1.305 #1:170@ca
setattr a color #55DB55DBF063 #1:170@ca
setattr a name '|| .KPR' #1:170@ca
setattr a radius 2.176 #0:171@ca
setattr a color #76EE75450000 #0:171@ca
setattr a name '|| N.F' #0:171@ca
setattr a radius 2.177 #1:171@ca
setattr a color #0E3C15A5F8D1 #1:171@ca
setattr a name '|| R.K' #1:171@ca
setattr a radius 0.626 #0:172@ca
setattr a color #AA949E3A2663 #0:172@ca
setattr a name '|| .GLA' #0:172@ca
setattr a radius 0.625 #1:172@ca
setattr a color #1A598EC1401A #1:172@ca
setattr a name '|| .SK' #1:172@ca
setattr a radius 1.216 #0:173@ca
setattr a color #6B816B816B81 #0:173@ca
setattr a name '|| .VLAI' #0:173@ca
setattr a radius 1.221 #1:173@ca
setattr a color #F10CF10C0000 #1:173@ca
setattr a name '|| S.' #1:173@ca
setattr a radius 1.241 #0:174@ca
setattr a color #C073C073B781 #0:174@ca
setattr a name '|| .PL' #0:174@ca
setattr a radius 1.235 #1:174@ca
setattr a color #49045EA27441 #1:174@ca
setattr a name '|| .HDS' #1:174@ca
setattr a radius 1.267 #0:175@ca
setattr a color #160161409C31 #0:175@ca
setattr a name '|| .RSDQ' #0:175@ca
setattr a radius 1.267 #1:175@ca
setattr a color #46EA486F40D5 #1:175@ca
setattr a name '|| .VLI' #1:175@ca
setattr a radius 0.946 #0:176@ca
setattr a color #A851D4280000 #0:176@ca
setattr a name '|| .GN' #0:176@ca
setattr a radius 0.941 #1:176@ca
setattr a color #EDAAE160C7C1 #1:176@ca
setattr a name '|| .P' #1:176@ca
setattr a radius 1.186 #0:177@ca
setattr a color #B717B7177935 #0:177@ca
setattr a name '|| .PSM' #0:177@ca
setattr a radius 1.182 #1:177@ca
setattr a color #7F767B6E1400 #1:177@ca
setattr a name '|| .GITV' #1:177@ca
setattr a radius 0.999 #0:178@ca
setattr a color #45A045A02C3C #0:178@ca
setattr a name '|| .LVTP' #0:178@ca
setattr a radius 0.996 #1:178@ca
setattr a color #5AD85E02032A #1:178@ca
setattr a name '|| .SI' #1:178@ca
setattr a radius 1.799 #0:179@ca
setattr a color #B00162170000 #0:179@ca
setattr a name '|| .ENGH' #0:179@ca
setattr a radius 1.804 #1:179@ca
setattr a color #F0B5F0B5FD4A #1:179@ca
setattr a name '|| A.PV' #1:179@ca
setattr a radius 0.831 #0:180@ca
setattr a color #769ABB4C5BCD #0:180@ca
setattr a name '|| .ATH' #0:180@ca
setattr a radius 0.832 #1:180@ca
setattr a color #F04DF04DF672 #1:180@ca
setattr a name '|| .P' #1:180@ca
setattr a radius 1.030 #0:181@ca
setattr a color #10D110D10AA3 #0:181@ca
setattr a name '|| .FV' #0:181@ca
setattr a radius 1.027 #1:181@ca
setattr a color #01E301E30142 #1:181@ca
setattr a name '|| .LI' #1:181@ca
setattr a radius 1.770 #0:182@ca
setattr a color #2B8B2B8B14B2 #0:182@ca
setattr a name '|| .LIVT' #0:182@ca
setattr a radius 1.776 #1:182@ca
setattr a color #F0DDF0DDDF74 #1:182@ca
setattr a name '|| .APS' #1:182@ca
setattr a radius 2.168 #0:183@ca
setattr a color #E3C5E3C5C2CD #0:183@ca
setattr a name '|| V.TPLG' #0:183@ca
setattr a radius 2.173 #1:183@ca
setattr a color #000000000000 #1:183@ca
setattr a name '|| I:.' #1:183@ca
setattr a radius 0.445 #0:184@ca
setattr a color #000000000000 #0:184@ca
setattr a name '|| .ML' #0:184@ca
setattr a radius 0.447 #1:184@ca
setattr a color #F125F1250000 #1:184@ca
setattr a name '|| .G' #1:184@ca
setattr a radius 1.218 #0:185@ca
setattr a color #1F511F511342 #0:185@ca
setattr a name '|| .LAIT' #0:185@ca
setattr a radius 1.222 #1:185@ca
setattr a color #002000200000 #1:185@ca
setattr a name '|| F.' #1:185@ca
setattr a radius 2.142 #0:186@ca
setattr a color #126212621262 #0:186@ca
setattr a name '|| L.V' #0:186@ca
setattr a radius 2.142 #1:186@ca
setattr a color #A2F6A2F650AB #1:186@ca
setattr a name '|| .ASIT' #1:186@ca
setattr a radius 1.658 #0:187@ca
setattr a color #000000000000 #0:187@ca
setattr a name '|| I.F' #0:187@ca
setattr a radius 1.657 #1:187@ca
setattr a color #858A858A858A #1:187@ca
setattr a name '|| V.L' #1:187@ca
setattr a radius 1.387 #0:188@ca
setattr a color #44CB44CB39CA #0:188@ca
setattr a name '|| .VLIC' #0:188@ca
setattr a radius 1.391 #1:188@ca
setattr a color #357D357D0000 #1:188@ca
setattr a name '|| .FGT' #1:188@ca
setattr a radius 1.473 #0:189@ca
setattr a color #FFFFFFFF9896 #0:189@ca
setattr a name '|| .VGS' #0:189@ca
setattr a radius 1.471 #1:189@ca
setattr a color #079C079C0712 #1:189@ca
setattr a name '|| .LMA' #1:189@ca
setattr a radius 2.088 #0:190@ca
setattr a color #000000000000 #0:190@ca
setattr a name '|| I.L' #0:190@ca
setattr a radius 2.082 #1:190@ca
setattr a color #FAA1F3956552 #1:190@ca
setattr a name '|| .VGN' #1:190@ca
setattr a radius 2.755 #0:191@ca
setattr a color #FFFFFFFF170B #0:191@ca
setattr a name '|| G:.SV' #0:191@ca
setattr a radius 2.756 #1:191@ca
setattr a color #0000000098CE #1:191@ca
setattr a name '|| H:.I' #1:191@ca
setattr a radius 2.315 #0:192@ca
setattr a color #9274927425A3 #0:192@ca
setattr a name '|| .TVISGCM' #0:192@ca
setattr a radius 2.324 #1:192@ca
setattr a color #000000000000 #1:192@ca
setattr a name '|| L:.' #1:192@ca
setattr a radius 1.827 #0:193@ca
setattr a color #FFFFFFFF0000 #0:193@ca
setattr a name '|| S.' #0:193@ca
setattr a radius 1.824 #1:193@ca
setattr a color #5CB75CB75AC6 #1:193@ca
setattr a name '|| .ALIFV' #1:193@ca
setattr a radius 1.935 #0:194@ca
setattr a color #000000000000 #0:194@ca
setattr a name '|| .LMF' #0:194@ca
setattr a radius 1.935 #1:194@ca
setattr a color #FFFFFDBF4A38 #1:194@ca
setattr a name '|| .TGAS' #1:194@ca
setattr a radius 2.560 #0:195@ca
setattr a color #FFFFFFFF0000 #0:195@ca
setattr a name '|| G:.N' #0:195@ca
setattr a radius 2.558 #1:195@ca
setattr a color #015D015D0074 #1:195@ca
setattr a name '|| I:.T' #1:195@ca
setattr a radius 1.343 #0:196@ca
setattr a color #4DEE4DEE0000 #0:196@ca
setattr a name '|| .GMFL' #0:196@ca
setattr a radius 1.337 #1:196@ca
setattr a color #DA41CD7B9FED #1:196@ca
setattr a name '|| .PADNY' #1:196@ca
setattr a radius 1.774 #0:197@ca
setattr a color #987498742AA6 #0:197@ca
setattr a name '|| .NLTPAM' #0:197@ca
setattr a radius 1.777 #1:197@ca
setattr a color #149F149F0473 #1:197@ca
setattr a name '|| .IFYV' #1:197@ca
setattr a radius 0.725 #0:198@ca
setattr a color #FFFFFFFF4F18 #0:198@ca
setattr a name '|| .NSVA' #0:198@ca
setattr a radius 0.725 #1:198@ca
setattr a color #FFFFD0AE0000 #1:198@ca
setattr a name '|| .TD' #1:198@ca
setattr a radius 0.112 #0:199@ca
setattr a color #FFFFFFFF0000 #0:199@ca
setattr a name '|| .' #0:199@ca
setattr a radius 0.113 #1:199@ca
setattr a color #F248F2480000 #1:199@ca
setattr a name '|| .N' #1:199@ca
setattr a radius 1.791 #0:200@ca
setattr a color #917391730D77 #0:200@ca
setattr a name '|| Y.FP' #0:200@ca
setattr a radius 1.793 #1:200@ca
setattr a color #F048F04842E4 #1:200@ca
setattr a name '|| .TAGC' #1:200@ca
setattr a radius 2.823 #0:201@ca
setattr a color #FFFFFFFFF7EC #0:201@ca
setattr a name '|| A:.' #0:201@ca
setattr a radius 2.817 #1:201@ca
setattr a color #F679F6790000 #1:201@ca
setattr a name '|| SG.' #1:201@ca
setattr a radius 0.864 #0:202@ca
setattr a color #018A018A018A #0:202@ca
setattr a name '|| .LI' #0:202@ca
setattr a radius 0.866 #1:202@ca
setattr a color #000000000000 #1:202@ca
setattr a name '|| .M' #1:202@ca
setattr a radius 0.045 #0:203@ca
setattr a color #FFFFFFFF0000 #0:203@ca
setattr a name '|| .' #0:203@ca
setattr a radius 0.045 #1:203@ca
setattr a color #FFFFFFFF0000 #1:203@ca
setattr a name '|| .' #1:203@ca
setattr a radius 0.045 #0:204@ca
setattr a color #FFFFFFFFFFFF #0:204@ca
setattr a name '|| .' #0:204@ca
setattr a radius 0.045 #1:204@ca
setattr a color #FFFF00000000 #1:204@ca
setattr a name '|| .' #1:204@ca
setattr a radius 0.531 #0:205@ca
setattr a color #FFFFFFFF3035 #0:205@ca
setattr a name '|| .S' #0:205@ca
setattr a radius 0.531 #1:205@ca
setattr a color #FFFFFFFFEBBF #1:205@ca
setattr a name '|| .A' #1:205@ca
setattr a radius 0.463 #0:206@ca
setattr a color #00000000FFFF #0:206@ca
setattr a name '|| .R' #0:206@ca
setattr a radius 0.461 #1:206@ca
setattr a color #FEE3FEE3E838 #1:206@ca
setattr a name '|| .V' #1:206@ca
setattr a radius 3.759 #0:207@ca
setattr a color #FFFF00000000 #0:207@ca
setattr a name '|| D>:.' #0:207@ca
setattr a radius 3.754 #1:207@ca
setattr a color #FFA7FFA71AF7 #1:207@ca
setattr a name '|| S>:.A' #1:207@ca
setattr a radius 1.042 #0:208@ca
setattr a color #000000000000 #0:208@ca
setattr a name '|| L.' #0:208@ca
setattr a radius 1.039 #1:208@ca
setattr a color #061206120000 #1:208@ca
setattr a name '|| .FTI' #1:208@ca
setattr a radius 0.778 #0:209@ca
setattr a color #FFFFFFFF7FFF #0:209@ca
setattr a name '|| .SP' #0:209@ca
setattr a radius 0.779 #1:209@ca
setattr a color #F053F37E490C #1:209@ca
setattr a name '|| .GA' #1:209@ca
setattr a radius 1.693 #0:210@ca
setattr a color #FFFFFFFFF590 #0:210@ca
setattr a name '|| P.' #0:210@ca
setattr a radius 1.690 #1:210@ca
setattr a color #F6DBF6DBCA1F #1:210@ca
setattr a name '|| A.SV' #1:210@ca
setattr a radius 3.698 #0:211@ca
setattr a color #00000000FFFF #0:211@ca
setattr a name '|| R>:.K' #0:211@ca
setattr a radius 3.694 #1:211@ca
setattr a color #FB4FFB4FF031 #1:211@ca
setattr a name '|| A>:.' #1:211@ca
setattr a radius 1.833 #0:212@ca
setattr a color #000000000000 #0:212@ca
setattr a name '|| L.IF' #0:212@ca
setattr a radius 1.833 #1:212@ca
setattr a color #F83CF83CF83C #1:212@ca
setattr a name '|| V.' #1:212@ca
setattr a radius 1.363 #0:213@ca
setattr a color #01B201B201B2 #0:213@ca
setattr a name '|| .FMLA' #0:213@ca
setattr a radius 1.361 #1:213@ca
setattr a color #0FED0FED0D1C #1:213@ca
setattr a name '|| I.V' #1:213@ca
setattr a radius 0.857 #0:214@ca
setattr a color #07CC69D412B8 #0:214@ca
setattr a name '|| .HTL' #0:214@ca
setattr a radius 0.852 #1:214@ca
setattr a color #66F5647C04F2 #1:214@ca
setattr a name '|| .SYFM' #1:214@ca
setattr a radius 1.303 #0:215@ca
setattr a color #80BB80DF45E0 #0:215@ca
setattr a name '|| .ASF' #0:215@ca
setattr a radius 1.303 #1:215@ca
setattr a color #B7E2D4300000 #1:215@ca
setattr a name '|| .NGR' #1:215@ca
setattr a radius 0.782 #0:216@ca
setattr a color #000000000000 #0:216@ca
setattr a name '|| .LFI' #0:216@ca
setattr a radius 0.775 #1:216@ca
setattr a color #3FD578780787 #1:216@ca
setattr a name '|| .KN' #1:216@ca
setattr a radius 1.073 #0:217@ca
setattr a color #668366835BA4 #0:217@ca
setattr a name '|| .AL' #0:217@ca
setattr a radius 1.065 #1:217@ca
setattr a color #2B9916D500FA #1:217@ca
setattr a name '|| .FWD' #1:217@ca
setattr a radius 0.926 #0:218@ca
setattr a color #FFFFFFFF5C0F #0:218@ca
setattr a name '|| .GP' #0:218@ca
setattr a radius 0.913 #1:218@ca
setattr a color #357C6AF90C87 #1:218@ca
setattr a name '|| .KTSH' #1:218@ca
setattr a radius 0.424 #0:219@ca
setattr a color #0000054D0584 #0:219@ca
setattr a name '|| .W' #0:219@ca
setattr a radius 0.417 #1:219@ca
setattr a color #BF4184C574F7 #1:219@ca
setattr a name '|| .PEA' #1:219@ca
setattr a radius 0.659 #0:220@ca
setattr a color #CFD3E7E90000 #0:220@ca
setattr a name '|| .G' #0:220@ca
setattr a radius 0.661 #1:220@ca
setattr a color #00000000F0C5 #1:220@ca
setattr a name '|| .R' #1:220@ca
setattr a radius 0.256 #0:221@ca
setattr a color #6E39664B1960 #0:221@ca
setattr a name '|| .' #0:221@ca
setattr a radius 0.266 #1:221@ca
setattr a color #F9FEF9FE0000 #1:221@ca
setattr a name '|| .T' #1:221@ca
setattr a radius 0.265 #0:222@ca
setattr a color #679A4E714E71 #0:222@ca
setattr a name '|| .' #0:222@ca
setattr a radius 0.274 #1:222@ca
setattr a color #F8B0F8B00000 #1:222@ca
setattr a name '|| .G' #1:222@ca
setattr a radius 0.440 #0:223@ca
setattr a color #B600B6009C7E #0:223@ca
setattr a name '|| .VA' #0:223@ca
setattr a radius 0.001 #1:223@ca
setattr a color #000000000000 #1:223@ca
setattr a name '|| .' #1:223@ca
setattr a radius 0.612 #0:224@ca
setattr a color #024B03550000 #0:224@ca
setattr a name '|| .FLI' #0:224@ca
setattr a radius 0.001 #1:224@ca
setattr a color #000000000000 #1:224@ca
setattr a name '|| .' #1:224@ca
setattr a radius 0.371 #0:225@ca
setattr a color #AAFDB14632EB #0:225@ca
setattr a name '|| .T' #0:225@ca
setattr a radius 0.001 #1:225@ca
setattr a color #000000000000 #1:225@ca
setattr a name '|| .' #1:225@ca
setattr a radius 0.432 #0:226@ca
setattr a color #B58ECBD53459 #0:226@ca
setattr a name '|| .G' #0:226@ca
setattr a radius 0.001 #1:226@ca
setattr a color #000000000000 #1:226@ca
setattr a name '|| .' #1:226@ca
setattr a radius 0.768 #0:227@ca
setattr a color #BBFFDDFF0855 #0:227@ca
setattr a name '|| .G' #0:227@ca
setattr a radius 0.769 #1:227@ca
setattr a color #000000000000 #1:227@ca
setattr a name '|| .F' #1:227@ca
setattr a radius 0.339 #0:228@ca
setattr a color #6FD566640000 #0:228@ca
setattr a name '|| .' #0:228@ca
setattr a radius 0.346 #1:228@ca
setattr a color #F548F5480000 #1:228@ca
setattr a name '|| .N' #1:228@ca
setattr a radius 0.717 #0:229@ca
setattr a color #B4B48FD70000 #0:229@ca
setattr a name '|| .NS' #0:229@ca
setattr a radius 0.718 #1:229@ca
setattr a color #D895D895B524 #1:229@ca
setattr a name '|| .AG' #1:229@ca
setattr a radius 0.928 #0:230@ca
setattr a color #012E012E0000 #0:230@ca
setattr a name '|| .WI' #0:230@ca
setattr a radius 0.925 #1:230@ca
setattr a color #B39535CA0000 #1:230@ca
setattr a name '|| .DNH' #1:230@ca
setattr a radius 0.242 #0:231@ca
setattr a color #90CB90CB4799 #0:231@ca
setattr a name '|| .' #0:231@ca
setattr a radius 0.246 #1:231@ca
setattr a color #636F723F0000 #1:231@ca
setattr a name '|| .' #1:231@ca
setattr a radius 1.508 #0:232@ca
setattr a color #592B592B0000 #0:232@ca
setattr a name '|| .YWFI' #0:232@ca
setattr a radius 1.501 #1:232@ca
setattr a color #EA1871E80778 #1:232@ca
setattr a name '|| .DNQG' #1:232@ca
setattr a radius 2.261 #0:233@ca
setattr a color #45F645F6317B #0:233@ca
setattr a name '|| .WAFSV' #0:233@ca
setattr a radius 2.265 #1:233@ca
setattr a color #00002EAFB6F6 #1:233@ca
setattr a name '|| H.QLN' #1:233@ca
setattr a radius 1.320 #0:234@ca
setattr a color #000000000000 #0:234@ca
setattr a name '|| .ILF' #0:234@ca
setattr a radius 1.320 #1:234@ca
setattr a color #020202020157 #1:234@ca
setattr a name '|| W.' #1:234@ca
setattr a radius 0.985 #0:235@ca
setattr a color #FFFFFFFFE472 #0:235@ca
setattr a name '|| .VG' #0:235@ca
setattr a radius 0.982 #1:235@ca
setattr a color #000000000000 #1:235@ca
setattr a name '|| .IL' #1:235@ca
setattr a radius 3.345 #0:236@ca
setattr a color #FFFFFFFFFFFF #0:236@ca
setattr a name '|| P>:.' #0:236@ca
setattr a radius 3.345 #1:236@ca
setattr a color #2E332E330000 #1:236@ca
setattr a name '|| FY.' #1:236@ca
setattr a radius 3.454 #0:237@ca
setattr a color #3E863E863E86 #0:237@ca
setattr a name '|| VI.M' #0:237@ca
setattr a radius 3.452 #1:237@ca
setattr a color #001500150000 #1:237@ca
setattr a name '|| W>:.' #1:237@ca
setattr a radius 0.625 #0:238@ca
setattr a color #06D106D10000 #0:238@ca
setattr a name '|| .FIMG' #0:238@ca
setattr a radius 0.626 #1:238@ca
setattr a color #F846F846F846 #1:238@ca
setattr a name '|| .VA' #1:238@ca
setattr a radius 2.025 #0:239@ca
setattr a color #FFFFFFFFA51A #0:239@ca
setattr a name '|| A.S' #0:239@ca
setattr a radius 2.024 #1:239@ca
setattr a color #F71AF71A08B2 #1:239@ca
setattr a name '|| G.' #1:239@ca
setattr a radius 0.103 #0:240@ca
setattr a color #FFFFFFFF0000 #0:240@ca
setattr a name '|| .T' #0:240@ca
setattr a radius 0.101 #1:240@ca
setattr a color #FB0690279027 #1:240@ca
setattr a name '|| .' #1:240@ca
setattr a radius 1.079 #0:241@ca
setattr a color #063406340634 #0:241@ca
setattr a name '|| .IV' #0:241@ca
setattr a radius 1.080 #1:241@ca
setattr a color #013001300000 #1:241@ca
setattr a name '|| F.' #1:241@ca
setattr a radius 1.389 #0:242@ca
setattr a color #0B3C0B3C0498 #0:242@ca
setattr a name '|| L.SA' #0:242@ca
setattr a radius 1.389 #1:242@ca
setattr a color #2D902D902B9F #1:242@ca
setattr a name '|| .IV' #1:242@ca
setattr a radius 0.590 #0:243@ca
setattr a color #FFFFFFFFFFFF #0:243@ca
setattr a name '|| .A' #0:243@ca
setattr a radius 0.590 #1:243@ca
setattr a color #FFFFFFFF0000 #1:243@ca
setattr a name '|| .G' #1:243@ca
setattr a radius 0.563 #0:244@ca
setattr a color #FFFFFFFF0000 #0:244@ca
setattr a name '|| .S' #0:244@ca
setattr a radius 0.561 #1:244@ca
setattr a color #FF16FF165CC2 #1:244@ca
setattr a name '|| .GA' #1:244@ca
setattr a radius 1.066 #0:245@ca
setattr a color #592159211718 #0:245@ca
setattr a name '|| .IYVLC' #0:245@ca
setattr a radius 1.074 #1:245@ca
setattr a color #F152F1529C50 #1:245@ca
setattr a name '|| .AG' #1:245@ca
setattr a radius 0.719 #0:246@ca
setattr a color #931693167451 #0:246@ca
setattr a name '|| .VAF' #0:246@ca
setattr a radius 0.722 #1:246@ca
setattr a color #003F003F0000 #1:246@ca
setattr a name '|| .LI' #1:246@ca
setattr a radius 2.666 #0:247@ca
setattr a color #C231C2310000 #0:247@ca
setattr a name '|| G:.LI' #0:247@ca
setattr a radius 2.667 #1:247@ca
setattr a color #F017F017F017 #1:247@ca
setattr a name '|| A:.V' #1:247@ca
setattr a radius 0.690 #0:248@ca
setattr a color #FFFFFFFFF776 #0:248@ca
setattr a name '|| .V' #0:248@ca
setattr a radius 0.687 #1:248@ca
setattr a color #B671B6719B1F #1:248@ca
setattr a name '|| .AW' #1:248@ca
setattr a radius 0.549 #0:249@ca
setattr a color #5EC05EC01546 #0:249@ca
setattr a name '|| .WT' #0:249@ca
setattr a radius 0.550 #1:249@ca
setattr a color #0AF00AF001D2 #1:249@ca
setattr a name '|| .L' #1:249@ca
setattr a radius 0.949 #0:250@ca
setattr a color #C2E1C2E19C47 #0:250@ca
setattr a name '|| .VAF' #0:250@ca
setattr a radius 0.956 #1:250@ca
setattr a color #B212B2120000 #1:250@ca
setattr a name '|| .YL' #1:250@ca
setattr a radius 1.218 #0:251@ca
setattr a color #6760676014D5 #0:251@ca
setattr a name '|| .YF' #0:251@ca
setattr a radius 1.218 #1:251@ca
setattr a color #00000988AB68 #1:251@ca
setattr a name '|| .HW' #1:251@ca
setattr a radius 0.459 #0:252@ca
setattr a color #1F80486B54BB #0:252@ca
setattr a name '|| .KL' #0:252@ca
setattr a radius 0.457 #1:252@ca
setattr a color #D63877200AF8 #1:252@ca
setattr a name '|| .QE' #1:252@ca
setattr a radius 0.636 #0:253@ca
setattr a color #19F630091EDE #0:253@ca
setattr a name '|| .LK' #0:253@ca
setattr a radius 0.633 #1:253@ca
setattr a color #6636604B0000 #1:253@ca
setattr a name '|| .YF' #1:253@ca
setattr a radius 1.376 #0:254@ca
setattr a color #16D616D60C0F #0:254@ca
setattr a name '|| .FMATL' #0:254@ca
setattr a radius 1.380 #1:254@ca
setattr a color #88D388D38B91 #1:254@ca
setattr a name '|| .VI' #1:254@ca
setattr a radius 1.213 #0:255@ca
setattr a color #35C935C91D68 #0:255@ca
setattr a name '|| .IVY' #0:255@ca
setattr a radius 1.211 #1:255@ca
setattr a color #028E030E0080 #1:255@ca
setattr a name '|| .LF' #1:255@ca
setattr a radius 0.879 #0:256@ca
setattr a color #54FA84E50E07 #0:256@ca
setattr a name '|| .GK' #0:256@ca
setattr a radius 0.877 #1:256@ca
setattr a color #046A15B84D86 #1:256@ca
setattr a name '|| .RI' #1:256@ca
setattr a radius 0.810 #0:257@ca
setattr a color #1A9C584C0000 #0:257@ca
setattr a name '|| .CRGM' #0:257@ca
setattr a radius 0.811 #1:257@ca
setattr a color #E4A4C2EAB745 #1:257@ca
setattr a name '|| .AP' #1:257@ca
setattr a radius 0.658 #0:258@ca
setattr a color #18E657ECD0EB #0:258@ca
setattr a name '|| .H' #0:258@ca
setattr a radius 0.647 #1:258@ca
setattr a color #4F04452D13FC #1:258@ca
setattr a name '|| .IRD' #1:258@ca
setattr a radius 0.299 #0:259@ca
setattr a color #036F083F0EE3 #0:259@ca
setattr a name '|| .LI' #0:259@ca
setattr a radius 0.304 #1:259@ca
setattr a color #F4465F09283D #1:259@ca
setattr a name '|| .E' #1:259@ca
setattr a radius 0.308 #0:260@ca
setattr a color #E8B9B6DA7B02 #0:260@ca
setattr a name '|| .P' #0:260@ca
setattr a radius 0.314 #1:260@ca
setattr a color #0000884F0C77 #1:260@ca
setattr a name '|| .RN' #1:260@ca
setattr a radius 0.171 #0:261@ca
setattr a color #BAF1C65B16D5 #0:261@ca
setattr a name '|| .' #0:261@ca
setattr a radius 0.184 #1:261@ca
setattr a color #F37500000000 #1:261@ca
setattr a name '|| .E' #1:261@ca
setattr a radius 0.183 #0:262@ca
setattr a color #816881685259 #0:262@ca
setattr a name '|| .' #0:262@ca
setattr a radius 0.191 #1:262@ca
setattr a color #FB2FFB2FFB2F #1:262@ca
setattr a name '|| .P' #1:262@ca
setattr a radius 0.175 #0:263@ca
setattr a color #99794A072503 #0:263@ca
setattr a name '|| .' #0:263@ca
setattr a radius 0.183 #1:263@ca
setattr a color #FBB2FBB2FBB2 #1:263@ca
setattr a name '|| .P' #1:263@ca
setattr a radius 0.257 #0:264@ca
setattr a color #B329B329527F #0:264@ca
setattr a name '|| .' #0:264@ca
setattr a radius 0.270 #1:264@ca
setattr a color #00000000FCF2 #1:264@ca
setattr a name '|| .H' #1:264@ca
setattr a radius 0.140 #0:265@ca
setattr a color #9C298771263E #0:265@ca
setattr a name '|| .' #0:265@ca
setattr a radius 0.150 #1:265@ca
setattr a color #00000000F248 #1:265@ca
setattr a name '|| .H' #1:265@ca
setattr a radius 0.114 #0:266@ca
setattr a color #3E674C965C30 #0:266@ca
setattr a name '|| .' #0:266@ca
setattr a radius 0.120 #1:266@ca
setattr a color #74C5C11E023E #1:266@ca
setattr a name '|| .' #1:266@ca
setattr a radius 0.694 #0:267@ca
setattr a color #7FC616752E21 #0:267@ca
setattr a name '|| .EDK' #0:267@ca
setattr a radius 0.691 #1:267@ca
setattr a color #F340F3408FC9 #1:267@ca
setattr a name '|| .PAT' #1:267@ca
setattr a radius 0.559 #0:268@ca
setattr a color #FFFFB29A7B30 #0:268@ca
setattr a name '|| .AE' #0:268@ca
setattr a radius 0.560 #1:268@ca
setattr a color #000023FC3B8A #1:268@ca
setattr a name '|| .FH' #1:268@ca
setattr a radius 0.707 #0:269@ca
setattr a color #8536492225DF #0:269@ca
setattr a name '|| .EGA' #0:269@ca
setattr a radius 0.703 #1:269@ca
setattr a color #25857EE278BF #1:269@ca
setattr a name '|| .KTD' #1:269@ca
setattr a radius 0.618 #0:270@ca
setattr a color #BAEDAA09396E #0:270@ca
setattr a name '|| .TV' #0:270@ca
setattr a radius 0.613 #1:270@ca
setattr a color #AB69B4BC4B4E #1:270@ca
setattr a name '|| .SA' #1:270@ca
setattr a radius 0.334 #0:271@ca
setattr a color #FFFF855C4FB6 #0:271@ca
setattr a name '|| .EP' #0:271@ca
setattr a radius 0.336 #1:271@ca
setattr a color #00B90173E5E0 #1:271@ca
setattr a name '|| .K' #1:271@ca
setattr a radius 0.230 #0:272@ca
setattr a color #FFFFD379011C #0:272@ca
setattr a name '|| .N' #0:272@ca
setattr a radius 0.228 #1:272@ca
setattr a color #D13FE7828A03 #1:272@ca
setattr a name '|| .P' #1:272@ca
setattr a radius 0.191 #0:273@ca
setattr a color #C60EA7676FFF #0:273@ca
setattr a name '|| .' #0:273@ca
setattr a radius 0.194 #1:273@ca
setattr a color #9CD39CD30000 #1:273@ca
setattr a name '|| .' #1:273@ca
setattr a radius 0.316 #0:274@ca
setattr a color #613689B83A54 #0:274@ca
setattr a name '|| .' #0:274@ca
setattr a radius 0.322 #1:274@ca
setattr a color #2BF82BF82BF8 #1:274@ca
setattr a name '|| .LV' #1:274@ca
setattr a radius 0.255 #0:275@ca
setattr a color #1DCE28090000 #0:275@ca
setattr a name '|| .L' #0:275@ca
setattr a radius 0.258 #1:275@ca
setattr a color #F302E9E9E9E9 #1:275@ca
setattr a name '|| .PV' #1:275@ca
setattr a radius 0.182 #0:276@ca
setattr a color #A4BD906E67CF #0:276@ca
setattr a name '|| .' #0:276@ca
setattr a radius 0.184 #1:276@ca
setattr a color #016802D0F167 #1:276@ca
setattr a name '|| .' #1:276@ca
setattr a radius 0.281 #0:277@ca
setattr a color #AB458B625775 #0:277@ca
setattr a name '|| .E' #0:277@ca
setattr a radius 0.284 #1:277@ca
setattr a color #F2BAF2BAD271 #1:277@ca
setattr a name '|| .V' #1:277@ca
setattr a radius 0.255 #0:278@ca
setattr a color #1FE932E882B0 #0:278@ca
setattr a name '|| .K' #0:278@ca
setattr a radius 0.260 #1:278@ca
setattr a color #232123210000 #1:278@ca
setattr a name '|| .I' #1:278@ca
setattr a radius 0.205 #0:279@ca
setattr a color #519792A45FEC #0:279@ca
setattr a name '|| .' #0:279@ca
setattr a radius 0.209 #1:279@ca
setattr a color #F4FD00000000 #1:279@ca
setattr a name '|| .D' #1:279@ca
setattr a radius 0.204 #0:280@ca
setattr a color #4A097FFF5B9B #0:280@ca
setattr a name '|| .' #0:280@ca
setattr a radius 0.202 #1:280@ca
setattr a color #43B843B847BE #1:280@ca
setattr a name '|| .' #1:280@ca
setattr a radius 0.286 #0:281@ca
setattr a color #11C61D553638 #0:281@ca
setattr a name '|| .L' #0:281@ca
setattr a radius 0.289 #1:281@ca
setattr a color #B50400000000 #1:281@ca
setattr a name '|| .D' #1:281@ca
%</AQPsublogo>
%    \end{macrocode}
%    \begin{macrocode}
%<*Standard>
%%%%%%%%%%%%%%%%%%%%%%%%%%%%%%%%%%%%%%%%%%%%%%%%%%%%%%%%%%%%%%%%%%%%%%%%%%%%%%%%
%%%%%                                                                      %%%%%
%%%%%                   Standard genetic code definitions                  %%%%%
%%%%%                                                                      %%%%%
%%%%%      (The last codon of each list is used for backtranslations       %%%%%
%%%%%        from protein to DNA sequences---therefore the wobbles)        %%%%%
%%%%%                                                                      %%%%%
%%%%%              These definitions are default in TeXshade.              %%%%%
%%%%%     There is no need to load them. This is an example file only.     %%%%%
%%%%%                                                                      %%%%%
%%%%%%%%%%%%%%%%%%%%%%%%%%%%%%%%%%%%%%%%%%%%%%%%%%%%%%%%%%%%%%%%%%%%%%%%%%%%%%%%

\codon{A}{GCA,GCG,GCC,GCT,GCU,GCN}
\codon{C}{TGC,TGT,UGC,UGU,TGY}
\codon{D}{GAC,GAT,GAU,GAY}
\codon{E}{GAA,GAG,GAR}
\codon{F}{TTC,TTT,UUC,UUU,TTY}
\codon{G}{GGA,GGG,GGC,GGT,GGU,GGN}
\codon{H}{CAC,CAT,CAY}
\codon{I}{ATA,ATC,ATT,AUA,AUC,AUU,ATH}
\codon{K}{AAA,AAG,AAG,AAR}
\codon{L}{CTA,CTG,CTC,CTT,TTA,TTG,CUG,CUG,CUC,CUU,UUA,UUG,YTN}
\codon{M}{ATG,AUG,ATG}
\codon{N}{AAC,AAT,AAU,AAY}
\codon{P}{CCA,CCG,CCC,CCT,CCU,CCN}
\codon{Q}{CAA,CAG,CAR}
\codon{R}{AGA,AGG,CGA,CGG,CGC,CGT,CGU,MGN}
\codon{S}{TCT,TCC,TCG,TCA,AGT,AGC,UCU,UCC,UCG,UCA,AGU,WSN}
\codon{T}{ACT,ACC,ACG,ACA,ACU,ACN}
\codon{V}{GTA,GTG,GTC,GTT,GUA,GUG,GUC,GUU,GTN}
\codon{W}{TGG,UGG,TGG}
\codon{Y}{TAC,TAT,UAC,UAU,TAY}
\codon{.}{TAA,TAG,TGA,UAA,UAG,UGA,TRR}
%</Standard>
%    \end{macrocode}
%    \begin{macrocode}
%<*Ciliate>
%%%%%%%%%%%%%%%%%%%%%%%%%%%%%%%%%%%%%%%%%%%%%%%%%%%%%%%%%%%%%%%%%%%%%%%%%%%%%%%%
%%%%%                                                                      %%%%%
%%%%%            Ciliate macronuclear genetic code definitions             %%%%%
%%%%%                                                                      %%%%%
%%%%%    Only exchanges compared to the standard code must be defined.     %%%%%
%%%%%                                                                      %%%%%
%%%%%      (The last codon of the list is used for backtranslations        %%%%%
%%%%%       from protein to DNA sequences---therefore the wobbles)         %%%%%
%%%%%                                                                      %%%%%
%%%%%                                                                      %%%%%
%%%%%    Activate these definitions for your alignment by the following    %%%%%
%%%%%    command in the texshade environment:                              %%%%%
%%%%%                                                                      %%%%%
%%%%%                         \geneticcode{ciliate}                        %%%%%
%%%%%                                                                      %%%%%
%%%%%%%%%%%%%%%%%%%%%%%%%%%%%%%%%%%%%%%%%%%%%%%%%%%%%%%%%%%%%%%%%%%%%%%%%%%%%%%%

\codon{Q}{TAA,TAG,UAA,UAG,YAR}
%</Ciliate>
%    \end{macrocode}
%    \begin{macrocode}
%<*meme>
%!PS-Adobe-3.0 EPSF-3.0
%%Creator: GraphicConverter 7
%%Title: meme.eps
%CreationDate: 2021-05-13
%%Pages: 1
%%BoundingBox: 0 0 897 882
%%DocumentData: Clean7Bit
%%EndComments
%%BeginProlog
%%EndProlog
%%Page: 1 1
%ImageData: 897 882 8 3 0 2691 2 } exec
/languagelevel where {pop languagelevel 2 lt} {true} ifelse {
 (JPEG picture requires Postscript level 2
) dup print flush
 /Helvetica findfont 20 scalefont setfont 100 100 moveto show showpage stop
} if
save
/RawData currentfile /ASCIIHexDecode filter def
/Data RawData << >> /DCTDecode filter def
597.00 582.00 scale
/DeviceRGB setcolorspace
{ << /ImageType 1
     /Width 897
     /Height 882
     /ImageMatrix [ 897 0 0 -882 0 882 ]
     /DataSource Data
     /BitsPerComponent 8
     /Decode [0 1 0 1 0 1]
 >> image
 Data closefile
 RawData flushfile
 showpage
 restore
} exec
FFD8FFE000104A46494600010100004800480000FFE100404578696600004D
4D002A00000008000187690004000000010000001A000000000002A0020004
0000000100000381A0030004000000010000037200000000FFED003850686F
746F73686F7020332E30003842494D04040000000000003842494D04250000
00000010D41D8CD98F00B204E9800998ECF8427EFFE207E84943435F50524F
46494C45000101000007D86170706C022000006D6E74725247422058595A20
07D900020019000B001A000B616373704150504C000000006170706C000000
000000000000000000000000000000F6D6000100000000D32D6170706C0000
00000000000000000000000000000000000000000000000000000000000000
00000000000000000000000000000B64657363000001080000006F6473636D
000001780000059C637072740000071400000038777470740000074C000000
147258595A00000760000000146758595A00000774000000146258595A0000
078800000014725452430000079C0000000E63686164000007AC0000002C62
5452430000079C0000000E675452430000079C0000000E6465736300000000
0000001447656E65726963205247422050726F66696C650000000000000000
0000001447656E65726963205247422050726F66696C650000000000000000
00000000000000000000000000000000000000000000000000000000000000
00000000000000000000006D6C7563000000000000001F0000000C736B534B
00000028000001846461444B0000002E000001AC6361455300000024000001
DA7669564E00000024000001FE707442520000002600000222756B55410000
002A0000024866724655000000280000027268754855000000280000029A7A
68545700000016000002C26E624E4F00000026000002D86373435A00000022
000002FE6865494C0000001E0000032069744954000000280000033E726F52
4F0000002400000366646544450000002C0000038A6B6F4B52000000160000
03B67376534500000026000002D87A68434E00000016000003CC6A614A5000
00001A000003E2656C475200000022000003FC7074504F000000260000041E
6E6C4E4C000000280000044465734553000000260000041E74685448000000
240000046C7472545200000022000004906669464900000028000004B26872
485200000028000004DA706C504C0000002C00000502727552550000002200
00052E617245470000002600000550656E5553000000260000057600560161
0065006F006200650063006E00FD0020005200470042002000700072006F00
660069006C00470065006E006500720065006C0020005200470042002D0062
00650073006B0072006900760065006C007300650050006500720066006900
6C0020005200470042002000670065006E00E800720069006300431EA50075
0020006800EC006E006800200052004700420020004300680075006E006700
500065007200660069006C0020005200470042002000470065006E00E90072
00690063006F0417043004330430043B044C043D043804390020043F044004
3E044404300439043B002000520047004200500072006F00660069006C0020
006700E9006E00E900720069007100750065002000520056004200C1006C00
740061006C00E1006E006F00730020005200470042002000700072006F0066
0069006C901A75280020005200470042002082725F6963CF8FF00047006500
6E0065007200690073006B0020005200470042002D00700072006F00660069
006C004F006200650063006E00FD0020005200470042002000700072006F00
660069006C05E405E805D505E405D905DC0020005200470042002005DB05DC
05DC05D900500072006F00660069006C006F00200052004700420020006700
65006E0065007200690063006F00500072006F00660069006C002000520047
0042002000670065006E00650072006900630041006C006C00670065006D00
650069006E006500730020005200470042002D00500072006F00660069006C
C77CBC1800200052004700420020D504B85CD30CC77C666E901A0020005200
470042002063CF8FF065874EF64E00822C0020005200470042002030D730ED
30D530A130A430EB039303B503BD03B903BA03CC002003C003C103BF03C603
AF03BB002000520047004200500065007200660069006C0020005200470042
002000670065006E00E9007200690063006F0041006C00670065006D006500
65006E0020005200470042002D00700072006F006600690065006C0E420E1B
0E230E440E1F0E250E4C002000520047004200200E170E310E480E270E440E
1B00470065006E0065006C0020005200470042002000500072006F00660069
006C00690059006C00650069006E0065006E0020005200470042002D007000
72006F006600690069006C006900470065006E006500720069010D006B0069
0020005200470042002000700072006F00660069006C0055006E0069007700
65007200730061006C006E0079002000700072006F00660069006C00200052
00470042041E04310449043804390020043F0440043E04440438043B044C00
200052004700420645064406410020062A06390631064A0641002000520047
004200200627064406390627064500470065006E0065007200690063002000
5200470042002000500072006F00660069006C00657465787400000000436F
707972696768742032303037204170706C6520496E632E2C20616C6C207269
676874732072657365727665642E0058595A20000000000000F35200010000
000116CF58595A20000000000000744D00003DEE000003D058595A20000000
0000005A750000AC730000173458595A20000000000000281A0000159F0000
B83663757276000000000000000101CD0000736633320000000000010C4200
0005DEFFFFF326000007920000FD91FFFFFBA2FFFFFDA3000003DC0000C06C
FFC00011080372038103012200021101031101FFC4001F0000010501010101
010100000000000000000102030405060708090A0BFFC400B5100002010303
020403050504040000017D0102030004110512213141061351610722711432
8191A1082342B1C11552D1F02433627282090A161718191A25262728292A34
35363738393A434445464748494A535455565758595A636465666768696A73
7475767778797A838485868788898A92939495969798999AA2A3A4A5A6A7A8
A9AAB2B3B4B5B6B7B8B9BAC2C3C4C5C6C7C8C9CAD2D3D4D5D6D7D8D9DAE1E2
E3E4E5E6E7E8E9EAF1F2F3F4F5F6F7F8F9FAFFC4001F010003010101010101
0101010000000000000102030405060708090A0BFFC400B511000201020404
03040705040400010277000102031104052131061241510761711322328108
144291A1B1C109233352F0156272D10A162434E125F11718191A262728292A
35363738393A434445464748494A535455565758595A636465666768696A73
7475767778797A82838485868788898A92939495969798999AA2A3A4A5A6A7
A8A9AAB2B3B4B5B6B7B8B9BAC2C3C4C5C6C7C8C9CAD2D3D4D5D6D7D8D9DAE2
E3E4E5E6E7E8E9EAF2F3F4F5F6F7F8F9FAFFDB004300020202020202030202
030503030305060505050506080606060606080A0808080808080A0A0A0A0A
0A0A0A0C0C0C0C0C0C0E0E0E0E0E0F0F0F0F0F0F0F0F0F0FFFDB0043010202
0204040407040407100B090B10101010101010101010101010101010101010
10101010101010101010101010101010101010101010101010101010101010
FFDD00040039FFDA000C03010002110311003F00FDFCA28A2800A28A2800A2
8A2800A28A2800A28A2800A28A2800A28A2800A28A2800A28A2800A28A2800
A28A2800A28A2800A28A2800A28A2800A28A2800A28A2800A28A2800A28A28
00A28A2800A28A2800A28A2800A28A2800A28A2800A28A2800A28A2800A28A
2800A28AF927F6B2FDAAB40FD9AFC2B008228F54F17EB4AE34CB0727CB0130
1AE2E3690C2142400010D237CAA400CC98D7AF0A50739BB244CE6A2AECFACE
49238A36965608880B3331C000724927A0159D61ADE8BAA48D1699A85BDDBA
0DCCB0CA92103D485278AFE6CF51BFFDA1FF006A2D4A6D4BC45A9DF6BB1962
EB6E1996CE1E720456E988D71D321727B9279AF19F13F82FC4FF000E35546B
812E9B7F6CD94922668E446E99565C107E95F1D538D60A768D3D3D4F29E6CA
F6513FAC6A2BF063F65BFF00828778B7C07AD587817E3A5FC9AEF85AE5D204
D5A72CF7BA7EE200799F969E11FC5BB3201F3066C6C3FBC504F0DCC31DCDB4
8B2C52A8747421959586410470411C835F5380CC29E221CF03D1A35A335744
B45145771B0514514005145140051451400514514005145140051451400514
51400514514005145140051451400514514005145140051451400514514005
14514005145140051451400514514005145140051451400514514005145140
05145140051451400514514005145140051451401FFFD0FDFCA28A2800A28A
2800A28A2800A28A2800A28A2800A28A2800A28A2800A28A2800A28A2800A2
8A2800A28A2800A28A2800A28A2800A28A2800A28A2800A28A2800A28A2800
A28A2800A28A2800A28A2800A28A2800A28A2800A2BC47E3D7ED01F0F7F676
F05BF8C7C7974732931D958C386BABD9C0CF971212381C177242A8C64E4807
F06BE31FFC140BF685F8B97D3DB7873547F03E82F958ECF497297054F432DE
71316EC7618D71FC39E6B8B178FA747E2DCEBC360A7576D8FE94E8AFE38355
8B5DD66E64D535C9AE2FEE65C979AE1DA5918939E59C927924F26B73C1FF00
143E287C31B95BAF0078AB53F0F3A7F0D9DD4B146C0316DAF186D8EBB893B5
948CF38AF3639F45BD63F89DD2C9DA5A48FEC0E8AFC4FF00D99BFE0A8B793E
AD6BE0BFDA4E38562BA711C3E21B58C42B131C01F6D8106D084E732C4142F1
98F1B9C7ED3595ED9EA5676FA8E9D711DD5A5D46B2C3344C1E392370195D19
490CAC082083823915EC50C442A2BC59E5D6A1283B48B54515F317C56FDB17
F678F83CF2D9F89FC5B6F79AA4276B69FA69FB6DD2B75DAEB112B11C7FCF56
4EDEA2B4A952305793B1118B7B23E9DAFE58FF006ABF8A373F143F697F1BEB
F753196D2C351974BB35390A969A7B18136AB72A1CAB48471F3393819C57E8
178CFF00E0A55F103C6FA837873E04F84A2D38CA485BCD45BED37023FEFF00
948561888E3EF34A3B7522BE5FD07F66DB7D66FEEFC45E36BC6BDD4B549A5B
AB9F21562479A762F21381DD893C0503B002BF31E32E35C052A6A94A77F43C
5C6E2A325CB13E89FD973F697F87BF0B3C172D8EA16E8D72CBF7B8C935F1DF
ED25F18EC3E2678A2E753B28D638DD890057BBCDFB387C3D861291E9D24648
EA2E27FEAF8AF01F1F7ECDE2DA292EBC2D7CEB2282443707729F60E0023F10
6BF3CC0F17606BC942ED7AAFF26CF352BD95CF8EF57BA5646C9AFE90FF00E0
9A7F14751F89BFB2DE950EAF2BDC5DF83EF6E34132B924BC56CB1CD6E327B4
704F1C63D97D6BF997F14C5AA6877F3693AB40F6B7511C323FA7A83D083D88
E2BFA2CFF824BF86354D07F6569B56D450A43E26F10DFEA36B918CC0914166
4FFDFCB67AFD9B86E972BBC5E8D1EC6061667E9CD14515F5C7A41451450014
51450014514500145145001451450014514500145145001451450014514500
14514500145145001451450014514500145145001451450014514500145145
00145145001451450014514500145145001451450014514500145145001451
450014514500145145007FFFD1FDFCA28A2800A28A2800A28A2800A28A2800
A28A2800A28A2800A28A2800A28A2800A28A2800A28A2800A28A2800A28A28
00A28A2800A28A2800A28A2800A28A2800A28A2800A28A2800A28A2800A28A
2800A28A2800AC3F13F88F47F07F87354F167882E05AE97A35ACD79752B744
86DD0C9237E0A0D6E57E74FF00C151BC7177E0EFD956F34DB3768DBC59AB58
E92ECBD4467CCBB719F465B62A7D41C77A8AB3E58B9174E3CD248FC39F8F9F
1D7C51FB46FC54D47C7DE209244B591CC3A6D9B1CA59592B1F2E250323711F
34847DE724F4C63E99FD96FF006728BE255FC4F7A0796482735F9D9A55C88A
5473D88AFD0FFD9FFF006874F8751218DC29502BE1F117949B67D6D1B46364
7DBBF1A3F64CF08F82FC2CD7706CDEA9E9ED5F8CDE3CD2EDF4BD526820C6D5
622BEF9F8D1FB605F78CF4F7B3136548C6335F9B3E22D79F54BB92E1DB2589
358AA7AE869CEEDAEE7297B820D7DEBFB397FC146FE26FC09F86B71F0A9B4B
B7F12FD95B768D75A8CF2E2C62C1DD6ED1A61A5881C18D4491EC1B94120A85
F820C4D7009DC147BD65CFA1B4AE244BB11B29C83B73FD6BD2C3D6E4D9D99C
55A973EEAE7D59F17BF6C9F8EFF1785C5B78CBC5F72BA64FC1D3AC8FD8ECF6
FF0075A28B1E60EFFBD2E7DEBE526D7A4BBBA8AC2C10CB3DC3AC71A2F56773
8503EA4E2BDABE0BFC01B3F88DA8CBA9F8B7549E3F0FDA3F9728B46482E259
301B6C72489322E01C9251B191C1AFBDB58F0FFECFBF0CBE1E5C685F0F3E1D
D85A6ADA94D6F0FF006D6A172FA96A6811C48CD1C93284899D536B7909182A
4FCB5E163B89B2FA55254AA55E6A8BA5A4F5B68AF6B7E27938ACDA349BA715
668A7F02FC036BE14D0ED6C7896FA7DAD712F7673D87FB2BD00FC7A935FAC7
F0A7E0CE9536930EA5AB2798F280429E8057E60F82B5D8ED24825CFDD20D7E
8FF807F684D134FD0E1B4BE38689401835F80B746A661EDB314E506EEEC7C9
51941D572ADFD33ABF8A7F0D341D3F4792F2CE310B20ED5F9F3E261123C91F
A66BE94F8ADF1E63F10C0D636076427DF935F186BDAF094BB16EB5E5D6C352
962A52C2C5C617D13EC6388E49546E9AB23E5EF8DBE06D13C510453DE878DE
D640C64842F9BE5E7E755DDC72B9C678CE0D7F483F053FE15E8F849E118FE1
414FF84462D36DE3D3767516F1A0501FBF9808224CFCDBF76EF9B35FCEAF8E
7568DE09413D8D7D39FF0004ECFDA6A4F00F8B23F823E33B909E1BF13DC31D
2A690F169A948D810E4F48EE4F0076976F1F3B11FD07E1E6652A54FD8557A3
DBC8EFC2D5E5D1F53F77A8A28AFD6CF4C28A28A0028A28A0028A28A0028A28
A0028A28A0028A28A0028A28A0028A28A0028A28A0028A28A0028A28A0028A
28A0028A28A0028A28A0028A28A0028A28A0028A28A0028A28A0028A28A002
8A28A0028A28A0028A28A0028A28A0028A28A0028A28A0028A28A0028A28A0
0FFFD2FDFCA28A2800A28A2800A28A2800A28A2800A28A2800A28A2800A28A
2800A28A2800A28A2800A28A2800A28A2800A28A2800A28A2800A28A2800A2
8A2800A28A2800A28A2800A28A2800A28A2800A28A2800A28A2800AFC9FF00
F82C06892DE7ECEFE1AF1040AEC749F12DB8930405586E2D6E50B303C93E60
8D460FF11C8EE3F582BE7EFDA9FE0E2FC7CF803E32F85D1AA9BFD52C8C9A7B
310A16FED984F6B963F754CA8AAE7FB8587426B3AB0E68B45D395A499FC825
A5E0C0E6B7A0D56688611C8AF3CB94D4743D4EEF43D66DDECF50D3E692DEE2
0946D9229A2628E8EA790CAC0823B1AB91DFF1D6BE6EAE15DCF7A9E20EE65D
52597EFB9359B2DDFBD7386FF8EB5525BECF7A88E18B95724D4755D42DE6FF
004743246DE9D8D41A7EA1E21D5F51B5D234CB392E6F6FA5482089065E4965
60A88A3B9662001EB5425B92DD4D7E81FF00C1333E05DFFC61FDA5B48F155C
C27FE11FF876F1EB377291F29BA89B363083D9CCC04BFEEC4DD0E2BD1A3422
EC9A386AD66B54CD1F87DFB33FEDC7E1BD1DECADBE166A4905CBFDA577C96F
1B65D146195A50CA70A382011D0F35A9E2FF00849FB5578474397C57F13BC0
779A3F87B4B78E4B8BB92585D22323885090923372D205E01EBCF15FD43572
9E3BF06689F113C19ADF817C4719974CD7AD26B3B80A70C12652A594F665CE
54F6201AE4C5F0AE0EA3954E4F79F5F33C5AB494DB93DCFE6BBC35E2D430A1
0FDBD6BD36D7C60C1401271F5AF917E24F86BC65FB3EFC49D5FE16F8F22686
FB4A9488A6DA562BBB624F9573093D639146477072AD8656012D3C7E8CA3F7
9FAD7E3F9A7094D4DE87993C3BB9F5C5D78B3703993F5AE1759F15AEC6F9FF
005AF089BC725D70AF9AE72F7C472DC6773E07D6B9B09C33CAEED131A074DE
26D7CDDEE50DF2D7ACFC12F8476BF12F446B8D5EECDA0B57796DD216026980
39FBDFC001E9DFD31C1AE0FE01FC1DF137ED1BF14B4CF877E1E5923B491D66
D4EF51772D958AB0F36524F1BB1F2C6A7EF3903A648FD30FDA7FE08C3FB37F
8DB4BF895F0EAD3EC7E09D6DE2B5B9B7887EEEC35055C2B60748AE55739E82
50D93F3A2D7BD9DE4D8D8E5752AE0B471FBEDD6C7653C3C64BDE3F407F66BF
8B12FC41F089F0FF0088272FE26F0EAA4174D21F9EEA20311DD7B97C624F47
04E0065AFA4ABF1DFC3BE32D43C2BAD695F16BC180492C1F2DD5BEEDAB3C2F
8F361623380E0654E0ED60AD8256BF597C23E2AD17C6FE1AD3FC59E1E9BCFD
3F52884B131E187386461CE1D181561D9811DABE93C33E33FED4C1FB2AEFF7
D4F4979AE92FF3F3F53B4E8E8A28AFD30028A28A0028A28A0028A28A0028A2
8A0028A28A0028A28A0028A28A0028A28A0028A28A0028A28A0028A28A0028
A28A0028A28A0028A28A0028A28A0028A28A0028A28A0028A28A0028A28A00
28A28A0028A28A0028A28A0028A28A0028A28A0028A28A0028A28A00FFD3FD
FCA28A2800A28A2800A28A2800A28A2800A28A2800AC6D7FC45A1785B4B9B5
BF125FC1A6D84032F35C388D07A0C9EA4F60392780335C47C5BF8B3E18F83D
E1497C4DE237323B9F2AD2D50FEF6EA7232117D077663C28E793807F1E3C7B
F12FC6BF1835E3AE78BEE8B46AC7ECD671922DAD94F68D093CFAB1CB1EE718
03E6788389A96063CBBCDF4FF33AB0F8673F43EF0F1BFEDBFE17D3E592CBE1
F68D36BAEB902EAE58DADBE73C15420C8E3D88435F3F6AFF00B62FC70D41D9
AC4E9BA5A13F2882D4B90067826679327D4E07B015E23A3784EFB50C0B784B
67D0575D3FC37D5E087CC78180FA57E538CE34C7D595F9ECBB2D3FE09EAC30
54D743AFB4FDB2BE3C698774F2E9BA960938B8B42A082318FDCBC4703AFAE7
A9C715EC7E11FDBFB4BF392D7E2378626B15240375A74827419EED0C9B1801
DF0EE7D057C59ABE8B35992B22E08AF3ED4AD5483915BE038BF1D07AD4BFAE
A4D4C1C1F43F7EFC15E3EF06FC46D197C41E08D5E0D5EC58ED2F0B7CC8DD76
48870F1B639DAE01C738AEBEBF9BEF0AF8F3C69F09FC4B178BBC05A8BE9D7D
16048A3261B88C1C98A78FA3A1F43D3AA90C011FB57FB367ED27E19FDA07C3
2F244134EF13E98AA352D3B77DC27813439E5A173D0F553F2B762DFA9E47C4
30C5AE56AD2FCFD0F2EBE1DC35E87D2D451457D11CC1451450014514500145
212141663803A9AF9FBC6DFB43F84FC39349A6F87A33AFDFC670DE4B85B643
DC19B0DB88F440C3B120D70E3F33A18587B4C44D457F5B772A10727647D055
0CF716F6B199AE64589071B9D828E7DCD7E7C6B9F173E2878A19C36A874AB7
73911588F2703B7EF3265FFC7BF0AF38B8D224BF97ED1A84AF752E31BE5632
37E6D935F9E63BC52C2C1DA85372F576FF0033AE38297567E9AFFC257E17FB
47D93FB62CFCFCE3CBFB447BF3D71B7766B62DEE6DAEE3F36D65499338DC8C
186476C8AFCA897C3D6E46360FCAB25B4792C64FB458C8F6D28C61E2628DC1
C8E5483C1E6B868F8AA9BF7E87FE4DFF0000B781F33F5D28AFCBAD13E34FC5
CF06C8A6D75A7D52D93AC1A80FB4AB7FC0C9128FC1C57D31F0F7F6AEF07788
A58749F1B47FF08D6A4E4289646DD64EDFF5D4E3CBCFFD3401474DE4D7D965
3C6982C5351BF2CBB3FF003D8E6A986944FAB68A6A3A488B246C1958020839
041E841A757D698051451400514514005148CCA8A5DC85551924F00015F0C7
C6AFF82857ECF5F0824B8D22C7527F196BB0165369A46D9628DC7696E98885
7904304323A9EA9401F74515FCF0FC49FF0082A9FC77F13B496FF0FF004ED3
BC156A5B28EB18D42EC0F4325C2F927F08057C73E27FDA8BF687F184A64D7F
E236BB32920F971DF4D6F0E4743E542C91E7DF6D007F5C9457F197278F3C67
348D2CBAFDFBBB92CCCD7529249E4924B724D5BB0F899F10B4A91A5D33C51A
A59BB8DACD0DECF1923AE09571C5007EC57FC1463FE09DF7FF0014AEEF7E3D
7C07B2F33C5A543EB1A3C5B57FB495171F69B70703ED2147CE9FF2D80CAFEF
7897F9DCBC8B53D1EFA7D2B56B796CAF6D1DA29A0991A3963910E195D1802A
C0F0411915FA1DE10FDB77F6A6F064B1C9A6FC45D4EF51082535275D495864
9209BB595B07246410718C11818F4BD6BF6ADF843F1E2E94FED69F08B4EF11
5E90A9FDBFE1F77D3357444185DFF3EDB8DBC80AF22A01D17239C674933585
568FCA4176C7F8A94DC0F5AFD9DF09FEC4FF00F04EFF008A7A65DF8ABC1DF1
A756F0F59DAC7E75CD8EAF71656F73681B92089A18B7AAE42E50C83381BD89
AF2BBDF057FC137FE11DFB7F66DBF8AFE346A36CE46CB8B94D334866463C89
228A1B823A74122B633C29C1C7EAE6BEDCF8ABF67FFD9D3E2B7ED31E3387C1
FF000CF4B6991597EDBA8CC192C6C2227992797040E39545CBBE3E553CD7F5
93FB31FECE3E0BFD983E16D8FC3AF0928B8B8245C6A7A832ED96FEF5942BCC
C096DAB8015101215401C9CB1FC74D2FFE0A5BE28F0168917857E0B7C2DF0C
783742B53FB8B3DB34E8A3A12DE43DA8776006E7232C412724F12A7FC157BF
68C5914CBE1BF0AB2023701697CA48EE01FB71C1F7C1FA56F4E0A2633A8D9F
D06515F8A7E0DFF82B9CCB32C1F117E1D8F288F9AE34ABCF981CF68274C1E3
FE9B0FD78FD08F835FB65FECF9F1CE7B7D2FC21E254B3D6EE07CBA5EA4BF63
BC2C4E3620726395BDA277E3E86B4B99937ED3DFB287C35FDA93C291E8BE30
46D3F59B00C74ED5ED914DD5AB1EAA777FAC858FDF88900F552AD861F813F1
5BFE09DBFB53FC26BD9CE93E1F3E37D223DCD1DE6899B87651CE1AD78B80F8
ECA8C339018F5AFEA528AE5C460E153E244B8A67F1C70FC25F8EDF6CFECC5F
875E243781CC7E48D22F0C9E60382BB3CACEECF18C66BB1B1F813F116C3C5B
0785FE2769D75E0D924863BAF26F222975243212159236E80952096E847438
C57F5D75F107EDC7F0326F897F0EE3F1EF85ADFCCF16781D64BB80203BEEAC
701AEADB03966DABE644304EF5DAB8F31ABC0CDB2497D5E6F0CFDFB6865529
7BAF9773E50FD8BF59F0D7C05F1CFF00C2236702C1A278CBC8B69AE6439945
EC5BBECD24921E4AB9768CA8C282E08C0073FAC7E3AF057877E23F84356F03
78B2D85DE93AD5BBDBCE9D1B0DD1D1BF85D1806461CAB00C3915F817E1AD5A
D3C55A0261838910739F6AFD12F0AFEDDFF0D7C19F0E74FB7F8AD757B3F8B2
C10DBCB6F6D6CD34D7822002DC090EC8419063706914EF0DC6304FC87875C4
53ABCF9762DDE6AED5FAAEABF5FBFB1CD82ACDDE12DCF8A34BD3F5FF0081FF
0012358F82BE3E7337D94836F7046D5BCB2949FB3DD20E7EF0043004ED7574
C92A6BEC1FD9DBE233FC32F1AB780B5B9FFE29DF12CA1AD6563F25BDEBF098
3D927E14FA3853C02C6BE25FDAAFF6B9F07FC768B47B8F0BF826E34ED5F419
F7DBEA97173189DAD9C7EF6DDA18D1832B1C32FEF7E565047DE615A7E0EF1C
E89E3AF03245A84B99E341E51C65981FE1E39E3F4AFCF78AF2BAFC399BC333
C12FDDC9EABA5BAC5FF5E67A07EF5515F983A67ED33F19A4F0D5A68825B48A
E2DE358CDFF9264BA90200033F98CD1973D58ECE4F61CE793BCF1EFC54D4E5
FB4DD78BB55573FF003C6EE48179FF0066228BFA57D8E65E3A65146CA94653
7E892F4D5FE96F303F5AA8AFC878BE217C57D2A433D9F8BF552F907F7D7525
C2F19FE198BAF7E78E7BD7A9F84BF6C2F1B787258EDBE2069F1EBD643EFDC5
B2ADBDD81EBB46217FA009F5AEAC97C6ACA715350A8A54DBEAF55F86BF801F
A4B45711E02F88DE0EF899A28D77C1BA8A5F5BA9D92AF292C2FF00DC9636C3
237A64608E412306BB7AFD6E85785482A94E49C5ECD6CC028A28AD4028A2B8
0F15F8EADB41BB4D16C105D6A92A07D873B22427019C8F5C1C28E4F720609F
0B89789B039460AA66198D450A50DDBFC92DDB7D12D4691DFD15E28DA978BA
E54DC497CE99E8A8AA8ABEC0633F9926B9FB8F89DE21F09389B565FED5B052
04830A93A28EA508C2B11E8C327FBC3AD7E03957D2BF86B138C8E16A53AB49
4B69CE0945F9E926EDF2F5B1D1570752094A48FA2E8ACBD135AD33C45A4DAE
B7A34E2E6CAF103C522E791EE0E0820F041E41183CD6A57F4C52AB19C54E0E
E9EA9AD9A39828A28AB00A28A2800A28A2800A28A2800A28A2800A28A2800A
28A2800A28A2800A28A2800A28A2800A28A2800A28A2800A28A2800A28A280
0A28A2800A28A2800A28A2800A28A2800A28A2800A28A2803FFFD4FDFCA28A
2800A28A2800A28A2800A28A2800A82EAEADAC6D66BDBC9561B7B74692491C
E15110659893D000324D364BCB38A65B79678D256C610B00C73D30339E6BE4
CFDB7FC7B71E06F807A9C1632986EFC4B710E931B29C3049F74938E3B3431B
A9FF007AB9F13898D3A72A9D91518DDA47E6CFC69F8BF77F1BBE23DDF89433
AE8B685ADF4B85B8D96CA78765ECF29F9DBB8E172428AB7E03F0FB6B37F140
0677102BE7FD0E655541F4AFAB3E10EAB6B65A9C12CA40C115FCFB9AD69D6A
B2A937AB3E8292515647E91FC30F847A2E91A541737902C933A83823815E9D
AC783344BBB3781ED517230085C62B3FC27E32D1AF74B8196E1119500209C7
41595E35F89BA268961298EE1649704000F02BD8A54F2C8E5EF9B5AAFF0033
9A4EABA9A6C7C01F197C3F69A55F4F1C200009AF907566556602BE80F8A9E3
75D62FA6757CEE26BE62D4AFB7B31CF5AF99A1499DB2672BAB10C1B35CDF82
7E28789BE0B7C40D33E237845FFD334D93F790B1C47736ED8F36093AFCB22F
19C654E187205686A9743079AF24F10CC191BF1AFACCA9CA12528EE7255D51
FD52F80FC6BA0FC47F06E8DE3AF0C4E2E34BD72D63BA81BB859064AB0ECE87
2AC3A86041E4575B5F947FF04AAF89971AE7C3EF16FC2DBE94B9F0B5EC7796
818F2B6DA8EF2E8BCF459A3773C7593DEBF572BF65C356F694D4FB9E34959D
828A28ADC90ACDD5F57D3340D32E759D66E52D2CAD10C92CB21C2AA8FE64F4
00724E00E6B4490A0B31C01C926BF357E38FC5C93E26788CF87B429B3E1AD2
A5F90AF4BB9D720CC7FD85C90833823E6EE00F0B8833CA780A0EACB56F65DD
FF0097735A549C9D8D8F891F1B35FF00893712695A21934BF0E8257603B66B
A1EB3107853D9071FDE2DC6386D374B48D4055ACDD22D000A00AF64F0B786A
6D4A545442735FCD79D6735F1755D5AD2BBFCBD0F629D3515646059E8F24B8
DA99ADC5F0E5C15CF967F2AFA5F41F0059DAC4A6E177BFA575FF00F08CE9E1
31E42E2B8F0F9362EB479E9C1B45B9C568D9F14DD68B2C40E54D72D7966572
08AFB37C41E0CB5920792DD3040E95F36F89349FB248CBB7A571384E12E59A
B31BB743C5350B45C1E2BCE35AD35240C0AE735EBBA920048AE0F544054D7A
B85A8D32248D1F84DFB4478B3E0C5E45A4EA3E66B5E142D87B366FDEDAA93C
BDB3374C75319F95BB6D24B57EAE785FC51A0F8CF41B3F137866F12FF4DBF4
12453467823A1041E5594E432900A9041008AFC3ED7AD95D5863AD74BF003E
3F5F7C08F198B6D625925F076AF205BF8465FECEE785BA8973C32F02403EF2
762C171FAFF08F134E16A15DDE3D1F6FF81F91E7D7A17D51FB7745436D736F
796F15E59CA93C13A2C91C91B0647461956561C1047208EB5357EAC7085787
FC77FDA13E1AFECEFE126F14FC40BFD924DB96CEC202AD797B22E32B0C648C
85C8DEE4845C8C9C900F27FB50FED43E09FD997C12DADEB656FF005FBF574D
2B4A47C4B7328FE373C9485091BDF1EC32C40AFE617E2BFC5AF1CFC67F19DF
78EBE206A4FA86A578C7009222822CE56181324471A67851EE492C49201F46
FED23FB72FC5FF00DA0EE2EB48374DE1AF08484AA68F6521C489FF004F5300
AF3B1EE0E23E3840724FC52D2135117CD3E385E53C0CD0030BD34935D05AE8
9713630A4D6C2F856E4AE761FCA80386C9A6EE35D7DCF87A7881250D73F736
3245D450050DE6A7883CA70B5544659F6D6F5AC214000526C0960B50304F26
B522B627A0AD5D27499AFA4548D73935F43783FE0E6A5AC046580907DAA00F
9C058CB8CEDA824B665EA2BEEAB8FD9F350B7B6F31ADC8C0F4AF10F167C3C9
F492C1A3DB8A00F9DE48AB3658B0432F0C0E411D4115D66A3646DDCA91D2B9
F9475A00FD01FD993FE0A37F13BE0DDDD9F85FE29CD71E35F05A058BF78C1F
53B2407EF4133906655071E54ADD0008E8060FF421E00F883E0CF8A5E14B1F
1BF807558759D17515DD0DC40D9191C32329C323A9E19180653C100D7F1AF7
080835F537EC7BFB59788BF65CF88714D732CD7BE06D66544D674E439014E1
7ED70A9E04D10E7031E628D8C47CACB4981FD5A57C33FB47FEDCDF0FBE0A4B
3785BC2C89E2CF17A7CAF6D0C98B5B33D3FD2665DDF38FF9E4996E3E629904
FC7FFB57FF00C14026F138B8F879FB3EEA2D0E912A05BBD7A1CA4B701D4EE8
AD3700C8982332F0E483B30A3737E5FDBC1CEE6E58F249EA49AA03D53C1FE3
F9DBC6DA8CF7B6B6FA743ACDCCB771DBDA218EDA13339778A2462C51149F91
7270B803A57ACFC44F0DC7E21D1C6A16ABBA7806F523A91DC7E35F304D632C
F106B63B2E2121E26F475E9F81E87D8D7D35F0D3C531F883475B7B8F9664CA
3A375565E083F435F8671D655532FC6C333C2E9ADFE7FF0004F2F154DC25ED
22780E93A24FAADF47636E30CE793D947726BEC7F00783F4ED02C63B3B0842
E49663DD998E4927DCD79BD9E836FA2789A6DABB56EBE743FF00A101FCFF00
FD55F47782628EE2F218DBA122BF3AF1478D2AE62E14A9E94D24EDDDB5ABF9
6CBFE09E852A8A515247B77817E1A6A3E222A628F09DC9E95EDCFF00021A3B
5DC1C1703A57B8FC38D36C6CF41B61080372824FBD7A2DD98162E315F13937
012C5E0678EAB5946DB23D1A7878D95FA9F99DE31F075C6852BC52AE315E1F
AC5BA10C08AFB43E37EA3626E1A38C82CA39AF8AB58BB52CD835F1B82A528C
AC8E492B368E57C31F107C55F087C590F8CBC1D3F977117CB3C0F930DD439C
B452A8EAA7B1EAA70410457ED27C25F8A1E1EF8C1E05D3FC71E1C6DB15D029
3C04E64B6B94C09617E072A7A1C0DCA430E1857E14788A65646AF65FD86BE2
E4FE07F8DC3E1EDF4FB746F1AA3C4A8DF752FE152F0B8E78DEA1A3381F3164
F4AFE95F0A388AAD1A8B0951DE12DBC9FF00C1EA49FB774514C9658A18DA59
9C468A32598E001EE4D7F44CA4A29CA4EC900FAFCEDF873F12EDBC55AF5E78
A751939D56E1E7C13F7119BF769F445C28F615F6AEA3E30693745A3202A073
71270A3FDD53D7EA78F635F9C7F19B4AF017C2AD2ECAFF00C1128B3FB1EC86
7B71216468FEEAB8DC4B6F071BB9E47279049FE1BFA4278859567F3A191E59
5BDA38C9B6E3AC39B45157FB5D75578ABDAF73B23879C23ED1E9D8FBDAE7C6
3A2269A15654C63B1E6BE63F887E32B37B7982B8DA735F1EAFC7481ADFFE3E
474F5AF24F19FC6637B1BC56F2EE273CE7815F8AD7E14CDB35C552963BFE5D
A5156564921D7C5CA6ACF63F4A3F62CF88336B57DE36F0349217B7D325B6BE
B619276FDABCC49947A006346033D598FAE7EF4AFC57FD8A7C4DA9F86DFC43
E35B7B8459357962B648E5019658ADF7163D8805DF00820E57D3AFEAE785FE
28E81E2031DADE9FECDBD7E0472366373FEC49800FB03827B035FDAFE1A789
993D384387EBD7E5AF4BDD5CDA29765196CDABF2DB4775A27B9C6D1E994514
57EF220A28A2800A28A2800A28A2800A28A2800A28A2800A28A2800A28A280
0A28A2800A28A2800A28A2800A28A2800A28A2800A28A2800A28A2800A28A2
800A28A2800A28A2800A28A2800A28A2803FFFD5FDFCA28A2800A28A2800A2
8AF99BE31FC48B8FED16F00F87E63190A3EDF321F98071910A9ED953973E84
0F515E0712F1150CAF073C5E2365B2EADF44BFADB535A345CE5CA8EB7C53F1
7ADADE6934AF0846B7F7284ABDC3E7ECE841C10B8C190FD085F73C8AF2FBB9
7C51E206F3358D4A79C37F02B148C7D1130BFA6693C1BE1E1745111708BD4D
7BD69FE1B86288048C1F72335FC779BF17F1071162254F0EDF2FF2C6E925E7
DFE7F81F414F0F4A8ABB3E64B9F07DA2211E4281F4AF2DF19FC3DD1B5FD386
97AB5AADD5AC4DE624720CAA3E08DCA3F85B04F239E4D7DD3A8E856EF1B2BC
601F502BE7EF15D8ADA4D2467B66BE2270C765F88F6756F1923A538C95D1F9
B1E31F843FD84CF75E1F6611AF3E4B9CFF00DF2C79FCFF003AE1F45F104FA5
CFB5B31BA1C1078208AFB6FC53047223823D6BE54D5ACB40D27C5D61AC6BB0
1974E8E61F6951FDC3C6E20677053C91DC0C57EADC3DC4552AC396BFBCFA77
392AAE55747D17F0DE5F1CF89EDA39ED99AD2C9BA4AF9CB0FF00657B8F73C5
7D016BF0E20BB50FAA33DE3F7F30923F21C576BE07D12DAEEDA07B1DAF03AA
B2321055948C8208E0823A62BDDEC7C35188C7C99AF9C78DCD735ACE960A2E
CBA2FD59E27B4AB5AEEF647CA97BF0ABC32D1147D2EDC8C639894FF4AF0FF1
97C07F085F44E61B2169276680F978FC07CBFA57E88EA7A0C42361B71815E1
1E27B4581DD0D78F2AD98606BFB3AD27197A9CD56152934EE7E407C4DF853A
FF008444B776A4DF5926492A3F7883FDA03A81EA3F215F23EBD7E0AB735FB4
BE35B18678A45600839AFCABF8F9E024D0EF9F5DD2936DB4AC7CE8C74563FC
43D8F7F7E7D6BF6DE0AE22FAC354ABFC5D1F73A68E2DCBDD91F24EAFAA6A9A
3DDDBEADA2DECDA7DF5BC9BA39EDE468A543EAAE84303F435F5FFC16FF0082
97FED37F08E4B7B0D77564F1EE8516D56B5D6B325C041D7CBBD5C4E1C81806
532A8EBB339CFC55AFCBE6154F7CD728EB5FBE65B754919D5DCFEB4BF664FD
B97E0B7ED37045A568576741F1708CBCDA25FB289FE5FBC6DE4184B841D72B
87039645AFB32BF85FB4BDBFD2AFA0D4F4BB996CEF2D5D6586785DA3963910
E5591D482AC0F208E41AFE82BF604FF828D9F88F7561F04FE3F5EC50F89DC2
41A46B2F88D3526E8B6F7273B56E8FF038C2CBF74812E3CCF4E32B999F787E
D5BF1325F04F81E2F0D6933797AB789D9EDD594E1A3B5403CF718E84E5507F
BC483915F01787A158D100ED8AE83F693F1B7FC265F1D3588A17DD69E1D09A
5C3FEF439698E3D7CE675CF70A3E839DD1E6002D7E13C71984ABE2A4AFA474
5FAFE27A98685A27B2F8720134A8BEA6BED2F877A1C56D64B7257E63C0AF8C
3C25708B7084FA8AFBA3C0B7D0CDA6468A4652BE0B0B4A0EBC154DAE7549E9
A1E9B042A8838E6A7201E2A08A642A326892E110715FD3786C5E128E1D72C9
28A478D28C9B32EF9170476AF94BE207949712EDF535F43789B5E834EB5766
61BC8E057C77E32D7C5CCD21DDD49AFE73E23C453AF8C94E92D0F5E945A8EA
799EAD28DED5C0EA528C1ADED46F0124E6B88D42E739E7A571E1E98DB392D6
1C1045789F8A235911C1F7AF58D56E7EF1CD791F8826055ABEA72E8D9A3199
F4E7EC9DFB5AEBBE13D42DFE13F8FB5663A1A04834CBA902B1B30485485C95
39873F2A93FEAF819D9F77F57D3C43AD45C89D25FF00AE918C7FE3BB6BF987
D52E96D35AB79DBEE3B796FF00EEBF07F2383F857EC97EC69F1A67F889E0DB
9F04F8867F375DF0988E20EC72F7164C36C32127AB2106373CF456272D5F9C
F8CCB3DCBA8C73CCA31952105653829CB955F452516ECB5B292B754FB9D997
7B29BF67522BC8F897FE0A0BFB3E7C63BDF156A3F1DA7D4A5F15E86E9FBF8D
136368F6E87091A45B9B36E3392EBC8259A41C973F94A5B35FD7F4D0C57113
C13A2C9148A5591802ACA460820F0411D457E0E7EDBDFB1DCBF09F54B8F8A7
F0E6D07FC215A84CBF69B5881CE973C9C631CFFA3C8DF70F4463B3006CCFB1
E057D206599D48E4F9E4D7B67F04F45CFF00DD97452ECF452DBE2F8B3CCF2A
E45ED296DD51F9DD6B6CD3B80066BD5FC31E0E96FDD004CE6B0FC2BA48B99D
38EA457DD3F0A7C0D1DCB43B9339C57F5C1E01C67857E104D74887C9CE7DAB
D34FC1478E1DDE4FE95F7F7807E19DAFD9E32621D076AF4BD5BC016505A93E
581C7A5007E32F8A3E199B356FDD631ED5F35F8AF414D39659651B5230493E
C2BF5BFE27786AD604970A3BD7E5E7C7D9E2D36DA2B088ED92F65C71FDC8F9
6FD768FC6803E74B63E639908C6E39FA574B61179B22AE2B99B520015D768E
CA265CFAD6607D55F077C191EAB7B0874C82457ED27C11F831A4B69F14D3C4
BC28278AFC91F829ADDB585CC0EE40C115FAE7F0CFE30E9BA669D1A348A30A
3BD007B1F8D7E19E836DA43C9146A0007B57E4FF00C75D06C2CE59C44A0633
5F7E7C48F8F3633584914532818EC6BF2B7E30FC428F549E62B26724F7A00F
8E7C551A25C4817D4D79BCE464D75BE20D405C4ECD9EA6B89964C934014E73
585760106B5E67AC5B961CD340773F0EB570D2CBA1CEDCA832439F4CFCCBF8
751F8D7B75A405C80057C91A6EA2DA56B965A829C7952AEEF74270C3F1526B
EDDF0FD81B899062AC0E87C3DE1B96F9D42AE735D56B3E0FD5FE1DDDDAF8D2
DE261A4DEBA41764748A66E2373E81FEE93D3701DDABE88F84FE075BE961DD
1E7915FA27A7FECFFA0F8D7E1EEB1E15D6EDD64B5D62CA6B560C4A732A10A4
32FCCA54E0865E41191C8AF3337CB29E330F2A153664548292B33F15BE237C
63F0A689A746B6528D435A8B0C9144788DBFE9ABF45047040C9C761D6BDCFE
1AF8EB4FD6B4BB0F10E9AF982E903804F2ADD191B1FC4AC083EE2BF21E08AE
749BF9F46BF0565B662A770DA48EC707A7B8EC78AFA2FE0678F1FC39AC3F85
EF64C596A4FBA027A25C6318F60E063EA07AD7E0FC55E1ED18E5FCF866E528
EB7EEBAD9796EBE65D2C2C69C3DD3F783C07F1C934BB14B6B860EAA3B9AE8B
C47FB41C4F6AF1D961091D73CD7E6A5AF8964451B5C8FC6AD49E2491C72E7F
3AFC39E06AA5C89BB7634556495AE7B5F8B7C6F3EB13BCB2C85B71F5AF20D4
B5404139AE6AE75CCE72D5C96A3AD8C1F9ABAF07943BEC663B5DD486D6E6BC
1756F166AFE17D7B4EF14F876E4DAEABA35D457B693001BCB9EDDC491B6D60
54E1941C1041E8462BA1F13789EDAD10F9D260B70AA39663E800E4D74FF01B
5BF871E1DF88563AF7C76D2209741B864487ED4E59ADE4661B677807CAF18F
E30C0E07CC32460FE97956071184C34F1B468CAA382BA8C6D776E8AED6BFD6
FA1AD3A4E4EC8FDE0F801FB45C3F1A3E10E8DE3C9348B8D3F57BA530DCDAC9
1B471F9F1001E589D861A07277230CFF0074F2A6B67C5BE34B2D3A26BDF12D
E2E1395814E1171EDDCFB9E6BE64F8A3FB4C783FC01A735968F2469E5AE102
63A01C631D0578CF83BE0DFC7CFDA7A68F5ED7E793C13E0D9C875B8B8426EE
EA3241FDC404A9DACBD247217905438C8AFC871B9AF1878835BD8422E8E11F
D95749AFEF37673F9DA3DA37D4ECF694A8FC1EF4BBBDBE48D6F897FB51EA1A
CEAB1F82BE1C59CFABEAD78C63B7B5B34696591B19E15012700127D0649C0A
862FD833E2E7C5BF08EA1ABFC52F178F0C6BB7B16EB0D3EDD05DADBBF040BB
9030539E54AC45B6F0DBDB94AFD12F843F01BE197C10D2DAC3C09A588AEA74
55BAD427225BEBADBCE659880719E422858C1FBAA2BD8EBFA43C3CF0232AC9
20A7523ED2AF56CE2A956537CD27767F393AF7EC03FB647867507B0D2F42B4
F14DB2B305B9B2D4EDA28D946304A5E496EE339E9B78C1F627E81F845FF04D
1F891AFDE25EFC72D761F0FE94339B2D2A51717D21E98333218621DF204A48
E30B9C8FDB5A2BF4C87076014FDA725FF2333F0D7E207ECAFF00B42FECDAD2
6B1E0667F1F783AD896CD9C65751B6881CFEF6D464B803ABC25B382CC882AC
7C31FDA5B45F10451DADECCA1CFCACAE7073DC106BF70ABE49F8EDFB197C1F
F8E124FAF4B6ADE18F16BA929AC6980472BC8790D731711DC0CF52E0498E16
45EA3F30E3DF0272CCDA32A9463C950699CF780FE325FD8C312595D2EA1640
002DE66FBA3FE99C9CB2F1C0072A07415F51F867C7BE1DF14621B49BC8BBEF
6F36164FAAF2430F75271DF15F88FE3AF86FFB45FECAF712DEF8A2D1BC49E1
0B76F975AD3C17895077B88B264808EE5814CF01DABD67E1BFED09A0789E08
43DC2EFE0839C303EA3B835F8DE0F88F8BF82AA2A1898BC4E197493774BFBB
3D5AF4775D921E87EC9D15F27782FE35DFC11C71DF49FDAF678EA587DA147B
39E1FE8DCFFB55F49681E28D0FC4D019B47BA598A637C67E59133FDE43823E
BD0F626BFA43813C5AC9B8823CB83ABCB57AD396935E8BED2F38DFCEC268DF
A28A2BF4B10514514005145140051451400514514005145140051451400514
51400514514005145140051451400514514005145140051451400514514005
14514005145140051451401FFFD6FDFCA28A2800A28A280399F19F896DFC1D
E13D5BC537437A6996D24E1338DECABF2267D59B0A3EB5F9BBE14D46E2FE79
352BF90CB7579234D339E4B49212CC4FD49AFAABF6BDD67FB23E0A5F47E6F9
6750BCB2B7E87E6FDF094AE474E23E7DB8EF5F0BF84B5A55863C37615FCE9E
37D7A952A52C3AF852BFCDBFF247AF96A49367E867C39F25AC4B0C6EC8CD7B
741342B10E715F0E782FC7A74A2B96CA9EA2BD907C57D38439C7CDF5E2BF35
E02E35FEC3F689D2E6E63B31587F696D4F60D62F6144691C80AA324D7CA5E3
7D6A2B8BB9590F19352789FE26CB7E8D144DB13D057856B1E21F30B33375AF
9AE20CCEAE6B8C962AA2B5CDA941423CA8CBF115FA957E6BE5EF1ECD1CB0CA
0F3906BD5FC41AE2ED6F9ABE6DF1AEB4A6393E6F5AFA6E1DC0494918D591FA
3BFB06FC401E2AF005F785AFE4DF7DE14B916C09EA6D2505ADF27FD9C3A0F6
415FA2914B108C60818AFE69FE087C68D73E1178B755BDD22EDAD5753B7512
6155B734326532181E81DBFAD7EB47C0FF00DA962F893E17BA9B5040351D36
5114ACA36AC8AEB947C7404E082071C67BE2BF43CBB399E4152BD7851E684E
CEFD9BFF008278D52705EECB43ECBD72FE08A277620002BE57F186B114D712
143C64D2F897E244B7E1915F6AFA0AF11D67C441F712D5F9467B98D5CCF192
C55456B9C788ADCF64B6462F8A2FD591F9AF8C3E2CC76FA869F756F300CB22
B023EB5EFDE25D7976BFCD5F2978FF005813AC918392D9AFAEE19C14A334D1
944FCE6D7ACA6B4BE9E09BAC4ECA33DC03C1FC6B9892BD2FE206D5D666C7F1
007F4C57994879AFEA5C04F9A8C25DD23A1BBEA557AA65A485D66858A49190
CACA70548E41047422ADB9AA9273C57688FD7CF0AEAB7B7110BFD66F25BEBE
B8264B8B899DA49A799FE6791D9B2CCEEC4B127249249AF58B3D5755741F63
8553DDF2C7F20463F335E3FF000E6CDAF2DED24994AC92AAE14F55DD8E3EB5
FADDF067F67EF0DCFE1CB7D635F8FCF96E5772AF4007BD7F3C62A956C4E27D
8D08F3499EE414631E691F0445E32F1868ECB2C0607DBFC3246707F2607F5A
F74F86BFB52E9DA7DEC3A578BA2FEC895C855B8DDBAD98FF00B4C7063FF816
57D58574DFB457C3BF0EF845D24D2808C48B9DBDC57E7AF8A3CA912453C8AF
229E1EF374EA2D51AB4AD747EE059FC4FD3E48433E0E47507AD666ADF152DA
38C8B7C03EB5F917F02BE2FDF4534DF0F756B92E6D50CB62ECC4B1894E1E2C
9FEE64151FDDC8E8A2BEC3F0FE87E2FF00186D7D3E031DBB7FCB6972A98F6E
E7F0E2B833BC64701072C655E582EADD97F5E44D3839BB456A771E28F1ECD7
CCE5A4EBEF5E2DAA6B265624B57D19A67C11B0C0935CBE92E5CF558FE45FEA
7F5AEBA2F849E0589769D395FDDCB37FE844D7E458DF19F27A32E5A4A53F34
ACBF169FE077C72DA8F7691F075EDF6ECF35C85FDEF04035FA2B79F06FC097
4853FB3235CFA654FE60D791F8A7F66FF0FDC46EFA549359376DAFBD7F10F9
3F9115D99578DD93D49A8D58CE1E6D26BF077FC02795D4B7BAD33E08D56F78
201AF29D76F061B9AFA2FE21FC1BF1AF85D64B8B78BFB4ED9324B42089001E
B1F5FF00BE49AF9135EBF2BBD5B20A92083D735FBCF0DE6984C75355707514
E3E5D3D56EBE67935E9CA0ED2563CD3C4D279ACF83D2BE86FD993E254BE02F
8F5E12D59A6F2AC75F74D2EF0670AD1DF6D542C7A0093796E4FA2D7CCDAB4B
BD5D8F7CD67FF69DDD8269FA85A4863B9B2F2E589FBABC64321FC0815F6B99
64B4F1F80AD81ABF0D48CA3F7AB7E1B9CD0AAE32525D0FEA7F22A9EA5A6E9F
AC69F73A4EAD6D1DE595E46D0CF04CA1E396371B591D5B20A907041EB4DD3E
F62D4AC2DB518032C775124AA1B0182BA8619C123383CF35F1C7ED4BFB6D7C
38FD9B6D5F434C7887C69345BE0D2E070161DDF764BB9067CA5EE179761D00
0770FF002E387787F30CCB1B0C165D49CEAB7A25D2DD5BD925DDB491F695AB
4211E69BD0FCFF00F8FBFB36CDF007C73F6FD3732F83B57919EC676FF9776E
59ADA563FC4839563F7939EA180C9F0FFED33F07BE1CC48DA96A6DA8CE9FF2
C6C23FB431C7FB7958FF003715F0B7C6DFDA6FE337ED07A87DA3E22EBAF269
E8DBA1D32D41834F84F6DB0827730ECF2177EDBB1C578308FD057FAD9C0F4B
35A19550A39D4E33C4455A528DECFB5EE95DDB77B37A9F0B897094DBA7B1FB
3EBFF0556F0A787D1A2F0DFC3EBDD4D53010DD5EC5665867A911C7718E3DCF
3C7BD559FF00E0ADB0EA7986FF00E183DAC454FCD16B0266CF6F95AD2318F7
CD7E377947D2985315F55ED198729FAB3ACFEDDBF0D7C680C57DA7EA1A2C8F
D5A58D268867FDA8999CFF00DFBFF0AF8BBE3278CF49F17F886C6E345BB4BB
B48E1660C8780CEDC820F20E1464100F4AF9CD96AEE9EDE5923D4D6909B6EC
C4D1DE5BC9C0ADFB2B8F2DC1AE3EDE6E056B45362988F79F0BF8BA4D3990AB
E315EF3A67C65BBB58022CE463DEBE2182F993A1AD25D5E50386A00FACB5FF
008BF797A8CA67273EF5E11AF78AA6BE762CE4E7DEBCFE4D4E57EAD59F25C9
6EA6802FDCDD1918926B32492A179AA9C935002CD275AC7B89073534D35645
C4DD6A920332F1F76457E8DF8000952D6490E5991093EE40AFCDD94B48DB54
6598E00F5CD7E87F846ED2CFC9881E230ABF9715407E9B7C16FB344F016C76
AFD3DF08DDDB4DA5451C4C32074AFC5BF86DE306B778922CBB9C602824FE42
BEF4D3A4F10DAF86BFB77E206ACDE0FF000E6DC950DB6FAE863EEAF78C1F6F
98FB0AF98E22E2DC265B1FDEBBCED75156BFABE915E6ECBA2BBD06A2D9F817
FB6EFC325F865FB4178C74BB0541691DF4B7B6661398CD95E1F3FCA5F7B767
319079051BAE41AF99ED2E0CF0A4D0B959130CACA70411C820FA8AFD49FDB7
3E2BFC39F895E18B1F06F84347834D4F0FCCD2E9D2AA837323B8DAEAEDD76C
A00C8E7E60A4938AFC9ED1E2BB4BD96D044CAAAC78618DBEAA47A83C62BE6B
86339FAF61E551C79757A6B6F95D26D7AA5D7446F49D9F29F68F827C7EDAEE
8B14F72C05DC3FBB9C0E3E71FC407A30E47A74ED5DA1F118C7DEFD6BE79D0B
E16FC423A3CBE25D2B4FB94B36037CC8195481D3F2CD72CFE28D7F4891A2B8
7370A9C157E187FC0BAFE79AF93CC38154AA4A741AB3D97E812A0F747D3779
E268E38DA49240AAA3249380057113EBFAA6B8B21D1C086D23FF00597B3FCB
120FF673F78FF9E6BCF8EBBE1BB7D3E0D6359BBFED6B89543C769192B0C67D
24EE581EA3F5AF2CF14FC45D5BC4132D9C44941C47042308A3D947F3AF5327
E09A54DF3D6D5F6E8546925AC8F4BD5FC6DA17869A47D20FF68EA4786BDB81
9C1FFA6687803D2BC5AFB5CD7BC5B7D2792ED3BB72F2B9F957EA7FA0A860D0
E49489F5A90B313F2C087BFA123F90FCEBBDB4B34B1B6081046C47DD5180A3
B015F7F9753A729F2D3574BAF4F444CEB7443F58D7F5ED434DB6B0D52F5EE0
5B42916493F36C50B93DCF4EF5FD2B7FC13B7E3149F177F667D123D4A5F375
7F073B685764B65996D510DBB9CF2775BBC60B1CEE756E7A81FCF57C3AF81F
F157E34DF4963F0E3C3D3EAAB0B059AE3E58AD612467124F2158D4E390BBB7
11D01AFD98FF00827AFC03F8C3FB336BFE30B7F8926C9741F12DB5AC882D6E
1AE1A3BCB5760B95DA36868E56DC4673B573D057958CE34E1DCA2AAC0D7C5D
2A53FE5728A7AF757D2FE614F09566B9A316D1FACB45416D756F7B02DCDAC8
B2C4FD194E47FF00AC77153D7DB52AB1A915383BA7AA6B54D77460D34ECC28
A28AD041451450035955D4A380CAC3041E4106BE04F8DDFF0004FF00F869E3
D9E7F14FC2D9BFE15F78A58B485AD133A75CC846409AD41023248FBF0EDEA4
B2B9C57DFB4572633034711074EBC5493EE07F3FDAEDEFC76FD997578B48F8
C3A2C90D83B048355B62D3E9D393D024E000AC707E4902498E76E306BE9CF8
79F1CF45D785BDE5A5E795729829246DB5D4FB11FAD7EAA6ADA4693AFE9B73
A36BB6506A3A7DE218E7B6B98D6686543D55E3705594F704115F9A1F197FE0
9DB6265B8F167ECDBAB0F0B6AE4EFF00EC8BD776D2E43DC46EAAF2C04FA7CE
9D82A0E6BF9C78EBE8F385C5CBEB595CBD9D45AAB6967E4C699E9F6BFB66E8
DE15F162784FC6D035ED888A32FA8DA80648647C9C4B0AFDE5DBB4929861FD
C6AFB3BC2DE2EF0CF8DF47875EF096A706AB6138CACB038600FF0075875461
DD5806078201AFE5BBC6175F16BE1078B5FC3BF1B742BCD0757B867756B950
62B90A70CF04C84C52A0271BA36651D38E95ED3F0C3E366BFE11D4935EF04E
B32E977876EF3137C92AA9C859633949173D98115F63906799A651429E1734
6EB72A49CFED7ADFAFCF5F33D258484E2B91EA7F497457C05F083F6E5F0E6B
E21D1FE2A5BA68B7ADF28D42DD59AC9FD3CC4CB49113D33F32F725457DCDA2
6BFA178974F4D57C3BA8DBEA964E4859AD6549A324751B909191DC57EA7966
7585C6479B0F34FCBAAF55B9C3568CA0ED246BD14515EA9905145140051451
40051451400514514005145140051451400514514005145140051451400514
514005145140051451400514514005145140051451401FFFD7FDFCA28AF993
F6A6FDA7FC1BFB2F7C3F3E29D7906A3ACEA0CD0695A5A48125BB9C0C924F25
218F8324983B72000599414D81EDDE37F1EF82FE1B787A7F15F8FB5AB4D074
8B7203DCDE4AB147B8E76A2EEE59DB076AAE58F606BF2CBE2B7FC15BBE1E68
93CDA6FC1EF0ADD78A5D72A2FAFE43A7DAE79C3A45B5E6917A70E213F973F8
E3F1A7E3C7C50FDA1FC5AFE2CF893AABDE3A9616B671E52CECA36FF9676F0E
48518032C72ED805D98F35C1E99E18D4750C791116CFB573CAB3E807DDDE35
FF008285FC65F8E661F0578C34CD134CD06E6E164D9636F30995A225A33E6C
D349F37F092028233C0AF57F0978AD5A08FE7EC2BF369FC13AD69C52E844CA
D190C0E0F51CD7B67833C76D128B6B86D9247C107A8AFCABC43C8DE25C6B2D
74B1DD83AD6D0FD1CD3FC5385187ADE1E2C3B71E67EB5F1EE9BE365651FBCF
D6BA34F19A95FF0059FAD7E1F5F871A7B1EA2AA7D1D75E28C8397AE2F54F13
0C1F9FF5AF19B9F18AE0FCFF00AD719A9F8C81046FFD6BA709C3CEFB132AA7
A07883C4E36B7CF5E05E22D65AE9DBE6F96AAEA7AFC9744E5B8AE0B55D5155
1896AFBACAF2854EDA1CD3A86EF85BC0DACFC44D66E6D748BF874F5B245696
59559B890E0050BD4F04F2474AFBEFE1368765F0BFC2FF00D870DE1BDBA9E4
335CDC15D9E64840002AE4ED550000327B9EF5F187C27D622D0F4B9EED982C
FA8C81DBD7627083F527F1AF6B4F1D295FF59FAD72F1152AF59BA09FB8BA5B
F5F53C6AF2BC99F4FDDF8AC6D277D70BAB78A86D6F9FF5AF0CBAF1C2E0FEF3
F5AE2353F1A34B9546CD78983E1D77D8CAC7A1F897C560AB00F5E01AFEA8D3
9791DBD6A3D435A7949691B26BCBFC53E238ACED24959BEE8FCCF6AFB9CAB2
9E56A315A9491B1E2DF829AF6BBE07B5F1FE8E8D2DC4C24636FDE4851D9432
7FB5F2E71DC74E7AFC8F2EE462AE08607041EA08AFD52F821F15740F10F86E
D7C3D72C3ECD0C6B12A9FBD19518FF00F5D64EAFF00FE1EF883C78FE2636FE
7C40E5E2538B79A4CE7CC651C9F7ECDDC1EFE8659C733C154A987CC20ECAFC
B6DFC97FC13F9CF03E3756C16371587CEE93514DB872AD52E917DFD7BDFA6D
F9EBE11F869E36F1DC98F0DE9924F0E70676FDDC20F7F9DB00E3B8193ED5F4
159FECA9ABE85A73789BC49ABC59D3C09DADA0899D5F610769918AF07A1F96
BF45346D02DAD208EDAD2158E38C055540000074000E00AD3F13783E6D67C2
9AB69B811FDA6D268D5D8ED5566421589ED83835F358FF0014F175AB28524A
10BFABB7ABFD123F3BC77D23331C463E9430F154A97346F65CD271BABDDBF2
EC91F3B78335F4B4F22556FB9822BEFF00F067ED617DE1DF0FC7A5ED593CA5
C296E715F9DFE0AF87FA85B68A97BE2A9DACE77504431952538FE33C827D87
E75CFEA3AE7F655ECB64671284E8C38C8FA76A1D28CEACBD8CB55D57F99FDE
77D353EAEF89DF18F53F1C5EC9757B3673D067B57CD3AD6AB25D38B7B75696
5958222282CCCCC7000039249E82B87BDF16AED3F3D7DCFF00B247C211AA34
1F133C4B016925E74E8DC70919E3CE23D5B9DBE8BCF718F0789B37A19260A7
8EC4EB6D97593E8BFCDF446B429BAB2E489D7FECE1FB2341A7EA567F123E23
2B49AAC60B5AD906222B712295265C7FAC72A48DBF71727863823F46EDEDA1
B68D618102AA8C000532CEDD6DE058D0638AD38A22C715FC33C47C4B98F106
37DB6264E4F64BA45764BFA6FADD9F430A50A51E5890842697CB35BD069C58
74A924D3CA8E95F4187F0BB1F2A3ED9C3439DE3217B1CD9522A3650E0AB74A
D39E0D86A830C57C0E6995CF0D37099D509DF5479FF893448678DB2A0835F0
1FC75F809A578A639B52D3156C75550489547CB211DA403AFA6EEA3DC715FA
4BAA28684E6BC33C576A92238C57D6F0371362B01888D6C3CDC5AFEB5EE878
9A31A90B491FCFD78C34DD47C3B7F73A46AF09B7BAB738753FA107B823A1AE
5EF9FCBB4443D5100FC857E8C7ED17F0BAD7C5360D7D6B184D46D3E68DC705
D41C98D8FA1EDE87F1CFE723F886DBC2DAB5AEBBAA5AB5D41A4DC453CD6FB8
A3482291498B70E5779F9777F0E73DABFD18E02E2B866D82F6F08FEF16F1F3
F2F27D0F8CC5E19D39D9EC7EC97ED93FB638FD9B3C17A6FC33F05B4573F116
FEC2252D95923D2A2D817CE91790D2B107C94231FC6D95015FF9E3D4752D4B
5BD4AEB59D66EA5BEBFBE95E79EE277324B2CB21DCCEEEC4966624924F26B5
3C67E33F127C46F16EABE39F185E35FEB3AD4ED717333F766E800FE155002A
A8E154003802B2ECADCCF2050339AEFF000BBC33C270DE07D941275E7AD49F
79765FDD5B25F37AB6638DC64AB4EEF65B166CB4F9AE9C2A2E735E99A2FC3A
D475001962273ED5E9FF000A7E1D1D6AE620C99C91DABF557E15FECE36F756
714925B8CE076AFD38E53F2126F84FA8C316F6848FC2BCFB58F095C5812190
8C57EFAF8DBE0158E9B60EDE40040F4AFCD8F8BBE09B6D3659822018CF6A00
FCF6B8B76898835551FCB707D6BB3F10DA2C33B803A1AE1E707071D45541D9
899D1DBDC7039AD68A7AE2ADAEB8EB5B315CFBD74D883A859FDEA613FBD73C
973EF530B8C52B01B867F7AD2D1B48D5FC477E9A5E896AF7773274541D07AB
13C28F738158FA269D7FE20D52DB47D353CCB8B960AA3B0F527D80E4D7E95F
C2AF877A4F8434B8AD2CE30F7120066988F9E47F7F61D876FCCD7C67187175
3CAE968B9AA3D97EAFCBF32E10B9E35E0FFD97C4F1C773E2DBF6673C982DBE
551EC5D864FE007D6BDA20FD9DBE1B410843A3AC98EEF248C4FE25ABED7F87
7F09356F16153044447DD8F4AF6DD4FF00672BBB3B069A365919464815FCE9
98F1C66D8993A8EAC92FEEDD25F77EA74469F647E4BEBBFB387C3BBB85961D
3DAD1F1C3C32B823F02597F306BE53F883FB3AEBFA02C979E1B9CEA76E9926
270166007A1FBAFF00A1F406BF5AFC59E1C97469E48255C1524578AEB96C8E
AC08AF5787F8FF0032A1357AAE4BB4B5FCF55F225C133F1E744B296E3C4765
613214713AEF56182021DCC083D0E01AFD26F83BF087C63F126FE15B154D3F
4F62375DDD3AC310527AEE72063DFF002C9E293C2FE09F8691FC56B2F1278D
34C92EA3314D18587003DC14FDC9914E0152C0231EA1589C1200AFD42D0BC6
5F0C3E1F7826DAE751B65BFD5F51664B7B48D844A8A8066476C1C282400319
3D3D48FD4B31E3DAF5A946182872B92D64EDA77496BF7B5F27D3CCC5E27D93
515AB67289AF7C10FD94F4A586C6787C5BE2E640C6E090D6B6ED8E0AB7F11F
A7FE3B5F9EDF16BF68FF001EFC61D7E68AD6E9EE5C920BB1DB040A7B01D147
B0193EF5F4AFC6FBEF07F8E7429175DB282342DBE14136D76DA72555C61C67
A123B1AE5ECBC21FB36E91A7C0BA769F75A809115FCA8669494C8CE1983A2E
ECF072C4E7AD7C865793D2AF59627194E73BBBD92BDDAEADB7ABF37E8ACB43
18669CD15783F91F3278634DD1BC332FF6B5D3FF00696B07937328E2327FE7
929C85FAFDEF7C715F2C6ABE238753F1F6A7AD4F1A8492F5C90005C846DA09
03B90327D4E4D7EAF69717ECEED32413E84D6527441785C863D31BBCC75CFD
4F3F5AD98FC21FB37E9E6478FC1FA33B48C599A5B58A5249EA72E18D7DE63B
8AA9508AA11C34E36E8D25FAB17F6BAA6F5833CCE6FDB67C3361F0453E1BE9
7A44293795B1A4DA3278AFCBCF106AE350BC9EECF1E6B16C7D6BF6CF45F137
C1BF0D6A967A869DA06956EF6AE3698ED21042F420617BAE47D2BECF87E2C7
C2BD0D1934F9ADAD558E5842891827D4EDC57E47C71E3457C9AA528D3CBA55
54D3B3E7B6AAD75F04BBA3E8F246B1B193F86DFA9FCAFDAF823C73AEDE7D8F
4CD175099641BF6C369348CC3B950AA73C57A6E81F01FE34DDC657C3BF0E3C
4370B9DAF2AE9576C32304867F2F19E41C7E95FD215CFED17F0DEDF39BE538
FF0068573D75FB547C38B7CEDB80D8FF006ABF3CC47D22F3DAF654B29B2ECE
4DFF00EDA8F61E450EB50FC3FD13F648FDA11E4563E03D564B971906480C28
BF4321500FD706BD33C29FB0E7ED11E22D72C2DF5AF0ACBA4E932DC2477773
25CDA2B430E46F758CCBBD885E985209E3D6BF542EFF006C1F01419F2C86C7
B9AE62F3F6D7F08420F9518344FC77E369D274B0B96D385D5AF69B6BCD7BE9
5FD53F98E3936193D67FD7DC7D8DF0CFE1FF00863C0DE1AD37C19E1AB38F4D
D334E896348D0019207CCCC7F89D8F2CC7927935DAEA96B6D6D26C81B70AF9
47C01FB46695E2BD386AB811472B30419C6429C67F315E849F11CEBD749A5E
811FDA6EE6070AA7A01D49278007A9AFE72C73A91C0D5C162B077C4CA6E52A
B26DBF3FC773D98D2BCD4D4BDDEC686A5F10E2F87BE33D36E269B6E99A9CD1
DB5EA7F08F348449BD8A310491D5723AE31F5757E52FED0DE0EF89B343A6AC
5189A3D5AF6DECD64B46677867B891638B208420B3B00A41C67032322BF56A
BFB83E8B78BAD2C867879D4E78425EEEB7B5F78F959EB6F3F33E6788231F6A
9A5BA0A28A2BFA60F0428A28A0028A28A0028A28A00E2FC7DF0EBC0DF14BC3
773E11F887A1DAEBFA45D021EDEEA30E01C603A37DE8DD7F85D08653CA9079
AFC77F8EFF00F04BEF12F8624BAF167ECC7AB35F5B2E643E1ED4A50261CFDD
B5BB7215FD927DA401CCAC4E2BF6F28AE5C560A9568F2D48DCA8CDC7547F36
7F02FE12FC65F885AEEBBA0EA7A3B787CF84D247D51B552D6C21F273B9000A
CECCD8F94AA95EF9C735F5EFEC57E37F14CDFB4B6A1E15D22231E8B77A55CD
C6A10A331823585D041201FDF576D8B9E76BB57A8685E2409F06FE2F7C5E2F
F3F8CF569E1B56F581E43B003E8237C7E15A3FF04DEF0D83A07C40F8872A0D
FAB6AB0E9B1B1CEEF2B4F87CC24678DA5EE48E3A95E7A0AFC678670D0AD9CA
9D1568C799E9DAED2FBD2FC4F671326A8D99FA5D451457EE07881451450014
51450014514500145145001451450014514500145145001451450014514500
1451450014514500145145001451450014514500145145007FFFD0FDE2F11F
88747F09787F53F14F886E92CB4BD1EDA6BBBA9DCE16282042F239F60A09AF
E473F694F8F5E21FDA3FE2F6AFF11359774B17736FA55A31CADA69F1B1F263
03A6E20EF908FBD2331181803F733FE0AA9F14AE3C0BFB36AF83F4E98C579E
3AD421B07DADB5859C00DC4E4739218A471B762AE41EB83FCDA5B1195AE7AD
2E807AB7807C2AFAFEA314017218815FB1DF00FF00653B1D634E8AEAE60072
01E457E5D7C17D4AD34FD5219662000C2BF6D3E137C7BD17C3FA1C7179AAA5
5477F6AE703CDFE33FECF3A2786F4C91A38954AA9ED5F8F3F13B467D07587B
BD2DB64919C328E8C07F51FE7B57EB27C7EFDA22CB59B69A28660720F435F9
01F10FC4C355BE9640D9C9359D5A519C5C64B41A76D515747F881708A23949
0C2BB783C752BAF0C7F3AF9B2494872C3AD598359920E37FE75F278CE1C8DE
F189D30AC7D267C533CC397C0FAD56935ADDCB364FBD78EE877FA8EBDAA5AE
89A447F69BDBD711C3187505DDBA0CB1007D4902BED69BF626F8D363E14B7F
166B97DA5D8C37519923804D24D30C750FB23F2C1FA3B57CAE653C2609C638
9A8A0DED7DDFCB73A21CD2F851F39DE6B8AA0FCD5E4BE28F19C71B8B584EF7
6FBD83D07F8D7ED37ECB9FF04D4D03C4BA3E89F13BE35EBC358B1D4228EEED
F45D3F7C50B238CAADD5C30591BFDA8E3098231E630EBF889F1F7C092FC2BF
8E5E3BF8792AB22683ACDEDB43B860B5B2CAC607C71C3C451871D0D7DCE559
04525527B74396A55E88EA743F1BDDC7126C977274CFA7B1AF41B4F1ADCC8A
3F795F487ECE7F03FC25FF000A2E4D5FC79A545793F89DFED685C6D9628106
D83648B865246E7CA9190F83C715F24FC53F0FE85E05D41CE897734903CA51
22930CE0753F30C0C0FA7E35F31471D84C562EAE16945DE2ED7E8EDBFE2725
D33B81E24925197933F8D324D7940FBD5E23A55EEAFAA0074FB69645FEF150
17FEFA240AEBA2F0E6B570A0DC5C2C39EC0163FD3F9D7B11E1F9B7B0729B5A
AF8A61823677930057886BFE209B589FA910AF41EA7D4D7A9C9F0F2DAE9B75
EDDCD2E3B02AA3F2C1FE7427C2FD22E244811E65690850770E33F857BF81CA
E9E1D7B496EBF01549C6117393B25A9E85FB34F81EF2EAEA6F16DEB3A5A731
41183C48C3EF3B7B2F41EF9F4AFD11D12C81D8882BC5BC09A5D9E8BA4DA699
649B21B68D6351ECA31CFBFAFBD7B9E8DA9C1A711294F3A41C840719FA9EC2
BF9DB8B332A98EC54AA25A745E47F0171AE2B1DC499CCA384A6E529BB462BB
2DAFF2D5B7A2F43D8AC34FD2F42D37FB575A904512F4CF2CCD8C8551DD8E3F
C8AF24F1978D85F6E6980B5B18CE638339271D19CF73DF1D076F53CB78BFC6
E632D7DAADC0925452147448C7A28EDFCCF7AF99E0BAF1F7C6DF17C5E03F86
964FA85F5C1CBB83B61822CE1A59A4E891AE7927AF450588076C8B872559A5
15EACFEA3F0C3C16C070ED28E3F30B4F12B5BBF861FE1BF5FEF7DD6EAEF883
F15669663A5E8AAF7373265523894BB13E8157249F615F218F1E5D5DDDCD70
F212643C64FA57EEB7C26FD99BC1BF01B4AFED1BC9535DF18DD478BAD45D7E
58B70F9A2B556E523E705BEFBF56C0C2AFC0DFB6AF837C0C9A9695ABD86970
59EAB7EF706E26B74F2DE61184E5F6E03105BA9193EBC57DDE031985C362A3
8050BB7D5775AEDD8F4B0DE3061F159FC325A349B52BDA77EA936EEBB596F7
BF976F9EBE0EE8D75F147E22693E16258DACB279B7454FDDB78F97E7B6EE14
1F522BFA20F87DA7DAE9B616D696B188A285151154602AA8C0007A015F8C7F
B0C7842D2CB56D6FC54773B9F2ECE266E7681FBC9003EFF27E55FB35E16BB0
B1A735FCADF48CCDE55F1EB094DFB94D7E2F56FF0025F23FA0725A768733DD
9EDF08040AE86C21048CD72F613ACB12906BA6B498262BF1CF0E1D0863232A
DD19D58D4ECD23B4B6863082A1BC58C29C5508EFC2AF5AA5737DB875AFEDCC
DB8EF2BA596FB28257B1F354F0B373B99979B726B164AB93CDB89AC5BBBA58
94F3CD7F077196650C46264E99F5186A6D2B197AB4E1508CD78CF88E604357
75ABEA2086E6BC7BC417E36B735C991E0DF323AAA3B2B1E0DE3C44962901E7
20D7E33FED3DA1CFA3F88E19AD976D8EA0ED238038F3D0639F4E0923D4E73D
ABF617C6378192419AFCE2FDA474B4D63C377B201992CF170A7AE361F9BFF1
D26BFAFBC16CC2585C6C2FF0CB47F3FF0027A9F399A439A0CFCFC8EBB1F0E4
4AF749BBD6B8C435D368D75E44CAD9EF5FD98CF974CFD2EF800B6304F6ED2E
3A8AFD90F86FE28D16C74C8C16518515FCF67C3CF1EFF6518C89318C77AFAD
748F8FD2DA5A041718C0F5A459FA6DF167E216946C654475E86BF1FF00E35F
8A2DAEE79F630E49AB3E35F8E93EA113AF9E4E73DEBE3CF1878CE4D4A4725F
39F7A00F3FF11DD2CB70E41EF5C3C879357EF6ECCCE4E6B299A825B284A5A0
9323EE9AB905E640E6A1900652A7A5775E33F835F133E1EF85BC35E37F1568
57163A0F8BAD96EB4DBC2BBA1951CB6C52EB90AEC8BE62A310C50860306954
CC28D39429D59A529BB453693934AF64BABB2BE9D0145B4DA4738975EF5605
D7BD72693B0E09AB4B3B7AD76927D97FB3AE8913C973E249D72EE7C88B3D94
60B11F5381F857E87F828452DE408E78C8AF84FE0A4F1DA78574E8D4F542C7
EAEC58FF003AFAC3C3DAE7D96449036318AFE64E3B94EBE36A4DF7B2F45A1D
30D8FDB7F83B6FA65B785ADBECDB7247CD8EB9AF4AD7751B1B2B292E2760A8
8A49CD7E57F827E3C6A5E1EB658619CA803A66A7F177ED03AC6B90342F7076
91D01AF0F0B9E55A181960614D7BDD7A9B296CC6FC61D72CAF357B87B7236E
E38C57CB7ABDD021B9AD2D6BC4B25ECAD248F926BCE354D5060F35E3E0B04E
E66D9C578A6EA6850DC599C4F11DE9FEF29C81F8D79DF893E325FEB2B617F6
121179A7AB433DB39DAE467395CF1907391DF3EDCED789353051F9AF9D3554
B6B8BF7925404B77E878F71CD7EC9C21838CDAA535E9F71C38AC3A9B52EA7A
8D87C4CBAF14F89C5ADDC8CADB54047C82A8A39001F7CFE35FAC5FB2CDFF00
C07D3FC3FA8DFF00C48F2A4BC500429282576E39DA07F11F7AFC74F02784B5
2F16789ACB48F0EC39BC91D4070A5D82B1C11EA739E067938AFAB7E367C31F
1BFECED7963A5F8AAFED2F9AEEDA39E4FB1BB335B171FEAE75651B5BDC165E
9CF38AFD6A9E1D538C60BA10A164923AFF008C5AE7862FB5FD4A4F0E4622D3
9E46F297D173C57C697DE23D79EF2782DB51658A26C00C0B71F5DC2AB6B7F1
04DCC6CB1BF5EF52D87C3BF894F6675293C2BAB7957004CAE6C6E3698CAE55
B3B31823907A60E6AFEA50A9A4E3717245FC47257DAEFC458AE03D95C5B4E8
AC08DECEA78F51C8FD6B7AD7C6BF112562FAC6A512AF6481589FFBED8FFECB
5977B15EE9F2FD9EFE092DA5C676C8A51B07BE0806B35EE3DEB96BF0FE0AAB
4EAD24EDDD1D34A4E0AD0D0EA26F156BCF9DDA84A7F103FA5644FAFEACFF00
7AF663FF00033FE358525C7BD67CB71EF574F22C147E1A31FF00C057F916EB
4FB9A571ABDFB6775D4A7FE06DFE358371A95D67779CC48E796355A7B9F7AC
A9652C735D71C1518E8A0BEE4473BEE7D09E0AF8EBE20D0208EC5AE311C602
81D303DABED6FD9CBF6A6BDD0FC76A2F55EF2CEF223F6828ACE618A3F99A46
C038419CB31E0704D7E4ACAFC1AFDB2FF82397C345BC9BE23FC55D4AD924B7
68EDF40B66650DBC3FFA4DE2107B605BE477CF3D39FCCB893C20CAB3285484
E3CAA6B5B2FC57CCF42866952163F58B48D5B4DF13E9306A9A6CC2586750EA
C8DCFA8208E411D88AF4DF0E6BAD78BFD9D7EC3ED918F94FFCF64503E6E800
607EF01F5180703E31D6B43BEFD98BC591359EF9BE1A6BB3AC76E599A46D2A
E5FA5BC8C73FB963FEA5C9FF00A66DCED2FF004741710EA16D16A1A74C4670
F1C919E41EC47F51D0F43C57F1C64F9AE6DE1CF104A8E21395197C4BA54874
9C7A73C7F1D53F2FA1AB4A9E2E97347FE19F63DB28AE7F42D723D562314DB6
3BD887EF2307A8CE03AE79DA7F1C1E327A9E82BFD09C933BC2E6584A78EC14
D4E9CD5D35DBF46B669EA9E8F53E42AD29424E325AA0A28A2BD5330A28A280
0A28A2800AF29F8E7E263E0FF83DE2FD7D1B64D069B3C70123205C4EBE4C39
03B798EB9AF56AF913F6C2BC9AFF00C21E18F87B66C4CFE2DD72D2DD903637
410132B123AB01288B8F5C1CF407CACF719F57C1D5ADDA2FEFB69F89AD0873
4D44F907E3A4C9F0FBF65EF0178317F772DF44FA84CBD0FCC372E7FEFBFD2B
EDCFD8DBC22DE0DFD9B3C116732325C6A568DAACBBFEF6ED4A46BA5C81D36A
48AB8EA00E79CD7E7BFEDC73CFE28F895A37C2DD05F6C9B6C345B60AA5B135
DBAA2E14724E5C0C0EB5FB2DA569965A2E9967A3E9B1086D2C218E08635E02
471284451EC0002BF3FF000DF056957ACFA5A3F72B3FC8EFCC27A457CCBF45
1457EAA7961451450014514500145145001451450014514500145145001451
45001451450014514500145145001451450014514500145145001451450014
5145007FFFD1D5FF0082CE6AF747C43F0A744240B686DB56B800672D24AF6C
877738380836F19196E79E3F17A097A57EE17FC1687C253BE8FF000BBC7F04
47C9B4B9D4B4CB8939C6EB94867817AE0710CC781CF73C0AFC29827E3AD73D
45A81E91A1EBF2E9CEAC8D8C57AE5A7C5BD4EDADC4493B0007AD7CD693FBD4
E2E5BB1AC7940F5ED7BE206A1A9EEF325273EF5E67777CF3B16639ACA69F3D
EA079A8E5027925ACD9A5A4926ACC9A6AD12036FC33AEB787BC57A46BAADB7
EC17704C4FFB28E09FCC66BFA87D2EE13C59F042DA70779B638FF80C8B9AFE
55ACEC351D5EE96CB4BB596F6E5F3B62851A4738F45504D7F4B5FB1FEA5AA6
B7F049741D7EDA6B3D4A2D3E21341711B452C7342A0306460181FA815F81F8
EB858FB2C3E213F7A2DAF3B743D6CADBBB47D6DFB1FEB0F7DF085B4494E5FC
39AA5F5867183867174A0FAE167038EC077C93F94DFF00053EFD96352F107E
D1FE07F897A05A15D2BC7C22D37569D067CABCD3D09F31CF18325920083BF9
27D457DF5FB3278C344F06F8F7C6DE13F105F41A745A90B4BEB569E45895A4
4DF14C3736067062C0CD5FFDABBE2259EA33699E12D26E63B8B4B44FB75C3C
4E1D1A5705615CA9FE14DCC7D43A9ED5F5D538DA9E1B869632335ED141452B
EBCCFDD5A7E3E8674F08E75F95AD2E7C05F12FC49A6F87BC2F1E93A6EDB78E
DE258638C7015546D007B015F9DDAA69567E21D57FB57508C5C18F2220FCA8
04F2DB7A127D7E95EE5F18FC4526A9AA7F67C0F9C9DBF9F53F80AF3282DD55
42A8C01C015C3E14E5B3587FACD55ABFEBFAF530C661E14E76832A436688A1
5570055D5B718E95AB059BBF415D6687E0CD5F5FB816BA6C0647EE7A2A8F52
6BF5AAB56308B94DD92390E07C8F6AB5A7C6A97F0B1FEF57BD5C7C08F12DBD
AF9EF343BF19DA377F3C7F4ADAF87DF0CB45B7B87D5BC4D3A4D7166D91699E
17D19C1FBC0F6EDF53C0F9ACDB3DC34F0B51539DDB56D3CCF333AC0D4C4E12
AE1E9BB3926AFEBA7E4745E07F0CDEEA16E97D760DBDA60104F05C7B7B7BD6
B78DFC59E1FF000E69AF0C3B2264070C3EF67DFD6B9BF88FF16B4FD06D5EDE
D5C02A080A0D7CDDA568FAC7C49BF1ABF88E57B5D1C9DCA809124E0F4DBFDD
53FDEEA7B75DC3F35C3E554695375B11A47F17E48F8ECBF2DCA384F072C449
D9BDE5F6A4FB25F925A2DDF566D784FC35E33F8FFE227B1D3271A66856F205
BBBF94128838CA46BC799260E42E401C6E61919FD57F861A1780BE0B785078
67C0F6C96C8F86B9B9721AE6EE419FDE4F260163C9DA3855070A00E2BE37D0
356B6D06C20D33498D6D2D6DD76C71C636AA8EBC01EA793EA7935BF2F8C6E9
D76990FE75F398CE21C437C987F721D175F99FCC5C75E2BE6B9B5671A6FD9D
15B457E727D5FE0BA777F4978B3E2046EAE2393731CF35F9ABFB53EAE352B8
D0D49CB46B74C7E8E6203FF41AF78BBD7A497259F39AF933E3DDF7DA6FF4F8
FBC70B9CFF00BCDFFD6A385A94A59842A4DDDEBF930F04F0739F1350AD5356
94DFFE4925FA9F45FEC8D2C765E100C0E0CF732B9FA8C2FF0025AFD27F0D6A
C0227CD5F935FB35EB896DE1F5B5DD8314F203CFA90DFD6BF41FC35E200513
E6AFC2BC58CA653CCABCDADE4CFF0045B2FA96823EC6D175B54DA09E2BD02D
B508655055EBE58D335F181F3576D69E23200F9EBF03AB9755A33E68687AAD
A92D4FA07ED7C7DEAAD35EC6BCB38AF1E5F13BE3EFFEB55E6F11161CBFEB4A
AE231935CB29684AA513D2EF3598901086B8AD4759CE7E6AE36EB5ECE7E6AE
4AFF005D183F354E1727949DD96E696C6E6ADAC0C37CD5E45E20D5C61BE6A3
56D7861BE6AF22D7F5E186F9ABF40C9725775A1C956A9CBF8BB561B1FE6F5A
F8F7E2315BFD2B528DFA4B04ABEBC1522BDABC4FAC998B46ADD6BC23C52977
A85B4B61A742D7377740C50C4832F248FF002AA28EE5890057F40F08601D29
C1F5BA3C7C4CEE7E712B55D86628722AB5DDA5E69D77369FA84125ADD5B3B4
72C52A9492375386565600820F0411915186AFEC8DCF973B3B1D725B7C6D6A
E897C5F72AB8121FCEBCB964C77A93CD6A9B05CEE2F3C49713820B935CC5C5
EBCA4E4E6B30C84D30B51CA04CCF9A84B530B575FE03F87DE37F8A3E26B5F0
6FC3DD16E75FD6AF33E5DB5AA6F6DA3AB31E151173F33B10ABDC8A76B6AC0F
56FD967E095DFED05F1B7C3FF0F8237F66349F6BD5245C8F2B4FB721A63918
20BE444A7B3BAD7F55BE28F01F83BC6FE13B9F01F8AB46B6D4FC3F7912C125
8C918F24C698D8154636EC20142B82A402A4102BE5AFF827EFEC43ACFECDFE
11D575AF88EF6CDE30F133C7E7ADAB9956D6D22198EDFCC2002FBD99A429F2
93B402C1431FD2AB5D32CACC6208829F5EA7F33CD7F337897E13677C599D53
ACF10B0F86A1A536AEE6E5A394D24D5B54927CD7F753B6A7B782C7D3C3D36A
D793DCFE783F680FF824678B62B9BBF12FECDF73FDA164E4BAE87A948219E3
CF3B2DEEE42239067851314207591CF35F929F103E157C4DF84DAB3689F127
C31A8F86EF14E025F5BBC21FBE63761B5C1EC54907B1AFEE67005676ADA3E9
1AF584DA56B96306A3657036C905CC4B344E0F66470548FA8AFE83E19CA311
81C2430D89C4CABC97DA9A8F33FF00C052BFABBB7D5B3CAAD514A578C6C7F1
D7F097C44BFD8569116C18C1423D369C7F2C57D33A5EBC368C357DA1FF0005
26FD96BC2DF0FF00C35A4FC6DF83FE18B2F0FE9DA7CBF64D76DB4BB68ED600
B3B2ADBDD18A20A8B893F74E42E58BA13D09AFCAFD1BC58AE8BF3D7C0714F0
FBF6D2925A3D4519687D616FE20C0E1EA77D7B2397FD6BE7FB7F13A91F7EAD
B789463EFD7C2CB2377D8D398F5CBAD7060FCD5C5EA9AE8DA7E6AE02EFC4A3
07E7AE4AFB5C79F203715E960F2477BB427234F5CD5CCC594371DEBEF9FF00
8275F803C2FAF3F8FBC7FE35D32CEEB4FD112D2382E2F608E5481956696E1D
5E407614408588C7079E95F993737724CEB0C0AD24B210A88A373331380001
C924F415FB6DF053E13EA5A27803C3FF00B345BCA2D266275FF1CDDA30DB6E
F36D71645C71BD2358D24C120955C12A5857EA7C1F85F635BDBDB48AFC5E8B
FAF238719794791753D4B46F1068F1A6A3FB49F8B2D16D7C3FA1092CFC27A6
BA8412484F3394C0019986E381C703276027F386DF44F1B7ED81F1D07832DE
E5F1A948F77ABDF637AD9D9291E63E3A64E4246BD0BB28381923D37F6BAF8E
51F8BB5BB4F03780E16FEC4D24AE9DA4D9C00933392114AA8E59A46C01DCF1
DEBF4EFF0063DFD9BE2F803F0C63B7D623593C5FE21D979ACCE30C56523E4B
656FEE40095F42E59870C2BEC60DD5A8E736613F7216823C23E13FFC1353E0
17C33D6EDFC49AAB5FF8C6FED1C490A6AB2466D23753F2B7D9E28D15C8F490
BAE79DA0E31F7A49A78F4AED0DB7B55592D47615EDD0C4C61A4558F22B46A3
D6679C6A3A2DA5F40F6D79024F0BE37248A194E0E4641C8382335E15E2CFD9
DBE0B78A04875BF046913C9272D2AD9C514C7AFF00CB58D55FB9FE2EBCF5AF
A9AE2DB1918AE7AEEDC60F15ECE1F12A5F12B9CF76B667E63F8F3FE09F1F04
35A8A47F0E2DFF0086A7272A6DEE1A7881C63E64B9F3188EF80EBCF7C715F0
E7C49FF827F7C4EF0D092E7C13AA5AF89A05E444E3EC5727D82BB3447D3265
5FA7A7EF85F5A820F15C16AB60AEAD915EBC329C262159C2CFBAD3FE07E06D
0C7558F5B9FCAEF8A3C2FE26F076A6FA3F8AF4CB8D2AF63CE62B88DA3247AA
E7861E84641EC6B95792BFA41F89BF0D3C2FE37D2E5D27C4FA643A8DB36484
95031538C6E46EAADE8CA411EB5F8BDFB447ECDBA9FC2A9A5F10F870C97DE1
C66F9B77CD2DA67A0723EF213C06EDC06E704FCCE6BC3D530E9CE0EF1FC51E
BE1B1D1A9A3D19F28B99257586152F2390AAAA32493C0000EA4D7F5F7FB1D7
C167F809FB3BF84BC017F00835916E6F75400E4FDBEF0F9B2AB11C13164440
8E3083EB5F867FF04C3FD98A5F8C7F1697E2D789AD8B784FE1FDC45326E1F2
5DEACB8920886461841C4D20EDFBB0461EBFA69AF9D3BCC8D7F41D1FC53A25
F7873C416AB7BA6EA50BC1710BE76BC720C30C8208E3A1041079041E6BE1ED
324D73F67CF19C3F0EFC577125EF85756763A1EA521DC768E4DACED8004D18
FC245F9D790CABF7C571FE3DF02F877E24785AF7C21E2885A5B1BC0A7746DB
258A4421A396271F75D180653C8EC410483F9D7897E1CE1388F2F785ADA545
AC25D62FFC9F53BB038D9509DD6DD51C64724AAD1EA5A6C855E3C3865E411D
4AB0EEADD08FC4104023D1BC3FE20B1F10D99B8B5216588EC9A227E689FD0F
00E0F5538191CD7C11E0EF1AF8C7E0F78F9FE137C539BED0C101D3EF82B241
7F6A381247BBA3A93891327637192A558FD4923DCE9975178ABC387CD0C3F7
B08384B88F9F95BAE08CE55BA83EC483FCCBE186698DE0BC64F2DCD1BF6537
EF2E907FCF1F27F692F55B1F518DCBA188A71A907ABD9FFEDAFCFB1EE14566
68FABD86B9611EA3A749BE293820F0C8C3AAB0ECC0F04569D7F6DD1AD0A905
529BBC5EA9AD9A7D4F8B9C1C5B8C959A0A28A2B42428A28A002BE35F88520F
187ED65E0AF0E3624B3F06E953EA92F4C24D74C460FB810C47DB771DEBECAA
FCF4D0FC451CFE2DF8F5F188C81A3B12DA45A38395FF0046416FF21F4668C3
6475CE7BD7C3F1F633D9E0943F9A4AFE91F79FE476E06179DFFAEC7CCFE035
6F8B9FB72E8F7720F3AD34FBEBDD624EFB52CD1BC83F84C62AFDA3AFC9AFF8
279787DB58F899F10FE21CEB91A75A5AE950BFA9BA91AE2700FB793113F515
FACB5B701E0DD1CB60DEF2BB163E77A8C28A28AFB238C28A28A0028A28A002
8A28A0028A28A0028A28A0028A28A0028A28A0028A28A0028A28A0028A28A0
028A28A0028A28A0028A28A0028A28A0028A28A00FFFD2FD38FDBCBE09DDFC
78FD98FC57E13D16DBED5AF69A89AB69680659AEEC49731A0FEFCB099225E9
CB8E6BF909B7BA181CD7F7755FCC37FC14E3F635D47E0A78FEEBE377802C1A
4F0178B6E1A6BB481309A4EA33365E360A30B04EC7742780AC5A2C2E23DF32
8DC0FCD04B9CF7A9C5C67BD72B1DE0F5AB22EBDEB2E503A0371D89A81EE3DE
B1FED24D2C5E7DCCC96F6E8D24B2B054450599998E000072493D052B582C76
FE1EF0E5EF896E62B3B14927B9B8711C30C4BBA4918F002819E735FA39F047
F60082F3C9D7BE2E5C3CA18864D32DDF6A81E93CABC93EAB1918FEF1E95EC1
FB22FECE36DF0DF418BC43E258566F12EA2A1DCB007EC91B01FB9439233FDF
61D4F1D00AFD25F0E684F74C91449D6BF8F7C58F1D310AACF03934F962B473
5BBFF0F65E7BBE9647D6E579246CA7596BD8F31F05FC25F0778274E5D33C2B
A35B6976C3194B78963DC477620658FB9C9AD9F01C63C37F192F34B71B21D5
EDE29875DA77A9848F4CE63C903D73DEBEA0D33C1D04718328DCD5E25F18F4
78FC2DE30F0678AED63F295E59ACA57C1F98B0596204F4E024847AE4FA57E1
1C372C756C4CAB57526A717793BBF3DFD51E9E3674B93963D0F9CBC6BA5AE9
3F18B4BF313E4BEFB4D939C740C8655FCDE351F8D713F123C3DADE97A7DDCD
A45AC97AAE59B6C2A59C67B6D1C9C741806BD93F694823D2754B3F15A8FDDD
85CDADF92067291BABB0FC4022BA2D620501B35F419BE3A542A61F1095D34D
35DDC5DFF268CF071E68CA27E4537C3AF186AF15E78A2EACE4B7F2E42890CC
8D1C8E3196750C071D00F5C1ACAD33479A794465086070411820D7E90789EC
63911C115F34DF69FA6E93E218AEAF408ED6670B2BF64CF463EDEBEDCD7F46
7867E292A92583C5C5460FE16BA793EEBCFA7A6DE2665965973C5EA73DE1CF
87F2DD2091D308A3249EC07535F4FF00ECD565E0EBBF11D845E206486C646D
EFB8E01F404FF3AFA3FE1A7C19B7D5EC7ECF2C6425C4650B2E32030C641E46
79E2BF2A13C51ABFC3AF126ABE05F116EB4D5BC3F752D95C2302A77C2C5770
07F8580DCA7A15208E08AFD8F8AA854A94A2A3B5F53C089FAFBFB46EAFF0BE
CF41B7B1F0EADB7DB41EB6F8C04C7722BF263E27EBCFA7ABEA3A749E5DC440
E083D41EAA7D8D54D5BE2AACD11692E371C7735F3FF8A7C5175E23BA4D3ECF
324972E238D7BB331C0FD6BE63098094EB3A8E295FA2565F714D95347FF89E
DF8D635F7F3C96DD1C2C72A3DDBB1F61D07BF6F6EB2D60607CD5F344CD7DE1
2D7AF3C37A99C4F63214247018755619ECCA411EC6BB7D3F5F0403BABCECF3
0152737CFD3F03F96F8E3058AC5E2653C4B6EDA25D12EC91EFF16AE31F7AAC
FF006B8C7DEAF1C875C047DEAB9FDB83FBD5F312CAF5D8FCD2A70FB4F63D3A
5D5863EF57CC5F197C4702F882DED2E32145A2B871CF2D238C1FCABD1E5D6C
7F7ABE69F8AF7A6F7C40AE4E425BA28FA6E63FD6BE978572B5F5A5CCBA33F4
CF0A329747358D5B6D197F91EBDF053C636B697F71A6C32E73B650391FECB7
5FC2BF40BC2DE2A568D087FD6BF1EFC0F75259F88EDE78C90551FA77E2BED5
F0778E06C45327EB5F2BE26F0929D77560AF747F59E5999465274AFAAB7DCF
6FC8FD17D2BC4C368F9EBB4B5F128C7DFAF8D347F188655F9FF5AEF6D3C5A0
81F3D7F3A661C2CD37A1F470C41F51A789063EF5364F118C7DEAF9DD3C5631
F7FF005A1FC5431F7FF5AF1BFD5977D8D7DB9EE175E23183F35725A8789383
F37EB5E4975E2A183F3D725A878A860FCF5EBE0B865DF6329D73D1757F120C
37CF5E45AEF888C85951B24D733A8F88A49890ADC572175A867249AFBDCB32
154ECDA392A56B92EA37B9DC49C935D67ECF9E1893C77F1BFC3FA7852D6FA6
BBEA5390376D5B45DE84FB197CB5CFBFAE2BC6F54D4D5558E6BF4DFF00610F
857368BE04D57E2CEB3114BAF1437D9EC43020AD8DBB10CE3A7FAE97F358D4
83835F7985A1ECA9B9B38E72BE87C8DFB737ECBD0789A1B9F891E0AB309E20
B252D75144BCDEC2A3D3BCA83EE9C6580DBCFCB8FC62CE0E0F0457F579F11A
CA3B88255619C835FCEFFED5FF000BE2F87DF10A4D574A87CAD2F5C6799554
6163B8073220F4072187D481C0AFD878233A95487D5AA3DB6FF23CBC552B7B
C8F99B752EE15543D383D7E85638CB1B8526EA837D7D39FB29FECD1E27FDA7
7E2443E15D30BD9E8764565D5B50500FD9E039C2A06E0CB210550738E5882A
A6A64ECAEC0D9FD94FF647F889FB5578B5B4CF0E0FECCF0F69EE8352D5A642
D1401883E5C4A31E6CC579099000E5994119FEA57F67BFD98FE13FECD9E154
F0EFC3AD2120B89557ED9A84A0497B78EBFC534D8C919C908308B93B5464D7
63F07FE13F82BE0E78274FF04780F4B8B4AD2AC1311C51E4924F2CEECD9677
63CB331249EA6BD5AB2A31753DE7B14F40A28A2BB890A28A2803235FD0746F
14E87A87867C45691EA1A5EAB6F2DADD5BCA3747341329491187A32920D7F3
05FB65FEC4FE38FD973C433F8A3C28971ADFC37BD9736D7D8DF2D8B3938B7B
BDBD08E892E0238C747CAD7F529556F6C6CB52B49B4FD46DE3BAB5B8431CB1
4A81E3911860AB2B020823A822B9B13858D58DA434CFE27ECFC4D390016ADA
4D7A59072F8AFE84FE337FC129FF0067CF88B7171AC7802E2EFE1E6AB3966D
B6405CE9DB9BA93692952A01E8B14B1A8191B7A63F2DB56FF8273FC518759B
BD3FC2BE2DD1F52B6B69E48927B91716BE62A1C0708914F8CFA6EE9DEBC687
0B56ACDAA30BDBB1352BC21F13B1F1C1D53232CD9AAB26A65D8451E5DDC855
5519249E800EE6BF45BC25FF0004BEF1C6A32C47C5DE3BB1B08BE5322D8DAC
B74DFED05329800E780C47BEDED5FA37F05BF640FD9FFF0067781BC671597F
68EABA5C4D7126B3ABB2CF2DBAC4A4BBC4A14470E0679440F8E0B1AECA3C11
884EF5972AF3392A6654D7C3AB3E03FD9DBF669D6FE1469961F1CFE2BE939F
145FBAC5E11F0F5C21338BD90FCB797709036F9208648DB90C4160AFB457D5
7F1BBC5F6DFB3CFC2D93E1A595EFDA3C61E251F6DF10DEEECC9BA5F9BCA2DD
79C9CFE27BD7A9DBF8CE38A1D57F6A3F1F45B2D6D84965E12D3A6E081C8131
5FEF363731E707201C2AD7E717847C1FE30FDB1BE3E27844DD4A2CA795AFF5
DBF5E7ECD64ADF3ED2723CC90911C4307E62091B5588EA54A11FDDD2D97E3E
6553527EF4B767D3BFF04F7FD9E6E7C77E27FF008695F1D5B6ED2B4E9248BC
3B04A3226B946D925EE0F056221923FF00A6819B828A4FED2563F87BC3FA2F
84F41D3FC31E1CB38F4FD2B4A823B6B5B788612286250A883BE001DF93DF9A
D8AEB8AB2B1D290546F1AB0F7A928AA4EC29453D198F730641E2B96BC8F19A
EEA75CAE6B91D4942935EAE06ADDD8F031947925A1C4DE20C1AE335088106B
B6BC3D6B91BFEF5F6B97B773819E61AD5A8756E2BE7FF1B787ED352B4B8B4B
B856686656474750CACAC3041078208EB5F49EA8A086AF23F1040195ABE82B
4138EA3A72B1ABFB0E6A7A678164D4FE0B5840965A6857D42C225FBA1F7013
80492496055BFE024F7AFD19AFC5FB9D6759F05788AC7C5FE1C93C8D474B94
4B193F75BB323018CABA92AC3B826BF5B7E1F78E347F88FE0ED33C65A1922D
B518F7146FBF148A4AC9137FB48E0A93D0E323820D7E5F9DE5DEC6A7347E16
7D1E0EBF3C6CF73B3A28A2BC43B0F24F8CDF08742F8C7E117D07527367A8DA
31B8D335041996CAEC290920008DC87A4919203AF19076B2FCA7F07FE256BD
E18D7EF7E127C508BEC7AEE92C2371F3797321E12781982978A4032AD8F504
06040FD07AF9F3E3F7C1087E2C68B06ABA04C9A678CF4205F4BBE6C843920B
DB4FB7930CB8C6704C6D87507E657FCDFC46E02A59CE1AF056AD1F85F7F267
B793E69EC5BA753583DFCBCD1AAB7375E12D43FB6F4E0F71A6DC73736E841D
C3B4880FF1A8FA6E1C1EC47B2D9DE5AEA16B15ED94AB3C13286474390C0FA5
7C25F07BE30EA9ACE9D73E10F116992DB788B4B93ECD776770363DB4ABF7BC
CEB91DD48C8604153B4EEAFA474BD135ED2F4C63677725B4723349B13EE2B3
F2DB41CE3279C7A927A926BF0FE00F10F17917B5CB31D467529C2ED28ABCA1
67AEF6F77AEFA7CCF7736CA635146A39A4DECFF997F9A3D968AF0AD37E2B49
A46BD67E1FF16B23437F20821BC00215998E116551C61DBE50CA0609191825
87BAD7F4B70C714E0F37C2AC5E0A578ECD3D1A7D9AFF00875E67CAE3F2FAB8
69F2555FE4C28A28AFA2384E73C61E22B6F07F84B5BF16DE297B7D12C6E6FA
403A94B689A561F92D7E47EBB7775E08FD8AB4F6BD949D4FC73AA4B7B3C878
7915092CC718EACA0FE35F74FED9DE216D17E02EADA640C05D788EE6D34B84
1CFCC67943C8A31D4986393DBEBD0FE7A7EDD7A80F0AF873C15F0CACB24F87
F4588346A324CF3000E00EE768FCEBE733A9DE4A1E5F9BFF0024CEFC22B2B9
F65FFC13BBC2EDA2FECE96FE239D40B8F16EA77FA8B139DFB124FB1C7BB3D8
ADB8618E30D9EA4D7DD55E77F08BC1BFF0AEFE167847C0ACBB64D074AB3B39
3E60D996185564624704B382491C73C715E895EF50A4A10504B638A72BB6C2
8A28AD490A28A2800A28A2800A28A2800A28A2800A28A2800A28A2800A28A2
800A28A2800A28A2800A28A2800A28A2800A28A2800A28A2800A28A2800A28
A2803FFFD3FDFCAC5F11F87341F17E837FE17F145843AA691AA42F6F756B70
8248A68A41864753C10456D51401FCD6FED79FF04A1F1EFC3EBFBBF1BFECDF
04FE2DF0BC84C8FA303E66AB6231CAC409CDDC79FBBB7F7C320147C192BF1F
AF2DAFB4BBC9B4ED4EDE4B4BBB6731CB0CC86392375382ACAD82083D411C57
F7A75E1BF167F668F809F1CC16F8ABE07D375FB929E58BB92230DEAA630156
EE131CEA07601C01D47349A03F899130AFB43F62EF86F0F8BFC7B278BF548B
CCB2F0F6C30861F2B5DBE769F43E5A82DEC4A9AFDB5F117FC11E7F64AD6EE8
CFA6DCF893404273E558EA30BA0180319BBB7B87C719FBD9C9EB8C01F3A787
7E0EF813E0078BBC45F0D7E1E4F7B73A46997ECA26D45E392EA494222CA5DA
28E143865217083E503393C9FCCBC5BCE6782C9AA2A4ED29FBBF27BFE0ADF3
3D2CAA8A9D657D96A7D47E1989156345AFADFC07A646964B391F3357C6FE1D
BE5468CE6BEBFF00016B96D2D92405C061D2BFCFA7461F5A8CAB7C37D4FB7A
ADFB36A27B7D959C61066BE5BFDB3B4EBF1F02B57D774391A2D43C3F3DADFC
4CABB8809288A63D0E310C921CE38C76EA3E938B510A9C1AF37F8AB6969E2C
F00F88FC277722A2EB5A75D5982DCED33C4C81BB742722BFA7B15C6D9250CB
E9D2A7157D2FFA9F28B0B565367E237C55D5B59F17F8434ED5352BFB8BC69A
DCC6C65959F94E3B9AFB0FE1F788C788FE1978635876DD2CFA75B898FF00D3
68D024BFF8FAB57C57A43B6B1F0BCC52A912D949CAB0C32EE182083C8C1AF5
BFD9D35DCFC3BBAD09CE0E8DA85CC4ABFEC4C45C03F42D237E20D7CB788795
A9E02328AD69CFF069A7F8A46B9256B55717D51EC1AF4C0AB57CE5E358D258
A453CE41AF67D72FC61B9AF06F15DE828FCD7CA70DE1E4A49A3D9C448FD1BF
D803E26C7E2EF06EA9E07D5A50FAC784DD1232C46F96C66CF92D8EA7CB2AD1
9F40132726B53F6B9FD84FC27FB49DCAF8D740D50F853C756F12C3F6D0865B
6BC8A30764775102A72BD16553B94704380A07E68FECA3F1266F877FB46695
72B2F9769AE433E9B71CE01122F991707A9F363403BF27E87F6AEF3E305A25
B92AEA1B15FDB9C258EFAC6021ED356B47FD7A1F198BA7CB3763F01BE277EC
11FB49FC29D0B53F126BB3E91AA69DA4A79B2BD85DCB2C8615237C8B1490C6
C42292CC0E0E01C678CD3F86BE09F0C59693F6A8BFD36FEE130F71201901BF
8507F08FD4F73DABF5C3E227C584BB8E5432860C0820F2315F92FA88B7F007
8EA6D3B4F1B343D49DA5B303EEC249CBC1F453CA7FB381C95268CFB032F65C
D4B4B6E608F28F8D5E10BCD502EA96AA4EADA7A6DCFF00CFCDB8E40FF7D7F5
FC463E6AD37C4EC984909523D6BF47BC47A4C5E20D2FCD8302651943EF5F0F
F8DFC0A9797735CD8A8B5BF427CC43C239EE78E8C7D7A1EFEB5C9809C7194A
D3F8D6FE7E67CD67990C311EFA5A94ADBC4C081F3D690F11823EF578B4BF6F
D32636F7D1344E3B30C7E5EA3DC56BD8DC5CDD616DA2794FA2296FE55CD572
98A7B1F9D62386E29EC7AAA6AF3DD48228016635E5DE395B88F5774B939708
9F9119AF5AF035A47F6C863BF062691C07DC30547A735FA25F173F663FD987
C69F0974AF17E99E299FC3DE2FB4B6D93792AB710DD775F361628432F40CAE
38E1836011DD9761A1467CCCFB0E1DC869613F7AFE26BEE3F20FC2824FED52
D1A96DB13138EC32057A05B6B53585C0785B827247BD6BC7E12D27C273DC0B
5BA6BE908F2FCD65080AE73C28271923B935E8DFB397C29F0BFC64F8C50784
FC577D3DAE9D15BC97AF15B8C35CF92E80C3E66731821B2580270080558861
19C52A5384AACF648F73EAB3FAD4713877EF6CD74688B40F1FF0A8EFB48EC4
D7A95878E54A8FDE7EB5ED3FB487EC25AE78756E3E207C0FB59755D18E64B9
D213325DDA03925A01F7A68874DA3322FF00B4325794FD96FF00636F17FED0
9E12BEF17FFC2409E1BB4B7B87B6B6135AB4FF00686880F31B8923D8AA4EDC
FCC4B0230319AFC7B36A396AA71AF39D94A4A2BBF34B6565AFE96BBDB53EE2
9CA77B189178D411FEB3F5A91FC6608E1EB99F8D1F043E20FC04D623D3FC5C
60BBB39C9105ED9B99217EB8560C15D1B1CE1971E84E0D79241ACAFF007ABC
7C1E5184C4D2588C3494E0F66B546B2AB24ECCF7197C532CBC2B566CBAB3C9
CBBE6BCC17595C7DEA47D6947F1575C32551D909D53D025D4C0079AE7EF757
0A0E5AB8ABAD7D141CBD7233F88ECAE2E7C8BABE8ED221CBBBB0C85FF657AB
1F415EB6132493D6C652AA77DA7EA7A5DDEBF66BADC53DCE951CAAD7696EC1
257881CB223B70ACFD0360EDCE70718AFD274FF828B5B59E9B67A2E9BF0C13
4FD3EC628EDE18A2D5F2B14312844545FB18E1540006474AFC86F0F78CF475
F1B471B66DF40BB9238199CE5A31C2F9EC4F7CFCCDFECF03A0AFB57E23FC0A
D7FC0408D4616423AE457D3D6C8294630F6D0E9DDFE8F73CAAB8B9F33E57A1
F575A7ED83F0C3C7056C7535B8F0EDE4BF28178AAD0163D84C84803DDC20AF
94BF6B6F0C5BF8BBC0F7B25A8134B6ABF6CB765C364C60B7CA475DC84818F5
AF99752B60A5918722ADF873E205F6800681AC4CD3E8736542BFCDF672DC6E
5CFF0007F797A770339CF6E5D97468D45528F4E811C5392B48F830494FF32B
F5B3E21FFC124BE285C783B49F893FB3FEB36FE2ED335BB0B5D423D2AF196C
B50892EA1594224AE45BCC06E3F333427181B49CD7E637C40F859F12FE146A
8345F897E16D4BC317AF9D91EA16B2DBF9801C168CBA8122E7F890907D6BF5
65AABA31317C3DA26AFE2AD7B4EF0CE816ED77A96AB7115ADB42B8CC92CCC1
1172700649EA781D4D7F5CBFB1AFECEDA07ECFFF000C74DF08E9C167BF71F6
9D46EF186BABC703CC7E9C20C0441D900CE4E49FC55FF825DFC1587C4DE36D
4BE306B56E25B6D04FD8B4EDC011F6B95434D20CF431C6C147FD743DC57F4A
DE1684456A00F4AF371B52EF90D20BA9D8A0C0A7D20E94B5E952568A4430A2
8A2B4105145140051451401C6FC41D75BC37E0BD5F588DB64D14056220E089
65FDDC678F46606BE2BF0C40B1A20C74C57D09FB45EA46DBC27A7E9AA70D7B
7AA48CF5489198F19E7E62BEDFA5781680C1556BF44E11A0A345CFBBFC8F07
349DE56EC7AD6998C2D79178C2EE4F8C1E34FF008549A5DC7D9FC33A16DBDF
12DF2B6102C5F3ADAEEE9C70CDFED6071B181E13F687F8E4DF073C0C91F879
05E78CFC46FF0061D0ECC2EF2F72F806665E729086DC7230CDB578DD91F96F
F16BE24EA7FB35782D3E12783F5891FC65AD3FF68F8B351121779AE25CB7D9
9D8E7781B8EFCF5624F5AF0B8C73654E7F568EADEFE4BFE0FE5EA74E5396BA
8BDABD96C7BD7ED67F1C27F889E2AB3F06F816D9DF4DB378F4BD16C2DD4969
599846815075791B000EBD057EB27EC8BFB3A58FECEBF0B60D1EF1239BC57A
DECBCD72E970C5EE4AFCB02BF5315B82513B125DC005C8AFCFEFF8269FC05D
47C71341FB537C47B20888B243E1EB6901DAD282639AFB69EC30521CE792ED
D5518FED457C9D28595CF4D46C1451456A505145140114C7086B8BD5651935
D4DECEB1A1E7A579DEA776198F35EC65945B773C2CC2A272B2E8615E49D6B9
2BD93AD6BDE4FD79AE56F67E0D7DE60689E636737A93F06BCBF5B20AB577FA
9CFC1AF33D626186AF6AAED60823C5FC551074715EB9FB187C4C9341F1CEA1
F0A352982D8EBAB25E5886FE1BC8547988A7FE9A44A5B1EB1F1D4E7C97C46E
0AB57CFD7BE24D47C13E28D2BC67A3B94BDD16EE2BB8F1DCC4C18A9F50C32A
477048AF97CDE82A94DC4F4B0B3E595CFE8668AE13E1BFC4CF047C5CF095A7
8DBE1F6AB16AFA4DDE57CC88FCD14A9F7E1993EF472C6787460194F515DDD7
E727BC145145007E7C7C5AF1BD9F853F689D420B9862B7925B1B1915C050D3
261807720024820A0CE4E140E9803D66E3E3CE9D3E94218CAAB6DEB9AE27F6
CDFD9BBC57F17B4AD37C73F0AA48E3F1BF8791A25B795C451EA366CDBBC832
310A8F1B12D13310BF33AB637065FC6ED6FE217C60F076A03C31E35F0E6A7A
16A858A2C1756B346EE4778C151BC1EA0AE411820E0D7F3971C700E39E3ABE
2B0D51C6357E2B3B5D767E47DDE579A61A7469D3AD0BB8EC7DC5F1E7E2C5AC
DA55C436F3ED9483B4AB618376208E847AD7EBFE8735D5CE8BA7DC5EE7ED12
DBC4F2646D3BD9016C8C0C1CF6C57F3EFF000B7C0BE22D5BC5DA0F8CFE2CD9
BDA785AD2F6DE6BA8266DB3CF08705B28394451CB86C3103681CEE1FD0DAB2
B286539079047422BEB7C22CAA8612856A746A2934D5EDD1EA707156267527
0728D95B4168A28AFD80F933E25FDA547FC265F19FE0FF00C2E4FDE42D7D36
B1789D825BED488E3DD7CE15F07FC52C7C68FDB7740F0B002E2D27F10DB472
46577ABD9E987CE9D703B3450BE4F400E4F02BED18F5A8F58FDABBE2278F6E
183D8FC37D016CE17208024F2BCE71CF70F2CAB9E3A7A609F937F61AD164F1
C7ED4BAE78DEF41923F0D69773387C9E2EF50904499FAC5E757CC4FF00798B
B79FE5FD33D05EED3FEBA9FB57451457D39E78514514005145140051451400
51451400514514005145140051451400514514005145140051451400514514
005145140051451400514514005145140051451401FFD4FDFCA28A2800A28A
2800AFC26F116B123FC56F1B3DC47F6795B5ED50BC79C946375212B9EF83C6
6BF766BF007E3AC13F833F68BF1EE8B72A63336A92DF27390C97E05D2907FE
DAFE078ED5F91F8C380957C042DD25FA1EAE5353966CF6DD1B580AAA735EB5
A1F8C65B2DA51F18F7AF90746F108DABF3577969E20181F357F1C661923E67
A1F5B4EB1F6347F14EF562DA263F9D71FAFF00C466304B737B72228A30599D
DB6AA81DC93C0AF9FD7C4191F7ABC43E21F8C25D5F5A8FC3304B8B6B60AF38
07EF48DF32A9F65186FA9F61519470ACB155D525A2DDBEC8CF158D54A0E473
FA6892D2EBC496B6110B8D3F53BAB87B73CA6D89E56643B48CF0A40C71F874
AADE14B9D43E1EDD6A973047F6B8B5331192263B36345B8654E0F2437391D8
57E9CFECD3F05BC03A9783878875E8A3BDB897236B1E10015F287ED2FA5F84
341F155C59F864A792A71853900F7AFDD7118494E828D549C27A5AFABB5B7B
6AB5F43E3E35651973C7467900F1E69DAEAC91C2CD0DCA0CB43260301EA319
047B8FC715E59E28D586D7F9BD6BCD7C5B7B25A482FACDFCBB880EE461EA3F
983D08EE2BE87F0DFC2CF0EF8A3C33A5F8AFC41E2BB2B08353B68AE3C9F377
489E62825595412194F0467A8AE6870E53C372D587C2FF0003D08E65CD1F7F
73E64D1355934EF1C681A9C4FB1E0D46D5F23D3CD5CFE62BEE9BCF8B33347B
7CE3F9D70E3C23FB387871D67BDD4E7D56788861E4C38019790434A4E39E87
14B73F16FE0A787C1FEC9F0D7DA997A3DDDC139FF80A6057E9191716BC2527
4E9D1726FCEDFA33CDC4494DDD19FABF8C756D5495B68E494B74C026BCFF00
58F869F10BC6B652A596933BB203246C1589122F2A4601E7B7B8E2BA3D47F6
B7B6D3014D074ED374C03A18A05661FF00026C9AF25F10FED7DE2BD433136B
139078091B6C1F805C57A92E27CDEBE94E8C52F46FF5FD0E7E54755E06D76F
E2697C3FE2085AD351B33E5CD0BF0CA7191F8104107B835178CBC17A45E6A3
1EA97D70F656841F3E58D3CC20019076820FB715F2468BE3ED793C59FDB12D
ADC340EC5649595B1B01382491838F6FF0AFBABC3FA85AF89F485E43A48BF5
CE6BB731E6C0CD622F64FF00AB19AF7A3667875C3FC0CD340FB4B5EEB4F19C
8CAA4499F519DC6A8CDF17BE1F68E36E89E16B51B7A35CC8D29FCB815E5FF1
47E175E685E2396686F1E1D36E9B7468173B4FF12EECE07B71D3E95E6A3C27
A627FAE926B83DF73E01FF00BE40AFAACBE34F154A35E12BA6727D575D59DA
EBFE3E8BC45E239F588E386D9A5540520511A028BB4600F602AE49E33BE96D
C40F72C5076DC715168BF03EFF005D960BC473A45B8FBC305DE45FF749183E
E4FE06BE87F0FF00C20F0A6911AE6D3ED9201CBDC7EF09C77C11B47E0057CB
677C659760E5C8A5CF25D17F9EDF75CE9865D291F27DF6B0F3EE588348D8CE
1464FE95EB7FB246B3368FFB46F866EAF2268A1B8177017704052F6F26DE71
D4B00BF8D7D350F87ECEDE311C10AC6A3A05000FC85326D1A120FC82BE3F11
E2753A94E74950D249AF8BBFC8ECA597F2B4EE7EC4E95E2FB1D1347B8D6752
9BCAB4B185E795FF00BA91A9663F90ACFD03E39F87BC33F0960F893E2ED35F
43835EBC944514444CEE2462EF3150A9C16CE71B8F7C926BF1F9A7D6EC6CE5
D3EC751B9B7B59976C90A4CEB1381D03203B58703822B4BE23FC5EF889E37F
0AE93E15D724825B5D0A0305AF951084EDC019709F293C0190A3DEBF36CC70
14B1CE3D1ABDB5B59C972B6BA5D26ECFA37A1EA42AB89ECBF1BFC19F14FF00
6BBBCD56EBE0ADBC3ACE97A23446466B98EDC4AD26E658E2694AA971B7E60C
CBB78C91915F13EB3FB297ED49E1A2EB7FF0E3589BCB009FB25BFDB7AB6DE0
DA9973CFA678E7A735FB17FB0C49E10F0B7C0FD3B48D075282F7569E492F75
811B1F323BCB83CA3AB00C3CB4558C1C6D6DA4A920E4FDB91F88108E706BE8
323AB1CA70F1C061E09C21A26F77E7F31557ED1F3BDD9FCC6E99FB37FED37A
A8FF0045F867E204F9827EFF004F9EDB93FF005D9538F7E9EF5EE9E0DFF827
A7ED4DE2C9E21ACE9965E15B57C1336A3791B90BD4911DA99DF3E8182F3D48
1CD7F40875E8F1C62A9CDE200070715E954E25AAFE0A6911EC9773F08BF69B
FD917C0FFB327C2AD2B51D5F5FB9F1478C7C437E96D136D16B676F0C4A659D
D210CEEEC3091E5DF04393B41C63E07D63C11178C6DA38EC9D2DF5287885DB
E54619E51C804E3D3D0FE35FA99FF0550D526B99BE17CA3718636D69598676
866FB11504F4C901B1F43EF5F981A66A8612AE8D822BEFB876A557858576FD
E777F8B5FA1E5631FBF647D27FB29FEC80FAC78E34BF16FC66D52CF4CF0BE8
F3C775259C52F9D737C633BD621B3E548C9003B16DD8E1579DCBF74FED77F1
8FC37E29B9961D24AB0031918C9AFCD3B0F897AD595B7910DCB2AE3B1AE6B5
6F155EEA6E5EE652C4FA9AF4F1356A557EF9C8DB63356BB59257607A935E75
ACC8AC8C0D6ADDDF8C1E6B88D56F810C01ADF0B4351A47F4CBFF0004D6F8AF
2FC4EFD97F48D3AFE6F3751F055C4BA14C491B8C36E164B538EB810489183D
CA1EF9AFB975DD0341F1369B2E8FE25D36DB56D3E71892DEEE149E1718C61A
390329EBDC57E3F7FC11BBEDE7E1F7C49790B7D88EAB6622C9F97CD10379B8
1EBB4C79FC2BF593E21EB7FF0008FF008235CD5D5B6BDBDA4BE59E3FD63295
4EBFED115F6B1AB6A49B358EA7C91F0F3C3DE0AF0FEA9A95BF80344B2F0FE8
725E4F25BDA584096D6EAACE7E658E30157775200C0CE0702BEBCF0EB8F242
FB57C7BF0EA558EDA103D057D4DE1EBC1B579AF1EABD6E6D13D254D3AA089C
300477A9EBDAC2D4E6899C90514515D0485145140051451401F247ED397BB6
FF00C316209185BA91876393105FC460FE75E1EFE24D27C31A1DE78875CB95
B4D3F4D85EE2795CF091C63731FC8703A9E82BD0BF6A2B8F2FC6FA2293C7D8
49C7FDB56AFCDBF88DADBFC6EF1DB7C23B3BD365E07F0981A878B2FC1C2379
243ADB06E87691C8EEE707050E7EEB0D9AC305972AB2DFA2EEEECF22584957
C438238ED4FE26CCBFDB1FB5E78FE1D9737024D37C0FA5CDFF002CA25C8FB4
953DC64B13DDC9C70057C77FB3BFC1CF127ED71F1D9349D5E699B49497FB43
5FBE04EE4B6DD928ADDA499BE44F4E5B0429ACBF8FFF0014B54F8DFF00116C
F40F07593FF66DBBC5A5685A6C03A26E11C4AABD3739C67DCD7EE9FEC9BF03
B47FD9DFE19D9F86600936BB7FB6EB58BB5E7CEBB23945279F2E207647D320
162033357C5E4595D5C7D7957ABAADDBF3EDFD743DACCF171C3D254E9FC8FD
10F075968DA0E8361E1ED0AD63B1D3B4D823B6B6B7886D8E286250A88A0740
AA0015D7D78AE91AD98CAE1ABD12CB5D8645018E2BD3C7E573849B4B43C4C2
63925CB23A6A2A8AEA16EC33B8535F51B741F785797EC67D8F43EB54FB9A15
56E2E5215249AC4BAD7A18C1D8735C6EA1AF1933F35776172CA937AA38B119
82B5A06AEADAA8248078AE0EF6F7713CD55BCD48B13CD739737B9CE4D7DA60
32CE5478B29DF525BABAEBCD72F7B75C1E69D75798CF35CA5F5EF5E6BE9695
250446E51D4AEB835E6FABDD0C3735BDA95E8E79AF39D5EF786E6B2AD3368A
387F105C021ABE70F1AB878A41F5AF6BD7AF7E56E6BE7EF15DCEFDCB9AF0F1
72D0ECA48F9B7E09FED5DE30FD8C7F68CD4F5BB049753F077881A26D634912
605C44E07EFE204ED5B889B71463F78650901891FD5A7803C7BE13F8A1E0BD
1FE20F81B508F54D0B5DB75B9B4B88FA3A370410795656055D4F2AC0A90082
2BF986F0F7EC45F103F6A7F1F43A8E8D7B1E89A0DB402DE6BD950CCED22BB3
6D861054B101BE6666551EA4822BF5FF00F654F819E28FD87747D53C29A9F8
A6EBC59E11D4E41751D9C96CB1496572701E4B72256056451F3C7C6480C083
B837E658EC552856945BEA7BF493714CFD2AA2B931E3AF08B40972BAAC2C92
20906092DB48CF2A06E0DFEC91BBB62A94DF11FC1F0E40BD691B19012194E7
F1DB8CFD4D43C4535BC97DE5599DCD7CB7FB5EFC389BC7BF082F352D261F37
5AF09B8D5ECC2E7738814FDA2218E4EF84BE179CB85E338AF589FE2B786A2C
F9715CCD800FCA8A33EDF3BAD655D7C5ED276BAC3A7CB22918C4AE88083D41
C17FEB9AE1C6D4C3D6A52A3524AD25635A529424A71DD1F9A3E0FD56DBC51E
184F983ACB1FD7A8AFD16FD9D3C60FE27F8736DA6DEC9BF51F0F31D3E7CF52
9100606E833988A827BB06E739AFCBDD36183C05F143C41E08B7416FA70B86
BAD3E3072AB65724BC48A70B911F31F4EA86BD2EE7E2D78CBE074B7BE2CF06
C36B729A8C71C37315E248F0FC8D94936C7246DB972CA3E6C618D7F3AF0763
DE4B9ECF0B55FB93D3FC9FF5DCFD0B36A2B19815561BAD7FCCFD70A8E69A2B
789E79DD638A352CCEC4055503249278000EA6BF15B52FDBEBE335CEFF0026
EF4BB0DCDB8791699DA3FBA3CD7938FAE4FBD79778ABF6CBF8CFAFE957BA5E
A7E316FB1DEABC6F1456D6B08D93028537244AFB76B11CB7BF5E6BFA11E794
BA26FE47C12C1CBAB3E8DD23C4B2DB7ECD1F13FE2C5C8686F3E23EBB379618
92C2296532ED04F650CCA0760315E85FF04CDF0B9B7F00F8DBC7F326D97C43
AC2DA21FEFDBE9D10DA73E9E64F28FC0D7C63F15BE2C783DBE07F84FE1DF85
75C9AEE0D3E27B8BB864812248AE5F248575019C7CC7A93DBE95FAC9FB1BF8
66DBC2BFB31FC3DB2B6C1FB769A9A9BB00C373EA4CD78C4EEE4E3CDC67A600
DBF2E2B8F265CF55CDF45F9EA6B8AD2363E99A28A2BE94F3C28A28A0028A28
A0028A28A0028A28A0028A28A0028A28A0028A28A0028A28A0028A28A0028A
28A0028A28A0028A28A0028A28A0028A28A0028A28A0028A28A00FFFD5FDFC
A28A2800A28A2800AFC75FF8297FC3DB8F0FF893C33F1B74E4C59EA08347D4
300FCB3C7BE5B790FF00BE9BD0938C6C51CE6BF62ABCFF00E29FC36F0D7C5E
F87FADFC38F16C464D335BB730BB2E37C4E0868A68C9E03C520575CF195190
46457999C65B1C5E1E5425D7F334A551C25CC8FE71340F16ABC6A43D7A3D9F
89C151F3D7CBBF143C1BE37FD9E7E246A3F0C7C7D118AF2C5B7DBDC0044379
6ACC44573093D524C1F7560CAD865202E9DE365651892BF9C739E0E9426D38
9F434717747D74BE2518FBF5F3DEA5E2409E35D583B61BCE53CFA14523F4AC
D8BC60A47DFF00D6BCA7C7FA9CD0DEC7E26B1CBAED11DC2AF5017EEBFD31C1
F4E3DEB3E1BC83D9D79464ADCCADF3BA6658FA9CF0D3A1F69E83F1C3C4FA1E
9874FD3EFDE28986085622BCDBC45E36B9D5257B9BC98C8ED924939AF95ED3
E20C724631275A86EFC6324EA42357D7C38664A5AA3C5E63B4F16788165474
53926B94D37C51E304D261B4D393FD122DCA9233A8046E39E339E0F1D2B889
F509F51B85B580EF964FC80EE4FA015D44B742DEDE1D32D0EE2AA141FE64FF
003AF75E094231A318A6DBEA694A09DDBD8905D78AF5467171A82C2AA7195C
BF3FF8EFB52FF602CA775EDFCF31F40428FE44FEB562DD9608963539C753EA
7B9A9FCFAFBDC1E554A9C127157EA73CA5AE8411E83A2C5C9B7121FF00A68C
CF9FC189157E28EDAD976DB44910F4450BFCAAA19FDEA269EBD28D34B64497
DE6AF45F85FE306D0B525D1EE1F16F29CC44F41EABF876F6FA5791BCFEF555
E76521E36DAE8432B0EA08E86BCCCEB29863284A8CBE5EA08FBDFC5FA0D878
BFC3B2C72900EC2CAFFDD20641FC2BC23C13F0FE1B4916FF0053027B807E41
D51476233DFDFB56AF853C53A9F883C3B6D6B3031A13FBCFF6F61C607FB39E
7DF8AF4FD32DC00AA057F3D63735C5E06954C042764DEBFD79F53BF0F457C4
CD6B0B00000AB5DB69FE1CBBBAC6C8C9FC2BAAF02783A6D6AEA38D5376E22B
F42FC01F01EC63B28AE7514C6E00E31CD7E792E7A93E582BB3D048FCE86F05
5FAA6E688FE55CCDFE8935AE43A915FB09A9FC1CD064B564B78F6B638C8AF8
B7E2978063D1E59542636E7B54D4A15294B96A2B319F105E5BEDCF15C95FDB
2904115E9BAE5B086465F4AE02F7BD7A5839B4CCD9C8E91E23F12F80B5D8BC
4DE10BE934ED42DFA3A1F95D7BA48A7E5743DD5811DFAE2BF533E09FED19A4
7C57D10EFDB61AF5928FB6596ECE3B0962CF2D1B7E6A7E53D8B7E55EA68083
5C8695E29D73C01E24B5F15F8766315DD9B64AE4EC9633F7E371DD58707D38
23900D7DAE062AAAE59193958FDF53E2718FBF5C7F8BFE2AF86FC17A636ADE
26D4A3B1B71C2EE39791BFBB1A282CEDECA09AF927C39F15BC67F1634F57F8
4DA56D83CB0D79AB6A7FB8D3EC091F3AB31C6F74E4601C6791B8715E33E32F
8C5F09FE115F4BA9DA5C8F89FE3E5186D56FC674EB361DADA03C10A7A1200E
F8AF7B09C3CE6FDED11CD5718A3A475627ED90DE3CF8C9F09878FF00FB0E2F
0E7873C3130BBB27D524116A5A889311482DE1C9C2056DE4739080EECFCB5F
959A7EB2ACA0EEAF67F1DFC5EF89FF001D3C52639A7BAF10EA770484897252
35F65E15107A9C003A9AF36F1F7C11F883F0DB4CB7D7EFA18EEACA55DD702D
58CBF64724FCB2703E5C63E619504EDCF42DFA3E5784A74692C3C9A5D91C2F
9A779304D5463EF5364D5463EF5794C3AD3B0EB563FB49DC7DEAF4BEA06763
B4BAD5720E0D73177765B258D65BDE8C649AFD3AFF00827D7EC41AAFC73F11
587C60F8936860F877A3DC978609787D5EEA02088D508FF8F656FF005AE787
C18D727794D1C234A3CD2D11518B6EC8FD73FF0082717C1EBDF83BFB2FE883
5B80DB6B1E2F9A4D7AEA36521E35BA544B646CF20FD9E38D8A9036B330C672
4FBD7ED19A99B0F84DAB6D2435CC96B1023B033A139F62A08FC6BDA4154508
802AA8C0038000EC057CF1FB51897FE14EEA77083E5B7B8B37739E8A6744FE
6C2BCB9E6FED2AC20B6BA3AD53B23C0FC09A80586219F4AFA67C3FA90DABCD
7C45E07D5C7931FCDD857D29E1FD5C6D5F9ABDAA91334CFA874DD415D42935
D0A3822BC5F4AD5F81F355CF197C55F0A7C31F06EA9E3BF1BDFAE9FA2E8F09
9AE256E4E070151472CEEC42A28E59880393534AB383BA1B49A3F3E7FE0A61
FB55F8B3E163F86BE07FC3AD61FC3DABF8AE1377A96A50B05B8B6D399DA045
89BEF46D2BAB9F31486511E14E4923F1EBE2EDCD87C17F18697E23F835F137
51D5B5D9E18AE2EF5282E258A61700676B481B7381D3E62411C1E2B96FDAB3
E37788FF0068FF008957DF16FC46069ED2AADA695603916DA642EED144CDDD
F2ECEEDD19D9880AA401F31E8F67E21F166AF6FE1EF0DE9F73AAEA978E2386
D6D6279A691CF002A202C4FD054D572AD3E74CF327794AE8FD55F0CFFC162F
F68DD1ED2D74ED7F40F0DEB7F6485627B87B7BA82E67750079B234773E56E6
C658244AB93C051C57AFE8DFF05ADF1042A3FB7FE155A5E36C009B5D5E4B60
64EEC03DB4F853CE17391FDE35F971AC7EC6BFB5A68C1E4BDF849E25902B05
3F66D3A6BB393E82059091EE38F7AF30D63E0EFC64F0F094F883C07AFE9820
6D927DAB4BBB8763671B5B7C63073D8D7A3194FA48DED23F7B347FF82D1FC3
4B807FB77E1CEAB66DB571F66BC82E46EFE219758781D8E39F415EB1A4FF00
C15EFF0065BD4722F74DF12E96576F33D8DB386DDD76F937521C2F7C81D780
79C7F2F3749776172F677D0C96D3C670F1C8A51D4F5E55B0453C5E7A3568A5
52DB9D142CD7BC7EEE7ED73FB7DFC22F1CB5BEA5F07350B8BFD492C8DAC6D3
DA4902C123BB1F34F9AA376D07200CE4E01E335F11FC49F8BBE0DF0A7C16D2
7E15FC2BD4DB51935B55D47C45A91578E5B8BA7E7C83E600E56339CE782790
4839AF80C5E1FEF57ABFC1CF0BC5F117E227877C1B7C924D677D76BF69119D
ADF668C6F97E6FE1F954F3FD6AAAAAD8874E937A2D12F5674D2942929491FA
33FB01FC08F2671F1E3C5F0032CAAF16870C8394539496EC83D0B0CA47FEC9
66C60A1AFD76B2D4718E6BC17429AD34FB2B7D3EC234B7B6B6448A28A30151
1100555503800018007415DDD9EA638E6BF67CA3070C2D15461F3F367C6E2E
B3AB37391ED569AA118C357496DADB2FF1578B5B6A78C7CD5B30EABEF5DB53
0D09EE72EA8F674F10381F7A87D7DC8FBD5E4E9AAFBD49FDA9EF5CBFD954C7
CC7A1CDACB37F15644FA916EA6B8D7D4FDEA9CBA97BD7453C14222BB3A89AF
FAF3589717FD79AE7E6D4BDEB12E752EBCD74AB476051B9B3777FD79AE4EFA
FF00AF3542EF52E0F35C95F6A5C1E6B19D5348C49352D4383CD79CEADA87DE
E6ADEA5A975E6BCDF57D4F86E6BCEAB54DE113035ED4386E6BC475BB9F3646
39E95D9EB9A896CA83C9AF33D466E18D78D889DCEB844FD18FD94BE33E83E0
BD06D20D8B22C7198E44C807713F31FAE7AD7A87C6AFDA9344B48ED2FA6D96
F68AC635CB67E6232327D4E0FE55F809E29F8C5A87C37F180B1B199D60B985
66900F99558B32F4EA321727FF00AF599E26FDA56DF59D2E4B5BC2FA8EF1FE
A59484247233B87AD7E5D8BCBE6AA4A0DDD5EE7BB4E69C533F5AB50FDB43C2
D06765E47C7A115C36A1FB7168899115C6EFA035F9F9FB18F86346FDA47E36
41F08BC51A9C3E189B58B4BB9B4EB98ACD2E775DDB2F9C2164774E0C4B2364
3839503BE47EABEA5FF04C1F8A9A36E3E0DF1CE89A83292C9F6BB29B4FC900
609311B9C1CFA671FA0CFF00B362BA15CE7CFF002FEDA779A871A4DADCDD96
6DA3CA8D9F2C7B0C0EBCF4AF3FF13FEDADAD685702D357B69F4EB874F3163B
90617642480C15F0482410081D41AFA1352FD8DBF6CCF0DB38D3F4BB1D7BCB
3C1B0D52250FF36DCAFDB3ECFDBE6E71C7BF15F9D5FB43FECD5FB62F88FC5C
35ED7FE12F88E58AD2DE3B489AD2C9F50FDDA1924CB1B3F3801B99B9CE00C6
7A8AAA78257D620E47A668BFB545AF8DBE22E8925EC8A93FCD007DD9DCADF3
0527D8838FA9F5AFD289E0B4F1878465B5980912E222A47B115FCDDEA7E1EF
1A780B53824F1068F7DA1DEDBC8AE8B7B6D2DB3875C38E24553D083F435FBA
BFB32FC46B7F1AF82AC2E84819A489770CE70C06083EE0F15F8978BD903A6E
9E61455ADA3F55B1F6FC258EBA961E67CD5E26FD957E27F877E1AEB3F19BC4
7F103C31A3784B48790C8935CDD49A9ED129862416B15B30334AE02A219002
792C13E6AF82F47D5BE2278C5651A0D95CDEAC632E6352401EE4035F60FF00
C1427E177D8B50D2FE26E9F17C927FA15D903D72D131FC77293EEA2BCAFF00
65FF00DA9ECFE0678475DD12EF4986FE7D40B08A49141645651C64FA1C9FC6
BF53E13CCA966197D2C525AB5AFAADFF00CCF98CD70F3A15E54BA74F43C4F4
CF89373A3F882C6D7C7FA7DD5FE956D7086FAD20B816B3CB02B0DE892B4720
462380C50E3D2BFB14FD9E7E2C7C31F8D1F08BC3FE36F845201E1B6B74B586
D8AF97258B5B28436B2A64ED78400B8048230CA594863FC54F8F3C52DE30F1
2DE6BA63117DAA467DA380326BF673FE0893E3DD6A2F1C7C45F85ED348FA45
D69D06B09110C628AE60996DDD81C6D5691255072416118C0214E3E9E95351
D91E6B937B9FD12D14515B9214514500145145001451450014514500145145
00145145001451450014514500145145001451450014514500145145001451
45001451450014514500145145007FFFD6FDFCA28A2800A28A2800A28A2803
E6AFDA63F65AF871FB4F7843FB0BC5D0FD8B59B2573A66AF0AE6E2CA47C678
CA89626206F898E18720AB8571FCD77C76FD9BBE327ECCFE206D2FC7FA5BB6
9524856CF56B60D2585D2F38DB201F23E064C6FB5C75C6DC13FD735666B3A2
E8FE22D2EE744F105841A9E9D78BB27B6BA892682543FC2F1B82AC3D88AF37
1D9653AEAEF73485468FE332D35DDC002F5B716B08C304820FAD7F415F14BF
E0975FB3778F26B8D47C2A97DE07BF98971FD9D2892D3711DEDE70E02E79DB
1BC63D302BE3DF10FF00C11F7C7566B2B784BE2558EA1B7945BCB096D09E3A
1292CE01CF03D7AF1D2BE5311C313BFBA8EA8E251F91B7FE1ED22E6433D9C8
D66EDC909829FF007C9E9F8115983408E2FF005FA8B3AFA2A053F992DFCAAB
F88135CF0B788B53F09F886DA4B1D5747B99ACEEE093878A781CC7221F7560
45518EEBCDE5D8B7E35D34329C5DB95CB4F3B7FC39129C37B1D2DB496F68A6
0D323E4F0CE7927EA7FA0FCAB4ED8888124EE76EAD5CDC57214003815712EB
DEBDDC0E570A2F99EB2EE613A8DE9D0E945C67BD3BCF1EB5CF8B9E3AD2FDA7
DEBD2B199BA6E3DEA26B8F7AC7FB4E7BD44D73C75A2C06ABDC536DC3DE5D45
6B1FDE9582FE7DEB11EE7DEBA3F06627D76363D2252DF8F4FEB5C599627D8E
1E755744CA82BBB1F52F866DA2B3B58A08861635007E15EB3A22879501F5AF
25D2260116BD2346BC11C8A735FCA79A46529B93DCF5E07E89FC00D32CA4BD
B76940EA2BF48AC96248D540C0038AFC84F85DE395D1E58983E318AFB9345F
8E7A77D8D05C90CC075CD6191E60B0588F6D28DCA92BAB1F4BDE4B12464920
00335F067C75D6ECE6B89C46C0F5AEAFC6BF1EA17B5921B4708083D0D7C31E
3BF1E49AACD231933927BD19E6672C757F6AE364115CAAC798F892E95E77C7
AD79BDE4A0935ABA96A0657639EB5C8DDDC75E6B0C2D0D41B3235090106BCE
35A2195857657F71C1AF3ED5A7C8615F639652B34632672979F147C69E1DF0
CCDE06B3D52E468573726E05923B08CDC480213B4752C15463D47A9A97C29F
0975FF00174B1EA7E32B96D234E6F9BC95FF008FA907D0E447F5604FFB3DEB
02D2FF00EC3E26B0BEFF009E732827D03FCA7F435F46C5AD0C7DEAFD4F0755
AA4B956BDCE2749735CF44F0C597873C17A70D2FC2B631D8C3FC6C066490FF
007A473F331FA9E3A0C0E2B4E7D57CE044A77061820F2083ED5E6635A5FEF5
4526B8A07DEAC5E1DB7766B7391F16FC13F875E229DEF20B77D26E5C925ACD
82213EF1B0283FE0205795DC7ECF9A2DB36EFEDE9CC60F4F29338F4CE7AFBE
3F0AF68BCF102A83F362BB5F847F09BE207C7DF11FF627832DCC76103A8BED
4A5045ADA231E4B1E37C98E56353B9BD86587A54F13569C6EE764BB90E29F4
38EFD9EFF651B7F8B1F10ECFC21E18B696ED632B2EA1A85CE248ACAD81F9A4
DA004DE7A46A4125BBE0330FE9D3C1DE14D07C05E16D2FC1BE17B6167A568F
025BDBC40E70883AB13CB331C9663CB3124F26BC77E077C23F057C05F0645E
14F09C7E6CF2ED92FAFE4502E2F670305DC8CE14648441C20E992598FB03EA
C3D6BE5B31CDA75A5A3D11B429A47485C0EA715E25FB43D9A6AFF05FC5B6A3
930D99BBEB8FF8F365B8F51FF3CFFF00AC7A1EEE6D5FFDAAE6F58B9B6D52C2
EB4DBCF9E0BB89E2917D52452AC3F235E7D0AD38D48CFB34FEE2DAD2C7E5B7
8175F1E5A2EEAFA67C3DAE8DABF357C0FA6DD5CF85FC457FE1DBD6FDFE9973
2DAC9DB2D0B9427E871915F457873C461953E6AFD72493574701F63E9FE218
6085A79E558E38D4B3331015540C9249E0003AD7E1E7EDA3FB5349F1D7C4E7
C35E1EBA61F0F3C2F3130ED381A9DF282A6E0FAC6B92B10FEEE5CF2FB57D13
F6B7FDA42EE7B59FE0C7826F0C4664FF0089E5E46C3F77091CDA291FC4FF00
F2D304607C9CEE603F2A7C43AD24B8B6B6F92DA1F9517FA9F735C153DE9724
4E3AF56FEEA2C5BD9788FE2078A2C7C33E1AB39352D63579D2DAD2D6119679
243B55467000F52480072481935FD4B7EC69FB38F877F658F8631F87E1315D
F8A757D971ADDFA6489AE003B628CB004430825506064EE7201622BE0DFF00
827EFECCB0FC2DD1E3F8C7E3DB4C78C35988FD8619305B4EB29075C7F0CF32
9CB73954213826415FA8F16BC31F7ABA1C6CB951BD0A6A27B6FF006DFF00B5
4C7D6BFDAAF1F1AE8C7DEA6BEBBFED567CA6FCC7A46A77965A95ABD96A30C7
756F2637472A8746C1C8CAB641C100FD6BC8B5FF0085FF000735FF003DB5DF
02E83A89B8004A6E74BB59BCC0A001BB7C673800633E9534DAEFFB558B73AE
F07E7AA4AC1CC7E7E7ED99FB267C01BFF0F695E24F0EF856D7C2D76976B6D3
49A2451D8AB44F1BB0DD0A279070CA3E6F2F7738CE2BC23E0DFC2EF01FC2DF
36E3C356EF25F5CA8492EEE5C493941CEC0405555CF242A8CF19CE063EF8F8
FF0027F6DFC2FD6E34E64B444BA538CE3C870EDFF8E061ED9AF80B40D7432A
FCD5F7BC2FECDD3E66BDE4F73CBC7395ED7D0FAAB4DD5B81F3576967AB0C0F
9ABE76D335AE07CD5DC596B2081F357DDD3AE791281EED6DAA8E39AD88B55F
7AF18B6D63A7CD5B50EAFF00ED575C6B99389EB69AA7BD4BFDA9FED579726A
FF00ED54DFDAFF00ED56BEDC9E43D21B54F7AA926A9EF5E7EDAB8FEF554935
7F7A4EB8729DC4DAA7BD62DCEA9D79AE3A7D5BAFCD58973ABF5F9AB19562D4
4EA2EF55EBCD7257DAAF5F9AB9FBCD5FAFCD5C85FEB1C1F9AB96A563450357
52D5460F35E6FABEADD7E6AAFA96B1D79AE26EEF0B9258D79F56B1BC62457B
74589627935C56A973B51B9AD6BBB9E0D799F8BB5D8348D2EEF52B83FBBB68
D9CFBED1D3F1E95E7559D95CE88C4F88BE2BEA6BAA78EF529236DC9015847D
6350187FDF59AF39A9EEAE25BCB996EE73BA49DD9D8FAB31C9FD6A0AF8EA93
E69391EAC55958F63FD9E7E213FC28F8EBE01F88C2E0DB43A06B7637370EA7
19B559945C213E8F09753EC4D7F7335FC01D7F76BF077C46BE31F845E07F17
23F98BAE687A65F06C96DC2E6D6397396E4E77753CD40CF47A28A28020BAB5
B6BDB792D2F2149E0941578E450C8CA7A82A7208FAD7E617ED6DF0E74EF867
F10BC39F137C2F61169DA478841D36FE3B78D22856FA3DD2C32955C7CF347B
C31031FBA04F2DCFEA25792FC72F86D17C59F859AFF824055BDBA83CDB191C
E0457B0112DBB16EA17CC501F1D50B0E84D789C4794C71D82A9869755A7AF4
3B72FC53A35A351743F287E32781AC7E2B7C2ED57C3B7001FB6DBB046EBB24
1F3230F75600FE15FCE76A5A7DE68FA8DD693A8C461BAB295E1950F5592362
AC0FD08AFE88FC0BF113451E1DFB3EB97020BA8418E5B73CCCB221DAC853A8
21810738C1EB5CF7C3CFD883F66FFDAE6FFC5573ACBEADE14F16E9B709399F
4EB88FCBBBB5B81849258678E55DEAEACADE594E0A6493927F16F08F339E1B
155B2BADA27AAF55BFDEBF23EBB8AE8C2A5386221BF5F43F037C33A8E95A57
882C352D6F4F5D52C2DA65926B5662AB32A9CED2570707BFAF4AFD0FF81DE3
4F127C1CF1649F1D3F642BF8EE99E1F2B59F0C5F22C8F25AEF123C58E1D937
282AF1957523AE0907F60FE1C7FC1223F64AF057993789ADF56F1B5C4B1B27
FC4CEF3CA8632E0AEE8E3B3583900F1BD9F0791838C7E63FED59FF0004E4F8
BDFB2C6AD71F16BE00DE5E7893C19665A773173AA6971A727ED088009A1039
F3635E003E62281B9BFA05C7B1F087ED8FECA7FB6FFC24FDA9B4DFECED266F
F8477C6F6684DF787AF5C0BA4DBF7A4B762145C4431CB20DC9C798A995CFD9
D5FC6FF873C7DE11F8BBA858EA57D7E7C03F12EC1D64B2D72C9CDBC53CE9F7
4C86320C6F9FE35C57EBD7ECDBFF000524D7BC17AAD9FC1EFDB4205D2EF805
8EC7C5912FFA25DA7010DE041B549FF9EF1FC87237A260BB0A5DC0FDA6A2AB
59DE59EA36706A1A7CF1DD5ADD46B2C52C4C1E392371B95D19721958104107
04559AB00A28A2800A28A2800A28A2800A28A2800A28A2800A28A2800A28A2
800A28A2800A28A2800A28A2800A28A2800A28A2800A28A2800A28A2800A28
A2803FFFD7FDFCA28A2800A28A2800A28A2800A28A2800A28A2803F09BFE0A
BDFB245CCAEDFB517C3EB2321448E1F134112E48540B1C37C1473850047363
A008F8C076AFC2FB6BEE9935FDCFDF58D96A96571A6EA56F1DDD9DDC6F0CD0
CC8248E58E4055D1D1810CAC090411823835FCC0FEDFDFB02EAFFB3AEB173F
147E18C12DFF00C36D466CBA0CBCBA34D2B710CA7A9B766388A53D388DCEED
AD2007E7BC579EF5792EFDEB848AED8706B423BC3EB52E20766B77EF4FFB50
F5AE516F0D4A2E89A5603A5375E86A36BAF7AC0FB431EF4BE69A2C06B35C93
D2BB5F01CFB752958FF747F3AF35F30574FE13BD106A646701D7F506BC8CFA
939E0EA45762E9BF791F5BE957836AF35DC595F6DC106BC5B49D44155E6BB6
B4BF181CD7F3CE3F02EFB1E94647B6E97E2296DB051B15DAC3E3EBB44DA243
F9D7CF30EA1EF57575238FBD5E14F0269CC7B25FF8CEEAE010643F9D71579A
BBCC4966CD71CDA8E47DEAA925FF00BD10C0F90391B371799CF35817779C1E
6A85C5F0E79AC0BBBFEBCD7AB86C110E43AFEF383CD701AB5D8C3735A17F7E
3079AE1AFEEBCC6233C57D665D83D4CA52302FA53E60901C10C08FC0D7A1C7
E2475FBD9AE3741D3E3D73C55A3E8B28DD1DF5E410B8048F91E4018E472303
3CD7EBBE87E1AF87F158C5A7AF8774E6B78D55423DAC4FC2F4C965249EF93C
93C9E6BEDF0B555386AAE62D1F981FF0939C77AF4BF09FC34F8C1F10248D3C
2DE17BE9E194E05C4B19B7B6F5FF005D2ED4E07BE7F315FA9DA05AF8634B94
5CE97A559D9CBC7CF0C11C6DF9A806BD32D35DE07CF535732B7C30FBC7C87C
93F0A3F611B3DF16AFF19F58FB69E1BFB334E7648BE935C10AEDEE230B8238
722BF4B3C3569A0784747B6F0FF85EC20D2F4DB45DB15BDBA08E351F41DCF5
24F24F279AF28835EFF6AB4935DE3EF57878AA952ABF7D9A46C8F65FED93FD
EA89F59FF6ABC9BFB7BFDAA89F5D1FDEAE4F605731EA12EB1EAD58F71AC819
F9ABCE65D77FDAAC6BAD7400599F00724935A4680731F13FED4DA3C9E1BF89
10F8C6D57165E228D77B0E8B7702846181D3720461EA771F5AF9BFC77F1C6E
3C19E1D161A038935FD494A5B0FBDE4A9E0CCC0F1F2FF083C16EC4035ED3FB
507ED1BF09AF7C1F7DE11B6BC7D73555747B76B201E186E533B4B4C485208C
AB6C2C704D7E53EB9AD5DC6D36A3A849E66A77A324F68A3ECABE9C7F9CD7DE
60B1135878C24AD2DBE47998AAB6765B993E22D5CC2B2DA2CC669A57325CCC
C72D2CAC72C49EA79FCEBEAFFD8CBE02C3E30F1043F163C6D06ED1349973A7
DBC806DBBBA8CFFAC607AC511EDFC4FC670AC0FC55A469571E23D43CBE45B4
4732BFF407D4FF00F5EBE93D1FC71E3CF0ED9C3A7E89E21BFB4B5B7458A285
2E64F2E345E02A2125540F615EA61B0D68DDEE73D185B567EF1DAF88B81F35
6D43E2218FBD5F88FA27ED27F1934164CEAE9A9451E3F7777023838F574092
1FFBEEBE80F097EDA36323476BE36D1E4B162406B8B46F362E7AB18DB0EA07
B173ED5A3A0CEB5347EA2AF88863EF535BC423FBD5F35F87BE24E81E29D3D7
54F0F6A315FDB37F1C4D9DA7AE187556F66008F4ADA6F128FEF565C851ED93
78847F7AB12E7C43C1F9ABCAA2D66E2FA6105A8323B7A741F5AF49D07C2AF3
ED9AF8F98DD707A0FC2BF3BE37F12701924792A7BF57A457EAFA7E2FC8EBC3
60E5536D8C7D4259B5DB3B8D3FC969A0B98DE27038055C15233EE0D7CBFA17
ECC1E2EB503ED9AC42A0671B2267FA7565AFD10D3FC3F1A28089FA56F0D0B0
BF73F4AFE76C67D2433C8CE5F5371A69F68A7FFA55FF00247ABFD8B45FC7A9
F027FC28FF0010E9E998B514988EC6229FFB335645D785FC51A2E4CF079A8B
DE339FD0E0D7DFF77A3A807295C46ADA242EAD9515EB645F49DE22A534EBCE
351769452FFD2794C2B70FE1DAF755BE7FE67C5706B2F1B6C932AC3A83C115
B50EB43FBD5E85E30F055A5D0691536483A32F0457CEFAA8BDD0AE4C375CA6
7E56EC6BFADFC39F1B7019EDA84D7B2ADFCADDD3FF000BFD1EBDAE7CC63F26
9D1F796A8F574D6BFDAA9FFB6BFDAAF178F5D1FDEAB235DFF6ABF685883C8F
667AEB6B3FED55493591FDEAF2B6D778FBD5524D73FDAA4F101ECCF4B9F5A1
FDEAC2B9D6BAFCD5E7B3EBBFED560DCEBBD406ACA5882D40EEAF35AEBF3572
17BAC972429AE5A7D4E494F5C0ACE7BAF7AE59D6B96A06ACF785892C726B1E
E2EBB9AA535DE0726B0EEAFB00F35CF299AA88FBFBD001E6BE52F8E5E2B02D
22F0E40D97B922497DA353F28FC587E95ECDE28F125AE8FA7CF7F7726C8E15
24FA9F403DC9E057C25AEEB177AFEAB71AADE1CC93B671D95470147B015E3E
6389B4795753AA853D6E64D14515E11D815FDBF7ECA2AC9FB2DFC1D4705597
C1BE1E041E08234E83835FC40D7F78FF000D7C3D75E11F875E15F0A5F1DD73
A2E95636529E065EDA048DBA161D57B123DCD0076B4514500145145007E11F
ED81F0EDBE137ED1B7BAB5A2797A2F8F95B5683030AB77B82DEA7BB1908989
FF00A6A3BE69BF02FE237FC2ACF8B7A078C2597CAD2EE9FEC1A89C80BF64BA
214B3648E2270929FF0073F0AFD11FDBABE144BF127E065F6B3A4C2D2EB9E0
B7FED8B4080979238548BA8801927742598003974415F8EFA1DE43AE688BC8
70E9FCC57F3BF1F60E79666B4F31A3B377F9F547D86535157A0E8C8FE9568A
F9ABF64DF890DF11BE0CE9525FCA65D5B40274ABD2C72CD25B01E5C84F53E6
4263627FBC58738CD7D2B5FBF607190C451857A7B4926BE67C9D5A6E12717B
A3F21BF6CCFF008258781FE31BDFFC46F8102DFC1FE379374D3D8FFABD2F52
909CB315507ECF337F7D06C63CBA866692BF111BC61E34F84DAA5E7C09FDA6
BC2F732DA58B79725A5F2ECBCB2DDF765B6979DC847CC8CA4A38E549041AFE
CD2BE64FDA8FF66FF817FB427802EACBE34DBC1650E970C92C1AE79896B73A
670732ADC3FCA231D5924CC67AB2E4023A5C6E667E147C00FDA3FE2E7EC8D6
91EB9F0CEFDBE27FC16B890BCFA4CB262E34E2E72C623866B771C920031392
495DC770FDFDF805FB487C24FDA57C20BE30F857AC2DEA46145DD94D88AFEC
646CFEEEE60C9287838604A3E094661CD7F24FA598BE0E7ED06DE08F857F11
2D7C53E1D6D423B4FED7B78658ECEEAD9D8072F04C00255739DACF1923E491
9486AFE917C2FF000B3F678F8B5F0B2E3C13F0FF004EB7D020BEB768EE6DAC
DBECB24DBC7CCED24251D998F25F39F5F4AF8BE2AE318654E9C254DCDCBA2D
1DBADAFA37DA374DF437A341CFA9EF9F12BF6BAFD99BE10CD3DA7C41F891A3
69B7B6A18CB671DC0BBBC8F6F50D6B6DE6CC0FA0D993CE33835F11F8BBFE0B
29FB29E813496DE1DD3BC47E26651F24B6D650DBDBB1E3AB5D4F14A3A9FF00
964791E9835F8E1FB56FEC03F123E005EDE6BDE1B826F10784E32CE6455CDD
5A20E7F7CABF7D00FF00968A31FDE55EA7F3EABDFC9B3CC2E614557C24F9A3
F8A7D9ADD3F26673838BB33FA19D6FFE0B8FE1C827DBE1CF843777D0EE71BA
EF5A8ED5B683F21DB1DA4E3246723771D8B75AC07FF82E6CA590C7F0540507
E707C47925707807FB3860E7073CF1918E723F0168AF5883FA27D0BFE0B87E
0AB8317FC24BF0A350B004B799F64D521BBDA003B76F9905BEEC9C6738C7BF
7FA77C11FF000571FD8E3C5D34706ADAA6AFE11794851FDADA7315049206E6
B27BA551C039240008C9186C7F27745007F759F0F7E367C1FF008B1009FE1A
78D348F136464A585EC33CA981921E2562E840E70CA0E39AF50AFE032CEF2F
34FBA8AFB4F9E4B6B98183C72C4C51D1872195970411D88AFD09F821FF0005
41FDAB7E0ECF05A6A9E201E3DD123015ACB5FDD7326DC8C94BC056E4360606
F7741D761A00FEBAA8AFCEEFD9ABFE0A63FB3B7ED0B2D9F876FEF4F817C5F7
65635D3356914453CADC04B5BC1B62949380AAE239189C2A1AFD11A0028A28
A0028A28A0028A28A0028A28A0028A28A0028A28A0028A28A0028A28A0028A
28A0028A28A00FFFD0FDFCA28A2800A28A2800A28A2800A28A2800A28A2800
ACED5F48D2B5FD2AF343D72CE2D434ED422782E2DE7459229A2914ABA3A302
195812082304568D1401FCF37ED9BFF04ADD7FC3D797FF00127F663B47D5B4
57DD35CF871497BDB539258D8E799E2C74889F34745F33385FC579E2BAB0B9
96CAF22782781DA3923914ABA3A1C32B29C10411820F22BFA0FF00F8280FFC
14367F0E5C6A1F03BE016A663D52167B7D6B5CB76E6DD87CAF6B68E3FE5A83
C4B28FB87E543BF253F022FB17B2BCF704BCD212CCE4E5989E4924F527DE80
32927AB2B3534E9176577C4038F4CE0D52922B8B73895193EA08ACA15A12D1
31B46A89AA412D62099853C5C7AD6A2367CEF7AB3677DF65BA8E7CE369E7E8
7AD73DF68A9224BBB9E20899FDC0E3F3ACEAC62E2E32D98D1F4668BAC06451
BABBEB3D501039AF9AB4ABDBCD3D238EF06D3D0739E07D2BD16C35B040F9AB
F25CD326B49DB54764267B7C3A97BD5C5D4B8FBD5E5506AE08EB5A0BAB0FEF
57CD54CADDF635E73D18EA5C7DEAA92EA3EF5C31D547F7AAACBAAFBD28659E
41CE75D3EA5D79AE7EEF52EBCD7373EAB9C806B226BC693A9AF4E865B6DC87
234AEEF8C84806B02E27C026A396E00158B77778071C93D2BDCC3E1BA221B3
D53E0FC4B2F8E21D564C18F4C46939FEFB828A3F224FE15FA0BA2789C10A43
D7CB1F0CFC33A5E9DE168AD6E706FEE98CD3303C866E8A0FFB2303D3393DEB
BE4FED1D1DB72933423B8EA07B8FF0AE5C2F13E0E7565869BE569D937B3F9F
4F994E9492B9F5BE9DE24040F9EBB2B3F10F030D5F23695E2A0C07CFFAD777
65E26C8187AF6EA6148523E9F83C45D3E6AD58FC443FBD5F385BF89B2065AB
5A3F128C7DEAE59614AE63E811E201FDEA63F8847F7ABC28789063EF7EB58D
E20F881A6F86F47BAD735497CBB6B442EC472C7D15477663C01EB52B08DBB2
41CC7AB78DFE29F87BC09A1CDAFF0088EEFC8B68C85503E692473D1235EACC
7D3B0C924004D7E60FC5CFDA23C65F138CF68F3368FE1E5CE2D226C1917D67
7182E7FD9FBA3D09E4F9EFC41F889AF7C4DD7DB57D5DCA5B44596D6D81F921
8C9F4EECD8059BB9F60007EB5F093C547C171F8CE58FFE25901325CC201F34
47C146C771D4B01CA8C13DF1EFC2861F0518CEBB5CCF457EFE470E231565B9
E3935DA283ACDDAE11722DA33DCFF788AE08AEA1E22D516CAD7324F3B724F4
03B93EC2AC6B3A9DC6A978905B217791847144832793855007524D7D95F0AF
E016A5A5E96B737D0EFD46EC0694819D83B460FB77F53EC057A586A2E4FDA4
CE4A70BBBB3C8B48F0FC3A35847656CBC2F2CD8E5DBBB1FF003C74AB8F011D
457D4D7BF08EFADE2DCD0918F6AF2CD73C2335896DC8462BD33A4F21788567
CD0822BA6BCB568588358B2AD00C8744F106BFE11D49757F0DDEC96372BD4A
1F95C7F75D4FCACBEC4115F6F7C27F8D4FF11A74D02EA1FB2EB60659132629
10632EA4F4C6795273E84F6F85275CE6BEF8FD93BC071695A549E2ABB8C7DB
3553F212395814FCA3FE0472C7D463D2BE478E388D657974F12BE3DA3EAFFC
B73A70749CE6A3D0FBCBC11E1F8ED2142C374870598F526BE8BF0E68CF74E9
1C6BD6BC9FC38AB84515F5778134F8D6D04F8F98F15FE7BE7588AF8CC4B94E
57949EEFBB3EBE9A5146D697E18B781177A6F6FD2B71B46842E0C431F4AEC6
C2D536826ADDD411AAD7ED79578154E7977D6EB4B5B1E74F336A7CA8F15D6B
408BCA69225C63B578BEAF6DE596523A57D3BA98411C99E9835F3778966413
381EB5F8166396FD57152A317748F5233BC6E7906B96EACADC57CEBE37D1A1
BB824475CE41AFA33579810D5E2DE2565657AFB1E1EAF3A738CA2ECD1CF552
68F88B57BA9F45BE7B39C9C0FBA7D45521E221FDEAE8FE2A58096DDEE23189
22F9811FA8AF9E63D409FE2AFF0044FC37E2F9E6797C6559FEF21A3F3ECFE7
D7CD1F138FC22A73D3667B1378878FBD5F637C0FFD93FC69F173458BC57AC6
A03C3FA35D64DB662335C4EBFDF11E542A1ECC4E4F50B8209FCECB2BD84DCC
5F683FBADEBBFF00DDCF3FA57F41FF0006FE3578517C23610C52222450A220
5C60281C003B62BC6F14B8D71396C70F4B0F270551BE69A8F338A5D93D2EDB
EAB4456030919B6E4B6E87C17F1A7F63BF14FC39D0E5F12F873581AFD9DB0D
D342D0F913A281C95019C3FAE383E99AF868DE8F5AFE89350B17F8A3A25C18
2736DA7CFB955D40677EC48078C03C7BD7E1DFED01F003C7BF04F5BB8BAD5A
DBED9A05D4CC6DB50814F92771C849073E5BFF00B2783FC24D7CFF00853E24
E271B3960B35A9793D69C9A517357B5B4B26D765AEE6D98606315CF4D69D4F
1A6BC1EB5525BDF7AE65F50F7ACE9F5203AB57EEEEA1E524743717E00EB5CA
6A9AC456F13CB2B855404924E0003B9AC1D53C4105A42F34F288D1064B31C0
03DCD7CC3E38F88373AF3BE9FA6B18EC7A33721A5FF05F6EFDFD2B8B118B51
46B4E95C8BE21F8DDFC4F7BF64B373F6080F1DBCC7FEF1F61DBF3AF36A28AF
06A4DC9DD9DA959590514515033E81FD94BE1E3FC55FDA4BE1BF80F6178753
D6ECCDC80769FB2DBC827B920E4608863723DEBFB7BAFE5CBFE0947F052E3C
53E38F137C63BF85C58F862DC69F652740D7D78332153EB142307D3CD535FD
11F82FE22C8B751786FC52FF00BF7223B7B920E243D02487B3FA3746EFCF24
1D8F6AA28A281051451400C9238E68DA19943A382ACAC32083C1041EA0D7F3
AFF127C00FF047E36F89FE192A347A6413FDAB4B2DD1B4FBAFDE4201EFE5E4
C44F76435FB55F1C3F697F85BF012C41F175F9B9D5E74DD6FA5DA6D92F251D
98AE408E327F8DC807042EE2315F869FB437ED13E21F8F7E3DD1FC6B77A0DA
E8D0E96AF650C76C6496E1ADA57DE8279588590A3E76ED8D00DEDC1AF3F8C3
C36CC335C92BE369517C949395F6D16F6EF65ABB69A19E0B89B0B86C6430F3
9FBD276B7AED7EC7D71FB21FC51B1F869F15E6D235EBC8ACB41F175BF932CD
3BAC5143776C1A4B777762000C0C91FBB3AF3C57E91DEFED31FB3DD84A219F
E22E84ECCBBB315FC332E0923EF46CC33C74CE7F315F89DA37C3EF147C45F0
CDCDD691A64D3D9409FBEBA2A5608BB7CD21F973EC327D057857887C28FE1D
D4E4D2E77599A300875180C0F7C738AE0FA3261F2ECEE73C831D8974F111F7
A9C6DF1C7572B37A5E36BDB769B6AF6768F11F1F89C04638DA1494A9BD24FB
3E97F5FEB747EDEFC74FDBE7E007C19F013F8BB4ED5D3C6FA8CE5E3B3D2F45
713CF2CAB8FF005CEB95B68C6725E4C1233E5AC846DAFE64FF006A2FDB3BE3
6FED55AEB4BE3DD47FB3FC3D6B296B3D06C8B47636F8C8567527334C075964
C904908114ED1F464BA746472B5C76BDE0BD0B5A8D9353B18AE33DD946E1F4
61C8FC0D7F5766BF47F9C22DE1711AF692FD53FD0FCC307E26C64ED5A97DCF
F47FE67E7823BC6EB246C5594E41070411DC1AFB3BF67FFDAF3C61F0B352B5
B7D52EE496CE261B650C77C78F5F6AF3BF16FC1010AC975E19948239FB3CA7
231E8AFD7FEFAFCEBC06EF4EBDD3EE5ECEFA168668CE195860FF009F7AFE7F
E36F0F2B518BC366946F17B3DD7C9FFC33F23F44CA73DA18A5CF8796BDBAFD
C7F5ABF053F6A9F877F1C7C3F6FA4F8B6686492650A97208EA463E6F43EFD3
D6BE26FDB0BFE0987A578885DFC43F82021D3B509B333DAAFCB65744F2480B
9F2643EABF231EA0125ABF127E1D7C48F177C39D4E3BFD02E9D630C0B4449D
8C3FA57EE6FECB1FB63F8B3E25E8F278134F798CED18F3D369778633C3143C
EDDDD33DB3C60E0D7F3367BC378FC86ABCC30737CABED79769ADA4BA5F75D1
F43E8E9D58D45CB247E0EF887E157C4BF0A5C5EDB788BC2DA958369CCC970D
25AC9E5C657A932052847A302411C838AE0706BFB2BF17DD7C30D2BC171DAD
9BA9D43628746043838F9B7A9E9F4AFE67BF6CBF027857C35F13BFB57C1F04
7696DACAC92CD044311ACE8C373281C2870C09038C827BD7D9F04F8933CCF1
5F53AF4946564D38CB993D2ED6DA3F9BD7431AF86E5574CF8D28AB2F091D6A
02B5FAB9C8368A28A002BF4F7F645FF829EFC60FD9FEEB4EF097C469EE3C77
F0FE22226B69DC3EA365174DD697121CB041D21958A606D468B3BABF30A8A0
0FEE9FE0E7C6CF867F1EFC176DE3DF859ADC3AD6953E15F61DB35BCB804C33
C47E68E45CF2AC3DC654827D56BF886FD9BBF698F89FFB2E78FE3F1E7C36BC
004A162BFD3E7CB59DFDB839F2E6404723928EA43A127070581FEBBFF664FD
A6FE1C7ED4FF000E2DFC7DE01B8F2A78B6C5A969B2B0375A75D119314A0755
38263900DAEBC8C10CAA01F455145140051451400514514005145140051451
40051451400514514005145140051451401FFFD1FDFCA28A2800A28A2800A2
8A2800A28A2800A28A2800AFCC6FF82927ED772FC0BF02C7F0C3C077C60F1C
F8B6124CD1315934ED358B234EAC305659594C7111C8C3BE4155CFE8DF8A7C
4DA3782FC33AB78BFC45702D74AD12D27BDBB948CEC82DD0C923607270AA78
1D6BF8D6F8EFF17BC43F1CFE2A7887E27F895CFDA75AB82F1C5FC36F6E9F24
10A0C9C2C71855F720939249A00F229E6691D9DD8B331C924E4927D69F6B17
98DB8F4154C9DC715ECDE06F839F117C6E91CBE1FD1667B57C62E250218707
B877C6E1FEEE6BC6CEF34A585A5CD566A2BBB76FCCD6953727A2386813A62B
B6D03409753996355DDBB8E95F55F867F632D7A511CDE24D72283A168EDA26
93F0DEE531FF007C9AFA27C2FF00B36785BC34C92C52DCDD3AE39959319FA2
AAD7E438EF13B27A4ECAAF33F24FF5B23BA382A8FA1F3B7817F66FB3F12A21
BAD2E29B77731293F9E335ED53FEC69E1BB7B3F39F428338EBE5D7D39A2AEB
9E1A454D12E05A84181FB989FF00F4346AEAAE7E22F8FE48041713DBCEA176
9DD6D18CF1D4EC0BCFD302B829F8BF96B764E6BE4BFF00923478099F977E35
F823A4786CB9B6D32080A71958954F1EE057CD7AE697F6395900C015FAC3E3
BD2F51F1547219E286191B3FEAD582FE4493FAD7C29F113E11F8D2C9A4BBB6
B2FB745C9FF473B987FC00E18FE00D7D2E5FC6D9762DDA1595FCF4FCCC2786
9AE87CA37502CA8637FF00F51F5AC44BAB9B093639C8EC7D6BA6BE8A6B69DE
0B88DA2910E195C15607DC1E4562DC2248BB5C6457D7D369AB3D51CED1A76D
AE1C0C9AD78F5A27BD79EBC5242728722A58EF59786ACAAE5F07AA41CC7A2F
F6B3B743486FDDBA9AE263BF1EB5645F0FEF5723C05B643E63AA377503DDFB
D734DA8003EF56968DA46BFE27B8FB36836525DB0386651845FF0079CFCA3F
135A53C036EC90730B717C14124D76FE17F05B6A301D4F59578D241981012A
DECE7D3FD91F8FA577DE19F845FD905350D788BBBC5C1540331467F1FBC47A
9181E9DEBB99ECD93A8AFA2C0652A1EF54FB8CA533CCADEEFC4DE1093CC8D9
B50B15FE251FBD41FED28EBF51F957B2F85BE2369FABC2A0C81B3D457252C7
8AE3754F0DC134C6FB4F736579D77A0F958FFB6BD0FD7AD7CEF1070261F149
CE92E591A53C435B9F4A4FA75AEA2BF69D3A5F226EB91CA9FA8EFF00CFDEB2
3FB5354D25FCABF42B8E8EBCA37E3DBE87F5AF0CD1BC7BABF87274B4D7D362
938595798DBF1EC7D8F35EF5A478A34CD72DC248CAE1C77C1AFCEA189CC727
9FB2AB1E6A7D9FE8FA7E4755A33D8D4B5F15F03E7ADB8BC520FF001D70FA8F
855581B9D21F1DF61E9F81ED5CA1B89ED2430DCA98DC7506BEF328CEB098E8
FEE65EF754F7309C1C773DB878A063EF57CBFF00193C7B71E22D553C376921
FB0E9EDBA5C1E249C8EFFEE0E3EA4FA0AEBF55F102693A4DD6A44E4C119650
7BB7451F89C0AF1CF86DE0BD4FE2078856D57718B779973377018E4F3FDE63
9C7E27B57BF4FD95084B1155DA3139ABD5518DD9E9FF00053E19C9E2DD4935
7D463C69B6AD91B87CB232FAFF00B2BDFD4F1EB5EF1F13FC7767A3D81D334F
60B0403000FE26F535D1788F54D2BE1F786D740D2F6C5E54615CAF6007DD07
F9D7C9B64CDE3CF1407BFF009B4F81FE653D242390A7DBD7DB8AFCEE8AAD9C
633DB4D5A0B65D97F5B9E0454ABCEEF63D83F671FD9FECFC53E265F883A8D8
2C524C4B59C4061407EB2ECE80B03F2FB1CF5208FD9EF007ECF56AF631C8F0
0E83A8AF9A7F67F1616E6DCBE06315FAA7E19D6B4A4D2E35122A6057EB9428
A845413D8F66314B43E41F1EFC13B2D3ECDD8420607A57E6BFC59F07C16124
A1500C66BF647E2C78CB4C5B19234753806BF22FE33F886DEE279B630E49AD
0763E05F11DA2C333803BD79FCFC1AF41F135D2C93B907BD79D4EF926802BC
56EF797515A45CBCCEA8BF563815FAC7F0D21B7D374AB3B2B71B628234451E
814003F4AFCB0F0D953E24D30B7417111FC98115FA5DE09D4C0822F9BB0AFC
27C6994A70A3496CAEFEFD3F43D6CAF4BB3EBFF0F5EAA9424FA57D6BE00D62
092D45B9601BA8AF83345D5428539AF60F0FF8B24B32A55F18AFE4CC5E1650
A9CF1DD1F4119687DEF6D7E1140CD36E751057AD7CCF65F142648C2B386FAD
54D4BE25DC4C85564DA0FA57DA4FC4ACCBEABF555B1CFF005385F98F52F167
896DEDA07863705CF5C57CDFADEAA25773BBAD676ABE2792E492CF9CD7097F
AB673F357E7D1C34EACDCE7AB6755D2D106AD7C30DCD7907886F0156E6BA2D
5755183F35791F88355186F9ABED327C03BA39AA48F1EF1E4CB2432A9EE0D7
C757B75F65BE9A0271B5B8FA1E6BE99F196A3E62BA03D6BC721F825F177C7F
6B77E2DF02784F50D7B4AB795ADA49AC6133ED9E28D6464D8997C847523E5E
7381935FD49E16E3E9E0A72F6F3518B5BB692BDD5B73C2CC60E6B447191EA4
3FBD5DBF87BE2978C7C2C9F67D26F8B5AF68D8FDDFF74FA7B76AE42E3E137C
62B19DED6EFC13AE4134670C9269974ACA7DC18F229B17C34F8B52CA9045E0
ED65E49085555D3AE4B33138000F2F2493D057EC59A62B2AC652F6789A906B
7F897F99E5D28548BF7533F48FF664FDB9AFBC3F3C5E1CF1B652CC3F969313
F28C9E03FF0074FBF43FA57EC169DABF81FE2DF86A48196DF52B1BF8B64D04
AAB246E8C390CA7835FCBE8F843F19FC33A65DF8BF5EF05EABA769164A0DD4
B7769240AA8582E4ACA1588C91920631C9E01AFA47E087C77F197C31B882EF
42B87BED281064B466CB20EE6227A8FF0060FE07A2D7F3DF1170C619CE73CB
2A4671BEB1524E37DF469FBB3FEB63D8A5527149544D1F45FED43FB066BDE1
6FB6F8DBE0A4326A5A626E967D1F25EE601D49B63D65403F80E5C76DDD07E3
CF897C7963A234B6D72C7ED5192A6103E70C3A861FC3F8D7F549F063F689F0
5FC5CD1E1686E916E88C3213860DDC1079047706BE5DFDB23FE09CBF0F7F68
882EBC6BE0A10F85BC7854BFDB234C5AEA0C070B78883EF1E82651BC770E00
15EFF0AF8A588C2AFAB662DCA0B4E67F1C3FC6BED2FEFAF9F9635B01197BD0
FB8FE5FBC41E2BD53C4527FA53F9700395894FCBF53EA6B99AF49F8ABF087E
227C13F175CF823E25E8B368DAA5B9385906639901C0921907CB221ECCA4FA
1C1C8AF36AFDA68E2A15E0AAD392945EA9A774FD0E2E5B6814514568015734
ED3AFF0057D42D749D2ADE4BBBDBD9520821894BC92CB230544451C9666200
03A9AA75FB93FF00048EFD8E6F7C45E2487F6A6F88361B744D1CCB1F86E195
4FFA55F2931C97801E0C76FF00324679CCB960418B900FD60FD96BF675B3FD
9BFE007873E1A1443ABAC5F6DD6254C112EA572034E77024308F0B1230EA91
A9EB5D2F8AB464991C15AFA57548832135E1DE3CD4F46F0CE87A8788FC4175
1D8E9BA64125C5CCF21C245144A59D8FB003EBE941A1A1F0AFE22CF7B723C1
BE2498BDF2026D666EB3A28C94624F3228E73FC4A093C824F57F117E357C25
F8496DF6AF895E2ED33C3A194B24777728934807FCF3873E63FD154D7F319F
1DBF6E7F893F10F5FBC83C01752784FC3D1CBFE8BF66FDDEA0EA87E59249D4
9647279C44542F032D82C7E21D4756D4355BC9B51D52E65BCBBB862F2CD33B
4923B1EA5998924FB93410CFE9D3C61FF055DFD93FC333341A3CFACF8A483B
77E9D61E5A743CE6F64B638C8C6403D78045786F88FF00E0ABA3C5BE1FBFB5
F855E0AB9D2AFA5554B6D43539A29123627E7616F18218A8C85CBEDCE09040
2A7F117E167C1EF1CFC5AD5A3B4F0DE9F23592B813DDB8290463A905CF56F6
5C9EF8C57EAFFC3CFD9F3C11F0D6C60BBF113A6A179028DA1C623423FBA87F
9B64F7E2BE7F1FE27F0EE4337573593AB523F0D186B293FEF3DA31EF7D5F44
F5439E47986322A9E1172A7BCDEC979777E9F368F13F0CFC36F887F163599F
C55E22B89E697517F36E350BE66925989E32371DCDC00074006003818AFA93
42F06F80BE185AF9E556E6F947CD2C8433E7D8F451EC315CAF8F7E38E8BE1C
B67B6B2758828C00A7935E0DA369BF147E386A0A9A706D2748908CDCCC0E5D
49FF009669C16E3A1385FF006BB57E619EF1C71DF89557EA1422E8E0FF00E7
D53BA8F2FF00D3C969CDFF006F350BED1477E0F24C8F86E9FD6AB494AAFF00
3CB577FEEAE9F2BBEED9EC5E37FDA866D3EC5FC336DA9BAE9E493F6756FDDE
7E95F3869979AA78F6FEEFC4A636166A3C9898F476072D8F5C7033EB91EB5F
A13E02FD86FC11FD9B13EADA52EAD39C334F7BFBD766FA1F940F6000F5CF5A
F41F147C01B6D0B4E10DA5AA43042A1511142AAAA8C0000E0003B0AFE8BF01
7C13CB785F3AA199E2F12A5529A768C6FBB4D6ADDB64F64B57D6DBFE51C7BE
27BCCB07530786A6D464F56FB277D179DBAFFC37E5D5DE9D2404865C560CF0
F5CD7D2DE38F08FF00674B22EDC633DABE7ED4A0F29D857FA351AD4EB53552
1B33F07A774ECCE2AF2DC1078AF27F19783F4ED7EDCADC461654FB9201F32F
FF0058F715ECF72060D52D27C31AC78BF5AB5F0F6836ED757B7AE11107EA49
ECA0724F615F01C6185C1BC2559636CA924DC9CB4492D5B6FA5B7BF43EA326
C4D68558BA37E6BE963E68F00FC14F13F8F7C576BE0EF0EDAFDA2FAE5F683D
11107DE91CF6503927F0009E2BF7CFE0A7C15F027ECB7E02315B849F5ABA41
25E5E3A8124D201FF8EA2E4EC5E83A9C9249B5F073E0EF84BF67AF093EA37B
B2E35DBC40F7372C30CC40C855CF445EC3BF53CD7CD5F14BE22789FE2AEB77
1E19F08BB0B684FF00A55D0E52253D87AB1EC3FA648FF1C78E7886A71466B2
CB3246FEA71968DE8E76EAFCBB27EAECF45FD4397C674E8C675D5A76D4F3AF
8EBFB465CDE6BEF6DA2EE976EE57740587E247F5AFCE6F8957BAAF8CB535BF
BF0DB6204206E4E5BA93F5C0AFBD353F84EB6B0955889C773C927D49EE6BC3
3C53E04FB387FDDE3F0AFD9785B82B0B96423ECD5E495AE6556B391F07DF69
2D113F2D7333DB9435F457893C3FE4337CB5E39AA59796C78AFB23138A65A8
EAECD1ED26AA11400DA28A2800AFA17F665FDA47C7BFB2E7C51B1F891E0794
4880791A8D8484F917F64CC0BC3201D0F0191C728E01E4641E7BE0B7C0EF1A
7C73F132786FC21100DB955E570C402E70AAAAA09673D80FC48AF5AF8EFF00
B1CF8FBE065BDC5D6A374B7EB6201BA8CC26DE7894F3BB61670CA3B90D9EF8
C648F36AE7186A75961E734A6FA7F5A18CB11052E56F53FAECF831F187C11F
1E7E1B68BF14BE1EDE7DAB48D6A10E15B026B794712DBCCA09DB2C4D957192
3232A4A904FA957F277FF04C9FDB07FE19CBE2BFFC209E35BDF27E1F78E268
E1BB691B11E9F7FF00720BCE78543C4739E3E4DAEC4F94057F5895E91B0514
5140051451400514514005145140051451400514514005145140051451401F
FFD2FDFCA28A2800A28A2800A28A2800A28A2800A28A2803F31FFE0AB3F159
BC0BFB37A7826C65297FE3BBE8ECC81C11676B8B8B820FBB2C484770E6BF9D
DF875F09BC65F14F54165E1CB52B6C8D89AEE5044110F76C72DE8A327E839A
FDC9FF008282F8053E387C6AF08E8177A803E1EF07E9D2BDCC50B02ED797B2
E5E1C82761F2E188B9201C10072495CBF08F83F48F0E69B6FA568F691DA5A5
B2858E38D76A803FAFA9EA7BD7E31E2178B5472CE6C2E0ED3ABD5F48FF009B
F2D975EC7A384C039FBD2D8F9D3E147EC9FE04F04AC37FAAC035ED5940267B
A406346FFA6709CAAFB13B987623A57D7165A1430AAA220000C6315BF69621
40E39AE86DEC49C715FC899FF15E331F55D6C55472979FE8B64BC91EFD2A11
8AB451CE45A6201D2ACFF67A81D2BB08B4D63D16AC1D2DF1F76BE61E319B72
9C0BD82FA5674FA78C74AF439B4E651C8AC5B8B52B9E2B6A78A62713CE2EB4
D520F15C8EA1A3A38395AF58B9B71C9AE6EF2D81CF15EDE131AD7533944F96
BC77F0A7C2FE2D88AEB16092C8A085940DB2AFD1C60FE1D3DABE10F88BF01F
C41E1332DFE8A5B53D3D39200FDFA0F751C30F75FCABF58F50B3520F15E6FA
D696922B065AFD578578DF15836A2A578F67B7CBB7C8E1AF868C8FC5C98919
07AD644CC41C838AFBA3E307C13B7D5FCED67C391ADB6A232CE83E58E6EE73
E8FEFDFBFA8F85B5082E2CEE25B4BB8DA29A162AE8C30CAC3A822BFA5F87B3
DA18FA5ED293D7AAEABFAEE78D5A9B8BB33D47E1C78374DF1769D7536A0F2A
490CDB15A260BF2ED07A10457B868BFB3D787F529133717A464701E3C1F6FF
0057587FB3CE926E742B998A9CCB76C07A61513FA935FA5DF0ABC002FE58B3
1E738ED5F674E841C5368E76CF9FFC13FB287834149DF4937927ADC33483F1
42761FFBE6BE98B2F819F62B3482DACD608631854440AAA3D001C0AFD16F87
1F07AD3ECB1CD3C614003922BDAE5F86BA2980C6883763D2B78C12D9127E26
7887E17CB648D98B18F6AF04F10F86DAD19815C62BF63FE287806DACD25C46
0633DABF3A7E2368D15BCB280BEB5407C697D6C63623D2B9F996BBFD761549
1B15C1DC601A00C4BB8629E368A640E8C30558641FC2B8FF00B26A3A14C6E7
4098841C981CFCBFF0163D3E878FA576B31AC79FBD72E2F054ABC392AC6E86
A4D6C763E13F8A492482CF50CC13AF051F835EC5BF47F11DB85942B13D08EA
3E86BE43D4ECE0BA00CABF32FDD61C32FD0D3F44F186AFE1AB844B990DCDAA
9FBDFC407B8FEA2BF28CFBC3D9425EDF02ECD7DE75C315A7BC7B8F8EBE1878
9EFF004B86D741DB731DCCF1EF0CC14AA67A9CF040383C738EC6BE82F0BF87
F48F843E0F8EDA1C1BE9577339FBC588E5CFF203B0AF3FF057C52D3353585C
48AFB4A9C1F6ED5E7FF16FE23CB732CB0C526E77E383FA0AF1EA55CC317C98
1AEF44F5F3F53E5EBD4A9565CAD68713F113C673EB7A83D941212A49DC7F99
A67842FA3B078D53E555AF258E6667691CE598E49ADFB1BE309041AFD6729C
AE185A4A096A7A7469282B23F463E1B7C491A588FF00798C63BD7D5763F1FD
A0B4082E31C7AD7E3DE97E2B9AD80C3918AEA478FAE426DF34FE75EA1B5CFB
FBC77F1C24BF8DD7CFCE7DEBE28F1B78D5F509243BF39CF7AF38BFF175D5D9
D81CB337000EA49AFA93E11FC094BCF23C41E358BCF9DF0F1DA3731C79E419
07F137B741EE7A73627150A51E690A523E66D0FE1E78EBC7B209341D32492D
D8FF00AF93F770FD4336377FC0726BD5ACBF645F1ADD441EFF0055B5B773FC
28AF2E3F1212BF4BB42F0C46B1A436B0801400001C5776BE0DBE116FF21B1F
4AF9BAD9F546FDC5632750FC9A97F64DF15E917D05F58EB305CFD9E44701E2
68890A727A17C7A57AB69165AE786F6C1AAC2502F1BD4E54FE3FE35F73EA7A
23419591318AF2FD7F458254656406BE6B3DC043328A55DEAB668E8C3E3254
F6385D1F5D0557E6AF40B2D77007CD5E157B67FD8BA8C4CF2B4562D2289485
2E634279655EF81CE3F2AFA135CFF8573A0785D2E20996E4491074BA59092D
9190C3071F86315FCEFC61C393CBEB429CE2E4A7B34B4FF87F23E930989556
3746945AFB63EF53DF5E247DEAF07B2F14C175047716F2078E40194FA83571
BC4431F7ABE76591B4ECD1D3ED4F59B8D6F20FCD5CCDF6B6307E6AF3BB8F10
8C1F9AB99BEF110C1F9ABBB0B923BEC44AA1D76ABAE0C1F9ABC8FC41AE0C37
CD59BAB788C60FCDFAD7986A9AB35C3124F15F7194E49669B4734EA19DADDF
9999DD8F15FB4DFB0568B3E99FB386997D7190358BEBEBB404630825300FC0
98891F5F4C57E1C3437FADEA569A1E93135C5F6A33C76D6F12F2D24B330445
1EECC40AFE97FE1D7836D3E1C7C3CF0EF80ECD83A6856105AB38E92491A012
49F577CB1FAF41D2B3F116AC28E5AA94B793FCBFE1C7834DCEE74528AE5350
61B8E2BA7B9902A9635C4DFCF9635FC5799DA555D8FB3CB29B72398D66C6C7
56B0BAD2B53812EACEF2278668A4019248E452AE8C0F04302411E95F853F17
3E19DFFC01F89F3F863E79340D44B5CE933B73BADCB73133779223F2B77236
B600615FBAD732F5AF9CBF689F84767F19BE1EDDF87D7643ACD9E6E74CB86F
F9657480E1588E44720F91FAF077609515FA87845C70F26C7F2577FB8AB653
F2ED2FFB77AF937D6C7AD9D644F1787F717BF1D57F97CFF33F337C3BADDD68
3A8C5E28D13536D1EE22C34B3282D1B20FF9EA830580F51F301D33D0FEACFC
18FDA96E20BCB6F057C4E896C7519230F0C9B83C17511E92C12F492339073D
4670403C57E03CBE26F10DBDC5D78535C81B4FBAD3E67B7BAB76E2459626DA
C8FF00461D2BED2F811E33D1BC75E1DFF854DE361E749640CDA54E58ACA814
7291C83E6578F92B8FE1CA9E062BFAAF8D30B3A108E362AF18EEE3F128BFB4
9FDA4BAC5E8D755677FCFB2AC146B37464F964F6BED7ECFF004F3EFD3F4EBF
6B2F09FC2DF8F5E0AB9F025D68D67E25BDB75132C85F6DCE9ECEB9592DDD30
E1F182467691C306048AFE667E377ECE1E33F8377B2DD4B1BEA5A097C25E22
1063C9C049D7F81B3C67EE9E30413B47E9B6A7AC7C45F823E247D72F2FA5D5
B44965DC350518961DC785B951C01D838F94F7DA485AFA9748F137803E3A68
A6D355105BEA7731942CC0182E430C15914F1CD71F0DF1662B2C6AB619A9D1
9EAD2D22DF5B7F24BBADBCB63F72C1787B94E3B2D54D5D544AEE6BE38CADAF
346F6942FB5B54BE6DFF00337457E947ED25FB0F6AFE17BCBAD7BE1CDA1551
9925D333904752D6CC7A8EFB0FFC04F45AF50FD88BFE0975E30F8D7269FF00
12FE3AC373E17F0192B34162730EA3AAA632BB411982DDB8CB901DD7FD5800
AC83FA3387F8970999D2F6B8696AB74F74FCD7F48FC2B89F84B1994D654F13
1F75FC325AC64BBA7FA6EBAA3C6FF607FD84BC4DFB5578BE1F14789E1934DF
865A1DC2FF00685D9DC8D7EE8431B2B52304B30E24901C44A739DE554FF5A3
E1FD0345F0AE85A7F867C3765169BA56936F15ADA5B40A1228208542471A28
E02AA8000A83C2FE16F0E7827C3D61E13F08E9B6FA3E8DA5C421B5B4B58D62
8618D7A2A22E00F53EA724F26B7ABDF3E64CFBF198EBF123FE0ACBF1BDBC3D
E17D17E0568F26DBAF12E352D4886C15B1B7931046463912CE8CD9CF1E4E30
7771FB5FACDE5AD8DA49777B325BC10A9679246088AA3A9663C01F5AFE713E
37FC1ED7FF0068AFDA2FC5BF133C69A89B4F0CFDB0DA6990C4C1A596C2CFF7
50B0246D89250A65C105FE7390A6BC4CF388F079753F6B8BA8A3D9757E8B77
F9773E9B87384F30CDAAFB1C0D272EEF68AF57B2FCDF44CFCBBF0EF86BC45E
30D4D347F0CD84BA8DE49D1225CE07AB13F2A8F762057DF1F0A3F641D234B9
20D77E2ACABA83A90DF60858AC2BDF1238C339F6185EDF30AFA5B49D33C01F
0AB461A6F87ACE0B18631F36D1F3391FC4EC72CEDEEC49AF04F881F1F123F3
2D34A6F31F9036F35F8866BC779AE733785CAA0E9C1E975F13F9ADBE5AF99F
D1B94F84D9264941E273DA8AA4EDB6D15E8B46FD65A7F751F5C6AFF10FC19E
00D163D3FC3F0C1636B026D8D23558D540ECAABC0FC2BE3DF19FC6DD7BC537
525968392A4E1A4270ABF53FD073ED5F3ECDAEEABE269BCFD5676642722307
8FC4FF0085753A798E24548D42A8E8070057EE3E0EFD0FA9CE71C6E7D3693D
7956B27EAF54BF17E8CFE46E3DF15A9D1AB528655EF4536949ED6E9EAFF0F5
3BFF0007F85EDEFF00538EFF005C7FED0B9241FDE731A9F653C1FA9CFB62BF
557E01E8960F2DBB4AA3A8AFCBDF0B6A0B6F708C4F422BEE8F85BF10A3D2FC
A6F3318C77AFEF6C4F076130393BC064F45538F68ADFCDBDDBF37767F34E61
9A623175FDB62E6E4FCFF4EDF23F5DF4BB7B486D95540000E2B86F1FCB6316
9537998E9C5784D87C76B68ECC07914903D6BC6FE21FC6BFB7C2E8B2F18F5A
FC2325E00CC278B4E49DAE7B18ECEE9CF0CA8C616B1E0FF16EE2D8DCCDB31D
4D7C67ADC8A666C7AD7AC78E3C5C6FE590EFCE735E310DAEA5E20D520D2348
B77BBBDBC904714518DCCECDD001FE71D4D7F59E11C30383FDF49251576DBB
2492D5B7D12EACF908D2739D92DCA5A5689ABF89F58B5D0341B67BCBFBD711
C514632598FF00203A92780393C57EA9FC25F821A3FC02F0D9D77C48893EB9
7699966382077F2D0F651F9B1E4F603ABFD9FF00E03685F03FC387C55E28D9
3F88EEE3CCB21C110A9E7CA8F3D00FE23FC47DB007CB7FB547ED246E45C786
B4398348E0A92A784AFF00283C75F1E730F107378F0B70B26F06A569495EF5
A49EFE54D6F14F7F8A5D12FDF383F86E9E594FEB98CB73DBFF00015FE7FF00
0CBCFCC3F694FDA2AE2EE5B8D13449B748721883C015F4CFECE7E14D175FF8
2DE19D7F4CB75517D6C5AE187264BA46314EE4F524C88DF4181D00AFC67D5A
E659DE49A7732492125989E49AFD99FD81F528E6FD9E21825555FB26A77B12
9EE412B264FE2E457EB782F0830BC3193D38AD6BB6B9DFC9E8BC977EBF71EA
65FC452C6E2A49690B69FE674DE2BF04C1144FFBB1D3D2BE38F885E1B8A2F3
3E5F5AFD06F885ACD869F672DC4EC001C003AB31E8A3DCD792F85FC0A7C577
CB73736A2E2590E563C6E0BFFD7F7AFC8FC41F13B2FE1DA49E26F2A92DA0B7
7E6FB2F3FB93D4FA448FC81F1A681727CC921B6775E790848FCF15F2DEBF6E
15DB02BFA8FBCF847796163E6CB64A100E700702BE32F8D5FB3C7803C77637
09A969A96F7C41DB7702AC770ADEBB80F9BE8D91ED5F8FE49F49FC2D7C42A5
8BC23845F552E66BE4E31FCFE454A0D6E7F3E9769863596D5EC1F17BE1AEB3
F0B7C512E81AA9F3E26CBDBDC2A9559A3CE338E70C3A30C9C7D0835E3EFDEB
FA73038EA389A31C450973424AE9AEA8823AEAFC17E0BF10F8FF00C436DE1A
F0D5B1B8BBB83EFB234CF2EE7B28FF00EB0C92054DE03F01F897E247896D7C
29E15B53757B74793CEC8D33F3492373851DCFE0324815FBA5F0BFE097817F
654F00C7A85DAC77BE2EB98C4A5E50325F1C3C839DAABFF2CD3B7539249AF8
DE31E3EC265328509BBD59ECBB7F7A5D97E7F26D7B381E1DC662E854AF8785
D457DEFB2EEEDAFF00C168F2AF0B7C2997F650F07E8FAFF87F5674F138905C
DC423FD6B02BC4A3FB847458FBAFFB5F7BD63E137C49F869F1B75AD5EE7E35
EA13CF73A98115B5D310F1C4FD19A58CFDE078503A019F623F3DFE3F7C7BBD
BDBDBA5370D717970CC793924FA9F451FAF415F20786BE24F89FC37AA36A56
D74CFE63EF7427824FA7A579785E1BAF8BA32AF565EFBD53EE7C6C3053926E
7B9FA67FB4EFFC13A352F0F2C9E32F856D14DA7DCE648D616DD65303CFEEDC
67CA63FDC6F97B02A2BF5FBFE09BDF1FB59F8BBF0413C0DE3D492DBC75F0D8
C3A46A514FC4D2DA84C595D11D4F991A142DCEE78D9BB8AFC9EFD9ABF6E6D4
34445D12FA64BCD3E7C2DC69F77F3C2E0F5C03D0FB8AFB21BE29E9C9E244F1
A7ECA37F6BA778A5621FDA9A35F2AC8B7316778B78C9237AE7E6203061C6D6
0735A60B88715809FB0C745CA3DFAAFF0035F89EB64781C562ABAC353577E7
D17A9FB4B457C7BF00FF006C4F03FC5DBB4F0678A603E0DF1DC67CB934BBC6
DA97120EA6D256C7999EBE590241CE0328DF5F6157E8984C652AF05528CAE8
EDC660AAE1E6E9D68D9A0A28A2BA4E50A28A2800A28A2800A28A2800A28A28
00A28A2800A28A2803FFD3FDFCA28A2800A28A2800A28A2800A28A2800AF99
7F687F8DABF0F74C1E16F0D4A1BC4DA9A70460FD8E06C8F38F51BC9188D4FB
B1E000DEC7F11BC75A5FC37F066A5E30D570C9651FEEE2CE0CD3B7CB144BC1
E5D8819C1C0C93C035F90D3EB1AC78BF5FBDF1578865F3F50D4A532CAC3851
9E8AA3270AA30AA3B00057E51E29F1B3CB30DF57C3BFDECD6FFCABBFABD97C
D9DD81C373CAEF646AE9767248ED7170ED2CD2B177772599D98E496279249E
493D6BBEB1B6000E2B0B4E8400B5DEE976A65650057F126618894A4DB67D24
51A7A6E9CD29000AF4AD1BC2935CEDC213F856D783FC306E5959D7815EF3A7
6950DB462389318EF5E5E1701571353D9D2576692924B53CCECFC0602832E0
55D93C0F16DC2E2BD712CB8E94925A6074AFB3FF0088718B54FDA339FEB71B
D8F9D756F084900276F15E61AAE94612DC57D8577689221475C835E27E2ED2
16067DA38AF8BC5606A61EA724CE88C935747CDF7906D24572F771F5AF40D5
A108ED5C25E7535D78764B38FBD8C722B87D4A0041E2BA8F136BFA1F876D1A
FF005FD42DF4EB619FDE5C4AB12FD0162327D8735F35F88BF694F841A7BBC3
0EACF7EEBDADA091C7D0330553F81AFBBC8B26C6E2B5C35194D7926D7DE8E5
AB5231F899BBAE592BAB022BE2BF8DFF000C175885F5ED222C6A36E3E6551F
EB90763FED01F74FE1E98F69D43F698F86D70DB638EFC86CF3E42E07D7E7CD
72975F17BE1DEB8A512FDAD9DFA2CF13A7E6C0151FF7D57EBDC3B95E6F80A9
1ADEC24ADE4FF1382AD4A7256B917ECF7A39B1F09E94932ED9260D33678FF5
8C597FF1D22BF5A7E08DA5AAB40580ED5F969E06D72C92FF00ECF6B711CF0E
7313C6C19481D4023238FE5F4AFBEFE18F8CE3B2F2B2F8C62BFA572DC6C311
423561D7A767D51E3CE36763F613C3BE48D2A110E300738AD97912352EE70A
3A9AF907C31F17E3B6B655F3874F5A93C47F1AD1EDD9566C71EB5DC4177E31
7882D19660AC3BD7E5EFC4DD4E29269707B9AF6CF88DF13FED9E67EF739CF7
AF89FC63E25FB5C927CDD6803CC35FB80D2362B81B87E4D6D6A579E63B735C
BCD2649A00AF33564CED57257ACA9DE8032EE9B835CB5FF20D74572FD6B9AB
D6E0D0073EFAB369F6AB6D661ADEE3CE763323105E37551B18671F295CAE07
F1367B61D1DF4F71833C8D211FDE39AC3D573F7C7F09CD32D2E320572FD5A1
17CC96A4B8ADCEC629AAFC7363A57350CFEF57D27AA158E8D2E98743537DB5
FD6B9F59EA41352B08FAB3F67BF04AEBDAC9F13EA49BEDEC9F6C0ADD1A5EA5
B1FECF6F7FA57E9DF8534F12B470A0EB815F20FC1AD3E2D13C37A7D8A8DAC9
1A97F776E58FE64D7D95E05BE8A2BD859CF008AF8CCD2AB9CD98CDEA7DDFF0
BFE1B69F169D16A5A8461CB8CAA9AF6D7D0B4978BC9FB326DC63EE8AE6FC09
ADD86A1A1DBC70C8BB91402335D65E6A565610B4D732AA2A8CF26BD4CBB0F8
45439AA59BEB73782563E53F8B9E16B3D324696DD42AB8C815F20EB7B5770A
FA77E2FF008E6DB53B978E06CA2F02BE44D67500E58E6BE6395733E5D8E77B
E8797F8A224911C1AF95BC5D36A092C3A3BDEBC5A6B4858A03D09E485F4DDF
9679C64D7D37E21BC055F9AF95FE2248248242A70C3907DC74AEBAB818E229
3A725E9E4FA337A151C649A3A8B0F11C36B6F15B4076C71284519E8146055F
3E2918FBFF00AD7CAB6FE2A9C1D84E08E3AD6B2789277182F8AFCCAB70C352
7747B8AB9F425C78A463EFFEB5CCDEF89CBE42B67F1AF25FED777FBCF9A6B6
A800EB574B228C7A03AA76973AA34A773357397BA90507E6AE72E75855072D
5ECDF05FE0278CFE354C759589F4FF000ADABED9EF9811E715CE62B7C8C3B6
4619BA275393853EA2C2C2943DA547648CA752DB9F5DFF00C13DFE07DE78B3
C64DF1C7C456E5744F0F3490E9624523ED37EC36B4AB9E0A40A48CE31E6118
398DB1FB2F3CDD49AF23F85DE20F04C5E15B4F0C7846CA3D160D122483FB39
06D1028EEBDDD58E4EF3CB1C96F9B35D55FEAE08214F15FC75E24F18E27178
D9D3A9070E5D145F45D3EFDFB6BA687D4E5797B9C5386A9F52CEA57C31B54D
71577739279A4BABE2E4F35853DCE4F5AFC8F91C9DCFBECBF2EE4413CD9CD6
1DCCC306A49EE00CF35CEDEDD800F35E861A836CFA8A147955CFCC8FDB9FE0
B6CB95F8E1E15B7C489B21D6A28C7DE51F2C5758F55E2393DB69C70C6BE22F
0BF882EEC6EAD755D36630DDDA3ACB138EAAE8720FFF005BBD7EE978816CF5
3B1B9D3750896E2D6EE3786689C6E4923901565607A82090457E19FC4BF07A
7C2BF891A9F862CA46974A13335948DD7CA383E5B1EED1E7693DF86C00457F
64783BC4B2C7605E558AD654D7BADF586D67FE1E9FDDF43F2DE32CB151ACB1
54B4527AFAF7F9FE7EA7E8B683E39D3BE20785A1D59022CD2298EE60EBB250
3E75C1EAA7A8F504679C8AF9D35CD2350F87FA9BEB7E066296DBB749A7E708
3D4C07F87FDC3F2FF771D2A97C0B4BD8ED75AD56E731E953402304F01EE158
152A7D5177648E9B856D5FEB1FDAFA789C9CB077463EA5188CFE2066951CA7
EA18BA94E83BD3BD9A7B6BAD9FA6D7DD1DD95712578B8D5A6F96A2EA8FA8BE
137ED13E1EF1AE949E1CF190FB4C2BF265F89E06E87AF3C7A1AFAB7C1BE36F
18FC182351F0CCC3C47E0EBA6DED6C5CE109EA636E4C4FEBC6D6EE33823C4B
F67EFD977C03FB47FECF2FA94327FC237E37F0FEA979676DACDAC7F338DB1D
C2457719205C443CEE395651C2B0F9837932F89BE327ECDBE2C6F875F1374D
FB3DD5D2B0B594E66D3B508C71E641260070B90590ED75E0305C8AFACC670C
E372F8D3CD32F9349ABDD6F1F27DD7A9FA9647C6197E7509E5F8B8454DEF09
7C137DE2F784FCD7FC13F6F3C03F11BC29F12B461ACF85AEFCD0BB44D03E16
E2DDDB3F24A993B4F070412AD8CA9239AF33F895FB467833C0AD2E97A51FED
ED65095305BB01144C3FE7ACDC81FEEA8660782075AFCFFB3F0F7C39B9F867
3F89F53D6FCCF114AE76DA851823DF8EFEDC0E98AF93BC49F146D7C1CF3D8C
92661705A2F5523AAFD3B8F4E6BD3ADE2A6675E8C6861E8C5556BE24EFF747
A3F26DFA1965FE0A64D4F13571389AF3746126B965171D57F7B7947A5E2A37
EE7D57F12FE30F897C6D3B5DF8C352FF00468C968ECA13E5DB47CE47C80FCC
C3FBCDB9BDF1C57C77E3EF8E5A668F1490DACAA5978C03C0AF1FF8BDAA7C5D
D3BE1B5BFC55D4F45B8D33C2FAA5EC5A7D94F71FB96BA9668A5983448D8768
8244D9900DB9200279C7C676B79A86BF7826BF90B96390BD87E14F23F0D715
8EA9F5BCD26DB7DF57FD7E06FC4BE3065F95D2FA8E494D7BBB59592FEBEFEE
7BA6B1E3AF1278EAF0C7148C9031FF00381FE35BBFF0804B69E16D53599232
C60B691CB1EBF77AE7DAB73E1878512E658B299CE2BEEF3F0C1754F85BE28B
3B6843DC4DA45EAC40AE7F78607D9C7A86C115FBBE4193E1B0328468C2C934
7F33711F126373272A98AA8DB77D3A1F957A44E36815DC59DC702BC9B49BBE
0735DCDA5D820735FDA191660B956A7F38E65857767A5D85F9858106BD2746
F17CD6606D7C63DEBC2E0BBE9CD6AC57A4746AFBDC2E62AC7CD56C29F4B27C
48BB11EDF34FE75CE6A9E37B9BA077484E7DEBC606A0DD3352DBBDDDFDCC56
7671BCF713BAA471A02CEEEC701540E4927A0AEE96630845C9B492DD9CDF54
6DD91D7ACFA86B77F0E9DA7C4F757574EB1C51460B3BBB1C0000E4926BF5A7
F671FD9EF4AF83FA10F1B78D9125F11DCC7939C30B643CF9687D7FBCC3A9E0
71587FB307ECDB63F0B7484F885F106246F10CE9BA38DB0C2D2361F757B798
47DE61F41C649E07F69FFDA5E3D2E29741D0E5DD70F945543D3B57F961F48C
FA416378D71EF83F851B786E6E59CE3FF2F9DFE156FF009769FF00E06F5F85
2BFECBC23C270C143EB98BF8ED749FD95DDF9FE463FED3FF00B4BFD984BE1E
D066CCEF95C29E1457E57EABA95C5F5C4B79772196694E598F5CD7747C21F1
23C617AD7D16897B7325C9DDE63C4C8873E8EFB571F8D75FA77ECCFF0011F5
301F527B4D290E33E6CBE638FA08C303FF007D0AFD5BC2ECA383FC3CC07366
78FA2B1925EF7BF1728FF7545372F5696BB2D37E3CE7158BCC6A725183F66B
CB7F37FA23E5FD46E320F35FA8DFB1A78A9B42F8222094EC13EA3772A72395
F9173F9A9AF12B0FD93FC256F87F14F8967BAC6094B644807B8CBF9848FC8D
7B7F87AD3E1E7C3CD06DFC3DA416FB25A6F2A2590B125D8BB127B9249FE5D2
BE3FC46FA5070EE269FB0CB954ACEFBC60E2BAF59F2BFC0F7787B25AB87973
D556D0EA7C59E35FF8487C636FA7799BA0B1412119E0CB267AFD1471FEF1AF
B6FF0067DD43465F305C3289C81B7763A57E48EB9E37D387C45B996C191609
A284A85E00C0DA7F515EFBE17F889259A23DBCE50FB1C57F127887531398E3
D668E16BD9A8BD6C95AC9FA753ECA94F96499FAFFE31F11E89A7E952BCD220
F948032324E2BF3C7C67ACDBCF2CEE840524E2B80BFF00899777B1E2E2E99F
EA6BCB7C47E378FCB73E657C4E3F0D8ACCF1DF5AAD149BB2B2565A1A56ADCE
7C95FB5EE8D67AF784E7BE5506E74D6F3E361D401C38FA6DC9FC07A57E71F8
27C15E21F887E26B3F09F85ED8DCDF5EBE07076C6BFC4EE403855EE7F01C90
2BF443E239D5FE2109FC27E1F8FED37DAA23DBC284E06E9148C93D828CB1F6
15F64FC16F82FE07FD933C0EBABEACB1DFF8BAFE3123BC806EDD8E1883F751
7F813F139249AFE93C978E5641922C35B9AB49BE48F6BF57E49F4EAF4EED7D
1F0B70C4F30A9CD2D29C777FA2F3FCB7EC9DAF83BF07BC0BFB2678112FEFD2
3BDF14DEA0919E403797C70EE3F8557F813B75392493F08FED1DFB475D5E5D
5CC71DC9B9BBB827033D7DCFA28FD7A0A77ED1DFB475D5E5D5CC71DC9B8BBB
927033D7DCFA28FF00EB0AFCE4BFBFBBD4EEE4BEBE90CB34A72CC7FCF4AEEF
0FBC3FAD8CACF37CDDB94E4EFAF5FF0081F96C8FACE28E28A787A7F51C0E89
69A74FEBFE0B0BFBFBBD4EEE4BEBE90CB34A72CC7FCF4AA74515FD1318A8AB
2D8FCA1B6DDD9341713DACAB3DBB98E45390CA70457AAF843E29EB7A0EA11D
DBDCC893236E132B10E1B39CE4735E494572E33034ABC792AAB9D582C754C3
CFDA527667EB7F83FE3F7827E2C69F6DA17C608BFD2E3016DB5CB6016EA123
EEF9B8FBE01FC6BF443E167ED59F143E074367A7FC54964F88BF0F67216D75
FB56F3AFAD50F4F358F3301DD642241CE1D80095F2C7EC59FB227C28FDAEFF
0063CB6BE8DDBC2BE3FF00086A9A8E949AC5A286171964BD885EC07689D02D
CEC560CB200A00936AECAF11F1169DFB43FEC53E2CFF008457E2569BBB46BF
631C53806E746D51075F2A46006EC725182C8A304A80413F9FE2F23C565F37
5F052D3AAFF347E8D84CFB0998C161F1B1D7A3FF00267F4C5E06F881E0BF89
7E1F83C53E03D62DF5AD2EE3A4D6ED9DAD8C94910E1E3719E51C061DC0AEC2
BF9E2F853E3753AC8F1EFECD5AF1F0A78A1806BCD02E1F759DE81C945563B6
45EB853865CFCA41E6BF51BE057EDA3E0FF88FA943E01F88B6A7C11E3A0444
6CEE895B5BB97A62DA66FE227A46F8639010C9826BE8726E2BA3896A955F72
7DBA3F4FF23E733BE12AD864EAD1F7E9F75BAF5FF33ED7A28A2BEACF910A28
A2800A2B175EF10695E1BB03A86AD37951E76A2819791CF44451C927FF00AE
70326BCFE3F15F89F5FCB69B126996EDF74B012CD8F539F907D307EB5F399F
715E0B2E4BEB33D5F45ABFEBD4E9C3E12753E15A1EB545796C8DE30440C9AB
B961EB0C383FF8E0ACB4F891A8E813887C596C25B5240FB4DBA90C9CF578F2
723D4AF3E8A6BC0CB7C4DCAF135153BB837FCCACBF06FF0013A2A659512BAB
3F43D9A8AAB657B67A95A457FA7CC9716D3A878E48C865653D0822AD57E829
A6AE8F3828A28A607FFFD4FDFCA28A2800A28A2800A28A2800A2ABDCDDDB59
42D71792AC312F5672147E66BCF7C47F153C33E1DD2AFB58904D736FA7C32C
F2B4698509121763972BD876CD7163332C3E1D5EB4D47D5951837B23E10FDB
07E22B788BC7B67F0EB4E9B7587875165BA0A787BD986403EBE5C4463D0BB0
3D2BC23488C2A8AF3A8759BDF136B9A8789B543BAF356B99AEA63FEDCCE5C8
1EC09C0F6AF48D35800057F1471DE713C7636AE225D5E9E496DF81F4985A7C
B148EFB4E504AD7AE785ECC4D32715E45A63F2A6BDCBC18C8654CFA8AFCAB1
51D4ED47D39E1AB04B6B28F039615DF5A42302B97D2194DAC4474C57596D20
02BF5CF0B307439B9AA1C38D93B1A8908C5569D0015389805AA37138C1E6BF
A2F3BC560E9E1ACAC7914E326CC7BA0066BC9FC6C50467D715EA1773AAA962
7815F347C5AF1EE81E13D26EF5FF00115EA59585A8F9E46C9EBC0555192CC4
F0140249E95FC75C51FED38DF6742376DD925AB6DF447BD09A841CA6EC91E4
7E2CD4AC34AB5B9D4752B88ED6D6D91A496595822222F25998F0001DCD7E5A
7C64FDB3A69EE6E341F84918112E50EAB3A64B1F582171800766901CFF0070
704C9F1DFE2AEA3F1BEF574B844D63E16B770E9685B6BDCC80F125C6D2738F
E1404A83C924E08F0C5F056930460258C631DCA027F335FD05E1D783746853
8E2F378F34DED0E91FF1777E5B2EB7E9F99E71E276129CDD2A17979AFD2FFE
4780EB5AD6BBE28D41F56F11EA13EA77927596E24691B1E80B13803B01C0EC
2B3841ED5EDFA8F8474C604A402361FDDF97F415CEDAF86EDA6BA164CC1247
384DDC0627A0FAFA0EF5FD0946308C5420AC97446795715E17172E5BB527DF
FCCF3168AAB3C22BD4B57F06DDE9F9F3108C7B5705756CD0B10456B73E99A3
2ACAF751D1EE96FB4BB87B59D3A3C6D83F8FAFE35F50FC36FDA4AE74F962D3
BC62020E02DE4630BFF6D1074FAAF1EC0735F2F48B59B320AD22EC49FB2BA1
7C42D47518124D30B5C46E0157439420F70DD2B7EE2FFC537E841F901F524D
7E6E7ECF1F1E1FE1BEAF07863C56E66F0C5EC817CC3CB58BB9E645F58B3CBA
F6FBCBCE437EC55AF852E2EAD23BCB68FCD86550E8E9F32B2B0C8208E0823A
115F9FF13F1263B09579229462F67BDFEFEBF23D1C2E1E9CD5CF95B5CF0D78
92ED5984EA49F506BC1BC4FE1CF13E9FBA49EDCCB18EAD1FCDFA75FD2BF413
52D124B7C891315E71ACE9513AB02B5E3E078E316A5FBC6A4BD3FCAC6D3C14
3A1F9D171392C735992495F4878FFC056B7624BBB4510DC8E770E037FBDFE3
5F325EA4D693BDBDC2EC910E08AFD2B2ACE2962E178E8FB1E6D5A2E0F52296
4ACA9E4EB524B2D65CD2D7AC6455B990735CE5E49C1AD3B897AD73B772E722
8030EEF0C4EE1907B7AD60E5ECEE5ED9CF28719F51D8FE239ADBB93C1A6F88
2C1A5816FE01978970E0775F5FC3F97D2BEA324E1AA98FC1E2AB51579D2E59
5BBC7DEE6B79AD1FA266156B284A29F52386E78AD18EE3DEB8B82EFB135AB1
5C92401C93DABE49A373B1B459EEE4F2E05DCDD7E82ADF972DADDC515CA756
1EE08CD6D6916634CB0F366E259065BDBDBF0AE2753D55EE350FDDB6162EFE
F59EE49FA7DE05D5545AC401EC2BE84D0F5DF24AB2B74AF82FE1EF8B12E2C6
DE40FF007D54FE62BE89D27C4436AFCD5F238AA0EE62D1F70F87BE29DFE971
858672B8F7AD4D5BE2EEA97F1147B8620FBD7C7B6FE22181F355C6F108C72F
5C0E811CA7AD6ADE2592E999E47CE7DEBCFB53D5C107E6AE42EBC40307E6AE
3753F100C1F9AB68501D8B7AFEAE36B7CD5F37F8CB50139640735D77883C42
08601B9AF1AD5AEA6B97F2E2569659485545059998F000039249E8057AB86A
36D4D228E8354FD96BF688B5F0BD87C46D2FC197BACF8735A805EDBDCE9ABF
6DC42E376648A12D2C7B47DE2E800F5C578236A17B64FE55E42F0BE33B5C15
38E9D0E2BFAD6834293E19FECD12680D913785FC24D6CC626F9B7D9586C255
8639CA641F5E6BF193C5CB12785DA39915D447D1864703DEBF3FE21E20FA9B
A1CF4D4BDA5DE8ED657D3A33D8A3479AFAEC7E675BEB335C4B1DBDBA34B2CA
C151154B33331C0000E4927A0AFA0FC05FB347ED21F13A6897C2DE03D496DA
6231777D17D82D429EAC25B9F2D5801C9D9B8F6009201FE897F63AF0F58787
7F66CF05A5ADAC56EF7968F7523463973713CB302CDD49C3F43D3A0E057BE5
C91935D5C4D8AFA861E1560937257D7A5C9A2B99B4CFC81F829FF04CCD2742
BAB7F117C76D6135EB888875D22C0BA596E0011E7CEC1649403D515631C725
94915F7A6AFA1E97A569B1693A45A4363656882386082358A2891461551140
55503A00302BD9AF650A18D794788AE0157AFC43159BE2317579AB4AF6D974
5E88E6C435CD647CA1E27B9D4FC2BAB2788B436D9736E7257242CA9FC51BE3
AAB7E870472057B7685E2DB2F1468B6BADE9EF98AE572573CA38E191BDD5B2
0FD38E2BC8FC75B248A407DEBCA7E0EF8ADB48F176A5E0DB9722DF5053756E
3B09A3C0900FF79307D3E43DCF3F2BE23F09C71B97BC7535FBCA5BF9C7AFDD
BFA5CFB0E08CCF9312B0D53E19EDE4FF00E0EDF71F5D4B759EF59935D0F5AC
89AFC0EF58B73A88E79AFE77A3836CFDBE308C4D4BBBE00119AE4EFAFF00AF
3556EF51EBCD6C781FC17AAFC40D585ADB068ECA361E74D8E83FBABEAC7F4E
A7DFE8B0196CA52518ABB7B25BB7D8E3C663A14A0E737648C1D3BC2BE2EF19
4A96FE19B17B8579562798F11445BF898FB0E78FC6B88FDA7F47F827E0DF02
A781A6B1B5D5BC4B6CFF0068B8D52440CF14D8E57D181E850E401CF0D823F4
03C4DE25F0A7C18F08B5BC0D1C02DE323AFE7CF5249EA7B9AFC43F883045F1
5BC7579E257DD6DE1B323482DC123ED12E79DBE91E7AE3E83FD9FDCF86F2CA
5816E1397BF1F8E4B5517FF3EE2FACFF0099AD16D7DD9F9763731AB8EA964B
DDE8BF567CDAFE3BF1578987F62E8761232A7C9FBB42B12A8E33BB01557DC9
AEC21B0FEC7D2A1D399C48F1025DC746762598FD3278F6C57A8EB1791C310B
7814471A00A8AA30AAA0600007402BCCB529BE5626BEEE3982AE9469D3508E
F6DDB7E6CAA781546EDCAECFD81FF8269DD993E0FF0089ECBCA60B17886590
49FC2C64B4B60547BAECC9FA8AFAABF686F8471FC6CF84FADF81A278ADB539
9166D3AE654DC2DEF2061244D91CA862BB1C8E763375E95F377FC137F4E96C
BE016A17926EDBA96BF793A654A8DAB05BC3C1FE21988F3EB91DABEFE2715F
D25C3F414B2EA54EA2BA71FC19F0D89C44A9E29D5A6ECD3BA7E68FE5CFC533
7C73F871AB49E0CF1BF82F56D33548DCC4AA6D9E58A66181982688347329C8
C346CC39EB5F7DFEC5FF00B296BD7FAB4FF18FE3FF0087162091AAE8BA4EA7
0E64590B2B9BC9607FF56C8005895C6EC9662AA5509FD81BAB9080815C9DFD
E75E6B8B2DE0BC1616B7B7A71D4FB0CE7C48CCF1D86FAB579E9D4FC4FF00F8
2C578B9CD97C35F04C47293CDA8EA130C7430AC30C5838EFE6499C1F4C8E98
FC7BF0A2AF9C99F6AFD01FF82B56B8D77F1FBC2FA406568EC7C390CBC1C90F
3DDDC8208EDF2C6A7D79FA57E75787AF445229CD7D69F9F33F403E13C90C6F
0E71DABF4C3E1B6A960B6A914A1591970CA402083C1047A57E3B780FC54B68
63F9F18C57D89E12F8A0B6B12012E31EF408FCDEF897E173F0DFE27F897C11
82B16937D34506E392D6C4EF81893FDE8995BF1AA5657DC0E6BE80FDAF6CED
B5FD42C3E25E98019822D9DFE3B853FB990FE6509FF7057C8D61A967033CD7
ED5C2FC43CD4A377AF5F53E0338CAED37647ADC17A0F7AD48EF3DEBCE2DF51
071CD6B457FEF5FA6E133A56DCF91AD979F5DFECDFF0EF4CF89DE3E8ACF5F6
3FD9563B659E30706524FCB193D81C1DD8E7031C6735FBEF17C39F037837C0
BFF12FB3B6B5B4B480B88A28D510055CF000C76AFE6E3E0D7C57FF008573E2
06BC99F6C173B4337F74AE719F6E6BF42A6FDAAEC754D01B4F9B5032C12C65
0A07FE161838AFF3B7E94CF37CC388674711173C3A84552576A316D26E765A
395F996BE9D0FD2F83F0B4696179A31F79DEEFFAF2B1F536B1E36D2BEC6D63
A96A8F35A8276C45F851FDD0DF7B1ED9FA578EDCFC43F859E179DEF74FB0B2
82E7926658D3CD24F525F1B8FE75F957E3AF88BE358EFA68A3D6D24B7663E5
94DC2429D8B0E80FD09AF1ABED7B55BD25AFB509A627AFCDC57C8F07FD1C73
0C7D0588A73B53975BD93EFA6EFEE3D2C6671428CB926F5EC7E9178F3F698D
0E0D5AEA44B956DC72B83DB1C7E55F3FEBBFB53990B2D8EF90FB57C7E3CABB
B98E138CC8C1773F2064F539AFABB44F82FE08F1059D80D3259654214CF264
AC85971B87F7406F40381EFCD7E9EBE8C99365597D7CDB37C54614E92BBDDB
6DED18ECDC9F456F99C584CEAA62B174B0583A2E539BD3C9756FB2460699F1
2BE297C45BC6B2F0AD83CA01C3484ED8D33FDE73803E9D4F615E8DA77C07F8
ADE22DB2788FC490D8236094B74799B1DC65BCB00FE7F8F7FA97C17E0ED2B4
1B0834ED2ED52DE08861510607D7DC9EE4F27BD7B7691E19B8BA03CB8CE3E9
5FC93C45E22D2A35A50CA282A705B3694A5EAF4B2F44B4EECFE9BC83C2DC2C
29A9639F34BAD9B4BFCFE77F923F3DFC4FFB346ABE1CF0F4DAEF84B50BCD63
5BB5C48D0CC531346B9DEB1A80087EE01639C60726BC63C3FF00169A21E45C
B3452C676B2B64104704106BF65E4F05DE451EE68C8FC2BE5DF8AFFB32780F
C797726AD77049A4EAEFCB5DDA108D21FF00A68841573EF80DEF5970F78914
AA37433A8B927B4D2D57935D57A6ABCFA67C4FE18D2AD15572B6935BC6FA3F
34F5D7D747E5D7E423F155644F92426B95D5BC753DDAB6E936AFD6BD2EEFF6
41D56CAE316DE2F56B6E797B23BD7D0604D86ED9391F4AEE7E1FFC2CF067C3
1D7A3D7FC4963FF0974B0237949780082297F865100CABEDFEEB961DC60815
F74B3BC8A9C79F0D3E77D128B4FF00F264923F3DA5E1EE67CFCB560A0BBB69
FE4D9E4FF07353B04D75FC43733E27886CB6C1C3212413203D43718523DFD4
57A87C7693E2278A3C3B75AD786253AD5EAA65E066C4B20F553C062073B782
7A0C9E2BE4CF8C9E34F0AF84FC6926A3E1D8FF00B364B994B4B69028485493
CB24630231EC063D075AF57F86DF1A2DAFA189659C32B01CE6BB733C8713CD
4B39A34EEB47CB25D174D3F34CF12A56C4E5D5AA50A553C9DB667E676AF71A
9DCEA5712EB1BC5E6F2251202ACAC382A54F4C74C76ACDAFD6CF89BF05BC0B
F1AAC9B53B72BA678802FEEEF2203E7C0E1665FE35F7E187638C83F99FF103
E1B78B7E19EB0747F1559984B64C33AE5A09D47F146FD0F5191D467902BF72
E0FE3BC1E6915462BD9D55BC1FFEDBDD7E2BAA3C3936DDD9C1D14515F7820A
28A2803FA2AFF821EEBE2E7C0BF157C2D919D3F52D32F71B4E47DB219A3CEE
E873F66E9DBF1AFDABF18F82FC25F10BC3B77E12F1CE8F6BAEE8D7EBB66B4B
C896689F1C83B581C329E55860A9C104119AFC28FF008218FF00CD6CFF00B9
6FFF007255FBFD401F82DFB4A7FC134BC6BF0D2E6E3E21FECBB35C6B7A4424
CD2787E490B6A36A0724D9CCC737080748D8F9C300032B1C0F94FC2FF1E3C3
9E3BB4FF008423E3AD83CB3DA335BAEA013CBD42CE443B4A4A1802DB48C32B
0C8C62BFA95AF8C7F69FFD87FE127ED2D6D26B1790FF00C237E348D3106B96
31A8958A8F952EA3E16E231D30C4381C2BA8CE7E5339E16A3894E50D247D6E
49C595B0CD467AC4F8EFE15FED41F153E04595943E2A9DBE277C316C2C3A95
BB799A8D947D977311BC28FF0096731CF40B22A80B5F497C60FF0082887C15
F037846C753F87D74BE34D735884C96D65031896D872375E961BA1208FF55B
7CC3C1C0521EBF1BBC6DE0AFDA63F62AF123E8DE29B5DFA56A25E1B5D42206
E748BF006704301B5F1C98A40ADC1C6546EAF0AD0A0B8D67514FB448659AE1
86E66E49EDFA761D857CB3CDF30C042542ABBF66F75FD799F5AB26CBB309C7
11495BBA5A27FD791F64F8D7F6E7FDA7FC7D70CF17897FE11CB46FBB6DA3C2
B6EABC83FEB5BCC9CF4EF263DB935C5E9FFB58FED41A1DDADED9FC45D55A55
1802E5D2EA3EA0F31CE9221E9DD7DBB9AFB9FF0067EFD8F340F14F861359D6
5816917205787FC62FD9EE4D07C6D6BA3F8434F9F55944C84C36D134CDB430
CE420381EA4F02BE5AAE718D72552551EBD6ECFA0582CBA119D28C2378ABB5
647DCDF0AFC67E3BF895A7E8DAEFC4DB986E7578AD91184118862527962101
2039E37918048E000001F75685A6429689C7000AFCD4F863AE4FA2EDD37508
9ED6EAD1BCB96290147475E19594F2083DABEE9F0C7C48D39ECD12E1F6B000
67AE6BC3C2E2E82CD1627324E71FBFD0FCBB13094A9DA9687AA5EDA4491EE4
18AF19F1AAC0F0491B60F5AE8B57F88160632207DC7F4AF03F16F8C6374909
93AE6BCDE28C4E17178DF6981872C7D2D7F90F034674E2DD424F833F1064F0
DFC41FF85777D2E74CD73CC7B407FE58DDA2991941ECB2A06E0FF1818E58E7
ED4AFC5CD77C61241F12BC2373652ED9E1D6B4F6439EE2E13838C6411C11DC
715FB475FBB701E26A4F02A153ECE8BD0F133182552EBA8514515F6A701FFF
D5FDFCA28A2800A28A2800AE2FC4FE3083436FB0D9A8B8D41C6427F0C60F42
F8E7E80727D8734BE36F14AF8674C0D061AFAE894814F2323EF391E8B9FCC8
1DEBC7349B7926737170C659A5259DD8E5998F524D7C3F177143C1C7D8D1F8
DFE1FF0004E9C3D0E6D5EC6C186FF58B8FB5EAB335C49DB3F7573D947402BC
37F6A6BE3E1AF82BAAF92BB64D525B7B253E8259033FE71AB0EDD7F0AFABB4
AD34040EE3AD7C47FF000504BEFB1781FC21A3F4177AB3DC631C7FA3C0E9D7
FEDB74EFF857E31569E271119E2257696ED9E8A69591F08F87DC2A20AF51D3
A6185AF1CD0EE0045E6BD2B4FB9181CD7E4F9AD16E4CEF833D474DB8DACBCD
7AFF0085F5110C88735F3FD8DD631CD779A4EA9E511CD7C7E230FA9BA67DD7
E14D7E19EDD6176C11D2BD1A1B918C835F14687E296836E1F18AF54D3FE203
C68019335AE5D9956C24AF4D84E0A5B9F451BBE3AD51B9BE48D4B48DB47BD7
8B4BF11582FCAC05717ABF8F25981064FD6BD5C6F15E32BC791B2234228F4D
F13F8BA18A278A16FC6BF097F684F8CB7BF17FC7D710D9CE4F87B4795E1B28
D5898E565255AE08E012FF00C3E8981DCD7DE9F1D3E20DD685F0CFC4FAC5AC
856E22B29522607959251E5A30FF0075981AFC74D0EE523F2CF618AFD87C09
E1AA752A56CD2B2BCA2F963E4DABC9FADAC97AB3F30F13B339C68C30B4DD94
B57FA23EF4FD9CFF0067D93E265E2C973F2C0BCB13E95F567C5AFD947C3FE1
4F0CBEA1664168D79C8C5782FECE1F1FB4FF0087B6E239F1B48E45777F1CBF
6B28BC55A53E99A79091B0E707AD7F4E250E577DCFC2952A7ECDF327CF7D3B
58FCE7F1669F1E9D7F35BA745245791EB51A32B577BE21D64EA17524EC7962
4D79AEA972086A989E8E5D4E49AB1EB7E14D5EDBC6DE13B9B5D4086D5B472A
9231EB2C2D9F2E4F73C156FA02796AF06F14592DBDC3AA8C60D3FC15AF3691
E31460D88EF2196093E8577AFF00E3CA2A2F14DF2CF72E41CE49ADCFDD328C
44AA61E2E5BAD0E025C0359F2D5C95B24D5090D5247A466DC2E41AFDF2FF00
8255FC75D33E24783B51FD9F7C68239B5BF09C3F6AD2A571F3CFA5960AF113
DDADA46007FD3375006109AFC0E98F15EBBFB357C5EBBF815F1EFC19F13A09
0C76DA5DFC697C01C07B0B8FDCDD29E403FBA762B9E03007B57166397D1C4D
270AD1BA5AFCFC8D29D471778B3FA58F8CDE0DB1D1E5692DD40561915F186B
21519857D5DF19FE20DAEB572E2D5C3440617073915F1B6B17E1D98E6BF03A
74FDEF7763DFBE9A9E73E2245747CD7C91F1234C54737D10C32F5C7715F546
B9760AB735F39F8E1D2582553C820D7D9E455A54AA467138EBA4D58F9EA59A
B2A79A8BA90C72321FE138AC79E7F7AFD7632BA4D1E431B73375AC499F249A
9A59371AA123550159A37B8952DE21979582A8F76381D6BA356ED5B7F0A345
5F107C42D2EDA51BA1B7737327D211B87E6DB45655D59B5A6AB369B21C3433
3447FE02DB4D7ED5E08E7D470F89C6539BB38C6127E9EFFE56FC4F3334A4DC
62C8B46F851A8789AF84F6120B4B27EA4A963BB3D1178C8FC78AF57B5F8216
DE1EB986FAF249E7F2FE6F2E401158F6E8335F54FECC3ABF80341F1B69979E
398566D2AD8A968CF438F6AFA03F6BBF8ADF0AFC697F6B0FC3CB18EDA28102
B3A204DC7E82BF9BB8AB8A258DC7D6C4E1A0A953949DA0BA2FEB7B595DE892
DA23566A366CFCAEF884D6DA65A799668D013F2942770E7B835E10B727AF73
C9AFA3BC6D6875D81B4DB58CCD3CE76C6A3A97ED5F3B6A5A0EBFA2311AB584
D6C07F13A1DBF83743F9D7665189E6A7CB37A9DF42775667A8F803C5EFA791
652B6021CAF3D8FF0081AFA8743F18ABC6A449FAD7E7E4374F13AC919C32F2
2BD2B41F164F18011F257AA93C8FF115D38AC2A96A6B289F7C5AF8AD4A8F9E
B40F8A463FD657C8561E3398A8E7F5ADD4F15CCE3EF7EB5E6CB0867CA7D197
5E2A5C1F9EB8CD4BC4E64CAAB66BCA1B5B793EFC95525D5D147DEA71C3D86A
27557BA997CB3B57D55FB11FC243F11BE38F87BC47E254F23C3FA248FA842D
29DA2F6F2CCAB430C59FBDB642246E3056365EB9C7C83E13D2AFFC537F35D2
E9F7779A3692167D4A6B68D9960833C9760085CF3CF50033745247ED2FC32D
7FC0BF14BC2FA6786F478A1D2B50B28A38AC7ECA04685631F22A84C6D231C1
1F5AF2B3BAAD509D3A72B369A6FF009742A32B34EC7DBFFB42EA6B67F03FC5
ACC0112DA2C1F31C0FDFCA917E7F3703B9E2BF17FE224BE4F879D7A7C98AFD
33F8C5178F17E09DCF85FC4DFF00133B99A4B432DD42992238AE125C4807F1
2EC197030467201E4FE657C59B195348F22126477C2AAAA92CC4F4000EA4D7
F3C4B38FED1AF45426A4A9B71BAF27E7667D253A32516ED6B9FB7DF05ADE3D
1FE0AF80F498B6E2D741D3518A2ED0CC2DA3DCD8F566C93EE6BB0BBBA55079
AF3B3E29834DB3B7D0B4606FA6B489210DD11762851BD877E3903F1C5422C3
58D58F99A9DCB2A37FCB38F28BFA727F126BE738CBC49A7525CB5257B74445
0C24E51F71597765DD675AB58731C93221F42C01AF26D6F5313AB7959707D0
1AF5CB7F0DE996E30B10FCAAD368D62463CA15F9AC7C427177852FBD9B2C92
1F6A67C29E339AE04523344E073C9535F205CEBADA37C44D07568D8AF93791
AB9079F2E53E5C9FF8E31AFD91BEF0A585C291E583F8578078F7E04F83BC48
0C9A869919994EE59506C9011C83B9707AF63C57D8E4FE29E16A42587C5D26
A324D3B3BEEADB3FF31D3CA654AA46AD396A9A7F71E7D36A7EF58973A9F5F9
A9DE23D0756D15998A99621DC0E47D4556F01783F5BF88BAE2E97A6A94B78C
83713E32B1A9EDEEC7B0FE99AFCD68E556B3B1FB0ACEA8CA0EA2968B73A5F0
3783B58F885AC0B1B1052D6323CF9F190A0F61EAC7B0FC4F15F62EB5ACF853
E0BF840C101483C88CE4E79CF724F727B9A75E5E784FE0A783FECF01484C08
49248DC4E3924F763DCD7E49FC5DF8B9ACFC61F104D0C33345A05B4843B648
F3883F71718E3D4F6FAF4FD2B29C9DE0DB845DAB7DA97FCFA4FECC7BD57B36
BE05A2F7AED7E7D98661531D55462BDDE8BBF9BF222F8A5F1335AF8CDAECCF
2CCF0787A07238254DC153F757D17D4FE039E479C6A9A8C7147E5C402220C2
A818000E8001515EEA515BC4208008E38C6D55518000E800AE16FAF9A52493
C57BB85C226A30846D08ECBF57DDBEACF568518D08D96ADEECAD7F746462C4
D711AADCE118E6B66F2E400699E0AF095F7C4CF883E1FF000069C584DAEDEC
36C5D7AC71337EF65E7B4718673EC3A1AFB4CAB04E738C22B738B135924DB3
FA03FD907C2CFE0BFD9B7C0DA65CC7E5CF756475093206E275091EE977600E
4248ABCF20000F4AFA0AE6F028C0358E92DB585A436366A2282DD1638D0745
4418503E805635D5FF005E6BFA8F0B4152A51A4B6492FB8FCE2A4B9A4E5DCB
9797DD79AE3EFAF7AF35CCF8D7C7BE19F04E8F36BDE2BD4E1D2EC21FBD2CCD
804F6551D598F655049EC2BF3AFE257EDE08D34BA7FC2DD1BED2A0902FB510
CA87B6520421883D417753EA95962F1F4A82BD4958145BD8FCD3FF0082956A
D26A1FB55EB10BC8AE2C34FD3A050319506112E0E3BE64279EC457C53A4C3A
94C54DA5ACB30E3EE23375E9D01AFBBFC61E67C44F1AEA3F10FC5B6F05EEBB
AAB44D3CE6245FF5312C31855030A1511578EB8C9C9C92D8F4D450005C015F
3B5F8B20B4A70BFA95EC3BB3E63D21FC456A4674EBA5C7AC2FFE15E856BE2E
D52C53F7F14B16DEBB95971F98AF5FFECF5F4A824D3D7D2B9A3C5D2BEB0FC4
7EC3CCF24D57C751EA563369F7A44B04E851D0F420F5AF9AEFA36D2EF1A389
B7C24928DEABEFEE2BECCD47C2BA4DFA9175691B93DF6E1BF31CD793F88BE1
25BDCA349A45CB40FD4249F3A7D33F787D79AFA3C9F8CE9427EF5D5CE4C560
79D1E356BAB7404D6E43A98C7DEAE435BD1B50F0DDFAE9FAAA88E571B90AB6
55D738CA9FAF6383ED5523B92BD78AFD7B2EE21E7829C25747CA6232D49D9A
3D1D752C8EB5622D667B71882568C7A06E2B808EEBDEAD2DCFBD7A18AAD431
492C4538CEDB5D27F99C90C3CA9FC0DAF43AE975292562CEC589EA49C9AAAD
72CDD4D608BAA1AEABBA38D518A8C744BA2D8C7EADADCEBF43D2EE3C45ABDA
E8B6A70F74E1491FC2BD59BF01CD7EA17C3DD0AD345D2ED74FB65C2408AA09
E49C0EA4F727A93EB5F037C04B34B9D72EF559067C855890FBB9CB7E3803F3
AFD0EF0E4EAA895FC45F495E2CAB8AC6472E84BF774D26D779357BFC93B2ED
AF73FA37C20C9214683C5C97BD3FC97F9BD7EE3E90F0368C351BA8E3C64645
7DFDE06F87BA75A58C53DC441DD803823815F0C7C2FD4E0B7BC8CC847515FA
2FE17F11D85CD84452550CAA0104D7F2E70CE0B09531ABEBAFDD3EF38FF1B8
9861E11A17B3DEC59D4FC27A5CB6ED1B40AB91D40AF8F3E236830E9F3CA883
1B49AFB0F5EF1569F636CEEF329600E0035F117C49F15457B712956C9626B5
E39A197AC45B2FF84F33C39FAE7B494EA37CBE67816ADB41615E53AFA2BA35
77BAA5E8666E6BCD75ABA186E6BCBC9E8C9491FA563EA2699F05FED2FE0387
58D21BC43671817DA602C48EAF0F5753F4FBC3F1F5AF86746D7B51D0AE05C5
8CA5467257F84D7EA9F8EBCBB8B59A27019594820F3906BF2A75ED3FFB2B5A
BED387DDB799D57FDD07E5FD315FD99E16660EBE0A584ADAA8ED7ECFA7DFF9
9FCEDC73828C310AB47ED6FEA7D85F0C3E3910D1DBCF2EC75C6518F3F87AD7
D86BA8783BE28787DF41F13DB457F67700651FA86ECCA7AAB0EC41C8AFC678
E49219165898A3A9C820E08AF71F017C5ED43439E38750908008C49DBF1FF1
AE6E2EF0D2355FD6B02F966B556D1A7E47C31DA7C68FD99B5CF00BCBAF7845
9F59D01B2E4019B9B65FF6C0FBEA07F181F503193F2BD7EAC782BE2D5B6B09
1ADC480E40039E315C0FC54FD9C3C3DE398E6F127804C5A5EAEC0BC96E3E5B
7B86FA0FF56E7D40C13D464935C9C37E23D6C2CD60B3C567B2A96FFD297EAB
E7DCF49E5EDC14A0EEFAA3F3968AD8D6BC3FADF877559743D72CA5B3BF85B6
B4322E1B27A63D41EC4641ED9ACA92292191A2990A3A1C1561820FA106BF6B
A756338A941DD3D55BB1E734D6E7F43BFF00043AD1C41E13F8B5AFF9720379
7BA45B6F20F96DF668AE5F6A9C6370F3F2DCF42BF8FEEE57E537FC11D7C18D
E1BFD92A4F11CD1157F166BD7F7A8E460BC36EB159A807032A24824F5E49E7
B0FD17F107C55F00F86CBC57DABC535C21C182D8FDA26DDE8563DDB4FF00BD
81EF5CB9866986C2537571556308F793497DECA853949DA2AE7A1515F39DD7
C79B89EEA31A4E81243A7EE024BABE95622173C958973BB8E465C7B8F4EE75
1F1468D6F6EB73AAEB02489807508E15483C8E13191F526BF3B978C192D494
E182A9ED651DF9745F7BB5FF00EDDB9DEB2AACADCCAC7E6CFF00C14EFC21F1
57E2AF887E1AFC3DF86BA05F6B90DB8D46F6EC5AA6E8565730450199C80911
55F336977190E7D2BC23E10FFC1383E30DCCF6FAA78EF55B2F0CC2A4318509
BCB918EC42158867D448DF4ED5FA5BE2CFDA5BE1CF83E1758668815CF52073
F415F15FC49FF82846996224B7D2E719E4000E3F4EA6BE2F33CE71D99D573A
34B953FEB77FFC89F4B80CC2AE1A92A54DA5F99FA07E0BF87DE1BF85FA047A
3DC6BD71731C4B8632C8109F5C08C2E07B66B2FC41F1BBE19F826193CA9615
23249185C9F53DCD7E14F8DFF6D0F1C7899E44B19DA28DB3CE71D7F5FE55F3
46BFF123C4DE20677D4F5195F7E720391FAE73FAD72D3E17A9369E26A7C96B
F9E8BE5637786C4E21F3CFAF56EDF86E7EACFC4DF8F1E18F147C4E975CF0FC
8882E618C4FB3A34B1E5437E28147E15DAE87F14E3312FEFBF5AFC3BB5F11D
D691722E2D2E194E79DCC5B3F524935ECFA07C63D46245495F07FDEE2BE927
C2D09D28FB377B77DCF2311CF426E13FF807EBFDC7C518FCBE66FD6BC9BC59
F1522113FEFB9FAD7C26BF15AF6E531E76DCFBD6AF84ADBC6DF153C450F853
C07A65CEBDABDC7221B752DB172019246E1634048CBB90A33C9AC30BC2B696
C633C5E87D55FB3EE89A97C5FF00DA0FC3163147E6D9E91749AB5E1237225B
D93AC9F38E98924D91F3DDBD335FBD75F2D7ECB5FB3669BFB3F784E5FB7CA9
A8F8B35908DA95E264C63664A4106402224C9C9232ED96381B557EA5AFD732
4CBBEAD4791EECF1711579E570A28A2BD7303FFFD6FDFCA28A2800A28AC8D7
F513A4685A96AC0806CADA69F27A7EED0B73F95294ACAEC0F993C51AF9F11F
8CAF675626DACD8DB423391B6338661FEF364FD31E95D968215E48D6BC07C3
3747CA46762CC79249C927D4D7B26897E237539E95FCEB9ED6956AD2A92EAC
F5E92B2B1F40D8C6BC0EC2BF38FF00E0A4B288B46F87E7A66F2FBFF45C55F6
1F89BE34FC3EF005AA4BE28D5A3B7B87195B68C196E1FD088901600E3EF1C2
FBD7E75FED6BF18345F8FBE1DD1743F0AD85DD8B68F7C6E4CD7891A8911A26
42008DDC8E483CF51D7078AFA0C366D86594CF0AF493FF003465283F68A47C
8BA25E0DABCD7A4D85DF0066BC7A0B3D434803ED285957BAF35D8697AA4532
6E89C30070707A1F71DABF17CDB2D9AF7ADA1E853A899EBD6979D066BA4B5B
F2B8E6BCBAD6FB81CD6FC17D8EF5F1B88C19D2A47AB5AEAE57186C56F45E20
900E1ABC7E3D400192718A26D66F82A2691A6DE6B3712C823486C61699D9D8
123A70071C926B0C2E4B5B115152A107293E88CF138BA7469BAB55DA2B767B
2BF88A523EF573DAC78BEC74CB7377AA5E456908E37CAE1173E9927AFB5733
75E03F8896F64356F89BAF697F0B34861BB17522DD6A4EBFECC23804FA1191
5E31AEFC6FFD98FE18CCD73E17D0EEBE23EBF170353F10484DB861DD2DFBAF
A035FADE47E06632ADA78EA8A9AECBDE97F92FBD9F9F665E245083E4C241CD
F77A2FBBE2FC17A977E296A9E23F899F0C3C5569E00F0F6A3AD5A5BD94B717
17A90325B4715AE2690867C173B50E140C9ED93815F999A4EA80C6B835F417
C59FDB77E2878FAD65D1AE756FB0E90CAC834FB0516D6C23C72A5131B863D7
35F1469DAEAC523440955C9DB9F4ED5FD03C2FC2785CA70CF0D84BD9BBB6DE
ADD92FD3A1F1F8BAD8CCC53AD898EDB59595BF1FCD9F44DAEBB342004908FC
69D71ADCB28F9DC9FC6BC920D70103E6A9DF5A18FBD5F43ECCF9D7943BEC76
B75A90C1E6B8CD4F511B48CD63DCEB1BB201CD73D737A5B258D5C627AB82CA
EDAB2E59DD32EAD0CC0E0A16627F03566FAF0CCE589EB5816ACD969DB8DC30
BF4FFEBD4AD266ADA3F47CB28F25249F515DEA9C8D4E67AA923D33D02199B8
AC3B939CD68CD275ACA94E7356901FB65E01F1BDD788BE1B785F53BB72D34F
A6DA7984F79044A1CF53C16071CE714DD4755183935E0FF07B517B1F855E1C
B691C122D437073C3B3301F8035D26A3AE0C1F9ABF17AF81B569A5B5DFE67B
4A7EEA2CEB7AA0DAC735E07E2BD43CDDC80F5AE9B5CD746186EAF24D5AF8BE
F773CD7BB97E16DA985491E5FAC4C12F6403D6B9F9252D526A173E75E4B267
20B7F2E2B39A4AFD2B0EAD4E29F63CE96E3DDEA84D2714E924AB9E1ED0EF7C
55AE5AE87603E7B86F99F1911C639673EC07E678EA6AEAD48C22E727648491
F4C7ECEFE1F36BA75FF8AEE130F7ADF67809EBE54672E47B33E07D56BCBBE2
8D83693E39D4B60C2DC482E50FAF9A371FFC7B22BEBAD2ACED344D2ED748B0
5D96F691AC683BE146327D49EA4FAD782FC73D2CCAB61E208C67666DA43EC7
2C9FAEEFCC5797E12711DB88E6EB7C15E2E16FCBF2B7AB2B30A37A3A743CFF
0044F18A889559F6BAF047A115AF77E2E56439933F8D789CB6EB2B7988C639
3FBC3BFD45224377D1A618F5C1CFF3AFBFCC7C16CCBDBB58387B4837A34E2B
EFBB567F8799E32AD4F79687BD7C32F8B72FC38F887A7F8E208EDE7B8D384B
E525DC4B3425A543136E46C83F233633D0F239E6BEEAB3FDA47F67BF88A9E5
FC4AF8776F693CBC3DE68929B66C9EA4C47286BF2BEDA28A3041F9CB752DCE
6A75B6801DD0E616F58C95FD3A7E95F6957E8B99857C2C2AC3110556DAC5DE
CBC9495EEFE495FADB539278E85ED667E9A6A5FB337ECA5F15019BE1FF008D
E0D2AF25FBB6DABC26CDF27A0F3A1F90FD48AF3D87FE0979F17B59F114161E
17BE83EC130771A80992EAD5422961F3C2430DC70A06DCF3CF7AF88ED754D7
AC0836B79E601DA4E0FE63FC2BED2FD96BF6B4D6BE0C6A9A9CFAACB320BC48
634C379918552C5FE53903276FE55F91F177859C5D91D09D55859CD2FE5F7D
7AE97B77D6C7A397D684EAC62E765E6723E38FD83FF691F8708EF77A13EB70
C5D6E34DC4E1BFED90C4BFF8E57CCD7FA5F89742BB7B0D56D24B5B98FEF453
234520FAAB8047E55FD15F80FF006FDF02F8A638E1D5CDB4E5B00F211BF235
ED97BABFECD9F196C459F896C2C6ED241C25DC31CAA09F4DC081F5AF81CB7C
66C0C5C68E7995AD346E94E7427EAD54F6F49BFF00B7A9A67D155C9A4D5E94
FF00292FC2CFF067F2C2B3EA647DD0BF56FF000AF64F84DE05D17C617AD378
9AEE4314122836F09D9BC119F99FAE0F4C000F1D6BF6DBC6DFF04D8FD9FBC7
11C97DE06BA9B439A4E41B2B8DF113EF1CBE6003D976D7CA379FF04E7F8CBF
0CF549F50F09EAB69E20B161CA3AB5ACE7072B819910E3D4BAFD074AFD5B3B
E35F0E713C3B889E5956BD0C6249C7DBABA7AAE64AA5252A5F0DED7927B6C7
90B2FC6AAA949271F2FF0027AFE07A97C333E1EF0B68D0689E1DB38AC2C63E
9144A0025BA963D598F72724F7359E9E09BBF865E2C3F127E125A9951B2F75
A3C233B7711BE5B2403E538CEE8870467660FCADCB787FC27F11741D56DB44
F1268D73A6CB2B6DDF22E62E3AE245CA138EC0D7E837C33F05C090C688BBDF
00B39E49AFE2AE28E20A78684A3097B4F6AACD277524FBB5B9EFE59964E53E
697BA91BFF0009BE33F85BE2C68518F391AE1D76B23773D0AB29E87B106BC7
357F835E01F12F8F8EAB63035C69F68F9585B06DDA753CB2AE3951DB3C13C8
18C1AD3F89FF000A7C1FA2FC48D3F5EF0DC93E9BAE5F46D26A50DB154B6B98
DB2A92CC98CF9C5B3F3A904A83BC1F948F69F086991C10C6AAB8C015F89E75
55E5D372A0E519496CF757E8FBF93DEDF72F7A92E6F43ACD0F45B7D3EDD02A
0181C715D2AAFA511A7402B5ADAD0BF6AF92C9B23AF8FAB682BB63AD594773
3446C690A115D4AE9A76E7154EE2CF676AFD131FE14E330F47DACE271C71D1
6EC73E4763542F2D9268C82326B5E58F69C55571C57E4F8EC24A8CECFA1DF0
9753C4FC51A2452A382A0E6BCB7C35E39D57E116A7712DB43F6AD26EC93340
00DE8E463CD8CF6600723A11E87047D09AF44A55ABE7AF18D92491C808F5AF
B0E18CD2A529C651766BAF55E6BB3ECF742AF4D3563E34F8E1F1A359F8A7E2
2B8D260924B7D1ED58899B952E7FB833CFD7D07D457875CEA70DB42B6F6E02
4718C2AAF00015E81F167447D3AE1F52B51819FDE01FCEBC065BB2D924E6BF
68CBF054DD182A4AD15D3CFAB7DD9E865F38A836B7EBFD7635AEAF9A424B1E
2B0EE2E78EB5526BBF7AC3BABD001E6BE830F84EC6D52A8B7F798079AFD03F
F8278FC3196F7C41ABFC69D550ADB69CAFA6E9B91F7E790037128CFF007108
4047077B8EAB5F9EFE14F0BEBFF127C61A6781FC311F9BA86AB308949CEC8D
3ABCAE47448D41663D703804E057F415E00F0AE85F0C7C13A47817C369B2C7
498044A4FDE91C9DD24AFF00ED48E59DBB649C715FAB703647CD57EB135A47
6F5FF81FE47CCE718CB47D9ADD9EB171A8F5C1AF917F686FDA9BC37F072D9F
46D3C26AFE2B9D018AC837C90061F2C972C3955EE101DCDC6300EE1CCFED39
FB4AC1F083464D0FC3ED1DCF8B354426DE36F996D62395FB448BDF904229E1
8824E42907F1EA7BBD475AD46E759D66E64BDBFBD91A59E6958BC923B1C966
279249AFBBCE739F62BD9D3F8BF23E7210B9D778DBC77E33F8A3AFB7893C6F
A8BEA1764158D4FCB14084E764518F95147B724F2C49C9AC782D00ED4B6D0F
02BA2B5B269301457E7788C44A4DCA4EECE948A11DB0F4AB4B6FED5D9D8786
EE2E31843F9574F1782AE8AE7CB3F9571399763C98C1EA2A178057A9DDF84E
E2107286B93BCD2A4849C8C50A407152C02B267B715D5CF095ED5933475BC2
64347CCFF1D3C2CDA878763D6AD9732E96C59FD4C526037E4707E99AF9360D
4AE61F949DEBE87FC6BF4AB51B282EEDE6B5B94124332323A9E8CAC3041FA8
AFCF2F1BF8626F08F892EF46901312B6F818FF001C2DF70FF43EE0D7E89C29
99CB95D14ECD6ABD0E3C4D24F568AB1EB317F1AB29F6E455C5D66DBFE7A63E
A0D723457DD4336AC8F3E582833AE6D76DD7A316FA0FF1C5675C6BB70E0AC2
367B9E4D61514AA66D5E4AD7B0E383A6BA1F647ECDB75E5E8F792C87323DE3
649EA408E3C57DE7A06A4362735F99DF0335B16897B605B1B255947BEF1B4F
FE822BEDFF000FEB8A514EEAFE53F1332A9CF1B526FAFF0091FB7F07639470
D08A3EB3D03C44D66EAEAD8C57BA68FF0015AEED2254129E3DEBE27B0D6F81
F3574F06B9C7DEAFC2F1593C933F4AA58D8CA3696A7D61AB7C52BCBC42BE69
E7DEBCAB54F114974C59DF39AF2F3ADF1F7EB3EE35A18FBD5C70CA1B7A9D1F
5D8C55A3A1D4DF6A8083CD79E6B3A90DADF3553BFD6860FCD5E75ACEB636B7
CD5F559564CEEB43C5C7660AC731E2ED44346E33C9AFCEDF88A1478CF52DA3
0374678F78D735F68EBBA919D98E7815F0E78BEF05FF0089B51B90720CA541
F64F93FA57F4A786B8274A727E5FAA3F1FE30C429C63EBFE6737451457EB87
C11D57873C5BA9F87675681CBC20F299E9F4AFB2FE1B7C688AE1634926CE30
082791F515F05575DE04BDD3EC7C57A7CBAB02D64CE525018AE43A951C823A
120FE15F21C53C2F86C6D09CA71F7926F4DD9D586C44A0D58FDE5FD98FC396
3F177E25D8DBC10A1105B4CF7373E58768A00A70371E9BA42A073D4F7E95DB
FED31F02BE1F6AAB79E19F1158DBEA1E4AFEEE7DA16688919051C7CCA47B1F
AF15F3A7C11F8ABA67C35D25E2F04CE74FFB581E73AB932498E9B989C90327
03A0AB9E39F8CC2FA39A7B8B9334D264924E4935FC5998E578AFED2854C0C6
70941E92BB4EDD345B59DFAB3EAE9D48F23E769A67D47F0E7C5DF0F3C1BF08
BC3DE0EB0BFBB8F44D2206862D292629042E2466943F3F3B190B36F6058EEC
F7AE43C4DFB50F867C391B43A2C56F64147054067FCCE4D27ECE9FB0F6B1F1
CFC1563F12FC71E2FBCD0341D6A5965B7D3AC604173240AE63F37CF94B2C7B
CA92A3C97CAE1BA115FA39F0EBF632FD9C7E19C90DEE8FE0FB7D4B528B07ED
BAA96D4272E3A3A89CB471B7BC4895FA9E07C06C466357EB99A62252E6D7DE
6E4ECFD6E79F3CE634D725389F92BA678DFE3EFC6F9C27C2FF00096ADE218A
62556E96268EC81F46BA9764087D9A41FA1AFAEFC1FF00B207ED19E2DD36D2
0F89FE33B4F0C5A4512A7D9EC035F5CE14602B3131C68474CAB4838E95FA9D
1C71C31A430A048D00555518000E000074029F5FB470D785B94E58AF4A9DDF
767955F32AB3DD9F0D49FF0004FBF81373E11D5B40D49F54D4B56D4ED26B65
D5EE6F1CDCDB492A90268628FCB803236197746C78C124139FE67FE3C7C17F
897FB34FC46BCF879F12ED4ADC2665B4BD4DCD6D7F6C49093C0E7AA9C6083F
32365580208AFED16BCB3E2F7C13F85FF1E3C28FE0BF8ADA041AF6985BCC8C
4B949A09718124132159227C71B91864654E54907EE2A65D49C7963148CF0D
8D9D39735CFE2FA1D7411D6A66D65987CA6BF70BE21FFC116B41B8BE6BBF84
FF0011EE34DB5627169AC59ADD328209E2E2078780700031671CEEC8E784F0
BFFC116BC6D34EA7C69F13EC2CA057F9974FD3E5BA76418E034B2C0149E467
0C070707381E4CB25D7447B90CF34D59F8E5035F6AB79069DA7C325DDDDDC8
91430C4A64924924215111172599890001C93C0AFE87FF00634FF8270F82B4
1F85371AAFED21E1F8F59F12F8A8472FD82679233A4DB004A44AF1323ADC3E
EDD3306F970A830558B7D51FB39FEC1FFB3FFECD7731EBDE17D325D73C508A
47F6C6ACC971731EEFBDF6755548A0EA465103953B59D8673F66D7A584CBA3
4F591E5E373275572AD8F8A2C7FE09DFFB22D85FBDF27825E504A14865D4F5
078A3DA3070A6E3E60DD4872C3D3038AFA9FC15F0F7C0BF0E34AFEC3F01681
65E1FB1C82D1594090876031B9CA805DBFDA6249F5AEC28AF4234E2B64798D
B0A28A2AC4145145007FFFD7FDFCA28A2800AE3BE226EFF857FE26D8096FEC
CBDC01C927C87AEC6AA5FD941A958DCE9D740986EA3789C0E0ED7054FE86A2
A479A2D0D1F9B9E1BD4C7949CF615CDFC4EF8CB73E1655F0D7861C1D6A750C
F290196D51BA1C1C82EDFC20F00727B03CAB6AD2783DB54B5D68F972688674
B81E86DB3BFF00F4135F2BE93ABDD6BFAA4FADEA4DBAE6FE532B9EC0B1E146
7B28C01E8057E0D9DE1DD26EE8F4A123D83C2BE0CD6BC5B7E5E2592F6EEE1B
73C8E4BBBB1EA598E493F5AF5ED5FE01F8A749D38DEDD5AB2A0193C74AFA63
F65DD3B448F487BBC21BA206338C81DF15F4678FF58D1F4BF0EDDC97CE815A
36001EE48E2BC1A394FB5C2D4C54AAA5CAF48F5657359DAC7E277887473692
3C332E0AF06BC335E82E34DB8FB7E9CDE5CCBF930F461DC7F915F4BFC46D52
D6E755B87B7C6D2C7A57CE3E23B84647A597C5BF765AA612674BE10D6A6F13
5B2B5840EF3AB797246AA4959000481EA0E722BEAFF067ECEBF147C4B02EA1
7B6A9A069B8DCD73A83792A17D769F9BF1C62BE52F817FB4F1F80B65E21D32
DB4BB29EEF51923B8B6BCB98F7C901552922AFFBC3611F43D6B2B5FF00DA17
E3CFED07ABBE93E14B6D57C55396C18AD23736F167FBFB408E31EEC40F7AA9
F0861E339549A7CBBEAD28A5E727FE5F32BEB32B58FBC6FAD7F665F84C864F
176BB2F8DF5587936D68765A861D8B0383F50DF857CEFF00173F6F1BBB7F0E
5EF863E185ADA7836D59088CD928F3F23A664C0C13D09001E6B98F0AFEC31F
1A7C5EC9A87C55D7A1F0DDA31CBD9D991757647753203E4A1F70641ED5ED10
7EC65F0BFC230A49A5594B717D1722EEF1FED1316ECD861E5A9F7445AE7871
7E53819A542A2725D20B4F9C9EEBD1B5E47357C3FB684A9D4D9AB33F2801F8
D1F17AF8EA305B5DDD2DC1CB5F5FBB244413CB0790E5FF00E0018FB57AB786
BF662D3CB25E78FF005A975193826DAD330C39EE1A46F9D87D021AFB2FC4FE
0AD73C38CCFB7ED102FF00128F980F71FE15E78DA81E79A79BF89199E295A8
354E2FF977FF00C09FE963872CE16C161D7BB1BBF3391D7FE1D784ECFE1E78
83C2FE13D1EDEC9AFEC27894C69991E4D84C7BDCE5DC8700FCC49AFC958CAC
C8013823A11D457EC749A87FB55F96FF00163C2537857C7BACDBD9C45AC1E4
FB546CA32B1C739DC14FA6D24AF3E80F7AFADF093329FB4AF87AD26E52B4AE
DDDBB68F7F91E8E61463CAACB43CFD6E6EEDF8724AFA8AB4BA9823E66AA914
F560189B96407F015FB733E6EAE5906EE894EA01BE55F98FA0E4D6D685650E
A577FE9EDB638CF31F76FAFB56423A2FDD000F6A7ADC3C32ADC4270E9FA8F4
A9F439B11953F66D537EF1E93AEE8C92C3F6AB2501907DD1C0207F8579EB3D
7A4E85AC457F028279E8456078A34468776A368BF21E5C0EDEFF00E359C5F4
678F916692A52FAAD7F97F91C6BC954E492A279AA9BC99AD523ED0591F26A9
3E58ED51B89E001D493523356CF85AD45E6BB6CCE3315BB095FD3E4E40FC4E
3F0CD2A93E58B931A5A9F73E8DA8A689E1FD3B475217EC56D14240F54400F4
F53CE6B2B50F1196C856C9AF316D66493EF3D549354503AD7C3AC0EB77B9DE
E674F79A8B484B3B570FAD5E5CBDADC4B6D13C8B028691954B08D490A1988E
83240C9EE40AE93C23E19F177C4BF1141E14F03E9D2EA9A8DC73B231858D07
592573F2A20EECC40E83A902BF51BC33FB37F87FE1C7C3C9FC2BA82C7AADFE
A899D4EE197E499C8236283C88D01C2F42796382703E2B8CFC45CBF2174E9D
6F7AAC9AF713D796FAC9F65DBBBDBAB5B50C2CEADDAD8FC45696ABBCB5EA1F
19FE1ADE7C2DF164BA512D269D73BA4B495BA94CF28C7FBCBD0FA8C1EF8AF2
0852E6FAE12CECA369E794E15106589F602BF65CB733A18BC3C3158795E125
74FF00AFC4F3270717CAF725065B8992DEDD0C92CAC15154659998E0003B92
6BEC2F861E0A4F06E966EAF803AADE80663C1F2D7A88C1F6EAD8EA7D4006B9
1F87BF0FADBC341358D6009F5523E55C829003FDDF56F56FC07A9F539B5158
FEF360FA77AF9ACEB173C57FB3D05EEF5F3FF806D4E2A3AB3A492F001D6B35
BC3137C4AB88FC05632225DEAD94899F3B519019031C027036F35C9DCEA934
99117CA3D4F5AF63FD9BADC5C7C54B4BC9396B3B7B8941C67965F2FAF6E1EB
CDAD96D4CBF0B531E9DA74D3947FC4B55F8D8B5353928773F3AB53B0BDD1F5
2BAD275285ADEF2CA5782689BEF2491B15753EE0820D551257DE5FB727C1F9
742F1243F16F458B3A6EBCCB0DF051C437A8B85738FE199173FEFAB1272C05
7E7F8931C57F63785BE2352CDF2DA38EA4FE25AAFE592F897C9FDEACFA9F39
8EC1BA737166AA4B56925AC41255949ABFA032ECF134B53C99D236965A984B
58EB2D4EB2D7D661B3AD3739DD23623B978983C6C518770706BBDD03E29F8E
7C34EADA56AD322AFF000B3161FAD7970969FE6D785C45C2991673171CD309
4EADFAB8AE6FFC0B497DCCDE8626B5277A7268FB9BC0FF00B6EFC45F0C491F
DB99A60B8F9A362AD5F737C3AFF82955BC9E55B6B7718CE015B85CFF00E3DF
FD7AFC34F37DA93CDAFE6DE29FA1E70B62A4EB657527869BFE5775F7692FFC
98F7B0FC51888E955292F33FA2FD6FE3AE9BF19BC45A56A1A4AA25869F011F
BB39579E56CBB71E8AAA07A735F567C36F1069961633DD5F48B1C704665763
D9506E63F90AFC1FFD957C5F158683F61327CD14EE0E4FA9CFF235FA0979E3
991FC353595BCB86BCDB09C1FE173F30FC5722BFCC7F13382EB64F9DD6C0B9
733A3371E6B5AFCAF7B79EE7E8580C546AE1F9AD6BA3DB745D6EEBC5FE21BE
F155FE7CDD466322AB73E5C7D238FE88802FE15F477870AED502BE55F034EB
1C1101E82BE93F0F5E81B39AFE7CE2BA952B579D49BBB6DB674E1D24AC8F56
B7505866BAFD3E25C0CD71B6B202030EF5D2DADD84039AFD53C1ACC30987AC
A55CF2F31849AB23A7C0E958F7FB79C539B501B7AD62DD5DEECF35FD01C7DC
758078374E0EECF2B0B859F35D997718C9ACC9982A1356A6933926B9BD46F9
514806BF8473DC4AAD5DB89F53421A1CE6B53821ABC3FC4EEA43D7A56B17C3
0DCD78BF88EF4156E6BD2C9B0ED3455567CC7F12AD22BAB59E37008208AF84
35091AD2E65B663CC6C457DDDE39B9568E4E7D6BE04F1DCA2D754798702424
7E22BFA1384E939D2E4672D1ADCB53D4CDB8BE001E6B9DB8BC9AE268ED6D11
A79E7658E38E352CEEEC70AAAA392493800753542DDB52D6F51B7D1B45B692
FAFEF1C470C10A9792473D828E4FF935FA87FB377ECDD63F0D441E39F1C04B
CF15C91E6287878B4FDD9C853C869483867E8BCAAF1966FD5723E1A9E227B5
97562C6660A0BCCF4EFD95FE032FC19D026F11789D124F176B4804FB487169
6FC30B7561C16CF3230E09000C8504FB47C55F8ADA3FC2FF00066A1E2FD61B
7ADB2ED821070D3DC3F11C4BD7EF1EA7076A82DD01AB72EADFED57E557ED73
F1427F1978FD3C1563396D2BC35F2B856CAC97B20CC8C71D4C6A44633CA9DF
EA6BF57ADC983C372D35B688F969CDCE5CD23C1B5CF126B9E39F125FF8B7C4
93FDA751D4A532CAFD00CF0154765518551D8002AE5AA56058A800574F6632
457E7389A8E4DB66C8E8F4DB432B8502BDB7C23E0F935091004CE4D70BE14D
3C4F3A0C77AFBDFE10783A2BA921DC99E9DABCC99A443C13F0665BD446F249
CE3B57B8C7F019D6DF71B7EDE95F5AF84BC2F65A558C5B631BC81CE3A576ED
6EBB7A57A783C8EB5787B48EC44AAA4EC7E60F8BFE129B2473E4E319ED5F28
78BFC29F63771B318AFD9AF883A2DA4FA73CC5006E735F9B3F1434F862965D
A3D6BC99D3706E2F7344CF85754B3F29DB8AE4AE170715E9BE238D525702BC
D2E880C6AE085230EE075AF0AF8CBE0F1E22D006A16A83EDFA71DC87A178DB
864FE4467B8C7735EEB39EB5E69E3ABBD965159A9E677C9FF753FF00AE457B
D92A9BC4C3937BFF00C3FE063536773E0DBBB0BDB16D97903C27FDB5233F4C
F5AA95F5A8823990C72A0746EA18641FC0D64DC7C3AF0D6AE7FE3D7ECCE78D
D01D9FF8EF2BFA57EA2711F30515F46BFECE7AF5F8DDE1EBF8E763D239D4C6
7FEFA5DC09FC05709E22F829F157C2D13DD6ADE19BCFB2C6096B8863371085
1DDA48B72A8FF7B1401CAF8435B6D0F5A8AE09C452FEEDFE87A1FC0D7D8DE1
BF148DAA37FEB5F08D7A67853C5134612D257C489C2E4FDE1E9F5AF8FE29C8
238A8FB44B547D06499A3A2F95EC7DFDA6F8914A8F9EBAB83C4431F7EBE3ED
2BC62400AEC548AED6D7C5CA40F9EBF18C7F09493D8FD130B9F26B73E993E2
11B7EF567DC788863EF57838F1582387CD412788E493EEB579B0E1669EA8EC
967575A33D5F50F110C1F9AB81D4B5979C91BB8AE525D50BF2EF9AE7F52D6E
1B585E69A408880924F402BDFC0E46A2D24B53CAC5664DABB654F1A7891347
D267B9DC3CCC6D8C7ABB74FF0013EC2BE4A24B12CC724F535D578B7C4B2F88
6FB7292B6B0E446A7BFAB1F73FCBF1AE52BF6BE1ECABEAB46D2F89EE7E719B
63BDBD4D3641451457BE796145145007A2681F11F5AD1635819DE445E010C4
1C7F5AFDDEFD8C7F6068BE2F780BC35F1AFE34EA7730695AE47F6CB5D0A15F
2E49AD4B1F2649EE4316093280E12350DB0A9DE09207E5D7EC15FB2A6A1FB5
3FC70B0D1B52B590F82BC3CC97DAFDC0DCAA6DD4E52D55C63125CB0D800218
26F71F72BFB18B5B5B6B1B686CACA14B7B7B7458E38E3508888830AAAA3000
006001C015E4CF22C24AA7B4953573555A56B5C8B4FD3EC749B0B6D2B4BB74
B4B3B289218218942471451A8544451C055500003802AE51457AC641451450
01451450014514500145145001451450014514500145145007FFD0FDFCA28A
2800A28A2803F297F6F1F06DE783BFB47C71A5C67FB37C490F97391D23BA5C
071EDE646370F521CD7E7AF86359558A321BA015FD0F7C5BF877A67C57F86F
E21F87FAAAA18F59B39618DDC64433953E54A31C831BE1B8F4C57F2F5A6EAD
A8E85A85DE81ADC2D67A96993C96B7303F0F14F0B14911873CAB020D7E73C6
394F335522B73A68CFA1F7F782BE306ABE1741F629DA3C7A1C56878B7E37EB
BE2288C77574EE08EE6BE32B5F132951F3D4D2F89411F7EBF37795BBEC74F3
1E91AB6BC662CEED926BC97C41AC2ED6F9AB2B52F12AED3F357976B5AFB4FB
951ABD9C0656EE44A65CD023B1D7BE24F8634BD4625B9B3BCD5AC609A37195
78E49D15948EE083835FD47689E13D0BC37A6C5A478774DB7D2EC20188EDED
6148214FF7510051F80AFE73BF62CF875A87C56FDA7BC1BA7DAC2D25968174
9ADDEC8012B143A7B0950B7B3CC238FEAFE99AFEA08E9F1EDE3AD7171A7879
88CE15274DDA304FE6D9A61F11085F991E6977A6A329E2BCEB5AD39406056B
DBAF600848AF32F110542DF4AFE62CDB26A980C4BA32E877CD269491F3478B
3448668DD5941CE6BE23F88FE116D3E59350D3D31824BA0E87DFEB5FA07E22
DA55ABE66F1CDAC72C52823AE6BEBB87B14D35196C72CB4D50EF0DFEC4BAEF
8CBE1CDC6B7AE78AC787AF350B659AC5ED3CB9E35570194C921E183A9EB19F
973B831E87F203E2B78EAD3C052EB9F0BBC3DA41D39A2925B4D52E6F42CD7B
76EA4870EE720231E46D38230727AD7E93787BE29FC5A8EDEDFE0BF832DE6D
5E6BC99FFB3A20E0794082F22966215635C17C9200F9BD8579478ABF60CF8E
9AE78C6EFC79AEE816BAE6A770C9F2C77B6E22CA0DA24C4AF1EE380065BA60
6172335FAC70B71A61323A95296635A1153578DE51527D3CAD1F5BEA9EAF52
2A509D5578267E348B851232A1F941E3E95712E3DEBEDFF8D9FB18FC71B392
6F1458F8175317182F731DBC06E51C004990183CC0ADF29DD9C03C1EA7E6F8
2EE22BCD3EEA6B1BF85EDEE6D9DA3962914A3C6E870CACA790C08C107906BF
76E1DE26C166947DAE0EAC65DD269DBEE382B50941DA48DC13FBD3BCFF007A
C259CD4825AF7EC6274BA7EA8FA7DC8994FCA7EF0FEB5EEFE188AE7C572C3A
6E9D1FDA269F8007403B927B01DC9AF997CCAF70F83BF11BFE116D4974EB9D
A904EC30E001CFA13FCAB87309CE149CE9ABB47CF67192C6BFEF23A334FE32
7C1DBDF87C60D5AC49B9D3265559580FF5531183F4463F77D3A7A67C059EBF
5C626D1FC6DA0C9A7DFC69736F75195756E410C2BF39BE257C21F11783BC5F
0681A459DC6A716AD288F4E1046D2CB33B9C2C215012D267800025BB0EC3C1
E1DCF5D6FDC57F8D7E3FF04ECCBB19CCBD9CF74793DBDBDD5FDD4361610BDC
DD5CBAC5145129792491CED55555C966627000E49E95FD0CFC0AFF00825D78
1F4CF83FA69F8B9673B78DF5353757B25BDD4919B4F317E4B5555631379439
662A73216C129815BBFF0004FAFF008274AFC206B3F8D5F1D2C926F1BE0BE9
9A5394962D241E934846E57BA23EEE095881CF32729FAFE56B2CFB32A8FF00
7741E8B77DCF7A9416ECFC2EF889FF0004C2B48166B8F87FE2BB9B46196587
50896E10FF00B3E645E5328F72AC7EBD6B81F85DFF0004D8BAD4B518BFE16E
78DE1B2556F9AC3498DDA5946780B737288AA7D844E7F9D7EFD5F69B0DCA10
5466BC77C55E11B6B88DD5E30C08F4AF82CFAA63F17849E1B0D8A7426F69C6
31935F2927F834FB3475D2E58CAF28DD1E2DE0DF839F0E7E0DF87FFE11CF87
9A345A65BE079B28F9EE2E187F1CD3365E43C9C64E17A280302B8BF165B874
7E3D6BB2D4AFF5BF0D9314C5EFAC471B58E6441FECB1EBF43F8115C6EAF7F6
9A95AFDAAD1C3C6F9FA83DC11D411DC1AFE15E27E03CEB27C6BAD99C9D5537
A55BB7CCFCDBD54BC9FC9B4AE7D2D1C553A91B434F23F3E3F68AF869A7F8EB
4196C2E4F91344C248660BB9A261D4819190464119AF04F0AFC1AF0A7856D3
162CF35D48B86B87DBBDB3F870BEC081F8D7DCFE3D81648A5046739AF98ECD
F63CF667AC0E40FF0075B91FD457FA3BF420CDB2BC7E22AE419C515539A3CD
4EEDE8D7C51B5EDAAF7969A72BEE7C6F14C2A422AB5276EE78A6A96777A7DC
BDADCE54AFA7423B11ED588CB5EFBAB6936BAB41E5DC0C3AFDC71D54FF0051
ED5E43AAE8979A64A5674CA13F2B8FBA7F1FE95FD1FE20785F5F24A8EAE1E3
CD876F492FB3E52EDEBB3F5D0F23059846AAB3F88E6596BE85FD9B17CAF1A5
FDDE71B2C9931EBBE58CE7FF001DAF0368EBDCFE014C2DFC4B7E09C335AF03
BF0EB5F84F1941CB2BAF15DBF547AD87F8D1F79788BC31A0FC45F086A3E0BF
12C5E7E9DAA4262900C6E43D51D09CE1D180653D8815F837F173E187887E10
F8DAFBC1DE204C980EFB79C0C25CDBB13E5CA9EC47047F0B02A7915FBB7A36
A3F2AF35C47C75F823E1FF008F3E10FECBBA74B2D6EC73269F7C54131B91CC
721037189FF880E8406192307F14F0ABC48A9C37983A5886FEAD51FBDFDD7D
2497E12EEBD123D2C760D5685D6E8FC0E1254825C74AD6F16F853C41E09D7E
F7C33E26B292C352B0731CD0C83907A8208E0A91CAB0C820820906B9AF308A
FF0046725E2A8CA119C6578B574D3D1A7D533E3EAD035926A9D66AC4127352
2CF5F7D83E214D6E724A89B826A7896B184F5209F15EE52CFD7732748D7F36
93CDACAF3A8F3ABA3FB7BCC5EC8F59F873E379FC29AA72C4433919F661FE35
FA09E00F8831F8924B781240E2321FAFE1FD6BF287CFAFA33F677F13CB6FE3
26B2B873B65B772BCF565653FCB35FC79F48BF0DA18D8D5CF708BDE51F7D7A
2B732F96FE97D753E9B25CC1C2D465F23F6FBC19A981147CFA57D07A1EAC00
5E6BE23F05EBE0C51FCDE95F42689AE8DABF357F97D9FE56F99E87DC51A87D
6BA46BA854248722BB186FA37194706BE63D3B5EC0043575B6DE2365030F5F
1B0A75A84AF4DD8DDA52DCF76377EAD5427BF89012CFCD7947FC248E47DFAA
53EBE48396A9C4627155749C8234E28F40BED6570421AE1F51D541CFCD5CCD
DEB6307E6AE4350D6C60FCD4B0B9636CA94CD0D6355186E6BC67C45AA8DADF
35686B1AE0C37CD5E33E22D7061BE6AFB9C9F2B775A1CB52A1C078D7530639
39F5AF99C7C38D6FE2B788A1F0E687710DA48BBA79A79C9DB1C0B85621472C
D965014633DC81935E97E2AD58DC3B203C77AD9F801697F7BE39D52EEC63F3
BECBA738655C17F9E58CE557A9C0524E338193D01AFDEF8370D4695587D61A
8C6FAB6ECBEF7DF63CAAF525BC773E8FF84BF07FC0DF07ED0B68911BCD5A65
DB3EA1380667071954ED1A64676AFE24919AF609358FF6ABCCDB5AFF006AAA
49ACFF00B55FD2F4A8C69C5460AC8F1A526DDD9D478BFC6B0F863C35AA788A
E0864D36DA59F6938DC6352C17EAC7007D6BF18A3BDBAD52F6E354BF90CD75
792BCD2BB75692462CCC7DC939AFBFBE3F6A1ABEADE026F0CE836F2DE5FEBB
776F691450A96763B8CC40C76C46724F18CE6BA2F819FB144090DB6B5F135B
ED57076B8B089B10C67AE2571CB9F50085EA3E615F9378A7E2265991528CB1
F52CED7515AC9FA2FD5D9799D783C24EABB411F15F84BC29E26F15DC7D93C3
7A5DC6A5202037911B385CF4DCC06147B922BEA5F0AFEC9BF15352DB2EA51D
B6969C656593CC7E7DA30CBC7FBD5FAB3E16F87BA1F87AC22B0D2ECA2B3B68
40091C4811547B28C015DEC5A75AC430A82BF87789FE9478EAB371CB68C611
EEFDE7FA2FC1FA9F4B472482F8DDCFCEBF0EFECABADE921259F584771D556D
CE3F32FF00D2BE90F06E83AA7825D0B422F153AE0EC3C7B106BE8DF221E9B0
5432595BC8305457C147E90DC4BCD775D5BFC10FFE44E9FECAA1DBF33434BF
8C3E198D521D5A1B9D3D80C1678B7A7E0632C7F3515E9DA7F88747D66D7ED7
A4DEC3770F19689C3019ECD8E87D8F35E01A8F87E0954E101AF30D53C39716
131BCD32692D2E07024858A363D323A8F51D0D7EC9C1BF4AEC7538AC3E3E94
651EF1D25FAA7F87A9C189C8A0F58B3DCBE2578AED2DEC9EDD240700E6BF36
BE266BD14F2CB86F5AF4DF1B6B1E2C8ED9C5D31BBDA0FCE061CFD40E0FE18A
F8B7C59E2692E25915C90C09041EA0D7F4570F714E0F37A5F58C24EFDD755E
ABFA5E678F5A8CA9BB491C2F886EC3CAD835E75752649AD5D4AF8C8CDCD731
3CD926BE9631306CAF33D78C78C2E8DC6B461CFCB6E8ABF89F98FF003AF5B9
64EF5E097D73F6BD46E2E7391248C47D33C7E95F5DC2D42F5A53ECBF330AEF
4B1340B922BBFF000F69FE7CAA319CD70D6632E2BD9BC1B02B4C99AFBB394F
A4FE1978356EE48B29D715FA43F0C7E164334711310EDDABE4FF008416F007
84B63B57EA4FC2D7B38A08B38E82803C8FC67FB0B7ECFF00F132DE49BC6BE0
7D3AEEEE6077DD4517D96ED89EE67B731CA48ED9638AFCD2F8D7FF000487F0
75979BA8FC28F16DEE8D2A92E2D75245BC8739E15258FCA92303D584A7F3E3
FA09334417796005785FC4BD66D3C890061D0D007F2A9E38FD973E3AFC2F91
A3D634B8B5EB68F8FB5695279E47A662758E53F5087DEBC756FAE2CE76B5BA
478664FBD1C8A5245FAA3608FC457EEFFC57D4E1769B691DEBF3DBE2259693
AC6F8752B48AE9013812A2BE3DC64707DC579589CA29D4D5687751C7CE3BEA
7C7F0EB2A0726ADFF6EC6A3EF7EB5735EF875A34ACCDA5DC5C69CC7B472332
7FDF2C49FC88AF29D5FE1DF88E2DCD6F723504F42C55CFE0C71FF8F5794F86
62DEB23BBFB69A5B1D3EA7E3AB0B4565593CD93B2A73F99E82BCAB5BF11EA1
ADBE266D908E9183C7E3EA6B36FB4CD474D7F2EFEDDE03DB72900FD0F43F85
51AF5B0593D1A1AC55DF7387139854A9A37A0514515EA1C21451450015E8BF
09FE1578DBE35FC40D1FE19FC3DB06D435BD6E658A2501BCB8973F3CD33286
D9146B9691F1F2A827DAAA7C36F867E3AF8BFE32D3FE1FFC38D1E6D735ED4D
8AC36D0019C28CB3BB310A88A0659D88551C935FD6D7EC35FB14F85FF644F0
248B71245ABF8EF5E8D1B58D4D57E550391696A580658233C9270646F9D800
111003D63F656FD9B3C1FF00B2D7C23D33E1B785D567BCC0B8D56FF6E24BFB
F703CC95BB8518091AFF000A00393927E90A28A0028A28A0028A28A0028A28
A0028A28A0028A28A0028A28A0028A28A0028A28A00FFFD1FDFCA28A2800A2
8A2800AFC4BFF828F7EC8FE22B7D62E7F68FF84DA7C97F15C007C45A75B216
951D0605FC68BCB29181305195C799C82E57F6D28AC6BD08D48F2C869D8FE3
374FF1D8741F3F35ACDE312EB85635FBF7FB407FC1323E057C63D46EFC51E1
0966F87BE23BB3BE4974F8965B096427E6792C98A0DC47FCF292204FCCC189
24FC27E1DFD8A7C4FF00B2AFC60F08EB9F11752D1BC55A46A97CB1D9189253
22BC0E8CC658268F62E558630EF8391C753F199B65B0C3529576AE91BC2777
63F372E35F926CEE7E2A858A6ADE22D56D3C3FE1EB49B52D4F509561B7B6B7
4324D34AE70A888B92493D857EC07FC141FC11A75F78FF00FB39F6D8DA5FA5
901222022043B622CAA303E503803D315EFDFB37FC03F84BF00A43A8785AC5
AF75B994A49AA5F159AEF630E52360AAB121EE1146EE3716C0AFCFF3FE38C3
E5B4A4E54DB9DE4925B3B6976FA2FBD9B46936CF4DFD843F65F8FF00673F87
B36A5E28092F8DFC54B14DA930C30B48D4662B3460483E59626465E19CF565
5435F75CB711A29E6BCCECBC4B652A03E6053525D7896CE3524CB9AF9EC2F8
DBC98570943DF3A1E0A2DDEFA1B7A8DD200CC4E00AF16F126A8B248F83C558
D77C5C2505236C2D792EADAD6EDD96AFC2733C654C762655E7D4E89495AC8C
7D7EF410DCD7CF5E2FBA5647E7D6BD175DD5C61BE6AF00F166AC0ABFCDEB5F
4D92E0DDD1CF26785DFF008C6EBC05E3AD0FC71A7E4CDA2DE4570557AC91AB
7EF23FA489B90FB1ED5FBAB04F0DDDBC5756EDBE29955D1877561907F115FC
F1F8E6F05CCAEB9C819AFDCFF819A97F6C7C17F02EA25FCC79744D3C48C46D
CC896E8AFC0C0FBC0F418F4AF96F1E3278FD57098B4BDE5CD17E9A35F76BF7
9E964F535944F4F2011835F117ED35FB14FC1DFDA0BCDD6B5AB26D1BC4FE5E
D4D5EC00499B03E513A1F92751803E71BC2F0AEB5F6F9EB59DA8C61E1AFE7C
E1FE20C6E578B8E2F0155D39AEA9DBE4FBAEE9E8FA9EDD6A519C7964AE8FE5
F3E2D7FC13DBE3CFC37BA9A5F0F5BC3E32D290E527B0611DC05EDE65B48430
6F68DA418EFD71F19EB9E19F14785A616DE26D22F3499892025DDBC9036578
23122A9E3BD7F5D5E24B757952D93FD64EDB57E9DCFE02BD97C15E1FB34D3C
59CD0AC90C8851D1C065656182083C1047506BFD0AF09FC52CDB33C0BC4E65
4E2E29D9497BAE4D6EDAD57DC92BDF4D0F93C76069C25CB067F12BE6D75BE1
4F0478EFC7578B63E06F0E6A5E21BA2DB445A75A4D7726EE380B0AB1CF22BF
B734F02F82BFE85FD3FF00F0122FFE26BB0B1B2B6B2B78EDAD224820880548
E350A8A0760A3802BF6ECBF389621DA30B7CFF00E01E6CE9DBA9FCB7681F09
BF68AF81DE0ED27C4DF193C2577A0691A95C7D9609A778D995F686413A2333
425C64289029255B8E067E82F03F8CB50F0E788349F1D7869E31AA691289A1
322EE46C82AC8C38255D4953820E0F041C11FBA1F19ADFC0B77F0B7C4B6BF1
2EC9751F0D4B68CB776CDD65048F2D10820AC864DBE5B020ABED2082323F10
3E12FC0AD6ADF54306ADA9ADC69224FF004744C9B878FF0085653B55548E84
AE738C80B9C0F8CE30A54309515653B49EB63E633550A52534ECCFD084FDA6
FC7DE38D26DF51F0CE9F0787A0B8439DE3ED33875255B0CE026DC838FDD9E3
BD73B73F14BE3344ED3AF89E50ECBB7FE3DAD4AF1E8A61DA0FBE335EFF00E0
2F82D64BA540F7482345401500C050060015078FFE185969762F736A31B466
BF3AC6E33359C3EB92E654DECF65F81E656AB8F70F6EDB51F5B7E078E68DFB
5378DBC39388FC69A6C3AD596E1BA5B602DEE517B90BCC6F81D170993D5857
D53E1BF1BF843E26683FDBDE11BD5BC801D92A60A4B0498C98E543CAB0FC88
E4120835F9D7E27B58D4BA11D335E33A3FC48F117C1BF1745E2FF0DB968B72
ADEDA6711DDDB86CB46DD83019D8F8CA9E791907AF2BCCA757DDA9A9E8E539
ED44D46B3BAFC4FD37F15E931C8AE0AF5AF92FC5B6579A05D4B7FA70CAB7FA
C8FA2C8A3F911D8FF4AFAFADBC49A2F8DFC31A7F8B7C3D37DA34ED5A04B881
C8C1D8E33861D994F0C0F20820D785F8D6C92449323D6BE8EBE068E2E8CB0D
8A829425A34FFAFB9EE9EA7DA466D5A5167C99E29BE82FAD4DC40DB91C123D
7DC1F71DEBE5DBB7F235F6C7DD9B2A7EA391FCABDEFC64AFA2DEC89FF2ED72
D83FECBF407F1E87F0AF9E35A90ADF4738EAB229FC33FE15F13E1A60AAF0B7
15E12AC65EE46A45A7DE0DD9AF5B5D3FBF6674E39AAF8792F237375452C714
E8629943A37504641A4DD4DDD5FEDC52C54671E596A9F43F2D71B1C95FF842
CE725ECDCC04FF0009F997FC4568F80B4FD43C37E2A82E2401A0995A17656E
81B91C707EF01DAB6F71A6872A432F04735F97715F82D9166946AC21074A52
4D5E2ECAEF6F75DD5BC9247A187CD6B536AEEFEA7D41A3EA582066BD374DD4
B81CD7CD9A26AFE7411CA0F2473F5EF5EA1A5EA8085E6BFC84E2EE19AB86C4
54C3D68DA706E2D7669D9AFBCFD1B0F5949292D8CAF8EFF007C29F1DB43DD7
016C3C4967115B2D400E9CE445301CBC44FE28492BD486FC51F88DF0D3C5DF
0C3C473F863C6362D6777164A3758A68F381244FD1D0F62391D0804103F7FA
C753181CD62FC41F879E07F8B9A07FC23DE37B05BC8572D0CAA764F6EE4637
C520E54F4C8E55B0030238AF6FC36F16B17C3F2583C5A7530DDBED43CE37E9
DE3F35677BE78CC046AFBD1D19FCEC1520F34DE457D8FF001A7F640F1FFC37
9E6D53C350C9E27F0F0CB09EDE3CDC42BE9342B96E3FBE995C0C9DBD2BE417
8596BFB878678C30798D0589C0565383EDBAF26B74FC9D99F335B0F283B491
5F7353848682B8F6A4C57DBE1F33F3399C094487D69C1EA0E69D835EA53C7F
999B81637D75FE03D73FB07C5BA6EA45B6A2CA11CE78DB20D849FA6735C480
7D697F1A8CC953C5E1AA616A7C338B8BF46AC385E32525D0FD81F02F8AC189
06FE95F49685E250557E7AFCACF85BF105A7D3E013CBFBF8408E4C9E491DFF
0011CD7D79E1CF18ABA2FEF3F5AFF3138CF836A61713528548EB16D33EE70D
89528A68FB974FF118C0F9EBA9B7F110C0F9ABE4CD33C58081F3D7616DE291
81F3D7E518BC81DF63D08D63E931E2218FBD50CBE211FDEAF045F140C7DFFD
6A297C5031F7FF005AE0590BEC5FB63D96EFC42307E6AE3F51F110C1F9ABCB
EEFC50307E7AE3352F148C1F9EBD7C2642EFB19CAA9DD6B3E23186F9BF5AF1
9F10F888BEE556C9AC3D57C4AD292A8D9AE0EF751CE599B26BEE32CC9942CD
A3927548B56BEE1998F26BE9EFD9B3C1AF2782F57F18CC1A39351BB115BB83
8609680FCEA4723F78EC383D56BE3D86D754F146B565E1AD0E3F3F50D4E648
214EC5DCE324F3803A93D8026BF5C74CF0EE9DE05F06697E0FD3DBCC874AB7
48779E0C8E065DC81DDD8963EE6B7E28C52A1878D05BCDFE0BFE0DBF1161E3
777385D3ED74AF1E5F378635398699E280ACD6D778FDCDF84192B328E7CE03
92EBCB0CB302C39F2DF10E83E28F0E6B2BA16AB66E9712B6222398E51EA8DD
08F5EE3BE2A97C45D42EB4BB88F59D2A536F7B6322CF0C83AAC919DCA7DF91
C8EF5F485878AE6F893FD97AB3DBB5B5B08229842E3E659A540D267FDD24A8
FA13DEBC09F89F9C70CC632A76AB869C649464F5A734B4E57BF23BAF71E8AC
D45C558D3EA54EB3D747F9937C36F87D0D9AC7797A825BB6192E4676E7B2FA
0FE75F53691A6456F129DBD2B93F0E592A22002BD2E24C28515FC43C6DC558
DCDF1D3C562EA39CE4F57FA2EC9744B447D1E1E8C69C39628503B0A788D8D5
D82D8B56BC7A7923A575F0FF0087D8AC6AE6844CEAE2A3139C31B0A6104574
92D8ED1D2B2A68769A33FE03AF82579A0A58A52D8CF232306B9ED56C95D49C
75AE8D860D53BB506235F051E6A733AD6A7827897468E557057AD7C41F17FE
1A0BD49751D317CABB404FA0703B1F7F435FA1FAE420AB578178BAC1248E40
466BF68E00E28C4607110AD42566BF1F27E47062E829AB33F23EF1E6866786
75292212194F0411D8D65492D7BBFC68F090B1BA6D72CD36F38980EE3B37E1
DEBE72967F7AFF0041786B39A798E12189A7D775D99F255A9384ACC8754BCF
B3D95C4F9C18D19BF102BC3A26AF49F145D94D22700F2FB57F3233FA5797C4
D5FAB70CD1E5A529777F97FC39C159EA74366F8615EBBE14BC58A5439AF15B
7976906BB4D2350F29D4E6BE94C4FD07F86DE284B53165B18C57DDBE09F89B
1DB4283CDC631DEBF1D7C37E2C6B52BF3E315EE1A3FC4B920451E6F4F7A00F
D70B8F8C318B6C79DDBD6BE7FF001E7C545B8490097D7BD7C5D71F156531E3
CEFD6BCDF5EF8892DC06FDEF5F7A00EDFC7FE32FB5349F3E739AF947C47AAF
9D239CF5AB7AE7899EE0B65FAD799DFDF995892680295DCDB98F35912353E5
9726A93BD00473047528E0329EA0F20D71DA9784B40BECB35A885CFF00145F
27E83E5FCC5756ED54A47A00F27BFF000049192DA7DD061D964183FF007D0C
FF00215C7DDE89AAD913E7DBB6D1FC4BF32FE62BDDE66AC8B86EB401E155F5
37ECC9FB1F7C67FDAA3C471E9DE00D29EDF4286611DFEB7748C96168382C37
E3F792852088932C7233B57E61CA78234EF0E5FF00C45F0C45E27B25BDD2E4
D52C96F2220FEF6DCCE8254254AB7CC991C107D08EB5FDB4787BC3BA0784B4
4B2F0D785B4DB7D2349D3A310DB5A5A4490410C6BD1638D005503D00A00F9A
7F656FD8E7E11FEC9BE186D33C0D6A6FF5EBE8C2EA5ADDD2837977820EC047
11420805624E38058B37CC7EAFA28A0028A28A0028A28A0028A28A0028A28A
0028A28A0028A28A0028A28A0028A28A0028A28A00FFD2FDFCA28A2800A28A
2800A28AE0BC7DE301E17D3922B4C36A37995854F3B40FBD211E83B0EE7DB3
5C79863E96168CB1159DA315765462DBB219E2DF1CDBE84CDA6E9CA2E75220
7CA7EE459E85C8EFFECF5FA719FCEAFDAD9B59960F05F8B2FEE1E736DAB985
98E30A6685994019F947EECF0063FAFD49A55B33319A663249212CCCC72CCC
792493D49AF00FDA67C57E09D43E1B6A3E1D8E737FAA5BCB6B7311B701D219
22994B1690FCA3E4DEA40CB73D057F2F665C6D8FCDB318C29C5FB34FE15B25
B5E5FE6FE47AD1C3C610F33CEFFE0A0164DA9E93E1EF11DB92A6FF00495657
5382197E60411DC6E18AEC7C31E2F5D474BB3D451801750C728C1C8C3A86E0
FE3547F69884788FF672F87BAF81BCADA1B763F48C0FE686BE68F861E2F2DE
0BD1637932F05AC50B12724B42BE5939F72B4F8C303F5882979DFF00F038A6
651763EE4B7F143AA8C3D4B2F89DD872F5F38C1E29181F3D587F148C7DFAFC
D5E4AEFB1A731ECD79E21C8396AE2753F100C1F9ABCDAEFC50307E7AE2754F
140C1F9EBD2C264AEFB09C8EAB5EF100DADF3578078AFC400AB80D92699AF7
8A461806AF20D575569D99DCD7DE6539472D9B463291CD788AF72923B1EB9A
FE833E0F6912681F08BC13A34F188A6B2D174E8A5551B479AB6E81CE3D4B64
9EF9EB5F837F0B3C1773F15FE2CF86BC096F1B4B0DFDE46D77819096709F32
E18F61FBB52067A9207522BFA2B60AA02280147000E800AFCA7C7FC7423430
D83EBAC9FA6CBF53D6C9E0EF29113550BE2045578F5AE4FC4FAA2699A5DE5F
B72B6B0C9291ED1A96FE95FCAB87A12AB5A34E0AEDBB2F567D03765767995A
DC2EAFE2CB9901CC566DE427FBCA7E73FF007D71F857D2DE1A50B1281E95F2
47C3991BECF1CB2B6E91FE6627A963C93F9D7D55E1EB901539AFF4B729C9A9
E5D82A382A5B42297ABEAFE6EEFE67C754A8E72727D4F418864D688E056644
C3835A6A722BF51E179C5D37DCE2AE8F847F6E3F1749A7E8FE12F0644FB575
7BB9EEE600F252CD55541F62D367DCAFB57CF5F0E755B7B6BFB6925236A906
BA8FF828919B4AF137C3BD6C9FDC5C43A8DB9E38578DA07193FED06381FEC9
AF96FC39E290A88CAFC8AFCF78FB0F3A98B937D97E47C46774DCAAB3F6B3C3
7E33D126D262613A801477AF21F8B5F1274C363258DA4A1D88C706BE0AB3F8
8B7D0C2234B82063D6B0F53F18CB73969652C4FA9AF9BC56718CAD868E0E6F
DC5E473D7CCEB54A5EC5EDF99B9E24D5D65691F3D735F33F8EAFD1E19413D8
D75DAE789576B7CD5F3AF8CFC462457456C935B651807CC8CB0B41DCFD12FD
873C6973ACFC2FD77C2976E5C787B527F232490B6F76BE60400F0312890F1F
DEE83A9F76F16302AF5F187FC13FED658749F1F6B873E55DDC585B838E375B
A4CED83EB89867F0AFAE3C5176087E6BECDC2D51D8FD0B06DFB28DCF92FE27
D9A5C5BCCA7B835F226AB3996DC17FBE0ED6FA83835F6178FA6568E4CFA1AF
8C359252EA48470249011F52707FA573E7B95FB7A74AA457BD0947EE6D27FA
7DC75539D9B4763BA9A5B8A66EF4A66EAFF50F0D9969B9F032A63F7526E351
EEA6EE15EBD3CCBCCC9D337F46D41ADA63131C2B723EB5E99A5EADB4852D5E
27BF07238AE8AC35338018F22BF8A3E925E1E2FAD7F6D61A3EED4D27E52B68
FD24BF15E67D4E478DF77D94BA6C7D1561AAF03E6AEBAD354E9CD7CFFA76B3
D32D5DAD96AE300EEAFE22CD72269BD0FA98553DBADF53E9F35781FC4EFD99
3E11FC5269750B8B0FEC4D5E5C937B61B626763CE658F1E5BE4F524063FDEA
ED2DB56E9F356E41AAFF00B55F3F80AB8DCBABFD63035654E7DE2EDF27DD79
3D0D64A335692B9F93BF103F640F889E16D5AEAD3C36D1788ED61395788882
62A46E04C52363383D159B9AF99B5DF0AF887C33766CBC41A6DCE997033F25
CC4F131C752030191EE2BF7235ABC2BAF1909FF5888DF97CBFD2BB7D3FEC1A
9DA7D8F52B78AEEDDC61A3951644208C72AD9078AFD9705F48ECD30308AC7D
08D556DD3E497ABD1C7EE48F3A593C24FDD763F9E33111DA9361AFE80B52FD
9B3E0078B1B7EADE0BB18DDB9CD987B2E718CFFA33443FC7BD7307F606FD9D
6FE51225BEA76AB8036457A4AFD7F78AE7F5AFAAC37D30787E2BFDA68D683F
F0C64BE4D4AFF82307C3F57ECB47E12ED34E0B5FBF9A37FC13E7F66AB69C4B
7763A96A0BC7EEE6BE755E0E7AC2236E7A75FF001AFA2FC11FB307ECE7E079
239FC3DE00D2C4D1FDC96EA237D2A9F557BA32B03EE0E7F335E3677F4E3E1E
C341FD570D5AA4BD2315F7B937FF0092B34A7C33565F14923F9E8F83FF0008
FE3378FF00528DFE1B784F51D6E195B63CD0C0CB6A3071F35C3ED85483FDE7
15EAD6FAF6B7E10D6AEBC35E25B6974ED4B4E95A0B88261B5E3910E19587B1
AFE976CAE62B6863B7B7558E28942A2280AAAAA3000038000E82BE46FDA97F
64AF0C7ED0569FF0926892C7A2F8D6D630915D91FB9BB45FBB15D05049C0E1
241F328E0865000FE7AADF4B5A39E671FF000AD838D0A1256528B72927D1CD
D926ADA6914D79ADBD78E44E953FDDCAEFFAD8FCB6D1FC74AEAB893F5AEF2D
3C68081FBCFD6BE45F17784FC79F09FC4737853C75A6CFA46A36FC98E5190E
B9C078DC129221C70C84A9F5A96C3C4D3103F7B5FB3D4C82857A71AF424A50
92BA69DD35DD35A33815669D99F67278C948FF0059FAD325F18AE3FD67EB5F
2DDBF8865603321AD48F59DDC97CD7972E1B8A669EDCF76B9F1733F08D9AE7
EE35C966E59B8AF321AC281F7AA1975A503EF56F4B2851D909D43B99F53001
F9AB96D4B5A58D5896E95CACDAD4971325ADAA34F3CAC15234059D98F40AA3
924FA0AFAAFE0F7C17B5B4B88BC59F122349E74DAF6DA7B1DC91B0390F3E38
66F44E547F164F036C4FB2C2D3F6957E4BAB146F27647B57ECBFF0A66F09C0
FF00137C5F0341ABDEC663B0B79061ADEDE41F348C3A89241C01C155C83CB1
03E81F127885763FCD5C4EA1E325DA7F795E49E25F19A047FDE7EB5F9ACF0D
5B198975EAAD5FE0BB1DCA4A31B234F5DD2F53D6AE2D2E27B49469B7336C13
32911B95E4AAB1E0F4C71EF5F51F822D5208225518000AF9FAD7E2CD9FC4EB
5D0AC74A554D3F4047802AAEDFDE8DAA430F550BFA9AFA23C2D28544AFC4BC
53C6569CD50A8ADC9756FEBAEC7A1828ADD1F426838C20AEF604048AF33D06
E402BCD7A75BB0386F5AFE77C2D15F5A5CE7AB37EE9D2D8C20E2BA78A040B5
CBD9CE1715B6B7802F5AFEF1F0BF1F9761F0ABDA5AF63E5B1B09B9682DDC6A
01AE52EC0C9ADBBABB041E6B9BB9977135F9B78C19EE0EADD52B1D980A725B
99B275ACDBD902C78ABF2B85C935CB6A5780E79AFE40AAB9EAE87D0474472B
AC4830D5E2DE24C157AF4ED5EEC61B9AF20F105C82AD5F77905069A39EAB3E
60F891A7C579693C32286575208F635F9E1ABC2FA7DFCF64FF00F2C98819EE
3B7E95FA3DE349034720FAD7E7D7C4B885B6B26E17812641FA8AFED2F0671D
28C9E1DED25F8A3E6F3285D5CF22F175C8FB0C7183CB483F200D71113D6B78
AAEB77D993B12C7F2C7F8D7391495FD799353E5C3AF99F3F537372392B52DE
E4A1EB5CEA4956D25AF54CCEEED35678F186AE8EDFC472A0E1EBCB12723BD5
85BA6F5A00F537F134C47DFF00D6B1EE75D924CFCD5C3FDADFD6986E09EF40
1B9717ECE4E4D64C9396EF54CCB9A89A4A009DA4AAECF51349503C9400E77A
A523D2BC954A492802295EB22E1FAD5B9A4AC89E4CE450050964646DE84AB0
E410790477AFEEB2BF883F875A3CDE24F88DE15F0EDB46B2CDAA6AD636A88E
0B2B34F3A2004004904B60800FD0D7F6F94005145140051451400514514005
1451400514514005145140051451400514514005145140051451401FFFD3FD
FCA28A2800A28A2801090064F0057C65AF78993C4BE22BEF104F26CB38B2B0
9738096F1E704FA67963F5AFA4FE276AEDA1F80F59BF8DB649E4F9284750D3
B08811F4DD9AFCC8F1E78A1E511F84EC1F6A2857BA2A7AF758FE9D18FE1EF5
F96788B1AB8B9D1CB693B27EF49F92D3FCFE763B30B657933A0F187C4CD43C
4AD2691A0C8D69A48CAB3AE564B81DF27AAA1FEEF523EF75C0F2ED7F45FB67
85756B48A3DCD25A4E178070DB0ED201EE0E08F7AD3D2ADB7B2C6A2BEA0F86
FF000C17C46B8BB8C340EA43861C1523047E35C381CB28E1692A1878D97E2D
F77DD9A39393BB3C9B5723C4DFB13E8F723E76D2AECA1F6059C7FECE2BF317
C2BE2BFEC88A6D21DF06D6799793D99CB8FD1ABF57BE10E8693FECAFE2FF00
085D8334DA3DC481F78DA77C3B7276E4E3943C678AFC4BF8A1A3F89340F13E
ADAE699672BE8F14A8B34C837245238E03E3254363827009C80720D7854A8D
3AD254A6ED78AB5FBA6E16F5D0527D8FA92D7C6EA40C49FAD5F6F1A2E3892B
E1FD3BC6D72CA3E7FD6BA48BC577120E64C5455E17B3D89F687D4977E345C1
F9FF005AE3750F16BCD908D5E2DFDBA5F967CFE35149ADA81F7AB4A392463D
05CC77B77AB3392CED935CF6AADA9C7A536B46D261A7EFF2BED1E5B793E675
D9E6636EEF6CE6B43E16EADE0CBCF88DA1DB78ECF99A335C299A2DC5166C72
B1B30208563C120838E0104E6BFA1FD5FC0DF0FF00E2BFC35B4B4D0F4EB35B
0B78E3F26D5618CC1179443201110530080718C5736658C7847CB0A7CD24B9
ADB5E2BE2E5D1DE51DDAED76B62A11B9F16FEC19F026F3C13E19B9F8B7E2DB
66835CF13C423B2864186B7D3490E188ECD70C0310738454E84B0AFD05635C
AF86BC4726AF14D63A9A0B7D5B4F212EA21C0C9CED9133FC0F8247A1CAE4E3
2766E2F12307079AFE0CE3BE20C5E6398D5AD8A56937B744BA25E56EBD77EA
7D7612946304A3B0FB99C44A79E6BC53E2BEA4D0F8335664232D10539F4760
A7F435DE6A1A97079AF08F8B7A893E09D5B6E4908A78E780EA49FC05570260
6F9AE11C97FCBC87FE948AC54BDC97A321F015F2882219EC2BE97D0352002F
35F10F80F5A06DE221BB0AFA5742D606D53BABFD1DC552D4F904CFA874EBE4
9900279ADE8E6C5789699ADE36FCD5D55C78CB4DD1B4D9F54D5EE16DED2D57
7C9239E140FD493D001924F039ACF078EA987973442693576784FEDC7F0C5B
E26FC03D5AE74E21758F09B7F6CD964EDDE6DD584D167B9921670A3BBECAFC
1EF0AFC428E5853F79CF707D6BF553E3BFC7D9FC645ACED19ECB40B63B9216
3879DC7492500F6FE15E83A9E7A7E3AF8BBC21ABF8CFE25BDA7C1FD367D4B5
4D51A495F4EB542FB9D417924403EE8C02CD9E0727BE2BA1E690CCAAB8CA36
B2DFA1F275EB53AF55C60B45D7B9F435BF8E10A0FDE7EB50DDF8DD769FDE7E
B5F1A5C7893C43A0EA13E89E20B59F4CD46D1B64D6D731BC3344E3F85E3701
94FB11563FE12C9261F3CBFAD43E1EB3B9CEF0163E80D6BC68640CA8F93F5A
F27D575796E1F6AE649643B55546492780001C926B99B1BCD4F5FD461D1B41
B49B52D42E9B6456F6D1B4D2C8C7B2A202C4FD057E98FECD3FB285EF83F50B
5F897F16D233AB41FBCB1D28ED916D9F1959A76E54CABFC0ABC21F9892D809
DB0C24282BB3B70B81727A1F4AFC00F004BF093E0F695E1DD463116AD79BAF
F5003A8B9B803E43EF1C6A919C704AE6ADF89752043F35D56BBAC64360D786
789B581B5F9ACE941C9DD9F489249451E4DE38D403248335F306A7119B58B5
51CE24DC7FE03F37F4AF65F14EA1E7C8CA0E40AF1E91F7EB1173F7439FD31F
D6BEC387B0B0A98AA309EDCD1BFC9A6615A5EEB68DA269A5A9A4D309AFEBCC
266E9F53E6A5487934D2D4C27B5349CD7B74B35F33274C793EB424851B3DAA
126826B3CC674717879E1B10B9A12566BFAFC0704E2D496E6D437EF09073C7
AD75561ADE71F3579EA4817E56E54FE952E5E23BE33C1FCABF88F8F7C3F965
D5DC5EB4DFC32EFE4FCD7E3BA3EAB098C535E67B55A6B238E6B7E0D5F38F9A
BC22D75595080D5D0DBEB278C9AFC771DC3FAE88F46354F46D66FB7DC5B4C0
F66527F223FAD75FE1FD4F85E6BC5EE754135BF5C942187E1FFD6AE9343D53
6B28DDD6BE0788F256A9DAC75D2A9A9F4F697A97039AEFEC350E9CD7CFFA4E
A8085F9ABD0B4FD4C60735F8567394EAF43D2A750F6DB3D4381CD7496DA870
39AF1FB3D4BA735D35B6A5D39AFCEB1F95EBB1D7099EAF06A1EF5A1FDAD1C1
1F9923607EA4FA0F7AF34835218EB5D57856D64D76F56F251FB843FBA07BFF
00B5F8F6F6AF1A8640AACDB9E915BFF97A9ABAB62D7893E11F833E31691FD9
DF137458756B001BC98E5043C25C60B472A90E8E401964607B57E7AFC4FF00
F825E3F9D3EA5F06FC5623562596C3575385E7385BA8549C7601A2CFAB1EB5
FAF5044B0C4B1A8C002A7519AFD2784F8DB33CA6D0CBAAB8C3F9378FDCFABE
AD59BEE71D7C3C2A6B35A9FCDBF88FF635FDA9FC1EE45CF826E352841F964D
3E586F03719E12276907A7CC82B848FE0B7ED0C3AFC31F147FE09AF7FF008D
57F51D14258D69476648E95FD0991789F9C62E2B9F0B06FBAE65FAB3CAA983
A71FB47F325E1FFD98FF0069EF146C1A67C39D5A1DFD3EDD1AE9FDC8E4DDB4
58E9DFDBD45607C61F817F14BE066A1A3DA7C51B64D3EDF5A8DA482E2DA417
116F438784B8C289146091C8C30209E71FD457950339895D4C8BD5411918F6
AF3CF89FF0B3C17F16BC217BE07F1EE9A9A96957B8254FCAF1C8BF72589C72
922F661EE0E4120FA72F10331A15E33C4D08FB2EBCB7E6F936ED75DACAFB69
D27EAB06AD17A9FCED780355F0AF86D45C697129BA71F35C4843CA7E8DFC20
F70A003E95EC2BF11E329FEB7F5A83E367FC13EFE33FC2ABEB8D57E17097C7
1E1BC96448001A940BD964B718F348E81A1C96EA510715F165EEABAE787EF1
B4CF11DA5CE99791FDE86EA378645C123947008E411D3B57DEE0E9603345ED
F09594FE7AAF55BAF99CCE528692563EC0D43E22AB02164C9F6AF3AD5BC593
5DE4B3E16BC574AD6351D72E8586896B3EA374C3222B68DA6908C81C2A027A
91DABECDF84DFB157C66F8852C3AA78EED24F0568070CC6ED717F2A9E76C76
C7E68CF62660B8ECADC8AC733A997E554FDAE32A282F37ABF45BBF95C70E79
BB451B5F066ED0E996F3023F78598FFDF47FA57DB1E19BE01139AF91EFB40D
1BE1D78F353F05682AF1D8692E90C42462EE4796AC5998F52C4963D064F000
E2BDE7C35AB0D89CD7F2971FD258AAD3C4C3E19EABD1EA8F730AECAC7D51A3
6A206DE6BD7349D4E39635466E6BE60D2756185F9ABD134ED68A630D5F8166
1974A32E689EA42773E848AE3156BED671D6BCA2CBC48CAA016C8AD75F1221
1DABA30BC498CA11E48B25D08B3B792E4B77ACD9EE51012E6B919BC4248E08
158775AD16EAD5E463F1D89C53F7D9A429A89D25F6A630403815C5EA1A8641
E6B2EEF55C83F35723A86ABD7E6A781CADDF61CA61AB5F8C1E6BC8F5FBF043
735B9AB6AA307E6AF23D7F55186F9ABF4AC8B2B775A1C556679CF8C2F8147E
6BE1FF008A6C1ADFED1FDD947EB915F5178B754CAB807935F247C54B909A39
C9E5A4503F9D7F50F86B8474F134BCDA3C4C6CAF167CE7E21B8DD3C1CF407F
A565C52D53D62E18DCC793FC3FD6A08A6AFEC7C0C6D4A28F9D9EE74492D5A5
92B0E396ADA4B5D649B0B2D4A25ACA596A512D0069F99479959C25A5F36802
F1929864AA5E6D4665A00B6D2557792ABB4B55DE5A0099E5AA524BD6A3925A
A32CB4004D2D66C8D9E69EEF9AAAED401F62FF00C13FBC0937C41FDAF7E1DD
82C7BE0D1EF8EB13B70446BA6235CA31CFACC91A8F761F5AFEBAABF073FE08
D5F08247B8F1BFC78D461223554D034E6206189297378C33CF1881411EAE3B
57EF1D00145145001451450014514500145145001451450014514500145145
001451450014514500145145007FFFD4FDFCA28A2800A28A2803C0BF691D40
69BF0D9EE1DB6C42EE1F338CFC8A19CFEAA0D7E56E9F7B2DFDCCBA85C7FADB
973237D58E71F857E9DFED71129F811AF5DB100DA4968E339FF9693A42718F
690D7E54E8F72022F35F1D9CE16D88757AB497DD73A29BD2C7B77864A35D47
BBD6BF473E12BDA2E8A123C07E09AFCBDD1F51F2645607A57D21E0BF89F368
CAA124C63DEBC5A7374EA29AE869E4759F086C7CBD63E377825976A9BFBF96
28F6950B1CD24CC8003D82BAE0F42304715F057C38782CBE21F8AF45BD459A
0BF82D9CC4E032308DA547054F5077A83DABEDEF817E224D5FF688F1D1E047
AD5AC728519E58410863CFAB063C7AD7C23AF6EF0BFC74B881FE559A2BBB7C
71F7925461EFC056E95F9371B61BDB612B417F7D7FE03253FF00DB8EAC2BB4
D1A1E36FD8C3E1478EDDF51F0BCD2F8435094E4FD9544B684E3AFD9D8AEDED
C23A0F6C9CD7CE9AC7EC15F1B74E667F0F6ABA4EB508CED1E7496F31C74CAC
89B013FEF9FEB5FA09A16BC0AAFCD5E9DA76B6303E6AFC430BE2571065CBD9
C2B73C5749AE6FC7E2FC4F5A582A53D5A3F25AD7F622FDA6A791233A558C4A
E402ED7F0ED504F53B49381DF009F415EDBE0CFF00826FFC40D4A5497E20F8
CF4FD22DCE18C7A7452DECA41C12A5A51022375191BC7D6BF4AAD75B181F35
6D45ADF1F7AB8F33F1BF88EA47929B841F78C35FFC99C97E05432BA2B7BB3C
EFE127EC9FF01FE104467D2F455D73557431C97FAB6DBA98ABAED608A54451
82320EC404824124543E18D66FFF00674F1C41A23CB25C78275A936584B231
736EE4126D6466249650098D8F2C8083965627D4BFB6F8FBF5CBF8B6CB46F1
968179E1AD7E33358DF26C70AC51D4E72AC8CB82ACA402AC390466BE2F21E3
3CCE8661F5CC55694F99ABB6EED5B671BECE3D168AD75B337AF8584A1CB156
25F8F7E34934DF0D58FC5EF84B6B06BD0DB4ED15DDC433652340C03C4E155B
2AF8C6ECFCA70403C54BE11F897A37C40F0D5B789B4490F9538DB246DC3C33
2FDF89C7F794FE04608C820D7C19F08BE216B9FB2EFC48BBF853F1109D5FC1
FADAEDDEEBFBABAB573B04CAB921644E9226723E8549D8F1EEA5A6FC02F8AB
703E196A906BFA16B96E975259A4BF2C69282D1798C036C9533903A943F301
9523FA378E7C2DC0E7581866194A5ED5FBC9AD14D3DD6BA47C968A2EF176BE
9E26131F3A3370A9B7E47DB5A86ADC1F9ABC97C713B6A7E1ED534F8CE5EE6D
A645E33F33210A71EC715F17EBFF0012BE2278B2567BAD524B181BA41664C0
807A1653BDBDF7311ED5E7B77A29BA6F36ECB4CFFDE7258F3CF535F33C3DE0
8E228CA15ABE22319269D9272DB5DEF136AD9CC5DD46373DFF00E1EF8955AD
E2F9FB0AFA87C3FE201B57E6AFCC9FB25F696FE7699732DABAF398DCA8FC40
E0FE35DEF85BE387887C39731DAF8961FED1B4638F36201275FF00808C2B0F
6F94FB9E95FD413946A6C78F1AA7E9EC9E33D3744D367D5F56BA4B5B4B542F
2CAE701547F327A00392780335F08FC40FDA627F187894C776CD65E1A8032D
AC2C705A4CF13CD8E371190074407FDE27C7FE29FC5ABEF15B794EED6FA4DB
1CC301382EC3FE5A49EADE83A28F7C93F2633F89FE2778A6DFC15E0B80DD5F
5DB63AE238D07DE9246FE1451D4FE032480718E5AAB2717B7730C525562E9F
467BDEB5ADF8BBE2CF8B20F00FC3EB76BED42F091F29C24683EF4923F45451
D58FB0192403FA9DFB387C0AF0C7C00F0DB08D9352F14EA68A752D488E5BBF
930E795854F6EAE7E66EC17CAFE00FC28F0B7C0BF0BFF66698C2F758BD0ADA
86A0E30F3B8FE151CEC89493B13F139624D7BE9F118C7DEFD6B9FD8469C7D9
52DBF32703808515E66DFC41F00FC32F89D6FF0067F1FF008674FD776A9557
BAB7479901FEE4B81227FC05857CE937EC71FB2C43766FA3F03C424C96DBF6
EBE31E4FFD33371B31EDB715EC5378887F7AB06EFC45C1F9AB4A5ED22AD193
B7A9DD28C5EE8B1E1CF09FC3EF8776CF6BE05F0FD86849270E6D2DD227939C
FCEEA373FF00C089ED54F55D7786F9AB93D43C43D72D5E79AB788C007E6ADA
145B77617EC6CEBBAF001BE6AF0BF126B9BB72AB7269DAE78889DC0364D797
6A3A83484B31E4D7A7428D886CCAD5AEF8624E4D711A6A4973A8DC5CA82521
4099ED9739C67D7E5AB5ACEA01118935FB05FB25FECF7E19D53F661834FF00
1F69CB72DE369A4D60EE5DB34092208AD5E27C641F25448A7FE9A11C8273F4
D91B50AF1A8FA18D5D558FC8C6041A8C9AFAC3E3DFECB7E32F83F773EAB651
C9AC78599B31DF46B968549E16E557EE11D37FDC6E3182768F955E220D7EBB
83CD5DAF17A1E6CA99589A613EB5294A8CA9AFA0A39ACBB983A4337537269F
B68D86BD0866527D49E423C9F4A963959383CAFA526DA5DB5863A14B1549D1
C4479A2FA3FEBF11C6F17744E515C6E8CE6952564386A8D7729C83CD4FF2C9
F7860D7E11C49E1E56A0DD4C27BF0EDF697F9FCBEE3D6A18C4F4968CBD15C0
23079157F4DD44C2FE596E50FE9DAB9F31BA72B556E27781D2E3B0F95BFA1A
FC673AC9FDA536ADAA3D2A552CCF7FD1B58042FCD5E93A76AA38F9ABE60D23
59DAC016E2BD3B4CD6781F357E139F64366F43D3A554FA16CB54E9CD74F6BA
A7039AF0CB1D5F81F35763A65EDC5E4D1DB5AA34B2C870AA83249F602BF30C
CB27B5DB3B6150F5FB6BB7BE9E2B143FEB8FCC47651D7F3E95F50F84E3B5D3
6C9649DD62451D58E3FC9AF16F017C35D6564FED2D5C880B01B507240F73D3
F9D7D0165A359DA0076F98E06373727F5AFC9388F8BB078787D5F0CF9DF5B6
D7F53D0A38793D59AADAD8907FA0DBBCDE85BE45FD79FD2A2371ABCA399920
1FEC2E4FE6D9FE55300074A2BF36C4715E326FDC972FA7F9EE762A112A1B69
24FF005F753480F632301F90E2AB9D1F4E27734218FA9E4D69D15E457CCF13
57F8951BF56CD1412D9196746D34F1E42FE550FF0060D82FFA9531739F94E3
9F5E3BD6D51594319562EF1935F31B8A660B69FAA43CDA6A33AE3A0672C07E
0D91597A84FACC90ADBEAF676BABC08DBF6CF10CEE0080470541C13CEDAECA
90A8618619AF630BC518DA4D3551BF5339508BE8737A3789FC37A7A1B2FB12
E8A19B71558C2C2588C13B9063381D580AE9AFA782EACC4F6F22CB1BF2AC84
3291EC4715897DA35B5DA9F9464D79F5EE8BA8E93234DA4CEF6E4F254728DF
553C1FAE335F5B9771542B497B6D19CF3A16D8FCE6FDA6207F0F7C70B9BC20
AA6AF6B6D720F382517C823392323CA19031D471CE4BBC31AF0289F35753FB
63E9BA8DFE85A3F8C64B6DB71A2CCD05C3C63E4305CED019BB8DB22A80391F
39E6BE5CF0978A018D017AFEA2CBF0F4F31CA28D7A4EF65CAFD569F959FCCF
1A5270A8D33EE5D235B185F9ABBEB1D6C607CD5F2A68DE230557E6AF45B0F1
00207CD5F9BE69C3CEEF43B2158FA36DF5BE07CD5A89ADF1F7ABC1EDF5F181
F356A26BE31F7ABE46B640EFB1D0AA9ECADAD71F7AA84DAD0FEF579636BC3F
BD5426D787F7AA69E42EFB03AA7A25DEB5C1F9AB90D435A183F3571977AF0C
1F9AB8FD475F183F357D160387DDF6319553A1D5B5B183F35790F8835C186F
9AAB6AFE21186F9ABC9F57D61AE1880DC57E9591E41669B471D5AA676B7A89
9DD9B3C0AF983E2D6A43CAB5B3079790BFE0A31FD6BDBB54BE0A8DCD7C83F1
035C1A8F8825546CA5A8F2C7D472DFAF1F857F41F87D94B962A32B691D4F27
1752D13CF75B9489616F50DFA62A9C33F4A66B0E64855D79F2CFE86B1E1B8C
77AFE8CC27F0D23C796E75B1CD57126AE6A2B8ABC93FBD7492740B354A26AC
359FDEA613FBD006CF9D479B595E751E75006A19AA33356719AA3335005F69
AABBCD545A7AAED29340169E6AA8EF9EB5117F7A8CB500399AAD68DA36AFE2
7D6AC3C37E1FB57BFD4F549E3B5B5B78865E59A660888A3D59881596F262BF
6C7FE0923FB2F43AF6B579FB4CF8BE0496D746924B1D0A162AC7ED85479F74
CBCE3CB47091E7BB3301F2A9A4D81FB2BFB36FC18D37F67FF827E15F857602
379B48B45FB6CD1AE05C5F4BFBCB9979E486959B6E7242855E8057B9514530
0A28A2800A28A2800A28A2800A28A2800A28A2800A28A2800A28A2800A28A2
800A28A2800A28A2803FFFD5FDFCA28A2800A28A2803C0BF6A3D3CEA7F00BC
6302A6FF002AD63B8C038E2DA68E627F0D99AFC66D1F500635E6BF7FBC43A2
DA789340D4BC3BA80DD6BAA5B4D6B2823398E7431B71DF826BF9D6BAB6D47C
27AFEA5E15D657CABFD1EE65B49D7D2485CA363D464707B8AF1B36A3CD6669
4D9EC569A8E00E6B7135D36EA3E6249E140EA4FA0AF24B4D5410326B7BC397
D1EA1AC3B39CA40422FD7F88FE7C7E15F099EE2BEAB41D54B5E9EA744753DA
BC1F6DE31B6D765F13786EE67B1D42787C92D6FF00782608E09070707A8C7B
73CD727E31F065DDC6A675AD5DA57D4D599FCF909326F7CEE393EB9E6BF4E3
E04697E16B4F06C177108DAE6419958E0B03E95F38FED29ACF8724D5845A59
8DA445C485318DDF857E4D8C85774A189A9553E76FDDEABA36D79D8DD1F126
9DE369BC3F7F1E9FAE90B1C8C152E070A09E81C76FA8E3D71D6BDD34CF1102
07CD5F22FC429E09ADA556C1041A7FC2DF1ECFA9680B1DE485E7B095AD9D8F
56D982A73DCED2327B9CD7C7F12F08539D058AA31B746BF268EDC2E29DF959
F755AF880607CD5B716BE303E6AF9A6CBC4C081F3D7430F89063EF57E5988E
1E77D8F49553DFC6BC31F7AA0935E1FDEAF131E23E3EFD412789000497E07B
D72472077D8AF6A71DFB496BBE1CD5B44B1F0DDEDBADD6AAD32DCDB3862AF6
AA870D2647387194DA786E4F5515E0BA06926429144B966FE75CBDF788A5F1
6F89AFBC4133165B8908881FE1853E541FF7C8C9F7CD7BDFC28B1B6BCD66DD
67C6DDC339AFEB8E11E1EFECBCBA9E15B77DDF937BA5D92DBF1EA7CAE32BFB
4A8E47B07C3DF803AF78AE259A2B762879CE2BA4F1AFECF3AB7866D0CF3424
003AE2BF4EBE1E69DA669FE18B38EC1542940491DCD725F1A354D2ACFC2D34
776CBE630F941EBD2BECA795F2E19621CD6BD0E7B1F87FE20D29AC26789C60
AD797EA8041BA709BCA0240EB5EE7F10AFA0975198C58C64F4AF0CD4E51835
CB43724F99BC61E2FBBD56EDB4CD3892C720919C281D49C761DEBF42FE0678
53C29F0D3C3719D1254BEBED49125B9BFC7CD3E46542FF0076319F957F1392
49AF8B3C430DA86966489564907CCC000C7EA7BD775F06FC792DBE9773E1DB
99493A74998B27FE594B9200FA306FA02057D2463CF4F963A1A537A9FA3717
8A548FBFFAD583E285C7DEFD6BE5883C64A47DFF00D6AF7FC260B8FBFF00AD
733C21BF31F46CDE2818FBD58377E29183F3D782CDE305C7DFFD6B0EE7C5AC
F908D9AA8E145CC7B26A3E28E0FCF5E7BA9F891A42551B35E7D71AD4B37DE7
E2B2E5D4BDEBA6142C26CE86EF502D9666C9AE56FF0051001E6B32EF54001C
B573F6506BBE2ED72CBC2DE15B29754D5F5394436D6D08DD24923761D801D4
93800024900135D54E95C4D9EC3F01BE13EA3F1F7E2CE97E078158693130BA
D56719022B18981906E1D1E4E234FF006981E80D7F48569696B61690D8D944
B05BDB22C71C6836AA220C2AA81D00030057CD1FB2A7ECEBA6FECF7F0F974E
B82977E28D6765C6AF76A07CD200764119EBE542090BCFCCC59F8DDB47D3F5
F4585A1C91F33193B91CD0C57113C1708B24522957460195948C1041E0823A
8AF84FE317EC39E0FF00161B9D73E1ACEBE1DD524CBFD8DC66C2463D94282D
0E79FBBB947408057DE14577D1AF3A6EF164349EE7F3C1F10FE0C7C47F8637
4D078CB439ECA1DE512E42F996D21EDB264CA1C8E719C8EE01AF2C6888AFE9
AEEAD6DAF6DE4B4BD852E20946D78E450E8C0F62A7208AF98BC6FF00B1DFC0
EF19979E0D25FC3D76E49F374C7F257279C792C1E103D9501C719E98F6F0F9
DDBE34632A5D8FC2D31FE14DD95FA4FE27FF00827A788E16924F07F8AAD2F1
33944BE864B76033D0BC7E70271DF68CFA0AF0DD5BF62FFDA034D67F27C3F1
DFC6993BEDEF2DC8383D95E44739EA06DCFE35ECD1CE69BFB466E933E4AD87
D29767B57BE4FF00B357C73B795A193C15A8965EA522DEBF832E41FC0D5AB2
FD97FE3C5F67C8F065F2ED207EF4243D7FEBA32E6BB166D0FE623D91F3E04F
6A916335F60E8BFB10FC7AD524097DA6DA690B9C6EBABC898638E7FD1CCC71
CFA6783C74CFBAF85BFE09E975E6472F8D3C591A4631BE2D3E02C49EE04B29
18FAF967E95855CFA9AFB435459F99E919AF6DF86BFB387C4EF8BCC87C3DA4
341A64BC35FDD830DA053C121C8264C77118623D2BF5D7C09FB2BFC13F00C9
1DDD8682BA9DF45D2E7516FB53FB108C04408EC56307DEBE8700280AA30070
00AF91CE6BE1B17F1D24DF7EBF81D34A2E3D4FE627C69E16F107C30F186A3E
0AF12A84BDD365788B29252408C577A12012A71C1C55DD275DCE016E457E8F
FEDE7F061356D46DBC6563188DF525D825EC97B12E36B1ECB346A3F142C7D2
BF3CBE0BFC25F157C48F153694F0CB65A7E9F205BE9C8C6CC7FCB35CF591BF
41C9EC0FF3C788395E172EA3531988928D24AF77F979BECBA9EC612729B515
B9ECDF0E3C37E20F1F6A22C7468CF948479B3B03E5C60FAFA93D80FE5CD7E9
07C3DF857A27842D11D23F36E580DF33F2EC7FA0F615A9F0FBC03A2F83345B
7D2F49B65B782118007527BB313C927B93C9AF4AE95FE69F889E2555CD2B4A
8E1BDCA2BA757E6FFCB65E7B9F6384C1A82BCB71AAAA802A8C014EA28AFC94
EF0A28C669DB6AE34DB15C6D14FDB46DAD3D8483986514EDB4DACE50682E14
514540C2ABCF6F1CEA558558A29A6D6A80F19F883E02D33C59E1FD4BC3BAA4
7BED352824B7931D42C8A572A7B30CE41EC4022BF0675FD3B5DF855E37D4FC
0BE23F96EF4B94A6F030B346798E54EBF2BA10C3B8CE0F20D7F4857702CD19
04735F017ED81FB35BFC5CD063F11784E348BC5DA2AB7904E105DC1D5ADDDB
D73CC649C0248380C48FE87F03BC43A580C53C06612B50ABA5DFD997497A3D
9FC9F43C8CCB08E51E686E8F82340F16ABAAFCFF00AD7A8E9DE28040F9EBE0
6D37C4D7DA45D4961A946F6F3C0ED1C91C80AB23A9C32B03C820F041E41AF5
AD27C68ACABFBCFD6BFACF3BE0C77E68ABA67874F127D9D6BE26181F3D6BA7
8907F7ABE54B3F178207CF5BB178B011F7EBE0B11C22EFB1D6B107D22DE241
8FBD54A7F120C7DEAF026F158C7DFAA3378AF3C07CD634F84DDF61BAE7B5DE
789060FCD5C66A5E25183F3D7974FE21965CE0E2B265D449C966AF7B07C371
86E8CA558EB2FB58927272D815CBDE5F800F358F71A9000F359DA6D9EBDE2F
D76C7C2BE15B29754D5F5394416D6D00DD24923761D8003924E02804920026
BEB70194372518A39E750E3FC6DE25934FB265B4469EEE73E5C11202CCF237
4000E4E3AD5DF853FB0AFED33F17AE23B987C332786F4C9B0CD7FAE6EB24C3
1EAB13033BE46482B195E9961919FE86BE03FECDDA27ECEBF086682E961BEF
16EA2629F53BF0809F319828821623708A2562A3FBC4B3E06EDA3D52C2F010
326B1F113C41C6F093A797E0E945D5A90537295DDAEE492495B556BDDBB5DE
C184C2C6BDE727A23E00F80DFF0004D3F82DF0CD60D67E2367C7FAFA00C45D
2797A744DC1F92D413E66391995981E084535F94BFB777EC63A97ECDDE2B3E
31F06C2F75F0F35F9DFECAC03B9D3663CFD926739E08CF92E4E58020FCCA4B
7F4EA082322B27C45E1DD03C61A15EF863C53A7C1AAE93A94661B9B5B9412C
32C6DD432B0C1F51E8704735F8F70778EB9C60B34FAF636ABAB196928B7656
FEEADA2D6EAC9754F73BF119653943962AC7F1431DC153CD5F8EE7DEBF48FF
006C6FF827578C3E0C4F7DF107E11C13F88FC09979A5810196FB4A5C92448A
32D2C0A3A4A32CA01F300037B7E62AB915FDFDC31C5580CE30B1C665F514A2
FEF4FB49747FF0EB4D4F96AD4254E5CB34746971EF53ADC7BD73C929F5AB0B
29F5AFA231377ED03D697ED158C243EB4F127BD006A19FD0D34CC4D500F4BB
E802D97269A5EAB6FA6992802C97A85A4AAED2815F5AFC05FD94FC51F1445B
F8AFC58B2689E13CABABB0DB717CBD71029E8847FCB53C73F206E76F959CE7
785CBE83C462E7CB15F7B7D92EACD29D294DDA28E1FE06FC15D47E2A6A8FAA
DF836FE1BD3645FB54BC83337DEF22323B91F7987DD073D4807FA35FF827F5
D5BD968DE2FF000959A2C169A7C96334112E02AAC89246428EB8022507F0AF
89AF74DD1BC35A341A0787ED23B1D3ECD0470C310C2AA8E3F13EA4E493C939
AFAA3F60FBB917C7BE29B1006C9B4E8A427BE6398018F6F9CE7F0AFCA387B8
BEB6679BD3AB25CB0D5463E567ABF37FF00EEAB8750A6D753F50A8A28AFDB0
F3428A28A0028A28A0028A28A0028A28A0028A28A0028A28A0028A28A0028A
28A0028A28A0028A28A00FFFD6FDFCA28A2800A28A2800AFC9AFDBFBE0ADDE
91AAC3F1D3C35016B3BBF2EDB5955FF9673280905C631F75C011B9ECC13A96
26BF596B3F57D274CD7B4ABCD135AB68EF6C35085E0B88255DD1CB148A55D1
81EA08241A8A94D49598D33F99CB1D755941DD573C31E2A5B4D56F2DA47DAE
B313F837CC3F435F417ED45FB1FF008AFE09EA177E31F01C336B1E05918C99
5CCB3E9A0F2527EE621FC32FA70F8382DF9FFAE5D6A70BFF00C241A52191A0
4C5C46BD59074603B91DC7523A74AF8CE23C9BDB50717D0DA133F42F48F8B5
ABE9D67F67B4BD68D08C60362B8CF1078EA4BC2F2DC4C5D8F726BE26D3BE2B
C52C43F79CD32FFE21CB72A446FC1F5AFCCE1C29353D51B7B43D37C6FE2C59
237447C939AE1FE1F788D74D17A0BE3CD9B3F928E6BCE52EB51F12EA29A6E9
C3CD9E53CB1FBA8BDD98F603F5E839AD9F13F81358F0781A9E9B335ED8705E
4030D19FF6D46783D8F4EC71C67DCC464F0FABBA0FA8E9C9A773EA3D3FC60A
CA3E7FD6BAAB7F16838F9EBE22D2FC5772B805F915DC5A78AE7207CDFAD7E7
F8DE11B3D8EE8E20FAD478A811F7EB1BC45E2D78B41D4A489FE74B698AF3DC
21C578043E259D872D45FEAB3DEE9D756C1F99A27418FF006948AF3F0FC351
8558B92D2E8A956D0BBE1CBB510A007B0AF75F077899B49BA8E756C1520D7C
9BE1BD615A14C37615E9B65AB60021ABF6CC4E1F53C13F4F7C2DFB4D6A5A36
9AB6B1DC1000C75AF2DF88BF1DB53F13EF12CE5837BD7C691EBAE171BF8FAD
413EB0CC396AE254181D5EAFAC35D48D23B649AE0352BD183CD55BBD5060FC
D5C76A5AA0C1E6BAE8E1C0C7D7AF06D6E6BCF740D5E5D3B5C9248DF689A320
FB90411FD6ADEB5A8990940697C09E09D4FC73AA5E47A6E57EC7102CFB7203
3B7CA0FA642B7E55F4D9761A526A315AB0E6B6ACF4EB4F14DC103F795BD0F8
86661CC95E7BABFC3CF1FF008672F71A7C97102F3E6420C831F41F30FC4573
11EB9340C6399591978208C115DB570728BB4958D233BEC7BBA6AFBBAB66AC
0D547F7ABC4A3F12281CBD4C7C4E98FBF5CFF572B98F627D5D40E5AB2AE75C
5507E6AF16D47C79A758A9FB55D2467D0B73F80EB5E5BAE7C5B2E1A2D22332
37677E17F01D4FE38AE9A197549FC28895548F7BF1278DEC348B57BABC9C46
8BEFC93E807735FA65FF00049AF891F023C57A878874CBC892D7E2B09A47B4
FB615669B4A11AE7EC448C2BAB6F33A2FCE5769CB20213F9EDD4F55D4358B8
373A8CCD33F6CF403D00E82B4FC15E33F147C3AF16695E39F05EA32E95AE68
970973697309C3C7221FC8A9195653956525581048AF768E551A6AEF73075A
ECFEF228AF977F640FDA5741FDA97E0AE93F112C0C706B308167AD59A1FF00
8F5D46251E680A492239322488927E46009DC1B1F51566D58D028A28A40145
145001451450014514500145145001451514F3C36B049757322C50C2A5DDD8
E1555464924F40075A52924AEC0F25F8E565E1DD5BE1DEA1A2F88339BE016D
027FAC174BF346CBECA465BD5723BE0F85FC3EF07D9E876115BDBC7B428049
C72C7B927B9356758F104DF103C52FAC904584198ACD08C111679723FBCE79
3E8303B57A3E9F6EB040A00E4D7F967F494F18E79E639E03093FF66A4DA8DB
ED3D9CDFE51ECBCDB3EDF27CBD528F34BE265E550A001D052D1457F281EE85
382D2A8AB091E6BD2C1602551D919CA646A84D4CB11AB91C39ABA96D9ED5FA
4E51C1552AA4EC7254C4A464F92690C55B7F65F6A8DEDF15F4188E009C237E
53258A461B475115AD5922C55374C57C0E6D903A4DE874C2ADCA446292A561
8A8C8C57C5D7A3CACE98B128A28AE6282B9CD5ED432920574754AF537446B7
C3D4719262923F34BF6A1FD93B45F8A7E6F8AFC2822D27C551A9DEFB76C37D
8E8B3ED190E3A2C98271C364636FE3FF00887C3FE2BF87DADCDE1FF14D8CDA
6DEC0798E418C8ECCAC38653D994907D6BFA60D72DC10DC57CC1F15BE1BF84
FC7BA6B69BE2AD363BE89725198159233EA8EB8653F43CF7E2BFAE3C27F1A3
1380A51C0E3D7B5A2B6FE68FA37BAF27B7469687838ECBA327CD1D19F89163
E209703F786BA7B7D76520664AFB7B5CFF00826CF8A353D02D7C49F0DFC470
192E95A4FB06A6AD1154DC7614B8883062CB83868D71FDE39E3E51F167ECC1
FB47F805E41AD781B519E189883358C7F6F8B0067716B6326D5C776C63A1C1
E2BFA472BE3CE1CCCDB8E17170E7DB964F9657F495AFEAAE8F22786AD0F8A2
CC58F5766EAF9AB8BAA0FEF5797EA09E20D024F275CD36E74F932576DC42F1
1DC3A8C380723BD6B69563E2DD6847FD8DA2DF5FF9A094FB3DB4B2EE03A91B
14E715F47572DA6A3ED39972F7BAB7DE62A6F63BD3AA8C7DEAA73EB2AA3EF5
773E18FD99BF68FF001794363E08BEB085B932EA4A2C155738C959CA39FA2A
924720639AEAB59FD927E2778497ED9E2C83EDB6E80338D398C8ABDC86242B
F1DF0807FB55F315388725A757D8CB174DCFB2926FD347A3F53E9786B85F17
9AE2A184C3F2C5CBACE4A11FBE5BFA2BB7D11E63E08F0E6BFF0013BC536DE1
1F0CF97F6AB83F34B33EC8614CF2EEDC9C0F40093D81AFD26D73F64CBBF821
E0A83C67F0CBC4BA88F165B447ED77F6B7125A3CF13805E38562605138C85D
CC5BF889E31F0AE92AFE189EDAF7C36A2CAEAC1F7C4578C9EEAE7A90DD0E72
7BF515FA89F0DBE26789FC6DE07B68AFEC5A29AE500DD38E621D0ED5EF9ED9
E3BF22BE6F3D8E6B88C6D0AF80C47B2C3D377925F149F67E56FB9F46EC7EA7
C5FE1DE3F81F1D86C4D7E5ABCDAC656BC1B5F141A7EBBE974EEACEF6F8D34B
F8C1FB456BC1748D33C5DAE6A6136E6369E49FEE9047985F764647F11E7BD7
D03E0CD73F6B7B468E69B53B49A1041F27518E37C83D72D0207FFC7C57D3BE
13F00E99A6C3B2D2D522DEC5DB6A81966E4938EE7D6BD5ED7C2E360C45FA55
67188A18A97FB5C15476B5E4AEEDD937AA3B78ABC61A18DBC3059550A50F38
294BEFB452FBAFE6797F853E2AFC42B458E1F1C6856920FE39F4E9A41F9433
29FF00D1B5EB7A67C4BF086A52081EF3EC331380974A62CE7A618FC873D806
CD51BBF0EA2A90D1E3F0AF3CD77C316F2A3068C1FC2BE071BC0395625B74E2
E9BFEEBD3EE77FC2C7E3B89C7CA73E7E54BC96DF99F4723ABA87421948C820
E4106BF377F695FF00826BFC24F8D135EF8AFC02EBE05F165C1F31DADE3074
DB99339669AD971B1DFBC9115E72CC8EC4E7D46CBC49E28F87571E669121BA
D38125ECA527CB23BEC3C98CFD38F506BDF7C13F173C1FE3A2F6BA4DD88F52
8177CD633612E2319C6EDBFC4B9FE35257B139E2BE46581CFB85ABFD7B2DAA
F97F9A3B5BB4E2EEBEFBAF3B8A34E1884E2E37B6AFCBCCFE5AFE35FEC8BF1F
3E014B3CDE3BF0CCCFA442D81AB58E6EB4F61D89950663CF612AA31F4AF9B1
6422BFB6A6BB85D4A380CAC3041E4106BE59F899FB1C7ECBDF1559EE7C4DE0
3B0B6BE7258DDE9A1B4E9CB9EACED6A631237FD740FF00A0AFDCF85BE95BEE
AA79CE11DFF9A9FF00F2127F94BE478F5F23EB4E5F79FCA0ACB52896BF797C
61FF000495F83F7EEF27827C6BAC688581C25E4705FA2B1CE30145BB6D1C70
589E0FCDCF1E277DFF00048ED7E1B875B0F89D6B34031B5A5D2E489CF1CE51
6E1C0E73FC47D7DABF63C07D20B856BC6EF12E0FB4A13FD22D7E279D2CA6BA
E87E4479A3D697CD1EB5FA27F15BF60AF097C0CD08F88FE2A7C61B4D36D9C3
0B782DB4B6B9BDBA9063E4820373196EBF336E0ABC6E201CD7E76DFC761F6C
9174B795ED14811B4CAA923003059954B05C9C9DA19B68E3736327F44E1CE2
DC066D4DD6CBE6E705D796497C9C924FCED7B75392B509537698C330151995
9BEE8A411A8EB5285F6AFA3B989ECDF08BE217C36F8777E759F16F80BFE135
D4508307DAAFC41690E3BFD9BECF2891B3DDD8AF70A0F35F64FF00C3C2ECEF
894BEF03C9691F00793A82CC71DF8304638EDCFE55F9A3B6936D7C8673C0B9
666157DBE2E9B94BA3E79E9E8B9ACBE48E8A78A9C15A2CFD3687F6BAF861E2
39563BD17BA43375373086407EB0B4871EE5457E9D7FC13B35BF0EF8AFC53E
2ED7340D462D42382C6DE2CC2EAC17CD94B10E072ADF20C038EFC57F31A56B
67C39E25F12783B59B6F11784755BAD1355B360F0DD594CF6F3C6C0E415923
2AC391EB58E57C0983C162215F0EDAE5E8F55B7DE39E2A525667F76B457F38
BFB34FFC15DFE20F846EAD7C35FB46D99F16E8876C7FDAD671C70EA76EA060
3491A848AE071CFDC7EAC59CFCA7FA00F86FF133C07F177C2365E3AF86FAD5
BEBDA25F8CC7716EFB82B60168E45FBD1C8B91B91C0653C102BED8E63BAA28
A2800A28A2800A2BCE7C4FF183E127823531A2F8CFC6FA1E81A8B00C2DB50D
4ADAD6621B041F2E5915B07231C577B657B67A8DA437FA7CF1DD5ADC2878A5
8983C6E8C321959720823A1149496C059A28A2980514514005145140051451
400514514005145140051451401FFFD7FDFCA28A2800A28A2800A2B2357D7F
45D06113EB37B15A29FBBBD8066C7F757AB1F600D79C5C7C6CF07C2EC228AF
6E107474830A7E81CAB7E60578B99711E03072E5C5578C1F66D27F76E75D0C
056AAAF4E0DFC8F5B90C623633604601DDBBA63BE73DABF98FFDABBC6BE02B
FF0089DE24BCF86BA45BE8BA299CC5125B2EC49DA3F95E7099DAA246CB2AA8
5017195DD927F5F7F69BFDA6342B2F86F73E1FF0A4D345AAEB00C330923689
E0B523E720FDD264FB9C13C6EE9C67F9F0F1E799E24F1038D3C3492CD20558
D14B9924638015472493D86726BE6731CFE8E32AC68E1E4A505AB6B54DF45F
E7FF0000E8797548439A6ACFB163E0BFECF7F123F688F13EA3A4FC37D21F50
974C83ED77AE92C56EA88CDB157CC99963DEEC4ED527242B1190A6BE8D4FF8
27AFED0915D7D9A4F066A72957D841BBB3F2C9CE3EFAB818F70D8C739EF5FB
6BFB157ECEB07ECE7F05EC743D42203C51AE11A86B327048B875C2400FF760
4C263246FDEC3EF57D775F5D86C0454173AD4F39B3F96EF1DFC12F883FB3A6
A5A7697E3DD03FB0BFB6A379AD809E1B81208582B83242F202CB952416C80C
BEB4BA76B49247B5C86561839EE3DEBF7F3F6ABF80565FB42FC28BDF0AC7B2
0D7AC0FDB348B96C0F2EED14808CDD44728CA3F6190D825457F3533CDAD785
75ABDF0CF896D25D3755D3267B7B9B699764914B19C32B0F506BC3CDF2AF7B
9A3B1B5399D2F88BE1A68FAC3B5F68520D3AE5B931E3F72C7E8394FC323DAB
C8B52D27C49E1A7C6A96CCB18E922FCD19FF00810E3F0383ED5EC367E20520
7CD5B91EB8ACBB4B641EC6BE6A78692D1AB9B5CF01B6D7C0C6E38ADA875E43
FC55E897FA2F8575225EE2C230E7F8A3FDD9CFAFC8464FD6B989BC11E19273
1CB7117B2C8BFD54D724F2E84BA15CE799A6A834CD56580362290978FD30DD
BF03C577D61AFAB28C3564EB7E03D1EE6D196CEEE64BA41989A465650DE840
50707BFE75E40BABEA3A35D369FAA46D0CD1F63D08F507B83EB5F5785A4AAD
349EE8E3A91B33E958F5C18FBD4D935B18FBD5E150789C1030D561BC42EE38
269FD44CCF52BBD6C60FCD5C7DFEB064C8535C8CBAA3C9CB3566CFA8281D6B
7A785480D0BDBCC0249AFD83FD8EBF677BAD23E1841E22F105B18B52F11B7D
B191C61A3808C40A476CAFCF83C8DF83D2BE14FD97BE07C7F10BC4F67E2DF1
AC5B7C396322C91DBC839BE914E54107FE5883CB13C3FDD1C648FDFAF0FEB7
A7C162A323A57D8E498071FDECBE473D69F447CD3E34F863676B6AE4440601
ED5F9E5F16BC23650B4A5A252573824035FAA7F123C5363F66902B0E86BF30
7E306BD6F2BCC15877AFA16AFB9CE7E7BF8BE29ACA5716F2344067EEB11FCA
BC5B52BEBF772B25CCAE3A72EC7FAD7B1F8DAF5649A420FAD787DF3EE726A1
C22BA14A4CCB3D2994F3D2995320EA30F5A8CF5A7D31BAD6323489FA2FFF00
04C8FDA46E7E037ED17A7787356BB31784BE21490E91A82336238EE5D8AD95
C9CF03CB95F6331202C723B1E82BFAD8AFE04125920952785CC7246432B29C
3291C8208E841AFED7FF00643F8C0FF1DFF66EF017C4DBB97CED4752D3921D
41BA16BFB3636D74D8EC1E6899D47F758751C9F2F111D6E6F067D214514573
961451450014514500145145001451450015F3CFC79F16FD9B4FB6F04D84B8
BAD53125CED3CA5B29E01F4F31863DC2B0EF5F41C92470C6D34CC123405999
8E000392493D00AFCF4BBD7A4F1B78D753F13B92D0DCCC56DC1FE1B78FE58C
63B65464FB935F81FD2278DA594E432A145DAA57BC57947ED3FBAD1F99EAE5
186F6956EF647A4F852C1628D001E95EA0A30A07A571DE1F882AAD7655FE48
E6955CEAB6CFBCA6B40A50292A4515C9421CCC7264A8B9AD1862CD55896B62
DD3A57EBFC1992AAB357470622A58B50419ED5A71DB7145BA0EB5AB1A0C57F
66706F06D274D368F9EC4621DCA06DC63A554961ADE6518ACF980C57D27107
0BD18527A18D2AEDB39D9A3ACA956B7AE3A1AC69ABF9338E32E841BB1EEE1A
66638A81AACBD573D2BF9E733A69499EAC761945145786CD82A1B81FBB3535
41707119A70DC19C06B0990D5E4BAA5835EDD47671FDF9DD631DF96381FCEB
D775620835CEF86AC7ED9E2681D8656DC34A7F0181FA906BEDF2CC43A54DD4
EC8E59C6ECF5EB5B78ACEDA2B4806238115147FB2A303F4A9E9F8146057C8B
526EECD86514FDA29A462A2C031955861864573DAA6816B7B1B7C80E6BA3A2
B4A55A50778B2949A3E05F8C5F0CBC10FE21D3E06B08C6AB33F9EEC8366620
719900E1B730C0CFA1E7D7D83C17A4436F0C51C68000000057879F130F1A78
EF58F130712413DC325B11D3ECF11D9111D7195018FB926BE90F06C8865881
F515FE84F0CE4F5B0194D0C35793752C9CAFD1BD6DE91DBE57EA799996778A
C5F253AD5652842EA29B6D4575B5F6BF91F42F857C3686049A65CE7A0AF4B8
B4855400002A0D096316F16DE9B463F2AE96BF75E03E05C162709F59C4AE66
FF0003E57198B92972C4E2B52D262910A328CFAD78B7882C842EE8474AFA17
5295003EC2BE7BF16DFC6F7126C3DEBF37E2ECA686131F2A5867A23BF09525
285E4786F89ED63911C119AF8B3E2A683751489AD68B3C963A9D8379B6F710
318E58DC77575C11F81AFB4BC4172A55ABE65F1E32490CA3AF06B6CA35F764
AE99DD84C7D5C2D7862683B4A2EEBFAEA9ECD6CD68CBFF00053F6C65BB7B7F
09FC5A75B6BBC88A2D502ED8A43D00B850308C4FF18013FBC17049FBA46AE9
2207470CAC320839041EE2BF037C4B6C916AF75128E33B87E35EF3F03FF690
D63E1E4D6FE15F16CCF7BE1A6211246CBCD65D86CEED10EE9D40E57A6D3F9A
71978594DCE75B2F8DA4B78F47E71ECFCBEEECFF00A3F887C35A59A6554788
72385BDA479A54D6D75A4943D249AB75E96D9FEB8C9AAFFB55F27FED43FB57
F85BF676F0A79F36CD4BC53A9237F66E9BBB973D3CE9B1CA4287A9EAE7E55E
7256FF00C58F8E7E17F859F0DAFBE256A53ADDD94312B5AAC4C0FDAE69B886
38C8C83BC9C923385CB7406BF9BFF891F117C53F163C6BA9F8F3C6173F68D4
B5390B1032238631C245129276C71AE154649EE492493BF83FE0FACE310F15
8D4D61E9BD56DCD2FE5F24BED3DFA2DEEBF9A332C73A4B917C5F90EF88BF12
7C6DF16BC5775E34F1FEA726A9AA5D900BBE152341C2C7122E151147455007
7392493C585A4515A56766F70E028AFEEFC361A9D1A71A3462A318AB249592
5D9247CB4A4DBBB2AA44CDD055D4B099FA29AF50F0E7822E75065023273ED5
EE3A47C1ABCB88837907F2AD847C84DA74CA395AA8F6EE9D457D93AB7C21B9
B58CB184F1ED5E29E20F07CB62CDB9318A00F192BEB5191D8D6CDED9B42E41
1D2B25862802BB0AFA23F66AFDA87E287ECB7E398BC5BE01BC32E9F3BA0D4F
49998FD9350854F2922F3B5C027CB9546E43EAA595BE7A6F5A89855203FB54
FD9AFF0068DF017ED3FF000CACFE23F816430927C8BFB09581B8B0BB500BC3
263A8E43238003A90700E547BFD7F181FB277ED39E2DFD957E2C59F8F742F3
2F348B9DB6FAC6981F6A5F5993C8E72A258CFCD13E32ADC7DD6607FB15F03F
8D7C35F11FC1FA3F8F3C1D7A9A8E89AF5AC577693A7478A55DC323AAB0E8CA
7956054804114C0EAABE79FDABFE28EADF05FF00677F1C7C49D007FC4D34AB
254B46DBBFCAB8BB952DA294AE181113CA1C8231F2F3C64D7D0534D0DBC6D3
4EEB1C6832CCC42803DC9AF1BF8A43E1CFC4CF01F883E1BF89247D434CF10D
A4D6570B6AA5982CAB80F1C98D8190E194E78600D7959A67384C1C39B155A3
4EFF00CCD2FCCD29D2949FBAAE7F333A2EBFF06757F859E21D5FC6EDA96ABF
12F53BA69A2BA790BA3190EE77959B259D892492724D7DCBFF00049EF8D9E2
74F891AF7C08BDBA92EBC3975A6CDAB59C4C4B0B3BA8268D6411FF00712659
4971D37AA90016627D5BE0E7FC13DFE0EF83BC37369DF12F4CBBF1B6AAF34C
16EBED135842B007221291412A90C5002E1D9C0624024015A7FB3DFECADE27
FD997F699D43E28F81E582E7C09AC594FA7DC585C967D4A1B79764A0452ED0
8C5678A33CB0CA641CB00D5F8FE51E2BE4AF1D3A32ADCA95FDE7F0CB5B68D3
6FEF48F4AA65D5796F63F6028AE120F88DE197758EE5E5B366FF009ED1301F
8B2EE51F89C575F65A8586A508B9D3AE23BA88F1BE270EB9FA826BF5ECB33F
C0E355F095E33FF0C93FC11E754A3387C4AC5CA28A2BD7320A28A2800A28A2
800A28A2800A28A2800A28A2803FFFD0FDFCA28A2800AF31F895F1013C1763
0DA58059F58D4322DE26E55547DE95C0C7CAB9E071B8F1D3247A757E78EA9E
2D1E31F895AE6B4F2F9B04572F696DE82DED98A2EDF6620BFD58D7E79E25F1
4D4CAF2D73A1FC4968BCBBBF97E6D1EE6419747115D29FC2B567AC68BE1ED4
7C4776DA9EA92BDEDDCDCBCB29CFE03B003B01C0EC2BA4D7FC2F6BA168B77A
C6A0EB1DAD8C2F3CAC47DD48D4B31FC00AF42F01FD8C69B138C1CF5FAD69F8
FB4ED1FC47E16D4FC3DA9B15B5D46DA5B7959080CAB2295CA920FCC3391EF5
FCCB97F0751C66595337C6E2AD3BED7BB7EBD6ECFA8C5E795A9E2150A31B45
69B1FCEC7ED01F1366F106B37B78A761B8722341FC118E1578F41D4F739359
1FB22DA681A47C5BD1FE29F8F74E6BFF000CF86AEC3B00377FA5953E54C14F
DFFB3B15948C8390B8CF4AEC3E247EC8BF16FF00E12EB9B6B2B9B1BDD33CC2
22BD33F96BE593C3491905C301F7828619E84F5AF5FBDF0A68DE01F045A784
34BC34565161E42306595B99246F76624FB0C01C015FA460B88B0B82C3D1A7
8092949DBCEC977FCAC63570352ACE52AC9A47EF1D9DE5AEA16905FD8CCB71
6D728B2C52210C8E8E372B291C1041C83566BE15FD80FE235D78BFE12DFF00
84B5191A5B9F06DFBDA44CC4B136738F360049FEE13246A3A045503D07DD55
FD1581C52AF4615975573E16B53E4938BE815F9D9FB6FF00EC57A4FC76D366
F895E0A9ADF46F1D6970665967710DB6A16D0AFDC9DCF092228C4729E31F24
876ED68FAAFDA9BF6EEF86BFB3A2CBE1AD3157C55E362BC69D04A162B427A3
5E4C036CF511A82E78C84521ABE1CF00780FF68CFDB76EE0F1AFED05AFDD69
3E049184B67A258E6D20BA43C822252488FF00BB24A5E4607E521486AE3CDB
36C3E169B9567A053A6E4F43F357E1D781BE23FC45D79BC35E07D1E6D66F62
2449E415314601C65E62446AB9E84B007B66BF4D7E17FF00C13A3C4FA84305
F7C4CF100B3DC016B4D3977B0CF633C83683EA04647A13D6BF53BE1E7C2AF0
77C38D06DBC3FE15D2E0D36CADC616285028CF7663D598F762493DCD7A6A46
AA30A302BF1FCC78B2B57935878F2AEFBBFEBFAB9E842825B9F1CF857F627F
823E1B4563A0A6A130C65EF5DEE77639E51C94FC9457AB5B7C00F861648896
7E19D3211172812CE15DA473C6178E6BDD36D1B6BC1A8F133D6536CD74EC7C
E9AB7ECF3F0D7504905D78634D9C4A72DBAD216C9CE7272BC9CF35F2FF00C4
EFD837E0878DAD1A1B8D07FB3A65C949ACA4685D09FEE8E53F02A457E9495A
A5736915C2957519ACE9E6189A12E68CDAF983845F43F9B3F8A1FF0004E1F8
8FE15967BCF875ACC1AE59AE592DEEC7D9AE801D1438CC6E7DCF963DABF3C7
C450EB7E0BD6AE3C39E2BB19B4BD4ED0ED9609D0A3AFA119E083D986411C83
8AFEC4F5FF000F4522B656BF3F3F6A5FD953C25F1C7C392DADE44B63AEDA2B
1B0D451019216E4847E85E263F7909F7186C1AFBEE1FE34E69A863758F75BA
FF00338EB6174BC0FE7547885243B623926BBCF0843A75C5EC773A98170148
211BEE7E23BFE3C7B579F78EBC09E28F85DE2EBFF0578C2D0D9EA7A7BED619
CA3A9E5648DBF891C72A7F0201C80ED1F576B66183D2BF6CC26068595486A9
ECCF267396CCFD48F873F1122D3A28911C22A00001C00057D2D6BF1A963B5D
A26C71EB5F8FFA378DE5B655024E9EF5D87FC2CB9C478F34FE75EAF318D8FB
DBC6FF0018BED3138F3BD7BD7C4BE3DF1C9BD793F799CE7BD79A6B1E3E9EE1
483213F8D794EADAFC970C496CE68E60B0CD7B5437123126B8499F731AB173
726463CD5063DEA5B2921ADE94C26949A8FAD66D820A61E4D389C0A6565265
C46B57F495FF000450F1DC9AB7C1AF1EFC3B9A4323786F5A86FA30724A43AA
41B028E70177DABB600EACC79CF1FCDA1EB5FB75FF000446D56587E29FC4AD
102E63BBD1AD2E0B67A35BDC94031DF3E71E7B63DEB82BEC6B03FA39A28A2B
8CD428A28A0028A28A0028A28A0028A28A00F1DF8F1E24FF00846FE19EA8D1
B62E3530B61176E6E387E7B6230E47B8AF903C1B088E18C7A015E9BFB596B2
CFA9F85BC371B7CA3CFBB957D49DB1C67D3FBFDABCEFC2E408D3F0AFF3CFE9
499DCB139C3C327EED28A8FCDFBCFF003B7C8FACC9295A9DFB9EF3A270A2BA
BAE2F469385AECD4E4035FC398F8DA67D343616A65EB50D4ABD69E0FE2148B
F0D6C5B9E958911AD385F18AFDCF8271B1A72573CDC4C6E74703015A4920C5
73D14D8AB82E38AFEB9E1AE2CA74E9AD4F02B506D9AAF20C567CD2542F71C5
50967CD65C4DC6B0941AB8E8E19DC8E77EB5912B558965CD67C8D5FCA5C639
E2A927667B987A76207350374A7B1CD444F35F8763EB7348F4A284A28A2BC8
66A154AF5F6A62AE93819AC1D426E0D6D421790A4CE475593EF56A780ACB7F
DB6FC8CEE65894FA606E3FCC5733AA4DC1AF5DF05E966DFC3D6A587CD38329
FF00819C8FFC7715FA3F0FE47531B7A30471D5AAA3AB2F792690C445743F63
F6A824B6C57D562FC38AB4E1771318E2D3300A114CAD1962C55375C57E6B9A
E4F2A2DA68EB854B9011DEB93F1D6A4745F04F88758127926C74EBB9F7E09D
BE542CDBB039E319AEBAB96F1AE9326BBE0ED7B428504926A361756CAADD18
CD132007D8E6BCACB39238AA5ED7E1E657F4BAB953BD9D8FCC3F86DA8A2DAC
201EC2BEAFF0C6B4226460DD2BF3C7E1CF8954C11FCFE95F53681E221B17E6
AFF51B3AC03E767C95299FA25E12F1D5A9B6482E5F1B4706BD05FC5DA708F7
7DA0115F9FDA6F8A1900DAF8FC6BA45F16CACB8F33F5AF230B8FC661A2E9D1
A8D265CE9C24EED1F4DF88FC750B46D15B375EF5E1DABEB7E6B3316C935C2D
CF890B0397AE4B51F100C1F9ABCF5859CE4E53776CD6E92B22FEBDAA8DADF3
57CE7E34D514C727CD5D66BDE211B5BE6AF9D3C67E22051C06EB5F5394E01F
3239AA4CF0FD7E5171AE4CC3A6D3FA11FE35CADDC60835ABE61B8D4E57CE76
A1CFD588FF000AC7D7AE869DA65E6A04022D6192539E9F2296FE95D59AE98B
E48EF647FA0FE01D3B70642AD6D22A551AF44DDFF1B9F2F7C5DF887AF6BD25
BF811B50965D0B40964786DC9CC6B732002561DF8C6003C03B88C6E6CF8DA8
A88C8F33BCD2B167725989EA49E49A996BFA0F2CCBE9E16846853564BF17D5
FCDEA7F9E5C479CCB30C7D6C6C95B9E4DDBB5DE8BE48B50A6F602BD6BC19A0
FDB6E231B73922BCC2C1419457D2FF000CA2885D445B1D45769E29F65FC1DF
8551DFF925A2CE71DABF433C31F02ED0D92398074F4AF0FF00819358C2B017
C76AFD15D0F5FD321D3D4657EED2607C45F10BE115A595B484440601ED5F9B
BF157C250D9C9280806335FB13F15FC4FA7B5B4A1587435F941F18B58B79A5
9B691DE803F3EFC4B64B0CCE00AF3E94609AF4CF154EAF3BE3D4D799CC7E63
4C0AC7A544D52355ED2347D53C41AADA687A25B3DE5FDFCAB0C10C632F248E
70AA07B9A529A8A7293B2048CD82D6E6F6E62B3B389E7B89DD638E38D4B3BB
B9C2AAA8C92493800724D7F4A3FF0004DDD2FF00684F845F05AF3C11E3BB28
6DF4EB9BBFB668F0DC48EF756093826E237880D8AAEF891137655DA42EB96C
0F3FFD8DFF00628D07E145BDBF8BBC59147AAF8C2E10132B2868AC830E63B7
C8FBD838693A91C0C2920FEACE93A4C1630A85500E2BF963C49F1CEA4252C1
E4CECB673EAFFC3D979EFDADD7E8305952B73D5FB8CE5D16EF54956EF5FB97
BD941C8DE7E55CFF007547CABF80AE820B0B5B700471818F6ABA07614F0B5F
CCF89AF88C5D47571137293DDB6DB3D7BA8AB2D088228E8290A21EA2A7D9ED
4856B39609A5B0B9CCDB8D3AD6E14878C1CD71D7DE177B59BEDFA44CF6970B
D1E1628DF4C8EA3D8F06BD048C53480452A188AD879AA9464D35D8BE6BAB33
8AD37E296B5A04C2D7C5D6E6F6D7A7DA6150254F77418561EA57047A1E95EE
7A66A9A7EB36316A5A55C25D5ACE329221C83EA3D883C10790783CD78BEBBA
3C5750B1DA0D78E5BEAFAFFC37D59F55D07F796F291F68B5627CB9941E7FDD
7C70180E3BE4707FA47C39F1BEBC251C366AF9E1B737DA5EBFCCBF1F37B1E4
E372A8B5CD4B43EDAA2B98F08F8BB44F1AE8D1EB7A1CDE644C76C88DC490C8
3EF4722F661F91182320835D3D7F58D0AF0AB05529BBC5EA9AEA8F9D69A766
1451456A20A28A2800A28A2800A28A2803FFD1FDFCA28A2800AFC5CF06F8A0
DAEA37B6D39F2E686E6647523690CB21041079041EC6BF68EBF0DBF6A9D075
1F825F1E753BA7431E83E3291F54B0973F2996420DDC47D192662C076474F7
AFCC7C51C8A58DC1C2515F0B7F8FFC31F43C3B8C54AAB4FA9F66786BE285C6
9B0848A4C8F4CD5BD77E29DDDFC451E5C2FA0AFCFED27E2644F129137EB5A1
75F11E3D87F7BFAD7F32CB85EA27CB6D0FBC58BA77E7B2B9EF7E25F1729572
5FF5AF92BE22F8AD1A29515F24E475AC9F11FC465656549324FBD7CFBE24F1
3B4FBE699FD6BED787B869C249B479B8DC7DD591FA9FFF0004BE8E77B8F8A5
A832FEE656D1A356C8E5E3178CC31D780EBEDCF1DEB5BF6FCFDB8A4F84505C
7C1CF84F76BFF099DE45FF00130BF8D813A54320C844C74B9753904FFAB521
BEF3295C9D07C5537EC1DFB141F12EB702C5F103C6B7125C5A594E177C7797
5185844887E6296D046B24AA47121319C17CD7E3DFC13F875AEFED0BF17E2B
2D6EE27BC4BB99EFF58BC762D2B216DD21673CEF958ED0793925B04035FD1F
42AFD530518CF4B2BBF2EA7E7F5DF3D56D1F57FEC67FB2ADCFC5BD607C5BF8
A71C97BA4BCED35B417259DB509F7EE79E62C72D1EECF0D9F31B24FCA3E6FD
F4D0746B7D32D238A240814000018000AF3BF86FE17D3B41D26CB4AD36DD2D
AD2CE248A28D06D548E30155540E8001815ECC838C0AFC4F37CCEA636BDE5B
74476D3828A1EAB56163269D0C5B8D6A246AA2BEBF87B861D75CD2D11CF56B
58CF16EDE94D6848AD6A42A0F5AFB3A9C234396D17A9CEB10CC364C540C3B5
6A4F1E2B39C62BF31CF72BF61271676D29DCC9D42DC4B11E39AF1FF11E9CAC
AFC735ED730CA3579AEBEA811D9C85500924F402BE4E8BB4AC8DD9F94DFB6A
7ECDB6BF18BC1D26A7A1DBA2F8B3444692C64C0569D07CCF6CCDE8FF00C19E
15F1C805B3FCFCB2DD59DCBDA5C46F14F0B9478DC157575382A54F2083C115
FBD7FB47FEDCBE05F095E5D785BE19471F8AB578F7249761FF00E25F038EC1
D799C8EE1085FF006F3915F8EBE22D56FBC59E24D47C59ACAC4DA8EAB334F3
BC51244ACEDD4854000F73D49E4924927FA0B81ABE2E8E1DC2BC7DDDE37DFE
EEDFD753C7C5C23295D1E7B691EA4E0108507FB5C7E9D6B612D6EC8F9E5C7D
2B6D6202A411D7D8CF1B37B6873AA48E7DB4DDFF007DD8FE35564D0A07EACF
F98FF0AEAFCBA698EB2FAC4FB95C8BB1C3CBE1C5EB1CAC3EA33FE15933E877
B1E4A6241EC707F5AF4968EABBC42AE38B9AEA274D1E4B2C52C2DB664287DC
5455EA1716B1CAA56450C0F622B96BED05797B53B4FF0074F4AE886313F88C
DD23952734D27152CB1490B98E552AC3B1A809CD6B29084AFD89FF00823178
874ED13E3E78AACB517F2BFB63455B481CE021B8FB424AB1927A17546DBEA4
63A915F8ED5FA51FB087876E61F0F78B7C676EED0CFF006BB4B7B79509578E
4B65698B2B0E873221C8E46057C9F16E7AB2EC14B14D5ECE2ADDEF249FE173
A70D4B9E5CA7F5994578D7C07F8A117C57F87B67AF4CC3FB56D0FD93514002
E2EA251B9828E02C80875C740D8EA0D7B2D77613150AF4A35A9BBC64AE8251
69D98514515D020A28A2800A28A2800A28A2803F353F68ED57EDBF1BAE2D4B
03FD9B676B6F8041C6E533608EDFEB7383F5EF53786A6C2A579CFC66BE966F
8F3E2B79DB732DC4483A0F95208D5471E80015D5F872E70ABCD7F985E33B75
B39C5CDFF3CBF06D2FC0FB5CBB4A715E47D07A3DC70BCD7A05B481E315E41A
45D70BCD7A369D75C019AFE5FCD70CD49B3D9848E8A9E0D46A430C8A7038AF
1A94ECCD19691AAE472566838A955EBEC329CDDD37B9CF3A7736126C54FF00
68F7AC61253BCD3EB5FA1E138CE518DAE72CB0E6A34FEF55DE6AA465A88C95
C398F184E6AD72A187B133C99AAACD485AA326BF3CCC7359547B9D5085809A
651457CE549DD9BA41451514B208D6B34864573284522B91D42E7AD68DF5D8
E79AE2F50BAEBCD7B597E16ECCA7232AE7CCBCB98AD22E5E775451EEC703F9
D7D6B65611DBC11C118F96250A3E806057CBDE06B6FED4F1958A1E52DCB4ED
EDE58CAFFE3D8AFADA1C015FD99F47CE18A5569D5C4D65BB497C95FF005FC0
F9ECD6BB4D24406D801D2B3AE2102B79D862B22E58735FBBF1864D85A741D9
23CCC3D49367397098AC89456D5CB0E6B1A535FC35C774A9C6A4B94FA5C33D
0A87AD31A9E7AD31ABF15AEBDE3D03F0ABE35F87EEFE0EFC71D7FC393298EC
2FE63A8D83636AB5ADD31750BED1B6E8BEA95D4F873C62AC8BFBCFD6BEE5FD
B23F67CBAF8D5E028F58F09C6BFF000987860493D88E86EE161996D09E997C
068C9E038C6543B30FC4FF000FF8EEEF4EB9934CD511ED6EED9DA2962954A4
91BA1DACACAD821811820F20F15FE99783DC5B4789323A72E6FDFD24A15175
BA5652F49257BF7BAE87C8E3A8BA351F67B1FA41A7F8B46D1F3D7429E2A523
EFD7C4BA5F8F91947EF7F5AEAE2F1DA6DFF59FAD7DC56C8DA7B18AAC7D513F
8AC63EFD729A978B1769F9EBE7D9FC72BB4FEF3F5AE5AFFC6524B911B6734E
8E4AEFB09D63D57C41E2D015BE7AF10D735A6B82D23B71DAB1AFB5B2E4B48F
935CE4066D7AF7ECD192204C195C761E83DCFF00F5EBDDA787A787A6EA5476
48F4B21C8F179AE36965F8287354A8EC97EAFB25BB7D16A745A446C6DA4BC7
EB70D91FEEAF03FA9AE6FC743FE291D73FEBC6E7FF004535778EA91C6B1A0D
AAA00007602B8DF15DBFDB3C3FAA598FF96F6B3C7C1C7DE422BF3FA18D7571
8ABCBAC93F95FF00C8FF0055A1C254F28E19593E1F5F6749C6FDE4D36DFCE4
DBF99F9EE9D2AC29AAA86A753DABFA94FF001F0D5B27DB2035EDFE0AD696D2
58CEEC608AF048DCA9CD74BA6EA6D6EC30718A903F4FBE1D7C4E5D3D23FDEE
318EF5F4B5BFC75F2ED7689FB7AD7E3C693E329AD80C3918F7AEB3FE1635C8
8F6F9A7F3A00FBDFC75F194DE4720F3B39CF7AF88FC75E3337F24877E739EF
5E79AAF8DEE2E01CC84FE35E71A8EB325C124B67345808F56BEF3A4639CE6B
9876C9A92594B924D55273400139AFDA0FD823F66B8BC3BA741F14BC576B9D
73568F368922F3696AFD0E0F49251C93D4290BC65857E6D7ECD7F0C87C52F8
A9A6E95791F99A5E9C45EDEE47CAD1C4C36C67FEBA390A4673B7711D2BFA53
F0069D1DBC1100300015F81F8DBC613C3D05975076725797A745F3DDF95BB9
ECE5585E67ED19EF1E16D3238225F97A0AEF1476AC4D1902C0315D046B935F
C694E0EAD6B1F4355D896388B56847685BB54F6700622BA48ADD11464735FD
1FC03E17AC64154A9B1E262B1BCAEC8E64D9103A5549602B5DAB448C318C56
25DC2066BE8F8BFC2EA586A0EA53D8C70F8D727639865C543D2AECCB8354DB
AD7F3166B84F653713DAA6EE432A86520F7AF2AF14E9C922B8C57AC3570DE2
14055ABC7C24DC6A6874C763E72F0F78CEF7E1378C175C8833E9378562D421
193BA20789140EAF1E495F504AF1BB23F42ECAF2D351B38350B09567B6BA8D
658A4439578DC065653DC10720D7E74F8E6CD668A4047506BD4FF64CF88525
E5B6A7F0B7536CCFA2837764C5B2CD6923E244C1E7F75230C1E987038C73FD
83E0AF1649AFECEACF47AC7C9F55F3DFD7D4F9FCD70DF6D1F66514515FD167
861451450014514500145145007FFFD2FDFCA28A2800AF20F8DFF04BC13F1F
3C0975E04F1B40DE539F36D6EA2C0B8B2B950424F0B1E8CB92083C3292AC08
26BD7E8A99C14959EC34EDAA3F991F8D5FB38FED0DFB38EA17035DD326D7BC
331B3791ADE9F13CB6CD18E86755DCD6ED8EA24F9739DACC066BE798FE22CD
70BF35C607B1AFEBD2BE4EFDA53E09FC32D6BE1AEB7E265F0768EFAD69423B
F5BC3A75BB5C2A432A3CFF00BCD9BF0D12B0639C7427A57C6E6DC374634E75
E11D9376EF6D6C7AB86C7C9C9425D4FE6F63D7AFF53590E9B6F35E98865FC9
469368F53B41C0FAD7EBF7EC8FFB01EB1A3EB7A77C56F8FA912DD58BADC69D
A123ACC23950868E7BB9109462A46E5890B2E76976EA94FF001468F617DE02
369636F1C1079242C71A844191D946057DC9F0F3E284369FB2E58FC51BF613
FF00C23FE1D9EE2E8962DBE5D26174B8C9EB92F0B64751D327AD7C7F8719D6
1F32A955BA7CAE16B2BDFE7B23D4CF7093C3C63695EE7E11FF00C149BE354B
F14FF689BFF0C58CBBB45F002B6910203F29BB56CDEC98E70DE68111F6894E
01CD7D4FFB00FC358BC39F0FD7C57771017FE2493CF2C46196DD0958573DC1
1B9C7FBF5F8DB17F6AF8D3C568B773B5CEA7AEDE8124ADCB3CF75272C7D496
6C9AFE913E1169967A1685A769562823B7B386386351FC291A8551F8015ED7
1A631AA5C8BABFC8F070D1D6E7D7DE1B8D52218F4AEC631C8AE23C3F38DAA2
BB788F22BF29C1DBDB6A77CB6362D9703356AAADBB023156ABFA2B24E5FAB4
794F22AFC41451484851935EACA492BB332ADC74ACA90F357679326B0F52D4
2CB4BB1B9D4F529D2D6D2D2379A69A460A91C7182CEECC780AA01249E82BF1
6E2FC7C6A557CA7A5878D9181E32F167877C0DE19D47C5BE2DD422D2F48D2E
269EE6E676DA91C6BDCF7249C050012C48001240AFE707F6B9FDBA3C51F1EE
F2EBC1DE0233E81E0304A3213B2EB521C7373B49DB1E41DB129C11CBEE380B
47F6E3FDB1353FDA3FC64FE17F08DD4B6FF0EF43948B4841283509D0906F25
52AADC838891BEE2F3C33301F0A4698AFA6E11E0F8D04B158A579BD52FE5FF
0083F97A9857C45F48EC11C4055A54A74699E95AD6B6124A400335FA21C866
AC44F41530B76EC2BBAB0F0CCF3E0843CD74B1782AE0AE7CB3F9500790181B
D2A268C8EB5EB375E129A2049435C95EE8EF0E72B8A00E3196A165AD59E028
4E6A83AD005075AA72460D69B0AAAEB401CDDFD8457285641CF63DC570D776
92DA49B24E41E87B1AF4F956B16F6D5278CA38E0D690A8D1328DCE02BF76BF
668F87EFE03FD9F3C39677717957BABC4DAA5C02003BAF3E78F23A8221F2D4
83CE4738E83F273F67CF83F79F17FE30E8BE083196D3C3FDAF5171D12C6DC8
69724720BF11A9ECCE2BFA02F112C71C4638D42220C28030001D0015FCF7E3
A712C54B0F95D37ADF9E5E9AA8FDFABF923D6CAE86F37E869FECA5F103FE10
8F8C9FF08B5ECBE5E9BE2E8FECC431C2ADE4597B76FAB65A3C772E3D2BF5A2
BF9F0F135E5EE93A95B6B1A64A60BCB09A39E191786496260C8C3DC100D7EF
5F837C4B69E32F0968DE2DB1E20D62CE0BB45CE7689903ED3EEB9C1F422BEB
FC2CCD5D5C13C3C9EB1DBD1FFC1FCCC71D4ED2B9D251542FB55D2F4C50DA95
E43680F796458C7A7F1115CEC9F113E1FC2ED14BE27D2D1D4E0AB5EC0083EE
0BD7EA2709D8D15CB5AF8E3C157D8FB178834FB8CB6C1E5DDC4F963DB863CF
238AE9A39639A3596170E8C32194E411EC45003E8ACED4F58D2B45B6377ABD
E45650F4DF338404FA0C9193EC39AF27D5BE3C78174F252C4DC6A4C38CC11E
D5CFFBD214E3E80D007B4D15F2DDDFED1974C76E9FA0003FBD2DC67FF1D541
FCEB14FED0FE2CFF00A0559FE727FF0015401F18FC6CBAF2FE3DF8B97D2ED7
FF0044C75D27872F7E54E7D2BC83E2FF00896EB5EF8C7AD6BD790476B25FFD
964F2E3CEC1B6DE38C919FEF1424F5E4D757E1BD4328BCD7F9C7E2C65AD667
8ABAFB73FF00D299F5F819FB91F43E98D22F46179AF45D3AF781CD78468F7F
C2F35E8FA75F8C0E6BF9AB37CBF567B14E47B259DD820026B5810C322BCEAC
AFF81CD7536B7C30326BE1715837177474C646F038A70355926471D6A6AE48
D4712AC4BB8D2EEA87269726BA638C685CA49B8D216F5A8E8A89E29B0511C5
A9B4515CD29B652414531A455EA6A84F781460511837B0365B966541583777
9D79AAD757DD79AE6AF2FBAF35EAE1702DB339487DEDE7079AE2F50BCEBCD4
B7B7DC1E6B8CD42FBAF35F6197603C8E79C8F70F8350992FF52D518708890A
9F52C7730FC302BE8B8EE4015E15F0960369E128EE1BADECD24DF803B07FE8
39AF545B9E2BF7FE07E338E5D41508BDAE79789C3F3BB9D13DD0C75ACC9EE3
3541AE6A9C93E7BD7A3C4DE267B5835CC451C1598E9A4C9ACC76C9A91E4CD5
6639AFE6BE22CE9D693773D8A34EC36987AD389C532BE1A72BB3A02BE00FDA
BBF630F0C7C699E6F1B783A58BC3FE34DA3CC94822D6FF0068C28B80B92AE3
8025504E38656E36FDF334A235AE3755BE4404BB05FA9AFABE0AE27CC727C7
471D965470A8B4D366BAA6B669F67EBBA4C2A61A1520E351687F339E33F067
C4AF843AC1D0BC7DA45CE9530384771BA1947AC52AE63907AED638E870722B
F7EBF63EF0A693A1FECEB6BA8CB630F9FA9D924D70E6352D2B4ABBCEF38F9B
1918CFA0F4AD3F13D8D9EBB613585FE9EBA9DA4A3E6864844D1BE3900AB02A
7DB35EB3E117F0C681E05B3F0DC72C56D1416E88F112621B954020038EFC00
3F0AFEC4C478FAF31C0C638DC2B8545BB83D25A3D527AAD7A5DFA9F3DFD96A
9CFDD95D1F837FB485EE8DA47C5AD56CF47B78ACA1F959A3854226F39C90A3
819EF81EF5E0B2F892341F7EBED0F89BF0B3C19E26F14EB5E27D5A29A6B89E
EE655FDF3A011C4E5146176FF77FFAF5E223C2FE1FD12F07F6658C7095070C
46F718E3EFB65BF5AFDD32BF107032A1185384A528A4B5B257F5BB7F81EBF0
5F06CB39CDF0F963AAA9FB595B9AD7B68DED757FBD1E69A768FADEB6564995
ACED4FF1B8C3B0FF00654F3F89E3EB5E9963636DA6DB2DADAAED45FC493DC9
3DC9ABA4E6A376DA2BC3CDB3DAF8C769E91E896DFF00059FE9DF865E0E651C
2D4DCB09173AD25695497C4FC9748C7C96FA5DBB22B4EDD6B0EEB6B2956190
7822B4A77AC5B97EB5CD8789F659AD74D33F3D751B43A76A577A7B75B69A48
8E7FD862BFD2A0535DCFC52D3FFB3BC6D7A47DCBB0B3AFFC0C61BFF1E06B81
535FD4F97627DB61E9D5EE93FC0FF1AB8A729780CCB13827FF002EE728FC93
76FC0B20D4E9215AA80D6B68FA2EB5E20BD5D3741D3EE352BB704886DA279A
420752110135D52928ABC9E878290A976EBD0D4FFDA127AD5FD77C15E34F0B
451CDE26D0350D22394E11AF2D65B70C47606455CFE15CCE48A9A556135CD0
775E4369ADCD07BB76EA6AB3485AABEE3464D588793EB4C269B914D26803F5
8BF608F07C7A7784EF3C573C604FAD5D1547EE60B6F907FE44327E95FB17E1
021638C7D2BF377F663B18F44F871E17B145D8458C3230F47997CC7FFC798D
7E857856F0044E6BF883C4CC54B138FAD55F56FEE5A2FC0FACC04796091F4C
68D206880AE9A0EB5E73A1DF0C2F35DFC3202030EF5F8FE0A6A9574D9DB595
D1D558B018AE81482322B8CB7B8DB5AD1DF6075AFEC1F0FB8EF0B470CA9546
7CE62B0ADBBA3759828C9AC2BC941CD325BD2475ACA9AE3754F1E788987A94
1D2A4F40C2E11A7765699B26A99EB523BE4D404E2BF91B38C72A951B3DEA71
B219236149F4AE035E9C10D5D5DFDDAC684035E61AE5E8C3735E760A9B94EE
75256479178B5C1471F5AF00F0978C3FE15EFC63F0CF89DA431DA8BC5B5BA2
3A7D9EEBF72E4FA850DBFEAB5ECBE29BC055F9AF8D7E27CDBEDE500E0F623A
E6BF69E08AB3A188A7563BA699E762D5E2D1FBEB4572BE05D71BC4FE09F0F7
895D83B6ADA75A5D961C02678564C8E075DDE82BAAAFEE084949292EA7C930
A28A2A8028A28A0028A28A00FFD3FDFCA28A2800A28A2800ACDD674BB6D734
8BED16F466DF50825B7907AA4AA51BF435A5452693566099F943E1E8AE25F0
BCBA4DF8C5D5917B7987A491128C3F306B1AF7C5577A77EC2FF1BFC311B913
684F3451A9C616D3546889C63079779FAFEBD07AB78B34CFF847FE2E78CF43
DB88AE6EBEDD1FFB42F504EE7FEFE338FC2BE54F8953DC697E02F8D9E16B75
1B7C43E185BBC74DD2695771C98C9E06219666C77C6073807F99B80DBCBF88
ABE09ECDC97F97E47DFE72BDBE02157D19F94DF042D05FFC5AF0C40C3705BB
12E063FE58A9933CFA6DCD7F421E02BC0B6F1007B0AFE7BFE055DA59FC59F0
F5C49C0124ABF8B43228FD4D7EEAF817560608B9EC2BF44E30A6DCE3E87C76
1CFB4342BF01579AF4FB3B959901079AF9C342D586D5E6BD474CD5F6E0835F
9955A6E32E6476A67AAC52E2B416E4639AE3ADB558A40371C1AD35BA8C8C86
15F5994F1755C3C794C2A504CDF370A3A55496E09EF598D729DD855496FE14
1D735BE63C6D5AAC7950A18648BEF2679AFC58FF0082A3FED4B269B66BFB35
7826E4A5CDF24773E21990E0A40D8920B307AFEF462497A7C9B1790EC07E9D
7C62F8B7A47C26F86FE22F88DACFCD69A0D9C971E58214CB20F9628949E374
921545F7615FC8DF8B3C5BAF78FF00C59ABF8DFC53726EF56D72EA5BBB9909
3CC92B162141270ABF7557A2A800702BAB81F2A78BC4BC5D65EEC36F397FC0
DFD6C4E267CAB9518F12E055F8D6AAC62B52D63DEE057ECC79C6AE9D62D3B8
0066BDBBC25E0B92F5D3E4CE7DAB99F0868E2E66418EA457DEFF0009BC0A97
2D0931E738ED401CC784BE104974884C39CFB57ADA7C137587260FD2BEF0F8
79F0BED9A08D8C43A0ED5ECD7BF0DED21B527CA1D3D2803F18BC4FF0B4DA46
DFBAC63DABE65F15F853EC8CFF00262BF63FE25F842DADE39708075AFCEBF8
97A4C50BCB81EB401F076AF61E53B0C571D32609AF5AF1340A92BE2BCB6EC6
18D0064BD5671569EAB35005390567CAB5A2F54641401EC9FB39FC66B8F81F
F12EDBC45312DA2DFAFD8F538C2EE26D9D81DEA3FBF1B00E31C900AFF157ED
8EB57F6F796A9796922CD04E8248DD0E55918655811D411C8AFE75AE1720D7
E9FF00EC8DF14A6F147C389FC11AACE64BEF0BB2C709624B359499F2864F5F
2C86418E8A1057E05E32F0746A2866F457BD1B467E6BA3F93D3D1AEC7AB976
22D7A6CF61F16B072C2BE90F827E28F14EABF0E6DF428F59BB4B1D2669AD96
0499D10066F370402323F79C6738E9DABE62F114BE64AC2BE9AFD9A74A6FF8
44352BE71C4F7EEABF448D39FCC91F851E15CE51C4F2AEB177FC031DF09E8E
3C3D1B3191D773B1C9279249EA49A53A0443F83F4AF491623D298D643D2BF7
D3CA3CAA7F0FC447DCFD2A94169A9E8D2FDA345BC9F4F941CEFB795E16CFAE
50835EAB2D90F4AC8B8B11CF1401E67A96BBE2DB8BC6BFD4AFA4D4E53D4DD1
321C7A06CEE03D06703D2AE5878CF4D2C21D5A26B26FEFFDF8C9FA8191F88C
7BD6EDE69E0E78AE3352D291C3656803D5205B7B98567B7759637195642194
8F623834E6B61E95F3F5BDE6B5E189CDC68D36C4272F0B0DD13FD57D7DC60F
BD7B07857C6DA5789C0B53FE89A8A83BADDCF271D4A1FE21FA8EE3BD4B03E6
7F8F36A74CF18E97A881B52F6D0A67B1685C93F8E1D7F4A87C31AA82A9F357
A47ED29A21B8F0441AF4631268F751BB1EFE54FF00BA603FE0650FE15F33F8
5B58C055DD5FC7FE33E41CB98D4A896924A5F859FE299F419755F7123EBDD1
B52E179AF49D3751E0735F38689AAE42FCD5E9DA66A7D3E6AFE52CE729D5E8
7B94EA1EF165A8F4E6BAAB4D43A735E2F63A98E39AEAED352E9CD7E7D8DCAF
C8EB8CCF5CB7D43DEB5E2D4381CD7965BEA5D3E6AD78752F7AF9BC4659E46C
A67A525EA9EF530BA435C047A8FBD5A5D47DEBCE9E5ACBE73B7FB42534DCA0
AE37FB447A9A69D47DEA165CC39CEBDAF5455392FC7AD728FA8FBD519751FF
006AB7A796F909CCE9E6D43DEB1E7BFE0F38AE7A6D47DEB1AE351EBCD7A987
CB7C8CDCCDCB9D43AF35CCDE6A1D79ACBBAD47AF35CCDE6A5D79AFA3C1E59E
4652997AFB50E0F35C66A1A867201C9355EFB52EBCD57F094635CF19E91A61
F9965B942E3AE510EF7FFC741AFACC360953A72AB25A4537F71CEE57763EED
D02D3FB2B44B0D3BA1B68238CFD55403FAD6C894D5407B8A5C9AFCA69E7552
EDDCEFF668B265351B3E6A2C9A4ACABE6D397505042939A6938A42DE94C2C0
7535E4D4AAE4CB485EB552F2F6DAC607B9BA90471C632589E0547757F05B44
F34AE1238D4B33138000E4927D0579AE9F7571E35D416E640534F89B30C678
DDFEDB0F53D8761EF9AF5329C9E5896E52D20B77FA2F362A93505AEE74B15C
EA7E21908B1536B6B9FF0058C3E761EA01FBA3EBCFD0D74367E18D36D8F992
279D29EAEFF31FCCD6E5ADB476D12C518C015792067E82BEBB078472FDDE1A
365E5BBF56734E6DEB3667AD9DAA8C2C4A07D2A09F4AB0B852B242A73ED5BA
6D5C0CE2AB34656BAF139562286B34D12A516789F8DBE11F86BC4B6B225E59
472820F38C30FA30E47E06BF387E2F7ECE5ABF86E47D57C32CF796D1E4B40D
CCAA3FD923EF7D3AFD6BF61480460D79AF8C74386E607CA8208AFA7E12E34C
5E06B28F35E3D99E9E459B54CB31F4331A09395292924F66D7476E87E0B383
192AE3695E083D411546592BEBAF8FDF0B12D669BC47A345B255C9991470E3
FBD8FEF0FD6BE34965AFEACC9331A78CA2AB53EBF81FE98F02789D83E22CBD
62F0DEEC969383DE32FD53E8FAAECD34A39E4EB58B7125599A5AEA741F87DA
C788712B9FB2C0DC82C32C7E83FC6BEC72FC054AD2B53573C1E3DE3FCAF24C
3FD6734AEA9C5ED7DDBEC92BB7F24CF8FBE3769BBE1B0D690731B181CFB37C
CBF9107F3AF9FD5ABF5C755FD9DB40D7B4A974DD566B89239B1B8AB2A9CA90
463E5E3915F3AF8D3F63A4B2B77B8F086AD279AB9222BC0195BDB7C6AA57F1
56AFDB786717EC30D1A15DEAB6F43FCC5F14B8EF27CDB3AA98ECBA4F96695E
F1B7BCB4FB9A49FADCE0FF00645FD9AF54FDA73E28C3E125B86B0D0AC02DC6
A97680178E12D811C79C8F324C10A4E4000B1071B4FF0055BF087E07FC3AF8
41E1B8BC2BF0C3C3F6DA2D8C6AA2430A7EFA765180F3CC7324AFFED3B13DBA
57E417FC131B4A93E1AD87886CBC4500B5D5AE751224071930A46823C11D57
25883EE7BE6BF73748D6E25877C243AB0CD7F3E7889C41471BC45F50CD6ACA
1848745B376BDDF7D76ECB6EA79D83A76A0AA52576CC4F11786B49D734EB9D
0FC456306A363748639ADEE235962911BAAB2B02083E8457F39BFF00050CFD
8CF42F837A88F89DF0BADFEC9E1CBB902DEE9EB931DA4921C2C9164F11B31C
14E8A48DBF29C2FF00485A96A5192D3CEC140AFCD5FDB735FD23C41F0E35FD
12E590C73DA4C9F3630BF21C37FC04F39ED8AF88E06E21797F102595CDBA0E
5669ECE3E6BBF667657A1CF45FB4DCFE6633EF499A8B751BABFBACF9424DD4
D2DEB4C2D4C2DE9401FB9DF05AF52DFC35A242920754B3B750C3A1C46BC8FA
D7DB7E17D506C4F9BD2BF31FE05F8992EFC1BA05C0619FB2428D8E819102B0
EFD0835F72F8575C0553E6F4AFE32E31CB24AB4D3EEFF33EA30D53447D8FA2
EA830BCD7A8E99AAAB285635F2EE89AD0C2FCD5E9BA6EB5C0F9ABF1CC7E5ED
3B9E9C6575A9EFB15C2B0CA9AB22722BCB6CF5D200F9ABA08B5C040C9AE3A5
8DC452D1325D14CEC4CC7D6A16726B9DFED94F5A824D6460F358D7C757A9B8
2A363A379157A9AC9BAD41114E0D73973ACF079AE62F7581CFCD5CF4F09293
BB3549235354D53AF35E5DAE6A630DCD49A9EB2307E6AF2DD735A186F9ABE9
B2ECBDDD194E6731E28D4C6D7F9ABE47F8817C26DC99CE4D7B3F8A75B015FE
6AF987C51A879CD2484F4CE2BF5BE1CC072B4CF3ABCCFDEAFD9D6E2E2EBE04
780A5B93971A359A0C803E448C2A74C745039EFD6BD9EB81F855A3DC787BE1
7F83F40BB3BA7D3346D3EDA438C65E1B7446E327B8F5AEFABFB0709171A508
BE897E47CCCB76145145741214514500145145007FFFD4FDFCA28A2800A28A
2800A28AF2CF11FC56D174A91ACF494FED4B95C8251B6C2A7D0BF39FF8083F
515CF88C553A31E6A92B22A316F63E72FDA3F4DFECAF8A3E1CF112AE1358B0
92D1C8E9BECE4DE33EE44F81ECBED5F2778FB428352F15E97A75C48B05B788
85CE853C8EC56358B5BB7934F2EE4764F3C3E7B1507B66BDEBF68CF1E6B7A9
78674FF11EA36F12DBE877F14A560521C24F98082CC4F1B9D49E9D057C73F1
5FE216A3AB7878DC68D6CB6971001245296DEC9227CCAC0600C820119CD7F3
57145454B8929E3A87C2DC5FDDA3FC8FB8CBEBC65807466F5D4FC89F0ADC5C
F85BC7BA71BF536F3E9F7C91CEADC142AFB2407E9CE6BF683E1F78895EDA23
BFB0AFCCDFDAF7C3AF69F16A7F8876562963A47C43862D7AD162CF96B25DA2
C97080F4C8958B60138575AFA37E0B78F86ABA158DD3BFCEE8038CF475E1BF
506BF61E20C3AAD4E3561AA3E4293B3B33F50340D746D5F9ABD4B4ED6F81F3
57C75E1CF128655F9ABD6F4CF1082061ABF36C4E0F53AE323E9AB5D70607CD
5B51EB7C7DFAF9FAD3C40081F356D45AF0C7DFAF3278534E63DB1B5BE3EF55
09B5A1FDEAF293AF0C7DFAA536BC31F7AA1610398FCF4FF82A4FC589EDFC25
E16F851A7CE53FB6AE1F51BD556C660B4C2C28E3BABC8E587FB5183DABF17A
1AFAA7F6E3F193F8C3F68FD7544865B7D0E0B5D3A13BB2008E312C807A6259
5C63D727BD7CAD11AFDEF85F02B0F81A71EAD5DFCF5FCB43CAAF2BC99A11D6
FE9A81A415CF466BA0D35F6B8AFA0323E8EF87D0219E3DDED5FA5DF0712DE3
3096C76AFCBDF04EA2B04A849C74AFB9BE1C78C92D047F3E318A00FD85F01D
FD9C36D1E481802BB7D77C416696AC030E95F047877E2AC76F028F37A0F5A9
B5EF8BEAD0301376F5A00D7F8ADE21B6649B0C3BD7E69FC4DD5639249707D6
BD8FE207C4817224C4B9CE7BD7C57E34F13FDA5E4F9F39CD00792F89AE43CA
E735E5B76F9635D3EB17DE6BB735C6CCF926802B39AACE6A56355D8D0040E6
A9C9565CD5390D00519B9CD7BBFECA5E20B7D13E3C786F4FBF94C361E2297F
B26E08217FE3F30909C9E0059FCB273D811EF5E0D29AA515DDCD85DC17F652
186E2D9D658DD7AABA1DCAC3DC100D7266183862284E85457524D150934D34
7EED7C49F861E23F08CB35CC886EAC813FBD41CA8FF697B7D7A7D2BEC3F81F
E147D1BE1A68A8C9896F23376E71D7ED04BAE7DC2151F857A0D90B1F88DF0F
3C3DE358A2530F88F4BB3D4540185D97902CC300E4E30DEB59FE11F15C3E17
D4ADFC37E2AE34B9196282E8F1F6724E156524FF00AB1FDEFE1EFF002F2BF8
FF0002D4C3617152E7D1B565D93BFF00C03D0C55DC743AEFECB7C7DDAAD2E9
ECA3915F441F05A32064DA7238C573BA9784E48413B2BF6B3CD3C127B4C678
AC7B883AF15EA1A9690D0E72B8AE26F2DB6678A00E1AEA0183C572D7B6C0E7
8AEF2EA3C66B9ABC8C73401E65A959020F15E69AB58490C82E6D99A19A33B9
1D0956523B823906BDA6FE21CD705AA5B821B3401E5DF167E35EB117C2CD43
40B9B1FB46A97845B3DCED5F285BBE7748CBDA4E028C0C027764118AF973C2
DAF0648E40DC1C77AFA5FC4BA545730C914A8195C104119041ED5F1FEB3A4C
DE0ED64C299FB0DC313113FC27BA1FA76F6FA1AFCCBC4DE1F58AC22AF15AC3
7F47FE4CEDC155E5958FAA741D6B2ABF357AC697ABE42FCD5F207873C41C2A
96AF66D235CC85F9ABF8E33FC81A93D0FA0A554FA5AC355E9F3575D67AAF03
E6AF00D3F59C81F3576367ABE71F357E659864AF5D0ED8D43DC2DF54FF006A
B661D4FA735E376DAB74F9AB6E1D5BFDAAF96C464FE46EAA1EB51EA7EF56D7
53F7AF2D8F55FF006AAE2EABFED57975329F22D4CF4AFED3FF0068D34EA7EF
5E783541EB4D6D547F7AB2594F90FDA1DEBEA5EF54A5D4FF00DAAE1DF55FF6
AA94BAAFFB55D34F29F21399D8CDA9E3BD63DC6A7D79AE4A6D57FDAAC5B8D5
7AF35EAE1F29F233750E9EEB54EBCD73179AA75E6B9EBAD5BAFCD5CBDEEADD
7E6AFA4C1651E4632A86DDF6A839E6BBFF00817FE9BE3993506194D3EDDDC1
C670F21083E84A96AF9D6FF56EBF357D03FB3F4822D3355D5DB83733AC209F
EEC4BBB8F6CBFE95DBC4781F6396D5696AD5BEFDFF000B8F0D2BD447DBF0EA
60F7ABC9A829EF5E5306ADC0F9AB45355FF6ABF9F6AE54FB1EEF3267A4FDB5
3FC9A61BD415C00D57DE9ADAAFFB558FF6630F74EE64D41477ACD9F52C6706
B8D9355FF6AB2E7D57FDAAEAA3953EC1CE918FE3FF0010BDD5C5A785EDDB9B
B3E6CF8ED129E17FE04C3F20477AF56F055A4705B20518C0AF92B4FD5FFB5F
C77AA5E31C88A6FB3AFB087E43FF008F027F1AFACFC2774BE5A2E7B57E9198
E57F54C153A115ADAEFD5EBF86DF23C6955E7A8D9E95047BDABADD3EC15802
4572D6AC03035D8595E2A28AFD9FC0EC165AEB29E36D63833394EDEE962E6C
1113A572579105635D55DEA0ACBD6B92BA983B1AFADF1B71194F2F2E0D230C
B554FB4661EB5CEEBE14C041F4AE85881C9AE335DB96903244A5CFA019AFE4
7A2B9AAE87B751E87CB9F12AC22B9B7991941C835F93FE3FD24683E22B9B54
1885C9741E809E47E06BF647C49E11F10EB4ACB6D004DDDE4751FE35F06FC7
FF00817AC68F15A6B1AAEA96504D73232430217791C2E0B9FBA000B919E7B8
AFE96F0B334A6EBC706A7794F4496AEFBFFC3F65A9F77E15F1DC387F35589C
4CF9684938CDF44B74ECBB4ADE766D2DCF953C1BA4A6ABA909A71BA3888C03
D09AFBF3E0E7C31B9F1BEA50E9D68BB53B9F415F0EF81A55B3792DD98168E4
752477DAD8AFD12FD9DFE29E9DE0AD4966BCC146183F4AFECCC0E5F1C3D28D
3FBFD7A9F8078C1C755B88B3DAD8DAB36E95DC60BB413D2CBBB5EF3F36CFA7
35DFD972C74ED0A4B9B5B8F3278932463D057C15E30D186997735A3F542457
E8BF8EBF696F0CFF00614D6FA49DD34CA4649E9915F9A3E30F128D52F66BA2
7EF926BBF150A4A7FB9D8FCD3194E8AA9FECEDD8F1DB8D7F53F036B7178934
462B2427F7880E1654EEA7FA1EC7F115F67F823F6A4D3EEB4E8A45BC032A0E
09C11EC457C13E30BF4686404F635F05F89FE236B9E0CF155C5BD9812DB481
645524A919241E467B83DABE5F893802867094A4BDF5D7BA3EE7857339D24E
9B7A1FBF5E27FDA5ADA4B6648AE77B11D01AE2FE17FC307FDA8B57D526F18B
CC9E0FB68E4B7B8F29DA392E25990A88A271CAED0773B0E47CAA3EF647E165
EFED09E2B9ADCC3650476EC4637331931F4185AFA4BF658FDBBBC67F06F558
B49F144ED7BA2CB2125B1CC7BCE4E40E08C9AF8DC5F84F8DC160E7532E4BDA
ADB5D6DD6DE7D8FB0598C6724A7B127ED7DFB00FC43FD9BAE6EBC5BE1BF37C
4FE002F95BE44CDC58AB1E12F117A01D04CBF231C6761216BF3E3757F65DF0
ABE347803E3778623BAD2AE60BC8AF62DB242C55D5D5C60A953C1041C106BF
2C3F6CAFF825EC5A81BDF897FB345B25BDC92D35DF86F2122933CB3D8331C2
377F258853FF002CCAE0237770278C8F9D65D9F7B95169CEF45E935D1FF797
BAFADBAE38BCB7EDD2FBBFC8FC1FDD499CD58D474ED4747BFB9D2756B596CA
F6CE46867827468A58A443B591D1806565230410083D6A957F444649ABA3C7
3EE5FD9BBC5FBBC38748924FDE69F332804FF03FCEBFA961F857E84F843C4A
1913E6AFC56F861E2A3E17F1344EEDB6DEF31149E8093F2B7E078FA135FA31
E10F1685D837F0718E6BF12E3CE1FF00DF4AA25A4B5FF33D6C1D6D2C7E8968
7E211B57E6AF50D375F040F9ABE2BF0FF8A95957E7AF5BD2FC4C081F3D7E11
99648D37A1EBD3AA7D5B69AF0C0F9AB7E1D78607CD5F3659F894607CD5D0C3
E2318FBD5F2B5B2777D8E8554F7F1AF0C7DEA8A4D7463EF57890F118C7DFA8
E4F118FEF5732C9FC8AF6A7ADDCEBA307E6AE62FB5E183F3579ADCF88C60FC
DFAD72D7FE24183F357A187C9DF622554EDF54D7C007E6AF29D77C42006F9A
B9ED5BC4C30DF3D7936B5E22694B2A357D765B936BB1CD52A8EF126B867668
D5BEB589F0DBC2AFF123E2AF85BC0CA8658F56D42149C28C9FB321F32E1B18
3F761573E9C73C571FA8EA030496AFD0AFF82787C2E9EF752D6FE346A91E2D
E157D2F4D0C3EF3B157B89477F94058D48E0EE71D457EAFC2B937B5C4429DB
4DDFA23CEC455B45B3F56A8A28AFE833C50A28A2800A28A2800A28A2803FFF
D5FDFCA28A2800A28AE53C71AFFF00C231E14D47595389618F6C5FF5D64211
3F26209F6A99CD462E4F64348F21F89BE3E9EEAF66F09E892148613B2EA553
F33B77894F603A377278E80E781D17439EFDD62B74CE7D2B8CD2C96219D8B3
31C924E4927A926BE9EF873696CB6866C02FC0AFCA332C554C557D5EFA2F23
D0A7151478E7C4EF847A8EBFF0DBC45A5C51EFB89EC2730A8E4F9C885E2E3F
DF0B5F952AC9AB78595C7CC1E3FE62BFA0B98C4B193C57E0FEB3A247E1AF17
78AFC1F18C43A56A5770C23FE980958C271EF19535F19C7F94FD59539295DA
3D1CB6B73368F817E2AEA1AEF8CB4BB3F0DEA73075F0BA3DBD84606D544DDB
88F72E3009EF85EC00AF3DF843E319BC3FA9C9A35CB14591B7A03C61C70CBF
8819FC0D7AFF00C46B23A7F8AAE0630B3A871F51C1FE95E11E25F0ECB34BFD
B1A40DB7519DCCABD588E770F7FE7F5EBFB2E49C98BCB6949758FE28F1B109
C2A347E8D784FC66B24519127A77AF70D27C54085F9EBF2D7C05F11A40890D
CB6C953E5753D88FE95F4FE81E3949117127EB5F318FCA9C5BD0D233B9F745
9F89F207CF9ADD8BC4A31F7ABE49B0F1882A3E7FD6BA783C5A081FBCAF0EA6
5FA9AA91F4BFFC248B8FBDFAD5597C463FBD5F3F8F16023EFD5697C5631FEB
2B2581F21F31F8F7F15F567D73E2BF8C757724FDAB58BF75DD8C84F3DF68E3
D1702B908CD49E239E79FC53AC4F76A12792F6E1A403A0732312075E87DCD5
58DABF6AA11E58452EC8F35EE692356A5ACBB58562A355A8DF15A88F52D0F5
530329DD8AF72F0EF8CDED82E24C63DEBE51B6BB3191CD74B69ACBC607CD40
1F70D97C4D9238C0F37F5AA7A97C4D964423CDFD6BE444F124A0637D413F88
A461F7E803D97C41E3692E377EF339F7AF1CD5F5B69D98EEEB5CF5D6ACF267
E6AC29AE99FBD004D73725C9E6B2DDA91E4CD40CD40033540CD4ACD5033500
46ED55246A95DAAA48D4015A5359927249ABB2B5507A00FEAD7F626BF4F137
EC75F0CB525612797A749699552A3365732DB1183DC18F04F42464715D378E
7428AE23954A839CD711FF0004E640DFB0FF00C3A3FF00618FFD3BDE57B9F8
AAD9591EBF9BF1F2E4C7D78AE9397E6CF621AC51CBFC03F8BF3689ACC1F0AF
C5D396B7B96D9A55CC87EE3638B5627B1C7EE89E87E4EE807DB524692A9471
906BF23FE21E9AEB9B9B676867818491C8876B23A9CAB291C8208C835FA03F
B3E7C541F15FE1F5BEA77AC06B5A63FD8B51418E678C022500745954871C60
12CA33B735FB0709E6EEB52F6337AADBD0F3EBD3B3BA3A4F1468A8819D0706
BC3358B6F2D9862BEA5D7E357B124F6AF9C3C4215647AFB039CF2DBC5C135C
C5D8EB5D3DFB0C9AE56EDBAD0072F7A060D713A8A641AED6F5BAD795EB7E30
F09D83B457FAD58DB38382B25CC48738CE30CC2AE9D294DDA2AE454AB18ABC
9D8E5F5884106BC03C7BA045AB584D6CE36B1E55BBAB0E8457A9EA3F11FC07
203B35CB56FA482B84D4F5FF000F6A487EC5A95B4E5867093231E467900E6B
4C465B55C1C6A53767BDD331A78EA32768CD37EA8F94F4AD627B0BA7B4BAF9
26818A38F71FD3D2BD9742F110655F9ABC77E27698FA7EA29ACDA0C6FF0092
403BE3A1FAF6FCAB0743F13142A0B6057F37718706FB2A8D25A3D8FA0C3E22
E8FB4B4CD77207CD5DCD8EB6303E6AF95347F12060BF3FEB5E89A7F880103E
6AFC4734E1C69BD0F4A158FA42D75ACE3E6ADE8359071F357805A6BA38F9AB
A2B7D73A7CD5F178BC83C8E88D53DCA2D633FC5575757E3EF578C45ADFFB55
7D35B1FDEAF1AA647E468AA9EBBFDAE077A43ABFFB55E5435AFF006A83AC8F
EF561FD88FB0FDA1E9CFABFF00B55465D5FAFCD5E6EFAD7FB55465D6BAFCD5
D14F24F213A87A14FAB8E7E6AC4B9D5C60FCD5C1CFAD75F9AB0AE75BEBF357
AF86C8F5D8CE554EDAEF58EBCD72B7BAC75F9AB94B8D59E4E14D64C9745B96
35F4985C9D4774652A86D5CDFBC84E4F15F5C7C2EB81A678374E84FCAF2869
5BDFCC62C0FF00DF3815F1199CBB845E598E07D4D7D57A4EAC96B6B0DAA37C
B0A2A0FA28C578DC6381E7C3C2925D6FF77FC39D18395A4D9F43C1AC8FEF56
9A6B3FED5786DBEBBD3E6AD48F5DFF006ABF26AD91791EA2AA7B30D647F7A9
1B5818FBD5E4835CC7F15236B9FED572FF00617915ED8F509759E3EF565CFA
C8FEF579C4BAE7FB55953EB9D7E6AECA1917912EB183E05D647F6C5F316C93
753127FE066BEC5F0B6B4A113E6F4AFCD8D0B5BFECAF176A960CDB42DCBB28
3FDD73BD7AE7F858735F5AF857C4EA6343BFD2BF41E2EC8DC9466968D2FC8F
1A954B33EE0D33558A78D43360D7451DD103835F35E93E260157E7FD6BB9B5
F1410A3E7AFCB211C56125FBA763BB9D3DCF5B7BA66EF54659D579635C0378
9F23EFD64DDF89060FCD5C98AA98BC4BFDE3294A2B63B4D47568E34214D7C8
DF1DFE36CBF0A60B4D5EE637FECEB9263699549549472AAC474DC3EEE7AE0D
7A7EABE24183F3578EF89F5EB7B886486E02C91B0E5580607EA0F15F4BC2F9
3D3A75E33AF4F9E3D55ED7F9EBF91CF5AA36B467CB377FB735A5C398EC04B3
BF27088CC71F80AF9C7E32FED17AFF008E92C6FADB4EBB26C4480E62600A49
B791F423F5AFA73C63E2086081ADED82C68A085550001F403815F1E7C47D48
0B1704FCD2B01FAE7FA57F527026130187C553C561709C925B372BEEACFA2E
8CE6C1653FDA389A581A92D2A4A29F95DAD7E5B9C0782BC58ED339B9568657
72C51F861B8E791F8D7D0FA478AB62A95931F8D7C47AACF736B32EA765CBC6
30EA3AB2FA8F715D2E89F10E378D4193047635FD35859AC5525523BF55D99F
2FE21F0154C9F309E1B570DE32EEBFCD6CFCFCAC7DB8FE2D79130D293F8D73
3A97895769F9EBE781E3C8CA7127EB58B7DE32926042371EB5B4704EE7C0C7
2F773B6F15F89048AE8ADC9AF92FC4FE18BAF17EBAA34C2CF78F88913190DD
4E38E4724E4D7A46A37BA85E5B4F756B0C93471101DD549552C70013EA7B0A
FA2BE10FC3D8745B65D57508C3EA37232C48CF96A79D83FA9EF5D15B1B0C25
2751EFD11F4B94E1356CF81FC49F087E23F8503CDAA685726D9064CF0219A1
03D4BC790BFF0002C579C6EF5AFDE9D26C776D4515DD587ECDDE0BF8892799
AD7856C3509643932C96B19949FF00AE806EFD6BE7E971E72E95A9FDCFF47F
E67BCF0BD99F89FF00053F681F1E7C10D7A1D4FC337AE6D03832DB331D8C33
CE3D0D7F47BFB2D7EDBFE05F8DBA441637F74B6BAAAA859239080C1BDC7F5A
F0BBCFF826BFC1A9E2FB6DC78336B6DFF96775791AFF00DF29301FA579CB7E
C87F0BFE1B5FBEA1E19D0E5D26F9410B7115E5D34887B15F325619079E4107
B822BE078EB0393E750E7E4942B2DA492FC75D51D985A9529E8F547D95FB59
FEC2DF0BFF006A3D324F10D9ECF0F78DE38C7D9F58B74044E146163BC45C79
C980006CEF4E36B6DCAB7F3BFE34FD8FBE3D7C3EF195D7833C5BA07D85ED5C
817AD20FB14F1F6960907322118FBABB97A32AB0207EE8FC29FDA43C4BF0CA
F20F0BFC4C90DEE90EC23835203006780B30FE06ED9FBADDB04E2BEEFD6F41
F007C69F0A7D8B538A2D4EC6E17723A91E644C470C8C39561FAF43915F9D70
FF001DE6FC352583C52F6943ECF5B7F81FFEDAFE56D4ECAD84A75973477FEB
73E09FD8AFFE09B1FB37683A1E93F163C49AB45F15B55902C91ACD12A69569
30E4A1B36DCCF2C79C379E48C80446A6BE53FDB6FF0067DBFF00D9D7E273F8
9BC3B6E57C0DE2B9DE6B1641F2595CB7CF2D9B60600072D08EF1FCBC94635F
5EDA691F11FF00626F1DB78C34E59BC43F0E754758F5258865D62DDF2BB264
2A4F1E4EC7FB8E09525770C7E9478D3C1BF0F3E3FF00C319FC39AF247ACF86
3C4F6B1CB14B137252402486785F076BAF0E8D8E0F5046457F4AE559B6133D
C12AB4249FE8FF0035E8CF16709529599FCD1F863C6EAC1417E78EF5EDDA3F
8C15957E7FD6BEA1F85BFF0004C6D1EC7C41E2A83E22F8967BDD26DE630E91
369AE2190C38C99A61223A890121360CA86563960457C63F1FFE13F887F66E
F19C7A0DF6A31EB1A45F067B1BD8C6C765523293C793E5C8B91D09561C839D
CABF92E32382AF88786A73F7AED2F371F8B97BDB66FBA76D99E941C92BB3DB
EC7C5A303E7FD6BA683C5431F7EBE26D37C74A40FDE7EB5D7DB78D9481FBCF
D6BC6C4F0D34F6368D73EBA5F150C7DFFD6A293C5431F7EBE5F4F1A0C7FACA
1FC60CFF00758D70FF00AB8FB15EDCFA26EBC54307E7AE4B50F160C1F9EBC5
A4F114D2FF001605517D573CB364D76D1C8947725D63BEBED7E5B827E6C0AE
5AEF5200124D73371AB051F7AA7F09F873C5FF0013BC5169E0BF0269D2EABA
B5EB00B1C43844C80D248DD1235C8DCEC428EE6BDDC2656E4D4608C6553A9D
8FC38F01F893E3378FF4DF87BE1518B9BF7CCD3952C96B6CA479B3C8063E54
07A646E62141C915FD197813C17A17C3AF07E93E08F0D43E469BA3C0B04438
DCD8E59D88C65DD89773DD893DEBC5FF00668FD9C7C3FF00B3D783DAC22917
51F11EA7B64D4F50DB8F31C0E228B3C8863E7683CB1258E09C0FA52BF67E1C
C8D60E9DE5F13DFCBC8F2EBD5E661451457D19805145140051451400514514
01FFD6FDFCA28A2800AF09FDA12F1ED7C13688870B71A8431B73D84723FF00
3515EED5E37F1E74B9B52F86D7F35BA97934E78AEB03FBB1B61CFF00C05198
9FA5736320E54A49762A0F53E53D3AEC00A735EB3E1BF194BA48011B8F4AF9
CF4ED4C141CD74716A3C70D5F976270AF98F42323E8CD57E25CF3DBB22B841
8E48E2BF1BED3E21FF00C2C0F8A1E2ED78B068EFEF59A1C7781144511FA944
04FBD7DA9F11BC45F61F07EA51A4C62B9BE85ED6DD87DE134CA5548FF7796F
A0AF3FF819F02A2D5923D3B43B648D6300CB330E4FA963D4935F13C619950A
385961E69CAACDAB25D3FE1F648F43014A4E7CFB247C3BF15FC0DAFEB5A85B
DE68D62F7053707230A36919EAC40EA057CFF3697A869977F66D46DDEDE41D
9D48CFD3D6BF72FE247C0DB8F07697F6F32ACF1F4240C60D7C0DF127C35617
D6B2A4918DCB920E3907D4563C19E2456C15B0189A5EEA7E6A4BEFFCACBD4D
71B96C6A7EF2123E421F0BECFC567ED7A64C34FD531C3E3314A7D24039FF00
810E47A1E2B0352D2FC73F0EE68A1F1258C96D14A488A5FBF0C98E4EC9172A
78EDD47702BD6BC2DA80B0BD7B690E1E07DA7FA1FC6BEA5D2755F0FF008834
59341F125AC5A8585CA8124328DCA7D08EE18750C0820F20835FD0CE14F134
D5483D1ABA3E6DDE2ECCF8A74AF1E3ED5CB115DBDAF8E0B01F3559F1FF00EC
EB269D349AAFC34BEFB65B31DDF61B860264CF68E4E15C0FF6B6903BB1AF9F
6E66D5F40BB361AE59CD6170BD63991A36FAE18038F7AF9EC46169AA9EC9B5
CD6BDBADBBDBB1AA93B5CFA453C64EC38269E7C51349D0E2BC02D7C40A7186
FD6B662D7971F7AB99E0EDD0AE63E51F16C46D7C61AD404E40BC9C827A9567
2467DF06B3A26AE97E26A2C7E319EED3EEDEA472FB671B0FEAB9AE3E27AFB4
C3CAF4E2FC8E67B9AE8D5615EB351EACABD6C234164C55859C8EF5981EA40F
401ABF693EB4D3704F7ACEDF46FA00B6D2935133D405E985E80252D5196A8C
BD44CF400E66A81DA919EABB3D000ED54E46A73BD5476A0089CE6AAC86A663
5A1E1FD0753F15F8874BF0B68B1F9DA8EB3750595B273F3CD71208E35E013C
B301C0A52924AEC0FEB37F62BD0E5F0D7EC81F0BB4E9D046D368F1DE001767
CB7CEF74A71EE25073DFAF7AF43F131055EBBAD2F45B4F07784746F08D836E
B6D12CADEC6238C663B68962538C9C70BEA6BCD7C477036B0AFE619D6F6B88
9D5FE66DFDEEE7B495958F9C3C71107493F1AE5FF65FF1FF00FC205F1CD3C3
9772F97A678C63FB1B83F745DC797B663EE4968C7BC833C735D3F8C650564A
F8DFC63AA5F683AE58F883497F2AFF004BB98AEADDFF00BB2C0E2446FC1803
5F73C3D5DD2AB19AE873565747EF4F89EE961B3284F26BE6AF10DD0676E6BB
1D43C6F6DE25D1ACB5CB36FF0047D46DE2B98F9CFC932075FD0D7C2FFB427E
D1BA47C3247D0F4809A9789A740C96E4E63B756E8F39041191CAA0E4FB039A
FDC70181AB89A8A9508DDB3C3C6E36961E9BAB5A5648F43F1CF8DFC31E09D3
9B56F14EA31585BF2177925DC8EC8832CE7D94135F0AF8E3F6B0D5F52796CB
C01A68B287902EEEC07948CF55887C8BFF00022DF406BE64F11788BC49E3AD
664F1078AEFA4BFBC978CBF088BD92341F2A28F4007AF524D4505A00318AFD
B387BC36A14D29E2FDF976E9FF0007E7F71F8BF10788D5EA370C2FB91EFD5F
F97CBEF2F6B9E28F1978B199BC49ACDD5F2BF26379088867AE225C20FC1457
3EBA64607DD15D225AFA0AB1F646F4AFD3B0D9452A71E5A71497923F36C466
D56A4B9AA49B7E6CE4DB4E4C7DD154E6D3233D56BB46B623B5529E248D1A47
215541249E800EB5557051B6A453C6CAFA1E59E2297FB3ACBC8DEC4CC7091E
E3B7239CE3DAB8886DA6C0789B0DEFDEAC6A1A8B6B5AA4B7A73E5676C40F64
078FC4F535A16E838AFE0DF1338AA198E612741254E3A474DFBC9FAF4F2B79
9FD53C2193CF07838C6B36E6F57E5E4BD3F327D3F5CBBB0754B90547627A7E
75E93A5F8A0100EFFD6B878A20F8461907B1AEA74AF0B5B6A0EAA81A227BA1
C7E9D3F4AFC831B86A357E25667D6C1B47A558F88C607CF5D4DAF884103E6A
C5D1FE0678BB5750DA0DEC52B119093868CFD372861F9815A9A9FC0BF8EBE1
F88DCCFE12BDBDB755DDE6D8A7DB176F72443B9940EFB80C75E9CD7CAE2F87
A12F84DE354E8A1D7FA7CD5A31EBE3FBD5E0526A97F6172F677F0C96D711E0
3C72A94752464655B04715762D79FD4D7815B86FC8D5563DE46BC31F7A94EB
A3FBD5E249AD4ADD0D581AACADD4D713E1E48AF6A7AEBEBA3FBD5465D773FC
55E682FD9BAB54A2F3DEAA392C5740F6A76D26AF23F4354DAF598F26B97FB6
7BD34DE7BD6F1C025B2279CE91AE8FAD567BBF7AE79EF47AD5392F80EF5D30
C089C8EAECEF00BE8189E15C37FDF2735EB965E24181F357CB77FAF2D96C72
F82CD81F956CE9FE2C040F9EB3CC387A5562A5635A55AC7D636FE22047DEAD
88BC423FBD5F32DA78A4103E7FD6B721F138C0F9EBE4311C2EEFB1D4B107D1
2BE2018FBD43788063EF57822F8987F7FF005A6CDE2A8614324D28451D4B10
07EB5C2B85DDF445FB73DC65F107FB559571E201FDEAF9FEEFE25F87EDC959
3518891D95B7FF00E839AE7E7F8B1E19C9537FD3D2394FEA16BD9C2F03E265
AC6949FC9994B171EE75FE36D58E97E22B5D762388EE8795211D03A7DD27DC
AF1FF01AF67F06F8E5248A3224F4EF5F256ABE3AF066BB63369D71A96D120E
18C722EC61F75812A39079FF00EB5723E18F1FBE9176D633DC2C8236DA1D1B
28C3B1047AD7DBAE13A95F06A954A6D4A1DD3575D3EED8F3E759295D3DCFD6
2D1BC64ACABF3D77769E2F5C0F9EBF3CFC3FF11A3745FDEFEB5E9B65E3F42A
3127EB5F96E65C112527EE9D50C41F689F16AEDFBF59777E2E183F3FEB5F2C
FF00C278B8FF0059FAD65DDF8F1707F79FAD7954B82E57F84B75CFA0B55F17
8DADF3D78DF88FC6400601F35E55A978D9E6C88DF39AE0AFF5D6909791F26B
ECB29E125069C91CF3AD736B5DD6DAE19A491ABE5EF1EEB82EEFC5AA36561E
4FD4FF00F5BF9D763E2DF1747636EDB5B323642AE7A9AF9F2E2EA4B891A695
B7339249F735FB270E64DC8B9DAF43F4DF0C32694B13FDA3517BB0BA8F9C9E
97F92FC7D09A59CB9AE7EEF4DB699CCD19314A7BAF43F5157DA5F4AAAF257D
DE1653A6EF0763F5ACEF0B86C6D374B170528F9FE9DBE467A457F0FCAB2AB0
F5391FE35A164E9F6C83FB56426D8B80E10E0E0FBFF3A84BD674AC656C0FBA
2BEAB2DAB5EB4AD27A75D0FC238C386B27C161E4E9C1F3CB48FBCF47DF7E87
DDBE1EB2D12EBC389A6E9902436FB7215077FEF13D49F73CD779E1790347E4
BFCB2C276B8F7EC47B1EDF976AF963E14789EEECC41A6DEE4A499F25BAEE0A
7047D47F515F5643A7CD6BB75D8FFBB8741C929D73F51DBF11DEBE6B8AB134
30F0B579D9B7A79FF5D7B1F88E551AD4710E16BC5EFF00E67BEF81B4F5BDBF
8A361D48AFD5FF00845E0FD3F4CD0E1BAF294CB20CE71D2BF24FC05ACC56D7
50CE18153820FA8ED5FA83F0AFE28692DA4C5677728428060E6BE4F06E9471
10957578DF53EBDDEDA1F4A496D1142315F20FC78D12C21569A340ACEB938A
F7FD57E24787AC6D5A5FB4AB903A0AF867E2FF00C4C4D6E79446FF002F200A
F578A31B84AF28FD563B6EED633A3192F88F94BC4B0DBC8D35BCA8B246F956
5600A907A820F51595F0FF00E28F8C3E0AEA4B3E812C9A86825BF7B62CC59E
15EE6127EF2FFB0791FC27A2D57D6350F3A566CF5AE50DA5F6A5916511907A
F41F9D7C8E36861AA5174F176E47DFFADFD0EBA4A6E5FBB5A9F7FAFED19E01
F8CDA6FF00C23F68EAD653201701BA34BDD194FF000A9EA0F7FA035EDFF01E
D357F0A787E7F87BA1C61B4286479AD27F30E6D5677DEF04631CA962CCA770
2993D46D03F176EBC11E24D2F5FB7D774894E9570D2289E55F9A374CF26441
8C903A30E7D7238AFD4FB1F8BFA4FC01F86FA3A6ABA8A788EF35041712F932
46561B77EBB1D0619BEA79C119039AFCFE961E594E279B2BC54BD9CD35251D
5DBB76E67B27BADFA1CB4726C4CF152AF55E9DBFAE8BFADCFA23E29FC43D17
E18F8526BA95C2F94B88D07577E8001DFD00AFCAFF00DA73F63BFDA6FE25E9
161F152C447AC4F2C2D7171A0A1F2AFACF712542073B26611E0B282AEAD955
573CD7D23E15F195DFC50F1CD9FC61BAB34D53C33A34DBB4DB2918A99980FF
008F9C741E59FF0056181CB02481856AFB6355F8B1A7EAFA3DADAF8165171A
DEACE618A2700496A40CC924C8738D80E40E431C632B935F6FE1DCB2BA7F58
C7E61512AD056505F622B68C57DA7D1DAEF9B7D598E3735A7CCE94257B6FE6
CFE4EAE7FB5FC3DA84FA46B769358DEDAB949A09E368A58DC75574701948F4
22B6ECF5C538F9ABF787F6A2D2BE0C68BE07FF008473C5BA158EBDA9441A49
AF6E230F76259396D938C4AACEDC9C373DF9AF03F047FC12C74CF1978120F1
1F887C417BE0ED7B532F3C564B125E456B6EE3F7292AB3239931F338F30633
B7A826BEDB86F3459BCA70A749A70DDE8D5FB5FAB5D74B5EF66F739A8E3549
B8DB547E5EDBEB4BC7CD5AD1EB4BFDEAFD09D4FF00E092FF0011ADA4C683F1
134DBB8F730CDCDA4D6E760FBA708D3727B8CF1EA6A0B2FF008251FC5C6B84
4BFF001E68D0C073B9E386E2471C7185654079FF00687AFB57D04F86AB37F0
9D6ABA3E04FEDC503EF5539FC451A024BE2BF5B7C1DFF04A2F0C5AB4537C40
F8837DAA00417874DB58EC971C7CBE64AD7048CE79DAA71D875AFB7BE16FEC
8DFB3EFC20923BDF09784ADE5D4E22185FDF937D74AEBD191E72C223FF005C
827E75B50E11A8DFBED213C4AE87E347C08FD8E7E32FC749ADF55BCB593C27
E15936B9D4AFE22AF34670736B6E76BCB90721CED8CFF7C9E2BF70FE0B7C05
F871F01BC3BFD83E04B02934C17ED77D39125E5DB2F4696400703270AA1517
B28C9CFB3515F5B9764F470CAF05AF739A755CB70A28A2BD5330A28A2800A2
8A2800A28A2800A28A2803FFD7FDFCA28A2800A82E6DA0BCB696D2EA312C33
A346E8C32ACAC30411E8454F45007E5AFC43F0A6A5F0CBC5971A15DAB1B294
B496531E92C04F1CF4DCBD187AF3D08AC18759523EF57E98F8FBE1F7877E23
686FA278821DC065A09938960908C0743FCC1E0F715F9D3E3EF81BF117E1F4
F2CA96AFAC694A495BBB452F841DE58C65A32075272BE8C6BC1C66549BE689
B46A1F3DFC56F1093AF68560CF889526948CF5625541FC067F3AFA5BE027C6
1D1FC1514B0EA2B94980E47518AF84BE38CD736BA6E9BE2E88E534B94C5391
C622B82AA18FAE1C28FF00817D6B88D1BE2242F029597B7AD7E0FC73916223
8E588A774D59A7E87D0602B45D3E467E9F7C69F8F963E2AB1FECBD286D8072
49EA4D7E7DF8D35D8CC527CDD735C86A1E3E8FCB3FBDFD6BC53C55E363721E
38DFAF5AF9CCBB23AF5ABBAD59B726EEDB3AA75A318F2C76286992C17FE2DB
C8DE7300316F040C825180C1FF00BEABD52C65BBB40041769263D72BFE35F2
70D56FD2F24BBB585E4461B772903BF3C6726B4A2F1A5FDBF3289A303BB2B6
3F3AFD9697F6A61E115879350B6CD26BF2B9E2CBD949BE6DCFB134CD7F5483
50B67B821A0571BC86CE17B9C75AF7CB7D67E1EEBD691E9BE286D3EEE17C85
8EF4465493C1DA25EFF4E6BF362D7E244A0E05C827D09AB7AA7C4569B4E686
46562C576F3DC1EBF957C87136518CCD6AC255DF2CA3A271D1FF005F717CF0
A509387E27E894DFB2EFC08F17C46E34DB5974E32609974EBB62391C6164F3
5075ECA07F5F16F8C5FB2DF80BC01E1B9AEFC2FAE6A379AE5DE23D3ECAE658
3616380D24CE9106F293963B5727851C9CD7C7967E39781C4B13B46E38DCA7
079F715D7D9FC56D5A3D422D4DF509A6B888A90D3399785390087C823D8F15
E660B86B3DC2D78CD66139535AF2BD6FE576DE9DEDADB6B1E4D6CD538351A4
94BB9C9F89BF678F186A3670DD5DEA76DF69841C2089D5483DB79627A8FEED
7CD5AE787B57F0B6A0DA6EB30F9528FBA41CAB8F553DC7EBEB5FBABF187E28
7C38F1DF807C21AB786B4C834ED56EB4E826BF481764693BA02EAA324ED0D9
C679C6335F96DF1BA3B3BEB079401E6C077A9EE31D7F315FA3F0E71362A588
F6188B357B5D2B7F91E4D0C5CDCF964EE7CC89255957AC88E5AB4B257E9E7A
66987A903D67AC952092802F6FA378AA7E652EFA00B5BE9A5EAB1929A64A00
B05EA267A80C95133D004ACF503BD46D255767A0073BD572D416A899A80119
ABF40BFE099BF091FE25FED3BA5788AF60F3348F01C4FACCECC3E5FB427EEE
CD01C603F9CC251EA226AFCF776AFE9DBFE09EFF0003D7E06FECF967A96AF6
C6DFC4FE372BAA6A01C6248E1208B480FB24477904021E4707A57C9F1A66BF
56C0C92F8A7EEAF9EEFE4BF1B1BE1E1CD23EE7D6AF81DC335E2FE22BD186E6
BB0D6753015B9AF18F10EA5C30CD7E2184A27A726796F8B6F0157E7D6BE40F
1CCA2699875EB5F44F8B7521B5F9AF97FC493F9AF23E7D6BED32DA76D4E69B
3DDB5CFDA3FF00E15E7C02F0BD8E9532DC789EFAD1ADAD2371B84115B4AF07
9D2293F75026D4078661D0A86C7E79B4F7FAB5FCFAB6AB3BDDDEDDBB4B34D2
B16791D8E4B313D49A8B58F327D726323B388F0A809242AFDEC283D0649381
DC93DEB46D23C015FDC7E1C64B4E8E029564BDE9C536FD56DE87F38F1E6775
2B62A745BF760DA4BD3A9A16D074E2BA5B1D3649C8DAB9CD45A558B4F22A81
D6BE91F01F8064D49E3FDDE738ED5FAD51824AE7E5D5AA36CF2ED3BC1F7570
0158CFE55BE7C0776A99319FCABF417C1FF0504D023B4381F4AEE750F82F14
76E4AC60903B0AF1AB7166069D4F66E7A98AA136AF63F27750F0C4F6F9DC86
BC27E2A5E1D1F475D3D0ED9B516318F5F2D705CFF21F8D7EA2F8F3E1D2D909
3F778C67B57E4A7C76BC597E234FA444418F4A892238E46F71E631FAE1803F
4AF90F15F89161321AB568CB5A9EE2FF00B7B7FF00C9533EDBC39CA7EB79AD
3535EEC3DE7F2DBF1B1E7164B8C57416E2B0AD3802B72035FC078967F58C4E
86C63DCEA2BDDFC0DA6ACD3C791DC578769A7E7535F4478066449A324FA579
1511713F417E0E7866095A12CA3B57E92785343B5D3F4F88AA0DC475AFCF3F
845AD416FE4E580C62BF40FC35E25B2B8B28D1A400815143954D39EC366C78
87C1BE13F14DBFD9BC51A2D96AF1608D9796F1DC2E08208C48AC3A123F135F
11FC62FD93BE03DD096F2C7C391E9170D96DD6123DBA8E3A0894F9407FC03F
AD7DBF7DE22D3ED612ED28240AF953E28F8EA0952550E3BD746633A5392F66
8504D6E7E5678F3E05687E1D9A4FEC5D4EE022E70270921FCD427F2AF9F351
D3EE74B9191A55900EE322BEBDF88FE2149E4970DEB5F22F88AF44923F35E6
3C241EE8BE630FFB5A243876C1F7A9935681BA48A7F115C5DD4996352D8787
7C45ADFF00C8174BBAD43271FE8F03CBCFA7C80D2FECA83D85ED0EDBFB4171
D6A36D440EF5CD4FF0DFE24C2AD23F85B568D47249B1B80001EA76570DA8AE
ADA6CCD6D7AB3DA4C392926E461F81C1AD2390A7D41D4B1EA726A600FBD593
71AC2A8396AF249F51BF00FF00A43FFDF46B99BFBEBC91583CEE41EDB8E2BD
2C370DA6F5643AC7A76B92EA17F68BA95BC323D95B48627982931AC840210B
0E0360E719EF59765AAC91900B115F45FEC1AFA7EBFE39D7BE1DEAF147736D
A95B2DDAC5300C8FE5308E552ADC36E591723D17D2BF473E287FC130FC2FE3
2D2C6BDF097565F0CEB0E9B8D8DD6F974E95BD9D774B0127A90245EC105774
E8C69D6FA9B8FA3EF712936B98FC86B3D6A4E3121ABB77E2FB7D2E30D71296
76FBA8BCB1FF0001EE6BA3F8B7FB397C71F80CD2CFF10BC317369A7459C5FC
005CD8B0CE01F3E2DC8A5BB2B956FF006457CD31CD2DE4ED777072EFFA0EC0
7D2B4A1C350A936EA2D103C434B43D22EBC79AFDF652D9FEC919FEEF2DFF00
7D1FE80566C706A1AA481E7924B87F5762C79FAD54D22C8DDCE91A8CE4D7DD
DF043E054BE2B9612D16E0D8ED5EE51C151A3A528246329B7BB3E398FC27A8
B26FF28E3E959379A45C5B64488462BF746E3F6458ACF45FB435B8CEDCF4AF
CFFF008C9F0B22F0E4F32AC7B76E7B56CA449F084D0E3208ACB70636DCB5D7
EAD6DE44CCBD306B979C57441DF460CD7D27C4D7964C1525381D89E457A669
BE38BA200F36BE7CBC5C72A7047714DD3754946E82462644E41EE57FFAD5E7
63F21A5523CF63D2CA703531757D8536949ED776BF97A9F57C1E2FB9900CCD
D6AD8F1196E5A4CFE35F32C5AEDD47C2B1A9CF896FB1856C57CDCB86A37D0F
A15C199A5EDECBFF00268FF99F45CFE2445192F5B3E04F077C4AF8CDAE8F0C
7C2FD0EE35BBCC812C8836DB5B83FC53CED88E3180719393D1413C573DF087
E125FF008E6F6D755F184CF6BA3390CB0A92B2DC2E7D472887D7A91D31906B
FA3AFD99F46F087867C036FA2785ACEDF4F86203F750AAA0E9C938EA4F7279
2792735EAE59C2D4DD48AA8F73871796AC226EAB5292E8B6F9BFCD2FBD1F97
1A6FFC1323C43A669E3C43F1435CFED1BB0373DA69D95813FD93338DEC3E8A
9F5AE62F7E107833C1529B6D3743B681E2CA8768C492E3DE47DCE7F135FBB1
E2DD5F4FD2F44BA96F64500A10013D6BF21FE29EAB6B75ABCEF6E460B1E95F
4F8AC0D3A13E4A7AA3BE8E7F8AC451519CAD15A24B456F447CD3AFD8432C4F
1B202A46304718AF8CBE28783EDEC5DB55D362116D3FBC4518047F7B03B8EF
5F6B6AF3A90D5E03E3748A7B6991C643022B6A54D4972C8CF0B98D4C2D4556
93FF0083E47C652396F9456DF877C35AAF892FD34ED2603348DD4E0ED45FEF
310381FF00EA1CD7DD1E03F057C1983C33A66B736950CB753C28D2F9E4CD89
40C3E04848FBC0E302BD024F889E07F0EC3E469F0DBDBA274540AA07E0302B
A2B66DF55A56A349C9FDDFE67959DD7AB8BA9EDEB4BD17647CC717C2FF0011
E91A6491C16B37956004C97457680F8CB363AEDC920F1C2FE75ECFF0E3C589
AD69E2D2EBE49E3CC722375565E08355F5AFDA074731C9690912AC80A1451F
781E08AF14B3D5DB44D4A1F105B1C44E552E547A74593F0E87DB07B1AFCDF8
A32AAB98D1954A90E59AD57F91E2B496C7D556D732E877DF66CE21949688FA
1EA57FA8FC6BD5745F1FDE58A811CA463DEBC92D67B7F1268EAC8DF3E01561
D411D0D3744D3356D4B784010C4C51F9E8C3DBDC735F13946794A341C3172B
3877EA8D6323DEAEFE266A37119433123EB5E7FA9788A6BB62CEF9CD3ADFC1
77AEA3CD99BF018A59FC0D301F2CCFFA566F8D32E52B733FB8AE639182E06A
BAB2E9FBBF7698693DF3D17FA9FC2BF433E0AFC03B5F16686357D425F2203C
22A81938AFCF6B7F0BEA5E1FD426D4118DC09583118C1180071F80AFAC7E1A
7ED39AA782B487D1A383CD8E2E487C02BEE06738FD2BE7734C661F178A8D49
DE549744ECFF001F33E8B015A3ECB928B4A7E6741F193E1C5AF822FCD9C720
96275CA93D715F116B3E1AB1D77C61E1CF0C6AB7F2D8E83AB6AF616DA8B46F
B0A5A4F7091CEE84860AC2362776D3D0641AF7EF1FFC5EB9F1CDD36AD77307
0F9DB8E800EDF857C6BF143C450CD633C7B81DC08FCEA784F0D3A78C8C947D
DBECFD7A9E863657A5693D6C7DE7AC7817E2CFECB5E21B0D13593FDB9E0EBC
945B69FAC409B110004AC1751027C99028E0F28F8F94E72A3EA9D4FC4DF0D9
7C136F7DA7DC0FED8701D99495657EA483DB1DB15F4078EFE1C69FF1AFE0E1
F06789677B79755B18244BB55CCB6D78A8AF1CEABF2E4A49C95C80C32A7835
F98171FB177ED796DA97FC23D6FAC68579A5962ABA8FDA668879790033C262
2EAF8392ABBC0208DC7827F4DE36F091D6C4FD632E824A4B5EE9DEF75DBFE1
CFC8F1181F79B825AFF5A1ECDFB3C7812DBE357C5AD57C77E30BA3A9695E09
9603676CE0B2CDA84BB9D65958F0C200A1957A9765638DB86FD3EAF11F801F
04F4AF811E008BC236776DA9DFDCCAD77A8DF3AED3737722AAB32AE4EC4555
5545C9C01924B124FB757EC5C2590C72DC053C2A4AE96B6EACEFC3D2E48288
514515F486E145145001451450014514500145145001451450014514500145
145007FFD0FDFCA28A2800A28A2800A424282CC7007249A5AF96FF006B7F89
537803E16CDA7E99318754F1239B28597EF242466771E9F27C991C82E08F5A
4DD80FCEBFDB23E27E93F16F52D5346F0E5B5B9D134D495239D2250F7B28C1
795DF1964257080F6F9BA9E3F24AE6D35FD1E52DA34A67B7ED1B361D3D813C
11F539FAD7DD8E4302AC320F515F2BF8874D6D2B54B9B16E9139DA4F753CA9
FC462BCEAD84A7884E3555CD23371D8F2C6D7F5F93F773DBDC67FDC6C7E7D2
848352BD39BCFDC45DD73963F974FE75D6482AAB0AE4A1C3B8684B9AD73496
264CA8A8B1A8441B55780053726A6615111DEBDD48E7209638E5E2540E3FDA
00FF003AF7BF863E10F879E2CD0BECFAEE896D35C5AC8637750627656C9525
A32A7B91D7B5784119AEFF00E1C6BCDA1F8812391B6C17B889BD9BF80FE7C7
E35F35C5797CEBE0A6A936A4B556D1E9E9E46D42494B5D8F6FF15FECF9F077
4FF0A6B3E2548AFB4E4D2ACAE6EC8B7B967188223274984A4E36F4F7FA63E0
0F066A1A56B31C526A57CF1489FEB62000CFD0F607E86BF447E29EB4E7E0EF
8B92390AB3E997099071C321047D08383ED5F8F104F35BB8921728C3B8AF86
F0FF000B88C461EBFD62AB6D4925777B69E7DEE5661463A28AB1F7E5DFC41B
48ACD2DADDC08E34088A0F015460015E7975A26BBF10EDEFBEC27CA863461E
6303B4BF641EA4F7F41F803F3245E25D5176E640CA08C823AE2BED0F857E3E
D3B58D2E2B28C2C2D10DA631C053F4F7AF5F30CAE780A7EDA8C6EFBF63C4A9
49D35747C67756D79A5DE4BA7DFC661B8818ABA375047F9E0D3925AFB17E2D
7C3087C5566759D1D426A702FCB8E04AA39DA7FA1F5F6AF8A9FCEB699EDAE1
0C72C4C55D5860861C10457D764B9C43194B9E3A496E8ECA15D4D5CD8596A5
125642CD53096BD937353CCA5F32B3BCDA5F345005FF0033D69A64AA465A69
96802D992A2692AB1909A6EE340129726984FAD465A9A5A801C5AA266A6B35
7BE7ECE5FB3CF8C3F68CF1EC1E18D051AD748B568E4D575265CC5676C4F279
E1A570088A3EAC793850CCB8E23110A50752A3B25B8D26DD91F437FC13FF00
F65E9BE377C488BC79E2DB2DFE06F08CEB2CFE60C25F5EA0DF0DB2E7865538
79874D9853F7C1AFE8B354D5800DF35795F807C1FE0FF843E07D33E1EF812D
059691A4C7B235273248C7979656E37492312CCD8E49E001814DD575B18396
AFC1F3FCD6798625D47A456915E5FE6FA9EA528282B136B5AB70C4B578BF88
B57E1CEEAD0D6F5C186CB5789F88F5DE1BE6A9C1E14252394F156ABBCB229A
F0FD767011C935D9EAB7A6466763C9AF26F135F88ADE47278009AFABC15077
5146139595D9E453B09751B893D5CFE9C56DD9AE4815CB5AC85D8BB75639FC
EBA8B161B857FA03C3D85FABE169507F66297DC923F92B3CAFEDB1152AFF00
336FEF67B178334F5B8B98C119C915FA43F063C2D6F2180B28ED5F9E1E04B8
48EE2324F715FA3FF087C436F6C21DCC3B57A19ECEA7D4E7ECB73E6E31F7B5
3EDED1B4A82285515400055CD42CE258CF1587A5EBD6CF0AB4720231EB5475
CF14DA5AC0CF248381D335FCE7F55AF3AD6D6E7E991CC3030C172A5AD8F9EF
E2DD9DB2248C000715FCDD78E7553AD7C41F11EA85830B8D46E8A90770D824
60801EE028001AFDF6F8A9E368AE0CB871DF15FCE85B5CCB732BDC4EDBA495
8BB1E996639278F7A5E2F46AD1CB30787A9D5C9FDCA2BF53DAF0868C5D7C4D
55D1457DEDFF0091D75AB56CC4DC8AE76DDF18AD985F8AFE66AF13F778B3AA
B0982B0AF5EF0B6AC2DA4439AF0BB79B6915D6E9BA91848E6BCB9C0A3EFBF0
3F8E7EC823FDE6318AFA7F43F8B86085479D8FC6BF2BF49F14490630F5DE5A
F8F268D40F30FE75CEE05291FA39ABFC6267848F3BF5AF9D3C67F129AE83FE
F739F7AF9AAEFC7B33A91E61FCEB85D53C552CF9F9F39A9E41F31D478A7C4C
6E5DCEEEB517C3DF83FE32F8B57DFF0012A8FEC9A6AB624BC954EC07B845E0
BB7D381DC8AD2F837F0D6EFE29F887CCBE0CBA359B0F3D81C191BA88D4FEAC
7B0F722BF613C05E0CD3B46B0B7B1B0B748218542AA22855007600568D28AE
66694E1CC7807C33FD917E1DF86238AE6F34F1AC5F00099AF0097E61DD508D
8BED819F7AFABAC3C01616D12C70DBAA2A80000A000057AD68BE1F0C8AC570
B5D9C7A242ABF7735CD175AAFF000E2D9BB708E8CF9DEE7C156A54FEE87E55
E53E32F84DE1CF1059BD96B1A6C17B0B6729346B22FE4C0D7DA577A2C654ED
1CD79EEB1A604DCA45446B5484AD2D194AD25A1F8C3F18BF62FD16749B50F0
248749BA193E4392F6CE7F1CB267D4647FB35F98BE32F0C6BDE0ED566D13C4
766F6777173B5BA32F40CA470CA71D457F4E5E25D263911C15AF83BE3FFC1B
D13C7BA2CD677B16CB88B2D04EA3E789FD47A83DC7423DF06BEA32BCC2ED29
1C95A8F547E6AFEC7BE263E15FDA4FC237ACFB20BA967B4941214309A07551
93FF004D36903B918AFEA8FC23E32D2EF74D855A65575503935FC8BF8634CD
67E1EFC6BF0F586A29E55DD86AF6809FE16469546E5F62A723D3BF35FB43A2
FC60B9B385544C463DEB3E2672A75E9D5A6F5B7EA61476699D1FFC1577E2FA
E95F077C3FF0D347BA659BC5DA9192E7637CAF67A728768D877CCD240C3FDD
AFC18B4C6DAFACBF6DFF001C5D78C3C6FE1BF3E4DF1DA5849B47A34929DDCF
FC0457C8B6B2702BE9729E69612139BBB777F898D4F8B43D63C15E57DBE232
74C8AFDA1FD96FC41A1E9896ED7054631D6BF0D746D40DA4CAE0E306BEA0F0
37C5EBBD091047315DBEF5A4E207F45DE2AF89DE1A83429764CA4B274C8E2B
F1A7F68BF19586A577706060724F4AF2FD67F686D4AEED0C3F693D31D6BE64
F16F8E6E75895DA490B67DEA5A6DEA0715AEDC2CB70E47AD7233364D5AB9B9
323124D654B26056D4E026CCFBB3C1AE61E5F22F22907009DA7E86B72EA4EB
5CADFBFCD91DB9AF568D3BAE566B83C4BA3561563BC5A7F71D409ABD3FE19E
8169AC6A7F6FD514496B6A46D8CF477EBCFB0EA477FA66BC783D7AC78035B4
B4416F9C15639F7CF35C5430A9CB53F6AE23CEEA53C2BF64F56ED7ECBFAD0F
BC340D5C41B1D0E318AFA7FC13F1AB55F0E42A96D394C7A1AF81747D7D4A2E
1ABBAB5D786D1F355D4A3A9F99C2AE87DAFE2AF8EBAD6BD098A6B96607DEBE
7CD5F5E7BB7691DB24D79CB6BA08FBD59377AE0C1F9AB38D22DD546C6ADA92
ED6E6BC3FC5BA8A98A419F5ADED5F5D015B2D5E33E20D5BCF2DF3702BBF0F4
7538ABD53B4F07F865B5BD192E2E751B88E0DF22AC3110AA30C73C907A9C9E
D5F487C1DFD996FF00E2CF889744F09E8DF6B542A6EEF2E8B496F6D1B1C6E9
4B92B9C676A01B9B0700E0D7AFFEC13FB3A683F1DFC2975ACF88B59FB2D978
7AF4DBCF616E3FD266128132BB48DF2C71B6E2A080CC4A37DDE09FDC7F08F8
33C2DE03D162F0F783F4C874AD3E1E9142B8C9EECEC72CEC7BB3124F735F47
4E8C525A1E054A8DBDCFCF4FDA03F61DF0B5BFECB37FE13F86D646E7C53E1D
93FB6A19D23027BF9E142B343B101E2484B08A35FE309C92589FC27F08CD75
AFDCA693691F9EF3020A9FBBB7A12DEDEB5FD8557F3E5E2FF0F7C3597E36F8
D7C4BF0CAC05968FAA6A2F24614E62661812BC23036452CA1E444E80360614
055F87E3FCE2965F84F6EF59BD12FD7D17F90E95EE79CFC36F00F8B347B87D
3EE2547D3A3DA6298B1DE548FBA5793953C649C11CFB0FAA343F0E416D1848
23E58E58E3963EA6A9787AC8394451D6BEBAF869F0D3FB5D56E6E57118C64D
7F15E7F9D56C5556DEEFB2DCEB843A23C360F0ECECA0888E3E955EEB4578C1
0C98FC2BF43ADBE1CE870C0105B83C7535E5DE3BF875696F6AF75689803A8F
4AF0F1B966330F1552BC1A4CDDD2695CF85751D2D0A9056BC9BC47E1D8A556
651B587423822BE8FD76CBECF2BA11D2BCBF588410D5DD94E3A71926998BD3
53E2BF10DCEA1E0BB69D199E6B7925670C79DBBBB7D2BB9FD953E0CF88FF00
693F8BDA74725ABB784B41B88AEF59B965FDC88A36DEB6D93C19272BB028E4
2EE6C616B77C6FA3DBDF5ACD04C8191C1041F435FB05FB127897C03AE7C0BD
374DF03E8F6BA04BA239B3D4ACED942837A88BBAE18925DFCF5DAFBDC939CA
EE3B335FD41E1A7D5F1D53F7BF1AD7D7FAEA5D6C74F93959F5DD14515FD007
94145145001451450014514500145145001451450014514500145145001451
4500145145007FFFD1FDFCA28A2800A28A2800AFC97FDB93C4F26A9F14ECBC
388E4DBE876118D9D84D724C8EDF8A7963F0AFD68AFC3EFDA96E65B8F8F5E2
E9266DCCB3C280E00F952DE3551C7A0005675761A3C04B5798FC4AD04CF6D1
6BD6C99317EEE6C7F749F91BF3E0FD457A416A79105C4325A5D20921994ABA
9EE0F5AE48BB3B94D1F20CA9549D6BB8F14787E5D0B5192D8E5A16F9A273FC
4BFE23A1AE3DD2BB93BEA419ECB5111E957196A12B4C0AC56906E52194E08E
411C62A72B4DDA2803DA24D55FC61F0E75BD21BE6BAB9B0B8B72BD7F78D1B0
460323A9C1FAF15F982A6BEF0D07579745BE13A12637E245F51EBF515F1EF8
FB435F0EF8BF52D3A21FE8C6532C047430CBF3A60F7C0383EE0D7CFE5596AC
257AB08AF765AAFD57E46D527CC91CA838ADCD035EBDF0FEA11EA162E4153F
32E7861E86B001F5AF67F851F00BE29FC64D4EC6D3C19A0DD4F63777096EFA
8342E2CE1DC70CCF2E36E106490093D80C900FA98FC4D0A54A53C44928756F
444D3A12A8F920AED9F54FC3CF1E58F89F4D8DB78DC46194F507B835E7FF00
193E142EAF1BF89BC3B1017B18CCB1A8FF005CA3FF00661DBD7A7A57B6FC6C
FD9A9FF639F08E9D65E28325FF0089F58BAF3ADAF2DB72D9ADA460E639323F
D6BE7250F2BB46095EB8DE0CF17D878974E46570DB8608F435F9AE2694B075
238EC14AF07FD59FF57F99E4E2284E84EE8FCE80EC84AB6411C106A6590D7D
B7F173F65DF155EF852EFE2E7856C59AD91B335B2A9DF3463EFCF1AE39DBDF
1F7B92391CFC2E24C57E859466F471B4BDA517B68D767D8F4212BC548D2121
A7F986A8092A40F5EA145CDE6977555DF4BBFDE802CEEA4DD55F7D34C82802
C17A633E2BD3FE1CFC13F8A5F162754F046833DDDB1386BB9310DA2E0E0E66
930848C7DD04B7A0AFD36F825FB077823C1F2DBF883E2DDD47E29D4A3C3AD8
460AE9F1B63A3EEC3CF83FDE0AA7A1435E3E639E61F0C9F3CAEFB2DFFE07CC
D2149B3E29FD9BFF00648F1EFED017A9AABEED07C210BE27D4E64E65C1F9A3
B5438F31FD5BEE2F739C29FDF5F875E09F037C1CF07DAF823C01A7C7A769D6
A32C401E6DC4B801A69E4C03248D81963DB006140032C6B56D676D1D9D9AA4
10408B1C71C602A2228C2AAA8E00006001D05615EF89073F3E6BF2DCEB37AF
8E97BDA416CBFCFBB3BA9D3513D0F50D7F8386AF3DD57C40006F9AB8BD47C4
A30DF3D79D6ADE25CE42B570D0C114E4749AE78871BBE6E6BC9753D4DA762C
CDC554BED4DE62599AB95BDBFE0F35ED50C3A466E443A95E800F35E1BE36D5
C6C16CADF34A71F80EB5DBEB9AB2431B92D802BE6DBED73FB635492E11B312
9DA9F41DFF001AFD33C3CC93EB38F84A4BDD87BCFE5B2FBFF0B9F2FC5798FB
0C2492F8A5A2FD7F03AAB393815D25A4D8239AE2AD26E057416F3E315FD7F8
4AFA1FCE18BA3A9EBFE1DD57ECF22B03D2BEA0F05F8FCD9041BF18C77AF886
CEF8C64106BB1B0F10C90E30D5ECC2AA6ACCF1AA51D4FD30D3BE3018A10BE7
638F5AC0F107C5E79A261E767F1AF8562F184EAB8121FCEA95DF8AA6941CB9
FCEB2860F0F1973A82B98B833D8FC5BE3A92F4BFEF339F7AFC9A83104F2400
EE11B1507D70715F6BDE6B0F29396CD7C63E218CD8F8A353808C01712301FE
CB9DCBEBD88AFC43C78C3BA986C3555B45C97DE93FFDB4FD83C23AAA35ABD2
EE93FB9BFF003366DE4C815B10CB5CADB4DD2B66297BD7F2862296A7EEA99D
1C72568437254F5AE7639AAEA4B5E74E916A475906A0CBD0D68A6AEE3F8AB8
A597D0D4A2635CD2A451D7BEACEC3EF532D0DDEAB7F6FA75A0DF3DD48B1A0F
567381FAD72A6626BD9BE0369A9A9F8FEDEE26194B146947A6F3841FA127F0
A29D0BBB01FA9DF03BC1963E14F0F58E956AA3F74837B6305DCF2CC7DC9E6B
EDCF0858ACD24698E2BE5EF044A91C3181ED5F557812EE313A0638CD7063B5
67743447BAD8DAA4718C0E0569605416CE1A300558AFDAF20C251A58587B25
BA3C2AB26E4EE54B9894AEEAF38F12448B96F5AF46BA995508CFD6BC87C51A
9A33B2A9E9C57E6FC6FEC5E252A7BF53D2C0276D4F2AD78290D5F3A78CEDD2
48E40475CD7BAEB77808639AF05F15DCA957E6BC1C12773AE67E6BFC6CF86D
26B7E34F0DEB7A3C0D2EA569A85B284894BC932995484551C93BB18FA9C75A
F6BBCF835F1A74BF086A5E339BC2F75F62D2EDDAE648F282768D065CAC4583
1DAB962300E01C0278ACFF008917D7361326A9A7C9E55D594A93C2E3AAC91B
06461F4201AFA8FC49FB74E86DA644B11065B9854BC2838CBA8DCB8EFD718A
F6731C2D5ACA9B8C6E91CB1E5D6E7E2BFC47BE6F1F6AF6DAACCA6C56D6130A
8670C5BE62D93C0C75E9CFD6BC8627685DA263CA120FD45751E23D33C585F5
5D62C74C96DF4AB77964469484710863B7E4621F85E4F1C579BDB5F34A77C8
72CDC9AFB2C0E1D469F241AB2F9D8E09CB53B786E318AD586FE441F2B571B1
5CFBD5E4BAF7AA9D10B9D63EA7330C6EAA725CB3753589F69F7A6B5CFBD66A
88731A2F30ACE9A7AAB25CD67CD700724D6F0A220B99B00D66E9DA75CEBFAC
D8E8B678F3F51B88ADA3CF4DF33845FD4D41713EFE01AF43F83568B3FC41D3
F53987EE7486FB59FF007D3FD5F3EA1C83F406BD4C3D2D5221B3EBF83F613F
10821AF7C69A6C7160E4C714D2367B7CA420FD6B9FF15FECB67C0FA25D6ABA
3F8A7FB6B53B74574B65B5F215CAF2EA18CAE49C7DDE064F5EBC7B44BF142E
4C5B7CE3F9D79FEB7E399EEC3664273EF5E9C7074D6A91E8D7CEF13522E139
69E88F9C340F178C08E46DAEBC1078208EB5E9169E2A52A3E7AF34F18F87A1
D5EF24D5F4C616F7AFCB8E8921F538E8DEFDFBFAD79D1D4756D2A4F22FD1A2
3D89E87E87A1AE5AB8333A58B3EA23E285DBF7EB2AF3C50B8387AF038FC432
38E65A7B6B6BFDFC935CEB0C6EF127A36A1AEBCF9F9B02B94694DF4DB07FAB
4E58FF004ACE861BEBEF9A40608BD5B863F41FE35B4A91C1188E218515DB43
0D6D59C75B11D11FAA5FF04A4F18C9A67C5FF167825C9106BDA42DD003A19A
C260133FF009E4E7F0EF5FBD55FCCEFF00C13AAFE6B4FDAE7C1B6F1052B7B1
6A90B939C851A7CF2FCBEF941D7B67EA3FA62AEF388F9CBF6B0F1FCDF0EBE0
4F89356B293CBBFD4235D3AD88E0892F0F96CC0FAAC65D87B8AFC64F082AC7
0C607615FA03FF00052AD71AD3C03E0BD0BCC655BED5DEE0A8030DF6681972
4F5E3CEE0743DFA0AFCEAF0C5E288D39AFE74F186B4EA62553E915F9EA7451
D8FA7FC1C636BA88374C8AFD2EF86CB6CBA0C1E563DFEB5F947E1DD57ECF22
383D2BEC2F877F15C6950AC133E53B826BF9D694DE1F131AEE37B3B9D94A76
7A9F799922115798F8EAFADADF499CC840DC302B849BE33692B06E4E5B1EB5
E05E3AF8A12EAE5915F0BD80AFAAE34E3879A518D0853E548B4E314ECF73CD
FC59771BDD4854F526BC8B549860D6DEABAB79CCCC5BAD703A9DF0C1E6BE2B
018577316CE0FC48CACAF5EBBFB07FC4797C1FFB40C9E099E709A778DAD248
4A370A6F2CD5A781B3EBB04A807192E3A9C0AF06F105E8DADCD71FF08B5B93
4BFDA2BE1ADEC12888FF00C247A644ECC7004535CA45264E4606C63FD6BF73
F0F273A18BA735DFFE1CC2A6A8FE9B28A28AFEAD39428A28A0028A28A0028A
28A0028A28A0028A28A0028A28A0028A28A0028A28A0028A28A00FFFD2FDFC
A28A2800A28A2800AFC57FDB07493A57C79D7A50BB63D423B4B9418F581118
FBE5D18D7ED457E727EDF1E0692683C3BF116D222C21DDA6DD3019DAA49960
27D06E320C9EE40A89ABA03F3449A8F750C6A32715CAE25A6676B3A5DAEB76
4D67763DD1C7546F51FD457CFDACE8977A45D35B5D2F4E5587DD61EA0D7D1E
4FAD67EA3A7D9EA96E6DAF230E87A1EEA7D41EC6AA9CDA1347CC0F1D40C95E
93ADF82AFAC0B4D660DCC1EAA3E75FA8FEA2B8692020E08C575269EC499452
9BB2B40C54CF2BDA9814B65798FC52F0D4BAC6931EAB689BEEB4E0723BB41C
923DF69E47D5BDABD6BCBF6A708A93480F913E1B68561E2AF1F7877C39AA39
8ECB51BFB782720E0F94F200F83D895CE0D7F5C5F0EB50F845E16F85B06916
70416CD6B008A182350AA8AA30A140E0015FCADF8D7C05A8F87B515F17F848
144B7916731C63E68644208741DD7232476FA74FA27C31FB69EAB1E8D1D8EB
2CD04E8A14B00CCA48EE300FEB5F97788BC318BC7CA94E82525152566AFBE9
7B77ECFA1F4990E3E8D28CA351B4DB4EEB4DBA7A1F67FEDB7E36D635FF000C
5EE95A4C4FAE388CFF00A31533111C63EF81C9C44A377B01C115F917F09FE2
045E05F155A5F6AB0FDB34B32033C449C7FBD81D71DC7715FB2DFB38EBFE11
F1269D37892EAF23D4B54D5902B392088A2273E5A679E4E0B7A903D39F9F3F
6A6FD8A9352FB57C42F8456EA976D996EB4E4C2A4DD4B3C3D9643DD780DEC7
AFCC70967D83C13964D8D8F2A7F69ECDF5BF6F27F79E9E7795D4C547EB54F5
F2F23F407E187C43F0EFC48F0BDB496B2C73413C436EDC11B48E00AFCB2FDB
3FF6509FC0D7F75F13BC056A5B44B86696FEDA319FB3331C99500FF9667AB0
FE13C8F973B7C1FF0067AF8F1AE7C1AF138D1B5B6963D2DA5F2E58A4055ADA
407072A79033F7876FCEBF76BC2BE26F0F7C4CF0C08E431DD5BDDC5820E195
9587208EE08AE6C5E1F15C398EF6F47DEA52FB9AFF00347C87C4ACCFE60779
1D38ABFA6B5949791C7A9BBC56EC70CE80165CF438EE077AFB37F6BCFD96AF
7E0EEB7378BFC2B6E64F09DFC9CAA027EC523FF09FFA66C7EE9EC7E53DB3F1
C681A06B3E29D62DB41D02D5EF2FAED82471A0FCC93D028EA589000E49C57E
E195E69471B42388A12BC5FE1E4FCD183563DD749F82FA56B3B1ADFC466157
008CDB09320F7044ABDABD1AC3F654B2BD008F196CCFFD3867FF006E0575F6
DF074F80BC15676EDAA9BAD7A22649D437EE17763F7511C0384FEF1FBC49E0
0C62B68DE32BDB3710DC120AF1E86B8B1B2C541FBB3FC17F91AC5459D6681F
B1FF00C3BF3449ADF88350BD4073B21114008F424AC871EB8C1F4C57D25E0A
F81DF037C1AC93E9FE19B7BCB94FF96D7F9BC6CFA8594B229FF754578B697E
3CC85DCF5DADA78E0103127EB5F338B9E2A7A4A6FF002FC8DA2A3D8FAF21F1
3430C490C24247180AAAB80140E0000740052BF8A971F7FF005AF96E3F1A82
3024A94F8B9DFEEB578CF2E7D8D398FA2EE3C54307E7AE6AF7C54307E7C7E3
5E24FE219E5EAF8AACDA996E5989AD2381484E47A3DE788A49890A6B9D9F51
2C492D935C93EA4A3F8AB32E356001F9ABAA1872798E9AEB51001F9AB90D4F
575453F35606A1AEAA2925B15E49E24F15C8EAD0DAB7278DDFE15DD470C4B6
7AFF00826F74CD7FC5534372DE6FF67462611E32ACDBB033FEEF5C7AE2BE58
8A37D3750B9D3A4397B595E23F5462A7F957A9FC1CBA783C41A8DC339DBF66
C367A64B82093EBC1AF14F89BE2BD1B4CF1B6AAFA2CA97EB3C9E69313828B2
3806405864677E738CFBF35FAC7056369E11373764FF0043E5389B2F9E2211
54D5DA3D1ED2E3815A12EBDA5580FF004FBD86D88FF9E922A7F335F24EA3E3
1F106A394372D0447F821CA0FC48E4FE26B99C124927935F715BC458D3D28D
3BF9B76FC0F92A5E1FB9EB5AA5BC92BFE27DB11F8FFC223AEAF6FF00F7D8AD
1B7F883E1276D8BACDA83FED4AAA3F3240AF85B68A368AC61E28E297FCBA8F
E2693F0D30AFFE5E4BF03F446D75BB4BC5DD67709381DE370C3F4CD596BC26
BF39D0BC6C1E362AC3A11C11F8D759A6F8FBC61A511F66D526651FC329F357
E987CE3F0C57B584F1569BD2BD16BD1DFF00076FCCF2315E184D6B46AA7EAA
DF8ABFE47DC8D724F7AF9FFE26DAFD9B5E835241F25E47863FEDC7C1FF00C7
4AD61E8FF1B6E9711EBD62241FF3D2DCED6FFBE18907FEFA15D46B9ADE87E3
8D0644D1EE564BBB7FDF244DF2CBF283B86D3C9C827A646714F8AB33C1E739
6D4A342779AF7927A3BAF27BDD5D69DCC320CA7179563A156BC3DC7A36B556
7E7D35B3D4E26D67E00ADC867E95C2D9DCF0066BA082E33C57F2CE2B0DA9FB
9267591CC0D5B49B15CE473D5E4B8F7AF22A50344CDE59AA5135622CC2A513
7BD733A23363CEAFA2BF675B948B5EBF95BA81081F425F3FCABE5FF34FAD7A
E7C1CD6869BE239232D813203D7BA1FF00EBD69428FBE869EA7EC7F83B5551
146735F4478735DF2191C374C57C29E0CF122B451FCDD857D01A378846D5F9
ABC8C6E15DCEE848FBB744F1A5B490AACED823BD747278AEC1533E6E6BE31B
2F1295030FFAD6CFFC24CC5705FF005AC69637134A3C909B4852A506EED1F4
0EB5E324646481B02BC9755D6CC85896AE1EE7C439CFCD5CAEA1AF8C1F9AB9
151949DE5B9A2B256469EB5AB0DADF3578578A7551B1FE6AD8D6BC4030DF35
78478ABC42BB5CEEAF63078677329CCF1EF88DA80955E3CF5AF9D629D22DDB
542924E70319E7BD7A3F8AF54FB4BC9213C0CD7879BF059886EE6BE8EB61EF
4944E294B53AC92E239A368A501D1C156079041E0835F1B78CBC3D3783F5C7
B619363704BDB3F5F933F749FEF2F43EBC1EF5F4F7DBBDFF005AC3D7ACB4EF
10E9F269BA926F8DF9047DE46ECCA7B11FFD63C55E5727879EDEEBDCCE7A9F
34417BC75AD14BC06B1B5FF0EEABE17B9293FEFAD89F92651F291D81F43EC7
F0CF5ACA8EF89AFAEF63192E68EC6173B4FB50F5A6B5D0F5AE616F09A5FB41
359FD587CC6E49783D6A8C93B37D2A8198524665B99560B7469247E02A8C93
5A46892D92B39621101666380072493D857D03E0AD34F8734CDAFC5D5C90F3
11DBD17FE03FCF35C8F85BC2834D29A8EA3F35DF3B501CAC79FE6DFA7F3AEF
7CC35E950A3CBAB25B3A36D524C7DEAA72DF3BF53591E61A6992BA492DBCC4
F5355E4D92294900653D41E41A84BD30BD00539349D25CEE6B48B3ECA07F2A
7C56D696DFF1EF0A45FEEA81FCAA467A859E9590EE48CF555DE919EA0CE4D3
11F77FFC13834A9351FDACFC2F769BB1A5DB6A772DB57230D672C1F31EC332
8E7D703BD7F4B55F84DFF0498F03C97DF113C69F11654CC1A4699169B19600
8F36FA51292B9E772A5BE323B3E0F5AFDD9A00FCD0FF00829EE97237C27F09
F892340DFD9DAE2C0ED8F9952E6DE539CF652D1283EFB6BF2CBC31AC031A7C
D5FBEDFB55FC309FE2FF00C02F17F82F4F88CBA9BDA7DAEC155773B5E59B09
E245EE0CA53CBC8ECE6BF9A1F0A788B6858DCE194E083C118AFC73C48C99D4
A8AAA5BAFC8DA9C8FB374BD5F007CD5DDD8788248C0DAF5F33695AF82ABF37
EB5DBDA6BA303E6AFE7CC7E4AD3D8E8523E801E279CAE0C87F3ACFB8D6DA4F
BCD5E4E9AE0C7DEA64BAE0C7DEAF2164EEFB0EE77777AA8C1F9AB8CD4F5518
3F3573977AE0C1F9AB89D535D001F9ABDAC064CEEB41390BE21D586D6F9AB2
7E05E9F2F8ABF68AF873A5C7189D5FC43A6C922302434305C24B2E40078F2D
1BDBD481CD79DEB9AD34C5955ABEDCFF0082687C389BC5FF001D2F3E21CE87
EC1E09B37656C641BCBF568235CF4E22F35BD4103D6BF61E10C9DFB782B754
63367EFBD14515FD0473851451400514514005145140051451400514514005
1451400514514005145140051451401FFFD3FDFCA28A2800A28A2800AE3BE2
078274AF88BE0DD57C17ACF16DAA4263DE065A29010D1C8A3B9470180EE460
F15D8D1401FCED78D7C21ADF813C4FA8F84BC450182FF4D94C6E3B30EAAEA7
BABA90CA7B820D7207DEBF6C3F696FD9E2CBE31E86357D0923B6F1669A98B7
95BE517310C9FB3C8DD0649CA31FBA78E01247E33EBBA1EAFE1BD5AEB42D7A
D24B1D42C9CC73432AED7461D88F7EA0F423047158CA0062139A6938A7918A
8F6D4728EE264D62EA1E1FD2B53CB5CC00487F8D7E56FC48EBF8D6DEDA50B4
D211E6777F0F55896B2BAE3B2C8BFD47F8561CDE02D7109F2D125C7F75C0FF
00D0B15ED816AD4313C8EB1C6A5998800019249E800AD14981E0D17803C4F3
C8B0C3626477215555D0924F0001BB926BDA6C3F64EF8A97371F67BC4B2B02
02EE32DC070A48C95FDD07E57A1ED9EF8E6BEDBF03FC14D43E1FD9E9FE29F1
AC26DB5DD417CDB2B2700BDAC1FF003DE51CED958F11AF54C1270F8DBEC5A6
5A35CC8B120E4D6752AB5A194EA5B447C4FE18FD8D74F8D84DE2CD764B81DE
2B38C4633FF5D24DF91FF001F5F4FC6FF885E0FD266F13EAEDA229B345BB9C
468791B048DB41C7438C648FCABFAF8F0AFC2B4BBB25BABB21430E335FC9C7
8D631178AF5A88744BDB85FCA4614E8CA4F5614E4DEACE07E1E78F3C71F087
5A4D434CDF25A6ECCB129CC6C3D463EE9FAFE35FB3BF023F68BF0E7C48D1E2
5FB42997015E3638656EE0835F8D92568681E21D57C2DAA47ABE8B31B7B88C
8391D180ECDEA2BE3F8BB8270F99D3724AD516CFBFA9F4194675530D2B6F13
F527F69CFD90B41F8B567378CFC0C23D3FC4E8BB8903115DE070B2E3A36060
38E4743918C7C55FB3C7C67F18FC10F1A37C34F1EDB4F691C32797B2507740
7D0F5CC67AAB0E31D0E0D7DA9FB3FF00ED3361E27862D1F5A90417C8006463
D7DC7A8AF00FDBB3E2DF80E2D7F46F0E784EDE1BCF115B1F3B51B9423F730B
03B2D895EAEC4EF20FDC01783B8E3F3CE1A963956790E634DCE1AD9F58DBAD
FB7F4BB1ED67586C355A5F5CA32B3EDDFF00E09FA8513F863E2BF84A5D3751
8A2BDB4BF84A491B80CAEAE304107A835F9D5E2AF855A5FECE134F6BE1FD34
43677ECD8BD39795C13911BC879017B28C038CE33935E7DFB327ED24342BEB
7D0F52B93F629982C65CF31B1FE06F6F4AFD7AB1F02685F1EB41FEC8BFB617
7A7DC80256F6EBF2B0E437A10723AD7252C56278671AD54D68CB7F35DD79FF
00C31F256E647E20F897C6F2DD96FDE673EF5E4375E209639CC830E3B83DEB
FA50BEFD8B3E0BD8E8C6D4782F4D640B82EF6E924BD3FE7A302FF8E6BF2F3F
69BFD8B345F0FDADCF88BE192B5A4B00677B166678A40392232C4B2B7A0C95
3D38EB5F6B94F8A797632B7D5EA4250BEDCCB4F9F626546513E16D27C516F3
1558E4DAFF00DC6E0FE1EB5DF58F8807019B06BE616574728E0AB29C107820
8AE9347D6AE203E44D21743D371CE3DABEE31397AB73449533EA3B6D750E3E
6ADA875A5C7DEAF21D2F46F13EA009D3F4CBC98AE33B2DE47C67A670A715DF
DA7C30F8BD3A33C1E15D5195064EEB39D491EC19467F0CD7CED6AB423F14D2
F566B76758BAD0C7DEA56D7001F7ABE7E9BC636D6B23433DD2C72212ACA564
CA91C1046DC822A8BF8EEC3BDC39FF0076327FF422B5A7D5E23E63E809FC40
8A397AE2B56F1EE9968C6292E5778FE107247D71D2BC5359F1A457566D069F
E7095CE0BC9B570BDF6ED2793EF5F51F82BF679F871AE7ECF577F17351F164
316A8ACEB1D8E7126E4241041E7391FD7A53F6715AD84D9E3B75E215D46332
C73AC89D3E4391F8D715AA6A505AC2F7776E1224EA4FF21EF5E777172FA6DD
34F68FF71BF0600F43F5AE335DD6AE75ABB32C994850911C7D947A9F735E86
0F09CEEFD0894AC6CEA9E37D5EEE2B9B0B099ED2CAEB689114E1A4099C0623
9C73F7471EB9AE342FAD2814E009AF752495918894B8356A2B6793A0AD6874
79A4C7CA68039FC1A306BAB3A0CC06769AA13E97247D568030A8AB5240C879
15588C500464628059583A12ACA7208EA08A7D308C55C581BFA6DF315D921C
B2FEB5D65B5D820735E6A8ED1B075EA2BA0B5BDE0303C57918DC1DF5469191
E8515CF1D6AFA5C035C6417990306B4E3BA15E0D5C21B291D52CFEF5289C81
5CE25CFBD4C2E4E2B8A5852948E83CFAD3D17587D2754B7BE53C46DCFF00BA
7835C80BAA43739E9531C3D9DD0EE7E8F7803C7092C117EF3B0EF5F4BE87E2
D0557E7AFC8FF05F8DE7D2655B795CED5E9F4F4FC2BEB0F0C7C438E58D0897
3F8D5E2B01CCB9A26B0AA7E835978A5481F3D6CAF8A063EFFEB5F1DE9FE384
6507CCFD6BA18FC6A847FACAF16597BBEC6EAA1F4CCFE27183F3D72BA8F8A5
403F3D7844FE345C1FDE5727A8F8D570407C9ABA79789D43D575CF150DADF3
D7867887C44D74CCAADC5739A9788E6B9272F815C36A7ACAA29C1C93D3D4D7
AF87C1D8C65333BC55A94A2D2616EA6470A490A3381DCFD0578DA5FBD7D5DA
3F83A7B0D25EF7564C5EDE0C943D634ECA7DFBB7E5DABC2B57F87BAADD6B13
C3E1F8D6538327925D51B19E76EE201C7A673E83D33CBB39A15EB4A82E9B3E
FDFF00E011569B4AE7182F5CD2FDA58F7AB937843C616B2F932E8D765B381B
216704E71C150456B5AFC38F89579218ACFC27AB4EE06E2B1D84EC40E99C04
3C735EF4E8C63BB4733A896ECE6E468E6468E650E8C305586411EE2BCFF54F
87DA35DB34B60ED64E7B2FCC9FF7C9E7F238AFAB348FD99FE37EAD119C786A
4B2857AB5D491C4DD71FEAF7190FE087F967BCD03F67AB5D26E15FC6B34971
32609B68D5A18FE8CC70EC3E9B6AF08D49BF6534FD190AAC65B3B9F0D785FE
01FC40F195EB5AF86521B98E3FF5933B347147DC6E62A464F60093ED8C9A4F
1AFC09F883E00BBB5B4F102DAAFDB10BC724529743B4E186768391C1231D08
AFD6FD2934FD22CE3D3F4C812D6DA21858E350AA3F01FA9EF5C4FC57F08C7E
3EF0A4B65001FDA3684CD6AC4E32E07299F471C7A6707B57B14DBBAE61B3F2
A2CBC02C486D46F38EEB10FF00D99BFC2BBDD3349D37494DB630842782E797
3F56EBFD2A57F32091A1994A3A12ACAC30411C1041E8452892BBD412D88B9A
3E652EFAA024A5F32A845DDF4864AA7E652192802D1929864AAA64A8CC9401
65A4A85A4AAED2546589A0098B134F5EB508ED5F69FEC3DFB38CFF00B417C5
DB75D62D8BF847C3463BDD59CAE524193E4DAFD6765208FEE2B9EA06403F6A
7F604F84527C25FD9CB441A8C060D63C54CDAD5E2B8C3AFDA55440872011B6
058C953D1CB7A9AFB5291555142200AAA3000E000296800AFE743FE0A1BFB3
7EA3F03FE264DF16BC2364C3C11E30B86965F294F9761A949969626ECA931C
C91740096400055CFF0045F5CA78E7C0FE15F893E12D53C0BE37D3A3D5B43D
661305D5B4A0ED75C86041182ACAC0323290CAC032904035C398606188A6E1
21A67F261A1F8B91D54EFAF45B2F13A903E7AEFF00F6B2FD82FE25FECD97B7
9E31F03ADC78A3E1E025FED6837DD69E9804ADEC6806141C81328D840F9B61
201F88EC3C59380031E6BF25CD7846D27A1AA99F5CC7E2518FBF492789463E
FD7CDB0F8AA561D6AD7FC241349FC78AF9A7C2F67AA2B98F6BBCF13000FCF5
C65FEBAF3E406C0AE05F56CF2CF9FC6B7BC17E15F1AFC4FF00125B783FE1EE
8F73AEEB176711DBDB26E20776763858D07F13B95551C92057A783C8926945
5D89C89B4EB1D6BC55ADD8F867C376926A3AAEA9325BDADBC437492CB21DAA
AA3DC9EFC0EA78AFE9E3F659F80D61FB3C7C21D33C100A4FAC5C137BAB5CA6
489AFA651BC293FC11A858D3819550C40666AF0CFD8D7F622D1FF679813C75
E3499357F883790B46D246C5AD34F8A4FBD15BE4296761C49291C8F95005DC
5FEFFAFD5787F24FAB479E7F13FC0CE52B8514515F4A405145140051451400
51451400514514005145140051451400514514005145140051451401FFD4FD
FCA28A2800A28A2800A28A2800AF19F8B9F023C03F192C045E25B5306A50AE
DB7D42DF097310EA14920874C93F23023938DA4E6BD9A8A00FC67F891FB197
C56F064D2DCF87201E2AD3172564B318B80BDB7DB925B3EC85FEB5F2A6A5A3
EA7A35DBD86AF67358DD47F7A29E368A45FAAB00457F48759FA8E93A56B10F
D9B57B286FA1FEE4F1ACABF9302297281FCDE79753456D2CD22C30A1777215
5546492780001D49AFE8224F83FF00092691A597C13A1BBB92CCCDA6DB1249
E4924C7C935D4691E16F0CF87F2741D22CF4DC80BFE8D6F1C3C0E83E4038A5
CA07E2B7803F65AF8C7E3D9A3687449347B17EB75A9036C807A8461E6B7B6D
423DC57E92FC18FD95BC09F09E5875BBB3FDBFE218BE65BB9D02C70B7FD308
B2C148FEF12CDE8474AFA868A6901F9FDF15B5A9357F88DAC3BB652D2416A8
3D0423691FF7D6E3F8D53F0CDD4505E46EFD0115CF78E84969F107C470CBC3
1D42E5FF000924675EBEC4566DB5E34641538C571545A9CB3DCFBEF42F1A68
EBA5C4B2C814A2D7F24FF1FF00448FC39F1B3C79A2C09B2DED75CD45611800
7906E1CC44019032854E3B57EFDA788AE91360738AFC6AFDB6BC313695F192
7F1185FF0047F11DB437208181E6C2A20917EBF22B1FF7BD735AD09BD99A53
93D8F8B641551EAF482A9B8AE9360B7BBBBB19D6EAC667B79939578D8AB29F
5045799EB3617905CC97571235C79CE59A562599998E497279C93DFBD7A230
A8248D24529200CADC107A1A9E457BDB51F33D887E0A69117883E2B785F44B
A6C5ADDDFC2B3A9380F12B6F743FEF2823FC3AD7F5CBF01AFF00C3DA558416
10A456B0C71AAC688022281D801802BF911F0F5CCFE06F17695E2DB10CF169
D731CCCA3960A0FCC07FC0722BF75BE177C7AD3F51D1ED6EADAED596445604
3750457E37E29603113AB42BD2D54355DAF7BEABEE37A3248FD8FF0018F8B3
47B5D35D44AACECA40008EF5F9D9F18B5DB49ACAE4311C835CAEABF19A07B7
25AE33C7AD7C9FF133E2A8D4639618A5F94E72735F9DFD5B1998E3BEB7895E
F3B6CACB4347256B23B4F85F71F0820F0DCCDAFF008674893514BBB8592792
CE032CC0B97567629962036DC9E78F5E6BD623F8AFF0F747DFFD9B6D696DBB
19F2A344CE3A6768AFC64D6FE29EBABA9DD45A3C88B6BE61DA48249EC4F5EF
DAB9C97E21F8BA7E0DF15FA015FA454F0E6AD797B4A93DF5B36CC7DA1FB5D7
1FB42F8721CF96EBC7A0AE76EFF697D2A3CF96D5F8C7278AFC473FFACD4253
F438FE559579AE5F47199AEEFE4083AEE91B1F977ADA1E16D15F1483DA9DAF
ED05E2DD27C43F183C4BADE8EAA905E4D1C8C13A79C614F38FD4C9B8B7B935
E2FF00DA65BEE8CFD2B78DCA6DDE88A9BB927033CD645D6A2A9900F35F6F85
C3C29538D182BF2A4BEE06556D4261C95207D0D74365E21D420D34D80B9748
1D8B18C13B727BE2B374BD0757F104C2386091D4FF000C685E461ECA013F8D
6EF8C3E1EEB7E0ED1ADF59D52092C63BC97CA8629D944AE40DCCDB321C2A8C
6491D48AEC95183B45EE2B9CC5EDE8604035CCFDE627D682ECFF0078D28E95
DB468A82224EE380C9C56C5869EF70E001552D21323815ECFE0CF0E9BC9906
DCE6B4246F873C1535E95C479CFB57B968DF09AE278C1F24FE55F41FC2EF86
2B79E51316738ED5F73F85BE0C44D6CA4C1DBD2803F2BAF3E12CD0C64984FE
55E55E21F023DA06CC7D3DABF6AFC4BF0960B6B663E48181E95F167C4AF044
569E6E1318CD007E5FEB1A39B766F97A57173C45188AFA3BC69A42C124802E
3AD7836A50EC73401CF118A6B74A91C60D467A5340329D1CAD13E4743D4536
9AD54D5F4606DC17191B90F15A515E11D6B910EF19DC8706BA9D074BD57C44
6E22D22D9AEA7B58FCE78E3F99CC60852CABD5B048C8192339E9923CEC4505
15CCF62E2CD44BC1EB5656EFDEB9A0EC8C558608E08356126AE1961D17CC74
22ECD2FDA8D62ACD5289AB27410EE6A0B9704329C115DA687E2B9ED1D57CCF
2DBFF1D3FE15E75E68A3CD15508B885CFA934AF1D5D050A5F9FAD7656FE35B
9703E7FD6BE34B7D5EEED31E549F28EC7915D5E8DE26D6F52BFB7D274AB19B
50BEBA758E182DD5A596591B80A91A82CC4F6032689508357D8A533EAEFF00
849E697EF4B50BEBCB8C97AFA47E12FF00C13EBE3B78E3C3CBE24F1E5E59FC
3F826C18ADAFC34F7CC8573B9E088E2319206D7712039DC83033E9EDFB0E78
13C287778A7C6379E2091003B2D614B188B770DB9A772BDB8653DF8E95F238
FE2DCAF0EDC6556ED748A6FF0015A7E274C28CDF43E118F53BDD56F22D334A
85EEEEE73B638A2059D8FB01E9D4FA7535F45780BE12BE8E63F1078BC2CBA8
A9DD0DB021A380F6663D19C76C70BDB27047D0D61E17F08782AD9AC3C27A6C
56111C066505A47C74DF2365DBFE044D62EA3723079AFCFB3AE38A98A4E8E1
A3C907BBEAFF00CBFAD4EAA78751D64707AE0055BF1AF9F7C4D3CFA65E45AA
DA1C4D6AE1C7BE3AA9F62320D7BAEB57030D5E0FE2E915A2907D69640DC669
A155D8FAF7E05F86D7C5BACC1E2F9909D2B4FDB2C458712CEC0320FF008002
1CFA1DBEB5F6749719EF5F327ECB1AE0D4BE0CE956E461F4D9AE6D98FAE252
EBDBB23A8FC2BE8369FDEBEAF35AF2AF59CA5D3447C363A4E555DFA171E6AC
9BD8AD6F13CABB85274FEEBA861F91CD0F3553926AC28D269DD1CF1472F7BE
0DF0BDC12CD62B193DE3664FD01C7E95C0EBBE15D22C2E2316C1D5240782D9
E457AB4B35711E286DD0C530FE06C7E0C3FF00AD5F5F93E26AFB551949B5EA
7A784AB2E749BD0F877E3CFC0D7BC497C6DE0C85A5B851BAF6D546E6900FF9
6A807561FC4A3AF51CE73F11F9841C1AFD958E6EE2BE5CF8BBFB3C5AF89A59
FC49E08096BA9C84BCD6A4858A763C9643D11CF7CFCA4F270724FDAD1C45BD
D99EAB47C2025A779BEF52EADA3EA9A1DF4BA66AF6D259DD4270F1CAA5587E
07B7A1E86B37E6AEF24BBE6D37CCAA996A5A00B064A66F278A8A9E0628017A
D4A0535457D63FB387EC7BF16BF68ED421B9D02D0E91E1759365CEB776845B
2053875857833C839F950E01E1DD01CD007937C1DF83DE39F8E5E39B2F00F8
06C8DD5FDD1DD248D9105AC0A4079E7700EC8D33C9E49242A82C403FD4AFC0
3F81FE12FD9F7E1BE9FF000F3C26A64587335E5D38025BCBB900F326930075
C0551FC28AAB938CD53F809FB3C7C38FD9D7C24BE17F01D99F3E70AD7BA84E
15AEEF645CE1A57000C2EE211000AA3A0C924FBA5001451450014514500359
55D4A380CAC3041E4106BE13F8D1FF0004E8FD9A3E315C4DABC7A349E0ED6A
77323DE686CB6C2463C9F32DD95E0393CB1545724F2D5F775151529466AD25
703F063C47FF000479F1CDA5DFFC519F12AC2FAD5893FF00130B196D64419E
07EE9E70C40EA7E5CFA0E95C7E97FF000490FDA1E498AEB1E2EF0C5AC38186
826BD9DB39E7E56B58874FF6ABFA17A2B82594507BC47CC7E417C3BFF8246F
80F4B952EFE28F8E2FFC425486FB369D0269D0F4E51DDDAE24719EEBE59AFD
33F863F077E187C1AD13FE11EF863E1CB4F0FD9363CCF21499662BC069A672
D2CAC071BA4763EF5E954574D0C1D2A7F046C170A28A2BA441451450014514
50014514500145145001451450014514500145145001451450014514500145
145007FFD5FDFCA28A2800A28A2800A28A2800A28A2800A28A2800A28A2800
A28A2803E0FF00DA4BC333687E328BC5112FFA26B68A188E8B710A8523DB28
148F53BBD2BC1E2BDE3AD7E9D78E7C19A578F7C3775E1CD586D49B0D14A002
D0CABF7245F71DFA64123A135F989E33F0A7887E1DEB92685E2184A3024C33
283E54F1F678DBB8F51D41E0F358CE9DF5339C4B82F063AD7CA5FB5DF814F8
D7E1A9D72C937DFF00861DEEC7A9B6618B851F4015CFB257BE8D4063AD4535
CA5C44F0CA8248E405595864303C1041EA08ACE31B3B99A563F04A55E6A8BA
D7BF7C76F85D71F0D7C653DB5BC67FB1F502D3D8BF2408C9E6224E7E68C9DB
D72461BBD7844895D47419EC2A123B55C65A8196802B11D8D6E787BC61E20F
0664E8F23BDAE72600795F5D9FE1F956415A66DACABD08548F24D5D0D33D71
7F68A492DFFD29AE0B0E36639FE78FD6B82D6BE305D7891DAC14359DBC9C6E
63F337B1C74FD7EB5C06ABA1ADDE67B6C24DDC740DF5F7F7AE312DA77B85B5
D844ACC1769EB935E750C970F4A5CF18957B9E94F736B09DB2CA887D19803F
AD5846575DC8C181EE3915ECBAC7EC79F10740F85307C55D55A38AC6E57722
13F31079AF91FCC9620F0AB9556E1803C1C7AD7A971729DB6A1E21B7B5CC56
B89A4F5FE11FE3F8571975797179279B72E5CF6F41F41DAA08D2499C450A97
76E800C935DAE97E13DDB66D4CFBF96A7F99FF000FCE9D83628E9CFA86ACAB
04084EC00339FBBF89F5AF41D13C35045346587DA6E5C8009E809E981FD6AE
DB5BC70A2C50A0445E800C015FAC9FF04F3FD8B749F8D16DAC7C4EF8A9673F
FC22D1452D8E9688CD0B5C5DB82925CA38EAB6DD1782A653CE7CB6539C6946
3B215CF9E7C21243A168967A542C31026091D0B13963F8926BE32FDA27C60F
E26F1D9D3637DD6BA2C620400F1E6BE1E53F5CE14FFBB5FA51FB467ECC7F12
3F66CBFB9D4B508DF57F076E26DF57853F7617B25C28CF93263D7E56FE163C
81F8CF7F7B36A7A85CEA570732DDCAF2B9FF006A462C7F535E7E1F09CB51CE
5B95296857A91474A8EA54EA2BD0641D268F087957EB5F5AFC30D36379A2DC
3B8AF9474470B2A9F7AFAAFE1C6A71C12C449F4A803F54BE0D6916BB20240E
D5FA05E1CB3B18ECD4E074AFCC2F863E3486D238BE7C631DEBEB5D37E28C11
5A81E68E9EB401EB3E3B6B24B69318E86BF37BE2F4F6F99B6E3BD7BFF8DFE2
9452C32012F6F5AF827E2678DD6E8CB87CE73DE803E58F88124666931EF5F3
6EAC4798D5EB9E2FD604F23E1B3D6BC52FE6DEE68031E4EB511E94F639351B
7A53403691BA52D35AAC061E95D6780BC4D37843C5FA6EBD1B1548650B28CF
0D0C9F2B83FF0001248F700F6AE4CF4A8DBA5454A4A717096CF41A76773F41
BC5FE1AF09F89774DA9D8C6D3B7FCB78FE497D8EE5C6EFF816457CFF00AB7C
2958DD9B46BFCA7649C73F8BAF5FFBE4577DE13D725D57C2BA65CCAFB9BC95
463EAD1FC84FE256AF5CDD800F35F0F855568B74D4B63B6567A9E5DE0EF817
F157E21789ADFC1DE02D0E4D7F57B9591E382D5D3256252EEC4B9500003A9C
73C75201F44BEFD8BFF6B5D36716D71F09FC44EE65F2730D8493AEEE79DF10
6509C7DFCECE9CF22BF7E3FE09C5FB36DCFC36F04CBF187C636860F1278BE1
55B38A41F3DAE96C55D323F85EE180908EC823CE1B701FA695F5B87C339413
9EE7249EBA1FC3CEA1F0AFE26691ADDDF86F54F0E5E59EA7632B433C1347E5
BC522FDE57DD8C11EF5D8687FB3CFC51D6D94BDA41A723747B99D40FCA2F31
87E22BF493F6C5D125F047ED55E32B66188357920D4E13CFCCB750AB39E49E
928907A71C01D0705A2EAC0AAFCD5F9D67DC458AC3D4953A715A3DF57FA9D9
4A8C5ABB3C7BC19FB1D694D2473F8DFC4324EB905A0B14118E0F4F364DC483
EC8A7DEBF42BE0F7873E1B7C1D5F37C07A0DB69B76C9B1EEF6F9974EBDC34F
2167C13C95042FB578CD8EAA303E6AE92DF582070D5F96E719B63B1578D6A8
DAEDB2FB97EA77D2A708EC8FADAFFE2D5FDC4450CE4FE35E57ACF8AE7BD625
DC9CFBD7949D6188FBD54E6D538E5ABE6E38266DCC7477BA8EEC926B8FD42F
860F3546EB54183CD723A86A8307E6AF530B827733948A9ACDF0C3735E1DE2
ABD0CAEB9EB5DA6B5AA8DADCD78CEB77C667639E057DCE4F83B34CE6A923EF
EFD8E2623E1A6AEA4F035A9FF0FF0046B7AFABDA7AF9A3F663D3A5D23E10E9
B2CC0ABEA335C5D60FF759CA21FC5501FC6BDF1AE3DEBE86A61AF26CF8FC52
BD493349E7AA72DC5507B9AA325C57452C299C605B9AE2B075253776D2403A
91C7D474A91E62DDEA1DD5EC61A8F2B4CE9A71B3B9C0A4CD1B147E194E0835
716707BD77BA5FC32F14F8F8DFDC7846C5AF26D361F3A74520165CE005CF05
CF242F5600E32460F94BB4D6D2BC17086392325595860A90704107A115F591
8A9C5491EAC277456F13F84BC31E33B3FB0F8974F8AF5003B5986244CFF71C
6197F03CF7AF973C53FB2B233BDC783B57080F482F41E3E92A0FC814FC6BEA
F5BAF7A7FDA7DE9C79E3B32AC7E70EABF023E28E9458B68CD75183C35BC91C
B9FA286DFF009AD728FF000DFE2046E51BC35A912A4838B4988E3D085C1FC2
BF520DC13D298652D5BAAF2EA1CA7E625A7C2FF88778DB62F0E5F21C81FBC8
1E2EBEEE16BD1344FD9C7E206A2CA75216FA5C6792659448FF0082C5B867EA
C3F3E2BEF70DEB5D4F83FC1DE21F1DEB90787BC3368D757531E480764499C1
9246C615067927E83248155EDA4F61591E59FB3AFEC67E19F16F8D6C34AD53
CCD65613E75E4B22ED8228148DD88C1392DF7543330C9CE38AFE8034FD3EC3
49B0B7D2F4BB78ED2CED2358A186250891C6830AAAA38000E0015E73F097E1
668DF0A7C34347B06FB45EDC9125E5D11869A403000F445E422F6C93D4927D
4AB78A7D44C28A28AA10514514005145140051451400514514005145140051
45140051451400514514005145140051451400514514005145140051451400
514514005145140051451401FFD6FDFCA28A2800A28A2800A28A2800A28A28
00A28A2800A28A2800A28A2800AE7FC4BE16F0FF008BF4C7D23C49631DF5AB
FF000B8E54F4DC8C3E656F7520D74145007C6DE25FD922C2595A7F086B6F6A
A7FE585DA79807D244C103EAA4FBFAF003F64FF88C1B1FDA5A56DCF5F367CE
3E9E47F5AFD08A2958563F21BE27FECBFA578DF4AD53C11E2FBEF2AF2D5DC5
B4F147FEA2750424A371059483CAF1B94F507047E1AFC48F875E27F861E2CB
EF0778B2D4DB5F5931C1C1F2E68C93B2589BF891C720FE0704103FABDF8CFE
10B8E3C67A521731284BC45EBB147CB2FF00C0470DED83D0135F0FFC63F847
E0CF8E7E163A0F8953C8BC832D657F1A833DB487B8CE37237F1A1386F66018
79D2AD3A751F36B1391D49425EF6C7F3BCF1E2ABB257BB7C60F819E39F837A
CB69FE27B42F632B116B7F082D6D703AFCADFC2F8EA8D861E84609F1578ABD
18C9357475A69EA8CE2B4C2B574C74C294C654DBED5E77A9CB2DA6B724EBC3
A32B0FC862BD37CBF6AE6F5FD0A4BE02EAD00F39060A9E370FAFA8A18D33B9
F12FED1BF11FC4FE0FB7F04EA1A83B69B6CA1563CF1815E02D962493D6A496
19EDDB64F1B46DE8C08AD8D2F40BFD4A45250C5067E67618E3D8773FA52486
D9D77856D234D3527D803CA5B2D8E480703F9575D1C7496B6890449044B848
C0007B0AFD12FD927F603F883FB405CDA78B7C5B1CDE16F0106573792A15B9
BF4EA56CA371CAB0E3CF61E58CFCA242ACB4C93CEFF63EFD92BC53FB4D78DD
2231CB61E0DD2A543ABEA43E5DABD7ECF0120869DC74182101DCDFC21BFA95
F0B786341F05F8734DF09785ECE3D3F49D2208EDADA08861638A318503D4F7
24F24E49E4D647C3CF877E0DF855E10D3FC0BE02D323D2745D3136430C7939
27967763967773CB3B1258F24D76B401F31FEDA5ADC3E1FF00D92FE2E6A13E
DDB2786752B51BCED01EF206B7420FAEE906DF538AFE2B874AFEC6FF00E0A2
5FF265BF14FF00EC1D0FFE95435FC720E9532025A7A9A8C74A7038A9606E58
4FE5B8ED5EC3E18D78DB321DDD2BC2A29369CD74165A8B44460D4580FBA3C2
FF0010DAD5131263F1AF548FE2D48B163CEFD6BF3CECBC492C4000F5B3FF00
096CDB71BE901F5E7887E28C93A30F3BF5AF9E7C4FE317BA2F97CE6BCCAF3C
4D2C808DF5C95E6AAF29396A00BDAB6A66663CD72134858E69D34E5CD542DD
CD002138A651D68AB4802987934E27B5329808DD2A26A79EB4C6A680FD0FFD
953E09C5F143E1CDC6B526A52DBA59EA13DA18D235E0AC71CB90C4F3C49E95
F627817F66DF0DF833C5DA7F8A758B7FF849E1D3A4F3469F7C316B2C8BF73C
E54C17553CECCED6C61B2B95395FF04DED1667F811AB5D5C44556E3C4174F1
93FC482DAD5091EDB9587E15F7DCBA3263EED7C363F14E1889F2F73AA2B43D
6EC3F6C2B985C26BFE10658C7DE92D6EC3B7E11BC6A3FF001FAF69F087ED1F
F09BC5EF1DAC7ABFF645EC9D2DF525FB3364F41E6126124FA2C84FE95F0E5E
68A841CAD705AC786E195581406BAE867D553F7B525D246F7FC1537E165C5C
689E18F8EFA343E67F6291A56A6EB83B6D677DF6B21FF6566674EBD655E3A9
AFCA6F0FF8895E35F9ABF4A21F1AF8BBC39E1CD4BC0D7CC35FF07EAD6D2D9D
E68F7CCCF03412A95610BFDF818039468C80AD86C122BF22BC4D657DF0EFC5
775A15D2CAB67E63B59CB2E332DBEEF9092005DE060381D0F4E304F9D9D612
9E2FF7B05AF54694A4E3A33E9CB0D70103E6AE9A0D68103E6AF98F4BF142B2
8F9EBB0B7F11823EFD7E7B89C9DA7B1D91A87BC8D6463AD57975918EB5E403
C4231F7AABCBE2118FBD5C6B297D8AF687A5DDEB4307E6AE3B51D6C60FCD5C
45DF88860FCD5CB5DEAEF313CF15EAE1729EAD19CAA1B3AAEACD292AA78AE6
B4ED3350F14EBB63E1BD2977DDEA332C29E8371E59B1FC2A32CC7B004D63DE
6A0AA0926BECEFD9BBE1ACBA35B1F883AFC263BFBE42967138C18A06EB2107
A349DBD17FDE38FABC0607B2392BD5E5573EBED16CACB40D1AC742D39765AE
9F0C70443FD88D428CFBE07357DAEAB0FED0C693CC63D6BD55833C1F6669BD
CD40642D5503D2EEADA187B14A058DD5A7A2E8FAA788B54B7D1745B76BABCB
A6D91C69D49EE49E8001C92780393537863C2FAFF8CB568F44F0E59BDE5D49
C90BF7514756763C2A8F53F4EA457E8FFC23F839A47C34B1FB54C56F35DB84
DB3DC8CED507931C40E30BEA7196C64E3803D3C2609C9DDEC6F0A773A1F85F
F0F6C3E1C785A1D1ADF125E4B896EE6EF2CC4738FF00657A28F4E7A926B8AF
8ABFB3D7827E2709751D9FD91AE30C8BD8147EF0F41E747C0938EF90DD3E6C
715EF7457D028A4AC8EA47E3F7C45FD9FBE22FC3856BDBFB45D434CDFB56EE
D0F98BCF4DE9C3A13EE319E0135E28C9244C51D4A91D41E0D7EF1DE59DAEA1
6B2D8DF44B3DBCEA51D1C65594F50457C4BA8687E11D266D7CD9C097166933
5ADBADC2ACBC8E1C82C3903076FB11DF9AC2BFBAAE95C691F9EC0D5BB6B5BA
BC7F2ED61795BD114B7F2AFAE61F869AF6A68BA8691E1B925B49798DD20CAB
00719071D322BB9F0F7C0EF1DEACCA2E2D1349B6FEFCEC01C77C46B96CFD40
1EF4946FAD8773E5DF04FC25F11F8C35AB4D1D00B76B960307E6655EACC71C
0007273CFB13C57EADF813E1F785BE1CE8C346F0BD9ADBA360CD29F9A599C0
FBD239E49EB81D067000150F81BE1EE85E04B368B4E5335DCC009AE64FBEF8
E7007455CF403F124F35DE56F08584D8514515620A28A2800A28A2800A28A2
800A28A2800A28A2800A28A2800A28A2800A28A2800A28A2800A28A2800A28
A2800A28A2800A28A2800A28A2800A28A2800A28A2800A28A2803FFFD7FDFC
A28A2800A28A2800A28A2800A28A2800A28A2800A28A2800A28A2800A28A28
00A28A2801ACAAEA51C06561820F20835F3AF8DFE06477B712EABE0E956D65
90966B57E2224F5F2D87DDFF00748C7A102BE8CA2A6504D5989C53DCFCE5F1
5FC34D76E2C2EB42F1778724BDB0B805258E483CF81C75EA0321F51CF1D6BF
3B3E257EC17A76B2D2EA9F0AAEDF4C91B2DF62BB124B6F9EB84980691001D9
83FD40AFE8B6BE4FBD83CBB59EDB03F74CE9C74F9491C57978CACF0F671EA7
0E22A7B2B35D4FE6EFE21FEC9BFB42FC3189AF3C59E05D4A3B00A1FED96F09
BAB508464334B0EF54E39C3ED23B815F3B980FA57F6C3E199FED3E1CD2EE0F
592D6163CE792833CD717E2EF827F07BC7C647F1A782746D6A594EE696EAC2
0965DC4E4912326F0724F20E793EB5EAA77477A3F8CFF24D28879AFEB3EEBF
61AFD936F2E1EE66F86BA6ABB9C91199A35FC11240A3F002ADE9BFB13FECA7
A4C8D2DAFC33D25CB633E7A3DC0E3D04CCE07BE073DE981FC99D8E977BA8DD
4765A7DBC9737131DA91C4A5DD8FA2AA8249FA57DA7F08FF00E09F1FB4CFC5
678AE47865BC2BA5B9E6F35DDD6400EB9580A9B86047DD222DA7FBC0735FD3
8785BC01E04F03C26DFC15E1CD37C3F11054A69F670DA2E09048C44AA31902
BADA00FCE0FD9EFF00E09A7F067E11C969E23F1E1FF84F3C49015915AEA309
A74120C1CC76B96F30A9CFCD2B303C108A7A7E8F2AAA2844015546001C0005
2D140051451401F2B7EDC36E975FB227C5A89D03E3C3D7AE030C8DC89B94FD
41008F435FC5CA9AFEEA3E31F8567F1D7C22F1C7822D73E7788742D4F4E8F1
D775DDAC908EC7BB7A1FA57F0B5344D6F3496EFF007A262A7F038A9900A0E2
9F51039A7038A9024071532C8474A8339A29580BEB72C3A1A97ED6F8EB5999
34BB8D2B0175AE58F7AAED21350EE3494EC038B5368A29D802909C504FA532
800A42714BD2984E68012987AD389C0AFA6BF639F81B2FED11FB45F83FE1AC
B196D2A6BA179AAB00485D3ACFF7D7009FE1322AF94A4F01DD720D35B5C0FD
F6FD91FE105F7C3FFD9C7C0FA1DE5B98AEE6B05BFB85618612DFB1B92ADEE8
24087FDDAF77BCD0E4881CAE2BED1FEC3D3D63112408AAA300018000AE27C4
5E10B792169605C11CE2BF3AC7E06BC24EA545BB3B6325B23E38BDB1DB9045
71BA8598C1E2BDB7C41A59B691971D2BCC351880C8AE1848A68F1AD6B4D475
6E2BE61F8ABF0DF48F18E93369BA945CF2D148BC3C4FD994FF004E84706BEC
4D4E20430AF24F10DA2B2B715E961EA34C868FC81F117833C65E039AE5AEE0
92E74DB67553771A93101212137919D85B0700F7E84D67D9F89E5207CDFAD7
E96D9788EF3E1E78B2D7C5967690EA1145BA3BBB2B84592DEF6D24E26B6991
815649178E41C1C3632057D5DF143F607FD90BE287806DFE34F80EF6EFC036
DABDBA5DC5FD9CCB3DA3BCDCF97F6398FCB206CA18E2923542A4103071ED43
011AF1BADCCF9EC7E2047E23958633FAD4A75991C72F5F40EAFF00B1678F2C
EE193C3FE24D3AFA00C42BDD24D6AE579C12A8B380718E371FA9EF8169FB24
7C6199CA4F73A55B28FE27B894E7E81226E7EB8AE4964524F4894AA9E30DA9
0EA5AABA5DDC5EDCC7656313DCDCCEC1238A252EEEC780AAAB9249EC057D79
E1FF00D8CA42EB278BBC4ECC9DE2B18769FC25949FFD175F51781FE12F803E
1D0F33C33A524574C30D752932DC1E307F78F92A0F70BB41F4AE9A190CAFEF
68899563E63F845FB3A6A16B716FE2CF88F6E032624B7D3980600F5579C8C8
C8EA23E7B6E39CAD7D881ABA01251B617FBE80FE15ED472D8455A272CD393B
B308353C357AE780B52F04E9DA80B7F18E8916A16131E65F984B09F51B586E
5F51D7B8F43F70689E0CF85F15AC17DA3F87F4E963914347298526C83C8219
C31FC7AD73D5C3723D48544FCDED1B40D7BC413791A1E9F3DFBE704431B381
F520607E35EF3E19FD9D75EBA0B75E2DB95D361CF304444B311FEF0CA2FF00
E3DF4AFB60CF1431086DD1628D7A2A80AA3E80715C9EBFAED9693652DF5F4A
228631924F73D801DC9EC29429ABEC68A8A5B9A3F0A746D13C2F1DE681A2DB
25BC785972397723E525D8F2DDBAF4ED577C4FF193C0BE1891EDA4BC37F749
D62B5024C1F42F9083DC6EC8F4AF8E7C49F10356D6669E2B095ECACA4050A2
36D69109CE1C8EC71F77A7D6B825EB5EC41592433EA3D4BF697BB2EC348D0D
1107469E52C4FD5540C7E66B0D3F692F180958CBA6D834673B4059430F4C9F
30E78F615F3DD308C555C0FA7BFE1A52EE6D3AE2DA7D1C4176F132C73452E5
5642B807632F4079FBD5E52971178924D1FC31A23E649D9233B811FBD95B0C
C47A2F73E82BCD1AA30EF1B8963628E841520E0823A107D6981FAA3A6D85B6
95A7DB69966BB20B48D2241FECA0C0CFBF1576BE30F86BF1EEFB4A921D13C6
D235DD89F956F0E5A68876DF804C8BEFF787BF4AFB26DAE6DEF2DE3BBB4916
686650E8E8432B2B0C8208E08228026A28A2800A28A2800A28A2800A28A280
0A28A2800A28A2800A28A2800A28A2800A28A2800A28A2800A28A2800A28A2
800A28A2800A28A2800A28A2800A28A2800A28A2800A28A2800A28A2803FFF
D0FDFCA28A2800A28A2800A28A2800A28A2800A28A2800A28A2800A28A2800
A28A2800A28A2800A28A2800AF99B5E8BCBD535684E7FE3E6661918E198B7F
5AFA66BCBF55855359D494281B8AB0F5E635CFEB9AF9FE23ABC9414EDD4C6B
613DB5A37B1BBF0FA533783B4D63FC08E9D31FEADD93FA57675C47801C9D0E
484E7F73712AF5CF53BFF0FBD5DBD7B385A9CD4A32EE91B72DB40A28A2B700
A28A2800A28A2800A28A2800AFE36BF6FDF82D3FC0DFDA93C61E1E8E2F2F4B
D62E1B58D3885DA86D6FD8C81507A44FBA2CFAA1AFEC96BF2FBFE0A95FB2E2
FC74F828BE3EF0C5B193C65E04266B644197BCB29982CF6D81D5C1C49175F9
94A0199321A8B6EC81B3F95007BD480E6A0069F9ACDA024A76EA606A5A403F
229722A3A28024A4C8A651400EDD4849349450014138A696F4A6D002939A4A
2984E69A4004E6BFA56FF8239FECF12782BE17EB1F1F7C456823D57C70DF64
D30B82248F4AB57F9D867181713AE71839589181C357E177ECB1F00B5CFDA5
7E377877E15E90AF1DA5DCA27D4EE5067EC9A6C241B8989E80ED2113380646
45CF35FDAA786FC39A1F83FC3DA67853C33671E9FA468D6D0D9D9DB443090D
BC08238D17393855000CF3549AD80DAAAF72A1A320F7AB158DAC6A31595A3C
8ED8C038FAD79B9BD584684B9BA974D6A7CE1E3B8E38EE240BEA6BC0F5461B
9ABD63C65ABADC4F2107A935E29A95C862C6BF3C8A3B1B394D458735E61AE0
055ABD03519860D79AEB33021ABB69220F10F16401E3707DEB73E0878F7547
B6BBF867A85E48F65A7BC97F610336523F3881701476F9B6B63D598F1CE723
C4B2028D5E39E18D5DB42F89DA1DEA36D49AE05B49CE015B8FDDF3EC0B03F8
57D1E5B539648CA68FBC8B1635B3A769135EB00AB9CD65DA279928535F457C
3ED0EDEE190B81DABE95981C143E06BA68B798CF4F4AE6755F0F4B679DCB8C
57DF4BE1CB08EC73B474AF9DBC7D636B0B384005007CBB2A18DB06A20D5A1A
A05594E3D6B2775080BC8F5F437C17F885369774BE14D46426CEE98FD9C939
F2E53FC23FD97FFD0BEA6BE7056ABB04AF1BABA3156539041C1047A529C149
59823F46EE75644467770AAA09249C00057CA3E38F17CDE28D4888D88B1B72
444BD37762E7DCF6F41F8D6DEB9E309EFBC1365329226D4C18DCE31FEAF897
1F538FC0D7942B5654A972EA36CB2BD715B163A64D764051D6B3ACE3F36455
AFA1BC01E1A8AF1D378F4AD4479D41E0EBB78F76C3F95636A3E1F9AD33B971
8AFBE6DFC1FA7C76C10A8DD8F4AF0DF1EE876D6BBF601DE901F264D198D883
5518F6ADFD62311CAC07AD738CD4C0631AFA4FE017C4F7D27508FC0FAD4BFE
8378D8B4763FEAA663FEAFFDD90F4F46FA9C7CD0CD5033B290CA4A91C823A8
35407EBA515E79F0B3C5C7C6BE08D3F5999B75DAA982E7FEBB45C31FF810C3
7E35E874005145140051451400514514005145140051451400514514005145
14005145140051451400514514005145140051451400514514005145140051
451400514514005145140051451401FFD1FDFCA28A2800A28A2800A28A2800
A28A2800A28A2800A28A2800A28A2800A28A2800A28A2800A28A64922451B4
B23054404927A00286C07D79CEB853FB7A608C1B7C284E3B302C307DF8AD0B
8D4AE755731C04C56BD063867F73E83DBF3A922D36345C6D15F9E711F11D2A
91787A6AFE7FE47651A2D6ACC6F065F5BE991EA16F7CE20125C798858119CA
283CF4C7CB5E8F14D14E82485D6443D0A9047E62B8AB9D3D369F96B9598DE6
9337DAB4E94C4E0E48FE16F661D0FF009C53CA38B5251A5523A2D055287547
B1D15CDF873C476FAF40CB811DDC38F363F4FF00697D54FE9D0D7495F7D4AA
C671528BBA672B560A28A2B41051451400514514005717F112C9B50F04EAF6
E832561F3703FE989127FECB5DA53258A39A2786550C9202AC0F420F0455D3
9F2C94BB0A4AEAC7F21DFB777ECDD37C2FF1A49F137C296A7FE113F13CECF2
AC6B84B1BF932CF160702397978FA0077260055CFC000D7F5A1F153E1C685A
CC3AEFC35F1A58A6A1A55D068648A51C3C4DF346EA7F8580C32B0E558020E4
57F389FB4BFECD1E29FD9E3C53F679F7EA3E18D45DBFB3752DB80E3AF93363
85994751D187CCBDC2FAB98E095BDB53F85FE071E1ABFF00CBB96E8F9B01CD
2E7151039A76E35E2F29DA4BB8D1BA99B851914AC03F70A375372293228B00
EDC6929370A6EE345807D3775368AAE50149CD0AA58855192780077A4AFD03
FD89FF0067F5F146BD07C59F195B6745D264CE9D048A36DD5D21FF005A41EB
1C2791D9A4C73846070C4E2234A0E721C55DD8FD56FF00826F7C17B7FD9DFC
032F89BC41085F17F8C1219AF03A8DD6B6AB9686D41EA0FCDBA51DDF00FDC0
6BF58ECFC6DA5CF18667DA6BF3BB4FD75E1C6D6C575F6FE2F9D171E61FCEBE
0A58EADED1D452B3675595AC7DB97FE3AD32DE3251F71AF12F1678F1EF3728
7E3D2BC3AE3C593C83973F9D73577ADBCB9CB573D6AD52A3BCDDCA565B1BFA
BEAE666625AB85BDBBCE79AAD75A8673935CD5EDF8C1E69420222D4AEC60F3
5E6DAC5DE77735B3A96A03079AF38D5EFC61B9AECA50136711E24BA1B5B9AF
9CF5FBDFB2EA50DEA939B6916418EB946078F7E2BD7BC47A80C3735E15AA86
BEB85B753869DD501C67EF1C0AF670B13391FA9965288E504D7BAF83FC4A96
3B096C62BE7A0E54F15A76DAA4B0630715F50607D9D71F1113ECBB03F6F5AF
0AF16F89FEDACF86CE6BCD24F104ECBB779AC3B9BE926272734005E4FE6C84
D5207B8A6124F2681D6802C29AB51B5525AB494C0EBAD2EAE1EC62B69242D0
C4CEC8A4F0A5F1BB1F5C0AB21AB2ED1B1020FF003D6AE06A00DED365092A93
5F48F80F5F86CB61661C62BE588A628720F4AEA74FD765B6C61B1480FBE5FC
7768B6B90C3763AD78078E3C551DE17C36739AF227F175C18F6EF3F9D72F7F
AD4B704E5B39A4045AA5D79B2B1AC166A2598B9C9AAECD4EC0399AA076A466
A81DE981F617ECABACB1FEDFF0FBB7CA3C9BA8D7D09CA487FF0040AFB06BE0
BFD96A47FF0084F75140C769D32524678244D0E0FE1935F7A5001451450014
51450014514500145145001451450014514500145145001451450014514500
14514500145145001451450014514500145145001451450014514500145145
00145145007FFFD2FDFCA28A2800A28A2800A28A2800A28A2800A28A2800A2
8A2800A28A2800A28A2800A28A2800AE2BC597E7CC8349438F33F7927FBA0E
147D09CFE55DAD78F788EE76F8AEE159BEEA4607B0DB9FEB5E2710D69430B2
E5EBA1AD15791D8E9912EC18AE95204D9935C5E977E8A00635D27DBE309F7C
57E759362F0F45CDD78DDF43AEA26F61B741555AB87D48AE0D6F5F6A2A410A
78AE2B51BC183CD796BDFAAE515645F43945D61FC3DADDBEAD112123602503
F8A263F38C77E391EE057D2CACAEA1D082AC3208E4106BE3FF00105C870C07
7AFA97C2F23CDE1AD26590EE77B48189F5263526BF51E1C9BF64E2CE2ADB9B
9451457D1988514514005154AF751B1D3A3125ECEB0A9E9B8F27E83A9FC2B9
B93C71A22310A26900FE209C1FCC835CD5B194A9BB4E497CC9734B76763457
2D69E33F0FDDB88CDC7D9DCF4128D83FEFAFBBFAD6F5DDEDBD9DB1BA95B29F
C3B792C4F40BEA4F6AB86269CA3CD192B7A8D496E787FC6FF07AEA1A747E2A
B2005DD9623957A1962278C7AB293C0EE09F402BE2BF18784BC2BF10FC3779
E11F19E9D16ABA4DF2ED961946464746523055D4F2ACA41079041AFB77C4BE
247B32BA9EB8ABF6B5C9B5B2CE56007A492FABFA57875E7C3AF14EB9A6DC78
C34BB2DD04AC5FCA1C4920392D2468072B9EC393FC208AFB1A34E587A2A539
269FF5A775E7F71E2BACAB546A2B6EA7F3A3FB4AFEC4DE32F83F3DDF8ABC0E
93788FC1A3748D22A86BBB15C93B6E11796451FF002D946DEBB82719F86726
BFAC1F35D18AB8208E083D735F137C6BFD873E13FC5496E75DF0EA9F087882
E09769AD1035A4D21EA65B6CAAE4F768CA124966DC6BCDC4E0212F7A8BF91D
B4B10D6933F064114B5F57FC4CFD8AFE3C7C399259E2D1BFE126D353245D69
3BAE0ED1CE5A0C0994E3AFC85473F311CD7CAB776979A7DCC9657D0BDB5C42
76BC72A947561D994E083F5AF2674E5176923B2324F622A29BBA8DC2A0A1D4
526E14AA19D8246A5998E001C924F61400515F5F7C17FD843F6A5F8EB2C727
84BC1175A7696E158EA5AC2B69D67B5BA323CC03CC3FEB8A487DABD0FC6BFB
204DFB3B7C5197C09F15258B5ABEB686DEEA178372D8CF14C9B83A070AEEAA
E1A33B8282C8C369159D5A8A0AEC691E27F00BE01DEFC4DD522D6FC469259F
862DDB2CDCABDD953831C47B2E461DFB741CF23F643419ACB4BB1B6D374E89
2DAD6D6358A28A30151110615540E0003815F35F877598628228610B1C71A8
5555002AA8E00007000F4AF4FD3F59040F9ABE5330AD2AB2D76378248F79B6
D5781F356BC7AB7FB55E336DAC703E6AD78F5818FBD5E54A91A1EA6755E3AD
559754F7AF3D3AC71F7AAACBAC7FB5495203B4B8D5383F357357BA98E7E6AE
6AE357E0FCD5CD5EEAE067E6AD61484D9ABA8EA63924D79BEB3AAF0DCD45A9
EB2307E6AF33D6358CE403CD76D2A4436666BBA8F98C541AC6F03E9EFAEFC4
2D034D41B81BB8E571EA901F35FF00F1D5358F7F77D4935EE5FB35786DAFB5
8D53C63711E61B34FB2C0C47065930D2107D55703E8D5EBE1695E4919C99F5
FB77A889C54EC2A123B57B8663324D464E69F498148045A78EB494F029A01E
B56A31CD57515B1A569D77AA5FDBE9B63199AE6EA458A341D59DCE140FA934
01A696F736F6B6F34D1948EE14B46C7A32862A48FA10453C3D7D69F1A7E17C
7A07C30D025D357CC7F0DA886E59470CB7182F27FDFDE9ECC73D2BE3F12500
5EDF4F1291D0D52DF4BE65005D33B5466426AAEFF7A697A00B05EA32F5097A
88BD004CCF55D9B34C6929A393401F5B7ECA5A5BC9ADEB9AD153B20B68EDC3
638CCCFBC8073FF4CC678FCBBFDB95E21FB3F784DFC31F0F6DA7BA42975ABB
1BB707A84600443FEF801BEAC6BDBE800A28A2800A28A2800A28A2800A28A2
800A28A2800A28A2800A28A2800A28A2800A28A2800A28A2800A28A2800A28
A2800A28A2800A28A2800A28A2800A28A2800A28A2800A28A2803FFFD3FDFC
A28A2800A28A2800A28A2800A28A2800A28A2800A28A2800A28A2800A28A28
00A28A2800AF0AF88AADA6F892DEF8F115EC4067FDB8CE08FC8AD7BAD72FE2
FF000DC5E29D165D38B08E753E64121FE0917A67D8F43EC7D6B8F1F85F6D49
C0A84ACEE793596AA368F9AB5C6ADF2FDEAF1496FAFF0043BE934AD5E36B7B
A80E195BF420F707A835A0BE215DBF7BF5AFCDF1393494B63B1543D36E3541
83CD725A8EAA307E6AE527D7770E1B358771A8349CB1E2B6C2E5567AA14A65
9BCB9FB44C14B050CC064F4193D6BE9BD2356D145B430E9D726158955542B7
18518191D0FE55C3FC2CF034AD20F13EB706D5C116D148BD73C19083EDC2E7
EBE95EA17DE03F0BDF31916CC5A4A7F8EDC984E7D70B8527EA0D7D155C8AA5
4845C27CAD192AA93D51A316A1740655A3B95FFBE5BF3191FA534789B4E599
A0BADF6EEA71965CA93EC573FAE2B8D97C0BAF581DDA26AFE681D23BA5C1FC
644FFE22B8EB93E39D2199B54D2E59E3C925E10275F7385CB01F502B183CD3
0FBFBEBFAF9986227A27047BEDBDD5B5D27996B2ACABEA8C187E95CEF89BC4
69A244B04003DE4E0EC53D147F79BDBD077FCEBC76CFC55A45C4993FB89978
2509461FD4544D78DA85FCB7524AD364ED567393B5781457E24A9ECE51E4E5
91C72C43B6C6DC56F75A8CE6E2E19A799FAB31C9FF00F57B56949A2CC91EE2
9C574DE19B684C4243826BA8BD58442D9C74AF948E12A55A72C4396C2A786E
68F3367875EDA28C822A969DE22B9F0F4B9753730206F2D18922263FC4A3A7
D456FEAEE9E7305F5AE135420AB5197D595EC73D393B9DAE81E0D9FC69A847
E28F122E2C49DD1C47EF4F8E85C7F73D077FA75F785554508802AA8C003800
0AF21F841ACB5DE9979A2CAD96B0903479ED1CB9381F4607F315EC15FAB60E
7FB984536D25A5CF429C125A23CE3C5DF0B7C29E2F2F737307D92F9F9FB441
85627FDB1D1BF119F715F3F6B5F00BC5962CCFA3CD06A517F08DDE4C87EAAF
F2FF00E3F5F64515DB4F1128ECC72A69EE7E7F4DF0CBC7D6F218DF44B8247F
71438FCD491FAD72FE20FD9BA5F8891ADB78BBC096DAD04528A750B681CA06
FEE34DCAFD548C1F7AFD29A2BA1E3E76B19AC3C4FC9C1FF04A6F811E259DAE
FC43E1CB7D14C8727EC177741C7392046B288541F60703B702BA8B1FF8244F
EC6D6B3492DC697AC5E2C814049754902A11D4AF9611B9EFB89F6C57E9DD15
C939B93B9AC6363F3E341FF825B7EC47A1C82793C0526A72AB6E56BCD53509
00E318D8B3A230EFF329E6BEA6F017ECF7F02BE17491DC7C3CF87FA1787EE6
3000B9B3D3E08EE4E092374E13CD6C1271B98E3B57B0D1505057E79FFC143F
F67DBBF8ABF0C23F887E12B569FC53E07596E0471FDFBAD3986EB888281977
4DA248C75E1D5416702BF4328A9945356607F245E19F15AC88997AF60D33C4
408186AFA0FF006F7FD8C358F85DAE6A1F1CBE0FE9CF71E0FBE2F73ABD8C00
1FECA9891BE645CEE36F2125885044273D23DA17F3B743F18AC8ABF3D78988
C258D548FB0ED3C40303E6ADB8B5E047DEAF9A2C7C52AC07CF5D1C3E25523E
FD79F2C31573DF7FB7463EF556975D18FBD5E2C3C46B8FBFFAD412F88D71F7
EA561C773D62E75E1CFCDFAD72D7DAEF07E6AF38B8F11E7386CD605C6AF24B
9CB6056B1A02B9D76A3AD97CAA9C9AE3EEAF49C9639ACA9AFB00F358CD7173
7B750D85844F737572EB145144A5E492473B551157259989C000649E0574C2
992D9BFA2E8BADF8DFC4BA77843C3501B9D4F569D2DE04EDB98FDE63D95465
99BA2A824F02BF61755F807A6FC18F09691A6F85AF7FB574C8A254B89CA6C9
3ED47991DD7B2C8C495E4E3EEE7819F4BFD89FF634B3F843E1A3E3AF89B669
71E37D761DAD6F26D75D2ED98E4428C33FBE71832B83C7DC5E0333FD93AB7C
2FD1AFE0960B599E08E5055A39079D1953DB0C437FE3D5E9D1A128D9A21B3F
2B1D6ABB2D7D3BF103F675F16681E6EA7E1F886A762B9629092D34607FB070
587D327F9D7CDD2C2E8C51C15653820F041F7AEE24CF2BEB49B455A294DD86
80200BE94F0B5308CD5EB1D3AF350B98ECEC207B8B894E1238D4BBB1F40A01
24FD28029C695F6EFECE3F09E6B464F883E2180C6E548D3E27183B5860CE47
B8384F6CB63EE9A67C25FD9C9EDA687C45F106352D190F0D86430CF50673D0
E3FB83FE047AAD7D8C00000030050054D474FB3D5AC2E34CD422135B5DC6D1
4887A3238C11F957E5C7C49F026A3F0EFC4D3E8B77BA4B6725ED6723896227
83E9B87461D8FB62BF54EB88F1F780743F887A1BE8DACA6D75CB413A81E641
263EF2FB1FE25E847BE0800FCA71252F995D7F8F7E1CF8A3E1DEA46CB5CB72
6073FB9B9404C330EBF2B63A8EEA7047D304F07E61A00BBE6534C954FCC349
BCD005A3254664A84127934F03BD003C726BD9BE0BFC369FE20789E337719F
EC7D3D965BB7E8180E56207D5C8C1F45C9EB8CE5FC35F84FE23F88F7C059A1
B4D3236C4D78EB944C75551C6F7F61F8915FA43E12F0A68FE0BD0ADFC3FA1C
463B780649639791CFDE773DD8FE5D860002803A3555450880055180070001
4B451400514514005145140051451400514514005145140051451400514514
00514514005145140051451400514514005145140051451400514514005145
1400514514005145140051451401FFD4FDFCA28A2800A28A2800A28A2800A2
8A2800A28A2800A28A2800A28A2800A28A2800A28A2800A28A280396F13F83
740F175BAC3AC5BEE78FFD5CC876CA99FEEB7A7B1C8F6AF0ED47E046AF1484
E89AAC53479E05C2B4640F4CA07CFD702BE9AA2B2A94632DD0D33E59B7F823
E2D7702E6F6CE243D4AB48EC3F0D8A3F5AF4FF000C7C23D0343952F35173AA
5D21CA99142C4A7D447CE4FF00BC4FE06BD5E8A886161177486E4C28A28AE8
2428A28A00C0D7740D1B57B599B50B18AE6408DB59901704038C37DE1F81AF
9A749BA1E5AF3DABEB5AF8F75DB493C31E23BCD2240523472F0FA189F9423D
70383EE0D7CEF10E15CE09A30AF1BA3D434AD79ECD40078ABF7DE2796742A0
E335E4D0EA431F7AA76D4463EF57C14B072D8E3D6D637EEEF371249E4D721A
8DD0C1E6A3B9D4460F35CADEDF190900D7A581C134CA840F61F828EEFAFEA6
C33B0DB8CFA6778C67DFAE3F1AFA4ABC67E0BE86F63A15C6B33A95935271B3
23FE594790A7F125BEA306BD9ABF43C1C3969A4CEF8EC14514575141451450
0145145001451450014514500453430DCC325BDC46B2C52A94747019595860
820F0411D457E29FED69FF0004C47B8B9BBF889FB32C496F23032DCF86D982
2337566B1918ED5CFF00CF172075D8C0612BF6C68A99453DC2E7F181A8AF89
BC21AC5C7877C556173A46A964C639ED6F227827898767470181FA8AD2B6F1
0CA40F9EBFAC7F8BBFB3DFC1BF8ED62967F14FC2D69ADBC086386E58345770
293BB115CC452545CF25436D27A835F98FF11BFE090BA14D2CB79F087C7973
A729395B3D6205B9419EC2E20F2D95476CC4C71D493C9E59E1BB17CC7E43C7
ADCCDFC75606A8EDF79EBEBCD6FF00E098FF00B5BE872B45A7D968FAEAA824
3D96A2A8AC40E805D240727A0C803D48AE56D7FE09F1FB68CB3A452F8016DD
5CE0C8FAB6965547A9D974CD8FA026B07867D8773E6EFED11EB55A5D4C01CB
57DE9E13FF00825BFED2FAEBA3F89750D0FC37067E712DD4973381923E5482
268CF4CF328E08EF903ED7F859FF0004ABF843E189E2D4BE286BB7DE36B98F
07ECC83FB3AC49E09DC91B3CCD83C0FDF28233953D9C70B262B9F8D9F0C3E1
87C4AF8DFE245F0AFC31D126D62F3832C8B84B7B64FEFCF3361235F4C9CB1E
1416C03FBDFF00B277EC3BE11FD9E843E31F13CF1F88FC7AF1B2B5E053F65B
2120C325A2380724654CCC03B024008ACCA7EC9F08F833C25E00D0A0F0CF82
747B5D0F4AB61FBBB6B389618C1C005885032C71F331CB13C924D74D5D74E8
28EA4B61451456E20AF3BF187C2CF0478D5259758D2E26BC7538B88F314BBB
18059D305B1FED647B57A251401F9C379F07EC4CAF1DBDF496EF192AC92A07
218704120AF7ACE1F06E6C8CEA898FFAE47FF8AAFB2BC77E03BCBABA7D7340
4F3249399A007059BFBC99E327B8FC4726BC52696E6DA4305CC6D148BD55D4
AB0FA835C753DA27A303CDAD3E12E8D6E775F5ECB738EC8A2307FF00423FAD
7D07F05BC31A6699ADDCDD6976AB0C7040519C0CB167618058E4F404F5AE63
47D0F5AF114EB0E9D6ED2293CB9188D7DCB74FC3AFA0AFA6BC2FE1CB6F0CE9
8B6509F32573BA593182EFFE03B0FEB9AAA5193779303A3A28A2BA8028A28A
00CFD5349D335CB1974CD62D63BCB49861E29543A9FC0F71D8F507915F2DF8
CBF657D1AFE57BBF065F9D3598E7ECF700CB08F6571F3A8FAEEAFAD28A00FC
D1D5BF676F8A7A5C8445A625FC60E37DB4C8C0FF00C05CABFF00E3B5CCA7C1
EF89CF20897C3778189C731E07E64E3F5AFD54A2803F36B46FD9CBE28EA8CB
F69B18B4D8D8FDFB999077C1F9632EDFA57D05E0CFD97FC35A43C77BE2DBB6
D6275C1F25018ADC1F7C1DEF8FAA83DD6BEA3A2802BDA59DA585B4565630A5
BDBC2A15238D42A2A8E8001C015628A2800A28A2800A28A2800A28A2800A28
A2800A28A2800A28A2800A28A2800A28A2800A28A2800A28A2800A28A2800A
28A2800A28A2800A28A2800A28A2800A28A2800A28A2800A28A2800A28A280
3FFFD5FDFCA28A2800A28A2800A28A2800A28A2800A28A2800A28A2800A28A
2800A28A2800A28A2800A28A2800A28A2800A28A2800A28A2800A28A2800AF
3AF887E048FC6360B2DA32C1A9DA83E4C87EEB0EF1BE3F84F63D41FC41F45A
2A6514D5981F02DF1D5B41BC7D3759B67B5B88FAAB8C647AA9E847A11C1A84
EB058715F76EABA2E93AE5BFD9757B48AEE2EC2450D8F704F20FB8AF3BB8F8
2DE0399F7C76D35B8CE76A4CE47D3E62C71F8D79553298B77464E923E4D92F
99F926BD3FC03F0DF50F145C47A96AB1B5B692A77648DAD3E3F8501FE13DDB
F01CF4FA0348F86BE09D16413DA698924ABC8798B4C41EC40724023D4015DD
56B432E8C5DD8D53B11C31456F12410208E38D42AAA8C055030001E8054945
15E91A05145140051451400514514005145140051451400514514005145140
051451400514514005145140051451400514514005452410CD8F3A357C74DC
01C7E752D140080050001803B52D1450014514500145145001451450014514
50014514500145145001451450014514500145145001451450014514500145
14500145145001451450014514500145145001451450014514500145145001
451450014514500145145001451450014514500145145007FFD6FDFCA28A28
00A28A2800A28A2800A28A2800A28A2800A28A2800A28A2800A28A2800A28A
2800A28A2800A28A2800A28A2800A28A2800A28A2800A28A2800A28A2800A2
8A2800A28A2800A28A2800A28A2800A28A2800A28A2800A28A2800A28A2800
A28A2800A28A2800A28A2800A28A2800A28A2800A28A2800A28A2800A28A28
00A28A2800A28A2800A28A2800A28A2800A28A2800A28A2800A28A2800A28A
2800A28A2800A28A2800A28A2800A28A2800A28A2800A28A2800A28A2800A2
8A2800A28A2800A28A2800A28A2800A28A2800A28A2800A28A2800A28A2800
A28A2800A28A2803FFD9>
%%EOF
%</meme>
%    \end{macrocode}
%    \Finale
\endinput
